% Archive-normalized complete cumulative translated source, numbered Papers 1--43.
% Derived mechanically from the 219 exact producer-frozen v014 source units in manifest order.
% No translated wording was edited. Every unit marker records its original byte count and SHA-256.
% Producer manifest SHA-256: F13DE01A4568708CA7A8F22FBA4CE2E08F9D784AFEE3BF7F8E0B4F293829503C
% Language surface: russian
\documentclass[11pt]{article}
\usepackage{fontspec}
\IfFontExistsTF{Times New Roman}{\setmainfont{Times New Roman}}{\setmainfont{TeX Gyre Termes}}
\IfFontExistsTF{Arial}{\setsansfont{Arial}}{\setsansfont{TeX Gyre Heros}}
\IfFontExistsTF{Consolas}{\setmonofont{Consolas}}{\setmonofont{Latin Modern Mono}}
\usepackage{amsmath,amssymb,mathtools,array,booktabs}
\usepackage{adjustbox,diagbox,enumitem,pdflscape}
\usepackage{geometry,microtype}
\geometry{a4paper,margin=2.25cm}
\setlength{\parindent}{1.25em}
\setlength{\parskip}{0pt}
\allowdisplaybreaks
\emergencystretch=24em
\raggedbottom
\DeclareMathOperator{\Div}{Div}
\begin{document}


\clearpage
\begingroup
\setcounter{footnote}{0}
% BEGIN PRODUCER UNIT 001
% Source: C:/Users/Floris/Documents/interlanguage/03_projects/noether/02_slavic_working_corpus/translations/paper01/russian/v001/Noether_Paper01_Russian_v001.tex
% Identity: 11629 B / SHA-256 3E1C5D12E8694E181CDD33CB4A97BC9F4A100415DCE6669A4649C805285C7551
\setlength{\parindent}{1.25em}
\setlength{\parskip}{0pt}
\emergencystretch=4em
% BEGIN UNIT BODY
\begin{center}
{\Large\bfseries 1. Построение системы форм для тернарной\\ формы четвёртой степени}\par
\vspace{1.5em}
Sitzungsberichte Физико-медицинского общества в Эрлангене 39 (1907), с. 176--179
\end{center}

\vspace{4em}
\begin{center}
(Отрывок из диссертации авторки.)
\end{center}

Систему форм тернарной биквадратичной формы, то есть тернарной формы четвёртой степени, изучали Гордан, Майзано и Паскаль.\footnote{P. Gordan, \emph{Über das volle Formensystem der ternären biquadratischen Form} $f=x_1^3x_2+x_2^3x_3+x_3^3x_1$ (Math. Annalen, т. XVII, с. 217--233, 1880); G. Maisano, 1. \emph{Sistemi completi dei primi cinque gradi della forma ternaria biquadratica e degl' invarianti, covarianti e contravarianti di sesto grado} (Giorn. di Battaglini XIX). 2. \emph{Sui covarianti indipendenti di $6^\circ$ grado nei coefficienti della forma biquadratica ternaria} (Rend. Circ. Mat. di Palermo I, 1887); E. Pascal, \emph{Contributo alla teoria della forma ternaria biquadratica e delle sue varie decomposizioni in fattori} (Memoria premiata dalla R. Accademia delle scienze fisiche e matematiche di Napoli, 1905).}
Гордан строит полную систему форм, состоящую из 54 построений, для специальной автоморфной формы
\[
f=x_1^3x_2+x_2^3x_3+x_3^3x_1,
\]
исходя из принципов, подобных тем, которые он дал для систем форм в бинарной области.

У Майзано для общей биквадратичной формы построены все формы до пятого порядка включительно,\footnote{Под "порядком" здесь понимается степень по коэффициентам, а под "степенью" -- степень по переменным.}
а также некоторые инварианты, коварианты и контраварианты более высокого порядка. Он использует метод, применённый Горданом в первом томе \emph{Mathematische Annalen} к тернарной кубической форме. Паскаль, пользуясь результатами Майзано, занимается главным образом вопросом о разложении биквадратичной формы на множители.\footnote{Во второй короткой заметке Майзано пытается доказать линейную зависимость трёх ковариантов порядка 6 и степени 6. В противоположность этому не вполне полному доказательству Паскаль считает доказанной линейную независимость этих трёх ковариантов порядка 6, а также трёх инвариантов порядка 9. Однако в ходе моих вычислений выяснилось, что линейное соотношение действительно существует как между тремя ковариантами, так и между тремя инвариантами. В обозначениях Паскаля явные формулы имеют вид:
\[
0=20\Omega_1+6\Omega_2-3\Omega_3+4fC_1-12fC_2-4A\cdot\Delta,
\]
\[
0=4A^3-15AB+30C-90D+9E.
\]}

\textbf{Цель моих исследований состоит в построении системы форм для общей тернарной формы четвёртой степени. Сначала я даю только основные основания и строю так называемую "относительно полную систему".\footnote{О термине см. Gordan--Kerschensteiner, \emph{Vorlesungen über Invariantentheorie}, т. II, с. 227.}}

Основная идея построения систем тернарных форм та же, что и в бинарной области. Начиная с первой относительно полной системы, то есть с системы форм, перенесённых из бинарного случая, по определённому закону переходят к системам со всё более высоким модулем. Этот процесс продолжается до тех пор, пока система некоторого модуля, если сам модуль взять за основную форму, не станет конечной и известной, либо пока модуль не удастся свести к формам, имеющим инвариант множителем. В первом случае абсолютно полная система получается трансвекцией относительно полной системы с системой модуля; во втором случае относительно полная и абсолютно полная системы совпадают. Благодаря общей теореме Гильберта о конечности систем форм эта процедура необходимо должна завершиться.

В нашем случае модуль $(abc)$ первой относительно полной системы можно свести к модулям
\[
\Delta=(abc)^2a_x^2b_x^2c_x^2\quad\text{и}\quad \nu=(abu)^4;
\]
отсюда легко получается относительно полная система по модулю $\nu$. В качестве ряда модулей теперь выбираем формы
\[
\nu;\qquad \nu^{(\nu)}=(\nu\nu_1x)^4=s_x^4,
\]
\[
\nu^{(s)}=(ss'u)^4\ldots
\]
и присоединяем как дальнейшие модули две квадратичные формы, возникающие при построении относительно полной системы по модулю $s$, а именно $u_\sigma^2$ и $t_x^2$. Затем можно показать, что модуль, следующий за $s$, то есть $(ss'u)^4$, сводим к модулю $(\varrho,t)$. Но одновременная система двух квадратичных форм конечна и известна; тем самым достигнуто описанное выше завершение. Относительно полная система по модулю $(\varrho,t)$ даёт 331 построение. ---

Добавлю ещё несколько замечаний о применяемом методе.

Основной процесс образования форм, процесс свёртки, в тернарной области можно определить следующим образом.

Если в символическом произведении
\[
s_x^m t_x^n u_\sigma^\mu u_\tau^\nu
\]
заменить пары множителей
\[
\begin{array}{c|c|c|c|c}
 & s_xt_x & u_\sigma u_\tau & s_xu_\sigma\ \text{или}\ s_xu_\tau & t_xu_\tau\ \text{или}\ t_xu_\sigma\\
\text{соответственно на} & (stu) & (\sigma\tau x) & s_\sigma\ \text{или}\ s_\tau & t_\tau\ \text{или}\ t_\sigma\\
\text{свёртка} & \mathrm{I} & \mathrm{II} & \mathrm{III} & \mathrm{IV}
\end{array}
\]
то возникающие формы получены из исходной формы свёрткой.\footnote{Gordan, Math. Annalen, т. XVII, с. 219.}

При этом всегда ещё можно положить $s_\sigma=0$ и $t_\tau=0$. Для связи между отдельными свёртками справедлива следующая теорема.

\textbf{Свёртки I и II являются основными свёртками. Из них, независимо от порядка композиции, можно составить свёртки III и IV. Другими словами: чтобы построить все формы, возникающие из данного выражения путём свёртки, достаточно применять только свёртки I и II.}\footnote{Теорема имеет исключение в случае, когда $n$ и $\mu$, соответственно $m$ и $\nu$, одновременно обращаются в нуль; тогда ряд форм сводится к диагональному члену.}

Под "рядом форм" мы понимаем начальную форму вместе со всеми формами, возникающими из неё повторной свёрткой с самой собой. По этой теореме ряд форм
\[
s_x^m t_x^n u_\sigma^\mu u_\tau^\nu
\]
можно расположить в прямоугольной схеме. При продвижении в этой схеме получают:
\begin{enumerate}
\item на одну колонку вправо -- все формы, возникающие из соседних форм посредством свёртки I;
\item на одну строку вниз -- все формы, возникающие из стоящих выше форм посредством свёртки II.
\end{enumerate}
Тем самым получаются все формы, возникающие посредством свёртки.

При этих предположениях теоремы о редуцентах можно сформулировать в самой общей форме, используя определение
\begin{center}
\textbf{редуцент -- это редуцируемый ряд форм.}
\end{center}
Тогда общая теорема имеет вид:

\begin{center}
\textbf{Если начальная форма ряда форм редуцируема потому, что один из её членов возник посредством свёртки с редуцентом, и если конечная форма этого ряда возникла из того же самого члена посредством свёртки, то весь ряд форм редуцируем.}
\end{center}

В качестве дальнейших методов редукции следует назвать:
\begin{enumerate}
\item так называемую двойную редукцию, то есть редукцию одной формы двумя способами для получения соотношения между формами более высокого порядка;
\item свёртку с формами, разлагающимися на множители; этот метод отчасти имеет характер редукции при помощи редуцентов, а отчасти характер двойной редукции.
\end{enumerate}
% END UNIT BODY
% END PRODUCER UNIT 001
\endgroup


\clearpage
\begingroup
\setcounter{footnote}{0}
% BEGIN PRODUCER UNIT 002
% Source: C:/Users/Floris/Documents/interlanguage/03_projects/noether/02_slavic_working_corpus/translations/paper02/russian/v001/Noether_Paper02_Introduction_Russian_v001.tex
% Identity: 10422 B / SHA-256 74031BEB51C843447192DFF5C16F9A606A094D3327BF63FB1C04BCB6EA4DE335
\setlength{\parindent}{1.25em}
\setlength{\parskip}{0pt}
\emergencystretch=4em
% BEGIN UNIT BODY
\begin{center}
{\Large\bfseries 2. Построение системы форм для тернарной\\ биквадратичной формы}\par
\vspace{1.5em}
\emph{Journal für die reine und angewandte Mathematik} 134 (1908), с. 23--90 и две таблицы
\end{center}

\vspace{3em}

\begin{center}
\textbf{Введение.}
\end{center}

Систему форм тернарной биквадратичной формы изучали Гордан, Майзано и Паскаль.\footnote{Краткий отрывок из введения и главы I появился в \emph{Sitzungsberichte der physikalisch-medizinischen Sozietät Erlangen 1907}, с. 176--179.}\footnote{P. Gordan, \emph{Über das volle Formensystem der ternären biquadratischen Form}
\[
f=x_1^3x_2+x_2^3x_3+x_3^3x_1.
\]
(\emph{Math. Annalen}, т. XVII (1880), с. 217--233.)
G. Maisano, 1) \emph{Sistemi completi dei primi cinque gradi della forma ternaria biquadratica e degl' invarianti, covarianti e contravarianti di sesto grado}. (Giorn. di Battaglini XIX.)
2) \emph{Sui covarianti indipendenti di $6^\circ$ grado nei coefficienti della forma biquadratica ternaria}. (Rend. Circ. Mat. di Palermo I, 1887.)
E. Pascal, \emph{Contributo alla teoria della forma ternaria biquadratica e delle sue varie decomposizioni in fattori}. (Memoria premiata dalla R. Accademia delle scienze fisiche e matematiche di Napoli. 1905.)}
Гордан строит полную систему форм, состоящую из 54 построений, для специальной автоморфной формы
\[
f=x_1^3x_2+x_2^3x_3+x_3^3x_1,
\]
исходя из принципов, подобных тем, которые он дал для систем форм в бинарной области.

У Майзано для общей биквадратичной формы построены все формы до пятого порядка включительно,\footnote{Под "порядком" здесь понимается степень по коэффициентам, а под "степенью" -- степень по переменным.}
а также некоторые инварианты, коварианты и контраварианты более высокого порядка. Он использует метод, применённый Горданом в первом томе \emph{Mathematische Annalen} к тернарной кубической форме. Паскаль, опираясь на результаты Майзано, занимается главным образом вопросом о разложении биквадратичной формы на множители.\footnote{Во второй короткой заметке Майзано пытается доказать линейную зависимость трёх ковариантов порядка 6 и степени 6. В противоположность этому не вполне завершённому доказательству Паскаль считает, что он доказал линейную независимость этих трёх ковариантов порядка 6, а также трёх инвариантов порядка 9. Однако то, что линейное соотношение действительно существует между тремя ковариантами, с одной стороны, и тремя инвариантами, с другой, показывают формулы (13), § 11, и (22), примечание к § 17, настоящей работы, где эти соотношения даны явно. По сообщению Паскаля, противоречие объясняется численной ошибкой в выражении его контраварианта $p$ (здесь обозначенного $\varrho$), с. 46 его статьи.}

\emph{Цель следующих исследований состоит в построении системы форм для общей тернарной биквадратичной формы. В этой работе даны только основные основания и построена так называемая "относительно полная система".\footnote{О термине см. § 3a.}}

Работа тесно примыкает к работе Гордана; однако принципы, которые там были даны лишь в самом общем виде, сначала требовалось разработать в деталях. С помощью теоремы о связи между свёртками и через введение "ряда форм" (§ 1) теоремы редукции формулируются точно и дополняются (§ 3), тогда как рекурсивное построение специальных разложений в ряды (§ 2) даёт вычислительное средство для фактического проведения редукций.

Основная идея построения систем тернарных форм та же, что и в бинарной области. Начиная с первой относительно полной системы --- системы форм, перенесённых из бинарного случая, --- по определённому закону переходят к системам со всё более высоким модулем. Этот процесс продолжается до тех пор, пока система некоторого модуля, если сам модуль взять за основную форму, не станет конечной и известной, либо пока модуль не удастся свести к формам, имеющим инварианты множителями. В первом случае абсолютно полная система получается трансвекцией относительно полной системы с системой модуля; во втором случае относительно полная и абсолютно полная системы совпадают. Благодаря общей теореме Гильберта о конечности систем форм эта процедура необходимо должна завершиться.

В нашем случае модуль $(abc)$ первой относительно полной системы (§ 4) сводится к модулям
\[
\Delta=(abc)^2a_x^2b_x^2c_x^2\quad\text{и}\quad \nu=(abu)^4
\]
(§ 5), а затем находится относительно полная система по модулю $\nu$ (§ 6). Эти результаты уже приведены в работе Гордана для общей биквадратичной формы; однако наш способ их вывода легче переносится на формы более высокой степени.

В качестве ряда модулей теперь выбираем формы (§ 7):
\[
\nu,\qquad \nu(\nu)=(\nu\nu_1x)^4=s_x^4,\qquad \nu(s)=(ss'u)^4,\ldots .
\]

При построении относительно полной системы по модулю $s$ две квадратичные формы, возникающие в системе, $u_\varrho^2$ и $t_x^2$, присоединяются как модули (§§ 9 и 10), и соответственно строится относительно полная система по модулю $(s,\varrho,t)$. Затем показано (§ 17), что модуль $(ss'u)^4$ сводится к модулям $\varrho$ и $t$. Благодаря этому следующая более высокая система переходит в относительно полную систему по модулю $(\varrho,t)$. Но одновременная система двух квадратичных форм конечна и известна, а в общем случае состоит из 20 построений;\footnote{Ср. Clebsch--Lindemann, \emph{Vorlesungen über Geometrie} (т. I, с. 288 и далее), система двух когредиентных форм. Gordan, \emph{Über Büschel von Kegelschnitten} (\emph{Math. Ann.}, т. XIX, с. 530), система двух контрагредиентных форм.}
тем самым с построением относительно полной системы по модулю $(\varrho,t)$ достигнуто описанное выше завершение. Трансвекция этой системы с системой $(\varrho,t)$ для построения абсолютно полной системы остаётся на будущее.

В качестве относительно полной системы по модулю $(\varrho,t)$ найдено 331 построение; в приложенной таблице они упорядочены по степени в переменных $x$ и $u$.

В заключение дадим обзор введённых символов:
\begin{align*}
f&=a_x^4,\\
\theta&=\theta_x^4u_\vartheta^2=(abu)^2a_x^2b_x^2,\\
K&=K_x^6u_k^3=(a\theta u)u_\vartheta^2a_x^3\theta_x^3,\\
j&=u_j^6=(a\theta u)^4u_\vartheta^2,\\
\Delta&=\Delta_x^6=a_\vartheta^2a_x^2\theta_x^4=(abc)^2a_x^2b_x^2c_x^2,\\
N&=N_x^8u_\eta=(a\Delta u)a_x^3\Delta_x^5,\\
\nu&=u_\nu^4=(abu)^4,\\
H&=H_x^2u_\eta^4=(\nu\nu_1x)^2u_\nu^2u_{\nu_1}^2,\\
L&=L_x^3u_l^6=(\nu\eta x)H_x^2u_\nu^3u_\eta^3,\\
g&=g_x^6=(\nu\eta x)^4H_x^2,\\
\sigma&=u_\sigma^6=H_\nu^2u_\nu^2u_\eta^4,\\
s&=s_x^4=(\nu\nu_1x)^4,\\
Z&=Z_x^4u_\zeta^2=(ss'u)^2s_x^2s_x'^2,\\
\varrho&=u_\varrho^2=\theta_\nu^4u_\vartheta^2,\\
t&=t_x^2=a_\eta^4H_x^2,\\
i&=a_\nu^4,\\
J&=s_\nu^4.
\end{align*}
% END UNIT BODY
% END PRODUCER UNIT 002
\endgroup


\clearpage
\begingroup
\setcounter{footnote}{0}
% BEGIN PRODUCER UNIT 003
% Source: C:/Users/Floris/Documents/interlanguage/03_projects/noether/02_slavic_working_corpus/translations/paper02/russian/v001/Noether_Paper02_Section01_Russian_v001.tex
% Identity: 9970 B / SHA-256 CDB5E334CC4E590171241DE5F4958A62AE4EDFD4E3ED005919F744147BFADE53
\setlength{\parindent}{1.25em}
\setlength{\parskip}{0pt}
\emergencystretch=4em
% BEGIN UNIT BODY
\begin{center}
{\Large\bfseries Глава I. \quad Общие теоремы о тернарных формах.}
\end{center}

\begin{center}
\textbf{§ 1. \quad Процесс свёртки. \quad Ряды форм.}
\end{center}

Основным процессом образования форм является процесс свёртки, который в тернарной области можно определить следующим образом. Пусть дано символическое произведение
\[
s_x^m t_x^n u_\sigma^\mu u_\tau^\nu .
\]
Если пары множителей
\[
s_xt_x\qquad u_\sigma u_\tau\qquad s_xu_\sigma\ \text{или}\ s_xu_\tau
\qquad t_xu_\tau\ \text{или}\ t_xu_\sigma
\]
заменить соответственно на
\[
(stu)\qquad (\sigma\tau x)\qquad s_\sigma\ \text{или}\ s_\tau
\qquad t_\tau\ \text{или}\ t_\sigma,
\]
то есть на свёртки
\[
\text{I}\qquad\text{II}\qquad\text{III}\qquad\text{IV},
\]
то возникающие формы получены из исходной формы посредством свёртки.\footnote{Gordan, \emph{Math. Annalen}, т. XVII, с. 219.}

При этом выражение $s_x^m t_x^n u_\sigma^\mu u_\tau^\nu$ можно рассматривать либо как настоящее произведение двух форм $S\cdot T$, либо как единственную форму с несколькими рядами символов:
\[
s_x^m t_x^n u_\sigma^\mu u_\tau^\nu=A_x^{m+n}u_x^{\mu+\nu}.
\]
В первом случае мы говорим о свёртке формы $S$ с формой $T$, во втором -- о «свёртке формы $A$ с самой собой». Для краткости в дальнейшем всегда предполагаем (что действительно имеет место для всех форм, рассматриваемых ниже), что
\[
\tag{a}
s_\sigma^\lambda=0,\qquad t_\tau^\chi=0
\]
для любых $\lambda$ и $\chi$, отличных от нуля, и вместе с любыми другими свёртками. Это означает, что форма $s_x^m t_y^n u_\sigma^\mu u_\tau^\nu$ является так называемой «нормальной формой», то есть при применении $\Omega$-процесса она обращается в нуль как по переменным $x$ и $u$, так и по переменным $y$ и $v$.

Следовательно, условие (a) не является ограничением; его всегда можно добиться.\footnote{Ср. Gordan, \emph{Math. Annalen}, т. V, с. 104.}

Для связи между отдельными свёртками справедлива следующая теорема:

\emph{Теорема I. Свёртки I и II являются основными свёртками; из них, независимо от порядка композиции, можно составить свёртки III и IV. Другими словами: чтобы построить все формы, возникающие из данного выражения посредством свёртки, надо применять только свёртки I и II.}

Доказательство. По теореме тождества, соответственно по теореме произведения для матриц, с учётом (a) получаем
\[
(\widehat{st}\sigma x)=-t_\sigma,\qquad
(\widehat{\sigma\tau}su)=-s_\tau,
\]
\[
(\widehat{st}\tau x)=s_\tau,\qquad
(\widehat{\sigma\tau}tu)=t_\sigma,
\]
\[
(stu)(\sigma\tau x)=s_\tau+t_\sigma-u_x\,s_\tau t_\sigma .
\]
Следовательно, посредством композиции свёрток I и II возникают свёртки III и IV; это происходит как при применении свёртки II к свёртке I, то есть к множителям $(stu)u_\sigma u_\tau$, так и при применении свёртки I к свёртке II, то есть к множителям $(\sigma\tau x)s_xt_x$.

Под \emph{рядом форм}\footnote{Ср. Clebsch, \emph{Über eine Fundamentalaufgabe der Invariantentheorie} (Abh. der Gött. Ges. d. Wiss., т. XVII): § 17 и конец § 18. «Ряд форм» отличается от введённой там «собственно редуцированной эквивалентной системы» принципом упорядочения по более высоким формам.}
мы понимаем начальную форму вместе со всеми формами, возникшими из неё посредством свёртки с самой собой, и обозначаем ряд форм его начальной формой. «Более высокой формой» здесь считается всякая форма с большим числом свёрток с самой собой; при этом среди равноправных свёрток $s_\tau$ и $t_\sigma$ одну необходимо нормировать как более высокую. Из закона построения ряда форм следует:

1) Каждая форма ряда форм линейна по символам начальной формы.

2) Ряд форм однозначно определён для начальных форм, каждая из которых имеет два ряда контрагредиентных символов; для начальных форм с большим числом рядов символов его надо однозначно нормировать, выделяя специальные свёртки среди тех, которые связаны теоремой тождества.

По теореме I ряд форм $s_x^m t_x^n u_\sigma^\mu u_\tau^\nu$ ($m>n;\ \mu>\nu$) можно расположить по следующей прямоугольной схеме:
\begingroup\small
\[
\begin{array}{c|c|c|c}
s_x^m t_x^n u_\sigma^\mu u_\tau^\nu
& (stu)
& (stu)^2
& \cdots\\[0.3em]
(\sigma\tau x)
& t_\sigma,\ s_\tau
& t_\sigma(stu),\ s_\tau(stu)
& \cdots\\[0.3em]
(\sigma\tau x)^2
& t_\sigma(\sigma\tau x),\ s_\tau(\sigma\tau x)
& t_\sigma^2,\ t_\sigma s_\tau,\ s_\tau^2
& \cdots\\[0.3em]
\vdots&\vdots&\vdots&\ddots\\[0.3em]
(\sigma\tau x)^\nu
& t_\sigma(\sigma\tau x)^{\nu-1},\ s_\tau(\sigma\tau x)^{\nu-1}
& t_\sigma^2(\sigma\tau x)^{\nu-2},\ t_\sigma s_\tau(\sigma\tau x)^{\nu-2},\ s_\tau^2(\sigma\tau x)^{\nu-2}
& \cdots\\[0.3em]
\vdots
& t_\sigma(\sigma\tau x)^\nu
& t_\sigma^2(\sigma\tau x)^{\nu-1},\ t_\sigma s_\tau(\sigma\tau x)^{\nu-1}
& \cdots
\end{array}
\]
\endgroup
и соответственно продолженную нижнюю часть\footnote{Здесь, и подобным образом далее, нижнюю часть, отмеченную звёздочкой, следует присоединять к так же отмеченному месту справа вверху схемы.}
\begingroup\small
\[
\begin{array}{c|c|c}
(stu)^n & \cdots & \cdots\\[0.3em]
t_\sigma(stu)^{n-1},\ s_\tau(stu)^{n-1}
& s_\tau(stu)^n & \cdots\\[0.3em]
t_\sigma^2(stu)^{n-2},\ t_\sigma s_\tau(stu)^{n-2},\ s_\tau^2(stu)^{n-2}
& t_\sigma s_\tau(stu)^{n-1},\ s_\tau^2(stu)^{n-1}
& s_\tau^2(stu)^n .
\end{array}
\]
\endgroup

Здесь получают, если двигаться

1) на одну колонку вправо, все формы, возникающие из соседних форм посредством свёртки I;

2) на одну строку вниз, все формы, возникающие из стоящих выше форм посредством свёртки II;

а следовательно, все формы, возникающие посредством свёртки.

Произвольная форма полного ряда форм, $t_\sigma^\kappa s_\tau^\lambda(stu)^\mu$ или $t_\sigma^\kappa s_\tau^\lambda(\sigma\tau x)^\nu$, образует начало нового ряда форм. Он состоит из всех тех форм прямоугольной схемы, которые имеют форму $t_\sigma^\kappa s_\tau^\lambda(stu)^\mu$ или $t_\sigma^\kappa s_\tau^\lambda(\sigma\tau x)^\nu$ в качестве левого верхнего угла и имеют множитель $t_\sigma^\kappa s_\tau^\lambda$.

Дополнительное замечание. Теорема I имеет исключение в случае, когда числа $n$ и $\mu$ (соответственно $m$ и $\nu$) равны нулю, так что свёртки I, II и IV больше не существуют, но свёртка III (соответственно IV) ещё существует. В этом случае рядом форм мы называем диагональный член схемы:
\[
\begin{array}{ccccc}
s_x^m u_\tau^\nu &&&& t_x^n u_\sigma^\mu\\
& s_\tau && t_\sigma &\\
& s_\tau^2 && t_\sigma^2 &\\
&&\ddots&&\\
&&s_\tau^\nu && t_\sigma^\mu .
\end{array}
\]
В соответствии с отсутствием свёрток I и II разложение в ряд по полярам ряда форм в этом случае всегда сводится к начальному члену; однако теоремы о редуцентах остаются в силе.
% END UNIT BODY
% END PRODUCER UNIT 003
\endgroup


\clearpage
\begingroup
\setcounter{footnote}{0}
% BEGIN PRODUCER UNIT 004
% Source: C:/Users/Floris/Documents/interlanguage/03_projects/noether/02_slavic_working_corpus/translations/paper02/russian/v001/Noether_Paper02_Section02_Russian_v001.tex
% Identity: 9228 B / SHA-256 BC65D909577D444C52343ACDB59169594CC96F778C99EF4C93F1F44FCF9088BC
\setlength{\parindent}{1.25em}
\setlength{\parskip}{0pt}
\emergencystretch=4em
% BEGIN UNIT BODY
\begin{center}
\textbf{§ 2. \quad Разложения в ряд по полярам ряда форм.}
\end{center}

Для форм с $n$ переменными явно построены два вида разложений в ряд\footnote{Ср. Gordan, \emph{Über Kombinanten}, § 2 и 5 (\emph{Math. Ann.}, т. V).}:

1) для форм, каждая из которых имеет один ряд контрагредиентных переменных, $s_x^m u_\tau^\nu$, -- ряд, идущий по степеням $u_x$, коэффициенты которого являются «нормальными формами»;

2) для форм с двумя рядами когредиентных переменных $s_x^m t_x^n$ -- разложение по полярам элементарных ковариантов (полярам ряда форм $s_x^m t_x^n$), которое в тернарном случае имеет следующий вид\footnote{Ср. Study, \emph{Methoden zur Theorie der ternären Formen} (Teubner 1889) II. § 7. В отличие от приведённого там вывода, наш даёт метод быстрого вычислительного определения числовых коэффициентов $C_{pr\varrho}$. Разложение Клебша--Штуди при специализации a), a') имело бы вид:
\[
(stu)^\mu u_\tau^{\nu-\kappa}v_\tau^\kappa s_x^m t_x^{\,n-\lambda}(tuv)^\lambda
=\sum_{p,r,\varrho} C_{pr\varrho}
\bigl[s_\tau^\varrho(stu)^{\mu+r-\varrho}\bigr]_{\widehat{uv}^{\lambda-r+\varrho}v^{\nu-\varrho},\,u^p,\,(vw=\widehat{xw})},
\]
и поэтому не допускает упрощения, которое в нашем разложении возникает уже в ходе вычисления, если известно, что все члены с множителем $u_x$ следует опустить.}:
\[
\tag{I}
s_x^m t_x^{n-\lambda}t_y^\lambda
=\sum_{\varrho}
\frac{\binom{m}{\varrho}\binom{\lambda}{\varrho}}
{\binom{m+n-\varrho+1}{\varrho}}
\bigl[(st\widehat{xy})^\varrho s_x^{m-\varrho}t_x^{n-\varrho}\bigr]_{y^{\lambda-\varrho}}.
\]

Мы объединяем оба разложения, строя для форм, каждая из которых имеет два ряда контрагредиентных переменных,
\[
s_x^m t_x^{n-\lambda}t_y^\lambda u_\sigma^\mu u_\tau^{\nu-\kappa}v_\tau^\kappa
\]
разложения по степеням $u_x$, коэффициентами которых являются составные поляры ряда форм $s_x^m t_x^n u_\sigma^\mu u_\tau^\nu$, взятые по переменным $x$ и $u$.

Достаточно построить ряд для двух контрагредиентных рядов символов (соответственно переменных), так как, объединяя каждые два ряда символов (переменных) в один новый ряд, формы с большим числом рядов символов (переменных) можно свести к этому случаю.

Чтобы добиться того, что все формы ряда форм содержат только по два ряда символов или переменных, специализируем:
\[
\begin{array}{ll}
\text{a) } \sigma=\widehat{st}, & \text{a') } y=\widehat{uv},\\[0.3em]
\text{b) } s=\widehat{\sigma\tau}, & \text{b') } v=\widehat{xy},
\end{array}
\]
и выполняем разложение для специализации a), a').

Прямым переносом разложения I получаются составные поляры для форм $(stu)^\mu s_x^m t_x^{n-\lambda}t_y^\lambda$, тогда как для форм $(stu)^\mu u_\tau^{\nu-\kappa}v_\tau^\kappa s_x^m t_x^{n-\lambda}t_y^\lambda$ вместо составных поляр возникают члены такого типа, вычисление которых требует введения всё новых вспомогательных переменных, приравниваемых лишь в конце; из-за этого вычисление излишне усложняется.

Поэтому мы выбираем косвенный путь: представляем составные поляры всех форм ряда форм, то есть выражения
\[
\bigl[s_\tau^\varrho(stu)^{\mu-\varrho+r}\bigr]_{y^\lambda}^{v^\kappa}
\]
через начальные члены этих поляр и, обращая возникающую рекурсивную систему уравнений, получаем требуемое разложение по полярам.

По принципу переноса для $\lambda$-й поляры формы $s_x^m t_x^n:[s_x^m t_x^n]_{y^\lambda}$ ($\lambda\le n$) справедливо следующее разложение (случай $\lambda>n$ сводится к первому посредством разложения $[s_y^m t_y^n]_{x^{m+n-\lambda}}$)\footnote{Gordan, \emph{Die Resultante binärer Formen} (Kap. I, § 5). Rend. Circ. Mat. di Palermo XXII. 1906.}:
\[
\bigl[s_x^m t_x^n\bigr]_{y^\lambda}
=\sum_r (-1)^r\,
\frac{\binom{m}{r}\binom{\lambda}{r}}{\binom{m+n}{r}}\,
(st\widehat{xy})^r s_x^{m-r}t_x^{n-\lambda}t_y^{\lambda-r},
\]
и соответственно для $\kappa$-й поляры формы $u_\sigma^\mu u_\tau^\nu:[u_\sigma^\mu u_\tau^\nu]_{v^\kappa}$
\[
\bigl[u_\sigma^\mu u_\tau^\nu\bigr]_{v^\kappa}
=\sum_\varrho (-1)^\varrho\,
\frac{\binom{\mu}{\varrho}\binom{\kappa}{\varrho}}{\binom{\mu+\nu}{\varrho}}\,
(\sigma\tau\widehat{uv})^\varrho u_\sigma^{\mu-\varrho}u_\tau^{\nu-\kappa}v_\tau^{\kappa-\varrho}.
\]

Умножением, с введением специализации a) и a'), получаем
\[
\begin{aligned}
\bigl[s_x^m t_x^n(stu)^\mu u_\tau^\nu\bigr]_{\widehat{uv}^{\lambda}}^{v^\kappa}
={}&
\sum_{r,\varrho}c_{r\varrho}\,
(st\widehat{uv})^r(st\widehat{uv})^\varrho
s_x^{m-r}t_x^{n-\lambda}(tuv)^{\lambda-r}(stu)^{\mu-\varrho}
u_\tau^{\nu-\kappa}v_\tau^{\kappa-\varrho},
\end{aligned}
\]
а отсюда, разлагая по степеням $u_x$ с выделением первого члена ($\sum'_{r,\varrho}$ означает, что член $r=0$, $\varrho=0$ следует опустить) и учитывая (a) на с. 26,
\[
\tag{II}
\begin{aligned}
&s_x^m t_x^{n-\lambda}(tuv)^\lambda(stu)^\mu u_\tau^{\nu-\kappa}v_\tau^\kappa\\
&\quad =
\bigl[s_x^m t_x^n(stu)^\mu u_\tau^\nu\bigr]_{\widehat{uv}^{\lambda}}^{v^\kappa}\\
&\qquad
-\sum_{r,\varrho}' c_{r\varrho}\,s_\tau^\varrho(stu)^{\mu-\varrho+r}
u_\tau^{\nu-\kappa}v_\tau^{\kappa-\varrho}
s_x^{m-r}t_x^{n-\lambda}(tuv)^{\lambda-r+\varrho}v_x^r\\
&\qquad
+u_x\sum_{\varrho}\sum_{r=1}^{\lambda} c_{r\varrho}\,
s_\tau^\varrho(stu)^{\mu-\varrho+r-1}(stv)
u_\tau^{\nu-\kappa}v_\tau^{\kappa-\varrho}
s_x^{m-r}t_x^{n-\lambda}(tuv)^{\lambda-r+\varrho}v_x^{r-1}\\
&\qquad
\vdots\\
&\qquad
-(-1)^\lambda u_x^\lambda\sum_\varrho c_{\lambda\varrho}\,
s_\tau^\varrho(stu)^{\mu-\varrho}(stv)^\lambda
u_\tau^{\nu-\kappa}v_\tau^{\kappa-\varrho}
s_x^{m-\lambda}t_x^{n-\lambda}(tuv)^\varrho,
\end{aligned}
\]
где
\[
c_{r\varrho}=(-1)^{r+\varrho}\cdot
\frac{\binom{m}{r}\binom{\lambda}{r}}{\binom{m+n}{r}}\cdot
\frac{\binom{\mu}{\varrho}\binom{\kappa}{\varrho}}{\binom{\mu+\nu}{\varrho}}
\qquad
\begin{array}{l}
\text{для }\lambda\le n,\\
\kappa\le\nu.
\end{array}
\]
и соответственно для остальных сочетаний чисел
\[\lambda,n;\kappa,\nu.\]
Следовательно, выражение
\[s_x^m t_x^{n-\lambda}(tuv)^\lambda(stu)^\mu u_\tau^{\nu-\kappa}v_\tau^\kappa\]
сведено к составной поляре и, с применением тождества
\[
(stv)u_\tau=(stu)v_\tau-s_\tau(tuv),
\]
к сумме аналогично построенных выражений, соответствующих более высоким формам ряда форм.

Итерацией приходим к разложению:
\[
\tag{III}
s_x^m t_x^{n-\lambda}(tuv)^\lambda(stu)^\mu u_\tau^{\nu-\kappa}v_\tau^\kappa
=
\sum_{p,r,\varrho}u_x^p\,C_{pr\varrho}\cdot
\bigl[s_\tau^\varrho(stu)^{\mu-\varrho+r}\bigr]_{\widehat{uv}^{\lambda-r+\varrho}v^{\kappa-\varrho+p},\,v_x^{r-p}}.
\]
Числовые коэффициенты $C_{pr\varrho}$ однозначно и рекурсивно вычисляются из коэффициентов $c_{r\varrho}$, $c'_{r\varrho}$ системы уравнений, возникающей при итерации (II.) (см. пример на с. 42). Аналогичные разложения получаются при остальных специализациях.
% END UNIT BODY
% END PRODUCER UNIT 004
\endgroup


\clearpage
\begingroup
\setcounter{footnote}{0}
% BEGIN PRODUCER UNIT 005
% Source: C:/Users/Floris/Documents/interlanguage/03_projects/noether/02_slavic_working_corpus/translations/paper02/russian/v001/Noether_Paper02_Section03_Russian_v001.tex
% Identity: 12653 B / SHA-256 601B188641B5CDAB4E0B55370EBD107CC16C0E45A6DFCF9CDB81C2F2FFCEA0A7
\setlength{\parindent}{1.25em}
\setlength{\parskip}{0pt}
\emergencystretch=4em
% BEGIN UNIT BODY
\begin{center}
\textbf{§ 3. \quad Редукционные теоремы.}
\end{center}

Известные теоремы о редукции форм и систем форм теперь надо связать с параграфами 1 и 2; для этого нужны некоторые определения, заимствованные из теории бинарных форм.

a) \emph{Относительно полной системой} по модулю заданного ряда форм мы называем систему форм со следующим свойством: все образования, возникающие при свёртке произвольного произведения этих форм, можно выразить как целые рациональные функции форм системы, с точностью до добавочного члена из выражений, возникших при свёртке системных форм с системой модуля\footnote{Gordan--Kerschensteiner, \emph{Vorlesungen über Invariantentheorie}, т. II, с. 227.}.

b) Форму будем называть \emph{редуцируемой}, если она выражается через формы, имеющие инварианты множителями, или через «более высокие формы»; то есть через более высокие формы полного ряда форм, которому она принадлежит, либо через формы, содержащие символы модуля в более высоком порядке.

Теоремы о редуцентах\footnote{Gordan, \emph{Math. Annalen}, т. XVII, с. 222.} в наиболее общей форме можно сформулировать так:

Определение: \emph{Редуцент есть редуцируемый ряд форм.}

\emph{Теорема II. Пусть начальная форма ряда форм редуцируема потому, что один из её членов возник при свёртке с редуцентом (то есть имеет редуцент множителем), и пусть конечная форма этого ряда форм возникла из того же самого члена при свёртке. Тогда весь ряд форм редуцируем}\footnote{Упрощение, которое даёт теорема II, можно видеть на следующем примере. Для специальной формы $f=x_1^3x_2+x_2^3x_3+x_3^3x_1$ выражение $a_y^2a_x^2u_y^2$ является редуцентом. Отсюда по теореме II следует редукция ряда форм
\[
\theta_\nu(\vartheta\nu x)\theta_x^3u_\vartheta u_\nu^2
=
a_\nu^2(abu)-a_\nu b_\nu(aba),
\]
так как $\theta_\nu^4=a_\nu^2b_\nu^2(abu)^2$ возникла из члена $a_\nu^2(abu)$ при свёртке. У Гордана, там же, с. 231, напротив, редукция форм
\[
\theta_\nu(\vartheta\nu x),\quad \theta_\nu(\vartheta\nu x)^2,\quad
\theta_\nu^2,\quad \theta_\nu^2(\vartheta\nu x),\quad
\theta_\nu^2(\vartheta\nu x)^2
\]
проводится по отдельности.}.

Доказательство: ряд форм возникает из начального члена при свёртке с самим собой, то есть, с одной стороны, при более высоких свёртках с редуцентом, а с другой -- при свёртке редуцента с самим собой. В обоих случаях, по § 2, все формы ряда форм можно представить как поляры (перекрывания) над рядом форм редуцента.

Переход от начальной формы ко всему ряду форм больше не действует для остальных методов редукции. Это следующие методы:

1) так называемая «двойная редукция».

Под этим мы понимаем редукцию выражения, имеющего множителями два независимых друг от друга редуцента, двумя способами; вследствие этого возникают соотношения между более высокими формами, к которым выражение было сведено.

Систематическим применением «двойной редукции», с одной стороны, надо получить все соотношения\footnote{Так, например, с помощью «двойной редукции» были найдены: редукция формы $a_\eta^4$ (§ 10), единственной формы, которую Maisano излишне включил в систему; затем -- упомянутые в примечании к введению соотношения между 3 ковариантами и 3 инвариантами.}; с другой стороны, если «двойная редукция», применённая к различным выражениям, не даёт новых соотношений, мы имеем одновременно критерий независимости более высоких форм и контроль вычисления.

2) \emph{Свёртка с распадающимися формами.}

Под «распадающимися формами», с небольшим обобщением обычного понятия, будем понимать формы, которые выражаются через произведения форм меньшей степени и через «более высокие» формы. Произведения форм при однократной свёртке переходят в произведения; так же и произведения нелинейных форм при двукратной свёртке в специализациях a') и b') (§ 2), по теореме произведения:
\[
a_yb_y a_xb_x=\frac12a_y^2b_x^2+\frac12b_y^2a_x^2-\frac12(ab\widehat{yx})^2.
\]
Следовательно, первые поляры (перекрывания) и некоторые вторые поляры распадающихся форм снова являются распадающимися формами. Отсюда следует:

Член однократного (соответственно двукратного) перекрывания над распадающейся формой, возникший при однократной (соответственно двукратной) свёртке с такой формой, сам становится распадающейся формой:

a) если следующие более высокие формы в ряду форм распадающейся формы распадаются или становятся редуцируемыми, либо если при однократной свёртке наступают специализации $y=\widehat{uv}$ или $v=\widehat{xy}$;

b) если остальные однократные свёртки приводят к более высоким формам (ср. § 12, B. $H_k$ и $H_k(k\eta x)$).

Такой член перекрывания также может распасться:

c) если остальные однократные свёртки приводят к более низким формам, которые, независимо от свёртки с распадающейся формой, сами распадаются, то есть выражаются через произведения и через редуцируемую форму. При этом каждый раз надо сначала вычислением убедиться, что числовой коэффициент соответствующего члена не исчезает тождественно. Это не что иное, как «двойная редукция» низшей формы (ср. § 23, C. $H_q(\theta H u)(\theta s u)$).

Распадающиеся формы возникают при всех однократных свёртках с функциональными детерминантами\footnote{Ср. Gordan, там же, § 3.}, а кроме того -- при некоторых более высоких свёртках. Имеем
\[
(st\widehat{yx})s_y=\frac12\,t\,s_y^2-\frac12\,s\,t_y^2+\frac12(st\widehat{yx})^2,
\]
\[
(st\widehat{yx})^2s_y^2=\frac13\,t\,s_y^3-\frac13\,s\,t_y^3+s_y(st\widehat{yx})^2-\frac13(st\widehat{yx})^3.
\]
Аналогичные формулы справедливы для дуалистических форм
\[
(\sigma\tau\widehat{\nu u})v_\sigma
\quad\text{и}\quad
(\sigma\tau\widehat{\nu u})v_\sigma^2.
\]

Из теорем этого параграфа для построения всех нередуцируемых форм, возникающих из свёртки $S$ с $T$, получается следующее правило:

Начиная с самых низких свёрток, надо продвигаться в построении ряда форм $ST$ по любому заранее зафиксированному пути (например, построчно или симметрично к диагональному члену), пока не дойдём до формы, имеющей редуцент множителем. Ряд форм, определённый этой формой, следует опустить, как только конечная форма удовлетворяет условию, установленному в теореме II. После исключения всех редуцируемых рядов форм нужно, применяя двойную редукцию, удалить все остальные редуцируемые формы, а также все распадающиеся формы. Места полного ряда форм, дважды покрытые независимыми друг от друга редуцируемыми рядами форм, дают повод к редукции более высоких форм (ср. § 11, D).

\clearpage
\noindent\textbf{Теорема III.}
\emph{Формы, заимствованные из бинарной области, то есть те формы, которые возникают из системных форм соответствующей бинарной формы по принципу переноса Клебша путём окаймления, образуют относительно полную систему по модулю $(abc)$.}

\noindent\textbf{Доказательство:}
Бинарному соотношению, утверждающему, что формы $Q'_1,\ldots,Q'_n$ образуют полную систему бинарной основной формы:
\[
Q'-F(Q')=0=\sum_\lambda\bigl\{(ab)c_x+(bc)a_x+(ca)b_x\bigr\}\varphi_\lambda^{*}
\]
путём окаймления соответствует тернарное соотношение:
\[
Q-F(Q_i)=\sum_\lambda u_x\cdot(abc)\varphi_\lambda .
\]
Но это и есть определяющее равенство для относительно полной системы по модулю $(abc)$. Для биквадратичной формы система заимствованных форм состоит из следующих форм:
\[
\begin{gathered}
f=a_x^4,\qquad
\theta=\theta_x^4u_\vartheta^2=(abu)^2a_x^2b_x^2,\qquad
K=K_x^6u_k^3=(a\theta u)u_\vartheta^2a_x^3\theta_x^3,\\
j=u_j^6=(a\theta u)^4u_\vartheta^2,\qquad
\nu=u_\nu^4=(abu)^4,
\end{gathered}
\]
среди которых первые четыре образуют относительно полную систему по модулю $((abc),\nu)$.
% END UNIT BODY
% END PRODUCER UNIT 005
\endgroup


\clearpage
\begingroup
\setcounter{footnote}{0}
% BEGIN PRODUCER UNIT 006
% Source: C:/Users/Floris/Documents/interlanguage/03_projects/noether/02_slavic_working_corpus/translations/paper02/russian/v001/Noether_Paper02_Section05_Russian_v001.tex
% Identity: 9411 B / SHA-256 86D23CF5D810E53CA1CB43D79939F70AC8D794F013FA240829A4028232D8A797
\setlength{\parindent}{1.25em}
\setlength{\parskip}{0pt}
\emergencystretch=4em
\let\A\relax
\newcommand{\A}{\mathfrak A}
% BEGIN UNIT BODY
\begin{center}
\textbf{§ 5. \quad Сведение модуля $(abc)$ к модулям $\A$ и $\nu$.}
\end{center}

Модуль $(abc)$, заданный лишь одним скобочным множителем, в соответствии с определением (§ 3, a) надо свести к формам системы. Иными словами: ряд форм $(abc)$ надо однозначно нормировать (ср. § 1).

\[
a_s^2a_xu_x=\frac13 a_s^3u_x=\frac13 S\cdot u_x. \tag{2}
\]
\[
\left.
\begin{aligned}
(a\A u)^2&=a_s-\frac{S}{6}u_x^2,\\
(a\A u)^3&=0.
\end{aligned}
\right\} \tag{3}
\]
\[
\left.
\begin{aligned}
\theta(s)=(ss,x)^2&=2a_t+\frac13 S\cdot\theta-\frac13 Tu_x^2,\\
\theta(t)=(tt,x)^2&=-\frac13 S\cdot a_t+\frac23 T\cdot a_s+\frac1{18}S^2\cdot\theta-\frac1{18}u_x^2\cdot S\cdot T.
\end{aligned}
\right\} \tag{4}
\]

Ряд модулей: $(abc)$; $\A$; $s$; $t$ (система модуля $t$ состоит из единственной формы $t$).\footnote{Gordan--Kerschensteiner, там же, с. 134.}

Окажется, что форм
\[
\A=(abc)^2a_x^2b_x^2c_x^2,\qquad \nu=(abu)^4
\]
как модуля достаточно, то есть что формы
\[
\begin{gathered}
\A=(abc)^2a_x^2b_x^2c_x^2,\qquad
a_\nu=(abc)(bcu)^3a_x^3,\\
a_\nu^2=(abc)^2(bcu)^2u_x^2,\qquad
i=a_\nu^4=(abc)^4
\end{gathered}
\]
задают ряд форм $(abc)$. В ряду форм $a_\nu$ форма $a_\nu^3$ отсутствует в силу тождества
\[
a_\nu^3=(abc)^3a_x(bcu)=\frac13 u_x\cdot(abc)^4=\frac13 i\cdot u_x.
\]

Чтобы свести модуль к более высокому, по определению надо свести самые общие свёртки или все специальные рассматриваемые свёртки с этим модулем к свёрткам с более высоким модулем.

В нашем случае самые общие свёртки возникают из выражений
\[
(abc)a_x^3b_y^3c_z^3,\qquad
(abc)a_x^2b_y^2c_z^2a_zb_xc_x
\]
(заимствованных из кубических форм),
\[
(abc)a_xb_yc_za_x^2b_x^2c_x^2
\]
(заимствованного из квадратичных форм) и путём поляризации этих выражений.

Чтобы вычислить эти выражения через поляры $\A$ и ряда форм $a_\nu$, соответственно через их начальные члены, мы, как и в § 2, выбираем косвенный путь. Разлагаем полярные члены
\[
(abc)^2a_x^2b_y^2c_z^2,\qquad
a_\nu(\nu xy)(\nu yz)(\nu zx)a_xa_ya_z
=(abc)(bc\widehat{x}u)(bc\widehat{y}z)(bc\widehat{z}x)a_xa_ya_z\quad\text{и т. д.}
\]
как линейные комбинации выражений с символическим множителем $(abc)$ и, обращая систему уравнений, получаем искомые соотношения, а также соотношение между полярами, взятыми по различным комбинациям переменных. Для вычисления служат теорема тождества и теорема произведения для детерминантов.

Система уравнений для $(abc)a_x^3b_y^3c_z^3$ имеет вид:
\begingroup\small
\[
\begin{aligned}
1)\quad
\sum_3 a_\nu(\nu yz)^3a_x^3
&=6(abc)a_x^3b_y^3c_z^3
-6\sum_3(abc)a_x^3b_y^2b_zc_z^2c_y^{*},\\[0.4em]
2)\quad
3(abc)^2a_x^2b_y^2c_z^2\cdot(xyz)
-\frac12\sum_3 a_\nu^2(\nu yz)^2\nu_x^2(xyz)
&=2\sum_3(abc)a_x^3b_y^2b_zc_z^2c_y\\
&\quad-4\sum_3(abc)a_xa_ya_zb_y^2b_zc_z^2c_x,\\[0.4em]
3)\quad
a_\nu(\nu xy)(\nu yz)(\nu zx)a_xa_ya_z
&=2\sum_3(abc)a_xa_ya_zb_y^2b_zc_z^2c_x,\\[0.4em]
4)\quad
\sum_3 a_\nu(\nu yz)^3u_x^3
-\frac32\sum_3 a_\nu^2(\nu yz)^2\nu_x^2(xyz)
+\frac16 i\cdot(xyz)^3
&=6\sum_3(abc)a_xa_ya_zb_y^2b_zc_z^2c_x.
\end{aligned}
\]
\endgroup

Отсюда для
\[
(abc)a_x^3b_y^3c_z^3,
\]
и аналогичным вычислением для
\[
(abc)a_x^2b_y^2c_z^2a_zb_xc_x,\qquad
(abc)a_xb_yc_za_z^2b_x^2c_x^2
\]
получаем:
\begingroup\small
\[
\begin{aligned}
(abc)a_x^3b_y^3c_z^3
&=\frac32(abc)^2a_x^2b_y^2c_z^2(xyz)
+\frac32 a_\nu(\nu xy)(\nu yz)(\nu zx)a_xa_ya_z\\
&\quad-\frac1{12}i\cdot(xyz)^3,\\[0.4em]
(\text{IV.})\quad
(abc)a_x^2b_y^2c_z^2a_zb_xc_x
&=\frac23(abc)^2a_x^2b_y^2c_z^2c_x\cdot(xyz)\\
&\quad+\frac13a_\nu(\nu xy)(\nu yz)(\nu zx)a_x^3
-\frac16 a_\nu^2(\nu xy)(\nu zx)a_x^2\cdot(xyz),\\[0.4em]
(abc)a_xb_yc_za_z^2b_x^2c_x^2
&=\frac16(abc)^2a_x^2b_z^2c_z^2\cdot(xyz).
\end{aligned}
\]
\endgroup

Тем самым мы получили относительно полную систему по модулю $(\A,\nu)$, состоящую из четырёх форм: $f,\theta,K,j$. Отсюда следует: все перекрывания $f$ над $\theta$, не заимствованные из бинарной области, так же как перекрывание $(a\theta u)^2$, редуцируемое в бинарной области к $(ab)^4$, выражаются через символы $\A$ и $\nu$.\footnote{$\sum_3$ относится к сумме трёх членов, возникающих из начального члена при циклической перестановке переменных. Уравнения справедливы также отдельно для каждого члена суммы.}

Основные для дальнейшего редукционные формулы мы получаем специализацией разложения IV или, короче, прямым вычислением (перестановкой символов) для $a_\vartheta(a\theta u)^2$, $a_\vartheta^2(a\theta u)^2$, $(a\theta u)^2$ по следующим исходным формулам:
\[
\begin{aligned}
a_\vartheta(a\theta u)^2
&=(abc)(bcu)(abu)^2-\frac13(abc)(bcu)\{(bcu)a_x-u_x(abc)\}^2,\\
a_\vartheta^2(a\theta u)^2
&=(abc)^2(abu)^2-\frac13(abc)^2\{(bcu)a_x-u_x(abc)\}^2,
\end{aligned}
\]
для $((a\theta u)^2)^{*}$:
\[
\begin{aligned}
(abu)a_y^3b_x^3&=\frac32u_\vartheta(\vartheta yx)\theta_y^2+\frac14(\nu yx)^3,\\
((abu)(acu))b_x^3&=\frac32u_\vartheta(\vartheta\widehat{a}ux)(\theta au)^2+\frac14(\nu\widehat{a}ux)^3.
\end{aligned}
\]
\begingroup\small
\[
\left.
\begin{aligned}
(a\theta u)^2&=\frac16\nu\cdot f-\frac23u_x\cdot a_\nu
+\frac23u_x^2\cdot a_\nu^2-\frac{i}{18}u_x^4,\\
a_\vartheta&=\frac13u_x\cdot\A,\\
a_\vartheta(a\theta u)&=0,\\
a_\vartheta^2&=\A,\\
(a\theta u)^3&=0,\\
a_\vartheta(a\theta u)^2&=-\frac56a_\nu+\frac76u_x\cdot a_\nu^2-\frac{i}{9}u_x^3,\\
a_\vartheta^2(a\theta u)&=0,\\
(a\theta u)^4&=j,\\
a_\vartheta(a\theta u)^3&=0,\\
a_\vartheta^2(a\theta u)^2&=\frac23a_\nu^2-\frac{i}{9}u_x^2.
\end{aligned}
\right\} \tag{1}
\]
\endgroup

К этому добавим формулу, также возникшую из редукции модуля $(abc)$:
\[
a_\nu^3a_xu_\nu=\frac13 i\cdot u_x. \tag{2}
\]

\emph{Из формул (1.) и (2.) и из аналогично построенных формул для основной формы $\nu$ посредством поляризации выводятся все последующие редукционные формулы}, за исключением соотношений для распадающихся форм. Из формул (1.) видим:

Ряд форм приводит:
\[
\begin{array}{rcl}
a_\vartheta(a\theta u)^2 &\text{к}& \text{одним символам }\nu,\\[0.2em]
a_\vartheta &\text{к}& \text{символам }\A\text{ и }\nu,\\[0.2em]
(a\theta u)^2 &\text{к}& \text{символам }j,\A\text{ и }\nu,\\[0.2em]
(a\theta u)^2\text{ при свёртке I} &\text{к}& \text{символам }j\text{ и }\nu,\\[0.2em]
(a\theta u)^2\text{ при свёртке II} &\text{к}& \text{символам }\A\text{ и }\nu.
\end{array}
\]

Итак, можно заменять:
\begin{enumerate}
\item по модулю $(\A,\nu)$: свёртки с $j$ соответствующими свёртками с $(a\theta u)^2$ или $(a\theta u)^3$;
\item по модулю $(\nu)$: свёртки с $\A$ соответствующими свёртками с $a_\vartheta$ или $a_\vartheta(a\theta u)$, а также специальными свёртками с $(a\theta u)^2$ (которые ведут только к однократной свёртке I).
\end{enumerate}
В частности:
\begingroup\small
\[
\begin{aligned}
(a\theta u)^2\theta_y^2\theta_x^2&=\frac13(jyx)^2\pmod{\nu},&
(a\theta u)^2u_y\theta_y&=-\frac16(jyx)^2\pmod{\nu},\\
(a\theta u)^2\nu_y^2&=\frac13(\A u\nu)^2\pmod{\nu},&
(a\theta u)(a\theta\nu)u_\vartheta\nu_\vartheta&=-\frac16(\A u\nu)^2\pmod{\nu},\\
(a\theta u)^2u_\nu^2\vartheta_\nu^2&=\frac13(\A u\nu)^2\A_\nu\pmod{\nu}.
\end{aligned}
\]
\endgroup
% END UNIT BODY
% END PRODUCER UNIT 006
\endgroup


\clearpage
\begingroup
\setcounter{footnote}{0}
% BEGIN PRODUCER UNIT 007
% Source: C:/Users/Floris/Documents/interlanguage/03_projects/noether/02_slavic_working_corpus/translations/paper02/russian/v001/Noether_Paper02_Section06_Russian_v001.tex
% Identity: 11136 B / SHA-256 BEAA14315D2B34CDB8EDF1946225144A90B302E7FAC95014652C62F314F6A02F
\setlength{\parindent}{1.25em}
\setlength{\parskip}{0pt}
\emergencystretch=4em
\allowdisplaybreaks
\let\A\relax
\newcommand{\A}{\mathfrak A}
% BEGIN UNIT BODY
\begin{center}
\textbf{§ 6. \quad Сведение модуля $(\A,\nu)$ к модулю $(\nu)$.}
\end{center}

Редукционные формулы предыдущего параграфа дают нам средство прямо построить относительно полную систему по модулю $\nu$; мы увидим, что к системе заимствованных форм нужно добавить только формы $\A$ и $(a\A u)=N$ (аналогично кубической форме и вообще как общее правило).

Для редукции модуля $\A$ рассмотрим все специальные свёртки, которые здесь нужно учитывать, то есть свёртки относительно полной системы по модулю $(\A,\nu)$ с системой $\A$, и начнём с форм, самых низких по коэффициентам, а именно со свёрток $f$ с $\A$.

Чтобы редуцировать формы $(a\A u)^2$, $(a\A u)^3$, $(a\A u)^4$, по § 5 заменяем символы $\A$ более низкими формами ряда форм, то есть раскладываем выражения
\[
a_\vartheta b_\vartheta(b\theta u)^2,\qquad
a_\vartheta(a\theta u)b_\vartheta(b\theta u)^2,\qquad
a_\vartheta(a\theta u)b_\vartheta(b\theta u)^3
\]
1) по полярам ряда форм $a_\vartheta$ (символы $\A$ и $\nu$),

2) по полярам ряда форм $b_\vartheta(b\theta u)^2$ (символы $\nu$)

и получаем редукционные формулы (вычисление см. ниже):
\begin{align}
\text{a)}\quad (a\A u)^2
&=\frac12(\vartheta\nu x)^2-\frac25 f\cdot a_\nu^2
-\frac45u_x\cdot\theta_\nu(\vartheta\nu x)^2 \notag\\
&\quad+u_x^2\left\{\frac16 i\cdot f-\frac1{40}S\right\},\notag\\
\text{b)}\quad (a\A u)^3
&=-\frac3{10}\theta_\nu(\vartheta\nu x),\tag{3}\\
\text{c)}\quad (a\A u)^4
&=\frac{21}{10}\theta_\nu^2-\frac3{10}\Pi
-\frac65u_x\cdot\theta_\nu^3+\frac1{10}u_x^2\cdot\theta .\notag
\end{align}

(К тем же результатам приводит также двойная редукция выражений $a_\vartheta^2(b\theta u)^2$, $a_\vartheta^2(b\theta u)^3$, $a_\vartheta^2((a\theta u)^2(b\theta u)^2)$ и $a_\vartheta^2((a\theta u)(b\theta u))^3$ после исключения $a_\vartheta^2$.)

Из формул (3.) следует, что $(a\A u)^2$ является редуцентом относительно модуля $\nu$. Отсюда получаем:

1) Редукция системы $\A$: ряд форм $(\A\A' u)^2=a_\vartheta^2((a\A u)^2)$ имеет редуцент $(a\A u)^2$ множителем.

2) Редукция системы, возникшей свёрткой с модулем $\A$:

ряд форм $(\theta\A u)^2$ имеет редуцент $(a\A u)^2$ множителем.

Для форм $(\theta\A u)$ и $\A_\vartheta$ получаем:
\[
(a\A u)^2(abu)=(\theta\A u)-u_x\cdot\A_\vartheta(\theta\A u),
\]
\[
(a\A u)(a\A b)=-\frac12\A_\vartheta+\frac12u_x\cdot\A_\vartheta^2.
\]

Из редуцента $(\theta\A u)$ следует редукция системы, возникающей при свёртке с модулем $\A$.

\bigskip

Итак, относительно полная система по модулю $\nu$ состоит из шести форм:
\[
f,\ \theta,\ K,\ j,\ \A,\ N.
\]

В качестве примера разложения в ряд из § 2 подробно выведем формулу (3.)a; в дальнейших вычислениях будем сразу записывать конечную формулу III. (Поляры исчезающих форм уже опущены в ходе вычисления; первая строка даёт значения, подставленные по разложению II, вторая строка -- конечные значения, найденные рекурсивно, начиная с последнего уравнения системы.) По разложению II следует:

\begingroup\scriptsize
\begin{align*}
\text{1)}\quad&m=3,\quad n=4,\quad \lambda=2,\quad \mu=0,\quad \nu=\chi=1,\\
a_\vartheta b_\vartheta(b\theta u)^2
&=[a_\vartheta]_{b^2}b+\frac67\{a_\vartheta b_\vartheta(a\theta u)(b\theta u)-u_xa_\vartheta b_\vartheta(a\theta b)(b\theta u)\}\\
&\quad-\frac17\{a_\vartheta(a\theta u)^2b_\vartheta-2u_xa_\vartheta(a\theta u)(a\theta b)b_\vartheta+u_x^2a_\vartheta(a\theta b)^2b_\vartheta\}\\
&=\frac23(a\Delta u)^2+\frac15[a_\vartheta(a\theta u)^2]b+\frac1{30}f[a_\vartheta^2(a\theta u)^2]-\frac25u_x[a_\vartheta(a\theta u)^2]b^2\\
&\quad-\frac1{15}u_x[a_\vartheta^2(a\theta u)^2]b+\frac15u_x^2[a_\vartheta(a\theta u)^2]b^3+\frac1{30}u_x^2[a_\vartheta^2(a\theta u)^2]b^2;\\[.6em]
\text{2)}\quad&m=2,\quad n=3,\quad \lambda=1,\quad \mu=1,\quad \nu=\chi=1,\\
a_\vartheta(a\theta u)b_\vartheta(b\theta u)
&=\frac25\{a_\vartheta(a\theta u)^2b_\vartheta-u_xa_\vartheta(a\theta u)(a\theta b)b_\vartheta\}+\frac12a_\vartheta^2(b\theta u)^2\\
&\quad-\frac15\{a_\vartheta^2(b\theta u)(a\theta u)-u_xa_\vartheta^2(a\theta b)(b\theta u)\}\\
&=\frac12(a\Delta u)^2+\frac25[a_\vartheta(a\theta u)^2]b+\frac1{15}[a_\vartheta^2(a\theta u)^2]f\\
&\quad-\frac25u_x[a_\vartheta(a\theta u)^2]b^2-\frac1{10}u_x[a_\vartheta^2(a\theta u)^2]b+\frac1{30}u_x^2[a_\vartheta^2(a\theta u)^2]b^2;\\[.6em]
\text{3)}\quad&m=2,\quad n=3,\quad \lambda=1,\quad \mu=1,\quad \nu=1,\quad \chi=2,\\
a_\vartheta(a\theta b)b_\vartheta(b\theta u)
&=\frac25\{a_\vartheta(a\theta b)(a\theta u)b_\vartheta-u_xa_\vartheta(a\theta b)^2b_\vartheta\}\\
&=\frac25[a_\vartheta(a\theta u)^2]b^2+\frac1{30}[a_\vartheta^2(a\theta u)^2]b-\frac25u_x[a_\vartheta(a\theta u)^2]b^3-\frac1{30}u_x[a_\vartheta^2(a\theta u)^2]b^2;\\[.6em]
\text{4)}\quad&m=1,\quad n=2,\quad \lambda=0,\quad \mu=2,\quad \nu=\chi=1,\\
a_\vartheta(a\theta u)^2b_\vartheta
&=[a_\vartheta(a\theta u)^2]b+\frac23a_\vartheta^2(a\theta u)(b\theta u)\\
&=[a_\vartheta(a\theta u)^2]b+\frac16[a_\vartheta^2(a\theta u)^2]f-\frac16u_x[a_\vartheta^2(a\theta u)^2]b;\\[.6em]
\text{5)}\quad&m=1,\quad n=2,\quad \lambda=0,\quad \mu=2,\quad \nu=1,\quad \chi=2,\\
a_\vartheta(a\theta u)(a\theta b)b_\vartheta
&=[a_\vartheta(a\theta u)^2]b^2+\frac13a_\vartheta^2(a\theta b)(b\theta u)\\
&=[a_\vartheta(a\theta u)^2]b^2+\frac1{12}[a_\vartheta^2(a\theta u)^2]b-\frac1{12}u_x[a_\vartheta^2(a\theta u)^2]b^2;\\[.6em]
\text{6)}\quad&m=1,\quad n=2,\quad \lambda=0,\quad \mu=2,\quad \nu=1,\quad \chi=3,\\
a_\vartheta(a\theta b)^2b_\vartheta&=[a_\vartheta(a\theta u)^2]b^3;\\[.6em]
\text{7)}\quad&m=2,\quad n=4,\quad \lambda=2,\quad \mu=\nu=\chi=0,\quad (\text{по разложению I}),\\
a_\vartheta^2(b\theta u)^2&=[a_\vartheta^2]_{ub^2}+\frac1{10}\{[a_\vartheta^2(a\theta u)^2]f-2u_x[a_\vartheta^2(a\theta u)^2]b+u_x^2[a_\vartheta^2(a\theta u)^2]b^2\};\\[.6em]
\text{8)}\quad&m=1,\quad n=3,\quad \lambda=1,\quad \mu=1,\quad \nu=\chi=0,\quad (\text{разложение I}),\\
a_\vartheta^2(a\theta u)(b\theta u)&=\frac14[a_\vartheta^2(a\theta u)^2]f-\frac14u_x[a_\vartheta^2(a\theta u)^2]b;\\[.6em]
\text{9)}\quad&m=1,\quad n=3,\quad \lambda=1,\quad \mu=\chi=1,\quad \nu=0,\quad (\text{разложение I}),\\
a_\vartheta^2(a\theta b)(b\theta u)&=\frac14[a_\vartheta^2(a\theta u)^2]b-\frac14u_x[a_\vartheta^2(a\theta u)^2]b^2.
\end{align*}
\endgroup

Из 1) и 4) следует:
\begingroup\scriptsize
\[
\begin{aligned}
(a\Delta u)^2
&=\frac65[a_\vartheta(a\theta u)^2]b+\frac15f[a_\vartheta^2(a\theta u)^2]+\frac35u_x[a_\vartheta(a\theta u)^2]b^2-\frac3{20}u_x[a_\vartheta^2(a\theta u)^2]b\\
&\quad-\frac3{10}u_x^2[a_\vartheta(a\theta u)^2]b^3-\frac1{20}u_x^2[a_\vartheta^2(a\theta u)^2]b^2,\\
(a\Delta u)^2
&=-a_\nu b_\nu+\frac35fa_\nu^2+\frac45u_xa_\nu^2b_\nu-\frac3{20}u_x^2a_\nu^2b_\nu^2-\frac1{12}u_x^2if.
\end{aligned}
\]
\endgroup

(О пересчёте в символы $\theta$ см. \S\ 8 и \S\ 9.) Формулу (3.)a можно было бы вычислить ещё короче прямым переносом разложения I с введением вспомогательных переменных $w=x\widehat{u}b$, тогда как уже для формул (3.)b и (3.)c выбранный нами путь становится проще; для (3.)b заранее известно, по \S\ 9, что все члены с множителем $u_x$ надо опустить. Для вычисления (3.)b и (3.)c получаем:
\begingroup\scriptsize
\[
\begin{aligned}
(3.)\mathrm{b)}\quad 1)&\quad a_\vartheta(a\theta u)b_\vartheta(b\theta u)^2=\frac12(a\Delta u)^3+\frac45[a_\vartheta(a\theta u)^2]_{ub}b+\frac7{360}[a_\vartheta^2(a\theta u)^2]_{ub^2},\\
2)&\quad a_\vartheta(a\theta u)b_\vartheta(b\theta u)^2=[a_\vartheta(a\theta u)^2]_{ub}b+\frac{13}{36}[a_\vartheta^2(a\theta u)^2]_{ub^2};\\[.5em]
(3.)\mathrm{c)}\quad 1)&\quad a_\vartheta(a\theta u)b_\vartheta(b\theta u)^3=\frac12(a\Delta u)^4+\frac65[a_\vartheta(a\theta u)^2]_{ub^2}b+\frac{53}{120}[a_\vartheta^2(a\theta u)^2]_{ub^3}\\
&\qquad-\frac65u_x[a_\vartheta(a\theta u)^2]_{ub^3}b^2-\frac35u_x[a_\vartheta^2(a\theta u)^2]_{ub^2}b+\frac{19}{120}u_x^2[a_\vartheta^2(a\theta u)^2]_{ub^3}b^2,\\
2)&\quad a_\vartheta(a\theta u)b_\vartheta(b\theta u)^3=\frac34[a_\vartheta^2(a\theta u)^2]_{ub^2}.
\end{aligned}
\]
\endgroup

Примечание. Аналогично вычисляются редуцируемые по § 5 перекрывания $f$ над $j$. Получаем
\[
\tag{4}
\begin{aligned}
g_\nu&=-\theta_\nu+u_x\left(\frac92\theta_\nu^2-\frac12H\right)-3u_x^2\theta_\nu^3+\frac12u_x^3\theta,\\
g_\nu^2&=\frac{21}{10}\theta_\nu^2-\frac3{10}H-2u_x\theta_\nu^3+\frac12u_x^2\theta,\\
g_\nu^3&=-\frac7{10}\theta_\nu^3+\frac12u_x^2\theta,\qquad g_\nu^4=\frac35\theta.
\end{aligned}
\]

\[
\tag{4}
g_\nu^3=-\frac7{10}\theta_\nu^3+\frac12u_x^2\theta,
\qquad
g_\nu^4=\frac35\theta.
\]

Из (3.) и (4.) следует, переносом из бинарной области,
\[
(\theta\theta'u)^2=\frac13\,j\cdot f\pmod{\nu}.\footnotemark
\]
Остальные формы ряда форм $(\theta\theta'u)^2$ и форма $\theta'_\nu$ редуцируемы к символам $\nu$. Значит, можно заменять:
\[
\bmod\ \nu:\quad \text{свёртки с } f\cdot j
\text{ соответствующими с }(\theta\theta'u)^2.
\]
Далее имеем:
\[
(\vartheta\vartheta'x)^2=\frac43 f\cdot\A,
\qquad
\theta_{g\nu}(\vartheta\vartheta'x)=-N.
\]
\footnotetext{Gordan--Kerschensteiner, там же, с. 181.}
% END UNIT BODY
% END PRODUCER UNIT 007
\endgroup


\clearpage
\begingroup
\setcounter{footnote}{0}
% BEGIN PRODUCER UNIT 008
% Source: C:/Users/Floris/Documents/interlanguage/03_projects/noether/02_slavic_working_corpus/translations/paper02/russian/v001/Noether_Paper02_Section07_Russian_v001.tex
% Identity: 7497 B / SHA-256 F074473FF9AC70BB94995253FEADAFAF649AF13D423F1E5CF19059EFFA294C37
\setlength{\parindent}{1.25em}
\setlength{\parskip}{0pt}
\emergencystretch=4em
\allowdisplaybreaks
\let\A\relax
\newcommand{\A}{\mathfrak A}
% BEGIN UNIT BODY
\begin{center}
\textbf{Глава III. \quad Относительно полная система по модулю $(s,\varrho,t)$.}
\end{center}

\begin{center}
\textbf{\S\ 7. \quad Обзор построения системы.}
\end{center}

Чтобы перейти от относительно полной системы по модулю $\nu$ к относительно полной системе со следующим более высоким модулем, согласно определению \S\ 3 надо построить систему формы $\nu$ относительно этого более высокого модуля и свернуть ее с формами относительно полной системы по модулю $\nu$. Но $\nu=u_\nu^4$, как контравариант четвертой степени, дуалистически противопоставлен форме $f$. Поэтому, рассматривая $\nu$ как основную форму, мы получаем относительно полную систему из шести форм по модулю $\nu(\nu)=(\nu\nu_1x)^4=s_x^4$.

Следовательно, для построения относительно полной системы по модулю $s$ (система III) надо перекрывать
\[
\begin{array}{ll}
\text{Система I:} & f,\ \theta,\ K,\ j,\ \A,\ N\quad \text{над}\vphantom{\dfrac11}\\[0.3em]
\text{Система II:} & \nu,\ H=(\nu\nu_1x)^2u_\nu^2u_{\nu_1}^2,\ L=(\nu\eta x)H_x^2u_\nu^3u_\eta^3,\\
& g=(\nu\eta x)^4H_x^2,\ \sigma=H_\eta^2u_\nu^2u_\eta^4,\ (\nu\sigma x).
\end{array}
\]
Ниже будет видно (формула (13.)), что форма $g$ редуцируема; значит, система II фактически состоит только из пяти форм.

В ходе вычисления квадратичные формы
\[
\varrho=u_\varrho^2=\theta_\nu^4u_\vartheta^2,
\qquad
 t=t_x^2=a_\eta^4H_x^2
\]
присоединяются как модули, так что получается относительно полная система по модулю $(s,\varrho,t)$; при этом модуль $(\varrho,t)$ считается более высоким, чем модуль $s$.

Отдельные формы упорядочиваются по порядку в коэффициентах. Будем считать свертку
\[
\theta_\nu(K_\nu,\theta_\nu,\ldots)
\]
более высокой, чем свертка
\[
H_y(H_y,L_y,\ldots).
\]
Тем самым, по \S\ 1 и \S\ 3b, порядок отдельных форм одного и того же порядка определяется однозначно.

Построение форм одинакового порядка проводится таблично:

I. Указывается, какие формы систем I и II надо свертывать между собой, чтобы получить заданный порядок в коэффициентах.

II. Выписываются нередуцируемые формы по схеме ряда форм.

III. Редуцируются редуцируемые или распадающиеся на множители ряды форм или формы, то есть те места, где полный ряд форм обрывается; этим доказывается, что все нередуцируемые образования исчерпываются указанными.

IV. Делаются выводы, а именно: 1) указываются вновь возникшие редуценты; 2) указываются те формы, которые не входят в произведениях с формами той же системы в свертку с формами другой системы.

Будет видно, что возникают только:

1) свертки одной формы системы I с одной формой системы II;

2) свертки степеней $\theta$ или $H$, либо двух форм одной системы с одной формой второй системы.

Получим следующую схему свертки:
\[
\begin{array}{r@{\ }c@{\ }l@{\qquad}c@{\qquad}r@{\ }c@{\ }l}
\multicolumn{3}{c}{\text{Система I}} & \text{свертывается с} & \multicolumn{3}{c}{\text{Система II}}\\[0.25em]
1. & \text{порядок} & f & & 2. & \text{порядок} & \nu\\
2. & \text{порядок} & \theta & & 4. & \text{порядок} & H\\
3. & \text{порядок} & K,j,\A & & 6. & \text{порядок} & L,\sigma\\
4. & \text{порядок} & N,\theta^2,f\cdot j & & 8. & \text{порядок} & (\nu\sigma x),H^2\\
5. & \text{порядок} & \theta\cdot K & & 10. & \text{порядок} & H\cdot L\\
6. & \text{порядок} & \theta^3,\A\cdot j & & 12. & \text{порядок} & H^3.
\end{array}
\]

Для контроля вычислений, кроме «двойной редукции», служат две специальные формы:
\[
\begin{array}{ll}
\text{I.}&\begin{array}{rlrlrl}
f&=\displaystyle\sum_3x_1^3x_2,& u_\nu^2&=\frac{i}{6}u_x^2,& s&=\frac{i}{3}f,\\
a_\eta^2&=-\frac19 i\theta,& H_\nu^2&=\frac{i}{6}\theta,& \sigma&=-\frac{i}{6}j,\\
g&=-\frac23i\Delta,& \varrho&=0,& t&=0,
\end{array}\\[1.2\baselineskip]
\text{II.}&\begin{array}{rlrlrl}
f&=\displaystyle\sum_3x_1^4,& s&=\frac43if,& a_\eta^2&=\frac29i\theta,\\
\Delta_\nu^2&=\frac1{15}i\theta,& \sigma&=\frac43ij,& g&=\frac43i\Delta,\\
\varrho&=0,& t&=0.
\end{array}
\end{array}
\]

Примечание. Формулам (1.), (2.), (3.), (4.) для основной формы $\nu$ соответствуют следующие:
\[
\tag{5}
\begin{array}{c|c|c}
(\nu\eta x)^2=A & H_\nu(\nu\eta x)=0 & H_\nu^2=\sigma\\[0.4em]
(\nu\eta x)^3=0 & H_\nu(\nu\eta x)^2=B & H_\nu^2(\nu\eta x)=0\\[0.4em]
(\nu\eta x)^4=g & H_\nu(\nu\eta x)^3=0 & H_\nu^2(\nu\eta x)^2=C,
\end{array}
\]
\[
A=\frac16s\nu-\frac23u_xs_\nu+\frac23u_x^2s_\nu^2-\frac{J}{18}u_x^4,
\quad B=-\frac56s_\nu+\frac76u_xs_\nu^2-\frac{J}{9}u_x^3,
\quad C=\frac23s_\nu^2-\frac{J}{9}u_x^2.
\]
\[
\tag{6}
s_\nu^3u_xs_\nu=\frac13u_xs_\nu^4=\frac13Ju_x.
\]

\[
\tag{7}
\begin{aligned}
\mathrm{a)}\quad (\nu\sigma x)^2
&=\frac12(Hsu)^2-\frac25\nu\cdot s_\nu^2-
\frac45u_x\,s_\eta(Hsu)^2
+u_x^2\left\{\frac16J\cdot\nu-\frac1{40}(ss'u)^4\right\},\\
\mathrm{b)}\quad (\nu\sigma x)^3&=-\frac3{10}s_\eta(Hsu),\\
\mathrm{c)}\quad (\nu\sigma x)^4&=\frac{21}{10}s_\eta^2-\frac3{10}Z-
\frac65u_xs_\eta^3+\frac1{10}u_x^2s_\eta^4.
\end{aligned}
\]
\[
\tag{8}
\begin{aligned}
\mathrm{a)}\quad g_\nu&=-s_\eta+u_x\left(\frac92s_\eta^2-\frac12Z\right)-3u_x^2s_\eta^3+\frac12u_x^3s_\eta^4,\\
\mathrm{b)}\quad g_\nu^2&=\frac{21}{10}s_\eta^2-\frac3{10}Z-2u_xs_\eta^3+\frac12u_x^2s_\eta^4,\\
\mathrm{c)}\quad g_\nu^3&=-\frac7{10}s_\eta^3+\frac12u_xs_\eta^4,\\
\mathrm{d)}\quad g_\nu^4&=\frac35s_\eta^4.
\end{aligned}
\]
% END UNIT BODY
% END PRODUCER UNIT 008
\endgroup


\clearpage
\begingroup
\setcounter{footnote}{0}
% BEGIN PRODUCER UNIT 009
% Source: C:/Users/Floris/Documents/interlanguage/03_projects/noether/02_slavic_working_corpus/translations/paper02/russian/v001/Noether_Paper02_Section08_Russian_v001.tex
% Identity: 1156 B / SHA-256 57C552C93DF872FDEEEC9769A30A853FEE3EE3B857C132C8C2A22CC32FF6C174
\setlength{\parindent}{1.25em}
\setlength{\parskip}{0pt}
\emergencystretch=4em
\allowdisplaybreaks
% BEGIN UNIT BODY
\begin{center}
\textbf{\S\ 8. \quad Формы третьего порядка (система III).}
\end{center}

\noindent\textbf{Свертка $f$ с $\nu$.}

\noindent\textit{Нередуцируемые формы:}
\[
\begin{array}{cccc}
a_\nu&\cdot&\cdot&\cdot\\
\cdot&a_\nu^2&\cdot&\cdot\\
\cdot&\cdot&\cdot&\cdot\\
\cdot&\cdot&\cdot&a_\nu^4=i.
\end{array}
\]

\noindent\textit{Редукция:}
\[
a_\nu^3=\frac13 i u_x\qquad\text{(формула (2.)).}
\]
\noindent\textit{Следствия:} 1) $a_\nu^3$ является редуцентом. 2) $f$ входит в свертку с $\nu$ только в произведениях с контрагредиентными формами $(j)$. То же самое относится к $\nu$ по отношению к $f$.
% END UNIT BODY
% END PRODUCER UNIT 009
\endgroup


\clearpage
\begingroup
\setcounter{footnote}{0}
% BEGIN PRODUCER UNIT 010
% Source: C:/Users/Floris/Documents/interlanguage/03_projects/noether/02_slavic_working_corpus/translations/paper02/russian/v001/Noether_Paper02_Section09_Russian_v001.tex
% Identity: 2247 B / SHA-256 EE382A74BD8F14D984D48822F99CF6C8D7074B7FD7436BC14055E30A4FC142B4
\setlength{\parindent}{1.25em}
\setlength{\parskip}{0pt}
\emergencystretch=4em
\allowdisplaybreaks
% BEGIN UNIT BODY
\begin{center}
\textbf{\S\ 9. \quad Формы четвертого порядка (система III).}
\end{center}

\noindent\textbf{Свертка $\theta$ с $\nu$.}

\noindent\textit{Нередуцируемые формы:}
\[
\begin{array}{cccc}
(\theta\nu x)&\theta_\nu&\cdot&\cdot\\
(\theta\nu x)^2&\theta_\nu(\theta\nu x)&\theta_\nu^2&\cdot\\
\cdot&\theta_\nu(\theta\nu x)^2&\cdot&\theta_\nu^3.
\end{array}
\]

\noindent\textit{Редукции:} Из редуцента $a_\nu^3$ посредством двойной редукции выражений $a_\nu^3(abu)$, $a_\nu^3b_\nu$, $a_\nu^3b_\nu(abu)$ по следующей исходной схеме:
\[
\begin{aligned}
\theta_\nu^2(\vartheta\nu x)
&=a_\nu^2(\widehat{a}b\nu x)(abu)-\frac13(\nu\nu_1x)^3\\
&=a_\nu b_\nu(\widehat{a}b\nu x)(abu)+\frac16(\nu\nu_1x)^3\\
&=a_\nu^3(abu)-a_\nu^2b_\nu(abu)
=2a_\nu^2b_\nu(abu)=\frac23a_\nu^3(abu),\\[0.4em]
\theta_\nu^2(\vartheta\nu x)^2
&=a_\nu^2(\widehat{a}b\nu x)^2-\frac13(\nu\nu_1x)^4\\
&=a_\nu b_\nu(\widehat{a}b\nu x)^2+\frac16(\nu\nu_1x)^4\\
&=if-2a_\nu^3b_\nu+a_\nu^2b_\nu^2-\frac13s\\
&=2a_\nu^3b_\nu-2a_\nu^2b_\nu^2+\frac16s\\
&=\frac23if-\frac23a_\nu^3b_\nu-\frac16s,\\[0.4em]
\theta_\nu^3(\vartheta\nu x)
&=a_\nu^2b_\nu(\widehat{a}b\nu x)(abu)=a_\nu^3b_\nu(abu).
\end{aligned}
\]

Отсюда получаются формулы:
\[
\tag{9}
\begin{array}{c|c|c}
\theta_\nu^2(\theta\nu x)=0
& \theta_\nu^3(\theta\nu x)=0
& \theta_\nu^4u_x^2=u_\varrho^2=\varrho\quad(\text{присоед. модуль, § 7}),\\[0.6em]
\theta_\nu^2(\theta\nu x)^2=\dfrac49 i f-\dfrac16s
& &
\end{array}
\]

\noindent\textit{Следствия:} 1) $\theta_\nu^2(\theta\nu x)^2$ является редуцентом. 2) $\nu$ входит в свертку с $\theta$ только в произведениях с контрагредиентными формами $(g)$.
% END UNIT BODY
% END PRODUCER UNIT 010
\endgroup


\clearpage
\begingroup
\setcounter{footnote}{0}
% BEGIN PRODUCER UNIT 011
% Source: C:/Users/Floris/Documents/interlanguage/03_projects/noether/02_slavic_working_corpus/translations/paper02/russian/v001/Noether_Paper02_Section10_Russian_v001.tex
% Identity: 5912 B / SHA-256 3E5BC867F51F7FBC44CB946B5BB9682EAB2B485F11A7D7163682617902F4E6DF
\setlength{\parindent}{1.25em}
\setlength{\parskip}{0pt}
\emergencystretch=4em
\allowdisplaybreaks
\let\A\relax
\newcommand{\A}{\mathfrak A}
\newcommand{\R}{\mathfrak R}
\newcommand{\hata}{\widehat{a}}
% BEGIN UNIT BODY
\begin{center}
\textbf{\S\ 10. \quad Формы пятого порядка (система III).}
\end{center}

\[
\text{Свертка }\left\{\begin{array}{l}
K,j,\A \text{ с }\nu,\\
f\text{ с }H.
\end{array}\right.
\]

\noindent\textit{Нередуцируемые формы:}
\[
\begin{array}{c@{\qquad}c@{\qquad}c@{\qquad}c}
\mathrm{A)} & \begin{array}{cccc}
\cdot&K_\nu&\cdot&\cdot\\
(k\nu x)^2&\cdot&K_\nu^2&\cdot\\
(k\nu x)^3&K_\nu(k\nu x)^2&\cdot&\cdot\\
\cdot&K_\nu(k\nu x)^3&\cdot&\cdot
\end{array}
&
\mathrm{B)}\ \begin{array}{c}(j\nu x)\\(j\nu x)^2\\(j\nu x)^3\\ \cdot\end{array}
&
\mathrm{C)}\ \begin{array}{c}\A_\nu\\\cdot\\\cdot\\\cdot\end{array}
\end{array}
\]
\[
\mathrm{D)}\qquad
\begin{array}{cccc}
(aHu)&(aHu)^2&\cdot&\cdot\\
a_\eta&a_\eta(aHu)&a_\eta(aHu)^2&\cdot\\
\cdot&a_\eta^2&\cdot&\cdot
\end{array}
\]

\noindent\textit{Редукции:}

\noindent A) Ряд форм $K_\nu^3$: редуцент $a_\nu^3$.

Форма $K_\nu^2(k\nu x)^2$: редуцент $\theta_\nu^2(\vartheta\nu x)^2$.
\[
(k\nu x),\quad K_\nu(k\nu x),\quad K_\nu^2(k\nu x)
\]
являются распадающимися на множители формами по \S\ 3.

\noindent B) Форма
\[
\begingroup\small
\begin{aligned}
(j\nu x)^4
&=(\hata\theta\nu x)^4\\
&=a_\nu^4\theta+\theta_\nu^4 f-4a_\nu^3\theta_\nu
-4a_\nu\{a_\nu\theta_x-(\hata\theta\nu x)\}^3\\
&\quad+6a_\nu^2\{a_\nu\theta_x-(\hata\theta\nu x)\}^2\\
&=\R f+\frac53 i\theta-6a_\nu^2(\hata\theta\nu x)^2
+4a_\nu(\hata\theta\nu x)^3,
\end{aligned}
\endgroup
\]
и поэтому
\[
(j\nu x)^4=\R f+\frac53 i\theta\pmod{(\A,\nu)},\qquad
(\text{ср. конец § 5}).
\]

\noindent C) Форма $\A_\nu^4$: редуцент $\theta_\nu^4$.

Формы $\A_\nu^2$ и $\A_\nu^3$: редуцент $a_\nu^3$ или $\theta_\nu^4$ благодаря замене формы $\A$ нижними формами ряда форм (конец \S\ 5), то есть посредством двойной редукции выражений
\[
a_\vartheta a_\nu^3,
\qquad a_\vartheta a_\nu^3\theta_\nu,
\qquad a_\vartheta\theta_\nu^4.
\]
Двойное вычисление $\A_\nu^3$ дает основание для редукции более высокой формы $a_\eta^3$.

\noindent\textit{Формулы редукции:}
\[
\tag{10}
\begin{aligned}
\mathrm{a)}\quad \A_\nu^2
&=\frac35a_\eta^2-\frac3{10}(asu)^2+\frac13 i\theta
+\frac15u_x^2(2a_\R^2-t),\\
\mathrm{b)}\quad \A_\nu^3
&=-\frac3{10}a_\R+\frac35u_x\left(\frac{13}{10}a_\R^2-\frac25t\right),\\
\mathrm{c)}\quad \A_\nu^4&=a_\R^2-\frac25t.
\end{aligned}
\]

\noindent\textit{Вывод формул:}
\begingroup\small
\[
\begin{aligned}
\mathrm{a)}\quad
 a_\vartheta a_\nu^3
&=\frac13 i\theta
=[a_\vartheta]_{\nu^3}+\frac67[a_\vartheta^2]_{\nu^2}
+\frac65[a_\vartheta(a\theta u)^2]_{\nu^2}\,
u x^2\\
&\quad+\frac{11}{30}[a_\vartheta^2(a\theta u)^2]_\nu\,\nu x^2
-\frac67u_x[a_\vartheta^2]_{\nu^3}\\
&\quad-\frac16u_x[a_\vartheta^2(a\theta u)^2]_{\nu^2}\,\nu x^2,\\
\frac13 i\theta
&=\A_\nu^2-\frac23u_x\A_\nu^3-a_\nu a_{\nu_1}(\nu\nu_1x)^2\\
&\quad+\frac25a_\nu^2(\nu\nu_1x)^2
+\frac15u_x^2a_\nu^2a_{\nu_1}(\nu\nu_1x)^2;
\end{aligned}
\]
\[
\begin{aligned}
\mathrm{b)}\quad
 a_\vartheta a_\nu^3\theta_\nu
&=0=[a_\vartheta]_{\nu^4}+\frac9{14}[a_\vartheta^2]_{\nu^3}
+\frac35[a_\vartheta(a\theta u)^2]_{\nu^2}\,\nu x^2\\
&\quad+\frac{17}{120}[a_\vartheta^2(a\theta u)^2]_{\nu}\,\nu x^2
-\frac9{14}u_x[a_\vartheta^2]_{\nu^4}\\
&\quad-\frac1{24}u_x[a_\vartheta^2(a\theta u)^2]_{\nu^2}\,\nu x^2,\\
0&=\frac56\A_\nu^3-\frac12u_x\A_\nu^4-\frac14a_\eta^3+\frac1{20}u_xa_\eta^4;
\end{aligned}
\]
\[
\begin{aligned}
\mathrm{b')}\quad
 a_\vartheta\theta_\nu^4
&=a_\R=a_\vartheta\theta_\nu\{a_\nu-(a\theta\nu x)\}^3\\
&=-\frac32[a_\vartheta]_{\nu^3}+\frac35[a_\vartheta(a\theta u)^2]_{\nu^2}\,\nu x^2
-\frac{37}{120}[a_\vartheta^2(a\theta u)^2]_{\nu}\,\nu x^2\\
&\quad+\frac32u_x[a_\vartheta^2]_{\nu^4}
+\frac{49}{120}u_x[a_\vartheta^2(a\theta u)^2]_{\nu^2}\,\nu x^2,\\
a_\R&=-\frac32\A_\nu^3+\frac32u_x\A_\nu^4-\frac{11}{20}a_\eta^3+\frac7{20}u_xa_\eta^4;
\end{aligned}
\]
\[
\begin{aligned}
\mathrm{c)}\quad
 a_\vartheta^2\theta_\nu^4
&=a_\R^2=[a_\vartheta^2]_{\nu^4}+\frac35[a_\vartheta^2(a\theta u)^2]_{\nu^2}\,\nu x^2,\\
a_\R^2&=\A_\nu^4+\frac25a_\eta^4.
\end{aligned}
\]
\endgroup

\noindent D) Редукция $a_\eta^3$ посредством двойной редукции $\A_\nu^3$ (C).

Остальные редукции получаются вычислением, дуалистически противоположным вычислению в \S\ 9.

\noindent\textit{Формулы редукции:}
\[
\tag{11}
\begin{array}{rl|rl}
a_\eta^2(aHu)&=-\dfrac16(asu)^3,&
 a_\eta^2(aHu)^2&=\dfrac49 i\nu-\dfrac16(asu)^4,\\[0.7em]
a_\eta^3&=-a_\R+\dfrac35u_xa_\R^2+\dfrac15u_xt,
& a_\eta^3(aHu)&=0,\\[0.7em]
&&a_\eta^4&=t\quad(\text{присоединено как модуль}).
\end{array}
\]

\noindent\textit{Следствия:}

1) $(j\nu x)^4$, $\A_\nu^2$, $a_\eta^2(aHu)$ являются редуцентами.

2) $\nu$ входит только в произведениях с контрагредиентными формами в свертку с $K,j,\A$; то же относится к $f$ относительно $H$ и к $j$ относительно $\nu$, поскольку $\theta_\nu(\theta\cdot j,\nu,x^3)$ по \S\ 11 B становится редуцируемой.
% END UNIT BODY
% END PRODUCER UNIT 011
\endgroup


\clearpage
\begingroup
\setcounter{footnote}{0}
% BEGIN PRODUCER UNIT 012
% Source: C:/Users/Floris/Documents/interlanguage/03_projects/noether/02_slavic_working_corpus/translations/paper02/russian/v001/Noether_Paper02_Section11_Russian_v001.tex
% Identity: 18036 B / SHA-256 EF43B6B1F8C66BF0E38D31F0CE80B3ABA71B074F33884C54DC2EF8766281423D
\setlength{\parindent}{1.25em}
\setlength{\parskip}{0pt}
\emergencystretch=4em
\allowdisplaybreaks
\let\A\relax
\newcommand{\A}{\mathfrak A}
\newcommand{\R}{\mathfrak R}
\newcommand{\hata}{\widehat{a}}
% BEGIN UNIT BODY
\begin{center}
\textbf{\S\ 11. \quad Формы шестого порядка (система III).}
\end{center}
\[
\text{Свертка }\left\{\begin{array}{l}
\theta^2,
 f\cdot j,
 N \text{ с }\nu,\\
\theta\text{ с }H.
\end{array}\right.
\]

\noindent\textit{Нередуцируемые формы:}
\[
\begin{array}{c@{\quad}c@{\quad}c}
\mathrm{A)}\ (\vartheta^2\nu x)^3,\quad \theta_\nu(\vartheta^2\nu x)^3
&
\mathrm{B)}\ a_\nu(j\nu x),\quad a_\nu^2(j\nu x)
&
\mathrm{C)}\ N_\nu
\end{array}
\]
\[
\mathrm{D)}\quad
\begin{array}{c|c|c|c}
 & (\theta Hu)&(\theta Hu)^2&\cdot\\
(\vartheta\eta x)&\theta_\eta&H_\vartheta(\theta Hu),\ \theta_\eta(\theta Hu)&\theta_\eta(\theta Hu)^2\\
(\vartheta\eta x)^2&H_\vartheta(\vartheta\eta x),\ \theta_\eta(\vartheta\eta x)&\theta_\eta^2&\theta_\eta^2(\theta Hu)\\
\cdot&\theta_\eta(\vartheta\eta x)^2&\cdot&\cdot
\end{array}
\]

\noindent\textit{Редукции:}

\noindent A) $(\vartheta^2\nu x)^4$ распадается по формуле (3.)\mbox{a}:
\[
(a\A\vartheta x)^4=\frac12(\vartheta^2\nu x)^4-\frac35 f\cdot a_\nu^2(\nu\vartheta x)^2
=\frac13\A^2+f\A_\vartheta^2.
\]

\noindent B) $a_\nu(j\nu x)^2$ распадается по формуле (9.)\mbox{a}, если по концу \S\ 6 заменить $f\cdot j$ на $(\theta\theta'u)^2$ и учесть редуцируемость $\theta_\vartheta'$:
\[
\frac13a_\nu(j\nu x)^2=(\theta\theta'\nu x)^2\theta_\nu
=\theta_\nu^3\theta-\theta_\nu^2\theta_\nu'
=\theta_\nu^3\theta+\theta_\nu^2(\widehat{j}\nu\theta'u)
=\theta_\nu^3\theta-\frac23\theta_\nu^3\theta.
\]
Формы $a_\nu(j\nu x)^3$ и $a_\nu^2(j\nu x)^2$: редуценты $a_j$ и $(j\nu x)^4$:
\[
a_\nu(j\nu x)^3=a_j(j\nu x)^3-(j\nu x)^3(j\nu\widehat{a}u),
\]
\[
a_\nu^2(j\nu x)^2=a_j^2(j\nu x)^2-2a_j(j\nu x)^2(j\nu\widehat{a}u)+(j\nu x)^2(j\nu\widehat{a}u)^2.
\]

\noindent C) Ряд форм $N_\nu^2$; редуцент $\A_\nu^2$. $(n\nu x)$ и $N_\nu(n\nu x)$ распадаются как перетаскивание над функциональными детерминантами по \S\ 3.

\noindent D) \textit{Обзор редукций:}

\noindent\mbox{a)} Ряд форм $\theta_\eta^2H_\vartheta$: редуцент $a_\eta^2(aHu)$. (Приводит к формам с более высокими символами.)

\noindent\mbox{b)} Ряд форм $\theta_\eta H_\vartheta$: редуцент $a_\eta^2(aHu)$ посредством двойной редукции. (Приводит к формам с более высокими символами и к более высоким формам общего ряда форм, то есть к ряду форм $\theta_\eta^2$.) Вместо $\theta_\eta H_\vartheta$ рассматриваем ряд форм
\[
\theta_\eta(\vartheta\eta x)(\theta Hu)=\theta_\eta H_\vartheta+\theta_\eta^2,
\]
в нем заменяем символы $\theta$ нижней формой $(abu)$ и таким образом приходим к двойной редукции ряда форм
\[
a_\eta^2(aHu)(abu)=\theta_\eta H_\vartheta+\frac32\theta_\eta^2\pmod{s}
\]
вплоть до конечной формы
\[
a_\eta^3b_\eta(bHu)(abu)=\theta_\eta^3H_\vartheta+\frac12\theta_\eta^4.
\]

\noindent\mbox{c)} Ряд форм $\theta_\eta^3$: посредством двойной редукции тех форм, которые одновременно принадлежат рядам форм $\theta_\eta H_\vartheta$ и $\theta_\eta^2H_\vartheta$, то есть форм ряда $\theta_\eta^2H_\vartheta$, к формам с более высокими символами с одной стороны, и к формам с более высокими символами и к более высокому ряду форм $\theta_\eta^3$ с другой стороны.

\noindent\mbox{d)} Формы $\theta_\eta^2(\vartheta\eta x)$, $\theta_\eta^2(\vartheta\eta x)^2$, $\theta_\eta^3(\vartheta\eta x)$: посредством двойной редукции при помощи редуцента $a$ аналогично вычислению в \S\ 9.

\noindent\mbox{e)} Формы $\theta_\eta^2(\theta Hu)^2$ и $\theta_\eta^2(\vartheta\eta x)^2$: посредством двойной редукции выражений $\A_\nu^2(a\A u)^4$ и $(a\A\nu x)^4$ при помощи редуцентов $\A_\nu^2$ и $(a\A u)^4$.

\noindent\mbox{f)} Формы $H_\vartheta$ и $H_\vartheta^2$: распадающиеся формы. Для $H_\vartheta^2$ это получается посредством двойной редукции выражения $\A_\nu^2(a\A u)^2$ или также прямым вычислением произведений $(a_\nu^2)^2$, $a_\nu^2b_{\nu_1}$ формами системы. (Аналогичное вычисление $(a_\nu)^2$ дает соотношение между распадающимися формами.)

\noindent\mbox{g)} Редукция более высоких форм за счет того, что формы ряда $\theta\cdot H$ допускают две возможности редукции (принадлежат двум редуцируемым рядам форм), а именно:
\begin{quote}\small
$g$: посредством двойной редукции $\theta_\eta^2(\vartheta\eta x)^2$; допускает редукции \mbox{d)} и \mbox{e)}.\par
$s_\vartheta^2(\theta su)$: посредством двойной редукции $\theta_\eta^3(\vartheta\eta x)$; допускает редукции \mbox{c)} и \mbox{d)}.\par
$s_\vartheta^2(\theta su)^2$: посредством двойной редукции $\theta_\eta^2H_\vartheta^2$; допускает редукцию \mbox{a)} и имеет редуцент $\theta_\nu^2(\vartheta\nu x)^2$ как множитель.\par
$s_\nu^2$: посредством двойной редукции $\theta_\eta^3H_\vartheta$; допускает редукции \mbox{a)} и \mbox{c)}.
\end{quote}
Доказательство независимости остальных форм получается посредством двойной редукции ряда форм $\A_\nu^2(a\A u)^2=\theta_\eta^2$.

\noindent\textit{Выполнение редукций:}

Существование редукций \mbox{a)}, \mbox{b)}, \mbox{c)}, \mbox{d)} априори ясно, поскольку по указанному закону построения соотношения, возникающие при двойной редукции, независимы друг от друга. Для редукций \mbox{e)}, \mbox{f)}, \mbox{g)} нужно вычислением доказать, что тождества не возникают.

Симметрично диагональному члену в ряду форм $\theta\cdot H$, продвигаясь к форме $\theta_\eta$, мы приводим, частично лишь с указанием вычисления, также формулы, полученные редукциями \mbox{a)}, \mbox{b)}, \mbox{c)}, \mbox{d)}, поскольку они частично применяются при редукциях \mbox{e)}, \mbox{f)}, \mbox{g)}, а частично позднее.

\noindent 1) $H_\vartheta$ и $H_\vartheta^2$ по \mbox{f)}. По исходной схеме\footnote{Знак $\sim$ вместо знака равенства означает, что произведения из $u_x$ с формами более высокими, чем возникающие посредством сверток III и IV, опущены.}
\begingroup\small
\[
\begin{aligned}
a_\nu b_{\nu_1}^2
&=\nu a_\nu b_\nu^2-(aHu)(bHu)b_\eta
\sim \nu a_\nu b_\nu^2-f a_\eta(aHu)^2-(abu)(aHu)b_\eta+(abu)^2b_\eta,\\
a_\nu^2b_{\nu_1}^2
&=b_\nu^2\{a_\nu u_\nu-(a\nu\nu_1u)\}^2
\sim \nu\{a_\nu^2b_\nu^2-2a_\nu b_\nu(aHu)(bHu)
-\frac16(asu)^2(bsu)^2\}\\
&\sim \nu a_\nu^2b_\nu^2-2a^2(aHu)(abu)-\frac16(asu)^2(bsu)^2+(\vartheta\eta x)^2(\theta Hu)^2
\end{aligned}
\]
\endgroup
при подстановке значения $\theta_\eta H_\vartheta$ возникают формулы:
\[
\tag{12}
\begin{aligned}
\mathrm{a)}\quad H_\vartheta&\sim 2a_\nu a_\R^2+2\nu b_\nu(\vartheta\eta x)^2+2f a_\eta(aHu)^2-2\theta_\eta,\\
\mathrm{b)}\quad H_\vartheta^2&\sim (a^2)^2+2\theta_\eta^2+\frac23(\theta su)^2
+\frac1{12}s\nu-\frac19if\nu-\frac16f(asu)^4.
\end{aligned}
\]

\noindent 2) $\theta_\eta H_\vartheta$ по \mbox{b)}:
\begingroup\small
\[
\begin{aligned}
a_\eta^2(aHu)(abu)&\sim [a_\eta^2(aHu)]_{bu}+\frac12 f[a_\eta^2(aHu)^2]
=a_\eta b_\eta(aHu)(abu)\\
&\quad+a_\eta(a\widehat{b}\eta x)(abu)(aHu)
\sim \theta_\eta(\vartheta\eta x)(\theta Hu)+\frac12\theta_\eta^2+\frac14(\nu\eta x)^2.
\end{aligned}
\]
\endgroup
По формулам (5.) и (11.) следует:
\[
\theta_\eta H_\vartheta\sim -\frac32\theta_\eta^2-\frac14(\theta su)^2-\frac1{12}s\nu+\frac29if\nu.
\]

\noindent 3) $\theta_\eta H_\vartheta(\theta Hu)$ вычисляется по \mbox{b)} аналогично $\theta_\eta H_\vartheta$:
\[
\theta_\eta H_\vartheta(\theta Hu)\sim-\theta_\eta^2(\theta Hu)-\frac16(\theta su)^3.
\]

\noindent 4) $\theta_\eta H_\vartheta(\vartheta\eta x)$ по \mbox{b)}; $\theta_\eta^2(\vartheta\eta x)$ по \mbox{d)} (ср. \S\ 9):
\[
\theta_\eta H_\vartheta(\vartheta\eta x)\sim-\theta_\eta^2(\vartheta\eta x)-\frac16s_\vartheta(\theta su)
\sim-\frac13(\vartheta\R x)-\frac16s_\vartheta(\theta su),
\]
\[
\theta_\eta^2(\vartheta\eta x)\sim\frac23a_\eta^3(abu)\sim-\frac13(\vartheta\R x).
\]

\noindent 5)
\[
\scriptsize
\begin{array}{ccc}
\theta_\eta H_\vartheta^2\text{ по } \mathrm{b)},&
\theta_\eta^2H_\vartheta\text{ по } \mathrm{a)},&
\theta_\eta^2H_\vartheta\text{ по } \mathrm{b)}\text{, а тем самым }\theta_\eta^3\text{ по } \mathrm{c)},\\
\text{из }a_\eta^2(aHb)(bHu);&
\text{из }a_\eta^3(Hab)(abu);&
\text{из }a_\eta^3(bHu)(abu).
\end{array}
\]
\[
\theta_\eta H_\vartheta^2\sim-\frac32\theta_\eta^2H_\vartheta-\frac14s_\vartheta(\theta su)^2+\frac16s_\nu-\frac49ia_\nu
\sim-\theta_\R-\frac16s_\nu-\frac19ia_\nu,
\]
\[
\theta_\eta^2H_\vartheta\sim\frac23\theta_\R-\frac16s_\vartheta(\theta su)^2+\frac29s_\nu-\frac29ia_\nu,
\]
\[
\theta_\eta^3H_\vartheta\sim\frac23\theta_\R-\frac23\theta_\eta^3+\frac5{36}s_\nu,
\qquad
\theta_\eta^3\sim\frac14s_\vartheta(\theta su)^2-\frac18s_\nu+\frac13ia_\nu.
\]

\noindent 6) $\theta_\eta^2(\theta Hu)^2$ по \mbox{e)}:
\begingroup\small
\[
\A_\nu^2(a\A u)^2=[(a\A u)^4]_{\nu^2}
=-\frac{12}{5}\theta_\eta^2(\theta Hu)^2-\frac3{10}\sigma-\frac1{20}(\theta su)^4+\nu\R
\quad(\text{по формуле (3.)}\mathrm{c}),
\]
\[
\A_\nu^2(a\A u)^4=[\A_\nu^2]_{aa}^4=-\frac35\theta_\eta^2(\theta Hu)^2+\frac15\sigma-\frac3{10}(\theta su)^4+\frac13ij
\quad(\text{по формуле (10.)}\mathrm{a}),
\]
\endgroup
\[
\theta_\eta^2(\theta Hu)^2=-\frac16\sigma+\frac1{12}(\theta su)^4-\frac19ij+\frac13\nu\R.
\]

\noindent 7) $\theta_\eta^2(\vartheta\eta x)^2$ по \mbox{d)} и \mbox{e)}. Отсюда редукция более высокой формы $g$.

По вычислению, аналогичному \S\ 9, следует:
\[
\begin{aligned}
\theta_\eta^2(\vartheta\eta x)^2
&=\frac12 i f-\frac12a_\eta^2b_\eta-\frac1{12}g
=\frac23tf-\frac23a_\eta^3b_\eta-\frac16g\\
&=-\frac16g-\frac13(\vartheta\R x)^2+\frac4{15}fa_\R^2+\frac8{15}ft
\quad(\text{по формуле (11.)}).
\end{aligned}
\]
По соответствующему вычислению, как выше, следует:
\begingroup\small
\[
((a\A\nu x)^4)=[(a\A u)^4]_{\nu x^4}
=\frac{21}{10}\theta_\eta^2(\vartheta\eta x)^2+\frac7{10}s_\vartheta^2-\frac3{10}g
\quad(\text{по формуле (3.)}\mathrm{c}),
\]
\[
\begin{aligned}
(a\A\nu x)^4&=-\frac13i\A+6[\A_\nu^2]a^2-4[\A_\nu^3]a+\A_\nu^4f\\
&= -\frac{36}{5}\theta_\eta^2(\vartheta\eta x)^2-\frac95s_\vartheta^2-\frac35g-\frac35(\vartheta\R x)^2\\
&\quad+\frac53i\A+\frac{37}{25}fa_\R^2+\frac{74}{25}ft\quad(\text{по формуле (10.)}).
\end{aligned}
\]
\[
\begin{aligned}
\frac{93}{10}\theta_\eta^2(\vartheta\eta x)^2
&=-\frac3{10}g-\frac52s_\vartheta^2-\frac35(\vartheta\R x)^2\\
&\quad+\frac53i\A+\frac{37}{25}fa_\R^2+\frac{74}{25}ft,
\end{aligned}
\]
\endgroup
и, сравнивая два значения $\theta_\eta^2(\vartheta\eta x)^2$ по \mbox{g)}, получаем:
\[
\tag{13}
0=g-2s_\vartheta^2+2(\vartheta\R x)^2+\frac43i\A-\frac45fa_\R^2-\frac85ft.
\]

\noindent 8) $\theta_\eta^2H_\vartheta(\theta Hu)$ по \mbox{a)} и \mbox{b)}, отсюда $\theta_\eta^3(\theta Hu)$ по \mbox{c)} (формы с множителем $u_x$ исчезают):
\[
\theta_\eta^2H_\vartheta(\theta Hu)=-\frac16s_\vartheta(\theta su)^3-\frac13(\nu\R x),
\]
\[
\theta_\eta^2H_\vartheta(\theta Hu)=-\frac13\theta_\eta^3(\theta Hu)-\frac13(\nu\R x),
\]
\[
\theta_\eta^3(\theta Hu)=\frac12s_\vartheta(\theta su)^3.
\]

\noindent 9) $\theta_\eta^2H_\vartheta(\vartheta\eta x)$ по \mbox{a)} и \mbox{b)}, отсюда $\theta_\eta^3(\vartheta\eta x)$ по \mbox{c)}. $\theta_\eta^3(\vartheta\eta x)$ по \mbox{d)}, отсюда редукция более высокой формы $s_\vartheta^2(\theta su)$ по \mbox{g)} (формы с множителем $u_x$ исчезают):
\[
\theta_\eta^2H_\vartheta(\vartheta\eta x)=\frac13(atu),
\]
\[
\theta_\eta^2H_\vartheta(\vartheta\eta x)=\frac13(atu)+\frac13\theta_\eta^3(\vartheta\eta x)-\frac16s_\vartheta^2(\theta su),
\]
\[
\theta_\eta^3(\vartheta\eta x)=\frac12s_\vartheta^2(\theta su),
\]
\[
\theta_\eta^3(\vartheta\eta x)=\frac25\theta_\R(\vartheta\R x)+\frac15(atu),
\]
\[
\tag{14}
s_\vartheta^2(\theta su)=\frac45\theta_\R(\vartheta\R x)+\frac25(atu).
\]

\noindent 10) $\theta_\eta^2H_\vartheta^2$ по \mbox{a)} и через редуцент $\theta_\nu^2(\vartheta\nu x)^2$, $\theta_\eta^3H_\vartheta$ по \mbox{a)} и \mbox{c)}, $\theta_\eta^4$ по \mbox{c)}, отсюда редукция более высоких форм $s_\vartheta^2(\theta su)^2$ и $s_\nu^2$ по \mbox{g)}:
\begingroup\scriptsize
\[
\begin{aligned}
\text{По }\mathrm{a)}\quad
\theta_\eta^2H_\vartheta^2&=[a_\eta^2(aHu)^2]b^2+\frac13[a_\eta^4]_{bu^2}-\frac13H_\nu^2(\nu\eta x)^2\\
&=\frac49ia_\nu^2-\frac16s_\vartheta^2(\theta su)^2-\frac5{18}s_\nu^2+\frac13(atu)^2+\frac1{54}Ju_x^2.
\end{aligned}
\]
\[
\begin{aligned}
\text{Через }\theta_\nu^2(\vartheta\nu x)^2:
\quad\theta_\eta^2H_\vartheta^2
&=[\theta_\nu^2(\vartheta\nu x)^2]_{\nu_1}+\frac13[\theta_\nu^4]_{\nu_1}\hat\nu_1x^2-\frac13s_\vartheta^2(\theta su)^2\\
&=\frac49ia_\nu^2-\frac16s_\nu^2-\frac13s_\vartheta^2(\theta su)^2+\frac13(\nu\R x)^2.
\end{aligned}
\]
\[
\begin{aligned}
\text{По }\mathrm{a)}\quad
\theta_\eta^3H_\vartheta&=-[a_\eta^2(aHu)]_{bu}b^2-\frac12[a_\eta^2(aHu)^2]b^2-\frac13[a_\eta^3]b+\frac1{12}[a_\eta^4]_{bu^2}\\
&\quad+\frac12u_x[a_\eta^2(aHu)^2]b^3\\
&=-\frac29ia_\nu^2+\frac1{12}s_\nu^2+\frac4{15}\theta_\R^2-\frac7{90}(\nu\R x)^2+\frac1{30}(atu)^2+\frac2{27}i^2u_x^2-\frac1{36}Ju_x^2.
\end{aligned}
\]
\[
\begin{aligned}
\text{По }\mathrm{c)}\quad
\theta_\eta^3H_\vartheta:
 a_\eta^3(Hab)(bHu)&=\frac12(atu)^2
=\theta_\eta^2H_\vartheta(\theta Hu)(\vartheta\eta x)+\frac14H_\nu^2(\nu\eta x)^2\\
&\quad+\frac12\theta_\eta^2H_\vartheta^2
=\frac32\theta_\eta^2H_\vartheta^2+\theta_\eta^3H_\vartheta-u_x[a_\eta^2(aHu)^2]b^3\\
&\quad+\frac14H_\nu^2(\nu\eta x)^2,
\end{aligned}
\]
\endgroup
и поэтому
\begingroup\small
\[
\begin{aligned}
\theta_\eta^3H_\vartheta&=-\frac23ia_\nu^2+\frac5{12}s_\nu^2+\frac12(\nu\R x)^2-\frac12(atu)^2\\
&\quad+\frac4{27}i^2u_x^2-\frac1{12}Ju_x^2.
\end{aligned}
\]
По \mbox{c)}
\[
\begin{aligned}
\theta_\eta^4&=\frac49ia_\nu^2-\frac16s_\nu^2+\frac{16}{15}\theta_\R^2\\
&\quad-\frac{14}{45}(\nu\R x)^2+\frac2{15}(atu)^2.
\end{aligned}
\]
\endgroup
По \mbox{g)} из двух значений для $\theta_\eta^2H_\vartheta^2$:
\[
\tag{15}
\begin{aligned}
s_\vartheta^2(\theta su)^2
&=\frac23s_\nu^2-2(atu)^2+2(\nu\R x)^2-\frac19Ju_x^2\\
&=\frac89ia_\nu^2+\frac8{15}\theta_\R^2-\frac{14}{15}(atu)^2+\frac{38}{45}(\nu\R x)^2-\frac4{27}i^2u_x^2.
\end{aligned}
\]
По \mbox{g)} из двух значений для $\theta_\eta^3H_\vartheta$:
\[
\tag{16}
s_\nu^2=\frac43ia_\nu^2+\frac45\theta_\R^2+\frac85(atu)^2-\frac{26}{15}(\nu\R x)^2+\frac16Ju_x^2-\frac29i^2u_x^2.
\]
Формула (16.) дает основу для редукции модуля $(ss'u)^4$, формула (13.) -- для редукции системы $s$. (\S\ 17.)

\noindent\textit{Следствия:}

1) $a_\nu(j\nu x)^3$, $\theta_\eta H_\vartheta$, $\theta_\eta^2(\vartheta\eta x)$, $\theta_\eta^2(\theta Hu)^2$ являются редуцентами; $a_\nu(j\nu x)^2$, $H_\vartheta$ и $H_\vartheta^2$ являются распадающимися формами.

2) Форма системы II: $j(\nu)=g=(\nu\eta x)^4H_x^2$ является редуцируемой.

3) Поскольку единственная контрагредиентная форма $g$ становится редуцируемой, $\nu$ вообще не входит ни в степенях, ни в произведениях с формами той же системы в свертку с системой I.

$\theta$, рассматриваемая как контравариант, входит в свертку только в произведениях или степенях, и соответственно $H$, рассматриваемая как ковариант (ср. \S\ 13 B).
% END UNIT BODY
% END PRODUCER UNIT 012
\endgroup


\clearpage
\begingroup
\setcounter{footnote}{0}
% BEGIN PRODUCER UNIT 013
% Source: C:/Users/Floris/Documents/interlanguage/03_projects/noether/02_slavic_working_corpus/translations/paper02/russian/v001/Noether_Paper02_Section12_Russian_v001.tex
% Identity: 8570 B / SHA-256 2509365B706DFB983E490C4E71A7BFEFEC9D3521CACB0BF21932C690DAB0CEDE
\setlength{\parindent}{1.25em}
\setlength{\parskip}{0pt}
\emergencystretch=4em
\allowdisplaybreaks
\let\A\relax
\newcommand{\A}{\mathfrak A}
\newcommand{\R}{\mathfrak R}
% BEGIN UNIT BODY
\begin{center}
\textbf{\S\ 12. \quad Формы седьмого порядка.}
\end{center}
\[
\text{Свертка }\left\{\begin{array}{l}
\theta\cdot K \text{ с }\nu,\\
K,j,\A \text{ с }H,\\
f \text{ с }L,\sigma.
\end{array}\right.
\]

\noindent\textit{Нередуцируемые формы:}
\[
\mathrm{B)}\quad
\begin{array}{@{}c@{\qquad}c@{\qquad}c@{\qquad}c@{\qquad}c@{}}
\cdot&\cdot&(KHu)^2&\cdot&\cdot\\
\cdot&K_\eta&\cdot&K_\eta(KHu)^2&\cdot\\
(k\eta x)^2&\cdot&K_\eta^2&\cdot&\cdot\\
(k\eta x)^3&H_k(k\eta x)^2,\ K_\eta(k\eta x)^2&\cdot&\cdot&\cdot\\
\cdot&K_\eta(k\eta x)^3&\cdot&\cdot&\cdot
\end{array}
\]
\[
\mathrm{C)}\quad
\begin{array}{cc}
(j\eta x)&H_j\\
\cdot&H_j(j\eta x)\\
\cdot&H_j^2
\end{array}
\qquad
\mathrm{D)}\quad
\begin{array}{c}
(\A Hu)\\\A_\eta
\end{array}
\]
\[
\mathrm{E)}\quad
\begin{array}{cccc}
a_i&(aLu)^2&(aLu)^3&\cdot\\
\cdot&\cdot&a_i(aLu)^2&a_i(aLu)^3\\
\cdot&a_i^2&\cdot&\cdot
\end{array}
\]

\noindent\textit{Редукции:}

\noindent A) $(\vartheta\cdot k,\nu,x)^4$ распадается как однократное перекладывание над распадающейся формой $(\vartheta^2\nu x)^4$.

\noindent B) Ряд форм $H_kK_\eta$: редуцент $H_\vartheta\theta_\eta$.

Ряд форм $K_\eta^2(KHu)$: редуцент $a_\eta^2(aHu)$.

Ряд форм $K_\eta^2(k\eta x)$: редуцент $\theta_\eta^2(\vartheta\eta x)$.

Получающаяся отсюда двойная редукция ряда форм $K_\eta^3$ дает основание для редукции более высокого ряда форм $(j\eta x)^3$.

$(KHu)$, $(k\eta x)$, $K_\eta(k\eta x)$ являются распадающимися формами по \S\ 3.

$H_k(KHu)$, $K_\eta(KHu)$ распадаются как возникшие через однократную свертку с распадающейся формой $K_\nu(k\nu x)$ и функциональным детерминантом
\[
K_\nu^2=\theta_\nu^2(a\theta u)\equiv a_\nu^2(a\theta u)\pmod{\nu^2}.
\]
Двойному представлению $K_\nu^2$ соответствует соотношение между распадающимися формами. Получаем:
\[
\begin{aligned}
H_k(KHu)&=2K_\nu(\nu\nu_1k)=2\theta_{\nu_1}(\nu\nu_1a\theta)\\
&=2\theta_\nu^2a_\nu-a_\nu^2\theta_\nu
+f\theta_\eta(\theta Hu)^2-f\nu\theta_\vartheta^3
\pmod{\nu^3},
\end{aligned}
\]
\[
\begin{aligned}
K_\eta(KHu)&=-K_\nu^2(\nu\nu_1x)
=-a_\nu^2(a\theta u)(\nu\nu_1x)\\
&=a_\eta(aHu)^2\theta+a_\nu^2\theta_\nu
=-\theta_\nu(\theta Hu)^2f-\theta_\nu^2a_\nu.
\end{aligned}
\]
Следовательно
\[
\tag{17}
0=a_\nu^2\theta_\nu+\theta_\nu^2a_{\nu_1}
+f\theta_\eta(\theta Hu)^2+\theta a_\eta(aHu)^2\pmod{\nu^3}.
\]
Формы $H_k$, $H_k(k\eta x)$, $H_k^2$, $H_k^2(k\eta x)$ распадаются как возникшие через однократную свертку с распадающимися формами $H_\vartheta$ и $H_\vartheta^2$ (ср. \S\ 3. 2), a) и b)).

Имеем
\[
H_k=H_\vartheta(a\theta u)\sim [H_\vartheta]_{ua}-fH_\vartheta(\theta Hu)
\]
(специализация $y=uv$ из 2)a)).
\[
H_k(k\eta x)=H_\vartheta a_\eta-H_\vartheta\theta_\eta f,
\qquad
H_\vartheta a_\eta=\frac54[H_\vartheta]a-\frac14H_\vartheta a_\vartheta
\quad(\S\ 3. 2)b),
\]
\[
H_\vartheta^2(a\theta u)=[H_\vartheta^2]_{ua},
\qquad
H_\vartheta^2a_\eta=[H_\vartheta^2]a=H_k^2(k\eta x)+H_\vartheta^2\theta_\eta f.
\]

\noindent C) Ряд форм $(j\eta x)^3$: редуцируем через двойную редукцию ряда форм $K_\eta^3$ по B), по схеме:
\[
0=a_\eta^3(a\theta u)=\theta_\eta^3(a\theta u)
=\theta_\eta+(a\theta\eta x)^3(a\theta u)-\theta_\eta^3(a\theta u)
=\frac12(j\eta x)^3\pmod{(\nu^3,s,\R,t,i,J)}.
\]
Аналогичная схема действует вплоть до конечной формы ряда форм:
\[
0=H_\vartheta^2(a\theta\eta x)a_\eta^3
=H_\vartheta^2(a\theta\eta x)\theta_\eta^3
-\frac12H_\vartheta^2(j\eta x)^4\pmod{(\nu^3,s,\R,t,i,J)}.
\]
Форма $(j\eta x)^2$ распадается: 1) через двукратную поляризацию соотношения для распадающейся формы $H_\vartheta^2$ ((12.) b);

2) через прямое вычисление произведения $a_\nu\theta_\nu^3$; отсюда распад более высокой формы $(\A Hu)^2$. Получаем:
\[
\tag{18}
\left.\begin{array}{ll}
\mathrm{a)}&\dfrac13(\A Hu)^2+\dfrac13(j\eta x)^2=\theta_\nu^2a_{\nu_1}^2,\\[0.6em]
\mathrm{b)}&(j\eta x)^2=\dfrac43\theta_\nu^3a_{\nu_1},
\end{array}\right\}\pmod{(\nu^3,s,\R,t,i,J)}.
\]
по схеме:
\[
H_\vartheta^2(a\theta u)^2-\frac13(\A Hu)^2
=2\theta_\eta^2(a\theta u)(aHu)+\theta_\nu^2a_{\nu_1}^2
=(a\theta u)(aHu)\{a_\eta-(a\theta\eta x)^2\}+\theta_\nu^2a_{\nu_1}^2,
\]
\[
\begin{aligned}
a_{\nu_1}\theta_\nu^3
&=-(a\widehat{\nu}\nu_1u)\theta_\nu^2(\theta\widehat{\nu}\nu_1u)
-\frac12(u\nu\nu_1u)\theta_\nu\theta_{\nu_1}(\vartheta\nu\nu_1u)\\
&=-\frac32\theta_\eta^2(\theta Hu)(aHu)
=-\frac32\theta_\eta^2(\theta Hu)(a\theta u)\\
&=-\frac32\{a_\eta-(a\theta\eta x)^2\}\{(aHu)-(a\theta u)\}(a\theta u).
\end{aligned}
\]

\noindent D) Ряд форм $\A_\eta(\A Hu)$: редуцент $\A_\nu^2$.

$(\A Hu)^2$ является распадающейся формой по C).

\noindent E) Ряд форм $a_i^2(aLu)$: редуцент $a_\eta^2(aHu)$. Формы $(aLu)$ и $a_i(aLu)$: распадаются как перекладывания над функциональными детерминантами по \S\ 3.

\noindent F) Ряд форм $a_\sigma^3$: редуцент $a_\R^3$.

Формы $a_\sigma^2$ и $a_\sigma^3$: редуценты $a_\nu^3$ и $a_\eta^3$ через двойную редукцию соответствующих выражений
\[
H_\nu a_\nu^3\quad\text{и}\quad H_\nu a_\eta^3,
\qquad
H_\nu a_\nu^3a_\nu\quad\text{и}\quad H_\nu a_\eta^4
\]
(ср. \S\ 10. C) редукцию $\A_\nu^2$ и $\A_\nu^3$).

Отсюда редукция для более высоких форм $a_\nu s_\nu(asu)^2$, $a_\nu^2s_\nu(asu)^2$ и для форм в символах $\R,t(a_\nu,a_\eta^2)$ с учетом (16.):
\[
\tag{19}
\left.\begin{array}{rcl}
a_\nu s_\nu(asu)^2&=&\dfrac13 i\theta_\nu^2-\dfrac1{18}iH,\\[0.5em]
a_\nu^2s_\nu(asu)^2&=&0,
\end{array}\right\}\pmod{(\R,t)}.
\]
Форма $a_\sigma$ распадается через поляризацию (12.)b или, короче, через прямое разложение произведения $a_\nu^2\theta_\nu^3$ по схеме:
\[
a_\nu^2\theta_{\nu_1}^3=(a_\nu^2\theta_{\nu_1})a_{\nu_1}-(a\theta\nu_1x)^2
=-2a_\eta(aHu)(a\theta H)(a\theta\eta x)+(a\theta\nu_1x)^2a_\nu^2\theta_{\nu_1}
\]
и с учетом редукции $H_j(j\eta x)^2$ (C):
\[
\tag{20}
a_\sigma=2a_\nu^2\theta_{\nu_1}^3\pmod{(s,\R,t,i)}.
\]

\noindent\textit{Следствия:}

1) $(j\eta x)^3$, $\A_\eta(\A Hu)$, $a_\sigma^2$ являются редуцентами; $(j\eta x)^2$, $(\A Hu)^2$, $a_\sigma$ являются распадающимися формами.

2) $f$ входит только в произведения с контрагредиентными формами в свертке с $L$; $f$, а потому также $K$ и $N$, вообще не входят в свертку с $\sigma$ и, следовательно, также с $(\nu\sigma x)$. $j$ входит только в произведения с контрагредиентными формами; $\A$ вообще не входит в произведения в свертке с $H$ и $L$ (это следует из $\A_\eta(\A Hu)$ и $(\A Hu)^2$). Равным образом $\sigma$, а значит и $(\nu\sigma x)$, не входит в произведения в свертке с какой-либо формой системы I. Из произведений в системе II, с учетом следствий в \S\ 11, остаются только степени $H$ и произведения $H$ с $L$.
% END UNIT BODY
% END PRODUCER UNIT 013
\endgroup


\clearpage
\begingroup
\setcounter{footnote}{0}
% BEGIN PRODUCER UNIT 014
% Source: C:/Users/Floris/Documents/interlanguage/03_projects/noether/02_slavic_working_corpus/translations/paper02/russian/v001/Noether_Paper02_Section13_Russian_v001.tex
% Identity: 6277 B / SHA-256 6FC10A9CDCA32E2A851380E4104DD6553D2FF905DD40B2E65241C938F1D5CA7A
\setlength{\parindent}{1.25em}
\setlength{\parskip}{0pt}
\emergencystretch=4em
\allowdisplaybreaks
\let\A\relax
\newcommand{\A}{\mathfrak A}
\newcommand{\R}{\mathfrak R}
% BEGIN UNIT BODY
\begin{center}
\textbf{\S\ 13. \quad Формы восьмого порядка (система III).}
\end{center}
\[
\text{Свертка }\left\{\begin{array}{l}
\A\cdot j \text{ с }\nu,\\
\theta^2,\ f\cdot j,\ N \text{ с }H,\\
\theta \text{ с }L,\sigma.
\end{array}\right.
\]

\noindent\textit{Нередуцируемые формы:}
\[
\begin{array}{c@{\quad}l@{\qquad}c@{\quad}l}
\mathrm{A)}&\A_\nu(j\nu x)&
\mathrm{B)}&\begin{array}{c}(\vartheta^2\eta x)^3\\(\vartheta^2\eta x)^4\end{array}
\quad H_\vartheta(\vartheta^2\eta x)^3,
\ \theta_\eta(\vartheta^2\eta x)^3,\\[1em]
\mathrm{C)}&(aHu)H_j,\quad a_\eta(aHu)H_j&
\mathrm{D)}&H_n,\quad N_\eta.
\end{array}
\]
\[
\mathrm{E)}\quad
\begin{array}{c|c|c|c}
\cdot&(\theta Lu)^2&(\theta Lu)^3&\cdot\\
\theta_i&\cdot&L_\vartheta(\theta Lu)^2,
\ \theta_i(\theta Lu)^2&\theta_i(\theta Lu)^3\\
(\vartheta lx)^2&\theta_i^2&\cdot&\cdot\\
\cdot&\theta_i(\vartheta lx)^2&\cdot&\cdot
\end{array}
\]

\noindent\textit{Редукции:}

\noindent A) $\A_\nu(j\nu x)^3$: редуцент $a_\nu(j\nu x)^3$.

$\A_\nu(j\nu x)^2$ редуцируема через распадающуюся форму $a_\nu(j\nu x)^2$ по \S\ 3. 2)b с учетом редуцента $\theta_{\vartheta j}$. Действительно, по \S\ 11 B) получаем:
\[
\begin{aligned}
&\frac3{10}\A_\nu(j\nu x)^2+\frac35a_\nu(j\nu x)a_\vartheta(j\nu x)
+\frac1{10}a_\nu(j\nu x)^2\\
&\qquad=\frac35\theta_\nu^3\theta_{\vartheta j}
+\frac25\theta_\nu^3\theta_{\vartheta j}\theta_{\vartheta j},
\end{aligned}
\]
(редуценты $a_j$ и $\theta_{\vartheta j}$).

\noindent B) Формы $H_\vartheta(\vartheta'\eta x)^2$, $H_\vartheta^2(\vartheta'\eta x)$, $H_\vartheta^2(\vartheta'\eta x)^2$ распадаются как возникшие через последовательную однократную свертку с распадающимися формами $H_\vartheta$ и $H_\vartheta^2$, соответственно $H_\vartheta a_\eta$ и $H_\vartheta^2a_\eta$, по \S\ 3. 2)b) и по соотношению:
\[
H_\vartheta(\vartheta'\eta x)^2=2H_\vartheta a_\eta^2f-2H_\vartheta a_\eta b_\eta.
\]

\noindent C) Ряд форм $a_\eta(j\eta x)^2$: редуценты $a_j$ и $(j\eta x)^3$.

Форма $a_\eta(j\eta x)$ распадается: редуцент $a_j$ и однократная свертка с распадающейся формой $(j\eta x)^2$ по \S\ 3, 2)a (специализация $y=uv$).

Форма $a_\eta H_j$ распадается: 1) через однократную свертку с распадающейся формой $a_\nu(j\nu x)^2$;

2) через специальную двукратную свертку с распадающейся формой $(j\eta x)^2$; отсюда соотношение между распадающимися формами.

Из $a_\nu(j\nu x)^2$:
\[
0=\theta_\nu'\theta_\nu+2a_\eta H_j\pmod{(s,\R,t,i,J)}.
\]
Из $(j\eta x)^2$:
\[
0=\theta_\nu'\theta_\nu-\frac32a_\eta H_j\pmod{(s,\R,t,i,J)}
\]
(произведения с контравариантами опущены), по схеме:
\[
\begin{aligned}
0={}&\frac12a_\nu(abu)^2\theta_\nu^3
+\frac12u_x(ubu)\theta_{\nu_1}^3(\theta bu)
-\frac34(\widehat{j}\eta bu)^2\\
&+\frac34H_j(\widehat{j}\eta bu)-\frac18fH_j^2
\end{aligned}
\]
и перестановкой символов в $a_\nu(ubu)(\theta bu)\theta_{\nu_1}^3$.

\noindent D) Ряд форм $N_\eta(NHu)$: редуцент $\A_\nu^2$.

$(NHu)$, $(n\eta x)$, $N_\eta(n\eta x)$, $H_\eta(NHu)$ являются распадающимися формами (ср. \S\ 12 B); $(NHu)^2$ распадается через однократную свертку с формой $(\A Hu)^2$.

\noindent E) Ряды форм $\theta_iL_\vartheta$, $\theta_i^2(\theta Lu)^2$, $\theta_i^2(\vartheta lx)$:

редуценты: $\theta_\eta H_\vartheta$, $\theta_\eta^2(\theta Hu)^2$, $\theta_\eta^2(\vartheta\eta x)$.

Формы $(\theta Lu)$, $(\vartheta lx)$, $L_\vartheta$, $L_\vartheta(\theta Lu)$, $\theta_i(\theta Lu)$, $L_\vartheta(\vartheta lx)$, $\theta_i(\vartheta lx)$, $L_\vartheta^2$, $L_\vartheta^2(\theta Lu)$, $\theta_i^2(\theta Lu)$ распадаются. (Они дуалистически противоположны распадающимся формам из \S\ 12 B).)

\noindent F) Ряд форм $\theta_\sigma(\vartheta\sigma x)$: редуцент $a_\sigma^2$.

Формы $(\vartheta\sigma x)$ и $(\vartheta\sigma x)^2$ распадаются как возникшие через однократную свертку с распадающейся формой $a_\sigma$; $\theta_\sigma$, возникшая через двукратную свертку, распадается, поскольку ряд форм
\[
a_{\nu_2}^2(a\theta u)\theta_{\nu_1}^3\equiv H_j^2(j\eta x)\equiv0\pmod{u_x(s,\R,t,i,J)}:
\]
\[
\theta_\sigma=a_\sigma(abu)^2=a_\nu^2(abu)^2\theta_\nu^3.
\]

\noindent\textit{Следствия:}

$\theta^3$ или $\theta^2K$ не входит в свертку с $H$; $\theta$, рассматриваемая только как контравариант, входит в степени или произведения в свертке с $L$, но вообще не входит в свертку с $\sigma$ и $(\nu\sigma x)$.

$N$ не входит в произведения в свертке с какой-либо формой системы II, равно как и со степенями или произведениями форм системы II. Из произведений в системе I, с учетом следствий последних параграфов, остаются только произведения $f\cdot j$, $\A\cdot j$, степени $\theta$ и произведения $\theta\cdot K$.
% END UNIT BODY
% END PRODUCER UNIT 014
\endgroup


\clearpage
\begingroup
\setcounter{footnote}{0}
% BEGIN PRODUCER UNIT 015
% Source: C:/Users/Floris/Documents/interlanguage/03_projects/noether/02_slavic_working_corpus/translations/paper02/russian/v001/Noether_Paper02_Section14_Russian_v001.tex
% Identity: 3624 B / SHA-256 9691C1A284C7B4B1F2DAECC56D4192AF5ED25B5D7195D350AD8298073AF6646D
\setlength{\parindent}{1.25em}
\setlength{\parskip}{0pt}
\emergencystretch=4em
\allowdisplaybreaks
\let\A\relax
\newcommand{\A}{\mathfrak A}
\newcommand{\R}{\mathfrak R}
% BEGIN UNIT BODY
\begin{center}
\textbf{\S\ 14. \quad Формы девятого порядка (система III).}
\end{center}
\[
\text{Свертка }\left\{\begin{array}{l}
\theta\cdot K \text{ с }H,\\
K,j,\A \text{ с }L,\\
j,\A \text{ с }\sigma,\\
f \text{ с }H^2.
\end{array}\right.
\]

\noindent\textit{Неприводимые формы:}
\[
\begin{array}{c@{\quad}l@{\qquad}c@{\quad}l}
\mathrm{A)}&(\vartheta\cdot k,\eta,x)^4,\quad
H_k(\vartheta\cdot k,\eta,x)^4&
\mathrm{B)}&\begin{array}{l}
(KLu)^3,\quad K_i(KLu)^3,\quad K_i^2,\\
(klx)^3,\quad K_i(Klx)^3
\end{array}\\[0.7em]
\mathrm{C)}&L_j,\quad L_j^2&
\mathrm{D)}&\A_i,\\[0.7em]
\mathrm{E)}&(j\sigma x)&
\mathrm{G)}&(aH^2u)^3,\quad (aH^2u)^4,\quad a_\eta(aH^2u)^3.
\end{array}
\]

\noindent\textit{Редукции:}

\noindent A) Формы $H_\vartheta(k\eta x)^3$, $H_\vartheta^2(k\eta x)^2$, $H_\vartheta^2(k\eta x)^3$ разлагаются как возникшие через последовательную однократную свертку с разлагающимися формами.

\noindent B) Ряды форм $K_iL_k$, $K_i^2(KLu)$, $K_i^2(klx)$:

редуценты: $\theta_\eta H_\vartheta$, $a_\eta^2(aHu)$, $\theta_\eta^2(\vartheta\eta x)$.

Формы $L_i$, $K_i(KLu)$, $K_i(klx)$, $(KLu)^2$, $(klx)^2$ разлагаются как трансвекции над функциональными детерминантами (\S\ 3).

Формы $L_k$, $L_k(KLu)$, $L_k(KLu)^2$, $L_k(klx)$, $L_k(klx)^2$, $L_k^2$, $L_k^2(KLu)$, $L_k^2(klx)$, $L_k^3$, $K_i(KLu)^2$, $K_i(klx)^2$ разлагаются как возникшие через однократную свертку с разлагающимися формами:
\[
\begin{gathered}
L_\vartheta,\quad H_k(KHu),\quad L_\vartheta(\vartheta lx),\quad H_k^2,\\
L_\vartheta^2,\quad L_\vartheta^2(\theta Lu),\quad
K_\eta(KHu),\quad \theta_i(\vartheta lx),
\end{gathered}
\]
по \S\ 3. 2), a) и b).

\noindent C) Ряд форм $(jlx)^3$: редуцент $(j\eta x)^3$.

Формы $(jlx)^2$ и $L_j(jlx)$ возникли через однократную свертку с разлагающейся формой $(j\eta x)^2$.

\noindent D) Ряд форм $\A_i(\A Lu)$: редуцент $\A_\nu^2$.

Формы $(\A Lu)^2$, $(\A Lu)^3$: однократная свертка с разлагающейся формой $(\A Hu)^2$.

\noindent E) Ряд форм $(j\sigma x)^2$: редуцент $a_\sigma^2$ через замену $j$ низшими формами ряда форм (\S\ 5):
\[
a_\sigma^2(a\theta u)^2=\frac13(j\sigma x)^2,\qquad
a_\sigma^2(a\theta\sigma x)^2=\frac13(j\sigma x)^4.
\]

\noindent F) Ряд форм $\A_\sigma^2$: редуцент $a_\sigma^2$.

Форма $\A_\sigma$: редуцент $a_\sigma^2$ через замену $\A$ низшими формами ряда форм:
\[
a_\sigma a_\R^2=\frac23\A_\sigma\pmod{\nu}.
\]

\noindent\textit{Следствия:} Только форма $j$ входит в свертку с $\sigma$ и $(\nu\sigma x)$.

$N$ не входит в свертку с $L$, поскольку соответствующая $\A_i$ форма $N_i$ разлагается.
% END UNIT BODY
% END PRODUCER UNIT 015
\endgroup


\clearpage
\begingroup
\setcounter{footnote}{0}
% BEGIN PRODUCER UNIT 016
% Source: C:/Users/Floris/Documents/interlanguage/03_projects/noether/02_slavic_working_corpus/translations/paper02/russian/v001/Noether_Paper02_Section15_Russian_v001.tex
% Identity: 1756 B / SHA-256 FB054E7479574EA2DD749945DA29406918DC3952D87AF94D8904A6859E14B0C0
\setlength{\parindent}{1.25em}
\setlength{\parskip}{0pt}
\emergencystretch=4em
\allowdisplaybreaks
\let\A\relax
\newcommand{\A}{\mathfrak A}
\newcommand{\R}{\mathfrak R}
% BEGIN UNIT BODY
\begin{center}
\textbf{\S\ 15. \quad Формы десятого порядка (система III).}
\end{center}
\[
\text{Свертка }\left\{\begin{array}{l}
\theta^2,\\
f\cdot j \text{ с }L,\\
\theta \text{ с }H^2.
\end{array}\right.
\]

\noindent\textit{Неприводимые формы:}
\[
\begin{array}{c@{\quad}l@{\qquad}c@{\quad}l}
\mathrm{A)}&(\vartheta^2lx)^4,\quad
\theta_i(\vartheta^2lx)^4&
\mathrm{B)}&(aLu)^2L_j,\quad a_i(aLu)^2L_j,\\[0.8em]
\mathrm{C)}&(\theta H^2u)^3,\quad
(\theta H^2u)^4,\quad
H_\vartheta(\theta H^2u),\quad
\theta_\eta(\theta H^2u)^3.
\end{array}
\]

\noindent\textit{Редукции:}

\noindent A) и B): остальные формы разлагаются посредством однократной свертки с разлагающимися формами.

\noindent C) Остальные формы разлагаются по \S\ 13 B) посредством дуалистически противоположного рассмотрения.

\noindent\textit{Следствия:}

$\theta^2K$ не входит в свертку с $L$; соответственно $H^2L$ не входит в свертку с $K$. $H^3$ и $H^2L$ не входят в свертку с $\theta$. Тем самым схема свертки из \S\ 7 доказана.
% END UNIT BODY
% END PRODUCER UNIT 016
\endgroup


\clearpage
\begingroup
\setcounter{footnote}{0}
% BEGIN PRODUCER UNIT 017
% Source: C:/Users/Floris/Documents/interlanguage/03_projects/noether/02_slavic_working_corpus/translations/paper02/russian/v001/Noether_Paper02_Section16_Russian_v001.tex
% Identity: 2394 B / SHA-256 7D15601A994308815AA483D0A07A829548B4854F8E874F77539441025E2B35ED
\setlength{\parindent}{1.25em}
\setlength{\parskip}{0pt}
\emergencystretch=4em
\allowdisplaybreaks
\let\A\relax
\newcommand{\A}{\mathfrak A}
\newcommand{\R}{\mathfrak R}
% BEGIN UNIT BODY
\begin{center}
\textbf{\S\ 16. \quad Формы 11-го, 12-го, 13-го и 15-го порядков (система III).}
\end{center}

\noindent\textit{11-й порядок.}
\[
\text{Свертка }\left\{\begin{array}{l}
\theta\cdot K \text{ с }L,\\
K,j \text{ с }H^2,\\
j \text{ с }(\nu\sigma x),\\
f \text{ с }H\cdot L.
\end{array}\right.
\]
\noindent\textit{Неприводимые формы:}
\[
\begin{array}{c@{\quad}l@{\qquad}c@{\quad}l}
\mathrm{A)}&(\vartheta\cdot k,l,x)^5,\quad
\theta_i(\vartheta\cdot k,l,x)^5&
\mathrm{B)}&(KH^2u)^4,\quad K_\eta(KH^2u)^4,\\[0.8em]
\mathrm{C)}&(\nu\sigma j),&
\mathrm{D)}&(a_\eta H\cdot L,u)^4.
\end{array}
\]

\noindent\textit{12-й порядок.}
\[
\text{Свертка }\left\{\begin{array}{l}
\theta^3 \text{ с }L,\\
\theta \text{ с }H\cdot L.
\end{array}\right.
\]
\noindent\textit{Неприводимые формы:}
\[
\mathrm{E)}\ (\vartheta^3lx)^6,
\qquad
\mathrm{F)}\ (\theta,H\cdot L,u)^4,
\qquad L_\vartheta(\theta,H\cdot L,u^4).
\]

\noindent\textit{13-й порядок.}
\[
\text{Свертка }K\text{ с }H\cdot L.
\]
\noindent\textit{Неприводимые формы:}
\[
\mathrm{G)}\ (K,H\cdot L,u)^5,
\qquad K_\eta(K,H\cdot L,u)^5.
\]

\noindent\textit{15-й порядок.}
\[
\text{Свертка }K\text{ с }H^3.
\]
\noindent\textit{Неприводимая форма:}
\[
\mathrm{H)}\ (KH^3u)^6.
\]

\noindent\textit{Редукции:}

Формы $H_j^2$, $H_j^2$. Редуцент $L_j^3$ из соотношения:
\[
(\nu lx)L_x^3=-\frac12H^2\pmod{(\sigma,s)}.
\]
Разложение или приводимость остальных форм следует из рассмотрений предыдущих параграфов.

\noindent\textit{Следствия:}

По схеме свертки из \S\ 7 и следствиям \S\S\ 8--15 тем самым исчерпаны неприводимые формы системы III (117 форм).
% END UNIT BODY
% END PRODUCER UNIT 017
\endgroup


\clearpage
\begingroup
\setcounter{footnote}{0}
% BEGIN PRODUCER UNIT 018
% Source: C:/Users/Floris/Documents/interlanguage/03_projects/noether/02_slavic_working_corpus/translations/paper02/russian/v001/Noether_Paper02_Section17_Russian_v001.tex
% Identity: 8197 B / SHA-256 7E9BEAB53686AFEA9D85F542B3183880D71F6C8886BB736E1F3106E326EB36DE
\setlength{\parindent}{1.25em}
\setlength{\parskip}{0pt}
\emergencystretch=4em
\allowdisplaybreaks
\let\A\relax
\newcommand{\A}{\mathfrak A}
\newcommand{\R}{\mathfrak R}
% BEGIN UNIT BODY
\begin{center}
{\Large\bfseries Глава IV.\quad Относительно полная система по модулю $(\varrho,t)$.}
\end{center}

\begin{center}
\textbf{\S\ 17. \quad Редукция модуля $(ss'u)^4$ и системы от $s$.}
\end{center}

Для построения следующей более высокой относительно полной системы надо, по аналогии с \S\ 7, построить систему от $s$ относительно следующего модуля
\[
\nu(s)=(ss'u)^4
\]
и сворачивать ее с формами относительно полной системы по модулю $s$. Нужно показать, что при введении модуля $(\varrho,t)$ как более высокого модуля
\[
\begin{array}{ll}
\mathrm{A)}&\text{модуль }(ss'u)^4\text{ становится приводимым,}\\
\mathrm{B)}&\text{система от }s\text{ состоит из единственной формы }s.
\end{array}
\]

\noindent A) Редукция модуля $(ss'u)^4$ сразу получается из редуцента $s_\nu^2$, найденного по формуле (16), \S\ 11. Из
\[
 s_\nu^2=\frac43 i a_\nu^2+\frac45\theta_\varrho^2+\frac85(atu)^2-
 \frac{26}{15}(\nu\varrho x)^2+\frac16J\cdot u_x^2-
 \frac29i^2\cdot u_x^2
\]
посредством поляризации $x$ по $\nu_1$ следует:
\[
\tag{21}
\begin{aligned}
(ss'u)^4&=-\frac23J\cdot\nu+6s_\nu^2s_{\nu_1}^2\\
&=\frac{24}{5}\theta_\nu^2\theta_\varrho^2+
\frac{48}{5}a_\nu^2(atu)^2-
\frac{52}{5}H_\varrho^2\\
&\quad+\frac43i\cdot(asu)^4+\frac13J\cdot\nu-
\frac49i^2\cdot\nu,
\end{aligned}
\]
а тем самым и редукция модуля $(ss'u)^4$.

\footnotetext{Из формулы (21) следует упомянутое в примечании к введению соотношение между тремя инвариантами 9-го порядка при применении формулы (10)c.}

\noindent B) Редукция
\[
Z=(ss'u)^2=\theta(s)
\]
и отсюда редукция системы от $s$ получается заменой формы $Z$ более низкой формой $g_\nu^2$ по формуле (8)b, \S\ 7, с учетом соотношения
\[
 s_\eta^2=s_\nu^2(\nu\nu_1x)^2-\frac13Z.
\]
Форма $g_\nu^2$ по формуле (13), \S\ 11 имеет редуцент $s_\nu^2$ множителем. По формулам (8) и (16) получаем:
\[
 g_\nu^2=-Z+\frac{14}{5}ia_\eta^2+\frac{14}{15}i(asu)^2
 \pmod{(\varrho,t)},
\]
а по формулам (13), (14), (15), (16):
\[
\begin{aligned}
g_\nu^2&=2[s_\vartheta^2]_{\nu^2}-\frac43 i\A_\nu^2\\
&=2s_\vartheta^2s_\nu^2-\frac65[s_\vartheta^2(\theta su)^2]\widehat{\nu x}^{\,2}
        -\frac43i\A_\nu^2\\
&=\frac45i\cdot a_\eta^2-\frac25i\cdot(asu)^2+\frac13J\cdot\theta
\pmod{(\varrho,t)}.
\end{aligned}
\]
Отсюда следует
\[
\tag{23}
 Z=2i\cdot a_\eta^2+\frac43i\cdot(asu)^2-\frac13J\cdot\theta
 \pmod{(\varrho,t)}.
\]

Итак, относительно полная система по модулю $((ss'u)^4,\varrho,t)$ переходит в относительно полную систему по модулю $(\varrho,t)$, а из нее посредством трансвекции над известной системой двух квадратичных форм возникает абсолютно полная система. Для построения относительно полной системы по модулю $(\varrho,t)$ нам надо, поскольку по формуле (23) система от $s$ редуцируется к форме $s$, перекрывать относительно полную систему по модулю $(s,\varrho,t)$ с $s$ и степенями $s$. Ниже будет показано, что степени $s$ входят в свертку только с системой I.

Объявленное соотношение таково:
\[
\tag{22}
\begin{aligned}
(ass')^4&=\frac{24}{5}\theta_\nu^2\theta_\varrho^2a_\nu^2a_\varrho^2+
\frac{48}{5}a_\nu^2b_\nu^2(tab)^2\\
&\quad-\frac{52}{5}H_\varrho^2a_\nu^4+
\frac53J\cdot i-\frac49i^3,\\
(ass')^4&=\frac{24}{5}a_\varrho^2a_{\varrho'}^2-
\frac{12}{5}t_\varrho^2+\frac53J\cdot i-\frac49i^3.
\end{aligned}
\]

Как \emph{высшие формы} среди форм одинакового порядка и одинаковой степени мы определяем те формы, которые возникли из более высоких форм относительно полной системы по модулю $(s,\varrho,t)$ посредством свертки с $s$, рассматривая свертку с $s$ только как поляризацию исходных форм. Тем самым порядок отдельных форм определен однозначно (ср. \S\ 1. ряды форм 2). Например:
\[
(\theta_\eta(\theta Hu),s,u)^2\text{ выше, чем }s_\eta((\theta Hu)^2,s,u)^*),
\]
\[
(\theta_\eta(\theta Hu)^2,s,u)\text{ выше, чем }(\theta_\eta(\theta Hu),s,u)^2.
\]

Возникающую систему делим на четыре подсистемы:

\noindent\begin{tabular}{@{}l p{0.72\textwidth}@{}}
Система $s_I$: & возникает посредством свертки $s$ и степеней $s$ с системой I,\\
Система $s_{II}$: & возникает посредством свертки $s$ с системой II,\\
Система $s_{III}$: & возникает посредством свертки $s$ с отдельными формами (без произведений) системы III и с произведениями по одной форме из систем I и II,\\
Система $s_{IV}$: & возникает посредством свертки $s$ с произведениями по одной форме из систем I и III, II и III, III и III.
\end{tabular}

\medskip

Отметим еще, что отныне все формулы следует понимать по модулю $(\varrho,t)$, а начиная с формулы (27) -- по модулю $(\varrho,t,i,J,\text{высшие формы})$, не указывая этого каждый раз явно.

Для редукции соберем ранее полученные формулы:
\[
\tag{24}
\begin{aligned}
s_{\nu}^2&=\frac43i\,a_{\nu}^2+\frac16J\cdot u_x^2-\frac29i^2\cdot u_x^2,\\
s_{\nu}^3&=\frac13J\cdot u_x,
& s_{\nu}^4&=J,\\
g&=2s_\vartheta^2-\frac43i\Delta,\\
s_\vartheta^2(\theta su)&=0,\\
s_\vartheta^2(\theta su)^2&=\frac89i\cdot a_{\nu}^2-\frac4{27}i^2\cdot u_x^2,\\
Z&=2i\cdot a_\eta^2+\frac43i(asu)^2-\frac13J\cdot\theta,\\
g_{\nu}&=-s_\eta+u_x\left\{2ia_\eta^2-\frac23i(asu)^2+\frac23J\cdot\theta\right\},\\
g_{\nu}^2&=\frac45ia_\eta^2-\frac25i(asu)^2+\frac13J\cdot\theta,
&g_{\nu}^3&=0,\quad g_{\nu}^4=0.
\end{aligned}
\]
\footnote{Для более краткой записи мы берем член трансвекции вместо полной трансвекции $[(\theta Hu)^2]_{\widehat{s}u}s$ как форму системы. При составлении формул все трансвекции заменяются их начальными членами; например, $(\theta_\eta(\theta Hu),s,u)^2$ заменяется на $\theta_\eta(\theta Hu)(\theta su)^2$.}

\noindent\textit{Следствия:}
\[
 s_{\nu}^2,\quad s_\vartheta^2(\theta su),\quad s_\vartheta^2s_{\nu},\quad s_\vartheta^2\theta_{\nu}
\]
являются редуцентами.
% END UNIT BODY
% END PRODUCER UNIT 018
\endgroup


\clearpage
\begingroup
\setcounter{footnote}{0}
% BEGIN PRODUCER UNIT 019
% Source: C:/Users/Floris/Documents/interlanguage/03_projects/noether/02_slavic_working_corpus/translations/paper02/russian/v001/Noether_Paper02_Section18_Russian_v001.tex
% Identity: 4407 B / SHA-256 3B6A742D9A6EFE182E98A7FC128E8D099275F2BEC651DA793A9E37FE314C9D79
\setlength{\parindent}{1.25em}
\setlength{\parskip}{0pt}
\emergencystretch=4em
\allowdisplaybreaks
\let\A\relax
\newcommand{\A}{\mathfrak A}
\newcommand{\R}{\mathfrak R}
% BEGIN UNIT BODY
\begin{center}
\textbf{\S\ 18. \quad Система $s_I$.}
\end{center}

\[
\text{Свертка }\left\{\begin{array}{l}
s\text{ с } f;\ \theta;\ K,j,\Delta;\ f\cdot j,N;\ \Delta\cdot j,\\
s^2\text{ с }\theta,K.
\end{array}\right.
\]

\noindent\textit{Неприводимые формы.}

\noindent 5-й порядок:
\[
(asu);\qquad (asu)^2;\qquad (asu)^3;\qquad (asu)^4.
\]

\noindent 6-й порядок:
\[
\begin{array}{cccc}
(\theta su)&(\theta su)^2&(\theta su)^3&(\theta su)^4\\
s_\vartheta&s_\vartheta(\theta su)&s_\vartheta(\theta su)^2&s_\vartheta(\theta su)^3\\
\cdot&s_\vartheta^2&\cdot&\cdot
\end{array}
\]

\noindent 7-й порядок:
\[
\begin{array}{c@{\quad}c@{\quad}c@{\quad}c@{\qquad}l}
\cdot&(Ksu)^2&(Ksu)^3&(Ksu)^4&\\
\text{A)}\ s_k&s_k(Ksu)&s_k(Ksu)^2&s_k(Ksu)^3&
\text{B)}\ s_j,\\
\cdot&s_k^2&\cdot&\cdot&
\text{C)}\ (\Delta su),\ (\Delta su)^2.
\end{array}
\]

\noindent 8-й порядок:
\[
\begin{array}{ll}
\text{A)}\ s_j(asu),\ s_j(asu)^2,\ s_j(asu)^3,
&
\text{B)}\ \begin{array}[t]{cc}
\cdot&(Nsu)^2\\
s_n&s_n(Nsu).
\end{array}
\end{array}
\]

\noindent 10-й порядок:
\[
\text{A)}\ s_j(\Delta su),\qquad \text{B)}\ s_\vartheta(\theta s'u)^4.
\]

\noindent 11-й порядок:
\[
\begin{array}{cc}
(Ks^2u)^6&\cdot\\
s_k(Ks^2u)^5&s_k(Ks^2u)^6.
\end{array}
\]

\noindent\textit{Редукции:}

Ряды форм
\[
 s_\vartheta^2(\theta su),\quad s_k^2(Ksu),\quad s_j^2,\quad (\Delta su)^3,\quad (Nsu)^3
\]
имеют редуцент $s_\vartheta^2(\theta su)$ множителем, а $s_j^2$ и $(\Delta su)^3$ получаются возвращением от $j$ и $\Delta$ к более низким формам ряда форм:
\[
\begin{aligned}
s_\vartheta^2(\theta su)(a\theta u)^3&=-\frac12s_j^2,\\
s_\vartheta^2(\theta su)(a\theta u)^2&=\frac13(\Delta su)^3,\\
s_\vartheta^2(\theta su)^2(a\theta u)^2&=\frac13(\Delta su)^4+\frac13s_j^2.
\end{aligned}
\]
Формы $s_\vartheta^2s_{\vartheta'}$ и $s_\vartheta^2s_{\vartheta'}^2$: редуценты $s_\vartheta^2(\theta su)$ и $\theta_{\vartheta'}$:
\[
\begin{aligned}
s_\vartheta^2s_{\vartheta'}&=s_\vartheta^2\theta_{\vartheta'}-s_\vartheta^2(\theta s\widehat{\vartheta'}x),\\
s_\vartheta^2s_{\vartheta'}^2&=s_\vartheta^2s_{\vartheta'}\theta_{\vartheta'}-s_\vartheta^2s_{\vartheta'}(\theta s\widehat{\vartheta'}x).
\end{aligned}
\]
Ряд форм $s_j(\Delta su)^2$: редуценты $\Delta_j$ и $(\Delta su)^3$.

\noindent\textit{Следствия:}

$s$ в степенях не входит в свертку с $f,j,\Delta,N$. Формы $f,j,\Delta,N$ входят в свертку только в произведениях с контрагредиентными формами; однако в произведениях с $f$ эти формы еще могут быть квадратичными по $x$, а в произведениях с $\Delta$ -- линейными по $x$; то есть нужно рассматривать следующие произведения форм:
\[
 f\cdot(0,\lambda),\quad f\cdot(1,\lambda),\quad f\cdot(2,\lambda),\quad
 j\cdot(\mu,0),\quad N,\quad
 \Delta\cdot(0,\lambda),\quad \Delta\cdot(1,\lambda)\quad(\lambda>0),
\]
где $(\mu,\lambda)$ означает форму $\mu$-й степени по $x$ и $\lambda$-й степени по $u$.

$\theta$ или $K$ не входят в степенях или произведениях с произвольными формами в свертку с $s$ или $s^2$. Для произведений с функциональным детерминантом $K$, $K\cdot\varphi_x$ ($\varphi\ne f$ или $\theta$), выполняется соотношение между разлагающимися формами:
\[
K\cdot\varphi_x=(a\theta u)\varphi_x=f\cdot(\varphi\theta u)-\theta\cdot(\varphi a u).
\]
Тем самым приведенная выше схема свертки системы $s_I$ доказана.
% END UNIT BODY
% END PRODUCER UNIT 019
\endgroup


\clearpage
\begingroup
\setcounter{footnote}{0}
% BEGIN PRODUCER UNIT 020
% Source: C:/Users/Floris/Documents/interlanguage/03_projects/noether/02_slavic_working_corpus/translations/paper02/russian/v001/Noether_Paper02_Section19_Russian_v001.tex
% Identity: 2203 B / SHA-256 DDFF9BAFE6DC7CFEB092A0B47F8F2A8EDD7FCB0485C42D7566E8E14867D2C155
\setlength{\parindent}{1.25em}
\setlength{\parskip}{0pt}
\emergencystretch=4em
\allowdisplaybreaks
\let\A\relax
\newcommand{\A}{\mathfrak A}
\newcommand{\R}{\mathfrak R}
% BEGIN UNIT BODY
\begin{center}
\textbf{\S\ 19. \quad Система $s_{II}$.}
\end{center}

Свертка $s$ с $v;\ H;\ L;\ H^2$.

\noindent\textit{Неприводимые формы:}
\[
\begin{array}{c@{\qquad}c@{\qquad}c@{\qquad}c}
\text{6-й порядок} & \text{8-й порядок} & \text{10-й порядок} & \text{12-й порядок}\\[0.3em]
s_{\nu} & (Hsu),\ (Hsu)^2 & (Lsu)^2,\ (Lsu)^3 & (H^2su)^3\\
\cdot & s_\eta & s_l & \cdot\\
\cdot & \cdot & \cdot & \cdot\\
s_{\nu}^4=J & \cdot & \cdot & \cdot
\end{array}
\]

\noindent\textit{Редукции:}

Ряды форм $s_\eta(Hsu)$ и $s_l(Lsu)$: редуцент $s_{\nu}^2$.

Ряд форм $s_\sigma$: редуцент $s_{\nu}^2$ из $H$, $s_{\nu}^2=\dfrac23s_\sigma$.

$(H^2su)^4$ разлагается по формуле (7.)a. (ср. \S\ 11. A). Отсюда: разложение $(H\cdot L,s,u)^4$.

\noindent\textit{Следствия:}

$s^2$ не входит в свертку с системой II.

$\nu$ входит только в произведениях с контрагредиентными формами, а $H$, рассматриваемая как ковариант, входит в свертку в произведениях с формами $(1,\lambda)$, $(2,\lambda)$, $(3,\lambda)$, где $\lambda$ произвольно.

$\sigma$ и $(\nu\sigma x)$ вообще не входят, а $L$ как функциональный детерминант не входит в произведениях в свертку (полные трансвекции над ковариантами системы III становятся приводимыми).

Как произведения систем I и II остаются:
\[
f\cdot \nu,\qquad \Delta\cdot \nu.
\]
% END UNIT BODY
% END PRODUCER UNIT 020
\endgroup


\clearpage
\begingroup
\setcounter{footnote}{0}
% BEGIN PRODUCER UNIT 021
% Source: C:/Users/Floris/Documents/interlanguage/03_projects/noether/02_slavic_working_corpus/translations/paper02/russian/v001/Noether_Paper02_Section20_Russian_v001.tex
% Identity: 3817 B / SHA-256 592C4B3724753E9250F55495077B3ACF7C0BE5873A31C22D14CAB5A47523B11C
\setlength{\parindent}{1.25em}
\setlength{\parskip}{0pt}
\emergencystretch=4em
\allowdisplaybreaks
\let\A\relax
\newcommand{\A}{\mathfrak A}
\newcommand{\R}{\mathfrak R}
% BEGIN UNIT BODY
\begin{center}
\textbf{\S\ 20. \quad Формы 7-го порядка (система $s_{III}$).}
\end{center}

Свертка $s$ с $f\cdot \nu$, $a_{\nu}$, $a_{\nu}^2$.

\noindent\textit{Неприводимые формы:}
\[
\begin{gathered}
s_{\nu}(asu),\quad a_{\nu}(asu),\quad s_{\nu}(asu)^2,\quad a_{\nu}(asu)^2,\quad
s_{\nu}(asu)^3,\quad a_{\nu}(asu)^3,\\
a_{\nu} s_{\nu},\qquad a_{\nu} s_{\nu}(asu),\qquad a_{\nu}^2(asu),\qquad a_{\nu}^2(asu)^2.
\end{gathered}
\]

\noindent\textit{Редукции:}

Самоочевидные редукции, возникающие из прямой свертки с редуцентами $s_{\nu}^2$ и $s_\eta(Hsu)$, отдельно упоминаться не будут; равно как и разложение трансвекций с функциональными детерминантами.

Форма $a_{\nu}s_{\nu}(asu)^2$ и $a_{\nu}^2s_{\nu}(asu)^2$: см. формулу (19.), \S\ 12;
\[
a_{\nu}s_{\nu}(asu)^3=a_{\nu}(a\widehat{\nu}_1\nu_2u)^3(\nu_1\nu_2\nu)=0.
\]
Ряд форм $a_{\nu}^2s_{\nu}$: редуцент $s_\vartheta^2(\theta su)$ получается возвращением $a_{\nu}^2$ к более низким формам, то есть двойной редукцией выражений
\[
s_\vartheta^2(a\theta s)(a\theta u),\qquad
s_\vartheta^2(a\theta s)(a\theta u)^2
\]
по следующему разложению:
\[
\begin{aligned}
s_\vartheta^2(a\theta s)(a\theta u)
&=[(a\theta u)^2]s^3+\frac15[a_\vartheta(a\theta u)^2]s^2+\frac1{60}[a_\vartheta^2(a\theta u)^2]s\\
&=\frac12a_{\nu}^2s_{\nu}-\frac23a_{\nu}s_{\nu}^2,\\
s_\vartheta^2(a\theta s)(a\theta u)^2
&=\frac45[a_\vartheta(a\theta u)^2]_{\widehat{us}}s^2+\frac{23}{180}[a_\vartheta^2(a\theta u)^2]_{\widehat{us}}s\\
&=-\frac12a_{\nu}^2s_{\nu}(asu).
\end{aligned}
\]
При этом выполняется $a_{\nu}^2b_{\nu}(abu)=0$.

Получаются формулы:
\[
\tag{25}
\begin{array}{rcl@{\qquad}rcl}
\underline{a_{\nu}^2s_{\nu}}&=&\dfrac43i\cdot a_{\nu}^2b_{\nu},
& \underline{a_{\nu}s_{\nu}(asu)^2}&=&\dfrac13i\cdot\theta_{\nu}^2-\dfrac1{18}iH,\\[0.7em]
\underline{a_{\nu}s_{\nu}(asu)^3}&=&0,
& \underline{a_{\nu}^2s_{\nu}(asu)}&=&0,\\[0.7em]
\underline{a_{\nu}^2s_{\nu}(asu)^2}&=&0.
\end{array}
\]

\noindent\textit{Следствия:}

1) $a_{\nu}s_{\nu}(asu)^2$, $a_{\nu}^2s_{\nu}$ являются редуцентами.

2) Отсюда получаются значения форм $(\Delta su)^3$, $(\Delta su)^4$, которые в \S\ 19 были признаны приводимыми; аналогично для соответствующего ряда форм $(agu)^2$, который при свертке с $\nu$ через редуцент $g_{\nu}$ приводит к двойной редукции.
\[
\tag{26}
\begin{aligned}
\mathrm{a)}\quad (\Delta su)^3&=-\frac65a_{\nu}^2(asu)+3a_{\nu}s_{\nu}(asu)+2i\theta_{\nu}(\vartheta\nu x),\\
\mathrm{b)}\quad (\Delta su)^4&=-\frac25a_{\nu}^2(asu)^2+\frac43i\theta_{\nu}^2+\frac19iH.
\end{aligned}
\]

Из формул (24.) и (26.), по модулю $(i,J)$ (ср. с. 71):
\[
\tag{27}
\begin{aligned}
\mathrm{a)}\quad (agu)^2&=\frac23(\Delta su)^2-\frac43a_{\nu}s_{\nu}+\frac35s\cdot a_{\nu}^2,\\
\mathrm{b)}\quad (agu)^3&=2a_{\nu}s_{\nu}(asu)-a_{\nu}^2(asu)^2,\\
\mathrm{c)}\quad (agu)^4&=-a_{\nu}^2(asu)^2.
\end{aligned}
\]
% END UNIT BODY
% END PRODUCER UNIT 021
\endgroup


\clearpage
\begingroup
\setcounter{footnote}{0}
% BEGIN PRODUCER UNIT 022
% Source: C:/Users/Floris/Documents/interlanguage/03_projects/noether/02_slavic_working_corpus/translations/paper02/russian/v001/Noether_Paper02_Section21_Russian_v001.tex
% Identity: 5485 B / SHA-256 BFEFF9265642867093F9700043273C206D2998C1C4EA3E36C24EDFDC7D834871
\setlength{\parindent}{1.25em}
\setlength{\parskip}{0pt}
\emergencystretch=4em
\allowdisplaybreaks
\let\A\relax
\newcommand{\A}{\mathfrak A}
\newcommand{\R}{\mathfrak R}
% BEGIN UNIT BODY
\begin{center}
\textbf{\S\ 21. \quad Формы 8-го порядка (система $s_{III}$).}
\end{center}

Свертка $s$ с формами 4-го порядка (система III). (\S\ 9.)

\noindent\textit{Неприводимые формы:}

\[
\begin{array}{c@{\quad}c|c|c}
\cdot
& (\vartheta\nu x)s_{\nu},\ ((\vartheta\nu x)^2,s,u)
& ((\vartheta\nu x),s,u)^2,\ (\theta_{\nu}su)
& \bigstar\\[0.45em]
& (\vartheta\nu x)^2s_{\nu}
& ((\vartheta\nu x)^2,s,u)^2,\ \theta s_{\nu}
& \\[0.45em]
&& ((\vartheta\nu x)^2,s,u)s_{\nu},\ (\theta_{\nu}(\vartheta\nu x)^2,s,u)
&
\end{array}
\]
\[
\begin{array}{c|c|c|c}
\bigstar
& ((\vartheta\nu x),s,u)^3,\ (\theta_{\nu}su)^2
& ((\vartheta\nu x),s,u)^4,\ (\theta_{\nu}su)^3
& \cdot\\[0.45em]
& ((\vartheta\nu x)^2,s,u)^3,\ (\theta_{\nu}(\vartheta\nu x),s,u)^2,\ (\theta_{\nu}^2su)
& (\theta_{\nu}(\vartheta\nu x),s,u)^3,\ (\theta_{\nu}^2su)^2
& (\theta_{\nu}(\vartheta\nu x),s,u)^4\\[0.45em]
& \cdot
& (\theta_{\nu}^3su)
& \cdot
\end{array}
\]

\noindent\textit{Редукции:}

$((\vartheta\nu x),s,u)s_{\nu}$ разлагается как трансвекция над функциональными детерминантами по \S\ 3.

Ряд форм $((\vartheta\nu x),s,u)^2s_{\nu}$: редуценты $s_{\nu}^2$ и $s_\vartheta^2\theta_{\nu}$:
\[
(\widehat{\vartheta\nu su})s_{\nu}(\theta su)=s_{\nu}s_\vartheta(\theta su)=-\frac12(\widehat{\vartheta\nu su})^2(\theta su),
\]
\[
\theta_{\nu}(\widehat{\vartheta\nu su})s_{\nu}=\theta_{\nu}s_{\nu}s_\vartheta=-\frac12\theta_{\nu}(\widehat{\vartheta\nu su})^2.
\]
Ряд форм $\theta_{\nu}s_{\nu}(\theta su)$: редуцент $a_{\nu}s_{\nu}(asu)^2$ через двойную редукцию выражений
\[
a_{\nu}s_{\nu}(asu)^2(abu)\quad\text{с}\quad a_{\nu}s_{\nu}(asu)^2(abu)(sbu);
\]
$\theta_{\nu}s_{\nu}(\theta su)^3$ через двойную редукцию $(agu)^3(abu)(bgu)^3$.

Ряд форм $\theta_{\nu}(\vartheta\nu x)s_{\nu}$: редуцент $a_{\nu}^2s_{\nu}$ из равенства $\theta_{\nu}(\vartheta\nu x)s_{\nu}=a_{\nu}^2(abu)s_{\nu}$.

Ряд форм $(\theta_{\nu}(\vartheta\nu x)^2,s,u)$: двойная редукция форм ряда $\theta_{\nu}(\widehat{\vartheta\nu su})s_{\nu}$, которые одновременно принадлежат к рядам форм $((\vartheta\nu x)su)^2s_{\nu}$ и $\theta_{\nu}(\vartheta\nu x)s_{\nu}$.

Форма $(\theta_{\nu}(\vartheta\nu x),s,u)$ разлагается как однократная трансвекция над функциональным детерминантом $\theta_{\nu}(\vartheta\nu x)=a_{\nu}^2(abu)$.

Форма $((\vartheta\nu x)^2,s,u)^4$: двойная редукция выражения $(a\A u)^2(\A su)^4$, или также $(\theta gu)^4$, из формул (27.) и аналогично редукции $(j\eta x)^4$ (\S\ 12 C)).

Двойная редукция дает формулы:

\[
\tag{28}
\begin{aligned}
0&=\theta_{\nu}s_{\nu}(\theta su)-\frac16(\widehat{\vartheta\nu su})^2(\theta su)-\frac1{12}(Hsu),\\
0&=\theta_{\nu} s(\theta su)^2-\frac14(\widehat{\vartheta\nu su})^2(\theta su)^2-\frac1{24}(Hsu)^2,\\
0&=(\widehat{\vartheta\nu su})^2(\theta su)^2+2\theta_{\nu}(\widehat{\vartheta\nu su})(\theta su)^2
  +3\theta_{\nu}^2(\theta su)^2-\frac13H(su)^2,\\
0&=\theta_{\nu}s_{\nu}(\theta su)^3+\frac12\theta_{\nu}(\widehat{\vartheta\nu su})(\theta su)^3,\\
0&=\theta_{\nu}s_{\nu}s_\vartheta-\frac14s_\eta,
\quad 0=\theta_{\nu}s_{\nu}s_\vartheta(\theta su),
\quad 0=\theta_{\nu}s_{\nu}s_\vartheta(\theta su)^2.
\end{aligned}
\]
Последняя группа соответствует $((\theta_{\nu}(\vartheta\nu x)^2,s,u)^2,\ldots)$.

\noindent\textit{Следствия:}

1) $\theta(\vartheta\nu x)s_{\nu}$, $(\widehat{\vartheta\nu su})(\theta su)s_{\nu}$, $\theta_{\nu}(\widehat{\vartheta\nu su})^2$, $(\widehat{\vartheta\nu su})^2(\theta su)^2$ являются редуцентами.

2) Формы четвертого порядка $(5,\lambda)$, $(6,\lambda)$ и т. д. не входят в свертку с $s^2$.

3) Для ряда форм $(\theta gu)^2$ из формул (27.), с учетом редукций этого параграфа, получаем формулы:

\[
\tag{29}
\begin{aligned}
\mathrm{a)}\quad (\theta gu)^2&=\frac43\theta_{\nu}s_{\nu}+\frac23s_{\nu}s_\vartheta-\frac13s\cdot\theta_{\nu}^2-\frac19s\cdot H\\
&\quad-\frac13f\cdot a_{\nu}^2(asu)^2+\frac13a_{\nu}^2\cdot(asu)^2,\\
\mathrm{b)}\quad (\theta gu)^3&=-\frac13(\widehat{\vartheta\nu su})^2(\theta su)-\theta_{\nu}(\widehat{\vartheta\nu su})(\theta su)\\
&\quad-\theta_{\nu}^2(\theta su),\\
\mathrm{c)}\quad (\theta gu)^4&=2\theta_{\nu}^2(\theta su)^2-\frac13(Hsu)^2,\\
g\vartheta(\theta gu)&=(\widehat{\vartheta\nu su})(\vartheta\nu x)s_{\nu}-\theta_{\nu}(\vartheta\nu x)(\widehat{\vartheta\nu su}),\\
g\vartheta(\theta gu)^2&=\frac23s\cdot\theta_{\nu}^3,
\quad g\vartheta(\theta gu)^3=\frac13\theta_{\nu}^3(\theta su),
\quad g\vartheta(\theta gu)^4=0.
\end{aligned}
\]
% END UNIT BODY
% END PRODUCER UNIT 022
\endgroup


\clearpage
\begingroup
\setcounter{footnote}{0}
% BEGIN PRODUCER UNIT 023
% Source: C:/Users/Floris/Documents/interlanguage/03_projects/noether/02_slavic_working_corpus/translations/paper02/russian/v001/Noether_Paper02_Section22_Russian_v001.tex
% Identity: 5143 B / SHA-256 E22190D5DAEDD45FC707D6C1BF9A7E5CBCB95ACDD75B37FF23B7D69EF73DC353
\setlength{\parindent}{1.25em}
\setlength{\parskip}{0pt}
\emergencystretch=4em
\allowdisplaybreaks
\let\A\relax
\newcommand{\A}{\mathfrak A}
\newcommand{\R}{\mathfrak R}
% BEGIN UNIT BODY
\begin{center}
\textbf{\S\ 22. \quad Формы 9-го порядка (система $s_{III}$).}
\end{center}

Свертка $s$ с формами 5-го порядка (система III) и с произведением $\A\cdot \nu$ (\S\ 10).

\noindent\textit{Неприводимые формы:}
\[
\begin{array}{c@{\,}c|c|c|c}
\text{A)}& \cdot & \cdot & (K_{\nu}su)^2 & \bigstar\\[0.45em]
& \cdot & ((k\nu x)^2,s,u)^2,\ K_{\nu}s_{\nu} & ((k\nu x)^2,s,u)^3 & \\[0.45em]
& (k\nu x)^2s_{\nu},\ ((k\nu x)^3,s,u) & ((k\nu x)^3,s,u)^2 & \cdot & \cdot\\[0.45em]
& (k\nu x)^3s_{\nu} & ((k\nu x)^3,s,u)s_{\nu},\ (K_{\nu}(k\nu x)^3,s,u) & \cdot & \cdot\\[0.45em]
& \bigstar\ (K_{\nu}su)^3 & (K_{\nu}su)^4 & \cdot & (K_{\nu}^2su)^4
\end{array}
\]
\[
\begin{array}{c@{\,}c|c@{\quad}c@{\,}c}
\text{B)}& (j\nu x)s_{\nu},\ ((j\nu x)^2,s,u) & ((j\nu x)^2,s,u)^2 && \\
& ((j\nu x)^3,s,u) & \cdot && \text{C)}\quad s_{\nu}(\A su),\ (\A_{\nu}su),
\end{array}
\]
\[
\begin{array}{c@{\quad}c|c|c}
\text{D)}& \cdot & ((aHu),s,u)^2,\ ((aHu)^2,s,u) & \bigstar\\[0.45em]
& (aHu)s_\eta,\ (a_\eta su) & (aHu)^2s_\eta,\ (a_\eta su)^2 & \\[0.45em]
& \cdot & (a_\eta^2su) & \\
\bigstar& ((aHu),s,u)^3,\ ((aHu)^2,s,u)^2 & ((aHu),s,u)^4 & \\
& (a_\eta su)^3,\ (a_\eta(aHu),s,u)^2,\ (a_\eta(aHu)^2,s,u) & (a_\eta(aHu),s,u)^3 &
\end{array}
\]

\noindent\textit{Редукции:}

A) Редукции непосредственно следуют из редукций \S\ 21 посредством однократной свертки. Ряд форм $K_{\nu}s_{\nu}(Ksu)^2$ имеет редуценты $a_{\nu}s_{\nu}(asu)^2$ и $\theta_{\nu}s_{\nu}(\theta su)^2$ множителями. Отсюда разложение высших форм $K_{\nu}^2(Ksu)^2$, $K_{\nu}^2(Ksu)^3$ по следующему разложению:
\[
0=a_{\nu}s_{\nu}(asu)^2(a\theta u)=K_{\nu}s_{\nu}(Ksu)^2=\theta_{\nu}s_{\nu}(\theta su)^2(a\theta u).
\]

B) Ряд форм $(j\nu x)^2s_{\nu}$: редуцент $a_{\nu}^2s_{\nu}$ из
\[
(j\nu x)^2s_{\nu}=3a_{\nu}^2s_{\nu}(a\theta u)^2.
\]
$((j\nu x)^3,s,u)^2$: редуцент $a_{\nu}^2s_{\nu}$ из
\[
((j\nu x)^3,s,u)^2=-6a_{\nu}^2s_{\nu}(a\theta u)^2(\theta su).
\]

C) Формы $\A_{\nu}(\A su)^2$ и $s_{\nu}(\A su)^2$: редуцент $g_{\nu}$ из формулы (27.)a.

Форма $\A_{\nu}s_{\nu}$: 1) редуцент $a_{\nu}^2s_{\nu}$ из $a_ga_{\nu}^2s_{\nu}=\frac23\A_{\nu}s_{\nu}$;

2) разлагается по формуле (27.)a через двойную редукцию выражения $(\A s\widehat{\nu}x)^2$.

Отсюда: разложение высшей формы $a_\eta s_\eta$.

D) Ряд форм $(aHu)^2s_\eta(asu)$: редуцент $a_{\nu}s_{\nu}(asu)^2$ через двойную редукцию ряда форм $[a_{\nu}s_{\nu}(asu)^2]_{\widehat{\nu_1}x}$; редукция $(aHu)^2s_\eta(asu)^2$ из $a_{\nu}s_{\nu}(asu)^2$ и $a_{\nu}s_{\nu}(asu)^3$. Отсюда редукция $(a_\eta,s,u)^4$.

Ряд форм $a_\eta(aHu)s_\eta$: редуценты $a_{\nu}^2s_{\nu}$, $a_\eta^2s_\eta$ через двойную редукцию выражения $(\A s\widehat{\nu}x)^3$, так как $a_\eta^2s_\eta$ не возникает из редуцированного члена $a_{\nu}^2s_{\nu}(\nu\nu_1x)$ формы $a_\eta(aHu)s_\eta$ посредством свертки.

Ряд форм $(a_\eta^2su)^2$: редуцент $a_{\nu}^2s_{\nu}$ из
\[
a_{\nu}^2s_{\nu}s_{\nu_1}=-\frac12a_\eta^2(Hsu)^2.
\]
Форма $(a_\eta(aHu),s,u)$ разлагается как трансвекция над функциональным детерминантом.

Получаются редукционные формулы:
\[
\tag{30}
\begin{aligned}
\mathrm{a)}\quad 0&=(aHu)^2(asu)s_\eta-a_\eta(asu)(Hsu)^2-a_\eta(aHu)(asu)(Hsu),\\
\mathrm{b)}\quad 0&=(aHu)^2(asu)^2s_\eta-a_\eta(aHu)(asu)^2(Hsu),\\
\mathrm{c)}\quad 0&=a_\eta(asu)^2(Hsu)^2,\\
0&=a_\eta s_\eta+\frac14a_{\nu}^2\cdot g-\frac14s\cdot a_\eta^2,\\
0&=a_\eta(aHu)s_\eta-\frac12a_\eta^2(Hsu),
\quad 0=a_\eta^2(Hsu)^2,\ldots,
\quad 0=a_\eta^2s_\eta.
\end{aligned}
\]

\noindent\textit{Следствия:}

1) $(j\nu x)^2s_{\nu}$, $\A_{\nu}s_{\nu}$, $\A_{\nu}(\A su)^2$ и, следовательно, $N_{\nu}(Nsu)^2$, $(aHu)^2(asu)s_\eta$, $a_\eta(aHu)s_\eta$, $a_\eta^2(Hsu)^2$ являются редуцентами; $a_\eta s_\eta$ является разлагающейся формой, которая при всех одно- и двукратных свертках переходит в разлагающиеся формы.

2) Формы 5-го порядка $(5,\lambda)$, $(6,\lambda)$ и т. д. не входят в свертку с $s^2$.
% END UNIT BODY
% END PRODUCER UNIT 023
\endgroup


\clearpage
\begingroup
\setcounter{footnote}{0}
% BEGIN PRODUCER UNIT 024
% Source: C:/Users/Floris/Documents/interlanguage/03_projects/noether/02_slavic_working_corpus/translations/paper02/russian/v001/Noether_Paper02_Section23_Russian_v001.tex
% Identity: 6294 B / SHA-256 BE9D3BB24AF58916B9A9B0D90A0427655F179F8631242430A33618D7321E131E
\setlength{\parindent}{1.25em}
\setlength{\parskip}{0pt}
\emergencystretch=4em
\allowdisplaybreaks
\let\A\relax
\newcommand{\A}{\mathfrak A}
\newcommand{\R}{\mathfrak R}
% BEGIN UNIT BODY
\begin{center}
\textbf{\S\ 23. \quad Формы 10-го порядка (система $s_{III}$).}
\end{center}

Свертка $s$ с формами 6-го порядка (система III) (\S\ 11).

\noindent\textit{Неприводимые формы:}
\[
\begin{array}{r@{\quad}c}
\text{A)}&
\begin{array}{c|c}
(\vartheta^2\nu x)^3s_{\nu}&\cdot\\[0.45em]
\cdot&(\vartheta^2\nu x)^3s_{\nu}s_\vartheta,\ \theta_{\nu}(\vartheta^2\nu x)^3s_\vartheta
\end{array}
\end{array}
\]
\[
\begin{array}{r@{\quad}c}
\text{B)}&
\begin{array}{c|ccc}
\cdot&(a_{\nu}(j\nu x),s,u)^2&(a_{\nu}(j\nu x),s,u)^3&(a_{\nu}(j\nu x),s,u)^4\\[0.45em]
a_{\nu}(j\nu x)s_{\nu}&\cdot&(a_{\nu}^2(j\nu x),s,u)^2&(a_{\nu}^2(j\nu x),s,u)^3
\end{array}
\end{array}
\]
\[
\begin{array}{r@{\quad}c}
\text{C)}&
\begin{array}{c|c|c}
\cdot&\cdot&((\theta Hu),s,u)^2,\ ((\theta Hu)^2,s,u)\\[0.35em]
\cdot&((\vartheta\eta x),s,u)^2,\ (\theta Hu)s_\eta,\ (\theta_\eta su)&((\vartheta\eta x),s,u)^3,\ (\theta_\eta su)^2\\[0.35em]
((\vartheta\eta x)^2,s,u)&\cdot&(\theta_\eta^2su)\\[0.35em]
\cdot&(\theta_\eta(\vartheta\eta x)^2,s,u)&\cdot
\end{array}
\end{array}
\]
\[
\begin{array}{c|c}
((\theta Hu),s,u)^3,\ ((\theta Hu)^2,s,u)^2&((\theta Hu),s,u)^4\\[0.45em]
(H_\vartheta(\theta Hu),s,u)^2,\ (\theta_\eta(\theta Hu),s,u)^2,\ (\theta_\eta(\theta Hu)^2,s,u)&(\theta_\eta su)^4,\ (\theta_\eta(\theta Hu),s,u)^3\\[0.45em]
\cdot&(\theta_\eta^2(\theta Hu),s,u)^2
\end{array}
\]

\noindent\textit{Редукции:}

A) формы $((\vartheta^2\nu x)^3,s,u)^2$, $((\vartheta^2\nu x)^3,s,u)^3$ разлагаются:
\[
(\vartheta\nu x)(\widehat{\vartheta\nu su})(\widehat{\vartheta'\nu su})
=(\vartheta\nu x)s_\vartheta s_{\vartheta'}-(\vartheta\nu x)s_{\nu}s_\vartheta
=-2\theta(\vartheta\nu x)s_{\nu}s_\vartheta+(\vartheta\nu x)s_\vartheta^2.
\]
Остальные формы: либо имеют редуценты множителями, либо являются однократными свертками с разлагающимися формами.

B) Редуцент $a_{\nu}^2s_{\nu}$ и $(\widehat{j\nu su})s_{\nu}$.

C) 1. Ряд форм $(\theta Hu)^2s_\eta$: a) редуцент $(aHu)^2(asu)s_\eta$; b) двойная редукция рядов форм $(\theta gu)^3g_{\nu}$ и $g_{\nu}(agu)^3(bgu)^2(abu)$.

Отсюда редукция высших форм $(\theta_\eta su)^3$, $(H_\vartheta(\theta Hu),s,u)^3$ и $(a_{\nu}\cdot b_{\nu_1}^2,s,u)^4$ [последняя форма вместо $(H_\vartheta su)^4$].

2. $((\vartheta\eta x),s,u)^4$: редуцент $(a_\eta su)^4$.

3. $((\vartheta\eta x)^2,s,u)^3$: редуцент $s_\vartheta^2s_{\nu}$ из $(\widehat{\vartheta\eta}su)^2(Hsu)$. Формы $(\vartheta\eta x)s_\eta$, $(\vartheta\eta x)^2s_\eta$, $((\vartheta\eta x)^2,s,u)^2$, $\theta_\eta s_\eta$ разлагаются посредством свертки с разлагающейся формой $a_\eta s_\eta$.

4. Ряды форм $H_\vartheta(\theta Hu)s_\eta$, $\theta_\eta(\theta Hu)s_\eta$, $H_\vartheta(\vartheta\eta x)s_\eta$, $\theta_\eta(\vartheta\eta x)s_\eta$: редуценты $a_\eta(aHu)s_\eta$ и $a_\eta^2s_\eta$.

Ряд форм $(\theta_\eta(\vartheta\eta x),s,u)^2$: редуцент $a_\eta^2(Hsu)^2$ [за исключением $(\theta_\eta^2(\theta Hu),s,u)^2$, которое возникает из члена $a_\eta^2(bsu)(Hsu)$ посредством свертки].

5. $(H_\vartheta(\vartheta\eta x),s,u)$, $(H_\vartheta(\vartheta\eta x),s,u)^2$: двойная редукция $a_\eta(aHu)s_\eta(abu)$ и $g_\vartheta(\theta gu)g_{\nu}$.

Ряд форм $(H_\vartheta(\vartheta\eta x),s,u)^3$: редуцент $((\vartheta\nu x)^2,s,u)^2s_{\nu}$.

6. $(H_\vartheta(\theta Hu),s,u)$ разлагается посредством однократной свертки с разлагающейся формой $(\theta_{\nu}(\vartheta\nu x),s,u)$ по \S\ 3. 2), c); то есть через двойную редукцию более низкой разлагающейся формы $(H_\vartheta su)^2$: \footnote{Знак $\equiv$ означает, что произведения форм более низкой степени опущены.}
\[
\begin{aligned}
0&\equiv \theta_{\nu}(\vartheta\nu\nu_1)(\theta su)+\frac23\theta_{\nu}^2(\vartheta\nu_1x)(\theta su)+\theta_{\nu}(\vartheta\nu x)(\theta su)s_{\nu_1}\\
&\equiv -\frac12H_\vartheta(\theta Hu)(\theta su)+\theta_\eta(\theta su)(Hsu)+\frac12H_\vartheta(\theta su)(Hsu)\\
&\equiv -\frac12H_\vartheta(\theta Hu)(\theta su)+\theta_\eta(\theta su)(Hsu)+\frac12[H_\vartheta]_{\widehat{su}^2}-\frac12\cdot\frac35H_\vartheta(\theta Hu)(\theta su)\\
&\equiv -\frac35H_\vartheta(\theta Hu)(\theta su).
\end{aligned}
\]
Через двойную редукцию по 1. и 5. получаются формулы:
\begin{align*}
0&=\theta_\eta(\theta su)(Hsu)^2+\frac14H_\vartheta(\theta Hu)(\theta su)^2-\frac34\theta_\eta(\theta Hu)(\theta su)(Hsu)\\
&\quad-\frac54\theta_\eta(\theta Hu)^2(\theta su)-(asu)^3\cdot b_\eta(bHu)^2-a_{\nu}(asu)^3\cdot b_{\nu_1}^2,\\
0&=H_\vartheta(\theta Hu)(\theta su)^3-7\theta_\eta(\theta Hu)(\theta su)^2(Hsu),\tag{31}\\
0&=a_{\nu}(asu)^2b_{\nu_1}^{2}(bsu)^2-\theta_\eta(\theta su)^2(Hsu)^2+6\theta_\eta(\theta Hu)(\theta su)^2(Hsu),\\
0&\equiv (H_\vartheta(\vartheta\eta x),s,u),\qquad
0\equiv H_\vartheta(\widehat{\vartheta\eta su})(Hsu)-4\theta_\eta^2(Hsu).
\end{align*}

\noindent\textit{Следствия:}

1) $(\theta Hu)^2s_\eta$, $H_\vartheta(\theta Hu)s_\eta$, $\theta_\eta(\theta Hu)s_\eta$, $H_\vartheta(\vartheta\eta x)s_\eta$, $\theta_\eta(\vartheta\eta x)s_\eta$, $(H_\vartheta(\vartheta\eta x),s,u)$, $(\theta_\eta(\vartheta\eta x),s,u)^2$ являются редуцентами.

2) $s^2$ не входит в свертку с формами $(5,\lambda)$ и т. д. 6-го порядка.
% END UNIT BODY
% END PRODUCER UNIT 024
\endgroup


\clearpage
\begingroup
\setcounter{footnote}{0}
% BEGIN PRODUCER UNIT 025
% Source: C:/Users/Floris/Documents/interlanguage/03_projects/noether/02_slavic_working_corpus/translations/paper02/russian/v001/Noether_Paper02_Section24_Russian_v001.tex
% Identity: 3322 B / SHA-256 620B6A8F82CC5DD55E1FEA7C77E2FFEEF116643BAE25D0B21F8FF151ED530C8A
\setlength{\parindent}{1.25em}
\setlength{\parskip}{0pt}
\emergencystretch=4em
\allowdisplaybreaks
\let\A\relax
\newcommand{\A}{\mathfrak A}
\newcommand{\R}{\mathfrak R}
\newcommand{\D}{\mathfrak D}
% BEGIN UNIT BODY
\begin{center}
\textbf{\S\ 24. \quad Формы 11-го порядка (система $s_{III}$).}
\end{center}

Свертка $s$ с формами 7-го порядка (система III) (\S\ 12).

\noindent\textit{Неприводимые формы:}
\[
\text{A)}\quad
\begin{array}{c|c|c|c}
\cdot&\cdot&\cdot&((KHu)^2su)^2\\[0.35em]
\cdot&\cdot&(K_\eta su)^2&\cdot\\[0.35em]
\cdot&\cdot&\cdot&\cdot\\[0.35em]
((k\eta x)^3,s,u)&\cdot&\cdot&\cdot\\[0.35em]
\cdot&(K_\eta(k\eta x)^3,s,u)&\cdot&\cdot
\end{array}
\]
\[
\begin{array}{c|c}
((KHu)^2,s,u)^3&((KHu)^2,s,u)^4\\[0.45em]
\cdot&(K_\eta(KHu)^2,s,u)^3
\end{array}
\]
\[
\text{B)}\quad ((j\eta x),s,u)^2,\ (H_jsu)\ \big|\ ((j\eta x),s,u)^3,
\qquad \text{C)}\quad (\D_\eta su),
\]
\[
\text{D)}\quad
\begin{array}{c|c|c}
\cdot&((aLu)^2,s,u)^2,\ ((aLu)^3,s,u)&((aLu)^2,s,u)^3\\[0.5em]
(aLu)^2s_\eta,\ (a_lsu)^2&(a_lsu)^3&(a_l(aLu)^2,s,u)^2
\end{array}
\]

\noindent\textit{Редукции:}

A) Редукция или разложение посредством однократной свертки с соответствующими формами в символах $\theta\cdot H$.

$(K_\eta(KHu)^2,s,u)$ и $(K_\eta(KHu)^2,s,u)^2$: разложение по \S\ 3. 2), a), посредством свертки с $K_{\nu}^2(Ksu)$, $K_{\nu}^2(Ksu)^2$.

$(K_\eta(KHu),s,u)^4\equiv(a_{\nu}^2\cdot\theta_{\nu_1},s,u)^4$: разложение по \S\ 3. 2), b), из разлагающейся формы $K_{\nu}^2(Ksu)^3$.

B) Ряд форм $(j\eta x)s_\eta$: редуцент $(j\nu x)^2s_{\nu}$.

Форма $H_j(\widehat{j\eta su})(Hsu)$: редуцент $(j\nu x)(\widehat{j\nu su})^2$.

Форма $H_j(j\eta x)(Hsu)$: разлагается посредством редукции разлагающейся формы $((j\eta x)^2,s,u)^2$ редуцентом $(j\nu x)^2s_{\nu}$ (формула (18.)a):
\[
0=-2(j\nu x)^2s_{\nu}s_{\nu_1}=[(j\eta x)^2]_{\widehat{su}^2}+H_j(j\eta x)(Hsu).
\]
C) Редуценты $\D_{\nu}s_{\nu}$ и $\D_{\nu}(\D su)^2$.

D) Редукция и разложение посредством однократной свертки с соответствующими формами в символах $f\cdot H$.

E) Из (30.)c получается редукция разлагающегося ряда форм $(a_\sigma su)^2$:
\[
\begin{aligned}
0&=a_\eta(Hsu)^2(as\widehat{\nu}x)^2\\
&=a_\eta(Hsu)^2(a\widehat{\nu}\eta u)^2
  +2a_\eta(a\widehat{\nu}\eta u)(\widehat{\nu}\eta su)(Hsu)^2\\
&=\frac13a_\sigma(asu)^2,\\
0&=a_\eta a_{\nu}(Hsu)^2(as\widehat{\nu}x)(asu)\\
&=a_\eta(a\widehat{\nu}\eta u)^2(Hsu)^2(asu)
  +a_\eta(a\widehat{\nu}\eta u)(\widehat{\nu}\eta su)(Hsu)^2(asu)\\
&=\frac13a_\sigma(asu)^3.
\end{aligned}
\]
\noindent\textit{Следствие: $s^2$ не входит в свертку с формами 7-го порядка.}
% END UNIT BODY
% END PRODUCER UNIT 025
\endgroup


\clearpage
\begingroup
\setcounter{footnote}{0}
% BEGIN PRODUCER UNIT 026
% Source: C:/Users/Floris/Documents/interlanguage/03_projects/noether/02_slavic_working_corpus/translations/paper02/russian/v001/Noether_Paper02_Section25_Russian_v001.tex
% Identity: 4018 B / SHA-256 FB7B0E63CFDC4D17BBD5B10D4C6F5624ED509F9A2FA8770E134504D607B27621
\setlength{\parindent}{1.25em}
\setlength{\parskip}{0pt}
\emergencystretch=4em
\allowdisplaybreaks
\let\A\relax
\newcommand{\A}{\mathfrak A}
\newcommand{\R}{\mathfrak R}
\newcommand{\D}{\mathfrak D}
% BEGIN UNIT BODY
\begin{center}
\textbf{\S\ 25. \quad Формы 12-го, 13-го, 14-го, 15-го порядков (система $s_{III}$).}
\end{center}

Свертка $s$ с формами 8-го, 9-го, 10-го, 11-го порядков (система III) (\S\S\ 13--16).

\noindent\textit{Неприводимые формы:}

\noindent 12-й порядок.
\[
\begin{adjustbox}{max width=0.98\linewidth}
$\displaystyle
\begin{array}{c@{\qquad}c@{\qquad}c}
\mbox{A)}&
\begin{array}{c|c}
((\vartheta^2\eta x)^4,s,u)&\cdot\\[0.35em]
\cdot&\theta_\eta(\vartheta^2\eta x)^3s_\vartheta
\end{array}
&
\mbox{B)}\quad ((aHu)H_j,s,u)^2\;\big|\;((aHu)H_j,s,u)^3
\\[1.0em]
\mbox{C)}&
\begin{array}{c|c|c}
\cdot&((\theta Lu)^2,s,u)^2,\ ((\theta Lu)^3,s,u)&((\theta Lu)^2,s,u)^3\\[0.35em]
(\theta_lsu)^2&\cdot&(\theta_l(\theta Lu)^2,s,u)^2
\end{array}
&
\end{array}$
\end{adjustbox}
\]
\noindent\textit{Редукции:} Редукции, возникающие посредством прямой свертки с редуцентами или разложимыми формами, далее уже не будут упоминаться.

\mbox{B)} Разложение $((aHu)H_j,s,u)$ и $(a_\eta(aHu)H_j,s,u)$ как трансвекция через функциональные детерминанты.

Разложение $(a_\eta(aHu)H_j,s,u)^2$: двойная редукция разложимой формы $[\theta_\nu\cdot\theta_{\nu_1}^{\prime 3}]_{\widehat{su^3}}$ (\S\ 13. \mbox{C)}):
\begin{align*}
0&\equiv[\theta_{\nu_1}^{\prime 3}\cdot\theta_\nu]_{\widehat{su^3}}\\
&\equiv-\frac34\theta_\nu(\theta su)^2(\theta\theta'u)\theta_{\nu_1}^{\prime 3}\\
&\equiv-\frac98(\theta\theta'\eta su)(\theta\theta'u)(\theta Hu)(\theta'Hu)
  \theta_\eta'(\theta su)\\
&\equiv-\frac38a_\eta(aHu)H_j(asu)^2.
\end{align*}

\mbox{C)} Из формул (31.) следует редукция разложимых форм
\[
 [L_\vartheta(\theta Lu)]_{\widehat{su^3}}\quad\hbox{и}\quad [\theta_l(\theta Lu)]_{\widehat{su^3}},
\]
то есть произведений $[\theta_\nu^2\cdot H]_{\widehat{su^3}}$ и $[a_\nu^2\cdot(bHu)^2]_{\widehat{su^3}}$:
\begin{equation}
\tag{32}
\begin{aligned}
0&\equiv\theta_\nu^2(\theta su)^2(Hsu),\\
&0&\equiv [L_\vartheta(\theta Lu)]_{\widehat{su^4}}-7[\theta_l(\theta Lu)]_{\widehat{su^4}},\\
0&\equiv a_\nu^2(asu)^2(bHu)^2(bsu)+6\theta_l(\theta Lu)^2(\theta su)(Lsu),\\
0&\equiv\theta_\nu^3(\theta su)(Hsu)^2
\quad\text{из }0\equiv\theta_l^2(\theta Lu)(\theta su)(Lsu)^2\\
&\quad(\text{редуцент }a_\eta^2(Hsu)^2).
\end{aligned}
\end{equation}

\noindent\textit{Неприводимые формы.}

\noindent 13-й порядок.
\[
\begin{array}{rl}
\mbox{A)}&((KLu)^3,s,u)^2;\ ((KLu)^3,s,u)^3,\\[0.25em]
\mbox{B)}&(L_jsu)^2,\\[0.25em]
\mbox{C)}&((aH^2u)^3,s,u)^2.
\end{array}
\]
\noindent 14-й порядок.
\[
\begin{array}{rl}
\mbox{A)}&((aLu)^2L_j,s,u)^2,\\[0.25em]
\mbox{B)}&((\theta H^2u)^3,s,u)^2.
\end{array}
\]
\noindent 15-й порядок.
\[
 ((KH^2u)^4,s,u)^2.
\]
\noindent\textit{Редукции:} Разложение $(K_l(KLu)^3,s,u)^2$, $(H_\vartheta(\theta H^2u)^3,s,u)$, $((K,H\cdot L,u)^5,s,u)$ по \S\ 3. \mbox{2)}, \mbox{c)}; то есть при одновременном рассмотрении разложимых форм $(K_l(KLu)^2,s,u)^3$, $(H_\vartheta(\theta H'u)^2,s,u)^2$, $((K,H\cdot L,u)^4,s,u)^2$.

\noindent\textit{Следствие:} $s^2$ не входит в свертку с отдельными формами системы III.
% END UNIT BODY
% END PRODUCER UNIT 026
\endgroup


\clearpage
\begingroup
\setcounter{footnote}{0}
% BEGIN PRODUCER UNIT 027
% Source: C:/Users/Floris/Documents/interlanguage/03_projects/noether/02_slavic_working_corpus/translations/paper02/russian/v001/Noether_Paper02_Section26_Russian_v001.tex
% Identity: 23497 B / SHA-256 0D0510A2F9B55CAE693A02B56749EB8CEE058F3FCBE7F228DE6EC76FE9222F27
\setlength{\parindent}{1.25em}
\setlength{\parskip}{0pt}
\emergencystretch=4em
\allowdisplaybreaks
\let\A\relax
\newcommand{\A}{\mathfrak A}
\newcommand{\R}{\mathfrak R}
\newcommand{\D}{\mathfrak D}
\newcommand{\mcell}[1]{\ensuremath{\begin{gathered}#1\end{gathered}}}
% BEGIN UNIT BODY
\begin{center}
\textbf{\S\ 26.\quad Система $s_{IV}$.}
\end{center}
\noindent 1) \textit{Свертка $s$ с системами I--III.}

\noindent\textit{Редукции:} После редукций последних параграфов $(a_\nu^2s_\nu,\ (aHu)^2(asu)s_\eta$ и т. д.) формы $(0,\lambda)$, $(1,\lambda)$, $(2,\lambda)$ не допускают одновременно сверток I и III; следовательно, по \S\ 18, для образования системы $s_{IV}$ формы $f,\D,N$ не входят в произведениях в свертку.

$j$ не входит в свертку, поскольку по тем же редукциям сверткам, контрагредиентным к $j$,
\[
(K_\nu(k\nu x)^3,s,u),\qquad ((\vartheta^2\eta x)^4,s,u)\quad\hbox{и т. д.}
\]
соответствуют свертки, когредиентные к $j$:
\[
K_\nu(k\nu x)^2s_k,
\qquad (\vartheta^2\eta x)^3s_\vartheta\quad\hbox{и т. д.}
\]

\noindent 2) \textit{Свертка $s$ с системой II--III}, то есть $s$ с $H\cdot(1,\lambda)$, $H\cdot(2,\lambda)$, $H\cdot(3,\lambda)$.

\noindent\textit{Неприводимые формы:}
\[
\begin{array}{rl}
\mbox{A)}&(a_\nu\cdot H,s,u)^4,
((aHu)^2\cdot H,s,u)^3,
((aHu)^2\cdot H,s,u)^4,\\[0.25em]
&((aLu)^2\cdot H,s,u)^4,
((aLu)^3\cdot H,s,u)^3,
((aH^2u)^3\cdot H,s,u)^4,\\[0.45em]
\mbox{B)}&(a_\nu^2\cdot H,s,u)^3,
(a_\eta(aHu)^2\cdot H,s,u)^3,
(a_l(aLu)^2\cdot H,s,u)^4,\\[0.45em]
\mbox{C)}&(\theta_\nu\cdot H,s,u)^4,
((\theta Hu)^2\cdot H,s,u)^3,
((\theta Hu)^2\cdot H,s,u)^4,\\[0.25em]
&((\theta Lu)^2\cdot H,s,u)^4,
((\theta Lu)^3\cdot H,s,u)^3,
((\theta H^2u)^3\cdot H,s,u)^4,\\[0.45em]
\mbox{D)}&(\theta_\eta(\theta Hu)^2\cdot H,s,u)^3,
(\theta_l(\theta Lu)^2\cdot H,s,u)^4,\\[0.45em]
\mbox{E)}&((KLu)^3\cdot H,s,u)^4,
((KH^2u)^4\cdot H,s,u)^4,\\[0.45em]
\mbox{F)}&((j\nu x)^2\cdot H,s,u)^3,
((j\nu x)^2\cdot H,s,u)^4,
((j\eta x)\cdot H,s,u)^4,\\[0.25em]
&(H_j\cdot H,s,u)^3,
(L_j\cdot H,s,u)^4.
\end{array}
\]
\noindent\textit{Редукции:} $\nu$ не входит в свертку аналогично $j$.

\mbox{B)} и \mbox{D)}. $(a_\nu^2\cdot H,s,u)^4$ и $(\theta_\nu^2\cdot H,s,u)^4$ разлагаются из формул (27.)c) и (29.)c), с учетом $(gHu)^2=-\frac13s\cdot\sigma$:
\[
(a_\nu^2\cdot H,s,u)^4=\frac13(asu)^4\cdot\sigma,
\qquad
(\theta_\nu^2\cdot H,s,u)^4=-\frac16(\theta su)^4\cdot\sigma;
\]
отсюда разложение $(a_\eta(aHu)\cdot H,s,u)^4$, $(\theta_\eta(\theta Hu)\cdot H,s,u)^4$.

$(\theta_\nu^2\cdot H,s,u)^3$, $(\theta_\nu^3\cdot H,s,u)^3$ и отсюда $(\theta_\eta^2(\theta Hu)\cdot H,s,u)^4$ разлагаются по \S\ 25, формы 12-го порядка \mbox{C)}.

\noindent 3) \textit{Свертка $s$ с системой III--III}, то есть $s$ с формами системы III:
\[
(3,\lambda)(3,\lambda),\quad (3,\lambda)(2,\lambda),\quad (3,\lambda)(1,\lambda),\quad
(2,\lambda)(2,\lambda),\quad (2,\lambda)(1,\lambda),\quad
(1,\lambda)(1,\lambda).
\]
\noindent\textit{Неприводимые формы:}
\[
\begin{array}{rl}
\mbox{A)}&(a_\nu^2\cdot a_\nu^2,s,u)^3,
(a_\nu^2\cdot a_\nu^2,s,u)^4,
(a_\nu^2\cdot a_\eta(aHu)^2,s,u)^3,\\[0.45em]
\mbox{B)}&(a_\nu\cdot H_j,s,u)^4,
((aHu)^2\cdot H_j,s,u)^3,
((aLu)^2\cdot H_j,s,u)^4,\\[0.25em]
&((aLu)^3\cdot H_j,s,u)^2,
((aH^2u)^3\cdot H_j,s,u)^3.
\end{array}
\]
\noindent\textit{Редукции:}

Произведения II--III при одинаковом порядке в символах $\nu$ и $f$ следует считать более высокими, чем произведения III--III.

Редукции возникают частично из соотношений между произведениями (сизигий), частично путем замены произведений соответствующими разложимыми формами и редукции свертки этих форм с $s$. Произведения, содержащие контраварианты, всегда опускаются, поскольку они переходят в произведения.

\mbox{A)} $f$ квадратично, символы $\nu$ произвольного порядка. По \S\ 11, формулам (12.) и аналогичным вычислением имеем:
\begin{align*}
1)\quad 0&=a_\nu\cdot b_{\nu_1}+\frac12f\cdot(aHu)^2-\frac14\theta\cdot H+\frac12u_xH_\vartheta-\frac14u_x^2\cdot H_\vartheta^2,\\
2)\quad 0&=a_\nu\cdot b_{\nu_1}^2+f\cdot a_\eta(aHu)^2-\theta_\eta-\frac12H_\vartheta,\\
3)\quad 0&=a_\nu^2\cdot b_{\nu_1}^2-H_\vartheta^2+2\theta_\eta^2.
\end{align*}
Отсюда следует:
\begin{align*}
4)\quad 0&=a_\nu^2\cdot b_{\nu_1}^2(j\nu_1x)\quad(\text{из }3)),\\
5)\quad L_\vartheta(\theta Lu)&=-a_\nu\cdot b_\eta(bHu)^2+\frac34a_\nu^2\cdot(bHu)^2+\frac34\theta_\nu^2\cdot H\quad(\text{из }2)),\\
6)\quad L_\vartheta(\theta Lu)&=-6a_\nu\cdot b_\eta(bHu)^2+2a_\nu^2\cdot(bHu)^2-\frac12\theta_\nu^2\cdot H\quad(\text{из }1)),\\
7)\quad 0&=a_\nu\cdot(bHu)^2+f\cdot(aLu)^3+\theta_\nu\cdot H\quad(\text{из }1)),\\
8)\quad 0&=a_\nu\cdot(bLu)^3+\frac34(\theta Hu)^2\cdot H,\\
9)\quad 0&=(aHu)^2\cdot(bHu)^2+2(\theta Hu)^2\cdot H,\\
10)\quad 0&=a_\eta(aHu)^2\cdot b_\eta(bHu)^2\quad(\text{из }5)\text{ и }6)),\\
11)\quad 0&\equiv [a_\nu^2\cdot b_\eta(bHu)]_{\widehat{su^4}}\quad(\text{из формулы (30.)b)}).
\end{align*}
Редукции
\[
(H_\vartheta su)^4,
\qquad (L_\vartheta(\theta Lu),s,u)^3,
\qquad (L_\vartheta(\theta Lu),s,u)^4
\]
(формулы (31.) и (32.)) вместе с формулами 1) до 11) дают редукцию всех сверток $s$ с произведениями, квадратичными в символах $f$ и произвольными по $\nu$.

\mbox{B)} Произведения каждых двух из символов $f,\theta,j$; $\nu$ произвольного порядка. По формулам (17.), (18.), (20.) и прямым вычислением получаем соотношения:
\[
\begin{gathered}
\left.
\begin{aligned}
1)\quad 0&=a_\nu\cdot\theta_{\nu_1}+\frac14\theta\cdot(aHu)^2+\frac14f\cdot(\theta Hu)^2,\\
2)\quad 0&=\theta_\nu\cdot\theta_{\nu_1}+\frac12\theta\cdot(\theta Hu)^2
\end{aligned}
\right\}\\[-0.2em]
\text{(прямое вычисление аналогично A 1).}
\end{gathered}
\]
Отсюда следует:
\begin{align*}
3)\quad 0&=3a_\nu\cdot(\theta Hu)^2+\theta\cdot(aLu)^3+2f\cdot(\theta Lu)^3,\\
4)\quad 0&=3\theta_\nu\cdot(aHu)^2+f\cdot(\theta Lu)^3+2\theta\cdot(aLu)^3,\\
5)\quad 0&=a_\nu\cdot(\theta Lu)^3,\quad 0=\theta_\nu\cdot(aLu)^3,\quad 0=(aHu)^2\cdot(\theta Hu)^2,\\
6)\quad 0&=a_\nu\cdot\theta_\nu^2+\theta_\nu\cdot a_\nu^2+f\cdot\theta_\eta(\theta Hu)^2+\theta\cdot a_\eta(aHu)^2\quad(\text{формула (17.)}),\\
7)\quad 0&=\theta_\nu\cdot\theta_\nu^2+\theta\cdot\theta_\eta(\theta Hu)^2+\frac16f\cdot H_j\\
&\qquad [\text{аналогичное выведение, как для }6)\text{ из }\theta_\nu^2(\theta\theta'\nu_1x)],\\
8)\quad 0&=a_\nu\cdot(j\nu x)^2+f\cdot H_j\quad(\text{из }6)),\\
9)\quad 0&=(aHu)^2\cdot(j\nu x)^2-2a_\nu\cdot H_j\quad(\text{из }8)),\\
10)\quad 0&=\theta_\nu\cdot(j\nu x)^2,\quad 0=\theta\cdot H_j\quad(\text{из }7)\text{ и }8)),
\end{align*}
\[
\left.
\begin{aligned}
11)\quad 0&=a_\nu\cdot\theta_\nu^3-\frac34(j\eta x)^2,\\
12)\quad 0&=a_\nu^2\cdot\theta_\nu^2-\frac13(j\eta x)^2-\frac13(\D Hu)^2
\end{aligned}
\right\}\quad\text{(формула (18.)),}
\]
\begin{align*}
13)\quad 0&=a_\nu^2\cdot\theta_\nu^3-\frac12a_\sigma\quad(\text{формула (20.)}),\\
14)\quad 0&=\theta_\nu\cdot\theta_\nu^3,
&0&=a_\eta H_j\quad(\S\ 13. C)).
\end{align*}
Формулы доведены до того места, где произведения становятся поляризуемыми; то есть при свертке возникает только одно новое произведение, потому что \mbox{a)} свертка произведения с самим собой становится приводимой, \mbox{b)} остальные произведения, возникающие при свертке, становятся приводимыми или ведут к контравариантам (ср. \S\ 3. \mbox{2)}, \mbox{a)} и \mbox{b)}).

Формулы 1) до 5) дают редукцию всех сверток $s$ с произведениями, возникающими из $a_\nu\theta_\nu$ при свертке. Формулы 6) до 10) дают редукцию тех произведений, которые возникают из $a_\nu\theta_\nu^2$ при свертке. По \S\ 24 \mbox{A)}, с учетом формул 6) до 10), свертки $s$ с произведениями, возникающими из $a_\nu^2\theta_\nu$, разлагаются.

Имеем:
\begin{align*}
0&\equiv a_\nu^2(asu)^2\theta_{\nu_1}(\theta su)^2,
&0&\equiv a_\nu^2(asu)\theta_{\nu_1}(\theta su)^3-K_\eta(KHu)^2(Ksu)^3,\\
0&\equiv a_\nu^2(asu)^2\theta_{\nu_1}(\theta su),
&0&\equiv a_\nu^2(asu)\theta_{\nu_1}(\theta su)^2,\\
0&\equiv a_\nu^2(asu)\theta_{\nu_1}(\theta su).
\end{align*}
Редуценты:
\[
 ((j\eta x)^2,s,u)^3\ [\hbox{из }(\vartheta\eta\widehat{su})^2(Hsu).\ \S\ 23.\ \mbox{C\ 3})],
\]
\[
(H_j(j\eta x),s,u)^2\ [\S\ 24,\ \mbox{B}],
\quad ((\D Hu)^2,s,u)^2\ [\S\ 24,\ \mbox{C}],
\quad (a_\sigma su)^2\ [\S\ 24,\ \mbox{E}]
\]
вместе с формулами 11) до 14) дают редукцию всех произведений, возникающих из $a_\nu\theta_{\nu_1}^3$, $a_\nu^2\theta_{\nu_1}^2$, $a_\nu^2\theta_{\nu_1}^3$.

Наши предшествующие вычисления можно суммировать в следующем \textit{выводе}:

Относительно полная система по модулю $(\R,t)$ состоит из следующих 331 форм и подсистем; ниже мы также приводим их в табличном порядке.
\[
\left.
\begin{array}{lr}
u_x&1\ \text{форма}\\
\text{система I}&6\ \text{форм}\\
\text{система II}&5\ \text{форм}\\
\text{система III}&117\ \text{форм}
\end{array}
\right\}
\quad
\begin{gathered}
\text{Относительно полная система по модулю }(s,\varrho,t)\\
{}[\text{см. таблицу I}].
\end{gathered}
\]
\[
\left.
\begin{array}{lr}
s&1\ \text{форма}\\
\text{система }s_I&35\ \text{форм}\\
\text{система }s_{II}&9\ \text{форм}\\
\text{система }s_{III}&125\ \text{форм}\\
\text{система }s_{IV}&32\ \text{формы}
\end{array}
\right\}
\quad [\text{см. таблицу II}].
\]

\clearpage
\begin{landscape}
\thispagestyle{plain}
\begin{center}\bfseries Таблица I (см. с. 90)\end{center}
\vspace{-0.6em}
\begingroup
\setlength{\tabcolsep}{2pt}
\renewcommand{\arraystretch}{1.35}
\tiny
\begin{center}
\begin{adjustbox}{max width=0.98\linewidth,max totalheight=0.78\textheight,center}
\begin{tabular}{|c|c|c|c|c|c|c|c|c|c|c|c|c|c|c|c|c|c|c|c|c|c|c|}
\hline
\diagbox{$u$}{$x$} & 0 & 1 & 2 & 3 & 4 & 5 & 6 & 7 & 8 & 9 & 10 & 11 & 12 & 13 & 14 & 15 & 16 & 17 & 18 & 19 & 20 & 21 \\ \hline
0 & \mcell{i} &  &  &  & \mcell{f} &  & \mcell{\A} &  & \mcell{K_{\nu}(k\nu x)^3} &  & \mcell{K_{\eta}(k\eta x)^3} &  &  &  & \mcell{(\vartheta^3\eta x)^4} & \mcell{H_k(k\cdot\vartheta,\eta x)^4} &  & \mcell{\theta_l(\vartheta\cdot k,lx)^5} &  &  &  & \mcell{(\vartheta^2lx)^6} \\ \hline
1 &  & \mcell{u_x} &  &  &  & \mcell{\theta_\nu(\vartheta\nu x)^2} &  & \mcell{\theta_\eta(\vartheta\eta x)^2} & \mcell{N} & \mcell{(k\nu x)^3} & \mcell{\theta_\nu(\vartheta^2\nu x)^3} & \mcell{(k\eta x)^3} & \mcell{H_{\vartheta}(\vartheta^2\eta x)^3\\ \theta_\eta(\vartheta^2\eta x)^3} &  & \mcell{\theta_l(\vartheta^2lx)^4} &  & \mcell{(\vartheta k,\eta,x)^4} &  & \mcell{(\vartheta k,l,x)^5} &  &  &  \\ \hline
2 &  &  & \mcell{a_\nu^2} &  & \mcell{\theta\\ a_\eta^2} &  & \mcell{(\vartheta\nu x)^2} & \mcell{K_\nu(k\nu x)^2} & \mcell{(\vartheta\eta x)^2} & \mcell{H_k(k\eta x)^2\\K_\eta(k\eta x)^2} &  & \mcell{(\vartheta^2\nu x)^3\\ K_l(klx)^3} &  & \mcell{(\vartheta^2\eta x)^3} &  & \mcell{(\vartheta^2lx)^4} &  &  &  &  &  &  \\ \hline
3 &  & \mcell{\theta_\nu^3} &  & \mcell{a_\nu} & \mcell{\theta_\nu(\vartheta\nu x)} & \mcell{\A_\nu\\a_\eta} & \mcell{K\\H_{\vartheta}(\vartheta\eta x)\\\theta_\eta(\vartheta\eta x)} & \mcell{\A_\eta} & \mcell{(k\nu x)^2\\\theta_l(\vartheta lx)^2} &  & \mcell{(k\eta x)^2} &  & \mcell{(klx)^3} &  &  &  &  &  &  &  &  &  \\ \hline
4 & \mcell{\nu} &  & \mcell{H\\\theta_\nu^2} & \mcell{(j\nu x)^3\\a_\eta(aHu)} & \mcell{\theta_\eta^2} & \mcell{(\vartheta\nu x)\\a_l^2} &  & \mcell{N_\nu\\(\vartheta\eta x)} &  & \mcell{H_\eta\\N_\eta\\(\vartheta lx)^2} &  &  &  &  &  &  &  &  &  &  &  &  \\ \hline
5 &  & \mcell{a_\eta(aHu)^2} & \mcell{\theta_\eta^2(\theta Hu)} & \mcell{\theta_\nu} & \mcell{\frac{K_\nu^2}{(aHu)}} & \mcell{\theta_\eta} & \mcell{K_\eta^2\\(\A Hu)\\a_l} &  & \mcell{\A_l} &  &  &  &  &  &  &  &  &  &  &  &  &  \\ \hline
6 & \mcell{j\\\sigma} &  & \mcell{(j\nu x)^2\\(aHu)^2} & \mcell{L\\a_\nu^2(j\nu x)\\H_{\vartheta}(\theta Hu)\\\theta_\eta(\theta Hu)} & \mcell{\frac{K_\nu}{\theta_l^2}} &  &  & \mcell{K_\eta} &  &  &  &  &  &  &  &  &  &  &  &  &  &  \\ \hline
7 &  & \mcell{\theta_\eta(\theta Hu)^2} & \mcell{H_j(j\eta x)\\a_l(aLu)^2} &  & \mcell{a_\nu(j\nu x)\\(\theta Hu)} &  & \mcell{\A_\nu(j\nu x)\\\theta_l} & \mcell{K_l^2} &  &  &  &  &  &  &  &  &  &  &  &  &  &  \\ \hline
8 & \mcell{H_j^2\\a_l(aLu)^3} & \mcell{(\nu\sigma x)\\(j\nu x)} & \mcell{(\theta Hu)^2} & \mcell{K_\eta(KHu)^2\\(j\eta x)\\(aLu)^2} &  &  &  &  &  &  &  &  &  &  &  &  &  &  &  &  &  &  \\ \hline
9 &  & \mcell{H_j\\(aLu)^3} & \mcell{a_\eta(aHu)H_j\\L_{\vartheta}(\theta Lu)^2\\\theta_l(\theta Lu)^2} &  & \mcell{(KHu)} &  &  &  &  &  &  &  &  &  &  &  &  &  &  &  &  &  \\ \hline
10 & \mcell{\theta_l(\theta Lu)^3} & \mcell{L_j\\(j\sigma x)\\a_\eta(aH^2u)^3} &  & \mcell{H_j(aHu)\\(\theta Lu)^2} &  &  &  &  &  &  &  &  &  &  &  &  &  &  &  &  &  &  \\ \hline
11 &  & \mcell{(\theta Lu)^3} & \mcell{K_l(KLu)^3\\L_j\\(aH^3u)^3} &  &  &  &  &  &  &  &  &  &  &  &  &  &  &  &  &  &  &  \\ \hline
12 & \mcell{(aH^2u)^4} & \mcell{a_l(aLu)^2L_j\\H_{\vartheta}(\theta H^2u)^3\\\theta_\eta(\theta H^2u)^3} &  & \mcell{(KLu)^3} &  &  &  &  &  &  &  &  &  &  &  &  &  &  &  &  &  &  \\ \hline
13 & \mcell{(\nu\sigma j)} &  & \mcell{(aLu)^2L_j\\(\theta H^2u)^3} &  &  &  &  &  &  &  &  &  &  &  &  &  &  &  &  &  &  &  \\ \hline
14 & \mcell{(\theta H^2u)^4} & \mcell{K_\eta(KH^2u)^4\\(a,H\!\cdot L,u)^4} &  &  &  &  &  &  &  &  &  &  &  &  &  &  &  &  &  &  &  &  \\ \hline
15 & \mcell{L_9(\theta,H\!\cdot L,u)^4} &  & \mcell{(KH^2u)^4} &  &  &  &  &  &  &  &  &  &  &  &  &  &  &  &  &  &  &  \\ \hline
16 & \mcell{(\theta,H\!\cdot L,u)^5} &  &  &  &  &  &  &  &  &  &  &  &  &  &  &  &  &  &  &  &  &  \\ \hline
17 & \mcell{K_\eta(K,H\!\cdot L,u)^5} &  &  &  &  &  &  &  &  &  &  &  &  &  &  &  &  &  &  &  &  &  \\ \hline
18 &  & \mcell{(K,H\!\cdot L,u)^5} &  &  &  &  &  &  &  &  &  &  &  &  &  &  &  &  &  &  &  &  \\ \hline
19 &  &  &  &  &  &  &  &  &  &  &  &  &  &  &  &  &  &  &  &  &  &  \\ \hline
20 &  &  &  &  &  &  &  &  &  &  &  &  &  &  &  &  &  &  &  &  &  &  \\ \hline
21 & \mcell{(KH^3u)^6} &  &  &  &  &  &  &  &  &  &  &  &  &  &  &  &  &  &  &  &  &  \\ \hline
\end{tabular}
\end{adjustbox}
\end{center}
\vspace{-0.2em}
\begin{center}\footnotesize (Числа в верхней горизонтальной строке обозначают степень по $x$; числа в первой вертикальной строке обозначают степень по $u$.)\end{center}
\endgroup
\end{landscape}
\clearpage
\begin{landscape}
\thispagestyle{plain}
\begin{center}\bfseries Таблица II (см. с. 90)\end{center}
\vspace{-0.6em}
\begingroup
\setlength{\tabcolsep}{2pt}
\renewcommand{\arraystretch}{1.35}
\tiny
\begin{center}
\begin{adjustbox}{max width=0.98\linewidth,max totalheight=0.78\textheight,center}
\begin{tabular}{|c|c|c|c|c|c|c|c|c|c|c|c|c|c|c|c|c|c|}
\hline
\diagbox{$u$}{$x$} & 0 & 1 & 2 & 3 & 4 & 5 & 6 & 7 & 8 & 9 & 10 & 11 & 12 & 13 & 14 & 15 & 16 \\ \hline
0 & \mcell{J} &  &  &  & \mcell{s} &  & \mcell{s_\vartheta^2} &  &  &  &  & \mcell{s_\nu} & \mcell{(k\nu x)^3s_\nu} & \mcell{(\vartheta^3\nu x)^3s_\nu s_\vartheta\\\theta_\nu(\vartheta^2\nu x)^3s_\vartheta} &  & \mcell{\theta_\eta(\vartheta^2\eta x)^3s_\vartheta} &  \\ \hline
1 &  &  &  &  &  &  & \mcell{(asu)} & \mcell{s_\vartheta} & \mcell{s_k^2\\(\Delta su)} & \mcell{s_\nu(Nsu)\\(\vartheta\nu x)^2s_\nu} & \mcell{((k\nu x)^3,s,u)s_\nu\\(K_\nu(k\nu x)^3,s,u)} &  & \mcell{(K_\eta(k\eta x)^3,s,u)} &  & \mcell{(\vartheta^2\nu x)^3s_\nu} &  & \mcell{((\vartheta^2\eta x)^4,s,u)} \\ \hline
2 &  &  &  &  & \mcell{(asu)^2} & \mcell{s_\vartheta(\theta su)} & \mcell{(\Delta su)^2\\a_\nu s_\nu} & \mcell{((\vartheta\nu x)^2,s,u)s_\nu\\(\theta_\nu(\vartheta\nu x)^2,s,u)} &  & \mcell{s_k\\(\theta_\eta(\vartheta\eta x)^2,s,u)} &  & \mcell{(k\nu x)^2s_\nu\\((k\nu x)^3,s,u)} &  & \mcell{((k\eta x)^3,s,u)} &  &  &  \\ \hline
3 &  &  & \mcell{(asu)^3} & \mcell{s_\vartheta(\theta su)^2\\s_\nu} & \mcell{a_\nu s_\nu(asu)\\a_\nu^2(asu)} & \mcell{s_\eta} & \mcell{(\theta su)\\(a_\eta^2su)} & \mcell{s_k(Ksu)} & \mcell{(Nsu)^2\\(\vartheta\nu x)s_\nu\\((\vartheta\nu x)^2,s,u)} & \mcell{((k\nu x)^3,s,u)^2} & \mcell{((\vartheta\eta x)^2,s,u)} &  &  &  &  &  &  \\ \hline
4 & \mcell{(asu)^4} & \mcell{s_\vartheta(\theta su)^3} & \mcell{a_\nu^2(asu)^2} & \mcell{(\theta_\nu^3su)} & \mcell{(\theta su)^2} & \mcell{s_k(Ksu)^2\\a_\nu(asu)\\s_\nu(asu)} & \mcell{((\vartheta\nu x)^2,s,u)^2\\\theta_\nu s_\nu} & \mcell{s_\nu(\Delta su)\\(\Delta_\nu su)\\(aHu)s_\eta\\(a_\eta su)} &  & \mcell{(\Delta_\eta su)} &  &  &  &  &  &  &  \\ \hline
5 &  &  & \mcell{(\theta su)^3} & \mcell{s_k(Ksu)^3,s_j\\s_\vartheta(\theta s^\prime u)^4\\s_\nu(asu)^2\\a_\nu(asu)^2} & \mcell{(Hsu)\\((\vartheta\nu x)^2,s,u)^3\\(\theta_\nu(\vartheta\nu x),s,u)^2\\(\theta_\nu^2su)} & \mcell{((j\nu x)^3,s,u)\\(aHu)^2s_\eta\\(a_\eta su)^2} & \mcell{(Ksu)^2\\s_l\\(\theta_\eta^2su)} &  & \mcell{((k\nu x)^2,s,u)^2\\K_\nu s_\nu} &  &  &  &  &  &  &  &  \\ \hline
6 & \mcell{(\theta su)^4} & \mcell{s_\nu(asu)^3\\a_\nu(asu)^3} & \mcell{(Hsu)^2\\(\theta_\nu(\vartheta\nu x),s,u)^3\\(\theta_\nu^2su)^2} & \mcell{(a_\eta(aHu),s,u)^2\\(a_\eta(aHu)^2,s,u)} & \mcell{(Ksu)^3} & \mcell{((\vartheta\nu x),s,u)^3\\(\theta_\nu^2su)} & \mcell{((k\nu x)^2,s,u)^3} & \mcell{((\vartheta\eta x),s,u)\\(\theta Hu)s_\eta\\(\theta_\eta su)} &  &  &  &  &  &  &  &  &  \\ \hline
7 & \mcell{(\theta_\nu(\vartheta\nu x),s,u)^4} & \mcell{(a_\eta(aHu),s,u)^3} & \mcell{(Ksu)^4\\(\theta_\eta^2(\theta Hu),s,u)^2\\(a_\nu^2\cdot a_{\nu_1}^2,s,u)^3} & \mcell{s_j(asu)^2\\s_k(Ks^2u)^3\\((j\nu x),s,u)^3\\(\theta_\nu su)^2} & \mcell{(j\nu x)s_\nu\\((j\nu x)^2,s,u)\\((aHu),s,u)^2\\((aHu)^2,s,u)} & \mcell{((\vartheta\eta x),s,u)^3\\(\theta_\eta su)^2} & \mcell{(aLu)^3s_l\\(asu)^3} &  &  &  &  &  &  &  &  &  &  \\ \hline
8 & \mcell{(a_\nu^2\cdot a_\nu^2,s,u)^4} & \mcell{s_j(asu)^3\\s_k(Ks^2u)^6\\((\vartheta\nu x),s,u)^4\\(\theta_\eta su)^3} & \mcell{((j\nu x)^2,s,u)^2\\((aHu),s,u)^3\\((aHu)^2,s,u)^3} & \mcell{(Lsu)^2\\(a_\nu^2(j\nu x),s,u)^2\\(H_\vartheta(\theta Hu),s,u)^2\\(\theta_\eta(\theta Hu),s,u)^2\\(\theta_\eta(\theta Hu)^2,s,u)} & \mcell{(a_lsu)^3} & \mcell{(K_\nu su)^2} &  & \mcell{(K_\eta su)^2} &  &  &  &  &  &  &  &  &  \\ \hline
9 & \mcell{(K_\nu^2su)^4\\((aHu),s,u)^4} & \mcell{(Lsu)^2\\(a_\nu^2(j\nu x),s,u)^3\\(\theta_\eta su)^4\\(\theta_\eta(\theta Hu),s,u)^3} & \mcell{(Ks^2u)^6\\(a_l(aLu)^2,s,u)^2\\(a_\nu^2\cdot H,s,u)^3} & \mcell{(K_\nu su)^3} & \mcell{(a_\nu(j\nu x),s,u)^2\\((\theta Hu),s,u)^2\\((\theta Hu)^2,s,u)} &  & \mcell{(\theta_lsu)^3} &  &  &  &  &  &  &  &  &  &  \\ \hline
10 &  & \mcell{(K_\nu su)^4\\(a_\nu^2\cdot a_\eta(aHu)^2,s,u)^3} & \mcell{(a_\nu(j\nu x),s,u)^3\\((\theta Hu),s,u)^3\\((\theta Hu)^2,s,u)^2} & \mcell{((j\eta x),s,u)^2\\(H_jsu)\\((aLu)^2,s,u)^2\\((aLu)^3,s,u)} &  &  &  &  &  &  &  &  &  &  &  &  &  \\ \hline
11 & \mcell{(a_\nu(j\nu x),s,u)^4\\((\theta Hu),s,u)^4} & \mcell{(K_\eta(KHu)^2,s,u)^3\\((j\eta x),s,u)^3\\((aLu)^2,s,u)^4\\(a_\nu\cdot H,s,u)^4} & \mcell{(H^2su)^3\\(\theta_l(\theta Lu)^2,s,u)^2} &  & \mcell{((KHu)^2,s,u)^2} &  &  &  &  &  &  &  &  &  &  &  &  \\ \hline
12 &  & \mcell{(a_\eta(aHu)^2\cdot H,s,u)^3} & \mcell{((KHu)^2,s,u)^3} & \mcell{((aHu)H_j,s,u)^2\\((\theta Lu)^2,s,u)^2\\((\theta Lu)^3,s,u)} &  &  &  &  &  &  &  &  &  &  &  &  &  \\ \hline
13 & \mcell{((KHu)^2,s,u)^4} & \mcell{((aHu)H_j,s,u)^3\\((\theta Lu)^2,s,u)^3\\(\theta_\nu\cdot H,s,u)^4} & \mcell{(L_jsu)^2\\((aH^2u)^3,s,u)^2\\((j\nu x)^2\cdot H,s,u)^3\\((aHu)^2\cdot H,s,u)^3} &  &  &  &  &  &  &  &  &  &  &  &  &  &  \\ \hline
14 & \mcell{((j\nu x)^2\cdot H,s,u)^4\\((aHu)^2\cdot H,s,u)^4} & \mcell{(\theta_\eta(\theta Hu)^2\cdot H,s,u)^3} &  & \mcell{((KLu)^3,s,u)^2} &  &  &  &  &  &  &  &  &  &  &  &  &  \\ \hline
15 & \mcell{(a_l(aLu)^2\cdot H,s,u)^4} & \mcell{((KLu)^3,s,u)^3} & \mcell{((aLu)^2L_j,s,u)^2\\((\theta H^2u)^3,s,u)^2\\((\theta Hu)^2\cdot H,s,u)^3} &  &  &  &  &  &  &  &  &  &  &  &  &  &  \\ \hline
16 & \mcell{((\theta Hu)^2\cdot H,s,u)^4\\(a_\nu\cdot H_j,s,u)^4} & \mcell{((j\eta x)\cdot H,s,u)^4\\(H_j\cdot H,s,u)^3\\((aLu)^2\cdot H,s,u)^4\\((aLu)^3\cdot H,s,u)^3} &  &  &  &  &  &  &  &  &  &  &  &  &  &  &  \\ \hline
17 & \mcell{(\theta_l(\theta Lu)^2\cdot H,s,u)^4} &  & \mcell{((KH^3u)^4,s,u)^3} &  &  &  &  &  &  &  &  &  &  &  &  &  &  \\ \hline
18 &  & \mcell{((\theta Lu)^2\cdot H,s,u)^4\\((\theta Lu)^3\cdot H,s,u)^3\\((aHu)^2\cdot H_j,s,u)^3} &  &  &  &  &  &  &  &  &  &  &  &  &  &  &  \\ \hline
19 & \mcell{((aH^3u)^3\cdot H,s,u)^4\\(L_j\cdot H,s,u)^4} &  &  &  &  &  &  &  &  &  &  &  &  &  &  &  &  \\ \hline
20 &  & \mcell{((KLu)^3\cdot H,s,u)^4} & \mcell{((aLu)^3\cdot H_j,s,u)^2} &  &  &  &  &  &  &  &  &  &  &  &  &  &  \\ \hline
21 & \mcell{((\theta H^2u)^3\cdot H,s,u)^4\\((aLu)^2\cdot H_j,s,u)^4} &  &  &  &  &  &  &  &  &  &  &  &  &  &  &  &  \\ \hline
22 &  &  &  &  &  &  &  &  &  &  &  &  &  &  &  &  &  \\ \hline
23 & \mcell{((KH^2u)^4\cdot H,s,u)^4} & \mcell{((aH^2u)^3\cdot H,s,u)^3} &  &  &  &  &  &  &  &  &  &  &  &  &  &  &  \\ \hline
\end{tabular}
\end{adjustbox}
\end{center}
\vspace{-0.2em}
\begin{center}\footnotesize (Числа в верхней горизонтальной строке обозначают степень по $x$; числа в первой вертикальной строке обозначают степень по $u$.)\end{center}
\endgroup
\end{landscape}
% END UNIT BODY
% END PRODUCER UNIT 027
\endgroup


\clearpage
\begingroup
\setcounter{footnote}{0}
% BEGIN PRODUCER UNIT 028
% Source: C:/Users/Floris/Documents/interlanguage/03_projects/noether/02_slavic_working_corpus/translations/paper03/russian/v001/Noether_Paper03_Russian_v001.tex
% Identity: 13311 B / SHA-256 67EC541EFB17335299516884AF397C443F624733E624D1E305372498CA3875C7
\setlength{\parindent}{1.25em}
\setlength{\parskip}{0pt}
\emergencystretch=4em
\allowdisplaybreaks
\newcommand{\eps}{\varepsilon}
\newcommand{\pair}[2]{\left(#1\,\middle|\,#2\right)}
% BEGIN UNIT BODY
\begin{center}
{\Large\bfseries 3. К инвариантной теории форм от $n$ переменных\footnote{Доклад, прочитанный на Зальцбургском собрании естествоиспытателей 1909 года.}}\par
\vspace{1.2em}
\emph{Jahresbericht d. DMV} 19 (1910), с. 101--104
\end{center}

\vspace{1em}
В проективной инвариантной теории форм от $n$ переменных главная задача решена: доказана конечность системы форм. Однако ряд дальнейших вопросов еще не рассмотрен, а именно вопросы, относящиеся к методам исследования связи между формами. Результаты, которые я получила по таким вопросам, я хочу кратко сообщить, но ради ясности сначала намечу постановки задач в тернарной области, где они уже известны.

Я прежде всего имею в виду две фундаментальные теоремы символического метода: теорему о том, что все инвариантные образования можно представить символически, и вторую теорему, по которой все отношения между инвариантами получаются последовательным применением небольшого числа символических тождеств; следовательно, символический метод дает все отношения, не выходя из области инвариантов.\footnote{Работа: \emph{Methoden zur Theorie der ternären Formen}. II. \S\ 6. (Leipzig, Teubner, 1889.)} На этом основаны процессы образования форм, символический процесс свертки и несимволические процессы дифференцирования, теория разложений в ряды и т. п. Сюда же я хотела бы добавить еще одну теорему, доказанную в моей диссертации,\footnote{\emph{Über die Bildung des Formensystems der ternären biquadratischen Form}. \S\ 1. (\emph{Crelle's Journal}, Bd. 134, 1908.)} а именно: все свертки можно породить последовательным применением двух возможных ``когредиентных'' сверток; под этим для символического произведения $s_x t_x u_\sigma u_\tau$ я понимаю две свертки $(stu)$ и $(\sigma\tau x)$. На этой теореме можно построить теорию редуцентов и редукции систем форм, аналогично бинарной области, как я показала в своей диссертации.

Если теперь перейти к $n$-арной области, то уже первая теорема символического метода верна лишь условно. Как известно, уже для кватернарных форм, кроме точечных координат $x$ и координат плоскостей $u$, возникают координаты прямых $p_{ik}$, то есть миноры из двух строк точечных координат. Аналогично в $n$-арной области мы имеем строки переменных $p_\rho$, отдельные элементы которых являются минорами из $\rho$ строк $x$, и строки символов $S^\rho$, элементы которых ведут себя как миноры из $\rho$ строк $s$, когредиентных к $u$. Символическое представление через детерминанты, как известно, действительно только тогда, когда строки символов $S^\rho$ разложены на строки $s$ и эти строки распределены между различными детерминантами; тогда реальный смысл имеют только такие агрегаты детерминантов, которые зависят от $S^\rho$. Но невозможно заранее обозреть, какие именно агрегаты детерминантов это обеспечивают; тем самым теряются все остальные теоремы, названные для тернарной области.

Теперь я хочу указать простой принцип, позволяющий получить символическое представление, явное относительно строк $S^\rho$, и тем самым вернуть также остальные теоремы. Это применение известной теоремы о соответствующих матрицах: ``Каждой строке переменных $p_\rho$ взаимно однозначно соответствует другая строка $q_{n-\rho}$, отдельные элементы которой являются минорами из $n-\rho$ строк $u$, связанных с $\rho$ строками $x$ уравнениями $(u\mid x)=0$.'' Соответственно каждой строке символов $S^\rho$ поставлена в соответствие другая строка $S_{n-\rho}$, ведущая себя так, как если бы она была составлена из $n-\rho$ строк, когредиентных к $x$.

Эта теорема допускает также вторую формулировку, удобную для приложений: ``Каждому матричному произведению $(v^{(1)}\cdots v^{(\rho)}\mid y^{(1)}\cdots y^{(\rho)})$,\footnote{То есть сумме по произведениям соответствующих детерминантов, образованных из матрицы $\rho$ строк $v$ и матрицы $\rho$ строк $y$.} где строки второй матрицы контрагредиентны строкам первой, соответствует детерминант; и обратно, каждый детерминант можно преобразовать в матричное произведение с заранее заданным числом строк $\rho$.'' Тем самым детерминант и матричное произведение выступают как полностью эквивалентные, а первая теорема символического метода принимает форму: ``Все инвариантные образования можно символически представить матричными произведениями.'' Из двух строк символов $S^\sigma$, $T^\tau$ при помощи строк переменных теперь очень легко образовать матричное произведение, явно содержащее эти строки символов и потому представляющее инвариантное образование формы $(S^\sigma\mid p_\sigma)(T^\tau\mid p_\tau)$, а именно:
\begin{equation}
 \pair{q_{n-\sigma-\lambda}S^\sigma}{T_{n-\tau}p_{\tau-\lambda}},
 \qquad (\lambda=1,2,\ldots,\tau\ \text{или}\ n-\sigma).
\tag{1}
\end{equation}

Важнее обратное утверждение: все инвариантные образования можно явно представить матричными произведениями, так что указанный процесс является расширением бинарного и тернарного процесса свертки. Для доказательства этого обратного утверждения необходимо перенести вторую теорему символического метода на $n$-арную область, то есть установить всю совокупность независимых неразложимых тождеств, в которых все строки $S^\rho$, $p_\rho$ рассматриваются как неразложимые. Эти тождества получаются из известных тождеств, содержащих только строки $x$ и $u$,\footnote{E. Pascal в \emph{Memorie d. R. Acc. d. Lincei}. (4) 5 (1888), p. 375.} если заменить теорему о произведении для детерминантов системой теорем о произведении для матриц и стягивать различные матричные произведения в одно именно с помощью этих теорем о произведении. Так приходят к двум дуалистическим формулам, из которых приведу только одну:
\begin{equation}
\small
\begin{aligned}
 \pair{S_1^{\rho_1}S_2^{\rho_2}\cdots S_k^{\rho_k}}{x p_{\rho-1}}
 &=\sum_{i=1}^{k}\eps\,\pair{S_1^{\rho_1}\cdots S_{i-1}^{\rho_{i-1}}S_{i+1}^{\rho_{i+1}}\cdots S_k^{\rho_k}q_{n-\rho_i+1}}{S_{n-\rho_i}x},\\
 \rho&=\sum_{i=1}^{k}\rho_i<n,\qquad \eps=\pm1.
\end{aligned}
\tag{2}
\end{equation}

Единственная специализация, содержащаяся в этих формулах, а именно появление строки $x$, неизбежна; для совершенно общих строк символов и строк переменных мы приходим к ``тождеству разложения'', то есть к тождеству, в котором строка $p_\rho$ разлагается на отдельные строки $x$. Простейший частный случай тождества разложения уже использовал Клебш в своей ``фундаментальной задаче инвариантной теории'' для вывода элементарных ковариантов в теории разложений в ряды. С помощью этих тождеств, которые осуществляют переход от неявных детерминантных выражений к матричному произведению, действительно можно доказать, что матричное представление дает желаемое явное символическое представление.

Для процесса свертки, однако, верна теорема, вполне аналогичная тернарному случаю; на ней точно так же аналогично можно построить теоремы редукции. Если в (1) назвать число $\lambda$ дефектом свертки, а свертки, для которых $\sigma=\tau$, назвать когредиентными, то все свертки можно породить последовательным применением когредиентных сверток дефекта 1. Доказательство этой теоремы существенно опирается на то, что наиболее общие формы можно заменить ``нормальными формами'', то есть формами, которые тождественно обращаются в нуль при применении всех инвариантных дифференциальных операций. В кватернарной области этот последний факт доказал Мертенс;\footnote{\emph{Über invariante Gebilde quaternärer Formen}. \emph{Wiener Berichte}. Bd. 98, 1889.} наше знание наиболее общих дифференциальных операций -- которые, как известно, возникают из процессов свертки заменой строк символов дифференциальными символами -- позволяет, применяя тождество разложения, почти дословно перенести доказательство Мертенса на $n$-арную область.
% END UNIT BODY
% END PRODUCER UNIT 028
\endgroup


\clearpage
\begingroup
\setcounter{footnote}{0}
% BEGIN PRODUCER UNIT 029
% Source: C:/Users/Floris/Documents/interlanguage/03_projects/noether/02_slavic_working_corpus/translations/paper04/russian/v001/Noether_Paper04_Introduction_Russian_v001.tex
% Identity: 15164 B / SHA-256 6EA6562172E6BC373F9EE854A0C5E1E1A8A6910FD367359AB47C9E86CADE163E
\setlength{\parindent}{1.25em}
\setlength{\parskip}{0pt}
\emergencystretch=6em
\allowdisplaybreaks
% BEGIN UNIT BODY
\begin{center}
{\Large\bfseries 4. К инвариантной теории форм от $n$ переменных}\par
\vspace{1.2em}
\emph{Journal f. d. reine u. angew. Math.} 139 (1911), с. 118--154
\end{center}

\section*{Введение.}
В проективной инвариантной теории форм от $n$ переменных вопросы, примыкающие к главной проблеме -- доказательству конечности системы форм, -- почти не рассматривались, если выйти за пределы бинарной и тернарной областей. Это вопросы о взаимной зависимости форм и о методах овладения этой связью между формами. Эти методы существенно опираются -- как это полностью доказано в бинарной и тернарной областях -- на две теоремы, названные Study ``фундаментальными теоремами символического метода'', а именно на теорему о символической представимости форм и теорему о символических тождествах.\footnote{E. Study, \emph{Methoden zur Theorie der ternären Formen} (Leipzig, Teubner, 1889). II. \S\ 6. -- О появлении этих фундаментальных теорем также в инвариантной теории специальных групп преобразований ср. Study, \emph{Das Apollonische Problem} (\emph{Math. Ann.} Bd. 49 [1897]) и \emph{Zur Differentialgeometrie der analytischen Kurven} (\emph{Transactions of the American Math. Society}, Bd. 10 [1909]; в оригинальном издании ошибочно напечатано Bd. 1).} Последняя теорема утверждает, что все отношения между целыми инвариантными образованиями получаются последовательным применением небольшого числа символических тождеств; следовательно, символический метод, и только он, дает знание связи между формами, не выходя из области инвариантов.

На причину трудности при переходе к $n$ переменным уже указал Gordan в своей Эрлангенской программе;\footnote{P. Gordan, \emph{Über das Formensystem binärer Formen}. S. 46. (Leipzig, Teubner, 1875).} она состоит в том, что первая теорема символического метода в ее общей формулировке уже не верна. Как известно, в $n$-арной области появляются строки переменных более высокого уровня $p_\rho$\footnote{Об обозначении ср. \S\ 1.} и аналогично строки символов более высокого уровня $S^\rho$; к этим строкам приходят как тогда, когда, следуя первоначальному подходу Clebsch,\footnote{A. Clebsch, \emph{über eine Fundamentalaufgabe der Invariantentheorie}. (\emph{Abh. der Ges. d. Wiss. zu Göttingen}, Bd. XVII, 1872).} исходят из форм с произвольно многими строками $p_\rho$, так и тогда, когда исходят из форм, содержащих только произвольно много когредиентных строк $x$, по подходу F. Deruyts и A. Capelli.\footnote{Ср. A. Capelli, \emph{Lezioni sulla teoria delle forme algebriche} (Napoli, 1902), \S\S\ 22--24, и приведенную там литературу.} Символического представления, которое явно\footnote{Более точное определение ``явного представления'' в \S\ 2.} содержало бы строки $p_\rho$, $S^\rho$, теперь не существует; необходимо разложить $p_\rho$, $S^\rho$ на отдельные строки $x$ и контрагредиентные им $s$ и распределить их между различными детерминантами. Правда, с помощью дифференциальных условий (ср. формулу (19.)) можно проверить, что данный агрегат детерминантов обладает свойством зависеть только от строк $p_\rho$, $S^\rho$; но таким способом невозможно получить обзор всей совокупности инвариантных образований и процессов их порождения.\footnote{Также вычислительно очень ценное дальнейшее развитие символики, восходящее к Study (сообщено F. Engel, \emph{Ber. d. K. Sächs. Ges. d. Wiss., math.-phys. Kl.}, 1900, S. 126, а для кватернарной области -- E. Study, \emph{Geometrie der Dynamen}, S. 135), не является явным представлением в нашем смысле. Например, оно едва ли привело бы к несимволическим процессам дифференцирования для порождения форм.}

Теперь мы покажем в \S\S\ 2 и 6, как, применяя ``теорему о соответствующих матрицах'',\footnote{Если не считать Graßmann-овой \emph{Ausdehnungslehre} 1862 года, Nr. 112, это впервые встречается у A. Brill, \emph{über zwei Berührungsprobleme}, \S\ 2 (\emph{Math. Ann.} Bd. 4, 1871).} можно вновь получить первую фундаментальную теорему во всей ее общности, заменив представление через детерминанты представлением через ``матричные произведения''. В \S\ 2 показывается, что всякое матричное произведение, явно содержащее строки $S^\rho,p_\rho$, является инвариантным образованием основной формы; обратное утверждение, данное в \S\ 6, а именно что всякое инвариантное образование можно явно представить матричными произведениями, удается только через перенос второй фундаментальной теоремы (\S\S\ 3--5).

Эта теорема известна для форм, содержащих только строки переменных $x$.\footnote{E. Pascal, \emph{Giorn. di mat.} 26 (1888), S. 33, 102; \emph{Rendic. d. Accad. d. Linc.} (4) 4 (1888), S. 119; ib. \emph{Mem.} (4) 5 (1888), S. 375.} Общая проблема -- установить всю совокупность неразложимых тождеств, в которых все строки $S^\rho,p_\rho$ рассматриваются как неразложимые, -- была сформулирована Study;\footnote{\emph{Methoden \dots}, S. 204, примечание 17).} одновременно он видит в числе возникающих тождеств меру трудности методов в $n$-арной области. Поскольку теперь удается объединить всю совокупность тождеств в одну формулу (36.), эта трудность по крайней мере сведена к минимуму. В приложениях, однако, часто будут возникать некоторые выделенные специальные случаи (36.), такие как (20.), (21.), характеризующиеся тем, что они допускают явное представление без подстановки новых строк; или формула (31.), называемая ``тождеством разложения'', поскольку строка $p_\rho$ в ней как будто разложена на строки $y$.

На двух фундаментальных теоремах теперь легко построить дальнейшую теорию. Первая теорема непосредственно дает процессы порождения форм, символический процесс свертки и несимволические процессы дифференцирования,\footnote{В работе \emph{Invariante Gebilde von Nullsystemen} (\emph{Sitzungsber. d. Ak. d. Wiss. Wien}, Bd. 97, Abt. IIa, 1888) Mertens другим, не непосредственно обобщаемым путем пришел к кватернарному процессу свертки. Возникающий там альтернированный инвариант от 6 строк $S^2$ отличается от нашего, возникающего при однократном применении подстановки (11.), агрегатом четных инвариантов. Для одного специального случая кватернарный процесс свертки встречается также у P. Gordan, \emph{Das simultane System von zwei quadratischen quaternären Formen} (\emph{Math. Ann.} Bd. 56, 1903).} тогда как тождество разложения устраняет среди этих процессов все эквивалентные (\S\ 6). Знание процессов дифференцирования далее дает перенос понятия ``нормальной формы'' на $n$-арную область и, с помощью способа доказательства, указанного Mertens\footnote{F. Mertens, \emph{Über invariante Gebilde quaternärer Formen}. (\emph{Sitzber. d. Ak. d. Wiss. Wien}. Bd. 98, Abt. IIa, 1889.)} в кватернарной области, теорему о том, что систему инвариантных образований можно заменить эквивалентной системой нормальных форм (\S\ 7). При этом последнем предположении я показала в своей диссертации,\footnote{\emph{Über die Bildung des Formensystems der ternären biquadratischen Form}. (Этот журнал, Bd. 134, 1908.)} что тернарный процесс свертки можно породить последовательным применением двух основных сверток. На основании этой теоремы и теории разложений в ряды я затем, введя понятие ``ряда форм'', развила основные теоремы редукции для систем форм. Совершенно аналогично в $n$-арной области процесс свертки можно породить из $n-1$ основных сверток (\S\ 8), и можно установить теоремы редукции (\S\ 9).

\paragraph{Послесловие.} По поводу Зальцбургского сообщения проф. Е. Мюллер из Вены обратил мое внимание на то, что вычисление с матричными произведениями, лежащее в основе этой работы, в принципе тождественно исчислению Graßmann-овой теории протяженности.\footnote{\emph{Hermann Graßmanns gesammelte mathematische und physikalische Werke. Ersten Bandes zweiter Teil: Die Ausdehnungslehre von 1862. In Gemeinschaft mit H. Graßmann d. Jüngeren herausgegeben v. Fr. Engel} (Leipzig, Teubner, 1896).} Действительно, Graßmann-ово ``комбинаторное произведение'' $[S^{\rho_1}S^{\rho_2}\cdots S^{\rho_i}]$ можно рассматривать как матрицу, заданную этими строками (ср. \S\ 2), причем ``прогрессивное произведение'' является матрицей самих строк, ``регрессивное'' -- матрицей сопоставленных строк $S_{n-\rho_1},S_{n-\rho_2},\ldots,S_{n-\rho_i}$, а матрица, соответствующая ``смешанному произведению'', возникает, когда несколько строк $S$ объединяются подстановкой (9.) в новые строки $P,Q$. Нашей ``сопоставленной строке'' при этом соответствует Graßmann-ово ``дополнение'', а нашему ``матричному произведению'' -- ``смешанное произведение уровня 0''.

При таком сопоставлении разложение, данное в \emph{Ausdehnungslehre} 1862 года, Nr. 113, становится тождественным формуле, дуалистической к нашей формуле (32.). В специальном случае $\rho=1$ формула (32.) переходит в (17.), а дуалистическая формула -- в (16.). Из этих специальных случаев Graßmann-ова разложения Мюллер недавно вывел упомянутые выше тождества (20.), (21.).\footnote{E. Müller, \emph{Beiträge zur Graßmannschen Ausdehnungslehre. 1. Mitteilung: Einige allgemeine Sätze}. (\emph{Sitzungsber. d. Ak. d. Wiss. Wien}; Bd. 118, Abt. IIa, Juli 1909), Satz 16 и 18.}

В отличие от вывода Graßmann--Müller, наш вывод одновременно дает доказательство полноты тождеств, что для рассматриваемых здесь вопросов представляется самым существенным; от (20.), (21.) он далее переходит к наиболее общим неразложимым тождествам (36.), которые, по-видимому, до сих пор еще не были замечены.
% END UNIT BODY
% END PRODUCER UNIT 029
\endgroup


\clearpage
\begingroup
\setcounter{footnote}{0}
% BEGIN PRODUCER UNIT 030
% Source: C:/Users/Floris/Documents/interlanguage/03_projects/noether/02_slavic_working_corpus/translations/paper04/russian/v001/Noether_Paper04_Section01_Russian_v001.tex
% Identity: 8709 B / SHA-256 2A8214E53246EA17B4C297089D22983479ACCC81850291D8A497F0C587A82FC9
\setlength{\parindent}{1.25em}
\setlength{\parskip}{0pt}
\emergencystretch=6em
\allowdisplaybreaks
% BEGIN UNIT BODY
\section*{\S\ 1. Обозначения и определения. Сводка известных результатов.}

Как обычно, будем обозначать через строку переменных $x$ строку элементов $x_1,x_2,\ldots,x_n$, а аналогично через строку переменных $p_\rho$ ($\rho=1,2,\ldots,n-1$) -- строку из $\binom n\rho$ элементов $p_{i_1i_2\ldots i_\rho}$, где $p_{i_1i_2\ldots i_\rho}$ означает $\binom n\rho$ детерминантов, образованных из матрицы $\rho$ строк $x^{(1)},x^{(2)},\ldots,x^{(\rho)}$:
\[
\left|\begin{array}{cccc}
 x_{i_1}^{(1)}&x_{i_2}^{(1)}&\cdots&x_{i_\rho}^{(1)}\\
 x_{i_1}^{(2)}&x_{i_2}^{(2)}&\cdots&x_{i_\rho}^{(2)}\\
 \vdots&\vdots&&\vdots\\
 x_{i_1}^{(\rho)}&x_{i_2}^{(\rho)}&\cdots&x_{i_\rho}^{(\rho)}
\end{array}\right|,
\qquad (i_1<i_2<\cdots<i_\rho=1,2,\ldots,n).
\]
По теореме о соответствующих матрицах\footnote{Ср., например: Gordan--Kerschensteiner, \emph{Vorlesungen über Invariantentheorie}, Bd. I, S. 99.} каждой строке $p_\rho$ взаимно однозначно сопоставлена вторая строка $q_{n-\rho}$ посредством соотношения, справедливого для каждого элемента строки:
\begin{equation}
 p_{i_1i_2\ldots i_\rho}=(-1)^{\frac{\rho(\rho+1)}{2}+i_1+i_2+\cdots+i_\rho}q_{j_1j_2\ldots j_{n-\rho}},
\tag{1}
\end{equation}
где $i$ и $j$ по отдельности упорядочены и вместе образуют полную перестановку чисел от 1 до $n$. Строка $q_{n-\rho}$ состоит из $\binom n\rho$ детерминантов степени $n-\rho$, $q_{j_1j_2\ldots j_{n-\rho}}$, образованных из матрицы $n-\rho$ строк $u^{(1)},u^{(2)},\ldots,u^{(n-\rho)}$, контрагредиентных к $x$. Эти строки удовлетворяют $\rho(n-\rho)$ уравнениям условий:
\begin{equation}
 (u^{(k)}\mid x^{(l)})=\sum_{\alpha=1}^{n}u_\alpha^{(k)}x_\alpha^{(l)}=0;
 \qquad (k=1,2,\ldots,n-\rho;\ l=1,2,\ldots,\rho).
\tag{2}
\end{equation}
Вместе с Clebsch\footnote{a. a. O. \S\ 5.} будем считать их нормированными так, что множитель пропорциональности, который иначе входил бы в (1.), положен равным единице.

По Clebsch\footnote{a. a. O. \S\ 12.} (и так же по более поздним исследованиям F. Deruyts и A. Capelli, ср. Введение), самые общие формы можно свести к таким, которые содержат не более одной из каждой из $n-1$ строк $p_\rho$, то есть символически представляются следующим образом:
\begin{equation}
 (S^1\mid p_1)^{m_1}(S^2\mid p_2)^{m_2}\cdots(S^{n-1}\mid p_{n-1})^{m_{n-1}},
\tag{3}
\end{equation}
где, по аналогии с (2.), положено
\[
 (S^\rho\mid p_\rho)=\sum S^{i_1i_2\ldots i_\rho}p_{i_1i_2\ldots i_\rho}.
\]
Из-за нелинейных тождественных соотношений между элементами строки $p_\rho$ строки символов $S^\rho$ можно нормировать так, чтобы они удовлетворяли тем же самым тождествам, то есть вели себя как детерминанты $\rho$-строчной матрицы, строки которой $s^{(1)},s^{(2)},\ldots,s^{(\rho)}$ контрагредиентны к $x$.\footnote{Clebsch, a. a. O. \S\ 4.} Поэтому каждой строке $S^\rho$ можно взаимно однозначно сопоставить другую строку $S_{n-\rho}$ посредством соотношения:
\begin{equation}
 S_{i_1i_2\ldots i_{n-\rho}}=(-1)^{\frac{(n-\rho)(n-\rho+1)}2+i_1+i_2+\cdots+i_{n-\rho}}S^{j_1j_2\ldots j_\rho},
\tag{4}
\end{equation}
где снова $i$ и $j$ по отдельности упорядочены и вместе в точности исчерпывают числа от 1 до $n$. Элементы $S_{i_1i_2\ldots i_{n-\rho}}$ строки $S_{n-\rho}$ ведут себя как детерминанты, образованные из $(n-\rho)$ строк $\sigma^{(1)},\sigma^{(2)},\ldots,\sigma^{(n-\rho)}$, когредиентных к $x$, а строки $s$ и $\sigma$ связаны $\rho(n-\rho)$ условиями:
\begin{equation}
 (s^{(k)}\mid \sigma^{(l)})=0,\qquad(k=1,2,\ldots,\rho;\ l=1,2,\ldots,n-\rho).
\tag{5}
\end{equation}
Из соотношений (1.) и (4.) следует, что каждый множитель в (3.) допускает эквивалентное четырехкратное символическое представление:
\begin{equation}
\begin{aligned}
 (S^\rho\mid p_\rho)&=(S^\rho q_{n-\rho})=(S_{n-\rho}p_\rho)\\
 &=(-1)^{\rho(n-\rho)}(q_{n-\rho}\mid S_{n-\rho}),
\end{aligned}
\tag{6}
\end{equation}
где, как и всюду далее, $(S^\rho q_{n-\rho})$ и $(S_{n-\rho}p_\rho)$ служат сокращениями для детерминантов $(s^{(1)}\cdots s^{(\rho)}u^{(1)}\cdots u^{(n-\rho)})$ и $(\sigma^{(1)}\cdots\sigma^{(n-\rho)}x^{(1)}\cdots x^{(\rho)})$; разложение Лапласа показывает, что они зависят только от строк $S^\rho,q_{n-\rho}$ соответственно $S_{n-\rho},p_\rho$, на что и должна указывать приведенная запись. Выражения $(S^\rho\mid p_\rho)$ и $(q_{n-\rho}\mid S_{n-\rho})$ означают, согласно сказанному выше, произведения матриц строк $s$ и $x$, соответственно строк $u$ и $\sigma$, причем под матричным произведением понимается сумма произведений соответствующих детерминантов, то есть значение детерминанта матрицы произведения.

Характеристическое число $\rho(n-\rho)$, принадлежащее строке символов (соответственно строке переменных), мы, вслед за Clebsch, называем весом строки;\footnote{a. a. O. \S\ 11.} число $\rho$, задающее число строк $s$, -- верхним весовым индексом, а число $n-\rho$, задающее число строк $\sigma$, -- нижним весовым индексом. Из соотношений (4.) следует, что строка символов может входить в одно и то же соотношение один раз с верхним и другой раз с нижним весовым индексом; то же верно для появления соответствующих друг другу строк $p_\rho$ и $q_{n-\rho}$. Отметим еще, что вообще мы обозначаем: строки переменных более высокого уровня, соответствующая матрица которых состоит из строк $x$, через $p$; те, матрица которых состоит из строк $u$, через $q$; строки символов, когредиентные к $x$, малыми греческими буквами $\sigma,\tau,\alpha,\beta$; строки символов, когредиентные к $u$, малыми латинскими буквами $s,t,a,b$; все остальные строки символов -- большими латинскими буквами с весовым индексом: $S^\rho$, $T^\tau$, $A^\rho$ или $S_{n-\sigma}$, $T_{n-\tau}$, $A_{n-\rho}$. В соответствии с верхним или нижним весовым индексом мы также пишем отдельные элементы строки с верхними или нижними индексами, как, например, в (4.).
% END UNIT BODY
% END PRODUCER UNIT 030
\endgroup


\clearpage
\begingroup
\setcounter{footnote}{0}
% BEGIN PRODUCER UNIT 031
% Source: C:/Users/Floris/Documents/interlanguage/03_projects/noether/02_slavic_working_corpus/translations/paper04/russian/v001/Noether_Paper04_Section02_Russian_v001.tex
% Identity: 13795 B / SHA-256 A042DE7DAA078CEEE123B884A8A79389691396FBA01C18573E8CFC70492C6A00
\setlength{\parindent}{1.25em}
\setlength{\parskip}{0pt}
\emergencystretch=6em
\allowdisplaybreaks
\newcommand{\pair}[2]{\left(#1\middle|#2\right)}
\newcommand{\eps}{\varepsilon}
% BEGIN UNIT BODY
\section*{\S\ 2. Представление матричными произведениями.}

В дальнейшем под ``матричным произведением'' мы всегда понимаем сумму по произведениям соответствующих детерминантов двух матриц с одинаковым числом строк, где строки второй матрицы контрагредиентны строкам первой. При этом строки каждой матрицы могут: 1. быть заданы по отдельности, 2. как в (6.) быть объединены в одну строку $S^\rho$, 3. быть объединены в различные строки $S^\rho,S^\sigma\ldots$. Матричное произведение будем обозначать символом $(\ \mid\ )$, использованным в (6.). Представление (6.) является специальным применением второй формулировки теоремы о соответствующих матрицах: ``Каждому $\rho$-строчному матричному произведению $(v^{(1)}v^{(2)}\cdots v^{(\rho)}\mid y^{(1)}\cdots y^{(\rho)})$ сопоставлен детерминант; и обратно, каждому детерминанту можно сопоставить матричное произведение заранее заданного числа строк $\rho$ ($\rho=1,2,\ldots,n-1$)''.\footnote{Для $\rho>n$ матричное произведение, как известно, равно нулю; для $\rho=n$ оно переходит в произведение двух детерминантов.} В самом деле, достаточно лишь объединить строки $y$ в одну строку $p_\rho$ (соответственно строки $v$ в одну строку $q_\rho$) и выполнить подстановки, заданные (1.), чтобы от данного матричного произведения перейти к детерминанту и обратно. Значение детерминанта при этом, однако, не задается его отдельными элементами, а получается только через разложение Лапласа. Итак, поскольку детерминант и матричное произведение оказываются эквивалентными, теорему о символической представимости форм (первую фундаментальную теорему) можно выразить также следующим образом:

\emph{Теорема I:} ``Все инвариантные образования можно символически представить матричными произведениями, и обратно, все матричные произведения являются инвариантными образованиями''.\footnote{Инвариантными образованиями заранее заданных основных форм являются, конечно, лишь такие агрегаты матричных произведений, которые зависят от строк символов $S^\rho\ldots$.}

Представление матричными произведениями позволяет преодолеть существенную трудность неявного представления, вызванную детерминантным представлением.\footnote{Ср. Введение.} Под ``явным представлением'' мы здесь понимаем такое, в которое входят только сами строки $S^\rho,T^\tau,\ldots$ или строки $R^\rho$, однозначно выводимые из них, но в котором отдельные частичные строки $s,t,\ldots$ уже не появляются. В \S\ 6 мы для всех инвариантных образований, которые зависят только от строк символов $S^\rho,T^\tau,\ldots$, получим явное представление через эти строки символов, а тем самым и перенос порождающих процессов, процесса свертки и процессов дифференцирования, на $n$-арную область. Возможность явного представления становится понятной из того, что только простейшие матричные произведения можно прямо сопоставить детерминантам, возникающим в детерминантном представлении; следовательно, все остальные матричные произведения являются объединениями разных детерминантов в одно выражение. Ведь если разложить строки символов и строки переменных на отдельные строки $s,u,x$, то по первой фундаментальной теореме в ее обычной форме необходимо должен возникнуть агрегат детерминантов. Доказательство того, что, обратно, все агрегаты детерминантов, представляющие инвариантные образования заранее заданных основных форм, можно объединить в явные матричные произведения, получается с использованием второй фундаментальной теоремы (\S\ 6).

Чтобы теперь из двух строк символов $S^\sigma,T^\tau$ с помощью строк переменных получить инвариантное образование основной формы $(S^\sigma\mid p_\sigma)(T^\tau\mid p_\tau)$, явно содержащее все строки (процесс свертки), достаточно образовать матричное произведение третьего типа по определению, данному в начале этого параграфа:
\begin{equation}
 \pair{q_{n-\sigma-\lambda}S^\sigma}{T_{n-\tau}p_{\tau-\lambda}},
 \qquad (\lambda=1,2,\ldots,\tau,n-\sigma).
\tag{7}
\end{equation}
В развернутом виде (7.) дает:
\begin{equation}
\begin{aligned}
&\pair{v^{(1)}\cdots v^{(n-\sigma-\lambda)}s^{(1)}\cdots s^{(\sigma)}}
       {\alpha^{(1)}\cdots\alpha^{(n-\tau)}y^{(1)}\cdots y^{(\tau-\lambda)}} \\
&\qquad =\sum_{i,k,l}\eps\,q_{k_1\ldots k_{n-\sigma-\lambda}}
 S^{k_{n-\sigma-\lambda+1}\ldots k_{n-\lambda}}
 T_{l_1\ldots l_{n-\tau}}
 p^{l_{n-\tau+1}\ldots l_{n-\lambda}},
\end{aligned}
\tag{8}
\end{equation}
то есть выражение, зависящее только от строк $S,T,q,p$. Здесь $\eps=\pm1$\footnote{Это обозначение будет сохраняться всюду далее.}, а сумму следует понимать так: индексы каждой строки $S\ldots$ упорядочены между собой по возрастанию, и для каждой комбинации $i_1,i_2,\ldots,i_{n-\lambda}$ как $k$, так и $l$ по отдельности пробегают все еще возможные перестановки чисел $i$. Число $\lambda$, входящее в (7.), мы называем дефектом свертки в случае $\sigma\geq\tau$; причина этого последнего ограничения выяснится в \S\ 6.\footnote{Число $\lambda$ всегда пробегает до меньшего из двух последних указанных чисел.}

Детерминанты двух матриц, входящих в (7.), можно рассматривать как элементы новых строк $Q^{n-\lambda},P_{n-\lambda}$, сопоставленные которым строки пусть будут $Q_\lambda,P^\lambda$. Эти строки $Q,P$ можно, по (8.), выразить через исходные $S$ и $q$, соответственно через $T$ и $p$. Мы обозначаем это уравнениями:
\begin{equation}
\begin{aligned}
 q_{n-\sigma-\lambda}S^\sigma&=Q^{n-\lambda}\sim Q_\lambda,
 &\quad\text{или также}\quad Q_\lambda&=[q_{n-\sigma-\lambda}S^\sigma]_\lambda,\\
 T_{n-\tau}p_{\tau-\lambda}&=P_{n-\lambda}\sim P^\lambda,
 &\quad\text{или также}\quad P^\lambda&=[T_{n-\tau}p_{\tau-\lambda}]^\lambda.
\end{aligned}
\tag{9}
\end{equation}
С учетом (4.) тогда, аналогично (6.), для матричного произведения (7.) справедливо эквивалентное четырехкратное символическое представление:
\begin{equation}
\begin{aligned}
\pair{q_{n-\sigma-\lambda}S^\sigma}{T_{n-\tau}p_{\tau-\lambda}}
&=(q_{n-\sigma-\lambda}S^\sigma P^\lambda)\\
&=(Q_\lambda T_{n-\tau}p_{\tau-\lambda})\\
&=(-1)^{\lambda(n-\lambda)}\pair{P^\lambda}{Q_\lambda}.
\end{aligned}
\tag{10}
\end{equation}
Еще одну подстановку новых строк можно выполнить в (7.), положив:
\begin{equation}
 \pair{q_{n-\sigma-\lambda}S^\sigma}{T_{n-\tau}p_{\tau-\lambda}}
 =\pair{R^{\sigma+\lambda}}{p_{\sigma+\lambda}}(R^{\tau-\lambda}p_{\tau-\lambda}).
\tag{11}
\end{equation}
И здесь (8.) показывает, что произведения элементов строк $R$ однозначно выражаются через строки $S$ и $T$. Подстановки (9.) и (11.) будут применяться в частности тогда, когда вместе с (7.) выполняется дальнейшая свертка с другими строками.

Матричное представление позволяет нам инвариантно записать квадратичные соотношения между элементами строки $p_\rho$ (из которых, как известно, следуют все остальные), то есть, по \S\ 1, линейные по коэффициентам основной формы соотношения, которым удовлетворяют строки $S^\rho$. Это тождества:
\begin{equation}
 \pair{q_{n-\rho-\lambda}S^\rho}{S_{n-\rho}p_{\rho-\lambda}}=0,
 \qquad (\lambda=1,2,\ldots,\rho,n-\rho),
\tag{12}
\end{equation}
где $p$ и $q$ следует рассматривать как произвольные величины. В \S\ 5 мы увидим, что тождества с нечетным числом дефекта $\lambda$ являются следствием тождеств более высокого дефекта. Соотношения (12.) являются прямым следствием соотношений (5.); действительно, имеем:
\[
\begin{aligned}
&\pair{q_{n-\rho-\lambda}S^\rho}{S_{n-\rho}p_{\rho-\lambda}}\\
&\quad
=(v^{(1)}\cdots v^{(n-\rho-\lambda)}s^{(1)}\cdots s^{(\rho)}
\mid \sigma^{(1)}\cdots\sigma^{(n-\rho)}y^{(1)}\cdots y^{(\rho-\lambda)}),
\end{aligned}
\]
и применение теоремы о произведении вследствие (5.) дает нулевой $(n-\lambda)$-строчный детерминант. Тождества (12.) утверждают, что для отыскания всех типов сверток достаточно сворачивать произведение форм, линейных в каждой строке переменных; или также, по \S\ 6, что нормированная форма $(S^\rho\mid p_\rho)^m$ не имеет инвариантного образования, линейного по коэффициентам.\footnote{Ср. дополнение Study в сочинениях Graßmann, первый том, вторая часть, с. 510 (условие, что величина $S^\rho$ является простой).}

Условия (12.) также достаточны для характеристики строк $p_\rho$. Поскольку свойство того, что отдельные элементы $p_\rho$ являются детерминантами матрицы, инвариантно, оно должно также иметь инвариантное выражение; а условия (12.), по теореме VI, задают наибольшее возможное инвариантное ограничение для $p_\rho$, а именно то, что линейная форма, образованная из строки $p_\rho$ как из коэффициентов, вообще не имеет инвариантных образований.
% END UNIT BODY
% END PRODUCER UNIT 031
\endgroup


\clearpage
\begingroup
\setcounter{footnote}{0}
% BEGIN PRODUCER UNIT 032
% Source: C:/Users/Floris/Documents/interlanguage/03_projects/noether/02_slavic_working_corpus/translations/paper04/russian/v001/Noether_Paper04_Section03_Russian_v001.tex
% Identity: 16026 B / SHA-256 95A5C13A93CAEF1AB5956A955E5F577465E9C17AC8611A2B3742E42F3C059CBC
\setlength{\parindent}{1.25em}
\setlength{\parskip}{0pt}
\emergencystretch=8em
\allowdisplaybreaks
\newcommand{\pair}[2]{\left(#1\middle|#2\right)}
\newcommand{\eps}{\varepsilon}
% BEGIN UNIT BODY
\section*{\S\ 3. Символические тождества.}

Теперь нам нужно перенести вторую фундаментальную теорему символического метода, а именно установить всю совокупность неприводимых неразложимых тождеств, в которых все строки $S^\rho,p_\rho$ рассматриваются как неразложимые. При этом условимся еще о следующем. Соответственно значению строк символов и строк переменных составными следует считать те тождества, которые возникают, когда строку переменных $p_\rho$ специализируют до $p_{\rho-1}y$ и к возникающим выражениям применяют одно из тождеств системы; с другой стороны, тождества, содержащие $k$ строк символов, должны считаться неразложимыми даже тогда, когда они возникают из тождеств с меньшим числом строк символов при помощи указанного процесса. Строки $S^\rho,p_\rho$ не обязаны входить в тождества явно; для их понимания как неразложимых величин достаточно доказать, что возникающие матричные произведения зависят только от этих строк. Но мы покажем, что в этом случае, отчасти с применением подстановки (11.), всегда можно получить явное представление. Общие тождества мы получаем при помощи принципа матричного представления из специальных тождеств, данных Pascal, которые содержат только строки $x$ и $u$ и являются прямым переносом тождеств, данных Study для тернарной области:\footnote{О формулировке задачи и литературе ср. Введение.}

\begin{equation}
\begin{aligned}
&\sum_{i=1}^n(-1)^{i+1}\pair{a^{(i)}}{x}
(a^{(1)}\cdots a^{(i-1)}a^{(i+1)}\cdots a^{(n)}u)\\
&\qquad =(-1)^{n+1}\pair{u}{x}(a^{(1)}\cdots a^{(n)}).
\end{aligned}
\tag{13}
\end{equation}
\begin{equation}
\begin{aligned}
&\sum\pm\pair{a^{(1)}}{x^{(1)}}\pair{a^{(2)}}{x^{(2)}}\cdots
\pair{a^{(n)}}{x^{(n)}}\\
&\qquad=(a^{(1)}\cdots a^{(n)})(x^{(1)}\cdots x^{(n)}).
\end{aligned}
\tag{14}
\end{equation}
\begin{equation}
\begin{aligned}
&\sum_{i=1}^n(-1)^{i+1}\pair{u}{a^{(i)}}
(a^{(1)}\cdots a^{(i-1)}a^{(i+1)}\cdots a^{(n)}x)\\
&\qquad =(-1)^{n+1}\pair{u}{x}(a^{(1)}\cdots a^{(n)}).
\end{aligned}
\tag{15}
\end{equation}

Еще два тождества получаются из (13.) и (15.) подстановками $(a^{(i)}\mid x)=(a^{(i)}q_{n-1})$ соответственно $(u\mid a^{(i)})=(p_{n-1}a^{(i)})$, так что в нашем понимании их не следует считать существенно отличными. (14.) представляет теорему о произведении для детерминантов; мы покажем, как можно заменить (13.) и (14.) (а дуалистически также (15.) и (14.)) системой подходящим образом образованных теорем о произведении для матриц. Теорема о произведении, образованная один раз для $\rho$-строчных и один раз для $(\rho-1)$-строчных матриц, имеет вид:
\[
\begin{aligned}
\pair{a^{(1)}a^{(2)}\cdots a^{(\rho)}}{xp_{\rho-1}}
&=(a^{(1)}a^{(2)}\cdots a^{(\rho)}\mid xy^{(2)}\cdots y^{(\rho)})\\
&=\sum\pm\pair{a^{(1)}}{x}\pair{a^{(2)}}{y^{(2)}}\cdots\pair{a^{(\rho)}}{y^{(\rho)}},
\end{aligned}
\]
\[
\begin{aligned}
&\sum\pm\pair{a^{(2)}}{y^{(2)}}\cdots\pair{a^{(\rho)}}{y^{(\rho)}}\\
&\quad=(a^{(2)}\cdots a^{(\rho)}\mid y^{(2)}\cdots y^{(\rho)})\\
&\quad=(a^{(2)}\cdots a^{(\rho)}p_{\rho-1})
=(a^{(2)}\cdots a^{(\rho)}q_{n-\rho+1});
\end{aligned}
\]
и отсюда следует система теорем о произведении, где $\rho=2,3,\ldots,n$:
\begin{equation}
\begin{aligned}
&\pair{a^{(1)}a^{(2)}\cdots a^{(\rho)}}{xp_{\rho-1}}\\
&\quad=\sum_{i=1}^{\rho}(-1)^{i+1}\pair{a^{(i)}}{x}
(a^{(1)}\cdots a^{(i-1)}a^{(i+1)}\cdots a^{(\rho)}\mid p_{\rho-1})\\
&\quad=\sum_{i=1}^{\rho}(-1)^{i+1}\pair{a^{(i)}}{x}
(a^{(1)}\cdots a^{(i-1)}a^{(i+1)}\cdots a^{(\rho)}q_{n-\rho+1}).
\end{aligned}
\tag{16}
\end{equation}
При $\rho=n$ (16.) переходит в (13.); если далее специализировать $p_{n-1}=y^{(2)}p_{n-2}$, $p_{n-2}=y^{(3)}p_{n-3}$ и т. д. и применить тождества (16.) к выражениям $(a^{(1)}\cdots a^{(i-1)}a^{(i+1)}\cdots a^{(n)}\mid y^{(2)}p_{n-2})$, $(a^{(1)}\cdots a^{(i-1)}a^{(i+1)}\cdots a^{(h-1)}a^{(h+1)}\cdots a^{(n)}\mid y^{(3)}p_{n-3})$ и т. д., то мы переходим от (13.) к (14.). Следовательно, тождество (14.) по нашему определению составлено из тождеств (16.); или система (16.) эквивалентна тождествам (13.) и (14.). Система (16.) удовлетворяет одному из условий, требуемых для общих тождеств: она содержит строку $p_{\rho-1}$ как неразложимую величину и одновременно в явной форме. Это единственная система, удовлетворяющая этому и только этому требованию. Ведь из строк $a,x,p$ мы не можем образовать никаких других детерминантов или матричных произведений; а между ними не может быть соотношений, независимых от (16.), поскольку иначе после исключения $(a^{(1)}\cdots a^{(\rho)}\mid xp_{\rho-1})$ возникло бы соотношение между $\rho$ ($\rho\leq n$) выражениями $(a^{(i)}\mid x)(a^{(1)}\cdots a^{(i-1)}a^{(i+1)}\cdots a^{(\rho)}q_{n-\rho+1})$\footnote{В оригинальной печати и в R823 в этом месте пропущено многоточие; формула (16.) однозначно дает $a^{(i+1)}\cdots a^{(\rho)}$. Здесь многоточие восстановлено редакционно.}, тогда как по (13.) в соотношение должны входить по меньшей мере $n+1$ таких выражений. Из аналогичных соображений получается система, эквивалентная (15.) и (14.):
\begin{equation}
\begin{aligned}
&\pair{u q_{\rho-1}}{a^{(1)}\cdots a^{(\rho)}}\\
&\quad=\sum_{i=1}^{\rho}(-1)^{i+1}\pair{u}{a^{(i)}}
(q_{\rho-1}\mid a^{(1)}\cdots a^{(i-1)}a^{(i+1)}\cdots a^{(\rho)})\\
&\quad=\sum_{i=1}^{\rho}(-1)^{i+1}\pair{u}{a^{(i)}}
(p_{n-\rho+1}a^{(1)}\cdots a^{(i-1)}a^{(i+1)}\cdots a^{(\rho)}).
\end{aligned}
\tag{17}
\end{equation}
Чтобы перейти от (16.) к тождествам между произвольными строками $A^{\rho_1},B^{\rho_2},\ldots,x,p$, заметим, что эти тождества должны становиться тождественными с (16.), как только мы разлагаем $A^{\rho_1}=a^{(1)}a^{(2)}\cdots a^{(\rho_1)}$, $B^{\rho_2}=b^{(1)}b^{(2)}\cdots b^{(\rho_2)}$ и т. д.; и что их также можно объединить в заключительные выражения только при помощи операций, определенных (16.), как только отдельные выражения имеют тип входящих в (16.). Поэтому, если еще положить
\begin{equation}
 \rho=\rho_1+\rho_2+\cdots+\rho_k\le n,
\tag{18}
\end{equation}
то получаем:
\[
\begin{aligned}
&\pair{a^{(1)}\cdots a^{(\rho_1)}b^{(1)}\cdots b^{(\rho_2)}
\cdots k^{(1)}\cdots k^{(\rho_k)}}{xp_{\rho-1}}\\
&\quad=\pair{A^{\rho_1}B^{\rho_2}\cdots K^{\rho_k}}{xp_{\rho-1}} \\
&\quad=\sum_{i=1}^{\rho_1}(-1)^{i+1}\pair{a^{(i)}}{x}
(a^{(1)}\cdots a^{(i-1)}a^{(i+1)}\cdots a^{(\rho_1)}
B^{\rho_2}\cdots K^{\rho_k}q_{n-\rho+1})\\
&\qquad+(-1)^{\rho_1\rho_2}\sum_{i=1}^{\rho_2}(-1)^{i+1}\pair{b^{(i)}}{x}
(b^{(1)}\cdots b^{(i-1)}b^{(i+1)}\cdots b^{(\rho_2)}
A^{\rho_1}\cdots K^{\rho_k}q_{n-\rho+1})\\
&\qquad+\cdots\\
&\qquad+(-1)^{(\rho_1+\cdots+\rho_{k-1})\rho_k}
\sum_{i=1}^{\rho_k}(-1)^{i+1}\pair{k^{(i)}}{x}\\
&\qquad\quad\times
(k^{(1)}\cdots k^{(i-1)}k^{(i+1)}\cdots k^{(\rho_k)}
A^{\rho_1}\cdots (K-1)^{\rho_{k-1}}q_{n-\rho+1}).
\end{aligned}
\]
Каждая отдельная сумма (но еще не часть суммы) действительно содержит строки $A^{\rho_1},B^{\rho_2}\ldots$ как неразложимую величину, поскольку она удовлетворяет необходимым и достаточным условиям:\footnote{Capelli, a. a. O. \S\ 23, дает достаточные условия: $\bigl(a^{(k)}\mid \frac{\partial}{\partial a^{(i)}}\bigr)=0$, $(k>i,\ i=1,2,\ldots,\rho_1-1)$. Из них необходимые и достаточные условия проще всего следуют применением процесса скобок Пуассона по Wellstein, \emph{Von den Differentialgleichungen der projektiven Invarianten}, \emph{Math. Ann.} 67 (1909), \S\ 3.}
\begin{equation}
 \left(a^{(i+1)}\middle|\frac{\partial}{\partial a^{(i)}}\right)=0;\ \ldots\ ;
 \left(k^{(j+1)}\middle|\frac{\partial}{\partial k^{(j)}}\right)=0.
 \quad (i=1,2,\ldots,\rho_1-1;\ j=1,2,\ldots,\rho_k-1)
\tag{19}
\end{equation}
Покажем, что она содержит все строки также явно. Действительно, если положить $B^{\rho_2}\cdots K^{\rho_k}q_{n-\rho+1}=q_{n-\rho_1+1}$, то по (16.) и (10.) получаем:
\[
\begin{aligned}
&\sum(-1)^{i+1}\pair{a^{(i)}}{x}
(a^{(1)}\cdots a^{(i-1)}a^{(i+1)}\cdots a^{(\rho_1)}q_{n-\rho+1})\\
&\quad=\pair{A^{\rho_1}}{xp_{\rho_1-1}}\\
&\quad=(A_{n-\rho_1}xp_{\rho_1-1})
=(-1)^{(\rho_1+1)(n+1)}\pair{q_{n-\rho_1+1}}{A_{n-\rho_1}x}\\
&\quad=(-1)^{(\rho_1+1)(n+1)}
\pair{B^{\rho_2}\cdots K^{\rho_k}q_{n-\rho+1}}{A_{n-\rho_1}x}.
\end{aligned}
\]
Если аналогично вычислить остальные суммы и после этого, поскольку строки уже охарактеризованы как различные весовыми индексами, положить $B^{\rho_2}=A^{\rho_2}$, $C^{\rho_3}=A^{\rho_3},\ldots$, то получаем искомые тождества:
\begin{equation}
\begin{aligned}
&\pair{A^{\rho_1}A^{\rho_2}\cdots A^{\rho_k}}{xp_{\rho-1}}\\
&\quad=\sum_{i=1}^{k}(-1)^{\delta_i}
\pair{A^{\rho_1}\cdots A^{\rho_{i-1}}A^{\rho_{i+1}}\cdots A^{\rho_k}q_{n-\rho+1}}{A_{n-\rho_i}x},\\
\delta_i&=(\rho_1+\rho_2+\cdots+\rho_{i-1})\rho_i+(\rho_i+1)(n+1).
\end{aligned}
\tag{20}
\end{equation}
А из (17.) получаются дуалистические тождества:
\begin{equation}
\begin{aligned}
&\pair{u q_{\sigma-1}}{A_{\sigma_1}A_{\sigma_2}\cdots A_{\sigma_k}}\\
&\quad=\sum_{i=1}^{k}(-1)^{\eps_i}
\pair{u A^{n-\sigma_i}}{p_{n-\sigma+1}A_{\sigma_1}\cdots A_{\sigma_{i-1}}A_{\sigma_{i+1}}\cdots A_{\sigma_k}},\\
\eps_i&=(\sigma_1+\sigma_2+\cdots+\sigma_{i-1})\sigma_i+(\sigma_i+1)(n+\sigma).
\end{aligned}
\tag{21}
\end{equation}
\emph{Теорема II.} Тождествами (20.) и (21.) исчерпывается вся совокупность неразложимых тождеств между строками $A^{\rho_i},x,p$, соответственно $A_{\sigma_i},u,q$, и вместе с тем вся совокупность тех неразложимых тождеств, которые допускают явное представление без выполнения подстановки (11.).

Первая часть теоремы следует из соображений, ведущих к (20.) и (21.); вторая часть -- из невозможности тождества между двумя строками символов:
\[
\begin{aligned}
\pair{q_\rho A^\sigma}{B_{\rho+\sigma-\tau}p_\tau}
&=\eps_1\pair{q_\rho B^{n-\rho-\sigma+\tau}}{A_{n-\sigma}p_\tau}\\
&\quad+\eps_2\pair{q_\rho q_{n-\tau}}{B_{\rho+\sigma-\tau}A_{n-\sigma}},
\end{aligned}
\]
где $\rho\le\tau\le\sigma$ и ни одна из строк не имеет веса $1\cdot(n-1)$. (Другие предположения о $\rho,\tau,\sigma$ можно свести к этому перестановкой обозначений строк.) Эта невозможность следует, например, из развертывания (8.) при специализации $q_\rho=lM^{\rho-1}$, $q_{n-\tau}=lN^{n-\tau-1}$, если положить $l_1=1$, $M^{2,3,\ldots,\rho}=1$, $A^{\rho+1,\ldots,\rho+\sigma}=1$, все остальные элементы строк $l,M,A$ равными нулю, тогда как $N$ и $B$ произвольны. Поскольку тождества между произвольными строками, после разложения строки $p_\tau$ на $y^{(1)}y^{(2)}\cdots y^{(\tau)}$, переходят в тождества, составленные из (20.), вторая часть утверждения будет доказана, как только (20.) можно будет получить из тождеств только с двумя строками символов. Действительно, из
\begin{equation}
\begin{aligned}
\pair{A^{\rho_1}A^{\rho_2}}{xp_{\rho-1}}
&=\eps\pair{A^{\rho_1}q_{n-\rho+1}}{A_{n-\rho_2}x}
+\pair{A^{\rho_1}}{xp_{\rho_1-1}},\\
&\quad \text{где }p_{\rho_1-1}\sim A^{\rho_2}q_{n-\rho+1},
\end{aligned}
\tag{20a}
\end{equation}
посредством специализации $A^{\rho_1}=B^{\sigma_1}B^{\sigma_2}$, то есть посредством
\[
\begin{aligned}
\pair{B^{\sigma_1}B^{\sigma_2}}{xp_{\rho_1-1}}
&=\eps\pair{B^{\sigma_1}q_{n-\rho_1+1}}{B_{n-\sigma_2}x}
+\pair{B^{\sigma_1}}{xp_{\sigma_1-1}},\\
&\quad \text{где } p_{\sigma_1-1}\sim B^{\sigma_2}q_{n-\rho_1+1},
\end{aligned}
\]
получаем тождество (20.) в трех строках символов; а дальнейшими специализациями $B^{\sigma_1}=C^{\tau_1}C^{\tau_2}$ и т. д. -- общие тождества (20.). Следовательно, несуществование явного представления тождества между двумя строками символов и двумя произвольными строками переменных влечет несуществование и при произвольно многих строках символов, если среди строк переменных нет строки $x$ или $u$.
% END UNIT BODY
% END PRODUCER UNIT 032
\endgroup


\clearpage
\begingroup
\setcounter{footnote}{0}
% BEGIN PRODUCER UNIT 033
% Source: C:/Users/Floris/Documents/interlanguage/03_projects/noether/02_slavic_working_corpus/translations/paper04/russian/v001/Noether_Paper04_Section04_Russian_v001.tex
% Identity: 10309 B / SHA-256 5B75154F9A29C030750A8F640EE7C45EE3C8381DC7855BB4DDC4E17F5CD21E28
\setlength{\parindent}{1.25em}
\setlength{\parskip}{0pt}
\emergencystretch=8em
\allowdisplaybreaks
\newcommand{\pair}[2]{\left(#1\middle|#2\right)}
\newcommand{\eps}{\varepsilon}
% BEGIN UNIT BODY
\section*{\S\ 4. Тождества разложения.}

Теперь рассмотрим самый общий случай тождеств между произвольными строками символов и строками переменных, в заключительных выражениях которых без выполнения подстановки (11.) строка $p_\tau$ появляется разложенной на отдельные строки $y^{(1)}\cdots y^{(\tau)}$, и будем называть такие тождества тождествами разложения. Опираясь на замечание в конце предыдущего параграфа, сначала определим их для случая только двух строк символов, то есть развернем выражение $\pair{q_\rho A^\sigma}{B_{\rho+\sigma-\tau}p_\tau}$. При этом положим:
\begin{equation}
\begin{aligned}
q_{\rho-\alpha}^{(i_1\ldots i_\alpha)}
&\sim p_{n-\rho}\,y^{(i_1)}\cdots y^{(i_\alpha)},\\
p_{\tau-\alpha}^{(i_1\ldots i_\alpha)}
&=y^{(i_{\alpha+1})}\cdots y^{(i_\tau)},\\
p_\tau&=y^{(1)}\cdots y^{(\tau)}.
\end{aligned}
\tag{22}
\end{equation}
Относительно чисел $\rho,\sigma,\tau$ надо различать четыре случая:
\[
\begin{array}{llll}
1)&\rho+\sigma\le n,& \tau\le\rho,& \tau\le\sigma;\\[2pt]
2)&\rho+\sigma\le n,& \tau\ge\rho,& \tau\le\sigma;\\[2pt]
3)&\rho+\sigma\le n,& \tau\le\rho,& \tau>\sigma;\\[2pt]
4)&\rho+\sigma\le n,& \tau\ge\rho,& \tau\ge\sigma.
\end{array}
\]
Выделим первый из них, а затем кратко охарактеризуем остальные. По (20.), (10.) и (9.), если разложить строку $y^{(1)}y^{(2)}\cdots y^{(\tau)}B_{\rho+\sigma-\tau}$ на $y^{(1)}$ и $y^{(2)}\cdots y^{(\tau)}B_{\rho+\sigma-\tau}$, получаем:
\begin{equation}
\begin{aligned}
&\pair{q_\rho A^\sigma}{y^{(1)}y^{(2)}\cdots y^{(\tau)}
B_{\rho+\sigma-\tau}}\\
&\quad=\pair{q_{\rho-1}^{(1)}A^\sigma}
{y^{(2)}\cdots y^{(\tau)}B_{\rho+\sigma-\tau}}\\
&\qquad+(-1)^\rho
\pair{q_\rho[A_{n-\sigma}y^{(1)}]^{\sigma-1}}
{y^{(2)}\cdots y^{(\tau)}B_{\rho+\sigma-\tau}}.
\end{aligned}
\tag{23}
\end{equation}
Последний член в (23.) имеет точно форму исходного выражения -- только $\sigma$ заменено на $\sigma-1$, а $\tau$ на $\tau-1$, -- значит, он снова допускает равенство (23.). Поэтому, так как $\tau\le\sigma$, мы можем применить операцию (23.) $\tau$ раз и с учетом (10.) приходим к соотношению:
\begin{equation}
\begin{aligned}
&\pair{q_\rho A^\sigma}{p_\tau B_{\rho+\sigma-\tau}}\\
&\quad=(-1)^{\rho\sigma+(\sigma+\tau)(n+1)}
\pair{q_\rho B^{n-\rho-\sigma+\tau}}{A_{n-\sigma}p_\tau}\\
&\qquad+\sum_{i_1=1}^{\tau}(-1)^{\rho(i_1-1)}
\pair{q_{\rho-1}^{(i_1)}
[A_{n-\sigma}y^{(1)}\cdots y^{(i_1-1)}]^{\sigma-i_1+1}}
{y^{(i_1+1)}\cdots y^{(\tau)}B_{\rho+\sigma-\tau}}.
\end{aligned}
\tag{24}
\end{equation}
Каждый член под знаком суммы снова имеет тип исходного выражения: в нем $\rho$ заменено на $\rho-1$, $\sigma$ на $\sigma-i_1+1$, $\tau$ на $\tau-i_1$, так что условие 1) выполнено еще сильнее. Поэтому, так как $\tau\le\rho$, мы можем применить операцию (24.) $\tau$ раз и получаем нужное тождество разложения:
\begin{equation}
\begin{aligned}
&\pair{q_\rho A^\sigma}{p_\tau B_{\rho+\sigma-\tau}}\\
&\quad=\eps_{i_0}\pair{q_\rho B^{n-\rho-\sigma+\tau}}
{A_{n-\sigma}p_\tau}\\
&\qquad+\sum_{i_1=1}^{\tau}\eps_{i_1}
\pair{q_{\rho-1}^{(i_1)}B^{n-\rho-\sigma+\tau}}
{A_{n-\sigma}p_{\tau-1}^{(i_1)}}\\
&\qquad+\sum_{i_2=i_1+1}^{\tau}\eps_{i_2}
\pair{q_{\rho-2}^{(i_1i_2)}B^{n-\rho-\sigma+\tau}}
{A_{n-\sigma}p_{\tau-2}^{(i_1i_2)}}+\cdots\\
&\qquad+\eps_{i_\tau}
\pair{q_{\rho-\tau}^{(i_1\ldots i_\tau)}B^{n-\rho-\sigma+\tau}}
{A_{n-\sigma}}\\
&\quad=\sum_{\alpha=0}^{\tau}
\sum_{i_\alpha=i_{\alpha-1}+1}^{\tau}\eps_{i_\alpha}
\pair{q_{\rho-\alpha}^{(i_1\ldots i_\alpha)}B^{n-\rho-\sigma+\tau}}
{A_{n-\sigma}p_{\tau-\alpha}^{(i_1\ldots i_\alpha)}}.
\end{aligned}
\tag{25}
\end{equation}
Для модуля $\eps$ из (24.) получается:
\begin{equation}
\begin{aligned}
\eps_{i_\alpha}&=(-1)^{\delta_{i_\alpha}},\\
\delta_{i_\alpha}
&=\sum_{\beta=1}^{\alpha}(\rho-\beta+1)(i_{\beta-1}+i_\beta+1)\\
&\quad+(\rho-\alpha)(\sigma+i_\alpha+\alpha)
+(\sigma+\tau+\alpha)(n+1),
\end{aligned}
\tag{26}
\end{equation}
где надо положить $i_0=0$. Знак суммы в (25.) следует понимать так, что к каждой комбинации $i_1i_2\ldots i_{\alpha-1}$, где $i_1<i_2<\cdots<i_{\alpha-1}$, присоединяются все значения $i_\alpha$, для которых $i_{\alpha-1}<i_\alpha\le\tau$. Отдельные суммы в (25.), как легко убедиться с учетом (26.), удовлетворяют дифференциальным уравнениям, аналогичным (19.):
\begin{equation}
 \left(\frac{\partial}{\partial y^{(i)}}\middle|y^{(i+1)}\right)=0,
 \qquad (i=1,2,\ldots,\tau-1),
\tag{27}
\end{equation}
и потому зависят только от строки $p_\tau$. Значит, тождество разложения является неразложимым тождеством между строками $A,B,q,p$; к явному представлению мы вернемся в следующем параграфе.

В случае 2) операцию (23.) также можно применить $\tau$ раз, но операция (24.) обрывается после $\rho$-го раза. Поэтому в заключительном выражении (25.) первый знак суммы надо заменить на $\sum_{\alpha=0}^{\rho}$. В случаях 3) и 4) операцию (23.) можно выполнить только $\sigma$ раз; вместо (24.) тогда возникает соотношение:
\begin{equation}
\begin{aligned}
&\pair{q_\rho A^\sigma}{p_\tau B_{\rho+\sigma-\tau}}\\
&\quad=(-1)^{\rho\sigma}
\pair{q_{\rho-(\tau-\sigma)}^{(\sigma+1\ldots\tau)}B^{n-\rho+\tau-\sigma}}
{A_{n-\sigma}p_{\tau-(\tau-\sigma)}^{(\sigma+1\ldots\tau)}}\\
&\qquad+\sum_{i_1=1}^{\sigma}(-1)^{\rho(i_1-1)}
\pair{q_{\rho-1}^{(i_1)}
[A_{n-\sigma}y^{(1)}\cdots y^{(i_1-1)}]^{\sigma-i_1+1}}
{y^{(i_1+1)}\cdots y^{(\tau)}B_{\rho+\sigma-\tau}};
\end{aligned}
\tag{28}
\end{equation}
$(\tau-\sigma)$-кратное повторение (28.), при котором знак суммы переходит в $\sum_{i_\beta=i_{\beta-1}+1}^{\sigma+\beta-1}$, -- как показывает $(\sigma-i_1+1)$-кратное применение (23.) к членам под знаком суммы в (28.), -- дает все члены отдельной суммы в (25.), соответствующей индексу $\alpha=\tau-\sigma$, с соответствующим знаком. После этого случаи 3) и 4) переходят в случаи 1) и 2), потому что числа, соответствующие $\sigma,\tau$, становятся такими: $\sigma'=\sigma-i_{\tau-\sigma}+\tau-\sigma$, $\tau'=\tau-i_{\tau-\sigma}=\sigma'$; а при продолжении процедуры будет $\tau'<\sigma'$. Поэтому в заключительном выражении (25.) первый знак суммы надо заменить на $\sum_{\alpha=\tau-\sigma}^{\tau,\rho}$.

Аналогично из (21.) получается дуалистическое тождество разложения, в котором строка $q_\rho$ разлагается на отдельные строки $v^{(1)},\ldots,v^{(\rho)}$. Положим
\begin{equation}
\begin{aligned}
 \dot p_{\tau-\beta}^{(i_1\ldots i_\beta)}
&\sim v^{(i_1)}\cdots v^{(i_\beta)}q_{n-\tau};\\
 \dot q_{\rho-\beta}^{(i_1\ldots i_\beta)}
&=v^{(i_{\beta+1})}\cdots v^{(i_\rho)};\\
 q_\rho&=v^{(1)}\cdots v^{(\rho)}.
\end{aligned}
\tag{29}
\end{equation}
Тогда имеем:
\begin{equation}
\begin{aligned}
&\pair{q_\rho A^\sigma}{p_\tau B_{\rho+\sigma-\tau}}\\
&\quad=\sum_{\beta=\max(0,\tau-\sigma)}^{\min(\tau,\rho)}
\sum_{i_\beta=i_{\beta-1}+1}^{\rho}\eps\\
&\qquad\quad\pair{\dot q_{\rho-\beta}^{(i_1\ldots i_\beta)}
B^{n-\rho-\sigma+\tau}}
{A_{n-\sigma}\dot p_{\tau-\beta}^{(i_1\ldots i_\beta)}} ,
\end{aligned}
\tag{30}
\end{equation}
где индекс суммирования $\beta$ пробегает от большего из нижних значений до меньшего из верхних. Все члены отдельной суммы в (25.) и (30.) являются полярами одной формы, возникшей через полярные процессы
\[
\left(\frac{\partial}{\partial p_{\tau-\alpha}}\middle|
p_{\tau-\alpha}^{(i_1\ldots i_\alpha)}\right),\qquad
\left(q_{\rho-\alpha}^{(i_1\ldots i_\alpha)}\middle|
\frac{\partial}{\partial q_{\rho-\alpha}}\right).
\]
Чтобы это выразить, обозначим через $\Pi$ агрегат таких полярных операций, умноженный на постоянные, и тогда вместо (25.) и (30.) получаем соотношение:
\begin{equation}
\begin{aligned}
&\pair{q_\rho A^\sigma}{p_\tau B_{\rho+\sigma-\tau}}\\
&\quad=\sum_{\alpha=\max(0,\tau-\sigma)}^{\min(\tau,\rho)}
\Pi\pair{q_{\rho-\alpha}B^{n-\rho-\sigma+\tau}}
{A_{n-\sigma}p_{\tau-\alpha}}.
\end{aligned}
\tag{31}
\end{equation}
% END UNIT BODY
% END PRODUCER UNIT 033
\endgroup


\clearpage
\begingroup
\setcounter{footnote}{0}
% BEGIN PRODUCER UNIT 034
% Source: C:/Users/Floris/Documents/interlanguage/03_projects/noether/02_slavic_working_corpus/translations/paper04/russian/v001/Noether_Paper04_Section05_Russian_v001.tex
% Identity: 12293 B / SHA-256 C6CB5A58949335A40815AA3B5DEEF935600F8CCE153E1CB28C5BECC95BFF5774
\setlength{\parindent}{1.25em}
\setlength{\parskip}{0pt}
\emergencystretch=8em
\allowdisplaybreaks
\newcommand{\pair}[2]{\left(#1\middle|#2\right)}
\newcommand{\eps}{\varepsilon}
% BEGIN UNIT BODY
\section*{\S\ 5. Тождества разложения в явной форме.}

Если в случае 4) в (25.) специализировать $\tau$ до $\sigma+\rho$, то получается:
\begin{equation}
\begin{aligned}
 \pair{q_\rho A^\sigma}{y^{(1)}\cdots y^{(\rho+\sigma)}}
 &=\sum_{1\le i_1<\cdots<i_\rho\le \rho+\sigma}
 \eps_{i_\rho}\pair{q_\rho}{y^{(i_1)}\cdots y^{(i_\rho)}}\\
 &\qquad\cdot
 \pair{A^\sigma}{y^{(i_{\rho+1})}\cdots y^{(i_{\rho+\sigma})}},
\end{aligned}
\tag{32}
\end{equation}
где $\eps_{i_\rho}=(-1)^{\delta_{i_\rho}}$, $\delta_{i_\rho}=\rho(\rho+1)/2+i_1+\cdots+i_\rho$. Это более общая форма теоремы о произведении для матриц, чем (16.) и (17.); она получается из разложения Лапласа детерминанта продуктовой матрицы
\[
 \sum\pm(v^{(1)}y^{(1)})\cdots(v^{(\rho)}y^{(\rho)})
 (a^{(1)}y^{(\rho+1)})\cdots(a^{(\sigma)}y^{(\rho+\sigma)}).
\]
Следовательно, более общая теорема о произведении составлена из тождеств (20.), соответственно (21.), и этим объясняется выделение специальной формы (16.), (17.).

Соотношение (32.) теперь дает средство для явного представления тождеств разложения. Для этого рассмотрим формы, входящие в (31.), как основные формы; то есть выполним подстановку (11.), которая ничего не меняет в явном представлении через строки $A,B$, поскольку по (8.) вновь введенные строки $R$ выражаются через сами строки $A,B$, а не через их частичные строки $a^{(1)}a^{(2)}\ldots$, $b^{(1)}b^{(2)}\ldots$. Пусть, следовательно,
\begin{equation}
 \pair{q_{\rho-\alpha}B^{n-\rho-\sigma+\tau}}{A_{n-\sigma}p_{\tau-\alpha}}
 =\pair{R_{\rho-\alpha}p_{n-\rho+\alpha}}{R^{\tau-\alpha}p_{\tau-\alpha}},
\tag{33}
\end{equation}
тогда для отдельной суммы из (25.), соответствующей индексу $\alpha$, получаем:
\begin{equation}
\begin{aligned}
&\sum \eps_{i_\alpha}
\pair{R_{\rho-\alpha}p_{n-\rho}y^{(i_1)}\cdots y^{(i_\alpha)}}
{R^{\tau-\alpha}y^{(i_{\alpha+1})}\cdots y^{(i_\tau)}}\\
&\quad=\sum \eps_{i_\alpha}
\pair{[R_{\rho-\alpha}p_{n-\rho}]^\alpha
y^{(i_1)}\cdots y^{(i_\alpha)}}
{R^{\tau-\alpha}y^{(i_{\alpha+1})}\cdots y^{(i_\tau)}}\\
&\quad=\eps\pair{[R_{\rho-\alpha}p_{n-\rho}]^\alpha R^{\tau-\alpha}}{p_\tau}
 =\eps\pair{R^{\tau-\alpha}q_{n-\tau}}{R_{\rho-\alpha}p_{n-\rho}}.
\end{aligned}
\tag{34}
\end{equation}
Тем самым действительно задано явное заключительное выражение, а симметрическая форма, в которой входят $q_\rho$ и $p_\tau$, показывает, что (30.) приводит к тому же заключительному выражению, которое поэтому не зависит от исходного разложения. Различие между (25.) и (30.) существует лишь до тех пор, пока сами строки $y$ и $v$ имеют самостоятельное значение, то есть пока тождество разложения по нашему определению переходит в разложимое тождество. Чтобы перейти к тождествам с большим числом строк символов, нужно по \S\ 3 (конец) лишь разложить строку $q_\rho$ в $q_{\rho_1}C^{\sigma_1}$, соответственно $p_\tau$ в $p_{\tau_1}C_{\tau-\tau_1}$. Тогда возникают выражения:
\begin{equation}
 \pair{q_{\rho_1}C^{\sigma_1}}
 {[R^{\tau-\alpha}q_{n-\tau}]_\lambda R_{\rho_1+\sigma_1-\alpha-\lambda}},
 \qquad
 \pair{[R_{\rho-\alpha}p_{n-\rho}]^\lambda R^{\tau-\alpha}}
 {p_{\tau_1}C_{\tau-\tau_1}},
\tag{35}
\end{equation}
а они имеют точно тот же тип, что и в случае двух строк символов, и потому при подстановке, соответствующей (33.), снова допускают явное представление; повторением получается явное представление для произвольного числа строк. Стоит еще отметить, что первая и последняя суммы в (25.), соответственно (30.), дают явное представление без применения (33.), только по формуле (32.), или вообще сводятся к одному члену, явно содержащему все строки. Соотношение, возникающее из (32.) при разложении $q_\rho$ в $q_{\rho_1}C^{\sigma_1}$ и т. д., можно рассматривать как теорему о произведении в ее самой общей форме.

Подводя итог, покажем:

\emph{Теорема III:} Тождествами
\begin{equation}
 \pair{q_\rho A^\sigma}{p_\tau B_{\rho+\sigma-\tau}}
 =\sum_{\alpha=\max(0,\tau-\sigma)}^{\min(\tau,\rho)}
 \eps\pair{R^{\tau-\alpha}q_{n-\tau}}{R_{\rho-\alpha}p_{n-\rho}},
\tag{36}
\end{equation}
где строки $R$ определяются через (33.), и тождествами, возникающими из них при разложении $q_\rho,p_\tau$ в $q_{\rho_1}C^{\sigma_1}$ и т. д., исчерпывается вся совокупность неразложимых тождеств.

Действительно, возникающие тождества являются самыми общими в своем роде, поскольку отдельные строки уже не подчиняются никаким ограничениям по весу и числу, как это еще было в случае (20.), (21.). Но между этими строками не существует и никаких других тождеств: при разложении строки $p_\tau$ в $y^{(1)}\cdots y^{(\tau)}$ самые общие тождества составлены из (20.), следовательно, по \S\ 3 (конец), из (20a.); причем композиция, обсуждавшаяся после (16.), есть именно та, которая проведена в \S\ 4. Переход к произвольному числу строк, как показывает обсуждение (20a.), дает применение одной и той же операции к выражениям, возникающим при разложении $q_\rho$ в $q_{\rho_1}C^{\sigma_1}$ и т. д. Отсюда также следует, что прямой переход от (20.) вместо (20a.) не дает ничего нового; это легко подтвердить вычислением, если один раз специализировать строку $q_\rho$ в (25.), а другой раз выполнить переход от (20.) методом \S\ 4. В обоих случаях возникают одни и те же суммы, которые при применении общей теоремы о произведении (см. выше) переходят в тождества, вытекающие из (36.). Тождества (20.), (21.) соответствуют специальным случаям $\tau=1$, $\rho=1$; тогда появляются только по две отдельные суммы, а значит, согласно замечанию выше, явное представление без подстановки (33.), что согласуется с прежним.

В заключение кратко упомянем еще два специальных случая тождества разложения. Если в (31.) положить $\tau=\sigma$, $\rho=n-\sigma$ и обозначить строку $p_\tau$ как $p'_\sigma\sim q'_{n-\sigma}$, то получается:
\begin{equation}
\begin{aligned}
\pair{A^\sigma}{p_\sigma}\pair{B^\sigma}{p'_\sigma}
-\pair{A^\sigma}{p'_\sigma}\pair{B^\sigma}{p_\sigma}
=\sum_{\alpha=1}^{\min(\sigma,n-\sigma)}
\Pi\pair{q_{n-\sigma-\alpha}B^\sigma}{A_{n-\sigma}p_{\sigma-\alpha}}.
\end{aligned}
\tag{37}
\end{equation}
Это формула Клебша\footnote{Clebsch, a. a. O., S. 25--32. Клебш доказывает, что отдельные члены правых частей зависят только от строк $A,B$; явное представление получалось бы применением (32.) к его выражениям $\Phi_h$.} для разложения формы с двумя когредиентными строками $p_\sigma,p'_\sigma$ по элементарным ковариантам. Элементарные коварианты, принадлежащие форме $\pair{A^\sigma}{p_\sigma}^m\pair{B^\sigma}{p'_\sigma}^n$, тогда следующие:
\begin{equation}
 \left\{\sum_{\alpha=1}^{\min(\sigma,n-\sigma)}
 \eps\pair{q_{n-\sigma-\alpha}B^\sigma}{A_{n-\sigma}p_{\sigma-\alpha}}
 \right\}^{\rho}
 \pair{A^\sigma}{p_\sigma}^{m-\rho}\pair{B^\sigma}{p'_\sigma}^{n-\rho},
 \qquad (\rho=0,1,\ldots,m,n).
\tag{38}
\end{equation}
Они, как мы коротко будем говорить, возникают из основной формы через ``когредиентную свертку с дефектом $\rho$'' (ср. \S\ 2).

Во-вторых, если положить $\tau=\sigma-\lambda$, $\rho=n-\sigma-\lambda$; $B^\sigma=A^\sigma$, то получается:
\begin{equation}
\begin{aligned}
\pair{q_{n-\sigma-\lambda}A^\sigma}{A_{n-\sigma}p_{\sigma-\lambda}}
&=(-1)^\lambda
\pair{q_{n-\sigma-\lambda}A^\sigma}{A_{n-\sigma}p_{\sigma-\lambda}}\\
&\quad+(-1)^{(n-\sigma)(\sigma-\lambda)}
 \sum_{\alpha=1}^{\min(\sigma-\lambda,n-\sigma-\lambda)}
 \Pi\pair{q_{n-\sigma-\lambda-\alpha}A^\sigma}
 {A_{n-\sigma}p_{\sigma-\lambda}}.
\end{aligned}
\tag{39}
\end{equation}
Следовательно, как отмечено в \S\ 2, условия (12.), характеризующие строки символов, при нечетном дефекте $\lambda$ являются следствием условий более высокого дефекта. Для линейной формы $\pair{A^\sigma}{p_\sigma}=\pair{B^\sigma}{p_\sigma}$ из (39.) следует, что ее неприводимые инвариантные образования, квадратичные по коэффициентам, сводятся к образованиям четного дефекта; это согласуется с известным фактом, что линейная форма в бинарной и тернарной областях не имеет инвариантных образований.

Заметим, что общие тождества первоначально были выведены при предположении, что отдельные строки удовлетворяют условиям (12.). Однако, поскольку эти отдельные строки входят в тождества только линейно, последние справедливы и для более общих строк, не удовлетворяющих условиям (12.).
% END UNIT BODY
% END PRODUCER UNIT 034
\endgroup


\clearpage
\begingroup
\setcounter{footnote}{0}
% BEGIN PRODUCER UNIT 035
% Source: C:/Users/Floris/Documents/interlanguage/03_projects/noether/02_slavic_working_corpus/translations/paper04/russian/v001/Noether_Paper04_Section06_Russian_v001.tex
% Identity: 15360 B / SHA-256 4BCD007C28226970CE627897581284610490B7157AEA7F8097DCAEEF953A3B04
\setlength{\parindent}{1.25em}
\setlength{\parskip}{0pt}
\emergencystretch=8em
\allowdisplaybreaks
\newcommand{\pair}[2]{\left(#1\middle|#2\right)}
\newcommand{\eps}{\varepsilon}
% BEGIN UNIT BODY
\section*{\S\ 6. Явное представление инвариантных образований.\\Процессы свертки и дифференцирования.}

Результаты предыдущих параграфов позволяют завершить сформулированную в \S\ 2 теорему о символической представимости (первую фундаментальную теорему), добавляя явную представимость; а именно мы покажем:

\emph{Теорема IV:} ``Все инвариантные образования произвольных основных форм можно явно представить матричными произведениями; то есть в заключительные выражения, кроме строк переменных $p_\rho$, входят только строки $S^\sigma,\ldots,T^\tau$ исходных форм или строки $P,Q,R$, полученные из них подстановкой (9.) или (11.).''

Для доказательства предположим, что инвариантное образование, кроме строк переменных, может зависеть от строк $S^\sigma,\ldots,T^\tau$; причем эти строки достаточно далеко разложены на отдельные строки:
$S^\sigma=S^{\sigma_1}\cdots S^{\sigma_i},\ldots,T^\tau=T^{\tau_1}\cdots T^{\tau_k}$, чтобы сделать возможным представление детерминантами. Пусть строки упорядочены в произвольном самом по себе, но фиксированном порядке $S^\sigma,\ldots,T^\tau$. Выделим такой агрегат детерминантов, который зависит от первой полной строки $S^\sigma=S^{\sigma_1}\cdots S^{\sigma_i}$, а не только от отдельных ее частичных строк. По \S\ 3--\S\ 5 эта зависимость является следствием общих тождеств. Если теперь разложить строку переменных $p_\rho$ в строки $y,v$ или одну из последующих строк $T$ в отдельные строки $t$, соответственно $\tau$, то самые общие тождества составлены из (20.), (21.). Но эти последние тождества всегда приводят к выражению, явно содержащему строку $S^\sigma=S^{\sigma_1}\cdots S^{\sigma_i}$, то есть к левой стороне (20.), (21.); через (19.) действительно убеждаемся, что только полная сумма справа, а не ее часть, соответствующая члену (20a.), обладает свойством зависеть только от $S^\sigma$. При продолжении процедуры все строки, однажды вошедшие явно, остаются явно сохраненными; применение тождеств может потребовать лишь объединения нескольких строк в одну, то есть подстановки (9.). Если наконец остается привести к явной форме только одну строку $T^\tau$, возможны два случая. Может еще появляться свободная строка переменных, то есть не связанная с другими строками подстановкой (9.); ее разложением на частичные строки $y$ или $v$ процедуру можно завершить, и возникают, как в (31.), поляры одной или нескольких форм, явно содержащих все строки. Если же свободной строки переменных больше нет, то по способу их возникновения мы распознаем в совокупности выражений, которые содержат отдельные строки $t$ или $\tau$, но зависят от строки $T^\tau$, члены тождества разложения; а они по \S\ 5 при применении подстановки (11.) всегда допускают явное представление во всех строках. Тем самым указанная выше теорема доказана в общем случае. Стоит еще отметить, что если появляются только две строки символов, подстановка (11.) никогда не возникает: действительно, ввиду $\sigma,n-\sigma,\tau,n-\tau<n$ строки переменных могут не входить только тогда, когда обе строки символов объединены в одном детерминанте, то есть появляются явно; как только входят строки переменных, (34.) и (33.) показывают, что подстановка (11.) излишня.

Из сказанного следует, что для порождения всех инвариантных образований достаточно, добавляя строки переменных, объединять строки символов во все возможные матричные произведения. Самое общее матричное произведение имеет форму:
\begin{equation}
 \pair{P^{\mu_1}\cdots P^{\mu_i}}{Q_{\nu_1}\cdots Q_{\nu_k}},
 \qquad \sum\mu=\sum\nu,
\tag{40}
\end{equation}
где $P$ и $Q$ могут быть строками исходных форм, или возникать из исходных строк через последовательность подстановок (9.) или (11.), или, наконец, быть строками переменных. Но выражения (40.) можно породить последовательным объединением попарно строк символов в матричное произведение с помощью строк переменных; ведь из
\[
 \pair{q_{n-\mu-\lambda}P^\mu}{Q_{\nu_1}p_{n-\mu-\nu_1-\lambda}}
 =\pair{R_\lambda Q_{\nu_1}}{p_{n-\mu-\nu_1-\lambda}}
\]
при объединении со строкой $Q_{\nu_2}$ получаем выражение
\[
 \pair{[R_\lambda Q_{\nu_1}]^{n-\lambda-\nu_1}}
 {Q_{\nu_2}p_{n-\mu-\nu_1-\nu_2-\lambda}}
 =\pair{q_{n-\mu-\lambda}P^\mu}
 {Q_{\nu_1}Q_{\nu_2}p_{n-\mu-\nu_1-\nu_2-\lambda}},
\]
и, продолжая процедуру, получаем выражение (40.). Если при этом $P$ и $Q$ не являются строками исходных форм, то (9.) и (11.) вместе с только что доказанным показывают, что они также возникли через последовательность объединений попарно строк символов. Отсюда следует:

\emph{Теорема V:} Процесс порождения всех инвариантных образований, процесс свертки, состоит в последовательном объединении попарно строк символов в матричное произведение с помощью строк переменных.

Из двух строк символов $S^\sigma,T^\tau$, где $\sigma\ge\tau$, можно образовать следующие матричные произведения:
\begin{align}
&\pair{q_{n-\sigma-\lambda}S^\sigma}{T_{n-\tau}p_{\tau-\lambda}},
&&(\lambda=1,2,\ldots,\min(\tau,n-\sigma)),\tag{41}\\
&\pair{q_{n-\tau-\mu}T^\tau}{S_{n-\sigma}p_{\sigma-\mu}},
&&(\mu=\sigma-\tau+\lambda),\tag{42}\\
&\pair{q_{n-\tau-\nu}T^\tau}{S_{n-\sigma}p_{\sigma-\nu}},
&&(\nu=1,2,\ldots,\sigma-\tau).
\tag{43}
\end{align}
Но из (31.) для (42.) и (43.) получаются значения:
\begin{align}
&\pair{q_{n-\sigma-\lambda}S^\sigma}{T_{n-\tau}p_{\tau-\lambda}}
 +\sum_{\alpha}\Pi
 \pair{q_{n-\sigma-\lambda-\alpha}S^\sigma}
 {T_{n-\tau}p_{\tau-\lambda-\alpha}},\tag{42a}\\
&\Pi(q_{n-\sigma}S^\sigma)(T_{n-\tau}p_\tau)
 +\sum_{\beta}\Pi
 \pair{q_{n-\sigma-\beta}S^\sigma}{T_{n-\tau}p_{\tau-\beta}}.
\tag{43a}
\end{align}
Следовательно, формы (42.) линейно выражаются через (41.) и поляры от (41.), а формы (43.) -- через поляры основной формы и форм (41.). Так же из (31.) следует линейная зависимость форм (41.) от (42.) и поляров от (42.). Поэтому по определению Клебша\footnote{Clebsch, a. a. O. \S\ 7.} формы (42.) эквивалентны формам (41.), тогда как (43.) сводит обратно к основной форме. Следовательно, порождающий процесс равноправно задается (41.) или (42.); процесс (41.), более простой по числу $\lambda$, мы будем определять как процесс свертки. Число $\lambda$, дефект свертки (ср. \S\ 2), дает число строк, недостающих до детерминантного произведения, причем именно в том матричном произведении, которое при заданной свертке содержит максимальное число строк. Поскольку $\lambda$ пробегает только до меньшего из двух значений $\tau,n-\sigma$, где $\sigma\ge\tau$, то:
\begin{equation}
 \lambda\le\frac n2,\ n\text{ четное};\qquad
 \lambda\le\frac{n-1}{2},\ n\text{ нечетное}.
\tag{44}
\end{equation}
Сведение (42.) и (43.) к (41.) остается справедливым и тогда, когда при дальнейшей свертке строки $p,q$ перешли в строки символов; ведь по (36.) получается:
\[
 \pair{A^{n-\tau-\varkappa}T^\tau}{S_{n-\sigma}B_{\sigma-\varkappa}}
 =\sum_{\alpha=\max(0,\varkappa-\sigma+\tau)}^{\min(\tau,n-\sigma)}
 \eps\pair{R^{\tau-\alpha}B^{n-\sigma+\varkappa}}
 {A_{\tau+\varkappa}R_{n-\sigma-\alpha}},
\]
где строки $R$ возникли из $S$ и $T$ по (33.). Значит, эти формы, содержащие только строки символов, получают сверткой $\pair{q_\tau A^{n-\tau-\varkappa}}{B_{\sigma-\varkappa}p_{n-\sigma}}$ соответственно $\pair{q_{\tau-\alpha}B^{n-\sigma+\varkappa}}{A_{\tau+\varkappa}p_{n-\sigma-\alpha}}$ -- в зависимости от того, $\sigma-\tau\gtreqless 2\varkappa$ -- с основной формой и формами (41.).

Из символического процесса свертки известным способом получают инвариантные процессы дифференцирования: нужно лишь заменить строки символов $S^\sigma,T_{n-\tau}$ строками дифференциальных символов $\frac{\partial}{\partial p_\sigma}$, $\frac{\partial}{\partial q_{n-\tau}}$, причем, как известно, каждому произведению $\frac{\partial}{\partial p_{i_1\ldots i_\sigma}}\frac{\partial}{\partial q_{j_1\ldots j_{n-\tau}}}$ соответствует реальное значение $\frac{\partial^2}{\partial p_{i_1\ldots i_\sigma}\partial q_{j_1\ldots j_{n-\tau}}}$. Поскольку строки $\frac{\partial}{\partial p_\sigma}$, $\frac{\partial}{\partial q_{n-\tau}}$ когредиентны строкам $S^\sigma,T_{n-\tau}$, они удовлетворяют тем же тождествам, что и эти строки, в частности также тем, которые существуют между (42.) и (42a.), (43.) и (43a.). Подводя итог, получаем:

\emph{Теорема VI:} Все инвариантные образования произвольных основных форм возникают из них повторным применением символического процесса свертки (41.) или также инвариантных процессов дифференцирования:
\begin{equation}
 \left(q_{n-\sigma-\lambda}\frac{\partial}{\partial p_\sigma}
 \middle|\frac{\partial}{\partial q_{n-\tau}}p_{\tau-\lambda}\right),
 \qquad(\sigma\ge\tau,\ \lambda=1,2,\ldots,\min(\tau,n-\sigma)).
\tag{45}
\end{equation}
Остается заметить, что при каждой свертке (41.) общий вес формы, то есть сумма весов отдельных строк переменных (ср. \S\ 1), уменьшается на положительное число $2(\lambda^2+\lambda(\sigma-\tau))$. Отсюда, независимо от общего доказательства конечности, следует конечность ряда форм, то есть совокупности инвариантных образований данной основной формы, линейных по коэффициентам.\footnote{Ср. Clebsch, a. a. O. \S\ 11 и 12. Конечность системы элементарных ковариантов.}

Далее заметим, что теорема VI остается справедливой также для форм, не удовлетворяющих условиям (12.). Действительно, каждый множитель из (3.), $\pair{S^\rho}{p_\rho}^{m_\rho}$, можно заменить произведением $m_\rho$ линейных множителей $\pair{A^\rho}{p_\rho}\pair{B^\rho}{p_\rho}\cdots\pair{M^\rho}{p_\rho}$; при этом инвариантные образования переходят в образования системы линейных форм, которые лишь впоследствии нужно приравнять друг к другу. Но для каждой отдельной линейной формы условия (12.), которые при введении дифференциальных символов содержат только частные производные 2-го порядка, всегда выполнены.
% END UNIT BODY
% END PRODUCER UNIT 035
\endgroup


\clearpage
\begingroup
\setcounter{footnote}{0}
% BEGIN PRODUCER UNIT 036
% Source: C:/Users/Floris/Documents/interlanguage/03_projects/noether/02_slavic_working_corpus/translations/paper04/russian/v001/Noether_Paper04_Section07_Russian_v001.tex
% Identity: 17886 B / SHA-256 8B2B7E08B97625D4A6647DF51DFEAC597B058B752EC4D1E135C506567223FC1C
\setlength{\parindent}{1.25em}
\setlength{\parskip}{0pt}
\emergencystretch=10em
\allowdisplaybreaks
\newcommand{\pair}[2]{\left(#1\middle|#2\right)}
\newcommand{\eps}{\varepsilon}
% BEGIN UNIT BODY
\section*{\S\ 7. Разложение общих форм по нормальным формам.}

В дальнейшем, в \S\ 8, мы выведем главное свойство процесса свертки (41.), а именно его составимость из основных сверток. Для этого нужно перенести на $n$-арную область важную теорему, которую Mertens\footnote{F. Mertens, \emph{\"Uber invariante Gebilde quatern\"arer Formen} (см. Введение), далее цитируется как ``Mertens''.} дал в кватернарной области. Эта теорема утверждает, что произвольную конечную систему форм можно заменить эквивалентной системой нормальных форм. Под ``нормальной формой'' мы, обобщая бинарное и тернарное определение\footnote{Study, указ. соч., с. 54.}, понимаем форму, все инвариантные образования которой, линейные по коэффициентам, тождественно исчезают, за исключением самой основной формы. Поэтому по Теореме VI (\S\ 6) нормальная форма $\varphi$ удовлетворяет системе дифференциальных уравнений:
\begin{equation}
 \left(q_{n-\sigma-\lambda}\frac{\partial}{\partial p_\sigma}\middle|\frac{\partial}{\partial q_{n-\tau}}p_{\tau-\lambda}\right)\varphi=0,
 \qquad(\sigma\ge\tau,\ \lambda=1,2,\ldots,\min(\tau,n-\sigma)),
\tag{46}
\end{equation}
где строки $q_{n-\sigma-\lambda},p_{\tau-\lambda}$ следует рассматривать как произвольные величины. В кватернарной области, с учетом (39.) при $\sigma=\tau=2$, эти дифференциальные уравнения полностью совпадают с десятью дифференциальными уравнениями, которые Mertens дал без введения произвольных строк $p,q$ (указ. соч., формула 28)). Знание этих уравнений вместе с Теоремой VI позволяет почти дословно перенести его способ доказательства на $n$ переменных. Добавить нужно только проверку, получаемую с помощью тождества разложения (30.), что построенные формы являются нормальными формами; в кватернарной области эта проверка еще получается без дальнейшего преобразования. Поэтому достаточно кратко изложить доказательство, причем для простоты в единственном случае, который понадобится далее, а именно для одной основной формы $f$ и ее ряда форм (ср. конец \S\ 6), поскольку для инвариантных образований произвольной степени по коэффициентам произвольного числа основных форм доказательство остается тем же.

Основная идея доказательства такова: введением преобразованных коэффициентов заменить форму с $n-1$ строками $p_\rho$ формами с $n$ строками $\xi$, когредиентными с $x$, то есть строками величин преобразования. Из этих форм искомые нормальные формы непосредственно выводятся инвариантными процессами дифференцирования. Разложение общих форм по нормальным формам опосредуется рядовым разложением для форм с $\rho$ ($\rho\le n$) когредиентными строками $\xi$; инвариантные процессы дифференцирования возвращают к формам в исходных строках $p_\rho$.

Пусть $\xi^{(1)},\ldots,\xi^{(n)}$ -- строки величин преобразования, так что имеем:
\begin{equation}
\begin{aligned}
x_i&=\xi_i^{(1)}x'_1+\xi_i^{(2)}x'_2+\cdots+\xi_i^{(n)}x'_n,\\
p_{i_1\ldots i_\rho}&=\sum_j(\xi_{i_1}^{(j_1)}\cdots \xi_{i_\rho}^{(j_\rho)})p'_{j_1\ldots j_\rho}.
\end{aligned}
\tag{47}
\end{equation}
Преобразованные коэффициенты форм $f,\varphi,\ldots$ будем обозначать через $(f),(\varphi),\ldots$, и пусть
\[
 f=\pair{S^1}{p_1}^{m_1}\pair{S^2}{p_2}^{m_2}\cdots\pair{S^{n-1}}{p_{n-1}}^{m_{n-1}}.
\]
Тогда
\begin{equation}
\begin{aligned}
(f)&=\pair{S^1}{\xi^{(1)}}^{m_{11}}\cdots
\pair{S^1}{\xi^{(n)}}^{m_{1n}}
\pair{S^2}{\xi^{(1)}\xi^{(2)}}^{m_{21}}\cdots\\
&\quad\cdot
\pair{S^{n-1}}{\xi^{(2)}\cdots\xi^{(n)}}^{m_{n-1,n}}
\pair{s^{(1)}}{\xi^{(1)}}^{k_1}
\pair{s^{(2)}}{\xi^{(2)}}^{k_2}\cdots
\pair{s^{(n)}}{\xi^{(n)}}^{k_n},
\end{aligned}
\tag{48}
\end{equation}
где $m_{11}+\cdots+m_{1n}=m_1$ и так далее. Из каждого коэффициента $(f)$, для которого
\begin{equation}
 k_1\ge k_2\ge k_3\ge\cdots\ge k_n,
\tag{49}
\end{equation}
процесс дифференцирования, точно исчерпывающий строки $\xi$,
\begin{equation}
\begin{aligned}
&\left(\frac{\partial}{\partial\xi^{(1)}}\middle|p_1\right)^{k_1-k_2}
\left(\frac{\partial}{\partial\xi^{(1)}}\frac{\partial}{\partial\xi^{(2)}}\middle|p_2\right)^{k_2-k_3}\cdots\\
&\quad\cdot
\left(\frac{\partial}{\partial\xi^{(1)}}\cdots\frac{\partial}{\partial\xi^{(n-1)}}\middle|p_{n-1}\right)^{k_{n-1}-k_n}
\left(\frac{\partial}{\partial\xi^{(1)}}\cdots\frac{\partial}{\partial\xi^{(n)}}\right)^{k_n}
\end{aligned}
\tag{50}
\end{equation}
дает форму
\begin{equation}
\begin{aligned}
 \varphi&=\bigl(s^{(1)}\mid p_1\bigr)^{k_1-k_2}
 \bigl(s^{(1)}s^{(2)}\mid p_2\bigr)^{k_2-k_3}\cdots\\
 &\quad\cdot
 \bigl(s^{(1)}\cdots s^{(n-1)}\mid p_{n-1}\bigr)^{k_{n-1}-k_n}
 \bigl(s^{(1)}s^{(2)}\cdots s^{(n)}\bigr)^{k_n}.
\end{aligned}
\tag{51}
\end{equation}
Формы $\varphi$ и есть искомые нормальные формы. Действительно, они обладают следующими свойствами, определяющими нормальные формы.

1. Они являются инвариантными образованиями основной формы $f$, линейными по коэффициентам. Это следует из того, что они возникают из $f$ с помощью инвариантных процессов дифференцирования
$\left(\frac{\partial}{\partial p_\rho}\middle|\xi^{(i_1)}\xi^{(i_2)}\cdots\xi^{(i_\rho)}\right)$ и (50.). Следовательно, по концу \S\ 6, они линейно выражаются через поляры конечного ряда форм $f$. Эти поляры получают, рассматривая переменные $p_\rho$, входящие в $\varphi$, как коэффициенты формы, разложенной на линейные формы со строками переменных $p_\rho$:
\begin{equation}
\begin{aligned}
 H={}&(q_{n-1}\mid p_{n-1})^{k_1-k_2}
(q_{n-2}\mid p_{n-2})^{k_2-k_3}\cdots\\
&\cdot(q_1\mid p_1)^{k_{n-1}-k_n}.
\end{aligned}
\tag{52}
\end{equation}
Симультанные инварианты рядов форм $f$ и $H$ дают все поляры, которые здесь входят в рассмотрение, так как последние должны иметь в отдельных строках $p_\rho$ ту же размерность, что и сама $\varphi$.

2. Они удовлетворяют дифференциальным уравнениям (46.). Если применить к $\varphi$ дифференциальный процесс (45.), то получается:
\[
 \varphi'=\bigl(q_{n-\sigma-\lambda}s^{(1)}\cdots s^{(\sigma)}\mid [s^{(1)}\cdots s^{(\tau)}]_{n-\tau}p_{\tau-\lambda}\bigr),\qquad (\sigma>\tau).
\]
Из тождества разложения (30.) следует, если строки $q_\sigma=s^{(1)}\cdots s^{(\sigma)}$ и $p_{n-\tau}=[s^{(1)}\cdots s^{(\tau)}]_{n-\tau}$, образованные из строк $s$, отождествить с входящими там $q_\rho,p_\tau$, что
\[
 \varphi'=\sum_{\beta=\sigma-\tau+\lambda}^{n-\tau,\sigma}\sum_{i_\beta}\eps\,
 \pair{\dot{q}_{\sigma-\beta}^{(i_1\ldots i_\beta)}q_{n-\tau+\lambda}}{p_{\sigma+\lambda}\dot{p}_{n-\tau-\beta}^{(i_1\ldots i_\beta)}}.
\]
При этом
\[
 \dot{p}_{n-\tau-\beta}^{(i_1\ldots i_\beta)}\sim s^{(1)}\cdots s^{(\tau)}s^{(i_1)}\cdots s^{(i_\beta)},
\]
где $i_1,\ldots,i_\beta$ означают произвольные $\beta$ попарно различные числа, выбранные из чисел от 1 до $\sigma$. Так как $\beta>\sigma-\tau$, среди этих чисел $i_1,\ldots,i_\beta$ необходимо есть такие, которые лежат между 1 и $\tau$; поэтому $\dot{p}_{n-\tau-\beta}^{(i_1\ldots i_\beta)}$ тождественно исчезает, а следовательно исчезает каждое отдельное матричное произведение и вместе с ним $\varphi'$.

Для эквивалентности системы нормальных форм с рядом форм $f$ остается только показать, что все формы ряда форм раскладываются по полярам нормальных форм в смысле, зафиксированном в пункте 1. Пусть $\Pi$ -- агрегат полярных операций
\begin{equation}
 \left(\frac{\partial}{\partial\xi^{(i)}}\middle|\xi^{(k)}\right),
 \qquad (i\ne k;\ i=1,2,\ldots,\rho;\ k=1,2,\ldots,\rho),
\tag{53}
\end{equation}
и пусть $\Pi^\sigma$ -- агрегат тех специальных операций (53.), для которых $i=\sigma$ и $k<\sigma$. Для формы $F$ от $\rho$ строк $\xi^{(1)},\ldots,\xi^{(\rho)}$, имеющей в строке $\xi^{(\rho)}$ степень $k_\rho$, справедливо разложение\footnote{F. Mertens, \emph{\"Uber eine Formel der Determinantentheorie}. Sitzungsber. d. Ak. d. Wiss., Wien, math.-naturw. Kl., Bd. 91. 2. Abt. (1885). Формула 20).}:
\begin{equation}
 F=\sum_{\alpha=0}^{k_\rho}\Pi F_\alpha,
 \qquad(\rho\le n),
\tag{54}
\end{equation}
где $F_\alpha$ имеет степень $\alpha$ в последней строке $\xi^{(\rho)}$ и возникает из $F$ с помощью инвариантных дифференциальных операций
\begin{equation}
 \left(\frac{\partial}{\partial\xi^{(1)}}\cdots\frac{\partial}{\partial\xi^{(\rho)}}\middle|\xi^{(1)}\cdots\xi^{(\rho)}\right)^\alpha
\tag{55}
\end{equation}
и $\Pi^\rho$. Если при построении $F_\alpha$ заменить процесс (55.) на
\begin{equation}
 \left(\frac{\partial}{\partial\xi^{(1)}}\cdots\frac{\partial}{\partial\xi^{(\rho)}}\middle|p_\rho\right)^\alpha,
\tag{56}
\end{equation}
то возникает форма $F_\alpha(p_\rho)$ только с $\rho-1$ строками $\xi^{(1)},\ldots,\xi^{(\rho-1)}$, к которой снова можно применить (54.). Продолжение процедуры приводит к формам $G(p_\rho,p_{\rho-1},\ldots,p_1)$. Чтобы из них получить разложение $F$ по (54.), нужно отдельные строки $p_\rho$ заменить на $\xi^{(1)},\ldots,\xi^{(\rho)}$ и к возникающим формам применить полярный процесс $\Pi$ (53.); это дает, ср. (48.), сумму преобразованных коэффициентов $(G)$. При $\rho=n$ на первом шаге сохраняется процесс (55.). Если $c$ означает числовые постоянные и, для краткости, положено
\begin{equation}
\begin{aligned}
(\xi^{(1)}\xi^{(2)}\cdots\xi^{(n)})&\sim\Delta,\\
\left(\frac{\partial}{\partial\xi^{(1)}}\frac{\partial}{\partial\xi^{(2)}}\cdots\frac{\partial}{\partial\xi^{(n)}}\right)&=\nabla,
\end{aligned}
\tag{57}
\end{equation}
то в случае $\rho=n$ получается:
\begin{equation}
 F=\sum c\Delta^\alpha(G),
\tag{58}
\end{equation}
тогда как при $\rho<n$ справа входит только $\alpha=0$, то есть $\sum c(G)$.

Применим (58.) к преобразованному коэффициенту $(f)$ (48.). В силу коммутативности полярных процессов $\Pi^\sigma$ со всеми операциями (56.), для которых $\rho\ge\sigma$, и потому что поляры преобразованных коэффициентов снова дают преобразованные коэффициенты, получаем:
\[
\begin{aligned}
G={}&\left(\frac{\partial}{\partial\xi^{(1)}}\middle|p_1\right)^{\beta_1}
\left(\frac{\partial}{\partial\xi^{(1)}}\frac{\partial}{\partial\xi^{(2)}}\middle|p_2\right)^{\beta_2}\cdots\\
&\cdot
\left(\frac{\partial}{\partial\xi^{(1)}}\cdots\frac{\partial}{\partial\xi^{(n-1)}}\middle|p_{n-1}\right)^{\beta_{n-1}}
\nabla^{\beta_n}\Pi(f)=\sum c\varphi,
\end{aligned}
\]
и из (58.) следует:
\begin{equation}
 (f)=\sum c\Delta^\alpha(\varphi)\qquad\text{(Mertens, формула 23)).}
\tag{59}
\end{equation}
Чтобы перейти от (59.) к разложению самой $f$, снова рассмотрим переменные $p_\rho$ как коэффициенты формы $H$ (52.) со степенями $m_1,\ldots,m_{n-1}$, соответствующими $f$, и пусть $m=m_1+\cdots+m_{n-1}$. Тогда по определению инвариантов
\[
 f\cdot\Delta^m=\sum(f)(H)=\sum c\Delta^\alpha(\varphi)(H),
\]
а применение процесса $\nabla^m$, точно исчерпывающего строки $\xi$, дает:
\begin{equation}
 f=\sum c\Phi,
\tag{60}
\end{equation}
где формы $\Phi$ выводятся из $\varphi$ и $H$ инвариантными процессами дифференцирования по переменным, то есть являются симультанными инвариантами рядов форм $\varphi$ и $H$.\footnote{Вместо прямого заключения Mertens в \S\ 7 вводит ряд дальнейших дифференциальных процессов, которые по существу служат построению ряда форм $H$; у нас это выполняет Теорема VI.} Но так как $\varphi$ является нормальной формой, ряд форм $\varphi$ сводится к самой форме $\varphi$; ряд форм $H$ содержит все вообще возможные инвариантные соединения переменных. Поэтому симультанные инварианты $\varphi$ и $H$ являются полярами $\varphi$ в смысле, указанном в пункте 1; то есть (60.) дает разложение $f$ по полярам нормальных форм. Теперь нужно еще доказать такое же разложение для произвольной формы $\theta$ из ряда форм $f$. Пусть $\Gamma$ -- нормальные формы, выведенные из $(\theta)$ с помощью (50.), так что
\begin{equation}
 (\theta)=\sum c\Delta^\beta(\Gamma).
\tag{61}
\end{equation}
Формы $\theta$ характеризуются отношением, справедливым для каждого преобразованного коэффициента:
\[
 (\theta)\Delta^\sigma=\sum c(f)\qquad\text{(Mertens, формула 9)).}
\]
Но каждая форма $F$ от $n$ строк $\xi$, содержащая множитель $\Delta^\sigma$, содержит также второй выраженный член $\nabla^\sigma\Pi F$ (Mertens, формула 20)); отсюда
\[
 (\theta)=\sum c\nabla^\sigma(f),\qquad\text{следовательно}\qquad \Gamma=\sum c\varphi,
\]
и поэтому из (61.) имеем отношение, соответствующее (59.):
\[
 (\theta)=\sum c\Delta^\beta(\varphi),
\]
которое для $\theta$ влечет разложение (60.) по полярам нормальных форм $\varphi$. Подводя итог, получаем:

\emph{Теорема VII.} Каждый ряд форм $f$ можно заменить эквивалентной системой нормальных форм.
% END UNIT BODY
% END PRODUCER UNIT 036
\endgroup


\clearpage
\begingroup
\setcounter{footnote}{0}
% BEGIN PRODUCER UNIT 037
% Source: C:/Users/Floris/Documents/interlanguage/03_projects/noether/02_slavic_working_corpus/translations/paper04/russian/v001/Noether_Paper04_Section08_Russian_v001.tex
% Identity: 21844 B / SHA-256 59E06CC47AD5EFADA47DC3391F74C0BBA99547C61EF601ECA5C5AFDE9C95AC54
\setlength{\parindent}{1.25em}
\setlength{\parskip}{0pt}
\emergencystretch=10em
\allowdisplaybreaks
\newcommand{\pair}[2]{\left(#1\middle|#2\right)}
\newcommand{\eps}{\varepsilon}
% BEGIN UNIT BODY
\section*{\S\ 8. Сведение процесса свертки к основным сверткам.}

Опираясь на Теорему VII, мы теперь проводим упомянутое в начале \S\ 7 сведение процесса свертки к основным сверткам. При этом (ср. \S\ 5) всякую свертку двух когредиентных рядов $S^\rho,T^\rho$ будем называть ``когредиентной сверткой'', а именно когредиентную свертку дефекта 1 двух рядов $S^\rho,T^\rho$ -- ``основной сверткой типа $\rho$''. Тогда, в дополнение к Теореме VI, мы докажем следующую теорему.

\emph{Теорема VIII.} Общий процесс свертки (41.) можно составить из $(n-1)$ основных сверток типа $\rho$, где $\rho=1,2,\ldots,n-1$. Иными словами, для порождения всех инвариантных образований достаточно последовательного применения этих $n-1$ основных сверток.

В бинарной области процесс свертки (41.) прямо совпадает с единственно возможной основной сверткой; в тернарной области я доказала Теорему VIII в \S\ 1 своей диссертации (ср. Введение). Там, вследствие (44.), общие когредиентные свертки еще тождественны основным сверткам, и существенно, что Теорема VIII для $n$ переменных дает сведение именно к основным сверткам, а не только гораздо более слабое сведение к когредиентным сверткам. Доказательство по идее построено по тернарному образцу: самые общие свертки конструируются из основных сверток с помощью символических тождеств. Мы будем порождать:
\begin{enumerate}
\item произвольные свертки дефекта 1 с помощью когредиентных сверток дефекта 1, то есть с помощью основных сверток (аналог тернарного случая);
\item произвольные свертки дефекта $\lambda$ с помощью когредиентных сверток дефекта $\lambda$ и произвольных сверток дефекта 1;
\item когредиентные свертки дефекта $\lambda$ с помощью когредиентных сверток дефекта $\lambda-1$ и произвольных сверток дефекта 1.
\end{enumerate}
При этом существенно, что по Теореме VII мы предполагаем обе формы $S$ и $T$, которые нужно свернуть, нормальными формами. Следовательно, по (46.) имеют место соотношения:
\begin{equation}
\begin{aligned}
\pair{q_{n-\sigma_1-\lambda}S^{\sigma_1}}{S_{n-\sigma_2}p_{\sigma_2-\lambda}}&=0,
&& (\sigma_1\ge\sigma_2,\ \lambda=1,2,\ldots,\sigma_2,n-\sigma_1),\\
\pair{q_{n-\tau_1-\lambda}T^{\tau_1}}{T_{n-\tau_2}p_{\tau_2-\lambda}}&=0,
&& (\tau_1\ge\tau_2,\ \lambda=1,2,\ldots,\tau_2,n-\tau_1).
\end{aligned}
\tag{62}
\end{equation}
Далее, по \S\ 2 (конец), достаточно считать формы $S$ и $T$ линейными по всем строкам переменных; то есть построенные формы принадлежат следующему ряду форм:
\[
 f=\pair{S^1}{p_1}\pair{S^2}{p_2}\cdots\pair{S^{n-1}}{p_{n-1}}
   \pair{T^1}{p_1}\pair{T^2}{p_2}\cdots\pair{T^{n-1}}{p_{n-1}}.
\]
\noindent 1) Порождение свертки $\pair{q_{n-\sigma-1}S^\sigma}{T_{n-\tau}p_{\tau-1}}$, где $\sigma>\tau$, из основных сверток типа $\tau+\alpha$; $\alpha=0,1,\ldots,\sigma-\tau$. Из основной свертки типа $\tau$
\[
 \pair{q_{n-\tau-1}S^\tau}{T_{n-\tau}p_{\tau-1}}
\]
с помощью подстановки $w\sim T_{n-\tau}p_{\tau-1}$ получаем форму ряда форм $f$ с тремя строками верхнего весового индекса (ср. \S\ 1) $\tau+1$:
\begin{equation}
 \pair{wS^\tau}{p_{\tau+1}}\pair{S^{\tau+1}}{p_{\tau+1}}\pair{T^{\tau+1}}{p_{\tau+1}}.
\tag{63}
\end{equation}
Основная свертка типа $\tau+1$, примененная к (63.), дает:
\[
 \pair{q_{n-\tau-2}wS^\tau}{S_{n-\tau-1}p_\tau};
\]
отсюда по (31.), после подстановки $q_{n-\tau-2}w=q'_{n-\tau-1}$, получается значение:
\[
 \eps\pair{q_{n-\tau-2}wS^{\tau+1}}{S^\tau p_\tau}
 +\sum_{\alpha=1}^{\tau,n-\tau-1}\Pi\pair{q'_{n-\tau-1-\alpha}S^{\tau+1}}{S_{n-\tau}p_{\tau-\alpha}}.
\]
Выражение под знаком суммы тождественно исчезает по (62.); тем самым мы приходим к форме
\[
 \pair{wS^{\tau+1}}{p_{\tau+2}}\pair{S^{\tau+2}}{p_{\tau+2}}\pair{T^{\tau+2}}{p_{\tau+2}},
\]
которая получается из (63.) заменой $\tau$ на $\tau+1$. Поэтому $(\sigma-\tau)$-кратное повторение процесса приводит к искомой свертке:
\[
 \pair{wS^\sigma}{p_{\sigma+1}}=\eps\pair{q_{n-\sigma-1}S^\sigma}{T_{n-\tau}p_{\tau-1}}.
\]
Проведенный процесс кратко обозначим формулой:
\begin{equation}
 T^\tau S^\tau,\\ S^{\tau+1},\ S^{\tau+2},\ldots,S^\sigma;
\tag{64}
\end{equation}
мы видим, что используется только свойство $S$ быть нормальной формой, но не соответствующее свойство $T$. Дуалистический к (64.) процесс можно провести в обратном порядке, используя лишь свойство $T$ быть нормальной формой, а именно:
\begin{equation}
 S^\sigma T^\sigma,
\quad T^{\sigma-1},\quad T^{\sigma-2},\ldots,T^\tau,
\tag{65}
\end{equation}
по соотношениям:
\begin{align*}
\pair{q_{n-\sigma-1}S^\sigma}{T_{1-\sigma}p_{\sigma-1}}&=\eps\pair{T_{n-\sigma}y}{p_{\sigma-1}},\quad y\sim q_{n-\sigma-1}S^\sigma,\\
\pair{q_{n-\sigma}T^{\sigma-1}}{T_{n-\sigma}yp_{\sigma-2}}&=\eps\pair{T_{n-\sigma+1}y}{p_{\sigma-2}},\\
&\vdots\\
\pair{q_{n-\tau-1}T^\tau}{T_{n-\tau-1}yp_{\tau-1}}&=\eps\pair{q_{n-\sigma-1}S^\sigma}{T_{n-\tau}p_{\tau-1}}.
\end{align*}
Следует еще указать на связь с убыванием общего веса, приведенным в конце \S\ 6: для сверток дефекта 1 это убывание имеет значение $2(1+(\sigma-\tau))$, а для основных сверток -- значение $2\cdot1$. Поэтому, если составимость из основных сверток вообще возможна, необходимо требуются именно $\sigma-\tau+1$ таких сверток, как и было проведено выше.

\noindent 2) Порождение общих сверток из когредиентных и основных сверток. Убывание веса для общей свертки (41) равно
\[
2(\lambda^2+\lambda(\sigma-\tau))
=2[\lambda^2+(\sigma-\tau)(1+(\lambda-1))],
\]
и поэтому дает возможность разложения на когредиентную свертку дефекта $\lambda$ (убывание веса $2\lambda^2$) и $\sigma-\tau$ сверток дефекта 1 с разностью весовых индексов $\lambda-1$ (убывание веса $2\{1+(\lambda-1)\}$). В случае $\lambda=1$ это разложение совпадает с указанным в пункте 1); мы покажем, что оно действительно реализуется.

Из когредиентной свертки дефекта $\lambda$
\[
 \pair{q_{n-\tau-\lambda}S^\tau}{T_{n-\tau}p_{\tau-\lambda}}
\]
с помощью подстановки $R^\lambda\sim T_{n-\tau}p_{\tau-\lambda}$ получаем форму с двумя строками верхних весовых индексов $\tau+\lambda$ и $\tau+1$:
\begin{equation}
 \pair{R^\lambda S^\tau}{p_{\tau+\lambda}}\pair{S^{\tau+1}}{p_{\tau+1}}.
\tag{66}
\end{equation}
Строка с меньшим весовым индексом $\tau+1$ принадлежит нормальной форме; следовательно, (66.) допускает свертку дефекта 1, составленную из $\lambda$ основных сверток, в порядке (65.):
\[
 (R^\lambda S^\tau)S^{\tau+\lambda},\quad S^{\tau+\lambda-1},\quad S^{\tau+\lambda-2},\ldots,S^{\tau+1}.
\]
Это дает, с учетом (31.) и (62.):
\[
 \pair{q_{n-\tau-\lambda-1}R^\lambda S^\tau}{S_{n-\tau-1}p_\tau}
 =\eps\pair{R^\lambda S^{\tau+1}}{p_{\tau+\lambda+1}},
\]
то есть форму, которая получается из (66.) заменой $\tau$ на $\tau+1$, аналогично тому, как в пункте 1) это произошло с (63.). $(\sigma-\tau)$-кратное повторение процесса приводит к искомой свертке:
\[
 \pair{R^\lambda S^\sigma}{p_{\sigma+\lambda}}
 =\eps\pair{q_{n-\sigma-\lambda}S^\sigma}{T_{n-\tau}p_{\tau-\lambda}}.
\]
Во всем этом процессе использовалось только свойство $S$ быть нормальной формой; рассмотрение $T$ как нормальной формы приводит к дуалистическому процессу с полным обращением порядка, по формулам:
\begin{align*}
\pair{q_{n-\sigma-\lambda}S^\sigma}{T_{n-\sigma}p_{\sigma-\lambda}}&=\eps\pair{T_{n-\sigma}R_\lambda}{p_{\sigma-\lambda}},\quad R_\lambda\sim q_{n-\sigma-\lambda}S^\sigma,\\
\pair{q_{n-\sigma}T^{\sigma-1}}{T_{n-\sigma}R_\lambda p_{\sigma-\lambda-1}}&=\eps\pair{T_{n-\sigma+1}R_\lambda}{p_{\sigma-\lambda-1}},\\
&\vdots\\
\pair{q_{n-\tau-1}T^\tau}{T_{n-\tau-1}R_\lambda p_{\tau-\lambda}}&=\eps\pair{q_{n-\sigma-\lambda}S^\sigma}{T_{n-\tau}p_{\tau-\lambda}}.
\end{align*}
\noindent 3) Порождение когредиентных сверток дефекта $\lambda$ из когредиентных сверток дефекта $\lambda-1$ и основных сверток. Убывание веса для когредиентных сверток дефекта $\lambda$ равно
\[
2\lambda^2=2((\lambda-1)^2+2\lambda-1),
\]
поэтому оно допускает возможность разложения на когредиентную свертку дефекта $\lambda-1$ и $2\lambda-1$ основных сверток. Это разложение проще всего найти, если сначала специально положить $n=2\lambda$. Единственно возможной когредиентной сверткой дефекта $\lambda$ будет $\pair{S^\lambda}{T_\lambda}=\pair{S^\lambda}{T_{n/2}}$; она содержит на одну строку $u$ и одну строку $x$ меньше, чем свертка дефекта $\lambda-1$ тех же самых рядов $\pair{uS^\lambda}{T_\lambda x}$. Обратно, мы составляем свертку дефекта $\lambda$ из основных сверток, которые вызывают исчезновение одной строки $u$ и одной строки $x$ (убывание веса $2(n-1)=2(2\lambda-1)$) и одновременно порождают новую символическую строку $R^\lambda$, а также из когредиентной свертки дефекта $\lambda-1$. Простейшая свертка, выполняющая это, есть следующая свертка дефекта 1:
\[
 \pair{sT^\lambda}{\sigma p_\lambda}
 =\eps\pair{S^{2\lambda-1}}{R_{\lambda-1}p_\lambda}
 =\eps\pair{R^\lambda}{p_\lambda},
\]
где положено:
\[
 R^{\lambda+1}=sT^\lambda,
 \quad s=S^1,
 \quad \sigma=S_1,
 \quad R^\lambda\sim\sigma R_{\lambda-1}.
\]
По (65.) и (64.) она составлена из следующих $2\lambda-1$ основных сверток:
\begin{equation}
 T^\lambda S^\lambda,S^{\lambda-1},\ldots,S^1;
 \qquad R^{\lambda+1}S^{\lambda+1},S^{\lambda+2},\ldots,S^{2\lambda-1}.
\tag{67}
\end{equation}
С учетом (21.) и (62.) свертка дефекта $\lambda-1$ дает:
\[
 \pair{uS^\lambda}{\sigma R_{\lambda-1}x}=\eps\pair{R^{\lambda+1}}{S_\lambda x}
 =\eps\pair{sT^\lambda}{S^\lambda x}=\eps\pair{S^\lambda}{T_\lambda}.
\]
Общий случай, для двух рядов $S^\rho,T^\rho$ и произвольного $n$, теперь прямо переносится из специального случая. Вместо (67.) появляются следующие $2\lambda-1$ основных сверток:
\begin{equation}
 T^\rho S^\rho,S^{\rho-1},\ldots,S^{\rho-\lambda+1};
 \qquad R^{\rho+1}S^{\rho+1},S^{\rho+2},\ldots,S^{\rho+\lambda-1}.
\tag{68}
\end{equation}
Первая половина сверток дает форму:
\[
 \pair{q_{n-\rho-1}T^\rho}{S_{n-\rho+\lambda-1}p_{\rho-\lambda}},
\]
из которой с помощью подстановок
\[
 S_{n-\rho+\lambda-1}p_{\rho-\lambda}\sim w,
 \qquad T^\rho w=R^{\rho+1}
\]
и применения второй половины сверток получается:
\[
 \pair{q_{n-\rho-\lambda}S^{\rho+\lambda-1}}{R_{n-\rho-1}p_\rho}.
\]
Если подставить:
\[
 q_{n-\rho-\lambda}S^{\rho+\lambda-1}\sim y,
\]
то через когредиентную свертку дефекта $\lambda-1$ получаем:
\begin{equation}
 \pair{q_{n-\rho-\lambda+1}S^\rho}{yR_{n-\rho-1}p_{\rho-\lambda+1}}.
\tag{69}
\end{equation}
Покажем, что это и есть искомая свертка дефекта $\lambda$. Если подставить
\[
 R_{n-\rho-1}p_{\rho-\lambda+1}=R_{n-\lambda}
\]
и учесть, что после разрешения относительно $y$ получается
\[
 \pair{q_{n-\rho-\lambda+1}R^\lambda}{S_{n-\rho}y}
 =\eps\pair{q_{n-\rho-\lambda}S^{\rho+\lambda-1}}{S_{n-\rho}p'_{\rho-1}}=0,
\]
то по (20.) для (69.) получается значение:
\[
 \eps\pair{S^\rho R^\lambda}{p_{\rho+\lambda-1}y}
 =\eps\pair{q_{n-\rho-\lambda}S^{\rho+\lambda-1}}{[S^\rho R^\lambda]_{n-\rho-\lambda}p_{\rho+\lambda-1}}.
\]
Но эта свертка имеет тип (43.), а значит, эквивалентна
\[
 \pair{q_{n-\rho-\lambda}S^\rho}{R_{n-\rho-1}p_{\rho-\lambda+1}}
 =\eps\pair{wT^\rho}{R_\lambda p_{\rho-\lambda+1}},
 \qquad R_\lambda\sim q_{n-\rho-\lambda}S^\rho.
\]
Здесь $R_\lambda$ уже имеет требуемую конечную форму; выражение соответствует форме $\pair{sT^\lambda}{S_\lambda x}$ в специальном случае. Если теперь учесть, что после разрешения относительно $w$ и $R_\lambda$ получается
\[
\begin{aligned}
 \pair{wR^{n-\lambda}}{T_{n-\rho}p_{\rho-\lambda+1}}
 &=\eps\pair{q'_{n-1}R^{n-\lambda}}{S_{n-\rho+\lambda-1}p_{\rho-\lambda}}\\
 &=\eps\pair{q''_{n-\sigma-1}S^\rho}{S_{n-\rho+\lambda-1}p_{\rho-\lambda}}=0,
\end{aligned}
\]
то по (21.) получается значение:
\[
\begin{aligned}
 \pair{wq_{n-\rho+\lambda-1}}{T_{n-\rho}R_\lambda}
 &=\eps\pair{q_{n-\rho+\lambda-1}[T_{n-\rho}R_\lambda]^{n-\lambda}}{S_{n-\rho+\lambda-1}p_{\rho-\lambda}}\\
 &\quad\text{эквивалентное}\quad
 \pair{q_{n-\rho-\lambda}S^\rho}{T_{n-\rho}p_{\rho-\lambda}}.
\end{aligned}
\]
Во всем процессе использовалось только свойство $S$ быть нормальной формой, но не соответствующее свойство $T$; следовательно, когредиентную свертку дефекта $\lambda-1$, нужную для порождения (69.), можно заменить $2\lambda-3$ основными свертками и когредиентной сверткой дефекта $\lambda-2$, которая, в свою очередь, снова допускает разложение на основные свертки, и т. д. Аналогично получается порождение при использовании свойства $T$ быть нормальной формой и применении дуалистического процесса, то есть основных сверток:
\[
\begin{aligned}
 S^\rho T^\rho,\quad T^{\rho-1},\ldots,T^{\rho-\lambda+1};\\
 \qquad R^{\rho-1}T^{\rho-1},\quad T^{\rho-2},\ldots,T^{\rho-\lambda+1},
\end{aligned}
\]
и свертки, дуалистической к (69.). Теперь возникает свертка, эквивалентная предыдущей:
\[
 \pair{q_{n-\rho-\lambda}T^\rho}{S_{n-\rho}p_{\rho-\lambda}}.
\]

Добавляя разложение, полученное в пункте 2), мы действительно получаем порождение самых общих сверток из основных сверток. Остается еще доказать, что это порождение, не считая перестановки порядка, однозначно, то есть что каждой свертке соответствует одно и только одно число $\alpha_\rho$ основных сверток типа $\rho$, где $\rho=1,2,\ldots,n-1$.

Пусть $m_\rho$ обозначают степени в строках переменных $p_\rho$ исходной формы $f$, а $l_\rho$ -- степени формы, возникшей через свертку. Каждая основная свертка типа $\rho$ преобразует степени $m_{\rho-1},m_\rho,m_{\rho+1}$ соответственно в $m_{\rho-1}+1,m_\rho-2,m_{\rho+1}+1$. Следовательно, числа $\alpha$ удовлетворяют системе уравнений:
\begin{equation}
 m_\rho-l_\rho=-\alpha_{\rho-1}+2\alpha_\rho-\alpha_{\rho+1},
 \qquad (\rho=1,2,\ldots,n-1),
\tag{70}
\end{equation}
где следует положить $\alpha_0=\alpha_n=0$, и определитель этой системы имеет значение $n$.\footnote{Если $\Delta_n$ обозначает $n$-строчный определитель, то справедлива рекуррентная формула $\Delta_n=2\Delta_{n-1}-\Delta_{n-2}$; то есть $\Delta_n-\Delta_{n-1}=\Delta_{n-1}-\Delta_{n-2}=\cdots=\Delta_2-\Delta_1=1$, так как $\Delta_2=3$, $\Delta_1=2$.} Тем самым числа $\alpha$ определены однозначно, причем по Теореме VIII как положительные целые числа; это дает условия для разностей $m_\rho-l_\rho$.

Результаты этого параграфа справедливы в силу соотношений (62.); однако если соотношения (62.) не выполняются, вовсе не следует считать, что конечные выражения эквивалентны исходным. Данное в \S\ 6 определение эквивалентности допускает также такую формулировку:\footnote{Ср. с рассуждением в конце следующего параграфа.} ``Две формы тогда и только тогда эквивалентны, когда они различаются на формы более высокой свертки, то есть на формы с большим убыванием веса''. Это большее убывание веса всегда имеет место, когда обе входящие в свертку строки $q_{n-\sigma-i},p_{\tau-\lambda}$ действительно являются строками переменных, как это было в (41.), (42.), в свертках этого параграфа, выступающих как эквивалентные, или также в эквивалентных формах Теоремы VII. Но оно не обязано иметь место, когда эти строки получены из символических строк подстановкой, как это было с формами этого параграфа, которые в силу (62.) выражались друг через друга. Легко видеть, что исчезающие формы отчасти даже имеют больший общий вес.
% END UNIT BODY
% END PRODUCER UNIT 037
\endgroup


\clearpage
\begingroup
\setcounter{footnote}{0}
% BEGIN PRODUCER UNIT 038
% Source: C:/Users/Floris/Documents/interlanguage/03_projects/noether/02_slavic_working_corpus/translations/paper04/russian/v001/Noether_Paper04_Section09_Russian_v001.tex
% Identity: 11430 B / SHA-256 EBDF1756CF42ED1C6BD9029EF610F2B3B63D7F13C0AA132B08C81E753CE0D784
\setlength{\parindent}{1.25em}
\setlength{\parskip}{0pt}
\emergencystretch=10em
\allowdisplaybreaks
\newcommand{\pair}[2]{\left(#1\middle|#2\right)}
\newcommand{\eps}{\varepsilon}
% BEGIN UNIT BODY
\section*{\S\ 9. Ряды форм. Редукционные теоремы.}

Введенный в \S\ 6 (в конце) ряд форм теперь определим несколько точнее как ``совокупность инвариантных образований данной исходной формы, линейных по коэффициентам, то есть совокупность форм, возникающих из исходной формы через свертку ее с самой собой, в порядке более высоких форм''.

Для объяснения более высоких форм пусть исходная форма, по Теореме VII, задана как произведение двух нормальных форм $S\cdot T$; более высокой формой мы будем называть форму с более высокой сверткой с самой собой, а среди форм с одинаковой такой сверткой -- специально нормированную форму. Точнее: пусть форма $A$ возникла посредством $\alpha_\rho$ основных сверток типа $\rho$, а форма $B$ -- посредством $\beta_\rho$ основных сверток того же типа. Тогда $B$ выше, чем $A$:
\begin{enumerate}
\item если в системе неравенств $\beta_\rho\ge\alpha_\rho$, $\rho=1,2,\ldots,n-1$, хотя бы для одного значения $\rho$ неравенство строгое;\footnote{Если в системе неравенств отчасти стоит знак $>$, а отчасти знак $<$, то формы несравнимы между собой, поскольку форме, возникшей из другой посредством свертки, всегда соответствует положительная система значений основных сверток, значит в нашем случае положительные или нулевые значения $\beta_\rho-\alpha_\rho$.}
\item если для всех значений $\rho$ стоит знак равенства, но меньше символических строк объединено в одно матричное произведение. Эта нормировка оказывается необходимой из-за редукций в \S\ 8. Соответственно, например,
\[
\pair{S^{n-1}}{T_{n-1}}
\quad\text{выше, чем}\quad
\pair{S^{n-1}}{[S^1T^\rho]_{n-\rho-1}p_\rho}.
\]
\item Если для всех значений $\rho$ стоит знак равенства и число символических строк, объединенных в одно матричное произведение, одинаково, то $B$ становится выше, чем $A$, благодаря нормировке, самой по себе произвольной, но фиксированной раз и навсегда. Такая нормировка всегда возможна из-за конечного числа форм, принадлежащих одной системе значений $\alpha$; ее можно, например, осуществить выделением одной формы $S$ (большее число символических строк $S$, большая сумма верхних весовых индексов строк $S$ и так далее).
\end{enumerate}
Удобно представлять ряд форм упорядоченным как $(n-1)$-мерное множество решеточных точек, где $n-1$ направлений соответствуют $n-1$ основным сверткам. Тогда каждой форме соответствует одна и только одна система значений $\alpha$, то есть положительная целочисленная решеточная точка, тогда как различные формы, принадлежащие одной системе значений $\alpha$, однозначно определены в своем порядке указанной выше нормировкой. Произвольная форма ряда форм задает исходную форму нового ряда форм, который содержится как часть в общем ряду форм; именно как та часть, которая состоит из форм, получаемых из новой исходной формы движением в $n-1$ положительных направлениях.

В силу Теоремы VIII и общей теории рядовых разложений для рядов форм имеют место в точности те же редукционные теоремы, что и в бинарной и тернарной областях (ср. диссертацию, \S\ 3); главную из них мы еще кратко выведем. Сначала дадим определение:

``В заданной системе форм пусть некоторое конечное число форм объединено в модуль; редуцируемой будем называть форму, которую можно выразить через формы, имеющие инварианты множителями, или через ``более высокие формы''; то есть через более высокие формы общего ряда форм, которому эта форма принадлежит, или через формы, содержащие символы модуля в более высоком порядке. Редуцируемый ряд форм назовем редуцентом.'' Тогда:

\emph{Теорема IX.} Если исходная форма ряда форм редуцируема вследствие того, что она возникла посредством свертки с редуцентом, то весь ряд форм, возникающий из этой исходной формы, редуцируем.

Для доказательства учтем, что форму $h$, возникшую посредством свертки формы $f$ со второй формой $g$, ввиду (45.) можно рассматривать как форму $f$, содержащую несколько когредиентных строк переменных $p_\rho,p'_\rho$. Такие формы по общей теории рядового разложения можно представить как поляры элементарных ковариантов $f$, последовательно применяя рядовое разложение к каждой паре когредиентных строк переменных $p_\rho,p'_\rho$. При этом по (38.) возникающие элементарные коварианты являются формами, порожденными когредиентной сверткой. Но так как порядок применения рядового разложения еще произволен, мы можем, в частности, выбрать такой порядок, который соответствует порождению общих сверток из основных сверток. Тогда получающаяся система элементарных ковариантов тождественна ряду форм $f$; именно формы, возникающие посредством когредиентной свертки более высокого дефекта, упорядочиваются среди форм, порожденных основными свертками. К полярам ряда форм $f$ приходят также при произвольном порядке рядового разложения, которому может соответствовать вторая система элементарных ковариантов. Поскольку ряд форм исчерпывает совокупность инвариантных образований, всякую форму второй системы можно линейно выразить через поляры ряда форм, причем через поляры тех форм, которые имеют такое же или большее убывание веса. Последнее следует из того, что (ср. \S\ 7) поляры можно рассматривать как симультанные инварианты ряда форм $H$ (52.) с соответствующими форме $h$ степенями и ряда форм $f$. Каждой свертке $H$ с самой собой соответствует убывание веса, которое после выполнения подстановки (11.) переносится на символические строки, а тем самым и на формы из $f$.\footnote{Это рассуждение лежит также в основе второй формулировки понятия эквивалентности (\S\ 8, конец). Более подробные объяснения различных систем элементарных ковариантов в тернарной области имеются у Study, там же, II, \S\ 7, теоремы 7) и 8). В силу нашей Теоремы VII то же самое справедливо и для $n$ переменных.}

Произвольная форма ряда форм $h$ возникла посредством свертки $h$ с самой собой; то есть посредством более высокой свертки $g$ с $f$, с одной стороны, и посредством свертки $f$ или $g$ с самими собой, с другой. В обоих случаях рядовое разложение, так же как выше было показано для исходной формы $h$, приводит к полярам ряда форм $f$. Если, в частности, $f$ является редуцентом, то получается разложение по полярам редуцента; следовательно, ряд форм $h$ становится редуцируемым.

Приложения теоремы к редукции полных систем форм получаются аналогично бинарной и тернарной областям.
% END UNIT BODY
% END PRODUCER UNIT 038
\endgroup


\clearpage
\begingroup
\setcounter{footnote}{0}
% BEGIN PRODUCER UNIT 039
% Source: C:/Users/Floris/Documents/interlanguage/03_projects/noether/02_slavic_working_corpus/translations/paper05/russian/v001/Noether_Paper05_Russian_v001.tex
% Identity: 13473 B / SHA-256 E8F378A0DB715BF391EBAF9C111A04B67496637B3EDC93DFA8E9183412640CC1
\setlength{\parindent}{1.25em}
\setlength{\parskip}{0pt}
\emergencystretch=10em
\allowdisplaybreaks
\providecommand{\Kfield}{\mathfrak{K}}
\providecommand{\Ssys}{\mathfrak{S}}
\providecommand{\Mfield}{\mathfrak{M}}
\providecommand{\tq}{\tau}
\providecommand{\ds}{\displaystyle}
% BEGIN UNIT BODY
\begin{center}
{\Large\bfseries 5. Поля рациональных функций}\par
\vspace{1.2em}
J. Ber. d. DMV 22 (1913), S. 316--319
\end{center}

\vspace{2em}

Следующие вопросы первоначально восходят к беседам с Э. Фишером. Некоторые из них, впрочем, для специального случая одной неопределенной -- причем при более общих предположениях о поле коэффициентов -- уже были поставлены и решены Э. Штайницем.\footnote{E. Steinitz: Algebraische Theorie der Körper (особенно \S\ 24). Crelles Journal 137. 1910.}

Под ``полем рациональных функций'' я понимаю поле, элементы которого являются рациональными функциями от $n$ неопределенных с коэффициентами из заранее заданного числового поля, которое, в частности, может содержать и все комплексные числа. Примерами таких полей служат поле симметрических функций от $n$ величин или, более общо, совокупность рациональных функций от $n$ неопределенных, допускающих перестановки некоторой группы (лагранжевы области родов); далее -- поле инвариантов.

Между двумя первыми названными полями и последним имеется характерное различие: первые содержат $n$ алгебраически независимых функций, элементарные симметрические функции; тогда как число алгебраически независимых инвариантов всегда меньше числа неопределенных. Поэтому важно, что вообще такие поля второго типа можно взаимно однозначно сопоставлять с полями первого типа, заменяя некоторые из неопределенных числами. Поэтому далее я ограничусь полями первого типа, то есть полями с $n$ алгебраически независимыми функциями; в силу этого сопоставления результаты затем имеют общий характер.

Вопросы группируются вокруг трех базисных понятий: рациональная база, минимальная база и интегральная база.

1. Под ``рациональной базой'' я понимаю конечное число функций поля такое, что всякую функцию поля можно представить как рациональную комбинацию этого конечного числа функций с коэффициентами из заданной области коэффициентов. Простыми рассуждениями -- которые в случае одной неопределенной имеются также у Э. Штайница, там же -- доказывается существование рациональной базы для всякого поля рациональных функций. Например, по теореме Лагранжа элементарные симметрические функции и одна функция, принадлежащая группе, образуют рациональную базу для упомянутых выше лагранжевых областей родов.

2. Всякая рациональная база должна (для полей первого типа) содержать по меньшей мере $n$ функций; если она содержит ровно $n$ функций, которые тогда должны быть алгебраически независимыми, я называю ее ``минимальной базой''. Вопрос о существовании минимальной базы частично решается теоремами Люрота, Кастельнуово и Энрикеса.\footnote{J. Lüroth: Beweis eines Satzes über rationale Kurven. Math. Ann. 9. 1875. --- G. Castelnuovo: Sulla razionalità delle involuzioni piane. Math. Ann. 44. 1893. --- F. Enriques: Sopra una involuzione non razionale dello spazio. Rend. Acc. Linc. vol. 21. 21. Jan. 1912.} Из них следует, что для полей от одной неопределенной минимальная база всегда существует: это функция Люрота данного поля, единственно определенная с точностью до дробно-линейного преобразования (ср. Штайниц, \S\ 24). Для полей от двух неопределенных минимальная база существует всегда и, вообще говоря, только тогда, когда допускается алгебраическое расширение поля коэффициентов, тогда как для трех и более неопределенных минимальная база вообще не существует. Охарактеризовать специальные поля с минимальной базой еще не пытались.

Теперь я продолжила исследовать вопрос о минимальной базе для лагранжевых областей родов. Здесь он приобретает следующий смысл: ``Если лагранжева область родов, принадлежащая группе $G$, обладает минимальной базой, то всю совокупность уравнений с заданной группой $G$ можно рационально построить посредством параметрического представления; и именно для всякого числового поля, содержащего область коэффициентов минимальной базы -- следовательно, в частности, для любого числового поля, если эта область коэффициентов есть поле рациональных чисел.''

Самый простой пример дает симметрическая группа; здесь элементарные симметрические функции образуют минимальную базу с рациональными числовыми коэффициентами. Отсюда в качестве параметрического представления получается просто уравнение с неопределенными коэффициентами. Как от него перейти к сколь угодно многим уравнениям ``без аффекта'' над каждым числовым полем, показывает теорема Гильберта о неприводимости.

Теперь непосредственно из теорем Люрота и Кастельнуово следует существование минимальной базы для всех групп, возникающих при уравнениях 3-й и 4-й степеней; при этом далее оказывается, что эта минимальная база имеет рациональные числовые коэффициенты. Значит, над произвольным числовым полем можно рационально построить совокупность уравнений 3-й и 4-й степеней с заданной группой. Для циклических и диэдральных групп существует минимальная база с полем $n$-х корней из единицы в качестве области коэффициентов; то же справедливо и для некоторых других метациклических групп. Вопрос, является ли минимальная база вообще характерной для некоторых групп, должен остаться нерешенным.

3. Чтобы перейти к последнему базисному понятию, я рассматриваю совокупность полиномов, содержащихся в поле. Под ``интегральной базой'' я понимаю конечное число этих полиномов такое, что всякий полином поля можно представить как целую рациональную комбинацию этого конечного числа полиномов с коэффициентами из заданной области коэффициентов. Вопрос о существовании интегральной базы -- который, например, как специальный случай включает конечность системы инвариантов -- был поставлен Гильбертом в его ``Математических проблемах''\footnote{D. Hilbert: Mathematische Probleme. Доклад, Париж 1900. Проблема 14. Göttinger Nachr. 1900.} в несколько иной формулировке, как проблема относительно целых функций.

Пока я могу указать лишь один класс полей, для которых интегральная база существует. А именно, если поле содержит систему из $n$ полиномов такую, что однородная результанта членов наивысшей степени этих полиномов отлична от нуля, то интегральная база существует; ее дают коэффициенты всех $u$ в неприводимом над полем уравнении для линейной формы:
\[
  u_0=u_1x_1+u_2x_2+\cdots+u_nx_n.\footnotemark
\]
\footnotetext{Это неприводимое уравнение в случае одной неопределенной встречается у Э. Штайница в доказательстве теоремы Люрота.}
Если, в частности, значение результанты равно единице области коэффициентов, то обеспечена также интегральность представления. Пример этого класса полей дают лагранжевы области родов, поскольку результанта элементарных симметрических функций имеет значение $\pm 1$. Еще один пример (без интегральности) дают все поля, содержащие только один алгебраически независимый полином; интегральная база здесь тождественна функции Люрота наименьшего подполя, содержащего все полиномы. Так, как известно, все целые рациональные проективные инварианты квадратичной формы от $n$ переменных являются целыми функциями дискриминанта.

\enlargethispage{2\baselineskip}
Указанное условие является только достаточным, отнюдь не необходимым для существования интегральной базы. Легко построить поля, в которых результанта любых $n$ полиномов обращается в нуль и которые тем не менее обладают интегральной базой, заданной коэффициентами того неприводимого уравнения. Возникает вопрос, дают ли, быть может, и в самом общем случае коэффициенты того уравнения интегральную базу.
% END UNIT BODY
% END PRODUCER UNIT 039
\endgroup


\clearpage
\begingroup
\setcounter{footnote}{0}
% BEGIN PRODUCER UNIT 040
% Source: C:/Users/Floris/Documents/interlanguage/03_projects/noether/02_slavic_working_corpus/translations/paper06/russian/v001/Noether_Paper06_Introduction_Russian_v001.tex
% Identity: 11129 B / SHA-256 C2331538827F4680DAD2A7036D0395159270BD5A4CDCA014054C25972A590F39
\setlength{\parindent}{1.25em}
\setlength{\parskip}{0pt}
\emergencystretch=10em
\allowdisplaybreaks
\providecommand{\Kfield}{\mathfrak{K}}
\providecommand{\Ssys}{\mathfrak{S}}
\providecommand{\Mfield}{\mathfrak{M}}
\providecommand{\tq}{\tau}
\providecommand{\ds}{\displaystyle}
% BEGIN UNIT BODY
\begin{center}
{\Large\bfseries 6. Поля и системы рациональных функций}\par
\vspace{1.2em}
Math. Ann. 76 (1915), S. 161--196
\end{center}

\vspace{2em}

Настоящая работа рассматривает вопросы о базах для произвольных систем рациональных и целых рациональных функций. Используемые здесь методы теории полей показывают, что существенным является решение этих вопросов для полей из рациональных функций --- полей рациональных функций\footnote{Ср. предварительное сообщение по случаю Венского съезда естествоиспытателей: Rationale Funktionenkörper --- Jahresber. d. D. Math.-Ver. 22 (1913) ---, где дан обзор постановок вопросов и результатов, относящихся к функциональным полям.} ---, тогда как обобщение результатов на произвольные системы получается как следствие.

Из вопросов о базах для общих систем до сих пор было известно лишь существование модульной базы, обеспеченное теоремой Гильберта (Math. Ann. 36). Ниже вопрос о рациональной представимости полностью решается через существование рациональной базы для всякой произвольной системы (\S\ 7); рациональная база полей получается уже в \S\ 4. Это существование рациональной базы позволяет всюду исходить из абстрактно определенного поля или системы и тем самым избегать трудностей, которые вызваны только специальным выбором рациональной базы, а не самой системой, например появлением специальных знаменателей или фундаментальных точек базисных функций.

Для полей будет также рассмотрен вопрос о минимальной базе, то есть о рациональной базе из алгебраически независимых функций (\S\ 6). Методы теории полей далее приводят к выделенной рациональной базе, инволюционной базе\footnote{Название выбрано потому, что рациональное функциональное поле в геометрическом толковании задает самую общую инволюцию в линейном пространстве.} (\S\ 5), которая становится особенно существенной для интегральных областей из полиномов (\S\ 8). Здесь она дает представление с фиксированным знаменателем (степенью функции, определенной интегральной областью), как это известно, например, в специальном случае типического представления инвариантов.

Из рациональной представимости теперь выводятся как можно более далеко идущие следствия для целой рациональной представимости, то есть для вопроса о конечности в более узком смысле.

Все известные до сих пор теоремы о конечности опираются на тот факт, что теорема Гильберта о модульной базе гарантирует конечность для всех таких систем, в которых представление
\[
  F=A_1f_1+A_2f_2+\cdots+A_kf_k
\]
влечет второе представление
\[
  F=B_1f_1+B_2f_2+\cdots+B_kf_k,
\]
где $B_i$ принадлежат системе и имеют меньшую степень, чем $F$.

Напротив, теорема о рациональной базе и методы теории полей позволяют заключать о конечности при предположениях совсем другого рода.

Так, как частичный ответ на одну из проблем Гильберта (Mathematische Probleme 14), получается класс относительно целых функций (относительно целые области первого типа), которые являются конечными интегральными областями и чья интегральная база задается упомянутой выше инволюционной базой (\S\ 10). По аналогии с гильбертовым доказательством теоремы Кронекера о фундаментальной системе алгебраически целых величин (Volle Invariantensysteme, \S\ 2; Math. Ann. 42) далее можно показать, что все регулярные системы из полиномов --- то есть все системы, которые можно отобразить в системы без фундаментальных точек, --- обладают интегральной базой (\S\ 12), тогда как нерегулярные системы без интегральной базы легко указать. В качестве простого примера этого факта упомянем теорему: ``Во всякой произвольной системе из бесконечно многих полиномов от одной неопределенной существует конечное число этих полиномов такое, что всякий полином системы можно представить как целую рациональную комбинацию этого конечного числа.'' Дальнейшие примеры находятся в \S\ 13.

Наконец, последние параграфы (\S\S\ 14 и 15) дают усиление теорем о базах в отношении целоалгебраичности.

Существенным вспомогательным средством исследования служит принцип переноса из \S\ 3, согласно которому достаточно рассматривать все поставленные здесь вопросы о базах для таких систем, в которых число неопределенных совпадает с числом алгебраически независимых функций системы.

Толчком к настоящей работе стали беседы с Э. Фишером, в частности его вопрос о минимальной базе лагранжевых областей родов. Относящиеся к этому исследования, которые привели к построению уравнений с заданной группой (ср. цитированное сообщение о ``Полях рациональных функций'', часть 2), оставлены для дальнейшей публикации.

Существенное расширение этой работы по сравнению с первоначальным сообщением состоит в том, что теоремы о базах, сформулированные там только для полей или специальных интегральных областей, здесь перенесены на произвольные системы. Этот перенос удается простым образом благодаря понятию наименьшего поля, содержащего произвольную систему, как обобщению поля частных интегральной области (\S\ 7). К образованию этого понятия меня привело то, что К. Генцельт, узнав рациональную базу полей, смог доказать существование рациональной базы для произвольных систем обходным путем через гильбертову модульную базу. Также к теореме об интегральной базе регулярных систем меня привело замечание К. Генцельта, что рассуждения, примененные мной к интегральным областям, справедливы для произвольных систем, подчиняющихся тем же предположениям; кроме того, я обязана Генцельту еще несколькими отдельными небольшими замечаниями.

Наконец, отметим, что в случае одной неопределенной рациональная база и минимальная база полей имеются у Э. Штайница в его ``Algebraische Theorie der Körper'', \S\ 24 (Crelles Journ. 137); при этом область коэффициентов там берется как поле в самом общем смысле. Вопрос, насколько при этих более общих предположениях сохраняются данные здесь теоремы, ниже не затрагивается; во всяком случае доказательства требуют модификации, поскольку, например, уже средства из теории функциональных детерминант отказывают для полей характеристики $p$.
% END UNIT BODY
% END PRODUCER UNIT 040
\endgroup


\clearpage
\begingroup
\setcounter{footnote}{0}
% BEGIN PRODUCER UNIT 041
% Source: C:/Users/Floris/Documents/interlanguage/03_projects/noether/02_slavic_working_corpus/translations/paper06/russian/v001/Noether_Paper06_Section01_Russian_v001.tex
% Identity: 8441 B / SHA-256 A6B7095F2F09F3D5E664D5AFE7D84CA2B744EF0110E247ABC8AB5B22C827BA52
\setlength{\parindent}{1.25em}
\setlength{\parskip}{0pt}
\emergencystretch=10em
\allowdisplaybreaks
\providecommand{\Kfield}{\mathfrak{K}}
\providecommand{\Ssys}{\mathfrak{S}}
\providecommand{\Mfield}{\mathfrak{M}}
\providecommand{\tq}{\tau}
\providecommand{\ds}{\displaystyle}
% BEGIN UNIT BODY
\section*{\S\ 1. Поля $\Kfield_{n\rho}$ и системы $\Ssys_{n\rho}$.}

В дальнейшем речь идет об исследовании систем рациональных функций от $n$ неопределенных, в особенности о полях рациональных функций.

Под $f(x_1,\ldots,x_n)$, $g(x_1,\ldots,x_n)$, \ldots\ или сокращенно $f(x)$, $g(x)$, \ldots\ поэтому всегда понимаются рациональные функции от $x_1,\ldots,x_n$; они предполагаются заданными в сокращенном виде, то есть с взаимно простыми числителем и знаменателем. Целые рациональные функции (полиномы) будут обозначаться большими буквами $F(x)$, $G(x)$, \ldots. Область коэффициентов $\Omega$ предполагается некоторым числовым полем, следовательно, в частности может содержать и все комплексные числа; $\Omega$ может также содержать еще конечное число параметров.

Величины из $\Omega$, то есть все функции нулевой степени, всегда считаются принадлежащими системе.

После этих соглашений поле из рациональных функций --- рациональное функциональное поле --- можно определить абстрактно.

\emph{Определение I: Система рациональных функций называется полем, если она удовлетворяет условиям:}
\begin{enumerate}
\item \emph{Вместе с $f(x)$ она содержит также $c\cdot f(x)$ для каждой величины $c$ из $\Omega$.}
\item \emph{Вместе с $f(x)$ и $g(x)$ она всегда содержит также $f(x)+g(x)$, $f(x)\cdot g(x)$ и, при $g(x)\ne0$, частное $f(x):g(x)$.}\footnote{При этом $f(x)$ и $g(x)$ могут быть также нулевой степени, то есть величинами из $\Omega$.}
\end{enumerate}
Специальные поля суть те, которые возникают присоединением к $\Omega$ конечного числа рациональных функций
\[
  f_1(x),\ f_2(x),\ldots, f_k(x);
\]
они будут обозначаться через
\[
  \Omega\bigl(f_1(x),\ldots,f_k(x)\bigr)
  \quad\text{или}\quad
  \Omega(f_1,\ldots,f_k).
\]
\footnote{То, что эти ``специальные поля'' уже исчерпывают все поля, следует в \S\ 4.}
Среди этих специальных полей упомянем еще поле, возникающее присоединением $x_1,\ldots,x_n$ к $\Omega$:
\[
  \Omega(x_1,\ldots,x_n),
\]
оно тождественно совокупности рациональных функций от $x_1,\ldots,x_n$ с коэффициентами из $\Omega$.

Введем еще, следуя Э. Штайницу, понятие промежуточного поля:
\begin{quote}
``Поле $\Mfield$ лежит между $\Omega_1$ и $\Omega_2$, если оно содержит $\Omega_1$ и содержится в $\Omega_2$'';
\end{quote}
тогда Определение I допускает и такую формулировку:

Поле рациональных функций от $n$ неопределенных есть промежуточное поле между $\Omega$ и $\Omega(x_1,\ldots,x_n)$, которое по меньшей мере в одной функции действительно содержит все $n$ неопределенных.

Наряду с числом неопределенных для систем рациональных функций возникает еще одно характеристическое число: число алгебраически независимых функций, или алгебраический ранг, по следующему определению.

\emph{Определение II: Система рациональных функций имеет алгебраический ранг $\rho$, если можно указать $\rho$ функций системы, являющихся алгебраически независимыми, тогда как любые $(\rho+1)$ функций системы алгебраически зависимы.}\footnote{$\rho$ функций $f_1(x),\ldots,f_\rho(x)$ называются алгебраически независимыми, если из каждого соотношения
\[
F\bigl(f_1(x),\ldots,f_\rho(x)\bigr)=0\quad\text{тождественно по }x_1,\ldots,x_n
\]
следует:
\[
F(\lambda_1,\ldots,\lambda_\rho)=0\quad\text{тождественно по }\lambda_1,\ldots,\lambda_\rho.
\]
В противоположном случае они называются алгебраически зависимыми.}

Системы из конечного или бесконечного числа рациональных функций от $n$ неопределенных и алгебраического ранга $\rho$ будут обозначаться двойным индексом
\[
  \Ssys_{n\rho}.
\]
В частности, под
\[
  \Kfield_{n\rho}
\]
понимается поле от $n$ неопределенных и алгебраического ранга $\rho$. Имеет место неравенство:
\[
  1\leq \rho\leq n.
\]

Приведем еще несколько примеров полей $\Kfield_{nn}$ и $\Kfield_{n\rho}$.

1. Полями $\Kfield_{nn}$ являются лагранжевы области родов, которые можно определить как
\begin{quote}
``совокупность рациональных (и целых рациональных) функций от $x_1,\ldots,x_n$, допускающих перестановки некоторой группы перестановок $G$ величин $x_1,\ldots,x_n$.''
\end{quote}
Они всегда содержат поле симметрических функций, а значит, и $n$ алгебраически независимых элементарных симметрических функций.

2. Полями $\Kfield_{n\rho}$ являются проективные поля инвариантов, образованные из совокупности рациональных (и целых рациональных) инвариантов одной или нескольких основных форм.\footnote{Целесообразно рассматривать только преобразования с детерминантом $1$; тогда всякая рациональная комбинация произвольного числа инвариантов снова допускает это преобразование, и действительно получается поле. Одни лишь однородные изобарные инварианты поля не образуют, но образуют систему $\Ssys_{n\rho}$.} Здесь $\rho<n$, так как число алгебраически независимых инвариантов всегда меньше числа коэффициентов.

Соответствующее утверждение справедливо для полей инвариантов относительно подгрупп проективной группы.
% END UNIT BODY
% END PRODUCER UNIT 041
\endgroup


\clearpage
\begingroup
\setcounter{footnote}{0}
% BEGIN PRODUCER UNIT 042
% Source: C:/Users/Floris/Documents/interlanguage/03_projects/noether/02_slavic_working_corpus/translations/paper06/russian/v001/Noether_Paper06_Section02_Russian_v001.tex
% Identity: 8639 B / SHA-256 24B59DC13FC0EE6E6A30EF364265091AFDFF15F99895A8E5EE13539F1B54F8F0
\setlength{\parindent}{1.25em}
\setlength{\parskip}{0pt}
\emergencystretch=10em
\allowdisplaybreaks
\providecommand{\Kfield}{\mathfrak{K}}
\providecommand{\Ssys}{\mathfrak{S}}
\providecommand{\Mfield}{\mathfrak{M}}
\providecommand{\tq}{\tau}
\providecommand{\ds}{\displaystyle}
% BEGIN UNIT BODY
\section*{\S\ 2. Вспомогательная лемма об алгебраически зависимых функциях.}

С помощью вспомогательной леммы этого параграфа в \S\ 3 выяснится, что алгебраический ранг системы является собственно характеристическим числом. А именно, лемма показывает, что при такой последовательной числовой специализации некоторых неопределенных, при которой $\rho$ алгебраически независимых функций остаются алгебраически независимыми, произвольная функция, алгебраически зависимая от этих $\rho$ функций, не может ни тождественно исчезнуть, ни стать неопределенной.

Пусть, следовательно, система
\begin{equation}
  g_1(x),\ g_2(x),\ldots,g_\rho(x)
\tag{1}
\end{equation}
имеет алгебраический ранг $\rho$, где $\rho<n$, а $h(x)$ является рациональной функцией, алгебраически зависимой от функций (1). По известным теоремам ранг функциональной матрицы системы (1) равен $\rho$. Поэтому неопределенные можно, например, так перенумеровать, что
\begin{equation}
 \frac{\partial(g_1,g_2,\ldots,g_\rho)}{\partial(x_1,x_2,\ldots,x_\rho)}
 =\frac{\Gamma(x)}{G(x)^{\rho+1}}\ne0,
\tag{2}
\end{equation}
где $G(x)$, как известно, означает общий знаменатель функций (1).

\emph{Число $a$ теперь выбирается так, чтобы произведение $G(x)\cdot\Gamma(x)$ не делилось на $(x_i-a)$}, где под $i$ понимается одно из чисел $\rho+1,\rho+2,\ldots,n$. Следовательно, при $x_i=a$ общий знаменатель функций (1) не может исчезнуть, и $\rho$ функций
\begin{equation}
  [g_1(x)]_{x_i=a},\ldots,[g_\rho(x)]_{x_i=a}
\tag{3}
\end{equation}
тоже являются алгебраически независимыми. Алгебраический ранг системы (1), как мы будем говорить, сохраняется подстановкой $x_i=a$.

Пусть неприводимое уравнение, которому $h(x)$ удовлетворяет относительно функций (1), имеет вид
\begin{equation}
\begin{aligned}
&A_0(g_1,\ldots,g_\rho)\,h(x)^\tau
 +A_1(g_1,\ldots,g_\rho)\,h(x)^{\tau-1}
 +\cdots \\
&\hspace{6em}+A_\tau(g_1,\ldots,g_\rho)=0,
\end{aligned}
\tag{4}
\end{equation}
причем $A_0(g_1,\ldots,g_\rho)$ и $A_\tau(g_1,\ldots,g_\rho)$ не исчезают тождественно. Чтобы перейти отсюда к тождеству между полиномами, положим:
\begin{equation}
  g_i(x)=G_i(x):G(x);\qquad h(x)=H_1(x):H_2(x),
\tag{5}
\end{equation}
и заметим, что вследствие предполагаемого сокращенного вида $h(x)$ функции $H_1(x)$ и $H_2(x)$ взаимно просты. В силу (5) уравнение (4) переходит в:
\begin{equation}
\begin{aligned}
&B_0(G,G_1,\ldots,G_\rho)\,G(x)^{\sigma_0}H_1(x)^\tau \\
&\quad+B_1(G,G_1,\ldots,G_\rho)\,G(x)^{\sigma_1}H_1(x)^{\tau-1}H_2(x)+\cdots \\
&\quad+B_\tau(G,G_1,\ldots,G_\rho)\,G(x)^{\sigma_\tau}H_2(x)^\tau=0,
\end{aligned}
\tag{6}
\end{equation}
где $B_k$ являются однородными формами от $G,G_i$, и по меньшей мере один из неотрицательных показателей $\sigma_k$ должен быть равен нулю.

Предположим теперь, что $H_1(x)$ делился бы на $(x_i-a)$. Тогда по (6) также
\[
  B_\tau(G,G_1,\ldots,G_\rho)\,G(x)^{\sigma_\tau}H_2(x)^\tau
\]
делился бы на $(x_i-a)$. Для $G(x)$ это исключено по предположению; значит, из делимости $B_\tau$ следовала бы также делимость
\[
  A_\tau=B_\tau:G^\lambda,
\]
и между функциями (3) существовало бы соотношение $A_\tau=0$, вопреки предполагаемой алгебраической независимости. Наконец, $H_2(x)$, как взаимно простой с $H_1(x)$, также не может делиться на $(x_i-a)$\footnote{Этот последний вывод, основанный на сокращенном виде, уже был бы невозможен при подстановке $x_i=a$, $x_k=b$ с сохранением прочих предположений. Тогда $H_1(x)$ и $H_2(x)$ могли бы одновременно исчезнуть, и частное из-за подстановки стало бы неопределенным, то есть $=\frac{0}{0}$; например,
\[
 h(x)=\frac{x_3(\alpha x_1+\beta x_2)+x_4(\gamma x_1+\delta x_2)}
 {x_3(\alpha' x_1+\beta' x_2)+x_4(\gamma' x_1+\delta' x_2)}
\]
при подстановке $x_3=0$, $x_4=0$.}; следовательно, предположение, что $H_1(x)$ делится на $(x_i-a)$, приводит к противоречию. Так же показывают, что и $H_2(x)$ не может делиться на $(x_i-a)$. Следовательно, функция $[h(x)]_{x_i=a}=h'(x)$ не может ни тождественно исчезнуть, ни стать неопределенной.

К системе (3) и функции $h'(x)$, которую снова предполагаем заданной в сокращенном виде, можно вновь применить приведенное выше рассуждение, и продолжение процедуры в конце концов приводит к следующей лемме.

\emph{Вспомогательная лемма. Пусть обе системы}
\[
  g_1(x),\ g_2(x),\ldots,g_\rho(x),
\]
\[
  g_1(x),\ g_2(x),\ldots,g_\rho(x),\ h(x)
\]
\emph{имеют алгебраический ранг $\rho$, где $\rho<n$, и пусть, например,}
\[
 \frac{\partial(g_1,\ldots,g_\rho)}{\partial(x_1,
\ldots,x_\rho)}
 =\frac{\Gamma(x)}{G(x)^{\rho+1}}\ne0.
\]
\emph{Пусть числа}
\[
  a_n,\ a_{n-1},\ldots,a_{\rho+1}
\]
\emph{выбраны так, что при подстановке}
\[
  x_n=a_n,\quad x_{n-1}=a_{n-1},\ldots,\quad x_{\rho+1}=a_{\rho+1}
\]
\emph{произведение $G(x)\cdot\Gamma(x)$ не исчезает, то есть алгебраический ранг системы (1) сохраняется. В последовательных подстановках}
\[
\begin{aligned}
 [h(x)]_{x_n=a_n}&=h^{(n-1)}(x),\qquad
 [h^{(n-1)}(x)]_{x_{n-1}=a_{n-1}}=h^{(n-2)}(x),\ldots,\\
 [h^{(\rho+1)}(x)]_{x_{\rho+1}=a_{\rho+1}}&=h^*(x_1,x_2,\ldots,x_\rho)
\end{aligned}
\]
\emph{пусть функции $h(x)$, $h^{(n-1)}(x)$, \ldots, $h^{(\rho+1)}(x)$ каждый раз приведены к сокращенному виду. При этих предположениях в функции $h^*(x_1,\ldots,x_\rho)$ ни числитель, ни знаменатель не могут тождественно исчезнуть, и сама функция также не может стать $\frac{0}{0}$.}\footnote{Последовательная подстановка дает, например, в функции $h(x)$ из предыдущего примечания:
\[
 [h(x)]_{x_4=0}=h^{(3)}(x)=\frac{\alpha x_1+\beta x_2}{\alpha' x_1+\beta' x_2},\qquad
 h^*(x_1,x_2)=\frac{\alpha x_1+\beta x_2}{\alpha' x_1+\beta' x_2}.
\]}
% END UNIT BODY
% END PRODUCER UNIT 042
\endgroup


\clearpage
\begingroup
\setcounter{footnote}{0}
% BEGIN PRODUCER UNIT 043
% Source: C:/Users/Floris/Documents/interlanguage/03_projects/noether/02_slavic_working_corpus/translations/paper06/russian/v001/Noether_Paper06_Section03_Russian_v001.tex
% Identity: 7688 B / SHA-256 D926AF2130844C30796807540AD685BE7B291F593336A90FD71831807F48B56A
\setlength{\parindent}{1.25em}
\setlength{\parskip}{0pt}
\emergencystretch=10em
\allowdisplaybreaks
\providecommand{\Kfield}{\mathfrak{K}}
\providecommand{\Ssys}{\mathfrak{S}}
\providecommand{\Mfield}{\mathfrak{M}}
\providecommand{\tq}{\tau}
\providecommand{\ds}{\displaystyle}
% BEGIN UNIT BODY
\section*{\S\ 3. Взаимно однозначное отображение систем $\Ssys_{n\rho}$ на системы $\Ssys^*_{\rho\rho}$.}

Из вспомогательной леммы мы непосредственно получим отображение систем $\Ssys_{n\rho}$ на системы $\Ssys^*_{\rho\rho}$ благодаря тому, что некоторые неопределенные заменяются числами; алгебраический ранг $\rho$ тем самым оказывается собственно характеристическим числом.

Так как $\Ssys_{n\rho}$ имеет алгебраический ранг $\rho$, существуют $\rho$ функций из $\Ssys_{n\rho}$,
\begin{equation}
  g_1(x),\ldots,g_\rho(x),
\tag{1}
\end{equation}
которые алгебраически независимы. Далее, если $f(x)$ обозначает произвольную функцию из $\Ssys_{n\rho}$, то система
\begin{equation}
  g_1(x),\ldots,g_\rho(x),\ f(x)
\tag{2}
\end{equation}
также имеет алгебраический ранг $\rho$; следовательно, системы (1) и (2) удовлетворяют условиям вспомогательной леммы.\footnote{При необходимости неопределенные следует перенумеровать.}

Поэтому определим числа
\[
  a_n,\ a_{n-1},\ldots,a_{\rho+1}
\]
по вспомогательной лемме так, чтобы алгебраический ранг $\rho$ системы (1) сохранялся при подстановке --- что возможно сколь угодно многими способами. Тогда в функции, определенной посредством
\[
\begin{aligned}
 [f(x)]_{x_n=a_n}&=f^{(n-1)}(x),\qquad
 [f^{(n-1)}(x)]_{x_{n-1}=a_{n-1}}=f^{(n-2)}(x);\ldots,\\
 [f^{(\rho+1)}(x)]_{x_{\rho+1}=a_{\rho+1}}&=f^*(x_1,
\ldots,x_\rho),
\end{aligned}
\]
а именно в функции
\begin{equation}
  f^*(x_1,\ldots,x_\rho),
\tag{3}
\end{equation}
ни числитель, ни знаменатель не могут тождественно исчезнуть.

Пусть теперь $f(x)$ пробегает все (конечно или бесконечно многие) функции из $\Ssys_{n\rho}$. Рассмотрим систему, состоящую из совокупности всех функций (3); она обладает следующими свойствами:

1. Она имеет алгебраический ранг $\rho$, ибо содержит функции, возникающие из (1),
\[
  g_1^*(x_1,
\ldots,x_\rho),\ldots,g_\rho^*(x_1,\ldots,x_\rho),
\]
которые вследствие выбора $a_n,a_{n-1},\ldots,a_{\rho+1}$ алгебраически независимы. Поэтому эту систему будем обозначать через
\[
  \Ssys^*_{\rho\rho}.
\]
\footnote{Под $f^*(x)$, $g^*(x)$, \ldots\ соответственно всюду следует понимать функции отображения, то есть функции от $x_1,\ldots,x_\rho$.}

2. Сопоставление функций $f(x)$ и $f^*(x)$ без исключений взаимно однозначно.

Действительно, если $f_1(x)$ и $f_2(x)$ принадлежат $\Ssys_{n\rho}$, то и разность
\[
  d(x)=f_1(x)-f_2(x)
\]
алгебраически зависит от системы (1); кроме того, имеет место:
\[
  d^*(x)=f_1^*(x)-f_2^*(x).
\]
Так как $d(x)$ вместе с системой (1) удовлетворяет условиям вспомогательной леммы, то из
\[
  d(x)\ne0
\]
необходимо следует:
\[
  d^*(x)\ne0;
\]
то есть различным функциям из $\Ssys_{n\rho}$ соответствуют различные функции из $\Ssys^*_{\rho\rho}$.

3. Из всякого соотношения между функциями из $\Ssys^*_{\rho\rho}$:
\[
  F\bigl(f_1^*(x),f_2^*(x),\ldots,f_k^*(x)\bigr)=0
  \quad\text{тождественно по }x_1,\ldots,x_\rho
\]
для взаимно однозначно соответствующих функций из $\Ssys_{n\rho}$ следует:
\[
  F\bigl(f_1(x),f_2(x),\ldots,f_k(x)\bigr)=0
  \quad\text{тождественно по }x_1,
\ldots,x_n.
\]
Действительно, вместе с $f_1(x),f_2(x),\ldots,f_k(x)$ всякая рациональная комбинация этих функций также алгебраически зависит от системы (1).

Далее из
\begin{equation}
\begin{aligned}
 h(x)&=F\bigl(f_1(x),\ldots,f_k(x)\bigr):\\
 h^*(x)&=F\bigl(f_1^*(x),\ldots,f_k^*(x)\bigr)
\end{aligned}
\tag{4}
\end{equation}
по вспомогательной лемме следует, что из
\[
  h(x)\ne0\quad\text{также}\quad h^*(x)\ne0;
\]
q. e. d.

Итак, получаем следующую теорему.

{\itshape Теорема I. Каждой системе $\Ssys_{n\rho}$ из (конечно или бесконечно многих) рациональных функций можно --- сколь угодно многими способами --- взаимно однозначно сопоставить систему $\Ssys^*_{\rho\rho}$ так, что все соотношения, существующие между функциями из $\Ssys^*_{\rho\rho}$, сохраняются также в $\Ssys_{n\rho}$, и наоборот. При этом полю $\Kfield_{n\rho}$ по (4) в частности снова соответствует поле $\Kfield^*_{\rho\rho}$.}

Из теоремы I мы выводим еще следствие, упрощающее все дальнейшие доказательства: {\itshape свойства систем $\Ssys_{n\rho}$, которые характеризуются конечным или бесконечным числом соотношений между функциями системы, справедливы для систем $\Ssys_{n\rho}$, как только они справедливы для системы отображения $\Ssys^*_{\rho\rho}$.}

Это следствие мы кратко называем ``принципом переноса''.

\providecommand{\Lsys}{\mathfrak{L}}
\providecommand{\Jdom}{\mathfrak{J}}

\clearpage
Простейший пример этого отображения дает теория алгебраических инвариантов. Действительно, линейным преобразованием всегда можно численно зафиксировать некоторые коэффициенты основной формы так, чтобы все соотношения, справедливые для специальной формы, сохранялись и в общем случае.
% END UNIT BODY
% END PRODUCER UNIT 043
\endgroup


\clearpage
\begingroup
\setcounter{footnote}{0}
% BEGIN PRODUCER UNIT 044
% Source: C:/Users/Floris/Documents/interlanguage/03_projects/noether/02_slavic_working_corpus/translations/paper06/russian/v001/Noether_Paper06_Section04_Russian_v001.tex
% Identity: 8270 B / SHA-256 5B73436660878432D076614B094A12C4B8FCC736776024F9916985DE8BADA7F8
\setlength{\parindent}{1.25em}
\setlength{\parskip}{0pt}
\emergencystretch=10em
\allowdisplaybreaks
\providecommand{\Kfield}{\mathfrak{K}}
\providecommand{\Ssys}{\mathfrak{S}}
\providecommand{\Mfield}{\mathfrak{M}}
\providecommand{\tq}{\tau}
\providecommand{\ds}{\displaystyle}
% BEGIN UNIT BODY
\section*{\S\ 4. Рациональная база полей $\Kfield_{n\rho}$.}

В дальнейших исследованиях, сгруппированных вокруг понятий базы, мы всюду пользуемся «принципом переноса» из \S\ 3. Сначала мы доказываем существование базы для полей, но определения даем в общем виде:

\emph{Определение III: Конечное число функций из $\Ssys_{n\rho}$ называется рациональной базой, если каждую функцию из $\Ssys_{n\rho}$ можно представить как рациональную комбинацию этого конечного числа функций с коэффициентами из $\Omega$.}

Рациональная база должна содержать $\rho$ алгебраически независимых функций, поскольку алгебраический ранг системы, полученной из рациональной базы, равен алгебраическому рангу самой рациональной базы.

Сначала покажем, что доказательство существования рациональной базы допускает принцип переноса. А именно, Определение III можно сформулировать так:

Функции из $\Ssys_{n\rho}$
\[
  g_1(x),\ g_2(x),\ldots, g_k(x)
\]
тогда и только тогда образуют рациональную базу, когда выполнены бесконечно многие соотношения:
\begin{equation}
  f(x)=\gamma\bigl(g_1(x),\ldots,g_k(x)\bigr),
\tag{1}
\end{equation}
где $f(x)$ пробегает все функции из $\Ssys_{n\rho}$, а $\gamma$ обозначает рациональную функцию от $k$ аргументов, меняющуюся вместе с $f$.

Следовательно, если в системе отображения $\Ssys^*_{\rho\rho}$ рациональная база задана функциями
\[
  g_1^*(x),\ g_2^*(x),\ldots, g_k^*(x),
\]
что влечет соотношения
\[
  f^*(x)=\gamma\bigl(g_1^*(x),\ldots,g_k^*(x)\bigr),
\]
то по Теореме I для $\Ssys_{n\rho}$ выполнены соотношения (1). Иными словами, \emph{принцип переноса для рациональной базы} гласит:

\emph{Из существования рациональной базы для системы отображения $\Ssys^*_{\rho\rho}$ следует существование рациональной базы для исходной системы $\Ssys_{n\rho}$.}\footnote{Соответственно формулируется принцип переноса для каждого вопроса о базе, рассматриваемого здесь.}

Чтобы доказать существование рациональной базы всех полей $\Kfield_{n\rho}$, мы, таким образом, можем ограничиться полями $\Kfield^*_{\rho\rho}$; здесь к цели ведет прямое обобщение рассуждения, примененного Э. Штайницем.\footnote{E. Steinitz, Algebraische Theorie der Körper, \S\ 24, Crelles Journ. 137 (1909).}

Пусть
\begin{equation}
  g_1^*(x_1,\ldots,x_\rho),\ldots,
  g_\rho^*(x_1,\ldots,x_\rho)
\tag{2}
\end{equation}
есть система из $\rho$ алгебраически независимых функций из $\Kfield^*_{\rho\rho}$. Элиминация $x_1,\ldots,x_\rho$ дает $\rho$ соотношений:
\begin{equation}
\begin{aligned}
  F_1\bigl(g_1^*(x),\ldots,g_\rho^*(x),x_1\bigr)&=0,\ldots,\\
  F_\rho\bigl(g_1^*(x),\ldots,g_\rho^*(x),x_\rho\bigr)&=0,
\end{aligned}
\tag{3}
\end{equation}
причем в каждом соотношении $F_i=0$ неопределенная $x_i$ действительно входит; иначе, вопреки предположению, между функциями (2) существовала бы алгебраическая зависимость. Соотношения (3) означают, что
\[
  x_1,x_2,\ldots,x_\rho
\]
являются алгебраическими элементами относительно
\[
  \Omega\bigl(g_1^*(x),\ldots,g_\rho^*(x)\bigr).
\]
Поэтому поле, возникающее присоединением $x_1,\ldots,x_\rho$,
\[
  \Omega\bigl(g_1^*,\ldots,g_\rho^*,x_1,
  \ldots,x_\rho\bigr)=\Omega(x_1,\ldots,x_\rho),
\]
есть конечное алгебраическое расширение над $\Omega(g_1^*,\ldots,g_\rho^*)$. По известной теореме\footnote{См., например, Weber, Lehrbuch d. Algebra 1, \S\ 144 (1-е изд.) или \S\ 55 (малое изд.).} $\Kfield^*_{\rho\rho}$ как промежуточное поле между $\Omega(g_1^*,\ldots,g_\rho^*)$ и $\Omega(x_1,
\ldots,x_\rho)$ само является алгебраическим расширением над $\Omega(g_1^*,\ldots,g_\rho^*)$, тогда как $\Omega(x_1,
\ldots,x_\rho)$ является алгебраическим расширением над $\Kfield^*_{\rho\rho}$. Значит, $\Kfield^*_{\rho\rho}$ получается из $\Omega(g_1^*,\ldots,g_\rho^*)$ присоединением конечного числа элементов из $\Kfield^*_{\rho\rho}$, или также присоединением \emph{одной} подходяще выбранной функции; то есть
\[
  \Kfield^*_{\rho\rho}
  =\Omega(g_1^*,\ldots,g_\rho^*,g_{\rho+1}^*,\ldots,g_k^*)
  =\Omega(g_1^*,\ldots,g_\rho^*,g_0^*).
\]

По принципу переноса мы тем самым получаем

\emph{Теорема II: Каждое поле $\Kfield_{n\rho}$ обладает рациональной базой алгебраического ранга $\rho$; в частности, рациональную базу можно выбрать так, чтобы она содержала не более $\rho+1$ функций.}

Пусть дана произвольная система из $\rho$ алгебраически независимых функций из $\Kfield_{n\rho}$:
$h_1(x),\ldots,h_\rho(x)$. По приведенному выше методу ее можно расширить до рациональной базы
\[
  h_1(x),\ldots,h_\rho(x);
  h_{\rho+1}(x),\ldots,h_\sigma(x).
\]
По определяющему свойству существуют соотношения:
\[
\begin{aligned}
  h_i(x)&=\chi_i\bigl(g_1(x),\ldots,g_k(x)\bigr)
    &&(i=1,2,\ldots,\sigma),\\
  g_j(x)&=\psi_j\bigl(h_1(x),\ldots,h_\sigma(x)\bigr)
    &&(j=1,2,\ldots,k),
\end{aligned}
\]
где $\chi_i$, $\psi_j$ обозначают рациональные функции своих аргументов. То есть:

Из любой рациональной базы всякая другая рациональная база получается бирациональным преобразованием. Кроме того:

если
\[
  h_1(x),\ldots,h_\rho(x)
\]
есть произвольная система из $\rho$ алгебраически независимых функций из $\Kfield_{n\rho}$, то $\Kfield_{n\rho}$ является конечным алгебраическим расширением над
\[
  \Omega\bigl(h_1(x),\ldots,h_\rho(x)\bigr).
\]
% END UNIT BODY
% END PRODUCER UNIT 044
\endgroup


\clearpage
\begingroup
\setcounter{footnote}{0}
% BEGIN PRODUCER UNIT 045
% Source: C:/Users/Floris/Documents/interlanguage/03_projects/noether/02_slavic_working_corpus/translations/paper06/russian/v001/Noether_Paper06_Section05_Russian_v001.tex
% Identity: 12219 B / SHA-256 4F940F1642422F46351BCA90761E6AD537F3431C137E549C8CB7C4D9286CB882
\setlength{\parindent}{1.25em}
\setlength{\parskip}{0pt}
\emergencystretch=10em
\allowdisplaybreaks
\providecommand{\Kfield}{\mathfrak{K}}
\providecommand{\Ssys}{\mathfrak{S}}
\providecommand{\Mfield}{\mathfrak{M}}
\providecommand{\tq}{\tau}
\providecommand{\ds}{\displaystyle}
% BEGIN UNIT BODY
\section*{\S\ 5. Инволюционная форма и инволюционная база полей $\Kfield_{n\rho}$.}

В то время как каждая рациональная база, построенная в \S\ 4, имела тот недостаток, что зависела от произвольно выбранных исходных функций, теперь мы показываем существование \emph{рациональной базы, однозначно определенной полем}.

По \S\ 4 поле $\Omega(x_1,\ldots,x_\rho)$ является алгебраическим полем -- скажем, степени $i$ -- над $\Kfield^*_{\rho\rho}$ и, ввиду
\[
  \Kfield^*_{\rho\rho}(x_1,
  \ldots,x_\rho)=\Omega(x_1,\ldots,x_\rho),
\]
возникает присоединением элементов $x_1,\ldots,x_\rho$, алгебраических относительно $\Kfield^*_{\rho\rho}$. Поэтому рассмотрим неприводимую относительно $\Kfield^*_{\rho\rho}$ целую функцию $\Phi^*(z,u)$ с нулем
\[
  z=u_1x_1+u_2x_2+\cdots+u_\rho x_\rho=u(x),
\]
имеющую степень $i$ по $z$. Кроме того, как произведение величин, сопряженных с $[z-u(x)]$, $\Phi^*(z,u)$ является однородной формой степени $i$ по $z,u_1,\ldots,u_\rho$; следовательно, коэффициентом при $z^i$ она имеет не зависящую от $u$ функцию из $\Kfield^*_{\rho\rho}$. Сделав этот старший коэффициент равным $1$ делением, мы однозначно определяем неприводимую функцию $\Phi^*(z,u)$. Определенная таким образом однородная форма
\begin{equation}
\begin{aligned}
  \Phi^*(z,u)
  &=z^i+
    \sum{}' g^*_{\alpha\alpha_1\ldots\alpha_\rho}(x_1,
    \ldots,x_\rho)
    z^\alpha u_1^{\alpha_1}\cdots u_\rho^{\alpha_\rho},\\
  &\hspace{8em}(\alpha+\alpha_1+\cdots+\alpha_\rho=i)
\end{aligned}
\tag{1}
\end{equation}
будет называться \emph{инволюционной формой}\footnote{Основание для этого названия выясняется в конце параграфа; $\sum'$ означает, что суммирование распространяется на все члены, кроме $z^i$.} поля $\Kfield^*_{\rho\rho}$; из (1) получается
\begin{equation}
\begin{aligned}
  \Phi(z,u)
  &=z^i+
    \sum{}' g_{\alpha\alpha_1\ldots\alpha_\rho}(x)
    z^\alpha u_1^{\alpha_1}\cdots u_\rho^{\alpha_\rho},\\
  &\hspace{8em}(\alpha+\alpha_1+\cdots+\alpha_\rho=i)
\end{aligned}
\tag{2}
\end{equation}
как \emph{инволюционная форма} поля $\Kfield_{n\rho}$.

Коэффициенты инволюционной формы $\Phi^*(z,u)$ окажутся рациональной базой поля $\Kfield^*_{\rho\rho}$; то есть имеет место соотношение
\begin{equation}
  \Omega\bigl(g^*_{\alpha\alpha_1\ldots\alpha_\rho}(x_1,
  \ldots,x_\rho)\bigr)=\Kfield^*_{\rho\rho}
  \quad(\alpha+\alpha_1+\cdots+\alpha_\rho=i).
\tag{3}
\end{equation}
Для доказательства заметим, что $\Phi^*(z,u)$ имеет коэффициенты из $\Omega(g^*_{\alpha\alpha_1\ldots\alpha_\rho}(x))$ и, вследствие неприводимости относительно $\Kfield^*_{\rho\rho}$, тем более неприводима относительно этого подполя. Значит, $\Phi^*(z,u)$ задает неприводимое уравнение для $u(x)$ относительно $\Omega(g^*_{\alpha\alpha_1\ldots\alpha_\rho}(x))$. Отсюда следует, что $\Omega(x_1,\ldots,x_\rho)$ является алгебраическим полем -- степени $i$ -- над $\Omega(g^*_{\alpha\alpha_1\ldots\alpha_\rho}(x))$. Поскольку $\Kfield^*_{\rho\rho}$ лежит между $\Omega(g^*_{\alpha\alpha_1\ldots\alpha_\rho}(x))$ и $\Omega(x_1,\ldots,x_\rho)$, из равенства степеней, как известно, следует равенство полей, а значит соотношение (3). По принципу переноса коэффициенты $\Phi(z,u)$ соответственно дают рациональную базу поля $\Kfield_{n\rho}$; то есть мы имеем

\emph{Теорему III: Коэффициенты инволюционной формы образуют определенную полем рациональную базу, инволюционную базу; для полей $\Kfield^*_{\rho\rho}$ эта база определена однозначно.}\footnote{Для полей $\Kfield_{n\rho}$ инволюционная форма может зависеть от отображения; однозначности можно достичь отображением с неопределенными коэффициентами.}

Добавим еще иррациональное толкование этого факта, устанавливающее связь с геометрической теорией инволюции.

Пусть разложение $\Phi^*(z,u)$ на линейные множители в некотором поле расширения задано так:
\begin{equation}
\begin{aligned}
 \Phi^*(z,u)&=(z-u_1x_1-\cdots-u_\rho x_\rho)
 (z-u_1x_1^{(1)}-\cdots-u_\rho x_\rho^{(1)})\cdots\\
 &\qquad\cdot
 (z-u_1x_1^{(i-1)}-\cdots-u_\rho x_\rho^{(i-1)}).
\end{aligned}
\tag{4}
\end{equation}
Сравнение с (1) показывает, что функции
\[
  g^*_{\alpha\alpha_1\ldots\alpha_\rho}(x_1,
  \ldots,x_\rho)\quad(\alpha+\alpha_1+\cdots+\alpha_\rho=i)
\]
тождественны элементарным симметрическим функциям рядов величин:
\begin{equation}
\begin{array}{cccc}
  x_1 & x_2 & \cdots & x_\rho\\
  x_1^{(1)} & x_2^{(1)} & \cdots & x_\rho^{(1)}\\
  \vdots & \vdots & & \vdots\\
  x_1^{(i-1)} & x_2^{(i-1)} & \cdots & x_\rho^{(i-1)}.
\end{array}
\tag{5}
\end{equation}
Следовательно, каждая функция из $\Kfield^*_{\rho\rho}$ является рациональной комбинацией этих элементарных функций, симметричной по рядам (5), и наоборот, каждая симметрическая функция рядов (5) принадлежит $\Kfield^*_{\rho\rho}$.

Итак, когда
\[
  f^*(x)=\frac{F^*(x)}{H^*(x)}
\]
пробегает все функции из $\Kfield^*_{\rho\rho}$, имеет место система бесконечно многих соотношений:
\begin{equation}
\begin{aligned}
  \frac{F^*(x)}{H^*(x)}
  &=\frac{F^*(x^{(1)})}{H^*(x^{(1)})}=\cdots\\
  &=\frac{F^*(x^{(i-1)})}{H^*(x^{(i-1)})},
\end{aligned}
\tag{6}
\end{equation}
или, в свернутой форме:
\[
  F^*(x^{(k)})H^*(x^{(l)})-
  F^*(x^{(l)})H^*(x^{(k)})=0
  \quad(k,l=0,1,\ldots,i-1),
\]
где ни $F^*(x^{(k)})$, ни $H^*(x^{(k)})$ -- как сопряженные с $F^*(x)$, $H^*(x)$ и формально симметрические -- не обращаются тождественно в нуль.

Обратно, для каждого ряда величин $\xi$, удовлетворяющего бесконечно многим соотношениям
\begin{equation}
\begin{aligned}
  F^*(x)H^*(\xi)-H^*(x)F^*(\xi)&=0,\\
  H^*(\xi)&\ne0,
\end{aligned}
\tag{7}
\end{equation}
получается принадлежность к рядам величин (5).

В самом деле, если $G^*(x)$ обозначает общий знаменатель коэффициентов $\Phi^*(z,u)$, то по предположению (7) целая также по $\xi$ функция
\[
  G^*(\xi)\cdot\Phi^*(z,u;\xi)
\]
не обращается тождественно в нуль -- то есть коэффициенты $\Phi^*(z,u;\xi)$ не становятся неопределенными -- и имеет нуль
\[
  z=u_1\xi_1+u_2\xi_2+\cdots+u_\rho\xi_\rho.
\]
Поэтому в силу (7) и $\Phi^*(z,u;x)$ имеет нуль $z=u(\xi)$; то есть $\xi$ принадлежит рядам величин (5). Соответствующее утверждение следует из (6), если $\xi$ удовлетворяет системе уравнений относительно сопряженного ряда величин:
\[
  F^*(x^{(k)})H^*(\xi)-H^*(x^{(k)})F^*(\xi)=0;
  \qquad H^*(\xi)\ne0,
\]
то есть ряды величин (5) определяются как решения $\xi$ системы уравнений
\[
  F^*(x)H^*(\xi)-H^*(x)F^*(\xi)=0;
  \qquad H^*(\xi)\ne0,
\]
где $F^*(x):H^*(x)$ пробегает все функции из $\Kfield^*_{\rho\rho}$; при этом ряд $x$ можно заменить любым сопряженным рядом.

Геометрически это означает следующее, если понимать каждый ряд $x$ как точку: решения (7) задают в $\rho$-мерном линейном пространстве $\infty^\rho$ группы по $i$ точек так, что любая точка группы однозначно определяет всю отдельную группу. Такая система групп точек называется «\emph{инволюцией в линейном пространстве}».\footnote{Исключенные решения (7), для которых $F^*(\xi)=0$, $H^*(\xi)=0$, -- а также исключенные решения системы уравнений, соответствующей инволюционной базе, -- называются фундаментальными точками инволюции; при этом следует учитывать и те, которые возникают при гомогенизации, то есть «лежащие на бесконечности».}

Соответственно, поле $\Kfield_{n\rho}$ характеризует инволюцию в линейном пространстве, состоящую из кривых, поверхностей и т. д.

Так как, по сказанному выше, степень $i$ поля $\Omega(x_1,\ldots,x_\rho)$ относительно $\Kfield^*_{\rho\rho}$ равна числу систем решений (7) с дополнительным условием $H^*(\xi)\ne0$ -- которые мы можем назвать системами решений поля $\Kfield^*_{\rho\rho}$, -- то рассуждение, использованное для доказательства существования инволюционной базы, допускает также иррациональную формулировку:

«Если $\mathfrak{L}^*$ является делителем $\Kfield^*_{\rho\rho}$ и каждая система решений поля $\mathfrak{L}^*$ является также системой решений $\Kfield^*_{\rho\rho}$, то $\mathfrak{L}^*$ и $\Kfield^*_{\rho\rho}$ тождественны».

Соответствующее утверждение по принципу переноса справедливо для делителей $\mathfrak{L}$ поля $\Kfield_{n\rho}$.
% END UNIT BODY
% END PRODUCER UNIT 045
\endgroup


\clearpage
\begingroup
\setcounter{footnote}{0}
% BEGIN PRODUCER UNIT 046
% Source: C:/Users/Floris/Documents/interlanguage/03_projects/noether/02_slavic_working_corpus/translations/paper06/russian/v001/Noether_Paper06_Section06_Russian_v001.tex
% Identity: 4955 B / SHA-256 27C829FF0BE7858313842302AC3BD764C643A98A5F1A3016D80AA476C13210F3
\setlength{\parindent}{1.25em}
\setlength{\parskip}{0pt}
\emergencystretch=10em
\allowdisplaybreaks
\providecommand{\Kfield}{\mathfrak{K}}
\providecommand{\Ssys}{\mathfrak{S}}
\providecommand{\Mfield}{\mathfrak{M}}
\providecommand{\tq}{\tau}
\providecommand{\ds}{\displaystyle}
% BEGIN UNIT BODY
\section*{\S\ 6. Минимальная база полей $\Kfield_{n\rho}$.}

Этот параграф дает обзор известных теорем, но в расширенной форме, которую делают возможной принцип переноса и существование рациональной базы.

\emph{Определение IV: Рациональная база поля $\Kfield_{n\rho}$ называется минимальной базой, если она содержит только минимальное число $\rho$ функций; тогда эти функции должны быть алгебраически независимыми.}

Из теорем Люрота, Кастельнуово и Энрикеса\footnote{J. Lüroth: Beweis eines Satzes über rationale Kurven, Math. Ann. 9 (1875). --- G. Castelnuovo: Sulla razionalità delle involuzioni piane, Math. Ann. 44 (1893). --- F. Enriques: Sopra una involuzione non razionale dello spazio, Rend. Acc. Linc, Vol. 21, 21. Jan. 1912.} следуют такие факты:

I. Каждое поле $\Kfield_{n1}$ имеет минимальную базу: это функция Люрота данного поля, однозначно определенная с точностью до дробно-линейного преобразования.

II. Для полей $\Kfield_{n2}$ минимальная база существует всегда тогда и, вообще говоря, только тогда, когда допускается алгебраическое расширение поля коэффициентов $\Omega$.

III. Для полей $\Kfield_{n\rho}$ ($\rho\ge3$) вообще говоря минимальной базы не существует. Характеризация специальных полей $\Kfield_{n\rho}$, имеющих минимальную базу, еще не была предпринята.

Кратко напомним доказательство теоремы Люрота для полей $\Kfield^*_{11}$, данное Э. Штайницем.\footnote{A. a. O. \S\ 24.}

Пусть $\Phi^*(z,u)$ является инволюционной формой поля $\Kfield^*_{11}$, а $\psi^*(x)$ -- каким-либо коэффициентом $\Phi^*(z,u)$, действительно содержащим неопределенную $x$. Оказывается, что степень $\Omega(x)$ относительно $\Omega(\psi^*(x))$ равна степени $\Omega(x)$ относительно $\Kfield^*_{11}$, откуда следует тождество полей. Далее, если $\Omega(\chi^*(x))$ равна $\Omega(\psi^*(x))$, то из взаимной рациональной зависимости $\chi^*(x)$ и $\psi^*(x)$ следует дробно-линейная зависимость этих двух функций.

Следовательно, по принципу переноса теорема Люрота допускает также такую формулировку:

\emph{Минимальная база полей $\Kfield_{n1}$ состоит из какой-либо функции инволюционной базы; и любые две функции инволюционной базы поля $\Kfield_{n1}$ дробно-линейно зависят друг от друга.}

Доказательство факта II Кастельнуово проводит с помощью геометрической теории инволюции для полей $\Kfield^*_{22}$, выведенных из рациональной базы из трех функций. Поскольку принцип переноса сохраняется также при конечном алгебраическом расширении поля коэффициентов $\Omega$, существование рациональной базы тем самым дает доказательство для всех полей $\Kfield_{n2}$.

Относительно III следует заметить, что Энрикес строит нерациональную инволюцию в трехмерном пространстве, то есть поле $\Kfield_{33}$, не имеющее минимальной базы.
% END UNIT BODY
% END PRODUCER UNIT 046
\endgroup


\clearpage
\begingroup
\setcounter{footnote}{0}
% BEGIN PRODUCER UNIT 047
% Source: C:/Users/Floris/Documents/interlanguage/03_projects/noether/02_slavic_working_corpus/translations/paper06/russian/v001/Noether_Paper06_Section07_Russian_v001.tex
% Identity: 10728 B / SHA-256 F04A6B97D8E040EF37BCEF1771A6F601BDBFB8DABB92333A0160E6C71256B18C
\setlength{\parindent}{1.25em}
\setlength{\parskip}{0pt}
\emergencystretch=10em
\allowdisplaybreaks
\providecommand{\Kfield}{\mathfrak{K}}
\providecommand{\Ssys}{\mathfrak{S}}
\providecommand{\Jdom}{\mathfrak{J}}
\providecommand{\Lsys}{\mathfrak{L}}
\providecommand{\Mfield}{\mathfrak{M}}
\providecommand{\tq}{\tau}
\providecommand{\ds}{\displaystyle}
% BEGIN UNIT BODY
\section*{\S\ 7. Произвольная система $\Ssys_{n\rho}$, линейное семейство $\Lsys_{n\rho}$ и область целостности $\Jdom_{n\rho}$ рациональных функций. Существование рациональной базы.}

Из существования рациональной базы полей $\Kfield_{n\rho}$ теперь будет выведена рациональная база для произвольных систем рациональных и целых рациональных функций. Мы упорядочим эти системы так, как они последовательно возникают одна из другой посредством элементарных операций сложения, умножения и деления.

I. Пусть, как всегда, $\Ssys_{n\rho}$ обозначает произвольную систему рациональных (или целых рациональных) функций от $n$ неопределённых, алгебраического ранга $\rho$. Элементы из $\Omega$ также считаются принадлежащими системе.

II. Система $\Lsys_{n\rho}$ называется \emph{линейным семейством}, если она удовлетворяет условиям:

1. вместе с $f(x)$ система $\Lsys_{n\rho}$ всегда содержит также $c\cdot f(x)$, где $c$ -- произвольный элемент из $\Omega$;

2. вместе с $f(x)$ и $g(x)$ система $\Lsys_{n\rho}$ всегда содержит также $f(x)+g(x)$.

III. Линейное семейство $\Jdom_{n\rho}$ называется \emph{областью целостности} при дополнительном условии:

3. вместе с $f(x)$ и $g(x)$ область $\Jdom_{n\rho}$ всегда содержит также $f(x)\cdot g(x)$.

IV. Область целостности $\Kfield_{n\rho}$ называется \emph{полем} при дополнительном условии:

4. вместе с $f(x)$ и $g(x)$ --- где $g(x)\ne0$ --- поле $\Kfield_{n\rho}$ всегда содержит также частное $f(x):g(x)$.

Условия 1--4 являются теми же условиями, которые в \S\ 1 были заданы для поля.\footnote{$f(x)$ и $g(x)$ могут также иметь нулевую степень, то есть быть элементами из $\Omega$.}

Теперь эти области должны быть последовательно выведены одна из другой.

a) Произвольная система $\Ssys_{n\rho}$ расширяется до линейного семейства $\Lsys_{n\rho}$ следующим требованием:

$\Lsys_{n\rho}$ содержит, кроме $\Ssys_{n\rho}$, все функции $h(x)$ и только те, которые можно линейно представить с коэффициентами из $\Omega$ через конечное число функций из $\Ssys_{n\rho}$:
\[
  h(x)=c_1f_1(x)+c_2f_2(x)+\cdots+c_kf_k(x),
\]
где $f_1(x),\ldots,f_k(x)$ пробегают все функции из $\Ssys_{n\rho}$, а $c_1,\ldots,c_k$ пробегают все элементы из $\Omega$.

Действительно, добавление рациональных комбинаций функций системы не меняет алгебраического ранга системы. $\Lsys_{n\rho}$ удовлетворяет условиям 1 и 2, следовательно, является линейным семейством. Кроме того, всякое линейное семейство, содержащее $\Ssys_{n\rho}$, по условиям 1 и 2 содержит также $\Lsys_{n\rho}$; поэтому $\Lsys_{n\rho}$ есть наименьшее линейное семейство, содержащее $\Ssys_{n\rho}$.

b) Соответственно, из линейного семейства $\Lsys_{n\rho}$ наименьшая содержащая его область целостности $\Jdom_{n\rho}$ возникает по требованию:

$\Jdom_{n\rho}$ содержит, кроме $\Lsys_{n\rho}$, все функции $h(x)$ и только те, которые можно представить как \emph{целую рациональную комбинацию} --- с коэффициентами из $\Omega$ --- конечного числа функций из $\Lsys_{n\rho}$, $f_1(x),\ldots,f_k(x)$, где $f_1(x),\ldots,f_k(x)$ пробегают все функции из $\Lsys_{n\rho}$.

c) Наконец, из области целостности $\Jdom_{n\rho}$ наименьшее содержащее её поле $\Kfield_{n\rho}$ возникает по требованию: $\Kfield_{n\rho}$ содержит, кроме $\Jdom_{n\rho}$, все функции $h(x)$ и только те, которые можно представить как частные двух функций из $\Jdom_{n\rho}$.

Если объединить эти три расширения в один шаг, получается:

\emph{Всякую систему $\Ssys_{n\rho}$ можно расширить до наименьшего содержащего её поля $\Kfield_{n\rho}$ следующим требованием:}

\emph{$\Kfield_{n\rho}$ содержит, кроме $\Ssys_{n\rho}$, все функции $h(x)$ и только те, которые можно представить как рациональные комбинации --- с коэффициентами из $\Omega$ --- конечного числа функций из $\Ssys_{n\rho}$, $f_1(x),\ldots,f_k(x)$, где $f_1(x),\ldots,f_k(x)$ пробегают все функции из $\Ssys_{n\rho}$.}

Из последнего факта непосредственно следует существование рациональной базы для произвольных областей:

Пусть $\Kfield_{n\rho}$ есть наименьшее поле, содержащее $\Ssys_{n\rho}$, и пусть рациональная база $\Kfield_{n\rho}$ задана функциями
\begin{equation}
  h_1(x),\ h_2(x),\ldots,h_\sigma(x).
\tag{1}
\end{equation}
По определению $\Kfield_{n\rho}$ существуют соотношения
\begin{equation}
\begin{aligned}
 h_1(x)&=r_1\bigl(g_1(x),\ldots,g_\lambda(x)\bigr);\ldots;\\
 h_\sigma(x)&=r_\sigma\bigl(g_\mu(x),\ldots,g_\nu(x)\bigr),
\end{aligned}
\tag{2}
\end{equation}
где $r_1,\ldots,r_\sigma$ являются рациональными функциями своих аргументов, и
\begin{equation}
  g_1(x),\ldots,g_\lambda(x),\ldots,
  g_\mu(x),\ldots,g_\nu(x)
\tag{3}
\end{equation}
все принадлежат $\Ssys_{n\rho}$. Благодаря соотношениям (2), наряду с (1), также (3) задаёт рациональную базу $\Kfield_{n\rho}$; а поскольку функции (3) принадлежат $\Ssys_{n\rho}$, она одновременно является рациональной базой $\Ssys_{n\rho}$. Итак, имеем:

\emph{Теорема IV: Произвольная система $\Ssys_{n\rho}$ рациональных (или целых рациональных) функций обладает рациональной базой алгебраического ранга $\rho$, состоящей из конечного числа функций самой системы.}

Пусть теперь специально $\Lsys_{n\rho}$ является линейным семейством, а
\begin{equation}
  g_1(x),\ldots,g_\rho(x),\ g_{\rho+1}(x),\ldots,g_\nu(x)
\tag{4}
\end{equation}
есть рациональная база $\Lsys_{n\rho}$. Пусть, например,
\[
  g_1(x),\ldots,g_\rho(x)
\]
алгебраически независимы. Если тогда положить
\[
  g(x)=c_1g_{\rho+1}(x)+c_2g_{\rho+2}(x)+\cdots+c_{\nu-\rho}g_\nu(x),
\]
то $c_i$ можно выбрать как элементы из $\Omega$ так, что
\begin{equation}
  g_1(x),\ldots,g_\rho(x),\ g(x)
\tag{5}
\end{equation}
станет рациональной базой наименьшего содержащего поля $\Kfield_{n\rho}$. Но так как $\Lsys_{n\rho}$ является линейным семейством, $g(x)$ принадлежит $\Lsys_{n\rho}$, и (5) задаёт рациональную базу $\Lsys_{n\rho}$; то есть:

\emph{Теорема V: Рациональную базу каждого линейного семейства $\Lsys_{n\rho}$ рациональных (или целых рациональных) функций, а следовательно, в частности, рациональную базу каждой области целостности $\Jdom_{n\rho}$ и каждого поля $\Kfield_{n\rho}$, можно специально выбрать так, чтобы она содержала не более $\rho+1$ функций (и не менее $\rho$).}

Следовательно, ограничение числа функций рациональной базы, найденное в \S\ 4 для поля, основано на свойстве поля быть линейным семейством; тогда как ограничение до $\rho$ функций в случае существования минимальной базы использует все предпосылки поля.

Заметим ещё, что согласно \S\ 3 (4) все характеристические свойства систем и наименьших содержащих их систем сохраняются при отображении.
% END UNIT BODY
% END PRODUCER UNIT 047
\endgroup


\clearpage
\begingroup
\setcounter{footnote}{0}
% BEGIN PRODUCER UNIT 048
% Source: C:/Users/Floris/Documents/interlanguage/03_projects/noether/02_slavic_working_corpus/translations/paper06/russian/v001/Noether_Paper06_Section08_Russian_v001.tex
% Identity: 8002 B / SHA-256 E494EF6ECCBF059837B21CDF5E9950BDA1E130BDD3DF733C356F503B8170EBC4
\setlength{\parindent}{1.25em}
\setlength{\parskip}{0pt}
\emergencystretch=10em
\allowdisplaybreaks
\providecommand{\Kfield}{\mathfrak{K}}
\providecommand{\Ssys}{\mathfrak{S}}
\providecommand{\Jdom}{\mathfrak{J}}
\providecommand{\Gdom}{\mathfrak{G}}
\providecommand{\Lsys}{\mathfrak{L}}
\providecommand{\Omegaint}{[\Omega]}
\providecommand{\eps}{\varepsilon}
% BEGIN UNIT BODY
\section*{\S\ 8. Инволюционная база областей целостности $\Jdom_{n\rho}$ из полиномов.}

В следующих параграфах речь идет о системах $\Ssys_{n\rho}$, элементы которых являются целыми рациональными функциями, то есть полиномами; прежде всего об областях целостности $\Jdom_{n\rho}$ из полиномов. Для этих областей $\Jdom_{n\rho}$ сначала нужно зафиксировать понятия делимости.

Пусть $F(x),G(x)$ --- два полинома из $\Jdom_{n\rho}$, а $L(x)$ --- общий делитель этих полиномов, так что имеем:
\[
  F(x)=L(x)\cdot F_1(x),\qquad G(x)=L(x)\cdot G_1(x).
\]
$L(x)$ тогда и только тогда называется \emph{общим делителем в $\Jdom_{n\rho}$}, когда три полинома $L(x),F_1(x),G_1(x)$ принадлежат $\Jdom_{n\rho}$.

Два полинома из $\Jdom_{n\rho}$ называются \emph{взаимно простыми в $\Jdom_{n\rho}$}, если они не имеют общего делителя в $\Jdom_{n\rho}$.

Пусть $\Kfield_{n\rho}$ --- наименьшее поле, содержащее $\Jdom_{n\rho}$. Тогда из определения следует, что каждая функция из $\Kfield_{n\rho}$ имеет представление
\[
  f(x)=F_1(x):F_2(x),
\]
где $F_1(x),F_2(x)$ принадлежат $\Jdom_{n\rho}$ и взаимно просты в $\Jdom_{n\rho}$. Соответственно, несколько функций из $\Kfield_{n\rho}$ можно привести к общему знаменателю в $\Jdom_{n\rho}$:
\[
  f_1(x)=\frac{F_1(x)}{F(x)},\ldots,
  f_k(x)=\frac{F_k(x)}{F(x)}.
\]
$F(x)$ тогда и только тогда называется \emph{наименьшим общим знаменателем в $\Jdom_{n\rho}$}, когда полиномы из $\Jdom_{n\rho}$:
\[
  F(x),F_1(x),\ldots,F_k(x)
\]
взаимно просты в $\Jdom_{n\rho}$.

Такое представление может быть возможно различными способами, поскольку в $\Jdom_{n\rho}$ вообще говоря не действуют законы единственного разложения на неприводимые функции.

Теперь исследуем значение инволюционной базы $\Kfield_{n\rho}$ для $\Jdom_{n\rho}$, перенося понятия инволюционной формы и инволюционной базы на области целостности. Пусть
\[
  \Phi(z,u)=z^i+\sum{}' g_{\alpha\alpha_1\ldots\alpha_\rho}(x)
  z^\alpha u_1^{\alpha_1}\cdots u_\rho^{\alpha_\rho},
  \qquad
  (\alpha+\alpha_1+\cdots+\alpha_\rho=i)
\]
является инволюционной формой $\Kfield_{n\rho}$, а $G(x)$ --- наименьшим общим знаменателем в $\Jdom_{n\rho}$ коэффициентов $\Phi(z,u)$. Умножением на $G(x)$ форма $\Phi(z,u)$ переходит в форму $\Psi(z,u)$, коэффициенты которой являются полиномами из $\Jdom_{n\rho}$ и взаимно просты в $\Jdom_{n\rho}$:
\[
  G(x)\cdot\Phi(z,u)=\Psi(z,u)
  =G(x)z^i+
  \sum{}' G_{\alpha\alpha_1\ldots\alpha_\rho}(x)
  z^\alpha u_1^{\alpha_1}\cdots u_\rho^{\alpha_\rho},
  \qquad
  (\alpha+\alpha_1+\cdots+\alpha_\rho=i).
\]

Всякую форму $\Psi(z,u)$, которая возникает из инволюционной формы $\Phi(z,u)$ поля $\Kfield_{n\rho}$ умножением на полином из $\Jdom_{n\rho}$ так, что коэффициенты $\Psi(z,u)$ являются полиномами из $\Jdom_{n\rho}$ и взаимно просты в $\Jdom_{n\rho}$, будем называть \emph{инволюционной формой} $\Jdom_{n\rho}$; коэффициенты $\Psi(z,u)$ будем называть соответствующей \emph{инволюционной базой}.

Прежде всего из значения инволюционной базы $\Kfield_{n\rho}$ ясно, что каждая инволюционная база $\Jdom_{n\rho}$ одновременно является рациональной базой. Для дальнейшего исследования рассмотрим область отображения $\Jdom^*_{\rho\rho}$, для наименьшего содержащего ее поля $\Kfield^*_{\rho\rho}$ неприводимое уравнение с нулем
\[
  z=u_1x_1+\cdots+u_\rho x_\rho
\]
задано как $\Phi^*(z,u)=0$.

По \S\ 5 коэффициенты формы
\[
  \Phi^*(z,u)=z^i+\sum{}' g^*_{\alpha\alpha_1\ldots\alpha_\rho}(x)
  z^\alpha u_1^{\alpha_1}\cdots u_\rho^{\alpha_\rho},
  \qquad
  (\alpha+\alpha_1+\cdots+\alpha_\rho=i)
\]
задают элементарные симметрические функции рядов величин, сопряженных относительно $\Kfield^*_{\rho\rho}$:
\[
  x,\quad x^{(1)},\ldots,x^{(i-1)}.
\]
Пусть $F^*(x)$ --- произвольный полином из $\Kfield^*_{\rho\rho}$. Тогда из равенства
\[
  F^*(x)=\frac1i\Bigl(F^*(x)+F^*(x^{(1)})+
  \cdots+F^*(x^{(i-1)})\Bigr)
\]
следует, что $F^*(x)$ является целой симметрической функцией этих рядов величин и потому выражается как целая рациональная комбинация элементарных симметрических функций, то есть функций
\[
  g^*_{\alpha\alpha_1\ldots\alpha_\rho}(x).
\]
Если теперь
\[
\begin{aligned}
  G^*(x)\cdot\Phi^*(z,u)
  &=\Psi^*(z,u)\\
  &=G^*(x)z^i+
    \sum{}'G^*_{\alpha\alpha_1\ldots\alpha_\rho}(x)
    z^\alpha u_1^{\alpha_1}\cdots u_\rho^{\alpha_\rho},\\
  &\hspace{5em}(\alpha+\alpha_1+\cdots+\alpha_\rho=i)
\end{aligned}
\]
является инволюционной формой $\Jdom^*_{\rho\rho}$, то всякий полином из $\Kfield^*_{\rho\rho}$, а значит в частности всякий полином из $\Jdom^*_{\rho\rho}$, является целой рациональной комбинацией частных
\[
  G^*_{\alpha\alpha_1\ldots\alpha_\rho}(x):G^*(x).
\]

С учетом принципа переноса отсюда получается:

\emph{Теорема VI: Для каждой области целостности из полиномов $\Jdom_{n\rho}$ существует по меньшей мере одна выделенная рациональная база, инволюционная база, такая, что в рациональном представлении всех полиномов из $\Jdom_{n\rho}$ через функции инволюционной базы в знаменателе появляется только степень одной фиксированной функции --- старшего коэффициента соответствующей инволюционной формы.}
% END UNIT BODY
% END PRODUCER UNIT 048
\endgroup


\clearpage
\begingroup
\setcounter{footnote}{0}
% BEGIN PRODUCER UNIT 049
% Source: C:/Users/Floris/Documents/interlanguage/03_projects/noether/02_slavic_working_corpus/translations/paper06/russian/v001/Noether_Paper06_Section09_Russian_v001.tex
% Identity: 12091 B / SHA-256 4B8B2E1E31E9AE7BF1E8AFA22EF2C0382095F12287B0B207DB972D46B32E0C29
\setlength{\parindent}{1.25em}
\setlength{\parskip}{0pt}
\emergencystretch=10em
\allowdisplaybreaks
\providecommand{\Kfield}{\mathfrak{K}}
\providecommand{\Ssys}{\mathfrak{S}}
\providecommand{\Jdom}{\mathfrak{J}}
\providecommand{\Gdom}{\mathfrak{G}}
\providecommand{\Lsys}{\mathfrak{L}}
\providecommand{\Omegaint}{[\Omega]}
\providecommand{\eps}{\varepsilon}
% BEGIN UNIT BODY
\section*{\S\ 9. Относительно целые области $\Gdom_{n\rho}$ из полиномов.}

Далее рассматриваем специальные области целостности из полиномов, промежуточные между областью целостности и полем.

Область целостности из полиномов $\Gdom_{n\rho}$ называется ``относительно целой областью'' при условии:

\emph{4a) $\Gdom_{n\rho}$ вместе с $F(x)$ и $G(x)$ --- где $G(x)\ne0$ --- содержит частное $F(x):G(x)$ тогда и только тогда, когда это частное цело по неопределенным, то есть является полиномом.}

Сначала нужно доказать, что относительно целые области совпадают с совокупностями полиномов некоторого поля $\Kfield_{n\sigma}$ ($\sigma\ge \rho$) рациональных функций.

Пусть $\Kfield_{n\rho}$ есть наименьшее поле, содержащее $\Gdom_{n\rho}$. Для каждой функции $f(x)$ из $\Kfield_{n\rho}$ существует представление в виде частного двух полиномов из $\Gdom_{n\rho}$:
\[
  f(x)=F_1(x):F_2(x),
\]
причем допускаются и полиномы нулевой степени. Как только $f(x)$ становится полиномом, она по условию 4a) принадлежит $\Gdom_{n\rho}$; а так как $\Gdom_{n\rho}$ не может содержать ничего сверх полиномов, то \emph{каждая относительно целая область $\Gdom_{n\rho}$ тождественна совокупности полиномов наименьшего содержащего ее поля.}

Обратно, пусть $\Kfield_{n\sigma}$ есть произвольное поле рациональных функций (то есть не обязательно наименьшее поле, содержащее данную систему), и пусть $\Jdom$ есть совокупность полиномов, содержащихся в $\Kfield_{n\sigma}$. $\Jdom$ удовлетворяет условиям 1, 2, 3 области целостности и условию 4a), следовательно является относительно целой областью; то есть \emph{совокупность полиномов, содержащихся в поле $\Kfield_{n\sigma}$, образует относительно целую область $\Gdom_{n\rho}$ ($\rho\le\sigma$).}\footnote{Возможен также случай $\rho=0$; то есть совокупность полиномов из $\Kfield_{n\sigma}$ состоит из элементов области коэффициентов $\Omega$.}

Далее покажем тождественность относительно целых областей введенным Гильбертом ``областям целостности из относительно целых функций''.\footnote{D. Hilbert: Mathematische Probleme. Доклад, Париж 1900. Problem 14 (Göttinger Nachrichten 1900).}

Пусть
\begin{equation}
  G_1(x),\ldots,G_k(x)
\tag{1}
\end{equation}
есть рациональная база $\Gdom_{n\rho}$. По условиям 1, 2, 3, 4a каждая рациональная комбинация полиномов (1), целая по неопределенным, то есть становящаяся полиномом, принадлежит $\Gdom_{n\rho}$. Обратно, каждый полином из $\Gdom_{n\rho}$ по понятию рациональной базы рационально выражается через полиномы (1); поэтому $\Gdom_{n\rho}$ \emph{тождественна совокупности тех рациональных комбинаций полиномов (1), которые целы по неопределенным, то есть являются полиномами.} Тем самым мы имеем определение Гильберта, исходящее из специальной рациональной базы, тогда как наше определение использует только абстрактные свойства области.

Для относительно целых областей определение общего делителя из \S\ 8 для общих областей целостности упрощается:

Пусть $F(x),G(x)$ --- два полинома из $\Gdom_{n\rho}$, а $L(x)$ --- общий делитель этих полиномов. Тогда $L(x)$ называется \emph{общим делителем в $\Gdom_{n\rho}$} тогда и только тогда, когда $L(x)$ принадлежит $\Gdom_{n\rho}$.

В самом деле, тогда принадлежность частных $F(x):L(x)$ и $G(x):L(x)$ следует из условия 4a).

Другое существенное свойство относительно целых областей состоит в том, что понятие ``алгебраически целого относительно $\Gdom_{n\rho}$'' можно определять через произвольное уравнение точно так же, как для алгебраических числовых полей или для поля $\Omega(x_1,\ldots,x_n)$.

\emph{Определение V: Величина $\xi$ называется алгебраически целой относительно $\Gdom_{n\rho}$ --- или также относительно алгебраически целой ---, если она удовлетворяет некоторому уравнению, старший коэффициент которого является единицей, а остальные коэффициенты являются полиномами из $\Gdom_{n\rho}$.}

В самом деле, пусть относительно алгебраически целая величина $\xi$ удовлетворяет, например, уравнению
\[
  L(z)=z^\lambda+G_1(x)z^{\lambda-1}+\cdots+G_\lambda(x)=0.
\]
$L(z)$ делится на целую функцию
\[
  M(z)=z^\mu+H_1(x)z^{\mu-1}+\cdots+H_\mu(x),
\]
неприводимую относительно наименьшего поля $\Kfield_{n\rho}$, содержащего $\Gdom_{n\rho}$. Но из этой делимости по известному обобщению теоремы Гаусса\footnote{Ср. J. König: Einleitung in die allgemeine Theorie der algebraischen Größen (Leipzig, B.~G. Teubner, 1903), Kap. IX, \S\ 2, или Weber (малое издание), \S\ 20.} следует, что рациональные функции $H_1(x),\ldots,H_\mu(x)$ являются полиномами из $\Kfield_{n\rho}$, а значит принадлежат $\Gdom_{n\rho}$. Тем самым определение относительно алгебраически целой величины $\xi$ получается из неприводимого уравнения. В частности, каждый полином $F(x)$ из $\Gdom_{n\rho}$ также алгебраически цел относительно $\Gdom_{n\rho}$, поскольку он удовлетворяет уравнению $z-F(x)=0$.

Далее отметим, что для области целостности $\Jdom_{n\rho}$ из полиномов, аналогично наименьшему содержащему полю, определена и наименьшая содержащая ее относительно целая область $\Gdom_{n\rho}$ следующим требованием:

$\Gdom_{n\rho}$ вместе с $\Jdom_{n\rho}$ содержит все полиномы $H(x)$, и только те, которые можно представить как частные двух полиномов из $\Jdom_{n\rho}$.

Из этого определения далее следует: из каждой инволюционной формы и инволюционной базы $\Jdom_{n\rho}$ соответствующая форма и база для $\Gdom_{n\rho}$ получаются так, что нужно самое большее выделить общий делитель в $\Gdom_{n\rho}$ всех полиномов базы.

Действительно, $\Jdom_{n\rho}$ и $\Gdom_{n\rho}$ имеют одно и то же наименьшее содержащее поле $\Kfield_{n\rho}$; и каждая инволюционная форма $\Jdom_{n\rho}$ возникает из инволюционной формы $\Kfield_{n\rho}$ умножением на полином из $\Jdom_{n\rho}$, поэтому имеет коэффициенты из $\Gdom_{n\rho}$, которые, однако, могут иметь общий делитель в $\Gdom_{n\rho}$.\footnote{Например, для области целостности $\Jdom_{22}$, состоящей из всех целых рациональных комбинаций $x_1^2$ и $x_1x_2$, инволюционная форма имеет вид:
\[
\Psi(z,u)=x_1^2z^2-x_1^4u_1^2-2x_1^3x_2u_1u_2-x_1^2x_2^2u_2^2.
\]
Так как $x_1^2$ принадлежит $\Jdom_{22}$, а тем более наименьшей содержащей ее относительно целой области $\Gdom_{22}$, и потому является общим делителем в $\Gdom_{22}$, инволюционная форма $\Gdom_{22}$ имеет вид:
\[
\Psi(z,u)=z^2-x_1^2u_1^2-2x_1x_2u_1u_2-x_2^2u_2^2.
\]}

Наконец отметим, что \emph{область отображения} относительно целой области $\Gdom_{n\rho}$, которая по \S\ 7 снова является областью целостности $\Jdom^*_{\rho\rho}$ из полиномов, не обязана сама быть относительно целой областью. Действительно, если $F(x)$ и $G(x)$ --- два полинома из $\Gdom_{n\rho}$, частное которых не является полиномом и потому не принадлежит $\Gdom_{n\rho}$, то частное $F^*(x):G^*(x)$ также не принадлежит $\Jdom^*_{\rho\rho}$. Но это частное, возникающее специализацией некоторых неопределенных, вполне может быть полиномом, и тогда для $\Jdom^*_{\rho\rho}$ условие 4a) не выполнено.\footnote{Пусть, например, $\Gdom_{22}$ состоит из всех полиномов, рационально выражаемых через $F=x_1^2x_2$ и $G=x_1^2(1+x_3)$. Допустимая специализация $x_3=0$ дает $F^*:G^*=x_2$, тогда как $F:G$ не полином. Следовательно, $\Jdom^*_{22}$ не является относительно целой областью.}
% END UNIT BODY
% END PRODUCER UNIT 049
\endgroup


\clearpage
\begingroup
\setcounter{footnote}{0}
% BEGIN PRODUCER UNIT 050
% Source: C:/Users/Floris/Documents/interlanguage/03_projects/noether/02_slavic_working_corpus/translations/paper06/russian/v001/Noether_Paper06_Section10_Russian_v001.tex
% Identity: 9607 B / SHA-256 3D7F33C5DBE3A0AA57B2464EC66238AF8A796E4CB455CBEB9FF49E5C69517B8F
\setlength{\parindent}{1.25em}
\setlength{\parskip}{0pt}
\emergencystretch=10em
\allowdisplaybreaks
\providecommand{\Kfield}{\mathfrak{K}}
\providecommand{\Ssys}{\mathfrak{S}}
\providecommand{\Jdom}{\mathfrak{J}}
\providecommand{\Gdom}{\mathfrak{G}}
\providecommand{\Lsys}{\mathfrak{L}}
\providecommand{\Omegaint}{[\Omega]}
\providecommand{\eps}{\varepsilon}
% BEGIN UNIT BODY
\section*{\S\ 10. Относительно целые области первого рода\\ и их интегральный базис.}

Понятие ``относительно алгебраически целого'', введенное в \S\ 9 (Определение V), позволяет указать класс относительно целых областей, которые являются конечными областями целостности, и тем самым хотя бы для этого класса положительно ответить на вопрос, поставленный Д. Гильбертом (Mathematische Probleme, проблема 14). Дадим следующее

\emph{Определение VI: Конечное число полиномов из $\Gdom_{n\rho}$ называется интегральным базисом, если каждый полином из $\Gdom_{n\rho}$ можно представить как целую рациональную комбинацию этого конечного числа с коэффициентами из $\Omega$. Область целостности из полиномов, имеющая интегральный базис, называется конечной областью целостности.}

Рассматриваемый класс относительно целых областей, для которых инволюционный базис окажется интегральным базисом, будем называть относительно целыми областями первого рода по следующему определению:

\emph{Определение VII: Относительно целая область $\Gdom_{n\rho}$ называется областью первого рода, если ее область отображения является относительно целой областью $\Gdom^*_{\rho\rho}$ такой, что каждый произвольный полином от $x_1,\ldots,x_\rho$ (с коэффициентами из $\Omega$) алгебраически цел над $\Gdom^*_{\rho\rho}$.}

Действительно, по этому определению неопределенные $x_1,\ldots,x_\rho$, а следовательно и линейная форма $u(x)=u_1x_1+\cdots+u_\rho x_\rho$, алгебраически целы над $\Gdom^*_{\rho\rho}$. Поэтому неприводимая целая функция с корнем $z=u(x)$ --- инволюционная форма наименьшего содержащего поля $\Kfield^*_{\rho\rho}$ --- имеет старшим коэффициентом единицу, а остальными коэффициентами полиномы из $\Gdom^*_{\rho\rho}$. Следовательно, она одновременно является инволюционной формой $\Gdom^*_{\rho\rho}$, и в $\Gdom^*_{\rho\rho}$ не может существовать другая инволюционная форма. По принципу переноса $\Gdom_{n\rho}$ также имеет инволюционную форму, старший коэффициент которой равен единице; а так как по теореме VI в рациональном представлении всех полиномов из $\Gdom_{n\rho}$ через соответствующий инволюционный базис в знаменатель входит только этот старший коэффициент, инволюционный базис становится интегральным базисом, то есть:

\emph{Теорема VII. Каждая относительно целая область первого рода является конечной областью целостности; интегральный базис задается инволюционным базисом, полученным из характеризующей области отображения. В частности, интегральный базис относительно целых областей первого рода $\Gdom_{\rho\rho}$ задается однозначно определенным инволюционным базисом $\Gdom_{\rho\rho}$.}

Дадим еще критерий для относительно целых областей первого рода. Пусть $\Gdom_{n\rho}$ является областью первого рода, и пусть интегральный базис $\Gdom_{n\rho}$ задан полиномами
\begin{equation}
  F_1(x),F_2(x),\ldots,F_k(x).
\tag{1}
\end{equation}
Тогда по вспомогательной теореме Гильберта\footnote{D. Hilbert: Über die vollen Invariantensysteme, Math. Ann. 42 (1893), \S\ 1. (Вспомогательная теорема, сформулированная там только для однородных полиномов, так же верна для неоднородных.)} существуют $\rho$ алгебраически независимых полиномов из $\Gdom_{n\rho}$:
\begin{equation}
  G_1(x),G_2(x),\ldots,G_\rho(x)
\tag{2}
\end{equation}
таких, что полиномы (1), а значит и каждый полином из $\Gdom_{n\rho}$, алгебраически целы над полиномами (2). То же верно для каждого полинома области отображения $\Gdom^*_{\rho\rho}$ относительно соответствующих полиномов
\begin{equation}
  G_1^*(x),G_2^*(x),\ldots,G_\rho^*(x).
\tag{3}
\end{equation}

По Определению VII тогда и каждый произвольный полином от $x_1,\ldots,x_\rho$ алгебраически цел над полиномами (3).

Обратно, пусть область отображения $\Jdom^*_{\rho\rho}$ произвольной относительно целой области $\Gdom_{n\rho}$ содержит $\rho$ полиномов (3), над которыми каждый полином от $x_1,\ldots,x_\rho$ алгебраически цел. Тогда покажем, что отсюда сама $\Jdom^*_{\rho\rho}$ получается как относительно целая область $\Gdom^*_{\rho\rho}$, и далее, что $\Gdom^*_{\rho\rho}$, а следовательно и $\Gdom_{n\rho}$, являются областями первого рода.

В самом деле, пусть $F^*(x)$ --- полином из наименьшего поля $\Kfield^*_{\rho\rho}$, содержащего $\Jdom^*_{\rho\rho}$; по предположению он алгебраически цел над полиномами (3). Тогда соответствующая функция из $\Kfield_{n\rho}$ алгебраически цела над $\rho$ полиномами из $\Gdom_{n\rho}$, значит является полиномом $F(x)$ и потому принадлежит $\Gdom_{n\rho}$; следовательно, $F^*(x)$ принадлежит $\Jdom^*_{\rho\rho}$. Область отображения $\Jdom^*_{\rho\rho}$ содержит все полиномы из $\Kfield^*_{\rho\rho}$ и является относительно целой областью $\Gdom^*_{\rho\rho}$. Для $\Gdom^*_{\rho\rho}$, а значит и для $\Gdom_{n\rho}$, из предположения по Определениям V и VII непосредственно следует, что они являются областями первого рода, то есть:

\emph{Относительно целые области первого рода $\Gdom_{n\rho}$ характеризуются тем, что в области отображения $\Jdom^*_{\rho\rho}$ существуют $\rho$ полиномов, над которыми каждый произвольный полином от $x_1,\ldots,x_\rho$ алгебраически цел. Из этого предположения сама область отображения получается как относительно целая область.}\footnote{Например, в области $\Gdom_{22}$, упомянутой в заключительном замечании к \S\ 9, допустимая специализация $x_1=1$ дает два полинома $F^*=x_2$, $G^*=1+x_3$, над которыми каждый произвольный полином от $x_2,x_3$ алгебраически цел. Следовательно, $\Gdom_{22}$ является областью первого рода, причем с $F(x)$ и $G(x)$ в качестве интегрального базиса.}
% END UNIT BODY
% END PRODUCER UNIT 050
\endgroup


\clearpage
\begingroup
\setcounter{footnote}{0}
% BEGIN PRODUCER UNIT 051
% Source: C:/Users/Floris/Documents/interlanguage/03_projects/noether/02_slavic_working_corpus/translations/paper06/russian/v001/Noether_Paper06_Section11_Russian_v001.tex
% Identity: 6140 B / SHA-256 36112F0E19F81A6C973DA31D8F2698561B9C1F00511D07561AE700D279D41B94
\setlength{\parindent}{1.25em}
\setlength{\parskip}{0pt}
\emergencystretch=10em
\allowdisplaybreaks
\providecommand{\Kfield}{\mathfrak{K}}
\providecommand{\Ssys}{\mathfrak{S}}
\providecommand{\Jdom}{\mathfrak{J}}
\providecommand{\Gdom}{\mathfrak{G}}
\providecommand{\Lsys}{\mathfrak{L}}
\providecommand{\Omegaint}{[\Omega]}
\providecommand{\eps}{\varepsilon}
% BEGIN UNIT BODY
\section*{\S\ 11. Вспомогательная теорема\\ об алгебраически целой зависимости.}

Теперь дадим критерий того, что каждый произвольный полином от $x_1,\ldots,x_\rho$ алгебраически цел над $\rho$ полиномами от $x_1,\ldots,x_\rho$
\begin{equation}
  G_1^*(x),\ldots,G_\rho^*(x).
\tag{1}
\end{equation}

Если полиномы (1) специально являются однородными формами, то при каждой специальной системе значений, обращающей (1) в нуль, одновременно обращаются в нуль и все алгебраически цело зависимые формы; в частности, также $x_1,\ldots,x_\rho$.

Следовательно, формы (1) не могут обращаться в нуль ни при какой системе значений, кроме
\[
  x_1=0,\ldots,x_\rho=0;
\]
их результант должен быть отличен от нуля.

Обратно, это условие также достаточно и допускает распространение на неоднородные полиномы согласно следующей

\emph{Вспомогательной теореме: Достаточным условием того, что каждый произвольный полином от $x_1,\ldots,x_\rho$ алгебраически цел над $\rho$ полиномами от $x_1,\ldots,x_\rho$}
\[
  G_1^*(x),\ldots,G_\rho^*(x),
\]
\emph{является неравенство нулю результанта, взятого по $x_1,\ldots,x_\rho$, старших однородных частей}
\[
  H_1^*(x),\ldots,H_\rho^*(x)
\]
\emph{полиномов (1):}
\begin{equation}
  R_0=\left[
  \begin{array}{ccc}
    H_1^*&\cdots&H_\rho^*\\
    x_1&\cdots&x_\rho
  \end{array}
  \right]\ne0.
\tag{2}
\end{equation}
\emph{Для однородных форм (1) условие (2) необходимо.}\footnote{О теории результантов ср. F. Mertens, Zur Theorie der Elimination (I. u. II. Teil), Sitzungsber. d. k. Ak. d. Wiss. Wien, Math.-naturw. Kl. Abt. IIa, Bd. 108 (1899). (Частично воспроизведено у J. König, Theorie d. algebr. Größen, Kap. VI).}

Для доказательства образуем полиномы
\begin{equation}
\begin{gathered}
  \bar G_1^*(x)=G_1^*(x)-\lambda_1,
  \ldots,
  \bar G_\rho^*(x)=G_\rho^*(x)-\lambda_\rho,\\
  \bar\Phi^*(x)=\Phi^*(x)-\mu,
\end{gathered}
\tag{3}
\end{equation}
где $\lambda_1,\ldots,\lambda_\rho,\mu$ обозначают неопределенные, а $\Phi^*(x)$ --- произвольный полином от $x_1,\ldots,x_\rho$. Пусть результант полиномов (3) по $x_1,\ldots,x_\rho,1$\footnote{То есть однородный результант по $x_1,\ldots,x_\rho,x_0$, если однородить полиномы (3) посредством дополнительной неопределенной $x_0$ так, чтобы степень по $x$ сохранялась.} задан как
\begin{equation}
  R(\lambda,\mu)=\left[
  \begin{array}{cccc}
   \bar G_1^*&\cdots&\bar G_\rho^*&\bar\Phi^*\\
   x_1&\cdots&x_\rho&1
  \end{array}
  \right];
\tag{4}
\end{equation}
это целая рациональная функция от $\lambda_1,\ldots,\lambda_\rho,\mu$. По Ф. Мертенсу в разложении (4) по степеням $\mu$ можно явно указать старший коэффициент;\footnote{A. a. O. \S\ 3 (доказательство Satz I) и явно \S\ 6 (определение совокупности всех членов результанта $R$, содержащих заданное степенное произведение $a_{\lambda\lambda}^{p_\lambda} a_{\mu\mu}^{p_\mu}\cdots a_{\sigma\sigma}^{p_\sigma}$ как множитель). В нашем случае положено $p_\lambda=0$, $p_\mu=0,\ldots,p_\rho=\tau$, $a_{\sigma\sigma}=(-\mu)$. Воспроизведено у König, Kap. VI, \S\ 9.} получается
\begin{equation}
  (-1)^\tau R(\lambda,\mu)
  =R_0^\sigma\mu^\tau+A_1(\lambda)\mu^{\tau-1}+\cdots+A_\tau(\lambda),
\tag{5}
\end{equation}
где $A_1(\lambda),\ldots,A_\tau(\lambda)$ обозначают целые целочисленные функции от $\lambda_1,\ldots,\lambda_\rho$ и коэффициентов полиномов (3). Это разложение прежде всего показывает, что при предположении (2), $R_0\ne0$, также $R(\lambda,\mu)$ не может тождественно обращаться в нуль. Далее тождество по $x_1,\ldots,x_\rho,\lambda_1,\ldots,\lambda_\rho,\mu$:
\[
  R(\lambda,\mu)\equiv 0\pmod{\bar G_1^*(x),\ldots,
  \bar G_\rho^*(x),\bar\Phi^*(x)}
\]
при подстановке
\[
  \lambda_i=G_i^*(x),\qquad \mu=\Phi^*(x)
\]
переходит в:
\begin{equation}
\begin{aligned}
  R(G_i^*(x),\Phi^*(x))
  &=R_0^\sigma\Phi^{*\tau}(x)
    +A_1(G_i^*(x))\Phi^{*\tau-1}(x)+\cdots \\
  &\qquad +A_\tau(G_i(x))=0,
\end{aligned}
\tag{6}
\end{equation}
чем вспомогательная теорема доказана, поскольку $R_0$ по предположению является числом из $\Omega$.
% END UNIT BODY
% END PRODUCER UNIT 051
\endgroup


\clearpage
\begingroup
\setcounter{footnote}{0}
% BEGIN PRODUCER UNIT 052
% Source: C:/Users/Floris/Documents/interlanguage/03_projects/noether/02_slavic_working_corpus/translations/paper06/russian/v001/Noether_Paper06_Section12_Russian_v001.tex
% Identity: 13769 B / SHA-256 11670A6BAF9A462B7F0B24D2ACCDCEA91C4F98ABA870D2D5F57660B75DC80A6B
\setlength{\parindent}{1.25em}
\setlength{\parskip}{0pt}
\emergencystretch=10em
\allowdisplaybreaks
\providecommand{\Kfield}{\mathfrak{K}}
\providecommand{\Ssys}{\mathfrak{S}}
\providecommand{\Jdom}{\mathfrak{J}}
\providecommand{\Gdom}{\mathfrak{G}}
\providecommand{\Lsys}{\mathfrak{L}}
\providecommand{\Omegaint}{[\Omega]}
\providecommand{\eps}{\varepsilon}
% BEGIN UNIT BODY
\section*{\S\ 12. Интегральный базис регулярных систем\\ $\Ssys_{n\rho}$ из полиномов.}

Вспомогательная теорема \S\ 11 вместе с заключением \S\ 10 непосредственно дает следующий факт:

\emph{Если в области отображения $\Jdom^*_{\rho\rho}$ области $\Gdom_{n\rho}$ существуют $\rho$ полиномов таких, что однородный результант их старших однородных частей не обращается тождественно в нуль --- $R_0\ne0$ ---, то $\Gdom_{n\rho}$ является областью первого рода и, следовательно, конечной областью целостности, причем ее инволюционный базис является интегральным базисом.}

Условие $R_0\ne0$ позволяет теперь, посредством вспомогательной теоремы, дать второе доказательство существования интегрального базиса, по существу принадлежащее Д. Гильберту;\footnote{\emph{Über die vollen Invariantensysteme}, \S\ 2. У Гильберта, при предположении уже иначе полученного существования интегрального базиса поля инвариантов, речь идет лишь о толковании этого базиса через кронекерову фундаментальную систему алгебраически целых величин поля.} это доказательство, правда, не позволяет распознать инволюционный базис как таковой, но зато равномерно действует для всех систем $\Ssys_{n\rho}$ с указанным свойством.

А именно, пусть $\rho$ полиномов из $\Jdom^*_{\rho\rho}$, для которых $R_0\ne0$ и которые, следовательно, алгебраически независимы, заданы как
\begin{equation}
  G_1^*(x),\ldots,G_\rho^*(x).
\tag{1}
\end{equation}
По вспомогательной теореме каждый произвольный полином от $x_1,\ldots,x_\rho$, а значит, в частности, каждый полином наименьшего поля $\Kfield^*_{\rho\rho}$, содержащего $\Jdom^*_{\rho\rho}$, алгебраически цел над полиномами (1). Если им в $\Ssys_{n\rho}$ соответствуют полиномы
\begin{equation}
  G_1(x),\ldots,G_\rho(x),
\tag{2}
\end{equation}
то по принципу переноса каждый полином наименьшего поля $\Kfield_{n\rho}$, содержащего $\Ssys_{n\rho}$, а следовательно, в частности, каждый полином из $\Ssys_{n\rho}$, алгебраически цел над полиномами (2). Обратно, всякая функция из $\Kfield_{n\rho}$, алгебраически целая над полиномами (2), также является полиномом. Поэтому, рассматривая $\Kfield_{n\rho}$ (по \S\ 4) как алгебраическое поле над $\Omega(G_1,\ldots,G_\rho)$, получаем, что совокупность полиномов из $\Kfield_{n\rho}$ --- относительно целая область $\Gdom_{n\rho}$ --- тождественна совокупности алгебраически целых величин этого поля. Пусть $G(x)$ есть полином из $\Kfield_{n\rho}$, дополняющий (2) до рационального базиса, а $D(x)$ --- дискриминант неприводимого уравнения $k$-й степени, которому $G(x)$ удовлетворяет относительно $\Omega(G_1,\ldots,G_\rho)$. Тогда каждый полином из $\Kfield_{n\rho}$, в частности каждый полином $F(x)$ из $\Ssys_{n\rho}$, как алгебраически целая величина допускает представление
\begin{equation}
\begin{split}
  F(x)=\frac{1}{D(x)}\big(&\Gamma_1(G_i)G(x)^{k-1}
    +\Gamma_2(G_i)G(x)^{k-2}\\
    &+\cdots+\Gamma_k(G_i)\big).
\end{split}
\tag{3}
\end{equation}

Пусть теперь $F(x)$ пробегает все полиномы из $\Ssys_{n\rho}$. Тогда применение гильбертовой теоремы I о модульном базисе (Math. Ann. 36) к системе, образуемой из соответствующих функций
\[
  v_1\Gamma_1(G_i)+v_2\Gamma_2(G_i)+\cdots+v_k\Gamma_k(G_i),
\]
показывает существование конечного числа полиномов из $\Ssys_{n\rho}$:
\begin{equation}
  F_1(x),F_2(x),\ldots,F_\lambda(x),
\tag{4}
\end{equation}
таких, что каждый полином из $\Ssys_{n\rho}$ допускает представление
\begin{equation}
\begin{aligned}
  F(x)&=A_1(G_i)F_1(x)+A_2(G_i)F_2(x)+\cdots\\
      &\qquad +A_\lambda(G_i)F_\lambda(x),
\end{aligned}
\tag{5}
\end{equation}
где $A_1,\ldots,A_\lambda$ обозначают полиномы от $G_i(x)$. Полиномы (2) и (4), следовательно, образуют интегральный базис $\Ssys_{n\rho}$; то есть:

\emph{Теорема VIII: Если в системе отображения $\Jdom^*_{\rho\rho}$ системы полиномов $\Ssys_{n\rho}$ существуют $\rho$ полиномов таких, что результант старших однородных частей не обращается тождественно в нуль --- $R_0\ne0$ ---, то $\Ssys_{n\rho}$ обладает интегральным базисом.}

Эта теорема допускает еще небольшое обобщение. Достаточно, чтобы $\rho$ полиномов (1), для которых $R_0\ne0$, принадлежали наименьшей области целостности $\Jdom^*_{\rho\rho}$, содержащей $\Ssys^*_{\rho\rho}$. В самом деле, рассуждения, ведущие к теореме VIII, показывают, что и тогда представление (5) еще остается в силе. Но по определению каждый полином $L(x)$ наименьшей содержащей области целостности $\Jdom_{n\rho}$ можно цело и рационально представить через конечное число --- меняющееся вместе с $L(x)$ --- полиномов из $\Ssys_{n\rho}$; следовательно, существуют $\sigma$ полиномов из $\Ssys_{n\rho}$
\begin{equation}
  K_1(x),K_2(x),\ldots,K_\sigma(x),
\tag{6}
\end{equation}
таких, что полиномы $G_i(x)$ становятся целыми рациональными комбинациями полиномов (6). Поэтому (6) и (4) дают интегральный базис. Вводя:

\emph{Определение VIII: Система $\Ssys_{n\rho}$ полиномов называется регулярной, если в области отображения $\Jdom^*_{\rho\rho}$ наименьшей содержащей области целостности существуют $\rho$ полиномов таких, что результант старших однородных частей не обращается тождественно в нуль --- $R_0\ne0$ ---}

получаем:

\emph{Теорема IX: Каждая регулярная система $\Ssys_{n\rho}$ из полиномов обладает интегральным базисом.}\footnote{Что существуют нерегулярные системы без интегрального базиса, Гильберт показал на примере (Gött. Nachr. 1891, с. 233). Еще немного проще следующий пример:
\[
  \Ssys_{22}=x_1,
  x_1x_2,
  x_1x_2^2,
  \ldots,
  x_1x_2^\nu,
  \ldots\quad\text{in inf.}
\]
Существование интегрального базиса здесь было бы равносильно существованию целого положительного числа $\nu$ такого, что выполняются тождества
\[
  x_1x_2^{\nu+\sigma}=\sum C\cdot x_1^{i_0}(x_1x_2)^{i_1}(x_1x_2^2)^{i_2}\cdots(x_1x_2^\nu)^{i_\nu}
  \quad(\sigma=1,2,\ldots).
\]
Сравнение размерностей по $x_1,x_2$ уже при $\sigma=1$ приводит к двум уравнениям условий для показателей:
\[
  1=i_0+i_1+\cdots+i_\nu,
  \qquad
  \nu+1=i_1+2i_2+\cdots+\nu i_\nu,
\]
которые очевидно не имеют решения в неотрицательных целых числах. Следовательно, $\Ssys_{22}$ не обладает интегральным базисом.}

Дадим еще, аналогично \S\ 5 для инволюции, геометрическое толкование регулярных систем. Для этого рассмотрим, как там в (7), систему уравнений
\begin{equation}
  L^*(x)-L^*(\xi)=0,
\tag{7}
\end{equation}
где $L^*(x)$ пробегает все полиномы из $\Jdom^*_{\rho\rho}$. Если однородить посредством $x_0,\xi_0$, то пусть (7) переходит в систему уравнений между однородными формами
\begin{equation}
  \xi_0^{m}M^*(x)-x_0^{m}M^*(\xi)=0,
\tag{8}
\end{equation}
причем, если $H^*(x)$ обозначает старшие однородные части $L^*(x)$, имеет место соотношение
\begin{equation}
  H^*(x)=M^*(x)\pmod{x_0}.
\tag{9}
\end{equation}

\emph{Фундаментальными точками} $\Jdom^*_{\rho\rho}$ будем называть (ср. примечание к \S\ 5) решения (8), не зависящие от $x$, то есть отличные от рядов, сопряженных с $x$. Следовательно, фундаментальные точки суть системы значений, отличные от $\xi_1=0,\ldots,\xi_\rho=0$, для которых одновременно выполняется
\[
  \xi_0^m=0,\qquad M^*(\xi)=0,
\]
или, ввиду (9),
\[
  \xi_0=0,
  \qquad H^*(\xi)=0.
\]

Если теперь существуют $\rho$ форм $H^*(\xi)$, для которых $R_0\ne0$, то есть не имеющих общей системы решений, то $\Jdom^*_{\rho\rho}$ не может иметь фундаментальной точки. Если, в частности, $\Ssys^*_{\rho\rho}$ состоит из однородных форм, то при $R_0\ne0$ и сама $\Ssys^*_{\rho\rho}$ не может иметь фундаментальной точки, поскольку такая точка была бы одновременно фундаментальной точкой $\Jdom^*_{\rho\rho}$. Но верно и обратное. Для этого рассмотрим систему $\Jdom(H^*)$, образованную из совокупности форм $H^*(x)$; она снова образует область целостности.

Из гильбертовой теоремы I о модульном базисе и гильбертовой вспомогательной теоремы (ср. цитату в \S\ 10) следует существование $\rho$ форм из $\Jdom(H^*)$, обращение которых в нуль влечет обращение в нуль всех форм из $\Jdom(H^*)$. Если теперь $\Jdom^*_{\rho\rho}$ не имеет фундаментальной точки, эти $\rho$ форм не могут обращаться в нуль одновременно; их результант $R_0\ne0$.

Итак:

\emph{Регулярные системы $\Ssys_{n\rho}$ из полиномов характеризуются тем, что область отображения $\Jdom^*_{\rho\rho}$ наименьшей содержащей области целостности не имеет фундаментальной точки. Для регулярных систем $\Ssys^*_{\rho\rho}$ из однородных форм уже сама система не имеет фундаментальной точки.}\footnote{Нерегулярная система $\Ssys_{22}$, указанная в предыдущем примечании, имеет фундаментальную точку:
\[
  \xi_1:\xi_2:\xi_0=0:1:0.
\]}
% END UNIT BODY
% END PRODUCER UNIT 052
\endgroup


\clearpage
\begingroup
\setcounter{footnote}{0}
% BEGIN PRODUCER UNIT 053
% Source: C:/Users/Floris/Documents/interlanguage/03_projects/noether/02_slavic_working_corpus/translations/paper06/russian/v001/Noether_Paper06_Section13_Russian_v001.tex
% Identity: 8009 B / SHA-256 D2B427916F8427817AF1CEBA966B7417030B25CEEA0A9DBF7B97ED3B79ED799A
\setlength{\parindent}{1.25em}
\setlength{\parskip}{0pt}
\emergencystretch=10em
\allowdisplaybreaks
\providecommand{\Kfield}{\mathfrak{K}}
\providecommand{\Ssys}{\mathfrak{S}}
\providecommand{\Jdom}{\mathfrak{J}}
\providecommand{\Gdom}{\mathfrak{G}}
\providecommand{\Lsys}{\mathfrak{L}}
\providecommand{\Omegaint}{[\Omega]}
\providecommand{\eps}{\varepsilon}
% BEGIN UNIT BODY
\section*{\S\ 13. Примеры относительно целых областей первого рода\\ и регулярных систем.}

1. В качестве первого примера относительно целой области первого рода $\Gdom_{nn}$ рассмотрим \emph{совокупность полиномов от $x_1,\ldots,x_n$, допускающих заданную группу перестановок $G$}, то есть совокупность полиномов лагранжевой области родов. Так как $x_1,\ldots,x_n$ в силу тождества
\begin{equation}
  x_k^n-\sigma_1(x)x_k^{n-1}+\sigma_2(x)x_k^{n-2}+\cdots\pm\sigma_n(x)=0
  \qquad(k=1,2,\ldots,n)
\tag{1}
\end{equation}
алгебраически целы над элементарными симметрическими функциями, принадлежащими $\Gdom_{nn}$, и потому по определению V алгебраически цело зависят от $\Gdom_{nn}$, область $\Gdom_{nn}$ имеет первый род. Далее, как область первого рода, поскольку ее элементы являются однородными формами и $n=\rho$, $\Gdom_{nn}$ образует (по заключению \S\ 10) регулярную систему. Действительно, в силу (1) результант $R_0$ элементарных симметрических функций не может тождественно обращаться в нуль; по правилам теории результантов он вычисляется как $\pm1$. \emph{Инволюционная форма} области $\Gdom_{nn}$, поскольку $\Gdom_{nn}$ имеет первый род и $n=\rho$, однозначно определяется как инволюционная форма соответствующей лагранжевой области родов. Величины, сопряженные с $x_1,\ldots,x_n$ относительно этой области, получаются перестановками группы $G$; следовательно, эта инволюционная форма задается равенством
\begin{equation}
\begin{aligned}
  \Psi(z,u)&=\prod_i
  (z-u_1x_{i_1}-u_2x_{i_2}-\cdots-u_nx_{i_n})\\
  &=z^i+\sum{}'G_{\alpha\alpha_1\ldots\alpha_n}(x)
  z^\alpha u_1^{\alpha_1}\cdots u_n^{\alpha_n},
\end{aligned}
\tag{2}
\end{equation}
где произведение берется по всем перестановкам из $G$. Формула (2) задает ``форму Галуа группы $G$'', и потому верен следующий факт:

\emph{Все полиномы от $x_1,\ldots,x_n$, допускающие группу перестановок $G$, являются целыми рациональными комбинациями коэффициентов $G_{\alpha\alpha_1\ldots\alpha_n}(x)$ формы Галуа этой группы.}

2. В качестве примера \emph{регулярных систем} рассмотрим системы $\Ssys_{n1}$ из полиномов, то есть системы, содержащие только один алгебраически независимый полином.

Действительно, пусть $\Ssys^*_{11}$ будет системой отображения для $\Ssys_{n1}$, а
\[
  F^*(x)=a_0x^k+a_1x^{k-1}+\cdots+a_k\qquad(a_0\ne0)
\]
будет полиномом из $\Ssys^*_{11}$. Тогда
\begin{equation}
  R_0=a_0\ne0.
\tag{3}
\end{equation}
Следовательно, $\Ssys_{n1}$ регулярна по определению VIII, и потому:

\emph{Каждая система $\Ssys_{n1}$ из полиномов, в частности каждая система из бесконечно многих полиномов одной неопределенной, обладает интегральным базисом.}\footnote{Самые простые возможные нерегулярные системы без интегрального базиса являются, таким образом, системами $\Ssys_{22}$. Что такие системы действительно существуют, показывают пример Гильберта и пример, приведенный в примечании к \S\ 12.}

3. Теперь специализируем системы $\Ssys_{n1}$ до относительно целых областей $\Gdom_{n1}$, которые как регулярные системы, по \S\ 12, являются относительно целыми областями первого рода. Их интегральный базис поэтому задается инволюционным базисом, который одновременно является инволюционным базисом наименьшего содержащего поля $\Kfield_{n1}$. По \S\ 6 любая функция инволюционного базиса $\Kfield_{n1}$ одновременно является минимальным базисом --- мы называли ее функцией Люрота этого поля --- и всякая функция инволюционного базиса линейно-дробно зависит от этой функции Люрота. В нашем случае функция Люрота, поскольку она принадлежит $\Gdom_{n1}$, является полиномом; всякий другой полином инволюционного базиса линейно-дробно зависит от нее, а в силу (3) эта зависимость также алгебраически цела, то есть фактически цела и линейна. Тем самым функция Люрота $L(x)$ становится интегральным базисом. Если еще положить $u_1=1$, инволюционная форма $\Gdom^*_{11}$ имеет вид
\[
  \Psi^*(z)=L^*(z)-L^*(x),
\]
откуда снова видно, что $L(x)$ является интегральным базисом. Следовательно:

\emph{Относительно целые области $\Gdom_{n1}$ имеют первый род; их интегральный базис состоит из одного полинома, функции Люрота наименьшего содержащего поля.}

4. В качестве еще одного примера относительно целой области $\Gdom_{n1}$ упомянем, в частности, целые рациональные проективные инварианты квадратичной формы от $s$ переменных. В согласии с предыдущим, они, как известно, представляются как целые рациональные комбинации детерминанта $|a_{ik}|$ этой формы; и этот детерминант одновременно является функцией Люрота соответствующего поля инвариантов.
% END UNIT BODY
% END PRODUCER UNIT 053
\endgroup


\clearpage
\begingroup
\setcounter{footnote}{0}
% BEGIN PRODUCER UNIT 054
% Source: C:/Users/Floris/Documents/interlanguage/03_projects/noether/02_slavic_working_corpus/translations/paper06/russian/v001/Noether_Paper06_Section14_Russian_v001.tex
% Identity: 10727 B / SHA-256 754F4A3C925E3E6E5AE36198DDB36305206D0F957BA4BD220B54A1B983878761
\setlength{\parindent}{1.25em}
\setlength{\parskip}{0pt}
\emergencystretch=10em
\allowdisplaybreaks
\providecommand{\Kfield}{\mathfrak{K}}
\providecommand{\Ssys}{\mathfrak{S}}
\providecommand{\Jdom}{\mathfrak{J}}
\providecommand{\Gdom}{\mathfrak{G}}
\providecommand{\Lsys}{\mathfrak{L}}
\providecommand{\Omegaint}{[\Omega]}
\providecommand{\eps}{\varepsilon}
% BEGIN UNIT BODY
\section*{\S\ 14. Системы из целочисленных полиномов.}

Полученные теоремы о базисах теперь нужно усилить в отношении целочисленности.

Поэтому далее коэффициентную область $\Omega$ будем считать (конечным) алгебраическим числовым полем; через $[\Omega]$ обозначим совокупность алгебраически целых чисел из $\Omega$, которые для краткости будем называть целыми числами. Под целочисленным полиномом будем понимать полином с коэффициентами из $[\Omega]$; отныне $F(x),G(x),\ldots$ будут означать только целочисленные полиномы.

По аналогии с \S\ 7 рассмотрим различные системы из целочисленных полиномов:

Произвольную систему из целочисленных полиномов будем называть целочисленной системой $\Ssys_{n\rho}$.

Целочисленное линейное семейство $\Lsys_{n\rho}$ есть целочисленная система, которая вместе с $F_1(x)$ и $F_2(x)$ всегда содержит также $c_1F_1(x)+c_2F_2(x)$, где $c_1,c_2$ --- произвольные целые числа из $[\Omega]$.

Целочисленная область целостности $\Jdom_{n\rho}$ есть целочисленное линейное семейство, которое вместе с $F(x)$ и $G(x)$ всегда содержит также $F(x)\cdot G(x)$.

Целочисленная относительно целая область $\Gdom_{n\rho}$ есть область целостности, которая вместе с $F(x)$ и $G(x)$ --- где $G(x)\ne0$ --- тогда и только тогда содержит частное $F(x):G(x)$, когда это частное является целочисленным полиномом.

Так же, как в \S\ 7, наименьшая целочисленная область целостности $\Jdom_{n\rho}$, содержащая целочисленную систему $\Ssys_{n\rho}$, состоит из всех полиномов $H(x)$, которые являются целыми рациональными целочисленными комбинациями конечного числа полиномов из $\Ssys_{n\rho}$;

наименьшая целочисленная относительно целая область $\Gdom_{n\rho}$, содержащая ее, состоит из всех целочисленных полиномов $H(x)$, которые являются рациональными комбинациями конечного числа полиномов из $\Ssys_{n\rho}$.

Теперь перейдем к переносу теорем о базисах. В отношении рационального базиса нужно отметить, что понятия поля, наименьшего содержащего поля и рационального базиса, основанные только на ``рациональном'', не изменяются. Кроме того, систему отображения $\Ssys^*_{\rho\rho}$ целочисленной системы $\Ssys_{n\rho}$ можно снова выбрать целочисленной системой сколь угодно многими способами. Так как и для целочисленного линейного семейства сохраняются рассуждения, ведущие к ограничению числа, то верно:

\emph{Теорема V$'$: Каждая целочисленная система $\Ssys_{n\rho}$ обладает рациональным базисом; рациональный базис целочисленного линейного семейства $\Lsys_{n\rho}$ можно, в частности, выбрать так, чтобы он содержал не более $\rho+1$ полиномов.}

Теперь исследуем аналог инволюционного базиса областей целостности из полиномов (\S\ 8). Для этого теорему о \emph{симметрических функциях рядов величин} нужно расширить в отношении \emph{целочисленности} следующим

\emph{вспомогательным утверждением. Целые целочисленные симметрические функции от $n$ рядов величин $(yz\cdots t)$:}
\begin{equation}
\begin{array}{cccc}
  y_1&z_1&\cdots&t_1\\
  y_2&z_2&\cdots&t_2\\
  \vdots&\vdots&&\vdots\\
  y_n&z_n&\cdots&t_n
\end{array}
\tag{1}
\end{equation}
\emph{являются целыми целочисленными функциями конечного числа среди них, а именно элементарных симметрических функций}
\begin{equation}
\begin{gathered}
  \sigma_1(y),\sigma_2(y),\ldots,\sigma_n(y);\quad
  \sigma_1(z),\sigma_2(z),\ldots,\sigma_n(z);\quad\ldots;\\
  \sigma_1(t),\ldots,\sigma_n(t);
\end{gathered}
\tag{2}
\end{equation}
\emph{и функций}
\begin{equation}
  \chi_1(y,z,\ldots,t),\ldots,\chi_k(y,z,\ldots,t),
\tag{3}
\end{equation}
\emph{которые алгебраически цело и целочисленно зависят от функций (2).}

Действительно, $n$ величин $y_1,\ldots,y_n$ алгебраически цело и целочисленно зависят от $\sigma_1(y),\ldots,\sigma_n(y)$; соответствующее утверждение верно для $z_1,\ldots,z_n$ в отношении $\sigma_1(z),\ldots,\sigma_n(z)$ и т.д. Следовательно, поле симметрических функций рядов величин (1) образует алгебраическое поле над полем, определенным функциями (2), так, что алгебраически целые и целочисленные величины этого поля тождественны целым целочисленным симметрическим функциям рядов величин (1). Поэтому, если выбрать (3) как кронекерову фундаментальную систему, вспомогательное утверждение доказано.

Пусть теперь $\Jdom_{n\rho}$ есть целочисленная область целостности, $\Jdom^*_{\rho\rho}$ --- целочисленная область отображения, и
\[
  \Psi^*(z,u)=G^*(x)z^i+
  \sum{}'G^*_{\alpha\alpha_1\ldots\alpha_\rho}(x)
  z^\alpha u_1^{\alpha_1}\cdots u_\rho^{\alpha_\rho}
\]
есть инволюционная форма $\Jdom^*_{\rho\rho}$; при этом $G^*(x)$ может также иметь нулевую размерность, то есть быть числом из $[\Omega]$. Элементарные симметрические функции, соответствующие функциям (2), тогда задаются некоторыми из частных
\begin{equation}
  G^*_{\alpha\alpha_1\ldots\alpha_\rho}(x):G^*(x).
\tag{4}
\end{equation}
Так как по вспомогательному утверждению функции (3) алгебраически цело и целочисленно зависят от (2), то --- по расширению теоремы Гаусса на полиномы с алгебраически целочисленными коэффициентами --- функции из $\Kfield^*_{\rho\rho}$, соответствующие (3), задаются в виде
\begin{equation}
  K_\nu^*(x):(G^*(x))^\sigma,
\tag{5}
\end{equation}
где $K_\nu^*(x)$ --- целочисленные полиномы из $\Jdom^*_{\rho\rho}$. Следовательно, если речь идет об относительно целых областях, $K_\nu^*(x)$ принадлежат области; вообще же они являются частными полиномов из $\Jdom^*_{\rho\rho}$.

Если теперь $F^*(x)$ есть произвольный целочисленный полином из $\Jdom^*_{\rho\rho}$, то имеем
\[
  F^*(x)=\frac1i\Bigl(F^*(x)+F^*(x^{(1)})+\cdots+F^*(x^{(i-1)})\Bigr)
\]
и поэтому $i\cdot F^*(x)$ становится целой целочисленной симметрической функцией рядов величин $x,x^{(1)},\ldots,x^{(i-1)}$.

Следовательно, по вспомогательному утверждению и принципу переноса верно:

\emph{Теорема VI$'$: Для целочисленных областей целостности $\Jdom_{n\rho}$ каждая инволюционная форма определяет выделенный рациональный базис так, что каждый полином $F(x)$ из $\Jdom_{n\rho}$ допускает представление}
\[
  F(x)=\frac{\Gamma(H(x),H_1(x),\ldots,H_s(x))}{i\cdot H(x)^i},
\]
\emph{где $\Gamma$ обозначает целую целочисленную функцию. Для целочисленных относительно целых областей $\Gdom_{n\rho}$, у которых целочисленная относительно целая область $\Gdom^*_{\rho\rho}$ является областью отображения, знаменатель $H(x)$ задается старшим коэффициентом соответствующей инволюционной формы.}
% END UNIT BODY
% END PRODUCER UNIT 054
\endgroup


\clearpage
\begingroup
\setcounter{footnote}{0}
% BEGIN PRODUCER UNIT 055
% Source: C:/Users/Floris/Documents/interlanguage/03_projects/noether/02_slavic_working_corpus/translations/paper06/russian/v001/Noether_Paper06_Section15_Russian_v001.tex
% Identity: 9375 B / SHA-256 6E1CBF8B7014C790B7138B04F605F7B658D0FBBDEAD7971399A976F07BDA4479
\setlength{\parindent}{1.25em}
\setlength{\parskip}{0pt}
\emergencystretch=10em
\allowdisplaybreaks
\providecommand{\Kfield}{\mathfrak{K}}
\providecommand{\Ssys}{\mathfrak{S}}
\providecommand{\Jdom}{\mathfrak{J}}
\providecommand{\Gdom}{\mathfrak{G}}
\providecommand{\Lsys}{\mathfrak{L}}
\providecommand{\Omegaint}{[\Omega]}
\providecommand{\eps}{\varepsilon}
% BEGIN UNIT BODY
\section*{\S\ 15. Целочисленные относительно целые области первого рода и регулярные системы.}

Вместо интегрального базиса для целочисленных систем выступает целочисленный интегральный базис, то есть конечное число целочисленных полиномов системы такое, что каждый целочисленный полином системы представляется как целая целочисленная комбинация этого конечного числа.

Теоремы о целочисленном интегральном базисе в существенном идут параллельно прежним, и только доказательства требуют небольших изменений.

Сначала определение V (\S\ 9) непосредственно переносится.

\emph{Определение $\mathrm{V}'$: величина $\xi$ называется алгебраически целой и целочисленной относительно целочисленной относительно целой области $\Gdom_{n\rho}$, если она удовлетворяет некоторому уравнению, старший коэффициент которого является единицей из $[\Omega]$, а остальные полиномы принадлежат $\Gdom_{n\rho}$.}

Действительно, расширение теоремы Гаусса на полиномы с алгебраически-целочисленными коэффициентами показывает, как и в \S\ 9, что отсюда можно получить определение на неприводимом уравнении.

Из определения $\mathrm{V}'$ далее получается перенос определения VII.

\emph{Определение $\mathrm{VII}'$: целочисленная относительно целая область $\Gdom_{n\rho}$ называется областью первого рода, если она имеет областью отображения относительно целую область $\Gdom^*_{\rho\rho}$ такую, что каждый произвольный целочисленный полином от $x_1,\ldots,x_\rho$ алгебраически цело и целочисленно зависит от $\Gdom^*_{\rho\rho}$.}

Чтобы показать конечность этих областей первого рода, сначала заметим, что, как в \S\ 10, из определения следует: старший коэффициент инволюционной формы является единицей из $[\Omega]$.

Следовательно, по теореме $\mathrm{VI}'$ (\S\ 14) каждый полином $F(x)$ из $\Gdom_{n\rho}$ допускает представление:
\[
  F(x)=\frac{\Gamma(H_1(x),\ldots,H_\sigma(x))}{i}.
\]
Если здесь к системе, образованной из целочисленных полиномов $\Gamma(H_1,\ldots,H_\sigma)$, применить гильбертову теорему II о целочисленном модульном базисе (Math. Ann. 36), сформулированную для целых рациональных числовых коэффициентов, в ее распространении на алгебраически-целочисленные полиномы\footnote{Если $w_1,\ldots,w_k$ обозначает фундаментальную систему алгебраически целых чисел из $\Omega$, и если положить
\[
\begin{gathered}
  \Gamma(H_1,\ldots,H_\sigma)
  =\Gamma_1(H)w_1+\cdots+\Gamma_k(H)w_k,
\end{gathered}
\]
то $\Gamma_1(H),\ldots,\Gamma_k(H)$ имеют целые рациональные числовые коэффициенты. Применение теоремы II к системе, образованной из форм $v_1\Gamma_1(H)+\cdots+v_k\Gamma_k(H)$, непосредственно показывает справедливость этого распространения.}, то отсюда выясняется существование целочисленного интегрального базиса, то есть:

\emph{Теорема $\mathrm{VII}'$: каждая целочисленная относительно целая область первого рода обладает целочисленным интегральным базисом.}

Далее переносим определение регулярных систем:

\emph{Определение $\mathrm{VIII}'$: целочисленная система $\Ssys_{n\rho}$ называется целочисленно-регулярной, если в области отображения $\Jdom^*_{\rho\rho}$ наименьшей содержащей целочисленной области целостности существуют $\rho$ полиномов такие, что результант членов наивысшей размерности является единицей из $[\Omega]$, то есть $R_0=\eps$.}

Чтобы перенести доказательство существования интегрального базиса, заметим, что вспомогательная теорема из \S\ 11 допускает следующую более острую формулировку; она получается из уравнений (5) и (6) в \S\ 11 и из того обстоятельства, что $A_1(\lambda),\ldots,A_\tau(\lambda)$ являются целыми целочисленными функциями от $\lambda_1,\ldots,\lambda_\rho$:

\emph{Достаточное условие для того, чтобы каждый произвольный целочисленный полином от $x_1,\ldots,x_\rho$ алгебраически цело и целочисленно зависел от $\rho$ целочисленных полиномов, состоит в том, что результант членов наивысшей размерности этих полиномов является единицей из $[\Omega]$, то есть $R_0=\eps$.}

Из этой вспомогательной теоремы точно так же, как в \S\ 12, для каждого полинома $F(x)$ целочисленно-регулярной системы следует представление (3) из \S\ 12, где теперь $\Gamma_1(G_i),\ldots,\Gamma_k(G_i)$ являются целочисленными полиномами от $G_i$. И далее, как в \S\ 12, применение гильбертовой теоремы II (Math. Ann. 36) в ее распространении на алгебраически-целочисленные полиномы, вместе с определением наименьшей содержащей целочисленной области целостности, дает существование целочисленного интегрального базиса; то есть:

\emph{Теорема $\mathrm{IX}'$: каждая целочисленно-регулярная система обладает целочисленным интегральным базисом.}

В качестве примера целочисленной относительно целой области первого рода, по аналогии с примером 1 в \S\ 13, укажем на совокупность $\Gdom_{nn}$ целочисленных полиномов лагранжевой родовой области. То, что $\Gdom_{nn}$ имеет целочисленно-первый род, следует из тождества:
\[
  x_k^n-\sigma_1(x)x_k^{n-1}+\cdots\pm\sigma_n(x)=0
  \qquad(k=1,2,\ldots,n);
\]
так как результант элементарных симметрических функций имеет значение $\pm1$, $\Gdom_{nn}$ является также целочисленной регулярной системой.

Еще один пример, по вспомогательной теореме из \S\ 14, дает совокупность целых целочисленных симметрических функций от $n$ рядов величин.

Эрланген, май 1914.
% END UNIT BODY
% END PRODUCER UNIT 055
\endgroup


\clearpage
\begingroup
\setcounter{footnote}{0}
% BEGIN PRODUCER UNIT 056
% Source: C:/Users/Floris/Documents/interlanguage/03_projects/noether/02_slavic_working_corpus/translations/paper07/russian/v001/Noether_Paper07_Russian_v001.tex
% Identity: 12700 B / SHA-256 146077299BC566EBF544A6DD6F97F6DD5D8AB91F5DE20CAC7895E0DAB8BA4585
\setlength{\parindent}{1.25em}
\setlength{\parskip}{0pt}
\emergencystretch=10em
\allowdisplaybreaks
\providecommand{\Hgrp}{\mathfrak{H}}
\providecommand{\GaloisRes}{\Phi}
\providecommand{\pol}[2]{\left(#1\frac{\partial}{\partial #2}\right)}
\providecommand{\Deltaop}{\Delta\!\left(\frac{\partial}{\partial x},\frac{\partial}{\partial y},\frac{\partial}{\partial z}\right)}
\providecommand{\Omegaop}{\Omega}
% BEGIN UNIT BODY
\begin{center}
{\Large\bfseries 7. Теорема конечности для инвариантов конечных групп}\par
\vspace{1.2em}
Math. Ann. 77 (1916), с. 89--92
\end{center}

\vspace{1em}

Ниже будет дано вполне элементарное доказательство конечности инвариантов конечных групп, основанное только на теории симметрических функций и одновременно дающее фактическое указание полной системы. Обычное же доказательство, опирающееся на гильбертову теорему о модульном базисе (Ann. 36), является только доказательством существования.\footnote{Ср., например, Weber, \emph{Lehrbuch der Algebra} (2-е изд.), т. 2, \S\ 57.}

Пусть конечная группа $\Hgrp$ состоит из $h$ линейных преобразований (с ненулевым детерминантом) $A_1,\ldots,A_h$, причем $A_k$ задает линейное преобразование
\[
  x_1^{(k)}=\sum_{\nu=1}^{n}a_{1\nu}^{(k)}x_\nu,\ldots,
  x_n^{(k)}=\sum_{\nu=1}^{n}a_{n\nu}^{(k)}x_\nu,
\]
или, сокращенно, $(x^{(k)})=A_k(x)$. Следовательно, группа $\Hgrp$ переводит ряд $(x)$ с элементами $x_1,\ldots,x_n$ в ряды $(x^{(k)})$ с элементами $x_1^{(k)},\ldots,x_n^{(k)}$. Так как среди $A_1,\ldots,A_h$ должна содержаться тождественность, среди рядов $(x^{(k)})$ содержится также ряд $(x)$. Под целым рациональным (абсолютным) инвариантом группы понимается такая целая рациональная функция от $x_1,\ldots,x_n$, которая при применении $A_1,\ldots,A_h$ остается тождественно неизменной. Поэтому для такого инварианта $f(x)$ имеем
\begin{equation}
  f(x)=f(x^{(1)})=\cdots=f(x^{(h)})=
  \frac1h\sum_{k=1}^{h}f(x^{(k)}).
\tag{1}
\end{equation}

\subsection*{1.}
Формула (1) выражает, что $f(x)$ есть целая рациональная симметрическая функция рядов величин $(x^{(1)}),\ldots,(x^{(h)})$. При этом, поскольку каждое слагаемое $f(x^{(k)})$ содержит только один ряд $(x^{(k)})$, речь идет о простейшем случае, обычно называемом \emph{одноформным}. По известной теореме о симметрических функциях рядов величин\footnote{Ср. примечание к 2.} функция $f(x)$ цело и рационально выражается через элементарные симметрические функции этих рядов, то есть через коэффициенты $G_{\alpha\alpha_1\ldots\alpha_n}(x)$ ``резольвенты Галуа''
\[
\begin{aligned}
  \Phi(z,u)
  &=\prod_{k=1}^{h}\bigl(z+u_1x_1^{(k)}+\cdots+u_nx_n^{(k)}\bigr) \\
  &=z^h+
    \sum G_{\alpha\alpha_1\ldots\alpha_n}(x)
    z^{\alpha}u_1^{\alpha_1}\cdots u_n^{\alpha_n},
\end{aligned}
\qquad
\left(\begin{array}{c}
\alpha+\alpha_1+\cdots+\alpha_n=h\\
\alpha\ne h
\end{array}\right),
\]
где $G_{\alpha\alpha_1\ldots\alpha_n}(x)$ являются инвариантами степени $\alpha_1+\cdots+\alpha_n$ по переменным $x$. Тем самым доказано:

\emph{Коэффициенты $G_{\alpha\alpha_1\ldots\alpha_n}(x)$ резольвенты Галуа образуют полную систему инвариантов группы в том смысле, что каждый инвариант можно цело и рационально выразить через эти конечные многие инварианты.}

\subsection*{2.}
Исходя из (1), можно провести и следующее еще более элементарное рассмотрение,\footnote{Оно сводится к доказательству теоремы о симметрических функциях рядов величин в указанном выше одноформном случае.} которое одновременно приводит ко второй полной системе. Пусть
\[
  f(x)=a+b x_1^{\mu_1}\cdots x_n^{\mu_n}
       +\cdots+c x_1^{\nu_1}\cdots x_n^{\nu_n},
\]
где $a,b,\ldots,c$ обозначают константы. Тогда по (1) получаем
\[
\begin{aligned}
 h\cdot f(x)
 &=h\cdot a+b\sum_{k=1}^{h}(x_1^{(k)})^{\mu_1}\cdots(x_n^{(k)})^{\mu_n}
     +\cdots
     +c\sum_{k=1}^{h}(x_1^{(k)})^{\nu_1}\cdots(x_n^{(k)})^{\nu_n} \\
 &=h\cdot a+b\cdot J_{\mu_1\ldots\mu_n}+\cdots+c\cdot J_{\nu_1\ldots\nu_n}.
\end{aligned}
\]
Значит, каждый инвариант можно составить как целую линейную комбинацию специальных инвариантов
\[
  J_{\mu_1\ldots\mu_n}
  =\sum_{k=1}^{h}(x_1^{(k)})^{\mu_1}\cdots(x_n^{(k)})^{\mu_n},
\]
и потому достаточно доказать утверждение для этих специальных инвариантов. Но $J_{\mu_1\ldots\mu_n}$, с точностью до числового множителя, является коэффициентом при $u_1^{\mu_1}\cdots u_n^{\mu_n}$ в выражении
\[
  S_\mu=\sum_{k=1}^{h}
  \bigl(u_1x_1^{(k)}+\cdots+u_nx_n^{(k)}\bigr)^{\mu},
\]
где положено $\mu=\mu_1+\cdots+\mu_n$, и это выражение есть $\mu$-я степенная сумма $h$ линейных форм
\[
  \xi_1=u_1x_1^{(1)}+\cdots+u_nx_n^{(1)},\ldots,
  \xi_h=u_1x_1^{(h)}+\cdots+u_nx_n^{(h)}.
\]
Как известно, бесконечно многие степенные суммы $S_\mu$ являются целыми рациональными комбинациями
\[
  S_1,S_2,\ldots,S_h,
\]
коэффициенты которых задаются инвариантами
\[
  J_{\mu_1\ldots\mu_n}\qquad(\mu_1+\cdots+\mu_n\le h).
\]
Тем самым получена вторая полная система:

\emph{Полная система инвариантов группы задается всеми инвариантами $J_{\mu_1\ldots\mu_n}$, для которых $\mu_1+\cdots+\mu_n\le h$, где $h$ означает порядок группы.}

Связь с 1. устанавливается замечанием, что степенные суммы $S_1,\ldots,S_h$ можно заменить элементарными симметрическими функциями от $\xi_1,\ldots,\xi_h$, что приводит к данной там полной системе $G_{\alpha\alpha_1\ldots\alpha_n}(x)$. Оба результата показывают, что все инварианты можно цело и рационально выразить через те из них, степень которых по переменным $x$ не превосходит порядка $h$ группы.

\subsection*{3.}
Из только что доказанного можно вывести следствия для рационального представления. Каждый рациональный абсолютный инвариант можно, как непосредственно видно из домножения числителя и знаменателя на сопряженные знаменателя относительно группы, представить в виде частного двух целых рациональных абсолютных инвариантов, не обязательно взаимно простых. Отсюда следует, что каждый рациональный абсолютный инвариант можно рационально выразить через коэффициенты $G_{\alpha\alpha_1\ldots\alpha_n}(x)$ резольвенты Галуа $\Phi(z,u)$. Эту теорему можно легко доказать разными способами и без прохода через теорему конечности; она уже сформулирована у Weber II, \S\ 58.\footnote{Однако данное там доказательство не является состоятельным: показано лишь, что функция $\Psi(t)$ в формуле (7) имеет инварианты коэффициентами, но не показано, что эти инварианты рационально выражаются через коэффициенты $\Phi(t)$. Этот пробел устраняется, если для представления $x_i^{(k)}$ через $\xi_k$ вместо интерполяционной формулы Лагранжа применить известную процедуру дифференцирования. Это соотношение, полученное дифференцированием тождества $\Phi(-\xi_k,u)=0$ по всем $u$:
\[
 \left[\frac{\partial\Phi}{\partial u_i}
      -x_i^{(k)}\frac{\partial\Phi}{\partial z}\right]_{z=-\xi_k}=0.
\]
Соответственно, вместо формул (7) и (8) у Weber для каждой рациональной функции $\omega(x)$ должна стоять формула
\[
 \omega(x_1^{(k)},\ldots,x_n^{(k)})=
 \left.
 \omega\!\left(
   \frac{\partial\Phi/\partial u_1}{\partial\Phi/\partial z},\ldots,
   \frac{\partial\Phi/\partial u_n}{\partial\Phi/\partial z}
 \right)\right|_{z=-\xi_k},
\]
которая содержит уже только коэффициенты $\Phi(z,u)$ и суммирование которой по всем $k$ для инвариантов $\omega(x)$ дает требуемое представление.}

\subsection*{4.}
Наконец, следует упомянуть, что выведенные здесь результаты неявно содержатся в моей работе ``Поля и системы рациональных функций'';\footnote{Math. Ann. 76, с. 161 (1915).} именно, полная система, данная в 1., содержится в теоремах VI и VII, а доказательство рационального представления, независимое от этой теоремы конечности, содержится в теореме III.\footnote{Для относительных инвариантов теоремы VIII и IX содержат доказательство конечности, которое также основано на иной основе, чем обычное, но тоже является только доказательством существования; причина в том, что относительные инварианты не образуют поля.} В данном специальном случае ход доказательства и результат, однако, существенно упрощаются по сравнению с общей теорией. К применению общих исследований к инвариантам конечных групп меня привели беседы с Э. Фишером; от него же по существу происходят и доказательство, данное в 2. (в общем случае возникающая там полная система не является рационально определимой), и примечание к 3.

Эрланген, май 1915 г.
% END UNIT BODY
% END PRODUCER UNIT 056
\endgroup


\clearpage
\begingroup
\setcounter{footnote}{0}
% BEGIN PRODUCER UNIT 057
% Source: C:/Users/Floris/Documents/interlanguage/03_projects/noether/02_slavic_working_corpus/translations/paper08/russian/v001/Noether_Paper08_Russian_v001.tex
% Identity: 35100 B / SHA-256 5E572F40A1137FED84435F198C7161703454D279DAF3A41763C76BEFECEB5881
\setlength{\parindent}{1.25em}
\setlength{\parskip}{0pt}
\emergencystretch=14em
\allowdisplaybreaks
\providecommand{\vdotswithin}[1]{\mathrel{\vdots}}
% BEGIN UNIT BODY
\begin{center}
{\Large\bfseries 8. О целом рациональном представлении инвариантов системы произвольно многих основных форм}\par
\vspace{1.2em}
Math. Ann. 77 (1916), с. 93--102
\end{center}

\vspace{1em}

В конце своего вклада в юбилейный сборник в честь Шварца ``Об инвариантах системы произвольно многих основных форм'' Д. Гильберт высказывает предположение, что эти инварианты можно цело и рационально представить через конечное число инвариантов системы $(J,PJ)$, где $J$ обозначает полную систему инвариантов для линейно независимых основных форм, а инварианты $PJ$ получаются из $J$ полярными процессами.

Ниже, в I, я показываю, что это предположение действительно верно, а именно как простое следствие редукционной теоремы: всякую форму от произвольно многих рядов по $n$ переменных в каждом можно представить как сумму поляр форм, содержащих только $n$ фиксированных рядов и сами выведенных из исходной формы полярными процессами. Эта редукционная теорема есть специальный случай, впервые сформулированный Ф. Мертенсом, обобщения разложения Клебша--Гордана по рядам на формы от произвольно многих рядов по $n$ переменных; это обобщение, данное А. Капелли, Ф. Мертенсом и Ж. Дерейтсом, доказано формальным преобразованием процессов дифференцирования.\footnote{F. Mertens: ``Über eine Formel der Determinantentheorie.'' (формула 5) Sitzb. d. Ak. d. Wiss. Wien, Bd. 91, II. Abt. (1885), и подробнее (с приложениями) в: ``Über ganze Funktionen von $m$ Systemen von je $n$ Unbestimmten''. Monatsh. f. Math. u. Phys. 4. Jahrg. (1893). Для некоторых доказательств Капелли (с 1882 г.) ср., например: Math. Ann. 37 или автографированные ``Lezioni sulla teoria delle forme algebriche'' (1902). J. Deruyts: \emph{Essai d'une theorie générale des formes algébriques}, Mém. d. 1. Soc. Roy. des Sciences de Liège (2) 17 (1892).}

В II я даю также более понятийное доказательство редукционной теоремы, рассматривая вопрос как проблему эквивалентности линейных семейств форм. Этот взгляд, равно как и существенную часть использованных соображений, я беру из работы ``Об алгебраических модульных системах''\footnote{E. Fischer: Über algebraische Modulsysteme und lineare homogene partielle Differentialgleichungen mit konstanten Koeffizienten, \S\ 1--4, J. f. M. 140 (1911).} и из примыкающих к ней неопубликованных теорем Э. Фишера, пользоваться которыми он любезно мне разрешил.

Наконец, в III я показываю, как соответствующие соображения приводят также к самым общим разложениям по рядам.

\subsection*{I.}
Пусть $\theta$ есть форма, однородная в каждом из $N>n$ рядов
\[
  A_1,A_2,\ldots,A_N,
\]
где вообще ряд $A_k$ состоит из $n$ элементов
\[
  A_k^{(1)},A_k^{(2)},\ldots,A_k^{(n)}.
\]
Коэффициенты $\theta$ могут быть неопределёнными или величинами заданной области рациональности. Далее пусть $P$ обозначает целую рациональную функцию -- с рациональными целыми коэффициентами -- от полярных процессов
\[
  P_{hk}=\left(A_h\frac{\partial}{\partial A_k}\right)
  =A_h^{(1)}\frac{\partial}{\partial A_k^{(1)}}+
   A_h^{(2)}\frac{\partial}{\partial A_k^{(2)}}+
   \cdots+
   A_h^{(n)}\frac{\partial}{\partial A_k^{(n)}},
\]
причём под произведением $P_{hk}P_{ij}\theta$ следует понимать применение $P_{hk}$ к $P_{ij}\theta$.

Тогда упомянутая редукционная теорема задаётся тождеством
\begin{equation}
  \theta=\sum PZ,
\tag{1}
\end{equation}
где формы $Z$ содержат уже только ряды
\[
  A_1,A_2,\ldots,A_n
\]
и выведены из $\theta$ полярными процессами $P$.

Рассмотрим теперь систему основных форм $(F')$, состоящую из $N$ форм одного порядка и с одним и тем же числом переменных
\[
  F_1,F_2,\ldots,F_N,
\]
коэффициенты которых соответственно берутся как $N$ рядов
\[
  A_1,A_2,\ldots,A_N.
\]
Пусть $(J)$ обозначает полную систему -- состоящую из конечного числа инвариантов -- форм $F_1,\ldots,F_n$, такую, что каждый инвариант форм $F_1,\ldots,F_n$ можно цело и рационально выразить через инварианты $(J)$. Тогда гильбертово предположение, подлежащее доказательству, состоит в том, что для симультанных инвариантов $S$ форм $F_1,\ldots,F_N$ полная система задаётся конечным числом инвариантов $(J,PJ)$, где $PJ$ означает все инварианты, выводимые из $J$ полярными процессами $P$.

В самом деле, каждый симультанный инвариант $S$ системы $(F')$ с рациональными целыми коэффициентами является формой, однородной в каждом из $N$ рядов $A_1,\ldots,A_N$, с рациональными целыми коэффициентами; следовательно, к нему применимо тождество (1). Так как применение полярных процессов $P$ к $S$, как известно, снова даёт инварианты, формы $Z$ также являются инвариантами, именно симультанными инвариантами форм $F_1,\ldots,F_n$, поскольку они содержат только ряды $A_1,\ldots,A_n$; значит, по предположению, они являются целыми рациональными функциями инвариантов $(J)$. Поэтому тождество (1) переходит в
\[
  S=\sum PY,
\]
где $Y$ обозначают произведения степеней инвариантов $(J)$. Но фактическое выполнение процессов $P$ над $Y$, по определению $P$, состоит в последовательном, конечное число раз повторённом выполнении простых полярных операций $P_{hk}$, которые по правилу дифференцирования произведений приводят к суммам произведений инвариантов $(J)$ и $(PJ)$. То же относится к применению $P_{hk}$ к произведению из $(J)$ и $(PJ)$; следовательно, и $PY$ даёт только целые рациональные комбинации из $(J,PJ)$. \emph{Тем самым доказана целая рациональная представимость всех инвариантов $S$ через $(J,PJ)$.}

\subsection*{II.}
Перед доказательством редукционной теоремы напомним некоторые факты о линейных семействах форм. Под \emph{линейным семейством форм над $K$} мы понимаем систему форм одной и той же степени, полную в том смысле, что вместе с $\theta_1$ и $\theta_2$ она всегда содержит также $c_1\theta_1+c_2\theta_2$, где $c_1,c_2$ -- величины предписанной области рациональности $K$, а $K$ содержит область коэффициентов форм $\theta$. Среди этих форм всегда есть конечное число -- скажем, $\rho$ -- линейно независимых над $K$, тогда как любые $\rho+1$ форм удовлетворяют линейному соотношению с коэффициентами из $K$; семейство имеет \emph{ранг $\rho$} относительно $K$.\footnote{Более общие линейные семейства форм -- ранг которых не обязан быть конечным -- получаются, если отбросить ограничение одинаковой степени или также ограничение, что $K$ содержит область коэффициентов форм $\theta$.} Мы будем пользоваться легко усматриваемым фактом: если $\mathcal L$ и $\mathcal T$ суть два линейных семейства над $K$ одного и того же ранга $\rho$, и $\mathcal T$ является подсемейством $\mathcal L$, то $\mathcal L$ и $\mathcal T$ тождественны (эквивалентны).

После этого переходим к редукционной теореме. Тождество из I
\[
  \theta=\sum PZ
\]
переходит в аналогичное тождество для $P\theta$, если к обеим частям применить один и тот же полярный процесс $P$; значит, редукционная теорема одновременно даёт представление всех поляр $\theta$ суммами специальных поляр $PZ$. С другой стороны, эти специальные поляры несомненно содержатся среди всех поляр, так что редукционная теорема утверждает эквивалентность этих двух линейных семейств форм; в частности, и эквивалентность тех подсемейств, формы которых в каждом отдельном ряду имеют ту же степень, что и $\theta$. Для доказательства этой эквивалентности введём промежуточное семейство, определяемое формами $Z$. Ради простоты сначала дадим доказательство в случае трёх бинарных рядов $(N=3,n=2)$, а затем кратко укажем полностью параллельное общее доказательство.

Итак, пусть $\theta(x,y,z)$ имеет соответственно степени $\alpha,\beta,p$ в рядах
\[
  x_1x_2;\quad y_1y_2;\quad z_1z_2;
\]
коэффициенты $\theta$ пока пусть будут неопределёнными $a$ (или также рациональными числами). Тогда рассмотрим следующие три линейных семейства форм.

1. Семейство $\mathcal L$, состоящее из всех форм $f(x,y,z)$, которые получены из $\theta$ полярными процессами $P$ и в отдельных рядах $x,y,z$ имеют ту же степень, что и $\theta$; в частности, $\mathcal L$ содержит и $\theta$. Если рассматривать $f(x,y,z)$ как форму от $x,y,z$ и неопределённых $a$, то область коэффициентов есть поле $R$ рациональных чисел; для $K$ также нужно взять $R$, поскольку только при рационально-целых $\theta_1,\theta_2$ выражение $c_1P_1+c_2P_2$ снова становится процессом $P$; пусть $\mathcal L$ имеет ранг $\rho$ относительно $R$.

2. Семейство $\mathcal S$, состоящее из всех форм $g(\lambda;x,y)$, определённых посредством
\[
  g(\lambda;x,y)=f(x,y,\lambda_1x+\lambda_2y),
\]
или, в выражении через полярные процессы,
\[
\begin{aligned}
  g(\lambda;x,y)
  &=\frac1{p!}\left[(\lambda_1x+\lambda_2y)\frac{\partial}{\partial z}\right]^p f(x,y,z) \\
  &=g_0(xy)\lambda_1^p+g_1(x,y)\lambda_1^{p-1}\lambda_2+\cdots+g_p(xy)\lambda_2^p,
\end{aligned}
\]
причём
\[
  i!(p-i)!\,g_i(xy)=
  \left(y\frac{\partial}{\partial z}\right)^i
  \left(x\frac{\partial}{\partial z}\right)^{p-i}f(xyz).
\]
\footnote{Как и в I,
\[
 \left(t\frac{\partial}{\partial x}\right)=t_1\frac{\partial}{\partial x_1}+t_2\frac{\partial}{\partial x_2},
\]
а значит, в частности,
\[
 \left[(\lambda_1x+\lambda_2y)\frac{\partial}{\partial z}\right]
 = (\lambda_1x_1+\lambda_2y_1)\frac{\partial}{\partial z_1}
 + (\lambda_1x_2+\lambda_2y_2)\frac{\partial}{\partial z_2},
\]
\[
 \left[z\left(\frac{\partial}{\partial\lambda_1}\frac{\partial}{\partial x}
 +\frac{\partial}{\partial\lambda_2}\frac{\partial}{\partial y}\right)\right]
 =z_1\left(\frac{\partial}{\partial\lambda_1}\frac{\partial}{\partial x_1}
 +\frac{\partial}{\partial\lambda_2}\frac{\partial}{\partial y_1}\right)
 +z_2\left(\frac{\partial}{\partial\lambda_1}\frac{\partial}{\partial x_2}
 +\frac{\partial}{\partial\lambda_2}\frac{\partial}{\partial y_2}\right).
\]}
Тем самым $g_i(xy)$ задают формы $Z$ из тождества (1): формы $g$ являются рационально-целыми формами от $x,y,\lambda$ и неопределённых $a$; пусть $\mathcal S$ имеет ранг $\sigma$ относительно $R$.

3. Семейство $\mathcal T$, получающееся из $\mathcal S$ обратным полярным процессом так же, как $\mathcal S$ получается из $\mathcal L$; формы $h$ семейства $\mathcal T$ определены через
\[
\begin{aligned}
 h(x,y,z)
 &=\frac1{p!}\left[z\left(\frac{\partial}{\partial\lambda_1}\frac{\partial}{\partial x}
 +\frac{\partial}{\partial\lambda_2}\frac{\partial}{\partial y}\right)\right]^p g(\lambda;xy)\\
 &=h_0(xyz)+h_1(xyz)+\cdots+h_p(xyz),
\end{aligned}
\]
где по 2. имеем
\[
  h_i(xyz)=
  \left(z\frac{\partial}{\partial x}\right)^{p-i}
  \left(z\frac{\partial}{\partial y}\right)^i g_i(xy).
\]
Отдельные $h_i$, а следовательно и $h$, являются формами $PZ$ и их суммами; пусть $\mathcal T$ имеет ранг $\tau$ относительно $R$.

Как показывает построение форм $h(x,y,z)$ семейства $\mathcal T$, эти формы получены из $\theta$ полярными процессами $P$ и в рядах $x,y,z$ по отдельности имеют ту же степень, что и $\theta$; поэтому семейство $\mathcal T$ образует подсемейство $\mathcal L$, и по упомянутому выше факту для доказательства эквивалентности остаётся доказать только равенство рангов. При этом доказательстве специальная структура форм $f(xyz)$ -- то есть то, что они являются полярами одной и той же формы $\theta$ -- далее уже не используется.

\emph{a) Для сравнения рангов $\mathcal L$ и $\mathcal S$ служит вспомогательная лемма a): из $f(x,y,\lambda_1x+\lambda_2y)=0$ необходимо следует $f(x,y,z)=0$.} Это непосредственно следует из тождественного соотношения между тремя бинарными рядами
\[
  (xy)z+(yz)x+(zx)y=0,
  \qquad (xy)=(x_1y_2-x_2y_1),\ldots,
\]
которое влечёт
\[
  f[x,y,(xy)x+(xz)y]=(xy)^p f(x,y,z).
\]
Из $f(x,y,\lambda_1x+\lambda_2y)=0$ (тождественно по $\lambda$), следовательно, как показывает подстановка $\lambda_1=(zy)$, $\lambda_2=(xz)$, и поскольку $(xy)\ne0$, необходимо следует $f(x,y,z)=0$.

По этой лемме между формами из $\mathcal S$ и формами из $\mathcal L$ существует взаимно однозначное соответствие. Действительно, из
\[
  f_1(x,y,\lambda_1x+\lambda_2y)=g(\lambda;xy),\qquad
  f_2(x,y,\lambda_1x+\lambda_2y)=g(\lambda;xy)
\]
необходимо следует
\[
  f_1(xyz)=f_2(xyz).
\]
Далее каждое соотношение в $\mathcal S$ сохраняется также между однозначно соответствующими формами в $\mathcal L$ (и обратно). Ибо из
\[
  c_1g^{(1)}(\lambda;xy)+c_2g^{(2)}(\lambda;xy)+\cdots+c_rg^{(r)}(\lambda;xy)=0
\]
по лемме необходимо следует
\[
  c_1f^{(1)}(x,y,z)+c_2f^{(2)}(x,y,z)+\cdots+c_rf^{(r)}(x,y,z)=0.
\]
\emph{Тем самым доказано равенство рангов $\mathcal L$ и $\mathcal S$; $\rho=\sigma$.}

\emph{b) Для сравнения рангов $\mathcal S$ и $\mathcal T$ соответственно служит вспомогательная лемма b): из $h(x,y,z)=0$ необходимо следует $g(\lambda;xy)=0$.} Для доказательства заметим, что по 3. равенство $h(x,y,z)=0$ равносильно исчезновению отдельных произведений степеней
\[
 \left(\frac{\partial}{\partial\lambda_1}\frac{\partial}{\partial x_1}
 +\frac{\partial}{\partial\lambda_2}\frac{\partial}{\partial y_1}\right)^{p_1}
 \left(\frac{\partial}{\partial\lambda_1}\frac{\partial}{\partial x_2}
 +\frac{\partial}{\partial\lambda_2}\frac{\partial}{\partial y_2}\right)^{p_2}g(\lambda;x,y),
 \qquad p_1+p_2=p,
\]
откуда дальнейшим дифференцированием следует также исчезновение произведений степеней
\[
\begin{aligned}
&\left(\frac{\partial}{\partial x_1}\right)^{\alpha_1}
 \left(\frac{\partial}{\partial x_2}\right)^{\alpha_2}
 \left(\frac{\partial}{\partial y_1}\right)^{\beta_1}
 \left(\frac{\partial}{\partial y_2}\right)^{\beta_2} \\
&\quad\cdot
 \left(\frac{\partial}{\partial\lambda_1}\frac{\partial}{\partial x_1}
 +\frac{\partial}{\partial\lambda_2}\frac{\partial}{\partial y_1}\right)^{p_1}
 \left(\frac{\partial}{\partial\lambda_1}\frac{\partial}{\partial x_2}
 +\frac{\partial}{\partial\lambda_2}\frac{\partial}{\partial y_2}\right)^{p_2}
 g(\lambda;xy)
\end{aligned}
\]
Следовательно, исчезает и всякая линейная комбинация этих произведений степеней, а значит, всякая форма
\[
  f\left(\frac{\partial}{\partial x},\frac{\partial}{\partial y},
  \frac{\partial}{\partial\lambda_1}\frac{\partial}{\partial x}
  +\frac{\partial}{\partial\lambda_2}\frac{\partial}{\partial y}\right)g(\lambda,xy).
\]
Если специально выбрать $f$ так, что
\[
  f(x,y,\lambda_1x+\lambda_2y)=g(\lambda;xy),
\]
то получаем также
\[
  g\left(\frac{\partial}{\partial x};
    \frac{\partial}{\partial x},
    \frac{\partial}{\partial y}\right)g(\lambda,x,y)=0;
\]
или -- если положить
\[
  g(\lambda,xy)=\sum G_i\lambda_1^{\nu_1}\lambda_2^{\nu_2}
  x_1^{\alpha_1}x_2^{\alpha_2}y_1^{\beta_1}y_2^{\beta_2},
\]
где $G_i$ являются рационально-целыми линейными формами от неопределённых $a$, -- имеем
\[
  \sum \nu_1!\nu_2!\alpha_1!\alpha_2!\beta_1!\beta_2!\,G_i^2=0
\]
и, следовательно, $g(\lambda;xy)=0$.

Из вспомогательной леммы b), точно так же как в a), следует \emph{равенство рангов $\mathcal S$ и $\mathcal T$; $\sigma=\tau$.}

Итак, по a) и b) имеем $\rho=\tau$; тем самым -- поскольку $\mathcal T$ есть подсемейство $\mathcal L$ -- доказана эквивалентность этих двух линейных семейств. Поэтому всякую форму из $\mathcal L$, в частности и $\theta$, можно представить как линейную рационально-целую комбинацию $\rho$ линейно независимых форм $h(x,y,z)$:
\[
  \theta=\sum c_i h^{(i)}(x,y,z)=\sum PZ.
\]
Это тождество, выведенное для неопределённых коэффициентов $a$ формы $\theta$, сохраняется, если заменить неопределённые $a$ величинами любой области рациональности, и \emph{редукционная теорема тем самым доказана вообще в случае трёх бинарных рядов.}

Общее доказательство для $N$ рядов по $n$ переменных в каждом проходит теперь совершенно параллельно. Пусть $\theta(A)$ имеет соответственно степень $\alpha_i$ в ряду $A_i$. Тогда снова рассмотрим три линейных семейства форм:

1. семейство $\mathcal L$, состоящее из всех форм $f(A)$, которые получены из $\theta(A)$ полярными процессами $P$ и в каждом отдельном ряду $A_i$ имеют ту же степень, что и $\theta(A)$;

2. семейство $\mathcal S$, состоящее из всех форм $g(\lambda;A_1\cdots A_n)$, которые получаются из $f(A)$ подстановками
\[
\begin{aligned}
 A_{n+1}&=\lambda_{n+1}^{(1)}A_1+\lambda_{n+1}^{(2)}A_2+
        \cdots+\lambda_{n+1}^{(n)}A_n,\\
 &\vdotswithin{=}\\
 A_N&=\lambda_N^{(1)}A_1+\lambda_N^{(2)}A_2+
        \cdots+\lambda_N^{(n)}A_n
\end{aligned}
\]
при этом коэффициенты отдельных произведений степеней $\lambda$ в $g(\lambda;A_1\cdots A_n)$ задают формы $Z$ тождества (1);

3. семейство $\mathcal T$, состоящее из всех форм $h(A)$, определённых посредством
\[
\begin{aligned}
 h(A)=&\left[A_{n+1}\left(
 \frac{\partial}{\partial\lambda_{n+1}^{(1)}}\frac{\partial}{\partial A_1}
 +\cdots+
 \frac{\partial}{\partial\lambda_{n+1}^{(n)}}\frac{\partial}{\partial A_n}\right)\right]^{\alpha_{n+1}} \\
&\cdots
\left[A_N\left(
 \frac{\partial}{\partial\lambda_N^{(1)}}\frac{\partial}{\partial A_1}
 +\cdots+
 \frac{\partial}{\partial\lambda_N^{(n)}}\frac{\partial}{\partial A_n}\right)\right]^{\alpha_N}
 g(\lambda;A_1\cdots A_n).
\end{aligned}
\]
$h(A)$ есть сумма форм $PZ$.

Теперь благодаря тождественному соотношению между любыми $n+1$ рядами вспомогательная лемма a) сохраняется; так же справедлива вспомогательная лемма b), в доказательстве которой число переменных вообще не играло роли. Следовательно, $\mathcal L$ и $\mathcal T$ имеют одинаковый ранг и -- поскольку $\mathcal T$ является подсемейством $\mathcal L$ -- тождественны. Поэтому тождество
\[
  \theta=\sum PZ
\]
имеет место для неопределённых коэффициентов, а следовательно, и для коэффициентов, являющихся величинами любой области рациональности; \emph{тем самым редукционная теорема доказана в общем виде.}

\subsection*{III.}
Коротко покажем ещё, как аналогичные соображения дают также \emph{общее разложение по рядам} (для $N\le n$). Пусть $Z$ есть отдельно однородная форма $n$ рядов $A_1,\ldots,A_n$ по $n$ величин в каждом, и пусть $\Delta$ есть детерминант этих рядов. Сначала докажем соответствующее (1) тождество:
\begin{equation}
  Z=\sum PH \pmod{\Delta},
\tag{2}
\end{equation}
где формы $H$ содержат уже только ряды $A_1,\ldots,A_{n-1}$ и выведены из $Z$ процессами $P$.

Доказательство состоит в превращении соотношений, выведенных в II, в конгруэнции по модулю $\Delta$. Ради простоты положим в основу три тернарных ряда $x,y,z$; коэффициенты предполагаются неопределёнными.

Три линейных семейства $\mathcal L,\mathcal S,\mathcal T$ пусть будут определены как в II для случая бинарных рядов; их ранги пусть будут $\rho,\sigma,\tau$, тогда как $\rho^*$ и $\tau^*$ обозначают ранги $\mathcal L$ и $\mathcal T$ по модулю $\Delta$ (то есть существуют $\rho^*$, но не $\rho^*+1$ форм из $\mathcal L$, линейно независимых по модулю $\Delta$). Тогда, как в II, имеет место

\emph{Вспомогательная лемма a): из $f(x,y,\lambda_1x+\lambda_2y)=0$ необходимо следует $f(xyz)\equiv0\pmod{\Delta}$ (и обратно).} Действительно, в силу соотношения между четырьмя тернарными рядами
\[
  \Delta_1x-\Delta_2y+\Delta_3z=\Delta t,
  \qquad \Delta_1=(yzt),\ldots,
\]
между тремя тернарными рядами всегда существует линейное соотношение по модулю $\Delta$:
\[
  \Delta_1x-\Delta_2y+\Delta_3z\equiv0\pmod{\Delta},
\]
и поэтому, как в II,
\[
  f(x,y,-\Delta_1x+\Delta_2y)=\Delta_3^p f(x,y,z)\pmod{\Delta}.
\]
Из $f(x,y,\lambda_1x+\lambda_2y)=0$ (тождественно по $\lambda$) отсюда -- поскольку $\Delta$ неприводим и $\Delta_3$ не делится на $\Delta$ -- необходимо следует $f(x,y,z)\equiv0\pmod{\Delta}$. Обратное ясно из $\Delta(x,y,\lambda_1x+\lambda_2y)=0$. Отсюда, как в II, заключаем: $\rho^*=\sigma$.

Вспомогательная лемма b) остаётся в силе вместе с доказательством, и поэтому имеем $\sigma=\tau$.

Для доказательства $\tau=\tau^*$ служит

\emph{Вспомогательная лемма c): из $h(x,y,z)\equiv0\pmod{\Delta}$ необходимо следует $h(x,y,z)=0$.}
Действительно, $h(x,y,z)$ как поляра исчезает при применении известного $\Omega$-процесса; это выражается дифференциальным уравнением
\[
  \Delta\!\left(\frac{\partial}{\partial x},\frac{\partial}{\partial y},\frac{\partial}{\partial z}\right)h(x,y,z)=0.
\]
Если теперь $h(x,y,z)=\Delta(x,y,z)\cdot k(x,y,z)$, то дальнейшим дифференцированием исчезает также
\[
  k\!\left(\frac{\partial}{\partial x},\frac{\partial}{\partial y},\frac{\partial}{\partial z}\right)
  \Delta\!\left(\frac{\partial}{\partial x},\frac{\partial}{\partial y},\frac{\partial}{\partial z}\right)
  =
  h\!\left(\frac{\partial}{\partial x},\frac{\partial}{\partial y},\frac{\partial}{\partial z}\right)h(x,y,z),
\]
а следовательно, и $h(x,y,z)$. Поэтому имеем $\rho^*=\sigma=\tau=\tau^*$, и, так как $\mathcal T$ есть подсемейство $\mathcal L$, тождество (2) доказано; соответствующее рассуждение даёт доказательство и при $n$ переменных.\footnote{По 3. имеется соответствующее, исчезающее из-за детерминантного множителя, соотношение для $\Omega(h)$; в операторной записи это уравнение, возникающее из произведений
\[
 \left(\lambda_1\frac{\partial}{\partial \xi}+\lambda_2\frac{\partial}{\partial \eta}\right),
 \qquad
 \left[z\left(\frac{\partial}{\partial\lambda_1}\frac{\partial}{\partial \xi}
 +\frac{\partial}{\partial\lambda_2}\frac{\partial}{\partial \eta}\right)\right]
\]
и оно исчезает из-за исчезновения детерминантного множителя.}

Применяя тождество (2) к множителю при $\Delta$, получаем разложение
\[
  Z=\varphi_0+\Delta\varphi_1+\Delta^2\varphi_2+\cdots+\Delta^\mu\varphi_\mu,
\]
где каждая из $\varphi_i$ исчезает при применении $\Omega$-процесса. $i$-кратное повторение процесса даёт, поскольку вообще $\Omega(\Delta\cdot\psi)=c_1\psi+c_2\Delta\cdot\Omega(\psi)$,
\[
  c\cdot \Omega^i(Z)\equiv\varphi_i\pmod{\Delta};
\]
с другой стороны, по (2) имеем
\[
  c\cdot\Omega^i(Z)\equiv\sum PH_i\pmod{\Delta},
\]
где $H_i$ выведены из формы $\Omega^i(Z)$ так же, как $H$ из $Z$. По вспомогательной лемме c) (в распространении на $n$ переменных) получаем поэтому $\varphi_i=\sum PH_i$, и, следовательно:
\begin{equation}
  Z=\sum PH+\Delta\cdot\sum PH_1+\Delta^2\cdot\sum PH_2+
  \cdots+\Delta^\mu\cdot\sum PH_\mu
\tag{3}
\end{equation}
\emph{как самое общее разложение формы от $n$ рядов по $n$ переменных по полярам форм только от $n-1$ рядов.}

Разложение по рядам для форм от $N$ рядов по $n$ переменных $(N<n)$ получается отсюда, как известно, простой подстановкой. Положим (для $N=3$, $n>3$):
\[
  u=\xi_1x+\xi_2y+\xi_3z,
  \quad v=\eta_1x+\eta_2y+\eta_3z,
  \quad w=\zeta_1x+\zeta_2y+\zeta_3z
\]
и разложим $H(u,v,w)=Z(\xi,\eta,\zeta)$ по (3). Тогда, поскольку $\xi,\eta,\zeta$ входят только в комбинациях $u,v,w$, имеем
\[
  \left(\eta\frac{\partial}{\partial \xi}\right)=
  \left(v\frac{\partial}{\partial u}\right)\cdots,
  \qquad
  \Omega=\sum_{ikl}\left(\frac{\partial}{\partial u}
  \frac{\partial}{\partial v}
  \frac{\partial}{\partial w}\right)_{ikl}(xyz)_{ikl}
  =\nabla_{uvw}(xyz).
\]
В силу специализации $\xi_1=\eta_2=\zeta_3=1$; $\xi_2=\xi_3=\cdots=\zeta_2=0$ форма $H(u,v,w)$ переходит в $H(x,y,z)$, а $\nabla_{uvw}(xyz)$ -- с точностью до числового множителя -- в $\nabla_{xyz}(xyz)=\nabla_{xyz}$, поскольку применение $(y\partial/\partial x)$ и т. д. к детерминантам $(xyz)_{ikl}$ обращает их в нуль. Так как соответствующее верно и для $N$ рядов, из (3) следует разложение
\begin{equation}
  H=\sum P\Phi+\sum P\Phi_1+\cdots+\sum P\Phi_\mu,
\tag{4}
\end{equation}
где $\Phi_i$ получены из $\nabla^i_{A_1\ldots A_N}H$ процессами $P$ и содержат уже только ряды $A_1,\ldots,A_{N-1}$, за исключением детерминантных комбинаций, входящих в $\nabla^i$, которые, как сказано, при поляризации не принимаются в расчёт.

Если поэтому заменить эти детерминанты в $\nabla^i$ неопределёнными $p$, то возникающие формы можно снова разложить по (4); так в конце концов получается разложение
\begin{equation}
  H=\left[\sum P\psi(p)\right]_{p=\text{детерминанты }A_1\cdots A_N},
\tag{5}
\end{equation}
где $\psi$ получаются из исходной формы процессами $P$ и $\nabla_{A_1\ldots A_N}(p)$, $\nabla_{A_1\ldots A_{N-1}}(p)$ и т. д. Этими разложениями и их комбинациями -- например, если применить к формам $Z$, входящим в (1), сначала (3), а затем (4) и (5), -- исчерпываются все известные разложения для форм от произвольно многих рядов по $n$ когредиентных переменных в каждом.

Эрланген, 5 января 1915 г.
% END UNIT BODY
% END PRODUCER UNIT 057
\endgroup


\clearpage
\begingroup
\setcounter{footnote}{0}
% BEGIN PRODUCER UNIT 058
% Source: C:/Users/Floris/Documents/interlanguage/03_projects/noether/02_slavic_working_corpus/translations/paper09/russian/v001/Noether_Paper09_Introduction_Russian_v001.tex
% Identity: 13405 B / SHA-256 D0470DF50EA822181D4CE7C256616587C08EAFCD8B7F529EE0C42E264CE0EFCE
\setlength{\parindent}{1.25em}
\setlength{\parskip}{0pt}
\emergencystretch=14em
\allowdisplaybreaks
\providecommand{\vdotswithin}[1]{\mathrel{\vdots}}
% BEGIN UNIT BODY
\begin{center}
{\Large\bfseries 9. Наиболее общие области целых трансцендентных чисел}\par
\vspace{1.2em}
Math. Ann. 77 (1916), с. 103--128
\end{center}

\vspace{1em}

В силу теоремы о вполне упорядочении Э. Цермело построил область «целых трансцендентных чисел», то есть числовую область $\mathfrak G$, обладающую следующими \emph{абстрактными} свойствами:
\begin{enumerate}
\item[I.] Сумма, разность и произведение двух чисел из $\mathfrak G$ снова являются числом из $\mathfrak G$.
\item[II.] Каждое действительное или комплексное число является частным двух чисел из $\mathfrak G$.
\item[III.] Каждое целое рациональное (соответственно алгебраическое) число является элементом $\mathfrak G$.
\item[IV.] Никакое нецелое рациональное (соответственно алгебраическое) число не принадлежит $\mathfrak G$.\footnote{Э. Цермело, «О целых трансцендентных числах», Math. Ann. 75, с. 434 (1914).}
\end{enumerate}

Конструкция основана на том, что теорема о вполне упорядочении позволяет установить существование \emph{алгебраической базы} всех чисел, то есть \emph{системы} $H$ \emph{чисел} $\eta$, \emph{между которыми нет алгебраических соотношений, но через которые можно алгебраически выразить все остальные числа}. Благодаря алгебраической независимости эти базовые числа $\eta$ можно рассматривать как \emph{неопределённые}; тем самым искомая область переходит в \emph{область целостности из рациональных и алгебраических функций от неопределённых} $\eta$,\footnote{По этой причине рассуждения настоящей работы отчасти идут параллельно рассуждениям прежней работы автора: «Поля и системы рациональных функций», Math. Ann. 76, с. 161 (1915).} коэффициенты которых принадлежат полю $K$ всех алгебраических чисел. Специальная область $\mathfrak G_\eta$, построенная Цермело, тогда характеризуется следующими \emph{базовыми свойствами}:

$\mathfrak G_\eta$ содержит:
\begin{enumerate}
\item[1)] все базовые числа $\eta$,
\item[2)] все многочлены от $\eta$ с целочисленными\footnote{Под «целочисленный» здесь и далее понимается, что коэффициенты принадлежат области целостности $[K]$ всех целых алгебраических чисел. Для определения $\mathfrak G_\eta$ в пунктах 1)--5) было бы достаточно рассматривать только целые рациональные числовые коэффициенты; однако для более общих областей базовые свойства тогда стали бы менее простыми.} коэффициентами,
\item[3)] все целые алгебраические функции от $\eta$ с целочисленными коэффициентами.
\end{enumerate}
$\mathfrak G_\eta$ исключает:
\begin{enumerate}
\item[4)] все рационально-дробные функции от $\eta$ с целочисленными коэффициентами,
\item[5)] все алгебраически-дробные функции от $\eta$ с целочисленными коэффициентами.\footnote{Цермело, цит. соч., \S\ 3.}
\end{enumerate}

Теперь легко видеть (ср. примеры в \S\ 2 и 3), что существуют области $\mathfrak G$, \emph{существенно отличные} от области Цермело $\mathfrak G_\eta$, то есть такие области $\mathfrak G$, которые ни при каком \emph{вполне упорядочении} не обладают всеми базовыми свойствами 1)--5) и, следовательно, ни при каком вполне упорядочении не могут быть получены конструкцией Цермело. Иными словами, существуют области $\mathfrak G$, которые \emph{нельзя} взаимно однозначно и изоморфно отобразить на $\mathfrak G_\eta$. Тем самым возникает вопрос о \emph{наиболее общих областях целых трансцендентных чисел}; ниже он полностью решается.

Для этого сначала надо установить, какие из базовых свойств являются следствиями абстрактных определяющих свойств, а какие обусловлены специальной конструкцией. Оказывается (\S\ 1), что для каждой заранее заданной области $\mathfrak G$ существует вполне упорядочение, при котором выполняются базовые свойства 1) и 2); \emph{следовательно, каждую область $\mathfrak G$ можно построить с помощью алгебраической базы всех чисел}. Из остальных базовых свойств ни одно не обязано выполняться (\S\ 2 и 3). Отсюда, в частности, следует, что ещё одно абстрактное свойство области Цермело не принадлежит всем областям $\mathfrak G$: \emph{алгебраически-целая замкнутость}. Его можно сформулировать так:
\begin{enumerate}
\item[V.] Каждое число, алгебраически-цело и целочисленно зависящее от конечного числа чисел из $\mathfrak G$, принадлежит $\mathfrak G$.
\end{enumerate}

Специальные области, которым наряду с I--IV присуще также V, будем называть областями $\mathfrak H$ из \emph{алгебраически-целых трансцендентных чисел}. Для наиболее общих областей $\mathfrak H$ получаются простые результаты (\S\ 4):
\begin{enumerate}
\item[a)] \emph{Пересечение всех областей $\mathfrak H$, которые можно построить с помощью одной и той же базы $H$, задаётся всеми целыми алгебраическими функциями этой базы, то есть областью Цермело $\mathfrak G_\eta$, которая сама является областью $\mathfrak H$.} (Тем самым одновременно получена абстрактная характеристика области Цермело.)
\item[b)] \emph{Каждая область $\mathfrak H$ возникает из $\mathfrak G_\eta$ посредством алгебраически-целого присоединения\footnote{Пояснение этого термина дано в \S\ 4.} произвольной системы $\mathfrak S(\eta)$; при этом $\mathfrak S(\eta)$ должна удовлетворять лишь одному условию: в результате присоединения в область не должно войти никакое алгебраически-дробное число.}
\end{enumerate}

Менее просты результаты для наиболее общих областей $\mathfrak G$, если за основу взята алгебраическая база. Здесь пересечение всех областей $\mathfrak G$ само уже не является областью $\mathfrak G$; поэтому и конструкцию наиболее общих областей труднее обозреть (\S\ 5).

Чтобы получить результаты, аналогичные результатам для областей $\mathfrak H$, алгебраическая база заменяется \emph{рациональной базой} $\Theta$, то есть \emph{системой} $\Theta$ \emph{чисел} $\vartheta$, \emph{которые позволяют рационально выразить все числа, тогда как по меньшей мере при одном вполне упорядочении никакое базовое число не может быть рационально представлено через конечное число предшествующих}. На эту рациональную базу $\Theta$ надо ещё наложить дополнительное условие: область целостности $[\Theta]$, состоящая из всех целочисленных многочленов от $\vartheta$, не должна содержать никакого алгебраически-дробного числа. (Построение такой базы с дополнительным условием показано в \S\ 7.) Тогда для наиболее общих областей $\mathfrak G$, которые можно было бы также назвать областями из \emph{рационально-целых трансцендентных чисел}, снова получаются простые результаты (\S\ 6):
\begin{enumerate}
\item[a)] \emph{Пересечение всех областей $\mathfrak G$, которые можно построить с помощью одной и той же рациональной базы $\Theta$, задаётся всеми целыми рациональными функциями этой базы, то есть областью $[\Theta]$, которая сама является областью $\mathfrak G$.}
\item[b)] \emph{Каждая область $\mathfrak G$ возникает из $[\Theta]$ посредством рационально-целого присоединения произвольной системы $\mathfrak S(\vartheta)$; при этом $\mathfrak S(\vartheta)$ должна удовлетворять лишь одному условию: в результате присоединения в область не должно войти никакое алгебраически-дробное число.}
\end{enumerate}

Чтобы иметь полную аналогию с областями $\mathfrak H$, следовало бы не включать алгебраические числа в определяющие свойства III и IV для областей $\mathfrak G$. При такой модификации III и IV конструкция целых элементов с помощью рациональной базы переносится на любое \emph{произвольное} поле (\S\ 10), тогда как конструкция Цермело предполагает алгебраическую замкнутость поля.

\emph{Когда $H$ пробегает все алгебраические базы, а $\Theta$ --- все рациональные базы, удовлетворяющие дополнительному условию, результаты a) и b) дают совокупность всех областей $\mathfrak H$ и $\mathfrak G$.} Если объединить все изоморфные области в один класс, то из этой совокупности можно выделить подмножество, представляющее совокупность всех классов --- по \S\ 9 по меньшей мере счётно бесконечно многих --- классов (\S\ 8).
% END UNIT BODY
% END PRODUCER UNIT 058
\endgroup


\clearpage
\begingroup
\setcounter{footnote}{0}
% BEGIN PRODUCER UNIT 059
% Source: C:/Users/Floris/Documents/interlanguage/03_projects/noether/02_slavic_working_corpus/translations/paper09/russian/v001/Noether_Paper09_Section01_Russian_v001.tex
% Identity: 3656 B / SHA-256 C3910B76AE8B0DB7580145C56F0A84E89F87A7AF5D84D04429FCAE7F6E805648
\setlength{\parindent}{1.25em}
\setlength{\parskip}{0pt}
\emergencystretch=14em
\allowdisplaybreaks
\providecommand{\vdotswithin}[1]{\mathrel{\vdots}}
% BEGIN UNIT BODY
\section*{\S\ 1. Доказательство базовых свойств 1) и 2) для всех областей $\mathfrak G$.}

Пусть $\mathfrak G$ есть произвольно заданная область целых трансцендентных чисел, то есть область, обладающая абстрактными свойствами I--IV. По процедуре, которую Э. Цермело применил к полю всех комплексных чисел, выберем из $\mathfrak G$ \emph{алгебраическую базу} $H$ \emph{всех чисел из $\mathfrak G$}. Так как по II \emph{каждое} действительное или комплексное число можно представить как частное двух чисел из $\mathfrak G$, то $H$ одновременно является алгебраической базой \emph{всех} чисел. Следовательно, для каждой произвольно заданной области $\mathfrak G$ алгебраическую базу всех чисел можно выбрать так, чтобы она принадлежала $\mathfrak G$; тем самым \emph{базовое свойство 1) выполнено}. Далее, по III область $\mathfrak G$ содержит область целостности $[K]$ всех целых алгебраических чисел, а потому по I также все многочлены от базовых величин $\eta$ с коэффициентами из $[K]$, то есть все целочисленные многочлены от $\eta$, образующие область целостности $[H]$. Значит, \emph{базовое свойство 2) также всегда выполнено}; то есть \emph{каждую область $\mathfrak G$ можно построить посредством алгебраической базы $H$ всех чисел так, что $\mathfrak G$ содержит область целостности $[H]$ как делитель.}

\textbf{Замечание.} Тем не менее, конечно, можно, исходя из базы $H$, построить область $\mathfrak G$, для которой базовые свойства 1) и 2) выполняются не именно относительно $H$, а относительно какой-нибудь другой базы $H'$. Пусть, например, $\mathfrak G'$ возникает из области Цермело $\mathfrak G_\eta$ так, что во всей конструкции некоторая базовая величина $\eta$ заменяется многочленом $f(\eta)$. Тогда $\mathfrak G'$ не содержит ни $\eta$, ни $[H]$; но если в базе $H$ заменить базовую величину $\eta$ на $\xi=f(\eta)$ --- причём $H$ снова переходит в алгебраическую базу $H'$ ---, то базовые свойства 1) и 2) для $\mathfrak G'$ выполняются относительно $H'$.
% END UNIT BODY
% END PRODUCER UNIT 059
\endgroup


\clearpage
\begingroup
\setcounter{footnote}{0}
% BEGIN PRODUCER UNIT 060
% Source: C:/Users/Floris/Documents/interlanguage/03_projects/noether/02_slavic_working_corpus/translations/paper09/russian/v001/Noether_Paper09_Section02_Russian_v001.tex
% Identity: 7750 B / SHA-256 17DB9C41AE602D02B3F0EB53C67A3FC6A9A829FAC23CA8CE781E40FFED42DE80
\setlength{\parindent}{1.25em}
\setlength{\parskip}{0pt}
\emergencystretch=14em
\allowdisplaybreaks
\providecommand{\vdotswithin}[1]{\mathrel{\vdots}}
% BEGIN UNIT BODY
\section*{\S\ 2. Исключение базовых свойств 4) и 5).}

Теперь нужно показать, что области $\mathfrak G$ не обязаны удовлетворять ни одному из дальнейших базовых свойств. Сначала свойства 4) и 5) будут исключены тем, что в конструкции Цермело область целостности $[H]$ заменяется областью из целых рациональных «функционалов».

По Веберу\footnote{Г. Вебер, \emph{Учебник алгебры}. Малое издание, \S\S\ 96 и 97.} под \emph{рациональным функционалом} понимают всякую рациональную функцию с рациональными числовыми коэффициентами в следующем однозначно определённом представлении:
\[
 A(\eta_1\cdots \eta_t)
   = a\cdot\frac{E_1(\eta_1\cdots \eta_t)}{E_2(\eta_1\cdots \eta_t)},
\]
где $E_1$, $E_2$ являются взаимно простыми примитивными многочленами с целыми рациональными числовыми коэффициентами, а $a$ --- положительным рациональным числом (абсолютной величиной $A$). Если, в частности, $a$ есть \emph{целое рациональное} число, то $A$ называется \emph{целым рациональным функционалом}. Как специальные случаи среди целых рациональных функционалов содержатся целые рациональные числа и многочлены с целыми рациональными числовыми коэффициентами, тогда как рационально-дробные числа и многочлены с рационально-дробными коэффициентами следует относить к дробным рациональным функционалам.

\emph{Пусть теперь область $\mathfrak G$ состоит из совокупности $\mathfrak F$ всех целых рациональных функционалов от каждый раз конечного числа базовых величин $\eta$ и из всех функций, алгебраически целых над $\mathfrak F$} --- то есть удовлетворяющих некоторому уравнению, старший коэффициент которого равен единице, а остальные коэффициенты являются целыми рациональными функционалами.

Доказательство свойств I--IV для $\mathfrak G$ идёт совершенно параллельно соответствующему доказательству у Цермело; коротко его наметим. Прежде всего II и III ясны, поскольку $\mathfrak G$ содержит область Цермело $\mathfrak G_\eta$ как делитель. Далее, по теореме Гаусса о многочленах с целыми рациональными числовыми коэффициентами сумма, разность и произведение целых рациональных функционалов снова являются \emph{целыми} рациональными функционалами; значит, $\mathfrak F$ образует область целостности, а по известным рассуждениям\footnote{Ср., например, Вебер, \S\ 97, 8.} областью целостности становится также $\mathfrak G$, чем выполнено I. Кроме того, из расширенной теоремы Гаусса\footnote{Вебер, \S\ 20, 5.} следуют две вспомогательные леммы: «Если $z$ алгебраически цел над $\mathfrak F$ посредством \emph{какого-либо} уравнения, то он алгебраически цел и посредством \emph{неприводимого} уравнения, которому $z$ удовлетворяет относительно поля всех рациональных функционалов»\footnote{Ср. Вебер, \S\ 97, 4.} и «уравнение, неприводимое относительно поля всех рациональных чисел, которому удовлетворяет алгебраическое число, неприводимо также относительно поля всех рациональных функционалов». Следовательно, $\mathfrak G$ не может содержать никакого дробного алгебраического числа; этим доказано IV, а значит, и все абстрактные свойства I--IV.

\emph{Напротив, ни при каком вполне упорядочении нельзя выбрать из $\mathfrak G$ базу так, чтобы $\mathfrak G$ исключала все рационально- и алгебраически-дробные функции от базовых величин.} Действительно, пусть $z(\eta)$ есть какая-либо функция области, которую следует выбрать в качестве базовой величины, то есть функция, определённая уравнением:
\[
 z^k+A_1(\eta)z^{k-1}+\cdots+A_k(\eta)=0,
\]
где
\[
 A_i(\eta)=a_i\cdot\frac{E_1(\eta)}{E_2(\eta)}
\]
являются целыми рациональными функционалами. Тогда функция $\frac{a_k}{z}$ принадлежит $\mathfrak G$, и точно так же принадлежат $\frac{a_k}{\sqrt z}$, $\frac{a_k}{\sqrt[3]z}$ и т. д.\footnote{Совершенно аналогично каждую алгебраическую функцию от $\eta$ можно, умножив на целое рациональное число, превратить в величину функциональной области $\mathfrak G$. Тем самым эта область обладает свойством: \emph{каждое комплексное число можно представить как частное величины из $\mathfrak G$ и целого рационального числа} --- по аналогии с алгебраическими числами.} Тем самым \emph{базовые свойства 4) и 5) для совокупности областей $\mathfrak G$ исключены.}

\textbf{Замечание.} Пример в \S\ 3 покажет, что $\mathfrak G$ при каждом вполне упорядочении может содержать рационально-дробные функционалы, не содержа рационально- или алгебраически-дробного числа.
% END UNIT BODY
% END PRODUCER UNIT 060
\endgroup


\clearpage
\begingroup
\setcounter{footnote}{0}
% BEGIN PRODUCER UNIT 061
% Source: C:/Users/Floris/Documents/interlanguage/03_projects/noether/02_slavic_working_corpus/translations/paper09/russian/v001/Noether_Paper09_Section03_Russian_v001.tex
% Identity: 10049 B / SHA-256 718901AE67C21991760EA9187E7A90DC92439CDB57EBB8E01F90B5C05BD6FFA9
\setlength{\parindent}{1.25em}
\setlength{\parskip}{0pt}
\emergencystretch=16em
\allowdisplaybreaks
\providecommand{\vdotswithin}[1]{\mathrel{\vdots}}
% BEGIN UNIT BODY
\section*{\S\ 3. Исключение базового свойства 3).}

Для исключения базового свойства 3) служит обобщённая
\emph{модульная область}: в ней аргументами многочленов области являются уже не
сами неопределённые, а все целые алгебраические функции от этих
неопределённых\footnote{Под «целыми алгебраическими функциями» здесь всегда
понимаются функции с \emph{целочисленными} коэффициентами, то есть функции,
удовлетворяющие некоторому уравнению, старший коэффициент которого равен
единице, а остальные коэффициенты являются многочленами от $\eta$ с
алгебраически-целыми числовыми коэффициентами, то есть многочленами из $[H]$.
В частности, все многочлены из $[H]$ принадлежат к целым алгебраическим
функциям.}; сами же неопределённые принимают на себя роль модульной базы.

Пусть $\mathfrak G$ состоит из \emph{всех величин}
\begin{equation}
 z=\alpha+g_1(\eta)\eta_1+g_2(\eta)\eta_2+\cdots+g_t(\eta)\eta_t,
\tag{1}
\end{equation}
\emph{где $\alpha$ пробегает все алгебраически целые числа,
$\eta_1\cdots\eta_t$ пробегают все базовые величины, а
$g_1(\eta)\cdots g_t(\eta)$ при каждой фиксированной комбинации
$\eta_1\cdots\eta_t$ по отдельности пробегают все целые алгебраические функции
от $\eta$.}\footnote{Модульное свойство принадлежит величинам $z-\alpha$.}

Легко проверить, что свойства I--IV для этой модульной области выполнены.
Сначала выполнено I, поскольку сумма, разность и произведение двух величин $z$
снова имеют форму (1). Кроме того, $\mathfrak G$ содержит все базовые величины
$\eta$, а также все величины $\eta\cdot g(\eta)$, где $g(\eta)$ пробегает все
целые алгебраические функции от $\eta$. Значит, все \emph{целые}
алгебраические функции от $\eta$, а потому и все алгебраические функции от
$\eta$ вообще, можно представить как частные двух величин из $\mathfrak G$;
этим доказано II. В представлении (1), в частности при
$g_1=0\cdots g_t=0$, содержатся все алгебраически целые числа; но $z$ не может
быть равным никакому дробному алгебраическому числу. Действительно, если $z$
представляет алгебраическое число $\beta$, то из тождества по $\eta$
\[
 \beta=\alpha+g_1(\eta)\eta_1+\cdots+g_t(\eta)\eta_t
\]
необходимо следует $\beta=\alpha$. Это видно из специализации
\[
 \eta_1=0\cdots \eta_t=0,
\]
при которой множители $g_1(\eta)\cdots g_t(\eta)$ как целые алгебраические
функции остаются конечными, поскольку переходят в целые алгебраические функции
или числа.\footnote{Это более подробное доказательство IV приведено из-за
расширенного примера в конце параграфа. Здесь хватило бы замечания, что $z$ по
конструкции является целой алгебраической функцией.} Итак, условия I--IV для
$\mathfrak G$ выполнены.

{\itshape Напротив, ни при каком вполне упорядочении нельзя выбрать из
$\mathfrak G$ базу так, чтобы все целые алгебраические функции от базовых
величин принадлежали $\mathfrak G$.} Пусть $z(\eta)$ есть любая функция
области, которую следует выбрать в качестве базовой величины, то есть функция,
которая действительно содержит неопределённые $\eta$. Мы покажем, что функции
\[
 \xi=\sqrt[\nu]{z-\alpha}
\]
начиная с некоторого конечного значения $\nu$, зависящего от $z$, уже не могут
принадлежать $\mathfrak G$.

Действительно, предположим, что $\xi$ принадлежит $\mathfrak G$. Тогда по (1)
оно имеет вид
\[
 \xi=\gamma+h_1(\eta)\eta_{i_1}+h_2(\eta)\eta_{i_2}
      +\cdots+h_\tau(\eta)\eta_{i_\tau},
\]
где снова $h_1(\eta)\cdots h_\tau(\eta)$ являются целыми алгебраическими
функциями от $\eta$, а $\eta_{i_1}\cdots\eta_{i_\tau}$ являются произвольными
базовыми величинами, которые в частности могут совпадать с
$\eta_1\cdots\eta_t$. Прежде всего нужно доказать, что алгебраически целое
число $\gamma$ равно нулю. Поскольку $h_1(\eta)\cdots h_\tau(\eta)$ являются
целыми алгебраическими функциями, $\xi$ становится равным $\gamma$ при
$\eta_{i_1}=0\cdots\eta_{i_\tau}=0$ и произвольных значениях остальных $\eta$;
в частности, имеем $\xi=\gamma$ при
\[
 \eta_{i_1}=0\cdots \eta_{i_\tau}=0; \qquad
 \eta_1=0\cdots \eta_t=0.
\]
Но при этой специализации функция $(z-\alpha)$ исчезает по (1), а значит
исчезает и $\xi$; следовательно, $\gamma=0$, и имеем
\[
 \xi=h_1(\eta)\eta_{i_1}+h_2(\eta)\eta_{i_2}
      +\cdots+h_\tau(\eta)\eta_{i_\tau}.
\]
Возведение в степень $\nu$ даёт
\[
 \xi^\nu=z-\alpha=F_\nu(\eta),
\]
где $F_\nu(\eta)$ обозначает однородную, не тождественно нулевую, форму степени
$\nu$ по $\eta_{i_1}\cdots\eta_{i_\tau}$, коэффициенты которой являются целыми
алгебраическими функциями от $\eta$. Теперь $z-\alpha$, которое по (1) является
\emph{целой} алгебраической функцией от $\eta$, удовлетворяет неприводимому
относительно $[H]$ уравнению
\[
 (z-\alpha)^\chi+B(\eta)(z-\alpha)^{\chi-1}+\cdots+C(\eta)=0,
\]
где последний коэффициент $C(\eta)$, произведение $(z-\alpha)$ с его
сопряжёнными, является многочленом; пусть его \emph{степень} равна $\lambda$.
С другой стороны, из $z-\alpha=F_\nu(\eta)$ умножением $F_\nu(\eta)$ на
соответствующие сопряжённые получаем
\[
 C(\eta)=G_{\nu\chi}(\eta),
\]
где $G_{\nu\chi}(\eta)$ обозначает однородную форму степени $\nu\chi$ по
$\eta_{i_1}\cdots\eta_{i_\tau}$, коэффициенты которой являются многочленами от
$\eta$. Значит, $G_{\nu\chi}$ имеет степень по $\eta$ по меньшей мере
$\nu\chi$, то есть $\lambda\geq\nu\chi$. Иными словами, {\itshape функции
$\sqrt[\nu]{z-\alpha}$, для которых $\nu>\lambda/\chi$, не принадлежат
$\mathfrak G$. Поэтому базовое свойство 3) исключено для совокупности областей
$\mathfrak G$.}

\textbf{Замечание.} Небольшое обобщение только что построенной области ведёт к
области $\mathfrak G$, которая при каждом вполне упорядочении содержит также
\emph{дробные функционалы}. Пусть $\mathfrak G$ состоит из всех величин
\[
 z=\alpha+\gamma_1(\eta)\eta_1+\gamma_2(\eta)\eta_2
      +\cdots+\gamma_t(\eta)\eta_t,
\]
где теперь $\gamma(\eta)$ пробегают все алгебраически целые, но \emph{не
обязательно целочисленные} функции от $\eta$. Доказательство свойств I--IV
остаётся тем же, что выше. Но поскольку $\mathfrak G$ наряду с каждой функцией
$z$ содержит также $a\cdot(z-\alpha)$, где $a$ обозначает произвольное
рациональное или алгебраически-дробное число, при каждом вполне упорядочении
возникают также дробные функционалы.
% END UNIT BODY
% END PRODUCER UNIT 061
\endgroup


\clearpage
\begingroup
\setcounter{footnote}{0}
% BEGIN PRODUCER UNIT 062
% Source: C:/Users/Floris/Documents/interlanguage/03_projects/noether/02_slavic_working_corpus/translations/paper09/russian/v001/Noether_Paper09_Section04_Russian_v001.tex
% Identity: 17545 B / SHA-256 1ADD5F4CBCB2704F1CE764ACFD440441D97D523955BA1FBDCBB51C40FB95AFBF
\setlength{\parindent}{1.25em}
\setlength{\parskip}{0pt}
\emergencystretch=18em
\allowdisplaybreaks
\raggedbottom
\providecommand{\vdotswithin}[1]{\mathrel{\vdots}}
% BEGIN UNIT BODY
\section*{\S\ 4. Самые общие области из алгебраически целых трансцендентных чисел.}

Для более удобной формулировки дальнейшего введём выражение:
«\emph{величина $z$ алгебраически цела над $\mathfrak T$}». Это означает, что
$z$ удовлетворяет некоторому уравнению, старший коэффициент которого равен
единице, а остальные коэффициенты являются целочисленными многочленами от
величин из $\mathfrak T$, где $\mathfrak T$ обозначает любую заданную
область.\footnote{Для области целостности $\mathfrak T$ старший коэффициент
равен единице, а остальные коэффициенты являются величинами из $\mathfrak T$.}

Область $\mathfrak T$ будем называть \emph{интегрально замкнутой}, если она
выполняет условие (ср. введение):
\begin{enumerate}
\item[V.] Всякая величина, алгебраически целая над $\mathfrak T$, принадлежит
$\mathfrak T$.
\end{enumerate}

Сначала покажем, что абстрактное свойство V принадлежит области Цермело
$\mathfrak G_\eta$. Пусть величина $z$ алгебраически цела над
$\mathfrak G_\eta$ посредством уравнения $f(z)=0$. Тогда умножение $f(z)$ на
его сопряжённые относительно $[H]$ показывает, что $z$ также алгебраически цела
над $[H]$ и, следовательно, принадлежит $\mathfrak G_\eta$. Точно так же
доказывается интегральная замкнутость функциональной области, рассмотренной в
\S\ 2.

То, что свойство V не принадлежит каждой области $\mathfrak G$, показывает
модульная область из \S\ 3, которая как раз не имеет базового свойства 3). Ещё
точнее, в \S\ 3 показано, что модульная область не интегрально замкнута
относительно \emph{ни одного трансцендентного} элемента области, тогда как по
III и IV это, конечно, имеет место относительно \emph{каждого алгебраического}
элемента области.

Поэтому из всей совокупности областей $\mathfrak G$ сначала выделим те, которым
принадлежит ещё и абстрактное свойство V; их будем называть
«\emph{областями $\mathfrak H$ из алгебраически целых трансцендентных чисел}».

Пусть $\mathfrak H$ есть такая область, а $H$ --- алгебраическая база всех чисел
из $\mathfrak H$; по II, $H$ одновременно является алгебраической базой
\emph{всех} чисел. По III, I и V область $\mathfrak H$ содержит, помимо $H$,
также все целые алгебраические функции от $\eta$, то есть область Цермело
$\mathfrak G_\eta$. Следовательно, по аналогии с \S\ 1:
\emph{каждую область $\mathfrak H$ из алгебраически целых трансцендентных чисел
можно построить посредством алгебраической базы всех чисел так, что
$\mathfrak H$ содержит область Цермело $\mathfrak G_\eta$ как делитель.}

Теперь пусть $H$ будет какой-либо алгебраической базой всех чисел, и пусть
$\mathfrak M(H)$ обозначает совокупность всех областей $\mathfrak H$, содержащих
$H$. Каждая область $\mathfrak H$ из $\mathfrak M(H)$ содержит, как только что
доказано, область Цермело $\mathfrak G_\eta$ как делитель; а поскольку сама
$\mathfrak G_\eta$ является областью $\mathfrak H$, то $\mathfrak G_\eta$ есть
\emph{наибольший общий делитель, или пересечение, всех областей $\mathfrak H$
из $\mathfrak M(H)$}. Тем самым $\mathfrak G_\eta$ определена абстрактно.
Одновременно прямо следует, что $\mathfrak M(H)$ замкнута относительно
образования пересечений, то есть пересечение всех областей $\mathfrak H$
произвольной подсистемы из $\mathfrak M(H)$ снова принадлежит $\mathfrak M(H)$.
Действительно, как пересечение оно выполняет I и V; II и III следуют из того,
что $\mathfrak G_\eta$ является делителем; IV следует, поскольку пересечение
является делителем области $\mathfrak H$.

Как показывает пример функциональной области в \S\ 2, области $\mathfrak H$ в
общем случае содержат и дальнейшие функции помимо $\mathfrak G_\eta$. Пусть
$\psi(\eta)$ есть такая рациональная или алгебраическая функция от $\eta$. Тогда
по I к $\mathfrak H$ принадлежат также все целые рациональные комбинации $\psi$
и каждый раз конечного числа функций из $\mathfrak G_\eta$. Возникающую таким
образом область целостности, состоящую из всех многочленов от $\psi$ с
коэффициентами из $\mathfrak G_\eta$, обозначим $[\mathfrak G_\eta,\psi]$, а
ведущий к ней процесс назовём «\emph{рационально-целым присоединением $\psi$ к
$\mathfrak G_\eta$}». По V к $\mathfrak H$ далее принадлежат все функции,
алгебраически целые над $[\mathfrak G_\eta,\psi]$; возникающую таким образом
область обозначим $\{\mathfrak G_\eta,\psi\}$, а соответствующий процесс
назовём «\emph{алгебраически-целым присоединением $\psi$ к
$\mathfrak G_\eta$}». Область $\{\mathfrak G_\eta,\psi\}$ состоит из всех
целых алгебраических функций от $\psi$ с коэффициентами из $\mathfrak G_\eta$ и
потому по известным рассуждениям\footnote{Ср., например, Вебер, \S\ 97, 6 и
7.} является областью целостности.

Покажем, что $\{\mathfrak G_\eta,\psi\}$ также интегрально замкнута, то есть
выполняет условие V. Действительно, пусть $z$ алгебраически цела над
$\{\mathfrak G_\eta,\psi\}$ посредством уравнения
\[
 f(z;a\cdots c)=z^\nu+a z^{\nu-1}+\cdots+c=0,
\]
где $a\cdots c$ как величины из $\{\mathfrak G_\eta,\psi\}$ удовлетворяют
уравнениям с коэффициентами из $[\mathfrak G_\eta,\psi]$:
\[
\begin{gathered}
 a^\lambda+A_1a^{\lambda-1}+\cdots+A_\lambda=0,\\
 \vdots\\
 c^\mu+C_1c^{\mu-1}+\cdots+C_\mu=0,
\end{gathered}
\]
и пусть их корни заданы как
$a,a'\cdots a^{(\lambda-1)}\cdots c,c'\cdots c^{(\mu-1)}$. Если теперь
образовать произведение по всем комбинациям корней,
\[
 g(z)=f(z;a\cdots c)f(z;a'\cdots c')\cdots
      f(z;a^{(\lambda-1)}\cdots c^{(\mu-1)}),
\]
то $g(z)$ имеет единицу старшим коэффициентом, а остальными коэффициентами ---
величины из $[\mathfrak G_\eta,\psi]$; следовательно, $z$ принадлежит
$\{\mathfrak G_\eta,\psi\}$.\footnote{Итак, если алгебраическая целость
определяется через \emph{какое-либо} уравнение, понятие неприводимости
относительно основного поля становится ненужным для доказательства I и V. Ср.
также \S\ 2, где только для доказательства IV нужна вспомогательная лемма,
переходящая от определения алгебраической целости через произвольное уравнение
к неприводимому. Поскольку эта вспомогательная лемма в настоящем случае вообще
уже не действует, IV приходится налагать на функцию $\psi$ как особое условие.}

Итак, область $\{\mathfrak G_\eta,\psi\}$ имеет свойства I и V; а поскольку
$\mathfrak G_\eta$ является её делителем, она имеет также II и III. Выполнено
ли IV, зависит от выбора $\psi$: например, IV не выполнено для
$\psi=\frac{1}{n\cdot\eta}$, потому что тогда $\frac1n$ принадлежит области;
зато по \S\ 2 оно выполнено для $\psi=\frac1\eta$, а также для
$\psi=\frac{\eta}{n}$. Действительно, в последнем случае, если функция из
$\{\mathfrak G_\eta,\psi\}$ представляет алгебраическое число, её значение
сохраняется при $\eta=0$; тем самым она переходит в функцию из
$\mathfrak G_\eta$, а значит, в алгебраически целое число.

\emph{Следовательно, $\{\mathfrak G_\eta,\psi\}$ становится областью
$\mathfrak H$, как только на $\psi$ наложено условие, что алгебраически-целое
присоединение не вводит в область никакого алгебраически-дробного числа.}

Совершенно аналогично определяется рационально-целое, соответственно
алгебраически-целое, присоединение произвольной системы $\mathfrak S$
рациональных или алгебраических функций от $\eta$ к $\mathfrak G_\eta$.
Рационально-целое присоединение ведёт к области целостности
$[\mathfrak G_\eta,\mathfrak S]$, состоящей из всех целочисленных многочленов
от каждый раз конечного числа величин из $\mathfrak G_\eta$ и $\mathfrak S$.
Алгебраически-целое присоединение даёт область
$\{\mathfrak G_\eta,\mathfrak S\}$, состоящую из всех величин, алгебраически
целых над $[\mathfrak G_\eta,\mathfrak S]$. Дословно как выше показывается, что
область целостности $\{\mathfrak G_\eta,\mathfrak S\}$ интегрально замкнута, а
также что выполнены II и III. Поэтому, как и выше:

\emph{$\{\mathfrak G_\eta,\mathfrak S\}$ становится областью $\mathfrak H$, как
только на $\mathfrak S$ наложено условие, что алгебраически-целое
присоединение не вводит в область никакого алгебраически-дробного числа} ---
то есть когда $\mathfrak S$ становится допустимой системой.\footnote{В общем
случае аналогично показывается: если $\mathfrak T$ есть произвольная область, а
$\{\mathfrak T\}$ обозначает совокупность величин, алгебраически целых над
$\mathfrak T$, то $\{\mathfrak T\}$ интегрально замкнута.}

Теперь важно, что верно и обратное утверждение: области
$\{\mathfrak G_\eta,\mathfrak S\}$ действительно исчерпывают все области
$\mathfrak H$ из $\mathfrak M(H)$, когда $\mathfrak S$ пробегает все допустимые
системы. Действительно, пусть $\mathfrak H$ будет произвольной областью из
$\mathfrak M(H)$, и пусть $\mathfrak L=\mathfrak H-\mathfrak G_\eta$ обозначает
систему, образованную из всех функций из $\mathfrak H$, не принадлежащих
$\mathfrak G_\eta$. По V область $\mathfrak H$ содержит
$\{\mathfrak G_\eta,\mathfrak L\}$ как делитель; но
$\{\mathfrak G_\eta,\mathfrak L\}$ также делится на $\mathfrak H$, потому что
она содержит $\mathfrak G_\eta$ и $\mathfrak L$, а так как они не имеют общего
элемента, содержит и $\mathfrak H$. Следовательно,
$\mathfrak H=\{\mathfrak G_\eta,\mathfrak L\}$; а поскольку IV выполнено для
$\mathfrak H$, система $\mathfrak L$, конечно, является допустимой.\footnote{Уже
присоединения некоторой подсистемы из $\mathfrak L$ может быть достаточно;
например, функциональная область из \S\ 2 возникает через алгебраически-целое
присоединение всех целых рациональных функционалов, за исключением многочленов,
к $\mathfrak G_\eta$.}

Когда $H$ пробегает все возможные алгебраические базы, то, как доказано в
начале, $\mathfrak M(H)$ пробегает совокупность всех областей $\mathfrak H$.
Следовательно, вопрос о самых общих областях $\mathfrak H$ из алгебраически
целых трансцендентных чисел полностью решён результатами этого параграфа:
\begin{enumerate}
\item[a)] \emph{Пересечение всех областей $\mathfrak H$ из $\mathfrak M(H)$
задаётся областью Цермело $\mathfrak G_\eta$, которая сама является областью
$\mathfrak H$; $\mathfrak M(H)$ замкнута относительно образования пересечений.}
\item[b)] \emph{Каждая область $\mathfrak H$ из $\mathfrak M(H)$ возникает через
алгебраически-целое присоединение допустимой системы $\mathfrak S$ к
$\mathfrak G_\eta$. Когда $H$ пробегает все возможные алгебраические базы,
$\mathfrak M(H)$ пробегает совокупность всех областей $\mathfrak H$.}\footnote{Допустимые
системы $\mathfrak S$ можно найти следующим образом. Подчинить все
алгебраически-дробные функции от $\eta$ некоторому вполне упорядочению $\Omega$
и взять в $\mathfrak S$ все функции, кроме тех, которые, будучи
алгебраически-цело присоединены к $\mathfrak G_\eta$ и ранее принятым функциям,
порождают алгебраически-дробное число. Эту процедуру можно прервать в любом
месте. Когда $\Omega$ пробегает все вполне упорядочения и для каждого
фиксированного $\Omega$ процедура прерывается в любом месте, этим
исчерпываются все допустимые системы $\mathfrak S$.}
\end{enumerate}
% END UNIT BODY
% END PRODUCER UNIT 062
\endgroup


\clearpage
\begingroup
\setcounter{footnote}{0}
% BEGIN PRODUCER UNIT 063
% Source: C:/Users/Floris/Documents/interlanguage/03_projects/noether/02_slavic_working_corpus/translations/paper09/russian/v001/Noether_Paper09_Section05_Russian_v001.tex
% Identity: 8470 B / SHA-256 AE608DBF09EFA287F32A61D320B17315CC73707610039ECB12CAB10E8311305B
\setlength{\parindent}{1.25em}
\setlength{\parskip}{0pt}
\emergencystretch=18em
\allowdisplaybreaks
\raggedbottom
\providecommand{\vdotswithin}[1]{\mathrel{\vdots}}
% BEGIN UNIT BODY
\section*{\S\ 5. Самые общие области из целых трансцендентных чисел при выборе алгебраической базы.}

Пусть $H$ есть какая-либо алгебраическая база всех чисел, и пусть
$\mathfrak M(\mathfrak G)$ обозначает совокупность всех областей
$\mathfrak G$, которые содержат $H$ и потому также $[H]$. Когда $H$ пробегает
все возможные базы, $\mathfrak M(\mathfrak G)$ пробегает совокупность всех
областей $\mathfrak G$; поэтому, как и в \S\ 4, можно ограничиться рассмотрением
областей $\mathfrak G$ из $\mathfrak M(\mathfrak G)$.

Все области $\mathfrak G$ из $\mathfrak M(\mathfrak G)$ содержат $[H]$ как
делитель; но, поскольку сам $[H]$ не является областью $\mathfrak G$, нельзя, как
в \S\ 4, заключить, что $[H]$ является наибольшим общим делителем. Что это всё же
так, показывает счётное множество областей $\mathfrak G$ из
$\mathfrak M(\mathfrak G)$, возникающих небольшим обобщением модульной области
из \S\ 3 и действительно не имеющих ни одного общего элемента, кроме $[H]$.

\emph{Пусть для каждого фиксированного $\nu$ области $\mathfrak G_\nu$
$(\nu=1,2,\ldots)$ состоят из всех величин}
\begin{equation}
z=f(\eta)+g_1(\eta)\eta_1^\nu+g_2(\eta)\eta_2^\nu+\cdots+g_t(\eta)\eta_t^\nu,
\tag{1}
\end{equation}
\emph{где $f(\eta)$ пробегает все целочисленные многочлены от $\eta$;
$\eta_1,\ldots,\eta_t$ пробегают все базовые величины, а
$g_1(\eta),\ldots,g_t(\eta)$ при каждой фиксированной комбинации
$\eta_1^\nu,\ldots,\eta_t^\nu$ по отдельности пробегают все целые
алгебраические функции от $\eta$.}

Как и в \S\ 3, видно, что областям $\mathfrak G_\nu$ принадлежат абстрактные
свойства I--IV; область $\mathfrak G_\nu$ содержит только целые алгебраические
функции и при $\nu=1$ совпадает с областью из \S\ 3. Теперь покажем, что области
$\mathfrak G_\nu$ не имеют никакого другого общего элемента, кроме многочленов
$f(\eta)$ из $[H]$. Действительно, пусть $z$ есть произвольная целая
алгебраическая, но не рациональная, функция от $\eta$, удовлетворяющая
неприводимому относительно $[H]$ уравнению
\begin{equation}
z^\chi+a_1(\eta)z^{\chi-1}+a_2(\eta)z^{\chi-2}+\cdots+a_\chi(\eta)=0,
\tag{2}
\end{equation}
где $\chi>1$; и пусть $\lambda$ есть наибольшая из степеней многочленов
$a_1(\eta),\ldots,a_\chi(\eta)$. Если $z$ принадлежит $\mathfrak G_\nu$, то
величина $\xi=z-f(\eta)$ удовлетворяет также неприводимому относительно $[H]$
уравнению
\begin{equation}
\xi^\chi+b_1(\eta)\xi^{\chi-1}+b_2(\eta)\xi^{\chi-2}+\cdots+b_\chi(\eta)=0,
\tag{3}
\end{equation}
причём $b_1(\eta),b_2(\eta),\ldots,b_\chi(\eta)$ как симметрические функции от
$\xi$ содержат, согласно (1), среди членов наименьшей степени члены по крайней
мере степеней $\nu,2\nu,\ldots,\chi\nu$. Сравнение (2) и (3) даёт
\[
 b_1(\eta)=a_1(\eta)+\chi f(\eta).
\]
Следовательно, если величина $\xi=z+\frac{a_1(\eta)}{\chi}$ удовлетворяет
уравнению
\begin{equation}
\xi^\chi+c_2(\eta)\xi^{\chi-2}+\cdots+c_\chi(\eta)=0,
\tag{4}
\end{equation}
то
\begin{equation}
 \xi-\zeta=\frac{b_1(\eta)}{\chi};
 \qquad
 c_i(\eta)\equiv b_i(\eta)\pmod{\frac{b_1(\eta)}{\chi}},
\tag{5}
\end{equation}
где $b_1(\eta)$ начинается членами по крайней мере степени $\nu$. Теперь
$c_i(\eta)$, как показывает сравнение (2) с (4), имеют степень $\chi$ по
коэффициентам $a_i(\eta)$ из (2), а значит, степень не выше $\chi\lambda$ по
$\eta$; с другой стороны, все $b_i(\eta)$ начинаются членами по крайней мере
степени $\nu$. Поэтому, если выбрать $\nu>\chi\lambda$, из (5) следует
\[
 c_i(\eta)=0;\qquad
 \xi^\chi=0\quad\text{или}\quad
 \left(z+\frac{a_1(\eta)}{\chi}\right)^\chi=0,
\]
то есть $z$, вопреки предположению, становится рациональной функцией от $\eta$.
Итак:

\emph{Если $z$ есть целая алгебраическая, не рациональная функция от $\eta$, то
$z$ не может принадлежать никакой области $\mathfrak G_\nu$ при
$\nu>\chi\lambda$, или:}

\emph{Пересечение всех областей $\mathfrak G$ из $\mathfrak M(\mathfrak G)$
задаётся областью целостности $[H]$, которая сама не является областью
$\mathfrak G$; следовательно, $\mathfrak M(\mathfrak G)$ не замкнута
относительно образования пересечений.}

Поэтому и конструкцию самых общих областей $\mathfrak G$ посредством
алгебраической базы труднее обозреть, чем конструкцию областей $\mathfrak H$.
Области $[H]$ принадлежат свойства I, III и IV; поэтому, если рационально-цело
присоединить произвольную систему $\mathfrak S$, то для возникающей области
целостности $[H,\mathfrak S]$ свойства I и III сохраняются, тогда как IV
становится особым условием на $\mathfrak S$. Чтобы возникла область
$\mathfrak G$, нужно также наложить на $\mathfrak S$ как условие свойство II;
или, иначе говоря: для каждого конечного над $(H)$\footnote{$(H)$ означает поле,
состоящее из всех рациональных функций от $\eta$ с алгебраическими числовыми
коэффициентами.} поля $\mathfrak K$ должно существовать конечное над
$\mathfrak K$ поле $\mathfrak L$ такое, что $[H,\mathfrak S]$ содержит
примитивный элемент из $\mathfrak L$.\footnote{Ср. по этому поводу \S\ 7.}
Обратно, всякую область $\mathfrak G$ из $\mathfrak M(\mathfrak G)$ также можно
представить в форме $[H,\mathfrak S]$; тем самым все области $\mathfrak G$ из
$\mathfrak M(\mathfrak G)$ получаются, когда $\mathfrak S$ пробегает все системы,
ограниченные таким образом. Структура систем $\mathfrak S$ проще выяснится в
следующем параграфе посредством «рациональной базы»; вместе с тем будет получена
полная аналогия с теоремами \S\ 4.
% END UNIT BODY
% END PRODUCER UNIT 063
\endgroup


\clearpage
\begingroup
\setcounter{footnote}{0}
% BEGIN PRODUCER UNIT 064
% Source: C:/Users/Floris/Documents/interlanguage/03_projects/noether/02_slavic_working_corpus/translations/paper09/russian/v001/Noether_Paper09_Section06_Russian_v001.tex
% Identity: 12025 B / SHA-256 86BBD4419D4CEFC6EDFFA3C54EA829B743B2ECD4C54DD571BBC4395ACA348186
\setlength{\parindent}{1.25em}
\setlength{\parskip}{0pt}
\emergencystretch=18em
\allowdisplaybreaks
\raggedbottom
\providecommand{\vdotswithin}[1]{\mathrel{\vdots}}
% BEGIN UNIT BODY
\section*{\S\ 6. Самые общие области из целых трансцендентных чисел при выборе рациональной базы.}

По аналогии с алгебраической базой рациональную базу следует определить так:
«Система $\Theta$ чисел $\vartheta$ называется рациональной базой всех чисел,
если $\Theta$ позволяет выразить каждое число рационально --- с коэффициентами
из $K$\footnote{Под рациональными функциями всегда понимаются функции с
алгебраическими числовыми коэффициентами; эти алгебраические числовые
коэффициенты включены в определение из-за III и IV. Более соответствующим было
бы требовать рациональные числовые коэффициенты как в III и IV, так и в
определении рациональной базы. Выводы \S\ 6 и \S\ 7 при этом дословно
сохраняются, если только поле $K$ всех алгебраических чисел заменить полем всех
рациональных чисел.} --- тогда как при по меньшей мере одном вполне
упорядочении никакое базовое число не выражается рационально --- с
коэффициентами из $K$ --- через конечное число предшествующих базовых
чисел».\footnote{В отличие от линейной и алгебраической базы, здесь порядок
элементов играет роль. В порядке $\eta,\sqrt[\nu]{\eta}$, например, надо взять
оба элемента, а в обратном порядке только $\sqrt[\nu]{\eta}$.}

Пусть область $\mathfrak G$ из целых трансцендентных чисел подчинена вполне
упорядочению $\Omega$. Тогда может случиться, что величина $z$ из
$\mathfrak G$ рационально выражается через конечное число предшествующих
величин из $\mathfrak G$, то есть удовлетворяет уравнению
$z=\varphi(\xi)$, где $\varphi(\xi)$ есть рациональная функция с
коэффициентами из $K$ от $\xi_1,\ldots,\xi_t$; в частности, каждое
алгебраическое число $\alpha$ области удовлетворяет уравнению $z=\alpha$.
Остальные величины $\vartheta$ из $\mathfrak G$ образуют систему $\Theta$,
которую надо доказать как рациональную базу всех величин из $\mathfrak G$.
Действительно, по предположению каждая величина $\vartheta$ рационально
независима от предшествующих величин из $\Theta$. Индукционный вывод\footnote{Ср.
Zermelo, указ. соч., \S\ 1.} далее показывает, что всякое отношение
$z=\varphi(\xi)$ влечет отношение $z=\psi(\vartheta)$; иначе должен был бы
существовать первый элемент
$z_0=\varphi_0(\xi_1,\ldots,\xi_t)$, который не выражается рационально через
$\vartheta$, тогда как предшествующие $\xi_1,\ldots,\xi_t$ являются
рациональными функциями от $\vartheta$; так получается противоречие. По II
система $\Theta$ вместе с тем является рациональной базой всех чисел; по III и
I область $\mathfrak G$ наряду с $\Theta$ содержит как делитель область
целостности $[\Theta]$, состоящую из всех целочисленных многочленов от
$\vartheta$, тогда как по IV $[\Theta]$ не содержит никакого
алгебраически-дробного числа. Следовательно, рациональная база всех чисел
$\Theta$ удовлетворяет дополнительному условию: выведенная из $\Theta$ область
целостности $[\Theta]$ не содержит никакого алгебраически-дробного
числа;\footnote{Что это действительно является условием, видно, например, из
того, что две величины $\eta,\frac1{2\sqrt{\eta}}$ недопустимы как базовые
элементы, а величины $\eta,\frac1{\sqrt{\eta}}$ допустимы.} $\Theta$ есть
\emph{рациональная база с дополнительным условием} для всех чисел. Поэтому по
аналогии с \S\ 2 верно: \emph{каждая область $\mathfrak G$ может быть построена
посредством рациональной базы с дополнительным условием так, что
$\mathfrak G$ содержит область целостности $[\Theta]$ как делитель.}

Пусть $\Theta$ есть какая-либо рациональная база с дополнительным условием для
всех чисел, существование которой тем самым обеспечено существованием областей
$\mathfrak G$. Обозначим через $\mathfrak N(\mathfrak G)$ совокупность всех
областей $\mathfrak G$, которые содержат $\Theta$ и, следовательно,
$[\Theta]$. Когда $\Theta$ пробегает любую базу с дополнительным условием,
$\mathfrak N(\mathfrak G)$ по сказанному выше пробегает совокупность всех
областей $\mathfrak G$; поэтому мы снова можем ограничиться областями
$\mathfrak G$ из $\mathfrak N(\mathfrak G)$.

Теперь сам $[\Theta]$ является областью $\mathfrak G$: ведь каждое число
выражается рационально через $\Theta$, то есть как частное двух многочленов из
$[\Theta]$; кроме того, $[\Theta]$ по своей конструкции удовлетворяет и
остальным условиям I, III, IV. \emph{Следовательно, $[\Theta]$ является
наибольшим общим делителем или пересечением всех областей $\mathfrak G$ из
$\mathfrak N(\mathfrak G)$; кроме того, $\mathfrak N(\mathfrak G)$ замкнута
относительно образования пересечений, то есть и пересечение всех областей
$\mathfrak G$ любой подсистемы из $\mathfrak N(\mathfrak G)$ принадлежит
$\mathfrak N(\mathfrak G)$.} Действительно, для пересечения выполнено I; II и
III следуют из того, что $[\Theta]$ является делителем; IV --- из того, что
область-пересечение содержится как делитель в некоторой области $\mathfrak G$.

Рационально-целое присоединение произвольной системы $\mathfrak S$
рациональных функций от $\vartheta$ к $[\Theta]$ --- каждая алгебраическая
функция от $\vartheta$ благодаря свойству рациональной базы рационально
выражается через некоторые другие $\vartheta$ --- дает область целостности
$[\Theta,\mathfrak S]$, которой принадлежат свойства I, II, III и которая
поэтому становится областью $\mathfrak G$, как только $\mathfrak S$ является
«допустимой системой», то есть удовлетворяет условию, что через
рационально-целое присоединение в область не входит никакое
алгебраически-дробное число. Обратно, каждая область $\mathfrak G$ из
$\mathfrak N(\mathfrak G)$ допускает представление $[\Theta,\mathfrak S]$,
если только $\mathfrak S$ означает остаточную систему $\mathfrak G-[\Theta]$.
Следовательно, для самых общих областей $\mathfrak G$ в полной аналогии с
\S\ 4 справедливы результаты:
\begin{enumerate}
\item[a)] \emph{Пересечение всех областей $\mathfrak G$ из
$\mathfrak N(\mathfrak G)$ задается областью целостности $[\Theta]$, которая
сама является областью $\mathfrak G$; $\mathfrak N(\mathfrak G)$ замкнута
относительно образования пересечений.}
\item[b)] \emph{Каждая область $\mathfrak G$ из $\mathfrak N(\mathfrak G)$
возникает через рационально-целое присоединение допустимой системы
$\mathfrak S$ к $[\Theta]$. Когда $\Theta$ пробегает каждую возможную
рациональную базу с дополнительным условием, $\mathfrak N(\mathfrak G)$
пробегает совокупность всех областей $\mathfrak G$.}\footnote{О самых общих
допустимых системах $\mathfrak S$ остается в силе сказанное в \S\ 4, если
только заменить «алгебраически» на «рационально».}
\end{enumerate}

\emph{Замечание:} Если из $\Theta$ выбрать алгебраическую базу $H$ всех величин
из $\Theta$, то $H$ вместе с тем становится алгебраической базой всех чисел;
обратно, всякую алгебраическую базу можно дополнить до рациональной, поставив
$H$ в начало вполне упорядочения. Поэтому построение областей $\mathfrak G$,
данное в \S\ 5 посредством алгебраической базы, можно истолковать так: систему
$\mathfrak S$, присоединяемую к $[H]$, разлагают на две подсистемы
$\mathfrak S_1$ и $\mathfrak S_2$; присоединение $\mathfrak S_1$ сводится к
дополнению $H$ до рациональной базы с дополнительным условием, к которой затем
еще рационально-цело присоединяется допустимая система $\mathfrak S_2$.
% END UNIT BODY
% END PRODUCER UNIT 064
\endgroup


\clearpage
\begingroup
\setcounter{footnote}{0}
% BEGIN PRODUCER UNIT 065
% Source: C:/Users/Floris/Documents/interlanguage/03_projects/noether/02_slavic_working_corpus/translations/paper09/russian/v001/Noether_Paper09_Section07_Russian_v001.tex
% Identity: 16631 B / SHA-256 6F2B08CA51B5A7F035FA1C5B475EB38AF76A7B59666FC18EDBA66BA3CB271D17
\setlength{\parindent}{1.25em}
\setlength{\parskip}{0pt}
\emergencystretch=18em
\allowdisplaybreaks
\raggedbottom
\providecommand{\vdotswithin}[1]{\mathrel{\vdots}}
% BEGIN UNIT BODY
\section*{\S\ 7. Построение самой общей рациональной базы с дополнительным условием.}

Результаты предыдущего параграфа нуждаются в дополнении в том отношении, что
там было доказано \emph{существование} рациональной базы с дополнительным
условием для всех чисел, но не метод ее построения --- при заданном вполне
упорядочении. Из-за дополнительного условия это построение становится сложнее,
когда вместо области $\mathfrak G$ дано поле всех комплексных чисел.\footnote{В
\S\ 6, b) было бы достаточно заставлять $\Theta$ пробегать только такие
специальные рациональные базы, которые состоят из алгебраической базы и
алгебраически целых функций базовых величин; тогда дополнительное условие
выполняется само собой. Чтобы это увидеть, надо лишь по \S\ 6 дополнить
алгебраическую базу из $\mathfrak G$ до рациональной и в полученной базе каждую
алгебраически-дробную функцию от $\eta$ умножить на подходяще выбранный
многочлен от $\eta$; при этом базовое свойство сохраняется. Это построение без
изменений переносится на поле всех комплексных чисел. Однако возникающий из
\S\ 6 вопрос о самой общей рациональной базе с дополнительным условием имеет и
самостоятельный интерес.}

Пусть поле всех комплексных чисел подчинено вполне упорядочению $\Omega$.
Тогда могут появляться величины трех типов:
\begin{enumerate}
\item[1.] Величины $y$, рационально-целое присоединение которых к
предшествующим --- рекурсивно определенным в 3. --- базовым
величинам\footnote{Как показывает индукционный вывод, такое рекурсивное
определение возможно.} вызывает вхождение алгебраически-дробного числа в
возникающую область целостности; если никакая базовая величина не предшествует,
эта область состоит из всех целочисленных многочленов от $y$. В частности, все
алгебраически-дробные числа являются величинами $y$, поскольку возникающая
область целостности содержит величины $y=\beta$.
\item[2.] Величины $z$, которым свойство 1. не принадлежит, но которые
выражаются рационально через конечное число предшествующих чисел, отличных от
$y$, то есть удовлетворяют отношению
$z=\varphi(\xi_1,\ldots,\xi_t)$, где $\varphi$ есть рациональная функция с
коэффициентами из $K$, а среди $\xi$ нет никаких $y$. В частности, сюда
относятся все целые алгебраические числа $\alpha$, поскольку им не принадлежит
1. и они удовлетворяют отношению $z=\alpha$.
\item[3.] Величины $\vartheta$, для которых не выполняется ни 1., ни 2. и
которые следует называть базовыми величинами. В частности, первое во вполне
упорядочении трансцендентное число $\vartheta_0$ является такой базовой
величиной: ведь целочисленные многочлены от $\vartheta_0$ не могут представить
алгебраически-дробное число, а $\vartheta_0$ также не может быть рационально
зависимым от предшествующих алгебраических чисел.
\end{enumerate}

\emph{Теперь надо доказать, что система $\Theta$ всех базовых величин
$\vartheta$ является рациональной базой с дополнительным условием для всех
чисел.}

Так как среди $\vartheta$ нет никакой величины $y$, область целостности
$[\Theta]$ не может содержать алгебраически-дробное число; следовательно,
дополнительное условие выполнено. Так как, далее, среди $\vartheta$ нет никакой
величины $z$, каждая $\vartheta$ рационально независима от предшествующих
базовых величин. Индукционный вывод, точно как в \S\ 6, далее показывает, что
всякое отношение $z=\varphi(\xi)$ --- где среди $\xi$ нет величин $y$ --- влечет
отношение $z=\psi(\vartheta)$, так что все величины $z$ рационально выражаются
через $\vartheta$.

Остается только более трудное доказательство того, что и величины $y$
рационально выражаются через $\vartheta$. Сначала покажем, что $y$ алгебраически
зависит от $\vartheta$. Действительно, по предположению для каждого $y$
существует по меньшей мере одно алгебраически-дробное число $\beta$, допускающее
представление
\begin{equation}
\beta=f_0(\vartheta)+f_1(\vartheta)y+\cdots+f_\chi(\vartheta)y^\chi,
\tag{1}
\end{equation}
где $f_0(\vartheta),\ldots,f_\chi(\vartheta)$ являются многочленами из
$[\Theta]$ и поэтому не могут представлять алгебраически-дробное число. Если бы
теперь $y$ было алгебраически независимым от $\vartheta$, то (1) выполнялось бы
тождественно по $y$, и мы имели бы $\beta=f_0(\vartheta)$, вопреки свойствам
$[\Theta]$.

Чтобы из алгебраической зависимости вывести рациональную, выберем из $\Theta$
алгебраическую базу $H$ всех величин из $\Theta$. Тогда каждое $y$, как
алгебраическая функция от $\vartheta$, является также алгебраически зависимым от
$\eta$; значит, через $y=y(\eta)$ определяется алгебраическое поле
$(H,y(\eta))$ над $(H)$. Доказательство рациональной зависимости --- поскольку
все $\eta$ вместе с тем являются величинами $\vartheta$ --- сводится к тому,
чтобы показать: для каждого такого поля существуют примитивные элементы,
которые являются рациональными функциями от $\vartheta$. Точнее, будет
показано, что после умножения $y(\eta)$ на некоторый многочлен от $\eta$ она
теряет свойство 1. и потому переходит в такой элемент.

Как известно, между $(H)$ и $(H,y,y',\ldots,y^{(\lambda_0-1)})$ лежит только
конечное число промежуточных полей. Пусть
$\Omega_0=(H),\Omega_1,\ldots,\Omega_\sigma$ суть те из них,
которые имеют примитивный элемент, рационально выражаемый через $\vartheta$, так
что каждый элемент поля становится рациональной функцией от $\vartheta$; это
условие по меньшей мере для $\Omega_0=(H)$ всегда выполнено. Пусть примитивная
функция, соответствующая $\Omega_i$, задана как
$\frac{g_i(\vartheta)}{h_i(\vartheta)}$, где $g_i,h_i$ --- многочлены из
$[\Theta]$. Тогда --- если под $\alpha_i$ понимать подходяще выбранное целое
число --- многочлен из $[\Theta]$
\[
 \xi_i=g_i(\vartheta)+\alpha_i h_i(\vartheta)
\]
становится примитивной функцией от $\Omega_i$ или от алгебраического поля над
$\Omega_i$,\footnote{Ср. об этом конец \S\ 5.} и поэтому каждая функция
$\psi(\eta)$ из $\Omega_i$ допускает представление
\begin{equation}
\psi(\eta)=
\frac{a_1(\eta)\xi_i^{\kappa-1}+a_2(\eta)\xi_i^{\kappa-2}+\cdots+a_\kappa(\eta)}
     {a(\eta)}
=\frac{A(\eta,\xi_i)}{a(\eta)},
\tag{2}
\end{equation}
где $a(\eta),a_1(\eta),\ldots,a_\kappa(\eta)$ являются целочисленными
многочленами от $\eta$, а значит $A(\eta,\xi_i)$ и $a(\eta)$ являются
величинами из $[\Theta]$.

Пусть неприводимое относительно $\Omega_i$ уравнение, которому удовлетворяет
$y$, благодаря (2) имеет вид
\[
 f^{(i)}(\eta)y^{\lambda_i}+f_1^{(i)}(\eta,\xi_i)y^{\lambda_i-1}
+\cdots+f_{\lambda_i}^{(i)}(\eta,\xi_i)=0,
\]
и все коэффициенты этого уравнения принадлежат $[\Theta]$. Тогда функция
$y_i=f^{(i)}(\eta)\cdot y$ удовлетворяет также неприводимому относительно
$\Omega_i$ уравнению
\[
 y_i^{\lambda_i}+f_1^{(i)}(\eta,\xi_i)y_i^{\lambda_i-1}
+\cdots+\bigl(f^{(i)}(\eta)\bigr)^{\lambda_i-1}
 f_{\lambda_i}^{(i)}(\eta,\xi_i)=0,
\]
благодаря которому она вместе с тем алгебраически цело зависит от
$[H,\xi_i]$; то же верно для произведения $y_i$ на произвольный многочлен от
$\eta$. В частности, функция
\begin{equation}
Y=f^{(0)}(\eta)f^{(1)}(\eta)\cdots f^{(\sigma)}(\eta)\cdot y
\tag{3}
\end{equation}
алгебраически цело зависит от всех областей $[H,\xi_i]$ благодаря всякий раз
неприводимому относительно $\Omega_i$ уравнению
\begin{equation}
Y^{\lambda_i}+F_1^{(i)}(\eta,\xi_i)Y^{\lambda_i-1}
+\cdots+F_{\lambda_i}^{(i)}(\eta,\xi_i)=0
\qquad (i=0,1,\ldots,\sigma).
\tag{4}
\end{equation}
Теперь покажем, что $Y$ представляет искомый примитивный элемент, рационально
выражаемый через $\vartheta$. Пусть через $Y$ и $[\Theta]$ представлено,
например, алгебраическое число $\gamma$:
\begin{equation}
\gamma=k_0(\vartheta)+k_1(\vartheta)Y+\cdots+k_\nu(\vartheta)Y^\nu,
\tag{5}
\end{equation}
где $k_0(\vartheta),\ldots,k_\nu(\vartheta)$ являются многочленами из
$[\Theta]$. Теперь определим неприводимое уравнение, которому $Y$ удовлетворяет
относительно поля
$\mathfrak K=(H;k_0(\vartheta),\ldots,k_\nu(\vartheta))$, возникающего из
области коэффициентов в (5). Как известно, оно задается неприводимым уравнением
$Y$ относительно поля-пересечения поля $\mathfrak K$ и
$(H,Y,Y',\ldots,Y^{(\lambda_0-1)})$, лежащего между $(H)$ и
$(H,Y,Y',\ldots,Y^{(\lambda_0-1)})$. Теперь
$(H,Y,Y',\ldots,Y^{(\lambda_0-1)})$ по (3) тождественно с
$(H,y,y',\ldots,y^{(\lambda_0-1)})$; кроме того, все
элементы из $\mathfrak K$, а значит и все элементы поля-пересечения,
рационально выражаются через конечное число $\vartheta$. Поэтому последнее поле
необходимо задается одним из полей $\Omega_0,\Omega_1,\ldots,\Omega_\sigma$,
скажем $\Omega_\tau$. Искомое неприводимое уравнение, следовательно, по (4)
имеет степень $\lambda_\tau$ и получается в виде
\begin{equation}
Y^{\lambda_\tau}+F_1^{(\tau)}(\eta,\xi_\tau)Y^{\lambda_\tau-1}
+\cdots+F_{\lambda_\tau}^{(\tau)}(\eta,\xi_\tau)=0
\tag{6}
\end{equation}
с многочленами из $[\Theta]$ в качестве коэффициентов. Посредством (6)
представление (5) можно преобразовать в
\begin{equation}
\gamma=K_0(\vartheta)+K_1(\vartheta)Y+\cdots+
K_{\lambda_\tau-1}(\vartheta)Y^{\lambda_\tau-1},
\tag{7}
\end{equation}
где $K_0(\vartheta),K_1(\vartheta),\ldots$ принадлежат $\mathfrak K$ и, с
другой стороны, являются многочленами из $[\Theta]$. Но поскольку $\gamma$ есть
алгебраическое число, отношение (7) по $Y$ не достигает степени
$\lambda_\tau$ и поэтому выполняется \emph{тождественно} по $Y$. Следовательно,
\[
 \gamma=K_0(\vartheta),
\]
то есть $\gamma$ принадлежит $[\Theta]$ и, значит, является целым
алгебраическим числом.

Когда теперь $\gamma$ пробегает все алгебраические числа, представимые через
$Y$ и $[\Theta]$, поле-пересечение $(H,Y,Y',\ldots,Y^{(\lambda_0-1)})$ и
соответствующего поля $\mathfrak K$ всегда пробегает только поля
$\Omega_0,\Omega_1,\ldots,\Omega_\sigma$, и все приведенные выше выводы
сохраняются; то есть \emph{через рационально-целое присоединение $Y$ к
$[\Theta]$ не может быть представлено никакое алгебраически-дробное число}.
Следовательно, $Y$ уже не имеет свойства 1.; по 2) или 3) оно является
величиной $z$ или $\vartheta$ и поэтому рационально выражается через конечное
число $\vartheta$; то же по (3) верно для $y$: \emph{$\Theta$ доказана как
рациональная база всех чисел.} Обратно, если поставить $\Theta$ в начало вполне
упорядочения, всякая рациональная база с дополнительным условием может быть
построена по 1), 2), 3). Тем самым доказано: \emph{построение, заданное в
1), 2), 3), приводит к самой общей рациональной базе с дополнительным условием
для всех чисел.}
% END UNIT BODY
% END PRODUCER UNIT 065
\endgroup


\clearpage
\begingroup
\setcounter{footnote}{0}
% BEGIN PRODUCER UNIT 066
% Source: C:/Users/Floris/Documents/interlanguage/03_projects/noether/02_slavic_working_corpus/translations/paper09/russian/v001/Noether_Paper09_Section08_Russian_v001.tex
% Identity: 12474 B / SHA-256 38A35111ABAD8DD23F979095CB6C5883A48E36DD5070B2F4B411387744C41CBA
\setlength{\parindent}{1.25em}
\setlength{\parskip}{0pt}
\emergencystretch=18em
\allowdisplaybreaks
\raggedbottom
\providecommand{\vdotswithin}[1]{\mathrel{\vdots}}
% BEGIN UNIT BODY
\section*{\S\ 8. Разбиение областей $\mathfrak G$ и $\mathfrak H$ на классы.}

Наряду с вопросом о наиболее общих областях $\mathfrak G$ и $\mathfrak H$
возникает дальнейший вопрос о \emph{существенно отличных} областях, которые ---
в смысле, подлежащем ниже более точному уточнению, --- нельзя получить одним и
тем же принципом построения. Это приводит к разбиению этих областей на классы.

\emph{Две области будем называть принадлежащими к одному и тому же классу, или
изоморфными, если их элементы можно взаимно однозначно поставить в соответствие
так, что сумме и произведению любых двух элементов одной области всегда
соответствуют сумма и произведение соответствующих элементов другой области.}
Из этого определения следует, что при всяком изоморфном соответствии нуль и
единица соответствуют сами себе, а все алгебраические соотношения между
конечным числом элементов сохраняются для соответствующих элементов, и
обратно. В частности, наряду с единицей все рациональные числа также
соответствуют сами себе, тогда как совокупность алгебраических чисел --- и
точно так же совокупность алгебраически целых чисел --- переходит в себя,
хотя отдельному алгебраическому числу вполне может соответствовать сопряжённое.

Прежде всего заметим, что всякая область из совокупности
$\mathfrak M(\mathfrak G)$ всех областей $\mathfrak G$, содержащих фиксированную
алгебраическую базу $H$, изоморфна некоторой области из совокупности
$\mathfrak M^*(\mathfrak G)$, принадлежащей другой алгебраической базе $H^*$.
Действительно, поле всех комплексных чисел получается из всякой алгебраической
базы алгебраическим расширением и потому имеет ту же мощность, что и
база;\footnote{Steinitz, \S\S\ 23 и 24.} следовательно, $H$ и $H^*$ имеют
одинаковую мощность, и требуемое изоморфное отображение получается
сопоставлением базовых величин $\eta_i$ и $\eta_i^*$.\footnote{Этот вывод,
разумеется, не верен для множеств $\mathfrak N(\mathfrak G)$ и
$\mathfrak N^*(\mathfrak G)$, принадлежащих двум рациональным базам $\Theta$ и
$\Theta^*$; также из взаимной рациональной выразимости $\Theta$ и $\Theta^*$
ещё нельзя заключить об изоморфности $\mathfrak N(\mathfrak G)$ и
$\mathfrak N^*(\mathfrak G)$. Однако этим устанавливается изоморфное
соответствие между полями $(\Theta)$ и $(\Theta^*)$.} Напротив,
$\mathfrak M(\mathfrak G)$ содержит неизоморфные, или существенно отличные,
области $\mathfrak G$, как следует из примеров в \S\S\ 2, 3 и 5. В самом деле,
при изоморфном отображении все алгебраические соотношения сохраняются, поэтому
алгебраической базе снова должна соответствовать алгебраическая база; но в
\S\S\ 2 и 3 было показано, что функциональная и модульная области не обладают
такой алгебраической базой, для которой выполнялись бы все базовые свойства
области Цермело. Поскольку функциональная область из \S\ 2 алгебраически цело
замкнута, подмножество $\mathfrak M(\mathfrak H)$ также содержит существенно
отличные области $\mathfrak H$.

Теперь, по аналогии с построением базы, покажем, как из
$\mathfrak M(\mathfrak G)$ можно выделить такое подмножество
$\mathfrak L(\mathfrak G)$, что все области $\mathfrak G$ из
$\mathfrak L(\mathfrak G)$ существенно отличны, но каждая область из
$\mathfrak M(\mathfrak G)$ --- а следовательно, и всякая область
$\mathfrak G$ вообще --- изоморфна некоторой области из
$\mathfrak L(\mathfrak G)$. Пусть $\mathfrak M(\mathfrak G)$ подчинено
вполне упорядочению $\Omega$; тогда область $\mathfrak G$ из
$\mathfrak M(\mathfrak G)$ может быть изоморфна какой-либо предшествующей
области, а может и не быть. Области $\mathfrak G$, не изоморфные никакой
предшествующей, образуют множество $\mathfrak L(\mathfrak G)$, которое по
своему построению содержит только существенно отличные области $\mathfrak G$.
Индукционное рассуждение далее показывает, что всякая область из
$\mathfrak M(\mathfrak G)$ изоморфна некоторой области из
$\mathfrak L(\mathfrak G)$. Действительно, пусть $\mathfrak G_1$ была бы
первой областью, для которой это не выполнено. Тогда $\mathfrak G_1$ либо
принадлежит $\mathfrak L(\mathfrak G)$, либо изоморфна предшествующей области
$\mathfrak G_0$, которая по предположению изоморфна некоторой области из
$\mathfrak L(\mathfrak G)$; значит, то же самое было бы верно и для
$\mathfrak G_1$. Так как всякая область $\mathfrak G$ принадлежит некоторому
множеству $\mathfrak M^*(\mathfrak G)$ и потому изоморфна некоторой области из
$\mathfrak M(\mathfrak G)$, \emph{совокупность классов областей
$\mathfrak G$ тем самым представлена множеством $\mathfrak L(\mathfrak G)$.}

Если в области $\mathfrak G$ из $\mathfrak L(\mathfrak G)$ заменить базовые
величины $\eta$ величинами $\eta^*$ какой-либо другой базы, то, как показано
выше, возникает область, изоморфная $\mathfrak G$. Обратно, пусть
$\mathfrak G^*$ --- произвольная область и, следовательно, изоморфна некоторой
области $\mathfrak G$ из $\mathfrak L(\mathfrak G)$. Если при этом базовым
величинам $\eta$ соответствуют величины $\eta^*$ из $\mathfrak G^*$, то они
снова образуют алгебраическую базу, а $\mathfrak G^*$ содержит те же
алгебраические функции от $\eta^*$, что и $\mathfrak G$ от $\eta$, --- или их
сопряжённые. \emph{Следовательно, всякая область, изоморфная $\mathfrak G$,
возникает так: в $\mathfrak G$ и в её сопряжённых относительно
$(H)$}\footnote{Элементы этих областей, сопряжённых с $\mathfrak G$, по II
заведомо можно рационально выразить через элементы $\mathfrak G$, но не
обязательно цело и рационально; поэтому сопряжённые могут отличаться от
$\mathfrak G$.} \emph{--- если они отличны от $\mathfrak G$ --- алгебраическая
база $H$ пробегает все возможные алгебраические базы.} Из каждой фиксированной
области $\mathfrak G_i$ из $\mathfrak L(\mathfrak G)$ таким образом возникает
класс $\mathfrak R_i$, который содержит все и только те области $\mathfrak G$,
которые изоморфны $\mathfrak G_i$, но, возможно, каждую область по нескольку
раз или даже бесконечно много раз. Из $\mathfrak R_i$ указанным выше способом
можно выделить подмножество $\mathfrak R_i^*$, содержащее каждую область,
изоморфную $\mathfrak G_i$, один и только один раз. Когда $\mathfrak R_i^*$
пробегает все классы, получается совокупность всех областей $\mathfrak G$,
каждая область один и только один раз. Для характеристики класса
$\mathfrak R_i$ вместо $\mathfrak G_i$ можно, конечно, взять любую другую
область этого класса, из которой, как выше, выводятся все области из
$\mathfrak R_i$. Итак, получаем:

\emph{Из $\mathfrak M(\mathfrak G)$ можно выделить подмножество
$\mathfrak L(\mathfrak G)$ так, что области $\mathfrak G$ из
$\mathfrak L(\mathfrak G)$ взаимно однозначно соответствуют совокупности
классов (без повторения). При этом из $\mathfrak G_i$ соответствующий класс
$\mathfrak R_i$ возникает тем, что в $\mathfrak G_i$ и в его сопряжённых
относительно $(H)$ база $H$ пробегает все возможные алгебраические базы. Если
$\mathfrak R_i^*$ означает все попарно различные области из
$\mathfrak R_i$, и $\mathfrak R_i^*$ пробегает все классы, то возникает
совокупность всех областей $\mathfrak G$, каждая область один и только один
раз.}\footnote{При этом $\mathfrak M(\mathfrak G)$ следует строить по
\S\S\ 6 и 7: по \S\ 7 алгебраическую базу $H$ дополняют до каждой
рациональной базы $\Theta$ с дополнительным условием, возникающей из неё, и для
каждой $\Theta$ по \S\ 6 b) строят совокупность $\mathfrak N(\mathfrak G)$
всех областей $\mathfrak G$, содержащих $\Theta$.}

Аналогичное утверждение верно для областей $\mathfrak H$ из алгебраически
целых трансцендентных чисел.
% END UNIT BODY
% END PRODUCER UNIT 066
\endgroup


\clearpage
\begingroup
\setcounter{footnote}{0}
% BEGIN PRODUCER UNIT 067
% Source: C:/Users/Floris/Documents/interlanguage/03_projects/noether/02_slavic_working_corpus/translations/paper09/russian/v001/Noether_Paper09_Section09_Russian_v001.tex
% Identity: 11347 B / SHA-256 BF1D8D426049D41F5DED9B56250278149CE967C381B43B9FB86D6A6F59E95A84
\setlength{\parindent}{1.25em}
\setlength{\parskip}{0pt}
\emergencystretch=18em
\allowdisplaybreaks
\raggedbottom
\providecommand{\vdotswithin}[1]{\mathrel{\vdots}}
% BEGIN UNIT BODY
\section*{\S\ 9. Пример счётно бесконечного числа классов областей $\mathfrak G$.}

Мощность конструкций, указанных в \S\ 4 b) и \S\ 6 b), из соответствующих
замечаний читается как равная мощности всех вполне упорядочений, то есть равная
$2^c$, если $c$ обозначает мощность континуума. Но отсюда нельзя сделать вывод
ни о мощности совокупности всех областей $\mathfrak G$, если каждую область
считать один раз и только один раз, ни тем более о мощности всех классов. Что
число классов во всяком случае не конечно, покажем теперь на примере счётно
бесконечного числа областей $\mathfrak G$ --- аналогичных областям из \S\ 5,
--- которые все окажутся существенно отличными и потому дадут счётно
бесконечно много классов.\footnote{Области $\mathfrak G$ из \S\ 6 также дают
счётно бесконечно много классов; однако доказательство несколько сложнее.}

Аналогично \S\ 5, (1), пусть область $\mathfrak G_\sigma$ состоит из
совокупности величин
\begin{equation}
 z=a_0+a_1(\eta)+\cdots+a_{\sigma-1}(\eta)
 \pmod{\mathfrak M_\sigma(\eta)},
\tag{1}
\end{equation}
\emph{где $a_i(\eta)$ обозначают однородные формы $i$-го измерения от $\eta$,
а модуль $\mathfrak M_\sigma$ содержит в качестве базы все степенные
произведения $\sigma$-го измерения от $\eta$\footnote{Разумеется, в каждом
отдельном $z$ в модуле встречается лишь конечное число степенных произведений
от $\eta$.} и в качестве множителей --- все целые алгебраические функции от
$\eta$.} Мы покажем, что при изоморфном соответствии между двумя областями
$\mathfrak G_\sigma$ и $\mathfrak G_\tau$ $(\sigma\ne\tau)$ модули
$\mathfrak M_\sigma$ и $\mathfrak M_\tau$ необходимо также изоморфно
соответствуют друг другу; отсюда уже легко следует невозможность.

Неопределённые в $\mathfrak G_\tau$ обозначим через $\xi$; при изоморфном
соответствии этим $\xi$ сопоставляются величины $f(\eta)$ из
$\mathfrak G_\sigma$. Поскольку изоморфное соответствие сохраняется и при
образовании частных, области $\mathfrak G_\xi$ всех целых алгебраических
функций от $\xi$ соответствует алгебраически цело замкнутая область
$\mathfrak T$, состоящая из всех функций, алгебраически целых над $f(\eta)$, а
следовательно также алгебраически целых по $\eta$.

Пусть (1) соответствует величине из $\mathfrak M_\tau$. Тогда по модульному
свойству $\mathfrak M_\tau$ произведение (1) с произвольной величиной
$\psi(\eta)$ из $\mathfrak T$ также должно принадлежать $\mathfrak G_\sigma$;
в формулах:
\begin{equation}
(a_0+a_1(\eta)+\cdots+a_{\sigma-1}(\eta))\psi(\eta)
\equiv b_0+b_1(\eta)+\cdots+b_{\sigma-1}(\eta)
\pmod{\mathfrak M_\sigma}.
\tag{2}
\end{equation}
Так как при обращении всех $\eta$ в нуль получается $a_0\psi(0)=b_0$, формула
(2) переписывается в виде
\begin{equation}
(a_0+a_1(\eta)+\cdots+a_{\sigma-1}(\eta))(\psi(\eta)-\psi(0))
\equiv c_1(\eta)+\cdots+c_{\sigma-1}(\eta)
\pmod{\mathfrak M_\sigma}.
\tag{3}
\end{equation}
Так как наряду с $\psi(\eta)$ также функции
\[
 \chi_\nu(\eta)=\sqrt[\nu]{\psi(\eta)-\psi(0)}
\]
при каждом $\nu$ принадлежат $\mathfrak T$, и так как $\chi_\nu(0)=0$, то
аналогично (3) имеют место формулы:
\begin{equation}
(a_0+a_1(\eta)+\cdots+a_{\sigma-1}(\eta))\chi_\nu(\eta)
\equiv c_1^{(\nu)}(\eta)+\cdots+c_{\sigma-1}^{(\nu)}(\eta)
\pmod{\mathfrak M_\sigma}
\quad(\nu=1,2,\ldots).
\tag{4}
\end{equation}
Пусть $A(\eta)$ степени $\lambda$ обозначает произведение
$\chi_1=\psi(\eta)-\psi(0)$ с его $\kappa-1$ сопряжёнными; ввиду
$\chi_1(0)=0$ выражение $A(\eta)$ должно начинаться с членов по крайней мере
первого измерения. Из (4) получается:
\begin{equation}
a_0\chi_\nu(\eta)\equiv0\pmod{\mathfrak M_1};\qquad
a_0^\nu\chi_1(\eta)\equiv0\pmod{\mathfrak M_\nu};\qquad
a_0^{\nu\kappa}A(\eta)\equiv0\pmod{\mathfrak M_{\nu\kappa}},
\tag{5}
\end{equation}
и потому, для $a_0\ne0$, так же как в \S\ 3, $\lambda>\nu\kappa$. Но
$\nu<\frac{\lambda}{\kappa}$ противоречит алгебраически целой замкнутости
$\mathfrak T$, поэтому необходимо $a_0=0$. Из (4) теперь получается:
\begin{equation}
a_1(\eta)\chi_\nu(\eta)\equiv c_1^{(\nu)}(\eta)\pmod{\mathfrak M_2};\qquad
[a_1(\eta)]^{\nu\kappa}A(\eta)
\equiv[c_1^{(\nu)}(\eta)]^{\nu\kappa}\pmod{\mathfrak M_{\nu\kappa+1}},
\tag{6}
\end{equation}
и отсюда, поскольку $A(\eta)$ начинается с членов по крайней мере первого
измерения, $c_1^{(\nu)}(\eta)=0$. Значит, в полной аналогии с (5), для каждого
$\nu$ имеем:
\[
 a_1(\eta)\chi_\nu(\eta)\equiv0\pmod{\mathfrak M_2};\qquad
 [a_1(\eta)]^{\nu\kappa}A(\eta)\equiv0\pmod{\mathfrak M_{2\nu\kappa}},
\]
и отсюда $a_1(\eta)=0$. Продолжая так, попеременное применение (5) и (6) даёт
\[
 a_0=0,\qquad a_1(\eta)=0,\quad \ldots,\quad a_{\sigma-1}(\eta)=0;
\]
то есть, поскольку аналогичное рассмотрение справедливо и относительно
$\mathfrak G_\tau$: \emph{при всяком изоморфном соответствии между
$\mathfrak G_\sigma$ и $\mathfrak G_\tau$ $(\sigma\ne\tau)$ модули
$\mathfrak M_\sigma$ и $\mathfrak M_\tau$ изоморфно соответствуют друг другу.}

Пусть, например, $\sigma>\tau$: неопределённым $\xi$ соответствовали величины
$f(\eta)$ из $\mathfrak G_\sigma$; и, поскольку все величины из
$\mathfrak G_\tau$ можно представить полиномами от $\xi$ по модулю
$\mathfrak M_\tau$, из соответствия $\mathfrak M_\tau$ и
$\mathfrak M_\sigma$ следует, что все величины из $\mathfrak G_\sigma$
получаются как полиномы от $f(\eta)$ по модулю $\mathfrak M_\sigma$. Так как
при $\sigma>\tau\ge1$ неопределённые $\eta$ заведомо не принадлежат
$\mathfrak M_\sigma$ и, значит, цело и рационально выразимы через $f(\eta)$ по
модулю $\mathfrak M_\sigma$, должно существовать такое $f(\eta)$ с
ненулевыми линейными членами. Итак, пусть соответствует:
\[
 \xi:f(\eta)\equiv a_0+a_1(\eta)\pmod{\mathfrak M_2}\quad[a_1(\eta)\ne0],
\]
\[
 \xi^\tau:[f(\eta)]^\tau\equiv a_0^\tau\pmod{\mathfrak M_1}\quad\text{для }a_0\ne0,
\]
\[
 \hphantom{\xi^\tau:[f(\eta)]^\tau}\equiv [a_1(\eta)]^\tau
 \pmod{\mathfrak M_{\tau+1}}\quad\text{для }a_0=0.
\]
$\xi^\tau$ принадлежит $\mathfrak M_\tau$; $[f(\eta)]^\tau$ начинается с
ненулевых членов нулевого или $\tau$-го измерения и при $\tau<\sigma$ не может
принадлежать $\mathfrak M_\sigma$, вопреки выведенному соответствию
$\mathfrak M_\sigma$ и $\mathfrak M_\tau$. Так как случай $\tau>\sigma$
одновременно покрывается перестановкой $\eta$ и $\xi$, тем самым показано, что
области $\mathfrak G_\sigma$ и $\mathfrak G_\tau$ \emph{неизоморфны, или
существенно отличны}. \emph{Следовательно, области $\mathfrak G_i$
$(i=1,2,\ldots\text{ in inf.})$ дают пример счётно бесконечного числа классов
областей $\mathfrak G$.}

\emph{Замечание.} Хотя любые две области $\mathfrak G_i$ неизоморфны, для двух
областей вполне может случиться, что каждая из них изоморфна части другой.\footnote{Такое
поведение может встречаться и для полей; ср. Steinitz, указ. соч., конец.}
Пусть, например, $\tau=1$, а $\sigma$ произвольно:
\[
\mathfrak G_1:a_0+\mathfrak M_1(\xi);\qquad
\mathfrak G_\sigma:a_0+a_1(\eta)+\cdots+a_{\sigma-1}(\eta)+\mathfrak M_\sigma(\eta).
\]
При сопоставлении $\xi=\eta$ область $\mathfrak G_\sigma$ становится
собственным делителем $\mathfrak G_1$; при взаимно однозначном сопоставлении в
$\mathfrak G_1$ и $\mathfrak G_\sigma$, $\xi=\eta^\sigma$,
$\eta=\sqrt[\sigma]{\xi}$, наоборот, $\mathfrak G_1$ становится собственным
делителем $\mathfrak G_\sigma$; следовательно, каждое $\mathfrak G_i$ также
изоморфно собственному делителю самого себя. Поэтому, в отличие от понятия
пересечения, понятие класса пересечения --- то есть такого класса, что каждая
его область была бы изоморфна собственному делителю каждой области всякого
другого класса, --- не существует, по крайней мере как однозначно
определённое.
% END UNIT BODY
% END PRODUCER UNIT 067
\endgroup


\clearpage
\begingroup
\setcounter{footnote}{0}
% BEGIN PRODUCER UNIT 068
% Source: C:/Users/Floris/Documents/interlanguage/03_projects/noether/02_slavic_working_corpus/translations/paper09/russian/v001/Noether_Paper09_Section10_Russian_v001.tex
% Identity: 8185 B / SHA-256 6D89A7D96E6B4363326ABE6281413BADA4F4C1A9E9AFF924F8A5665333A41048
\setlength{\parindent}{1.25em}
\setlength{\parskip}{0pt}
\emergencystretch=18em
\allowdisplaybreaks
\raggedbottom
\providecommand{\vdotswithin}[1]{\mathrel{\vdots}}
% BEGIN UNIT BODY
\section*{\S\ 10. Целые величины произвольного поля.}

Построения предыдущих параграфов --- если в определяющие свойства III и IV для
областей $\mathfrak G$ не включать алгебраические числа --- для областей
$\mathfrak G$ опирались только на тот факт, что совокупность всех комплексных
чисел образует поле;\footnote{Только в \S\ 7 была использована дополнительная
предпосылка, что ``простое поле'' рациональных чисел содержит бесконечно много
элементов; в случае простого поля лишь с конечным числом элементов этот
параграф, как мы увидим, просто отпадает.} для областей $\mathfrak H$ --- еще
и на то, что это поле алгебраически замкнуто. Поэтому полученные результаты
непосредственно допускают \emph{обобщение на самое общее поле}, которое вслед
за Вебером и Э. Штейницем\footnote{Там же, введение.} абстрактно определяется
как ``система элементов с двумя операциями, сложением и умножением, которые
подчинены ассоциативному и коммутативному законам, связаны дистрибутивным
законом и допускают неограниченные и однозначные обратные операции''.\footnote{Исключить
нужно только деление на нуль.}

Каждое такое поле содержит единичный элемент, а следовательно, и выведенное из
него операциями сложения и умножения и обратными к ним операциями ``простое
поле'', которое либо имеет тип поля рациональных чисел (характеристика $0$),
либо тип системы классов вычетов по простому числу $p$ (характеристика
$p$).\footnote{Штейниц, \S\ 4.} \emph{Целые величины простого поля} тогда
хорошо определены как величины, возникающие из единичного элемента сложением и
вычитанием; для простых полей характеристики $p$ целые величины уже исчерпывают
простое поле, и дробных величин простого поля не существует. Каждое
алгебраически замкнутое поле содержит ``абсолютно алгебраическое поле'',
возникающее из простого поля присоединением всех элементов, алгебраических над
простым полем; поэтому для полей характеристики $0$ оно имеет тип поля всех
алгебраических чисел. \emph{Алгебраически целые величины абсолютно
алгебраического поля} тогда хорошо определены как все величины, алгебраически
целые над целыми величинами простого поля (или, что то же самое, над единицей);
для полей характеристики $p$ поэтому также не существует алгебраически дробных
величин абсолютно алгебраического поля. После этих предварительных замечаний
определение ``целых'' и ``алгебраически целых'' величин можно перенести
соответственно на произвольные поля и на алгебраически замкнутые поля:
\emph{целые величины поля $\mathfrak R$ определяются как область целостности
$\mathfrak G$ величин из $\mathfrak R$, поле частных которой\footnote{То есть
поле, состоящее из всех частных любых двух величин из $\mathfrak G$.}
исчерпывает $\mathfrak R$ и которая содержит единичный элемент, но не содержит
дробной величины простого поля.}

\emph{Алгебраически целые величины алгебраически замкнутого поля $\mathfrak R$
определяются как алгебраически цело замкнутая область целостности
$\mathfrak H$ величин из $\mathfrak R$, поле частных которой исчерпывает
$\mathfrak R$ и которая содержит единичный элемент, но не содержит
алгебраически дробной величины абсолютно алгебраического поля.}

Для полей характеристики нуль буквально сохраняются результаты, полученные в
предыдущих параграфах, если только заменить ``алгебраические числа'' величинами
простого поля или соответственно абсолютно алгебраического поля. Для полей
характеристики $p$ результаты еще существенно упрощаются потому, что простое
поле не содержит дробных величин, а значит, дополнительное условие для
рационального базиса и для присоединяемых систем просто отпадает.\footnote{Ср.
первое примечание к этому параграфу.} Итак, результат таков: \emph{абстрактное
определение допускает в качестве целых величин поля характеристики $p$: каждую
область целостности, выведенную из произвольного рационального базиса, само
поле и каждую область целостности, лежащую между этими областями и полем.
Соответственно, алгебраически целыми величинами алгебраически замкнутого поля
простой характеристики считаются: каждая алгебраически цело замкнутая область
целостности, лежащая между областью, выведенной из алгебраического базиса, и
самим полем, --- включая обе граничные области.} Поэтому для самых общих полей
простой характеристики, как и для соответствующих простых полей, различие между
целым и дробным стирается.

\medskip
\noindent Эрланген, 30 марта 1915 г.

\clearpage
% END UNIT BODY
% END PRODUCER UNIT 068
\endgroup


\clearpage
\begingroup
\setcounter{footnote}{0}
% BEGIN PRODUCER UNIT 069
% Source: C:/Users/Floris/Documents/interlanguage/03_projects/noether/02_slavic_working_corpus/translations/paper10/russian/v001/Noether_Paper10_Intro_Section01_Russian_v001.tex
% Identity: 10610 B / SHA-256 6B578DE7D5EA6D46789DBED25F27716132AE7A9726F8CFE96405D87CE540E4D0
\setlength{\parindent}{1.25em}
\setlength{\parskip}{0pt}
\emergencystretch=18em
\allowdisplaybreaks
\raggedbottom
\providecommand{\Afield}{\mathfrak A}
\providecommand{\Bfield}{\mathfrak B}
\providecommand{\Cfield}{\mathfrak C}
\providecommand{\Sdomain}{\mathfrak S}
\providecommand{\Kfield}{\mathfrak K}
\providecommand{\Rfield}{\mathfrak R}
% BEGIN UNIT BODY
% Unit: Paper 10 introduction plus Part 1. Source segments P10-S0002--P10-S0008.
\begin{center}
{\Large\bfseries 10. Функциональные уравнения изоморфного\\ отображения}\par
\vspace{1.25em}
Math. Ann. 77 (1916), с. 536--545
\end{center}

\vspace{3em}

По Дедекинду\footnote{Ср., например: Dirichlet--Dedekind, \emph{Vorlesungen über Zahlentheorie} (4-е изд.), Supplement XI, \S\ 161. Название «изоморфное отображение» или «изоморфизм», образованное по образцу теории групп, встречается лишь в более поздней литературе; Дедекинд говорит о «перестановках поля».} под \emph{изоморфным отображением поля} \(\Afield\) \emph{на систему} \(\Bfield\) --- которая впоследствии сама также оказывается полем --- понимают \emph{такое однозначное соответствие, при котором каждому элементу из} \(\Afield\) \emph{соответствует один и только один элемент из} \(\Bfield\); \emph{и при котором сумме, разности, произведению и частному любых двух элементов из} \(\Afield\) \emph{снова соответствуют сумма, разность, произведение и частное однозначно соответствующих элементов из} \(\Bfield\).

Пусть такое отображение задается --- по сказанному выше, однозначной --- функцией \(f(z)\). Когда \(z\) пробегает все элементы \(x,y,\ldots\) поля \(\Afield\), функция \(f(z)\) характеризуется функциональными уравнениями:
\[
\begin{array}{ll}
(1)\quad f(x+y)=f(x)+f(y); & (2)\quad f(x-y)=f(x)-f(y);\\[0.45em]
(3)\quad f(x\cdot y)=f(x)\cdot f(y); & (4)\quad f\!\left(\dfrac{x}{y}\right)=\dfrac{f(x)}{f(y)}.
\end{array}
\]
\emph{Далее будут указаны самые общие однозначные решения} \(f(z)\) \emph{этих функциональных уравнений в случае, когда} \(x,y,\ldots\) \emph{пробегают все действительные и комплексные числа; то есть самая общая функция} \(f(z)\), \emph{осуществляющая изоморфное отображение поля всех комплексных чисел.}\footnote{Этот вопрос время от времени ставил передо мной Ландау. Как я позднее заметила, часть решений уже имеется у H. Lebesgue: Sur les transformations ponctuelles, transformant les plans en plans. Atti Acad. Torino, 1906/07; а также неявно у A. Ostrowski: Über einige Fragen der allgemeinen Körpertheorie. Crelles Journal 143 (1913), \S\ 2.} Совершенно аналогично получается решение для произвольных полей.

Самые общие решения функционального уравнения (1) Г. Гамель\footnote{G. Hamel: Eine Basis aller Zahlen und die unstetigen Lösungen der Funktionalgleichung: \(f(x+y)=f(x)+f(y)\), Math. Ann. 60, с. 459 (1905).} построил при помощи \emph{линейного} базиса всех чисел. Данная здесь конструкция решений уравнений (1)--(4) --- то есть функции отображения, относительно которой рациональные и алгебраические операции инвариантны, --- пользуется \emph{рациональным} и \emph{алгебраическим} базисами всех чисел. После того как в п. 1 изоморфное отображение будет обсуждено точнее, в п. 2 будут даны необходимые условия для \(f(z)\), которые в п. 3, будучи признаны достаточными, приведут к самой конструкции.\footnote{Ср. параллельные соображения в моей работе: Die allgemeinsten Bereiche aus ganzen transzendenten Zahlen. Math. Ann. 77, с. 103 (1915), особенно \S\ 8 и \S\ 10.} --- Решения (1)--(4), как известно, за исключением \(f(z)=z\) или \(\overline z\), разрывны; в п. 4 будет показано более сильное утверждение, что они «экстремально разрывны», --- факт, который Г. Гамель доказал для действительных решений (1), но который там в комплексной области не обязан выполняться.\footnote{Гамелев термин «полностью разрывный» в остальной литературе употребляется в другом смысле.}

\medskip
\noindent\textbf{1.}\quad Сначала выведем несколько следствий из функциональных уравнений (1)--(4), причем непосредственно для общих полей \(\Afield\). По (1) из однозначности следует: \(f(0)=0\). Если \(f(y)\) не обращается в нуль, то и исходный элемент \(y\) не может обращаться в нуль; поэтому функциональные уравнения (1)--(4) показывают, что \emph{система образов} \(\Bfield\) \emph{всех значений} \(f(z)\) \emph{снова образует поле.} Обратно, если задать начальное условие \(f(0)=0\), то из (2) следует однозначность функции \(f(z)\): \emph{требование однозначности и начальное условие} \(f(0)=0\) \emph{тем самым равносильны.} Из (4) следует, что \(f(z)\) не может тождественно обращаться в нуль; из (3) далее следует, что \(f(z)\) обращается в нуль только при \(z=0\). Действительно, из \(f(y)=0\) и \(y\ne0\) следовало бы, поскольку \(z/y\) также принадлежит \(\Afield\), что
\[
f(z)=f\!\left(\frac{z}{y}\cdot y\right)=f\!\left(\frac{z}{y}\right)\cdot f(y)=0,
\]
то есть, вопреки (4), функция \(f(z)\) тождественно обращалась бы в нуль. Следовательно, из \(f(x)=f(y)\) получается \(x=y\); то есть каждому значению \(f(z)\) соответствует также одно и только одно значение \(z\): \emph{\(f(z)\) есть взаимно однозначная функция от} \(z\). Как показывают функциональные уравнения, отображение, задаваемое обратной функцией, снова изоморфно. Так как каждая рациональная операция составляется из конечного числа сложений, умножений и обратных им операций, можно также сказать: \emph{изоморфное отображение устанавливает взаимно однозначную связь между элементами из} \(\Afield\) \emph{и из} \(\Bfield\), \emph{так что все рациональные соотношения, существующие между любыми конечными наборами элементов из} \(\Afield\):
\[
F(x,y,\ldots,t)=0
\]
\emph{сохраняются в} \(\Bfield\) \emph{и обратно.}

Следует еще заметить, что для поля функция \(f(z)\) уже характеризуется функциональными уравнениями (1) и (3), требованием, чтобы \(f(z)\) не обращалась тождественно в нуль, и начальным условием \(f(0)=0\). Действительно, (2) следует из (1), если заменить \(x\) принадлежащим \(\Afield\) элементом \((x-y)\); затем (2) вместе с \(f(0)=0\) дает однозначность \(f(z)\). Как было показано выше, из (3) далее следует, если заменить \(x\) принадлежащим \(\Afield\) элементом \(x/y\), что \(f(y)\) при \(y\ne0\) не может обращаться в нуль; значит, получаем справедливость (4) и одновременно взаимную однозначность. Если же вместо поля взять за основу область целостности \(\Sdomain\), то \(x/y\) не обязан принадлежать \(\Sdomain\), и из (3) нельзя заключить взаимную однозначность. Здесь, однако, достаточно следующей, наиболее привычной формулировки: изоморфное отображение есть \emph{взаимно однозначное} соответствие между элементами из \(\Sdomain\) и \(\Sdomain'\), такое, что сумме и произведению всегда соответствуют сумма и произведение.\footnote{Ср. к п. 1: Дедекинд, указ. место, \S\ 161.}

\clearpage
% END UNIT BODY
% END PRODUCER UNIT 069
\endgroup


\clearpage
\begingroup
\setcounter{footnote}{0}
% BEGIN PRODUCER UNIT 070
% Source: C:/Users/Floris/Documents/interlanguage/03_projects/noether/02_slavic_working_corpus/translations/paper10/russian/v001/Noether_Paper10_Section02_Russian_v001.tex
% Identity: 7857 B / SHA-256 9FC2B76C336E3D205F2F89DE3724C6DA546F3CB5BE332005C7EBFB23C908C952
\setlength{\parindent}{1.25em}
\setlength{\parskip}{0pt}
\emergencystretch=18em
\allowdisplaybreaks
\raggedbottom
\providecommand{\Afield}{\mathfrak A}
\providecommand{\Bfield}{\mathfrak B}
\providecommand{\Cfield}{\mathfrak C}
\providecommand{\Sdomain}{\mathfrak S}
\providecommand{\Kfield}{\mathfrak K}
\providecommand{\Rfield}{\mathfrak R}
% BEGIN UNIT BODY
% Unit: Paper 10 Part 2. Source segments P10-S0009--P10-S0011.

\medskip
\noindent\textbf{2.}\quad С этого момента, для простоты, вместо \(\Afield\) берём за основу поле \(\Cfield\) всех действительных и комплексных чисел. Прежде всего нужно обсудить системы значений, которые принимает \(f(z)\). Из~(3) или (4) следует: \(f(1)=1\), а отсюда, в силу функциональных уравнений, для каждого \emph{рационального числа} \(\alpha\): \(f(\alpha)=\alpha\); \emph{следовательно, каждое рациональное число при отображении соответствует самому себе}. Так как \emph{алгебраическое число} \(\beta\) удовлетворяет неприводимому уравнению с рациональными числовыми коэффициентами \(F(\beta,\alpha)=0\), то по п.~1 и только что сказанному для \(f(\beta)\) получаем: \(F(f(\beta),\alpha)=0\); \emph{следовательно, каждое алгебраическое число переходит в себя или в одно из своих сопряжённых}, а совокупности алгебраических чисел, ввиду взаимной однозначности, снова соответствует поле всех алгебраических чисел. Чтобы понять поведение \emph{трансцендентных} чисел и одновременно отделить сопряжённые значения, положим в основу \emph{рациональный базис всех чисел}\footnote{Ср. мою цитированную работу о «целых трансцендентных числах», \S\ 6.}, то есть систему \(\Theta\) чисел \(\vartheta\), которая позволяет выразить все числа рационально --- с рациональными числовыми коэффициентами\footnote{Под «рациональной функцией» здесь всегда следует понимать такую функцию, коэффициенты которой также являются рациональными числами.} --- тогда как по крайней мере при одном вполне упорядочении ни одно базисное число не выражается рационально через конечное число предшествующих. Существование такого базиса следует из теоремы о вполне упорядочении совершенно аналогично существованию линейных и алгебраических базисов.

Пусть \(\Cfield\) подчинено вполне упорядочению \(\Omega\). Тогда каждое число либо рационально выражается через конечное число предшествующих, либо рационально независимо от предшествующих. Именно эти последние числа образуют базис \(\Theta\), а индукция показывает, что действительно каждое число рационально выражается через конечное число \(\vartheta\).

Для каждого базисного числа \(\vartheta\) определено «поле начального отрезка \(\Kfield(\vartheta)\)»: оно состоит из всех рациональных комбинаций тех базисных чисел, которые предшествуют \(\vartheta\) во вполне упорядочении. Полем начального отрезка первого базисного числа следует считать поле \(\Rfield\) всех рациональных чисел. Число \(\vartheta\) может вести себя относительно \(\Kfield(\vartheta)\) двояко: оно может быть алгебраически зависимым от \(\Kfield(\vartheta)\) или алгебраически независимым от него. Так как по п.~1 все соотношения, справедливые в \(\Cfield\), справедливы также в поле образов, и обратно, совокупность значений \(f(\vartheta)\), соответствующих \(\vartheta\), образует базис области образов, который сам вполне упорядочен порядком, унаследованным от \(\Theta\). Следовательно, полю начального отрезка \(\Kfield(\vartheta)\) в частности соответствует поле начального отрезка \(\Kfield(f(\vartheta))\); и в зависимости от того, является ли \(\vartheta\) алгебраически зависимым или независимым от \(\Kfield(\vartheta)\), то же самое справедливо для \(f(\vartheta)\) относительно \(\Kfield(f(\vartheta))\); а в случае зависимости соответствуют друг другу также неприводимые уравнения для \(\vartheta\) и \(f(\vartheta)\). Совокупность тех значений \(\vartheta\), которые каждый раз алгебраически независимы от \(\Kfield(\vartheta)\), образует --- как снова показывает индукция --- \emph{алгебраический базис всех чисел}\footnote{Ср. E. Zermelo: \emph{Über ganze transzendente Zahlen}, \S\ 1, Math. Ann. 75 (1914).}, то есть систему \(H\) чисел \(\eta\), через которые все числа можно выразить алгебраически, но между которыми самими нет никаких алгебраических соотношений. Этому алгебраическому базису \(H\) в области образов снова соответствует алгебраический базис \(Z\), который ввиду взаимной однозначности имеет ту же мощность, что и \(H\), а именно мощность континуума, поскольку алгебраическое расширение, как известно, мощности не увеличивает.\footnote{Ср. E. Steinitz: \emph{Algebraische Theorie der Körper}. Crelles Journal 137 (1910), \S\ 24.} Но из знания значений \(f(\vartheta)\) непосредственно получается \(f(z)\) для всех систем значений, так как из \(z=\psi(\vartheta)\) всегда следует: \(f(z)=\psi(f(\vartheta))\).
% END UNIT BODY
% END PRODUCER UNIT 070
\endgroup


\clearpage
\begingroup
\setcounter{footnote}{0}
% BEGIN PRODUCER UNIT 071
% Source: C:/Users/Floris/Documents/interlanguage/03_projects/noether/02_slavic_working_corpus/translations/paper10/russian/v001/Noether_Paper10_Section03_Russian_v001.tex
% Identity: 13173 B / SHA-256 9D580E07BBEF458EF3460BC760043189D37F4DAC3AEE73C55C43898B51E5B494
\setlength{\parindent}{1.25em}
\setlength{\parskip}{0pt}
\emergencystretch=24em
\allowdisplaybreaks
\raggedbottom
\providecommand{\Afield}{\mathfrak A}
\providecommand{\Bfield}{\mathfrak B}
\providecommand{\Cfield}{\mathfrak C}
\providecommand{\Sdomain}{\mathfrak S}
\providecommand{\Kfield}{\mathfrak K}
\providecommand{\Rfield}{\mathfrak R}
% BEGIN UNIT BODY
% Unit: Paper 10 Part 3. Source segments P10-S0012--P10-S0019.

\medskip
\noindent\textbf{3.}\quad Теперь покажем, что найденные необходимые условия действительно дают и построение \(f(z)\). Базис \(Z\), взаимно однозначно соответствующий алгебраическому базису \(H\), построим и сопоставим с \(H\) следующим образом. Положим в основу второе вполне упорядочение \(\Omega'\), которое даёт алгебраический базис \(Z'\) всех чисел; пусть \(Z\), с элементами \(\xi\), будет подмножеством \(Z'\) той же мощности, что и \(Z'\), а следовательно, и \(H\). Каждому элементу \(\eta_i\) из \(H\) теперь взаимно однозначно сопоставим один элемент \(\xi_i\) с тем же номером по правилу: \(\xi_i\) есть \emph{первый} во вполне упорядочении \(\Omega'\) из всех тех элементов \(Z\), которые ещё не сопоставлены никакому элементу из \(H\), предшествующему \(\eta_i\).

Теперь \(f(z)\) \emph{определим так:}
\begin{enumerate}
\item[(a)] \emph{для всех рациональных чисел} \(\alpha\) \emph{по правилу} \(f(0)=0\), \(f(\alpha)=\alpha\);
\item[(b)] \emph{для всех величин алгебраического базиса} \(H\) \emph{по правилу} \(f(\eta_i)=\xi_i\);
\item[(c)] \emph{для всех прочих величин} \(\vartheta\) \emph{рационального базиса} \(\Theta\) --- \emph{тех, которые алгебраически зависят от поля начального отрезка} \(\Kfield(\vartheta)\) \emph{и потому удовлетворяют неприводимому в} \(\Kfield(\vartheta)\) \emph{уравнению} \(F(\vartheta;\vartheta_{i_1},\ldots,\vartheta_{i_\nu})=0\), --- \emph{как корень уравнения} \(F(f(\vartheta);f(\vartheta_{i_1}),\ldots,f(\vartheta_{i_\nu}))=0\);
\item[(d)] \emph{для всех прочих значений} \(z\), \emph{которые по определению} \(\Theta\) \emph{представимы в виде} \(z=\psi(\vartheta_{k_1},\ldots,\vartheta_{k_\nu})\), \emph{по правилу} \(f(z)=\psi(f(\vartheta_{k_1}),\ldots,f(\vartheta_{k_\nu}))\).\footnote{В (d) пункт (a) содержится как частный случай.}
\end{enumerate}
Чтобы показать, что так определённая \(f(z)\) действительно удовлетворяет функциональным уравнениям (1)--(4), по п.~1 достаточно доказать это для (1) и (3), так как \(\Cfield\) является полем, а \(f(z)\) вследствие (a) и (b) не может тождественно обращаться в нуль и, кроме того, удовлетворяет начальному условию \(f(0)=0\). С помощью рационального базиса \(\Theta\) пусть \(x,y,x+y\) выражаются так:
\[
x=\psi_1(\vartheta_1\cdots\vartheta_t),\quad
 y=\psi_2(\vartheta_1\cdots\vartheta_t),\quad
 (x+y)=\psi_3(\vartheta_1\cdots\vartheta_t).
\]
Тогда доказательство того, что \(f(z)\) удовлетворяет функциональному уравнению (1),
\[
f(x)+f(y)-f(x+y)=0,
\]
по (d) равнозначно доказательству, что из
\[
\psi_1(\vartheta)+\psi_2(\vartheta)-\psi_3(\vartheta)
 =\frac{H_1(\vartheta)}{H_2(\vartheta)}=0
\]
всегда следует
\[
\psi_1(f(\vartheta))+\psi_2(f(\vartheta))-\psi_3(f(\vartheta))
 =\frac{H_1(f(\vartheta))}{H_2(f(\vartheta))}=0,
\]
где \(H_1,H_2\) обозначают целые рациональные функции своих аргументов. К точно такой же форме условия приводит и функциональное уравнение (3); значит, выполнение всех функциональных уравнений равнозначно выполнению следующего --- по пп.~1 и 2 также необходимого --- условия: для \emph{каждой} целой рациональной функции \(H(\vartheta)\) из \(H(\vartheta)=0\) всегда следует также \(H(f(\vartheta))=0\), а из \(H(\vartheta)\ne0\) также \(H(f(\vartheta))\ne0\); последнее требуется из-за появления знаменателя.

Что это действительно так, показывает индукция. Пусть, например, в \(H(\vartheta_1\cdots\vartheta_t)\) базисная величина \(\vartheta_t\) является последней во вполне упорядочении среди \(\vartheta_1\cdots\vartheta_t\);\footnote{Выражения \(H(\vartheta)\), имеющие нулевую степень по \(\vartheta\), при отображении переходят в себя, поэтому их не нужно рассматривать дальше.} тогда \(H(\vartheta_1\cdots\vartheta_t)=0\) можно рассматривать как соотношение, которому \(\vartheta_t\) удовлетворяет относительно поля начального отрезка \(\Kfield(\vartheta_t)\), и так же \(H(f(\vartheta_1)\cdots f(\vartheta_t))=0\) как соотношение для \(f(\vartheta_t)\) относительно \(\Kfield(f(\vartheta_t))\). Если бы не все соотношения \(H(\vartheta)=0\) выполнялись также для \(f(\vartheta)\), и обратно, то должна была бы существовать первая базисная величина --- скажем, \(\vartheta_0\), --- для которой по крайней мере одно соотношение, справедливое относительно \(\Kfield(\vartheta_0)\), не сохранялось бы в поле образов, или наоборот, тогда как предыдущие поля начальных отрезков \(\Kfield(\vartheta_0)\) и \(\Kfield(f(\vartheta_0))\) были бы изоморфны. В частности, по (a) поле начального отрезка первой базисной величины из \(\Theta\), то есть поле \(\Rfield\) всех рациональных чисел, изоморфно своему полю образов. Пусть теперь \(H(\vartheta_0)\) имеет вид
\[
H(\vartheta_0)=a_0(\vartheta)\cdot\vartheta_0^\chi
+a_1(\vartheta)\cdot\vartheta_0^{\chi-1}+\cdots+a_\chi(\vartheta),
\]
где \(a_i(\vartheta)\) являются величинами из \(\Kfield(\vartheta_0)\). Тогда есть две возможности:

\noindent 1) \(\vartheta_0\) алгебраически независима от \(\Kfield(\vartheta_0)\). Тогда \(H(\vartheta_0)=0\) выполняется тождественно по \(\vartheta_0\),\footnote{В силу выразимости \(x,y,\ldots\) через \(\vartheta\) такая форма соотношения вполне может возникать.} то есть \(a_i(\vartheta)=0\) для всех коэффициентов, а значит, вследствие изоморфизма \(\Kfield(\vartheta_0)\) и \(\Kfield(f(\vartheta_0))\), также \(a_i(f(\vartheta))=0\). Поэтому из \(H(\vartheta_0)=0\) всегда следует \(H(f(\vartheta_0))=0\). Обратно, пусть \(H(f(\vartheta_0))=0\). Так как \(\vartheta_0\), будучи алгебраически независимой от \(\Kfield(\vartheta_0)\), принадлежит алгебраическому базису \(H\), то \(f(\vartheta_0)\) по (b) принадлежит алгебраическому базису \(Z\) и алгебраически независима от \(\Kfield(f(\vartheta_0))\). Следовательно, обращаются в нуль все \(a_i(f(\vartheta))\), а вследствие изоморфизма --- все \(a_i(\vartheta)\); из \(H(f(\vartheta_0))=0\) следует \(H(\vartheta_0)=0\), или из \(H(\vartheta_0)\ne0\) также \(H(f(\vartheta_0))\ne0\).

\noindent 2) \(\vartheta_0\) алгебраически зависима от \(\Kfield(\vartheta_0)\). Тогда \(H(\vartheta_0)\) делится на неприводимое относительно \(\Kfield(\vartheta_0)\) уравнение \(F(\vartheta_0)\), то есть тождественно по \(t\) имеем
\[
H(t;a_i(\vartheta))=F(t;\vartheta)\cdot G(t;\vartheta_i).
\]
А поскольку все входящие коэффициенты принадлежат \(\Kfield(\vartheta_0)\), в соответствующем изоморфном поле \(\Kfield(f(\vartheta_0))\) выполняется соответствующее соотношение:
\[
H(t;a_i(f(\vartheta)))=F(t;f(\vartheta))\cdot G(t;f(\vartheta)).
\]
Но по (c) \(f(\vartheta_0)\) было выбрано так, что оно является корнем \(F(t;f(\vartheta))\); поэтому из \(H(\vartheta_0)=0\) следует \(H(f(\vartheta_0))=0\).

Обратно, если \(H(f(\vartheta_0))=0\), то так же из делимости \(H(t;a_i(f(\vartheta)))\) на неприводимую относительно \(\Kfield(f(\vartheta_0))\) функцию \(F(t;f(\vartheta))\) для изоморфного поля \(\Kfield(\vartheta_0)\) следует делимость \(H(t;a_i(\vartheta))\) на \(F(t;\vartheta)\). А так как \(F(t;f(\vartheta))\) возникла из неприводимого уравнения для \(\vartheta_0\) посредством изоморфного соответствия, и это соответствие взаимно однозначно, то \(F(t;\vartheta)\) необходимо является той неприводимой функцией, корнем которой служит \(\vartheta_0\). Следовательно, из \(H(f(\vartheta_0))=0\) следует \(H(\vartheta_0)=0\), или из \(H(\vartheta_0)\ne0\) также \(H(f(\vartheta_0))\ne0\).

Итак, из изоморфных полей начальных отрезков \(\Kfield(\vartheta_0)\) и \(\Kfield(f(\vartheta_0))\) посредством присоединения \(\vartheta_0\), соответственно \(f(\vartheta_0)\), снова получаются изоморфные поля; предположение, что для некоторого \(\vartheta_0\) это не так, приводит к противоречию. Тем самым доказано, что \emph{функция} \(f(z)\), \emph{определённая по (a)--(d), действительно представляет решение функциональных уравнений (1)--(4), и притом, по п.~2, самое общее}.

Все рассуждения, использованные для построения \(f(z)\), опирались только на тот факт, что совокупность \(\Cfield\) всех чисел образует поле, и потому справедливы для каждого абстрактно определённого поля \(\Afield\). При этом надо заметить, что полю \(\Rfield\) рациональных чисел соответствует «простое поле» поля \(\Afield\), то есть поле, полученное из единичного элемента \(\Afield\) операциями сложения и умножения и обратными к ним операциями. Поэтому, если в (a)--(d) заменить рациональные числа элементами простого поля, то \emph{получающаяся таким образом функция} \(f(z)\) \emph{осуществляет самое общее изоморфное отображение произвольного абстрактно определённого поля}.
% END UNIT BODY
% END PRODUCER UNIT 071
\endgroup


\clearpage
\begingroup
\setcounter{footnote}{0}
% BEGIN PRODUCER UNIT 072
% Source: C:/Users/Floris/Documents/interlanguage/03_projects/noether/02_slavic_working_corpus/translations/paper10/russian/v001/Noether_Paper10_Section04_Russian_v001.tex
% Identity: 14125 B / SHA-256 A39500F9C993E697160CF887E6C40B329EA03F37D68461FA8570EA40641F2AED
\setlength{\parindent}{1.25em}
\setlength{\parskip}{0pt}
\emergencystretch=24em
\allowdisplaybreaks
\raggedbottom
\providecommand{\Afield}{\mathfrak A}
\providecommand{\Bfield}{\mathfrak B}
\providecommand{\Cfield}{\mathfrak C}
\providecommand{\Sdomain}{\mathfrak S}
\providecommand{\Kfield}{\mathfrak K}
\providecommand{\Rfield}{\mathfrak R}
\providecommand{\Xreal}{\mathfrak X}
\providecommand{\Yreal}{\mathfrak Y}
% BEGIN UNIT BODY
% Unit: Paper 10 Part 4. Source segments P10-S0020--P10-S0028.

\medskip
\noindent\textbf{4.}\quad Теперь остается показать, что функции \(f(z)\), определенные для поля \(\Cfield\) всех действительных и комплексных чисел, --- которые, как известно, за исключением \(f(z)=z\) или \(\overline z\), разрывны, --- являются \emph{экстремально разрывными}. Иными словами, в окрестности любого действительного или комплексного числа \(z_0\) существуют значения \(z\), для которых \(f(z)\) сколь угодно близко подходит к любому наперед заданному значению \(Z_0\). Эта экстремальная разрывность, как доказал Г. Гамель в указанной работе, присуща \emph{действительным} решениям \(\varphi(x)\), определенным для всех \emph{действительных} чисел, функционального уравнения (1). Но, как показывают приведенные ниже примеры, она не обязана иметь место для решений, определенных на комплексных числах. Рассуждение Гамеля таково. Если \(\varphi(x)\) отличается от \(C\cdot x\), то существуют по крайней мере два значения линейного базиса\footnote{Линейный базис всех (действительных) чисел задается системой (действительных) чисел, позволяющей линейно, с рациональными коэффициентами, выразить все (действительные) числа, тогда как между любым конечным набором базисных чисел нет линейного соотношения с рациональными коэффициентами.}, скажем \(x_1\) и \(x_2\), такие, что \(\varphi(x_1):x_1\ne\varphi(x_2):x_2\). Тогда две системы равенств
\[
\alpha x_1+\beta x_2=a;\qquad \alpha\varphi(x_1)+\beta\varphi(x_2)=b
\]
имеют решение в действительных числах \(\alpha,\beta\); следовательно, \(a\) и \(b\) можно сколь угодно хорошо аппроксимировать рациональными числами \(\alpha,\beta\). А поскольку \(\alpha,\beta\) рациональны, то \(\alpha\varphi(x_1)+\beta\varphi(x_2)\) действительно равно \(\varphi(\alpha x_1+\beta x_2)\). Для комплексных систем значений этот вывод уже не работает. Действительно, и в комплексной области можно задать решения (1), разрывные, но не экстремально разрывные, например:
\[
\varphi(z)=\varphi(x+iy)=\varphi(x)+i(y+\varphi(x)),
\]
где \(\varphi(x)\) означает действительное решение, экстремально разрывное на действительной прямой; или еще проще \(\varphi(z)=\varphi(x)\). В обоих случаях действительная часть \(\varphi(z)\) экстремально разрывна, тогда как мнимая часть определяется действительной.

Это поведение, а вместе с ним и поведение функции \(f(z)\), становится совершенно ясным, если ввести понятие ранга. Положим \(z=x+iy\), \(\varphi(z)=X+iY\), и будем называть \(\varphi(z)\) соответственно функцией ранга четыре, три, два или один, смотря по тому, существует ли ноль, одно, два или три линейно независимых соотношения
\[
c_1x+c_2y+c_3X+c_4Y=0,
\]
где \(c\), конечно, действительные числа.

\noindent 1) Пусть \(\varphi(z)\) имеет \emph{ранг четыре}. Тогда существуют по крайней мере четыре значения линейного базиса, скажем \(z_1,z_2,z_3,z_4\), такие, что определитель
\[
\begin{vmatrix}
x_1&x_2&x_3&x_4\\
y_1&y_2&y_3&y_4\\
X_1&X_2&X_3&X_4\\
Y_1&Y_2&Y_3&Y_4
\end{vmatrix}
\]
не обращается в нуль. Поэтому уравнения
\[
\begin{aligned}
\alpha x_1+\beta x_2+\gamma x_3+\delta x_4&=a,\\
\alpha y_1+\cdots+\delta y_4&=b,\\
\alpha X_1+\cdots+\delta X_4&=A,\\
\alpha Y_1+\cdots+\delta Y_4&=B
\end{aligned}
\]
имеют решение в действительных числах \(\alpha,\ldots,\delta\), а \(a,b,A,B\) можно сколь угодно хорошо аппроксимировать рациональными числами \(\alpha,\beta,\gamma,\delta\). Кроме того, имеем:
\[
\alpha(X_1+iY_1)+\beta(X_2+iY_2)+\gamma(X_3+iY_3)+\delta(X_4+iY_4)
 =\varphi(\alpha z_1+\beta z_2+\gamma z_3+\delta z_4).
\]
Следовательно, «в общем случае» решения \(\varphi(z)\) и в комплексной области экстремально разрывны; \emph{значения разрывности} \(\varphi(z)\) в окрестности \(z_0=a+ib\) \emph{образуют двумерное линейное многообразие}.

\noindent 2) \(\varphi(z)\) имеет \emph{ранг три}. Тогда существуют по крайней мере три системы значений линейного базиса, скажем \(z_1,z_2,z_3\), такие, что приведенный выше определитель имеет ранг три. Значения \(a,b,A,B\), удовлетворяющие соотношению
\[
c_1a+c_2b+c_3A+c_4B=0,
\]
и только они, могут быть сколь угодно хорошо аппроксимированы; \emph{значения разрывности} образуют \emph{одномерное линейное многообразие}, определенное числом \(z_0=a+ib\). Как показывают приведенные примеры, этот случай действительно может возникать.

\noindent 3) \(\varphi(z)\) имеет \emph{ранг два}. Так как всегда можно выбрать два комплексных числа \(z_1,z_2\) так, что
\[
\begin{vmatrix}x_1&x_2\\ y_1&y_2\end{vmatrix}\ne0,
\]
соотношения имеют вид
\[
X=\lambda_1x+\lambda_2y;\qquad Y=\mu_1x+\mu_2y.
\]
Получающееся отсюда выражение
\[
\varphi(z)=(\lambda_1x+\lambda_2y)+i(\mu_1x+\mu_2y)
\]
и есть самое общее \emph{непрерывное} решение функционального уравнения (1) в области комплексных чисел.

\noindent 4) \(\varphi(z)\) имеет \emph{ранг один}. В области комплексных чисел этот случай невозможен; в области \emph{действительных} чисел он означает \emph{непрерывное} решение \(\varphi(x)=C\cdot x\).

Теперь покажем, что \emph{решения} \(f(z)\) \emph{функциональных уравнений (1)--(4), определенные для всех действительных и комплексных чисел, могут иметь только ранг два или ранг четыре}. При этом ранг два дает только две \emph{непрерывные} функции \(f(z)=z\) или \(\overline z\), как показывает специализация \(f(1)=1\) и \(f(i)=\pm i\) вместе с п.~3. Так как ранг один уже исключен выше, остается исключить ранг три. Итак, предположим, что существует хотя бы одно соотношение
\[
c_1x+c_2y+c_3X+c_4Y=0,
\]
так что \(f(z)\) имеет ранг три, возможно более низкий, и покажем, что при этом предположении всегда наступает именно последний случай. Как снова показывает специализация \(f(1)=1\), \(f(i)=\pm i\), соотношение переходит в
\[
c_3(X-x)+c_4(Y\mp y)=0,
\]
где \(c_3\) и \(c_4\) не обращаются в нуль одновременно. Сначала пусть, например, \(c_4=0\). Соотношение \((X-x)=0\) означает, что чисто мнимому \(z\) соответствует также чисто мнимое \(f(z)\). Поэтому для каждого действительного \(y\) имеем \(f(iy)=i\cdot\Yreal\), где \(\Yreal\) снова действительно. Но из-за
\[
f(iy)=\pm i f(y)\quad\text{также}\quad f(y)=\pm\Yreal;
\]
каждому действительному \(z\) соответствует также действительное \(f(z)\); а из-за \((X-x)\) все действительные значения переходят в себя. Тем самым остаются только \(f(z)=z\) или \(\overline z\), то есть фактически ранг два. Совершенно аналогично рассуждаем, когда \(c_3=0\), \(c_4\ne0\).

Теперь пусть \(c_3\) и \(c_4\) отличны от нуля, так что соотношение можно записать в виде
\[
(Y\mp y)=c\cdot(X-x)
\]
с действительным, отличным от нуля и конечным \(c\). Это означает, что при \(y=\pm cx\) также \(Y=c\cdot X\). Поэтому, учитывая функциональные уравнения, для каждого действительного \(x\) выполняется
\[
f\{x\cdot(1\pm ic)\}=\Xreal(1+ic)=f(x)\cdot(1+i f(c)),
\]
где \(\Xreal\) снова действительно. Но здесь по п.~1 множитель при \(f(x)\) не может тождественно обращаться в нуль, и деление дает
\[
f(x)=\Xreal(\sigma+i\tau)
\]
с действительными \(\sigma\) и \(\tau\), не зависящими от \(x\) и \(\Xreal\). В частности, значению \(x=1\) также соответствует некоторое действительное \(X_0\), а условие \(1=X_0(\sigma+i\tau)\) показывает, что необходимо \(\tau=0\), значит \(f(x)=\sigma\cdot\Xreal\). Следовательно, каждому действительному \(z\) соответствует также действительное \(f(z)\); иначе говоря, из \(y=0\) необходимо следует \(Y=0\). Тогда линейное соотношение для действительных \(z\) переходит в \(X-x=0\);\footnote{Здесь используется факт, что \(c\ne0\), то есть что ни \(c_3\), ни \(c_4\) не обращается в нуль, тогда как выше использовалось только \(c_4\ne0\); поэтому случай, когда один из этих двух коэффициентов обращается в нуль, нужно было рассмотреть заранее.} каждое действительное значение соответствует самому себе. Снова остаются только \(f(z)=z\) или \(\overline z\), то есть фактически ранг два. Тем самым ранг три исключен во всех случаях. Но, как показывают исследования первых пунктов, разрывные решения действительно существуют и потому необходимо имеют ранг четыре; следовательно, по п.~1 получаем теорему: \emph{«Все решения \(f(z)\) функциональных уравнений (1)--(4), за исключением непрерывных решений \(f(z)=z\) или \(\overline z\), экстремально разрывны в действительной и комплексной областях».}\footnote{Lebesgue в указанной работе показывает, что действительная и мнимая части \(f(z)\) обе являются «неизмеримыми» функциями от \(x\) и \(y\); как показывает пример в начале п.~4, \(\varphi(z)=\varphi(x)+i(y+\varphi(x))\), отсюда еще не следует экстремальная разрывность в комплексной области.}

\medskip
\noindent Эрланген, 30 октября 1915 г.
% END UNIT BODY
% END PRODUCER UNIT 072
\endgroup


\clearpage
\begingroup
\setcounter{footnote}{0}
% BEGIN PRODUCER UNIT 073
% Source: C:/Users/Floris/Documents/interlanguage/03_projects/noether/02_slavic_working_corpus/translations/paper11/russian/v001/Noether_Paper11_Russian_v001.tex
% Identity: 37926 B / SHA-256 5AE5D0CA161D123BE968242E909EF80287CF1406D0037DD4217FE704B8AD7F2F
\setlength{\parindent}{1.25em}
\setlength{\parskip}{0pt}
\emergencystretch=30em
\allowdisplaybreaks
\raggedbottom
\providecommand{\Rfield}{\mathfrak R}
% BEGIN UNIT BODY
% Unit: Paper 11. Source segments P11-S0001--P11-S0027.

\begin{center}
{\Large\bfseries 11. Уравнения с заданной группой}\par
\vspace{1.5em}
Math. Ann. 78 (1918), с. 221--229
\end{center}

\vspace{4em}

Проблему построения уравнений с заданной группой можно атаковать в двух направлениях, которые коротко можно назвать «иррациональным», то есть характеризующим корни, и «рациональным», то есть характеризующим коэффициенты.

К «иррациональному» направлению, работающему функционально-теоретически и арифметически, относится теорема Кронекера о том, что все абелевы поля в области рациональных чисел являются циклотомическими полями, а также соответствующие теоремы для относительно абелевых полей относительно квадратичных числовых полей. Эти теоремы, с одной стороны, дают реализуемость всех абелевых групп относительно этих специальных областей рациональности; с другой стороны, они дают всю совокупность уравнений и вместе с тем более глубокое понимание структуры определяемых ими числовых полей. Однако каждая отдельная теорема ограничена специальной областью рациональности, и каждая новая область рациональности требует совершенно нового рассмотрения.

«Рациональное», алгебраическое направление опирается на теорему Гильберта о неприводимости, согласно которой в представлении коэффициентов можно ограничиться параметрическим представлением. Сюда относится, например, существование сколь угодно многих уравнений с альтернирующей группой относительно любой заданной области рациональности. Здесь, естественно, отсутствует описанное выше понимание структуры полей, задаваемых уравнениями; но преимущество этого «рационального» направления состоит в том, что реализуемость групп доказывается для \emph{всех областей рациональности одновременно}; однако известные до сих пор параметрические представления вообще не дают всей совокупности уравнений.\footnote{Для разрешимых уравнений параметрическое представление коэффициентов можно заменить представлением корней. Недавно Ф. Мертенс для некоторых групп 8-й степени дал наиболее общие представления корней: «Gleichungen 8ten Grades mit Quaternionengruppen», Sitzb. d. Ak. d. Wiss., Wien, Abt. IIa, Bd. 125, S. 735. Как я заключаю из этой заметки, еще раньше Г. Бухт для всегда осуществимых нормальных форм уравнений 3-й и 4-й степеней и уравнений 8-й степени с упомянутыми группами указал наиболее общие выражения корней: «Über einige algebraische Körper achten Grades», Arkiv for Math., Astron. och Physik, Bd. 6, Nr. 30. [Добавление при корректуре.]} 

Дальнейшее является вкладом в это «рациональное» направление: здесь указываются \emph{достаточные условия} для \emph{существования параметрических представлений, которые при любой области рациональности дают всю совокупность уравнений с заданной группой}. Эти условия состоят в том, что инвариантное поле, принадлежащее группе, -- лагранжев родовой домен, -- рационально представляется через \emph{независимые} функции этой области, то есть обладает «минимальной базой»; или, иначе говоря, изоморфно полю всех рациональных функций от $n$ неизвестных. Как будет показано еще в § 2, это, например, выполнено для всех групп, возникающих у уравнений 3-й и 4-й степеней. Фактическое построение и более точное обсуждение этих уравнений 3-й и 4-й степеней содержится в диссертации Ф. Зейдельмана,\footnote{Die Gesamtheit der kubischen und biquadratischen Gleichungen mit Affekt bei beliebigem Rationalitätsbereich. Erlangen 1916.} выдержка из которой непосредственно следует за этим сообщением.

Вопрос о минимальной базе лагранжевых родовых доменов был, как уже упоминалось в работе «Körper und Systeme rationaler Funktionen» (Math. Ann. 76), поставлен передо мной Э. Фишером независимо от описанной здесь связи и дал толчок настоящим исследованиям.\footnote{См. часть 2 предварительного сообщения о «Rationale Funktionenkörper», Jahresb. d. d. Math. Vereinig., Bd. 22, 1913.}

\begin{center}
§ 1.\\[0.75\baselineskip]
\textbf{Достаточные условия для параметрического представления.}
\end{center}

\noindent\textbf{1.}\quad Возможность свести построение уравнений с заданной группой к параметрическому представлению следует из теоремы Гильберта о неприводимости, которая утверждает следующее. Пусть уравнение
\begin{equation}
 f(x)=x^n+F_1(\lambda_1\cdots\lambda_r)x^{n-1}+\cdots+F_n(\lambda_1\cdots\lambda_r)=0
\tag{1}
\end{equation}
-- где $F_i(\lambda_1\cdots\lambda_r)$ обозначают рациональные функции независимых параметров $\lambda_1\cdots\lambda_r$ с коэффициентами из числового поля\footnote{Вместо числового поля могла бы фигурировать и общая область рациональности.} $\Omega$ -- имеет относительно поля $\Omega(\lambda_1\cdots\lambda_r)$ всех рациональных функций от $\lambda_1\cdots\lambda_r$ с коэффициентами из $\Omega$ группу $\Gamma$. Тогда параметры $\lambda$ можно заменить бесконечно многими рациональными числами так, что возникающее числовое уравнение
\begin{equation}
 g(x)=x^n+a_1x^{n-1}+\cdots+a_n=0
\tag{2}
\end{equation}
относительно $\Omega$ имеет именно группу $\Gamma$. Такое $g(x)$ точнее будем обозначать через $g_\Gamma$.

Теперь спрашивается, можно ли \emph{задать такое параметрическое представление} (1), \emph{чтобы возникающее числовое уравнение} (2) \emph{пробегало всю совокупность всех $g_\Gamma$, когда $\lambda_1\cdots\lambda_r$ пробегают все величины из $\Omega$}, -- за исключением тех величин из $\Omega$, которые вызывают редукцию группы.\footnote{Сюда следует отнести также появление кратного корня у $g(x)=0$.} Иными словами: нелинейная система уравнений
\begin{equation}
 F_1(\lambda_1\cdots\lambda_r)=a_1;\quad \cdots;\quad F_n(\lambda_1\cdots\lambda_r)=a_n
\tag{3}
\end{equation}
должна для каждой системы значений $a_1\cdots a_n$, соответствующей некоторому $g_\Gamma$, иметь решение, причем в величинах $\lambda_1\cdots\lambda_r$ из $\Omega$. Так как система (3) нелинейна, в требовании получить всю совокупность $g_\Gamma$ посредством параметрического представления с самого начала приходится сделать ограничение. Именно, в таких нелинейных системах всегда могут возникать сингулярные системы значений $a_1\cdots a_n$, удовлетворяющие алгебраическому соотношению $H(a_1\cdots a_n)=0$, для которых \emph{не} существует решения;\footnote{Это вскоре будет в общем виде выполнено в работе об «альтернативе в нелинейных системах уравнений».} тогда параметрическое представление (1) отказывает и необходимо дополнительное представление. Однако в данном случае эти сингулярные значения $a_1\cdots a_n$, как будет видно, можно просто охарактеризовать.

\noindent\textbf{2.}\quad Чтобы получить такое параметрическое представление, мы предъявляем более сильное требование: $\lambda_i$ должны быть \emph{естественными иррациональностями, принадлежащими группе, тождественно по корням, рассматриваемым как неизвестные}. Под этим мы понимаем следующее. Если в (1) заменить параметры $\lambda_1\cdots\lambda_r$ надлежаще выбранными \emph{рациональными} функциями $\varphi_1(x)\cdots\varphi_r(x)$ от неизвестных $x_1\cdots x_n$ -- с коэффициентами из $\Omega$ --, то (1) должно перейти в
\begin{equation}
 h(x)=x^n-\sigma_1(x)x^{n-1}+\sigma_2(x)x^{n-2}\cdots\pm\sigma_n(x),
\tag{4}
\end{equation}
где $\sigma_1(x)\cdots\sigma_n(x)$ обозначают элементарные симметрические функции от $x_1\cdots x_n$; и $h(x)$ относительно области рациональности $\Omega(\varphi_1(x)\cdots\varphi_r(x))$, возникающей при этом из $\Omega(\lambda_1\cdots\lambda_r)$, должно иметь группу $\Gamma$. В силу независимости параметров функции $\varphi_1(x)\cdots\varphi_r(x)$ также должны быть алгебраически независимыми.

Чтобы уточнить это требование, надо выяснить связь $\Omega(\varphi_1(x)\cdots\varphi_r(x))$ с группой $\Gamma$. Для этого сначала определим самую общую область рациональности $K$, относительно которой $h(x)$ становится $h_\Gamma$. Пусть $\mathfrak L_\Gamma$ обозначает поле всех рациональных функций с рационально-числовыми коэффициентами от $x_1\cdots x_n$, допускающих $\Gamma$, -- его также называют лагранжевым родовым доменом, принадлежащим $\Gamma$, или инвариантным полем $\Gamma$; а $\Omega(\mathfrak L_\Gamma)$ -- соответствующее поле, получающееся присоединением $\mathfrak L_\Gamma$ к $\Omega$. Итак, $\mathfrak L_\Gamma$ является подполем поля $\mathfrak R(x_1\cdots x_n)$ всех рациональных функций с рационально-числовыми коэффициентами от $x_1\cdots x_n$. Тогда по определению группы самая общая область $K$ задается как \emph{любое поле, которое содержит $\mathfrak L_\Gamma$ как подполе и пересечение которого с $\mathfrak R(x_1\cdots x_n)$ равно точно $\mathfrak L_\Gamma$}. Соответственно, для всякой области $K$, содержащей $\Omega$, ее пересечение с $\Omega(x_1\cdots x_n)$ задается полем $\Omega(\mathfrak L_\Gamma)$. Поэтому, в частности, так как по нашему требованию $\Omega(\varphi_1(x)\cdots\varphi_r(x))$ должно быть областью $K$, пересечение $\Omega(\varphi_1(x)\cdots\varphi_r(x))$ с $\Omega(x_1\cdots x_n)$ задается полем $\Omega(\mathfrak L_\Gamma)$; а так как, с другой стороны, $\Omega(\varphi_1(x)\cdots\varphi_r(x))$ является подполем $\Omega(x_1\cdots x_n)$, оно оказывается в точности равным $\Omega(\mathfrak L_\Gamma)$. Следовательно, требование означает, что \emph{лагранжев родовой домен $\Omega(\mathfrak L_\Gamma)$ должен возникать присоединением к $\Omega$ алгебраически независимых функций $\varphi_1(x)\cdots\varphi_r(x)$}; при этом необходимо $r=n$, потому что $\mathfrak L_\Gamma$ содержит $n$ алгебраически независимых элементарных функций, тогда как более чем $n$ функций не могут быть алгебраически независимыми. Функции $\varphi_1(x)\cdots\varphi_n(x)$, как мы еще будем говорить, образуют \emph{минимальную базу} $\Omega(\mathfrak L_\Gamma)$; или же $\mathfrak L_\Gamma$ обладает минимальной базой с коэффициентами из $\Omega$.

\noindent\textbf{3.}\quad Теперь существенно, что все это рассмотрение можно обратить; и так оно действительно приводит к достаточным условиям\footnote{Требование в пункте 2 было расширено по сравнению с пунктом 1, так что условия не обязаны быть необходимыми.} для искомого параметрического представления. Итак, предположим, что $\varphi_1(x)\cdots\varphi_n(x)$ является минимальной базой $\Omega(\mathfrak L_\Gamma)$; и пусть $\Omega$ будет наименьшим полем, содержащим коэффициенты $\varphi_1(x)\cdots\varphi_n(x)$. Тогда $\Omega$ необходимо является алгебраическим числовым полем и, в частности, может быть полем $\mathfrak R$ всех рациональных чисел. Так как $\sigma_1(x),\ldots,\sigma_n(x)$ принадлежат $\Omega(\mathfrak L_\Gamma)$, существуют тождества по $x_1\cdots x_n$ -- где $F_i$ обозначают рациональные функции своих аргументов:
\begin{equation}
\sigma_1(x)=F_1(\varphi_1(x)\cdots\varphi_n(x));\quad \cdots;\quad
\sigma_n(x)=F_n(\varphi_1(x)\cdots\varphi_n(x));
\tag{5}
\end{equation}
тогда как, наоборот, $\varphi_1(x)\cdots\varphi_n(x)$, а следовательно и $\varphi(x)=\mu_1\varphi_1(x)+\cdots+\mu_n\varphi_n(x)$, алгебраически зависят от $\sigma_1(x)\cdots\sigma_n(x)$ посредством уравнения, неприводимого относительно $\Omega(\sigma_1(x)\cdots\sigma_n(x))$:
\begin{equation}
G(\varphi(x);\sigma_1(x)\cdots\sigma_n(x))=0,
\tag{6}
\end{equation}
которое снова является тождеством по $x$.

В силу (5) это тождество (6) переходит в
\begin{equation}
G(\varphi(x);F_1(\varphi_1\cdots\varphi_n);\cdots;F_n(\varphi_1\cdots\varphi_n))
=H(\varphi_1(x)\cdots\varphi_n(x))=0;
\tag{7}
\end{equation}
и вследствие алгебраической независимости $\varphi_1(x)\cdots\varphi_n(x)$ равенство (7) выполнено не только тождественно по $x$, но и тождественно по $\varphi_i$; то есть тождественно в независимых параметрах $\lambda$ имеем
\begin{equation}
H(\lambda_1\cdots\lambda_n)=G(\mu_1\lambda_1+\cdots+\mu_n\lambda_n;F_1(\lambda_1\cdots\lambda_n);\cdots;F_n(\lambda_1\cdots\lambda_n))=0.
\tag{8}
\end{equation}
Из тождества (8) следует, что уравнение
\begin{equation}
f(x)=x^n-F_1(\lambda_1\cdots\lambda_n)x^{n-1}+F_2(\lambda_1\cdots\lambda_n)x^{n-2}\cdots\pm F_n(\lambda_1\cdots\lambda_n)=0
\tag{9}
\end{equation}
\emph{относительно} $\Omega(\lambda_1\cdots\lambda_n)$ \emph{имеет группу} $\Gamma$. Действительно, это тождество показывает, что рационально известное $\lambda=\mu_1\lambda_1+\cdots+\mu_n\lambda_n$ есть функция корней, принадлежащая $\Gamma$; но никакая дальнейшая функция, принадлежащая подгруппе $\Gamma$, не может быть рационально известной, как показывает подстановка $\lambda_i=\varphi_i(x)$; и при той же подстановке $\lambda$ также отличается от всех своих сопряженных. Если, кроме того, $\Omega^*$ есть произвольное числовое поле, содержащее $\Omega$, то $f(x)$ \emph{имеет группу} $\Gamma$ \emph{относительно} $\Omega^*(\lambda_1\cdots\lambda_n)$; ибо по пункту 2 группа не редуцируется при подстановке $\lambda_i=\varphi_i(x)$ даже после присоединения $\Omega^*$.

Остается доказать, что в модифицированном в пункте 1 смысле $f(x)$ действительно пробегает всю совокупность всех $g_\Gamma$. Для этого заметим, что в силу (5) функции $F_i(\lambda)$ являются алгебраически независимыми функциями своих аргументов $\lambda$. Поэтому система уравнений
\begin{equation}
F_1(\lambda_1\cdots\lambda_n)=-a_1;\quad F_2(\lambda_1\cdots\lambda_n)=a_2;\quad \cdots;\quad F_n(\lambda_1\cdots\lambda_n)=\pm a_n
\tag{10}
\end{equation}
имеет решения для каждой произвольной системы значений $a_1\cdots a_n$ -- за исключением еще подлежащих указанию сингулярных. Для каждой системы значений $a_1\cdots a_n$, соответствующей некоторому $g_\Gamma$ относительно $\Omega^*$, соответствующие $\lambda_i$ как естественные иррациональности, принадлежащие группе,\footnote{См. также ниже!} становятся величинами из $\Omega^*$. Таким образом, когда $\lambda_1\cdots\lambda_n$ пробегают все системы значений, $g(x)$ пробегает всю совокупность уравнений в модифицированном смысле; эта совокупность содержит всю совокупность $g_\Gamma$, которым соответствуют значения из $\Omega^*$. А так как, обратно, по теореме Гильберта о неприводимости значениям $\lambda$ из $\Omega^*$ «вообще говоря» также соответствуют уравнения $g_\Gamma$, то в модифицированном в пункте 1 смысле посредством параметрического представления (9) действительно получается вся совокупность всех $g_\Gamma$ относительно $\Omega^*$.

Теперь остается определить сингулярные системы значений, для которых параметрическое представление отказывает. Пусть функции минимальной базы, разделенные на числитель и знаменатель, заданы как
\[
\varphi_i(x)=\frac{\psi_i(x)}{\chi_i(x)}.
\]
Подставляя это в тождества (5), предположим, что они без сокращения переходят в
\begin{equation}
\sigma_k(x)=\frac{G_k(x)}{H_k(x)}\quad\text{или}\quad H_k(x)\cdot\sigma_k(x)=G_k(x),
\tag{11}
\end{equation}
так что произведение всех $H_k(x)$ делится на произведение всех знаменателей $\chi_i(x)$. Для всех систем корней $\alpha_1\cdots\alpha_n$, для которых произведение
\[
H_1(\alpha)\cdot H_2(\alpha)\cdots H_n(\alpha)\neq0
\]
и, следовательно, ни один знаменатель $\chi_i(\alpha)$ также не обращается в нуль, можно в (11) разделить на $H_k(\alpha)$ и получить
\begin{equation}
\sigma_k(\alpha)=\frac{G_k(\alpha)}{H_k(\alpha)}=F_k(\varphi_1(\alpha);\cdots;\varphi_n(\alpha)),
\tag{12}
\end{equation}
причем $\varphi_i(\alpha)$ вследствие необращения знаменателей в нуль принимают конечные значения, а для уравнений $g_\Gamma$ становятся величинами из $\Omega^*$. Уравнение (12) задает систему решений для (10), как только $\alpha_1\cdots\alpha_n$ определены так, что $a_k=(-1)^k\sigma_k(\alpha)$.\footnote{Тем самым решение уравнения сведено к решению системы независимых функций:
\[
\varphi_1(\alpha_1\cdots\alpha_n)=\lambda_1;\quad \cdots;\quad \varphi_n(\alpha_1\cdots\alpha_n)=\lambda_n.
\]}

\emph{Сингулярными} могут, следовательно, стать только такие системы коэффициентов $a_k=(-1)^k\sigma_k(\alpha)$, для которых
\[
H_1(\alpha)\cdot H_2(\alpha)\cdots H_n(\alpha)=0;
\]
то есть для которых (11) вместо представления переходит в $0=0$. Тогда необходимо специальное исследование, содержатся ли среди так определенных сингулярных $a_1\cdots a_n$ и такие, которым соответствуют уравнения $g_\Gamma$, и для каких числовых полей.\footnote{Например, как показал Ф. Зейдельман в указанном месте, только для альтернирующей группы уравнений 3-й и 4-й степеней возникают такие уравнения, соответствующие сингулярным системам значений, и притом только для тех числовых полей $\Omega^*$, которые содержат поле третьих корней из единицы. В каждом случае можно указать простое дополнительное представление.} В итоге получаем следующую

\emph{Теорему. Если лагранжев родовой домен, принадлежащий группе $\Gamma$, обладает минимальной базой}\footnote{Из существования одной минимальной базы сразу следует существование сколь угодно многих, выводимых из исходной обратимыми рациональными преобразованиями.} \emph{с коэффициентами из $\Omega$, то для каждого числового поля $\Omega^*$, содержащего $\Omega$, вся совокупность уравнений с заданной группой $\Gamma$ относительно $\Omega^*$ может быть рационально построена параметрическим представлением (2), возможно, за исключением сингулярных относительно этого параметрического представления уравнений, которые можно указать отдельно. При этом параметры должны пробегать все значения из $\Omega^*$, не вызывающие редукции группы. Построение возможно для всякого произвольного числового поля, если область коэффициентов $\Omega$ минимальной базы равна полю $\mathfrak R$ всех рациональных чисел.}

Так как по теореме Гильберта о неприводимости параметрическое многообразие не понижается исключением тех значений из $\Omega^*$, которые вызывают редукцию группы, то справедливо также: \emph{из существования минимальной базы следует, что многообразие уравнений с группой $\Gamma$ относительно $\Omega^*$ равно многообразию уравнений без аффекта; аффект не понижает многообразие.}

\begin{center}
§ 2.\\[0.75\baselineskip]
\textbf{Применение к специальным уравнениям.}
\end{center}

Теперь докажем существование минимальной базы для всех групп, возникающих у уравнений 3-й и 4-й степеней; для этого сначала совершенно общо покажем, что для уравнений $n$-й степени задачу можно свести к задаче от $n-2$ неизвестных. Первая редукция соответствует линейному преобразованию Чирнгауза для удаления второго по старшинству члена. В качестве первой базисной функции берем $\varphi_1(x)=\sigma_1(x)$; далее замечаем, что $\mathfrak L_\Gamma$ рационально выводится из всех \emph{однородных} форм от $x_1\cdots x_n$, допускающих $\Gamma$. Пусть $p(x_1\cdots x_n)$ -- такая форма; тогда также
\[
q(x)=p\left(x_1-\frac{\sigma_1(x)}{n};\cdots;x_n-\frac{\sigma_1(x)}{n}\right)=p\left(x_i-\frac{\sigma_1(x)}{n}\right)=p(y_i)
\]
принадлежит $\mathfrak L_\Gamma$, поскольку она допускает $\Gamma$, и кроме того имеем
\[
p(x)\equiv p\left(x_i-\frac{\sigma_1(x)}{n}\right)\operatorname{mod}\sigma_1(x)
\]
или
\[
p(x)=q(x)+\sigma_1(x)\cdot p_1(x)
\]
с $p_1(x)$, принадлежащим $\mathfrak L_\Gamma$ и имеющим меньшую степень, чем $p(x)$. Продолжение этого приема показывает, что все однородные формы из $\mathfrak L_\Gamma$, а следовательно вообще все рациональные функции из $\mathfrak L_\Gamma$, рационально выражаются через $\varphi_1(x)=\sigma_1(x)$ и через однородные формы
\[
q(x)=p\left(x_i-\frac{\sigma_1(x)}{n}\right)=p(y_i).
\]
Но между этими величинами $y_i$ существует линейное соотношение $\sum y_i=0$; следовательно, $q(x)$ \emph{зависят только от $(n-1)$ аргументов.}

Вторая редукция основана на том, что $q(x)$ являются однородными формами своих аргументов. В качестве второй базисной функции берем
\[
\begin{aligned}
\varphi_2(x)&=\frac{\tau_3(x)}{\tau_2(x)},\\
\tau_3(x)&=\sigma_3\left(x_i-\frac{\sigma_1(x)}{n}\right)=\sigma_3(y_i),\\
\tau_2(x)&=\sigma_2\left(x_i-\frac{\sigma_1(x)}{n}\right)=\sigma_2(y_i).
\end{aligned}
\]
то есть $\varphi_2(x)$ является однородно-дробной функцией первой степени от $(n-1)$ аргументов, именно от $y_i$ при $\sum y_i=0$. С помощью $\varphi_2(x)$ все $q(x)$ можно рационально выразить через функции
\[
r(x)=\frac{q(x)}{\varphi_2(x)^\lambda}=\frac{q(x)\cdot\tau_2(x)^\lambda}{\tau_3(x)^\lambda}=\frac{p(y_i)\cdot\sigma_2(y_i)^\lambda}{\sigma_3(y_i)^\lambda},
\]
также принадлежащие $\mathfrak L_\Gamma$, где $\lambda$ обозначает степень $q(x)$. Тогда $r(x)$ становится однородной функцией нулевой степени, а значит зависит только от $(n-2)$ аргументов, то есть от отношений $y_1:y_2:\cdots:y_n$ при $\sum y_i=0$. Следовательно, $\mathfrak L_\Gamma$ рационально выражается через
\[
\varphi_1(x)=\sigma_1(x);\quad \varphi_2(x)=\frac{\tau_3(x)}{\tau_2(x)}
\]
и через \emph{поле всех функций $r(x)$ из $\mathfrak L_\Gamma$, зависящих от этих $(n-2)$ аргументов},\footnote{Для уравнения $n$-й степени без второго члена
\[
f(x)=x^n+a_2x^{n-2}+a_3x^{n-3}+\cdots+a_n=0
\]
этой второй редукции соответствует подстановка, всегда возможная при $a_2\neq0$, $a_3\neq0$,
\[
y=\frac{x\cdot a_2}{a_3},
\]
посредством которой $f(x)$, с точностью до множителя, переходит в
\[
g(y)=y^n+\frac{a_2^3}{a_3^2}y^{n-2}+\frac{a_2^3}{a_3^2}y^{n-3}+\frac{a_4\cdot a_2^4}{a_3^4}y^{n-4}+\cdots+a_n\frac{a_2^n}{a_3^n}=0;
\]
то есть в уравнение, содержащее только $n-2$ параметра:
\[
g(y)=y^n+c\cdot y^{n-2}+c\cdot y^{n-3}+c_4y^{n-4}+\cdots+c_n=0.
\]} чем и выполнена желаемая редукция.

Для уравнений 3-й и 4-й степеней это, в частности, дает редукцию к функциональным полям от одного и соответственно двух аргументов. Но для них минимальная база всегда существует по теоремам Люрота и Кастельнуово.\footnote{J. Lüroth: Beweis eines Satzes über rationale Kurven, Math. Ann. 9 (1875); G. Castelnuovo: Sulla razionalità delle involuzioni piane, Math. Ann. 44 (1893). То, что для полей более чем от двух неизвестных минимальная база может не существовать, Ф. Энрикес показал на контрпримере, Rend. Acc. Linc., Vol. 21, 21 Jan. 1912.} В случае одного неизвестного минимальную базу можно вывести рациональными операциями, и потому специально для $\mathfrak L_\Gamma$ она содержит только рационально-числовые коэффициенты. В случае двух неизвестных по Кастельнуово еще могла бы идти речь о базе с коэффициентами из алгебраического числового поля $\Omega$; фактическое исследование (см. Ф. Зейдельман, цит. место) показывает, что и здесь для каждого $\mathfrak L_\Gamma$ получается база с рационально-числовыми коэффициентами. Это почти без вычисления видно уже из того, что для четверной группы можно выбрать такую рационально-числовую базу, а именно
\[
\varphi_1(x)=\sigma_1(x);\quad \varphi_2(x)=\frac{v w}{u};\quad \varphi_3(x)=\frac{u\cdot w}{v};\quad \varphi_4(x)=\frac{u\cdot v}{w},
\]
где
\[
u=x_1+x_2-x_3-x_4;\quad v=x_1-x_2+x_3-x_4;\quad w=x_1-x_2-x_3+x_4,
\]
так что альтернирующая группа переходит в циклическую группу трех базисных функций $\varphi_2,\varphi_3,\varphi_4$, а каждая группа восьмого порядка -- в симметрическую группу двух из этих функций; тем самым осуществляется редукция к группам уравнений 3-й и 2-й степеней. Для циклической группы рационально-числовая база уже известна от Абеля, хотя и не была сформулирована именно так.\footnote{См. Weber Algebra (малое издание), § 93, формула (12).} Итак, доказано: \emph{вся совокупность уравнений 3-й и 4-й степеней с заданной группой -- возможно, за исключением сингулярных относительно параметрического представления -- может быть рационально построена параметрическим представлением для любой заданной области рациональности; их многообразие равно многообразию уравнений без аффекта.}

Фактическое исследование показывает (см. Ф. Зейдельман, цит. место), что дополнительное представление требуется только для альтернирующей группы относительно областей рациональности, содержащих поле третьих корней из единицы; во всех остальных случаях параметрическое представление дает всю совокупность без оговорок.

Следует еще отметить, что минимальная база известна и для всех \emph{абелевых групп}.\footnote{E. Fischer: Die Isomorphie der Invariantenkörper der endlichen Abelschen Gruppen linearer Transformationen. Gött. Nachr. 1915, и Zur Theorie der endlichen Abelschen Gruppen. Math. Ann. 77 (1915).} Однако здесь речь идет о базе с коэффициентами из циклотомического поля, выводимого из характеров группы. Поэтому для построения уравнений с заданной группой в этом случае не удается получить новых результатов.

Гёттинген, июль 1916 г.

\clearpage
% END UNIT BODY
% END PRODUCER UNIT 073
\endgroup


\clearpage
\begingroup
\setcounter{footnote}{0}
% BEGIN PRODUCER UNIT 074
% Source: C:/Users/Floris/Documents/interlanguage/03_projects/noether/02_slavic_working_corpus/translations/paper12/russian/v001/Noether_Paper12_Russian_v001.tex
% Identity: 27442 B / SHA-256 D17508123567DD73325D98A4BC3461965E2B129A76470AE9B486ABDF3D32B563
\setlength{\parindent}{1.25em}
\setlength{\parskip}{0pt}
\emergencystretch=30em
\allowdisplaybreaks
\raggedbottom
\providecommand{\Rfield}{\mathfrak R}
\providecommand{\bdelta}{\bar{\delta}}
% BEGIN UNIT BODY
% Unit: Paper 12. Source segments P12-S0001--P12-S0019.

\begin{center}
{\Large\bfseries 12. Инварианты произвольных дифференциальных выражений}\par
\vspace{1.5em}
Nachr. v. d. Ges. d. Wiss. zu Göttingen 1918, с. 37--44
\end{center}

\vspace{3em}
\begin{center}
Представлено на заседании 25 января 1918 года Ф. Клейном.
\end{center}

Под дифференциальным выражением я понимаю функцию
\[
f(x,dx)=f(x_1,\ldots,x_n;\,dx_1,\ldots,dx_n)
\]
от \(n\) переменных \(x_1,\ldots,x_n\) и их дифференциалов
\(dx_1,\ldots,dx_n\), которая предполагается аналитической в обеих системах по \(n\) аргументов, но не обязательно вещественной. Теперь за основу я беру группу всех аналитических преобразований переменных, которой соответствует группа всех линейных преобразований дифференциалов:
\begin{equation}
x_i=x_i(y_1,\ldots,y_n);\qquad
dx_i=\sum_k\frac{\partial x_i}{\partial y_k}\,dy_k;\qquad
\delta x_i=\sum_k\frac{\partial x_i}{\partial y_k}\,\delta y_k,
\tag{1}
\end{equation}
при этом \(f(x,dx)\) пусть переходит в \(g(y,dy)\). Под инвариантом \(f(x,dx)\) я понимаю инвариант относительно этой группы и относительно индуцированной ею группы для высших дифференциалов \(d^2x,ddx,\ldots\) и для производных от \(f(x,dx)\); то есть такую аналитическую функцию \(J\), что для произвольного \(f\) и в силу (1) тождественно по переменным и всем входящим дифференциалам выполняется
\[
\begin{aligned}
&J\!\left(f,\frac{\partial f}{\partial dx}\cdots
\frac{\partial^{\rho+\sigma}f}{\partial x^\rho\partial dx^\sigma}
\cdots dx,\delta x,d^2x,\ldots\right)\\
&\qquad =
J\!\left(g,\frac{\partial g}{\partial dy}\cdots
\frac{\partial^{\rho+\sigma}g}{\partial y^\rho\partial dy^\sigma}
\cdots dy,\delta y,d^2y,\ldots\right).\footnotemark
\end{aligned}
\]
\footnotetext{\(\displaystyle
\frac{\partial^{\rho+\sigma}f}{\partial x^\rho\partial dx^\sigma}\)
всюду употребляется сокращенно для
\[
\frac{\partial^{\rho+\sigma}f}
{\partial x_{i_1}\cdots\partial x_{i_\rho}
 \partial dx_{k_1}\cdots\partial dx_{k_\sigma}},
\]
где индексы означают какие-либо из чисел от \(1\) до \(n\).}
Совершенно соответственно определяется инвариант одновременной системы дифференциальных выражений. Если такой инвариант, в частности, содержит только первые дифференциалы \(dx,\delta x\) и только производные по этим дифференциалам, а не по переменным, то появление переменных в функциях и линейное преобразование дифференциалов не принимаются во внимание; эти специальные инварианты являются инвариантами относительно группы всех линейных преобразований дифференциалов \(dx,\delta x\) с неопределенными коэффициентами и далее будут называться проективными инвариантами одновременной системы.

Для квадратичных однородных дифференциальных форм Кристоффель (Crelle 70) и Риччи (Math. Annalen 54) построили такую систему инвариантных образований, что все инварианты становятся проективными инвариантами этой одновременной системы; тем самым вопросы о всей совокупности инвариантов и об эквивалентности сводятся к вопросам линейной теории инвариантов. Это я хотела бы назвать «теоремой редукции». Полная система состоит из счетно бесконечного множества форм; но для инвариантов каждого конечного порядка \(\rho\), то есть таких, которые не содержат производных по переменным выше \(\rho\)-го порядка, она содержит только конечное число форм.

Ниже я показываю справедливость «теоремы редукции» для произвольных дифференциальных выражений; и здесь снова речь идет о полной системе счетно бесконечного множества инвариантных образований, но для каждого конечного порядка только о конечном их числе. Однако, тогда как Кристоффель и Риччи для порождения инвариантов и доказательства теоремы редукции опираются на трудно обозримые элиминации, я порождаю инварианты непосредственно инвариантными процессами дифференцирования и варьирования; последние просто можно понимать как процессы дифференцирования по другим параметрам. Алгоритм этих инвариантных процессов дифференцирования, в связи с Лагранжем (Mécanique analytique, 2-я часть, разд. IV), развили Риман (Werke XXII) и Липшиц (Crelle 70, 72) для квадратичных дифференциальных форм, однако не дойдя до теоремы редукции. Вспомогательное средство для доказательства теоремы редукции дают введенные Риманом для квадратичных форм «нормальные координаты» (Werke XIII), которые преобразуют экстремали соответствующей вариационной задачи, исходящие из одной точки, в прямые; но они существуют только для однородных дифференциальных выражений
\(f(x,\varkappa\cdot dx)=\Phi(\varkappa)\cdot f(x,dx)\footnote{Легко видеть, что \(\Phi(\varkappa)=\varkappa^c\), где \(c\) обозначает произвольную вещественную или комплексную величину.}\). Неоднородный случай, однако, сводится к этому.

\begin{center}
\textbf{I. Порождение инвариантов.}
\end{center}

Последовательность проективных инвариантов \(f(x,dx)\), вообще говоря не обрывающаяся, дается полярами. Из-за линейности преобразования дифференциалов из \(f(x,dx)=g(y,dy)\) также следует
\(f(x,dx+\lambda\delta x)=g(y,dy+\lambda\delta y)\); а отсюда, разлагая по \(\lambda\), получаются характерные для инвариантности соотношения:
\begin{equation}
\begin{gathered}
f(dx)=g(dy),\\
f_\delta=\sum\frac{\partial f(dx)}{\partial dx}\,\delta x
       =\sum\frac{\partial g(dy)}{\partial dy}\,\delta y,\\
\vdots\\
f_{\delta^\sigma}=
\sum\frac{\partial^\sigma f(dx)}{\partial dx^\sigma}\,\delta x^\sigma
=
\sum\frac{\partial^\sigma g(dy)}{\partial dy^\sigma}\,\delta y^\sigma,\\
\vdots
\end{gathered}
\tag{2}
\end{equation}

Каждая поляра через применение вариационного процесса \(\bdelta\) дает основание для необрывающейся последовательности общих инвариантов:
\begin{equation}
\begin{gathered}
\bdelta^\rho f(x,dx)=\bdelta^\rho g(y,dy),\\
\vdots\\
\bdelta^\rho f_{\delta^\sigma}
=\bdelta^\rho\sum\frac{\partial^\sigma f(dx)}{\partial dx^\sigma}\,
\delta x^\sigma
=\bdelta^\rho\sum\frac{\partial^\sigma g(dy)}{\partial dy^\sigma}\,
\delta y^\sigma,\\
\vdots\\[0.2em]
(\rho=0,1,2,\ldots).
\end{gathered}
\tag{3}
\end{equation}
Соотношения (2) и (3) в силу (1) позволяют считывать законы преобразования для производных от \(f(x,dx)\), но далее это не понадобится. Из инвариантов (3) теперь линейной комбинацией можно получить «нормальную форму \(\rho\)-й вариации», которая при \(\rho=1\) встречается у Лагранжа, а при \(\rho=2\) у Римана; она характеризуется тем, что смешанные дифференциалы порядка \(\rho+1\): \(d^\rho\delta\), \(d^{\rho-1}\delta^2,\ldots\), элиминированы. Для нее получаем:
\begin{equation}
\begin{gathered}
\Omega_1=\delta f(dx)-d f_\delta(dx),\\
\Omega_2=\delta^2 f-\delta d f_\delta+\frac12 d^2 f_{\delta^2},\\
\vdots\\
\Omega_\rho=\delta^\rho f-\delta^{\rho-1}d f_\delta
+\frac12\delta^{\rho-2}d^2 f_{\delta^2}
+\cdots+\frac{(-1)^\rho}{\rho!}\,d^\rho f_{\delta^\rho},\\
\vdots
\end{gathered}
\tag{4}
\end{equation}

Из инвариантов (4), обобщая подход Римана, можно перейти к инвариантам, содержащим только первые дифференциалы; такие инварианты будут называться «основными функциями». А именно, если в \(\Omega_1\) заменить дифференциалы дифференциальными отношениями, то \(\Omega_1=0\), тождественно по \(\delta x\), дает ровно лагранжевы уравнения вариационной задачи, принадлежащей
\[
f\!\left(\frac{dx}{dt}\right).
\]
Теперь я ставлю инвариантное условие, что направление \((d+\lambda\delta)\) удовлетворяет этим лагранжевым уравнениям:
\begin{equation}
\Omega_1(d+\lambda\delta)=
\bdelta f(dx+\lambda\delta x)
-(d+\lambda\delta)f_{\bdelta}(dx+\lambda\delta x)=0
\quad(\text{тождественно по }\delta x),
\tag{5}
\end{equation}
откуда подстановкой \(\bdelta=d+\lambda\delta\) для однородного \(f(dx)\) еще следует
\begin{equation}
(d+\lambda\delta)f(dx+\lambda\delta x)=0,
\tag{6}
\end{equation}
за исключением случая однородности первого порядка. В последнем случае (6) следует добавить как новое условие, поскольку тогда система уравнений (5) представляет только \((n-1)\) независимых условий. Если в (5), соответственно в (5) и (6), приравнять к нулю коэффициенты при нулевой, первой и второй степенях \(\lambda\), то получаются три инвариантные системы уравнений
\(\varphi=0\), \(\psi=0\), \(\chi=0\), тождественно по \(\delta x\), посредством которых \(d^2x,ddx,\delta^2x\) можно выразить через первые дифференциалы. Дальнейшие инвариантные системы уравнений
\begin{equation}
\begin{gathered}
d^\sigma\varphi=0,\qquad d^\sigma\psi=0,\qquad
d^\sigma\chi=0,\qquad d^{\sigma-1}\delta\chi=0,\\
d^{\sigma-2}\delta^2\chi=0,\ldots,\delta^\sigma\chi=0
\qquad(\sigma=0,1,2,\ldots)
\end{gathered}
\tag{7}
\end{equation}
тогда как раз достаточны, чтобы выразить все высшие дифференциалы через первые.\footnote{Только линейная форма \(\sum A_i\,dx_i\) требует особого рассмотрения, поскольку здесь система уравнений (5) не содержит вторых дифференциалов.} Следовательно, функции (4), соединенные с инвариантной системой уравнений (7), ведут к основным функциям \([\Omega_\rho]\), содержащим только первые дифференциалы; при этом \([\Omega_1]\) тождественно исчезает.

Так же из дифференциала \(\delta h(dx,\delta x)\) основной функции \(h\) посредством (7) снова можно получить основную функцию, которую следует назвать «ковариантной производной» \([h^{(1)}]\) от \(h\); вообще \([h^{(\nu)}]\) будет обозначать ковариантную производную от \([h^{(\nu-1)}]\). Эти два процесса таким образом ведут к дважды бесконечной последовательности основных функций:
\begin{equation}
\begin{array}{cccccc}
[\Omega_2], & [\Omega_2^{(1)}], & [\Omega_2^{(2)}], & \ldots,
& [\Omega_2^{(\nu)}], & \ldots\\
\vdots &&&& \vdots&\\
{}[\Omega_\rho], & [\Omega_\rho^{(1)}], & \ldots, & \ldots,
& [\Omega_\rho^{(\nu)}], & \ldots\\
\vdots & \vdots &&& \vdots&
\end{array}
\tag{8}
\end{equation}

Далее оказывается, что левые части уравнений, выражающих высшие дифференциалы через первые, когредиентны первым дифференциалам, то есть подчиняются тем же линейным преобразованиям. Поэтому, если вместо высших дифференциалов ввести эти левые части \(p,q,r,\ldots\) как аргументы инварианта, то, кроме производных, он зависит только от взаимно когредиентных величин. Указанное выше образование функций (8) просто сводится к введению этих новых величин; каждый инвариант тем самым переходит в сумму инвариантов, и из нее выбирается тот, который содержит именно \(dx,\delta x\), но не содержит \(p,q,r,\ldots\).

Ниже показывается, что при предположении однородности
\[
f(\varkappa\cdot dx)=\Phi(\varkappa)\cdot f(dx)
\]
основные функции (8) исчерпывают искомую полную систему.

\begin{center}
\textbf{II. Нормальные координаты, эквивалентность и теорема редукции.}
\end{center}

К нормальным координатам я прихожу интегрированием вариационной задачи, принадлежащей \(f\!\left(\frac{dx}{dt}\right)\):
\begin{equation}
\begin{gathered}
\Omega_1=\delta f\!\left(\frac{dx}{dt}\right)
-\frac{d}{dt}f_\delta\!\left(\frac{dx}{dt}\right)=0
\quad(\text{тождественно по }\delta x),\\
\frac{d}{dt}f\!\left(\frac{dx}{dt}\right)=0
\quad(\text{как следствие или дополнительное условие}).
\end{gathered}
\tag{9}
\end{equation}
Из-за предполагаемой однородности \(f(dx)\) эта система уравнений переходит в себя при подстановке параметра \(t=\varkappa\cdot\tau\); следовательно, если в силу (9) выразить \(\frac{d^\rho x_i}{dt^\rho}\) как функцию \(\varphi_\rho^{(i)}\!\left(\frac{dx}{dt}\right)\) первых производных, то \(\varphi_\rho^{(i)}\) становится однородной \(\rho\)-го порядка:
\begin{equation}
\varphi_\rho^{(i)}\!\left(\varkappa\cdot\frac{dx}{dt}\right)
=\varkappa^\rho\varphi_\rho^{(i)}\!\left(\frac{dx}{dt}\right).
\tag{10}
\end{equation}
Если теперь при интегрировании (9) для \(t=t_0\) задать начальные значения
\[
x=(x)_0,\qquad \frac{dx}{dt}=\left(\frac{dx}{dt}\right)_0
\]
и положить
\begin{equation}
u_i=(t-t_0)\left(\frac{dx_i}{dt}\right)_0,
\tag{11}
\end{equation}
где параметр \((t-t_0)\) в силу \(f\!\left(\frac{dx}{dt}\right)=\mathrm{const.}\) становится пропорциональным «длине экстремали», то разложение по степеням \((t-t_0)\) в силу (10) дает
\begin{equation}
x_i-(x_i)_0=u_i+\frac12\varphi_2^{(i)}(u)+\cdots+
\frac{1}{\rho!}\varphi_\rho^{(i)}(u)+\cdots .
\tag{12}
\end{equation}
Это разложение можно рассматривать как преобразование переменных \(x_i=x_i(u)\), где величины \(u\) как раз являются римановыми нормальными координатами. Значит, по (11) и (12) эти координаты преобразуют экстремали, проходящие через \((x)_0\), в прямые. Кроме того, произвольному преобразованию переменных \(x=x(y)\) по (11) соответствует линейное преобразование \(u=u(v)\) соответствующих нормальных координат, причем коэффициенты зависят только от произвольной системы значений \((x)_0\); следовательно, дифференциалы \(du,\delta u\) подчиняются тому же преобразованию, что и сами \(u\).

В силу (12) пусть \(f(x,dx)\) переходит в \(F(u,du)\), или, разложенная по степеням \(u\),
\begin{equation}
F(u,du)=F_0(u,du)+F_1(u,du)+\cdots+F_\rho(u,du)+\cdots
\tag{13}
\end{equation}
где \(F_\rho(u,du)\) обозначает однородную форму \(\rho\)-й размерности по \(u\). Из-за линейности преобразования \(u=u(v)\) из \(F(u,du)=G(v,dv)\) отдельные однородные составные части также переходят в соответствующие:
\begin{equation}
F_\rho(u,du)=G_\rho(v,dv).
\tag{14}
\end{equation}
Итак, если взять для \((x)_0\) какую-нибудь фиксированную постоянную систему значений, то (13) и (14) показывают, что эквивалентность \(f(dx)\) с \(g(dy)\) равносильна эквивалентности соответствующих систем функций \(F_\rho(u,du)\) относительно линейного преобразования. Кроме того, \(F_\rho(u,du)\), или соответствующие \(F_\rho(\delta u,du)\), представляют инварианты \(f(dx)\), взятые в точке \(x=(x)_0\), и именно они образуют полную систему основных функций для этой точки \(x=(x)_0\). А именно, имеем
\[
F_\rho(\delta u,du)=\frac1{\rho!}\sum
\frac{\partial F_\rho}{\partial \delta u^\rho}\,\delta u^\rho;
\qquad
\frac{\partial F_\rho}{\partial \delta u^\rho}
=\left(\frac{\partial F(du)}{\partial u^\rho}\right)_0,
\]
и эта последняя производная в силу (12) выражается через производные от \(f(dx)\), взятые при \(x=(x)_0\), совершенно так же, как соответственно образованная производная от \(G(dv)\) выражается через производные от \(g(dy)\). Поскольку
\[
\left(\frac{\partial^{\rho+\sigma}F(du)}
{\partial u^\rho\partial du^\sigma}\right)_0
=
\frac{\partial^{\rho+\sigma}F_\rho(\delta u,du)}
{\partial\delta u^\rho\partial du^\sigma}
\]
каждый инвариант \(f(dx)=F(du)\), взятый в точке \(x=(x)_0\), становится проективным инвариантом \(F_\rho(\delta u,du)\), как только высшие дифференциалы заменены на \(p,q,r,\ldots\), когредиентные к \(dx,\delta x\). Но \((x)_0\) можно рассматривать как параметр; каждому инварианту в точке \(x=(x)_0\) соответствует инвариант \(f(dx)\), если только снова заменить \((x)_0\) на \(x\). Следовательно, если \(F_\rho(\delta u,du)\) становятся инвариантами \(\Psi_\rho(\delta x,dx)\), взятыми при \(x=(x)_0\), ибо при \(x=(x)_0\) величины \(du,\delta u\) переходят в \(dx,\delta x\), то сами \(\Psi_\rho\) образуют основные функции полной системы. Остается выразить эту систему основных функций через явные инварианты \(f(dx)\), а именно через функции (8). Возможность этого показывает их порождение процессами дифференцирования. Для членов высшего порядка эти процессы переходят в простые процессы поляризации, и преобразование, аналогичное разложению в ряд Клебша--Гордана, вместе с индукционным выводом относительно отброшенных членов доказывает утверждение. Если речь идет об инвариантах одновременной системы, то к основным функциям (8) нужно еще добавить ковариантные производные остальных дифференциальных выражений, как показывает введение в эти выражения нормальных координат, принадлежащих \(f(dx)\).

Эквивалентность \(f(dx)\) с \(g(dy)\), которая была сведена к (14), следовательно, также равносильна эквивалентности \(\Psi_\rho\), или равным образом функций (8), относительно линейного преобразования дифференциалов, без необходимости в особых условиях интегрируемости, как следует из возможности сведения к (14). Тем самым теорема редукции доказана во всех частях. В частности, тождественное исчезновение функциональной последовательности \([\Omega_2],[\Omega_3],\ldots,[\Omega_\rho],\ldots\) необходимо и достаточно для того, чтобы \(f(dx)\) можно было преобразовать в выражение с постоянными коэффициентами.

Наконец, если вместо группы всех аналитических преобразований положить в основу подгруппу, то и группа соответствующих линейных преобразований \(u\) переходит в подгруппу проективной группы; инварианты \(f(dx)\) снова становятся инвариантами функций (8) относительно линейного преобразования, но теперь относительно этой подгруппы. Так, посредством однороднения, случай неоднородных функций \(f(dx)\) сводится к аффинной группе; полную систему здесь можно вывести из функций (8), образованных для одной дополнительной переменной.

Из теоремы редукции получается специальный случай квадратичных форм, более общо форм \(p\)-й размерности: здесь система функций обрывается на \([\Omega_2]\), более общо на \([\Omega_p]\) и ее ковариантных производных, поскольку все последующие основные функции становятся их проективными инвариантами. Этим объясняется выдающееся положение римановой формы кривизны \([\Omega_2]\) при квадратичных формах. Вопрос об эквивалентности квадратичных форм, соответственно форм \(p\)-й размерности, именно сводится к эквивалентности \([\Omega_2]\), соответственно \([\Omega_2]\ldots[\Omega_p]\), и их ковариантных производных относительно линейного преобразования; а тождественное исчезновение \([\Omega_2]\), соответственно \([\Omega_2]\ldots[\Omega_p]\), необходимо и достаточно для преобразования в формы с постоянными коэффициентами, чем для квадратичных форм высказан один из наиболее известных результатов.

Более подробное изложение должно вскоре появиться в Math. Annalen.

\clearpage
% END UNIT BODY
% END PRODUCER UNIT 074
\endgroup


\clearpage
\begingroup
\setcounter{footnote}{0}
% BEGIN PRODUCER UNIT 075
% Source: C:/Users/Floris/Documents/interlanguage/03_projects/noether/02_slavic_working_corpus/translations/paper13/russian/v001/Noether_Paper13_Russian_v001.tex
% Identity: 83745 B / SHA-256 9582F1D52E4E37521FD1E7BEDC023B2874E029733A5D4E5282487A85B2D64020
\setlength{\parindent}{1.25em}
\setlength{\parskip}{0pt}
\emergencystretch=30em
\allowdisplaybreaks
\raggedbottom
\providecommand{\G}{\mathfrak G}
\providecommand{\eps}{\varepsilon}
\providecommand{\D}{\Delta}
\providecommand{\p}{\partial}
\providecommand{\dd}{\mathrm d}
\providecommand{\bd}{\bar{\delta}}
\providecommand{\lg}{\log}
% BEGIN UNIT BODY
% Unit: Paper 13. Source segments P13-S0001--P13-S0044.

\begin{center}
{\Large\bfseries 13. Инвариантные вариационные задачи}\par
\vspace{1.5em}
Nachr. v. d. Ges. d. Wiss. zu Göttingen 1918, с. 235--257
\end{center}

\vspace{3em}
\begin{center}
Представлено Ф. Клейном на заседании 26 июля 1918 года.\footnote{Окончательная редакция рукописи была подана лишь в конце сентября.}
\end{center}

Речь идет о вариационных задачах, которые допускают непрерывную группу в смысле Ли; вытекающие отсюда следствия для соответствующих дифференциальных уравнений получают свое наиболее общее выражение в теоремах, сформулированных в \S{}1 и доказанных в следующих параграфах. Об этих дифференциальных уравнениях, возникающих из вариационных задач, можно делать гораздо более точные утверждения, чем о произвольных дифференциальных уравнениях, допускающих группу и составляющих предмет исследований Ли. Следовательно, дальнейшее основано на соединении методов формального вариационного исчисления с методами теории групп Ли. Для специальных групп и вариационных задач такое соединение методов не ново; я упомяну Гамеля и Герглотца для специальных конечных групп, Лоренца и его учеников (например, Фоккера), Вейля и Клейна для специальных бесконечных групп.\footnote{Гамель: \emph{Math. Ann.} т. 59 и \emph{Zeitschrift f. Math. u. Phys.} т. 50. Герглотц: \emph{Ann. d. Phys.} (4) т. 36, особенно \S{}9, с. 511. Фоккер, Verslag d. Amsterdamer Akad., 27.01.1917. О дальнейшей литературе см. вторую заметку Клейна: Göttinger Nachrichten, 19 июля 1918. В только что вышедшей работе Кнезера (\emph{Math. Zeitschrift} т. 2) речь идет о построении инвариантов сходным методом.} В частности, вторая заметка Клейна и настоящее изложение взаимно повлияли друг на друга; по этому поводу я могу сослаться на заключительные замечания заметки Клейна.

\section*{\S{}1. Предварительные замечания и формулировка теорем}

Все встречающиеся ниже функции следует считать аналитическими или по крайней мере непрерывными и конечное число раз непрерывно дифференцируемыми, а в рассматриваемой области однозначными.

Под «группой преобразований», как известно, понимают такую систему преобразований, что для каждого преобразования существует содержащееся в системе обратное преобразование, и что композиция любых двух преобразований системы снова принадлежит системе. Группа называется конечной непрерывной \(\G_\rho\), если ее преобразования содержатся в наиболее общем преобразовании, аналитически зависящем от \(\rho\) существенных параметров \(\eps\) (то есть эти \(\rho\) параметров не должны представляться как \(\rho\) функций от меньшего числа параметров). Соответственно под бесконечной непрерывной \(\G_{\infty\rho}\) понимают группу, наиболее общие преобразования которой аналитически или по крайней мере непрерывно и конечное число раз непрерывно дифференцируемо зависят от \(\rho\) существенных произвольных функций \(p(x)\) и их производных. Промежуточным звеном между ними является группа, зависящая от бесконечного числа параметров, но не от произвольных функций. Наконец, смешанной группой называют такую, которая зависит и от произвольных функций, и от параметров.\footnote{Ли в \emph{Grundlagen für die Theorie der unendlichen kontinuierlichen Transformationsgruppen} (Ber. d. K. Sächs. Ges. der Wissensch. 1891) определяет бесконечную непрерывную группу как группу преобразований, преобразования которой задаются наиболее общими решениями системы частных дифференциальных уравнений, если эти решения зависят не только от конечного числа параметров. Тем самым получается один из указанных выше типов, отличных от конечной группы; тогда как, наоборот, предельный случай бесконечного числа параметров не обязан удовлетворять системе дифференциальных уравнений.}

Пусть \(x_1,\ldots,x_n\) являются независимыми переменными, а \(u_1(x),\ldots,u_\mu(x)\) зависящими от них функциями. Если подчинить \(x\) и \(u\) преобразованиям некоторой группы, то вследствие предполагаемой обратимости преобразований среди преобразованных величин снова должно содержаться ровно \(n\) независимых --- \(y_1,\ldots,y_n\); остальные зависящие от них величины обозначим \(v_1(y),\ldots,v_\mu(y)\). В преобразованиях могут также появляться производные \(u\) по \(x\), то есть \(\p u/\p x,\ \p^2u/\p x^2,\ldots\).\footnote{Я, насколько возможно и при суммированиях, опускаю индексы; например, пишу \(\p^2u/\p x^2\) для \(\p^2u_\nu/\p x_i\p x_k\) и т. д.} Функция называется инвариантом группы, если существует соотношение:
\[
P\!\left(x,u,\frac{\p u}{\p x},\frac{\p^2u}{\p x^2},\ldots\right)
 =
P\!\left(y,v,\frac{\p v}{\p y},\frac{\p^2v}{\p y^2},\ldots\right).
\]
В частности, интеграл \(I\) является инвариантом группы, если существует соотношение:
\begin{equation}
\begin{aligned}
I&=\int\!\cdots\!\int f\!\left(x,u,\frac{\p u}{\p x},\frac{\p^2u}{\p x^2},\ldots\right)\dd x\\
 &=\int\!\cdots\!\int f\!\left(y,v,\frac{\p v}{\p y},\frac{\p^2v}{\p y^2},\ldots\right)\dd y,
\end{aligned}
\tag{1}
\end{equation}
интегрируемое по произвольной действительной \(x\)-области и соответствующей \(y\)-области.\footnote{Сокращенно пишу \(\dd x,\dd y\) для \(\dd x_1\cdots \dd x_n,\dd y_1\cdots \dd y_n\). Все аргументы \(x,u,\eps,p(x)\), входящие в преобразования, следует считать действительными, тогда как коэффициенты могут быть комплексными. Но так как в окончательных результатах речь идет о тождествах по \(x,u\), параметрам и произвольным функциям, они справедливы и для комплексных значений, если только все встречающиеся функции считать аналитическими. Впрочем, значительная часть результатов может быть обоснована без интегралов, так что ограничение действительным не необходимо и при обосновании. Напротив, рассуждения в конце \S{}2 и начале \S{}5, по-видимому, нельзя провести без интегралов.}

С другой стороны, для произвольного, не обязательно инвариантного интеграла \(I\) я образую первую вариацию \(\delta I\) и преобразую ее по правилам вариационного исчисления частным интегрированием. Как известно, если \(\delta u\) вместе со всеми встречающимися производными считать исчезающей на границе, а в остальном произвольной, получается:
\begin{equation}
\delta I=\int\!\cdots\!\int \delta f\,\dd x
=\int\!\cdots\!\int\left(\sum_i\psi_i\!\left(x,u,\frac{\p u}{\p x},\ldots\right)\delta u_i\right)\dd x,
\tag{2}
\end{equation}
где \(\psi\) обозначают лагранжевы выражения, то есть левые части уравнений Лагранжа соответствующей вариационной задачи \(\delta I=0\). Этому интегральному соотношению соответствует безинтегральное тождество по \(\delta u\) и ее производным, возникающее, если явно записать граничные члены. Как показывает частное интегрирование, эти граничные члены являются интегралами от дивергенций, то есть от выражений
\[
\Div A=\frac{\p A_1}{\p x_1}+\cdots+\frac{\p A_n}{\p x_n},
\]
где \(A\) линейно по \(\delta u\) и ее производным. Поэтому получается:
\begin{equation}
\sum_i\psi_i\,\delta u_i=\delta f+\Div A.
\tag{3}
\end{equation}
Если, в частности, \(f\) содержит только первые производные \(u\), то для простого интеграла тождество (3) совпадает с так называемым «центральным уравнением Лагранжа»:
\begin{equation}
\sum_i \psi_i\,\delta u_i
=\delta f-\frac{\dd}{\dd x}\left(\sum_i\frac{\p f}{\p u_i'}\,\delta u_i\right),
\qquad \left(u_i'=\frac{\dd u_i}{\dd x}\right),
\tag{4}
\end{equation}
а для \(n\)-кратного интеграла (3) переходит в:
\begin{equation}
\sum_i\psi_i\,\delta u_i
=
\delta f
-\frac{\p}{\p x_1}\left(\sum_i\frac{\p f}{\p \left(\frac{\p u_i}{\p x_1}\right)}\,\delta u_i\right)
-\cdots
-\frac{\p}{\p x_n}\left(\sum_i\frac{\p f}{\p \left(\frac{\p u_i}{\p x_n}\right)}\,\delta u_i\right).
\tag{5}
\end{equation}
Для простого интеграла и \(\kappa\) производных от \(u\) (3) задается формулой:
\begin{equation}
\begin{aligned}
\sum_i\psi_i\,\delta u_i
={}&\delta f
-\frac{\dd}{\dd x}\Bigg\{\sum_i\left[
\binom{1}{1}\frac{\p f}{\p u_i^{(1)}}\delta u_i
+\binom{2}{1}\frac{\p f}{\p u_i^{(2)}}\delta u_i^{(1)}
+\cdots+
\binom{\kappa}{1}\frac{\p f}{\p u_i^{(\kappa)}}\delta u_i^{(\kappa-1)}
\right]\Bigg\}\\
&+\frac{\dd^2}{\dd x^2}\Bigg\{\sum_i\left[
\binom{2}{2}\frac{\p f}{\p u_i^{(2)}}\delta u_i
+\binom{3}{2}\frac{\p f}{\p u_i^{(3)}}\delta u_i^{(1)}
+\cdots+
\binom{\kappa}{2}\frac{\p f}{\p u_i^{(\kappa)}}\delta u_i^{(\kappa-2)}
\right]\Bigg\}\\
&+\cdots
+(-1)^\kappa\frac{\dd^\kappa}{\dd x^\kappa}
\left\{\sum_i\binom{\kappa}{\kappa}
\frac{\p f}{\p u_i^{(\kappa)}}\delta u_i\right\}.
\end{aligned}
\tag{6}
\end{equation}
Соответствующее тождество имеет место и для \(n\)-кратного интеграла; \(A\), в частности, содержит \(\delta u\) вплоть до \((\kappa-1)\)-й производной. То, что формулами (4), (5), (6) действительно определены лагранжевы выражения \(\psi_i\), следует из того, что комбинациями правых частей устранены все высшие производные от \(\delta u\), тогда как выполнено соотношение (2), к которому частное интегрирование приводит однозначно.

В дальнейшем речь идет о двух теоремах:

I. Если интеграл \(I\) инвариантен относительно \(\G_\rho\), то \(\rho\) линейно независимых комбинаций лагранжевых выражений становятся дивергенциями; обратно, отсюда следует инвариантность \(I\) относительно \(\G_\rho\). Теорема остается верной и в предельном случае бесконечного числа параметров.

II. Если интеграл \(I\) инвариантен относительно \(\G_{\infty\rho}\), в которой произвольные функции входят вплоть до \(\sigma\)-й производной, то существуют \(\rho\) тождественных соотношений между лагранжевыми выражениями и их производными вплоть до \(\sigma\)-го порядка; и здесь также справедливо обратное.\footnote{О некоторых тривиальных исключительных случаях см. \S{}2, второе примечание.}

Для смешанных групп справедливы утверждения обеих теорем; следовательно, появляются как зависимости, так и независимые от них дивергентные соотношения.

Если от этих тождеств перейти к соответствующей вариационной задаче, то есть положить \(\psi=0\),\footnote{Несколько более общо можно также положить \(\psi_i=T_i\), ср. \S{}3, первое примечание.} то теорема I в одномерном случае, где дивергенция переходит в полный дифференциал, утверждает существование \(\rho\) первых интегралов, между которыми, однако, могут существовать нелинейные зависимости;\footnote{См. конец \S{}3.} в многомерном случае получаются дивергентные уравнения, которые в последнее время часто называют «законами сохранения»; теорема II утверждает, что \(\rho\) уравнений Лагранжа являются следствием остальных.

Простейший пример к теореме II, без обратного утверждения, дает параметрическое представление Вейерштрасса; здесь интеграл при однородности первого порядка, как известно, инвариантен, если независимую переменную \(x\) заменить произвольной функцией от \(x\), оставляя \(u\) неизменными \((y=p(x);\ v_i(y)=u_i(x))\). Значит, входит одна произвольная функция, но без производных; этому соответствует известное линейное соотношение между самими лагранжевыми выражениями:
\[
\sum_i \psi_i\,\frac{\dd u_i}{\dd x}=0.
\]
Другой пример дает «общая теория относительности» физиков; здесь речь идет о группе всех преобразований \(x\): \(y_i=p_i(x)\), тогда как \(u\) (обозначенные как \(g_{\mu\nu}\) и \(q\)) подчиняются индуцированным этим преобразованиям коэффициентов квадратичной и линейной дифференциальной формы, то есть преобразованиям, содержащим первые производные произвольных функций \(p(x)\). Этому соответствуют известные \(n\) зависимостей между лагранжевыми выражениями и их первыми производными.\footnote{См., например, изложение Клейна.}

Если, в частности, специализировать группу так, что в преобразованиях не допускаются производные \(u(x)\), а преобразованные независимые величины зависят только от \(x\), а не от \(u\), то из инвариантности \(I\) следует (как будет показано в \S{}5) относительная инвариантность \(\sum_i\psi_i\delta u_i\),\footnote{То есть \(\sum_i\psi_i\delta u_i\) при преобразовании приобретает множитель.} а также дивергенций, встречающихся в теореме I, как только параметры подчинены подходящим преобразованиям. Отсюда еще следует, что указанные выше первые интегралы также допускают группу. Для теоремы II так же получается относительная инвариантность левых частей зависимостей, собранных с помощью произвольных функций; и как следствие еще функция, дивергенция которой тождественно исчезает и которая допускает группу, та самая, что в теории относительности физиков посредствует связь между зависимостями и законом энергии.\footnote{См. вторую заметку Клейна.} Наконец, теорема II дает в группово-теоретической форме доказательство связанного с этим утверждения Гильберта об отказе собственных законов энергии при «общей относительности». С этими дополнительными замечаниями теорема I содержит все известные в механике и т. д. теоремы о первых интегралах, тогда как теорему II можно назвать максимально возможным группово-теоретическим обобщением «общей теории относительности».

\section*{\S{}2. Дивергентные соотношения и зависимости}

Пусть \(\G\) есть конечная или бесконечная непрерывная группа; тогда всегда можно добиться того, чтобы тождественному преобразованию соответствовали нулевые значения параметров \(\eps\), соответственно произвольных функций \(p(x)\).\footnote{См., например, Ли: \emph{Grundlagen}, с. 331. Если речь идет о произвольных функциях, специальные значения \(a^\sigma\) параметров следует заменить фиксированными функциями \(p^\sigma,\ \p p^\sigma/\p x,\ldots\); соответственно значения \(a^\sigma+\eps\) заменяются на \(p^\sigma+p(x),\ \p p^\sigma/\p x+\p p/\p x\) и т. д.} Наиболее общее преобразование поэтому имеет вид:
\[
y_i=A_i\!\left(x,u,\frac{\p u}{\p x},\ldots\right)=x_i+\D x_i+\cdots,\qquad
v_i(y)=B_i\!\left(x,u,\frac{\p u}{\p x},\ldots\right)=u_i+\D u_i+\cdots,
\]
где \(\D x_i,\D u_i\) обозначают члены наименьшей размерности относительно \(\eps\), соответственно относительно \(p(x)\) и ее производных; причем они будут считаться линейными относительно этих величин. Как выяснится позднее, это не ограничивает общности.

Пусть теперь интеграл \(I\) является инвариантом относительно \(\G\), то есть выполнено соотношение (1). В частности, тогда \(I\) инвариантен и относительно содержащегося в \(\G\) инфинитезимального преобразования:
\[
y_i=x_i+\D x_i,\qquad v_i(y)=u_i+\D u_i;
\]
для него соотношение (1) переходит в:
\begin{equation}
\begin{aligned}
0=\D I
&=\int\!\cdots\!\int f\!\left(y,v(y),\frac{\p v}{\p y},\ldots\right)\dd y\\
&\qquad-\int\!\cdots\!\int f\!\left(x,u(x),\frac{\p u}{\p x},\ldots\right)\dd x.
\end{aligned}
\tag{7}
\end{equation}
Первый интеграл следует брать по области \(x+\D x\), соответствующей \(x\)-области. Но это интегрирование можно также превратить в интегрирование по \(x\)-области с помощью преобразования, справедливого для инфинитезимальных \(\D x\):
\begin{equation}
\begin{aligned}
\int\!\cdots\!\int f\!\left(y,v(y),\frac{\p v}{\p y},\ldots\right)\dd y
&=\int\!\cdots\!\int f\!\left(x,v(x),\frac{\p v}{\p x},\ldots\right)\dd x\\
&\quad+\int\!\cdots\!\int \Div(f\cdot \D x)\,\dd x.
\end{aligned}
\tag{8}
\end{equation}
Следовательно, если вместо инфинитезимального преобразования \(\D u\) ввести вариацию
\begin{equation}
\bd u_i=v_i(x)-u_i(x)=\D u_i-\sum_\lambda\frac{\p u_i}{\p x_\lambda}\D x_\lambda,
\tag{9}
\end{equation}
то (7) и (8) переходят в:
\begin{equation}
0=\int\!\cdots\!\int\{\bd f+\Div(f\cdot\D x)\}\dd x.
\tag{10}
\end{equation}
Так как соотношение (10) выполнено при интегрировании по всякой произвольной области, интегранд должен тождественно исчезать; следовательно, дифференциальные уравнения Ли для инвариантности \(I\) переходят в соотношение:
\begin{equation}
\bd f+\Div(f\cdot\D x)=0.
\tag{11}
\end{equation}
Если здесь по (3) выразить \(\bd f\) через лагранжевы выражения, получается:
\begin{equation}
\sum_i\psi_i\bd u_i=\Div B,\qquad (B=A-f\cdot\D x),
\tag{12}
\end{equation}
и это соотношение, следовательно, для каждого инвариантного интеграла \(I\) является тождеством по всем встречающимся аргументам; это искомая форма дифференциальных уравнений Ли для \(I\).\footnote{(12) переходит в \(0=0\) в тривиальном случае, который может возникнуть только тогда, когда \(\D x,\D u\) также зависят от производных \(u\), именно когда \(\Div(f\cdot\D x)=0,\ \bd u=0\); такие инфинитезимальные преобразования всегда нужно отделять от групп, и при формулировке теорем считать только число остальных параметров или произвольных функций. Образуют ли остальные инфинитезимальные преобразования еще всегда группу, должно остаться открытым.}

Пусть сначала \(\G\) считается конечной непрерывной группой \(\G_\rho\). Так как, по предположению, \(\D u\) и \(\D x\) линейны по параметрам \(\eps_1,\ldots,\eps_\rho\), то по (9) то же верно для \(\bd u\) и ее производных; следовательно, \(A\) и \(B\) линейны по \(\eps\). Поэтому, если положить
\[
B=B^{(1)}\eps_1+\cdots+B^{(\rho)}\eps_\rho,\qquad
\bd u=\bd u^{(1)}\eps_1+\cdots+\bd u^{(\rho)}\eps_\rho,
\]
где \(\bd u^{(1)},\ldots\) являются функциями от \(x,u,\p u/\p x,\ldots\), то из (12) следуют искомые дивергентные соотношения:
\begin{equation}
\sum_i\psi_i\bd u_i^{(1)}=\Div B^{(1)},\ldots,\qquad
\sum_i\psi_i\bd u_i^{(\rho)}=\Div B^{(\rho)}.
\tag{13}
\end{equation}
Тем самым \(\rho\) линейно независимых комбинаций лагранжевых выражений становятся дивергенциями; линейная независимость следует из того, что по (9) из \(\bd u=0,\ \D x=0\) следовало бы и \(\D u=0,\ \D x=0\), то есть зависимость между инфинитезимальными преобразованиями. Но такая зависимость, по предположению, не выполняется ни при каком значении параметра; иначе \(\G_\rho\), вновь возникающая из инфинитезимальных преобразований интегрированием, зависела бы менее чем от \(\rho\) существенных параметров. Дальнейшая возможность \(\bd u=0,\ \Div(f\D x)=0\) была исключена. Эти выводы остаются верными и в предельном случае бесконечного числа параметров.

Пусть теперь \(\G\) есть бесконечная непрерывная группа \(\G_{\infty\rho}\). Тогда \(\bd u\) и ее производные, а значит и \(B\), снова линейны по произвольным функциям \(p(x)\) и их производным;\footnote{То, что не является ограничением считать \(p\) свободными от \(u,\p u/\p x,\ldots\), показывает обратное утверждение.} пусть после подстановки значений \(\bd u\), еще независимо от (12), имеем:
\[
\begin{aligned}
\sum_i\psi_i\bd u_i
={}&\sum_{\lambda,i}\psi_i\left\{
a_i^{(\lambda)}(x,u,\ldots)p^{(\lambda)}(x)\right.\\
&\left.
+b_i^{(\lambda)}(x,u,\ldots)\frac{\p p^{(\lambda)}}{\p x}
+\cdots+c_i^{(\lambda)}(x,u,\ldots)
\frac{\p^\sigma p^{(\lambda)}}{\p x^\sigma}\right\}.
\end{aligned}
\]
Теперь, аналогично формуле частного интегрирования, по тождеству
\[
\varphi(x,u,\ldots)\frac{\p^\tau p(x)}{\p x^\tau}
=(-1)^\tau\frac{\p^\tau\varphi}{\p x^\tau}\,p(x)
\quad \bmod \text{дивергенций}
\]
производные от \(p\) можно заменить самой \(p\) и дивергенциями, линейными по \(p\) и ее производным; следовательно:
\begin{equation}
\begin{aligned}
\sum_i\psi_i\bd u_i
={}&\sum_\lambda\left\{
a_i^{(\lambda)}\psi_i-\frac{\p}{\p x}\bigl(b_i^{(\lambda)}\psi_i\bigr)
+\cdots+(-1)^\sigma
\frac{\p^\sigma}{\p x^\sigma}\bigl(c_i^{(\lambda)}\psi_i\bigr)
\right\}p^{(\lambda)}+\Div\Gamma,
\end{aligned}
\tag{14}
\end{equation}
с суммированием по \(i\) в фигурной скобке. Вместе с (12) это дает:
\begin{equation}
\sum_\lambda\left\{
a_i^{(\lambda)}\psi_i-\frac{\p}{\p x}\bigl(b_i^{(\lambda)}\psi_i\bigr)
+\cdots+(-1)^\sigma
\frac{\p^\sigma}{\p x^\sigma}\bigl(c_i^{(\lambda)}\psi_i\bigr)
\right\}p^{(\lambda)}
=\Div(B-\Gamma).
\tag{15}
\end{equation}
Я теперь образую \(n\)-кратный интеграл от (15) по какой-нибудь области и выбираю \(p(x)\) так, чтобы они вместе со всеми производными, встречающимися в \(B-\Gamma\), исчезали на границе. Так как интеграл от дивергенции сводится к граничному интегралу, исчезает и интеграл от левой части (15) для произвольных \(p(x)\), исчезающих только на границе вместе с достаточно многими производными; отсюда по известным рассуждениям следует исчезновение интегранда для каждой \(p(x)\), то есть \(\rho\) соотношений:
\begin{equation}
\sum_i\left\{
a_i^{(\lambda)}\psi_i-\frac{\p}{\p x}\bigl(b_i^{(\lambda)}\psi_i\bigr)
+\cdots+(-1)^\sigma
\frac{\p^\sigma}{\p x^\sigma}\bigl(c_i^{(\lambda)}\psi_i\bigr)
\right\}=0,\qquad(\lambda=1,2,\ldots,\rho).
\tag{16}
\end{equation}
Это и есть искомые зависимости между лагранжевыми выражениями и их производными при инвариантности \(I\) относительно \(\G_{\infty\rho}\); линейная независимость видна так же, как выше, поскольку обратное утверждение возвращает к (12), и потому что от инфинитезимальных преобразований снова можно заключить к конечным, как подробнее изложено в \S{}4. Следовательно, при \(\G_{\infty\rho}\) уже в инфинитезимальных преобразованиях всегда появляются \(\rho\) произвольных преобразований. Из (15) и (16) еще следует \(\Div(B-\Gamma)=0\).

Если соответственно «смешанной группе» положить \(\D x\) и \(\D u\) линейными по \(\eps\) и \(p(x)\), то, полагая один раз \(p(x)\) нулем, а другой раз \(\eps\) нулем, видно, что существуют как дивергентные соотношения (13), так и зависимости (16).

\section*{\S{}3. Обратное утверждение в случае конечной группы}

Чтобы показать обратное утверждение, нужно сначала в существенном пройти предыдущие рассуждения в обратном порядке. Из существования (13) умножением на \(\eps\) и сложением следует существование (12); а посредством тождества (3) отсюда следует соотношение
\[
\bd f+\Div(A-B)=0.
\]
Если положить \(\D x=\frac{1}{f}(A-B)\), то тем самым мы пришли к (11); отсюда интегрированием наконец следует (7): \(\D I=0\), то есть инвариантность \(I\) относительно инфинитезимального преобразования, определенного \(\D x,\D u\), причем \(\D u\) определяется из \(\D x\) и \(\bd u\) по (9), а \(\D x\) и \(\D u\) становятся линейными по параметрам. Но \(\D I=0\) известным образом влечет инвариантность \(I\) относительно конечных преобразований, возникающих интегрированием одновременной системы:
\begin{equation}
\frac{\dd x_i}{\dd t}=\D x_i,\qquad
\frac{\dd u_i}{\dd t}=\D u_i,\qquad
\left\{\begin{array}{l}
x_i=y_i,\\
u_i=v_i
\end{array}\right\}
\quad\text{при }t=0.
\tag{17}
\end{equation}
Эти конечные преобразования содержат \(\rho\) параметров \(a_1,\ldots,a_\rho\), а именно комбинации \(t\eps_1,\ldots,t\eps_\rho\). Из предположения, что должно быть \(\rho\) и только \(\rho\) линейно независимых дивергентных соотношений (13), далее следует, что конечные преобразования, если они не содержат производных \(\p u/\p x\), всегда образуют группу. В противном случае по крайней мере одно инфинитезимальное преобразование, возникающее через процесс скобок Ли, не было бы линейной комбинацией остальных \(\rho\); а так как \(I\) допускает и это преобразование, существовало бы более чем \(\rho\) линейно независимых дивергентных соотношений; или же это инфинитезимальное преобразование имело бы специальный вид \(\bd u=0,\ \Div(f\D x)=0\). Но тогда \(\D x\) или \(\D u\) зависели бы от производных вопреки предположению. Может ли этот случай наступить, когда производные входят в \(\D x\) или \(\D u\), должно остаться открытым; тогда к определенному выше \(\D x\) нужно еще добавить все функции \(\D x\), для которых \(\Div(f\D x)=0\), чтобы снова получить групповое свойство; добавляющиеся этим параметры по соглашению не должны считаться. Этим обратное утверждение доказано.

Из этого обратного утверждения еще следует, что \(\D x\) и \(\D u\) действительно можно считать линейными по параметрам. Ведь если бы \(\D u\) и \(\D x\) были формами более высокой степени по \(\eps\), то из-за линейной независимости степенных произведений \(\eps\) следовали бы вполне соответствующие соотношения (13), только в большем числе; из них по обратному утверждению следует инвариантность \(I\) относительно группы, инфинитезимальные преобразования которой содержат параметры линейно. Если эта группа должна содержать ровно \(\rho\) параметров, то между дивергентными соотношениями, первоначально полученными из членов высших степеней по \(\eps\), должны существовать линейные зависимости.

Заметим еще, что в случае, когда \(\D x\) и \(\D u\) содержат также производные \(u\), конечные преобразования могут зависеть от бесконечного числа производных \(u\); ибо интегрирование (17) в этом случае при определении
\[
\frac{\dd^2x_i}{\dd t^2},\quad \frac{\dd^2u_i}{\dd t^2}
\]
приводит к
\[
\D\!\left(\frac{\p u}{\p x_\lambda}\right)
=\frac{\p\,\D u}{\p x_\lambda}
-\sum_\nu\frac{\p u}{\p x_\nu}
\frac{\p\,\D x_\nu}{\p x_\lambda},
\]
так что число производных \(u\) вообще возрастает на каждом шаге. Примером служит, например:
\[
f=\frac12 u'^2,\qquad
\psi=-u'',\qquad
\psi\cdot x=\frac{\dd}{\dd x}\{(u-u'x)\},\qquad
\bd u=x\eps,
\]
\[
\D x=-\frac{2u}{u'}\,\eps,\qquad
\D u=\left(x-\frac{2u}{u'}\right)\eps.
\]
Так как лагранжевы выражения дивергенции тождественно исчезают, обратное утверждение наконец показывает еще следующее: если \(I\) допускает \(\G_\rho\), то всякий интеграл, отличающийся от \(I\) только граничным интегралом, то есть интегралом от дивергенции, также допускает \(\G_\rho\) с теми же \(\bd u\), инфинитезимальные преобразования которой вообще будут содержать производные \(u\). Так, в соответствии с приведенным выше примером,
\[
f^*=\frac12\left\{u'^2-\frac{\dd}{\dd x}\left(\frac{u^2}{x}\right)\right\}
\]
допускает инфинитезимальное преобразование \(\D u=x\eps,\ \D x=0\), тогда как в инфинитезимальных преобразованиях, соответствующих \(f\), появляются производные \(u\).

Если перейти к вариационной задаче, то есть положить \(\psi_i=0\),\footnote{\(\psi_i=0\), или несколько более общо \(\psi_i=T_i\), где \(T_i\) являются вновь добавленными функциями, в физике называются «полевыми уравнениями». В случае \(\psi_i=T_i\) тождества (13) переходят в уравнения \(\Div B^{(\lambda)}=\sum_i T_i\bd u_i^{(\lambda)}\), которые в физике также называют законами сохранения.} то (13) переходит в уравнения:
\[
\Div B^{(1)}=0,\ldots,\qquad \Div B^{(\rho)}=0,
\]
которые часто называют «законами сохранения». В одномерном случае отсюда следует:
\[
B^{(1)}=\mathrm{const.},\ldots,\qquad B^{(\rho)}=\mathrm{const.};
\]
при этом \(B\) содержат самое большее \((2\kappa-1)\)-е производные \(u\) (по (6)), если \(\D u\) и \(\D x\) не содержат более высоких производных, чем входящие в \(f\) \(\kappa\)-е. Так как в \(\psi\) вообще входят \(2\kappa\)-е производные,\footnote{Если \(f\) нелинейна по \(\kappa\)-м производным.} мы имеем существование \(\rho\) первых интегралов. То, что между ними могут быть нелинейные зависимости, снова показывает приведенное выше \(f\). Линейно независимым \(\D u=\eps_1,\ \D x=\eps_2\) соответствуют линейно независимые соотношения:
\[
u''=\frac{\dd}{\dd x}u',\qquad
u''u'=\frac12\frac{\dd}{\dd x}(u')^2;
\]
тогда как между первыми интегралами \(u'=\mathrm{const.},\ u'^2=\mathrm{const.}\) существует нелинейная зависимость. Здесь речь идет об элементарном случае, когда \(\D u,\D x\) не содержат производных \(u\).\footnote{Иначе имеются еще \(u''u'^{\lambda-1}=\mathrm{const.}\) для каждого \(\lambda\), соответственно
\[
u''(u')^{\lambda-1}=\frac1\lambda\frac{\dd}{\dd x}(u')^\lambda.
\]}

\section*{\S{}4. Обратное утверждение в случае бесконечной группы}

Сначала покажем, что предположение линейности \(\D x\) и \(\D u\) не представляет ограничения; здесь это следует уже без обратного утверждения из того факта, что \(\G_{\infty\rho}\) формально зависит от \(\rho\) и только \(\rho\) произвольных функций. А именно оказывается, что в нелинейном случае при композиции преобразований, когда члены наименьшего порядка складываются, число произвольных функций увеличилось бы. В самом деле, пусть, например,
\[
\begin{aligned}
y&=A\!\left(x,u,\frac{\p u}{\p x},\ldots;p\right)\\
&=x+\sum a(x,u,\ldots)p^\nu
+b(x,u,\ldots)p^{\nu-1}\frac{\p p}{\p x}
+c\,p^{\nu-2}\left(\frac{\p p}{\p x}\right)^2+\cdots
+d\left(\frac{\p p}{\p x}\right)^\nu+\cdots,
\end{aligned}
\]
\[
p^\nu=(p^{(1)})^{\nu_1}\cdots(p^{(\rho)})^{\nu_\rho},
\]
и соответственно
\[
v=B\!\left(x,u,\frac{\p u}{\p x},\ldots;p\right).
\]
Тогда при композиции с
\[
z=A\!\left(y,v,\frac{\p v}{\p y},\ldots;q\right)
\]
для членов наименьшего порядка получается:
\[
z=x+\sum a(p^\nu+q^\nu)+
b\left\{p^{\nu-1}\frac{\p p}{\p x}+q^{\nu-1}\frac{\p q}{\p x}\right\}
+c\left\{p^{\nu-2}\left(\frac{\p p}{\p x}\right)^2
+q^{\nu-2}\left(\frac{\p q}{\p x}\right)^2\right\}+\cdots.
\]
Если здесь какой-либо отличный от \(a\) и \(b\) коэффициент не равен нулю, то есть действительно встречается член
\[
p^{\nu-\sigma}\left(\frac{\p p}{\p x}\right)^\sigma
+q^{\nu-\sigma}\left(\frac{\p q}{\p x}\right)^\sigma
\]
при \(\sigma>1\), то его нельзя записать как дифференциальное частное одной функции или степенное произведение такого частного; следовательно, число произвольных функций возросло вопреки предположению. Если все коэффициенты, отличные от \(a\) и \(b\), исчезают, то в зависимости от значений показателей \(\nu_1,\ldots,\nu_\rho\) второй член становится дифференциальным частным первого (как, например, всегда для \(\G_{\infty1}\)), так что фактически наступает линейность; либо число произвольных функций увеличивается и здесь. Поэтому инфинитезимальные преобразования вследствие линейности по \(p(x)\) удовлетворяют системе линейных частных дифференциальных уравнений; и так как групповое свойство выполнено, они образуют «бесконечную группу инфинитезимальных преобразований» по определению Ли (\emph{Grundlagen}, \S{}10).

Обратное утверждение получается теперь подобно случаю конечной группы. Существование зависимостей (16) умножением на \(p^{(\lambda)}(x)\) и сложением посредством тождественного преобразования (14) приводит к
\[
\sum_i\psi_i\bd u_i=\Div\Gamma;
\]
и отсюда, как в \S{}3, следует определение \(\D x\) и \(\D u\) и инвариантность \(I\) относительно этих инфинитезимальных преобразований, которые действительно линейно зависят от \(\rho\) произвольных функций и их производных вплоть до \(\sigma\)-го порядка. То, что эти инфинитезимальные преобразования, если они не содержат производных \(\p u/\p x,\ldots\), наверняка образуют группу, следует, как в \S{}3, из того, что иначе при композиции появилось бы больше произвольных функций, тогда как по предположению должно быть только \(\rho\) зависимостей (16); значит, они образуют «бесконечную группу инфинитезимальных преобразований». Такая группа (\emph{Grundlagen}, теорема VII, с. 391) состоит из наиболее общих инфинитезимальных преобразований некоторой тем самым определенной в смысле Ли «бесконечной группы \(\G\) конечных преобразований». Каждое конечное преобразование при этом порождается из инфинитезимальных (\emph{Grundlagen}, \S{}7),\footnote{Отсюда, в частности, следует, что группа \(\G\), порожденная инфинитезимальными преобразованиями \(\D x,\D u\) группы \(\G_{\infty\rho}\), снова приводит к \(\G_{\infty\rho}\). Ведь \(\G_{\infty\rho}\) не содержит инфинитезимальных преобразований, зависящих от произвольных функций и отличных от \(\D x,\D u\), и не может содержать независимых от них, зависящих от параметров, иначе она была бы смешанной группой. Но, согласно сказанному выше, конечные преобразования определены инфинитезимальными.} то есть возникает интегрированием одновременной системы:
\[
\frac{\dd x_i}{\dd t}=\D x_i,\qquad
\frac{\dd u_i}{\dd t}=\D u_i,\qquad
\left\{\begin{array}{l}
x_i=y_i,\\
u_i=v
\end{array}\right\}
\quad\text{при }t=0,
\]
причем может быть необходимо считать произвольные \(p(x)\) еще зависящими от \(t\). Следовательно, \(\G\) действительно зависит от \(\rho\) произвольных функций; если, в частности, достаточно считать \(p(x)\) независимыми от \(t\), то эта зависимость является аналитической по произвольным функциям \(q(x)=t\cdot p(x)\).\footnote{Вопрос, всегда ли наступает этот последний случай, был поставлен Ли в другой формулировке (\emph{Grundlagen}, \S{}7 и конец \S{}13).} Если появляются производные \(\p u/\p x,\ldots\), то может быть необходимо еще добавить инфинитезимальные преобразования \(\bd u=0,\ \Div(f\D x)=0\), прежде чем делать те же выводы.

В связи с примером Ли (\emph{Grundlagen}, \S{}7) укажем еще довольно общий случай, где можно дойти до явных формул, одновременно показывающих, что производные произвольных функций входят вплоть до \(\sigma\)-го порядка; здесь обратное утверждение тем самым полно. Это такие группы инфинитезимальных преобразований, которым соответствует группа всех преобразований \(x\) и тем самым «индуцированных» преобразований \(u\), то есть таких преобразований \(u\), при которых \(\D u\), а следовательно и \(u\), зависят только от произвольных функций, входящих в \(\D x\); при этом еще предполагается, что производные \(\p u/\p x,\ldots\) в \(\D u\) не входят. Итак:
\[
\D x_i=p^{(i)}(x),\qquad
\D u_i=
\sum_{\lambda=1}^{n}\left\{
a_i^{(\lambda)}(x,u)p^{(\lambda)}
+b_i^{(\lambda)}\frac{\p p^{(\lambda)}}{\p x}
+\cdots+
c_i^{(\lambda)}\frac{\p^\sigma p^{(\lambda)}}{\p x^\sigma}
\right\}.
\]
Так как инфинитезимальное преобразование \(\D x=p(x)\) порождает всякое преобразование \(x=y+g(y)\) с произвольным \(g(y)\), можно, в частности, определить \(p(x)\) зависящей от \(t\) так, чтобы порождалась однопараметрическая группа:
\begin{equation}
x_i=y_i+t\,g_i(y),
\tag{18}
\end{equation}
которая при \(t=0\) переходит в тождество, а при \(t=1\) в искомое \(x=y+g(y)\). Дифференцированием (18) действительно получается:
\begin{equation}
\frac{\dd x_i}{\dd t}=g_i(y)=p^{(i)}(x,t),
\tag{19}
\end{equation}
где \(p(x,t)\) определяется из \(g(y)\) обращением (18); и обратно, (18) возникает из (19) посредством побочного условия \(x_i=y_i\) при \(t=0\), которым интеграл однозначно фиксируется. Посредством (18) в \(\D u\) можно заменить \(x\) интеграционными постоянными \(y\) и \(t\); при этом \(g(y)\) входят точно вплоть до \(\sigma\)-й производной, поскольку в
\[
\frac{\p p}{\p x}=\sum_k\frac{\p g}{\p y_k}\frac{\p y_k}{\p x}
\]
\(\p y/\p x\) выражаются через \(\p x/\p y\), а вообще \(\p^\sigma p/\p x^\sigma\) заменяется своим значением в
\[
\frac{\p g}{\p y},\ldots,\frac{\p^\sigma g}{\p y^\sigma},\ldots,\frac{\p^\sigma x}{\p y^\sigma}.
\]
Для определения \(u\) получается поэтому система уравнений:
\[
\frac{\dd u_i}{\dd t}
=
F_i\!\left(g(y),\frac{\p g}{\p y},\ldots,\frac{\p^\sigma g}{\p y^\sigma},u,t\right),
\qquad (u_i=v_i\ \text{при }t=0),
\]
в которой переменными являются только \(t\) и \(u\), а \(g(y),\ldots\) принадлежат области коэффициентов, так что интегрирование дает:
\[
u_i=v_i+
B_i\!\left(v,g(y),\frac{\p g}{\p y},\ldots,\frac{\p^\sigma g}{\p y^\sigma},t\right)_{t=1},
\]
то есть преобразования, зависящие ровно от \(\sigma\) производных произвольных функций. Тождество содержится в них по (18) при \(g(y)=0\); а групповое свойство следует из того, что указанная процедура дает всякое преобразование \(x=y+g(y)\), благодаря чему индуцированное преобразование \(u\) однозначно фиксируется и группа \(\G\) исчерпывается.

Из обратного утверждения попутно следует еще, что не является ограничением считать произвольные функции зависящими только от \(x\), а не от \(u,\p u/\p x,\ldots\). В последнем случае в тождественном преобразовании (14), а значит и в (15), кроме \(p^{(\lambda)}\), появились бы еще \(\p p^{(\lambda)}/\p u,\ \p p^{(\lambda)}/\p(\p u/\p x),\ldots\). Если теперь последовательно брать \(p^{(\lambda)}\) как функции нулевой, первой, \(\ldots\), степени по \(u,\p u/\p x,\ldots\), с произвольными функциями от \(x\) в качестве коэффициентов, снова получаются зависимости (16), только в большем числе; но они по приведенному выше обратному утверждению путем объединения с произвольными функциями, зависящими только от \(x\), возвращают к прежнему случаю. Так же показывают, что одновременному появлению зависимостей и независимых от них дивергентных соотношений соответствуют смешанные группы.\footnote{Как в \S{}3, из обратного утверждения и здесь следует, что наряду с \(I\) всякий интеграл \(I^*\), отличающийся на интеграл от дивергенции, также допускает бесконечную группу с теми же \(\bd u\), причем \(\D x\) и \(\D u\) вообще будут содержать производные \(u\). Такой интеграл \(I^*\) Эйнштейн ввел в общую теорию относительности, чтобы получить более простую форму законов энергии; я указываю инфинитезимальные преобразования, которые допускает этот \(I^*\), точно придерживаясь обозначений второй заметки Клейна. Интеграл \(I=\int\!\cdots\!\int K\,\dd\omega=\int\!\cdots\!\int \mathfrak{R}\,\dd S\) допускает группу всех преобразований \(w\) и индуцированную ею группу для \(g_{\mu\nu}\); этому соответствуют зависимости ((30) у Клейна):
\[
\sum \mathfrak{R}_{\mu\nu}g_{\mu\nu}^{\alpha\beta}
+2\sum\frac{\p g_{\mu\nu}^{\alpha\beta}\mathfrak{R}_{\mu\nu}}{\p w^\sigma}=0.
\]
Теперь \(I^*=\int\!\cdots\!\int \mathfrak{R}^*\,\dd S\), где \(\mathfrak{R}^*=\mathfrak{R}+\Div\), и следовательно \(\mathfrak{R}_{\mu\nu}^*=\mathfrak{R}_{\mu\nu}\), где \(\mathfrak{R}_{\mu\nu}^*,\mathfrak{R}_{\mu\nu}\) каждый раз обозначают лагранжевы выражения. Указанные зависимости поэтому являются также зависимостями для \(\mathfrak{R}_{\mu\nu}^*\); и после умножения на \(p^\tau\) и сложения обратными преобразованиями дифференцирования произведения получается:
\[
\sum \mathfrak{R}_{\mu\nu}p^{\mu\nu}
+2\Div\!\left(\sum g^{\mu\sigma}\mathfrak{R}_{\mu\tau}p^\tau\right)=0,
\]
\[
\delta\mathfrak{R}^*+\Div\!\left(\sum 2g^{\mu\sigma}\mathfrak{R}_{\mu\tau}p^\tau
-\frac{\p\mathfrak{R}^*}{\p g_{\mu\nu}^{\sigma}}p^{\mu\nu}\right)=0.
\]
Сравнением с дифференциальным уравнением Ли \(\delta\mathfrak{R}^*+\Div(\mathfrak{R}^*\D w)=0\) отсюда следует:
\[
\D w^\sigma=\frac1{\mathfrak{R}^*}\left(\sum 2g^{\mu\sigma}\mathfrak{R}_{\mu\tau}p^\tau
-\frac{\p\mathfrak{R}^*}{\p g_{\mu\nu}^{\sigma}}p^{\mu\nu}\right),
\qquad
\D g^{\mu\nu}=p^{\mu\nu}+\sum g^{\mu\nu}_{\sigma}\D w^\sigma,
\]
как инфинитезимальные преобразования, допускаемые \(I^*\). Эти инфинитезимальные преобразования, следовательно, зависят от первых и вторых производных \(g^{\mu\nu}\) и содержат произвольные \(p\) вплоть до первой производной.}

\section*{\S{}5. Инвариантность отдельных составных частей соотношений}

Если специализировать группу \(\G\) к простейшему, обычно рассматриваемому случаю так, что в преобразованиях не допускаются производные \(u\), а преобразованные независимые переменные зависят только от \(x\), а не от \(u\), то можно делать выводы об инвариантности отдельных составных частей в формулах. Сначала по известным рассуждениям получается инвариантность
\[
\int\!\cdots\!\int\left(\sum_i\psi_i\delta u_i\right)\dd x,
\]
то есть относительная инвариантность \(\sum_i\psi_i\delta u_i\), где под \(\delta\) понимается какая-либо вариация. В самом деле, с одной стороны:
\[
\delta I
=\int\!\cdots\!\int \delta f\!\left(x,u,\frac{\p u}{\p x},\ldots\right)\dd x
=\int\!\cdots\!\int \delta f\!\left(y,v,\frac{\p v}{\p y},\ldots\right)\dd y,
\]
с другой стороны, для исчезающих на границе \(\delta u,\delta(\p u/\p x),\ldots\), которым вследствие линейного однородного преобразования \(\delta u,\delta(\p u/\p x),\ldots\) соответствуют также исчезающие на границе \(\delta v,\delta(\p v/\p y),\ldots\):
\[
\int\!\cdots\!\int \delta f\!\left(x,u,\frac{\p u}{\p x},\ldots\right)\dd x
=
\int\!\cdots\!\int\left(\sum_i\psi_i(u,\ldots)\delta u_i\right)\dd x,
\]
\[
\int\!\cdots\!\int \delta f\!\left(y,v,\frac{\p v}{\p y},\ldots\right)\dd y
=
\int\!\cdots\!\int\left(\sum_i\psi_i(v,\ldots)\delta v_i\right)\dd y;
\]
следовательно, для исчезающих на границе \(\delta u,\delta(\p u/\p x),\ldots\):
\[
\int\!\cdots\!\int\left(\sum_i\psi_i(u,\ldots)\delta u_i\right)\dd x
=
\int\!\cdots\!\int\left(\sum_i\psi_i(v,\ldots)\delta v_i\right)\dd y
\]
\[
=
\int\!\cdots\!\int\left(\sum_i\psi_i(v,\ldots)\delta v_i\right)
\left|\frac{\p y_i}{\p x_k}\right|\dd x.
\]
Если в третьем интеграле выразить \(y,v,\delta v\) через \(x,u,\delta u\) и приравнять его первому, получаем соотношение
\[
\int\!\cdots\!\int\left(\sum_i\chi_i(u,\ldots)\delta u_i\right)\dd x=0
\]
для исчезающей на границе, а в остальном произвольной \(\delta u\); отсюда, как известно, следует исчезновение интегранда для произвольной \(\delta u\). Значит, тождественно по \(\delta u\) справедливо соотношение:
\[
\sum_i\psi_i(u,\ldots)\delta u_i
=
\left|\frac{\p y_i}{\p x_k}\right|
\left(\sum_i\psi_i(v,\ldots)\delta v_i\right),
\]
которое выражает относительную инвариантность \(\sum_i\psi_i\delta u_i\) и, следовательно, инвариантность
\[
\int\!\cdots\!\int\left(\sum_i\psi_i\delta u_i\right)\dd x.
\]
\footnote{Эти выводы не работают, если \(y\) также зависит от \(u\), потому что тогда \(\delta f(y,v,\p v/\p y,\ldots)\) содержит и члены \(\sum(\p f/\p y)\delta y\), так что дивергентное преобразование не приводит к лагранжевым выражениям; так же если допустить производные \(u\), поскольку тогда \(\delta v\) становятся линейными комбинациями \(\delta u,\delta(\p u/\p x),\ldots\) и лишь после еще одного дивергентного преобразования приводят к тождеству \(\int\!\cdots\!\int(\sum\chi_i(u,\ldots)\delta u_i)\dd x=0\), так что справа снова не появляются лагранжевы выражения. Вопрос, можно ли из инвариантности \(\int\!\cdots\!\int(\sum\psi_i\delta u_i)\dd x\) уже заключить о существовании дивергентных соотношений, по обратному утверждению равносилен вопросу, можно ли отсюда заключить об инвариантности \(I\) относительно группы, которая не обязательно приводит к тем же \(\D u,\D x\), но приводит к тем же \(\bd u\). В специальном случае простого интеграла и только первых производных в \(f\), для конечной группы из инвариантности лагранжевых выражений можно заключить о существовании первых интегралов (см., например, Энгель, Gött. Nachr. 1916, с. 270).}

Чтобы применить это к выведенным дивергентным соотношениям и зависимостям, нужно сначала доказать, что \(\bd u\), полученное из \(\D u,\D x\), действительно удовлетворяет законам преобразования для вариации \(\delta u\), если только параметры, соответственно произвольные функции в \(\bd v\), определены так, как они соответствуют подобной группе инфинитезимальных преобразований в \(y,v\). Обозначим через \(T_q\) преобразование, переводящее \(x,u\) в \(y,v\), через \(T_p\) --- инфинитезимальное в \(x,u\); тогда подобное ему в \(y,v\) задается \(T_r=T_qT_pT_q^{-1}\), где параметры, соответственно произвольные функции \(r\), определяются из \(p\) и \(q\). В формулах это выражается так:
\[
T_p:\ \xi=x+\D x(x,p),\qquad u^*=u+\D u(x,u,p);
\]
\[
T_q:\ y=A(x,q),\qquad v=B(x,u,q);
\]
\[
T_qT_p:\ \eta=A(x+\D x(x,p),q),\quad
v^*=B(x+\D x(p),\,u+\D u(p),q).
\]
Но отсюда возникает \(T_r=T_qT_pT_q^{-1}\), то есть
\[
\eta=y+\D y(r),\qquad v^*=v+\D v(r),
\]
если посредством обращения \(T_q\) рассматривать \(x\) как функции от \(y\) и учитывать только инфинитезимальные члены; следовательно, имеем тождество:
\begin{equation}
\eta=y+\D y(r)=y+\sum\frac{\p A(x,q)}{\p x}\D x(p),
\tag{20a}
\end{equation}
\begin{equation}
v^*=v+\D v(r)=v+\sum\frac{\p B(x,u,q)}{\p x}\D x(p)
+\sum\frac{\p B(x,u,q)}{\p u}\D u(p).
\tag{20b}
\end{equation}
Если в этих формулах заменить \(\xi=x+\D x\) на \(\xi-\D \xi\), так что \(\xi\) снова переходит в \(x\), а значит \(\D x\) исчезает, то по первой формуле (20) и \(\eta\) снова переходит в \(y=\eta-\D\eta\); если этой подстановкой \(\D u(p)\) переходит в \(\bd u(p)\), то и \(\D v(r)\) переходит в \(\bd v(r)\), и вторая формула (20) дает:
\[
v+\bd v(y,v,\ldots,r)=v+\sum\frac{\p B(x,u,q)}{\p u}\,\bd u(p),
\]
\[
\bd v(y,v,\ldots,r)=\sum\frac{\p B}{\p u_\kappa}\,\bd u_\kappa(x,u,p),
\]
так что действительно выполнены формулы преобразования для вариаций, если только \(\bd v\) считать зависящим от параметров, соответственно произвольных функций \(r\).\footnote{Снова видно, что для справедливости выводов нужно считать \(y\) независимым от \(u\) и т. д. В качестве примера можно назвать указанные Клейном \(\delta g^{\mu\nu}\) и \(\delta q_\sigma\), которые удовлетворяют преобразованиям для вариаций, если \(p\) подчинены векторному преобразованию.}

Следовательно, в частности, получается относительная инвариантность \(\sum_i\psi_i\bd u_i\); поэтому также по (12), так как дивергентные соотношения выполнены и в \(y,v\), относительная инвариантность \(\Div B\); и далее по (14) и (13) относительная инвариантность \(\Div\Gamma\) и левых частей зависимостей, собранных с помощью \(p^{(\lambda)}\), причем в преобразованных формулах произвольные \(p(x)\) (соответственно параметры) всегда нужно заменять на \(r\). Отсюда еще получается относительная инвариантность \(\Div(B-\Gamma)\), то есть дивергенции нетождественно исчезающей системы функций \(B-\Gamma\), дивергенция которой тождественно исчезает.

Из относительной инвариантности \(\Div B\) в одномерном случае и при конечной группе можно еще заключить об инвариантности первых интегралов. Соответствующее инфинитезимальному преобразованию преобразование параметров по (20) становится линейным и однородным; а вследствие обратимости всех преобразований и \(\eps\) становятся линейными и однородными по преобразованным параметрам \(\eps^*\). Эта обратимость надежно сохраняется, если положить \(\psi=0\), потому что в (20) не входят производные \(u\). Приравнивая коэффициенты при \(\eps^*\) в
\[
\Div B(x,u,\ldots,\eps)=\frac{\dd y}{\dd x}\,\Div B(y,v,\ldots,\eps^*)
\]
получаем, что и \(\frac{\dd}{\dd y}B^{(\lambda)}(y,v,\ldots)\) являются линейными однородными функциями от \(\frac{\dd}{\dd x}B^{(\lambda)}(x,u,\ldots)\); так что из
\[
\frac{\dd}{\dd x}B^{(\lambda)}(x,u,\ldots)=0
\quad\text{или}\quad B^{(\lambda)}(x,u)=\mathrm{const.}
\]
следует также:
\[
\frac{\dd}{\dd y}B^{(\lambda)}(y,v,\ldots)=0
\quad\text{или}\quad B^{(\lambda)}(y,v)=\mathrm{const.}
\]
Следовательно, \(\rho\) первых интегралов, соответствующих \(\G_\rho\), каждый допускает группу, так что дальнейшее интегрирование также упрощается. Простейший пример здесь состоит в том, что \(f\) свободна от \(x\) или от некоторого \(u\), чему соответствуют инфинитезимальные преобразования \(\D x=\eps,\D u=0\) соответственно \(\D x=0,\D u=\eps\). Тогда \(\bd u=-\eps\,\dd u/\dd x\) соответственно \(\eps\), и так как \(B\) выводится из \(f\) и \(\bd u\) дифференцированием и рациональными соединениями, оно также свободно от \(x\) соответственно \(u\) и допускает соответствующие группы.\footnote{В тех случаях, где уже из инвариантности \(\int(\sum\psi_i\delta u_i)\dd x\) следует существование первых интегралов, эти интегралы не допускают полной группы \(\G_\rho\); например, \(\int(u''\delta u)\dd x\) допускает инфинитезимальное преобразование \(\D x=\eps_2,\D u=\eps_1+x\eps_3\), тогда как первый интеграл \(u-u'x=\mathrm{const.}\), соответствующий \(\D x=0,\D u=x\eps_3\), не допускает двух других инфинитезимальных преобразований, поскольку явно содержит и \(u\), и \(x\). Этому первому интегралу как раз соответствуют инфинитезимальные преобразования для \(f\), содержащие производные. Поэтому видно, что инвариантность \(\int\!\cdots\!\int(\sum\psi_i\delta u_i)\dd x\) во всяком случае дает меньше, чем инвариантность \(I\), что следует отметить по поводу вопроса, поставленного в предыдущем примечании.}

\section*{\S{}6. Утверждение Гильберта}

Из предыдущего наконец получается доказательство утверждения Гильберта о связи отказа собственных законов энергии с «общей относительностью» (первая заметка Клейна, Göttinger Nachr. 1917, ответ, 1-й абзац), причем в обобщенной группово-теоретической форме.

Пусть интеграл \(I\) допускает \(\G_{\infty\rho}\), а \(\G_\sigma\) есть какая-либо конечная группа, возникающая специализацией произвольных функций, то есть подгруппа \(\G_{\infty\rho}\). Бесконечной группе \(\G_{\infty\rho}\) тогда соответствуют зависимости (16), конечной \(\G_\sigma\) --- дивергентные соотношения (13); и наоборот, из существования каких-либо дивергентных соотношений следует инвариантность \(I\) относительно конечной группы, которая тогда и только тогда совпадает с \(\G_\sigma\), когда \(\bd u\) являются линейными комбинациями тех, что получаются из \(\G_\sigma\). Инвариантность относительно \(\G_\sigma\), следовательно, не может вести ни к каким дивергентным соотношениям, отличным от (13). Но так как из существования (16) следует инвариантность \(I\) относительно инфинитезимальных преобразований \(\D u,\D x\) группы \(\G_{\infty\rho}\) при произвольном \(p(x)\), отсюда, в частности, уже следует инвариантность относительно возникающих специализацией инфинитезимальных преобразований \(\G_\sigma\), а значит и относительно \(\G_\sigma\). Поэтому дивергентные соотношения
\[
\sum_i\psi_i\bd u_i^{(\lambda)}=\Div B^{(\lambda)}
\]
должны быть следствиями зависимостей (16), которые можно также записать как
\[
\sum_i\psi_i a_i^{(\lambda)}=\Div\chi^{(\lambda)},
\]
где \(\chi^{(\lambda)}\) являются линейными комбинациями лагранжевых выражений и их производных. Так как \(\psi\) входят линейно и в (13), и в (16), дивергентные соотношения, следовательно, в частности должны быть линейными комбинациями зависимостей (16); поэтому
\[
\Div B^{(\lambda)}=\Div\left(\sum \alpha_\mu\chi^{(\mu)}\right);
\]
а сами \(B^{(\lambda)}\) поэтому линейно составляются из \(\chi\), то есть из лагранжевых выражений и их производных, и из функций, дивергенция которых тождественно исчезает, как, например, появляющиеся в конце \S{}2 \(B-\Gamma\), для которых \(\Div(B-\Gamma)=0\) и где дивергенция одновременно обладает свойством инвариантности. Дивергентные соотношения, в которых \(B^{(\lambda)}\) можно составить из лагранжевых выражений и их производных указанным образом, я буду называть «несобственными», а все остальные --- «собственными».

Обратно, если дивергентные соотношения являются линейными комбинациями зависимостей (16), то есть «несобственными», то инвариантность относительно \(\G_\sigma\) следует из инвариантности относительно \(\G_{\infty\rho}\); \(\G_\sigma\) становится подгруппой \(\G_{\infty\rho}\). Следовательно, соответствующие конечной группе \(\G_\sigma\) дивергентные соотношения становятся несобственными тогда и только тогда, когда \(\G_\sigma\) является подгруппой бесконечной группы, относительно которой \(I\) инвариантен.

Специализацией групп отсюда получается первоначальное утверждение Гильберта. Под «группой сдвигов» будем понимать конечную группу:
\[
y_i=x_i+\eps_i,\qquad v_i(y)=u_i(x);
\]
то есть:
\[
\D x_i=\eps_i,\qquad \D u_i=0,\qquad
\bd u_i=-\sum_\lambda\frac{\p u_i}{\p x_\lambda}\eps_\lambda.
\]
Инвариантность относительно группы сдвигов, как известно, означает, что в
\[
I=\int\!\cdots\!\int f\!\left(x,u,\frac{\p u}{\p x},\ldots\right)\dd x
\]
\(x\) не входят явно в \(f\). Соответствующие \(n\) дивергентных соотношений
\[
\sum_i\psi_i\frac{\p u_i}{\p x_\lambda}=\Div B^{(\lambda)}
\qquad(\lambda=1,2,\ldots,n)
\]
будем называть «энергетическими соотношениями», поскольку соответствующие вариационной задаче «законы сохранения» \(\Div B^{(\lambda)}=0\) соответствуют «законам энергии», а \(B^{(\lambda)}\) --- «энергетическим компонентам». Итак: если \(I\) допускает группу сдвигов, то энергетические соотношения становятся несобственными тогда и только тогда, когда \(I\) инвариантен относительно бесконечной группы, содержащей группу сдвигов как подгруппу.\footnote{Законы энергии классической механики, а также старой «теории относительности» (где \(\sum \dd x^2\) переходит в себя), являются «собственными», поскольку здесь не появляются бесконечные группы.}

Пример таких бесконечных групп дается группой всех преобразований \(x\) и таких индуцированных преобразований \(u(x)\), в которых входят только производные произвольных функций \(p(x)\); группа сдвигов возникает специализацией \(p^{(i)}(x)=\eps_i\). Однако должно остаться нерешенным, заданы ли этим, вместе с группами, возникающими при изменении \(I\) на граничный интеграл, уже наиболее общие такие группы. Индуцированные преобразования указанного типа возникают, например, если подчинить \(u\) преобразованиям коэффициентов «тотальной дифференциальной формы», то есть формы
\[
\sum a\,\dd^2x_i+\sum b\,\dd x_i\dd x_k+\cdots,
\]
которая кроме \(\dd x\) содержит еще высшие дифференциалы; более специальные индуцированные преобразования, где \(p(x)\) входят только в первой производной, задаются преобразованиями коэффициентов обыкновенных дифференциальных форм \(\sum c\,\dd x_{i_1}\cdots \dd x_{i_\lambda}\), и именно их обычно рассматривали.

Еще одна группа указанного типа, которая из-за присутствия логарифмического члена не может быть преобразованием коэффициентов, например следующая:
\[
y=x+p(x),\qquad
v_i=u_i+\lg(1+p'(x))=u_i+\lg\frac{\dd y}{\dd x};
\]
\[
\D x=p(x),\qquad
\D u_i=p'(x),\qquad
\bd u_i=p'(x)-u_i'p(x).
\]
Зависимости (16) здесь становятся:
\[
\sum_i\left(\psi_i u_i'+\frac{\dd\psi_i}{\dd x}\right)=0,
\]
а несобственные энергетические соотношения:
\[
\sum_i\left(\psi_i u_i'+
\frac{\dd(\psi_i+\mathrm{const.})}{\dd x}\right)=0.
\]
Простейший инвариантный интеграл группы:
\[
I=\int \frac{e^{-2u_1}}{u_1'-u_2'}\,\dd x.
\]
Наиболее общий \(I\) определяется интегрированием дифференциального уравнения Ли (11):
\[
\bd f+\frac{\dd}{\dd x}(f\cdot \D x)=0,
\]
которое после подстановки значений для \(\D x\) и \(\bd u\), если считать \(f\) зависящей только от первых производных \(u\), переходит в
\[
\frac{\p f}{\p x}p(x)
+\left\{\sum_i\frac{\p f}{\p u_i}
-\sum_i\frac{\p f}{\p u_i'}u_i'+f\right\}p'(x)
+\left\{\sum_i\frac{\p f}{\p u_i'}\right\}p''(x)=0
\]
(тождественно по \(p(x),p'(x),p''(x)\)). Эта система уравнений уже для двух функций \(u(x)\) имеет решения, действительно содержащие производные, именно:
\[
f=(u_1'-u_2')\,
\Phi\!\left(u_1-u_2,\frac{e^{-u_1}}{u_1'-u_2'}\right),
\]
где \(\Phi\) обозначает произвольную функцию указанных аргументов.

Гильберт формулирует свое утверждение так: отказ собственных законов энергии является характерным признаком «общей теории относительности». Чтобы это утверждение было буквально верным, название «общая относительность» следует понимать шире, чем обычно, и распространять также на предшествующие группы, зависящие от \(n\) произвольных функций.\footnote{Этим снова подтверждается правильность замечания Клейна, что обычное в физике название «относительность» следует заменить на «инвариантность относительно группы». (Über die geometrischen Grundlagen der Lorentzgruppe. \emph{Jhrber. d. d. Math. Vereinig.} т. 19, с. 287, 1910; перепечатано в Phys. Zeitschrift.)}
% END UNIT BODY
% END PRODUCER UNIT 075
\endgroup


\clearpage
\begingroup
\setcounter{footnote}{0}
% BEGIN PRODUCER UNIT 076
% Source: C:/Users/Floris/Documents/interlanguage/03_projects/noether/02_slavic_working_corpus/translations/paper14/russian/v001/Noether_Paper14_Russian_v001.tex
% Identity: 99545 B / SHA-256 3FD6246F702A86117C52154C7AA67CFA3575494C71DDB684DDF3D8DC252D8DCC
\setlength{\parindent}{1.25em}
\setlength{\parskip}{0pt}
\emergencystretch=30em
\allowdisplaybreaks
\raggedbottom
\providecommand{\p}{\partial}
\providecommand{\frH}{\mathfrak H}
\providecommand{\frK}{\mathfrak K}
\providecommand{\frM}{\mathfrak M}
\providecommand{\frN}{\mathfrak N}
\providecommand{\frG}{\mathfrak G}
\providecommand{\frS}{\mathfrak S}
\providecommand{\frD}{\mathfrak D}
\providecommand{\frf}{\mathfrak f}
\providecommand{\frc}{\mathfrak c}
\providecommand{\fra}{\mathfrak a}
\providecommand{\frb}{\mathfrak b}
\providecommand{\frp}{\mathfrak p}
\providecommand{\frr}{\mathfrak r}
\providecommand{\frm}{\mathfrak m}
\providecommand{\calA}{\mathcal A}
\providecommand{\calB}{\mathcal B}
\providecommand{\calN}{\mathcal N}
\providecommand{\tmod}[1]{\;(\operatorname{mod} #1)}
% BEGIN UNIT BODY
% Unit: Paper 14. Direct Russian translation from the German final-audited source slice.

\begin{center}
{\Large\bfseries 14. Арифметическая теория алгебраических функций одной переменной в её связи с остальными теориями и с теорией числовых полей\par}
\vspace{1.5em}
\emph{Jahresbericht der Deutschen Mathematiker-Vereinigung} 38 (1919), с. 182--203
\end{center}

Следующий обзор возник из желания иметь связующее звено между отдельными теориями алгебраических функций; вместе с тем его можно рассматривать как дополнение к обзору Брилля--Нётер\footnote{\emph{Die Entwicklung der Theorie der algebraischen Funktionen in älterer und neuerer Zeit}. \emph{Jahresbericht der Deutschen Mathematiker-Vereinigung} т. 3 (1894) [далее цитируется как Брилль--Нётер].}, в котором рассмотрение арифметической теории по существу отсутствует.\footnote{За исключением обсуждения исследования Кронекера о дискриминанте. (\emph{J. f. M.} т. 91.)}

Помимо Брилля--Нётер, при обзоре отдельных теорий следует учитывать следующие энциклопедические статьи вместе с указанной там литературой:

для теории Римана: Виртингер (II B 2), № 1--19;

для теории Вейерштрасса: Виртингер (II B 2), № 23, 32, 33, 34;

для алгебраико-геометрической теории Брилля--Нётер: Виртингер (II B 2), № 21--31; Берцолари (III C 4), № 23--38;

для арифметической теории Дедекинда--Вебера и Гензеля--Ландсберга: Гензель (II C 5), [далее цитируется как Гензель]\footnote{Части, относящиеся к Дедекинду--Веберу, принадлежат А. Островскому.};

в более общем виде, для арифметической теории алгебраических элементов: Ландсберг (I C 5), а для арифметической теории алгебраических функций нескольких переменных: Юнг (II C 6).

Для теории числовых полей ср. «Zahlbericht» Гильберта\footnote{\emph{Die Theorie der algebraischen Zahlkörper}. \emph{Jahresbericht der Deutschen Mathematiker-Vereinigung} т. 4 (1894/95).} и лекции по теории чисел Дирихле--Дедекинда [далее цитируется как Дедекинд--теория чисел].\footnote{Брауншвейг, Vieweg, 4-е изд., 1894.}

\subsection*{Содержание}
\begin{enumerate}
\item Основания различных теорий.
\item Параллелизм между алгебраическими числами и алгебраическими функциями (сопряжённые элементы, теоремы о базисе, понятие модуля).
\item Параллелизм и различия в теории идеалов алгебраических чисел и функций.
\item Связь арифметического обоснования разложения в ряды с теорией идеалов.
\item Сравнение арифметической теории алгебраических функций с геометрической теорией. Различное понятие инвариантности.
\item Продолжение сравнения. Особые точки.
\item Сравнение между теоремой об остатке и теоремами теории идеалов.
\item Трансцендентные вопросы для алгебраических функций и алгебраических чисел.
\end{enumerate}

\subsection*{1. Основания различных теорий}
Теория Римана является единственной чисто трансцендентной теорией, основанной на теоремах существования (принцип Дирихле, методы Шварца--Неймана). Интегралы алгебраических функций здесь первичны, сами алгебраические функции вторичны. Функции на римановой поверхности характеризуются своими нулями и полюсами; это соответствует полигонам (соответственно дивизорам) арифметической теории, тогда как Вейерштрасс говорит также о самих функциях вместе с их нулями и полюсами. Главная задача, помимо теорем существования, состоит в построении всех функций с заданными полюсами; это выполняется с помощью интегралов первого и второго рода. Далее возникает вопрос о числе произвольных постоянных, которые ещё входят при заданных бесконечных точках. На этот последний вопрос отвечает теорема Римана--Роха; в арифметической теории ему соответствует вопрос о размерности класса полигонов (соответственно класса дивизоров), а в алгебраико-геометрической теории --- вопрос о размерности полной линейной системы, ответ на который также даёт теорема Римана--Роха. В арифметической теории фактическое построение функций осуществляется посредством теорем о базисе; в алгебраико-геометрической теории --- посредством теоремы об остатке. В частности, нулям подынтегральных функций первого рода соответствует дифференциальный класс в арифметической теории, а в алгебраико-геометрической теории --- учение о специальных группах, то есть о системе групп точек, высекаемой адъюнктными кривыми порядка \((n-3)\), \(\varphi\)-кривыми.

Как риманова теория исторически предшествовала остальным теориям, хотя строгое доказательство теорем существования было получено лишь позднее, так и позднее с помощью её трансцендентных средств были решены алгебраические вопросы, которые до сих пор недоступны чисто алгебраическому или арифметическому рассмотрению; прежде всего здесь следует упомянуть теорию особых соответствий Гурвица.\footnote{A. Hurwitz, Über algebraische Korrespondenzen und das erweiterte Korrespondenzprinzip. \emph{Math. Ann.} т. 28. (Ср. Брилль--Нётер, раздел 10.)}

Аналогии с трансцендентными постановками вопросов в теории алгебраических числовых полей обсуждаются в последнем пункте.

Для остальных теорий первичны алгебраические функции, а интегралы вторичны. При этом теория Вейерштрасса опирается на трансцендентные основания разложения в степенные ряды и теоремы о вычетах; затем, однако, она развивается алгебраически, через явное указание всех рассматриваемых функций.

Арифметическая теория в изложении Гензеля--Ландсберга также использует трансцендентные основания разложения в ряды. Но, сопоставляя на этой основе отдельным точкам дивизоры, она затем идёт чисто арифметическим путём и потому может рассматриваться как слияние Вейерштрасса и Дедекинда--Вебера. В отличие от последних, она одновременно даёт методы фактического построения всех рассматриваемых базисных функций. В деталях она имеет ряд упрощений по сравнению с Дедекиндом--Вебером, главным образом за счёт введения «дробных» дивизоров. Недавно Гензель\footnote{\emph{Mathematische Zeitschrift} т. 4; это обоснование положено также в основу его обзора.} арифметически обосновал теорию разложения в ряды (за исключением метода мажорант для доказательства сходимости), благодаря чему теория стала почти чисто арифметической.

Алгебраико-геометрическая теория Брилля--Нётер предполагает лишь понятие «бесконечно близких точек» одной ветви, которое сводится к разложению в ряд в обыкновенной точке\footnote{Для бесконечно близких «особых» точек ср. Брилль--Нётер, раздел 6, особенно № 10 и 11. Обыкновенные бесконечно близкие точки --- это, например, точки пересечения касательной с кривой или точки касания более высокого порядка (точки ветвления).} и соответствует вейерштрассову понятию элемента; затем она работает алгебраико-рациональными методами, по существу методами тернарной теории исключения, и также доходит до явных представлений.

Первоначальная, чисто арифметическая теория Дедекинда--Вебера\footnote{\emph{J. f. M.} 92 (цитируется как Д.--В.).} родственна римановой постольку, поскольку её утверждения снова имеют характер доказательств существования; она сохраняет полную чистоту метода.\footnote{По этой причине она также более пригодна для сравнения с остальными теориями, чем, например, теория Гензеля--Ландсберга, которая лишь изредка будет появляться в следующих сравнениях. --- Теория абелевых интегралов и их обращения, однако, насколько она предполагает непрерывность, у Д.--В. и также в дальнейшем развитии алгебраической теории (M. Noether, \emph{Ann.} 37) не рассматривается.} Арифметическое обоснование теории, где переменные играют лишь роль неопределённых, идёт параллельно теории алгебраических чисел, как будет подробнее объяснено в п. 2.\footnote{Сравнение теории Гензеля--Ландсберга с теорией чисел, в частности с теорией \(p\)-адических чисел, имеется у Гензеля.} Это обоснование даёт средство арифметически ввести понятие точки (ср. Гензель, № 10a) и тем самым построить понятие абсолютной римановой поверхности без каких-либо соображений непрерывности. Теория вследствие этого получает формальный характер; например, понятие дифференциала также вводится чисто формально. На основании, полученном через понятие точки, можно установить близкое родство с методами и понятиями алгебраико-геометрической теории, которая, если не считать понятия элемента, также свободна от трансцендентных средств; это будет сделано в следующих пунктах.

\subsection*{2. Параллелизм между алгебраическими числами и алгебраическими функциями. (Сопряжённые элементы, теоремы о базисе, понятие модуля)}
Алгебраическое числовое поле определяется как поле \(n\)-й степени над полем \(R\) рациональных чисел; поле алгебраических функций одной переменной --- как поле \(n\)-й степени над полем \(K(z)\) всех рациональных функций одной неопределённой \(z\) с произвольными комплексными коэффициентами. Следовательно, для обоих полей совместно верны все теоремы, полностью отвлекающиеся от специальной природы основного поля; это теоремы о базисе, нормах, следах и дискриминантах.\footnote{Д.--В., § 1 и 2; Гензель, № 4 (с примечанием 5).} Все эти понятия у Дедекинда--Вебера вводятся без использования сопряжённых элементов. Однако следует заметить, что понятие сопряжённого поля можно определить совершенно формально, как показал Штейниц, опираясь на Кронекера.\footnote{Steinitz, Algebraische Theorie der Körper. (\emph{J. f. M.} 137, § 6 и 8.)} Действительно, если \(K\) --- произвольное поле, а \(\varphi(x)\) --- неприводимый над \(K\) многочлен от \(x\), то система классов вычетов по модулю \(\varphi(x)\) образует поле. Если классу вычетов \(x\) поставить в соответствие элемент \(\xi\), то при изоморфном сопоставлении классу вычетов, представленному многочленом \(f(x)\), соответствует элемент \(f(\xi)\). Поскольку все многочлены, делящиеся на \(\varphi(x)\), представляют нулевой класс вычетов, \(\varphi(\xi)\) сопоставляется нулевому классу; следовательно, элемент \(\xi\) можно символически рассматривать как корень уравнения \(\varphi(x)=0\), и он порождает алгебраическое поле \(K(\xi)\). Повторение процедуры приводит к \(n\) сопряжённым полям \(K(\xi), K(\xi'),\ldots,K(\xi^{(n-1)})\). Значит, в частности и для алгебраических функций можно говорить о сопряжённых элементах, не покидая чисто арифметической теории.

Два основных поля \(R\) и \(K(z)\) обладают общим свойством: в них определена система всех целых элементов, соответственно целых рациональных чисел и целых рациональных функций одной неопределённой. Эти целые элементы имеют характерное свойство: любые два \(a\) и \(b\) имеют наибольший общий делитель \(d\), представимый в форме \(ma+nb\), где \(d,m,n\) также являются целыми элементами из \(R\) соответственно \(K(z)\); при этом \(d\) есть абсолютно наименьшее число соответственно многочлен наименьшей степени, представимый в такой форме. Только на этом свойстве, влекущем единственность разложения на простые множители для целых рациональных чисел соответственно функций одной неопределённой, основаны теоремы о целых элементах алгебраического поля или надполя. Рассуждение, доказывающее существование наибольшего общего делителя, даёт также существование фундаментальной системы, или базиса, целых элементов алгебраического поля. В самом деле, если рассматривать только такие базисы поля, которые состоят из целых элементов, то их дискриминант является целым рациональным числом соответственно целой рациональной функцией одной неопределённой. Каждый базис, отвечающий абсолютно наименьшему числу или функции наименьшей степени, образует фундаментальную систему.

Более точное понимание вопросов базиса, также для более общих полей, даёт введение понятия модуля, которое следует сформулировать так, чтобы оно охватывало различные определения, встречающиеся в литературе в зависимости от природы элементов. Модулем --- относительно основной области целостности \(\frS\) --- называется система элементов такая, что разность двух элементов снова принадлежит системе, равно как и произведение элемента на произвольный элемент из \(\frS\). (Если \(\frS\) состоит из целых рациональных чисел, то второе требование следует из первого.)\footnote{Понятие модуля восходит к Дедекинду, который впервые ввёл числовой модуль (\(\frS=\) целые рациональные числа) во 2-м изд. «Теории чисел»; модуль из целочисленных линейных форм здесь также содержится неявно. У Дедекинда--Вебера появляются «модули функций», элементы которых состоят из алгебраических функций, тогда как \(\frS\) состоит из всех целых рациональных функций одной неопределённой. --- Модуль из однородных форм от \(n\) неопределённых, в более общем виде из многочленов, впервые появляется у Гильберта (\emph{Annalen} 36); \(\frS\) состоит из всех многочленов, и, поскольку сами элементы также являются многочленами, этот модуль подпадает под наше определение идеала. Кронекеровы «модульные системы» из многочленов исходят из базиса, а не из абстрактного определения.}

Каждый модуль, имеющий базис из конечного числа элементов, через которые каждый элемент линейно представим с коэффициентами из \(\frS\), называется конечным (конечнопорождённым); в частности, каждый однопорождённый модуль имеет вид \(c\cdot\alpha\), где \(\alpha\) обозначает базисный элемент, а \(c\) пробегает все элементы из \(\frS\). \emph{Модуль, элементы которого сами принадлежат \(\frS\), называется идеалом в \(\frS\);} отсюда следует понятие конечного идеала, в частности однопорождённого (главного идеала).

Все теоремы о базисе для целых элементов и идеалов основаны на следующей теореме:

\emph{Если \(M\) --- модуль относительно \(\frS\), элементы которого являются линейными формами от \(n\) неопределённых с коэффициентами из \(\frS\), и если каждый идеал в \(\frS\) конечен, то \(M\) также конечный модуль. Если, в частности, каждый идеал в \(\frS\) главный, то \(M\) имеет базис из линейно независимых форм.}\footnote{В литературе эта теорема появляется отдельно для различных областей коэффициентов \(\frS\). Для целых рациональных чисел ср., например, Гильберт, Zahlbericht § 3 и 4; для многочленов от нескольких неопределённых: Гильберт, Über die vollen Invariantensysteme: конец § 2 (\emph{Math. Ann.} 42).}

В самом деле, если элементы \(M\) заданы как \(l(x)=a_1x_1+\cdots+a_nx_n\), то совокупность коэффициентов \(a_1\) пробегает идеал в \(\frS\), который по предположению имеет вид \(b_1a^{(1)}+\cdots+b_\rho a^{(\rho)}\). Принадлежащая \(M\) линейная форма
\[
m(x)=l(x)-b_1l^{(1)}(x)-\cdots-b_\rho l^{(\rho)}(x)
\]
поэтому зависит только от \((n-1)\) неопределённых; отсюда конечным повторением следует утверждение. Если каждый раз \(\rho=1\), то базис состоит из \(n\) форм соответственно от \(n,n-1,\ldots,1\) неопределённых, что даёт линейную независимость ненулевых форм.

Теорема о фундаментальной системе выводится отсюда в различных случаях следующим образом.

Пусть алгебраическое поле получается присоединением целого алгебраического элемента \(\delta\) к основному полю, и пусть \(\frS\) обозначает совокупность всех целых элементов основного поля. Тогда каждый целый элемент надполя допускает представление (ср., например, Zahlbericht § 3)
\[
\omega=\frac{a_1\delta^{\,n-1}+a_2\delta^{\,n-2}+\cdots+a_n}{\Delta(\delta)},
\]
где \(a_i\) --- элементы из \(\frS\), а \(\Delta(\delta)\) обозначает дискриминант \(\delta\). Посредством подстановки
\[
x_i=\frac{\delta^{\,n-i}}{\Delta(\delta)}
\]
модуль всех целых элементов переходит в модуль линейных форм; следовательно, существование фундаментальной системы получается во всех случаях, когда каждый идеал в \(\frS\) конечен. То же рассуждение показывает, что в более общем виде каждый идеал в области целых элементов надполя конечен. В частности, для целых рациональных чисел соответственно функций одной неопределённой каждый идеал является главным; отсюда следует, что фундаментальные системы в алгебраических числовых полях соответственно в полях алгебраических функций одной переменной состоят ровно из \(n\) элементов, если \(n\) обозначает степень. И поскольку в этих полях, как сказано выше, каждый идеал конечен, отсюда одновременно получается фундаментальная система для относительных полей, которая в общем случае будет состоять из большего числа элементов, чем указывает относительная степень.\footnote{По исследованиям Штейница о модулях из линейных форм, где \(\frS\) состоит из целых чисел (конечного) алгебраического числового поля, достаточно \(\nu+1\) базисных элементов, если \(\nu\) обозначает относительную степень (\emph{Math. Ann.} 71, § 4).} Далее, если \(\frS\) обозначает область всех многочленов от нескольких неопределённых, то в \(\frS\) каждый идеал также конечен;\footnote{По гильбертовой теореме о базисе модуля (\emph{Math. Ann.} 36). Модули из многочленов, которые здесь ради последовательности называются идеалами, обычно называются «модулями форм» или просто «модулями». И гильбертова теорема, что такой модуль \(N\) конечен, также сводится к теореме о линейных формах. Пусть \(f\) --- многочлен \(n\)-й степени от \(r+1\) неопределённых \(x_1,\ldots,x_r,x_{r+1}\), содержащийся в \(N\) и имеющий член \(x_{r+1}^{\,n}\) (этого всегда можно достичь линейным преобразованием). Тогда все многочлены из \(N\) по модулю \(f\) сравнимы с теми, которые представимы в виде \(a_1(x)x_{r+1}^{\,n-1}+\cdots+a_n(x)\); следовательно, посредством \(x_{r+1}^{\,n-i}=X_i\) они образуют модуль линейных форм, для которого каждый идеал в области от \(r\) неопределённых можно считать конечным, так как это известно для одной неопределённой.} значит, существование фундаментальной системы получается и для полей алгебраических функций нескольких неопределённых; она снова будет состоять из большего числа элементов, чем указывает степень.

\subsection*{3. Параллелизм и различия в теории идеалов алгебраических чисел и функций}
Доказательство основной теоремы теории идеалов, что каждый идеал однозначно представим как произведение степеней простых идеалов, можно вести совместно для алгебраических чисел и функций, опираясь только на свойство, что для целых рациональных чисел соответственно функций одной неопределённой каждый идеал становится главным.\footnote{Ср. замечание во введении Гурвица: Über die Theorie der Ideale (\emph{Gött. Nachr.} 1894). В учебнике алгебры Вебера доказательство, в связи с Кронекером и посредством введения функционалов, проведено параллельно для обоих полей. [Связь между кронекеровой теорией и теорией идеалов обеспечивает «расширенная теорема Гаусса»; Гурвиц, там же.]} Главный пункт доказательства состоит в установлении того, что из делимости идеала \(\frc\) на идеал \(\fra\) (то есть каждый элемент \(\frc\) содержится в \(\fra\)) следует произведительное представление \(\frc=\fra\frb\).\footnote{Особенно подчёркнуто у Дедекинда (Теория чисел); Supplement XI, Satz VII, § 177; ср. примечание на с. 554. Теорема использует тот факт, что речь идёт об алгебраических полях; она уже не верна для идеалов (модулей) из многочленов, где на место произведения становится наименьшее общее кратное.} У Гурвица это доказательство опирается на «расширенную теорему Гаусса», утверждающую, что если целый элемент \(\omega\), который делит каждый коэффициент произведения двух многочленов \(\varphi\) и \(\psi\), то он делит также произведение всех коэффициентов \(\varphi\) и \(\psi\); у Дедекинда --- на эквивалентную ей модульную теорему.\footnote{Теория чисел, Suppl. XI; Satz VI, § 173. На эквивалентность двух теорем указывает Дедекинд (Über die Begründung der Idealtheorie, \emph{Göttinger Nachr.} 1895).}

Далее совместно можно определить понятие комплементарного модуля (Гензель, № 12) и из него --- идеал ветвления (основной идеал, дифференту) \(\frD\) (норма которого становится дискриминантом поля \(D\)); и верна основная теорема, что идеал ветвления \(\frD\) есть \emph{наибольший общий делитель всех главных идеалов \(F'(\delta)\)}, где \(\delta\) пробегает все целые элементы поля, а \(F(\delta)=0\) каждый раз обозначает соответствующее неприводимое уравнение. Отсюда следует, что \emph{дискриминант поля содержит все и только те рациональные простые числа соответственно линейные множители по \(z\), которые делятся на квадрат простого идеала}. Для самого \(F'(\delta)\) получается представление
\[
F'(\delta)=\frf\cdot\frD,
\]
и при взятии норм
\[
\Delta(\delta)=R^2\cdot D;
\]
здесь \(\frf\) обозначает проводник кольца (порядка, области целостности), выведенного из \(\delta\), то есть \(\frf\) состоит из совокупности (целых) элементов, которые после умножения на каждый целый элемент делятся на модуль \([1,\delta,\ldots,\delta^{n-1}]\). \(R^2=\operatorname{Norm}\frf\) представляет квадрат детерминанта подстановки, переводящей базис \(1,\delta,\ldots,\delta^{n-1}\) в базис целых элементов.\footnote{Dedekind, Über die Diskriminante endlicher Körper (\emph{Abhandl. der Gött. Ges. d. Wissensch.} т. 29, 1882). Данные там для алгебраических чисел определения непосредственно переносятся на алгебраические функции; доказательства также переносятся почти дословно. Порядок доказательств у Д.--В. иной: см. ниже. Для указанных теорем ср. также Zahlbericht, § 31 и 32. (\(F'(\delta)\) там называется «дифферентой» целого числа \(\delta\).)}

\emph{Различия} в дальнейшем развитии обусловлены специальными свойствами основных полей. Так, в случае алгебраических чисел тот факт, что \emph{элементы основного поля одновременно являются показателями}, приводит к тому, что идеал ветвления может делиться на степень простого идеала \(\frp\), более высокую, чем \((e-1)\)-я, если \(\frp\) входит в \(p\) в \(e\)-й степени (именно когда \(e\) делится на \(p\)); это ведёт к дальнейшим понятиям поля инерции и поля ветвления,\footnote{Ср. Zahlbericht, § 39--45.} которые не имеют аналога в теории алгебраических функций. Но прежде всего вследствие \emph{существования бесконечно многих постоянных} в основном поле для алгебраических функций отпадает конечность числа классов идеалов, а вместе с ней все вопросы, связанные с определением числа классов. Впрочем, в известном смысле область всех (трансцендентных) прайм-форм\footnote{Klein, Zur Theorie der Abelschen Funktionen (\emph{Ann.} 36, 1890).} можно рассматривать как трансцендентный аналог поля классов, поскольку здесь каждая точка порождается одной-единственной формой, главным идеалом.

С другой стороны, существование бесконечно многих постоянных в основном поле \(K(z)\) для алгебраических функций ведёт к теоремам и постановкам вопросов, не имеющим аналога для алгебраических чисел. Сначала отсюда получается некоторое упрощение в построении доказательств; например, в данном Д.--В. доказательстве единственности разложения идеалов на простые идеалы можно обойти «расширенную теорему Гаусса» или эквивалентную ей теорему, используя тот факт, что каждая линейная форма \((z-c)\) имеет бесконечно много классов вычетов, а именно все постоянные (тогда как число классов вычетов по модулю простого числа конечно).\footnote{Д.--В., § 9. Ср. примечания на с. 212 и 213.}

К действительно новым результатам ведёт дальнейший факт: каждый многочлен с коэффициентами, являющимися целыми рациональными функциями от \(z\), по модулю \((z-\alpha)\) распадается на линейные множители (чего не происходит для целочисленного многочлена по модулю \(p\)); этот факт сохраняется даже тогда, когда вместо всех постоянных берутся только все алгебраические числа.\footnote{Д.--В. во введении замечают, что в этом случае вся теория сохраняется.} Отсюда следует, что каждый простой идеал есть идеал первой степени;\footnote{Д.--В., § 9, № 7.} далее следует, что дискриминант поля \(D\) становится пересечением всех дискриминантов \(\Delta(\delta)\), где \(\delta\) пробегает все целые элементы;\footnote{Д.--В., с. 224.} оба утверждения находятся в противоположности к случаю алгебраических чисел. Затем отсюда снова следует, что идеал ветвления становится пересечением всех главных идеалов \(F'(\delta)\).

Дальнейшие различия обусловлены тем, что для алгебраических чисел основное поле \(R\) рациональных чисел является абсолютно выделенным, а именно простым полем, содержащимся в данном поле (то есть полем, выведенным из единицы); тогда как основное поле \(K(z)\) относительно. Каждую непостоянную функцию \(\xi\) можно выбрать как независимую переменную; тогда поле становится алгебраическим полем над \(K(\xi)\). Теорема, что каждый простой идеал \(\frp\) есть идеал первой степени, то есть каждая функция по модулю \(\frp\) сравнима с постоянной, вместе с возможностью взять произвольную функцию в качестве независимой переменной ведёт к важнейшему понятию, выходящему за пределы алгебраических чисел: к понятию точки в чисто арифметической форме. (Ср. Гензель, 10a.) «Каждая точка при этом представлена полем постоянных, изоморфным функциональному полю.»\footnote{Два поля, как известно, называются «изоморфными», если они могут быть взаимно однозначно сопоставлены так, что также соответствуют сумма, разность, произведение и частное.} Совокупность так определённых точек образует «абсолютную риманову поверхность». Это определение точки зависит только от поля, а не от специальной переменной. Напротив, доказательство существования проводится с выделением какой-либо функции в качестве независимой переменной \(z\); точка \(P\) порождается простым идеалом \(\frp\), когда каждой функции сопоставляется её постоянный вычет по модулю \(\frp\). Следовательно, в частности, функции, делящиеся на \(\frp\), исчезают в \(P\), и наоборот. Когда \(\frp\) пробегает все простые идеалы по \(z\), тем самым получаются все точки, в которых \(z\) конечна; остальные точки возникают через введение независимой переменной \(\xi=1/z\) и соответствуют простым идеалам, входящим в главный идеал \(\xi\). Естественно, понятия, примыкающие к понятию точки, также не могут иметь аналога у алгебраических чисел: например понятие полигона, класса полигонов, размерности класса полигонов и т. д.

Следует ещё заметить, что свободный выбор независимой переменной даёт дальнейшее толкование разложения \(F'(\vartheta)=\frf\cdot\frD\). Если рассматривать \(\vartheta\) как независимую переменную, то соответственно получается
\[
\left(\frac{\p F}{\p z}\right)=\frf\cdot\fra,
\]
поскольку кольцо, выведенное из \(\vartheta\), было определено симметрично относительно \(z\) и \(\vartheta\); и \(\frf\) становится наибольшим общим делителем двух главных идеалов \(\frac{\p F}{\p z}\) и \(\frac{\p F}{\p\vartheta}\). Проводник \(\frf\) одновременно становится «идеалом двойных точек» относительно \(\vartheta,z\).\footnote{Д.--В., с. 263.}

\subsection*{4. Связь арифметического обоснования разложения в ряды с теорией идеалов}
У Гензеля--Ландсберга на место теории идеалов становится трансцендентное вспомогательное средство разложения в ряды и заимствованное из теории функций понятие точки. Недавнее арифметическое обоснование разложения в ряды для алгебраических чисел и функций, данное Гензелем,\footnote{\emph{Eine neue Theorie der algebraischen Zahlen} (\emph{Math. Zeitschr.} т. 2, 1918) и \emph{Neue Begründung der arithmetischen Theorie der algebraischen Funktionen einer Variablen} (\emph{Math. Zeitschr.} т. 4, 1919). Для более подробного изложения ср. Гензель, № 44 и далее.} поэтому можно рассматривать как эквивалент теории идеалов. В действительности это арифметическое обоснование сводится к тому, чтобы для каждого простого идеала перейти к такому (трансцендентному) полю расширения, где этот простой идеал становится главным, так что выполняются обычные теоремы разложения; достаточно провести это рассуждение для всех простых идеалов, входящих в простое число \(p\) соответственно в линейную функцию \(z-a\).

Сначала речь идёт о введении трансцендентного поля расширения, образованного всеми разложениями в ряды по целым степеням \(z-a\) соответственно \(p\) (\(p\)-адическими числами), каждый раз с не более чем конечным числом отрицательных степеней. В этом поле неприводимое уравнение, определяющее алгебраическое числовое или функциональное поле, распадается в произведение \(\rho\) неприводимых множителей; это соответствует разложению \(p\) соответственно \(z-a\) в произведение \(\rho\) «первичных» идеалов \(\mathfrak q_1,\ldots,\mathfrak q_\rho\). Дальнейшему представлению каждого первичного идеала как степени простого идеала
\[
\mathfrak q_i=\frp_i^{\,\ell_i},
\]
у Гензеля соответствует введение ещё одного поля, алгебраического относительно трансцендентного поля расширения. В случае алгебраических функций оно может быть определено чистым уравнением; в случае алгебраических чисел --- в более общем виде алгебраически разрешимым уравнением. Упрощение для алгебраических функций объясняется тем, что здесь не появляются поля инерции и ветвления (ср. п. 3).

Появляющийся в этих исследованиях «обрывающийся ряд» по простому элементу \(\pi\)\footnote{\emph{Neue Begründung \(\ldots\)}, формула (11).} (который ещё не требует трансцендентного доказательства сходимости)
\[
v=c_0+c_1\pi+\cdots+c_{\sigma-1}\pi^{\sigma-1}-v_\sigma\pi^\sigma,
\]
где \(c_0,\ldots,c_{\sigma-1}\) обозначают постоянные, а \(v_\sigma\) --- целый элемент поля, находится совершенно соответственно у Д.--В.\footnote{Д.--В., с. 231 и § 15.} Там он выводится из теории идеалов и даёт основу для определения порядка исчезновения (соответственно обращения в бесконечность) функции в точке.

\subsection*{5. Сравнение арифметической теории алгебраических функций с геометрической теорией. Различное понятие инвариантности}
«Инвариантным» является всякое понятие, характеризуемое одним только полем. В арифметической теории эта инвариантность вообще уже дана в определении, которое формулируется относительно всех элементов поля, а не относительно какого-либо базиса или выделенной переменной; примером служит указанное выше понятие точки. Напротив, остальные теории исходят из какой-либо специальной переменной или базиса и затем доказывают независимость понятия от этого специального выбора; инвариантность здесь становится инвариантностью относительно бирационального преобразования.

Пусть, например, поле с одной стороны получается присоединением алгебраического элемента \(s\) к \(K(z)\), а с другой --- присоединением \(\sigma\) к \(K(\xi)\), причём \(f(z,s)=0\) соответственно \(F(\xi,\sigma)=0\) обозначают неприводимые уравнения для \(s\) соответственно \(\sigma\); или ещё более общо --- присоединением конечного числа элементов к \(K(\eta)\), с соответствующими неприводимыми уравнениями
\[
g_1(\eta,\tau_1)=0,\qquad g_2(\eta,\tau_1,\tau_2)=0,\ldots .
\]
Тогда по определению все элементы рационально выражаются через \(z\) и \(s\), равно как через \(\xi\) и \(\sigma\), а также через \(\eta,\tau_1,\tau_2,\ldots\). В частности, \(z\) и \(s\) рационально выражаются через \(\xi\) и \(\sigma\), и наоборот; так же \(z\) и \(s\) рационально выражаются через \(\eta,\tau_1,\tau_2,\ldots\), и наоборот. Следовательно, плоская кривая \(f(z,s)=0\) переходит посредством «бирационального преобразования» в плоскую кривую \(F(\xi,\sigma)=0\) или также в кривую в пространстве переменных \(\eta,\tau_1,\tau_2,\ldots\), определённую уравнениями \(g_1(\eta,\tau_1)=0; g_2(\eta,\tau_1,\tau_2)=0,\ldots\). Инвариантность проявляется так: определение, сформулированное для специальной кривой, сохраняется для всех кривых, возникающих из неё бирациональным преобразованием (в пространстве любого числа измерений), то есть для всех кривых «класса» по Риману. «Точка» здесь определяется как согласованная система значений \(z=a,s=b\), так что \(f(a,b)=0\); эта точка при бирациональном преобразовании вообще снова переходит в согласованную систему значений \(\alpha,\beta\) на \(F(\xi,\sigma)=0\), но в особых случаях, когда \(a,b\) была «особой» точкой, может соответствовать нескольким «точкам» на \(F(\xi,\sigma)=0\). Всегда можно указать кривые, на которых «особая» точка \(f(z,s)=0\) соответствует конечному числу простых, раздельных или соседних точек (разрешение особенностей).\footnote{Ср., например, Брилль--Нётер, раздел 6. Подробнее об «особых» точках в следующем пункте.} Точно так же группа точек снова переходит в группу точек; число её точек инвариантно, если особую точку считать столькими точками, скольким простым точкам она соответствует при подходящем преобразовании. Так как каждую кривую можно бирационально преобразовать в такую, которая имеет «обыкновенные кратные» точки,\footnotemark[\value{footnote}] геометрическая теория развивает свои понятия только для таких специальных кривых и затем доказывает инвариантность относительно всякого бирационального преобразования, в том числе такого, которое ведёт к самым общим особым точкам.

В частности, как выделенную кривую «класса» --- за исключением гиперэллиптического случая --- можно ввести так называемую «нормальную кривую \(\varphi\)» в пространстве \((p-1)\) измерений, не имеющую особых точек; здесь \(\varphi\) --- это \(p\) «адъюнктных кривых порядка \((n-3)\)» для однороднённой \(f(z,s)=0\). Бирациональному преобразованию \(f(z,s)=0\) в \(F(\xi,\sigma)=0\) соответствует линейное преобразование \(\varphi\), так что здесь речь идёт об инвариантности в проективном смысле. Представление всех функций поля как рациональных функций от \(\varphi\) в геометрической теории называется «инвариантным представлением».\footnote{Ср., например, Брилль--Нётер, раздел 8. Для алгебраических функций нескольких переменных вопрос, можно ли свести инвариантность к инвариантности относительно линейного преобразования, не рассматривается.}

Итак, линейная система \(\varphi\), или также высекаемые ими «специальные группы», при каждом бирациональном преобразовании переходит в себя, а значит, характеризуется одним только полем. Ей соответствует дифференциальный класс в арифметической теории, где инвариантность снова непосредственно видна из определения, данного через свойства относительно каждой функции поля как независимой переменной.

\subsection*{6. Продолжение сравнения. Особые точки}
С различным пониманием понятия точки связано то, что «особые точки» кривых, создающие особые трудности в алгебраико-геометрической теории, в обосновании арифметической теории вообще не появляются; поэтому и «адъюнктные функции» не требуются для построения арифметической теории, как это имеет место в геометрической теории, а принадлежат к специальным более высоким частям теории. С другой стороны, в основании геометрической теории не появляются точки ветвления; фактически они исчезают и при представлении интегралов в аронгольдовой нормальной форме (Гензель, № 27, формула 80).

Данное в п. 5 определение особой точки можно сформулировать так: \(f(z,s)=0\) (соответственно кривая в пространстве более высокой размерности) имеет в \(z=a,s=b\) (соответственно \(z=a; s_i=b_i\)) особую точку, если эту систему значений принимает более чем одна точка абсолютной римановой поверхности (точка в арифметическом смысле), или если она принимается в одной или нескольких точках в более высоком порядке. В последнем случае речь идёт о точке полигона ветвления как для \(z\), так и для \(s\); это соответствует «бесконечно близким точкам» геометрической теории; точка становится (высшей) точкой возврата кривой \(f(z,s)=0\) (соответственно кривой в пространстве более высокой размерности).

Поскольку каждая произвольная функция \(\eta\) поля по предположению рационально выражается через \(z\) и \(s\), каждая точка абсолютной римановой поверхности, где \(z=a,s=b\), в силу изоморфии возникает так: каждой функции \(\eta\), представленной через \(z\) и \(s\), сопоставляется её значение при \(z=a,s=b\). Чтобы точка была особой, должна существовать по меньшей мере одна функция \(\eta\), которая при \(z=a,s=b\) принимает несколько значений или несколько раз одну систему значений. А вследствие рациональной выразимости через \(z\) и \(s\) это возможно только тогда, когда в таком представлении при \(z=a,s=b\) числитель и знаменатель одновременно обращаются в нуль. Отсюда следует, что приведённое определение тождественно обычному, по которому
\[
f,\quad \frac{\p f}{\p z},\quad \frac{\p f}{\p s}
\]
обращаются в нуль при \(z=a,s=b\). Ибо в этом случае
\[
\eta=\frac{\frac{\p f}{\p z}}{\frac{\p f}{\p s}}
\]
является функцией указанного типа, многозначной при \(z=a,s=b\), тогда как в противоположном случае всякая функция, даже если она становится \(0/0\), принимает только одно значение.

При арифметическом понятии точки особая точка не может появиться, поскольку две точки считаются различными, как только хотя бы одной функции соответствуют разные системы значений. Но и в теории идеалов, служащей для обоснования арифметического понятия точки, понятие особой точки не возникает, поскольку всегда одновременно рассматривается совокупность всех целых элементов поля. Все лежащие в конечной части особые точки \(F(z,\delta)=0\), где \(\delta\) алгебраически цела относительно \(z\), порождаются (по второму определению и концу п. 3) «идеалом двойных точек» \(\frf\). Так как по основной теореме идеал ветвления \(\frD\) является наибольшим общим делителем всех главных идеалов \(F'(\delta)\), для каждого простого идеала \(\frp\) можно указать такое \(\delta\), что соответствующий \(\frf\) не делится на \(\frp\); следовательно, поскольку \(\frf\) является наибольшим общим делителем \(F'(\delta)\) и \(F'(z)\), по меньшей мере один из двух главных идеалов \(F'(\delta)\) и \(F'(z)\) не делится на \(\frp\). Поэтому \emph{не существует системы значений, для которой одновременно обращаются в нуль все \(F'(\delta)\) и \(F'(z)\)}, и, следовательно, система значений \(z=a,\delta_i=b_i\) принимается только в одной точке абсолютной римановой поверхности; \emph{совокупность целых алгебраических функций поля} (которую можно толковать как кривую в пространстве бесконечно многих измерений) \emph{не имеет особой точки в конечной части}; в теории идеалов же рассматриваются только конечные значения.

Но и базисные функции \(\omega_1,\ldots,\omega_n\) всех целых элементов не имеют особых точек в конечной части. По сказанному выше, каждая точка в арифметическом смысле уже однозначно определена системой постоянных, изоморфно сопоставленной всем целым функциям. Если, следовательно, «пространственная кривая», определённая какими-либо целыми функциями \(s_i\), имеет особую точку в конечной части при \(s_i=\alpha_i\), то должна существовать по меньшей мере одна целая функция \(\eta\), принимающая при \(s_i=\alpha_i\) несколько систем значений. Если теперь \(\omega_1,\ldots,\omega_n\) обозначает базис всех целых элементов, то вследствие \(\delta=a_1\omega_1+\cdots+a_n\omega_n\) с целыми рациональными \(a_i\) каждой системе значений \(\omega\) соответствует только одна система значений \(\delta\); значит, пространственная кривая, определённая \(\omega_1,\ldots,\omega_n\), не может иметь особой точки в конечной части. И \emph{каждая точка абсолютной римановой поверхности, в которой \(z\) конечна, уже однозначно определена системой постоянных, изоморфно сопоставленной \(\omega\). Базисные функции суть именно функции \(\eta\), которые становятся многозначными при \(s_i=\alpha_i\).} Значит, введением базиса целых элементов, например вместо базиса \(1,\delta,\ldots,\delta^{n-1}\), обходится трудность особых точек.

В геометрической теории задача однозначного порождения всех точек абсолютной римановой поверхности --- в более общем виде построения всех функций, которые становятся бесконечными в заданных точках, --- выполняется посредством адъюнктных форм и их частных. Форма \(\psi\) называется адъюнктной, если \(\psi=0\) в каждой \(i\)-кратной точке основной кривой \(f(z,s)=0\) имеет по меньшей мере \((i-1)\)-кратную точку; более высокую особую точку при этом сначала надо разрешить в обыкновенные кратные точки. Выше было показано, что для самой общей функции \(\eta=\frac{\psi}{\chi}\) в особой точке \(f(z,s)=0\) числитель и знаменатель вообще должны обращаться в нуль; что они должны быть адъюнктными, показывает рассмотрение всюду конечных дифференциалов. Что, наоборот, условия адъюнктности достаточны для построения самой общей функции, показывает «теорема об остатке», к которой мы ещё вернёмся подробнее. Следовательно, адъюнктные функции занимают место базисных функций арифметической теории.

Арифметически адъюнктной форме соответствует числитель функции поля, делящейся на «идеал двойных точек» \(\frf\), как только особые точки --- чего всегда можно достичь преобразованием --- считаются все лежащими в конечной части. Отсюда формально видно представление базисных элементов \(\omega_1,\ldots,\omega_n\) как частных адъюнктных форм, а тем самым и связь различных способов их представления:

Пусть \(R(z)\) обозначает детерминант подстановки базиса \(1,\delta,\ldots,\delta^{n-1}\) в базис \(\omega_1,\ldots,\omega_n\). Тогда, наоборот, всякое \(\omega\) выражается как частное целой функции от \(z\) и \(\delta\) на \(R(z)\):
\[
\omega_i=\frac{G_i(z,\delta)}{R(z)}.
\]
Так как \(R(z)\omega_i\) принадлежит кольцу, выведенному из \(\delta\), то по определению проводника \(R(z)\) делится на \(\frf\) [из \((R(z))^2=\operatorname{Norm}\frf\) эта делимость следует только для \((R(z))^2\)]. Поскольку же \(\omega_i\), как целая функция, остаётся конечной для всех конечных \(z\), числительная функция \(G_i(z,\delta)\) также должна делиться на \(\frf\); числители и знаменатели всех \(\omega_i\) поэтому адъюнктны. Из этого представления \(\omega\) следует представление всех функций поля частными адъюнктных.

Согласно предыдущему, геометрическую теорию можно охарактеризовать как \emph{теорию кольца, выведенного из элемента \(\delta\)} (порядка), в противоположность арифметической теории, одновременно рассматривающей \emph{все целые элементы}. Поэтому ясно, что проводник кольца --- особые точки --- играет решающую роль. Можно указать и дальнейшие аналогии между арифметической теорией колец (порядков) и геометрической теорией. Так, геометрическая теория при пересечении с адъюнктными кривыми нигде не считает точки пересечения, попадающие в особые точки, а рассматривает только отличные от них группы точек. Этому соответствует то, что Дедекинд\footnote{Über die Anzahl der Idealklassen in den verschiedenen Ordnungen eines endlichen Körpers. (Festschrift Braunschweig 1877); § 3--5. (Теоремы, сформулированные в этих параграфах для алгебраических чисел, дословно верны для алгебраических функций одной переменной.)} рассматривает только такие идеалы кольца, которые взаимно просты с идеалом-проводником \(\frf\), и для этой области идеалов устанавливает все законы делимости, включая единственность разложения на простые идеалы.

Также лежащая в основании геометрической теории «основная теорема алгебраических функций»\footnote{M. Noether, Über einen Satz aus der Theorie der algebraischen Funktionen. \emph{Math. Ann.} т. 6 (1873).}, или по крайней мере непосредственное следствие из неё, в известном смысле имеет аналог в арифметической теории колец. Эта теорема даёт условия представимости
\[
\psi(z,s)=A(z,s)f(z,s)+B(z,s)\varphi(z,s),
\]
где все входящие величины целы рационально по \(s,z\) (в более общем виде целы и однородны по \(x_1,x_2,x_3\)). Из этих условий следует,\footnote{Брилль--Нётер, № 55.} что в силу \(f(z,s)=0\) частное \(\frac{\psi}{\varphi}\) всегда равно целой рациональной функции \(B(z,s)\), если оно адъюнктно к \(f\). Этому соответствует теорема из п. 3, лежащая в основании всех этих сравнений: проводник \(\frf\) кольца, выведенного из \(\delta\), равен идеалу двойных точек кривой \(f(z,\delta)=0\). Ибо по определению проводника всякий алгебраически целый элемент, делящийся на \(\frf\), принадлежит кольцу и потому цел и рационально представим через \(z\) и \(\delta\); тогда как упомянутая теорема утверждает, что такой элемент является адъюнктной функцией.\footnote{Здесь, следуя Дедекинду, всюду предполагается, что \(\delta\) является целым алгебраическим элементом (чего всегда можно добиться преобразованием). Случай, более точно соответствующий геометрии, \(f(z,s)=0\), где \(s\) может также алгебраически дробно зависеть от \(z\), рассматривают Гензель--Ландсберг; ср. Гензель, № 26 и 29.}

\subsection*{7. Сравнение между теоремой об остатке и теоремами теории идеалов}
Только что упомянутое следствие «основной теоремы алгебраических функций» в геометрической теории ведёт прямо к «теореме об остатке», а тем самым к основанию, на котором строится вся геометрическая теория. В известном отношении, однако, теорема об остатке снова имеет более элементарный характер, чем остальные основания, поскольку её можно арифметически сформулировать как прямое следствие единственности разложения идеалов на простые идеалы, или скорее фактов, лежащих в основе этой теоремы разложения; значит, она не опирается на высшую теорию колец. К такой арифметической формулировке приходят следующим образом.

Пусть основная кривая \(f(z,\delta)=0\) выбрана так --- чего всегда можно достичь линейно-дробным преобразованием переменных, --- что особые точки, как и конечное число рассматриваемых групп точек, лежат в конечной части; и пусть далее \(\delta\) алгебраически цела над \(z\), чего всегда можно добиться последующим линейным преобразованием, целым по переменным. Тогда простому идеалу \(\frp\), порождающему точку \(P\) абсолютной римановой поверхности, где \(z=a,\delta=b\), соответствует точка \(z=a,\delta=b\) кривой \(f(z,\delta)=0\); в более общем виде идеал \(\fra=\frp_1^{\,e_1}\cdots\frp_\rho^{\,e_\rho}\) порождает группу точек, состоящую из точек \(z=a_i,\delta=b_i\) с соответствующей кратностью; и так порождаются все точки \(f(z,\delta)=0\), лежащие в конечной части. Поскольку в точке \(P\) все функции \(g_1,g_2,\ldots\), делящиеся на \(\frp\), обращаются в нуль, и поскольку, наоборот, эта совокупность функций \(g\) порождает идеал \(\frp\), соответствующая точка на \(f\) является пересечением
\[
g_1(z,\delta)=0,\quad g_2(z,\delta)=0,\ldots
\]
с \(f=0\); то же верно для группы точек, соответствующей идеалу \(\fra\). Более того, по теореме, что каждый идеал является наибольшим общим делителем двух главных идеалов, всегда достаточно двух таких кривых, чтобы высечь группу точек. В частности, группа точек, соответствующая главному идеалу, высекается кривой \(g(z,\delta)=0\), и наоборот; главный идеал соответствует полному пересечению.

Две группы точек \(A\) и \(B\) в геометрической теории называются \emph{корезидуальными}, если их можно дополнить одной и той же группой точек \(R\) до полного пересечения. Если эти группы точек соответственно порождаются идеалами \(\fra,\frb,\frr\), то между этими идеалами имеется отношение
\[
\fra\cdot\frr=(\alpha),\qquad \frb\cdot\frr=(\beta),
\]
которое показывает, что идеалы \(\fra\) и \(\frb\) эквивалентны. Обратно, из эквивалентности двух идеалов, \(\fra=\eta\frb\), это отношение всегда следует на основании теоремы, что каждый идеал умножением на другой можно превратить в главный идеал;\footnote{Ср. Дедекинд, Теория чисел § 181; если \(\frr\) выбрано так, что \(\fra\frr\) становится главным идеалом \((\alpha)\), то \(\frb\frr=\eta^{-1}\cdot(\alpha)\); и поскольку \(\frb\frr\) содержит только целые элементы, то же верно для \(\eta^{-1}\cdot(\alpha)\), который поэтому также становится главным идеалом \((\beta)\).} \emph{следовательно, эквивалентные идеалы и корезидуальные группы точек являются тождественными понятиями.}

Теперь теорема об остатке утверждает: \emph{если \(A\) и \(B\) корезидуальны, то их можно высечь адъюнктными кривыми с одним и тем же остатком (вне особых точек).} Поскольку нули и полюсы функции всегда корезидуальны, тем самым одновременно доказано, что самые общие функции поля представимы как частные адъюнктных форм. На языке теории идеалов теорема об остатке просто сводится к тому, что наряду с \(\fra\) и \(\frb\) также \(\fra\frf\) и \(\frb\frf\) эквивалентны, где \(\frf\) обозначает идеал двойных точек. Ибо тогда, по сказанному выше, всегда существует идеал \(\frm\) такой, что
\[
\fra\frf\frm=(\alpha),\qquad \frb\frf\frm=(\beta)
\]
становятся главными идеалами; \(\alpha=0\) и \(\beta=0\) --- это адъюнктные кривые, высекающие группы точек \(A\) и \(B\), соответственно их подгруппы \(A'\) и \(B'\), отличные от особых точек.\footnote{Если \(\fra=\frf\fra_1,\ \frb=\frf\frb_1,\ \frr=\frf\frr_1\), где \(\fra_1,\frb_1,\frr_1\) взаимно просты, то уже \(\fra_1\frf\mathfrak n,\ \frb_1\frf\mathfrak n\) становятся главными идеалами, где \(\mathfrak n\) взаимно прост с \(\frf\). Ибо каждый идеал можно превратить в главный умножением на идеал, взаимно простой с данным (Д.--В., с. 214; или Теория чисел, § 178, теорема XI).} Поэтому доказательство теоремы об остатке арифметически содержится в возможности сопоставлять точки простым идеалам и в доказательстве, что эквивалентные идеалы можно умножением на один и тот же идеал превратить в главные.

Корезидуальные группы точек не обязаны состоять из одинакового числа точек, так же как эквивалентные идеалы не обязаны иметь одинаковую степень. Например, все главные идеалы эквивалентны, и все группы точек, порождённые полным пересечением, корезидуальны. Это различие с эквивалентными полигонами (дивизорами)\footnote{Полигон \(A\) состоит из конечного числа точек абсолютной римановой поверхности; два полигона \(A\) и \(B\) называются эквивалентными, если они являются нулями и полюсами некоторой функции \(\eta\); тогда \(\eta\) допускает представление полигональными частными \(\eta=\frac{A}{B}\). (Д.--В., § 17 и 18.)} объясняется тем, что точки, в которых \(z\) становится бесконечной, не порождаются простыми идеалами по \(z\). В представлении функции как идеальной дроби \(\eta=\frac{\fra}{\frb}\), выражающем эквивалентность \(\fra\) и \(\frb\), нули или полюсы \(\eta\), где \(z\) становится бесконечной, следовательно, не появляются. Так, например, главный идеал \((\alpha)\) становится произведением простых идеалов, соответствующих нулям целой функции \(\alpha\); полюсы \(\alpha\) не входят, поскольку \(\alpha\) становится бесконечной только в точках, где \(z\) становится бесконечной. В этом отношении теория идеалов является точным аналогом теории форм в геометрической теории, где также есть только нули и нет полюсов.

Переход от форм к функциям геометрически совершается посредством образования частных форм одного порядка; арифметически этому соответствует переход от классов идеалов к инвариантно определённым классам полигонов (дивизоров). Поскольку здесь больше не выделена никакая переменная, представление функции полигональными частными даёт все нули и полюсы. Отношение класса идеалов (класса корезидуальных групп точек) к классу полигонов таково. Пусть \(\calA\) и \(\calB\) --- полигоны, порождённые идеалами \(\fra\) и \(\frb\), а \(\calN\) --- полигон точек, в которых \(z\) становится бесконечной (полигон знаменателя \(z\) и целых функций от \(z\)). Тогда из эквивалентности \(\fra\) и \(\frb\) следует эквивалентность \(\calA\) и \(\calB\) тогда и только тогда, когда \(\fra\) и \(\frb\) имеют одинаковую степень; вообще же отсюда следует только эквивалентность \(\calA\calN^\nu\) с \(\calB\calN^{\nu'}\) при \((\nu\ne\nu')\). Поэтому разбиение идеалов (соответственно корезидуальных групп точек) на классы есть объединение бесконечно многих классов полигонов в один класс; его можно также понимать как разбиение полигонов на классы по модулю степени \(\calN\).\footnote{Гензель--Ландсберг, 25-я лекция, с. 438. --- Если представить две переменные \(z\) и \(s\) полигональными частными, это можно понимать как бинарную гомогенизацию каждой из этих переменных; если, в частности, выбрать \(z\) и \(s\) так, чтобы они имели один и тот же полигон знаменателя, возникает обычная в геометрии тернарная гомогенизация.}

Теорема об остатке для случая корезидуальных групп точек, состоящих из разного числа точек, появляется только в чисто геометрических вопросах. В применении к алгебраическим функциям фактически появляется только случай, соответствующий классам полигонов, то есть случай корезидуальных групп точек одинакового числа. «Полная линейная система», порождённая группой точек, то есть совокупность групп точек, корезидуальных этой группе и с тем же числом точек, точно соответствует классу полигонов, порождённому соответствующим полигоном; совпадают также понятия «сумма двух полных систем» и «произведение двух классов полигонов». То, что такая полная система содержит лишь конечное число линейно независимых групп точек, является прямым следствием теоремы об остатке, которая сводит размерность системы (число линейно независимых групп точек) к размерности всех адъюнктных кривых произвольного, но фиксированного порядка, проходящих через остаточную группу. Чтобы прийти к фактическому определению этой размерности, теореме Римана--Роха, геометрическая теория сначала выводит из теоремы об остатке «теорему редукции» (теорему о неподвижных точках);\footnote{Брилль--Нётер, № 58.} а из неё --- теорему Римана--Роха.

Совершенно соответственно арифметическая теория показывает, что каждый класс полигонов имеет конечную размерность, и определяет её с помощью специального базиса целых элементов, нормального базиса, определённого посредством комплементарного модуля; так она получает теорему Римана--Роха как следствие связи комплементарного модуля с полигоном ветвления. При этом в изложении Д.--В. также появляется теорема редукции,\footnote{Д.--В., § 30, № 2.} хотя и не как основание, как в геометрической теории. Вследствие упомянутого в п. 6 соответствия между базисом и адъюнктными функциями и здесь можно установить некоторое соответствие вспомогательных средств в двух теориях. В итоге можно сказать, что образования понятий в двух теориях по существу совпадают, хотя, помимо формулировки, они различаются своим расположением. Обе теории возникли независимо друг от друга; геометрическая --- примерно на десятилетие раньше. Следует ещё заметить, что теорема Римана--Роха первоначально, у Римана и Роха, а также у Вейерштрасса, появляется как теорема о ранге: речь идёт о ранге матрицы \((\varphi_i(x_j))\), где \(\varphi\) пробегают \(p\) адъюнктных кривых порядка \(n-3\), а \(x\) --- точки остаточной группы.

\subsection*{8. Трансцендентные вопросы для алгебраических функций и алгебраических чисел}
На аналогии между трансцендентными вопросами для алгебраических функций и для алгебраических чисел особенно указывал Гильберт;\footnote{Mathematische Probleme, проблема 12. (\emph{Göttinger Nachrichten} 1900).} по существу это вопросы существования. Так, римановой постановке вопроса о существовании алгебраической функции при заданной римановой поверхности (следовательно, также при заданных точках ветвления) соответствует вопрос о существовании алгебраических числовых полей с заданной группой и дискриминантом.\footnote{Подсказанное этим рассмотрение алгебраических функций, исходя из группы, Р. Кёниг провёл как специальный случай более общих вопросов; ср. особенно: Riemannsche Funktionen- und Differentialsysteme in der Ebene (\emph{J. f. M.} т. 148) и \emph{Math. Ann.} 78, 79. На месте алгебраического поля у Кёнига стоит специальный конечный модуль относительно \(K(z)\): имеется только «свойство класса».} Но тогда как в случае алгебраических функций доказательство существования можно провести в общем виде, в области алгебраических числовых полей оно удалось только для специальных основных полей (рациональных чисел, квадратичных числовых полей) и только для абелевых групп.\footnote{Только простой случай групп, возникающих в уравнениях третьей и четвёртой степени, и несколько вполне специальных групп в уравнениях 8-й степени можно обработать для произвольных основных полей: G. Bucht, Über einige algebraische Körper 8. Grades (\emph{Arkiv f. Math., Astron. och Physik}, т. 6, № 30.) Ср. также E. Noether, Gleichungen mit vorgeschriebener Gruppe и цитированную там диссертацию Зайдельмана. (\emph{Math. Ann.} 78).} Речь идёт о теореме Кронекера, что всякое абелево числовое поле над рациональными числами может быть составлено из полей корней из единицы, и о соответствующих теоремах Фютера\footnote{\emph{Math. Ann.} 75.} и Гекке\footnote{\emph{Math. Ann.} 76.} для относительно абелевых числовых полей. Значит, искомые числовые поля фактически даются трансцендентным средством --- экспоненциальной функцией, модулярными функциями, --- в противоположность чистому доказательству существования для алгебраических функций.

Доказательству существования алгебраических функций предшествует существование всюду конечных интегралов. Аналогом этого является существование «особых первичных чисел» алгебраического числового поля, определённых относительно простого числа \(\ell\), чьи \(\ell\)-е корни относительно неразветвлены и дают первые строительные блоки для поля классов.\footnote{Furtwängler, Reziprozitätsgesetze für Potenzreste mit Primzahlexponenten in algebraischen Zahlkörpern. \emph{Math. Ann.} 67 и 72.} Это доказательство существования проводится с существенной помощью теоремы, что в числовом поле всегда существуют простые идеалы с заданными вычетными характерами; поэтому эту теорему можно считать аналогом краевой задачи, на которой основано существование всюду конечных интегралов.

Всюду конечные интегралы дают трансцендентный критерий для определения того, принадлежат ли два полигона одному классу (две группы точек одинакового числа корезидуальны), а именно через теорему Абеля, которая утверждает, что интегралы первого рода, распространённые на «корезидуальные пути», исчезают; или через соответствующую дифференциальную теорему вместе с её обращением.\footnote{Ср., например, формулировку у M. Noether (\emph{Ann.} 37). Поскольку теорема об остатке исторически выросла из теоремы Абеля, геометрическая теория говорит о «сумме» групп точек и линейных систем, следуя интегральным суммам; в противоположность «произведению» классов в арифметической теории.}

Аналогом теоремы Абеля для алгебраических числовых полей следует считать «закон взаимности для \(\ell\)-х степенных вычетов», который для самого поля или по меньшей мере для подходящего надполя даёт критерий того, что два идеала принадлежат одному классу. Его содержание можно сформулировать и так: в этом надполе особые первичные числа имеют один и тот же вычетный характер относительно всех идеалов одного класса. Последний факт верен и в самом поле, но там он не совпадает с законом взаимности.\footnote{Фуртвенглер, там же.}

Наконец, теория \(p\)-адических чисел Гензеля даёт аналог разложения алгебраических функций в степенные ряды; по этому поводу см. обзор Гензеля.
% END UNIT BODY
% END PRODUCER UNIT 076
\endgroup


\clearpage
\begingroup
\setcounter{footnote}{0}
% BEGIN PRODUCER UNIT 077
% Source: C:/Users/Floris/Documents/interlanguage/03_projects/noether/02_slavic_working_corpus/translations/paper15/russian/v001/Noether_Paper15_Russian_v001.tex
% Identity: 68347 B / SHA-256 47F691B67090127AF816AF7E59D6B97B52A8CDC531C1228A6938B141BC00E42A
\setlength{\parindent}{1.25em}
\setlength{\parskip}{0pt}
\emergencystretch=30em
\allowdisplaybreaks
\raggedbottom
\providecommand{\p}{\partial}
\providecommand{\frH}{\mathfrak H}
\providecommand{\frK}{\mathfrak K}
\providecommand{\frM}{\mathfrak M}
\providecommand{\frN}{\mathfrak N}
\providecommand{\frG}{\mathfrak G}
\providecommand{\frS}{\mathfrak S}
\providecommand{\frD}{\mathfrak D}
\providecommand{\frf}{\mathfrak f}
\providecommand{\frc}{\mathfrak c}
\providecommand{\fra}{\mathfrak a}
\providecommand{\frb}{\mathfrak b}
\providecommand{\frp}{\mathfrak p}
\providecommand{\frr}{\mathfrak r}
\providecommand{\frm}{\mathfrak m}
\providecommand{\calA}{\mathcal A}
\providecommand{\calB}{\mathcal B}
\providecommand{\calN}{\mathcal N}
\providecommand{\tmod}[1]{\;(\operatorname{mod} #1)}
% BEGIN UNIT BODY
% Unit: Paper 15. Direct Russian translation from the German final-audited source slice.

\begin{center}
{\Large\bfseries 15. Конечность системы целочисленных инвариантов бинарных форм\par}
\vspace{1em}
\emph{Nachrichten von der Gesellschaft der Wissenschaften zu Göttingen} 1919, с. 138--156
\end{center}

\vspace{2em}
\begin{center}
Представлено Ф. Клейном на заседании 27 марта 1919 года.
\end{center}

Ниже, как расширение известной теоремы о конечности, доказывается следующая теорема:
\emph{все целые целочисленные инварианты бинарной системы основных форм являются целыми рациональными целочисленными функциями конечного числа из них.}
При этом, как обычно, под «инвариантом» всегда следует понимать целый рациональный инвариант основной формы или системы основных форм, а именно инвариант относительно преобразования коэффициентов, индуцированного группой всех линейных преобразований.

Доказательство опирается на целочисленное уточнение доказательства конечности Мертенса--Гильберта.\footnote{F. Mertens, Wiener Sitzungsberichte Jan. 1889; D. Hilbert, \emph{Math. Ann.} т. 33 (1889) [датировано мартом 1888]. Оба доказательства, возникшие независимо друг от друга вслед за Mertens, Crelle т. 100, совпадают по содержанию.}
Это доказательство конечности можно охарактеризовать так. Основные формы разлагают на их линейные множители, точнее: каждой основной форме сопоставляют произведение линейных форм, вследствие чего каждый инвариант \(I\) переходит в инвариант этих линейных форм, значит в функцию однородненных корней. Эта функция должна удовлетворять трем условиям: 1) условию инвариантности, 2) условиям веса, 3) условиям симметрии. Условия 2) и 3) необходимы для того, чтобы \(I\) стала целой и рациональной по коэффициентам основных форм. Условие 1) удовлетворяется тем, что \(I\) представляют как целую рациональную функцию от определителей линейных форм; условие 2) тем, что вместо отдельных определителей берут в качестве новых аргументов некоторые произведения их степеней. Условие 3) теперь просто говорит, что \(I\) становится целой рациональной симметрической функцией рядов величин, и что, обратно, всякая такая симметрическая функция рядов величин становится инвариантом основных форм. Следовательно, элементарные симметрические функции этих рядов величин дают полную систему.

Если теперь рассматривать только целочисленные инварианты, то условие 1) можно удовлетворить так же, как выше, как только для инвариантов линейных форм доказана более сильная теорема: \emph{все целочисленные инварианты бинарных линейных форм являются целыми целочисленными функциями определителей этих линейных форм}. Что это действительно так, я показываю в § 3, причем в формулировке, что совокупность всех соотношений между этими определителями по модулю любого простого числа совпадает с совокупностью всех тождественных соотношений; то есть модуль, соответствующий этим соотношениям, остается простым модулем по модулю каждого простого числа. Это доказывается нормированием классов вычетов относительно модуля квадратичных форм, соответствующих квадратичным соотношениям между определителями; эти квадратичные формы при этом оказываются целочисленным базисом модуля, притом по модулю каждого простого числа, тогда как прежде они были известны как базис модуля лишь без учета целочисленности.

Для удовлетворения условия 2) не требуется никакого изменения, поскольку введение произведений степеней не влияет на числовые коэффициенты. Условие 3) теперь говорит, что \(n_1!\cdots n_\mu! I\) (где \(n_1,\ldots,n_\mu\) означают степени отдельных основных форм) становится целочисленной симметрической функцией рядов величин, и что, обратно, всякая целочисленная симметрическая функция этих рядов становится целочисленным инвариантом. Но, как легко показать, эти целочисленные симметрические функции рядов имеют целочисленный базис, для которого одних элементарных функций уже недостаточно. Применение гильбертовой теоремы о целочисленном базисе модуля дает отсюда также целочисленный базис для самих \(I\).

В § 1 собраны специальные теоремы конечности, соответствующие условиям 1), 2), 3); из них затем в § 2 следует доказательство конечности. Специальные теоремы конечности, соответствующие условию 3), одновременно дают конечность целочисленных инвариантов конечных групп целочисленных или алгебраически-целых подстановок, как кратко показано в § 4.

\section*{§ 1. Специальные теоремы конечности}

Сначала приведу следующие специальные теоремы конечности; они будут служить для удовлетворения упомянутых во введении условий инвариантности, веса и симметрии.

\textbf{I.} Каждый целочисленный инвариант системы бинарных линейных форм является целой рациональной целочисленной функцией определителей этих линейных форм (доказательство будет дано в § 3).

\textbf{II.} Каждая система \(\frS\) произведений степеней от \(n\) переменных с неотрицательными показателями, которая вместе с любыми двумя произведениями степеней \(f\) и \(g\) содержит также их частное, если это частное является целым по переменным, имеет конечный мультипликативный базис \(f_1,\ldots,f_k\) из функций из \(\frS\), такой что для всякого \(f\) из \(\frS\) имеется представление
\[
        f=f_1^{\lambda_1}\cdots f_k^{\lambda_k},
\]
с неотрицательными показателями \(\lambda_1,\ldots,\lambda_k\).

Действительно, по гильбертовой теореме о базисе модуля существует конечное число функций из \(\frS\), \(f_1,\ldots,f_k\), каждая из которых действительно содержит \(x\), такое что всякое непостоянное \(f\) из \(\frS\) принадлежит порожденному ими модулю:
\[
        f\equiv0\pmod{(f_1,\ldots,f_k)}.
\]
Отсюда следует, что для каждого \(f=x_1^{i_1}\cdots x_n^{i_n}\) существует по крайней мере одно \(f_\nu=x_1^{\nu_1}\cdots x_n^{\nu_n}\), такое что \(\nu_1\le i_1,\ldots,\nu_n\le i_n\). В представлении
\[
        f=x_1^{i_1-\nu_1}\cdots x_n^{i_n-\nu_n} f_\nu = A f_\nu
\]
множитель \(A\) по предположению принадлежит \(\frS\), но имеет меньшую степень по переменным, чем \(f\), так что конечное повторение этого рассуждения дает утверждение. Чтобы представление сохранялось и для \(f=1\), нужно лишь положить все показатели \(\lambda\) равными нулю.\footnote{Это представление можно доказать и непосредственно, а из него обратно снова следует гильбертова теорема: ср. P. Gordan, Neuer Beweis des Hilbertschen Satzes über homogene Funktionen, Gött. Nachr. 1899. Другое прямое доказательство этого представления, из которого затем также выводится теорема II, см. у A. Ostrowski, Über die Existenz einer endlichen Basis bei gewissen Funktionensystemen, \emph{Math. Ann.} 78 (1916), § 2,1 и § 3,2. Специальный случай теоремы II, используемый в § 2, где показатели пробегают все неотрицательные решения системы диофантовых уравнений, уже содержится у Gordan--Kerschensteiner, \emph{Vorlesungen über Invariantentheorie}, т. I, с. 199.}

\textbf{III.} Все целые целочисленные симметрические функции от \(n\) рядов величин являются целыми целочисленными функциями конечного числа из них.

Пусть \(\tau_1(x),\ldots,\tau_n(x)\) обозначают элементарные симметрические функции неопределенных \(x_1,\ldots,x_n\). Тогда для каждого показателя \(k\) имеется тождество
\[
        x_i^k = G_0^{(k)}(\tau)+x_iG_1^{(k)}(\tau)+\cdots+x_i^{n-1}G_{n-1}^{(k)}(\tau),
        \qquad (i=1,2,\ldots,n),
\]
где \(G_0^{(k)},\ldots,G_{n-1}^{(k)}\) означают целые целочисленные функции от \(\tau(x)\). Пусть \(x_i,y_i,\ldots,z_i\) -- элементы \(i\)-го ряда; тогда посредством этих тождеств и соответствующих тождеств для \(y,\ldots,z\) все специальные (однорядные) симметрические функции
\[
        \sum_{i=1}^n x_i^{a}y_i^{b}\cdots z_i^{c}
\]
становятся целыми целочисленными функциями конечного числа функций
\[
        \sum_{i=1}^n x_i^{a}y_i^{b}\cdots z_i^{c}
        \qquad
        (0\le a<n,\;0\le b<n,\;\ldots,\;0\le c<n)
\]
и элементарных симметрических функций \(\tau(x),\tau(y),\ldots,\tau(z)\).\footnote{Hilbert в указанном месте включает в базис вместо \(\tau(x),\ldots,\tau(z)\) степенные суммы; благодаря этому базис также состоит только из однорядных функций, но целочисленность представления теряется.}
Если имеется случай наиболее общих симметрических функций \(n\) рядов, который ниже не встречается, то так же видно, что функции
\[
        \sum x_{i_1}^{a_1}y_{i_1}^{b_1}\cdots z_{i_1}^{c_1}
              x_{i_2}^{a_2}y_{i_2}^{b_2}\cdots z_{i_2}^{c_2}
              \cdots
              x_{i_n}^{a_n}y_{i_n}^{b_n}\cdots z_{i_n}^{c_n}
\]
\[
        (0\le a_\nu<n,\;0\le b_\nu<n,\;\ldots,\;0\le c_\nu<n),
\]
где \(i_1,i_2,\ldots,i_n\) пробегают все перестановки, вместе с \(\tau(x),\tau(y),\ldots,\tau(z)\) образуют целочисленный базис.

\textbf{IV.} Пусть \(F_1(x),\ldots,F_k(x)\) являются целыми целочисленными функциями от \(x_1,\ldots,x_n\), а \(a\) -- фиксированное целое число. Тогда все функции
\[
        \frac{G(F)}{a},
\]
где \(G\) означает целую целочисленную функцию, которые становятся целочисленными по \(x\), являются целыми целочисленными функциями конечного числа из них.

Действительно, по гильбертовой теореме о целочисленном базисе модуля для каждого такого \(G(F)/a\) имеется представление
\[
 \frac{G(F)}{a}
 =
 A_1(F_1\ldots F_k)\frac{G_1(F)}{a}
 +\cdots+
 A_\rho(F_1\ldots F_k)\frac{G_\rho(F)}{a},
\]
где \(A_1,\ldots,A_\rho\) означают целые целочисленные функции от \(F\), а
\[
        \frac{G_1(F)}{a},\ldots,\frac{G_\rho(F)}{a}
\]
принадлежат системе. Поэтому функции
\[
        F_1=\frac{aF_1}{a},\ldots,F_k=\frac{aF_k}{a};
        \quad
        \frac{G_1(F)}{a},\ldots,\frac{G_\rho(F)}{a}
\]
образуют требуемый базис.

\setcounter{equation}{0}
\section*{§ 2. Конечность целочисленных бинарных инвариантов}

Пусть \(x,y\) -- бинарные переменные, подвергаемые преобразованию с неопределенными коэффициентами
\begin{equation}
        x=s_{11}x'+s_{12}y',\qquad y=s_{21}x'+s_{22}y',
\end{equation}
и пусть \(f\) -- бинарная основная форма, коэффициенты которой являются неопределенными. Тогда из (1)
\begin{equation}
        f=a_0x^n+a_1x^{n-1}y+\cdots+a_ny^n
        =a_0'x'^n+a_1'x'^{\,n-1}y'+\cdots+a_n'y'^n.
\end{equation}
Под целочисленным инвариантом \(f\), как известно, понимают целую целочисленную функцию \(I(a)\), такую что
\begin{equation}
        I(a')=|s_{ik}|^\rho I(a),\qquad \text{(тождественно по \(a,s_{ik}\))};
\end{equation}
соответственно определяются целочисленные инварианты системы основных форм, причем в определение можно включить также однородность по коэффициентам отдельных основных форм, поскольку из таких по отдельности однородных инвариантов остальные составляются целочисленно-линейно.

Наряду с (2) рассмотрим теперь еще специальную бинарную форму \(g\), а именно произведение \(n\) линейных форм, коэффициенты которых снова являются неопределенными:\footnote{Если, как часто принято в теории инвариантов, записывать \(f\) с биномиальными коэффициентами: \(f=b_0x^n+\binom n1b_1x^{n-1}y+\cdots+b_ny^n\), то область всех целочисленных инвариантов \(I(b)\) расширяется по сравнению с областью \(I(a)\), и используемые здесь выводы о конечности перестают действовать; \(I(b)\) уже не будет целочисленной функцией корней.}
\begin{align}
        g&=(\alpha_1x+\beta_1y)(\alpha_2x+\beta_2y)\cdots(\alpha_nx+\beta_ny)\notag\\
         &=\sigma_0(\alpha,\beta)x^n+\sigma_1(\alpha,\beta)x^{n-1}y+\cdots+\sigma_n(\alpha,\beta)y^n\notag\\
         &=(\alpha_1'x'+\beta_1'y')(\alpha_2'x'+\beta_2'y')\cdots(\alpha_n'x'+\beta_n'y')\notag\\
         &=\sigma_0'x'^n+\sigma_1'x'^{\,n-1}y'+\cdots+\sigma_n'y'^n,
\end{align}
где \(\sigma_0(\alpha,\beta),\sigma_1(\alpha,\beta),\ldots,\sigma_n(\alpha,\beta)\) обозначают однородненные элементарные симметрические функции, а преобразованные коэффициенты \(\sigma_i'\) равны \(\sigma_i(\alpha',\beta')\).

В силу алгебраической независимости этих элементарных функций каждому соотношению
\begin{equation}
        G(\sigma_0,\sigma_1,\ldots,\sigma_n)=0
        \quad [\text{тождественно по \(\alpha,\beta\)}]
\end{equation}
соответствует еще одно:
\begin{equation}
        G(\sigma_0,\sigma_1,\ldots,\sigma_n)=0
        \quad [\text{тождественно по \(\sigma_0,\ldots,\sigma_n\)}],
\end{equation}
и тем самым для неопределенных \(a_0,\ldots,a_n\):
\begin{equation}
        G(a_0,a_1,\ldots,a_n)=0
        \quad [\text{тождественно по \(a_0,\ldots,a_n\)}].
\end{equation}

Теперь каждому целочисленному инварианту \(I(a)\) формы \(f\), посредством подстановки \(a_i=\sigma_i(\alpha,\beta)\), соответствует целочисленный инвариант \(I(\sigma)=H(\alpha,\beta)\) формы \(g\), где \(H\) становится целочисленной функцией от \(\alpha,\beta\). Для этого инварианта соотношение (3) переходит в
\begin{equation}
        I(\sigma')=|s_{ik}|^\rho I(\sigma)
        \quad [\text{тождественно по \(\sigma,s_{ik}\)}],
\end{equation}
а из \(\sigma'=\sigma(\alpha',\beta')\) по (8) следует
\begin{equation}
        H(\alpha',\beta')=|s_{ik}|^\rho H(\alpha,\beta)
        \quad [\text{тождественно по \(\alpha,\beta;s_{ik}\)}];
\end{equation}
следовательно, \(I(\sigma)=H(\alpha,\beta)\) становится целочисленным инвариантом \(n\) линейных форм \((\alpha_i x+\beta_i y)\). Пусть, обратно, \(H(\alpha,\beta)\) -- целочисленный инвариант этих \(n\) линейных форм, то есть удовлетворяет (9), и вместе с тем \(H(\alpha,\beta)\) равна целочисленной функции \(I(\sigma)\) от \(\sigma\). Так как тогда также
\[
        H(\alpha',\beta')=I(\sigma(\alpha',\beta'))=I(\sigma')
\]
то (9) можно записать и в виде
\begin{equation*}
        I(\sigma')=|s_{ik}|^\rho I(\sigma)
        \quad [\text{тождественно по \(\alpha,\beta;s_{ik}\)}],\tag{9a}
\end{equation*}
и отсюда по (5) и (6) следует также (8), а по (7) также (3). Поэтому каждому целочисленному инварианту \(I(a)\) соответствует целочисленный инвариант \(H(\alpha,\beta)=I(\sigma)\) \(n\) линейных форм, и наоборот. Кроме того, если \(I_1(a),\ldots,I_k(a)\) обозначают целочисленный базис всех \(I(a)\), то \(H_1(\alpha,\beta),\ldots,H_k(\alpha,\beta)\) становится целочисленным базисом всех \(H(\alpha,\beta)\), которые вместе с тем являются целочисленными функциями \(I(\sigma)\); и здесь также по (5), (6), (7) имеет место обратное утверждение. Значит, для доказательства существования целочисленного базиса всех \(I(a)\) достаточно доказать его для всех целочисленных инвариантов \(H(\alpha,\beta)\), которые одновременно являются целочисленными функциями \(I(\sigma)\); а базис
\[
        H_1(\alpha,\beta)=I_1(\sigma),\ldots,H_k(\alpha,\beta)=I_k(\sigma)
\]
дает для \(I(a)\) базис \(I_1(a),\ldots,I_k(a)\).

Пусть теперь \(H(\alpha,\beta)=I(\sigma)\) пробегает совокупность всех целочисленных инвариантов (4). Поскольку по (3) всякий инвариант становится однородным по \(\sigma\), а по (4) каждая \(\sigma\) однородна и линейна по коэффициентам каждой линейной формы, \(H(\alpha,\beta)\) однородна по каждому ряду \(\alpha_i,\beta_i\) и в каждом ряду имеет одинаковую размерность; кроме того, она, конечно, является симметрической функцией \(n\) рядов. Обратно, всякая целочисленная симметрическая функция \(H(\alpha,\beta)\) \(n\) рядов, по отдельности однородная и одинаковой размерности в каждом ряду, по основной теореме о симметрических функциях является целой целочисленной функцией от \(\sigma\); следовательно, по сказанному выше она является инвариантом \(I(\sigma)\), как только \(H(\alpha,\beta)\) инвариантна. Поэтому совокупность целочисленных инвариантов (4) характеризуется как совокупность целочисленных функций \(H(\alpha,\beta)\) со следующими свойствами:
\[
\begin{array}{ll}
1) & \text{целочисленный инвариант \(n\) линейных форм \(\alpha_i x+\beta_i y\);}\\
2) & \text{по отдельности однородна и одинаковой размерности в каждом из \(n\) рядов \(\alpha_i,\beta_i\);}\\
3) & \text{симметрична в \(n\) рядах \(\alpha_i,\beta_i\).}
\end{array}
\]
По специальной теореме конечности I, ввиду 1), \(H(\alpha,\beta)\) становится целой целочисленной функцией определителей
\[
        (ik)=\alpha_i\beta_k-\beta_i\alpha_k;
\]
следовательно
\[
        H(\alpha,\beta)=G((ik)).
\]
В этом представлении симметрию можно выразить и формально, применив к \(G((ik))\) все перестановки \(n\) рядов:
\begin{equation}
        n!\,H(\alpha,\beta)
        =
        G^{(1)}((ik))+G^{(2)}((ik))+\cdots+G^{(n!)}((ik))
        =G^*((ik)).
\end{equation}
Значит, \(G^*\) вместе с каждым произведением степеней \(P\) от \((ik)\) содержит также сумму \(Q=P^{(1)}+\cdots+P^{(n!)}\) по всем перестановкам \(P\); причем \(G^*\) становится целочисленной линейной комбинацией конечного числа \(Q\), \(G^*=\sum cQ\). Поскольку \(Q\) удовлетворяют условиям 1), 2), 3), то есть сами становятся инвариантами, достаточно рассматривать \(Q\) вместо \(G^*((ik))\).

Пусть теперь \(P\) пробегает систему \(\frS\) всех произведений степеней от \((ik)\), встречающихся в \(Q\), то есть удовлетворяющих условию 2). Если частное двух \(P\) является целым по \((ik)\), то оно также удовлетворяет условию 2), значит принадлежит \(\frS\). По теореме II существует конечное число таких \(P\), скажем \(P_1,\ldots,P_\rho\), такое что каждое \(P\) можно представить в форме
\[
        P=P_1^{\lambda_1}\cdots P_\rho^{\lambda_\rho}
\]
с неотрицательными показателями \(\lambda\). Применение всех перестановок дает отсюда
\[
        Q=
        P_1^{(1)\lambda_1}\cdots P_\rho^{(1)\lambda_\rho}
        +\cdots+
        P_1^{(n!)\lambda_1}\cdots P_\rho^{(n!)\lambda_\rho}.
\]
Следовательно, \(Q\) становится целочисленной симметрической функцией \(n!\) рядов величин, в каждом из которых \(\rho\) элементов:
\[
        P_1^{(1)}\ldots P_\rho^{(1)};
        \quad
        P_1^{(2)}\ldots P_\rho^{(2)};
        \quad \ldots;\quad
        P_1^{(n!)}\ldots P_\rho^{(n!)};
\]
поэтому по теореме III она является целой целочисленной функцией конечного числа симметрических функций этих рядов. Эти базисные функции сами являются величинами \(Q\) или элементарными симметрическими функциями \(n!\) элементов \(P_i^{(1)},P_i^{(2)},\ldots,P_i^{(n!)}\), следовательно удовлетворяют условиям 1) 2) 3) и становятся целочисленными инвариантами
\begin{equation}
        H_1(\alpha,\beta)=I_1(\sigma);\ldots;\quad
        H_l(\alpha,\beta)=I_l(\sigma).
\end{equation}
По (10), ввиду \(G^*=\sum cQ\), каждый \(H(\alpha,\beta)=I(\sigma)\) имеет форму
\[
        H(\alpha,\beta)=I(\sigma)=\frac1{n!}\Phi(I_1,\ldots,I_l),
\]
где \(\Phi\) означает целую целочисленную функцию; и обратно, \(\frac1{n!}\Phi\) становится целочисленным инвариантом, как только он становится целым по \(\alpha,\beta\). Поэтому по теореме IV к инвариантам (11) можно добавить еще конечное число \(I_{l+1}(\sigma),\ldots,I_k(\sigma)\) так, что
\[
        H_1(\alpha,\beta)=I_1(\sigma);\ldots;H_l(\alpha,\beta)=I_l(\sigma);
\]
\[
        H_{l+1}(\alpha,\beta)=I_{l+1}(\sigma);\ldots;H_k(\alpha,\beta)=I_k(\sigma)
\]
образуют целочисленный базис всех инвариантов \(H(\alpha,\beta)=I(\sigma)\), чем теорема для случая одной основной формы доказана.

Если речь идет о системе основных форм, то, как замечено выше, можно предполагать однородность инвариантов по коэффициентам каждой основной формы. Общие основные формы снова заменяют такими, которые соответственно являются произведениями линейных форм. Их инварианты тогда становятся: 1) инвариантами этих линейных форм; 2) по отдельности однородными по рядам, образованным из коэффициентов линейных форм, и одинаковой размерности в рядах каждой отдельной основной формы; 3) симметричными в рядах каждой отдельной основной формы. Обратно, каждому целочисленному инварианту всех линейных форм, удовлетворяющему этим условиям, снова соответствует целочисленный инвариант системы основных форм. Доказательство идет совершенно параллельно приведенному: к каждому произведению степеней \(P=P_1^{\lambda_1}\cdots P_\rho^{\lambda_\rho}\) отдельных определителей, удовлетворяющему условиям веса (теоремы I и II), надо применить все \(N=n_1!n_2!\cdots n_\mu!\) перестановок, где \(n_1,\ldots,n_\mu\) означают степени основных форм. Тем самым \(N\cdot I\) переходит в целочисленную симметрическую функцию \(N\) рядов, каждый из которых имеет \(\rho\) элементов; отсюда получается целочисленный базис для всех \(N\cdot I\), а следовательно, по теореме IV, целочисленный базис для всех \(I\).

\setcounter{equation}{0}
\section*{§ 3. Базис целочисленных инвариантов бинарных линейных форм}

Теперь остается добавить доказательство специальной теоремы конечности I. По известной теореме, впервые доказанной Clebsch, все целые рациональные инварианты \(n\) линейных форм \(\alpha_i x+\beta_i y\) являются целыми рациональными функциями определителей \((ik)\). В частности, все целочисленные инварианты этих линейных форм становятся целыми рациональными, но не обязательно целочисленными функциями этих определителей.

Ниже будет показано, что благодаря тождественным соотношениям между \((ik)\) это представление всегда возможно сделать целочисленным.

В формулах и на языке теории модулей это выражается так. Пусть \(\xi_{ik}\), где \(i<k\), -- неопределенные; тогда совокупность всех целочисленных многочленов \(F(\xi_{ik})\) со свойством, что \(F((ik))\) тождественно обращается в нуль по \(\alpha,\beta\), образует модуль \(\frM\), поскольку сумма двух таких \(F\), а также произведение такого \(F\) с произвольным целочисленным многочленом от \(\xi_{ik}\), снова принадлежит этой совокупности. В формулах: из
\[
        F((ik))=0\quad [\text{тождественно по \(\alpha,\beta\)}]
\]
следует
\[
        F(\xi_{ik})\equiv0\pmod{\frM}\quad [\text{тождественно по \(\xi_{ik}\)}],
\]
и наоборот.

Я теперь покажу, что совершенно аналогично имеет место следующее. Из
\[
        F((ik))\equiv0\pmod p\quad [\text{тождественно по \(\alpha,\beta\)}]
\]
\begin{equation}
        \text{следует}\qquad
        F(\xi_{ik})\equiv0\pmod{(\frM,p)}\quad [\text{тождественно по \(\xi_{ik}\)}],
\end{equation}
и обратно, причем для каждого простого числа \(p\). Тем самым \(\frM\) также тождествен модулю \(\frM_p\), соответствующему совокупности всех соотношений между \((ik)\) по модулю \(p\); этот модуль состоит из всех целочисленных многочленов \(F(\xi_{ik})\), для которых \(F((ik))\) по модулю \(p\) тождественно обращается в нуль по \(\alpha,\beta\).

Что (1) означает целочисленную представимость, видно так. (1) можно записать и в форме: из \(F((ik))=p\cdot H(\alpha,\beta)\) [тождественно по \(\alpha,\beta\)] следует
\[
        F(\xi_{ik})-p\cdot G(\xi_{ik})\equiv0\pmod{\frM}
        \quad [\text{тождественно по \(\xi_{ik}\)}],
\]
а значит, по определению \(\frM\),
\[
        F((ik))=p\cdot G((ik))\quad [\text{тождественно по \(\alpha,\beta\)}].
\]
\footnote{Этот последний вывод одновременно доказывает обратное направление (1).}
Следовательно, \(p\cdot H(\alpha,\beta)=p\cdot G((ik))\), или
\[
        H(\alpha,\beta)=G((ik)).
\]
Если, стало быть, дано представление
\[
        H(\alpha,\beta)
        =
        \frac1{p_1^{\lambda_1}\cdots p_\rho^{\lambda_\rho}}
        F((ik)),
\]
где \(p_1,\ldots,p_\rho\) означают простые числа, то отсюда также следует
\[
        H(\alpha,\beta)
        =
        \frac1{p_1^{\lambda_1-1}\cdots p_\rho^{\lambda_\rho}}
        G((ik)),
\]
и через конечное повторение получается целочисленность представления.

Для доказательства (1) обозначим через
\[
        f_{ik,rs}=\xi_{ik}\xi_{rs}-\xi_{ir}\xi_{ks}+\xi_{is}\xi_{kr},
        \qquad i<k<r<s
\]
формы от неопределенных \(\xi_{ik}\), соответствующие квадратичным соотношениям между \((ik)\):
\[
        (ik)(rs)-(ir)(ks)+(is)(kr)=0,
\]
а через \(\frN=(f_{ik,rs})\) -- порожденный ими модуль целочисленных многочленов. Из
\[
        F(\xi_{ik})\equiv0\pmod{\frN}
        \quad\text{следует значит:}\quad
        F((ik))=0\quad [\text{тождественно по \(\alpha,\beta\)}],
\]
\begin{equation}
\begin{array}{rcl}
        \text{из }F(\xi_{ik})&\equiv&0\pmod{(\frN,p)}\\
        \text{следует }F((ik))&\equiv&0\pmod p\quad[\text{тождественно по \(\alpha,\beta\)}].
\end{array}
\end{equation}
Я теперь покажу, посредством нормирования классов вычетов относительно \(\frN\), что имеет место и обратное к (2); тогда будет доказано \(\frN=\frM=\frM_p\), а вместе с тем и (1). При этом \(\frN=\frM\) получается как специальный случай \(\frN=\frM_p\) для \(p=0\); случай \(p=0\) всегда включен в дальнейшие рассуждения. Иначе говоря: \(f_{ik,rs}\) образуют целочисленный базис модуля \(\frM_p\), в частности также \(\frM_0=\frM\).

\subsection*{a) Случай двух или трех рядов \(\alpha_i,\beta_i\)}

Поскольку еще не встречаются четыре различных индекса, \(\frN\) является нулевым модулем. Поэтому надо показать, что из
\begin{equation}
\begin{aligned}
        &F((ik))\equiv0\pmod p\quad [\text{тождественно по \(\alpha,\beta\)}]\\
        &\Longrightarrow\quad
        F(\xi_{ik})\equiv0\pmod p\quad [\text{тождественно по \(\xi_{ik}\)}],
\end{aligned}
\end{equation}
чем (1) будет доказано при \(\frM=\frM_p=0\).

Поскольку условие \(F((ik))=p\cdot H(\alpha,\beta)\) должно выполняться отдельно для компонент, однородных в отдельных рядах, достаточно всюду предполагать, что \(F((ik))\) однороден в каждом отдельном ряду, а значит \(F(\xi_{ik})\) однороден в каждом индексе. В случае двух рядов имеется только один определитель \((12)\) и соответствующая неопределенная \(\xi_{12}\). Из принятой однородности получаем \(F=c\,\xi_{12}^\nu\). Из
\[
        c\,(12)^\nu\equiv0\pmod p\quad[\text{тождественно по \(\alpha,\beta\)}]
\]
ввиду \((12)\not\equiv0\pmod p\) [тождественно по \(\alpha,\beta\)] всегда следует \(c\equiv0\pmod p\); значит
\[
        F(\xi_{ik})\equiv0\pmod p\quad [\text{тождественно по \(\xi_{ik}\)}],
\]
чем доказано (3), а значит и (1).

В случае трех рядов имеются определители \((12),(13),(23)\), которым соответствуют неопределенные \(\xi_{12},\xi_{13},\xi_{23}\); поэтому
\[
        F=\sum c\,\xi_{12}^{\tau_{12}}\xi_{13}^{\tau_{13}}\xi_{23}^{\tau_{23}}.
\]
Пусть \(F\) имеет соответственно степени \(\sigma_1,\sigma_2,\sigma_3\) в индексах \(1,2,3\); тогда имеется система уравнений
\[
        \sigma_1=\tau_{12}+\tau_{13},\qquad
        \sigma_2=\tau_{12}+\tau_{23},\qquad
        \sigma_3=\tau_{13}+\tau_{23}
\]
с единственным обратным решением:
\[
        \tau_{12}=\frac12(\sigma_1+\sigma_2-\sigma_3),\quad
        \tau_{13}=\frac12(\sigma_1-\sigma_2+\sigma_3),\quad
        \tau_{23}=\frac12(-\sigma_1+\sigma_2+\sigma_3);
\]
следовательно, каждый по отдельности однородный \(F\) содержит не более одного члена,
\[
        F=c\,\xi_{12}^{\tau_{12}}\xi_{13}^{\tau_{13}}\xi_{23}^{\tau_{23}},
\]
откуда, как выше, следует выполнение (3), а значит и (1).

\subsection*{b) Случай 4 рядов \(\alpha_i,\beta_i\)}

Возникают шесть неопределенных \(\xi_{12},\xi_{13},\xi_{14},\xi_{23},\xi_{24},\xi_{34}\) и соответствующие определители; и \(\frN=(f_{12,34})\). В силу
\begin{equation}
        \xi_{14}\xi_{23}\equiv \xi_{13}\xi_{24}-\xi_{12}\xi_{34}\pmod{\frN}
\end{equation}
каждое произведение степеней от \(\xi_{ik}\) по модулю \(\frN\) конгруэнтно линейной целочисленной комбинации таких произведений степеней, в которых произведение \(\xi_{14}\xi_{23}\) не встречается. Поэтому для каждого многочлена \(F(\xi_{ik})\) имеем
\begin{equation}
        F(\xi_{ik})\equiv F_0(\xi_{ik})\pmod{\frN},
\end{equation}
где \(F_0\) не содержит произведения \(\xi_{14}\xi_{23}\); \(F_0\) будем называть нормированным многочленом класса вычета \(F\). При этом по отдельности однородному \(F\), поскольку (4) по отдельности однородно, снова соответствует по отдельности однородный \(F_0\).

Положу теперь
\begin{equation}
        F_0(\xi_{ik})=G(\xi_{ik})+\xi_{34}\,F_1(\xi_{ik}),
\end{equation}
где \(G\) не содержит произведения степеней, делимого на \(\xi_{34}\), и далее
\begin{equation}
        [G]_{\xi_{14}=\xi_{13}}=K(\xi_{12},\xi_{13},\xi_{23}).
\end{equation}
Из
\[
        F((ik))\equiv0\pmod p\quad[\text{тождественно по \(\alpha,\beta\)}]
\]
теперь по (5), с учетом (2), следует
\begin{equation}
        F_0((ik))\equiv0\pmod p\quad [\text{тождественно по \(\alpha,\beta\)}];
\end{equation}
а отсюда через специализацию \(\alpha_4=\alpha_3,\;\beta_4=\beta_3\), по (6) и (7),
\[
        K((ik))\equiv0\pmod p\quad [\text{тождественно по \(\alpha,\beta\)}].
\]
Поэтому по доказанному в пункте a)
\[
        K(\xi_{ik})\equiv0\pmod p\quad [\text{тождественно по \(\xi_{ik}\)}],
\]
и отсюда следует, как еще будет показано ниже,
\[
        G(\xi_{ik})\equiv0\pmod p\quad [\text{тождественно по \(\xi_{ik}\)}].
\]
Значит по (6)
\[
        F_0(\xi_{ik})\equiv\xi_{34}F_1(\xi_{ik})\pmod p\quad [\text{тождественно по \(\xi_{ik}\)}],
\]
а ввиду \((34)\not\equiv0\pmod p\) [тождественно по \(\alpha,\beta\)] из (8) также следует
\[
        F_1((ik))\equiv0\pmod p\quad [\text{тождественно по \(\alpha,\beta\)}],
\]
где \(F_1(\xi_{ik})\) снова нормирован и имеет меньшую степень по \(\xi_{34}\). Конечное повторение процедуры приводит к
\[
        F_0(\xi_{ik})\equiv\xi_{34}^{\lambda}F_\lambda(\xi_{ik})\pmod p,
\]
где \(F_\lambda\) имеет нулевую степень по \(\xi_{34}\), значит по только что доказанному обращается в нуль по модулю \(p\). Таким образом получаем
\begin{equation}
\begin{array}{rcl}
        F_0(\xi_{ik})&\equiv&0\pmod p\quad [\text{тождественно по \(\xi_{ik}\)}],\quad\text{или}\\
        F(\xi_{ik})&\equiv&0\pmod{(\frN,p)}\quad [\text{тождественно по \(\xi_{ik}\)}],
\end{array}
\end{equation}
чем доказано обратное к (2), а значит и (1); одновременно в более сильном утверждении для \(F_0(\xi_{ik})\).\footnote{Это более сильное утверждение означает, что каждый класс вычетов содержит только один нормированный многочлен.}

Остается добавить, что из \(K(\xi_{ik})\equiv0\pmod p\) следует также \(G(\xi_{ik})\equiv0\pmod p\); то есть из \(G(\xi_{ik})\not\equiv0\pmod p\) следует и \(K(\xi_{ik})\not\equiv0\pmod p\). Только для этой части индукционного шага нужно нормирование. Поскольку при подстановке \(\xi_{14}=\xi_{13}\) отдельные произведения степеней из \(G\) не исчезают, обращение \(K\) в нуль могло бы произойти лишь потому, что различным произведениям степеней в \(G\) соответствуют одинаковые в \(K\); а это, в свою очередь, возможно только если эти произведения степеней различаются лишь перестановкой индексов 3 и 4. Пусть такое, по предположению нормированное, произведение степеней, например, имеет вид
\[
        \xi_{12}^{\tau_{12}}\xi_{13}^{\tau_{13}}\xi_{14}^{\tau_{14}}\xi_{24}^{\tau_{24}}.
\]
В силу предполагаемой отдельной однородности по каждому индексу замене 3 на 4 в одном \(\xi_{ik}\) должна соответствовать замена 4 на 3 в другом; соответствующий член содержал бы вместо \(\xi_{13}\xi_{24}\) произведение \(\xi_{14}\xi_{23}\), что исключено нормированием. Совершенно так же против
\[
        \xi_{12}^{\tau_{12}}\xi_{13}^{\tau_{13}}\xi_{14}^{\tau_{14}}\xi_{24}^{\tau_{24}}
\]
мог бы сократиться лишь член, содержащий множитель \(\xi_{14}\xi_{23}\), что снова исключено нормированием. Следовательно, из \(K(\xi_{ik})\equiv0\pmod p\) следует также \(G(\xi_{ik})\equiv0\pmod p\).

\subsection*{c) Случай \(n\) рядов \(\alpha_i,\beta_i\)}

Случай \(n\) рядов \(\alpha_i,\beta_i\) теперь устраняется полной индукцией, в точном обобщении только что рассмотренного случая четырех рядов.\footnote{Доказательство остается верным и для \(n=3\).}
Сначала я снова покажу, что для каждого многочлена \(F\) имеет место конгруэнция
\begin{equation}
        F(\xi_{ik})\equiv F_0(\xi_{ik})\pmod{\frN}\quad [\text{тождественно по \(\xi_{ik}\)}],
\end{equation}
где \(\frN=(f_{ik,rs})\) и где \(F_0(\xi_{ik})\) нормирован. Нормирование должно состоять в следующем: если \(i,k,r,s\) -- четыре различных индекса и \(i<k<r<s\), то \(F_0\) не должен содержать произведение \(\xi_{is}\xi_{kr}\); то есть индексы одного \(\xi\) не должны оба лежать между индексами другого. Это нормирование еще не накладывает ограничения в случае двух и трех рядов, где каждый многочлен нормирован; и в случае четырех рядов, как показано, для каждого класса вычетов существует нормированный многочлен. Поэтому существование нормированного многочлена для каждого класса вычетов можно считать доказанным для \((n-1)\) рядов.

Произвольное произведение степеней из \(F\) запишем в форме
\[
        \xi_{i_1n}\cdots\xi_{i_\nu n}\,P,
\]
где \(P\) уже не содержит индекс \(n\). По предположению \(P\) конгруэнтно по модулю \(\frN\) некоторому целочисленному нормированному многочлену \(P_0\):
\[
        \xi_{i_1n}\cdots\xi_{i_\nu n}P
        \equiv
        \xi_{i_1n}\cdots\xi_{i_\nu n}P_0
        \pmod{\frN}.
\]
Если теперь \(P_0\) не содержит такого \(\xi_{rs}\), что хотя бы для одного индекса \(i_\lambda\) эти \(r,s\) лежат между \(i_\lambda\) и \(n\) (\(i_\lambda<r<s<n\)), то нормирование уже достигнуто, поскольку для произведения любых двух \(\xi\) из \(P_0\) условия выполнены. В противном случае пусть \(i_\lambda\) -- такой индекс, что \(i_\lambda<r<s<n\), а \(P_0^*\) -- произведение степеней из \(P_0\), содержащее \(\xi_{rs}\); из-за
\[
        \xi_{i_\lambda n}\xi_{rs}\equiv
        \xi_{i_\lambda s}\xi_{rn}-\xi_{i_\lambda r}\xi_{sn}
        \pmod{\frN}
\]
получаем
\begin{align*}
 &\xi_{i_1n}\cdots\xi_{i_\lambda n}\cdots\xi_{i_\nu n}P_0^*\\
 &\equiv
 \xi_{i_1n}\cdots\xi_{rn}\cdots\xi_{i_\nu n}Q
 +
 \xi_{i_1n}\cdots\xi_{sn}\cdots\xi_{i_\nu n}R
 \pmod{\frN}\\
 &\equiv
 \xi_{i_1n}\cdots\xi_{rn}\cdots\xi_{i_\nu n}Q_0
 +
 \xi_{i_1n}\cdots\xi_{sn}\cdots\xi_{i_\nu n}R_0
 \pmod{\frN},
\end{align*}
где \(Q_0\) и \(R_0\) снова обозначают нормированные многочлены классов вычетов \(Q\) и \(R\); \(Q_0\) и \(R_0\) существуют по предположению, так как \(Q\) и \(R\) не содержат индекс \(n\), и они заведомо являются целочисленными многочленами, поскольку \(Q\) и \(R\) -- просто произведения степеней. Из \(i_\lambda<r<s<n\) следует также
\[
 ((n-i_1)+\cdots+(n-r)+\cdots+(n-i_\nu))
 <
 ((n-i_1)+\cdots+(n-i_\lambda)+\cdots+(n-i_\nu)),
\]
\[
 ((n-i_1)+\cdots+(n-s)+\cdots+(n-i_\nu))
 <
 ((n-i_1)+\cdots+(n-i_\lambda)+\cdots+(n-i_\nu)).
\]
Таким образом, при переходе от \(P_0^*\) к \(Q_0,R_0\) эта сумма разностей уменьшилась; а так как эта сумма необходимо \(\ge \sigma\), процедура для каждого произведения степеней \(P_0^*\) должна завершиться после конечного числа шагов. Следовательно,
\[
        F(\xi_{ik})
        =
        \sum \xi_{i_1n}\cdots\xi_{i_\nu n}P^{(1)}
        \equiv
        \sum \xi_{j_1n}\cdots\xi_{j_\sigma n}S_0^{(j)}
        =
        F_0(\xi_{ik})\pmod{\frN},
\]
где целочисленный многочлен \(F_0(\xi_{ik})\) необходимо нормирован, поскольку иначе сумму разностей можно было бы еще уменьшить; этим доказано (10). Далее, так как по предположению при \((n-1)\) индексах каждому по отдельности однородному \(F\) соответствует по отдельности однородный нормированный \(F_0\), указанная процедура показывает, что это верно и для \(n\) индексов; поэтому ниже речь идет только о по отдельности однородных многочленах.

Положу теперь, аналогично пункту b),
\begin{equation}
        F_0(\xi_{ik})=G(\xi_{ik})+\xi_{n-1,n}F_1(\xi_{ik}),
\end{equation}
\footnote{Тождество, соответствующее (10) и (11),
        \[
          F((ik))=F_0((ik))=G((ik))+(n-1,n)\cdot F_1((ik)),
        \]
        является аналогом разложения Клебша--Гордана в ряд, если два ряда \((n-1)\) и \((n)\) заменить рядами переменных \(x,y\), то есть \((n-1,n)\) заменить на \((xy)\). Но тогда как у Клебша--Гордана речь идет о разложении по полярам, нормирование \(F_0\) соответствует «начальным членам» поляров, и тем самым позволяет добиться целочисленности, которой у Клебша--Гордана нет. Известное доказательство того, что без учета целочисленности \(f_{ik,rs}\) образуют базис \(\frM\), проводится именно посредством разложения Клебша--Гордана в ряд (Gordan--Kerschensteiner, \emph{Vorlesungen über Invariantentheorie}, с. 132).}
где \(G\) не содержит произведения степеней, делимого на \(\xi_{n-1,n}\), и далее
\begin{equation}
        [G]_{\xi_{in}=\xi_{i,n-1}}=K(\xi_{ik});
\end{equation}
как показывает определение нормирования, наряду с \(G\) нормирован и \(K\).

Тогда между \(G\) и \(K\) снова имеется взаимно однозначное соответствие; ведь только таким различным произведениям степеней из \(G\) могут соответствовать одинаковые в \(K\), которые отличаются лишь перестановкой индексов \(n\) и \((n-1)\). Пусть, например, \(G\) имеет степень \(\rho\) по индексу \((n-1)\), степень \(\sigma\) по индексу \(n\), и пусть
\[
        \chi_1\le\chi_2,\ldots,\le\chi_\rho\le\chi_{\rho+1},\ldots,\le\chi_{\rho+\sigma}
\]
являются индексами, соединенными в одном произведении степеней из \(G\) с \((n-1)\) и \(n\); их число должно быть \(\rho+\sigma\), поскольку по (11) неопределенная \(\xi_{n-1,n}\) в \(G\) не встречается. Тогда благодаря нормированию произведение степеней имеет вид
\[
        \xi_{\chi_1,n-1}\cdots\xi_{\chi_\rho,n-1}
        \xi_{\chi_{\rho+1},n}\cdots\xi_{\chi_{\rho+\sigma},n}\,P(\xi_{ik}),
\]
где \(P(\xi_{ik})\) уже не содержит индексов \(n\) и \((n-1)\). Всякая перестановка \(n\) с \((n-1)\) дает при \(i\ne j\) множитель \(\xi_{in}\xi_{j,n-1}\) с \(i<j\); значит \(j,n-1\) лежит между \(i\) и \(n\), и член не был бы нормирован. Поэтому при фиксированных \(\chi_1,\chi_2,\ldots,\chi_{\rho+\sigma}\) и фиксированном \(P\) многочлен \(G\) содержит только одно произведение степеней; тем самым показано, что различным произведениям степеней из \(G\) соответствуют различные в \(K\).
\begin{equation}
        \text{Из }K(\xi_{ik})\equiv0\pmod p\text{ следует поэтому }G(\xi_{ik})\equiv0\pmod p.
\end{equation}

Из
\[
        F((ik))\equiv0\pmod p\quad [\text{тождественно по \(\alpha,\beta\)}]
\]
теперь по (10), с учетом (2), следует
\begin{equation}
        F_0((ik))\equiv0\pmod p\quad [\text{тождественно по \(\alpha,\beta\)}],
\end{equation}
а отсюда через специализацию \(\alpha_n=\alpha_{n-1},\;\beta_n=\beta_{n-1}\), по (11) и (12),
\begin{equation}
        K((ik))\equiv0\pmod p\quad [\text{тождественно по \(\alpha,\beta\)}].
\end{equation}
Так как \(K(\xi_{ik})\) нормирован и содержит только \((n-1)\) индексов, по доказанному в пункте a) и в пункте b) (9), ввиду (15), можно предположить
\[
        K(\xi_{ik})\equiv0\pmod p\quad [\text{тождественно по \(\xi_{ik}\)}],
\]
и отсюда по (13) следует
\[
        G(\xi_{ik})\equiv0\pmod p\quad [\text{тождественно по \(\xi_{ik}\)}].
\]
По (11) получаем
\[
        F_0(\xi_{ik})\equiv\xi_{n-1,n}F_1(\xi_{ik})\pmod p\quad [\text{тождественно по \(\xi_{ik}\)}],
\]
а ввиду \((n,n-1)\not\equiv0\pmod p\) [тождественно по \(\alpha,\beta\)] из (14) также следует
\[
        F_1((ik))\equiv0\pmod p\quad [\text{тождественно по \(\alpha,\beta\)}],
\]
где \(F_1(\xi_{ik})\) снова нормирован и имеет меньшую степень по \(\xi_{n-1,n}\), чем \(F_0\). Конечное повторение процедуры приводит к
\[
        F_0(\xi_{ik})\equiv\xi_{n-1,n}^{\lambda}F_\lambda(\xi_{ik})\pmod p,
\]
где \(F_\lambda\) имеет нулевую степень по \(\xi_{n-1,n}\), значит по только что доказанному обращается в нуль по модулю \(p\). Таким образом получаем
\begin{equation}
\begin{array}{rcl}
        F_0(\xi_{ik})&\equiv&0\pmod p\quad [\text{тождественно по \(\xi_{ik}\)}]\quad \text{или}\\
        F(\xi_{ik})&\equiv&0\pmod{(\frN,p)}\quad [\text{тождественно по \(\xi_{ik}\)}].
\end{array}
\end{equation}
Этим доказано обратное к (2), а следовательно и (1); одновременно для \(F_0(\xi_{ik})\) доказано более сильное утверждение, использованное в индукционном шаге, и в случае \(n\) индексов, то есть вообще.

\setcounter{equation}{0}
\section*{§ 4. Конечность целочисленных инвариантов конечных групп}

Элементарное доказательство конечности, данное в моей заметке «Теорема конечности инвариантов конечных групп»,\footnote{\emph{Math. Ann.} 77, с. 89 (1915).} посредством специальных теорем конечности III и IV из § 1 допускает уточнение до целочисленности; тогда как обычное доказательство, опирающееся на гильбертову теорему о базисе модуля, этой целочисленности не дает. Даже привлечение целочисленного базиса модуля дает только представление, в знаменателе которого появляется сколь угодно высокая степень целого числа \(h\), порядка группы.

Я всюду пользуюсь обозначениями указанной заметки. Пусть конечная группа \(\frH\) состоит из \(h\) линейных преобразований (с ненулевым определителем)\footnote{Достаточно было бы предположить, что \(\frH\) просто изоморфна абстрактной конечной группе; то есть если определитель преобразования одного из \(A\) обращается в нуль, то \(\frH\) должна быть подобна группе \(\begin{pmatrix}\frS&0\\0&0\end{pmatrix}\), где определители подстановок из \(\frS\) не обращаются в нуль. Инварианты \(\frH\) в этом случае совпадают с инвариантами \(\frS\), так что условие ненулевости определителей фактически не ограничивает общности.} \(A_1,\ldots,A_h\), причем \(A_k\) задает линейное преобразование
\[
        x_1^{(k)}=\sum_{\nu=1}^{n}a_{1\nu}^{(k)}x_\nu,
        \quad\ldots,\quad
        x_n^{(k)}=\sum_{\nu=1}^{n}a_{n\nu}^{(k)}x_\nu
\]
или сокращенно \((x^{(k)})=A_k(x)\). Итак, группа \(\frH\) переводит ряд \((x)\) с элементами \(x_1,\ldots,x_n\) в ряды \((x^{(k)})\) с элементами \(x_1^{(k)},\ldots,x_n^{(k)}\). При этом коэффициенты \(a_{i\nu}^{(k)}\) предполагаются алгебраическими целыми числами; это не ограничивает общности, поскольку по теореме J. Schur\footnote{Über Gruppen linearer Substitutionen mit Koeffizienten aus einem algebraischen Zahlkörper. Satz VII. \emph{Math. Ann.} 71, с. 555 (1912).} всякая конечная группа линейных преобразований подобна такой, для которой \(a_{i\nu}^{(k)}\) становятся алгебраическими целыми числами.

Под целочисленным (абсолютным) инвариантом группы будем понимать такую целую рациональную функцию от \(x_1,\ldots,x_n\), с коэффициентами в целых числах алгебраического числового поля \(\frK\), порожденного \(a_{i\nu}^{(k)}\), которая при применении \(A_1,\ldots,A_h\) остается тождественно неизменной; для такого инварианта \(f(x)\) имеем
\begin{equation}
        f(x)=f(x^{(1)})=\cdots=f(x^{(h)})
        =\frac1h\sum_{k=1}^{h}f(x^{(k)}).
\end{equation}
Обратно, всякая целочисленная симметрическая функция \(h\) рядов величин \((x^{(k)})\) также становится целочисленным инвариантом группы.

Пусть \(\omega_1,\omega_2,\ldots,\omega_\rho\) обозначает базис всех целых чисел из \(\frK\), так что каждый инвариант можно представить как
\[
        f(x)=\omega_1 f_1(x)+\cdots+\omega_\rho f_\rho(x),
\]
где \(f_1(x),\ldots,f_\rho(x)\) имеют рационально-целые коэффициенты. По (1) получаем
\[
        h\,f(x)
        =
        \omega_1\sum_{k=1}^{h}f_1(x^{(k)})
        +\cdots+
        \omega_\rho\sum_{k=1}^{h}f_\rho(x^{(k)}).
\]
Здесь отдельные суммы являются симметрическими функциями \(h\) рядов величин \((x^{(k)})\), причем с рационально-целыми коэффициентами; значит по теореме III они выразимы цело и рационально, с рационально-целыми коэффициентами, через конечное число целорационально-числовых симметрических функций этих рядов. Эти базисные функции \(I_1(x),\ldots,I_\sigma(x)\) при этом, по сказанному выше, становятся целочисленными инвариантами группы, вообще говоря с алгебраически-целыми коэффициентами, поскольку \(\frH\) вместе с \(A\) не обязана содержать все алгебраически сопряженные преобразования.

Следовательно, существует представление
\begin{equation}
        h\cdot f(x)=\omega_1G_1(I)+\cdots+\omega_\rho G_\rho(I),
\end{equation}
где \(G_1,\ldots,G_\rho\) имеют рационально-целые коэффициенты.

Пусть теперь \(w_1,\ldots,w_\rho\) -- неопределенные, и пусть
\[
        L=w_1G_1(I)+\cdots+w_\rho G_\rho(I)
\]
пробегает все линейные формы от \(w\), которым при подстановке \(w_i=\omega_i/h\) соответствуют инварианты \(f(x)\) в представлении (2). Тогда по гильбертовой теореме о целочисленном базисе модуля получаем
\begin{equation}
        L=A_1(I)L_1+\cdots+A_\tau(I)L_\tau,
\end{equation}
где \(A_i\) имеют рационально-целые коэффициенты и где, ввиду линейности всех \(L\) по \(w\), функции \(A_i\) уже не содержат \(w\). Если теперь при подстановке \(\omega_i=w_i/h\) формам \(L_1,\ldots,L_\tau\) соответствуют инварианты \(\mathfrak I_1(x),\ldots,\mathfrak I_\tau(x)\), то (3) при этой подстановке переходит в
\begin{equation}
        f(x)=A_1(I)\mathfrak I_1(x)+\cdots+A_\tau(I)\mathfrak I_\tau(x).
\end{equation}
Это представление (4) показывает, что
\[
        I_1(x),\ldots,I_\sigma(x);\mathfrak I_1(x),\ldots,\mathfrak I_\tau(x)
\]
образуют базис всех целочисленных инвариантов группы \(\frH\), такой что каждый инвариант можно цело и рационально, с рационально-целыми коэффициентами, представить через это конечное число инвариантов.

Если группа \(\frH\), в частности, обладает свойством, что вместе с каждым преобразованием \(A\) она содержит также все \(A',A'',\ldots\) с алгебраически сопряженными числовыми коэффициентами, то целорационально-числовые симметрические функции \(h\) рядов \((x^{(k)})\) становятся функциями от \(x\) с рационально-целыми коэффициентами; это, в частности, верно и для базисных функций симметрических функций. Если поэтому ограничиться инвариантами \(f(x)\) с рационально-целыми коэффициентами, то и базисные функции
\[
        I_1(x),\ldots,I_\sigma(x);\mathfrak I_1(x),\ldots,\mathfrak I_\tau(x)
\]
имеют рационально-целые коэффициенты. Следовательно, при этом специальном предположении, которое вообще для групп \(\frH\) не выполняется, целочисленная конечность имеет место также для подобласти инвариантов с рационально-целыми коэффициентами.
% END UNIT BODY
% END PRODUCER UNIT 077
\endgroup


\clearpage
\begingroup
\setcounter{footnote}{0}
% BEGIN PRODUCER UNIT 078
% Source: C:/Users/Floris/Documents/interlanguage/03_projects/noether/02_slavic_working_corpus/translations/paper16/russian/v001/Noether_Paper16_Russian_v001.tex
% Identity: 24355 B / SHA-256 BD2E4D96BF8EEE18660DEA1B06A65C9D59205743816FB2B467A61BB2A967D75C
\setlength{\parindent}{1.25em}
\setlength{\parskip}{0pt}
\emergencystretch=30em
\allowdisplaybreaks
\raggedbottom
\providecommand{\p}{\partial}
\providecommand{\modu}[1]{\;(\operatorname{mod} #1)}
% BEGIN UNIT BODY
% Unit: Paper 16. Direct Russian translation from the German final-audited source slice.

\section*{16. К разложениям в ряды в теории форм}
\begin{center}
\emph{Mathematische Annalen} 81 (1920), с. 25--30
\end{center}

Под «разложением в ряды в теории форм», как известно, понимают, с одной стороны, обобщение разложения Клебша--Гордана: то есть разложения формы, содержащей два бинарных ряда, по степеням определителя этих рядов, где коэффициенты являются полярами форм только от одного ряда переменных, или, что то же самое, обращаются в нуль при применении \(\Omega\)-процесса. Это обобщается на формы от произвольно многих серий по \(n\) переменных. С другой стороны, как обобщение нормальных форм, возникающих для «комплексных координат» \(t_{ik}\), имеется в виду разложение по нормальным формам для форм, которые одновременно содержат величины различных уровней \(t_i,t_{ik},t_{ikl},\ldots\). Но этот второй вид разложения можно получить из первого посредством инвариантных процессов дифференцирования, как это прежде всего выполнили Мертенс и Дерейтс.\footnote{F. Mertens, Über invariante Gebilde quaternärer Formen (Sitzber. d. Ak. d. Wiss. Wien 98, Abt. IIa (1889)) для кватернарного случая. Общие нормальные формы имеются у J. Deruyts: Sur la réduction des fonctions invariantes dans le système des variables géométriques, Bulletins de l'Ac. R. de Belgique 24 (1892). Не зная этой работы, я, точно следуя Мертенсу, дала общее разложение по нормальным формам в \S~7 работы: Zur Invariantentheorie der Formen von \(n\) Variabeln, J. f. M. 139 (1910).}

Различные разложения первого вида все возникают, как это кратко намечено, например, в \emph{Math. Ann.} 77, цит. место III, как следствия разложения формы от \(n\) рядов по \(n\) переменных по степеням определителя этих рядов, причем коэффициенты являются полярами форм только от \((n-1)\) рядов и сами выводятся из исходной формы образованием поляр и \(\Omega\)-процессом.\footnote{Ср. библиографические указания в цитированном месте. С тех пор появилась также основанная на формальном методе работа: Schouten, Over reeksontwikkelingen van algebraische vormen met verschillende rijen van variabelen van verschillenden graad, Verslag Ak. Amsterdam 27 (1919), с. 1481.} Первоначальное разложение Клебша--Гордана \((n=2)\) шло еще на шаг дальше: специальные возникающие там полярные процессы заданы явно, даже вместе с числовыми коэффициентами.

То, что я теперь хочу добавить к указанной заметке, где разложение в ряды понимается скорее понятийно, как задача о линейных пучках форм, есть явное, с точностью до числовых коэффициентов, представление, получаемое специализацией исследований Э. Фишера\footnote{E. Fischer, Über die Differentiationsprozesse der Algebra, J. f. M. 148 (1917).} об общих модульных системах.

К такому включению приводит истолкование нормальной формы в \(t_{ik}\). Именно, если заменить «комплексные координаты» \(t_{ik}\) «линейными координатами»
\[
        p_{ik}=(x_i y_k-x_k y_i),
\]
то из-за тождественных соотношений между \(p_{ik}\) представление формы \(F(p_{ik})\) уже не является однозначным. Однозначности можно добиться, если требовать, чтобы между коэффициентами \(F\) выполнялись те же линейные соотношения, что и между произведениями степеней \(p_{ik}\), входящими в \(F\).\footnote{Ср. Clebsch, Über eine Fundamentalaufgabe der Invariantentheorie, \S~4, Abh. d. K. G. d. Wissensch. zu Göttingen 17 (1872).} Поэтому для этой нормальной формы \(\Phi\) имеем:
\begin{equation}
\begin{aligned}
        F(p_{ik})&=\Phi(p_{ik}) &&[\text{тождественно по }x,y],\quad\text{или}\notag\\
        F(t_{ik})&\equiv\Phi(t_{ik}) &&\modu{M},
\end{aligned}
\tag{1}
\end{equation}
где \(M\) обозначает модуль всех тождественных соотношений между \(p_{ik}\), то есть состоит из всех форм \(G(t_{ik})\), для которых \(G(p_{ik})\) тождественно по \(x,y\) обращается в нуль. \(\Phi(t_{ik})\) есть «нормальная форма» в комплексных координатах, тогда как базисные функции модуля \(M\) становятся квадратичными формами в \(t_{ik}\). Итерируя этот процесс относительно множителей этих базисных функций, получают полное разложение в ряды. Соответствующие рассуждения применимы и при одновременном появлении \(t_i,t_{ik},t_{ikl},\ldots\).

К аналогичной схеме сводится разложение по полярам. Если, например, в форме \(F(x,y,z)\) заменить три тернарных ряда \(x=(x_1,x_2,x_3),y,z\) такими рядами \(\xi,\eta,\zeta\), которые линейно составлены только из двух из них, например \(\xi=x,\eta=y;\zeta=\lambda_1x+\lambda_2y\), то представление \(F(\xi,\eta,\zeta)\) снова не является однозначным. И здесь можно предписать для коэффициентов те же линейные соотношения, что и для произведений степеней \(\xi,\eta,\zeta\), входящих в \(F\). Так как здесь модуль \(M\) всех тождественных соотношений имеет в качестве базисной функции определитель \(\Delta=(xyz)\) трех рядов (цит. место III, вспомогательная лемма a)), вместо [1] получается:
\begin{equation}
\begin{aligned}
        F(\xi,\eta,\zeta)&=\Phi(\xi,\eta,\zeta)&&[\text{тождественно по }x,y],\quad\text{или}\notag\\
        F(x,y,z)&\equiv\Phi(x,y,z)&&\modu{\Delta}.
\end{aligned}
\tag{2}
\end{equation}
Оказывается, что \(\Phi\) возникает образованием поляр (ср. формулу (2) в цитированном месте); полное разложение в ряды снова получается итерацией.

Итак, образование нормальных форм \(\Phi\) в обоих случаях подчиняется следующей задаче. Если в \(F(z)\) неизвестные \(z_i\) заменить многочленами \(f_i(\xi)\) от параметров \(\xi\), так что представление \(F(f_i(\xi))\) уже не является однозначным, то однозначность должна быть закреплена требованием, чтобы между коэффициентами выполнялись те же линейные соотношения, что и между произведениями степеней \(f_i(\xi)\), входящими в \(F\). Значит,
\begin{equation}
\begin{aligned}
        F(f_i(\xi))&=\Phi(f_i(\xi))&&[\text{тождественно по }\xi],\quad\text{или}\notag\\
        F(z)&\equiv\Phi(z)&&\modu{M},
\end{aligned}
\tag{3}
\end{equation}
где \(M\) обозначает модуль всех тождественных соотношений между \(f_i(\xi)\).

Для этой нормальной формы \(\Phi(z)\) Э. Фишер в \S~11, I своей цитированной работы дал явное, с точностью до числовых коэффициентов, разложение, а именно посредством линейных операций. Одновременно во введении III и в \S~6 II он дал для разности \(F-\Phi\), принадлежащей \(M\), при условии знания базиса модуля, явное выражение в том же смысле; и, соединяя эти разложения в \S~11 II, тем самым заменил формулу [3] явным представлением. При этом все еще открытый вопрос о числовых коэффициентах, согласно \S~12 той же работы, тождественен нахождению собственных значений и собственных функций линейных операций (формула 82).

Теперь в пункте 1 я специализирую формулы Фишера на случай разложения в ряды по полярам и получаю из них путем итерации разложение в ряды. Операция для определения \(\Phi\) при этом переходит в определенные полярные процессы, что существенно точнее приведенной выше, до сих пор известной формулировки. В операции для определения \(F-\Phi\) смешанно выступают умножение на определитель \(\Delta\) и \(\Omega\)-процесс. Чтобы получить обычную, неявную форму, нужно еще заметить, что \(\Omega\Delta\) выражается через полярные процессы.

Здесь следует добавить еще одно исправление. В указанной работе (с. 101, строка 6 сверху) я по ошибке приняла тождественное преобразование
\[
        \Omega(\Delta\psi)=c_1\psi+c_2\Delta\cdot\Omega(\psi),
\]
которое справедливо только в бинарной области, за общезначимое и тем самым сделала переход от фундаментальной формулы (2), дающей выражение для нормальной формы \(\Phi\) в [2], к разложению в ряды (3). Поэтому приведенные там рассуждения дают лишь формулу (2) и, как ее специальный случай, редукционную теорему (1); но, поскольку отсутствует явное выражение для разности \(F-\Phi\), они недостаточны для обоснования разложения в ряды.

В пункте 2 я точно так же, специализацией, получаю нормальную форму \(\Phi(t_{ik})\) из [1]; более общо, нормальную форму \(\Phi(t_i,t_{ik},t_{ikl},\ldots)\), что приводит к известным явным формулам. Но поскольку построение полного разложения в ряды требует знания базиса модуля \(M\), здесь предпочтителен инвариантно-теоретический путь (ср. указанную литературу), который этого знания не требует, а напротив, посредством разложения в ряды сам дает базис модуля.

1. По Фишеру, \S~11 (формула (19) и следующие строки), нормальная форма \(\Phi(z)\) из [3], там обозначенная \(N(\zeta)\), задается так:
\begin{equation}
        \Phi(z)=(a_1S+a_2S^2+\cdots+a_rS^r)F(z),
\tag{4}
\end{equation}
где постоянные \(a_1,\ldots,a_r\) принадлежат области коэффициентов \(f(\xi)\) и зависят только от степени \(F(z)\); оператор \(S\) определен формулами
\begin{equation}
\begin{gathered}
        S=AB;\qquad BF(z)=F(f(\xi));\\
        AF(f(\xi))=\left\{\sum z_i f_i\!\left(\frac{\p}{\p\xi}\right)\right\}^{p}F(f(\xi))\footnotemark.
\end{gathered}
\tag{5}
\end{equation}
\footnotetext{Ср. формулу (75). В случае [3] сумма \(\alpha\) пробегает все произведения степеней \(p\)-го измерения, и поэтому (75) специализируется в [5].}
Если в [2] рассматривать \(F\) как форму \(p\)-го измерения только по \(z\), то \(f_i(\xi)\) специализируется в \(\lambda_1x_i+\lambda_2y_i\); следовательно, получается:
\[
\begin{gathered}
        BF=F(x,y,\lambda_1x+\lambda_2y),\\
        SF=\left\{\sum z_i\left(\frac{\p}{\p\lambda_1}\frac{\p}{\p x_i}+
        \frac{\p}{\p\lambda_2}\frac{\p}{\p y_i}\right)\right\}^{p}
        F(x,y,\lambda_1x+\lambda_2y),
\end{gathered}
\]
или, в записи полярными процессами (ср. цит. место II, 2 и 3; также о сокращенном обозначении),
\begin{equation}
        SF(x,y,z)=\frac1{p!}\left[z\left(\frac{\p}{\p\lambda_1}\frac{\p}{\p x}+
        \frac{\p}{\p\lambda_2}\frac{\p}{\p y}\right)\right]^p
        \left[(\lambda_1x+\lambda_2y)\frac{\p}{\p z}\right]^p F(x,y,z).
\tag{6}
\end{equation}
\emph{Формулами} [4] \emph{и} [6] \emph{дано нужное явное выражение для} \(\Phi(x,y,z)\); вместе с тем \(\Phi(x,y,z)\) подчиняется общей форме, упомянутой в начале (\(\sum PH\) в (2), цит. место): полярам выражений, зависящих только от двух рядов и получаемых из \(F\) полярными процессами. Остающиеся неопределенными коэффициенты \(a_i\) являются рациональными числами, поскольку здесь \(f_i(\xi)\) имеют целые рациональные числовые коэффициенты. Если записать [2] в форме
\begin{equation}
        F=\Phi+\Delta F_1,
\tag{2a}
\end{equation}
то специализацией формулы, данной Фишером во введении III и в \S~6 II (с. 41), получается
\begin{equation}
        F_1=(c_1\Omega+c_2\Omega\Delta\Omega+\cdots+c_i\Omega\Delta\Omega\cdots\Delta\Omega)F\footnotemark,
        \qquad \left[\Omega=\left(\frac{\p}{\p x}\frac{\p}{\p y}\frac{\p}{\p z}\right)\right],
\tag{7}
\end{equation}
\footnotetext{Эту специализированную формулу Керим использует в своей скоро выходящей эрлангенской диссертации в применении к «формам инерции».}
где \(c_i\) снова рациональные числа, зависящие только от степени \(F\).

Итерация приводит к
\[
        F_1=\Phi_1+\Delta F_2;\quad F_2=\Phi_2+\Delta F_3;\quad\ldots;\quad
        F_i=\Phi_i+\Delta F_{i+1};\quad\ldots,
\]
где всякий раз \(\Phi_i\) являются нормальными формами, соответствующими \(F_i\). Так как \(F_i\) имеет по \(z\) степень \((p-i)\), из [4] и [6] получается:
\begin{equation}
\begin{gathered}
        \Phi_i=(a_{i1}S_i+a_{i2}S_i^2+\cdots)F_i,\\
        S_iF_i=\frac1{(p-i)!}
        \left[z\left(\frac{\p}{\p\lambda_1}\frac{\p}{\p x}+
        \frac{\p}{\p\lambda_2}\frac{\p}{\p y}\right)\right]^{p-i}
        \left[(\lambda_1x+\lambda_2y)\frac{\p}{\p z}\right]^{p-i}F_i.
\end{gathered}
\tag{8}
\end{equation}
Формуле [7] соответствует
\[
        F_i=(\alpha_{i1}\Omega+\alpha_{i2}\Omega\Delta\Omega+\cdots)F_{i-1},
        \quad\text{и через итерацию}
\]
\begin{equation}
        F_i=(\beta_1\Omega^i+\beta_2\Omega^i\Delta\Omega+\cdots
        +\beta_\rho\Omega^{k_1}\Delta\Omega^{k_2}\cdots\Delta\Omega^{k_\lambda}+\cdots)F,
        \quad(k_1+k_2+\cdots+k_\lambda-\lambda+1=i).
\tag{9}
\end{equation}
\emph{Посредством} [8] \emph{и} [9] \emph{в разложении}
\[
        F=\Phi+\Delta\Phi_1+\Delta^2\Phi_2+\cdots+\Delta^i\Phi_i+\cdots+\Delta^p\Phi_p
\]
\emph{достигнуто явное представление в указанном смысле.}

Если применить к [7] тождественное преобразование \(\Omega\Delta F=PF\), где \(P\) означает полярные процессы,\footnote{Ср., например, F. Mertens, Über ganze Funktionen von \(m\) Systemen von je \(n\) Unbestimmten, Monatsh. f. Math. u. Phys. 4, формула (5).} то из-за перестановочности этих процессов с \(\Omega\)-процессом получается известная формула
\[
        F_1=P\Omega F;\quad\ldots;\quad F_i=P\Omega^iF;
\]
и отсюда вместе с [8] или с рассуждениями в цитированном месте следует известная форма разложения в ряды, как, например, в (3) в цитированном месте. Если речь идет о \(n\) рядах по \(n\) переменных, \(x^{(1)},x^{(2)},\ldots,x^{(n-1)},z\), то соответственно
\[
        f_i(\xi)=\lambda_1x_i^{(1)}+\lambda_2x_i^{(2)}+\cdots+\lambda_{n-1}x_i^{(n-1)};
\]
и точно так же \(\Delta\) и \(\Omega\) нужно определить для \(n\) рядов.

2. Для нормальной формы \(F(t_i,t_{ik},t_{ikl},\ldots)\) функции \(f_i(\xi),f_{ik}(\xi),\ldots\) заданы одно-, двух-, вплоть до \((n-1)\)-рядными определителями
\[
        \xi_i^{(1)},\quad (\xi^{(1)}\xi^{(2)})_{ik},\quad\ldots,\quad
        (\xi^{(1)}\xi^{(2)}\cdots\xi^{(n-1)})_{i_1i_2\cdots i_{n-1}}.
\]
Тем самым, как специализация [5], получается:
\[
\begin{gathered}
        S=AB;\qquad
        BF=F\bigl(\xi_i^{(1)},(\xi^{(1)}\xi^{(2)})_{ik},\ldots\bigr),\\
        ABF=\left(\sum t_i\frac{\p}{\p\xi_i^{(1)}}\right)^{\lambda_1}
        \left(\sum t_{ik}\left(\frac{\p}{\p\xi^{(1)}}\frac{\p}{\p\xi^{(2)}}\right)_{ik}\right)^{\lambda_2}
        \cdots F\bigl(\xi_i^{(1)},(\xi^{(1)}\xi^{(2)})_{ik},\ldots\bigr),
\end{gathered}
\]
где \(\lambda_1,\lambda_2,\ldots\) обозначают степени по \(t_i,t_{ik},\ldots\).

Из инвариантно-теоретических соображений, а именно из того, что все инварианты \(F\), линейные по коэффициентам, можно линейно составить из \(\Phi\) и из форм из \(M\), следовательно в частности и не принадлежащее \(M\) выражение \(SF\),\footnote{Это следует, например, из \S~6 и \S~7, 2 моей цитированной работы «Zur Invariantentheorie der Formen von \(n\) Variabeln».} здесь следует конгруэнция \(\Phi\equiv a_1SF\modu{M}\); и в силу свойства нормальной формы отсюда вместо [4] получается простое соотношение
\begin{equation}
        \Phi=a_1SF=a_1S\Phi.
\tag{10}
\end{equation}
Второе равенство справедливо потому, что применение \(S\) обращает в нуль всякую форму из \(M\). Формула [10] показывает, что сама нормальная форма \(\Phi\) одновременно является собственной функцией; и действительно, по приведенному выше инвариантному рассуждению нет других собственных функций, которые одновременно были бы инвариантами \(F\) (ср. к [10] Дерейтс, формулы (10) и (10')).

\begin{center}
\textbf{Исправления.}
\end{center}

1. В работе «Die allgemeinsten Bereiche aus ganzen transzendenten Zahlen» (\emph{Math. Ann.} 77) на с. 121, строка 8 сверху, строка 9 сверху и строка 11 снизу, поле \((H,Y)\) нужно заменить на «поле \((H,Y,Y',\ldots,Y^{(\lambda_0-1)})\), где \(Y',\ldots,Y^{(\lambda_0-1)}\) обозначают сопряженные с \(Y\) относительно \(H\)». Соответственно на с. 120, строка 5 сверху, и с. 121, строка 9 сверху, \((H,y)\) нужно заменить на \((H,y,y',\ldots,y^{(\lambda_0-1)})\).

2. В работе «Gleichungen mit vorgeschriebener Gruppe» (\emph{Math. Ann.} 78) в формулировке теоремы Гильберта о неприводимости, с. 222, строка 3 снизу, вместо «числовое поле \(\Omega\)» точнее должно стоять: «конечное алгебраическое числовое поле \(\Omega\)», на что мне указывает господин Островский. Действительно, например, относительно поля всех алгебраических чисел группа всякого уравнения с алгебраическими числовыми коэффициентами является тождественной группой. Соответственно во введении под «произвольной областью рациональности» следует понимать конечное алгебраическое числовое поле, к которому еще можно присоединить конечное число параметров, или одну из возникающих таким образом промежуточных областей.

\begin{center}
(Принято в сентябре 1919 г.)
\end{center}
% END UNIT BODY
% END PRODUCER UNIT 078
\endgroup


\clearpage
\begingroup
\setcounter{footnote}{0}
% BEGIN PRODUCER UNIT 079
% Source: C:/Users/Floris/Documents/interlanguage/03_projects/noether/02_slavic_working_corpus/translations/paper17/russian/v001/Noether_Paper17_Through_Section02_Russian_v001.tex
% Identity: 23647 B / SHA-256 26C4977E794C632B803BDE5CC61448487FA79CC1F2041C531EE6DAC1245E9E43
\setlength{\parindent}{1.25em}
\setlength{\parskip}{0pt}
\emergencystretch=30em
\allowdisplaybreaks
\raggedbottom
\providecommand{\dd}{\mathrm d}
\providecommand{\p}{\partial}
% BEGIN UNIT BODY
% Unit: Paper 17 through Section 02. Direct Russian translation from the German final-audited source slice.

\setcounter{footnote}{0}
\section*{17. Совместно с В. Шмайдлером: модули в некоммутативных областях, в особенности из разностных выражений}
\begin{center}
\emph{Mathematische Zeitschrift} 8 (1920), с. 1--35
\end{center}

\begin{center}
\textbf{Содержание.}
\end{center}
\begin{center}
\begin{minipage}{0.88\linewidth}
Введение.\\
\S~1. Формальные сведения о дифференциальных и разностных выражениях.\\
\S~2. Общая полиномиальная область.\\
\S~3. Односторонние модули и их группы остатков.\\
\S~4. Приводимость групп остатков к двум подгруппам и их связи с модулем.\\
\S~5. Вычисления с единицами и приводимость к \(k\) подгруппам.\\
\S~6. Теоремы конечности.\\
\S~7. Один случай однозначного разложения.\\
\S~8. Аддитивное разложение вполне приводимых групп на простые группы. Теорема об изоморфии.\\
\S~9. Связь между понятиями «изоморфный» и «того же вида».\\
\S~10. Существование бесконечно многих разложений группы.\\
\S~11. Связи с системами дифференциальных уравнений и их интегралами.\\
\S~12. Примеры.
\end{minipage}
\end{center}

\begin{center}\textbf{Введение.}\end{center}

Формальная теория дифференциальных выражений от одной переменной приводит к разложению на неприводимые множители, которое в некотором смысле однозначно.\footnote{Ср. E. Landau, \emph{Ein Satz über die Zerlegung homogener linearer Differentialausdrücke in irreduzible Faktoren}, \emph{Journal für die reine und angewandte Mathematik} 124 (1902), с. 115--120, и вслед за этим A. Loewy, \emph{Über reduzible lineare homogene Differentialgleichungen}, \emph{Mathematische Annalen} 56 (1903), с. 549--584.} Но тогда как соответствующую теорему в алгебре можно перенести на многочлены от нескольких переменных, для частных дифференциальных выражений это уже не удается.\footnote{Ср. контрпример Ландау, опубликованный в H. Blumberg, \emph{Über algebraische Eigenschaften von linearen homogenen Differentialausdrücken}, диссертация, Göttingen 1912, 52 с., с. 51--52.} При дифференциальных выражениях от одной переменной, однако, наряду с представлением в виде произведения имеется еще представление как наименьшее общее кратное, причем действуют сходные законы; именно, при двух различных разложениях отдельные составные части можно попарно сопоставить друг другу: они «того же вида».\footnote{Ср., кроме названных работ, A. Loewy, \emph{Über vollständig reduzible lineare homogene Differentialgleichungen}, \emph{Mathematische Annalen} 62 (1906), с. 89--117; далее цитируется как Loewy.}

Это понятие наименьшего общего кратного теперь можно истолковать как модульное понятие и тем самым перенести на \(n\) переменных; в настоящей работе это дает естественное обобщение названных теорем. Кроме того, мы еще замечаем, что от дифференциальных выражений используется только то, что их можно понимать как многочлены с некоммутативным умножением, для которых степень произведения равна сумме степеней множителей. Эти условия выполняются, например, и для разностных выражений, так что теоремы справедливы и для них, где до сих пор они также были известны только для одной переменной.\footnote{Ср. G. Wallenberg--A. Guldberg, \emph{Theorie der linearen Differenzengleichungen}, Leipzig und Berlin 1911.}

Поэтому мы развиваем в общем виде теорию модулей, элементы которых являются многочленами с некоммутативным умножением, и при этом сводим рассмотрение модулей к рассмотрению классов остатков.\footnote{Для случая коммутативных областей ср. W. Schmeidler, \emph{Über Moduln und Gruppen hyperkomplexer Größen}, \emph{Mathematische Zeitschrift} 3 (1919), с. 29--42.} В отличие от коммутативного специального случая здесь можно говорить о сумме классов остатков, но не об их произведении; можно говорить только о произведении класса остатков на многочлен. Поэтому понятие «группы остатков» должно быть модифицировано по сравнению с коммутативным случаем. В специальном случае многочленов от одной переменной группа остатков, кроме того, «конечна», то есть имеет лишь конечное число линейно независимых классов остатков. В общем случае такого конечного «порядка» не существует; группы «бесконечны». Для наших методов, однако, это различие не существенно.

Представлению модуля как наименьшего общего кратного взаимно простых модулей, как и в коммутативном случае, соответствует более простой процесс аддитивного разложения группы остатков. При этом существенной оказывается изоморфия групп остатков, которая при специализации к дифференциальным выражениям от одной переменной переходит в понятие того же вида для соответствующих модулей, а при специализации к коммутативным модулям переходит в тождественность. Впрочем, понятие того же вида также можно перенести на \(n\) переменных и доказать его совпадение с понятием изоморфии. То, что в аддитивном разложении группы остатков всегда возникает лишь конечное число слагаемых, следует путем переноса гильбертовой теоремы о базисе модуля, основанной на методе доказательства Гордана; первоначальный метод доказательства Гильберта потребовал бы ограничений на законы умножения (в формуле (6) справа \(a\) вместо \(a_i\), что выполняется для дифференциальных, но не для разностных выражений).

Чтобы теперь прийти к теореме единственности, мы должны специализировать модули, как это делает и Лёви, до «вполне приводимых», то есть до таких, которые представимы как наименьшее общее кратное простых модулей. При двух различных разложениях составные части тогда либо попарно тождественны, либо по крайней мере каждой составной части одного разложения соответствует изоморфная составная часть другого разложения, и наоборот; вместе с тем каждое разложение тогда содержит по меньшей мере две изоморфные составные части. Далее, из появления двух изоморфных составных частей в одном и том же разложении следует существование бесконечно многих разложений, и это справедливо даже без ограничения на простые модули, что до сих пор не было известно даже в специальном случае дифференциальных выражений от одной переменной.

Наконец, для модулей из дифференциальных выражений мы переходим к связи с интегралами и показываем, что при конечных группах остатков порядок равен наибольшему числу линейно независимых интегралов; следовательно, произвольные функции возникают в общем интеграле тогда и только тогда, когда группа остатков не конечна. Понятие того же вида для двух модулей, кроме того, означает, как и в случае одной переменной, «бирациональное» родство интегралов. Завершают работу несколько примеров.

Первый толчок к настоящей работе был дан несколько лет назад вопросом Э. Ландау к Э. Нётер об обобщении теоремы о произведении. Тогда Э. Нётер лишь заметила, что речь идет о вопросе теории модулей в некоммутативной полиномиальной области, но подступиться к нему еще не удавалось. Известные тогда ласкеровы теоремы о модулях нельзя было непосредственно перенести на эту область, потому что метод доказательства Ласкера основан на теореме, что многочлен однозначно разлагается на неприводимые множители.\footnote{Недавно Э. Нётер заметила, что и ласкеровы теоремы можно обосновать без использования этой теоремы, методами, родственными применяемым здесь.} Только методы Шмайдлера дали возможность переноса, который мы затем осуществили совместно. В частности, следует еще упомянуть, что обобщение группы остатков, перенос и применение гильбертовой теоремы конечности и существование бесконечно многих разложений принадлежат Э. Нётер, а обобщение лёвиевского понятия полной приводимости, понятие изоморфии и его связь с понятием того же вида принадлежат В. Шмайдлеру.

\subsection*{\S~1. Формальные сведения о дифференциальных и разностных выражениях.}

1. Для дифференциальных выражений от одной переменной было введено символическое образование произведения. Пусть
\[
 a_0(x)\frac{\dd^n y}{\dd x^n}+a_1(x)\frac{\dd^{n-1}y}{\dd x^{n-1}}+\cdots+a_n(x)y=A(y),
\]
\[
 b_0(x)\frac{\dd^m y}{\dd x^m}+b_1(x)\frac{\dd^{m-1}y}{\dd x^{m-1}}+\cdots+b_m(x)y=B(y).
\]
Тогда символическое произведение \(AB(y)\) определяется как дифференциальное выражение
\[
 a_0(x)\frac{\dd^n B(y)}{\dd x^n}+a_1(x)\frac{\dd^{n-1}B(y)}{\dd x^{n-1}}+\cdots+a_n(x)B(y).
\]
Это образование произведения в действительности есть умножение двух операторов; проще всего его записывать, заменяя, как обычно, \(\frac{\dd^n}{\dd x^n}\) на \(\left(\frac{\dd}{\dd x}\right)^n\). Тогда \(AB\) переходит в
\[
 \left(a_0\left(\frac{\dd}{\dd x}\right)^n+\cdots+a_n\right)
 \left(b_0\left(\frac{\dd}{\dd x}\right)^m+\cdots+b_m\right)y,
\]
причем при выполнении произведения действуют обычные правила дифференцирования сумм и произведений. Каждый такой оператор теперь можно рассматривать как многочлен от одной переменной \(\frac{\dd}{\dd x}=\xi\): функцию, к которой применяется оператор, явно не записывают, а для переменной \(\xi\) фиксируют те правила умножения, которые соответствуют дифференцированию произведения:
\[
 \frac{\dd}{\dd x}\bigl(a(x)y\bigr)=a(x)\frac{\dd}{\dd x}(y)+a'(x)y,
\]
или, в нашей записи \((n=1, m=0)\),
\begin{equation}
 \xi\cdot a=a\cdot\xi+a'.
\tag{1}
\end{equation}

Совершенно так же поступают с линейными частными дифференциальными выражениями
\[
 \left(\sum a_{\alpha_1\ldots\alpha_n}(x_1,\ldots,x_n)
 \frac{\p^{\alpha_1+\cdots+\alpha_n}}
 {\p x_1^{\alpha_1}\cdots\p x_n^{\alpha_n}}\right)y
\]
от функции \(y\). Для \(n\) операторов \(\frac{\p}{\p x_1},\ldots,\frac{\p}{\p x_n}\) вводим переменные \(\xi_1,\ldots,\xi_n\) и тем самым получаем многочлен
\[
 \sum a_{\alpha_1\ldots\alpha_n}(x_1,\ldots,x_n)
 \xi_1^{\alpha_1}\cdots\xi_n^{\alpha_n}.
\]
Произведение двух многочленов тогда вычисляется по законам
\begin{equation}
 \xi_i a=a\xi_i+a_i \qquad (i=1,\ldots,n),
\tag{2}
\end{equation}
где \(a_i\) означает частную производную \(a\) по \(x_i\); произведения \(\xi\) между собой коммутативны. Кроме коммутативного закона умножения, для так определенных многочленов действуют все законы сложения и умножения, как в обычной алгебре. В частности, для произведения двух таких многочленов общая степень и степень по каждой отдельной переменной равны сумме соответствующих степеней двух множителей.

2. Для разностных выражений (ср. Wallenberg, цит. место) вида
\[
 a_x^{(0)}y_{x+n}+a_x^{(1)}y_{x+n-1}+\cdots+a_x^{(n)}y_x=A_x,
\]
\[
 b_x^{(0)}y_{x+m}+b_x^{(1)}y_{x+m-1}+\cdots+b_x^{(m)}y_x=B_x
\]
имеется соответствующее определение произведения
\[
 (AB)_x=a_x^{(0)}B_{x+n}+a_x^{(1)}B_{x+n-1}+\cdots+a_x^{(n)}B_x,
\]
которое можно выполнить с помощью правила умножения
\[
 (ay)_{x+1}=a_{x+1}y_{x+1}.
\]
Если для перехода от \(x\) к \(x+1\) ввести оператор \(\xi\), то получаем \(\xi\{ay\}=a^*\xi\{y\}\), где \(a^*=a_{x+1}\). Но и здесь можно снова опустить функцию \(y\), к которой применяется оператор, и получить закон
\begin{equation}
 \xi a=a^*\xi,
\tag{3}
\end{equation}
где наряду с \(a\) также \(a^*\) не обращается тождественно в нуль. Совершенно так же для \(n\) переменных имеем
\begin{equation}
 \xi_i a=a_i\xi_i\qquad (i=1,\ldots,n),
\tag{4}
\end{equation}
где \(\xi_i\) соответствует переходу от \(x_i\) к \(x_i+1\), а \(a_i\) обозначает функцию, полученную из \(a\) заменой \(x_i\) на \(x_i+1\), следовательно не обращающуюся тождественно в нуль, если это верно для \(a\). При этих соглашениях, за исключением коммутативного закона, также действуют все правила сложения и умножения, включая правило о степени произведения.

Оба случая в следующем параграфе будут подчинены вполне общо определенной полиномиальной области с произвольными коэффициентами.

\subsection*{\S~2. Общая полиномиальная область.}

Мы исходим из произвольно абстрактно определенного поля \(P\) (элементы которого обозначаем малыми латинскими буквами), для которого действуют все правила алгебры, в частности также коммутативный закон умножения с однозначным обращением. К этому полю присоединяем \(n\) неизвестных \(\xi_1,\ldots,\xi_n\), то есть рассматриваем область целостности \(J\), состоящую из произведений \(a\xi_1^{\nu_1}\xi_2^{\nu_2}\cdots\xi_n^{\nu_n}\) и их конечных сумм, причем должны действовать все законы сложения и умножения, за исключением коммутативного закона умножения. В частности, конечные суммы вида
\[
 F(\xi_1,\ldots,\xi_n)=\sum a_{\nu_1\ldots\nu_n}\xi_1^{\nu_1}\cdots\xi_n^{\nu_n}
\]
называем «многочленами» (которые всюду будем обозначать большими латинскими буквами). Теперь вместо коммутативного закона мы подчиняем величины из \(J\) таким законам произведения, чтобы \emph{произведение двух многочленов снова было многочленом}; тем самым область всех многочленов из \(J\) уже исчерпывает область \(J\). Для этого, очевидно, необходимо и достаточно, чтобы для величин \(\xi_1,\ldots,\xi_n\) и произвольной величины \(a\) из \(P\) действовали правила произведения вида
\begin{equation}
 \xi_i a=F_i(\xi_1,\ldots,\xi_n)\qquad (i=1,\ldots,n),
\tag{5}
\end{equation}
где \(F_i(\xi_1,\ldots,\xi_n)\) означает некоторый, естественно зависящий от \(\xi_i\) и \(a\), многочлен. (В частности, правила (2) и (4) для дифференциальных и соответственно разностных выражений имеют именно этот вид.) Действительно, поскольку \(\xi_i a\) само по себе не является многочленом, но по предположению принадлежит области, это равенство необходимо. С другой стороны, конечным повторным применением формулы (5) всякую величину области \(J\) можно превратить в многочлен.

Для единицы поля, которую мы обозначаем через 1, далее должно выполняться \(\xi_i\cdot1=1\cdot\xi_i=\xi_i\). Кроме того, установим, что умножение \(\xi_1,\ldots,\xi_n\) между собой коммутативно.\footnote{Это понадобится только начиная с \S~6 для теоремы конечности и ее следствий.} Тогда каждый многочлен в \(J\) записывается как обычный многочлен от \(\xi_1,\ldots,\xi_n\) с коэффициентами из \(P\), и два многочлена совпадают только тогда, когда совпадают все их коэффициенты.

Более специально, далее мы будем требовать,\footnote{Это также понадобится только начиная с \S~6.} чтобы вместо (5) действовало правило
\begin{equation}
 \xi_i a=a_i\xi_i+b_i\qquad (a_i\ne0)\quad(i=1,\ldots,n).
\tag{6}
\end{equation}
Сюда, в частности, также входят законы (2) и (4). Тогда степень произведения равна сумме степеней множителей, причем это верно как для общей степени, так и для степени по каждой отдельной переменной; следовательно, произведение двух многочленов, не обращающихся тождественно в нуль, также отлично от 0.
% END UNIT BODY
% END PRODUCER UNIT 079
\endgroup


\clearpage
\begingroup
\setcounter{footnote}{0}
% BEGIN PRODUCER UNIT 080
% Source: C:/Users/Floris/Documents/interlanguage/03_projects/noether/02_slavic_working_corpus/translations/paper17/russian/v001/Noether_Paper17_Section03_Russian_v001.tex
% Identity: 6785 B / SHA-256 0CD63E43009EDFD31F615E09C8D2267EA488DFC8FAEE2C09FEF9F3957A1A4845
\setlength{\parindent}{1.25em}
\setlength{\parskip}{0pt}
\emergencystretch=30em
\allowdisplaybreaks
\raggedbottom
\providecommand{\mM}{\mathfrak M}
\providecommand{\mx}{\mathfrak x}
\providecommand{\my}{\mathfrak y}
\providecommand{\mz}{\mathfrak z}
\providecommand{\me}{\mathfrak e}
% BEGIN UNIT BODY
% Unit: Paper 17 Section 03. Direct Russian translation from the German final-audited source slice.

\setcounter{footnote}{8}
\subsection*{\S~3. Односторонние модули и их группы остатков.}

Для нашей области \(J\), ввиду некоммутативности умножения, можно различать следующие виды модулей\footnote{Модули мы обозначаем большими готическими буквами.}:\footnote{Односторонние идеалы уже рассматривал А. Гурвиц; ср. \emph{Vorlesungen über die Zahlentheorie der Quaternionen} (Berlin, Springer 1919), лекция 6. Однако там по существу речь идет о главных идеалах; то же относится к работам Дю Паскье, цитированным там в примечании 1.}

1. \textit{Правосторонний модуль}, который вместе с \(M\) содержит также \(FM\) для произвольного многочлена \(F\), а вместе с \(M_1\) и \(M_2\) также \(M_1+M_2\).

2. \textit{Левосторонний модуль}, который вместе с \(M\) содержит также \(MF\), а вместе с \(M_1\) и \(M_2\) также \(M_1+M_2\).

Эти два вида мы называем общим именем \textit{односторонних модулей}; модули, которые являются одновременно правосторонними и левосторонними, будем называть \textit{двусторонними}. В дальнейшем мы ограничиваемся правосторонними модулями и замечаем, что теория левосторонних строится совершенно аналогично.

Назовем два многочлена \(F\) и \(G\) конгруэнтными по модулю \(\mM\), в обозначениях \(F\equiv G(\mM)\), если их разность делится на \(\mM\), то есть принадлежит \(\mM\). Совокупность всех многочленов, конгруэнтных данному многочлену, образует \textit{класс остатков} по модулю \(\mM\). (Классы остатков будем обозначать малыми готическими буквами.) Так как из
\[
 F_1\equiv F_2(\mM)\qquad\text{и}\qquad G_1\equiv G_2(\mM)
\]
следует конгруэнция \(F_1+G_1\equiv F_2+G_2(\mM)\), сумма двух классов остатков \(\mx\) и \(\my\) определена и снова является классом остатков \(\mx+\my\). Напротив, из-за односторонности модуля нельзя говорить о произведении двух классов остатков; действительно, из
\[
 F_1=F_2+M_1,\qquad G_1=G_2+M_2
\]
хотя и следует
\[
 F_1G_1=F_2G_2+F_2M_2+M_1G_2+M_1M_2,
\]
но поскольку \(M_1G_2\) не обязан принадлежать \(\mM\), когда \(M_1\) принадлежит \(\mM\), вообще говоря, \(F_1G_1\) и \(F_2G_2\) не обязаны быть конгруэнтными. Зато при \(M_1=0\), то есть \(F_1=F_2=F\), видно, что из \(G_1\equiv G_2\) следует и \(FG_1\equiv FG_2\); значит, произвольный класс остатков \(\my\) можно умножать спереди на многочлен \(F\), получая тем самым вполне определенный класс остатков \(F\my=\mz\).

Для вычислений с классами остатков действуют коммутативный и ассоциативный законы, а также закон однозначной обратимости сложения; кроме того, для умножения классов остатков на многочлены и для сложения действует дистрибутивный закон \(F(\my_1+\my_2)=F\my_1+F\my_2\), что непосредственно следует из вычислений с многочленами по модулю \(\mM\).

Класс остатков, которому принадлежит единица, будем называть единичным классом, или единицей, и обозначать \(\me\). Тогда \(F\me=\mx\) есть класс остатков, которому принадлежит многочлен \(F\). В отличие от общих классов остатков, здесь также произведение \(\mx\me\) вполне определено и равно \(\mx\); ибо из
\[
 F_1=F_2+M_1,\qquad G_1=1+M_2
\]
вследствие \(M_1\cdot1=M_1\) следует также
\[
 F_1G_1=F_2+M.
\]
Совокупность классов остатков по модулю \(\mM\) мы называем \textit{группой остатков модуля} \(\mM\). Следовательно, группа остатков имеет следующие свойства: 1. Вместе с каждым классом остатков \(\mx\) и для произвольного многочлена \(F\) она содержит также класс остатков \(F\mx\), а вместе с классами остатков \(\mx_1\) и \(\mx_2\) также класс остатков \(\mx_1+\mx_2\).\footnote{Это свойство в более общем смысле можно было бы назвать «модульным свойством» группы остатков.} 2. Она обладает единицей \(\me\), так что \(F\me\) для каждого многочлена \(F\) представляет класс остатков \(\mx\) многочлена \(F\). \textit{Итак, когда} \(F\) \textit{пробегает все многочлены,} \(F\me\) \textit{пробегает всю группу остатков.}
% END UNIT BODY
% END PRODUCER UNIT 080
\endgroup


\clearpage
\begingroup
\setcounter{footnote}{0}
% BEGIN PRODUCER UNIT 081
% Source: C:/Users/Floris/Documents/interlanguage/03_projects/noether/02_slavic_working_corpus/translations/paper17/russian/v001/Noether_Paper17_Section04_Russian_v001.tex
% Identity: 11743 B / SHA-256 2EE3179811A444DCA7C871965E4E9EA2276AA264A76E359F77EF07BA87C79717
\setlength{\parindent}{1.25em}
\setlength{\parskip}{0pt}
\emergencystretch=30em
\allowdisplaybreaks
\raggedbottom
\providecommand{\mM}{\mathfrak M}
\providecommand{\mN}{\mathfrak N}
\providecommand{\mL}{\mathfrak L}
\providecommand{\mG}{\mathfrak G}
\providecommand{\mA}{\mathfrak A}
\providecommand{\mB}{\mathfrak B}
\providecommand{\mbarA}{\bar{\mathfrak A}}
\providecommand{\mbarB}{\bar{\mathfrak B}}
\providecommand{\ma}{\mathfrak a}
\providecommand{\mb}{\mathfrak b}
\providecommand{\me}{\mathfrak e}
\providecommand{\bara}{\bar{\mathfrak a}}
% BEGIN UNIT BODY
% Unit: Paper 17 Section 04. Direct Russian translation from the German final-audited source slice.

\setcounter{footnote}{10}
\subsection*{\S~4. Приводимость группы остатков к двум подгруппам и ее связи с модулем.}

Для односторонних модулей понятия \emph{делимости}, \emph{наибольшего общего делителя} и \emph{наименьшего общего кратного} сохраняются так же, как в двустороннем случае; это показывают следующие определения:

Модуль \(\mM\) называется делимым на другой модуль \(\mN\), если из \(M\equiv0(\mM)\) следует также \(M\equiv0(\mN)\).

Для двух модулей \(\mM\) и \(\mN\) \emph{наибольший общий делитель} \((\mM,\mN)\) определяется как совокупность всех многочленов, представимых в виде \(M+N\), где \(M\equiv0(\mM)\), \(N\equiv0(\mN)\); это снова модуль. Соответствующее определение действует и для нескольких модулей.

Совокупность всех многочленов, делимых и на \(\mM\), и на \(\mN\), образует модуль \([\mM,\mN]\), \emph{наименьшее общее кратное} модулей \(\mM\) и \(\mN\). Это определение также соответственно применяется к нескольким модулям, даже к множеству модулей произвольной мощности.

Два модуля называются взаимно простыми, если их наибольший общий делитель есть единичный модуль, то есть состоит из совокупности всех многочленов.

Теперь исходим из модуля \(\mM\), который представим как наименьшее общее кратное двух взаимно простых модулей \(\mN\) и \(\mL\), причем оба предполагаются собственными делителями:
\begin{equation}
 \mM=[\mN,\mL].
\tag{7}
\end{equation}
Исследуем группу остатков \(\mG\) модуля \(\mM\). По предположению
\begin{equation}
 1\equiv N+L(\mM)
\tag{8}
\end{equation}
для двух многочленов \(N\equiv0(\mN)\) и \(L\equiv0(\mL)\). Отсюда покажем
\begin{equation}
 \mL=(\mM,L),\qquad \mN=(\mM,N),
\tag{9}
\end{equation}
где, например, под \((\mM,L)\) следует понимать модуль, состоящий из всех многочленов вида \(M+FL\). В самом деле, во-первых, \((\mM,L)\equiv0(\mL)\). С другой стороны, из (8) для произвольного многочлена \(F\) следует
\begin{equation}
 F\equiv FN(\mL).
\tag{10}
\end{equation}
Если теперь \(F\equiv0(\mL)\), то \(FN\equiv0(\mL)\), значит также \(\equiv0(\mM)\), и, следовательно, \(F\equiv FL(\mM)\). Соответственно доказывается \(\mN=(\mM,N)\). Отсюда, кроме того, следует: поскольку \(\mN\) и \(\mL\) являются собственными делителями \(\mM\), ни \(L\equiv0(\mM)\), ни \(N\equiv0(\mM)\) не может иметь места.

Пусть теперь \(\ma\) есть класс остатков многочлена \(N\) по модулю \(\mM\), а \(\mb\) -- класс остатков многочлена \(L\) по модулю \(\mM\). Тогда по (8)
\[
 \me=\ma+\mb\qquad(\ma\ne0,\ \mb\ne0).
\]
Так как каждый класс остатков имеет вид \(F\me\), то по дистрибутивному закону для каждого класса остатков получается
\begin{equation}
 F\me=F\ma+F\mb\qquad(\ma\ne0,\ \mb\ne0).
\tag{11}
\end{equation}
Однако это представление произвольного класса остатков как суммы двух классов остатков вида \(G\ma\) и \(H\mb\) является \emph{единственным}. Действительно, из соотношения \(0=G\ma+H\mb\), если \(K\) есть многочлен из класса остатков \(G\ma=-H\mb\), очевидно имеем \(K\equiv GN(\mM)\), \(K\equiv-HL(\mM)\), следовательно \(K\equiv0(\mN)\), \(K\equiv0(\mL)\) и потому \(K\equiv0(\mM)\), то есть \(G\ma=0\), \(H\mb=0\).

Пусть, обратно, соотношение (11) выполнено для каждого многочлена \(F\), причем \emph{единственно} в указанном выше смысле. Пусть \(N\) будет многочленом из \(\ma\), а \(L\) -- многочленом из \(\mb\). Образуем два модуля (9) и утверждаем (7) и (8). Последнее следует из (11) при \(F=1\). Далее, по определению \(\mN\) и \(\mL\), \(\mM\) является общим кратным этих двух модулей, притом наименьшим. Действительно, из \(F\equiv0(\mL)\) и \(F\equiv0(\mN)\) следует \(F\equiv HL\equiv GN(\mM)\), значит \(G\ma=H\mb\), \(G\ma-H\mb=0\), а потому, в силу единственности представления, \(G\ma=0\), \(H\mb=0\), \(F\equiv0(\mM)\). Тем самым показано, что \emph{представление модуля} \(\mM\) \emph{как наименьшего общего кратного двух взаимно простых модулей равносильно единственному аддитивному разложению группы остатков} \(\mG\) \emph{модуля} \(\mM\), \emph{как оно задается формулой} (11).

Теперь рассмотрим систему \(\mA\) классов остатков вида \(F\ma\), которую хотим связать с группой остатков \(\mbarA\) модуля \(\mL\). По (11) каждый многочлен \(F\) из \(\mL\) принадлежит классу остатков \(F\ma\), если \(\bara\) есть класс остатков \(N\) по модулю \(\mL\), то есть единичный класс \(\mbarA\). Если теперь сопоставить друг другу классы остатков \(F\ma\) в \(\mA\) и \(F\bara\) в \(\mbarA\), то каждому классу остатков из \(\mA\) сопоставляется класс остатков из \(\mbarA\), причем сумме двух классов остатков из \(\mA\) соответствует сумма соответствующих классов остатков из \(\mbarA\), а произведению класса остатков из \(\mA\) на многочлен соответствует произведение соответствующего класса остатков из \(\mbarA\) на тот же многочлен. Это сопоставление, кроме того, взаимно однозначно: из \(F\ma=0\) по (11) следует \(F\equiv0(\mL)\), значит \(F\bara=0\); обратно, \(F\bara=0\), то есть \(F\equiv0(\mL)\), по (10) означает также \(FN\equiv0(\mL)\), значит \(FN\equiv0(\mM)\), то есть \(F\ma=0\).

Тем самым мы приходим к фундаментальному понятию, которое определяем следующим образом:

\emph{Взаимно однозначное сопоставление между двумя системами \(\mA\) и \(\mbarA\) классов остатков называется изоморфным, если сумме двух классов остатков из \(\mA\) сопоставлена сумма соответствующих классов остатков из \(\mbarA\), а произведению произвольного класса остатков из \(\mA\) на многочлен -- произведение соответствующего класса остатков из \(\mbarA\) на тот же многочлен.}

Подсистема \(\mA\) классов остатков в \(\mG\), изоморфная группе остатков некоторого модуля, называется \emph{подгруппой} \(\mG\). Как показывает приведенное выше доказательство, здесь сопоставление специально устроено так, что каждый класс остатков из \(\mA\) возникает из соответствующего класса остатков из \(\mbarA\) присоединением некоторых многочленов, а именно всех многочленов, принадлежащих \(\mL\), но не принадлежащих \(\mM\).

Соответственно можно показать, что совокупность всех классов остатков вида \(F\mb\) образует подгруппу \(\mB\) группы \(\mG\), изоморфную группе остатков \(\mbarB\) модуля \(\mN\).

Группа остатков \(\mG\), классы остатков которой допускают единственное аддитивное разложение (11), будет называться \emph{приводимой} к двум подгруппам \(\mA\) и \(\mB\), в обозначениях \(\mG=\mA+\mB\); группа, не являющаяся приводимой, называется \emph{неприводимой}. Соответствующие модули мы также иногда называем приводимыми или неприводимыми. Итак, имеет место следующая теорема.

\textbf{Теорема 1.}\footnote{Ср. Schmeidler, там же, с. 39.} \emph{Модуль \(\mM\) тогда и только тогда представим как наименьшее общее кратное двух взаимно простых собственных делителей \(\mL\) и \(\mN\), когда его группа остатков \(\mG\) приводима к двум подгруппам \(\mA\) и \(\mB\). При этом \(\mA\) и \(\mB\) соответственно изоморфны группам остатков модулей \(\mL\) и \(\mN\).}
% END UNIT BODY
% END PRODUCER UNIT 081
\endgroup


\clearpage
\begingroup
\setcounter{footnote}{0}
% BEGIN PRODUCER UNIT 082
% Source: C:/Users/Floris/Documents/interlanguage/03_projects/noether/02_slavic_working_corpus/translations/paper17/russian/v001/Noether_Paper17_Section05_Russian_v001.tex
% Identity: 12167 B / SHA-256 8AAAA3C0D3955C0090F75A3443694B2AD016C6690B27882FD1DEB41DB24A4B32
\setlength{\parindent}{1.25em}
\setlength{\parskip}{0pt}
\emergencystretch=30em
\allowdisplaybreaks
\raggedbottom
\providecommand{\mM}{\mathfrak M}
\providecommand{\mN}{\mathfrak N}
\providecommand{\mG}{\mathfrak G}
\providecommand{\mU}{\mathfrak U}
\providecommand{\ma}{\mathfrak a}
\providecommand{\mb}{\mathfrak b}
\providecommand{\me}{\mathfrak e}
\providecommand{\mx}{\mathfrak x}
% BEGIN UNIT BODY
% Unit: Paper 17 Section 05. Direct Russian translation from the German final-audited source slice.

\setcounter{footnote}{11}
\subsection*{\S~5. Вычисления с единицами и приводимость на \(k\) подгрупп.}

Классы остатков \(\ma\) и \(\mb\) группы \(\mG\) имеют с единицей \(\me\) общее свойство: произведение произвольного класса остатков \(\mx\) с \(\ma\) и с \(\mb\) определено; поэтому \(\ma\) и \(\mb\) мы также будем называть единицами группы \(\mG\). Действительно, из (11) при \(F=M\) следует
\[
 0=M\me=M\ma+M\mb,
\]
и потому, в силу единственности, \(M\ma=0\), \(M\mb=0\), а значит
\[
 (F+M)\ma=F\ma=\mx\ma,
\]
где \(\mx\) означает класс остатков полинома \(F\). Поэтому вместо отношения (11) можно записать равносильное отношение между классами:
\[
 \mx\me=\mx\ma+\mx\mb.
\]
В частности, подставляя \(\mx=\ma\), получаем
\[
 \ma=\ma^2+\ma\mb,
\]
а отсюда в силу единственности \(\ma^2=\ma\), \(\ma\mb=0\). Точно так же получается \(\mb^2=\mb\), \(\mb\ma=0\).

Соответственно представлению группы остатков как суммы двух составляющих можно теперь определить и такое представление как сумму \(k\) составляющих. Пусть задано единственное представление
\[
 F\me=F\ma_1+F\ma_2+\cdots+F\ma_k\qquad(\ma_1\ne0,\ldots,\ma_k\ne0),
\]
причем из \(0=G_1\ma_1+\cdots+G_k\ma_k\) также следует \(G_i\ma_i=0\). Тогда, в частности при \(F=M\), получается отношение \(M\ma_i=0\), так что представление можно записать и в форме\footnote{Для аддитивного разложения групп остатков выполняются коммутативный и ассоциативный законы; однако, в отличие от коммутативного случая, из \(\mU_1+\mU_2=\mU_1+\mU_3\) не следует \(\mU_2=\mU_3\); ср. пример 1 в \S~12.}
\begin{equation}
 \mx\me=\mx\ma_1+\cdots+\mx\ma_k\qquad(\ma_1\ne0,\ldots,\ma_k\ne0).
\tag{12}
\end{equation}
Это представление можно также понимать как возникшее из последовательного разложения на две составляющие. Если его можно продолжать неограниченно, причем все \(\ma_i\ne0\), то мы говорим о сумме бесконечно многих составляющих.

Подставляя в (12) \(\mx=\ma_i\), получаем
\begin{equation}
 \ma_i^2=\ma_i,
 \qquad \ma_i\ma_j=0\quad(i\ne j),
 \qquad i,j=1,\ldots,k,
\tag{13}
\end{equation}
где \(\ma_i\ne0\). Пусть теперь, наоборот, в \(\mG\) имеется система классов остатков \(\ma_1,\ldots,\ma_k\), для которых выполнены условия ортогональности (13) и для которых
\begin{equation}
 \me=\ma_1+\cdots+\ma_k
\tag{14}
\end{equation}
имеет место. Мы покажем, что отсюда получается представление группы как суммы \(k\) составляющих. Прежде всего из предположения для каждого полинома \(A_i\) из класса \(\ma_i\), так как \(A_i\ma_i=(A_i+M)\ma_i\), следует \(M\ma_i=0\); следовательно, \(\ma_i\) можно умножать на любой класс. Поэтому из (14) следует представление (12), и оно единственно. В самом деле, если бы \(0=\mx_1\ma_1+\cdots+\mx_k\ma_k\), то умножением справа на \(\ma_i\), по (13), получили бы \(0=\mx_i\ma_i\). Тем самым доказаны единственность и само утверждение. Классы остатков \(\ma_1,\ldots,\ma_k\), возникающие при разложении единичного класса, мы называем единицами, принадлежащими этому разложению; они, таким образом, также задаются ортогональными соотношениями (13) и условием (14). Все дальнейшее развитие по существу сведется к вычислению с единицами.

Теперь остается установить связь между аддитивным разложением группы на \(k\) составляющих и соответствующим разложением модуля. Как и в случае \(k=2\), получится представление модуля как наименьшего общего кратного \(k\) взаимно простых модулей; доказывается это индукцией.

Пусть дано разложение (12), которое мы запишем в виде
\begin{align}
 \mx\me&=\mx\mb+\mx\ma_k,\tag{15}\\
 \mx\mb&=\mx\ma_1+\cdots+\mx\ma_{k-1}.\tag{16}
\end{align}
Из (15) по Теореме I следует, что классы остатков \(\mx\overline{\mb}\) изоморфны группе остатков модуля
\begin{equation}
 \mN=(\mM,A_k).
\tag{17}
\end{equation}
Для их классов остатков \(\mx\overline{\mb}\) поэтому также выполняется представление, соответствующее (16):
\begin{equation}
 \mx\overline{\mb}=\mx\overline{\ma}_1+\cdots+\mx\overline{\ma}_{k-1}.
\tag{18}
\end{equation}
Предположим, что уже доказано
\begin{equation}
 \mN=[\mM_1,\ldots,\mM_{k-1}],
\tag{19}
\end{equation}
где
\begin{equation}
\left\{
\begin{aligned}
 \mM_1&=(\mN,A_2+\cdots+A_{k-1}),\\
 \mM_2&=(\mN,A_1+A_3+\cdots+A_{k-1}),\\
 &\hspace{1.8cm}\vdots\\
 \mM_{k-1}&=(\mN,A_1+\cdots+A_{k-2}).
\end{aligned}\right.
\tag{20}
\end{equation}
Для двух слагаемых это утверждение совпадает с нашим прежним результатом. Кроме того, полиномы \(A_i\) можно здесь выбрать именно как такие полиномы из \(\overline{\ma}_i\), которые одновременно содержатся в \(\ma_i\), поскольку по \S~4 полиномы последнего класса остатков содержатся в первом. Тогда по (15)
\begin{equation}
 \mM=[\mN,\mM_k],
\tag{21}
\end{equation}
где
\[
 \mM_k=(\mM,A_1+\cdots+A_{k-1}).
\]
Из (19) и (21) следует
\[
 \mM=[[\mM_1,\ldots,\mM_{k-1}],\mM_k]=[\mM_1,\ldots,\mM_{k-1},\mM_k]
\]
по определению наименьшего общего кратного; а из (17) и (20) получается
\[
 \mM_1=(\mM,A_k,A_2+\cdots+A_{k-1}).
\]
Отсюда следует
\[
 \mM_1=(\mM,A_2+\cdots+A_k),
\]
ибо последний модуль, в силу (13), содержит также полином
\[
 A_k=A_k(A_2+\cdots+A_k)\pmod{\mM}.
\]
Соответствующие выражения получаются для \(\mM_2,\ldots,\mM_{k-1}\). Кроме того, модули \(\mM_1,\ldots,\mM_k\) взаимно просты, поскольку по (14)
\[
 \frac{1}{k-1}\Bigl((A_2+\cdots+A_k)+(A_1+A_3+\cdots+A_k)+\cdots+(A_1+\cdots+A_{k-1})\Bigr)\equiv1\pmod{\mM}.
\]
Более того, они образуют тотально взаимно простую систему, то есть наименьшее общее кратное любых из них взаимно просто с наименьшим общим кратным остальных.\footnote{Это обобщение понятия, введенного Лёви.} Это следует непосредственно, если в формуле (12) сгруппировать слагаемые в две подходящие группы.

Тем самым показано, что разложению группы остатков на \(k\) слагаемых соответствует представление модуля как наименьшего общего кратного \(k\) модулей, образующих тотально взаимно простую систему. Теперь докажем обратное утверждение, также индукцией.

Пусть
\[
 \mM=[\mM_1,\ldots,\mM_k]=[[\mM_1,\ldots,\mM_{k-1}],\mM_k]=[\mN,\mM_k],
\]
где \(\mM_1,\ldots,\mM_k\) образуют тотально взаимно простую систему. В сочетании с Теоремой I отсюда для группы остатков следует представление (15). Модули \(\mM_1,\ldots,\mM_{k-1}\) также образуют тотально взаимно простую систему. Ибо если бы при каком-нибудь разбиении модулей этой системы на две группы наименьшее общее кратное одной группы и наименьшее общее кратное другой имели общий делитель, то этот делитель сохранился бы и тогда, когда к одной из групп добавить еще модуль \(\mM_k\) и затем образовать наименьшее общее кратное.\footnote{Продолжая это рассуждение, получаем, что среди модулей \(\mM_1,\ldots,\mM_k\) каждые два взаимно просты; обратное утверждение, будто отсюда следует, что \(\mM_1,\ldots,\mM_k\) образуют тотально взаимно простую систему, вообще верно только в коммутативном случае; ср. контрпример 1 в \S~12.} Для \(k=2\) «тотально взаимно простые» означает то же, что и «взаимно простые»; поэтому наше утверждение можно считать доказанным до \(k-1\). Следовательно, для класса остатков \(\mx\overline{\mb}\) имеет место единственное представление (18), а в силу изоморфии также (16), что и доказывает утверждение.

\textbf{Теорема II.} \textit{Модуль \(\mM\) тогда и только тогда представим как наименьшее общее кратное \(k\) модулей \(\mM_1,\ldots,\mM_k\), образующих тотально взаимно простую систему, когда его группа остатков \(\mG\) приводима на \(k\) подгрупп \(\mU_1,\ldots,\mU_k\). При подходящем обозначении \(\mU_i\) изоморфна группе остатков модуля \(\mM_i\).}
% END UNIT BODY
% END PRODUCER UNIT 082
\endgroup


\clearpage
\begingroup
\setcounter{footnote}{0}
% BEGIN PRODUCER UNIT 083
% Source: C:/Users/Floris/Documents/interlanguage/03_projects/noether/02_slavic_working_corpus/translations/paper17/russian/v001/Noether_Paper17_Section06_Russian_v001.tex
% Identity: 11731 B / SHA-256 A0A7C3203256A4565DE206F24F7F108134A3856D4B9A1971CCB50B245E2E2EB2
\setlength{\parindent}{1.25em}
\setlength{\parskip}{0pt}
\emergencystretch=30em
\allowdisplaybreaks
\raggedbottom
\providecommand{\mM}{\mathfrak M}
\providecommand{\mG}{\mathfrak G}
\providecommand{\mS}{\mathfrak S}
\providecommand{\mP}{\mathfrak P}
\providecommand{\mQ}{\mathfrak Q}
\providecommand{\mU}{\mathfrak U}
\providecommand{\ma}{\mathfrak a}
\providecommand{\mx}{\mathfrak x}
% BEGIN UNIT BODY
% Unit: Paper 17 Section 06. Direct Russian translation from the German final-audited source slice.

\setcounter{footnote}{14}
\subsection*{\S~6. Теоремы конечности.}

В этом параграфе нужно показать, что каждую группу остатков можно представить как сумму конечного числа неприводимых составляющих. Для этого сначала перенесем гильбертову теорему о базисе модуля.

С этого места мы используем более специальные предположения, сформулированные в \S~2: умножение \(\xi_1,\ldots,\xi_n\) между собой коммутативно, а закон умножения полиномов имеет вид
\begin{equation}
 \xi_i a=a_i\xi_i+b_i,
 \qquad (a_i\ne0)\quad(i=1,\ldots,n).
\tag{6}
\end{equation}
Достаточно провести доказательство только для модулей. Действительно, если \(\mS\) -- какая-нибудь система полиномов, то мы образуем производный от нее модуль \(\mM\), присоединяя к \(\mS\) все полиномы, которые можно составить из конечного числа полиномов \(G_1,\ldots,G_k\) из \(\mS\) в виде \(A_1G_1+\cdots+A_kG_k\). Если теперь \(F_1,\ldots,F_l\) является базисом модуля \(\mM\) и
\[
 F_i=\sum_j A_{ij}G_j\qquad(i=1,\ldots,l),
\]
то конечное число полиномов \(G_j\), входящих в эти представления, образует базис системы \(\mS\).

Доказательство теоремы о базисе модуля мы проводим вслед за Горданом (\textit{Göttinger Nachrichten} 1899). Оно распадается на две части. Первая относится к системе \(\mP\) из бесконечно многих степенных произведений и, благодаря коммутативности \(\xi_1,\ldots,\xi_n\), по существу сохраняется. Степенные произведения
\[
 P=\xi_1^{a_1}\cdots\xi_n^{a_n},
\]
каждое из которых вносится в систему только один раз, упорядочим по возрастающей степени, а внутри одной степени -- каким-нибудь фиксированным образом. Рассмотрим теперь ту подсистему \(\mQ\) системы \(\mP\), которая содержит все такие степенные произведения \(P=Q\), не делящиеся ни на какое предшествующее \(P\). Тогда и никакое \(Q\) не делится на более позднее \(Q\), поскольку более поздние либо имеют большую степень, либо при той же степени отличаются от прежних. С другой стороны, каждое \(P\) из \(\mP\) делится хотя бы на одно \(Q\). В самом деле, если бы, напротив, \(P_m\) было первым \(P\), не делящимся ни на какое \(Q\), то \(P_m\) само не могло бы быть \(Q\); значит, оно делилось бы хотя бы на одно предшествующее \(P\), а это последнее по предположению делится на некоторое \(Q\). Следовательно, и само \(P_m\) делилось бы на \(Q\), что дает противоречие.

Теперь утверждаем, что число \(Q\) конечно. Действительно, пусть
\[
 Q_0=\xi_1^{h_1}\cdots\xi_n^{h_n}
\]
есть первое \(Q\). Так как произвольное \(Q_\lambda=\xi_1^{h_{1\lambda}}\cdots\xi_n^{h_{n\lambda}}\) не делится на \(Q_0\), то для каждого \(Q\) должен быть хотя бы один \(h_{\rho\lambda}<h_\rho\); пусть \(\lambda\) всякий раз обозначает наименьший индекс, для которого это выполнено. Объединим все \(Q\), для которых \(h_{\rho\lambda}\) имеет одно и то же фиксированное значение \(c_\lambda\), в один класс \(K(c_\lambda)\). Число этих классов конечно, поскольку \(c_\lambda<h_\lambda\). Но и в каждом классе содержится лишь конечное число элементов. В самом деле, если положить \(Q\) из класса \(K(c_\lambda)\) равным \(\xi_\lambda^{c_\lambda}Q'\), то \(Q'\) также образуют систему степенных произведений, ни одно из которых не делится на другое, и содержат только \(n-1\) переменных. При \(n=2\) произведения \(Q'\) содержат только одну переменную; тогда их число равно \(1\), следовательно конечно. Поэтому число можно считать конечным и для \(n-1\) переменных, что доказывает утверждение.

Во-вторых, пусть \(\mM\) -- модуль из полиномов, члены которых лексикографически упорядочены. Пусть \(\mP\) -- система степенных произведений, встречающихся каждый раз в старшем члене, каждое записанное только один раз. Тогда системе \(\mQ=(Q_1,\ldots,Q_\rho)\) соответствует система конечного числа полиномов \(F_1,\ldots,F_\rho\): каждому \(Q_i\) сопоставляется какой-нибудь полином из \(\mM\), старший член которого равен \(b_iQ_i\), так что можно написать \(F_i=b_iQ_i+\cdots\) с более низкими членами.

Пусть теперь \(F=bP+\cdots\) -- произвольный полином из \(\mM\), где \(P=\xi_1^{\beta_1}\cdots\xi_n^{\beta_n}Q_i\). Образуем произведение
\[
 \xi_1^{\beta_1}\cdots\xi_n^{\beta_n}F_i
 =\xi_1^{\beta_1}\cdots\xi_n^{\beta_n}(b_iQ_i+\cdots)
 =c\,\xi_1^{\beta_1}\cdots\xi_n^{\beta_n}Q_i+\cdots
 =cP+\cdots.
\]
Здесь \(c\ne0\) по (6), а все последующие члены являются более низкими членами лексикографического порядка, поскольку умножение, что касается показателей, коммутативно с точностью до более низких членов.\footnote{Отсюда видно, что вместо закона (6) можно было бы взять и несколько менее ограничительный закон
\[
 \xi_i a=a_i\xi_i+a_{i,i+1}\xi_{i+1}+\cdots+a_{in}\xi_n+b_i\qquad(a_i\ne0)
\]
при некоторой фиксированной нумерации \(\xi\).} Разность
\[
 F-\frac bc\,\xi_1^{\beta_1}\cdots\xi_n^{\beta_n}F_i
\]
содержит тогда только более низкие члены и снова принадлежит \(\mM\). Поэтому процесс обрывается через конечное число шагов, и \(F_1,\ldots,F_\rho\) являются базисом модуля.

\textbf{Теорема III.} \textit{Каждая система \(\mS\) из бесконечно многих полиномов имеет конечный базис модуля \(G_1,\ldots,G_k\), так что каждый полином \(G\) из \(\mS\) допускает представление \(G=A_1G_1+\cdots+A_kG_k\).}

Отсюда непосредственно следует

\textbf{Следствие.} \textit{Каждая система \(\mS\) из бесконечно многих классов остатков по модулю \(\mM\) имеет базис \(\mx_1,\ldots,\mx_k\), так что каждый \(\mx\) из \(\mS\) представим в форме \(\mx=A_1\mx_1+\cdots+A_k\mx_k\).}

А именно, каждому классу остатков поставим в соответствие какой-нибудь представитель и рассмотрим систему этих представителей вместе со всеми полиномами из \(\mM\).

Теперь перейдем к доказательству следующей теоремы.

\textbf{Теорема IV.} \textit{Каждая группа остатков представима как сумма конечного числа неприводимых составляющих; следовательно, по Теореме II каждый модуль представим как наименьшее общее кратное конечного числа неприводимых модулей, образующих тотально взаимно простую систему.}

Действительно, предположим противное: пусть \(\mG\) есть сумма бесконечно многих составляющих \(\mU_1,\mU_2,\ldots\) (ср. \S~5), и пусть \(\ma_1,\ma_2,\ldots\) -- соответствующие единицы, которые по определению все отличны от \(0\). Тогда образуем систему \(\mS\) всех этих единиц; по следствию из Теоремы III она имеет конечный базис \(\ma_{i_1},\ldots,\ma_{i_v}\) с \(i_1<\cdots<i_v\). Пусть теперь \(\mu>i_v\). Тогда имеем
\[
 \ma_\mu=\mx_1\ma_{i_1}+\cdots+\mx_v\ma_{i_v}.
\]
Это линейное отношение между единицами, которое по определению исключено. Поэтому при каждом аддитивном разложении группы остатков после конечного числа шагов появляется неприводимая составляющая; вследствие коммутативности разложения ее можно поставить в начало. Следовательно, без ограничения общности сами \(\mU_i\) можно считать неприводимыми. Но по приведенному выше рассуждению может появиться лишь конечное число составляющих \(\mU_i\), и тем самым утверждение доказано.
% END UNIT BODY
% END PRODUCER UNIT 083
\endgroup


\clearpage
\begingroup
\setcounter{footnote}{0}
% BEGIN PRODUCER UNIT 084
% Source: C:/Users/Floris/Documents/interlanguage/03_projects/noether/02_slavic_working_corpus/translations/paper17/russian/v001/Noether_Paper17_Section07_Russian_v001.tex
% Identity: 8567 B / SHA-256 790D93B2666EB02A61598D31EBE3C41DD104155DCAA75B3418C685C58BB35A26
\setlength{\parindent}{1.25em}
\setlength{\parskip}{0pt}
\emergencystretch=30em
\allowdisplaybreaks
\raggedbottom
\providecommand{\mM}{\mathfrak M}
\providecommand{\mG}{\mathfrak G}
\providecommand{\mU}{\mathfrak U}
\providecommand{\mB}{\mathfrak B}
\providecommand{\ma}{\mathfrak a}
\providecommand{\mb}{\mathfrak b}
\providecommand{\me}{\mathfrak e}
\providecommand{\mx}{\mathfrak x}
% BEGIN UNIT BODY
% Unit: Paper 17 Section 07. Direct Russian translation from the German final-audited source slice.

\setcounter{footnote}{15}
\subsection*{\S~7. Случай единственности разложения.}

Пусть
\[
 \mG=\mU_1+\cdots+\mU_k=\mB_1+\cdots+\mB_l
\]
есть группа остатков, составляющие которой \(\mU_1,\ldots,\mU_k\), \(\mB_1,\ldots,\mB_l\) неприводимы. Тогда имеем
\begin{equation}
 \mx\me=\mx\ma_1+\cdots+\mx\ma_k
 =\mx\mb_1+\cdots+\mx\mb_l,
\tag{22}
\end{equation}
где \(\ma_\chi\) есть класс остатков в \(\mG\), изоморфно соответствующий единичному классу \(\mU_\chi\), а \(\mb_\lambda\) -- класс остатков в \(\mG\), изоморфно соответствующий единичному классу \(\mB_\lambda\). Если теперь заменить \(\mx\) на \(\mx\mb_\lambda\), где \(\lambda\) -- фиксированный индекс из ряда \(1,\ldots,l\), то получаем
\begin{equation}
 \mx\mb_\lambda=\mx\mb_\lambda\ma_1+\cdots+\mx\mb_\lambda\ma_k.
\tag{23}
\end{equation}
Это представление группы \(\mB_\lambda\) единственно, но не является аддитивным разложением \(\mB_\lambda\), поскольку отдельные классы остатков в правой части не обязаны принадлежать \(\mB_\lambda\). Так как \(\mb_\lambda\) не обращается в нуль, не все произведения \(\mb_\lambda\ma_\chi\) могут быть нулевыми. Поэтому без ограничения общности пусть
\[
 \mb_\lambda\ma_1\ne0,
 \ldots,
 \mb_\lambda\ma_\rho\ne0,
 \qquad
 \mb_\lambda\ma_{\rho+1}=0,
 \ldots,
 \mb_\lambda\ma_k=0.
\]
Тогда
\[
 \mx\mb_\lambda=\mx\mb_\lambda\ma_1+\cdots+\mx\mb_\lambda\ma_\rho.
\]

Теперь при фиксированном \(\chi\) из ряда \(1,\ldots,\rho\) рассмотрим совокупность тех классов остатков \(\mx\mb_\lambda\), которые не равны нулю и для которых также \(\mx\mb_\lambda\ma_\chi\ne0\). Эти две системы, \(\mx\mb_\lambda\) и \(\mx\mb_\lambda\ma_\chi\), изоморфны. Действительно, во-первых, соответствие взаимно однозначно: из \(\mx\mb_\lambda\ma_\chi=0\) по предположению следует \(\mx\mb_\lambda=0\), а из последнего, в силу единственности представления (23), следует \(\mx\mb_\lambda\ma_\chi=0\). Кроме того, сумме соответствует сумма, а произведению с произвольным полиномом -- соответствующее произведение. При фиксированном \(\chi\) можно провести то же рассуждение для каждого \(\lambda\), для которого \(\mb_\lambda\ma_\chi\ne0\). Для каждого \(\ma_\chi\) должно существовать по крайней мере одно такое \(\mb_\lambda\), поскольку
\[
 \ma_\chi=(\mb_1+\cdots+\mb_l)\ma_\chi\ne0.
\]

Сначала предположим, что для каждого \(\chi\) существует только одно \(\lambda\), и для каждого \(\lambda\) только одно \(\chi\), такое что \(\mb_\lambda\ma_\chi\ne0\).

Из этого соответствия следует \(k=l\). Далее обозначения можно выбрать так, что \(\mb_\chi\ma_\chi\ne0\) и \(\mb_\lambda\ma_\chi=0\) при \(\lambda\ne\chi\). Тогда будем говорить, что произведения \(\mb_\lambda\ma_\chi\) образуют диагональную схему. В этом случае по (23) единица \(\mb_\chi=\mb_\chi\ma_\chi\), а из
\begin{equation}
 \ma_\chi=\mb_1\ma_\chi+\cdots+\mb_l\ma_\chi=\mb_\chi\ma_\chi=\mb_\chi
\tag{24}
\end{equation}
вообще следует \(\mU_\chi=\mB_\chi\). Значит, имеет место единственность разложения, и произведения \(\ma_\chi\mb_\lambda\) также образуют диагональную схему.

В частности, если для классов остатков \(\ma_\chi\mb_\lambda\) выполняется коммутативный закон, то указанное выше предположение всегда выполнено. Действительно, в (23) каждая составляющая \(\mx\mb_\lambda\ma_\chi=\mx\ma_\chi\mb_\lambda\) тогда содержится в \(\mB_\lambda\), так что (23) задает аддитивное разложение группы \(\mB_\lambda\). В силу неприводимости \(\mB_\lambda\) каждое \(\mb_\lambda\ma_\chi\), кроме одного, равно \(0\). Кроме того, по (24), при фиксированном \(\chi\) соответственно только одно \(\mb_\lambda\ma_\chi\) отлично от нуля, поскольку иначе (24), в силу перестановочности \(\ma_\chi\) и \(\mb_\lambda\), давало бы аддитивное разложение \(\mU_\chi\), которая должна была быть неприводимой.\footnote{Соответствующее утверждение верно в некоммутативном двустороннем случае, рассмотренном у \textit{Schmeidler} в цитированной работе, где произведение каждого класса остатков одной подгруппы с каждым классом остатков другой подгруппы обращается в нуль. Тогда из
\[
\mb_\lambda\ma_\chi=\mb_\lambda\ma_\chi\mb_1+\cdots+\mb_\lambda\ma_\chi\mb_l
\]
и из \(\mb_\lambda(\ma_\chi\mb_\mu)=0\) при \(\lambda\ne\mu\) следует также \(\mb_\lambda\ma_\chi=\mb_\lambda\ma_\chi\mb_\lambda\), то есть это произведение содержится в \(\mB_\lambda\). Поэтому из неприводимости \(\mB_\lambda\) следует, что только для одного \(\chi\), то есть ровно для одного \(\chi\), произведение \(\mb_\lambda\ma_\chi\ne0\). Так же показывают, что для каждого \(\ma_\chi\) существует один и только один такой \(\mb_\lambda\), что \(\ma_\chi\mb_\lambda\ne0\), откуда следует \(k=l\). Следовательно, \(\mb_\lambda\ma_\chi\) образуют диагональную схему, чем и доказана установленная там единственность.}

Те же рассуждения применимы и в случае, когда не для всех, а только для части \(\mb_\lambda\), а именно для \(\lambda=1,\ldots,\sigma\), выполнены условия \(\mb_\lambda\ma_\lambda\ne0\), \(\mb_\lambda\ma_\chi=0\) при \(\lambda\ne\chi\), \(\chi=1,\ldots,k\), и также для каждого \(\mu=1,\ldots,l\) условия \(\mb_\mu\ma_\chi=0\), \(\mb_\mu\ma_\mu\ne0\), \(\chi=1,\ldots,\sigma\). Тогда получаем \(\mU_\chi=\mB_\chi\) для \(\chi=1,\ldots,\sigma\), а если положить
\[
 \mU_{\sigma+1}+\cdots+\mU_k=\mU,
 \qquad
 \mB_{\sigma+1}+\cdots+\mB_l=\mB,
\]
то и для единиц \(\ma\) и \(\mb\) групп \(\mU\) и \(\mB\) выполнено условие \(\mb\ma_\chi=0\), \(\mb_\chi\ma=0\) для \(\chi=1,\ldots,\sigma\), а также \(\mb\ma\ne0\), откуда следует \(\mU=\mB\).
% END UNIT BODY
% END PRODUCER UNIT 084
\endgroup


\clearpage
\begingroup
\setcounter{footnote}{0}
% BEGIN PRODUCER UNIT 085
% Source: C:/Users/Floris/Documents/interlanguage/03_projects/noether/02_slavic_working_corpus/translations/paper17/russian/v001/Noether_Paper17_Section08_Russian_v001.tex
% Identity: 8558 B / SHA-256 40B9AF75A0CC2BE38BD6BDB28E838B474901B7DD42BE0D47564B374D93E51969
\setlength{\parindent}{1.25em}
\setlength{\parskip}{0pt}
\emergencystretch=30em
\allowdisplaybreaks
\raggedbottom
\providecommand{\mM}{\mathfrak M}
\providecommand{\mG}{\mathfrak G}
\providecommand{\mU}{\mathfrak U}
\providecommand{\mB}{\mathfrak B}
\providecommand{\ma}{\mathfrak a}
\providecommand{\mb}{\mathfrak b}
\providecommand{\me}{\mathfrak e}
\providecommand{\mx}{\mathfrak x}
% BEGIN UNIT BODY
% Unit: Paper 17 Section 08. Direct Russian translation from the German final-audited source slice.

\setcounter{footnote}{16}
\subsection*{\S~8. Аддитивное разложение вполне приводимых групп на простые группы. Теорема об изоморфии.}

Теперь рассмотрим другой случай: единственности в нем, правда, нет, зато есть изоморфия различных разложений. Здесь речь идет об обобщении случая, рассмотренного Лёви.

\textbf{Определение.} \textit{Модуль называется простым модулем, если он не имеет других делителей, кроме самого себя и единичного модуля, содержащего все многочлены. Группа остатков простого модуля называется простой группой.}

Всякая простая группа a fortiori неприводима. Существуют также бесконечные простые группы (см. определение \S~11 и пример \S~12, 2).

Пусть теперь данная группа остатков \(\mG\) двумя способами представима как сумма (в каждом случае конечного числа) простых групп \(\mU_1,\ldots,\mU_k\) соответственно \(\mB_1,\ldots,\mB_l\). Группу остатков, допускающую хотя бы одно такое представление как сумму простых групп, и соответствующий ей модуль \(\mM\) мы, следуя обозначению Лёви, будем называть вполне приводимыми.

Теперь рассмотрим произведения \(\ma_\chi\mb_\lambda\), где \(\chi=1,\ldots,k\), \(\lambda=1,\ldots,l\), и покажем, что либо \(\ma_\chi\mb_\lambda=0\), либо \(\mx\ma_\chi\mb_\lambda\) пробегает всю группу \(\mB_\lambda\), когда \(\mx\) пробегает все классы остатков \(\mG\). Действительно, модуль, принадлежащий \(\mB_\lambda\),
\[
 (\mM,\mb_1+\cdots+\mb_{\lambda-1}+\mb_{\lambda+1}+\cdots+\mb_l)=\mM_{\mb_\lambda},
\]
который получается из \(\mM\) присоединением всех многочленов класса остатков \(\mb_1+
\cdots+\mb_{\lambda-1}+\mb_{\lambda+1}+\cdots+\mb_l\), имеет делитель
\[
 (\mM,\mb_1+\cdots+\mb_{\lambda-1}+\mb_{\lambda+1}+\cdots+\mb_l,
 \ma_\chi\mb_\lambda).
\]
Но так как модуль \(\mM_{\mb_\lambda}\) по предположению является простым модулем, этот делитель либо равен \(\mM_{\mb_\lambda}\), либо является единичным модулем. В первом случае \(\ma_\chi\mb_\lambda\) принадлежит \(\mM_{\mb_\lambda}\); значит, существует соотношение
\[
 \ma_\chi\mb_\lambda=\mx(\mb_1+
\cdots+\mb_{\lambda-1}+\mb_{\lambda+1}+\cdots+\mb_l),
\]
откуда следует \(\ma_\chi\mb_\lambda=0\). В другом случае существуют два класса \(\mx_0\) и \(\mx_1\), такие что
\[
 \mx_0\ma_\chi\mb_\lambda
 +\mx_1(\mb_1+\cdots+\mb_{\lambda-1}+\mb_{\lambda+1}+\cdots+\mb_l)
 =\me=\mb_1+\cdots+\mb_{\lambda-1}+\mb_{\lambda+1}+\cdots+\mb_l+\mb_\lambda,
\]
и потому, в силу единственности, \(\mx_0\ma_\chi\mb_\lambda=\mb_\lambda\). Следовательно, \(F\mx_0\ma_\chi\mb_\lambda\), а значит и \(\mx\ma_\chi\mb_\lambda\), пробегает всю группу \(\mB_\lambda\), когда \(F\) соответственно \(\mx\) пробегает все многочлены соответственно все классы остатков \(\mG\). Тем самым утверждение доказано.

Кроме того, если \(\ma_\chi\mb_\lambda\ne0\), то \(\mU_\chi\) изоморфна \(\mB_\lambda\). Действительно, совокупность всех классов \(\mx\ma_\chi\ne0\), для которых также \(\mx\ma_\chi\mb_\lambda\ne0\), по \S~7 изоморфна системе классов \(\mx\ma_\chi\mb_\lambda\); но эти классы исчерпывают подгруппу \(\mB_\lambda\). Остается лишь показать, что названные классы \(\mx\ma_\chi\) исчерпывают группу \(\mU_\chi\). Для этого рассмотрим все классы остатков \(\mx\ma_\chi\), для которых \(\mx\ma_\chi\mb_\lambda=0\). Эта система обладает свойством, что сумма двух ее классов и произведение с произвольным многочленом снова принадлежат системе. Поэтому, если присоединить все эти классы остатков к модулю \(\mM_{\ma_\chi}\), снова возникает модуль, являющийся делителем \(\mM_{\ma_\chi}\), то есть либо равный \(\mM_{\ma_\chi}\), либо являющийся единичным модулем. Но в последнем случае сама \(\ma_\chi\) содержалась бы в нем, так что \(\ma_\chi\mb_\lambda=0\), вопреки предположению. Следовательно, нет таких классов \(\mx\ma_\chi\ne0\), для которых \(\mx\ma_\chi\mb_\lambda=0\), и изоморфия \(\mU_\chi\) с \(\mB_\lambda\) доказана.

Пусть теперь для некоторого \(\chi\), например,
\[
 \ma_\chi\mb_{\lambda_1}\ne0,
 \ldots,\ma_\chi\mb_{\lambda_i}\ne0\qquad(i>1).
\]
Тогда \(\mU_\chi\) изоморфна группам \(\mB_{\lambda_1},\ldots,\mB_{\lambda_i}\), следовательно, эти группы изоморфны между собой.

Теперь покажем, что тогда и по меньшей мере две из групп \(\mU\) изоморфны. Действительно, если бы в произведениях \(\mb_\lambda\ma_\chi\) каждой \(\mb_\lambda\) соответствовала только одна \(\ma_\chi\), так что \(\mb_\lambda\ma_\chi\ne0\), то есть если бы \(\mb_\lambda=\mb_\lambda\ma_\chi\), то было бы \(\mB_\lambda=\mU_\chi\), поскольку \(\mU_\chi\) как простая группа не может содержать собственной подгруппы. Тогда было бы \(k=l\), и при правильной нумерации \(\ma_\chi=\mb_\chi\). Но тогда и схема \(\ma_\chi\mb_\lambda\) была бы диагональной схемой, чего нет. Тем самым доказана следующая теорема:

\textbf{Теорема V.} \textit{Если вполне приводимая группа \(\mG\) двумя способами представлена как сумма простых групп:
\[
 \mG=\mU_1+\cdots+\mU_k=\mB_1+\cdots+\mB_l,
\]
то каждая группа \(\mU\) изоморфна по меньшей мере одной группе \(\mB\), и наоборот. Если каждая \(\mU\) изоморфна ровно одной \(\mB\), то разложения тождественны. В противном случае по меньшей мере две группы \(\mU\) изоморфны между собой, и точно так же по меньшей мере две группы \(\mB\) изоморфны между собой.}
% END UNIT BODY
% END PRODUCER UNIT 085
\endgroup


\clearpage
\begingroup
\setcounter{footnote}{0}
% BEGIN PRODUCER UNIT 086
% Source: C:/Users/Floris/Documents/interlanguage/03_projects/noether/02_slavic_working_corpus/translations/paper17/russian/v001/Noether_Paper17_Section09_Russian_v001.tex
% Identity: 9038 B / SHA-256 30A3FC807FE47C63267C2800CDB2300CB347E000D1D895DF6268FAC99A4F2180
\setlength{\parindent}{1.25em}
\setlength{\parskip}{0pt}
\emergencystretch=30em
\allowdisplaybreaks
\raggedbottom
\providecommand{\mM}{\mathfrak M}
\providecommand{\mN}{\mathfrak N}
\providecommand{\mA}{\mathfrak A}
\providecommand{\mB}{\mathfrak B}
\providecommand{\mP}{\mathfrak P}
\providecommand{\mD}{\mathfrak D}
\providecommand{\mS}{\mathfrak S}
\providecommand{\ma}{\mathfrak a}
\providecommand{\mb}{\mathfrak b}
% BEGIN UNIT BODY
% Unit: Paper 17 Section 09. Direct Russian translation from the German final-audited source slice.

\setcounter{footnote}{16}
\subsection*{\S~9. Связь между понятиями «изоморфный» и «того же вида».}

В этом параграфе мы показываем, что известное для дифференциальных выражений от одной переменной понятие «того же вида», если его естественно обобщить, приводит к изоморфии соответствующих групп остатков.

\textbf{Определение.}\footnote{Для дифференциальных выражений от одной переменной ср. формальное определение Блумберга, там же, с. 25.} \emph{Два модуля \(\mM\) и \(\mN\) называются модулями того же вида, если существуют два полинома \(P\) и \(Q\), причем \(P\) взаимно прост с \(\mM\), а \(Q\) взаимно прост с \(\mN\), такие, что для всякого \(M\equiv0(\mM)\) и \(N\equiv0(\mN)\)}
\begin{align}
 MQ&\equiv0(\mN),\tag{25}\\
 NP&\equiv0(\mM)\tag{26}
\end{align}
\emph{имеет место. Взаимная простота здесь означает существование двух полиномов \(X\) и \(Y\), таких что}
\begin{align}
 YQ&\equiv1(\mN),\tag{27}\\
 XP&\equiv1(\mM)\tag{28}
\end{align}
\emph{имеет место.}

Теперь справедлива следующая

\textbf{Теорема VI.} \emph{Два модуля тогда и только тогда имеют тот же вид, когда их группы остатков изоморфны.}

1. Пусть \(\mA\) есть группа остатков модуля \(\mM\), \(\mB\) -- группа остатков модуля \(\mN\), и пусть \(\mA\) изоморфна \(\mB\); пусть \(\ma\) и \(\mb\) -- единичные классы в \(\mA\) и \(\mB\). Далее пусть \(Q\mb\) есть класс в \(\mB\), соответствующий единице \(\ma\) в \(\mA\), и так же \(P\ma\) -- класс в \(\mA\), соответствующий единице \(\mb\) в \(\mB\). Это соответствие запишем в виде
\begin{align}
 \ma&\sim Q\mb,\tag{29}\\
 P\ma&\sim\mb.\tag{30}
\end{align}
Отсюда следует
\[
 0=M\ma\sim MQ\mb,
 \qquad\text{т. е.}\qquad MQ\equiv0(\mN),
\]
чем доказано (25). Аналогично имеем
\[
 NP\ma\sim N\mb=0,
 \qquad\text{т. е.}\qquad NP\equiv0(\mM),
\]
то есть соотношение (26). Далее, по (29), \(P\ma\sim PQ\mb\), а так как по (30) было предположено \(P\ma\sim\mb\), то из единственности соответствия следует \(PQ\mb=\mb\), следовательно \(PQ\equiv1(\mN)\) и соответственно \(QP\equiv1(\mM)\). Тем самым (27) и (28) доказаны при \(X=Q\), \(Y=P\).

2. Пусть, наоборот, \(\mM\) и \(\mN\) имеют тот же вид, так что выполнены соотношения (25)--(28). Тогда произвольному классу \(F\ma\) из \(\mA\) сопоставим класс \(FQ\mb\) из \(\mB\). При этом каждому классу в \(\mA\) соответствует только один класс в \(\mB\): действительно, из \(F\ma=0\), то есть \(F=M\), по (25) следует \(FQ\mb=0\). Кроме того, это соответствие исчерпывает всю группу \(\mB\), потому что специальному классу \(Y\ma\) соответствует класс \(YQ\mb\), равный \(\mb\) по (27). Итак, мы имеем однозначно определенное отображение \(\Phi\) из \(\mA\) на \(\mB\), но еще не доказано, что оно взаимно однозначно. Точно так же из равенств (26) и (28) следует однозначно определенное отображение \(\Psi\) из \(\mB\) на \(\mA\), которое также еще не доказано взаимно однозначным.

Если одно из этих двух отображений взаимно однозначно, то изоморфия доказана, поскольку требования относительно суммы и произведения выполнены. Но если бы оба отображения \(\Phi\) и \(\Psi\) не были взаимно однозначными, то существовали бы, например, в \(\mA\) классы \(G\ma\), которым в \(\mB\) соответствует класс \(GQ\mb=0\). Классы \(FQ\mb\), соответствующие всем остальным классам остатков \(F\ma\), тогда уже исчерпывали бы группу \(\mB\), и потому отображение \(\Psi\) возвращало бы в классах \(FQP\ma\) всю группу \(\mA\). Как мы знаем, существуют также в \(\mB\) отличные от нуля классы, которым посредством \(\Psi\) соответствует нулевой класс в \(\mA\); обозначим через \(G_1\ma\) все те классы, для которых \(G_1\ma\ne0\), но \(G_1QP\ma=0\). Все остальные полиномы \(F\) обладают свойством, что \(FQP\ma\) исчерпывает группу \(\mA\). Те полиномы \(F\), для которых \(FQP\ma=G_1\ma\), обозначим через \(G_2\); такие полиномы всегда существуют, как только существуют полиномы \(G_1\). Для каждого \(G_2\) далее имеем \(G_2QP\ma\ne0\), зато \(G_2(QP)^2\ma=0\). Соответственно для каждого целого \(\nu\) существуют классы \(G_\nu\ma\), такие что \(G_\nu(QP)^{\nu-1}\ma\ne0\), но \(G_\nu(QP)^\nu\ma=0\).

Теперь рассмотрим систему \(\mS\), состоящую из всех полиномов \(G_1,G_2,\ldots\). Она имеет конечный базис модуля \(G_{\lambda_1},\ldots,G_{\lambda_e}\). Теперь выберем \(\nu\) больше наибольшего из индексов \(\lambda_1,\ldots,\lambda_e\), который обозначим через \(\lambda\). Тогда
\[
 G_\nu=A_1G_{\lambda_1}+\cdots+A_eG_{\lambda_e}.
\]
Умножая это на \((QP)^\lambda\), получаем
\[
 G_\nu(QP)^\lambda\equiv0(\mM),
\]
что противоречит нашему предположению. Следовательно, получается изоморфия, а значит, по пункту 1 справедливы и более сильные соотношения (27), (28) при \(X=Q\), \(Y=P\).

В силу Теоремы VI Теорему V можно также сформулировать следующим образом:

\textbf{Теорема VII.}\footnote{Ср. Loewy, с. 100--101.} \emph{Если вполне приводимый модуль \(\mM\) двумя способами представим как наименьшее общее кратное простых модулей \(\mP_1,\ldots,\mP_k\) соответственно \(\mD_1,\ldots,\mD_l\), каждый раз образующих тотально взаимно простую систему, то каждый \(\mP\) имеет тот же вид по меньшей мере с одним \(\mD\), и наоборот. Если каждый \(\mP\) имеет тот же вид ровно с одним \(\mD\), то разложения тождественны. В противном случае по меньшей мере два из модулей \(\mP\) имеют тот же вид между собой, и точно так же по меньшей мере два из модулей \(\mD\) имеют тот же вид между собой.}
% END UNIT BODY
% END PRODUCER UNIT 086
\endgroup


\clearpage
\begingroup
\setcounter{footnote}{0}
% BEGIN PRODUCER UNIT 087
% Source: C:/Users/Floris/Documents/interlanguage/03_projects/noether/02_slavic_working_corpus/translations/paper17/russian/v001/Noether_Paper17_Section10_Russian_v001.tex
% Identity: 21735 B / SHA-256 4CDD48B4DAF539F4DAAAB70266F4AD183911B65D536D437B35035A9AF962D092
\setlength{\parindent}{1.25em}
\setlength{\parskip}{0pt}
\emergencystretch=30em
\allowdisplaybreaks
\raggedbottom
\providecommand{\mM}{\mathfrak M}
\providecommand{\mN}{\mathfrak N}
\providecommand{\mA}{\mathfrak A}
\providecommand{\mB}{\mathfrak B}
\providecommand{\mP}{\mathfrak P}
\providecommand{\mD}{\mathfrak D}
\providecommand{\mS}{\mathfrak S}
\providecommand{\ma}{\mathfrak a}
\providecommand{\mb}{\mathfrak b}
% BEGIN UNIT BODY
% Unit: Paper 17 Section 10. Direct Russian translation from the German final-audited source slice.

\setcounter{footnote}{19}
\subsection*{\S~10. Существование бесконечно многих разложений группы.}

Если представление вполне приводимой группы в виде суммы не единственно, то по Теореме V в каждом представлении существуют по меньшей мере две изоморфные между собой подгруппы. Рассмотрим теперь подгруппу, образованную такой системой изоморфных подгрупп,
\[
  \mathfrak H=\mathfrak A_1+\cdots+\mathfrak A_\alpha,
\]
где, следовательно, \(\mathfrak A_1,\mathfrak A_2,\ldots,\mathfrak A_\alpha\) все изоморфны, и утверждаем, что \(\mathfrak H\) тогда допускает бесконечно много различных таких разложений, если только в области рациональности \(P\) содержится бесконечно много ``констант'', то есть величин \(a\), для которых \(\xi_i a=a\xi_i\). Эта теорема справедлива даже тогда, когда о группах \(\mathfrak A_i\) больше ничего не предполагается, даже неприводимость. Итак, имеет место

\textbf{Теорема VIII.} \emph{Если \(\mathfrak G\) представима как сумма конечного числа попарно изоморфных подгрупп, то такое представление возможно бесконечно многими различными способами, как только область рациональности \(P\) содержит бесконечно много констант. При этом под константами следует понимать такие величины \(a\) области, для которых \(\xi_i a=a\xi_i\).}

\textbf{Доказательство.} Пусть, значит,
\[
  \mathfrak G=\mathfrak A_1+\cdots+\mathfrak A_\alpha,
\]
и \(\mathfrak A_i\) изоморфна \(\mathfrak A_\varkappa\) для \(i,\varkappa=1,\ldots,\alpha\) (\(\alpha\ge2\)). Чтобы задать определенное изоморфное соответствие, сначала выделим одну подгруппу, например \(\mathfrak A_1\). Пусть изоморфия между \(\mathfrak A_1\) и \(\mathfrak A_i\) задана соотношениями
\begin{equation}
 a_1\sim t_{1i}a_i,\qquad a_i\sim t_{i1}a_1,\qquad t_{11}a_1=a_1
 \quad (i,\varkappa=1,\ldots,\alpha),
\tag{31}
\end{equation}
где, следовательно, в группе \(\mathfrak A_1\) каждый класс остатков сопоставлен самому себе. Так как \(Ma_1=0\), \(Ma_i=0\), отсюда следует \(Mt_{1i}a_i=0\), \(Mt_{i1}a_1=0\); поэтому классы \(t_{1i}a_i\), \(t_{i1}a_1\), подобно единицам, допускают умножение не только на многочлены, но и на классы. Из определения изоморфии тогда следует, что для двух соответствующих классов \(\mathfrak y\) и \(\bar{\mathfrak y}\), допускающих это умножение классов, классы \(\mathfrak r\mathfrak y\) и \(\mathfrak r\bar{\mathfrak y}\) снова соответствуют друг другу. Поэтому соотношения \(a_ra_s=0\) \((r\ne s)\) (\(r=1,\ldots,\alpha\); \(s=1,
\ldots,\alpha\)) при \(s=1\) дают
\begin{equation}
\left\{
\begin{aligned}
 a_r t_{1\varkappa}a_\varkappa&=0 &&(\varkappa=1,
\ldots,\alpha,
\ r\ne1),\\
 \text{а при }s>1\qquad a_r t_{s1}a_1&=0 &&(r\ne s),\\
 \text{следовательно также}\qquad t_{1i}a_i&=a_1t_{1i}a_i,
 & t_{i1}a_1&=a_it_{i1}a_1.
\end{aligned}\right.
\tag{32}
\end{equation}
Далее из (31), как обобщение (27) и (28), из единственности соответствия следует
\begin{equation}
 t_{1i}t_{i1}a_1=a_1,
 \qquad
 t_{\varkappa1}t_{1\varkappa}a_\varkappa=a_\varkappa.
\tag{33}
\end{equation}
В качестве изоморфии между \(\mathfrak A_i\) и \(\mathfrak A_\varkappa\) теперь зафиксируем отношение, возникающее композицией (31); такая фиксация необходима, поскольку отдельные группы могут допускать изоморфии в себя. Из (31), (32) и (33) поэтому получаем
\begin{equation}
 a_i\sim t_{i1}t_{1\varkappa}a_\varkappa=t_{i\varkappa}a_\varkappa,
 \qquad
 t_{i\varkappa}a_i=a_i,
 \qquad
 a_rt_{s\varkappa}a_\varkappa=0\quad (r\ne s),
\tag{34}
\end{equation}
где положено \(t_{i1}t_{1\varkappa}=t_{i\varkappa}\). Из (33) и (34) далее следует
\[
 t_{i\varkappa}t_{\varkappa l}a_l=t_{i1}t_{1\varkappa}t_{\varkappa1}t_{1l}a_l
 =t_{i1}t_{1l}a_l=t_{il}a_l;
\]
тем самым выделение группы \(\mathfrak A_1\) снова устранено. Изоморфное отношение между \(\mathfrak A_i\) и \(\mathfrak A_\varkappa\) остается тем же самым, какая бы группа \(\mathfrak A_\varkappa\) ни использовалась вместо \(\mathfrak A_1\) для определения. Значит, \(t_{i\varkappa}\) являются классами, которые вместе удовлетворяют соотношениям
\begin{equation}
 Mt_{i\varkappa}a_\varkappa=0;
 \quad a_rt_{i\varkappa}a_\varkappa=0\ (r\ne i);
 \quad t_{i\varkappa}t_{\varkappa l}a_l=t_{il}a_l;
 \quad t_{\varkappa i}a_i=t_{\varkappa i}t_{i\varkappa}a_\varkappa=a_\varkappa.
\tag{35}
\end{equation}
Эти соотношения говорят в точности о том, что модули \(\mathfrak M_i\) и \(\mathfrak M_\varkappa\), принадлежащие \(\mathfrak A_i\) и \(\mathfrak A_\varkappa\), имеют тот же вид (ср. определение \S~9); вместо тамошних \(P\) и \(Q\) здесь стоят классы \(t_{i\varkappa}\) соответственно \(t_{\varkappa i}\), а первые два соотношения (35), взятые вместе, соответствуют формуле (25) соответственно (26). Обратно, из этих соотношений по Теореме VI следует изоморфия.\footnote{Причем без применения теоремы конечности, так как формулы уже каждый раз задают отображение \(\varphi_{i\varkappa}\) и обратное к нему \(\varphi_{\varkappa i}\). Действительно, из \(ra_i=ra_it_{i\varkappa}a_\varkappa t_{\varkappa i}a_i\) следует, что из \(ra_i\ne0\) также \(ra_it_{i\varkappa}a_\varkappa\ne0\).}

Из классов \(t_{i\varkappa}a_\varkappa\) теперь можно составить классы \(b_\rho\), которые снова окажутся единицами групп \(\mathfrak B_\rho\), откуда будет следовать представление
\[
 \mathfrak G=\mathfrak B_1+\cdots+\mathfrak B_\alpha.
\]
Пусть именно \(\lambda_{i\varkappa}\) являются константами, определитель которых \(|\lambda_{i\varkappa}|\ne0\), а \(\lambda^{\varkappa r}\) --- соответствующие миноры, деленные на \(|\lambda_{i\varkappa}|\), так что выполнены соотношения
\begin{equation}
\left\{
\begin{aligned}
 \sum_{r=1}^{\alpha}\lambda_{ri}\lambda^{r\varkappa}&=\delta_{i\varkappa}=\begin{cases}0,&i\ne\varkappa,\\1,&i=\varkappa,
 \end{cases}\\
 \sum_{r=1}^{\alpha}\lambda_{ir}\lambda^{\varkappa r}&=\delta_{i\varkappa}.
\end{aligned}\right.
\tag{36}
\end{equation}
Тогда определим
\begin{equation}
 b_\rho=\sum_{i,\varkappa}\lambda_{\rho i}\lambda^{\rho\varkappa}t_{i\varkappa}a_\varkappa
 \qquad(\rho=1,
\ldots,\alpha).
\tag{37}
\end{equation}
Из (35), (36) и (37) следует \(Mb_\rho=0\) и
\begin{equation}
 \sum_{\rho=1}^{\alpha}b_\rho
 =\sum_{i,\varkappa,\rho}\lambda_{\rho i}\lambda^{\rho\varkappa}t_{i\varkappa}a_\varkappa
 =\sum_{i,\varkappa}\delta_{i\varkappa}t_{i\varkappa}a_\varkappa
 =\sum_i t_{ii}a_i=\sum_i a_i=e,
\tag{38}
\end{equation}
а также
\[
\begin{aligned}
 b_\rho b_\sigma
 &=\sum_{\substack{i,\varkappa\\ \mu,\nu}}
  \lambda_{\rho i}\lambda^{\rho\varkappa}\lambda_{\sigma\mu}\lambda^{\sigma\nu}
  t_{i\varkappa}a_\varkappa t_{\mu\nu}a_\nu \\
 &=\sum_{i,\mu,\nu}\lambda_{\rho i}\lambda^{\rho\mu}\lambda_{\sigma\mu}\lambda^{\sigma\nu}t_{i\nu}a_\nu
 =\delta_{\rho\sigma}\sum_{i,\nu}\lambda_{\rho i}\lambda^{\sigma\nu}t_{i\nu}a_\nu
 =\begin{cases}0&(\rho\ne\sigma),\\ b_\rho&(\rho=\sigma),\end{cases}
\end{aligned}
\]
\((\rho,\sigma=1,
\ldots,\alpha)\). Следовательно, представление \(\mathfrak r e=\mathfrak r b_1+\cdots+\mathfrak r b_\alpha\) является аддитивным разложением группы \(\mathfrak G\). При этом группы \(\mathfrak B_\rho\) отличаются от групп \(\mathfrak A_\sigma\), как только \(\lambda_{i\varkappa}\) не образуют одну лишь диагональную матрицу. Действительно, если образовать \(a_\sigma b_\rho\), то по (35) получается:
\[
 a_\sigma b_\rho=
 \sum_{i,\varkappa}\lambda_{\rho i}\lambda^{\rho\varkappa}a_\sigma t_{i\varkappa}a_\varkappa
 =\lambda_{\rho\sigma}
 \sum_{\varkappa}\lambda^{\rho\varkappa}t_{\sigma\varkappa}a_\varkappa\ne0
 \quad\text{для }\lambda_{\rho\sigma}\ne0,
\]
ибо иначе в последней сумме должен был бы исчезать каждый отдельный член, чего нет. Значит, если при фиксированном \(\sigma\) существуют по меньшей мере два \(\rho\), для которых \(\lambda_{\rho\sigma}\ne0\), то \(\mathfrak A_\sigma\) не может быть тождественна никакой \(\mathfrak B_\rho\).

Но группы \(\mathfrak B_\rho\) также изоморфны между собой и группам \(\mathfrak A_\sigma\). Для этого сначала покажем существование классов \(r_{\rho\sigma}a_\sigma\) и \(r^{\sigma\rho}b_\rho\), которые осуществляют изоморфию \(\mathfrak A_\sigma\) и \(\mathfrak B_\rho\) (\(\rho,\sigma=1,
\ldots,\alpha\)). Именно, мы докажем, что для классов
\[
 r_{\rho\sigma}a_\sigma=\sum_i\lambda_{\rho i}t_{i\sigma}a_\sigma,
 \qquad
 r^{\sigma\rho}=\sum_\varkappa\lambda^{\rho\varkappa}t_{\sigma\varkappa}a_\varkappa
 =r^{\sigma\rho}b_\rho
\]
выполнены соотношения ``того же вида'' между модулями
\[
 (\mathfrak M,b_1+\cdots+b_{\rho-1}+b_{\rho+1}+\cdots+b_\alpha)
\]
и
\[
 (\mathfrak M,a_1+\cdots+a_{\sigma-1}+a_{\sigma+1}+\cdots+a_\alpha):
\]
\begin{equation}
\left\{
\begin{aligned}
 b_\varkappa r_{\rho\sigma}a_\sigma&=0 &&(\varkappa\ne\rho),&
 r^{\sigma\rho}r_{\rho\sigma}a_\sigma&=a_\sigma,\\
 a_\varkappa r^{\sigma\rho}b_\rho&=0 &&(\varkappa\ne\sigma),&
 r_{\rho\sigma}r^{\sigma\rho}b_\rho&=b_\rho,
\end{aligned}\right.
\tag{39}
\end{equation}
откуда, как и при (35), изоморфия следует посредством соответствия
\begin{equation}
 b_\rho\sim r_{\rho\sigma}a_\sigma,
 \qquad
 a_\sigma\sim r^{\sigma\rho}b_\rho.
\tag{40}
\end{equation}
В самом деле, получается:
\[
\begin{aligned}
 b_\varkappa r_{\rho\sigma}a_\sigma
 &=\sum_{\mu,\nu,i}\lambda_{\varkappa\mu}\lambda^{\varkappa\nu}t_{\mu\nu}a_\nu\lambda_{\rho i}t_{i\sigma}a_\sigma\\
 &=\sum_{\mu,i}\lambda_{\varkappa\mu}\lambda^{\varkappa i}\lambda_{\rho i}t_{\mu\sigma}a_\sigma
 =\delta_{\varkappa\rho}\sum_\mu\lambda_{\varkappa\mu}t_{\mu\sigma}a_\sigma
 =\delta_{\varkappa\rho}r_{\rho\sigma}a_\sigma=0\quad(\varkappa\ne\rho),
\end{aligned}
\]
далее
\[
 r^{\sigma\rho}r_{\rho\sigma}a_\sigma
 =\sum_{\nu,i}\lambda^{\rho\nu}t_{\sigma\nu}a_\nu\lambda_{\rho i}t_{i\sigma}a_\sigma
 =\sum_i\lambda^{\rho i}\lambda_{\rho i}a_\sigma=a_\sigma.
\]
Точно так же в обратную сторону находим
\[
 a_\varkappa r^{\sigma\rho}b_\rho
 =a_\varkappa\sum_\mu\lambda^{\rho\mu}t_{\sigma\mu}a_\mu=0
 \quad(\varkappa\ne\sigma),
\]
и далее
\[
 r_{\rho\sigma}r^{\sigma\rho}b_\rho
 =\sum_{i,\mu}\lambda_{\rho i}t_{i\sigma}a_\sigma\lambda^{\rho\mu}t_{\sigma\mu}a_\mu
 =\sum_{i,\mu}\lambda_{\rho i}\lambda^{\rho\mu}t_{i\mu}a_\mu=b_\rho.
\]
Из доказанных здесь соотношений того же вида теперь следует (опять без теоремы конечности) изоморфия \(\mathfrak A_\sigma\) и \(\mathfrak B_\rho\). Изоморфия между \(\mathfrak B_\rho\) и \(\mathfrak B_\tau\) находится отсюда композицией; соответствующие формулы мы еще укажем. Положим
\[
 \ell_{\rho\tau}b_\tau=r_{\rho\sigma}r^{\sigma\tau}b_\tau
 =\sum_{i,j}\lambda_{\rho i}t_{i\sigma}a_\sigma\lambda^{\tau j}t_{\sigma j}a_j
 =\sum_{i,j}\lambda_{\rho i}\lambda^{\tau j}t_{ij}a_j;
\]
тогда соответствие имеет вид \(b_\rho\sim\ell_{\rho\tau}b_\tau\). И здесь соотношения того же вида снова можно формально проверить; \(\ell_{\rho\tau}\), поставленные на место \(t_{i\varkappa}\), удовлетворяют соотношениям (35). В самом деле,
\[
 b_\sigma\ell_{\rho\tau}b_\tau=b_\sigma r_{\rho\sigma}r^{\sigma\tau}b_\tau=0\quad(\sigma\ne\rho),
\]
\[
 \ell_{\rho\sigma}\ell_{\sigma\tau}b_\tau=r_{\rho\varkappa}r^{\varkappa\sigma}r_{\sigma\lambda}r^{\lambda\tau}b_\tau
 =r_{\rho\varkappa}a_\varkappa r^{\varkappa\tau}b_\tau=\ell_{\rho\tau}b_\tau.
\]

Тем самым Теорема VIII доказана, так как нужно лишь выбирать \(\lambda_{i\varkappa}\) бесконечно многими способами так, чтобы они не образовывали диагональную матрицу и не могли быть преобразованы друг в друга диагональной матрицей.

Следствием Теоремы VIII, вместе с Теоремой V соответственно VII, является

\textbf{Теорема IX.}\footnote{Ср. Loewy, с. 106--107.} \emph{Пусть модуль \(\mathfrak M\) вполне приводим и является наименьшим общим кратным простых модулей \(\mathfrak P_1,
\ldots,
\mathfrak P_k\), образующих тотально взаимно простую систему, и пусть область рациональности \(P\) содержит бесконечно много констант. Тогда необходимое и достаточное условие того, что \(\mathfrak M\) можно бесконечно многими различными способами представить как наименьшее общее кратное простых модулей, состоит в том, что по меньшей мере два из модулей \(\mathfrak P_1,
\ldots,
\mathfrak P_k\) имеют тот же вид.}

Теперь выведем общие критерии того, что модуль \(\mathfrak M\) делится на бесконечно много простых модулей, и для этого сначала докажем следующую

\textbf{Лемма.}\footnote{Ср. Loewy, с. 96.} \emph{Если вполне приводимый модуль \(\mathfrak M=[\mathfrak P_1,
\ldots,
\mathfrak P_k]\) делится на простой модуль \(\mathfrak P'\), отличный от всех \(\mathfrak P_i\), то по меньшей мере два из \(\mathfrak P_i\) имеют тот же вид.}

Сначала выберем из \(\mathfrak P_1,
\ldots,
\mathfrak P_k\) тотально взаимно простую систему, последовательно вычеркивая каждый из этих простых модулей, который входит в наименьшее общее кратное еще не вычеркнутых. Поэтому теперь предположим, что \(\mathfrak P_1,
\ldots,
\mathfrak P_k\) сами образуют тотально взаимно простую систему. Пусть \(a'\) есть класс \(\mathfrak M\), сопоставленный единичному классу mod \(\mathfrak P'\). Тогда
\[
 Fa'=Fa'a_1+\cdots+Fa'a_k
\]
для каждого многочлена \(F\). Среди произведений \(a'a_i\) по меньшей мере два отличны от нуля, например \(a'a_1\) и \(a'a_2\); в самом деле, если бы, например, только \(a'a_1\ne0\), то было бы \(Fa'=Fa'a_1\), и классы остатков \(Fa'\) образовывали бы подгруппу \(\mathfrak A_1\), а значит, так как \(\mathfrak A_1\) есть простая группа, совпадали бы с \(\mathfrak A_1\), так что и модули \(\mathfrak P'\) и \(\mathfrak P_1\) совпадали бы. По способу рассуждения из \S~8 отсюда следует изоморфия \(\mathfrak A'\) с \(\mathfrak A_1\) и \(\mathfrak A_2\).

Пусть теперь \(\mathfrak M\) --- произвольный модуль. Тогда рассмотрим наименьшее общее кратное \(\mathfrak N\) всех простых модулей, на которые делится \(\mathfrak M\).\footnote{Если их бесконечно много, то этому представлению по Теореме IV не может соответствовать никакое аддитивное разложение группы остатков. Ср. также пример \S~12, 4.} Тогда возможны два случая. Либо \(\mathfrak N\) не представим как наименьшее общее кратное конечного числа простых модулей; в этом случае \(\mathfrak N\), а значит и \(\mathfrak M\), имеет бесконечно много простых модулей в качестве делителей. Либо \(\mathfrak N\) представим как наименьшее общее кратное конечного числа простых модулей; тогда \(\mathfrak N\) называется наибольшим вполне приводимым модулем модуля \(\mathfrak M\) (по аналогии с соответствующим построением у Loewy). Если среди этих конечных простых модулей два имеют тот же вид, то по Теореме IX модуль \(\mathfrak N\), а следовательно и \(\mathfrak M\), делится на бесконечно много простых модулей. Если же, наоборот, имеет место последнее, то существует по меньшей мере один простой делитель \(\mathfrak N\), отличный от всех тех конечных, как наименьшее общее кратное которых представим \(\mathfrak N\). По лемме тогда \(\mathfrak N\), а следовательно и \(\mathfrak M\), имеет два простых делителя того же вида. Тем самым доказана следующая теорема:

\textbf{Теорема X.} \emph{Если область рациональности \(P\) содержит бесконечно много констант, то необходимое и достаточное условие того, что модуль делится на бесконечно много различных простых модулей, состоит в том, что либо не существует наибольшего вполне приводимого модуля, либо, если такой существует, он делится по меньшей мере на два различных простых модуля того же вида.}
% END UNIT BODY
% END PRODUCER UNIT 087
\endgroup


\clearpage
\begingroup
\setcounter{footnote}{0}
% BEGIN PRODUCER UNIT 088
% Source: C:/Users/Floris/Documents/interlanguage/03_projects/noether/02_slavic_working_corpus/translations/paper17/russian/v001/Noether_Paper17_Section11_Russian_v001.tex
% Identity: 13324 B / SHA-256 431EF7693C53307C316BB498B41C07C966B031FEC5FD139D13D014201DA0CFB0
\setlength{\parindent}{1.25em}
\setlength{\parskip}{0pt}
\emergencystretch=30em
\allowdisplaybreaks
\raggedbottom
\providecommand{\mM}{\mathfrak M}
\providecommand{\mN}{\mathfrak N}
\providecommand{\mS}{\mathfrak S}
\providecommand{\mT}{\mathfrak T}
% BEGIN UNIT BODY
% Unit: Paper 17 Section 11. Direct Russian translation from the German final-audited source slice.

\setcounter{footnote}{23}
\subsection*{\S~11. Связи с системами дифференциальных уравнений и их интегралами.}

В силу теоремы конечности каждому модулю дифференциальных выражений можно сопоставить систему конечного числа дифференциальных выражений в качестве базиса так, что совокупность одновременных интегралов всех дифференциальных уравнений, возникающих при приравнивании к нулю дифференциального выражения модуля, совпадает с совокупностью интегралов конечной системы дифференциальных уравнений, получающейся при приравнивании к нулю базисных дифференциальных выражений. Тем самым дана связь с теорией систем линейных уравнений в частных производных; мы проследим эту связь еще немного в случае, когда группа остатков модуля является ``конечной''.

\textbf{Определение.} \emph{Группа остатков называется конечной, если существует число \(\mu\), порядок группы, такое, что между любыми \(\mu+1\) классами остатков существует линейное соотношение с коэффициентами из \(P\), тогда как существует по меньшей мере одна система из \(\mu\) классов остатков, между которыми соотношения нет и которую мы будем называть линейным базисом группы остатков. В противном случае группа остатков называется бесконечной.}

Сумма двух конечных групп остатков имеет конечный порядок, а именно сумму их порядков. Сумма двух бесконечных групп или одной конечной и одной бесконечной группы есть бесконечная группа. Примеры конечных групп дают все группы остатков модулей от одной переменной; но и для нескольких переменных порядок может стать конечным, например для модуля с базисом \(\xi_1^2\), \(\xi_2^2\), где все классы остатков можно линейно составить из классов \(1,\xi_1,\xi_2,\xi_1\xi_2\). Примеры бесконечных групп см. в \S~12.

Теперь имеет место следующая теорема.

\textbf{Теорема XI.} \emph{Если модуль имеет конечную группу остатков, то всякая система дифференциальных уравнений, возникающая при приравнивании к нулю базиса модуля, допускает ровно столько линейно независимых интегралов, сколько указывает порядок ее группы остатков. Обратно, для данной системы конечного числа линейных уравнений в частных производных, имеющих лишь конечное число линейно независимых интегралов, это число равно порядку группы остатков модуля, порожденного дифференциальными выражениями.}

Область рациональности \(P\) состоит здесь из аналитических функций от \(n\) переменных \(x_1,
\ldots,x_n\), у которых встречающиеся особенности внутри некоторой \(n\)-мерной области образуют не более чем многообразия меньших размерностей, так что существует сколь угодно много точек с регулярной \(n\)-мерной окрестностью.

Пусть теперь \(\mathfrak M\) есть модуль с конечной группой остатков. Тогда покажем, что существует линейный базис из степенных произведений с тем свойством, что \emph{при выражении произвольного класса остатков через этот базис в знаменателе могут встречаться только степени конечного числа различных величин из \(P\)}.

Для этой цели упорядочим все степенные произведения \(\xi_1^{a_1}\cdots\xi_n^{a_n}\) лексикографически и возьмем в качестве представителей \(Q_1,
\ldots,Q_m\) классов остатков все те, которые линейно независимы от предшествующих по модулю \(\mathfrak M\). По предположению таких только конечное число. Все остальные линейно выражаются через них.

Пусть \(\mathfrak S\) есть система всех степенных произведений, не включенных в линейный базис. Тогда, по первой части доказательства Теоремы III, \(\mathfrak S\) имеет конечную подсистему \(\mathfrak T\) таких степенных произведений, что каждое степенное произведение из \(\mathfrak S\) делится по меньшей мере на одно из \(\mathfrak T\). Пусть \(P_1,
\ldots,P_\rho\) будут степенными произведениями из \(\mathfrak T\), а \(M_1=0,
\ldots,M_\rho=0\;(\mathfrak M)\) -- соотношениями, посредством которых \(P_1,
\ldots,P_\rho\) соответственно выражаются по модулю \(\mathfrak M\) линейной комбинацией более низких степенных произведений; например,
\[
 M_i=b_iP_i+\text{низшие члены}.
\]
Если \(P\) есть произвольное степенное произведение из \(\mathfrak S\), то есть имеет вид
\[
 \xi_1^{h_1}\cdots\xi_n^{h_n}P_i,
\]
то отсюда следует
\[
\begin{aligned}
 \xi_1^{h_1}\cdots\xi_n^{h_n}M_i
 &=b_i\xi_1^{h_1}\cdots\xi_n^{h_n}P_i+\text{низшие члены}\\
 &=b_iP+\text{низшие члены}.
\end{aligned}
\]
Предположим теперь, что для всех более низких, чем \(P\), степенных произведений в знаменателе встречаются только степени \(b_1,
\ldots,b_\rho\), тогда как числители образованы из коэффициентов \(M_i\) сложением, умножением и дифференцированием; тогда это же утверждение верно и для самого \(P\). Такое предположение допустимо, поскольку оно выполнено для \(P_1\), первого элемента из \(\mathfrak S\) и \(\mathfrak T\).

Далее, по Теореме III, \(M_1,
\ldots,M_\rho\) является базисом модуля \(\mathfrak M\). Рассмотрим теперь такую систему значений \(\beta_1,
\ldots,\beta_n\) для \(x_1,
\ldots,x_n\), для которой \(b_1\ne0,
\ldots,b_\rho\ne0\), а остальные коэффициенты \(M_i\) регулярны, то есть регулярную точку системы дифференциальных уравнений \(M_1=0,
\ldots,M_\rho=0\). Степенным произведениям \(Q_i=\xi_1^{h_{i1}}\cdots\xi_n^{h_{in}}\) \((i=1,
\ldots,m)\) тогда соответствуют частные производные
\[
 \frac{\partial^{h_{i1}+\cdots+h_{in}}}{\partial x_1^{h_{i1}}\cdots\partial x_n^{h_{in}}}y,
\]
для которых мы задаем произвольные постоянные значения \((\gamma_1,
\ldots,\gamma_m)\) в точке \((\beta_1,
\ldots,\beta_n)\). Тогда дифференциальные уравнения \(M_i=0\) однозначно определяют значения всех высших производных \(y\) как конечные значения в этой точке, поскольку в знаменателе встречаются только \(b_i\), и, как известно, отсюда доказывается существование \(m\) линейно независимых аналитических интегралов.

С другой стороны, при наших предположениях не может существовать больше чем \(m\) линейно независимых интегралов, ибо иначе, как известно, в регулярной точке можно было бы произвольно задать значения \(m+1\) подходящим образом выбранных производных, что противоречило бы тому, что между любыми \(m+1\) степенными произведениями должны существовать линейные зависимости по модулю \(\mathfrak M\).

Если, обратно, модуль имеет бесконечную группу остатков, то тем же путем показывают, что существует бесконечно много линейно независимых интегралов; следовательно, в случае лишь \(m\) линейно независимых интегралов группа остатков конечна, а ее порядок по сказанному выше равен \(m\).

В заключение этого параграфа покажем, что понятие изоморфии или, что то же самое, того же вида для двух модулей дифференциальных выражений имеет для интегралов соответствующих систем дифференциальных уравнений то же значение, что и использованное Пуанкаре понятие вида для одной переменной (ср. Блумберг, цит. соч., Приложение I).

\textbf{Теорема XII.} \emph{Если два модуля \(\mathfrak M\) и \(\mathfrak N\) дифференциальных выражений имеют один и тот же вид, то существуют два дифференциальных выражения \(P\) и \(Q\), такие что \(P(y)\) и \(Q(z)\) пробегают все интегралы \(\mathfrak N\) и \(\mathfrak M\), когда \(y\) пробегает все интегралы \(\mathfrak M\), а \(z\) -- все интегралы \(\mathfrak N\).}

Действительно, по предположению (по \S~9) для двух подходящим образом выбранных многочленов \(P\) и \(Q\) существуют соотношения
\[
 MQ=0(\mathfrak N),\qquad PQ=1(\mathfrak N),
\]
\[
 NP=0(\mathfrak M),\qquad QP=1(\mathfrak M).
\]
Поэтому ввиду \(NP(y)=0\) функция \(P(y)\) является интегралом \(\mathfrak N\). Точно так же \(Q(z)\) является интегралом \(\mathfrak M\). Ввиду \(PQ(z)=z\) и \(QP(y)=y\) функция \(P(y)\) пробегает все интегралы \(\mathfrak N\), а \(Q(z)\) -- все интегралы \(\mathfrak M\). При этом каждому отличному от нуля \(y\) соответствует также отличное от нуля \(z\), и обратно.
% END UNIT BODY
% END PRODUCER UNIT 088
\endgroup


\clearpage
\begingroup
\setcounter{footnote}{0}
% BEGIN PRODUCER UNIT 089
% Source: C:/Users/Floris/Documents/interlanguage/03_projects/noether/02_slavic_working_corpus/translations/paper17/russian/v001/Noether_Paper17_Section12_Russian_v001.tex
% Identity: 19380 B / SHA-256 99F81EEBBC5885179E0A0F49FEC5DB7CB25166DFD52972678AA3EDBE4482C1BD
\setlength{\parindent}{1.25em}
\setlength{\parskip}{0pt}
\emergencystretch=30em
\allowdisplaybreaks
\raggedbottom
\providecommand{\mM}{\mathfrak M}
\providecommand{\mN}{\mathfrak N}
\providecommand{\mA}{\mathfrak A}
\providecommand{\mB}{\mathfrak B}
\providecommand{\mP}{\mathfrak P}
\providecommand{\mD}{\mathfrak D}
\providecommand{\mS}{\mathfrak S}
\providecommand{\ma}{\mathfrak a}
\providecommand{\mb}{\mathfrak b}
% BEGIN UNIT BODY
% Unit: Paper 17 Section 12 and end matter. Direct Russian translation from the German final-audited source slice.

\setcounter{footnote}{23}
\subsection*{\S~12. Примеры.}

1. В качестве первого примера рассмотрим обыкновенное линейное дифференциальное уравнение 2-го порядка, коэффициенты которого представлены как симметрические функции интегралов. Область рациональности \(P\) состоит здесь из всех рациональных функций от интегралов и их производных, причем мы ограничиваемся достаточно малой окрестностью регулярной точки дифференциального уравнения. Модуль \(\mathfrak M\) мы определяем как совокупность всех дифференциальных выражений, которые имеют \(y_1\) и \(y_2\) интегралами; для них мы снова пишем многочлены от одной неопределенной \(\xi\). Базис модуля тогда одноэлементен, а именно задается многочленом
\[
 M=
 \begin{vmatrix} y_1&y_2\\ y_1'&y_2'\end{vmatrix}\xi^2
 -\begin{vmatrix} y_1&y_2\\ y_1''&y_2''\end{vmatrix}\xi
 +\begin{vmatrix} y_1'&y_2'\\ y_1''&y_2''\end{vmatrix},
\]
который, с точностью до множителя из \(P\), определен однозначно. Теперь
\begin{equation}
 \mathfrak M=[\mathfrak M_1,
\mathfrak M_2],
\tag{41}
\end{equation}
где \(\mathfrak M_i\) состоит из всех дифференциальных выражений, имеющих интеграл \(y_i\), а его базис, с точностью до множителя из \(P\), определяется через
\[
 M_i=y_i\xi-y_i'\qquad(i=1,2).
\]
В самом деле, \(\mathfrak M\) делится на \(\mathfrak M_1\) и \(\mathfrak M_2\), поскольку он обращается в нуль для интегралов \(y_1\) и \(y_2\), следовательно, делится на \([\mathfrak M_1,
\mathfrak M_2]\); а так как порядок базисного многочлена для \([\mathfrak M_1,
\mathfrak M_2]\) должен быть по меньшей мере 2, то \(\mathfrak M=[\mathfrak M_1,
\mathfrak M_2]\). Кроме того, \(\mathfrak M_1\) и \(\mathfrak M_2\) взаимно просты, поскольку
\begin{equation}
 \frac{y_2}{D}(y_1\xi-y_1')-\frac{y_1}{D}(y_2\xi-y_2')=1,
 \qquad
 D=D(y)=\begin{vmatrix}y_1&y_2\\ y_1'&y_2'\end{vmatrix}.
\tag{42}
\end{equation}
К тому же они являются простыми модулями, так как их группа остатков имеет порядок 1.

Но \(\mathfrak M\), помимо \(y_1\) и \(y_2\), имеет также все интегралы
\[
 z_1=\alpha_{11}y_1+\alpha_{12}y_2,
 \qquad
 z_2=\alpha_{21}y_1+\alpha_{22}y_2,
 \qquad
 A=|\alpha_{ik}|\ne0,
\]
и потому тождественен наименьшему общему кратному двух взаимно простых модулей \(\mathfrak N_1\) и \(\mathfrak N_2\), которые соответственно имеют интегралы \(z_1\) и \(z_2\); их базисные многочлены заданы формулами
\[
 N_i=z_i\xi-z_i'=(\alpha_{i1}y_1+\alpha_{i2}y_2)\xi-(\alpha_{i1}y_1'+\alpha_{i2}y_2').
\]
Итак, \(\mathfrak M\) допускает бесконечно много различных представлений как наименьшее общее кратное взаимно простых простых модулей, и по теореме VII оба модуля \(\mathfrak N_1\) и \(\mathfrak N_2\) между собой и со всеми \(\mathfrak M_1\) и \(\mathfrak M_2\) имеют один и тот же вид. В самом деле, группы остатков всех этих модулей имеют порядок 1, а потому изоморфны.

Теперь выпишем аддитивное разложение группы остатков \(\mathfrak G=\mathfrak A_1+\mathfrak A_2\), соответствующее представлению (41). Для этого, как в \S~4, исходим из соотношения (42), выражающего взаимную простоту \(\mathfrak M_1\) и \(\mathfrak M_2\). Для единиц \(a_1\) и \(a_2\) получаем классы остатков, порожденные многочленами
\[
 -\frac{y_1}{D}(y_2\xi-y_2')
 \quad\text{и}\quad
 \frac{y_2}{D}(y_1\xi-y_1')
\]
соответственно. Отсюда
\[
 \mathfrak M_1=\left(\mathfrak M,\frac{y_2}{D}(y_1\xi-y_1')\right),
 \qquad
 \mathfrak M_2=\left(\mathfrak M,-\frac{y_1}{D}(y_2\xi-y_2')\right).
\]
Если теперь так же образовать единицы \(b_1\) и \(b_2\), соответствующие разложению \(\mathfrak M=[\mathfrak N_1,
\mathfrak N_2]\), заменяя \(y_i\) на \(z_i\), то получается
\[
 -\frac{z_1}{D(z)}(z_2\xi-z_2')
 =-\frac{\alpha_{11}y_1+\alpha_{12}y_2}{AD(y)}
 ((\alpha_{21}y_1+\alpha_{22}y_2)\xi-(\alpha_{21}y_1'+\alpha_{22}y_2')),
\]
\[
 \frac{z_2}{D(z)}(z_1\xi-z_1')
 =\frac{\alpha_{21}y_1+\alpha_{22}y_2}{AD(y)}
 ((\alpha_{11}y_1+\alpha_{12}y_2)\xi-(\alpha_{11}y_1'+\alpha_{12}y_2')),
\]
следовательно
\begin{equation}
\left\{
\begin{aligned}
 b_1&=\frac{\alpha_{11}y_1+\alpha_{12}y_2}{y_1}
 \left(\frac{\alpha_{22}}{A}\right)a_1
 +\frac{\alpha_{11}y_1+\alpha_{12}y_2}{y_2}
 \left(-\frac{\alpha_{21}}{A}\right)a_2
 =\frac{z_1}{y_1}\alpha^{11}a_1+\frac{z_1}{y_2}\alpha^{12}a_2,\\
 b_2&=\frac{\alpha_{21}y_1+\alpha_{22}y_2}{y_1}
 \left(-\frac{\alpha_{12}}{A}\right)a_1
 +\frac{\alpha_{21}y_1+\alpha_{22}y_2}{y_2}
 \left(\frac{\alpha_{11}}{A}\right)a_2
 =\frac{z_2}{y_1}\alpha^{21}a_1+\frac{z_2}{y_2}\alpha^{22}a_2.
\end{aligned}\right.
\tag{43}
\end{equation}
Здесь \((\alpha^{ik})\) есть матрица, обратная к матрице \((\alpha_{ik})\).

Обратно, переставляя \(z_i\) с \(y_i\), получаем
\begin{equation}
\left\{
\begin{aligned}
 a_1&=\frac{y_1}{z_1}\alpha_{11}b_1+\frac{y_1}{z_2}\alpha_{12}b_2,\\
 a_2&=\frac{y_2}{z_1}\alpha_{21}b_1+\frac{y_2}{z_2}\alpha_{22}b_2.
\end{aligned}\right.
\tag{44}
\end{equation}
Из этих формул видна изоморфия групп \(\mathfrak A\) и \(\mathfrak B\). Действительно, согласно нашим общим рассуждениям при \(a_1b_\sigma\ne0\) соответствующее сопоставление имеет вид \(a_1\sim a_1b_\sigma\), то есть по (44)
\[
 a_1\sim\frac{y_1}{z_\sigma}\alpha_{1\sigma}b_\sigma,
 \qquad\text{если }\alpha_{1\sigma}\ne0,
\]
и по (43) также
\[
 b_\sigma\sim\frac{z_\sigma}{y_2}\alpha^{\sigma2}a_2,
 \qquad\text{если }\alpha^{\sigma2}\ne0,
\]
следовательно
\[
 \frac{y_1}{z_\sigma}\alpha_{1\sigma}b_\sigma
 \sim
 \frac{y_1}{z_\sigma}\alpha_{1\sigma}\frac{z_\sigma}{y_2}\alpha^{\sigma2}a_2
 =\frac{y_1}{y_2}\alpha_{1\sigma}\alpha^{\sigma2}a_2.
\]
Поэтому
\[
 a_1\sim\alpha_{1\sigma}\alpha^{\sigma2}\frac{y_1}{y_2}a_2
 \quad\text{для }\alpha_{1\sigma}\alpha^{\sigma2}\ne0,
 \qquad
 \frac{1}{\alpha_{1\sigma}\alpha^{\sigma2}}\frac{y_2}{y_1}a_1\sim a_2,
\]
что представляет обратное сопоставление. Если \((\alpha_{ik})\) не является диагональной матрицей, то это условие \(\alpha_{1\sigma}\alpha^{\sigma2}\ne0\) всегда выполнено для некоторого \(\sigma\). Случай диагональной матрицы следует исключить, потому что тогда модули тождественны. Если, например, только \(\alpha_{12}=0\), то \(\mathfrak A_1=\mathfrak B_1\), но не \(\mathfrak A_2=\mathfrak B_2\). Значит, из \(\mathfrak A_1+\mathfrak A_2=\mathfrak B_1+\mathfrak B_2\), \(\mathfrak A_1=\mathfrak B_1\) не следует \(\mathfrak A_2=\mathfrak B_2\) (ср. примечание 12). Три модуля \(\mathfrak M_1\), \(\mathfrak M_2\) и \(\mathfrak N_2\) тогда попарно взаимно просты, но не образуют тотально взаимно простую систему, так как \([\mathfrak M_1,
\mathfrak M_2]\) делится на \(\mathfrak N_2\) (ср. примечание 15).

Далее, если положить
\[
 t_{12}=\alpha_{1\sigma}\alpha^{\sigma2}\frac{y_1}{y_2}a_2,
 \qquad
 t_{21}=\frac{1}{\alpha_{1\sigma}\alpha^{\sigma2}}\frac{y_2}{y_1}a_1,
\]
то соотношения (35), как подтверждается вычислением, выполнены.

2. Пример бесконечной группы дает каждый модуль от более чем одной переменной, базис которого состоит только из одного многочлена. Несколько сложнее следующий пример; он должен показать, что бесконечные группы могут возникать и у многоэлементных модулей, то есть у таких, чей базис модуля охватывает более одного многочлена.\footnote{Ср. пример Ландау, сообщенный Блумбергом (цит. соч., с. 51--52). Здесь и далее независимые переменные обозначаем через \(x\) и \(y\).}
\[
 \mathfrak M=[(\xi+x\eta),(\xi+1)]=[(P),(Q)]
 \qquad\left(\xi=\frac{\partial}{\partial x},\ \eta=\frac{\partial}{\partial y}\right).
\]
Правило умножения пусть будет правилом для дифференциальных выражений. Группа остатков здесь становится бесконечной, потому что группы остатков для \((\xi+x\eta)\) и \((\xi+1)\) обе бесконечны, тогда как модули \((\xi+x\eta)\) и \((\xi+1)\) взаимно просты ввиду
\[
 1=(-x\xi-x^2\eta+1)Q+(x\xi+x-1)P.
\]
Теперь очень легко показать, что \(\mathfrak M\) не содержит многочлена второй степени, положив
\[
 (u_1\xi+u_2\eta+u_3)(\xi+x\eta)=(v_1\xi+v_2\eta+v_3)(\xi+1).
\]
Отсюда получается \(u_1=u_2=u_3=v_1=v_2=v_3=0\). Если далее искать многочлены третьей степени в \(\mathfrak M\), то находятся два линейно независимых, например
\[
 (\xi^2+x\xi\eta+\xi+(2+x)\eta)Q=(\xi^2+2\xi+1)P
\]
и
\[
 (\xi^2+2x\xi\eta+x^2\eta^2+\eta)Q
 =(\xi^2+x\xi\eta+\xi+(x-1)\eta)P.
\]
Если бы теперь \(\mathfrak M\) был одноэлементным модулем, то оба должны были бы делиться на базисный многочлен; а так как \(\mathfrak M\) не содержит многочлена второй степени, они должны были бы делиться на многочлен третьей степени, то есть быть пропорциональными, чего нет. Многоэлементность \(\mathfrak M\) есть внутренняя причина отказа теорем Ландау о произведениях, действительных в случае одной переменной.

3. Еще один пример покажет, что бесконечные группы могут возникать и у простых модулей.\footnote{Наше определение ``простого модуля'' отличается от того, что в коммутативном случае называют простым модулем, поскольку там всегда существуют делители меньшей многообразности, кроме случая, когда многообразность \(=0\), то есть группа остатков конечна.} Пусть \(P\) будет областью всех рациональных функций от двух переменных \(x,y\). Модуль \(\mathfrak M\) с базисом
\[
 \xi-\frac{y}{x}=\xi-a
\]
имеет бесконечную группу, поскольку она содержит класс остатков каждой степени \(y\); с другой стороны, \(\mathfrak M\) является простым модулем. В самом деле, покажем, что каждый делитель
\[
 \mathfrak M_1=(\xi-a,
F(\xi,\eta))
\]
становится единичным модулем. Именно, каждый многочлен \(F(\xi,\eta)\) можно mod\((\xi-a)\) представить многочленом \(G(\eta)\):
\[
 F(\xi,\eta)=c_0\eta^\nu+c_1\eta^{\nu-1}+\cdots+c_\nu\quad(\mathfrak M).
\]
Без ограничения общности пусть \(c_0=1\). Но ведь по закону умножения \(\xi\eta^\nu-\eta^\nu\xi=0\); следовательно, ввиду
\[
 \xi\equiv a(\mathfrak M_1),
 \qquad
 \eta^\nu\equiv F_1(\mathfrak M_1),
\]
где положено \(F_1=-(c_1\eta^{\nu-1}+\cdots+c_\nu)\), получаем
\[
 \xi F_1-\eta^\nu a\equiv0(\mathfrak M_1).
\]
Разлагая
\[
 \eta^\nu a=a\eta^\nu+\nu\frac{\partial a}{\partial y}\eta^{\nu-1}+\cdots,
\]
и снова заменяя здесь \(\eta^\nu\) на \(F_1\), получаем
\[
 \xi(c_1\eta^{\nu-1}+\cdots+c_\nu)-a(c_1\eta^{\nu-1}+\cdots+c_\nu)
 +\nu\frac{\partial a}{\partial y}\eta^{\nu-1}+\cdots\equiv0(\mathfrak M_1),
\]
или
\[
 c_1\eta^{\nu-1}a+\frac{\partial c_1}{\partial x}\eta^{\nu-1}+\cdots
 -ac_1\eta^{\nu-1}+\cdots+
u\frac{\partial a}{\partial y}\eta^{\nu-1}+\cdots
 \equiv0(\mathfrak M_1),
\]
или, наконец,
\[
 \left(\frac{\partial c_1}{\partial x}+\nu\frac{\partial a}{\partial y}\right)\eta^{\nu-1}+\cdots
 \equiv0(\mathfrak M_1).
\]
Слева теперь стоит многочлен от \(\eta\) точной степени \(\nu-1\), который не обращается тождественно в нуль; ведь старший коэффициент
\[
 \frac{\partial c_1}{\partial x}+\nu\frac{\partial a}{\partial y}
 =\frac{\partial c_1}{\partial x}+\nu\frac{1}{x}
 \quad\text{есть }\ne0,
\]
так как иначе \(c_1=-\nu\log x+\varphi(y)\) не был бы рациональным. Следовательно, процедуру можно продолжать, и она приводит к многочлену нулевой степени, не обращающемуся тождественно в нуль; этим показано, что и единица содержится в \(\mathfrak M_1\).\footnote{Эта процедура сводится к доказательству того, что необходимые для системы дифференциальных уравнений условия интегрируемости не выполнены.}

4. Если, с другой стороны, присоединить к области рациональности еще функцию \(\log x\), то \(\mathfrak M\) имеет делители
\[
 \left(\xi-\frac{y}{x},\ \eta-\log x+\varphi(y)\right),
\]
где \(\varphi(y)\) пробегает все рациональные функции от \(y\). Все эти делители различны между собой; для каждого выполнено условие интегрируемости, а соответствующий интеграл равен
\[
 y\log x-\int\varphi(y)\,dy+c.
\]
Поэтому единица не может содержаться в модуле, так как иначе нулевая функция была бы единственно возможным интегралом. С другой стороны, \(\xi\) и \(\eta\) по отношению к модулю конгруэнтны некоторой величине из области; поэтому группа остатков конечна и имеет порядок 1. Кроме того, из бесконечно многих делителей любые два взаимно просты, а данный модуль \(\mathfrak M\) равен наименьшему общему кратному \(\mathfrak N\) всех этих делителей. В самом деле, он делится на все делители, следовательно, и на \(\mathfrak N\). Если бы теперь \(\mathfrak N\) был собственным делителем \(\mathfrak M\), то \(\mathfrak N\) должен был бы содержать еще по крайней мере один многочлен \(F\), скажем степени \(\rho\), зависящий только от \(\eta\). Если теперь заставить \(\varphi(y)\) пробегать функции \(y,
\ldots,y^{\rho+1}\), то между соответствующими интегралами
\[
 y\log x-\frac{y^{i+1}}{i+1}+c_i\qquad(i=1,
\ldots,
\rho+1)
\]
нет линейного соотношения с коэффициентами, не зависящими от \(y\). Но дифференциальное уравнение \(\rho\)-го порядка \(F=0\) не может иметь \(\rho+1\) линейно независимых интегралов.

Модуль \(\mathfrak M=\mathfrak N\), однако, не является вполне приводимым.\footnote{Тем самым мы имеем пример для первого случая теоремы X.} Иначе он был бы наименьшим общим кратным конечного числа этих простых модулей, а его группа остатков имела бы конечный порядок; между тем группа остатков для \(\xi-\frac{y}{x}\) содержит классы остатков всех степеней \(\eta\), следовательно, бесконечна.

Гёттинген, 1 августа 1919 г.

\begin{center}
(Поступила 4 августа 1919 г.)
\end{center}
% END UNIT BODY
% END PRODUCER UNIT 089
\endgroup


\clearpage
\begingroup
\setcounter{footnote}{0}
% BEGIN PRODUCER UNIT 090
% Source: C:/Users/Floris/Documents/interlanguage/03_projects/noether/02_slavic_working_corpus/translations/paper18/russian/v001/Noether_Paper18_Russian_v001.tex
% Identity: 4237 B / SHA-256 607A72E1304F6536F036E8F7EF75C286F9F1E3314147D98C7966DBD02C97F329
\setlength{\parindent}{1.25em}
\setlength{\parskip}{0pt}
\emergencystretch=30em
\allowdisplaybreaks
\raggedbottom
% BEGIN UNIT BODY
% Unit: Paper 18 complete. Direct Russian translation from the German final-audited source slice.

\setcounter{footnote}{0}
\section*{18. О работе погибшего на войне К. Генцельта по теории элиминации}
\begin{center}
J. Ber. d. DMV 30 (1921), с. 101
\end{center}

\begin{center}
\textbf{11-е заседание, пятница, 23 сентября 1921 года, 11 часов утра.}\\
(Председательствует Леви.)
\end{center}

\textbf{1. Э. Нётер: О работе погибшего на войне К. Генцельта по теории элиминации.}

Курт Генцельт с октября 1914 года числится пропавшим без вести под Диксмёйде, и его следует отнести к погибшим. Существеннейшие части его эрлангенской диссертации, которую он оставил построенной без пробелов, но изложенной почти непонятно, я намерена вскоре опубликовать в свободной переработке в \emph{Mathematische Annalen}.

Речь идет об общей задаче \textbf{теории элиминации}: определить общие нули всех полиномов идеала \(\mathfrak M\), состоящего из полиномов. Если предварительно подвергнуть переменные линейному преобразованию с неопределенными коэффициентами, то главный результат можно сформулировать так:

Идеалу \(\mathfrak M\) можно единственным образом сопоставить результирующую форму
\[
 R_{\mathfrak M}=R^{(1)}(x_1,
\ldots,x_n)\cdots R^{(i)}(x_1,
\ldots,x_n)\cdots R^{(m)}(x_n)=0(\mathfrak M),
\]
такую, что \(R_{\mathfrak M}\) обращается в нуль на всех нулях \(\mathfrak M\) и только на них. Если \(\mathfrak N\) является делителем \(\mathfrak M\), а \(R_{\mathfrak N}\) и \(R_{\mathfrak M}\) совпадают, то совпадают и сами \(\mathfrak N\) и \(\mathfrak M\).

Вторая часть теоремы показывает, что результирующая форма дает нули с характеристической кратностью. Она опирается на то, что \(\mathfrak M\) можно рассматривать как модуль линейных форм от степенных произведений \(x_1,
\ldots,x_{i-1}\); тогда отдельные множители \(R^{(i)}(x_i,
\ldots,x_n)\) формы \(R_{\mathfrak M}\) становятся нормами соответствующего ``основного модуля'' по этому модулю, когда \(x_{i+1},
\ldots,x_n\) присоединены к области коэффициентов. Если привлечь абстрактную теорию идеалов, получается следующая интерпретация: если
\[
 \mathfrak M=[\mathfrak Q,\mathfrak Q_1,
\ldots,
\mathfrak Q_r]=[\mathfrak Q,\mathfrak L]
\]
есть разложение \(\mathfrak M\) на примарные идеалы \(\mathfrak Q_i\), а \(\mathfrak L\) является дополнением к \(\mathfrak Q\), то идеалу \(\mathfrak Q\) соответствует примарный множитель \(Q^{(i)}(x_i,
\ldots,x_n)\) формы \(R^{(i)}(x_i,
\ldots,x_n)\), который в указанном выше смысле представляет норму \(\mathfrak L\) по \(\mathfrak Q\).

\clearpage
% END UNIT BODY
% END PRODUCER UNIT 090
\endgroup


\clearpage
\begingroup
\setcounter{footnote}{0}
% BEGIN PRODUCER UNIT 091
% Source: C:/Users/Floris/Documents/interlanguage/03_projects/noether/02_slavic_working_corpus/translations/paper19/russian/v001/Noether_Paper19_Introduction_Russian_v001.tex
% Identity: 21515 B / SHA-256 40BD3D5C19E144580B5D276D747CE6E6DDDC83E91E54D4443F065A59D9F6EA7F
\setlength{\parindent}{1.25em}
\setlength{\parskip}{0pt}
\emergencystretch=30em
\allowdisplaybreaks
\raggedbottom
\newcommand{\ideal}[1]{\mathfrak{#1}}
% BEGIN UNIT BODY
% Unit: Paper 19 title, contents, and introduction. Direct Russian translation from the German final-audited source slice.

\setcounter{footnote}{0}
\section*{19. Теория идеалов в кольцевых областях}
\begin{center}
Math. Ann. 83 (1921), с. 24--66
\end{center}

\begin{center}\textbf{Содержание.}\end{center}
\noindent Введение.\\
\S~1. Кольцевая область, идеал, условие конечности.\\
\S~2. Представление идеала как наименьшего общего кратного конечного числа неразложимых идеалов.\\
\S~3. Равенство числа компонентов в двух различных разложениях на неразложимые идеалы.\\
\S~4. Примарные идеалы. Единственность ассоциированных простых идеалов в двух различных разложениях на неразложимые идеалы.\\
\S~5. Представление идеала как наименьшего общего кратного наибольших примарных идеалов. Единственность ассоциированных простых идеалов.\\
\S~6. Единственное представление идеала как наименьшего общего кратного относительно-просто неразложимых идеалов.\\
\S~7. Единственность изолированных идеалов.\\
\S~8. Единственное представление идеала как произведения взаимно-просто неразложимых идеалов.\\
\S~9. Распространение исследования на модули. Равенство числа компонентов в разложениях на неразложимые модули.\\
\S~10. Специальный случай области полиномов.\\
\S~11. Примеры из теории чисел и из теории дифференциальных выражений.\\
\S~12. Пример из теории элементарных делителей.

\subsection*{Введение.}
Содержание настоящей работы составляет \emph{перенесение теорем о разложении целых рациональных чисел, соответственно идеалов в алгебраических числовых полях, на идеалы в произвольных областях целостности, а в более общем виде -- в кольцевых областях}. Для понимания этого перенесения сначала сформулируем теоремы о разложении для целых рациональных чисел в форме, несколько отличающейся от обычной.

Если в
\[
 a=p_1^{\varrho_1}p_2^{\varrho_2}\cdots p_\sigma^{\varrho_\sigma}=q_1q_2\cdots q_\sigma
\]
рассматривать степени простых чисел \(q_i\) как компоненты разложения, то этим компонентам присущи следующие характерные свойства:

1. Они \emph{попарно взаимно просты}; однако ни одно \(q_i\) не представимо как произведение попарно взаимно простых чисел, то есть в этом смысле имеет место неразложимость. Из попарной взаимной простоты следует также, что произведение \(q_1\cdots q_\sigma\) равно наименьшему общему кратному \([q_1\cdots q_\sigma]\).

2. Любые две компоненты, \(q_i\) и \(q_k\), \emph{относительно просты}; то есть если \(bq_i\) делится на \(q_k\), то \(b\) делится на \(q_k\). И в этом смысле также имеет место неразложимость.

3. Каждое \(q\) является \emph{примарным}; то есть если произведение \(b\cdot c\) делится на \(q\), но \(b\) не делится, то некоторая степень\footnote{Если эта степень всегда первая, то, как известно, речь идет о простых числах.} \(c\) делится. Кроме того, это представление есть представление через \emph{наибольшие примарные компоненты}, поскольку произведение двух различных \(q\) уже не является примарным. Также относительно разложения на наибольшие примарные компоненты эти \(q\) неразложимы.

4. Каждое \(q\) \emph{неразложимо} в том смысле, что его нельзя представить как наименьшее общее кратное двух собственных делителей.

Связь этих примарных чисел \(q\) с простыми числами \(p\) состоит в том, что для каждого \(q\) существует одно и, с точностью до знака, только одно \(p\), которое является делителем \(q\) и некоторая степень которого делится на \(q\): ассоциированное простое число. Если \(p^\varrho\) есть наименьшая такая степень -- \(\varrho\) есть экспонента \(q\), -- то здесь, в частности, \(p^\varrho\) равно \(q\). Теорему единственности теперь можно сформулировать так:

\emph{В двух различных разложениях целого рационального числа на неразложимые, наибольшие примарные компоненты \(q\) совпадают число компонентов, ассоциированные простые числа (с точностью до знака) и экспоненты. Поскольку \(p^\varrho=q\), отсюда следует также совпадение самих \(q\) (с точностью до знака).}

Неопределенность, задаваемая знаком, как известно, устраняется, если вместо чисел рассматривать выведенные из них идеалы (все числа, делящиеся на \(a\)); тогда формулировка точно так же справедлива для единственного разложения идеалов конечных алгебраических числовых полей на степени простых идеалов.

Далее (\S~1) за основу берется общая кольцевая область, которая должна удовлетворять только условию конечности: каждый идеал области имеет конечный идеальный базис. Без такого условия конечности неразложимые и простые идеалы вообще могли бы не существовать, как показывает область всех целых алгебраических чисел, в которой нет разложения на простые идеалы.

Оказывается, что -- соответственно четырем характерным свойствам компонентов \(q\) -- в общем случае существуют четыре отдельные разложения, возникающие одно из другого последовательным расщеплением. При этом разложение на взаимно-просто неразложимые идеалы является представлением в виде произведения, а остальные три разложения являются приведенными (\S~2) представлениями как наименьшие общие кратные. Сохраняется также связь между примарным идеалом -- неразложимые идеалы тоже примарны -- и ассоциированным простым идеалом: для каждого примарного идеала \(\ideal Q\) однозначно определяется ассоциированный простой идеал \(\ideal P\), который является делителем \(\ideal Q\) и некоторая степень которого делится на \(\ideal Q\). Если \(\ideal P^\varrho\) есть наименьшая такая степень -- \(\varrho\) есть экспонента \(\ideal Q\), -- то здесь \(\ideal P^\varrho\) уже не обязано совпадать с \(\ideal Q\). Теорема единственности выражается так:

Разложения 1 и 2 единственны; в двух различных разложениях 3 или 4 совпадают число компонентов и ассоциированные простые идеалы;\footnote{Предположительно, сверх этого совпадают также экспоненты, а еще более общо -- соответствующие компоненты изоморфны.} изолированные идеалы (\S~7), встречающиеся среди компонентов, определены однозначно.

Для доказательства теорем о разложении из условия конечности выводится впервые высказанная Дедекиндом для конечных числовых модулей ``теорема о конечной цепи'', а из нее -- представление 4 каждого идеала как наименьшего общего кратного конечного числа неразложимых идеалов. Преобразованием понятия разложимости компоненты отсюда получается основная теорема единственности для разложения 4 на неразложимые идеалы. Объединяя каждый раз конечное число компонентов, переходят к остальным разложениям, теоремы единственности которых получаются как следствие теоремы единственности 4.

Наконец показывается (\S~9), что представление конечным числом неразложимых компонентов имеет место и при более слабых предположениях; коммутативность кольцевой области не требуется, и достаточно вместо идеала рассматривать модуль относительно этой области. В этом более общем случае еще сохраняется равенство числа компонентов в двух различных разложениях, тогда как понятия простого и примарного связаны с коммутативностью и с понятием идеала; напротив, понятие взаимной простоты для идеалов в некоммутативных областях сохраняется.

Простейшая кольцевая область, в которой действительно возникают четыре отдельные разложения, -- это область всех полиномов от \(n\) переменных с произвольными комплексными коэффициентами. Отдельные разложения здесь можно иррационально истолковать через поведение алгебраических образований, а теорема единственности для ассоциированных простых идеалов соответствует основной теореме теории элиминации об единственной разложимости алгебраических образований на неразложимые. Дальнейшие примеры дают все конечные области целостности из полиномов (\S~10). Но даже простая область всех четных чисел, более общо -- всех чисел, делящихся на фиксированное число, уже дает пример частично разделенных разложений (\S~11). Пример теории идеалов в некоммутативных областях дает теория элементарных делителей (\S~12), где имеется единственное разложение на неразложимые идеалы, соответственно классы. Эти неразложимые классы полностью характеризуют неразложимые составляющие элементарных делителей и, возможно, для областей, где обычная теория элементарных делителей отказывает, могут рассматриваться как ее эквивалент.

О существующей литературе следует заметить следующее: разложение на наибольшие примарные идеалы для области полиномов с произвольными комплексными, соответственно целыми, коэффициентами было дано Ласкером и в отдельных пунктах далее развито Маколеем.\footnote{E. Lasker, Zur Theorie der Moduln und Ideale. Math. Ann. 60 (1905), с. 20, Satz VII und XIII. -- F. S. Macaulay, On the Resolution of a given Modular System into Primary Systems including some Properties of Hilbert Numbers. Math. Ann. 74 (1913), с. 66.} Оба опираются на теорию элиминации, то есть используют тот факт, что полином можно единственным образом представить как произведение неразложимых полиномов. В действительности теоремы о разложении идеалов не зависят от этой предпосылки, как подсказывает теория идеалов в алгебраических числовых полях и как показывает настоящая работа. Примарный идеал у Ласкера и Маколея также определяется на основе понятий из теории элиминации.

Разложение на неразложимые идеалы и разложение на относительно-просто неразложимые, по-видимому, не были замечены в литературе даже для области полиномов; только у Маколея имеется замечание о единственности изолированных примарных идеалов.

Разложение на взаимно-просто неразложимые идеалы для области полиномов было дано Шмейдлером,\footnote{W. Schmeidler, Über Moduln und Gruppen hyperkomplexer Größen. Math. Zeitschr. 3 (1919), с. 29.} с использованием теории элиминации для доказательства конечности. Однако там теорема единственности сформулирована только для классов идеалов, а не для самих идеалов. Последняя теорема единственности содержится в совместной работе,\footnote{E. Noether--W. Schmeidler, Moduln in nichtkommutativen Bereichen, insbesondere aus Differential- und Differenzenausdrücken. Math. Zeitschr. 8 (1920), с. 1.} где речь идет об идеалах в некоммутативных областях полиномов. Там используется только конечный идеальный базис; следовательно, теоремы и методы остаются действительными для общих кольцевых областей и в настоящей работе уточняются относительно равенства числа (\S~11). Настоящие исследования являются сильным обобщением и дальнейшим развитием тех понятийных образований, которые лежали в основе обеих работ. Существенным в обеих работах является переход от представления как наименьшего общего кратного к аддитивному разложению системы классов вычетов. Здесь, ради простоты изложения, мы снова остаемся при наименьшем общем кратном; аддитивному разложению тогда соответствует преобразование понятия разложимости в свойство дополнения (\S~3). Однако рассуждениями, существенно соответствующими рассуждениям совместной работы, все приведенные теоремы можно также понимать как аддитивные теоремы разложения для системы классов вычетов и некоторых подсистем. Эта система классов вычетов образует кольцо той же общности, что и исходно взятое за основу; ведь каждое кольцо можно рассматривать как систему классов вычетов того идеала, который соответствует совокупности тождественных соотношений между элементами кольца; или также подсистемы этих соотношений, если остальные соотношения предполагаются выполненными уже в области.

Это замечание дает также место работам Френкеля.\footnote{A. Fraenkel, Über die Teiler der Null und die Zerlegung von Ringen. J. f. M. 145 (1914), с. 139. Über gewisse Teilbereiche und Erweiterungen von Ringen. Habilitationsschrift, Leipzig, Teubner, 1916. Über einfache Erweiterungen zerlegbarer Ringe. J. f. M. 151 (1920), с. 121.} Френкель рассматривает аддитивные разложения колец, на которые налагаются такие ограничительные условия (существование регулярных элементов, деление на них, условие разложимости), что для соответствующего идеала четыре разложения совпадают. Из-за этого совпадения и его условие конечности, согласно которому идеал должен иметь только конечное число собственных делителей, -- от которого, впрочем, частично отступают, -- не означает более сильного ограничения, чем наше. Исходный пункт Френкеля обусловлен иными, по существу алгебраическими, целями его работ; посредством алгебраического расширения он затем приходит к более общим кольцам с менее ограничительными условиями.

\clearpage
% END UNIT BODY
% END PRODUCER UNIT 091
\endgroup


\clearpage
\begingroup
\setcounter{footnote}{0}
% BEGIN PRODUCER UNIT 092
% Source: C:/Users/Floris/Documents/interlanguage/03_projects/noether/02_slavic_working_corpus/translations/paper19/russian/v001/Noether_Paper19_Section01_Russian_v001.tex
% Identity: 11360 B / SHA-256 C54ABEAC994F11DD492C3AE95663D0BFC0A2019B1C8D91FE089C5BF60DEB4FC2
\setlength{\parindent}{1.25em}
\setlength{\parskip}{0pt}
\emergencystretch=30em
\allowdisplaybreaks
\raggedbottom
\newcommand{\ideal}[1]{\mathfrak{#1}}
% BEGIN UNIT BODY
% Unit: Paper 19, Section 1. Direct Russian translation from the German final-audited source slice.

\setcounter{footnote}{5}
\subsection*{\S~1. Кольцевая область, идеал, условие конечности.}

\textbf{1.} Положенная в основу область \(\Sigma\) пусть будет (коммутативным) кольцом в абстрактном определении;\footnote{Определение взято из габилитационной работы Френкеля, с опущением его ограничительных условий 6, I и II; зато пришлось включить коммутативный закон сложения. Итак, речь идет о законах, определяющих поле, с опущением обратимости умножения.} то есть \(\Sigma\) состоит из системы элементов \(a,b,c,\ldots,f,g,h,\ldots\), в которой как равенство определено отношение, удовлетворяющее обычным условиям; и в которой посредством двух операций (видов соединения), сложения и умножения, из каждых двух элементов кольца \(a\) и \(b\) всегда однозначно получается третий элемент как сумма \(a+b\) и как произведение \(a\cdot b\). Кольцо и эти вообще произвольные операции должны удовлетворять следующим законам:
\begin{enumerate}
\item Ассоциативный закон сложения: \((a+b)+c=a+(b+c)\).
\item Коммутативный закон сложения: \(a+b=b+a\).
\item Ассоциативный закон умножения: \((a\cdot b)c=a(b\cdot c)\).
\item Коммутативный закон умножения: \(a\cdot b=b\cdot a\).
\item Дистрибутивный закон: \(a(b+c)=ab+ac\).
\item Закон неограниченного и однозначного вычитания. В \(\Sigma\) существует единственный элемент \(x\), удовлетворяющий уравнению \(a+x=b\). (Его обозначают \(x=b-a\).)
\end{enumerate}
Из этих свойств следует существование нуля; однако кольцо не обязано иметь единицу; и произведение двух элементов может обращаться в нуль, хотя ни один множитель не обращается в нуль. Кольца, для которых из обращения произведения в нуль всегда следует обращение в нуль одного из множителей и которые, кроме того, имеют единицу, называются собственными областями целостности. Для конечной суммы \(a+a+\cdots+a\) введем обычное сокращенное обозначение \(na\), где целые числа \(n\) рассматриваются лишь как сокращающие знаки, а не как элементы кольца, и где рекуррентно определены \(a=1a\), \(na+a=(n+1)a\).

\textbf{2.} Под идеалом \(\ideal M\)\footnote{Идеалы обозначаются большими немецкими буквами. \(\ideal M\) должно напоминать пример идеала из полиномов, который обычно называется ``модулем'' или модулем форм. Впрочем, \S\S~1--3 используют только модульное свойство, а не свойство идеала; ср. \S~9.} в \(\Sigma\) понимается система элементов из \(\Sigma\), удовлетворяющая двум условиям:
\begin{enumerate}
\item Вместе с \(f\) идеал \(\ideal M\) содержит также \(a\cdot f\), где \(a\) -- произвольный элемент из \(\Sigma\).
\item Вместе с \(f\) и \(g\) идеал \(\ideal M\) содержит также разность \(f-g\); следовательно, вместе с \(f\) он содержит также \(nf\) для всякого целого числа \(n\).
\end{enumerate}
Если \(f\) является элементом \(\ideal M\), то, как обычно, мы выражаем это записью \(f\equiv0(\ideal M)\) и говорим, что \(f\) делится на \(\ideal M\). Если каждый элемент \(\ideal N\) одновременно является элементом \(\ideal M\), то есть делится на \(\ideal M\), мы говорим: \(\ideal N\) делится на \(\ideal M\); в обозначениях: \(\ideal N\equiv0(\ideal M)\). \(\ideal M\) называется собственным делителем \(\ideal N\), если он содержит элементы, отличные от элементов \(\ideal N\), то есть если обратно он не делится на \(\ideal N\). Из \(\ideal N\equiv0(\ideal M)\), \(\ideal M\equiv0(\ideal N)\) следует \(\ideal N=\ideal M\).

И остальные известные понятия также сохраняются буквально. Под наибольшим общим делителем двух идеалов \(\ideal A\) и \(\ideal B\), то есть \(\ideal D=(\ideal A,\ideal B)\), понимается совокупность элементов, представимых в форме \(a+b\), где \(a\) пробегает все элементы из \(\ideal A\), а \(b\) все элементы из \(\ideal B\); \(\ideal D\) снова является идеалом. Точно так же наибольший общий делитель бесконечно многих идеалов,
\[
 \ideal D=(\ideal A_1,\ideal A_2,\ldots,\ideal A_\nu,\ldots),
\]
определяется как совокупность элементов \(d\), представимых в виде суммы элементов каждый раз лишь из конечного числа идеалов:
\[
 d=a_{i_1}+a_{i_2}+\cdots+a_{i_n};
\]
и здесь \(\ideal D\) снова является идеалом.

Если идеал \(\ideal M\), в частности, содержит конечное число элементов \(f_1,f_2,\ldots,f_\varrho\) так, что
\[
 \ideal M=(f_1,\ldots,f_\varrho),\qquad
 f=a_1f_1+\cdots+a_\varrho f_\varrho+n_1f_1+\cdots+n_\varrho f_\varrho
\]
для каждого \(f\equiv0(\ideal M)\), где \(a_i\) являются величинами кольцевой области, а \(n_i\) целыми числами, то \(\ideal M\) называется конечным идеалом; \(f_1,\ldots,f_\varrho\) называются идеальным базисом.

Далее мы полагаем в основу только такие кольца \(\Sigma\), которые удовлетворяют условию конечности: \emph{каждый идеал в \(\Sigma\) конечен, то есть обладает идеальным базисом.}

\textbf{3.} Из условия конечности прямо следует лежащая в основе всех дальнейших рассуждений

\textbf{Теорема I (теорема о конечной цепи).}\footnote{Впервые сформулирована для числовых модулей Дедекиндом: Zahlentheorie, Suppl. XI, \S~172, Satz VIII (4-е издание); наше доказательство и обозначение ``цепь'' заимствованы оттуда. Для идеалов из полиномов у Ласкера, цит. соч., с. 56 (вспомогательная теорема). Однако в обоих случаях теорема имеет лишь отдельные применения. Наши применения всюду опираются на постулат выбора.}
\emph{Пусть \(\ideal M,\ideal M_1,\ideal M_2,\ldots,\ideal M_\nu,\ldots\) -- счетно бесконечная система идеалов в \(\Sigma\), каждый из которых делится на следующий. Тогда начиная с некоторого конечного индекса \(n\) все идеалы тождественны:
\[
 \ideal M_n=\ideal M_{n+1}=\cdots.
\]
Иными словами: если \(\ideal M,\ideal M_1,\ideal M_2,\ldots,\ideal M_s,\ldots\) образуют просто упорядоченную цепь идеалов так, что каждый идеал является собственным делителем непосредственно предыдущего, то цепь обрывается за конечное число шагов.}

Действительно, пусть
\[
 \ideal D=(\ideal M,\ideal M_1,\ideal M_2,\ldots,\ideal M_\nu,\ldots)
\]
есть наибольший общий делитель системы, а \(f_1,\ldots,f_\varrho\) -- базис \(\ideal D\), всегда существующий вследствие условия конечности. Тогда из предположения о делимости следует, что каждый элемент \(\ideal D\) одновременно является элементом одного из идеалов цепи; ведь из
\[
 f=g+h,\qquad g\equiv0(\ideal M_r),\qquad h\equiv0(\ideal M_s),\qquad (r<s)
\]
следует \(g\equiv0(\ideal M_s)\), а значит \(f\equiv0(\ideal M_s)\). То же самое верно, если \(f\) является суммой большего числа слагаемых. Следовательно, существует и конечный индекс \(n\) такой, что
\[
 f_1\equiv0(\ideal M_n),\quad f_2\equiv0(\ideal M_n),\quad\ldots,\quad f_\varrho\equiv0(\ideal M_n).
\]
Поскольку, обратно, \(\ideal M_n\equiv0(\ideal D)\), получаем \(\ideal M_n=\ideal D\); а далее, поскольку \(\ideal M_n\equiv0(\ideal D)\) и \(\ideal D=\ideal M_n\equiv0(\ideal M_r)\), получаем также \(\ideal M_r=\ideal D=\ideal M_n\) для каждого \(r>n\), чем теорема доказана.

Отметим, что из этой теоремы, обратно, снова следует существование идеального базиса, так что условие конечности можно было бы сформулировать и в этой безбазисной форме.

\clearpage
% END UNIT BODY
% END PRODUCER UNIT 092
\endgroup


\clearpage
\begingroup
\setcounter{footnote}{0}
% BEGIN PRODUCER UNIT 093
% Source: C:/Users/Floris/Documents/interlanguage/03_projects/noether/02_slavic_working_corpus/translations/paper19/russian/v001/Noether_Paper19_Section02_Russian_v001.tex
% Identity: 13459 B / SHA-256 3B28D609F80BBA6DFC1F6EFCED63AE047A9FCCB44530B05DA35FFAB226C5E8E2
\setlength{\parindent}{1.25em}
\setlength{\parskip}{0pt}
\emergencystretch=30em
\allowdisplaybreaks
\raggedbottom
\newcommand{\ideal}[1]{\mathfrak{#1}}
% BEGIN UNIT BODY
% Unit: Paper 19, Section 2. Direct Russian translation from the German final-audited source slice.

\setcounter{footnote}{8}
\subsection*{\S~2. Представление идеала как наименьшего общего кратного конечного числа неразложимых идеалов.}

Наименьшее общее кратное \([\ideal B_1,\ideal B_2,\ldots,\ideal B_k]\) идеалов \(\ideal B_1,\ideal B_2,\ldots,\ideal B_k\) определим, как обычно, как совокупность элементов, которые делятся как на \(\ideal B_1\), так и на \(\ideal B_2\), \(\ldots\), так и на \(\ideal B_k\); в обозначениях:
\[
 \text{из } f\equiv0(\ideal B_i),\quad (i=1,2,\ldots,k),
 \quad\text{следует}\quad
 f\equiv0([\ideal B_1,\ideal B_2,\ldots,\ideal B_k]),
\]
и обратно. Наименьшее общее кратное снова является идеалом; идеалы \(\ideal B_i\) мы также называем компонентами разложения.

\textbf{Определение I.} Представление \(\ideal M=[\ideal B_1,\ldots,\ideal B_k]\) называется редуцированным представлением, если ни один \(\ideal B_i\) не делит наименьшее общее кратное \(\ideal U_i\) остальных идеалов и если ни один \(\ideal B_i\) нельзя заменить собственным делителем.\footnote{Пример нередуцированного представления: \((x,xy)=[(x),(x^2,xy,y^\lambda)]\) для каждого показателя \(\lambda\ge2\); соответствующее \(\lambda=1\) представление \([(x),(x^2,y)]\) является соответствующим редуцированным. Это представление сообщил мне погибший на войне K.~Hentzelt как простейший пример неоднозначного разложения на примарные идеалы.} Если эти условия выполнены только для идеала \(\ideal B_i\), то представление называется редуцированным относительно \(\ideal B_i\). Наименьшее общее кратное \(\ideal U_i=[\ideal B_1,\ldots,\ideal B_{i-1},\ideal B_{i+1},\ldots,\ideal B_k]\) называется дополнением \(\ideal B_i\). Представления, в которых выполнено только первое условие, называются кратчайшими представлениями.

Итак, при представлении идеала как наименьшего общего кратного достаточно ограничиваться редуцированными представлениями, по следующей лемме.

\textbf{Лемма I.} \emph{Каждое представление идеала как наименьшего общего кратного конечного числа идеалов можно по меньшей мере одним способом заменить редуцированным представлением; такое представление, в частности, можно получить последовательным разложением.}

Действительно, пусть \(\ideal M=[\ideal B_1,\ldots,\ideal B_k]\) есть произвольное представление \(\ideal M\). Мы последовательно опускаем те \(\ideal B_i\), которые делят наименьшее общее кратное оставленных идеалов. Поскольку остальные идеалы всё ещё дают \(\ideal M\), в возникающем таким образом представлении
\[
 \ideal M=[\ideal B_1,\ldots,\ideal B_k]=[\ideal U_i,\ideal B_i]
\]
первое условие уже выполнено; значит, это кратчайшее представление, и это условие остаётся выполненным, если заменить какой-либо \(\ideal B_i\) собственным делителем. Второе условие, в свою очередь, всегда достижимо по теореме о конечной цепи (теорема I). В самом деле, пусть положено:
\[
 \ideal B_i\supset \ideal B_i'\supset \ideal B_i''\supset\cdots\supset \ideal B_i^{(\nu)}\supset\cdots,
 \qquad
 \ideal M=[\ideal U_i,\ideal B_i]=[\ideal U_i,\ideal B_i']=\cdots=[\ideal U_i,\ideal B_i^{(\nu)}]=\cdots,
\]
где каждый \(\ideal B_i^{(\nu)}\) является собственным делителем непосредственно предшествующего идеала. Тогда по этой теореме цепь \(\ideal B_i,\ideal B_i',\ldots,\ideal B_i^{(\nu)},\ldots\) должна оборваться за конечное число шагов; значит, в представлении \(\ideal M=[\ideal U_i,\ideal B_i^{(\nu)}]\) идеал \(\ideal B_i^{(\nu)}\) уже нельзя заменить его собственным делителем. Тем более это остаётся верным, если \(\ideal U_i\) заменить собственным делителем. Применяя, следовательно, эту процедуру последовательно к каждому \(\ideal B_i\) и каждый раз образуя дополнение с уже редуцированными \(\ideal B\), мы получаем редуцированное представление.\footnote{Что такое представление не определяется данным представлением однозначно, показывает предыдущий пример. Для \((x^2,xy)=[(x),(x^2,xy,y^\lambda)]\), где \(\lambda\ge2\), наряду с \([(x),(x^2,y)]\) также \([(x),(x^2,ux+y)]\) при произвольном \(u\) является редуцированным представлением.}

Чтобы получить такое представление последовательно, надо показать, что из отдельных редуцированных представлений
\[
 \ideal M=[\ideal B_1,\ideal C_1],\qquad
 \ideal C_1=[\ideal B_2,\ideal C_2],\quad\ldots,\quad
 \ideal C_{k-1}=[\ideal B_k,\ideal C_k]
\]
следует, что и возникающее из них представление \(\ideal M=[\ideal B_1,\ldots,\ideal B_k,\ideal C_k]\) редуцировано. Для этого достаточно показать, что вместе с редуцированными представлениями
\[
 \ideal M=[\ideal B,\ideal C],\qquad \ideal C=[\ideal C_1,\ideal C_2]
\]
редуцировано также \(\ideal M=[\ideal B,\ideal C_1,\ideal C_2]\). В самом деле, здесь по предположению ни один \(\ideal B\) не делит своего дополнения. Если бы это имело место для некоторого \(\ideal C_i\), то, вопреки предположению в первом представлении, \(\ideal C\) можно было бы заменить собственным делителем, поскольку в силу второго редуцированного представления \(\ideal C_1\) и \(\ideal C_2\) являются собственными делителями \(\ideal C\). Значит, представление становится кратчайшим. Далее, по предположению ни один \(\ideal B\) нельзя заменить собственным делителем; если бы это было возможно для некоторого \(\ideal C_i\), то это, вопреки предположению, соответствовало бы замене \(\ideal C\) собственным делителем, поскольку представление для \(\ideal C\) редуцировано. Тем самым лемма доказана.

\textbf{Определение II.} Идеал \(\ideal M\) называется разложимым, если он представим как наименьшее общее кратное двух собственных делителей; в противоположном случае \(\ideal M\) называется неразложимым.

Теперь, пользуясь редуцированным представлением, мы доказываем с помощью теоремы I о конечной цепи:

\textbf{Теорема II.} \emph{Каждый идеал представим как наименьшее общее кратное конечного числа неразложимых идеалов.}\footnote{Что такое представление в общем случае не единственно, показывает предыдущий пример: \((x^2,xy)=[(x),(x^2,ux+y)]\) при произвольном \(u\). Обе компоненты при произвольном \(u\) неразложимы. Все делители \((x)\) имеют вид \((x,g(y))\), где \(g(y)\) означает многочлен от \(y\); следовательно, наименьшее общее кратное двух таких идеалов тоже имеет этот вид и потому является собственным делителем \((x)\). Идеал \((x^2,ux+y)\) имеет только один делитель \((x,y)\), значит, с необходимостью также является неразложимым.}

Действительно, произвольный идеал \(\ideal M\) либо неразложим; тогда \(\ideal M=[\ideal M]\) есть представление, которого требует теорема II. Либо же \(\ideal M=[\ideal B_1,\ideal C_1]\), где \(\ideal B_1,\ideal C_1\) являются собственными делителями \(\ideal M\), и это представление по лемме I можно считать редуцированным. Для \(\ideal C_1\) действует та же альтернатива: либо он неразложим, либо существует редуцированное представление \(\ideal C_1=[\ideal B_2,\ideal C_2]\).

Продолжая так, получаем ряд редуцированных представлений
\begin{equation}
 \ideal M=[\ideal B_1,\ideal C_1]=[\ideal B_1,\ideal B_2,\ideal C_2]=\cdots=[\ideal B_1,\ldots,\ideal B_n,\ideal C_n]=\cdots .
\tag{1}
\end{equation}
В цепи \(\ideal C_1,\ideal C_2,\ldots,\ideal C_n,\ldots\) каждый \(\ideal C_n\) является собственным делителем непосредственно предшествующего идеала; следовательно, цепь обрывается за конечное число шагов. Поэтому существует индекс \(n\), для которого \(\ideal C_n\) становится неразложимым. По лемме I представление \(\ideal M=[\ideal U_n,\ideal C_n]\) также редуцировано; значит, \(\ideal C_n\) не может делить своё дополнение \(\ideal U_n\), а в представлении \(\ideal M=[\ideal U_n,\ideal C_n]\) идеал \(\ideal U_n\) нельзя заменить никаким собственным делителем. Если при необходимости заменить \(\ideal U_n\) собственным делителем,\footnote{На самом деле это представление также редуцировано относительно \(\ideal U_n\), как будет показано в \S~3 (лемма IV) в качестве обращения леммы I.} так, чтобы представление стало редуцированным, то тем самым показано, что каждый разложимый идеал допускает редуцированное представление как наименьшее общее кратное одного неразложимого и одного комплементарного ему идеала. Следовательно, в ряду (1) без ограничения общности все \(\ideal B_i\) можно считать неразложимыми; повторение приведённого выше рассуждения даёт существование неразложимого \(\ideal C_n\), чем теорема II доказана.

\clearpage
% END UNIT BODY
% END PRODUCER UNIT 093
\endgroup


\clearpage
\begingroup
\setcounter{footnote}{0}
% BEGIN PRODUCER UNIT 094
% Source: C:/Users/Floris/Documents/interlanguage/03_projects/noether/02_slavic_working_corpus/translations/paper19/russian/v001/Noether_Paper19_Section03_Russian_v001.tex
% Identity: 14879 B / SHA-256 0D70E61BE7B33C2985B0F45E847506745074C199F06674AA29A8A089D8935E2E
\setlength{\parindent}{1.25em}
\setlength{\parskip}{0pt}
\emergencystretch=30em
\allowdisplaybreaks
\raggedbottom
\newcommand{\ideal}[1]{\mathfrak{#1}}
% BEGIN UNIT BODY
% Unit: Paper 19, Section 3. Direct Russian translation from the German final-audited source slice.

\setcounter{footnote}{12}
\subsection*{\S~3. Равенство числа компонентов в двух различных разложениях на неразложимые идеалы.}

Для доказательства равенства числа компонентов сначала надо выразить разложимость, соответственно неразложимость, идеала через свойства его дополнения по следующей теореме.

\textbf{Теорема III.}\footnote{Теорема III соответствует переходу от модулей к фактор-группам в работах Шмейдлера и Нетер--Шмейдлера (ср. введение). Идеалу \(\ideal C\) соответствует фактор-группа, а \(\ideal N_1\) и \(\ideal N_2\) -- подгруппы, на которые разлагается фактор-группа.} \emph{Пусть кратчайшее представление \(\ideal M=[\ideal U,\ideal C]\) редуцировано относительно \(\ideal C\). Тогда необходимым и достаточным условием того, что \(\ideal C\) разложим, является существование двух идеалов \(\ideal N_1\) и \(\ideal N_2\), которые являются собственными делителями \(\ideal M\), таких, что}
\begin{equation}
 \ideal N_1\ideal N_2\equiv0(\ideal M),\qquad
 \ideal N_i\equiv0(\ideal U),\qquad
 [\ideal N_1,\ideal N_2]=\ideal M.
\tag{2}
\end{equation}
\emph{Отсюда ещё следует: если условия (2) выполнены, а \(\ideal C\) неразложим, то по меньшей мере один \(\ideal N_i\) не является собственным делителем \(\ideal M\); тогда \(\ideal N_i=\ideal M\).}

Пусть \(\ideal C=[\ideal C_1,\ideal C_2]\), где \(\ideal C_1,\ideal C_2\) являются собственными делителями \(\ideal C\). Тогда получается
\[
 \ideal M=[\ideal U,\ideal C]=[\ideal U,\ideal C_1,\ideal C_2]
       =[[\ideal U,\ideal C_1],[\ideal U,\ideal C_2]].
\]
Здесь идеалы \([\ideal U,\ideal C_i]\) являются собственными делителями \(\ideal M\), иначе \([\ideal U,\ideal C]\) не было бы редуцировано относительно \(\ideal C\). Поскольку делимость на \(\ideal U\) также выполнена, необходимость условия (2) доказана. (Представление (2) не редуцировано, так как один из \([\ideal U,\ideal C_i]\) можно заменить на \(\ideal C_i\).)

Пусть теперь, обратно, выполнено (2). Образуем идеалы
\[
 \ideal C_1=(\ideal C,\ideal N_1),\qquad
 \ideal C_2=(\ideal C,\ideal N_2),\qquad
 \ideal C^*=[\ideal C_1,\ideal C_2].
\]
Тогда \(\ideal C\) делится как на \(\ideal C_1\), так и на \(\ideal C_2\), а значит, и на наименьшее общее кратное \(\ideal C^*\). Чтобы показать делимость \(\ideal C^*\) на \(\ideal C\), пусть
\[
 f\equiv0(\ideal C^*);\quad f\equiv0(\ideal C_1);\quad f\equiv0(\ideal C_2),
 \quad\text{или также}\quad f=c+n_1=t+n_2,
\]
где \(c,t,n_1,n_2\) суть величины из \(\ideal C,\ideal C,\ideal N_1,\ideal N_2\); в частности, \(n_1\) делится на \(\ideal N_1\). Следовательно, разность
\[
 g=c-t=n_2-n_1
\]
делится как на \(\ideal C\), так и на \(\ideal U\), а значит, на \(\ideal M\). В силу \(n_1=n_2+m\) далее \(n_1\) (так же как \(n_2\)) делится как на \(\ideal N_1\), так и на \(\ideal N_2\), а значит, на \(\ideal M\). Поэтому
\[
 f=c+m,\qquad f\equiv0(\ideal C),\qquad \ideal C^*=\ideal C.
\]
При этом \(\ideal C_1\) и \(\ideal C_2\) являются собственными делителями \(\ideal C\). Ведь если бы \(\ideal C_1=(\ideal C,\ideal N_1)=\ideal C\), то \(\ideal N_1\) делился бы на \(\ideal C\); следовательно, из-за делимости на \(\ideal U\) он был бы равен \(\ideal M\), вопреки предположению. Тем самым \(\ideal C=[\ideal C_1,\ideal C_2]\) распознан как разложимый; теорема III доказана.

Заметим, что почти то же рассуждение показывает ещё следующее:

\textbf{Лемма II.} \emph{Если в кратчайшем представлении \(\ideal M=[\ideal U,\ideal C]\) идеал \(\ideal C\) можно заменить собственным делителем, то \(\ideal C\) разложим.}

Итак, пусть
\[
 \ideal M=[\ideal U,\ideal C]=[\ideal U,\ideal C_1],
\]
и пусть положено
\[
 \ideal C^*=[\ideal C_1,(\ideal U,\ideal C)].
\]
Тогда \(\ideal C\) снова делится на \(\ideal C^*\). Из \(f\equiv0(\ideal C^*)\) далее следует
\[
 f=c_1=a+c.
\]
Разность \(a=c_1-c\) делится как на \(\ideal U\), так и на \(\ideal C_1\), следовательно, на \(\ideal M\). Поэтому \(c_1=c+m\), \(f\equiv0(\ideal C)\), \(\ideal C^*=\ideal C\). Так как и \(\ideal C_1\), и \((\ideal U,\ideal C)\) по предположению являются собственными делителями \(\ideal C\), тем самым \(\ideal C=\ideal C^*\) распознан как разложимый.

Неразложимый \(\ideal C\), следовательно, нельзя заменить собственным делителем.

Пусть теперь заданы два различных кратчайших представления \(\ideal M\) как наименьшего общего кратного конечного числа неразложимых идеалов:
\[
 \ideal M=[\ideal B_1\cdots\ideal B_k]=[\ideal D_1\cdots\ideal D_l].
\]
По замечанию к лемме II эти представления одновременно редуцированы. Сначала докажем

\textbf{Лемму III.} \emph{Для каждого дополнения \(\ideal U_i=[\ideal B_1\cdots\ideal B_{i-1}\ideal B_{i+1}\cdots\ideal B_k]\) существует идеал \(\ideal D_j\) такой, что \(\ideal M=[\ideal U_i,\ideal D_j]\).}

Действительно, полагая \(\ideal M=[\ideal B_i,\ideal C_1]\), \(\ideal C_1=[\ideal D_1,\ideal C_2]\) и т. д., получаем
\[
 \ideal M=[\ideal U_i,\ideal B_i]=[\ideal U_i,[\ideal D_1,\ideal C_2]]=[\ideal B_i,[\ideal U_i,\ideal D_1],\ideal C_2].
\]
Здесь для \(\ideal N_1=[\ideal U_i,\ideal D_1]\), \(\ideal N_2=[\ideal U_i,\ideal C_2]\) выполнены условия (2) теоремы III, поскольку \(\ideal M=[\ideal U_i,\ideal B_i]\) редуцировано относительно \(\ideal B_i\), а представление является кратчайшим. Так как \(\ideal B_i\) был предположен неразложимым, необходимо один из этих \(\ideal N_i\) должен стать равным \(\ideal M\).

Если \([\ideal U_i,\ideal D_1]=\ideal M\), лемма доказана. Если же \([\ideal U_i,\ideal C_2]=\ideal M\), то соответственно получается \(\ideal M=[[\ideal U_i,\ideal D_1],[\ideal U_i,\ideal C_2]]\), где по тому же выводу снова одна компонента должна стать равной \(\ideal M\). Продолжая так, получаем либо \(\ideal M=[\ideal U_i,\ideal D_j]\), где \(j<l\), либо \(\ideal M=[\ideal U_i,\ideal D_1,\ldots,\ideal D_{l-1}]\); поскольку \(\ideal D_1\cdots\ideal D_{l-1}=\ideal D_l\), лемма доказана.

Отсюда теперь следует

\textbf{Теорема IV.} \emph{В двух различных кратчайших представлениях идеала как наименьшего общего кратного неразложимых идеалов число компонентов одинаково.}

Из леммы для \(i=1\) действительно следует:
\[
 \ideal M=[\ideal B_1\cdots\ideal B_k]=[\ideal U_1,\ideal D_{j_1}]
 =[\ideal D_{j_1},\ideal B_2,\ldots,\ideal B_k].
\]
Рассматривая теперь два разложения
\[
 \ideal M=[\ideal D_{j_1},\ideal B_2,\ldots,\ideal B_k]
          =[\ideal D_1,\ideal D_2,\ldots,\ideal D_l],
\]
и повторяя приведённый выше вывод относительно дополнения \(\ideal U_2=[\ideal D_{j_1},\ideal B_3,\ldots,\ideal B_k]\) идеала \(\ideal B_2\), получаем
\[
 \ideal M=[\ideal U_2,\ideal B_2]=[\ideal U_2,\ideal D_{j_2}]
          =[\ideal D_{j_1},\ideal D_{j_2},\ideal B_3,\ldots,\ideal B_k],
\]
а продолжение процедуры даёт
\[
 \ideal M=[\ideal D_{j_1},\ideal D_{j_2},\ldots,\ideal D_{j_k}].
\]
Так как по предположению представление \(\ideal M=[\ideal D_1\cdots\ideal D_l]\) является кратчайшим, то есть никакой \(\ideal D\) нельзя опустить, различные среди \(\ideal D_{j_i}\) должны исчерпывать все \(\ideal D\); следовательно, \(k\ge l\). Если в лемме и последующих выводах всюду поменять \(\ideal B\) с \(\ideal D\), соответственно получаем \(l\ge k\), а значит, \(k=l\), чем равенство числа компонентов доказано. Отсюда ещё следует, что идеалы \(\ideal D_{j_i}\) все попарно различны, иначе в кратчайшем представлении через \(\ideal D_{j_i}\) появилось бы меньше чем \(k\) компонентов; значит, обозначения можно выбрать так, чтобы \(\ideal D_{j_i}=\ideal D_i\). По той же причине все промежуточные представления \(\ideal M=[\ideal D_{j_1}\cdots\ideal D_{j_r},\ideal B_{r+1},\ldots,\ideal B_k]\) также кратчайшие и, по замечанию к лемме II, тем самым редуцированные.

Равенство числа компонентов приводит к обращению леммы I:

\textbf{Лемма IV.} \emph{Если в редуцированном представлении объединить компоненты в группы и образовать их наименьшие общие кратные, то возникающее представление будет редуцировано. Иными словами: из редуцированного представления}
\[
 \ideal M=[\ideal C_1\cdots\ideal C_k]
\]
\emph{следует, что также \(\ideal M=[\ideal N_1,\ideal N_2]\) редуцировано, если положено \(\ideal N_1=[\ideal C_1,\,\ldots,\ideal C_h]\), \(\ideal N_2=[\ideal C_{h+1},\ldots,\ideal C_k]\).}

Сначала надо заметить, что \(\ideal N_1\) не может делить своё дополнение \(\ideal N_2\), поскольку это не имеет места ни для одного из его делителей \(\ideal C_i\); следовательно, представление кратчайшее. Чтобы показать, что \(\ideal N_1\) нельзя заменить никаким собственным делителем, разложим \(\ideal C\) на их неразложимые идеалы \(\ideal B\),\footnote{Под этим всегда следует понимать кратчайшее, значит редуцированное, представление через \(\ideal B\).} так что возникают редуцированные (по лемме I) представления
\[
 \ideal N_1=[\ideal B_1\cdots\ideal B_\lambda],
 \qquad
 \ideal N_2=[\ideal B_{\lambda+1}\cdots\ideal B_\kappa].
\]
Пусть теперь \(\ideal M=[\ideal N_1,\ideal N_2]\) редуцировано относительно \(\ideal N_2\), а \(\ideal N_1^*\) является собственным делителем \(\ideal N_1\). По лемме II получается
\[
 \ideal N_1=[\ideal N_1^*,(\ideal N_1,\ideal N_2)],
\]
и это представление необходимо редуцировано относительно \(\ideal N_1^*\), поскольку иначе \(\ideal N_1\) можно было бы также заменить в \(\ideal M\) собственным делителем. Заменим теперь при необходимости и \((\ideal N_1,\ideal N_2)\) собственным делителем так, чтобы возникло редуцированное представление для \(\ideal N_1\). Если теперь разложить обе компоненты \(\ideal N_1\) на неразложимые идеалы, то число \(\lambda\) неразложимых идеалов в \(\ideal N_1\) аддитивно складывается из чисел компонентов; следовательно, число неразложимых идеалов, соответствующих собственному делителю \(\ideal N_1^*\), необходимо меньше \(\lambda\). Но тогда и разложение \(\ideal M=[\ideal N_1^*,\ideal N_2]\) на неразложимые идеалы приводит к менее чем \(\kappa\) идеалам, в противоречии с равенством числа компонентов. Как специальный случай \(\varrho=2\) ещё следует, что представление \(\ideal M=[\ideal N_1,\ideal N_2]\) также редуцировано относительно дополнения \(\ideal N_2\).

\clearpage
% END UNIT BODY
% END PRODUCER UNIT 094
\endgroup


\clearpage
\begingroup
\setcounter{footnote}{0}
% BEGIN PRODUCER UNIT 095
% Source: C:/Users/Floris/Documents/interlanguage/03_projects/noether/02_slavic_working_corpus/translations/paper19/russian/v001/Noether_Paper19_Section04_Russian_v001.tex
% Identity: 18490 B / SHA-256 A15DFCDA06F441A33C65D743EA5B8E8A4CD099A4535403BCFC8A0DC1CBBC1493
\setlength{\parindent}{1.25em}
\setlength{\parskip}{0pt}
\emergencystretch=30em
\allowdisplaybreaks
\raggedbottom
\newcommand{\ideal}[1]{\mathfrak{#1}}
% BEGIN UNIT BODY
% Unit: Paper 19, Section 4. Direct Russian translation from the German final-audited source slice.

\setcounter{footnote}{14}
\subsection*{\S~4. Примарные идеалы. Единственность ассоциированных простых идеалов при двух различных разложениях на неразложимые идеалы.}

В дальнейшем речь идет о связи между примарными и неразложимыми идеалами.

\textbf{Определение III.} Идеал \(\ideal Q\) называется примарным, если из \(a\cdot b\equiv0(\ideal Q)\), \(a\not\equiv0(\ideal Q)\) необходимо следует \(b^\varkappa\equiv0(\ideal Q)\), где показатель \(\varkappa\) является конечным числом.

Это определение можно сформулировать и так: если произведение \(a\cdot b\) делится на \(\ideal Q\), то либо один множитель делится на \(\ideal Q\), либо некоторая степень каждого множителя делится на \(\ideal Q\). Если, в частности, \(\varkappa\) всегда равно \(1\), то идеал называется простым идеалом.

Из определения примарного, соответственно простого, идеала благодаря существованию базисов следует формулировка, говорящая только о произведениях идеалов.\footnote{Если \(\ideal A=(a_1,\ldots,a_r)\), \(\ideal B=(b_1,\ldots,b_s)\), то под \(\ideal A\ideal B\) следует понимать идеал, имеющий базисом произведения \(a_i b_k\); что эти и только эти элементы образуют идеал-произведение, следует из базисного представления ввиду коммутативности.}

\textbf{Определение IIIa.} Идеал \(\ideal Q\) называется примарным, если из \(\ideal A\ideal B\equiv0(\ideal Q)\), \(\ideal A\not\equiv0(\ideal Q)\) необходимо следует \(\ideal B^\varkappa\equiv0(\ideal Q)\). Если \(\varkappa\) всегда равно \(1\), то идеал называется простым идеалом. Следовательно, для простого идеала \(\ideal P\) из \(\ideal A\ideal B\equiv0(\ideal P)\), \(\ideal A\not\equiv0(\ideal P)\) всегда следует \(\ideal B\equiv0(\ideal P)\).

Поскольку в IIIa при \(\ideal A=(a)\), \(\ideal B=(b)\) определение III содержится как частный случай, каждый идеал, примарный по IIIa, является примарным и по III. Обратно, пусть \(\ideal Q\) примарен по III, и пусть выполнена гипотеза IIIa: \(\ideal A\ideal B\equiv0(\ideal Q)\). Тогда либо следует \(\ideal A\equiv0(\ideal Q)\), либо существует по меньшей мере одна величина \(a\not\equiv0(\ideal Q)\) такая, что \(a\ideal B\equiv0(\ideal Q)\), \(a\not\equiv0(\ideal Q)\). Если теперь \(b_1,\ldots,b_s\) есть идеальный базис \(\ideal B\), то по определению III, поскольку \(a b_i\equiv0(\ideal Q)\), получается
\[
 b_i^{\varkappa_i}\equiv0(\ideal Q)\qquad (i=1,\ldots,s).
\]
Ввиду
\[
 b=n_1b_1+\cdots+n_s b_s+a_1b_1+\cdots+a_s b_s
\]
для \(\varkappa=\varkappa_1+\cdots+\varkappa_s\) произведение любых \(\varkappa\) величин \(b\) делится на \(\ideal Q\). Тем самым для идеалов, примарных по III, доказано выполнение определения IIIa. Специально для простых идеалов \(\ideal P\) из \(a\ideal B\equiv0(\ideal P)\), а значит и \(a b\equiv0(\ideal P)\) для всякого \(b\equiv0(\ideal B)\), и \(a\not\equiv0(\ideal P)\), следует \(b\equiv0(\ideal P)\), а потому \(\ideal B\equiv0(\ideal P)\). Тем самым показана эквивалентность двух определений.

Связь между примарными и простыми идеалами устанавливается замечанием, что совокупность \(\ideal P\) всех элементов \(p\), обладающих тем свойством, что некоторая степень \(p\) делится на \(\ideal Q\), образует простой идеал. Сначала ясно, что \(\ideal P\) есть идеал, поскольку вместе с \(p_1\) и \(p_2\) указанным свойством обладают также \(a p_1\) и \(p_1-p_2\). По тому же базисному рассуждению, которое применялось в определении IIIa, далее получается существование числа \(\lambda\) такого, что \(\ideal P^\lambda\equiv0(\ideal Q)\).

Пусть теперь
\[
 a b\equiv0(\ideal P),\qquad a\not\equiv0(\ideal P).
\]
Тогда по определению \(\ideal P\) получается
\[
 a^\lambda b^\lambda\equiv0(\ideal Q),\qquad a^\lambda\not\equiv0(\ideal Q),
\]
а значит, по определению \(\ideal Q\),
\[
 b^{\lambda\varkappa}\equiv0(\ideal Q),\qquad\text{и следовательно}\qquad b\equiv0(\ideal P),
\]
чем \(\ideal P\) доказан как простой идеал. \(\ideal P\) определяется также как наибольший общий делитель всех идеалов \(\ideal B\), обладающих тем свойством, что некоторая степень \(\ideal B\) делится на \(\ideal Q\). Ведь каждый такой \(\ideal B\) делится на \(\ideal P\); значит, и наибольший общий делитель \(\ideal D\) этих \(\ideal B\) делится на \(\ideal P\). Обратно, сам \(\ideal P\) является таким идеалом \(\ideal B\), значит, он делится на \(\ideal D\); этим доказано \(\ideal P=\ideal D\). Итак, \(\ideal P\) есть простой идеал, являющийся делителем \(\ideal Q\), и некоторая степень которого делится на \(\ideal Q\); этим свойством он определяется единственно. Действительно, из
\[
 \ideal Q\equiv0(\ideal P),\qquad \ideal P^\varrho\equiv0(\ideal Q),
 \qquad
 \ideal Q\equiv0(\ideal B),\qquad \ideal B^\sigma\equiv0(\ideal Q)
\]
следует
\[
 \ideal P^\varrho\equiv0(\ideal B),\qquad
 \ideal B^\sigma\equiv0(\ideal P),
\]
а потому по свойству простых идеалов
\[
 \ideal P\equiv0(\ideal B),\qquad \ideal B\equiv0(\ideal P),\qquad \ideal B=\ideal P.
\]
Итак, получаем:

\textbf{Теорема V.} \emph{Для каждого примарного идеала \(\ideal Q\) существует один и только один простой идеал \(\ideal P\), который является делителем \(\ideal Q\) и некоторая степень которого делится на \(\ideal Q\); \(\ideal P\) будет называться ``ассоциированным простым идеалом''.}\footnote{Что обратное неверно, показывает пример \(\ideal M=(x^2,xy)\). Простой идеал \((x)\) выполняет все условия, но \(\ideal M\) не является примарным.} \emph{\(\ideal P\) определяется как наибольший общий делитель всех идеалов \(\ideal B\), обладающих тем свойством, что некоторая степень \(\ideal B\) делится на \(\ideal Q\). Если \(\varrho\) есть наименьшее число такое, что \(\ideal P^\varrho\equiv0(\ideal Q)\), то \(\varrho\) называется экспонентой \(\ideal Q\).}\footnote{Что вообще не обязательно, как в области целых рациональных или алгебраических чисел, будет \(\ideal P^\varrho=\ideal Q\), показывает пример \(\ideal Q=(x^2,y)\); \(\ideal P=(x,y)\); \(\ideal P^2=(x^2,xy,y^2)\equiv0(\ideal Q)\); но \(\ideal Q\not\equiv0(\ideal P^2)\); следовательно, \(\ideal Q\) отлично от \(\ideal P^2\).}

Теперь докажем связь между примарностью и неразложимостью:

\textbf{Теорема VI.} \emph{Каждый непримарный идеал разложим; иначе говоря, каждый неразложимый идеал примарен.}\footnote{Что здесь обратное неверно, показывает, например, пример \(\ideal Q=(x^2,xy,y^\lambda)=[(x^2,y),(x,y^\lambda)]\), где \(\lambda\ge2\). Здесь \(\ideal Q\) примарен, но разложим. То, что \(\ideal Q\) примарен, следует из того, что он содержит все степенные произведения \(\lambda\)-го измерения от \(x,y\); значит, для всякого многочлена без постоянного члена некоторая степень делится на \(\ideal Q\). Но если в \(ab\equiv0(\ideal Q)\) многочлен \(b\) содержит постоянный член, то есть \(b^\varkappa\not\equiv0(\ideal Q)\) для каждого \(\varkappa\), то вследствие однородности базисных многочленов \(\ideal Q\) каждая однородная компонента \(ab\) делится на \(\ideal Q\), и потому \(a\) должен делиться на \(\ideal Q\).}

Пусть \(\ideal K\) есть непримарный идеал, так что по определению III существует по меньшей мере одна пара величин \(a,b\) такая, что
\begin{equation}
 a b\equiv0(\ideal K),\qquad
 a\not\equiv0(\ideal K),\qquad
 b^\varkappa\not\equiv0(\ideal K)\quad\text{для каждого }\varkappa.
\tag{3}
\end{equation}
Образуем теперь два идеала
\[
 \ideal K_1=(\ideal K,a),\qquad \ideal N_1=(\ideal K,b),
\]
которые по (3) являются собственными делителями \(\ideal K\) и для которых по (3) выполнено
\begin{equation}
 \ideal K_1\ideal N_1\equiv0(\ideal K).
\tag{4}
\end{equation}
Для элементов \(f\) наименьшего общего кратного \(\ideal K_2=[\ideal K_1,\ideal N_1]\) теперь имеет место альтернатива.

Либо из
\[
 f\equiv0(\ideal K_1),\qquad f\equiv0(\ideal N_1),
 \qquad\text{то есть}\qquad f\equiv a_1 b\;(\ideal K)
\]
всегда следует представление
\[
 f\equiv l_1 b\;(\ideal K),\qquad l_1\equiv0(\ideal K_1),
\]
а значит, по (4) \(f\equiv0(\ideal K)\). Поэтому \(\ideal K_2\equiv0(\ideal K)\), и ввиду \(\ideal K\equiv0(\ideal K_2)\) также \(\ideal K=\ideal K_2\), чем \(\ideal K\) распознается как разложимый.

Либо существует по меньшей мере одно \(f\equiv0(\ideal K_2)\), для которого такого \(l_1\) не существует. Тогда с соответствующим этому \(f\) элементом \(a_1\) образуем
\[
 \ideal K_2=(\ideal K,a,a_1),\qquad \ideal N_2=(\ideal K,b^2).
\]
Тогда вследствие \(a_1b\equiv0(\ideal K_1)\) по (4) также
\begin{equation}
 \ideal K_2\ideal N_2\equiv0(\ideal K),
\tag{4'}
\end{equation}
и \(\ideal K_2\) является собственным делителем \(\ideal K_1\).

Для элементов \(f\) из \(\ideal K_3=[\ideal K_2,\ideal N_2]\) имеет место та же альтернатива.

Либо из
\[
 f\equiv0(\ideal K_2),\qquad f\equiv0(\ideal N_2),
 \qquad\text{то есть}\qquad f\equiv a_2b^2\;(\ideal K)
\]
всегда следует
\[
 f\equiv l_2b^2\;(\ideal K),\qquad l_2\equiv0(\ideal K_2),
\]
и тем самым по (4') \(f\equiv0(\ideal K)\).

Либо для по меньшей мере одного \(f\) такого \(l_2\) нет; это приводит к образованию \(\ideal K_3=(\ideal K,a,a_1,a_2)\) и \(\ideal N_3=(\ideal K,b^3)\), причем \(\ideal K_3\ideal N_3\equiv0(\ideal K)\), где \(\ideal K_3\) является собственным делителем \(\ideal K_2\). Продолжая так, определяем в общем виде:
\[
 \ideal K_n=(\ideal K,a,a_1,\ldots,a_{n-1}),\qquad
 \ideal N_n=(\ideal K,b^n),
\]
\[
 \ideal K_n\ideal N_n\equiv0(\ideal K),
\]
где \(a_n\) определены тем, что существует \(f\) такое, что
\[
 f\equiv0(\ideal K_n),\quad f\equiv0(\ideal N_n),\quad
 f\equiv a_n b^n\;(\ideal K),\quad\text{но}\quad
 f\not\equiv l_n b^n\;(\ideal K),\quad l_n\equiv0(\ideal K_n).
\]
После этого в общем случае \(\ideal K_n\ideal N_n\equiv0(\ideal K)\), \(\ideal N_n\) по (3) является собственным делителем \(\ideal K\), а \(\ideal K_n\) -- собственным делителем \(\ideal K_{n-1}\). По теореме I о конечной цепи цепь \(\ideal K_n\) должна оборваться за конечное число шагов, скажем на \(\ideal K_n\). Для каждого \(f\equiv0(\ideal K_n)\), \(f\equiv0(\ideal N_n)\) тогда имеем \(f\equiv l_n b^n\;(\ideal K)\) с \(l_n\equiv0(\ideal K_n)\), и, следовательно, по приведенному выше выводу получается \(\ideal K=[\ideal K_n,\ideal N_n]\), чем \(\ideal K\) доказан разложимым.

Из только что доказанного \emph{единственность ассоциированных простых идеалов} получается следующим образом.

Пусть
\[
 \ideal M=[\ideal B_1\cdots\ideal B_k]=[\ideal D_1\cdots\ideal D_k]
\]
являются двумя кратчайшими, следовательно редуцированными, представлениями \(\ideal M\) как наименьшего общего кратного неразложимых идеалов, число которых по теореме IV совпадает. Тогда по этой теореме также возникающие там промежуточные представления (где, как там замечено, можно положить индекс \(j_i=i\))
\[
 \ideal M=[\ideal D_1\cdots\ideal D_{i-1}\ideal B_i\ideal B_{i+1}\cdots\ideal B_k]
       =[\ideal D_1\cdots\ideal D_{i-1}\ideal D_i\ideal B_{i+1}\cdots\ideal B_k]
       =[\ideal U_i,\ideal B_i]=[\ideal U_i,\ideal D_i]
\]
являются кратчайшими представлениями. Поэтому
\[
 \ideal U_i\ideal B_i\equiv0(\ideal D_i),\qquad \ideal U_i\not\equiv0(\ideal D_i),
 \qquad
 \ideal U_i\ideal D_i\equiv0(\ideal B_i),\qquad \ideal U_i\not\equiv0(\ideal B_i).
\]
Так как по теореме VI неразложимые идеалы \(\ideal B_i\) и \(\ideal D_i\) примарны, отсюда следует существование двух чисел \(\lambda_i\) и \(\mu_i\) таких, что
\begin{equation}
 \ideal B_i^{\lambda_i}\equiv0(\ideal D_i),
 \qquad
 \ideal D_i^{\mu_i}\equiv0(\ideal B_i).
\tag{5}
\end{equation}
Обозначим через \(\ideal P_i\), соответственно \(\bar{\ideal P}_i\), ассоциированные простые идеалы \(\ideal B_i\), соответственно \(\ideal D_i\); то есть \(\ideal P_i^{\varrho_i}\equiv0(\ideal B_i)\), \(\bar{\ideal P}_i^{\sigma_i}\equiv0(\ideal D_i)\). Тогда из (5) получается
\[
 \ideal P_i^{\lambda_i\varrho_i}\equiv0(\bar{\ideal P}_i),
 \qquad
 \bar{\ideal P}_i^{\mu_i\sigma_i}\equiv0(\ideal P_i),
\]
и отсюда по свойству простых идеалов
\[
 \ideal P_i\equiv0(\bar{\ideal P}_i),
 \qquad
 \bar{\ideal P}_i\equiv0(\ideal P_i),
 \qquad
 \ideal P_i=\bar{\ideal P}_i.
\]
Этим доказано:

\textbf{Теорема VII.} \emph{В двух различных кратчайших представлениях идеала как наименьшего общего кратного неразложимых идеалов ассоциированные простые идеалы совпадают; среди них могут встречаться и одинаковые\footnote{Это показывает, например, пример из сноски 19: \((x^2,xy,y^\lambda)=[(x^2,y),(x,y^\lambda)]\), где \(\lambda\ge2\). Оба ассоциированных простых идеала здесь: \((x,y)\).}, причем в каждом разложении одинаково часто. Следовательно, сами идеалы можно по меньшей мере одним способом попарно сопоставить так, что всякий раз некоторая степень одного идеала \(\ideal B_i\) делится на сопоставленный \(\ideal D_i\), и обратно. Их число совпадает по теореме IV.}\footnote{О единственности ``изолированных'' идеалов, содержащихся среди неразложимых идеалов, ср. \S~7.}
% END UNIT BODY
% END PRODUCER UNIT 095
\endgroup


\clearpage
\begingroup
\setcounter{footnote}{0}
% BEGIN PRODUCER UNIT 096
% Source: C:/Users/Floris/Documents/interlanguage/03_projects/noether/02_slavic_working_corpus/translations/paper19/russian/v001/Noether_Paper19_Section05_Russian_v001.tex
% Identity: 14868 B / SHA-256 B6DCC6D6691E2F41BA26EB757E0A0F62C921E335067137A21EF21F834453F4AF
\setlength{\parindent}{1.25em}
\setlength{\parskip}{0pt}
\emergencystretch=30em
\allowdisplaybreaks
\raggedbottom
\newcommand{\ideal}[1]{\mathfrak{#1}}
% BEGIN UNIT BODY
% Unit: Paper 19, Section 5. Direct Russian translation from the German final-audited source slice.

\setcounter{footnote}{20}
\subsection*{\S~5. Представление идеала как наименьшего общего кратного наибольших примарных идеалов. Единственность ассоциированных простых идеалов.}

\textbf{Определение IV.} \emph{Кратчайшее представление \(\ideal M=[\ideal Q_1\cdots\ideal Q_\alpha]\) будем называть представлением как наименьшего общего кратного наибольших примарных идеалов, если все \(\ideal Q\) примарны, но наименьшее общее кратное двух \(\ideal Q\) уже не примарно.}

То, что хотя бы одно такое представление всегда существует, следует из представления \(\ideal M\) как наименьшего общего кратного неразложимых идеалов. Эти идеалы примарны; либо уже имеется представление через наибольшие примарные идеалы, либо наименьшее общее кратное каких-либо двух идеалов снова становится примарным. Поскольку при этом число идеалов уменьшилось на единицу, повторение процедуры после конечного числа шагов приводит к желаемому представлению.

Это представление редуцировано по лемме IV. Обратно, каждое редуцированное представление через наибольшие примарные идеалы возникает именно таким образом, как показывает разложение \(\ideal Q\) на неразложимые идеалы.

Чтобы вывести здесь из теоремы VII соответствующую теорему единственности, нужно исследовать связь с ассоциированными простыми идеалами.

\textbf{Теорема VIII.} \emph{Если примарные идеалы \(\ideal N_1,\ideal N_2,\ldots,\ideal N_\lambda\) все имеют один и тот же ассоциированный простой идеал \(\ideal P\), то их наименьшее общее кратное \(\ideal Q=[\ideal N_1,\ideal N_2,\ldots,\ideal N_\lambda]\) также примарно и имеет \(\ideal P\) своим ассоциированным простым идеалом. Обратно, если \(\ideal Q=[\ideal N_1\cdots\ideal N_\lambda]\) является редуцированным представлением примарного идеала \(\ideal Q\), то все \(\ideal N_i\) примарны и имеют в качестве ассоциированного простого идеала ассоциированный простой идеал \(\ideal P\) идеала \(\ideal Q\).}

Для доказательства первой части утверждения заметим сначала, что из \(\ideal P^{\varrho_i}\equiv0(\ideal N_i)\) для каждого \(i\) также следует \(\ideal P^\tau\equiv0(\ideal Q)\), где \(\tau\) означает наибольший из индексов \(\varrho_i\). Так как \(\ideal P\) притом является делителем \(\ideal Q\), то \(\ideal P\) необходимо является ассоциированным простым идеалом, если \(\ideal Q\) примарен. Из
\[
 \ideal A\ideal B\equiv0(\ideal Q),\qquad
 \ideal B^k\not\equiv0(\ideal Q)\quad\text{(для каждого }k)
\]
следует тем самым
\[
 \ideal B\not\equiv0(\ideal P);
 \quad\text{значит}\quad
 \ideal B^k\not\equiv0(\ideal N_i)\quad\text{(для каждого }k),
\]
и поэтому \(\ideal A\equiv0(\ideal N_i)\); следовательно, \(\ideal A\equiv0(\ideal Q)\). Этим доказано, что \(\ideal Q\) примарен, а \(\ideal P\) является его ассоциированным простым идеалом.

Обратно, пусть сначала
\[
 \ideal Q=[\ideal N_1\cdots\ideal N_\lambda]=[\ideal N_i,\ideal N_i']
\]
есть кратчайшее представление \(\ideal Q\) примарными идеалами \(\ideal N_i\), и пусть \(\ideal P_i\) каждый раз обозначает соответствующий ассоциированный простой идеал. Из
\[
 \ideal N_i\ideal N_i'\equiv0(\ideal Q),\qquad
 \ideal N_i'\not\equiv0(\ideal N_i)
\]
(из-за кратчайшего представления) тогда получается
\[
 \ideal N_i'\not\equiv0(\ideal P_i),\qquad
 \text{но}\qquad
 \ideal N_i'\equiv0(\ideal P).
\]
Поскольку \(\ideal P_i\) одновременно является делителем \(\ideal Q\), \(\ideal P_i\) для каждого \(i\) совпадает с ассоциированным простым идеалом \(\ideal P\) идеала \(\ideal Q\).

Остается показать, что при каждом редуцированном представлении \(\ideal Q=[\ideal N_1\cdots\ideal N_\lambda]\) идеалы \(\ideal N_i\) примарны.\footnote{Что редуцированность представления здесь существенна, показывает пример нередуцированного представления: \(\ideal Q=[x^2,xy,y^\lambda]=[(x^2,xy,y^\lambda,yz),(x,y^\lambda)]\), где \(\lambda\ge2\). Здесь \((x^2,xy,y^\lambda,yz)=[(x^2,y),(x,y^\lambda,z)]\) по только что доказанному не является примарным; ибо последнее представление является кратчайшим представлением примарными идеалами, но ассоциированные простые идеалы \((x,y)\) и \((x,y,z)\) различны. (\(\ideal Q\) примарен по сноске 19.)} Для этого разложим \(\ideal N_i\) на их неразложимые идеалы \(\ideal B\); в возникающем таким образом кратчайшем представлении \(\ideal Q=[\ideal B_1\cdots\ideal B_s]\) каждый идеал примарен и, по только что доказанному, имеет \(\ideal P\) своим ассоциированным простым идеалом. Тогда по первой части теоремы, доказанной выше, то же самое верно для каждого \(\ideal N_i\), чем теорема VIII полностью доказана.

\textbf{Дополнение.} Отсюда еще следует, что простой идеал необходимо неразложим. Ведь из редуцированного представления \(\ideal P=[\ideal N,\ideal K]\) по теореме VIII получается \(\ideal N\equiv0(\ideal P)\); значит, поскольку \(\ideal P\equiv0(\ideal N)\), также \(\ideal P=\ideal N\), и точно так же \(\ideal P=\ideal K\). Неразложимость \(\ideal P\) получается и непосредственно: из \(\ideal P=[\ideal N,\ideal K]\) следует \(\ideal N\ideal K\equiv0(\ideal P)\), \(\ideal N\not\equiv0(\ideal P)\), \(\ideal K\not\equiv0(\ideal P)\), что противоречит определяющему свойству простого идеала.

Пусть теперь
\[
 \ideal M=[\ideal Q_1\cdots\ideal Q_\alpha]=[\ideal Q'_1\cdots\ideal Q'_\beta]
\]
-- два редуцированных представления \(\ideal M\) как наименьшего общего кратного наибольших примарных идеалов. Разлагая \(\ideal Q\) на их неразложимые идеалы \(\ideal B\), а \(\ideal Q'\) соответственно на неразложимые идеалы \(\ideal D\), получаем два редуцированных представления для \(\ideal M\) как наименьшего общего кратного неразложимых идеалов, у которых по теореме VII совпадают как число компонент, так и ассоциированные простые идеалы. По теореме VIII все неразложимые идеалы \(\ideal B\), встречающиеся при фиксированном \(\ideal Q_i\), имеют один и тот же ассоциированный простой идеал \(\ideal P_i\); тогда как принадлежащий \(\ideal Q_j\) идеал \(\ideal P_j\) необходимо отличен от него, потому что иначе по теореме VIII не было бы представления через наибольшие примарные идеалы. Следовательно, число \(\alpha\) идеалов \(\ideal Q\) равно числу различных ассоциированных простых идеалов \(\ideal P\) идеалов \(\ideal B\); эти различные \(\ideal P\) образуют ассоциированные простые идеалы идеалов \(\ideal Q\). То же самое верно для \(\ideal Q'\) относительно их разложения на \(\ideal D\). Поэтому из теоремы VII следует равенство числа \(\ideal Q\) и \(\ideal Q'\) и совпадение их ассоциированных простых идеалов. Вместе с тем видно, что указанное в начале параграфа объединение неразложимых идеалов в наибольшие примарные состоит в том, чтобы объединять все, и только те, которые имеют один и тот же ассоциированный простой идеал. Теорема VIII далее показывает свойство неразложимости наибольших примарных идеалов: они не допускают редуцированного представления как наименьшего общего кратного наибольших примарных идеалов.

Итак, доказано:

\textbf{Теорема IX.} \emph{В двух редуцированных представлениях идеала как наименьшего общего кратного наибольших примарных идеалов совпадают число компонент и ассоциированные простые идеалы, все попарно различные. Иными словами: каждому \(\ideal Q\) можно взаимно однозначно сопоставить один \(\ideal Q'\) так, что некоторая степень \(\ideal Q\) делится на \(\ideal Q'\), и обратно.}\footnote{Пример различных представлений -- пример, приведенный в сноске 12 к теореме II: \((x^2,xy)=[(x),(x^2,\mu x+y)]\) для произвольного \(\mu\). Так как ассоциированные простые идеалы \(\ideal P_1=(x)\); \(\ideal P_2=(x,y)\) различны, речь идет о наибольших примарных идеалах. -- О единственности ``изолированных'' наибольших примарных идеалов ср. \S~7.} \emph{-- Идеалы \(\ideal Q\) и \(\ideal Q'\) имеют свойство неразложимости относительно разложения на наибольшие примарные идеалы.}

\textbf{Дополнение.} Следует заметить, что теорема IX по существу сохраняется, если вместо редуцированного предполагать только кратчайшее представление. Если, например,
\[
 \ideal M=[\ideal Q_1\cdots\ideal Q_i^*\cdots\ideal Q_\alpha]
\]
редуцировано относительно \(\ideal Q_i^*\), и \(\ideal Q_i^*\) является собственным делителем \(\ideal Q_i\), а \(\ideal U_i\) -- дополнением \(\ideal Q_i\), то по лемме IV получается
\[
 \ideal Q_i=[\ideal Q_i^*,(\ideal U_i,\ideal Q_i)],
\]
и это представление редуцировано относительно \(\ideal Q_i^*\). По теореме VIII, при применении которой при необходимости следует заменить \((\ideal U_i,\ideal Q_i)\) собственным делителем, \(\ideal Q_i^*\) поэтому примарен и имеет тот же ассоциированный простой идеал \(\ideal P_i\), что и \(\ideal Q_i\). Продолжение процедуры показывает, что каждому такому представлению можно поставить в соответствие редуцированное представление через наибольшие примарные идеалы так, что число компонент и ассоциированные простые идеалы совпадают. Следовательно, и при кратчайшем представлении верно, что в двух различных представлениях совпадают число компонент и ассоциированные простые идеалы.

Определенные таким образом однозначно ассоциированные, попарно различные простые идеалы будем кратко называть ``ассоциированными простыми идеалами \(\ideal M\)''.
% END UNIT BODY
% END PRODUCER UNIT 096
\endgroup


\clearpage
\begingroup
\setcounter{footnote}{0}
% BEGIN PRODUCER UNIT 097
% Source: C:/Users/Floris/Documents/interlanguage/03_projects/noether/02_slavic_working_corpus/translations/paper19/russian/v001/Noether_Paper19_Section06_Russian_v001.tex
% Identity: 23379 B / SHA-256 C48A86D72031732B94660D002953D3AFD937A982F1F256E73D369BB4178BFA5F
\setlength{\parindent}{1.25em}
\setlength{\parskip}{0pt}
\emergencystretch=30em
\allowdisplaybreaks
\raggedbottom
\newcommand{\ideal}[1]{\mathfrak{#1}}
% BEGIN UNIT BODY
% Unit: Paper 19, Section 6. Direct Russian translation from the German final-audited source slice.

\setcounter{footnote}{22}
\subsection*{\S~6. Единственное представление идеала как наименьшего общего кратного относительно-просто-неразложимых идеалов.}

\textbf{Определение V.} \emph{Идеал \(\ideal R\) называется относительно простым к \(\ideal S\), если из
\[
 \ideal T\ideal R\equiv0(\ideal S)
\]
необходимо следует \(\ideal T\equiv0(\ideal S)\). Если и \(\ideal R\) относительно прост к \(\ideal S\), и \(\ideal S\) относительно прост к \(\ideal R\), то \(\ideal R\) и \(\ideal S\) называются взаимно относительно простыми.}\footnote{Условие относительной простоты не симметрично. Например, \(\ideal R=(x^2,y)\) относительно прост к \(\ideal S=(x)\); но \(\ideal S\) не относительно прост к \(\ideal R\), так как \(\ideal S^2\equiv0(\ideal R)\), но \(\ideal S\not\equiv0(\ideal R)\).}

\emph{Идеал называется относительно-просто-неразложимым, если он не может быть представлен как наименьшее общее кратное взаимно относительно простых собственных делителей.}

Если, в частности, вместо \(\ideal T\) взять наибольший общий делитель \(\ideal T_0\) всех \(\ideal T\), для которых \(\ideal T\ideal R\equiv0(\ideal S)\), то также \(\ideal T_0\ideal R\equiv0(\ideal S)\) и \(\ideal S\equiv0(\ideal T_0)\). Итак, \(\ideal T_0=\ideal S\), если \(\ideal R\) относительно прост к \(\ideal S\); а \(\ideal T_0\) является собственным делителем \(\ideal S\), если \(\ideal R\) не относительно прост к \(\ideal S\).\footnote{Так определенный \(\ideal T_0\) совпадает с ``остаточным модулем'' Ласкера; а в его расширении на числовые модули вместо идеалов -- с ``частным'' двух модулей у Дедекинда. Lasker, a. a. O. S. 49, Dedekind (Zahlentheorie), S. 504.}

В основе доказательства единственности лежит:

\textbf{Теорема X.} \emph{1. Если \(\ideal R\) относительно прост к идеалам \(\ideal S_1,\ldots,\ideal S_\lambda\), то \(\ideal R\) относительно прост и к их наименьшему общему кратному \(\ideal S\).}

\emph{2. Если идеалы \(\ideal S_1,\ldots,\ideal S_\lambda\) относительно просты к \(\ideal R\), то и их наименьшее общее кратное \(\ideal S\) относительно просто к \(\ideal R\).}

\emph{3. Если \(\ideal R\) относительно прост к \(\ideal S\) и \(\ideal S=[\ideal S_1\cdots\ideal S_\lambda]\) есть редуцированное представление \(\ideal S\), то \(\ideal R\) относительно прост также к каждому \(\ideal S_i\).}

\emph{4. Если \(\ideal S\) относительно прост к \(\ideal R\), то каждый делитель \(\ideal S_i\) идеала \(\ideal S\) также относительно прост к \(\ideal R\).}

1. Так как из \(\ideal T\ideal R\equiv0(\ideal S)\) необходимо следует \(\ideal T\ideal R\equiv0(\ideal S_i)\), то по предположению \(\ideal T\equiv0(\ideal S_i)\), а значит \(\ideal T\equiv0(\ideal S)\).

2. Положим
\[
 \widehat{\ideal S}_1=[\ideal S_2\cdots\ideal S_\lambda],\quad
 \widehat{\ideal S}_{12}=[\ideal S_3\cdots\ideal S_\lambda],\quad \ldots,\quad
 \widehat{\ideal S}_{12\ldots\lambda-1}=\ideal S_\lambda .
\]
Из \(\ideal T\ideal S\equiv0(\ideal R)\) тогда следует \(\ideal T\ideal S_1\widehat{\ideal S}_1\equiv0(\ideal R)\); значит, по предположению, \(\ideal T\widehat{\ideal S}_1\equiv0(\ideal R)\). Отсюда снова следует \(\ideal T\ideal S_2\widehat{\ideal S}_{12}\equiv0(\ideal R)\), значит \(\ideal T\widehat{\ideal S}_{12}\equiv0(\ideal R)\), и наконец \(\ideal T\widehat{\ideal S}_{12\ldots\lambda-1}=\ideal T\ideal S_\lambda\equiv0(\ideal R)\); следовательно, \(\ideal T\equiv0(\ideal R)\).

3. Пусть \(\widehat{\ideal S}_i\) обозначает дополнение \(\ideal S_i\). Из \(\ideal T_0\ideal R\equiv0(\ideal S_i)\) тогда, вследствие \(\widehat{\ideal S}_i\ideal R\equiv0(\ideal S_i)\), следует также
\[
 [\ideal T_0,\widehat{\ideal S}_i]\ideal R\equiv0(\ideal S).
\]
Если бы теперь \(\ideal T_0\) был собственным делителем \(\ideal S_i\), то из-за редуцированного представления \(\ideal S=[\ideal S_i,\widehat{\ideal S}_i]\) и \([\ideal T_0,\widehat{\ideal S}_i]\) был бы собственным делителем \(\ideal S\), вопреки предположению.

4. Из \(\ideal T\ideal S_i\equiv0(\ideal R)\), где \(\ideal T\) является собственным делителем \(\ideal R\), следует также \(\ideal T\ideal S\equiv0(\ideal R)\); значит, \(\ideal T\) был бы собственным делителем \(\ideal R\), вопреки предположению.

Определение V, в частности, показывает, что каждый идеал относительно прост к единичному идеалу \(\ideal D\), состоящему из всех элементов \(\Sigma\);\footnote{\(\ideal D\) играет роль единицы только относительно делимости и наименьшего общего кратного, но не относительно образования произведений. Например, \(\ideal D=(x)\) для области всех целочисленных полиномов от \(x\) без постоянного члена; \(\ideal D=(2)\) для области всех четных чисел.} однако следующие теоремы справедливы только для идеалов, отличных от \(\ideal D\).

Из только что доказанной теоремы X прежде всего получается:

\textbf{Лемма V.} \emph{Всякое представление идеала как наименьшего общего кратного взаимно относительно простых идеалов, отличных от \(\ideal D\), является редуцированным.}

Пусть
\[
 \ideal M=[\ideal R_1\ideal R_2\cdots\ideal R_\sigma]=[\ideal R_i,\ideal L_i]
\]
есть такое представление. По теореме X, 2 тогда и \(\ideal L_i\) относительно прост к \(\ideal R_i\); значит, \(\ideal R_i\) не может входить в \(\ideal L_i\), и представление кратчайшее. Ведь из \(\ideal L_i\equiv0(\ideal R_i)\), когда \(\ideal L_i\) относительно прост к \(\ideal R_i\), вследствие \(\ideal D\ideal L_i\equiv0(\ideal R_i)\) также следовало бы \(\ideal D\equiv0(\ideal R_i)\); значит, \(\ideal R_i=\ideal D\), что исключено предположением.

Пусть теперь \(\ideal R_i\) можно заменить собственным делителем \(\ideal R_i^*\). Тогда по лемме II \(\ideal R_i\) разложим и получается
\[
 \ideal R_i=[\ideal R_i^*,(\ideal R_i,\ideal L_i)].
\]
Если здесь при необходимости заменить \((\ideal R_i,\ideal L_i)\) собственным делителем \((\ideal R_i,\ideal L_i)^*\), а также \(\ideal R_i^*\) собственным делителем так, чтобы возникло редуцированное представление для \(\ideal R_i\), то теорема X, 3 показывает, что \(\ideal L_i\) также становится относительно простым к \((\ideal R_i,\ideal L_i)^*\), а этот идеал вследствие кратчайшего представления отличен от \(\ideal D\). Но так как \((\ideal R_i,\ideal L_i)^*\) входит в \((\ideal R_i,\ideal L_i)\), а \((\ideal R_i,\ideal L_i)\) входит в \(\ideal L_i\), этим доказано противоречие.

Из теоремы X далее получается теорема, передающая связь с ассоциированными простыми идеалами:

\textbf{Теорема XI.} \emph{Если \(\ideal R\) относительно прост к \(\ideal S\) и \(\ideal S\) отличен от \(\ideal D\), то ни один ассоциированный простой идеал идеала \(\ideal R\)\footnote{Определение ассоциированных простых идеалов идеала \(\ideal R\) в конце \S~5.} не делится на ассоциированный простой идеал идеала \(\ideal S\). Обратно, если такой делимости нет, то \(\ideal R\) относительно прост к \(\ideal S\), и, конечно, \(\ideal S\) отличен от \(\ideal D\).}

Пусть
\[
 \ideal R=[\ideal Q_1\cdots\ideal Q_\alpha],\qquad
 \ideal S=[\ideal Q_1^*\cdots\ideal Q_\beta^*]
\]
-- редуцированные представления \(\ideal R\) и \(\ideal S\) через наибольшие примарные идеалы, а \(\ideal P_1,\ldots,\ideal P_\alpha\) и \(\ideal P_1^*,\ldots,\ideal P_\beta^*\) -- соответствующие ассоциированные простые идеалы. Докажем утверждение в форме: если некоторый \(\ideal P\) делится на некоторый \(\ideal P^*\), то \(\ideal R\) не может быть относительно прост к \(\ideal S\), и обратно.

Итак, пусть \(\ideal P_\mu\equiv0(\ideal P_\nu^*)\), а потому также \(\ideal Q_\mu^{\varrho}\equiv0(\ideal P_\nu^*)\). Из определения \(\ideal P_\nu^*\) отсюда следует \(\ideal Q_\mu^{\varrho}\equiv0(\ideal Q_\nu^*)\), значит также \(\ideal R^{\varrho}\equiv0(\ideal Q_\nu^*)\). Пусть \(\ideal R^\tau\) -- наименьшая степень \(\ideal R\), делящаяся на \(\ideal Q_\nu^*\). При \(\tau=1\) идеал \(\ideal R\) делится на \(\ideal Q_\nu^*\); но так как по предположению \(\ideal S\) отличен от \(\ideal D\), то \(\ideal R\) не относительно прост к \(\ideal Q_\nu^*\). При \(\tau\ge2\) имеем \(\ideal R^{\tau-1}\ideal R\equiv0(\ideal Q_\nu^*)\), \(\ideal R^{\tau-1}\not\equiv0(\ideal Q_\nu^*)\), значит \(\ideal R\) не относительно прост к \(\ideal Q_\nu^*\);\footnote{\(\ideal R^0\) не определен, поскольку \(\Sigma\) не обязана содержать единицу; поэтому случай \(\tau=1\) пришлось рассмотреть заранее. \(\tau=0\) также в случае, когда \(\Sigma\) содержит единицу, исключается предположением о \(\ideal S\).} следовательно, в обоих случаях по теореме X, 3 и \(\ideal R\) не относительно прост к \(\ideal S\).

Обратно, если \(\ideal R\) не относительно прост к \(\ideal S\), то по теореме X, 1 \(\ideal R\) не относительно прост хотя бы к одному \(\ideal Q_\nu^*\). Поэтому
\[
 \ideal T_0\ideal R\equiv0(\ideal Q_\nu^*),\qquad
 \ideal T_0\not\equiv0(\ideal Q_\nu^*)
\]
и отсюда, так как \(\ideal Q_\nu^*\) примарен,
\[
 \ideal R^\varrho\equiv0(\ideal Q_\nu^*),\qquad
 \text{значит также}\quad
 \ideal Q_1^{\varrho}\cdots\ideal Q_\alpha^{\varrho}\equiv0(\ideal Q_\nu^*).
\]
Но отсюда для ассоциированных простых идеалов следует
\[
 \ideal P_1^{\varrho_1}\cdots\ideal P_\alpha^{\varrho_\alpha}\equiv0(\ideal P_\nu^*),
\]
и по свойству простых идеалов \(\ideal P_\nu^*\) тем самым входит по меньшей мере в один \(\ideal P\). Этим теорема XI доказана.

Из теорем X и XI теперь следуют существование\footnote{Существование разложения можно доказать и непосредственно, в точной аналогии с доказанным в \S~2 существованием разложения на конечное число неразложимых идеалов.} и единственность разложения на относительно-просто-неразложимые идеалы.

Пусть \(\ideal M=[\ideal Q_1\cdots\ideal Q_\alpha]\) -- редуцированное или по меньшей мере кратчайшее представление \(\ideal M\) через наибольшие примарные идеалы, а \(\ideal P_1,\ldots,\ideal P_\alpha\) -- ассоциированные простые идеалы. Объединим \(\ideal P\) в группы так, чтобы никакой идеал одной группы не делился на идеал другой группы, тогда как отдельная группа не могла бы распасться на две подгруппы, для обеих из которых выполнялось бы это свойство.

Чтобы построить такое разбиение на группы, заметим, что по определению отдельная группа \(G\) вместе с каждым идеалом \(\ideal P\) должна содержать также все его делители и кратные, встречающиеся среди \(\ideal P_1,\ldots,\ideal P_\alpha\) (то есть делящиеся на \(\ideal P\)). Пусть, например, \(\ideal P_j^{(i)}\) -- все кратные \(\ideal P\), \(\ideal P_{j_1}^{(i_1)}\) -- все делители \(\ideal P_j^{(i)}\), \(\ideal P_{j_1}^{(i_1 i_2)}\) -- все кратные \(\ideal P_{j_1}^{(i_1)}\) и так далее; вообще \(\ideal P_{j_1\ldots j_{\lambda-1}}^{(i_1\ldots i_\lambda)}\) -- все кратные \(\ideal P_{j_1\ldots j_{\lambda-1}}^{(i_1\ldots i_{\lambda-1})}\), а \(\ideal P_{j_1\ldots j_\lambda}^{(i_1\ldots i_\lambda)}\) -- все делители \(\ideal P_{j_1\ldots j_{\lambda-1}}^{(i_1\ldots i_\lambda)}\). Так как всего имеется только конечное число идеалов \(\ideal P\), процедура должна оборваться после конечного числа шагов, то есть уже не дать идеала, отличного от всех предыдущих. Полученная таким образом совокупность идеалов \(\ideal P\) действительно образует группу \(G\) с требуемыми свойствами.

Действительно, по определению \(G\) вместе с каждым \(\ideal P_{j_1\ldots j_{\lambda-1}}^{(i_1\ldots i_{\lambda-1})}\) содержит все кратные \(\ideal P_{j_1\ldots j_{\lambda-1}}^{(i_1\ldots i_\lambda)}\), а вместе с каждым \(\ideal P_{j_1\ldots j_{\lambda-1}}^{(i_1\ldots i_\lambda)}\) также все делители \(\ideal P_{j_1\ldots j_\lambda}^{(i_1\ldots i_\lambda)}\). Среди них, однако, содержатся и все делители \(\ideal P_{j_1\ldots j_{\lambda-1}}^{(i_1\ldots i_{\lambda-1})}\), тогда как кратные самих \(\ideal P_{j_1\ldots j_{\lambda-1}}^{(i_1\ldots i_\lambda)}\) снова являются \(\ideal P_{j_1\ldots j_{\lambda-1}}^{(i_1\ldots i_\lambda)}\). Поэтому идеалы, не содержащиеся в \(G\), не могут быть ни делителями, ни кратными содержащихся в \(G\).

Но \(G\) удовлетворяет и условию неразложимости. Ведь если имеется разбиение на две подгруппы \(G^{(1)}\) и \(G^{(2)}\), и \(G^{(1)}\) содержит, например, \(\ideal P_{j_1\ldots j_\lambda}^{(i_1\ldots i_\lambda)}\) или \(\ideal P_{j_1\ldots j_{\lambda-1}}^{(i_1\ldots i_\lambda)}\), то она содержит и все предшествующие, так как они попеременно являются кратными и делителями или делителями и кратными; следовательно, она содержит также \(\ideal P\), а тем самым всю группу \(G\). Если поступать аналогично с идеалами, не содержащимися в \(G\), то получается разбиение на группы \(G_1,\ldots,G_\sigma\) всех \(\ideal P\), обладающее требуемыми свойствами. Такое разбиение на группы единственно: ведь если \(G'_1,\ldots,G'_\tau\) -- второе разбиение, а \(\ideal P_{j_1\ldots j_\lambda}^{(i_1\ldots i_\lambda)}\) или \(\ideal P_{j_1\ldots j_{\lambda-1}}^{(i_1\ldots i_\lambda)}\) является элементом \(G'_i\), то по сказанному выше \(G'_i\) содержит всю группу \(G\) и потому в силу свойства неразложимости совпадает с \(G\).

Пусть теперь \(\ideal Q_{i\mu}\), где \(\mu\) пробегает от \(1\) до \(\lambda_i\), обозначают идеалы \(\ideal Q\), объединенные в одну группу \(G_i\). Так как все \(\ideal P\) попарно различны, примарные идеалы \(\ideal Q_{i\mu}\), соответствующие кратчайшему представлению, определены однозначно. Положим
\[
 \ideal R_i=[\ideal Q_{i1}\cdots\ideal Q_{i\lambda_i}],\qquad
 \text{тогда получается}\quad
 \ideal M=[\ideal R_1\cdots\ideal R_\sigma].
\]
Покажем, что этим достигнуто разложение \(\ideal M\) на относительно-просто-неразложимые идеалы. Сначала заметим, что теорема XI остается применимой и тогда, когда \(\ideal R_i=[\ideal Q_{i1}\cdots\ideal Q_{i\lambda_i}]\) является только кратчайшим представлением, поскольку по дополнению к теореме IX ассоциированные простые идеалы уже этим однозначно определены. Так как никакой ассоциированный простой идеал \(\ideal P_{i\mu}\) идеала \(\ideal R_i\) не делится на ассоциированный простой идеал \(\ideal P_{j\nu}\) идеала \(\ideal R_j\), и наоборот, то по теореме XI \(\ideal R_i\) и \(\ideal R_j\) взаимно относительно просты; а каждый отдельный \(\ideal R_i\) по теореме XI, вследствие свойства неразложимости группы \(G_i\), относительно-просто-неразложим, потому что при разложимости всегда рассматриваются идеалы, отличные от \(\ideal D\). По лемме V, кроме того, речь идет о редуцированном представлении, даже если первоначально исходили только из кратчайшего представления через \(\ideal Q\). Обратно, по теореме XI всякое разложение на относительно-просто-неразложимые идеалы приводит к указанному разбиению \(\ideal P_{i\mu}\) на группы.

Пусть теперь
\[
 \ideal M=[\bar{\ideal R}_1\cdots\bar{\ideal R}_\tau]
\]
-- второе представление \(\ideal M\) через относительно-просто-неразложимые идеалы, которое по лемме V редуцировано. Тогда, как показывает разложение \(\ideal R\) на наибольшие примарные идеалы, ассоциированные простые идеалы совпадают; следовательно, совпадают и разбиения этих простых идеалов на группы, единственность которых доказана выше. Значит, \(\tau=\sigma\); и обозначения можно выбрать так, что \(\ideal R_i\) и \(\bar{\ideal R}_i\) принадлежат одной и той же группе. Пусть тогда
\[
 \ideal M=[\ideal R_i,\ideal L_i]=[\bar{\ideal R}_i,\bar{\ideal L}_i]
\]
-- представление через идеал и дополнение. Поскольку \(\bar{\ideal R}_i\) отнесен к той же группе, что и \(\ideal R_i\), по теореме XI и \(\ideal L_i\) относительно прост к \(\bar{\ideal R}_i\), а \(\bar{\ideal L}_i\) относительно прост к \(\ideal R_i\). Из
\[
 \ideal R_i\ideal L_i\equiv0(\bar{\ideal R}_i),
 \qquad
 \bar{\ideal R}_i\bar{\ideal L}_i\equiv0(\ideal R_i)
\]
получается
\[
 \ideal R_i\equiv0(\bar{\ideal R}_i),
 \qquad
 \bar{\ideal R}_i\equiv0(\ideal R_i),
 \qquad
 \ideal R_i=\bar{\ideal R}_i.
\]
Тем самым доказано:

\textbf{Теорема XII.} \emph{Каждый идеал можно единственным образом представить как наименьшее общее кратное конечного числа взаимно относительно простых и относительно-просто-неразложимых идеалов.}
% END UNIT BODY
% END PRODUCER UNIT 097
\endgroup


\clearpage
\begingroup
\setcounter{footnote}{0}
% BEGIN PRODUCER UNIT 098
% Source: C:/Users/Floris/Documents/interlanguage/03_projects/noether/02_slavic_working_corpus/translations/paper19/russian/v001/Noether_Paper19_Section07_Russian_v001.tex
% Identity: 9975 B / SHA-256 4380BB55316686C9C315702B5AA925DFC8F72751C5FCE555DA8E55C407B02AB8
\setlength{\parindent}{1.25em}
\setlength{\parskip}{0pt}
\emergencystretch=30em
\allowdisplaybreaks
\raggedbottom
\newcommand{\ideal}[1]{\mathfrak{#1}}
% BEGIN UNIT BODY
% Unit: Paper 19, Section 7. Direct Russian translation from the German final-audited source slice.

\setcounter{footnote}{28}
\subsection*{\S~7. Единственность изолированных идеалов.}

\textbf{Определение VI.} \emph{Если кратчайшее представление \(\ideal M=[\ideal R,\ideal L]\) редуцировано относительно \(\ideal L\), то \(\ideal R\) называется изолированным идеалом, если ни один ассоциированный простой идеал идеала \(\ideal R\) не входит в ассоциированный простой идеал идеала \(\ideal L\), иначе говоря, если \(\ideal L\) относительно прост к \(\ideal R\).}

Согласно этому представление \(\ideal M=[\ideal R,\ideal L]\) удовлетворяет условиям представления \(\ideal M=[\ideal R_i,\ideal L_i]\) в лемме V, и леммой V доказано:

\textbf{Лемма VI.} \emph{Если \(\ideal R\) является изолированным идеалом кратчайшего представления \(\ideal M=[\ideal R,\ideal L]\), редуцированного относительно \(\ideal L\), то это представление редуцировано также относительно \(\ideal R\).}

Итак, поскольку для изолированных идеалов речь идет о редуцированном представлении, ассоциированные простые идеалы, возникающие при разложении \(\ideal R\) и \(\ideal L\) на неразложимые идеалы, дополняют друг друга до однозначно определенных ассоциированных простых идеалов, возникающих при соответствующем разложении \(\ideal M\). Поэтому никакой ассоциированный простой идеал идеала \(\ideal R\), относящийся к разложению на неразложимые идеалы, не входит в остальные ассоциированные простые идеалы, относящиеся к соответствующему разложению \(\ideal M\). Обратно, если это условие выполнено и \(\ideal R\) встречается хотя бы в одном представлении \(\ideal M=[\ideal R,\ideal L]\), редуцированном относительно \(\ideal R\), а значит и в редуцированном представлении \(\ideal M=[\ideal R,\ideal L^*]\), то \(\ideal R\) по определению VI изолирован. Отсюда получается не зависящее от специального дополнения \(\ideal L\) определение:

\textbf{Определение VIa.} \emph{\(\ideal R\) называется изолированным идеалом, если ассоциированные с разложением на неразложимые идеалы простые идеалы идеала \(\ideal R\) не входят в остальные простые идеалы, ассоциированные с соответствующим разложением \(\ideal M\), и если \(\ideal R\) встречается хотя бы в одном представлении \(\ideal M=[\ideal R,\ideal L]\), редуцированном относительно \(\ideal R\).}\footnote{Если бы вводились обычные ассоциированные простые идеалы (конец \S~5), то как особое условие надо было бы добавить, что все простые идеалы от \(\ideal L\) отличны от простых идеалов от \(\ideal R\). По определению VIa, следовательно, уже не надо предполагать, что представление редуцировано относительно дополнения.}

Пусть теперь
\[
 \ideal M=[\ideal R,\ideal L]=[\bar{\ideal R},\bar{\ideal L}]
\]
-- два представления \(\ideal M\) через изолированные идеалы \(\ideal R\) и \(\bar{\ideal R}\) и дополнения \(\ideal L\) и \(\bar{\ideal L}\), такие, что ассоциированные простые идеалы идеалов \(\ideal R\) и \(\bar{\ideal R}\) совпадают. Если заменить \(\ideal L\) и \(\bar{\ideal L}\) такими делителями \(\ideal L^*\) и \(\bar{\ideal L}^{*}\), чтобы представления стали редуцированными, то совпадают также ассоциированные простые идеалы от \(\ideal L^*\) и \(\bar{\ideal L}^{*}\); по теореме XI, следовательно, \(\ideal L^*\) относительно прост к \(\bar{\ideal R}\), а \(\bar{\ideal L}^{*}\) относительно прост к \(\ideal R\). Вследствие
\[
 \ideal R\ideal L^*\equiv0(\bar{\ideal R}),
 \qquad
 \bar{\ideal R}\bar{\ideal L}^{*}\equiv0(\ideal R)
\]
получается
\[
 \ideal R\equiv0(\bar{\ideal R}),
 \qquad
 \bar{\ideal R}\equiv0(\ideal R),
 \qquad
 \ideal R=\bar{\ideal R};
\]
таким образом, изолированные идеалы однозначно определены своими ассоциированными простыми идеалами. Это, в частности, дает усиление теорем VII и IX о разложении на неразложимые, соответственно наибольшие примарные, идеалы; при этом по дополнению к теореме IX достаточно предполагать лишь кратчайшее представление. В итоге получается:

\textbf{Теорема XIII.} \emph{В каждом кратчайшем представлении идеала как наименьшего общего кратного неразложимых, соответственно наибольших примарных, идеалов изолированные неразложимые, соответственно наибольшие примарные, идеалы однозначно определены; неоднозначность относится только к неизолированным неразложимым, соответственно наибольшим примарным, идеалам.}\footnote{Эта теорема для идеалов из полиномов в случае разложения на наибольшие примарные идеалы уже сообщалась без доказательства у Macaulay; его определение изолированных и неизолированных (imbedded) примарных идеалов можно рассматривать как иррациональную формулировку сообщаемой ниже.} \emph{Вообще изолированные идеалы однозначно определены своими ассоциированными простыми идеалами.}

Если в таком кратчайшем представлении через неразложимые, соответственно наибольшие примарные, идеалы \(\ideal B_i\), соответственно \(\ideal Q_j\), неизолированы, то по определению дополнения \(\ideal U_i\), соответственно \(\ideal L_j\), делятся на \(\ideal P_i\), соответственно на \(\ideal P_j\). Поэтому получается
\[
 \ideal U_i^{\varrho_i}\equiv0(\ideal B_i)
 \qquad \text{соответственно} \qquad
 \ideal L_j^{\sigma_j}\equiv0(\ideal Q_j).
\]
Из выполнения этих соотношений, обратно, следует, что \(\ideal P_i\), соответственно \(\ideal P_j\), входит по меньшей мере в один ассоциированный простой идеал дополнения, а значит, по дополнению к теореме IX, также в ассоциированный простой идеал того делителя \(\ideal L_j^*\) идеала \(\ideal L_j\), который дает представление, редуцированное относительно \(\ideal L_j^*\); поэтому \(\ideal B_i\), соответственно \(\ideal Q_j\), неизолированы. В частности, неразложимые идеалы \(\ideal B_i\), ассоциированный простой идеал которых встречается более одного раза при разложении \(\ideal M\), всегда неизолированы. \emph{Неизолированные примарные идеалы, следовательно, характеризуются также тем, что некоторая степень каждого дополнения становится делимой на них; изолированные примарные -- тем, что это не может выполняться.}
% END UNIT BODY
% END PRODUCER UNIT 098
\endgroup


\clearpage
\begingroup
\setcounter{footnote}{0}
% BEGIN PRODUCER UNIT 099
% Source: C:/Users/Floris/Documents/interlanguage/03_projects/noether/02_slavic_working_corpus/translations/paper19/russian/v001/Noether_Paper19_Section08_Russian_v001.tex
% Identity: 13707 B / SHA-256 511EED0B8D1E8A2B743F7054954A044587999CCA052D44D838BFA99D7ED6D95B
\setlength{\parindent}{1.25em}
\setlength{\parskip}{0pt}
\emergencystretch=30em
\allowdisplaybreaks
\raggedbottom
\newcommand{\ideal}[1]{\mathfrak{#1}}
% BEGIN UNIT BODY
% Unit: Paper 19, Section 8. Direct Russian translation from the German final-audited source slice.

\setcounter{footnote}{30}
\subsection*{\S~8. Единственное представление идеала как произведения взаимно-просто-неразложимых идеалов.}

Если лежащая в основе кольцевая область \(\Sigma\) обладает единицей, то есть элементом \(\varepsilon\), таким что \(\varepsilon a=a\) для каждого элемента из \(\Sigma\),\footnote{Как известно, ввиду коммутативности умножения \(\Sigma\) не может иметь более одной единицы, ибо для любых двух единиц \(\varepsilon_1\) и \(\varepsilon_2\) имеем: \(\varepsilon_1\varepsilon_2=\varepsilon_2=\varepsilon_1\).} то взаимно простые идеалы можно определить так:

\textbf{Определение VIII.} \emph{Два идеала \(\ideal R\) и \(\ideal S\) называются взаимно простыми, если их наибольший общий делитель равен единичному идеалу \(\ideal D=(\varepsilon)\), состоящему из всех элементов \(\Sigma\). Идеал называется взаимно-просто-неразложимым, если его нельзя представить как наименьшее общее кратное попарно взаимно простых идеалов.}

Отметим, что два взаимно простых идеала всегда взаимно относительно просты. В самом деле, по определению существуют два элемента \(r\equiv0(\ideal R)\), \(s\equiv0(\ideal S)\), такие что \(\varepsilon=r+s\). Из \(\ideal T\ideal R\equiv0(\ideal S)\) следует \(\ideal T r\equiv0(\ideal S)\), следовательно \(\ideal T\varepsilon=\ideal T\equiv0(\ideal S)\); и соответственно из \(\ideal T\ideal S\equiv0(\ideal R)\) также получается \(\ideal T\equiv0(\ideal R)\).\footnote{Обратное, однако, неверно; например, идеалы \(\ideal R=(x)\), \(\ideal S=(y)\) взаимно относительно просты, но не взаимно просты.} Следовательно, из леммы V следует, что всякое представление через попарно взаимно простые идеалы одновременно редуцировано.

В основе доказательства единственности лежит теорема, соответствующая теореме X:

\textbf{Теорема XIV.} \emph{Если \(\ideal R\) взаимно прост с каждым из идеалов \(\ideal S_1,\ldots,\ideal S_\lambda\), то \(\ideal R\) взаимно прост также с \(\ideal S=[\ideal S_1\cdots\ideal S_\lambda]\). Обратно, из взаимной простоты \(\ideal R\) и \(\ideal S\) следует также взаимная простота \(\ideal R\) с каждым \(\ideal S_j\). Если \(\ideal R=[\ideal R_1\cdots\ideal R_\mu]\) и каждый \(\ideal R_i\) взаимно прост с каждым \(\ideal S_j\), то \(\ideal R\) и \(\ideal S\) взаимно просты; и здесь также справедливо обратное.}

В самом деле, если \(\ideal R\) взаимно прост с каждым \(\ideal S_j\), то существуют элементы \(s_j\), такие что
\[
 s_j\equiv0(\ideal S_j),
 \qquad
 s_j\equiv\varepsilon(\ideal R).
\]
Следовательно, получается
\[
 s_1s_2\cdots s_\lambda\equiv0(\ideal S),
 \qquad
 s_1s_2\cdots s_\lambda\equiv\varepsilon(\ideal R),
 \qquad
 (\ideal R,\ideal S)=(\varepsilon).
\]
Но так как \((\ideal R,\ideal S)\) делится на каждый \((\ideal R,\ideal S_j)\), справедливо и обратное. Многократное применение этого рассуждения дает вторую часть утверждения. В самом деле, если \(\ideal R_i\) при фиксированном \(i\) взаимно прост с \(\ideal S_1,\ldots,\ideal S_\lambda\), то \(\ideal R_i\) взаимно прост с \(\ideal S\). Если это выполняется для каждого \(i\), то, поскольку понятие взаимной простоты является взаимным, \(\ideal S\) взаимно прост с \(\ideal R\). Обратно, из взаимной простоты \(\ideal R\) и \(\ideal S\) следует взаимная простота \(\ideal S\) с \(\ideal R_i\), а следовательно и \(\ideal R_i\) с \(\ideal S\).

Доказательство существования и единственности\footnote{Существование разложения снова можно доказать непосредственно по аналогии с \S~2; доказательство единственности также можно провести непосредственно (ср. сказанное во введении о Schmeidler и Noether--Schmeidler). Данное здесь доказательство вместе с тем дает представление о структуре взаимно-просто-неразложимых идеалов.} разложения на взаимно-просто-неразложимые идеалы, как и соответствующее доказательство для относительно простых идеалов, сводится к единственному разбиению на группы. Но поскольку относительно-просто-неразложимые идеалы \(\ideal R_1,\ldots,\ideal R_\sigma\) идеала \(\ideal M\) по теореме XII однозначно определены, здесь не нужно возвращаться к ассоциированным простым идеалам.

Объединим однозначно определенные относительно-просто-неразложимые идеалы \(\ideal R_1,\ldots,\ideal R_\sigma\) идеала \(\ideal M\) в группы так, что \emph{каждый идеал одной группы взаимно прост с каждым идеалом любой другой группы, тогда как отдельная группа не может быть разбита на две подгруппы так, чтобы каждый идеал одной подгруппы был взаимно прост с каждым идеалом другой подгруппы}. Такое разбиение на группы получается следующим образом: по определению отдельная группа \(G\) вместе с каждым идеалом \(\ideal R\) должна содержать также все идеалы, не взаимно простые с \(\ideal R\). Пусть они, скажем, заданы как \(\ideal R_{i_1}\); с ними не взаимно просты \(\ideal R_{i_1i_2}\); вообще пусть \(\ideal R_{i_1\ldots i_\lambda}\) не взаимно прост с \(\ideal R_{i_1\ldots i_{\lambda-1}}\). Так как всего имеется лишь конечное число идеалов, процедура после конечного числа шагов должна оборваться, то есть уже не давать идеалов, отличных от всех предыдущих; полученная таким образом совокупность идеалов образует группу \(G\) с требуемыми свойствами. Действительно, все идеалы, не содержащиеся в \(G\), по построению \(G\) взаимно просты со всеми идеалами, содержащимися в \(G\). Если далее имеется разбиение на две подгруппы \(G^{(1)}\) и \(G^{(2)}\), и пусть, например, \(\ideal R_{i_1\ldots i_\lambda}\) является элементом \(G^{(1)}\). Поскольку свойство взаимной простоты является взаимным, \(G^{(1)}\) тогда должна содержать также \(\ideal R_{i_1\ldots i_{\lambda-1}},\ldots,\ideal R_{i_1}\), следовательно также \(\ideal R\) и потому всю группу \(G\). Тем самым неразложимость доказана. Проделывая то же с группами, не содержащимися в \(G\), получаем разбиение на группы \(G_1,\ldots,G_\tau\) всех \(\ideal R\).

Это разбиение на группы единственно. Действительно, пусть \(G'_1,\ldots,G'_\tau\) есть второе разбиение и \(\ideal R_{i_1\ldots i_\lambda}\) является элементом \(G'_i\). Тогда \(G'_i\) содержит также \(\ideal R\), а следовательно \(G_1\), и ввиду условия неразложимости не может содержать элементов, отличных от элементов \(G_1\); поэтому \(G'_i\) равна \(G_1\).

Пусть теперь \(\ideal T_i\) есть наименьшее общее кратное идеалов \(\ideal R\), объединенных в одной группе \(G_i\). Покажем, что
\[
 \ideal M=[\ideal T_1\cdots\ideal T_\tau]
\]
становится представлением \(\ideal M\) через взаимно-просто-неразложимые идеалы. Прежде всего разложение \(\ideal T\) на \(\ideal R\) показывает, что \([\ideal T_1\cdots\ideal T_\tau]\) действительно представляет \(\ideal M\). Далее, по теореме XIV, идеалы \(\ideal T\) попарно взаимно просты, и каждый \(\ideal T\) взаимно-просто-неразложим. Тем самым существование такого представления доказано.

Для доказательства единственности пусть \(\ideal M=[\bar{\ideal T}_1\cdots\bar{\ideal T}_\tau]\) есть второе такое представление. Если разложить \(\bar{\ideal T}\) на их относительно-просто-неразложимые идеалы \(\ideal R\), то возникающие при разных \(\bar{\ideal T}\) идеалы \(\ideal R\) по теореме XIV взаимно просты друг с другом, следовательно также взаимно относительно просты; значит, они совпадают с однозначно определенными относительно-просто-неразложимыми идеалами \(\ideal R\) идеала \(\ideal M\). Кроме того, по теореме XIV идеалы \(\bar{\ideal T}_i\) порождают разбиение \(G'_i\) идеалов \(\ideal R\) на группы с указанными свойствами. Но так как это разбиение на группы единственно и каждый \(\bar{\ideal T}_i\) однозначно определяется группой \(G'_i=G_i\), получается \(\bar{\ideal T}_i=\ideal T_i\), чем единственность доказана.

Для попарно взаимно простых идеалов, далее, наименьшее общее кратное равно произведению. В самом деле, по теореме XIV для \(\ideal M=[\ideal T_1\cdots\ideal T_\tau]\) дополнение \(\ideal L_i\) также взаимно просто с \(\ideal T_i\); поэтому существуют элементы
\[
 t_i\equiv0(\ideal T_i),
 \qquad
 l_i\equiv0(\ideal L_i),
 \qquad
 \varepsilon=t_i+l_i.
\]
Из \(f\equiv0(\ideal T_i)\), \(f\equiv0(\ideal L_i)\) поэтому, ввиду
\[
 f=f\varepsilon=ft_i+fl_i
 \quad\text{также}\quad
 f\equiv0(\ideal L_i\ideal T_i).
\]
Так как, обратно, \(\ideal L_i\ideal T_i\) делится на \([\ideal L_i,\ideal T_i]\), получаем \([\ideal L_i,\ideal T_i]=\ideal L_i\ideal T_i\); и, продолжая процедуру на \(\ideal L_i\), в конце концов:
\[
 \ideal M=[\ideal T_1\cdots\ideal T_\tau]
       =\ideal T_1\ideal T_2\cdots\ideal T_\tau.
\]
Тем самым доказано:

\textbf{Теорема XV.} \emph{Каждый идеал единственным образом представляется как произведение конечного числа попарно взаимно простых и взаимно-просто-неразложимых идеалов.}
% END UNIT BODY
% END PRODUCER UNIT 099
\endgroup


\clearpage
\begingroup
\setcounter{footnote}{0}
% BEGIN PRODUCER UNIT 100
% Source: C:/Users/Floris/Documents/interlanguage/03_projects/noether/02_slavic_working_corpus/translations/paper19/russian/v001/Noether_Paper19_Section09_Russian_v001.tex
% Identity: 17368 B / SHA-256 4A8B18A3B9DCE3C2B067597A29ED539C9494EC356552E6D66461045AD4A01095
\setlength{\parindent}{1.25em}
\setlength{\parskip}{0pt}
\emergencystretch=30em
\allowdisplaybreaks
\raggedbottom
\newcommand{\ideal}[1]{\mathfrak{#1}}
\newcommand{\eqzero}[1]{\equiv 0\,(#1)}
% BEGIN UNIT BODY
% Unit: Paper 19, Section 9. Direct Russian translation from the German final-audited source slice.

\setcounter{footnote}{33}
\subsection*{\S~9. Распространение исследования на модули. Равенство числа компонентов при разложениях на неразложимые модули.}

Покажем теперь, что содержание первых трех параграфов, относящееся к неразложимым, а не к примарным и простым идеалам, сохраняется при более слабых предположениях. Эти параграфы действительно не используют коммутативный закон умножения и опираются только на то свойство идеалов, что они являются модулями; следовательно, они сохраняются для модулей относительно некоммутативных областей, которые теперь надо определить. В основу определения модулей следует положить двойную область \((\Sigma,T)\) со следующими свойствами:

\(\Sigma\) есть абстрактно определенное некоммутативное кольцо, то есть \(\Sigma\) есть система элементов \(a,b,c,\ldots\), для которых определены два вида операций, кольцевое сложение \((+)\) и кольцевое умножение \((\times)\), удовлетворяющие законам, установленным в \S~1, за исключением коммутативного закона 4 для кольцевого умножения.

\(T\) есть система элементов \(\alpha,\beta,\gamma,\ldots\), для которой в связи с \(\Sigma\) также определены две операции: сложение, которое из любых двух элементов \(\alpha,\beta\) однозначно образует третий \(\alpha+\beta\); умножение элемента \(\alpha\) из \(T\) на элемент \(c\) из \(\Sigma\), которое однозначно образует элемент \(c\alpha\) из \(T\).\footnote{Итак, здесь речь идет о ``правостороннем'' умножении, о ``правосторонней'' области \(T\) и, следовательно, о ``правосторонних'' модулях и идеалах. Если бы для \(T\) было положено в основу левостороннее умножение \(\alpha c\), то получилась бы соответствующая теория левосторонних модулей и идеалов; \(M\) вместе с \(\alpha\) содержит также \(\alpha c\). Ассоциативный закон здесь имел бы вид \((\gamma b)a=\gamma(b\times a)\).}

Для этих операций выполняются следующие законы:
\begin{enumerate}
\item Ассоциативный закон сложения: \((\alpha+\beta)+\gamma=\alpha+(\beta+\gamma)\).
\item Коммутативный закон сложения: \(\alpha+\beta=\beta+\alpha\).
\item Закон неограниченного и однозначного вычитания: в \(T\) существует один и только один элемент \(\xi\), удовлетворяющий уравнению \(\alpha+\xi=\beta\) (обозначают \(\xi=\beta-\alpha\)).
\item Ассоциативный закон умножения: \(a(b\gamma)=(a\times b)\gamma\).
\item Дистрибутивный закон:
\[
 (a+b)\gamma=a\gamma+b\gamma,\qquad c(\alpha+\beta)=c\alpha+c\beta.
\]
\end{enumerate}
Из этих условий, как известно, следует существование нулевого элемента, а также справедливость дистрибутивного закона для соединения вычитания и умножения:
\[
 (a-b)\gamma=a\gamma-b\gamma,\qquad c(\alpha-\beta)=c\alpha-c\beta,
\]
где \((-)\) означает вычитание в \(\Sigma\). Если \(\Sigma\) содержит единицу \(\varepsilon\), то должно выполняться \(\varepsilon\alpha=\alpha\) для каждого элемента \(\alpha\) из \(T\).

Под модулем \(M\) в \((\Sigma,T)\) будем понимать систему элементов из \(T\), удовлетворяющую двум условиям:
\begin{enumerate}
\item \(M\) вместе с \(\alpha\) содержит также \(c\alpha\), где \(c\) есть произвольный элемент из \(\Sigma\).
\item \(M\) вместе с \(\alpha\) и \(\beta\) содержит также разность \(\alpha-\beta\), следовательно вместе с \(\alpha\) также \(n\alpha\) для всякого целого числа \(n\).\footnote{Целые числа снова следует рассматривать как сокращающие знаки, а не как элементы кольца.}
\end{enumerate}
По этому определению само \(T\) образует модуль в \((\Sigma,T)\). Если, в частности, область \(T\) и заданные там операции совпадают с областью \(\Sigma\) и действующими там операциями, то модуль \(M\) переходит в правосторонний идеал \(M\) в \(\Sigma\). Если \(\Sigma\) еще предположить коммутативной, то получается обычное понятие идеала, которое тем самым оказывается частным случаем понятия модуля.\footnote{Простейший пример модуля дает модуль целочисленных линейных форм: \(\Sigma\) состоит здесь из всех целых рациональных чисел, \(T\) -- из всех целочисленных линейных форм. Несколько более общий модуль получается, если в \(\Sigma\) и \(T\) брать целые алгебраические числа вместо целых рациональных чисел, или, например, все четные числа. Если вместо линейных форм каждый раз рассматривать комплекс всех коэффициентов как один элемент, то операции в \(\Sigma\) и \(T\) действительно различны. Идеалы в некоммутативных кольцевых областях из полиномов составляют предмет совместной работы Noether--Schmeidler. Об идеалах в других специальных некоммутативных областях говорится в лекциях Hurwitz о теории чисел кватернионов (Berlin, Springer 1919) и в цитированных там работах Du Pasquier.}

Для модулей сохраняются все определения из \S~1. Так, \(\alpha\eqzero{M}\), соответственно \(N\eqzero{M}\), означает, что \(\alpha\), соответственно каждый элемент \(N\), является элементом \(M\); иначе говоря, \(\alpha\), соответственно \(N\), делится на \(M\). \(M\) становится собственным делителем \(N\), если \(M\) содержит элементы, отличные от элементов \(N\); из \(N\eqzero{M}\), \(M\eqzero{N}\) следует \(M=N\). Определения наибольшего общего делителя и наименьшего общего кратного также дословно сохраняются. Если модуль \(M\), в частности, содержит конечное число элементов \(\alpha_1,\ldots,\alpha_\rho\), таких что
\[
 M=(\alpha_1,\ldots,\alpha_\rho),
\]
то есть
\[
 \alpha=c_1\alpha_1+\cdots+c_\rho\alpha_\rho
       +n_1\alpha_1+\cdots+n_\rho\alpha_\rho
\]
для каждого \(\alpha\eqzero{M}\), причем \(c_i\) -- величины из \(\Sigma\), а \(n_i\) -- целые числа, то \(M\) называется конечным модулем, а \(\alpha_1,\ldots,\alpha_\rho\) -- базисом модуля.

В дальнейшем, аналогично \S~1, мы полагаем в основу только такие области \((\Sigma,T)\), которые удовлетворяют условию конечности: каждый модуль в \((\Sigma,T)\) конечен, то есть обладает базисом модуля.

Тогда для этой области \((\Sigma,T)\) теорема I о конечной цепи справедлива также для модулей, как показывает приведенное там доказательство, и тем самым выполнены все предпосылки для \S\S~2 и 3. Определение I и лемму I о кратчайшем и редуцированном представлении можно перенести непосредственно; так же и теорему II о представимости каждого модуля как наименьшего общего кратного конечного числа неразложимых модулей, причем лемма II показывает, что всякое такое кратчайшее представление одновременно редуцировано. Далее сохраняется теорема III, выражающая разложимость модуля через свойства его дополнения; из нее следует лемма III и, наконец, теорема IV, утверждающая равенство числа компонентов в двух различных кратчайших представлениях модуля как наименьшего общего кратного неразложимых модулей. Теорема IV дает еще лемму IV как обращение леммы I о редуцированном представлении.

\emph{Дополнение.} То же рассуждение показывает, что все эти теоремы и определения сохраняются, если для двусторонних областей \(T\), то есть таких, которые являются как правосторонними, так и левосторонними, под модулем всюду понимать двусторонний модуль, то есть такой, который вместе с \(\alpha\) содержит также \(c\alpha\) и \(\alpha c\), а вместе с \(\alpha\) и \(\beta\) также \(\alpha-\beta\).

В то время как все эти теоремы опираются только на понятие делимости и наименьшего общего кратного, дальнейшие теоремы единственности существенно основаны на понятии произведения и поэтому не допускают непосредственного переноса. Так, определение примарного и простого идеала нельзя перенести на модули, поскольку произведение двух величин из \(T\) не определено. Формально возможный перенос на некоммутативные кольца также теряет смысл, поскольку здесь не удается доказать существование ассоциированного простого идеала\footnote{В самом деле, из \(p_1^\rho\eqzero{\ideal Q}\), \(p_2^\sigma\eqzero{\ideal Q}\) здесь не следует \((p_1-p_2)^{\rho+\sigma}\eqzero{\ideal Q}\), и точно так же из \((ab)^\rho\eqzero{\ideal Q}\) не следует \(a^\rho b^\rho\eqzero{\ideal Q}\). Поэтому \(\ideal P\) нельзя установить ни как идеал, ни как обладающий свойством простого идеала. Для двусторонних идеалов \(\ideal P\) хотя и можно установить как идеал, но не как простой идеал.} и не работает также доказательство того, что неразложимый идеал примарен. Эти два факта, однако, образуют основу следующих теорем единственности. Зато, если некоммутативное кольцо обладает единицей, можно определить взаимно простые и взаимно-просто-неразложимые идеалы, и рассуждения, использованные при доказательстве теоремы II, показывают, что каждый идеал можно представить как наименьшее общее кратное конечного числа попарно взаимно простых и взаимно-просто-неразложимых идеалов.\footnote{Для специальных, ``вполне редуцируемых'' идеалов можно и здесь установить теоремы единственности; ср. совместную работу Noether--Schmeidler.}

Наконец, упомянем еще достаточный критерий того, что в \((\Sigma,T)\) выполняется условие конечности: если \(\Sigma\) содержит единицу и удовлетворяет условию конечности, а само \(T\) является конечным модулем в \((\Sigma,T)\), то каждый модуль в \((\Sigma,T)\) конечен.

В самом деле, из существования единицы в \(\Sigma\) для
\[
 \ideal M=(f_1,\ldots,f_e),\qquad
 f=\bar b_1f_1+\cdots+\bar b_ef_e+n_1f_1+\cdots+n_ef_e,
\]
где \(\bar b_i\) -- элементы из \(\Sigma\), а \(n_i\) -- целые числа, следует также представление
\[
 f=b_1f_1+\cdots+b_ef_e,
\]
где \(b_i\) -- элементы из \(\Sigma\). Ведь ввиду \(f_i=\varepsilon f_i\) получаем \(n_if_i=n_i\varepsilon f_i\), а так как \(n_i\varepsilon=(\varepsilon+\cdots+\varepsilon)\) принадлежит \(\Sigma\), величина \(b_i=\bar b_i+n_i\varepsilon\) является элементом из \(\Sigma\). Предположение о \(T\) теперь говорит, что каждый элемент \(\alpha\) из \(T\) допускает представление
\[
 \alpha=\bar a_1\tau_1+\cdots+\bar a_k\tau_k+n_1\tau_1+\cdots+n_k\tau_k,
\]
и, следовательно,
\[
 \alpha=a_1\tau_1+\cdots+a_k\tau_k,
\]
где второе представление ввиду \(\varepsilon\tau_i=\tau_i\) получается из первого так же, как выше.

Когда теперь \(\alpha\) пробегают элементы модуля \(M\) из \((\Sigma,T)\), коэффициенты \(a_k\) при \(\tau_k\) пробегают идеал \(\ideal M_k\) из \(\Sigma\); следовательно, по сказанному выше, для каждого \(a_k\eqzero{\ideal M_k}\) также
\[
 a_k=b_1a_k^{(1)}+\cdots+b_ea_k^{(e)}.
\]
Если \(\alpha^{(i)}\) обозначает элемент из \(M\), у которого коэффициент при \(\tau_k\) равен \(a_k^{(i)}\), то \(\alpha-b_1\alpha^{(1)}-\cdots-b_e\alpha^{(e)}\) становится принадлежащим \(M\) элементом, зависящим только от \(\tau_1,\ldots,\tau_{k-1}\). К совокупности этих уже не содержащих \(\tau_k\) элементов, образующих модуль \(M'\), можно снова применить процедуру; поэтому конечным числом повторений утверждение доказано.
% END UNIT BODY
% END PRODUCER UNIT 100
\endgroup


\clearpage
\begingroup
\setcounter{footnote}{0}
% BEGIN PRODUCER UNIT 101
% Source: C:/Users/Floris/Documents/interlanguage/03_projects/noether/02_slavic_working_corpus/translations/paper19/russian/v001/Noether_Paper19_Section10_Russian_v001.tex
% Identity: 18851 B / SHA-256 D48BB32F0A1E30FC612D7EDDFA0DBBD1A612032BB3214E87B008DAFC898A84EF
\setlength{\parindent}{1.25em}
\setlength{\parskip}{0pt}
\emergencystretch=30em
\allowdisplaybreaks
\raggedbottom
\newcommand{\ideal}[1]{\mathfrak{#1}}
\newcommand{\eqzero}[1]{\equiv 0\,(#1)}
% BEGIN UNIT BODY
% Unit: Paper 19, Section 10. Direct Russian translation from the German final-audited source slice.

\setcounter{footnote}{38}
\subsection*{\S~10. Специальный случай полиномиальной области.}

1. Пусть лежащая в основе кольцевая область \(\Sigma\) состоит из всех полиномов от \(x_1,
\ldots,x_n\) с произвольными комплексными коэффициентами; для нее, по теореме Гильберта о базисе модулей (Ann. 36), выполнено условие конечности. Речь идет о связи наших теорем с известными теоремами теории элиминации и теории модулей.

Эта связь устанавливается следующим специальным случаем известной теоремы Гильберта:\footnote{О полных системах инвариантов. Math. Ann. 42 (1893), \S~3, с. 313.}

Если \(f\) обращается в нуль при каждой конечной системе значений \(x_1,
\ldots,x_n\), которая является нулевой точкой всех полиномов простого идеала \(\ideal P\) (нулевой точкой \(\ideal P\)), то \(f\) делится на \(\ideal P\). Иными словами: простой идеал \(\ideal P\) состоит из всей совокупности полиномов, которые обращаются в нуль в его нулевых точках.\footnote{Как показал Lasker [Math. Ann. 60 (1905), с. 607], этот специальный случай для однородных форм можно доказать и непосредственно; обратно, отсюда снова следует теорема Гильберта в однородном и неоднородном случаях, которую можно сформулировать так: если идеал \(\ideal R\) обращается в нуль во всех нулевых точках \(M\), то некоторая степень \(\ideal R\) делится на \(M\). Однако эта теорема, как и специальный случай, справедлива только тогда, когда область значений \(x\) алгебраически замкнута; поэтому она не может следовать только из наших теорем, а должна использовать существование корней, например в том месте, что идеал, нулевые точки которого состоят только из \(x_1=0,
\ldots,x_n=0\), содержит все степенные произведения \(x\), начиная с некоторой степени. Остальное доказательство, с использованием наших теорем, можно несколько упростить по сравнению с Lasker. Ведь Lasker должен привлекать теорему Гильберта также для доказательства разложения идеала на наибольшие примарные.}

Следовательно, если произведение \(fg\) обращается в нуль во всех нулевых точках \(\ideal P\), то обращается в нуль по крайней мере один множитель; нулевые точки образуют неприводимое алгебраическое многообразие. Если, наоборот, положить это определение неприводимого многообразия в основу, то оказывается, что совокупность всех полиномов, обращающихся в нуль на неприводимом многообразии, образует простой идеал; тем самым простой идеал и неприводимое многообразие взаимно однозначно соответствуют друг другу. В силу \(\ideal P^\rho\eqzero{\ideal Q}\), \(\ideal Q\eqzero{\ideal P}\) нулевые точки примарного идеала совпадают с нулевыми точками его соответствующего простого идеала.\footnote{Это свойство примарного идеала, а именно обладать неприводимым многообразием, Macaulay (ср. Введение) берет за определение, тогда как Lasker включает в определение только понятие многообразия образования, а в остальном определяет абстрактно. Примарные идеалы, обращающиеся в нуль только при \(x_1=0,
\ldots,x_n=0\), у Lasker занимают особое положение.} Представление идеала как наименьшего общего кратного наибольших примарных идеалов тем самым дает разрешение всех нулевых точек идеала на неприводимые многообразия; и, как показал Lasker, верно также обратное. Доказательство единственности соответствующих простых идеалов здесь соответствует основной теореме теории элиминации об однозначной разложимости алгебраического многообразия на неприводимые многообразия; для специальных полиномиальных областей, где нет однозначного произведенного представления полиномов через неприводимые полиномы области, а следовательно нет и теории элиминации, его можно считать эквивалентом этой теоремы элиминации.

Неприводимые многообразия, соответствующие изолированным примарным идеалам, суть ровно те, которые появляются в ``минимальной резольвенте'';\footnote{Ср., например, J. König, Einleitung in die allgemeine Theorie der algebraischen Größen (Leipzig, Teubner, 1903), с. 235.} ведь нулевые точки каждого неизолированного наибольшего примарного идеала одновременно являются нулевыми точками по крайней мере одного изолированного, а именно такого, соответствующий простой идеал которого делится на простой идеал неизолированного идеала. Поэтому единственность изолированных примарных идеалов дает в показателях новые инвариантные числа кратности. Также теоремы единственности при разрешении примарных идеалов на неприводимые можно рассматривать как дополнение теории элиминации в смысле кратности.

После этих замечаний различные разложения можно истолковать по их поведению относительно алгебраических многообразий. Попарно взаимно простым идеалам соответствуют такие многообразия, которые не имеют общей нулевой точки; при взаимно относительно простых идеалах ни одно неприводимое многообразие одного идеала одновременно не является общей нулевой точкой; наибольшие примарные идеалы обращаются в нуль только на неприводимых многообразиях, все различающихся между собой; при разложении на неприводимые идеалы одни и те же неприводимые многообразия могут появляться и многократно.

Следует еще заметить, что вместо общей полиномиальной области можно положить в основу и область всех однородных форм, поскольку легко убедиться, что и при действующих там операциях -- сложение определено только для форм одинаковой степени -- общие теоремы сохраняются.\footnote{Что здесь, однако, в случае неоднозначности рядом с однородными могут существовать и неоднородные разложения, показывает пример \( (x^3,xy,y^3)=[(x^3,y),(y^3,x)]=[(xy,x^3,y^3,x+y^3),(xy,x^3,y^3,y+x^2)]\).}

Самый простой пример четырех различных разложений -- формулы для него приведены ниже -- по сказанному выше можно задать прямой и двумя скрещивающимися с ней прямыми, которые пересекаются между собой, причем одна из них содержит, отличную от точки пересечения, точку с большей кратностью. Разложению на взаимно простые неприводимые идеалы соответствует разложение на прямую и скрещивающееся с ней многообразие; это многообразие при разложении на относительно простые неприводимые идеалы распадается на две прямые; разложению на наибольшие примарные идеалы соответствует отделение точки большей кратности, тогда как разложение на неприводимые идеалы требует разрешения этой точки.

Если взять эту точку за начало координат, проходящую через нее прямую за ось \(y\), прямую, пересекающую ее, параллельной оси \(x\), а скрещивающуюся с ней прямую параллельной оси \(z\), то такая конфигурация, например, представляется следующими неприводимыми идеалами:\footnote{Из них первые три неприводимы как простые идеалы; \(\ideal B_4\) -- потому что имеет только делители \((x^2,y,z)\) и \((x,y,z)\); \(\ideal B_5\) -- потому что каждый делитель содержит полином \(xy\).}
\[
\begin{gathered}
 \ideal B_1=(x-1,y),\quad
 \ideal B_2=(y-1,z),\quad
 \ideal B_3=(x,z),\quad
 \ideal B_4=(x^3,y,z),\quad
 \ideal B_5=(x^2,y^2,z).
\end{gathered}
\]
Соответствующие простые идеалы будут:
\[
 \ideal P_1=\ideal B_1,
 \quad \ideal P_2=\ideal B_2,
 \quad \ideal P_3=\ideal B_3,
 \quad \ideal P_4=\ideal P_5=(x,y,z).
\]
Наибольшие примарные идеалы будут:
\[
 \ideal Q_1=\ideal B_1,
 \quad \ideal Q_2=\ideal B_2,
 \quad \ideal Q_3=\ideal B_3,
 \quad \ideal Q_4=[\ideal B_4,\ideal B_5]=(x^3,y^2,x^2y,z).
\]
Относительно простые неприводимые идеалы будут:
\[
 \ideal R_1=\ideal Q_1,
 \quad \ideal R_2=\ideal Q_2,
 \quad \ideal R_3=[\ideal Q_3,\ideal Q_4]=(x^3,x^2y,xy^2,z).
\]
Взаимно простые неприводимые идеалы будут:
\[
 \ideal S_1=\ideal R_1,
 \quad
 \ideal S_2=[\ideal R_2,\ideal R_3]=((y-1)x^3,(y-1)x^2y,(y-1)xy^2,z).
\]
Отсюда получается полный идеал:
\[
\begin{aligned}
 \ideal M&=[\ideal S_1,\ideal S_2]=\ideal S_1\ideal S_2 \\
 &=((x-1)(y-1)x^3,(y-1)x^2y,(y-1)xy^2,(x-1)z,y(y-1)x^3,yz),
\end{aligned}
\]
где
\[
 1=-(y-1)x^3+(y-1)(x^3-1)+y,
\]
\[
 (y-1)(x^3-1)+y\eqzero{\ideal S_1},
 \qquad -(y-1)x^3\eqzero{\ideal S_2}.
\]
При этом идеалы \(\ideal B_1\), \(\ideal B_2\), \(\ideal B_3\) изолированы, а потому и в разложениях на неприводимые и на наибольшие примарные идеалы определены однозначно. Идеалы \(\ideal B_4\) и \(\ideal B_5\), соответственно \(\ideal Q_4\), неизолированы; они определены не однозначно, а, например, могут быть заменены на
\[
 \ideal D_4=(x^3,y+\lambda x^2,z),
 \qquad
 \ideal D_5=(x^2+\mu xy,y^2,z),
\]
так же \(\ideal Q_4\) заменим на
\[
 \overline{\ideal Q}_4=(x^3,x^2y,y^2+\lambda xy,z).
\]

2. Так же как общая (и целочисленная) полиномиальная область, каждая конечная область целостности из полиномов -- как показывает теорема Гильберта о базисе модулей -- тоже удовлетворяет условию конечности,\footnote{Обратно, если для полиномиальной области выполнено условие конечности и каждый полином допускает по крайней мере одно представление, где множители имеют меньшую степень по \(x\), то область является конечной областью целостности.} причем коэффициенты можно брать из произвольного поля. Приведем еще пример для области всех четных полиномов как простейшей области, где из-за \(x^2y^2=(xy)^2\) не существует однозначного произведенного представления полиномов через неприводимые полиномы области. Речь идет о той же конфигурации, что и в примере выше; теперь она задается неприводимыми идеалами:
\[
\begin{gathered}
 \ideal B_1=(x^2-1,xy,y^2,yz),\quad
 \ideal B_2=(y^2-1,xz,yz,z^2),\\
 \ideal B_3=(x^2,xy,xz,yz,z^2),\quad
 \ideal B_4=(x^4,x^3y,y^2,xz,yz,z^2),\\
 \ideal B_5=(x^2,y^2,xz,yz,z^2).
\end{gathered}
\]
Соответствующие простые идеалы будут:
\[
 \ideal P_1=\ideal B_1,
 \quad \ideal P_2=\ideal B_2,
 \quad \ideal P_3=\ideal B_3,
 \quad \ideal P_4=\ideal P_5=(x^2,xy,y^2,xz,yz,z^2).
\]
То, что \(\ideal P_1\) становится простым идеалом, следует из того, что каждый полином области имеет вид
\[
 f=\varphi(z^2)+xz\,\psi(z^2)\pmod{\ideal P_1}.
\]
Итак, пусть
\[
 f_1=\varphi_1(z^2)+xz\,\psi_1(z^2),
 \qquad
 f_2=\varphi_2(z^2)+xz\,\psi_2(z^2),
\]
тогда из \(f_1f_2\eqzero{\ideal P_1}\) следует выполнение равенств
\[
 \varphi_1\varphi_2+z^2\psi_1\psi_2=0,
 \qquad
 \varphi_2\psi_1+\varphi_1\psi_2=0,
\]
и отсюда \(f_1\eqzero{\ideal P_1}\) или \(f_2\eqzero{\ideal P_1}\). Совершенно так же показывается, что \(\ideal P_2\) становится простым идеалом; \(\ideal P_3\) становится простым идеалом, поскольку каждый полином области по модулю \(\ideal P_3\) конгруэнтен некоторому полиному от \(y^2\); \(\ideal P_4\) состоит из всех полиномов области. Тем самым \(\ideal B_1\), \(\ideal B_2\) и \(\ideal B_3\) как простые идеалы тоже неприводимы; но \(\ideal B_4\) и \(\ideal B_5\) имеют каждый только единственный собственный делитель \(\ideal P_4\), а потому необходимо неприводимы.

Из неприводимых идеалов получаются наибольшие примарные:
\[
\begin{gathered}
 \ideal Q_1=\ideal B_1,
 \quad \ideal Q_2=\ideal B_2,
 \quad \ideal Q_3=\ideal B_3,
 \quad
 \ideal Q_4=[\ideal B_4,\ideal B_5]=(x^4,x^3y,y^2,xz,yz,z^2),
\end{gathered}
\]
относительно простые неприводимые идеалы:
\[
 \ideal R_1=\ideal Q_1,
 \quad \ideal R_2=\ideal Q_2,
 \quad
 \ideal R_3=[\ideal Q_3,\ideal Q_4]=(x^4,x^3y,x^2y^2,xy^3,xz,yz,z^2),
\]
взаимно простые неприводимые идеалы:
\[
 \ideal S_1=\ideal R_1,
 \quad \ideal S_2=[\ideal R_2,\ideal R_3]
 =((y^2-1)x^4,(y^2-1)x^3y,(y^2-1)x^2y^2,
 (y^2-1)xy^3,xz,yz,z^2).
\]
Как и в первом примере, будет:
\[
 [\ideal B_4,\ideal B_5]=[\ideal D_4,\ideal D_5],
 \qquad
 [\ideal Q_3,\ideal Q_4]=[\ideal Q_3,
\overline{\ideal Q}_4],
\]
где
\[
 \ideal D_4=(x^4,xy+\lambda x^2,\ldots),\qquad
 \ideal D_5=(x^2+\mu xy,\ldots),\qquad \lambda\mu\ne1,
\]
и
\[
 \overline{\ideal Q}_4=(x^4,x^3y,y^2+\lambda xy,\ldots).
\]
Остальные неприводимые, соответственно наибольшие примарные, идеалы \(\ideal B_1\), \(\ideal B_2\), \(\ideal B_3\) как изолированные идеалы определены однозначно.
% END UNIT BODY
% END PRODUCER UNIT 101
\endgroup


\clearpage
\begingroup
\setcounter{footnote}{0}
% BEGIN PRODUCER UNIT 102
% Source: C:/Users/Floris/Documents/interlanguage/03_projects/noether/02_slavic_working_corpus/translations/paper19/russian/v001/Noether_Paper19_Section11_Russian_v001.tex
% Identity: 10232 B / SHA-256 B798BFE7ADE965CF79D06A3D6F393181B6DB474743E9214D46131AEA7495B87A
\setlength{\parindent}{1.25em}
\setlength{\parskip}{0pt}
\emergencystretch=30em
\allowdisplaybreaks
\raggedbottom
\newcommand{\ideal}[1]{\mathfrak{#1}}
\newcommand{\eqzero}[1]{\equiv 0\,(#1)}
% BEGIN UNIT BODY
% Unit: Paper 19, Section 11. Direct Russian translation from the German final-audited source slice.

\setcounter{footnote}{45}
\subsection*{\S~11. Примеры из теории чисел и теории дифференциальных выражений.}

1. Пусть лежащая в основе область \(\Sigma\) состоит из всех четных целых рациональных чисел. Тогда \(\Sigma\) можно взаимно однозначно сопоставить с областью всех целых рациональных чисел, ставя каждому числу \(2a\) из \(\Sigma\) в соответствие число \(a\). Отсюда сразу следует, что каждый идеал в \(\Sigma\) является главным идеалом \((2a)\); при этом в базисном представлении \(2c=n\cdot2a\) для каждого элемента \(2c\) идеала нечетные \(n\) означают лишь сокращенные записи конечных сумм.

Простые идеалы области задаются как \(\ideal P_0=\ideal D=(2)\) и \(\ideal P=(2p)\), где \(p\) означает нечетное простое число; следовательно, каждый простой идеал делится на \(\ideal P_0\), но не делится ни на какой другой простой идеал. Примарные идеалы задаются как
\[
 \ideal Q_0=(2\cdot2^{e_0})
 \qquad\text{и}\qquad
 \ideal Q_p=(2p^e);
\]
они одновременно являются неразложимыми идеалами, и, согласно сказанному о простых идеалах, любые два, принадлежащие различным нечетным простым числам, взаимно относительно просты, но никакой \(\ideal Q\) не является относительно простым ни с каким \(\ideal Q_0\).\footnote{Действительно, из \(2b\cdot2p_1^{e_1}\eqzero{2p_2^{e_2}}\) для нечетных \(p_1\ne p_2\) всегда следует \(2b\eqzero{2p_2^{e_2}}\); из \(2b\cdot2p^e\eqzero{2\cdot2^{e_0}}\) же следует только \(2b=2\cdot2^{e_0-1}\pmod{2\cdot2^{e_0}}\).} Таким образом, \(\ideal Q_0\) являются единственными неизолированными примарными идеалами. Единственному разложению \(a\) на степени простых чисел соответствует единственное представление идеала \((2a)\) через наибольшие примарные, одновременно неразложимые идеалы:
\[
 (2a)=[(2\cdot2^{e_0}),(2p_1^{e_1}),
\ldots,(2p_\mu^{e_\mu})].
\]
Следовательно, здесь, в отличие от примеров из полиномиальной области, однозначно определены и неизолированные наибольшие примарные идеалы. При \(e_0=0\) это одновременно представление через взаимно относительно простые идеалы, тогда как при \(e_0>0\) идеал относительно-просто-неразложим.

Итак, тогда как в области всех целых чисел четыре различных разложения совпадают, здесь это имеет место лишь, с одной стороны, для двух разложений на наибольшие примарные и неразложимые идеалы, а при \(e_0>0\) -- с другой стороны, для взаимно-просто-неразложимых (каждый идеал взаимно-просто-неразложим, так как область не имеет единицы) и относительно-просто-неразложимых; при \(e_0=0\), напротив, взаимно-просто-неразложимые и относительно-просто-неразложимые разложения становятся различными. Вместе с тем здесь уже имеется пример того, что один простой идеал может делиться на другой, не будучи с ним тождественным; более общо, пример того, что из делимости не следует представление в виде произведения. Последнее -- следствие того, что область не содержит единицы, -- является также причиной того, что не существует единственного представления чисел области через неразложимые числа области в виде произведения, хотя каждый идеал становится главным; введение наименьшего общего кратного оказывается здесь необходимым. Заметим еще, что соотношения остаются в точности теми же, если вместо всех четных чисел взять за основу все числа, делимые на фиксированное простое число или степень простого числа.

Напротив, неразложимые и примарные идеалы также становятся различными, если \(\Sigma\) состоит из всех чисел, делимых на составное число \(g=p_1^{\nu_1}\cdots p_s^{\nu_s}\). Снова каждый идеал становится главным идеалом \((g\cdot a)\), а простые идеалы снова задаются как \(\ideal P_0=\ideal D=(g)\); \(\ideal P=(g\cdot p)\), где \(p\) означает простое число, отличное от входящих в \(g\) простых чисел. Напротив, неразложимые идеалы задаются как \(\ideal B_{\lambda_i}=(g\cdot p_i^{\lambda_i})\) и \(\ideal B_e=(g\cdot p^e)\), а примарные идеалы -- как \(\ideal Q_e=\ideal B_e\) и как отличные от неразложимых
\[
 \ideal Q_{\lambda_1\ldots\lambda_s}=(g\cdot p_1^{\lambda_1}\cdots p_s^{\lambda_s}),
\]
причем все \(\ideal B_{\lambda_i}\) и \(\ideal Q_{\lambda_1\ldots\lambda_s}\) имеют один и тот же ассоциированный простой идеал \(\ideal P_0=(g)\). Здесь также существует единственность разложения на неразложимые идеалы, а следовательно и для разложения на наибольшие примарные идеалы; тем самым вновь однозначно определены и неизолированные идеалы.

2. Пример некоммутативной кольцевой области дает теория идеалов в некоммутативных полиномиальных областях, рассмотренная в работе Нетер--Шмейдлера. Речь идет в частности о ``вполне редуцируемых'' идеалах, то есть таких, у которых компоненты разложения попарно взаимно просты и не имеют собственных делителей; следовательно, компоненты тем более неразложимы. Тем самым, в дополнение к доказанному там изоморфизму по \S~9, получается еще равенство числа компонентов при двух различных разложениях. Благодаря этому для разложения систем частных или обыкновенных линейных дифференциальных выражений, возникающего как специальный случай указанной работы, получен результат, который, по-видимому, не был замечен даже в известном случае обыкновенного линейного дифференциального выражения.

Вместе с тем система \(T\) всех классов вычетов фиксированного идеала \(M\) вместе с некоммутативной полиномиальной областью \(\Sigma\) дает двойную область \((\Sigma,T)\), где \(T\) обладает свойством модуля относительно \(\Sigma\). Действительно, разность двух классов вычетов снова является классом вычетов, равно как и произведение класса вычетов на произвольный полином, тогда как произведения двух классов вычетов не существует (ук. соч., \S~3). Поэтому системы классов вычетов, обозначенные там как ``подгруппы'', образуют примеры модулей в двойных областях \((\Sigma,T)\), где лежащая в основе кольцевая область \(\Sigma\) некоммутативна.
% END UNIT BODY
% END PRODUCER UNIT 102
\endgroup


\clearpage
\begingroup
\setcounter{footnote}{0}
% BEGIN PRODUCER UNIT 103
% Source: C:/Users/Floris/Documents/interlanguage/03_projects/noether/02_slavic_working_corpus/translations/paper19/russian/v001/Noether_Paper19_Section12_Russian_v001.tex
% Identity: 16507 B / SHA-256 D2B2AF3E8C40FC6BC8EC1639D02E187CD3879E41BC9149D047D080E49AD9EE01
\setlength{\parindent}{1.25em}
\setlength{\parskip}{0pt}
\emergencystretch=30em
\allowdisplaybreaks
\raggedbottom
\newcommand{\ideal}[1]{\mathfrak{#1}}
\newcommand{\eqzero}[1]{\equiv 0\,(#1)}
% BEGIN UNIT BODY
% Unit: Paper 19, Section 12. Direct Russian translation from the German final-audited source slice.

\setcounter{footnote}{46}
\subsection*{\S~12. Пример из теории элементарных делителей.}

Речь идет о таком понимании теории элементарных делителей, которое обусловлено общими разработками, но само предполагается известным.

Пусть \(\Sigma\) есть область всех целочисленных матриц из \(n^2\) элементов, для которых сложение и умножение определены в обычном для матриц смысле. Тогда \(\Sigma\) становится некоммутативной кольцевой областью; следовательно, идеалы вообще являются односторонними, а двусторонние содержатся среди них лишь как специальный случай.\footnote{Теория идеалов этих областей составляет предмет работ Du Pasquier: \emph{Zahlentheorie der Tettarionen}, Dissertation Zürich, Vierteljahrsschr. d. Naturf. Ges. Zürich, 51 (1906); \emph{Zur Theorie der Tettarionenideale}, ibid., 52 (1907). Содержание второй работы состоит в доказательстве того, что каждый идеал становится главным идеалом.}

Сначала покажем, что каждый идеал становится главным идеалом. Для этого каждой матрице \(A=(a_{ik})\) в случае правосторонних идеалов сопоставим модуль
\[
 A=(a_{11}\xi_1+\cdots+a_{1n}\xi_n,
 \ldots,
 a_{n1}\xi_1+\cdots+a_{nn}\xi_n)
\]
из целочисленных линейных форм. Обратно, этому модулю соответствует каждая матрица, задающая базис \(A\), то есть, наряду с \(A\), также \(UA\), где \(U\) унимодулярна. Более общо, произведению \(PA\) соответствует модуль \(B\), который становится кратным \(A\). Отдельная линейная форма из \(A\) задается такими \(P\), которые имеют только одну ненулевую строку.

Пусть теперь \(A_1,A_2,\ldots,A_\nu,\ldots\) -- все элементы идеала \(\ideal M\); пусть \(A_1,A_2,\ldots,A_\nu,\ldots\) -- сопоставленные им модули, \(A\) -- их наибольший общий делитель, а \(UA\) -- самая общая матрица, сопоставленная этому модулю \(A\). Тогда каждой отдельной линейной форме из \(A\), по определению наибольшего общего делителя, соответствует матрица \(P_1A_{i_1}+\cdots+P_\sigma A_{i_\sigma}\), где, по сказанному выше, \(P\) имеют только одну ненулевую строку. Отсюда следует, что и матрица \(A\), соответствующая базису \(A\), допускает такое представление, теперь уже с общим \(P\), а значит становится элементом \(\ideal M\). Так как, кроме того, каждый модуль \(A_i\) делится на \(A\), каждая матрица \(A_i\) делится на \(A\); \(A\) и, вообще, \(UA\) образует базис \(\ideal M\). Если речь идет о левосторонних идеалах, то соответственно столбцы каждой матрицы следует рассматривать как базис модуля; каждый идеал становится главным идеалом, для которого, наряду с \(A\), также \(AV\) является базисом, где под \(V\) понимается произвольная унимодулярная матрица.

Далее, для связи с теорией элементарных делителей, мы полагаем в основу двусторонние идеалы; для них, в частности, наряду с \(A\), также \(PAQ\) принадлежит идеалу. Тогда самый общий базис такого идеала, по сказанному выше, имеет вид \(UAV\), где \(U\) и \(V\) унимодулярны. Следовательно, базисные элементы исчерпывают класс эквивалентных матриц,\footnote{Базисным элементам односторонних идеалов соответствуют правые, соответственно левые, классы.} и существует взаимно однозначная связь между идеалом и классом, а следовательно также между идеалом и системой элементарных делителей \((a_1\mid a_2\mid\cdots\mid a_n)\) класса, где \(a_i\), как известно, являются неотрицательными целыми числами, каждое из которых делит следующее. Матрицу класса, появляющуюся в нормальной форме, обусловленной элементарными делителями, тем самым можно рассматривать как специальный базис идеала; из делимости идеалов, соответственно классов, следует делимость элементарных делителей, и обратно.

Однако Du Pasquier в указанном месте, \S~11, показал, что для каждого двустороннего идеала ранг равен \(n\) и все элементарные делители совпадают. Поэтому, чтобы рассматривать случай общих элементарных делителей, мы должны исходить не из идеалов, а непосредственно из двусторонних классов (которые также будут обозначаться большими готическими буквами). Класс \(\mathfrak A=UAV\) при этом делится на другой класс \(\mathfrak B\), если \(A=PBQ\).

В общем случае справедливо: \emph{наименьшее общее кратное (наибольший общий делитель) двух классов получается образованием наименьшего общего кратного (наибольшего общего делителя) соответствующих систем элементарных делителей}.\footnote{В работе \emph{Zur Theorie der Moduln}, Math. Ann. 52 (1899), S. 1, E. Steinitz определяет наименьшее общее кратное (наибольший общий делитель) классов через наименьшие общие кратные и наибольшие общие делители элементарных систем. Независимо от этого наименьшее общее кратное классов встречается как ``Kongruenzkomposition'' у H. Brandt, \emph{Komposition der binären quadratischen Formen relativ einer Grundform}, J. f. M. 150 (1919), S. 1.} Действительно, пусть
\[
 (a_1\mid a_2\mid\cdots\mid a_n),\quad
 (b_1\mid b_2\mid\cdots\mid b_n),\quad
 (c_1\mid c_2\mid\cdots\mid c_n),
\]
где \(c_i=[a_i,b_i]\), суть системы элементарных делителей классов \(\mathfrak A\), \(\mathfrak B\), \(\mathfrak C^*\), и пусть \(\mathfrak C=[\mathfrak A,\mathfrak B]\). Тогда \(\mathfrak C^*\) делится на \(\mathfrak A\) и \(\mathfrak B\), а значит делится на \(\mathfrak C\); обратно, элементарные делители \(\mathfrak C\) делятся на элементарные делители \(\mathfrak C^*\), следовательно \(\mathfrak C\) делится на \(\mathfrak C^*\), и потому \(\mathfrak C=\mathfrak C^*\). Соответственно получается доказательство для наибольшего общего делителя.

Итак, единственному представлению элементарных делителей \(a_i\) как наименьшего общего кратного степеней простых чисел соответствует представление \(\mathfrak A\) как наименьшего общего кратного классов \(\mathfrak Q\), элементарные делители которых заданы степенями одного простого числа, в символах
\[
 \mathfrak Q\sim(p^{r_1}\mid p^{r_2}\mid\cdots\mid p^{r_\rho}\mid0\mid\cdots\mid0),
 \qquad r_1\le r_2\le\cdots\le r_\rho.
\]
В частности, если ранг \(\rho\) равен \(n\), то здесь имеется разложение на взаимно простые и взаимно-просто-неразложимые классы, которое, несмотря на некоммутативность области, единственно.

Классы \(\mathfrak Q\) далее можно разложить на классы, соответствующие системе элементарных делителей:
\[
\begin{gathered}
 \mathfrak B_1\sim(p^{r_1}\mid\cdots\mid p^{r_1}),\quad
 \mathfrak B_2\sim(1\mid p^{r_2}\mid\cdots\mid p^{r_2}),\quad \ldots,\\
 \mathfrak B_\nu\sim(1\mid\cdots\mid1\mid p^{r_\nu}\mid\cdots\mid p^{r_\nu}),\quad
 \mathfrak B_{\rho+1}\sim(1\mid\cdots\mid1\mid0\mid\cdots\mid0),
\end{gathered}
\]
где в \(\mathfrak B_\nu\) каждый раз число \(1\) встречается \((\nu-1)\) раз. Если, например, \(r_1=r_2=\cdots=r_\mu=0\), то \(\mathfrak B_1,\ldots,\mathfrak B_\mu\) равны единичному классу и, следовательно, при разложении должны быть опущены; то же относится к \(\mathfrak B_{\rho+1}\), если \(\rho=n\). Если, далее, например, \(r_\nu=r_{\nu+1}=\cdots=r_{\nu+\lambda}\), то \(\mathfrak B_{\nu+1},\ldots,\mathfrak B_{\nu+\lambda}\) являются собственными делителями \(\mathfrak B_\nu\), а потому также должны быть исключены. Обозначим через \(\mathfrak B_{i_1},\ldots,\mathfrak B_{i_k}\) оставшиеся, которые теперь дают кратчайшее представление; тогда
\[
 \mathfrak Q=[\mathfrak B_{i_1},
\ldots,\mathfrak B_{i_k}]
\]
есть единственное разложение \(\mathfrak Q\) на неразложимые классы. Действительно, пусть
\[
 \mathfrak B_\nu\sim(1\mid\cdots\mid1\mid p^{r_\nu}\mid\cdots\mid p^{r_\nu})
\]
представима как наименьшее общее кратное
\[
 \mathfrak C\sim(1\mid\cdots\mid1\mid p^{s_\nu}\mid\cdots\mid p^{s_i})
 \quad\text{и}\quad
 \mathfrak D\sim(1\mid\cdots\mid1\mid p^{t_\nu}\mid\cdots\mid p^{t_i}).
\]
Тогда в системе элементарных делителей \(\mathfrak C\) и \(\mathfrak D\) на первых \((\nu-1)\) местах должно стоять число \(1\); на \(\nu\)-м месте показатель \(s_\nu\) или \(t_\nu\), скажем \(s_\nu\), должен стать равным \(r_\nu\). Но так как
\[
 r_\nu=s_\nu\le s_{\nu+1}\le\cdots\le s_i\le r_\nu
\]
то получается \(\mathfrak C=\mathfrak B_\nu\), следовательно \(\mathfrak B_\nu\) неразложим.\footnote{Напротив, \(\mathfrak B\) еще разложимы на односторонние классы; здесь уже нет взаимно однозначной связи между элементарными делителями и классом, а тем самым нет и единственного разложения на неразложимые односторонние классы. Это показывает, например, сообщенный мне H. Brandt следующий пример разложения на правосторонние классы (где классы представлены базисной матрицей, соответственно соответствующим модулем):
\[
\begin{gathered}
 \mathfrak B=[\mathfrak C_1,\mathfrak C_2]=[\mathfrak D_1,\mathfrak D_2],\quad
 \mathfrak B\sim\begin{pmatrix}p&0\\0&p\end{pmatrix},\quad
 \mathfrak C_1\sim\begin{pmatrix}1&0\\0&p\end{pmatrix},\quad
 \mathfrak C_2\sim\begin{pmatrix}p&0\\0&1\end{pmatrix},\\
 \mathfrak D_1\sim\begin{pmatrix}1&0\\0&p\end{pmatrix},
 \quad
 \mathfrak D_2\sim\begin{pmatrix}p&0\\(p-1)&1\end{pmatrix}.
\end{gathered}
\]
В самом деле, модули \((\xi,p\eta)\); \((p\xi,\eta)\); \((p\xi,(p-1)\xi+\eta)\) каждый имеют в качестве собственного делителя только модуль \((\xi,\eta)\), поэтому они неразложимы и, кроме того, различны между собой.} То же справедливо для \(\mathfrak B_{\rho+1}\), где \(p\) следует заменить на \(0\). Каждый из этих неразложимых классов, как показывает образование наименьшего общего кратного, задает определенный показатель и место, где этот показатель впервые появляется в системе элементарных делителей \(\mathfrak Q\), соответственно \(\mathfrak A\), тогда как \(\mathfrak B_{\rho+1}\) задает ранг. Так как эти числа однозначно определены системой элементарных делителей \(\mathfrak Q\), соответственно \(\mathfrak A\), и так как связь между элементарными делителями и классом взаимно однозначна, разложение \(\mathfrak Q\), а равно и произвольного класса, на неразложимые классы единственно.

В итоге можно сказать: \emph{каждый двусторонний класс \(\mathfrak A\) из целочисленных матриц с ограниченным числом элементов единственным образом представляется как наименьшее общее кратное конечного числа неразложимых двусторонних классов. Каждый неразложимый класс при этом представляет фиксированный простой делитель системы элементарных делителей \(\mathfrak A\), соответствующий показатель и место, где этот показатель впервые появляется. Неразложимый класс, соответствующий делителю \(0\), задает ранг \(\mathfrak A\).}

\begin{flushleft}
Эрланген, октябрь 1920.
\end{flushleft}
\begin{center}
(Поступила 16. 10. 1920.)
\end{center}
% END UNIT BODY
% END PRODUCER UNIT 103
\endgroup


\clearpage
\begingroup
\setcounter{footnote}{0}
% BEGIN PRODUCER UNIT 104
% Source: C:/Users/Floris/Documents/interlanguage/03_projects/noether/02_slavic_working_corpus/translations/paper20/russian/v001/Noether_Paper20_Russian_v001.tex
% Identity: 32167 B / SHA-256 EA6F48643EAE5F4EB94FFB648977E986CD3CF30D3432D039555AF31FB9F22D4A
\setlength{\parindent}{1.25em}
\setlength{\parskip}{0pt}
\emergencystretch=30em
\allowdisplaybreaks
\raggedbottom
\newcommand{\ideal}[1]{\mathfrak{#1}}
\newcommand{\eqzero}[1]{\equiv 0\,(#1)}
% BEGIN UNIT BODY
% Unit: Paper 20. Direct Russian translation from the German final-audited source slice.

\setcounter{footnote}{0}
\section*{20. Алгебраический критерий абсолютной неприводимости}
\begin{center}
Math. Ann. 85 (1922), с. 26--33
\end{center}

Многочлен -- с коэффициентами из произвольного, абстрактно определенного поля -- называется \emph{абсолютно неприводимым}, если он остается неприводимым в алгебраически замкнутом поле, до которого можно расширить область коэффициентов.\footnote{Что каждое поле, причем по существу единственным образом, можно расширить до алгебраически замкнутого поля, т. е. до такого, в котором каждый многочлен от одной переменной разлагается на линейные множители, показал E. Steinitz в своей \emph{Algebraische Theorie der Körper} (J. f. M. 137 (1910), с.~167); тем самым он дал рациональный эквивалент основной теоремы алгебры.} Ниже я покажу, что абсолютную неприводимость, в отличие от неприводимости относительно заданного поля, можно выразить алгебраически; точнее, справедлива следующая теорема:

\emph{Каждому многочлену от $n\ge2$ переменных с неопределенными коэффициентами можно сопоставить целую рациональную функцию с целочисленными коэффициентами от этих коэффициентов и дополнительных неопределенных, форму приводимости, так что для каждого специального многочлена той же степени необходимое и достаточное условие абсолютной неприводимости дается ненулевостью формы приводимости. Иными словами: для каждой специальной системы значений коэффициентов, для которой форма приводимости тождественно по неопределенным обращается в нуль и которая не вызывает понижения степени, специальный многочлен приводим, и обратно.}

Если вместо многочленов рассматривать однородные формы еще от одной переменной, то условие о степени заменяется более простым условием: система коэффициентов не должна тождественно обращаться в нуль.

Как прямое следствие теоремы получается также утверждение, впервые сформулированное \mbox{A.~Ostrowski}\footnote{\emph{Zur arithmetischen Theorie der algebraischen Größen}. Gött. Nachr. 1919, с.~279 (вспомогательная лемма, с.~296). Для случая полей, состоящих из обычных чисел, Ostrowski, насколько мне известно, также пришел к приведенной выше основной теореме; свои относящиеся к этому исследования он вскоре опубликует. То, что при моей организации доказательства основная теорема верна для произвольных, абстрактно определенных полей, основано на том, что форма приводимости, построенная для неопределенных, имеет числовые коэффициенты, не зависящие от специально положенного в основу поля.}:

\emph{Каждый абсолютно неприводимый многочлен с алгебраическими числами в качестве коэффициентов остается неприводимым по модулю каждого простого идеала конечного алгебраического поля $\Omega$, содержащего область коэффициентов, за исключением не более чем конечного числа простых идеалов из $\Omega$. В частности, $\Omega$ может быть также полем рациональных чисел.}

К дальнейшим применениям к теории относительно целых функций\footnote{D. Hilbert, \emph{Mathematische Probleme}. Gött. Nachr. 1900, с.~253, задача 14.} я намерен вернуться в другом месте.

1. Сначала покажем, что \emph{каждый многочлен от $n\ge2$ переменных с неопределенными коэффициентами абсолютно неприводим}. Именно, положим:
\begin{equation}
\begin{aligned}
F(x)&=\sum A_{i_1\ldots i_n}x_1^{i_1}\cdots x_n^{i_n} &&(i_1+\cdots+i_n\le l),\\
G(x)&=\sum B_{j_1\ldots j_n}x_1^{j_1}\cdots x_n^{j_n} &&(j_1+\cdots+j_n\le m),\\
H(x)&=F(x)G(x)=\sum C_{k_1\ldots k_n}(A,B)x_1^{k_1}\cdots x_n^{k_n} &&(k_1+\cdots+k_n\le l+m),
\end{aligned}
\tag{1}
\end{equation}
где $A$ и $B$ обозначают неопределенные, а $C(A,B)$ становятся целочисленными билинейными комбинациями $A$ и $B$. Поэтому частные
\[
D=C(A,B)/C_{0\ldots0}(A,B)
\]
являются рациональными функциями частных $A/A_{0\ldots0}$ и $B/B_{0\ldots0}$; но число этих функций больше числа их аргументов, ибо эти числа соответственно равны
\[
\binom{l+m+n}{n}-1,
\qquad
\binom{l+n}{n}-1,
\qquad
\binom{m+n}{n}-1,
\]
и
\[
\binom{l+m+n}{n}+1>
\binom{l+n}{n}+\binom{m+n}{n}
\quad\text{для } n\ge2,\ l>0,\ m>0.
\]
Следовательно, существует по меньшей мере \emph{одно} алгебраическое соотношение между $D(A,B)$, а потому также между $C(A,B)$:
\begin{equation}
 \Phi(Z)\ne0\quad[\text{тождественно по }Z_{k_1\ldots k_n}],
 \qquad
 \Phi(C(A,B))=0\quad[\text{тождественно по }A,B],
\tag{2}
\end{equation}
где $\Phi$ обозначает целочисленный многочлен.

Следовательно, многочлен с неопределенными $Z$ в качестве коэффициентов
\begin{equation}
 E(x)=\sum Z_{k_1\ldots k_n}x_1^{k_1}\cdots x_n^{k_n}
 \qquad(k_1+\cdots+k_n\le l+m)
\tag{3}
\end{equation}
при $n\ge2$ необходимо абсолютно неприводим.\footnote{Ибо неопределенные $Z$, в терминологии пункта 2, не являются нулями $\mathfrak P$.}

2. Совокупность тех многочленов $\Phi(Z)$, которые при подстановке $Z=C(A,B)$ тождественно по $A,B$ обращаются в нуль, образует \emph{идеал многочленов},\footnote{Такие совокупности обычно называют модулями; фактически же характеризующее их свойство -- вместе с $\Phi_1$ и $\Phi_2$ они содержат также $\Phi_1-\Phi_2$, а вместе с $\Phi$ также $\Psi\Phi$, где $\Psi$ означает произвольный целочисленный многочлен от $Z$, -- является свойством идеала.} точнее простой идеал $\mathfrak P$; ведь если в силу подстановки обращается в нуль произведение $\Phi_1\Phi_2$, то обращается в нуль по меньшей мере один множитель. Так как (2) выполнено тождественно по $A,B$, оно сохраняется для каждой специальной системы значений $A,B$, где $A,B$ и, следовательно, также $\Gamma=C(A,B)$ обозначают величины некоторого поля. \emph{Коэффициенты $\Gamma$ каждого многочлена, приводимого в поле расширения, удовлетворяют условиям $\Phi(\Gamma)=0$, то есть являются нулями $\mathfrak P$,} или также нулями конечного числа базисных многочленов $\mathfrak P$; теперь требуется доказать обратное.

Для этого нужно заметить, что все соотношения между $A,B,C$ сохраняются, если многочлены (1), введя новую переменную $y$, превратить в однородные формы $F(x,y)$, $G(x,y)$, $H(x,y)$ степеней $l,m,l+m$. Поэтому нулями $\mathfrak P$ становятся и такие коэффициенты $\Gamma$, при которых, например, $F(x,y)$\footnote{Многочлены и формы, возникающие при замене неопределенных в $F,G,\ldots,f,g,\ldots$ специальными системами значений, я всюду обозначаю чертой сверху: $\bar F,\bar G,\ldots,\bar f,\bar g,\ldots$.} становится степенью $y$; для них переход к неоднородным многочленам вызывает понижение степени $\bar H(x)$, а не разложение. Будет показано, что, за исключением этого понижения степени, обратное всегда верно; в частности, что для \emph{однородных форм обратное верно без исключений, если только не все $\Gamma$ обращаются в нуль; другими словами, каждый нуль $\mathfrak P$ задается параметрическим представлением $\Gamma=C(A,B)$ с величинами $A,B$ из некоторого поля расширения}.\footnote{То, что ``вообще'' нули $\mathfrak P$ задаются параметрическим представлением, следует из свойства простого идеала $\mathfrak P$ определять неприводимое алгебраическое многообразие. Но, как известно, отсюда не следует -- этого я первоначально не заметила, и на это \mbox{A.~Ostrowski} уже давно обратил мое внимание, -- что параметрическое представление не может отказать ни для какой специальной системы значений. Данное здесь доказательство опирается на более элементарные понятия.}

Это доказательство я проведу прямым построением конечного числа функций $\Phi(Z)$; именно посредством подстановки Kronecker'а
\[
 x_i=\xi^{d^{\,n-i}}
\]
я сопоставляю многочленам (1) многочлены от \emph{одной} переменной; показываю, что здесь возникают пробелы в показателях, и через образование нормы соответствующих коэффициентов прихожу к функциям $\Phi(Z)$ и к искомому критерию.

3. Положим:\footnote{\emph{Festschrift}, с.~11.}
\begin{equation}
\begin{aligned}
 d&=l+m+1,\\
 i&=i_1d^{n-1}+i_2d^{n-2}+\cdots+i_n,\quad
 j=j_1d^{n-1}+\cdots+j_n,\quad
 k=k_1d^{n-1}+\cdots+k_n.
\end{aligned}
\tag{4}
\end{equation}
Тогда по подстановке Kronecker'а многочленам $F,G,H$ из (1) соответственно отвечают многочлены
\begin{equation}
\begin{aligned}
 f(\xi)&=\sum A_{i_1\ldots i_n}\xi^i &&(i_1+\cdots+i_n\le l),\\
 g(\xi)&=\sum B_{j_1\ldots j_n}\xi^j &&(j_1+\cdots+j_n\le m),\\
 h(\xi)&=\sum C_{k_1\ldots k_n}(A,B)\xi^k &&(k_1+\cdots+k_n\le l+m).
\end{aligned}
\tag{5}
\end{equation}
Это соответствие, как известно, взаимно однозначно, поскольку для всех встречающихся в (5) индуцированных показателей из $k=0$ следует также $k_1=0,\ldots,k_n=0$; и потому из $i+j=i'+j'$ следует также $i_1+j_1=i'_1+j'_1;\ldots;i_n+j_n=i'_n+j'_n$. Значит, различные произведения $\xi^i\xi^j$ могут давать один и тот же показатель $\xi^k$ лишь тогда, когда и соответствующие произведения $x_1^{i_1}\cdots x_n^{i_n}x_1^{j_1}\cdots x_n^{j_n}$ приводят к тому же показателю. Поэтому получается:
\begin{equation}
 h(\xi)=f(\xi)g(\xi),
\tag{6}
\end{equation}
и из выполнения соотношения (6) такого рода, что в $f(\xi)$ и $g(\xi)$ не встречаются никакие показатели, отличные от индуцированных, в обратную сторону следует
\begin{equation}
 H(x)=F(x)G(x).
\tag{7}
\end{equation}
При этом \emph{среди индуцированных показателей в $f(\xi)$ и $g(\xi)$ действительно имеются пробелы}. В самом деле, наибольший показатель в $f(\xi)$ равен $ld^{n-1}$, следующий по величине равен $(l-1)d^{n-1}+d^{n-2}$; разность равна $d^{n-2}(d-1)>1$, так как $d=l+m+1\ge3$, $n\ge2$; то же верно и для $g(\xi)$.

Соотношения (6) и (7), справедливые для неопределенных $A,B$, сохраняются для каждой специальной системы значений $A,B$. Чтобы при этом сохранялись степени, я однородизую (6) и (7) посредством переменных $\eta$ и $y$. \emph{Итак, однородная форма $H(x,y)$ разлагается в алгебраическом поле расширения тогда и только тогда, когда в этом поле соответствующую бинарную форму $h(\xi,\eta)$ можно разложить на два множителя так, что в множителях не возникают никакие показатели, отличные от индуцированных.}

4. Чтобы выполнить образование нормы, упомянутое в конце пункта 2, пусть
\[
 a_1,b_1;\ldots; a_p,b_p;\quad a_{p+1},b_{p+1};\ldots; a_{p+q},b_{p+q}
\]
обозначают новые неопределенные. Положим
\begin{equation}
\begin{aligned}
 \varphi(\xi,\eta)&=(a_1\xi+b_1\eta)\cdots(a_p\xi+b_p\eta)
     =r_p\xi^p+r_{p-1}\xi^{p-1}\eta+\cdots+r_0\eta^p,\\
 \psi(\xi,\eta)&=(a_{p+1}\xi+b_{p+1}\eta)\cdots(a_{p+q}\xi+b_{p+q}\eta)
     =s_q\xi^q+s_{q-1}\xi^{q-1}\eta+\cdots+s_0\eta^q,\\
 \chi(\xi,\eta)&=\varphi(\xi,\eta)\psi(\xi,\eta)
     =t_{p+q}\xi^{p+q}+t_{p+q-1}\xi^{p+q-1}\eta+\cdots+t_0\eta^{p+q},
\end{aligned}
\tag{8}
\end{equation}
где $r,s,t$ каждый раз обозначают однородизированные элементарные симметрические функции, то есть те симметрические функции строк $a,b$, линейные в каждой отдельной строке, из которых все целые симметрические функции, однородные и одинаковой степени в каждой отдельной строке, составляются целым образом и с целочисленными коэффициентами -- причем только одним способом.

Теперь с дальнейшими неопределенными $u_{\mu\nu}$ образуем линейную форму
\begin{equation}
 \zeta(u,a,b)=\sum u_{\mu\nu}r_\mu s_\nu,
\tag{9}
\end{equation}
где суммирование распространяется на заданные индексы $\mu$ и $\nu$. Произведение $\zeta(u)$ со всеми его сопряженными, возникающими из перестановок строк $a,b$, отличными от $\zeta(u)$ и друг от друга, -- норма $N(\zeta(u))$ элемента $\zeta(u)$ относительно поля $t(a,b)$, если $\zeta(u)$ рассматривается как функция поля $r_\mu(a,b)$ и $s_\nu(a,b)$, -- тогда становится целочисленной симметрической функцией всех строк $a,b$, однородной и одинаковой степени в каждой строке; следовательно, по сказанному выше, это однозначно определенная целая целочисленная функция от $t(a,b)$ и $u_{\mu\nu}$:
\begin{equation}
 N(\zeta(u,a,b))=T(t(a,b),u)\qquad[\text{тождественно по }a,b,u].
\tag{10}
\end{equation}
То же рассуждение показывает, что также
\[
N(z-\zeta(u))=z^\delta+T_{\delta-1}(t,u)z^{\delta-1}+\cdots+T(t,u)
\qquad[\text{тождественно по }a,b,u,z]
\]
где все $T$ обозначают целые целочисленные функции от $t$ и $u$. Обе части этого тождества обращаются в нуль при $z=\zeta(u)$; именно, вследствие алгебраической независимости элементарных симметрических функций $r(a,b)$ и $s(a,b)$ они обращаются в нуль тождественно по $r,s$, если $t(a,b)$ согласно (8) положить равным целочисленной билинейной комбинации $w(r,s)$ от $r$ и $s$, а $\zeta(u)$ рассматривать как функцию от $r$ и $s$. Поэтому получается
\begin{equation}
\zeta(u)^\delta+T_{\delta-1}(w(r,s),u)\zeta(u)^{\delta-1}+\cdots+T(w(r,s),u)=0
\quad[\text{тождественно по }r,s,u].
\tag{11}
\end{equation}\footnote{Если суммирование в (9) распространяется на все индексы, то (11) представляет расширение Kronecker'а теоремы Gauss'а, по существу в Kronecker'овой организации доказательства (\emph{Zur Theorie der Formen höherer Stufen}, Berliner Ber. 1883, Werke, 2, с.~417).}

Тождества (10) и (11) сохраняются для всех специальных систем значений $\alpha,\beta$ величин $a,b$, соответственно $\varrho,\sigma$ величин $r,s$. Если $\varrho,\sigma$, в частности, выбраны так, что все произведения $\varrho_\mu\sigma_\nu$ обращаются в нуль -- где $\mu,\nu$ обозначают индексы, встречающиеся в (9), -- так что $\zeta(u)$ при подстановке обращается в нуль, то, полагая $w(\varrho,\sigma)=\tau$, из (11) получаем:
\begin{equation}
 T(\tau,u)=0\qquad[\text{тождественно по }u].
\tag{12}
\end{equation}

Обратно, пусть $\tau_0,\ldots,\tau_{p+q}$ -- такие системы значений, не все нулевые, для которых выполнено (12). Так как в алгебраическом поле расширения существует разложение
\begin{equation}
 \bar\chi(\xi,\eta)=\tau_{p+q}\xi^{p+q}+\cdots+\tau_0\eta^{p+q}
=(\alpha_1\xi+\beta_1\eta)\cdots(\alpha_{p+q}\xi+\beta_{p+q}\eta),
\tag{12'}
\end{equation}
где ни в одном множителе $\alpha$ и $\beta$ не обращаются одновременно в нуль,\footnote{Эта (рациональная) ``основная теорема алгебры'' является теоремой существования, на которой основана безысключительная справедливость обратного утверждения, высказанного в конце пункта 2.} хотя некоторые $\alpha$ или $\beta$ могут обращаться в нуль, имеем $\tau=t(\alpha,\beta)$. Подстановка этих значений в (10) с учетом (12) показывает, что $\zeta(u,\alpha,\beta)$ или один из сопряженных множителей тождественно по $u$ обращается в нуль. Следовательно, после подходящей нумерации $\alpha,\beta$, полагая $r(\alpha,\beta)=\varrho$, $s(\alpha,\beta)=\sigma$, мы утверждаем тем самым, что $\bar\chi(\xi,\eta)$ допускает по меньшей мере одно разложение
\begin{equation}
\bar\chi(\xi,\eta)=\bar\varphi(\xi,\eta)\bar\psi(\xi,\eta)
=\sum \varrho_\mu\xi^\mu\eta^{p-\mu}\cdot\sum \sigma_\nu\xi^\nu\eta^{q-\nu},
\tag{13}
\end{equation}
такое, что все произведения $\varrho_\mu\sigma_\nu$, встречающиеся в $\zeta(u)$, обращаются в нуль, но не все $\varrho$ или $\sigma$ обращаются в нуль.

5. Пусть теперь степени $p,q$ в (8) равны $ld^{n-1}$ и $md^{n-1}$, то есть степеням $f(\xi)$ и $g(\xi)$ в (5). Индексы $\mu,\nu$ разделяются на $\mu^{(1)},\nu^{(1)},\mu^{(2)},\nu^{(2)}$; при этом $\mu^{(1)}$ пробегает все показатели, отсутствующие в $f(\xi)$, $\nu^{(1)}$ -- все показатели $\xi$ в $\psi(\xi,\eta)$; соответственно $\nu^{(2)}$ пробегает все показатели, отсутствующие в $g(\xi)$, а $\mu^{(2)}$ -- все показатели $\xi$ в $\varphi(\xi,\eta)$. Далее пусть
\[
 e(\xi)=\sum Z_{k_1\ldots k_n}\xi^k\qquad(k_1+\cdots+k_n\le l+m)
\]
обозначает многочлен от одной переменной, соответствующий многочлену (3), и пусть
\[
 S(Z,u)
\]
обозначает целочисленный многочлен от $Z$ и $u$, возникающий из $T(t,u)$ -- однозначно определенной правой части (10), где суммирование в (9) распространяется на указанные индексы $\mu^{(1)},\nu^{(1)},\mu^{(2)},\nu^{(2)}$, -- посредством замены $t$ коэффициентами $Z$ и соответствующими нулями $e(\xi)$.

Если теперь заменить $r,s$ коэффициентами $A,B$ и соответствующими нулями $f(\xi)$ и $g(\xi)$, то при указанном суммировании $\zeta(u)$ обращается в нуль этой подстановкой. Далее $t=w(r,s)$ переходит в коэффициенты $C(A,B)$ и соответствующие нули $h(\xi)$; значит, $T(t,u)$ переходит в $S(C(A,B),u)$. Поэтому из тождества (11) получаем
\begin{equation}
 S(C(A,B),u)=0\qquad[\text{тождественно по }A,B,u].
\tag{14}
\end{equation}
Так как (14) сохраняется для каждой специальной системы значений, коэффициенты $\Gamma=C(A,B)$ таких специальных $H(x)$ или общих $H(x,y)$, которые в поле расширения разлагаются на два множителя, удовлетворяют условию
\begin{equation}
 S(\Gamma,u)=0\qquad[\text{тождественно по }u].
\tag{15}
\end{equation}

Если, обратно, для коэффициентов $\Gamma$ многочлена $\bar H(x)$, соответственно формы $\bar H(x,y)$, не все из которых обращаются в нуль, выполнено условие (15), то по сказанному выше это равносильно $T(\tau,u)=0$, где под $\tau$ понимаются все коэффициенты $\bar h(\xi,\eta)$. Следовательно, по (13) в поле расширения $\Gamma$ существует разложение:
\[
 \bar h(\xi,\eta)=\bar f(\xi,\eta)\cdot\bar g(\xi,\eta),
\]
где в $\bar f$ и $\bar g$ из-за заданного суммирования не встречается никаких показателей, отличных от индуцированных; следовательно, имеется соотношение
\[
 \bar H(x,y)=\bar F(x,y)\cdot\bar G(x,y).
\]
Отсюда, как только при $y=1$ ни один множитель не становится степени нуль, в частности если $\Gamma$ не вызывают понижения степени $H(x)$, следует также соотношение
\[
 \bar H(x)=\bar F(x)\cdot\bar G(x).
\]

\emph{Итак, условие (15) необходимо и достаточно для того, чтобы $\bar H(x,y)$, а при исключении понижения степени также $\bar H(x)$, разлагалась в поле расширения на два множителя.}

Отсюда еще следует, что $S(Z,u)$ не может тождественно по $Z$ и $u$ обращаться в нуль, поскольку в пункте 1 показано, что при $n\ge2$ многочлены с неопределенными в качестве коэффициентов, а следовательно и соответствующие однородные формы еще от одной переменной, абсолютно неприводимы. Поэтому, понимая под $U$ степенные произведения $u$, получаем
\[
 S(Z,u)=\sum \Phi_\lambda(Z)U_\lambda,
\]
где $\Phi_\lambda(Z)$ по (14) становятся многочленами из $\mathfrak P$. Тем самым показано, что $\mathfrak P$ не имеет иных нулей, кроме заданных параметрическим представлением $\Gamma=C(A,B)$, где $A$ и $B$ принадлежат алгебраическому полю расширения $\Gamma$. Обратное утверждение, высказанное в конце пункта 2, тем самым доказано.

6. \emph{Условие абсолютной неприводимости} теперь можно сформулировать непосредственно. Для этого нужно лишь степень $E(x)$ в (3) всеми возможными способами разложить на два положительных слагаемых $l+m$ и для каждого такого разложения образовать форму $S(Z,u)$. Произведение по этим формам
\begin{equation}
 R(Z,u)=\prod S(Z,u)
\tag{16}
\end{equation}
тогда образует форму приводимости для $E(x)$; ибо по сказанному выше каждому разложению $E(x)$, соответственно $E(x,y)$, соответствует обращение в нуль одного множителя из (16), и обратно. Тем самым доказана основная теорема, сформулированная во введении, и одновременно показано, как форму приводимости действительно можно построить за конечное число шагов.

7. Из существования формы приводимости (16) непосредственно следует упомянутая во введении теорема Ostrowski. Пусть
\[
 \bar E(x)=\sum \Gamma_{k_1\ldots k_n}x_1^{k_1}\cdots x_n^{k_n}
 \qquad(k_1+\cdots+k_n\le v)
\]
есть абсолютно неприводимый многочлен с целыми алгебраическими числами $\Gamma$ в качестве коэффициентов; пусть $\Omega$ обозначает конечное алгебраическое числовое поле, содержащее все $\Gamma$, а $\mathfrak p$ -- простой идеал из $\Omega$.

Тогда форма приводимости $R(Z,u)$, построенная для точной степени $\bar E(x)$, при подстановке $Z=\Gamma$ остается отличной от нуля; имеем
\[
 R(\Gamma,u)=\sum P_\lambda U_\lambda\ne0\qquad[\text{тождественно по }u].
\]
Значит, идеал $\mathfrak r=(P_1,P_2,\ldots)$ делится лишь на конечное число простых идеалов из $\Omega$. Если поэтому $\mathfrak p$ отличается от этих конечных, то $R(\Gamma,u)$ по модулю $\mathfrak p$ не обращается тождественно в нуль; а так как классы вычетов по модулю $\mathfrak p$ снова образуют поле, следует, что $\bar E(x)$ по модулю $\mathfrak p$ неприводим, точнее абсолютно неприводим. То же верно для $\bar E(x,y)$, так что и понижения степени по модулю $\mathfrak p$ не происходит.

Если коэффициенты $\bar E(x)$ являются алгебраически дробными числами, то $\mathfrak p$ нужно также брать отличным от конечного числа простых идеалов, входящих в общий знаменатель $\Lambda$; по модулю всех остальных простых идеалов $E(x)$ можно заменить на $\Lambda\bar E(x)$, так что приведенные выше предпосылки выполнены. Тем самым теорема доказана.

\begin{flushleft}
Erlangen, август 1921.
\end{flushleft}
\begin{center}
(Поступило 28. 8. 1921.)
\end{center}

\clearpage
% END UNIT BODY
% END PRODUCER UNIT 104
\endgroup


\clearpage
\begingroup
\setcounter{footnote}{0}
% BEGIN PRODUCER UNIT 105
% Source: C:/Users/Floris/Documents/interlanguage/03_projects/noether/02_slavic_working_corpus/translations/paper21/russian/v001/Noether_Paper21_Russian_v001.tex
% Identity: 14763 B / SHA-256 A435F1E25AB51ED2F23AA49DD97B5E35F63E5DDDFF2486941A371846B5C24A92
\setlength{\parindent}{1.25em}
\setlength{\parskip}{0pt}
\emergencystretch=30em
\allowdisplaybreaks
\raggedbottom
% BEGIN UNIT BODY
% Unit: Paper 21. Direct Russian translation from the German final-audited source slice.

\section*{21. Формальное вариационное исчисление и дифференциальные инварианты}
\begin{center}
Encyklopädie d. math. Wiss. II, 3 (1922), с. 68--71 (в: R. Weitzenböck, Differentialinvarianten)
\end{center}

\setcounter{footnote}{148}
\subsection*{28. Формальное вариационное исчисление и дифференциальные инварианты.\footnote{Этот номер принадлежит E. Noether.}}

Порождение всех дифференциальных инвариантов и вывод основных теорем можно наиболее единообразно провести формальными методами вариационного исчисления. Этот принцип восходит к Riemann\textsuperscript{72)} и получил свою главную формотеоретическую разработку у Lipschitz\textsuperscript{73)}.

Если исходить из вариационной задачи, принадлежащей однородной форме $f(dx)$, где
\[
\left|\frac{\partial^2 f}{\partial dx_i\,\partial dx_k}\right|\ne0
\]
предполагается, а именно
\begin{equation}
 \delta\int_{t_0}^{t_1} f(x')\,dt=0,
 \qquad \left(x_i'=\frac{dx_i}{dt}\right),
\tag{140}
\end{equation}
то интегрирование по частям приводит к инвариантной системе уравнений
\[
 2\psi_i=\frac{d}{dt}\frac{\partial f}{\partial x_i'}-\frac{\partial f}{\partial x_i}=0,
\]
где выражения Лагранжа $\psi_i$ контрагредиентны к $dx$. Формально это проще всего видно, если определить $\psi_i$ свободным от интегралов формальным тождеством, соответствующим интегрированию по частям (центральное уравнение Лагранжа\footnote{Для специального $f$ обозначение принадлежит Heun; см. его статью IV, 1 11, p. 447.})
\[
 \delta f-df_\delta=-2\sum\psi_i(d,d)\,\delta x_i,
 \qquad
 \left(f_\delta=\sum\frac{\partial f}{\partial dx_i}\,\delta x_i\right),
\]
где как $f$, так и $df_\delta$ являются инвариантами, а значит, инвариантна и правая часть. В частности, для квадратичной формы получается
\[
 \delta\sum g_{ik}dx_i dx_k-2d\sum g_{ik}dx_i\delta x_k
 =-2\sum\psi_\mu(d,d)\,\delta x_\mu.
\]
Если заменить здесь $d$ когредиентным $(d+\lambda\delta)$ и сравнить степени по $\lambda$, то получается соответствующая система тождеств:
\begin{equation}
\left\{
\begin{aligned}
D\sum g_{ik}dx_i dx_k-2d\sum g_{ik}dx_iDx_k
   &=-2\sum\psi_\mu(d,d)\,Dx_\mu,\\
D\sum g_{ik}dx_i\delta x_k-\delta\sum g_{ik}dx_iDx_k
   -d\sum g_{ik}\delta x_iDx_k
   &=-2\sum\psi_\mu(d,\delta)\,Dx_\mu,\\
D\sum g_{ik}\delta x_i\delta x_k-2\delta\sum g_{ik}\delta x_iDx_k
   &=-2\sum\psi_\mu(\delta,\delta)\,Dx_\mu.
\end{aligned}
\right.
\tag{141}
\end{equation}
Отсюда $\psi_\mu(d,\delta)$, а следовательно, в частности также $\psi_\mu(d,d)$ и $\psi_\mu(\delta,\delta)$, оказываются контрагредиентными векторами. Из (141) вычисляется
\[
 \psi_\mu(d,\delta)=\sum g_{i\mu}\,d\delta x_i+
 \sum\left[\begin{smallmatrix}ik\\ \mu\end{smallmatrix}\right]dx_i\delta x_k,
\]
и отсюда получается когредиентный вектор
\begin{equation}
 p^\sigma(d,\delta)=\sum g^{\sigma\mu}\psi_\mu
 =d\delta x_\sigma+
 \sum\left\{\begin{smallmatrix}ik\\ \sigma\end{smallmatrix}\right\}dx_i\delta x_k.
\tag{142}
\end{equation}
Если положить $p^\nu(d,d)=0$, то это, очевидно, дает уравнения геодезических линий.

Знание когредиентного вектора $p^\sigma(d,\delta)$ непосредственно приводит к ковариантной производной (ср. № 19). Пусть, например, $h(d,\delta)$ есть форма двух рядов $dx$ и $\delta x$; тогда выражение
\begin{equation}
 h^{(1)}(dx,\delta x)=\delta h-
 \sum\frac{\partial h}{\partial dx_\sigma}p^\sigma(d,\delta)-
 \sum\frac{\partial h}{\partial \delta x_\sigma}p^\sigma(\delta,\delta)
\tag{143}
\end{equation}
само является инвариантом как разность инвариантов и образовано так, что вторые дифференциалы взаимно уничтожаются. Поэтому $h^{(1)}(dx,dx)$ представляет форму только с первыми дифференциалами: ковариантную производную от $h$. Соответствующее утверждение верно, если $h$ содержит больше рядов $d^{(1)}x,\ldots,d^{(\lambda)}x$.\footnote{Lipschitz, J. f. Math. 72 (1870), p. 17.}

На том же принципе основано образование формы кривизны у Riemann и Lipschitz. От $\delta^2 f$ переходят к ``нормальной форме второй вариации''
\begin{equation}
 \Omega=\delta^2\sum g_{ik}dx_i dx_k
 -2d\delta\sum g_{ik}dx_i\delta x_k
 +d^2\sum g_{ik}\delta x_i\delta x_k,
\tag{144}
\end{equation}
которая, как совокупность инвариантов, сама снова является инвариантом и в которой теперь, в обобщение центрального уравнения Лагранжа, устранены третьи дифференциалы. Устранение вторых дифференциалов происходит, по аналогии с ковариантным дифференцированием, путем образования разности
\begin{equation}
 K=\Omega-2\sum\{p^\sigma(d,\delta)\psi_\sigma(d,\delta)-p^\sigma(\delta,\delta)\psi_\sigma(d,d)\},
\tag{145}
\end{equation}
что можно также записать как\footnote{Lipschitz, J. f. Math. 82 (1876), § 1.}
\begin{equation}
\begin{aligned}
 \frac12 K={}&d\sum\psi_\sigma(d,\delta)\delta x_\sigma
 -\delta\sum\psi_\sigma(d,d)\delta x_\sigma\\
 &-\sum\{p^\sigma(d,\delta)\psi_\sigma(d,\delta)-p^\sigma(\delta,\delta)\psi_\sigma(d,d)\}.
\end{aligned}
\tag{146}
\end{equation}
Закон образования $K$ можно выразить и так: в $\Omega$ вторые дифференциалы устраняются приравниванием к нулю $p(d,\delta)$ и $p(\delta,\delta)$, соответственно правых частей (141). Это определение формы кривизны по Riemann (ср. № 21), причем, однако, следует явно заметить -- что особенно ясно показывает (146), -- что это приравнивание к нулю может происходить только после выполненного дифференцирования; формальный способ образования Lipschitz обходит эту трудность. Соответственно ковариантную производную $h^{(1)}$ (143) можно было бы также определить так: в $\delta h$ вторые дифференциалы устранены приравниванием $p(d,\delta),\ldots$ к нулю.

Из формы кривизны последовательным образованием ковариантных производных получается ряд форм $[\Phi_4,\Phi_5,\ldots$ у Christoffel], проективные инварианты которого, после присоединения $f$, по Christoffel и Ricci исчерпывают совокупность дифференциальных инвариантов квадратичной формы, а эквивалентность этого ряда относительно линейных преобразований без какого-либо дальнейшего условия интегрируемости выражает эквивалентность двух квадратичных дифференциальных форм (ср. № 22). Внутреннее основание этой теоремы редукции состоит в существовании римановых нормальных координат, возникающих из вариационной задачи (140) (ср. № 20): они переводят геодезические линии, проходящие через точку, в прямые и потому при произвольном преобразовании $x$ преобразуются линейно. Тем самым образование всех инвариантов и вопрос эквивалентности сведены к соответствующему вопросу для отдельных однородных составляющих в разложении $f$ по нормальным координатам, и показывается, что эти составляющие линейно составляются из $f,\Phi_4,\Phi_5,\ldots$. Для квадратичных форм это подробно проведено у Vermeil\textsuperscript{92)} в связи с заметкой E. Noether\textsuperscript{93)}. Согласно этой заметке, теоремы и методы остаются справедливыми, если в основу положить произвольное дифференциальное выражение (ср. № 22 и 23); тем самым выявляется область действия методов формального вариационного исчисления по сравнению с методами чистой теории элиминации.\footnote{Ср. относительно этих изложений и дальнейших геометрических толкований семинарские доклады F. Klein, приведенные в разделе ``Литература''.}

Приравнивание $p(d,\delta)$ к нулю Levi-Civita\textsuperscript{134)} геометрически истолковал как ``параллельное перенесение вектора'' (ср. также Hessenberg\textsuperscript{133)}). Это понятие, аксиоматически сформулированное Weyl, по существу ограничено квадратичными формами, но допускает обобщение другого рода. Оно сохраняется при расширении группы (умножении $g_{ik}$ на произвольную функцию), как только, кроме квадратичной формы, в основу кладется еще и линейная форма. Тогда $p(d,\delta)$ определен однозначно, можно проводить образования инвариантов, соответствующие указанным выше, и определять геодезические координаты; но $p(d,d)=0$ теперь уже не возникает из вариационной задачи!\footnote{H. Weyl, Reine Infinitesimalgeometrie, Math. Ztschr. 2 (1918), p. 384--411, и Raum, Zeit, Materie (3-е и 4-е изд.).} Вопросы о теореме редукции и эквивалентности здесь еще не были рассмотрены в общем виде; только для инвариантов низшего порядка R. Weitzenböck дал теорему редукции.\footnote{Wien. Ber. 129 (1920), p. 683 и 697.}

Наконец, следует указать на общие ``инвариантные вариационные задачи'', т. е. на такие, в которых простой или кратный интеграл $J$ инвариантен относительно некоторой группы в смысле Lie. Специальными физически интересными случаями занимаются работы Hamel\footnote{Math. Ann. 59 (1904), p. 416--434; Ztschr. Math. Phys. 50 (1904), p. 1--54.}, Herglotz\footnote{Ann. d. Phys. (4) 36 (1911), p. 511, особенно § 9.}, Lorentz\footnote{Verslag der Amsterd. Akad., апрель, сентябрь 1916, октябрь 1916, май 1917.}, Fokker\footnote{Там же, январь 1917.}, Weyl\footnote{Ann. d. Phys. (4) 54 (1917), p. 117, § 2.} и Klein\footnote{Gött. Nachr. 1917, p. 469--482; там же 1918, p. 171--189.}. Принципиальная формулировка у E. Noether\footnote{Gött. Nachr. 1918, p. 235--257.} показывает, что инвариантности $J$ относительно $G_\rho$ (конечной группы с $\rho$ существенными параметрами) соответствуют $\rho$ линейно независимых дивергенций; инвариантности относительно бесконечной группы, содержащей $\rho$ произвольных функций вплоть до $\sigma$-й производной, соответствуют $\rho$ зависимостей между выражениями Лагранжа и их производными вплоть до порядка $\sigma$. В обоих случаях верно обратное. Так как выражения Лагранжа становятся относительными инвариантами группы, тем самым имеется и процесс, порождающий инварианты.

\begin{center}
(Завершено в марте 1921.)
\end{center}
% END UNIT BODY
% END PRODUCER UNIT 105
\endgroup


\clearpage
\begingroup
\setcounter{footnote}{0}
% BEGIN PRODUCER UNIT 106
% Source: C:/Users/Floris/Documents/interlanguage/03_projects/noether/02_slavic_working_corpus/translations/paper22/russian/v001/Noether_Paper22_Russian_v001.tex
% Identity: 115631 B / SHA-256 AD9A9C83F14B6929AFBF89C0F5CE00EF3B83C9A638C285B4BDBD3D3E5D76D7EC
\setlength{\parindent}{1.25em}
\setlength{\parskip}{0pt}
\makeatletter
\renewcommand\footnotesize{\@setfontsize\footnotesize{7.7}{8.7}\selectfont}
\makeatother
\emergencystretch=35em
\allowdisplaybreaks
\sloppy
\hbadness=10000
\vbadness=10000
\hfuzz=2pt
\vfuzz=10pt
\raggedbottom
\providecommand{\Amod}{\mathfrak A}
\providecommand{\Bmod}{\mathfrak B}
\providecommand{\Gmod}{\mathfrak G}
\providecommand{\Mmod}{\mathfrak M}
\providecommand{\Nmod}{\mathfrak N}
\providecommand{\Rmod}{\mathfrak R}
\providecommand{\Smod}{\mathfrak S}
\providecommand{\mideal}{\mathfrak m}
\providecommand{\nideal}{\mathfrak n}
\providecommand{\gideal}{\mathfrak g}
\providecommand{\hideal}{\mathfrak h}
\providecommand{\jideal}{\mathfrak j}
\providecommand{\tideal}{\mathfrak t}
\providecommand{\aideal}{\mathfrak a}
\providecommand{\bideal}{\mathfrak b}
\providecommand{\bb}{\mathfrak b}
\providecommand{\cc}{\mathfrak c}
\providecommand{\zero}[1]{\equiv 0\pmod{#1}}
\providecommand{\Norm}{N}
% BEGIN UNIT BODY
% Unit: Paper22 v001. Direct Russian translation from the German final-audited source slice.
\setcounter{footnote}{0}
\section*{22. Обработка K. Hentzelt:\\ к теории полиномиальных идеалов и результантов}
\begin{center}
Math. Ann. 88 (1923), с. 53--79
\end{center}

В дальнейшем речь идет об \emph{общей задаче теории исключения: определить общие нули всех полиномов идеала полиномов}; или, что то же самое, нули любого конечного числа полиномов, образующих базис идеала; при этом одновременно получится понятийное истолкование возникающих кратностей. Коэффициенты полиномов могут быть отнесены к произвольному полю; тогда нули принадлежат полученному из него алгебраически замкнутому полю. Кроме того, без ограничения общности идеал можно считать преобразованным (§~3). Тогда существенный результат таков:

\emph{Каждому идеалу $\mideal$ из полиномов можно сопоставить результантную форму:}
\[
 R_{\mideal}=R^{(1)}(x_1,\ldots,x_n)\cdots
 R^{(i)}(x_i,\ldots,x_n)\cdots R^{(n)}(x_n)\zero{\mideal},
\]
\emph{такую, что $R_{\mideal}$ обращается в нуль для всех нулей $\mideal$ и только для них. Если $\nideal$ есть делитель $\mideal$, а результантные формы $R_{\nideal}$ и $R_{\mideal}$ совпадают, то совпадают и идеалы $\nideal$ и $\mideal$.}\footnote{Kurt Hentzelt с октября 1914 года числится пропавшим без вести под Dixmuiden и должен быть причислен к погибшим. Данная работа представляет собой совершенно свободную переработку наиболее существенной части его диссертации ``Zur Theorie der Formenmoduln und Resultanten'', с которой он летом 1914 года получил степень доктора у E. Fischer в Erlangen. Эта диссертация, написанная целиком на основе его собственных идей, построена без пробелов; но лемма следует за леммой, все понятия описаны формулами с четырьмя и пятью индексами, текст почти полностью отсутствует, так что понимание встречает величайшие трудности. До запланированной переработки он сам уже не дошел. Я излагаю работу в чисто понятийной форме, благодаря чему достигается большое упрощение доказательств, основные мысли которых повсюду восходят к Hentzelt, и, как я надеюсь, становится явной красота работы. -- Части диссертации, которые -- при заданном базисе -- относятся к вопросу построения возникающих функций конечным числом шагов, должны быть оставлены для отдельной публикации. (E. N.)}

Вторая часть главной теоремы показывает, что результантная форма дает нули с характеристической кратностью. Эта вторая часть аналогична тому факту, что два идеала алгебраического числового поля, один из которых является делителем другого, совпадают тогда и только тогда, когда совпадают их \emph{нормы}; внутренняя причина обеих теорем также одна и та же.

А именно, $\mideal$ можно рассматривать как модуль линейных форм $\mathfrak M_{i-1}$, если каждый полином рассматривать как линейную форму от степенных произведений $x_1,\ldots,x_{i-1}$ и присоединить $x_{i+1},\ldots,x_n$ к области коэффициентов. Тогда результантный множитель $R^{(i)}$ становится равным \emph{норме} соответствующего основного модуля (§~1) $\Gmod_{i-1}$ относительно $\mathfrak M_{i-1}$; отсюда сразу получается и \emph{истолкование кратности -- произведения степени и показателя множителя $R^{(i)}$ -- через число линейно независимых классов вычетов}\footnote{Эти последние результаты проявляют свое настоящее значение, если привлечь разложение идеалов на первичные; кратность, соответствующая первичному идеалу, тогда определяется как число линейно независимых классов вычетов дополнения по этому идеалу; это будет рассмотрено в другом месте. (E. N.)}, факт, до сих пор известный только в специальном случае, когда результант можно определить для неопределенных коэффициентов\footnote{Ср. Macaulay, \emph{The algebraic theory of modular systems}, Cambridge Tracts 19 (Cambridge University Press, 1916), No. 67; также Lasker, Zur Theorie der Moduln und Ideale, Math. Ann. 60 (1905), с. 20--116; с. 98. Что в этом случае результант и результантная форма совпадают, показано в еще не опубликованных теоремах Hentzelt, упомянутых в примечании 1). (E. N.)}.

Для вывода этих результатов в §~1 и §~2 даются теоремы о модулях из линейных форм, коэффициенты которых суть целые рациональные числа или полиномы от одной неизвестной. Связь с идеалами из полиномов, введенными в §~3, устанавливает теорема изоморфизма §~4. §~5 дает вышеупомянутую вторую часть главной теоремы, §~7 -- теорему о нулях, которой в §~6 предшествует общая теория исключения. Там же, при введении новых неизвестных, доказывается разложение резольвенты на линейные формы от этих неизвестных; тем самым для данного здесь исключения получается безупречное доказательство одного утверждения Kronecker\footnote{Попытка доказательства у König, \emph{Algebraische Größen} (Leipzig 1903, Teubner), V, §~4 -- для кронекеровской теории исключения -- как известно, не удалась. Что и для кратностей, введенных König в кронекеровскую теорию исключения, вторая часть главной теоремы Hentzelt не имеет силы, показывает Macaulay, loc. cit., с. 23; там далее показано, что кронекеровская теория исключения не соответствует разложению на первичные идеалы. (E. N.)}.

\section*{§ 1. Модули из линейных форм.}

Под целыми величинами $a,b,c,\ldots$ в первых двух параграфах понимаются целые рациональные числа или полиномы от одной неизвестной с коэффициентами из произвольного, абстрактно заданного поля $P$; под линейными формами $a(\xi),b(\xi),\ldots$ понимаются такие формы от конечного числа неизвестных $\xi_1,\ldots,\xi_k$ с целыми величинами в качестве коэффициентов.

\emph{Модуль $\Amod$ из линейных форм} определяется условиями: вместе с $a(\xi)$ и $b(\xi)$ модуль $\Amod$ содержит также разность $a(\xi)-b(\xi)$; вместе с $a(\xi)$ он содержит также $c\cdot a(\xi)$, где под $c$ понимается произвольная целая величина.

Под \emph{рангом $\Amod$} мы, как обычно, понимаем максимальное число линейно независимых линейных форм из $\Amod$; под $\lambda$-м \emph{детерминантным делителем} $d_\lambda$ модуля $\Amod$ -- наибольший общий делитель всех детерминантов порядка $\lambda$ из $\Amod$ (то есть всех детерминантов порядка $\lambda$, которые можно образовать из матрицы коэффициентов любых $\lambda$ линейных форм из $\Amod$). Если $\rho$ есть ранг $\Amod$, так что $d_1,\ldots,d_\rho$ являются ненулевыми детерминантными делителями, то \emph{элементарные делители} модуля $\Amod$ определяются формулами $e_1=d_1$, $e_2=d_2/d_1,\ldots,e_\rho=d_\rho/d_{\rho-1}$, $e_{\rho+i}=0$. Как известно, $\Amod$ есть конечный модуль, имеющий модульный базис ровно из $\rho$ линейных форм $a_1(\xi),\ldots,a_\rho(\xi)$; через них каждая линейная форма из $\Amod$ представляется линейно, с целыми величинами в качестве коэффициентов:
$\Amod=(a_1(\xi),\ldots,a_\rho(\xi))$. Если $A$ обозначает матрицу коэффициентов этих линейных форм, то детерминантные и элементарные делители $\Amod$ совпадают с соответствующими делителями $A$; отсюда прежде всего следует, что \emph{каждый элементарный делитель $e_i$ делит следующий $e_{i+1}$}. Матричное равенство, даваемое теорией элементарных делителей,
\[
 P A Q=\begin{pmatrix}
 e_1&&&0\\
 &e_2&&\\
 &&\ddots&\\
 0&&&e_\rho
 \end{pmatrix}
\]
далее показывает -- если унимодулярным преобразованием $\xi=Q(\eta)$ ввести новые неизвестные $\eta_1,\ldots,\eta_k$ -- что имеет место равенство модулей:
\begin{equation}
 \Amod=(e_1\eta_1,e_2\eta_2,\ldots,e_\rho\eta_\rho).
\tag{1}
\end{equation}

Самое существенное для дальнейшего понятие -- это понятие основного модуля\footnote{Ср. Steinitz, Rechteckige Systeme und Moduln in algebraischen Zahlkörpern II. Math. Ann. 72 (1912), с. 297--345, No. 36. Впрочем, Hentzelt, по-видимому, во всяком случае независимо от Steinitz -- с которым и в других местах имеются точки соприкосновения -- пришел для своего более простого случая к этому понятию в форме формулы (4). (E. N.)}, согласно следующему

\medskip
\noindent\textbf{Определение I.} \emph{Модуль $\Gmod$ называется основным модулем, если он не имеет собственного делителя того же ранга}\footnote{Делитель понимается в модульном смысле: $\Amod$ делится на $\Bmod$, $\Amod\zero{\Bmod}$, если каждый элемент из $\Amod$ содержится в $\Bmod$; $\Bmod$ называется собственным делителем, если он содержит элементы, отличные от элементов $\Amod$.}.

Итак, $\Gmod$ является основным модулем тогда и только тогда, когда из
\begin{equation}
 c\cdot g(\xi)\zero{\Gmod};\quad c\ne0 \quad\text{всегда следует:}\quad g(\xi)\zero{\Gmod}.
\tag{2}
\end{equation}

Связь между произвольным модулем и основным модулем задается следующей

\medskip
\noindent\textbf{Теорема I.} \emph{Каждый модуль $\Amod$ делится на один и только один основной модуль $\Gmod$ того же ранга; $\Gmod$ определяется как наименьший делитель $\Amod$ того же ранга и представляется в виде}
\begin{equation}
 \Gmod=(\eta_1,\ldots,\eta_\rho);
\tag{3}
\end{equation}
\emph{$\Gmod$ будет называться основным модулем модуля $\Amod$.}

Условие, что $\Gmod$ должен быть наименьшим делителем того же ранга, выражается так: $\Gmod$ содержит все и только те линейные формы $g(\xi)$, которые удовлетворяют соотношению
\begin{equation}
 b\cdot g(\xi)\zero{\Amod};\quad b\ne0.
\tag{4}
\end{equation}
Отсюда прежде всего следует, что $\Gmod$ есть модуль, поскольку наряду с
\[
 b_1g_1(\xi)\zero{\Amod};\quad b_1\ne0;
 \qquad
 b_2g_2(\xi)\zero{\Amod};\quad b_2\ne0
\]
выполняется также
\[
 b_1b_2\bigl(g_1(\xi)-g_2(\xi)\bigr)\zero{\Amod};\quad b_1b_2\ne0,
\]
и, конечно, также $b_1\cdot c g_1(\xi)\zero{\Amod}$. Но $\Gmod$ является и основным модулем; действительно, из
\[
 c\cdot g(\xi)\zero{\Gmod};\quad c\ne0
\]
по (4) следует
\[
 b c g(\xi)\zero{\Amod};\quad bc\ne0,
\]
а значит, снова по (4), также $g(\xi)\zero{\Gmod}$.

Кроме того, не существует основного модуля $\Gmod_1$, отличного от наименьшего делителя $\Gmod$ того же ранга, который удовлетворял бы этим условиям. Ведь $\Gmod_1$ должен был бы делиться на $\Gmod$, но как основной модуль он не имеет собственного делителя того же ранга, а потому совпадает с $\Gmod$. Наконец, модуль, представленный правой частью (3), является основным модулем, так как всякий собственный делитель должен был бы содержать еще одну неизвестную $\eta$, а значит имел бы больший ранг. Тем самым (3) дает однозначно определенный основной модуль модуля $\Amod$, и \emph{Теорема I доказана во всех частях}.

Еще одно существенное для дальнейшего понятие -- дедекендово понятие частного двух модулей\footnote{У Dedekind речь идет о числовых модулях, для которых также определено умножение, так что дедекендово частное не совпадает с частным, определенным здесь. (E. N.)}, именно:

\medskip
\noindent\textbf{Определение II.} \emph{Частное $\cc=\Amod/\Bmod$ двух модулей из линейных форм определяется как совокупность целых величин $c$, удовлетворяющих условию}
\[
 c\cdot\Bmod\zero{\Amod}.
\]
\emph{$\cc$ представляет собой идеал из целых величин, следовательно главный идеал. Соответственно, частное $\mathfrak C=\Amod/\bb$ модуля $\Amod$ из линейных форм на идеал $\bb$ из целых величин определяется как совокупность линейных форм $c(\xi)$, удовлетворяющих условию}
\[
 c(\xi)\cdot\bb\zero{\Amod}.
\]
\emph{$\mathfrak C$ снова представляет собой модуль из линейных форм, а именно, как только $\bb$ отличен от нулевого идеала, модуль того же ранга, что и $\Amod$.}

Здесь нужно лишь заметить, что из определения ясно: частное каждый раз обладает модульным свойством; системы целых величин с модульным свойством суть идеалы.

Понятие частного приводит к связи между основным модулем и высшим элементарным делителем:

\medskip
\noindent\textbf{Теорема II.} \emph{Между основным модулем и высшим элементарным делителем модуля $\Amod$ существует взаимное соотношение}
\begin{equation}
 (e_\rho)=\Amod/\Gmod;\qquad \Gmod=\Amod/(e_\rho),
\tag{5}
\end{equation}
\emph{где $(e_\rho)$ означает главный идеал, порожденный $e_\rho$.}

Действительно, поскольку каждый элементарный делитель делит $e_\rho$, из (1) и (3) получаем
\begin{equation}
 e_\rho\Gmod\zero{\Amod}\quad \text{и следовательно}\quad (e_\rho)\cdot\Gmod\zero{\Amod};
\tag{6}
\end{equation}
значит, $(e_\rho)$ делится на частное $\Amod/\Gmod$. Но верно и обратное; ведь из
\[
 b\Gmod\zero{\Amod}
 \quad\text{следует}\quad b\eta_\rho\zero{\Amod}
 \quad\text{и потому}\quad b\zero{(e_\rho)},
\]
чем доказана первая формула (5)\footnote{Если положить $\Amod_0=\Amod,\ldots,\Amod_\lambda=(\Amod,\eta_\rho,\ldots,\eta_{\rho-\lambda+1})$, где каждый $\Amod_i$ имеет высший элементарный делитель $e_{\rho-i}$, то теми же рассуждениями получается:
\[
 (e_\rho)=\Amod_0/\Amod_1,\ldots,(e_{\rho-\lambda})=\Amod_\lambda/\Amod_{\lambda+1},\ldots,(e_1)=\Amod_{\rho-1}/\Amod_\rho.
\]
Более общо, пусть положено $\zeta_\lambda=b_{\lambda1}\eta_1+\cdots+b_{\lambda\lambda}\eta_\lambda$, пусть $(b_{\lambda\lambda},e_\lambda)=1$, и пусть
\[
 \Bmod_0=\Amod,\ldots,\Bmod_\lambda=(\Amod,\zeta_\rho,\ldots,\zeta_{\rho-\lambda+1}).
\]
Тогда и $\Bmod_i$ имеет высший элементарный делитель $e_{\rho-\lambda}$, и поэтому
\[
 (e_{\rho-\lambda})=\Bmod_\lambda/\Bmod_{\lambda+1}.
\]}. Формула (6) далее показывает, что $\Gmod$ делится на частное $\Amod/(e_\rho)$; это частное является делителем $\Amod$ того же ранга, значит, по определению основного модуля $\Amod$, наоборот, делится на $\Gmod$. Тем самым доказана и вторая формула (5).

Пусть $d$ обозначает какой-нибудь не тождественно нулевой детерминант порядка $\rho$ из $\Amod$. Тогда $d$ делится на $\rho$-й детерминантный делитель $d_\rho$, а следовательно и на $e_\rho$; значит, $d\Gmod\zero{\Amod}$, и по приведенному выше выводу получается $\Gmod=\Amod/(d)$. Для дальнейшего, однако, существенно, что это последнее представление $\Gmod$ как частного можно доказать и без возврата к элементарным делителям, а потому оно имеет более общую силу:

\medskip
\noindent\textbf{Теорема III.} \emph{Если под целыми величинами понимать полиномы от нескольких неизвестных}\footnote{На самом деле используется лишь то, что целые величины воспроизводятся сложением, вычитанием и умножением при сохранении обычных правил вычисления, то есть образуют кольцо; и далее то, что в этом кольце произведение обращается в нуль только тогда, когда обращается в нуль один из множителей, так что кольцо можно расширить до поля образованием частных (присоединением пар элементов). Напротив, базисное представление модуля не используется; следовательно, (7) справедлива, например, и для области всех целых алгебраических чисел в качестве коэффициентов. (E. N.)}, \emph{то сохраняется представление частным}
\begin{equation}
 \Gmod=\Amod/(d),
\tag{7}
\end{equation}
\emph{где под $d$ понимается какой-нибудь не тождественно нулевой детерминант порядка $\rho$ из $\Amod$.}

Прежде всего ясно, что понятия ранга, основного модуля и модульного частного -- которое теперь уже не обязательно является главным идеалом -- сохраняются при этом расширении; сохраняется и Теорема I, за исключением базисного представления.

Пусть теперь
\[
 a_1(\xi)=a_{11}\xi_1+\cdots+a_{1k}\xi_k,\ldots,
 a_\rho(\xi)=a_{\rho1}\xi_1+\cdots+a_{\rho k}\xi_k
\]
суть $\rho$ линейно независимых линейных форм из $\Amod$, которые вообще не будут образовывать модульный базис; обозначим неизвестные $\xi$ так, что
\[
 d=|a_{ij}|\ne0;\qquad i,j=1,2,\ldots,\rho.
\]
Отсюда следует, что $\rho$ линейных форм специального вида принадлежат $\Amod$:
\begin{equation}
 d\xi_\lambda+b_{\lambda,\rho+1}\xi_{\rho+1}+\cdots+b_{\lambda,k}\xi_k\zero{\Amod}
 \quad(\lambda=1,2,\ldots,\rho).
\tag{8}
\end{equation}
Но $\Amod$ не может содержать линейную форму $c_{\rho+1}\xi_{\rho+1}+\cdots+c_k\xi_k$, зависящую только от $\xi_{\rho+1},\ldots,\xi_k$, поскольку иначе по меньшей мере один детерминант порядка $\rho+1$, а именно $d\cdot c_\alpha$, был бы отличен от нуля.

Теперь частное $\Amod/(d)$, как делитель $\Amod$ того же ранга, делится на $\Gmod$; но верно и обратное. Действительно, для каждого $g(\xi)\zero{\Gmod}$ по определению имеем
\[
 c\cdot g(\xi)\zero{\Amod};\quad c\ne0
 \quad\text{и следовательно}\quad d\cdot c\cdot g(\xi)\zero{\Amod}.
\]
По (8), однако,
\[
 d\cdot g(\xi)=g_{\rho+1}\xi_{\rho+1}+\cdots+g_k\xi_k\zero{\Amod},
\]
следовательно,
\[
 c g_{\rho+1}\xi_{\rho+1}+\cdots+c g_k\xi_k\zero{\Amod}.
\]
Из только что замеченного следует $c\cdot g_{\rho+1}=0,\ldots,c\cdot g_k=0$, а потому, ввиду $c\ne0$, также $g_{\rho+1}=0,\ldots,g_k=0$; значит, $d\cdot g(\xi)\zero{\Amod}$, чем доказана (7).

Заметим, что $d\cdot g(\xi)\zero{\Amod}$ следует и непосредственно из того, что оба модуля $(a_1(\xi),\ldots,a_\rho(\xi))$ и $(g(\xi),a_1(\xi),\ldots,a_\rho(\xi))$ имеют тот же ранг $\rho$; значит, если положить $g(\xi)=c_1\xi_1+\cdots+c_k\xi_k$, то детерминант порядка $\rho+1$
\[
\begin{vmatrix}
 a_1(\xi)&a_{11}&\cdots&a_{1\rho}\\
 \vdots&\vdots&&\vdots\\
 a_\rho(\xi)&a_{\rho1}&\cdots&a_{\rho\rho}\\
 g(\xi)&c_1&\cdots&c_\rho
\end{vmatrix}
\]
обращается в нуль. Но приведенное выше доказательство снова встретится в §~4 после формулы (30).


\section*{§ 2. Теоремы о разложении норм и модулей.}

Теперь снова речь идет о модулях линейных форм относительно целых рациональных чисел или полиномов от одной неизвестной с коэффициентами из $P$.

\medskip
\noindent\textbf{Определение III.} \emph{Если, как обычно, все линейные формы от $\xi$, конгруэнтные между собой по модулю $\Amod$, объединить в один класс вычетов, то символ $\Gmod\mid\Amod$ будет обозначать систему классов вычетов модуля $\Gmod$ по модулю $\Amod$, то есть систему всех тех классов вычетов по $\Amod$, которые состоят из элементов $\Gmod$. Норма $\Gmod$ относительно $\Amod$ -- в обозначении $\Norm(\Gmod\mid\Amod)$ -- определяется как детерминант переходного преобразования от $\Gmod$ к $\Amod$, то есть как детерминант преобразования, выражающего какой-нибудь линейно независимый модульный базис $\Amod$ через такой же базис основного модуля $\Gmod$ модуля $\Amod$.}

Из (1) и (3), соответственно из замечания к Теореме II, получаем:
\begin{equation}
 \Norm(\Gmod\mid\Amod)=e_1e_2\cdots e_\rho=d_\rho;
 \qquad
 \Gmod=\Amod/\bigl(\Norm(\Gmod\mid\Amod)\bigr).
\tag{9}
\end{equation}

Из (1), (3) и (9) обычным образом следует, что в случае целых рациональных чисел $\Norm(\Gmod\mid\Amod)$ представляет число классов вычетов $\Gmod$ по $\Amod$, тогда как в случае полиномов от одной неизвестной -- где число классов вычетов вообще становится бесконечным -- степень $\Norm(\Gmod\mid\Amod)$ равна \emph{числу линейно независимых относительно $P$ классов вычетов}. Если $\Bmod$ является делителем $\Amod$ того же ранга, то вместе с представителем любого такого класса вычетов $\Bmod$ содержит и все элементы этого класса; значит, $\Bmod$ получается из $\Amod$ добавлением конечного числа классов вычетов, соответственно их линейных комбинаций. Отсюда сразу следует:

\medskip
\noindent\textbf{Теорема IV.} \emph{Если $\Bmod$ является делителем $\Amod$ того же ранга, так что основные модули совпадают, и если, кроме того, $\Norm(\Gmod\mid\Bmod)=\Norm(\Gmod\mid\Amod)$, то также $\Bmod=\Amod$.}

О связи между разложением норм и модулей справедливо следующее утверждение.

\medskip
\noindent\textbf{Теорема V.} \emph{Пусть}
\begin{equation}
 \Norm(\Gmod\mid\Amod)=r\cdot s,\qquad (r,s)=1;
\tag{10}
\end{equation}
\emph{тогда существует один и только один модуль $\Rmod$, а также один и только один модуль $\Smod$, оба делители $\Amod$ того же ранга, такие, что}
\[
 r=\Norm(\Gmod\mid\Rmod),\qquad s=\Norm(\Gmod\mid\Smod)
\]
\emph{выполняется. $\Rmod$ и $\Smod$ задаются как}
\[
 \Rmod=\Amod/(s);\qquad \Smod=\Amod/(r),
\]
\emph{и $\Amod$ равен наименьшему общему кратному $\Rmod$ и $\Smod$.}\footnote{Наименьшее общее кратное $[\Rmod,\Smod]$ здесь понимается в модульном смысле: как совокупность линейных форм, принадлежащих и $\Rmod$, и $\Smod$. Соответственно, наибольший общий делитель $(\Rmod,\Smod)$ в модульном смысле понимается как совокупность линейных форм, которые можно представить в виде суммы элемента из $\Rmod$ и элемента из $\Smod$.}
\emph{Если положить $r_\rho=(r,e_\rho)$, $s_\rho=(s,e_\rho)$, то имеют место соотношения взаимности}
\begin{equation}
 (s_\rho)=\Amod/\Rmod;\quad \Rmod=\Amod/(s_\rho);
 \qquad
 (r_\rho)=\Amod/\Smod;\quad \Smod=\Amod/(r_\rho).
\tag{11}
\end{equation}

Действительно, если вообще положить $r_i=(r,e_i)$, $s_i=(s,e_i)$, то из (9) и (10) получаем:
\[
 (r_i,s_i)=1;\qquad e_i=r_is_i;\qquad r=r_1\cdots r_\rho;\qquad s=s_1\cdots s_\rho,
\]
и каждый $r_i$, соответственно $s_i$, делит следующий $r_{i+1}$, соответственно $s_{i+1}$. Итак, если положить
\[
 \Rmod=(r_1\eta_1,\ldots,r_\rho\eta_\rho);\qquad
 \Smod=(s_1\eta_1,\ldots,s_\rho\eta_\rho),
\]
то $\Rmod$ и $\Smod$ являются делителями $\Amod$ того же ранга; вследствие взаимной простоты $r_i$ и $s_i$ модуль $\Amod$ равен наименьшему общему кратному $\Rmod$ и $\Smod$, и, кроме того,
\[
 \Norm(\Gmod\mid\Rmod)=r_1\cdots r_\rho=r;\qquad
 \Norm(\Gmod\mid\Smod)=s_1\cdots s_\rho=s.
\]
Далее имеем
\[
 s_\rho\Rmod\zero{\Amod};\qquad r_\rho\Smod\zero{\Amod};
\]
из $c\Rmod\zero{\Amod}$ следует $cr_\rho\eta_\rho\zero{\Amod}$, значит $c\zero{(s_\rho)}$, чем доказано $(s_\rho)=\Amod/\Rmod$ и, соответственно, $(r_\rho)=\Amod/\Smod$. Из
\[
 b(\eta)=b_1\eta_1+\cdots+b_\rho\eta_\rho\zero{\Amod/(s_\rho)}
\]
далее следует, поскольку все $r_i$ взаимно просты с $s_\rho$, что $b_i\zero{(r_i)}$; значит, $\Rmod=\Amod/(s_\rho)$ и соответственно $\Smod=\Amod/(r_\rho)$. Тем самым доказаны соотношения взаимности (11), которые в специальном случае $r=1$ переходят в (5).

Остается доказать \emph{единственность} $\Rmod$ и $\Smod$. Пусть $\overline{\Rmod}$ и $\overline{\Smod}$ -- модули, удовлетворяющие условиям Теоремы V, а $\bar r_i$ и $\bar s_i$ -- их элементарные делители. Так как $\bar r_i$ и $\bar s_i$ являются делителями $e_i$, то из
\[
 r=\bar r_1\cdots\bar r_\rho,
 \qquad
 s=\bar s_1\cdots\bar s_\rho,
 \qquad
 (\bar r_i,\bar s_i)=1
\]
следует также
\[
 e_i=\bar r_i\bar s_i,
 \qquad
 \bar r_i=(r,e_i)=r_i,
 \qquad
 \bar s_i=(s,e_i)=s_i.
\]
Следовательно, $\overline{\Rmod}$ и $\overline{\Smod}$ совпадают с $\Rmod$ и $\Smod$ по элементарным делителям и основному модулю. Поэтому, если с помощью унимодулярного преобразования $\eta=Q(\xi)$ ввести новые неизвестные $\xi$, то получаем
\[
 \overline{\Rmod}=(r_1\xi_1,\ldots,r_\rho\xi_\rho);
\]
\[
 \Amod=\bigl(r_1s_1(q_{11}\xi_1+\cdots+q_{1\rho}\xi_\rho),\ldots,
 r_\rho s_\rho(q_{\rho1}\xi_1+\cdots+q_{\rho\rho}\xi_\rho)\bigr).
\]
Из $\Amod\zero{\Rmod}$ отсюда получаем
\[
 r_i s_i(q_{i1}\xi_1+\cdots+q_{i\rho}\xi_\rho)\zero{\Rmod};
\]
значит, $r_is_iq_{ij}\zero{(r_i)}$. Поскольку $(r,s)=1$, это дает $r_iq_{ij}\zero{(r_i)}$, то есть $\overline{\Rmod}\zero{\Rmod}$. Но поскольку $\Norm(\Gmod\mid\Rmod)=\Norm(\Gmod\mid\overline{\Rmod})$, Теорема IV дает $\Rmod=\overline{\Rmod}$ и соответственно $\Smod=\overline{\Smod}$. Тем самым Теорема V доказана во всех частях.


\section*{§ 3. Идеалы из полиномов.}

Пусть $\overline{\mideal}$ есть \emph{идеал из полиномов} от $n$ неопределенных $y_1,\ldots,y_n$ с коэффициентами из произвольного, абстрактно заданного поля $P$; то есть вместе с $\Phi_1(y)$ и $\Phi_2(y)$ идеал $\overline{\mideal}$ содержит также разность $\Phi_1(y)-\Phi_2(y)$, а вместе с $\Phi(y)$ содержит также $A(y)\cdot\Phi(y)$, где под $A(y)$ понимается произвольный полином с коэффициентами из $P$\footnote{Само понятие приводит здесь к названию "идеал" вместо ранее общепринятых названий "модуль" или "модуль форм"; уже здесь это оказывается необходимым, чтобы отличать такие объекты от модулей из линейных форм. (E. N.)}.

Пусть $y$ подвергнуты преобразованию с неопределенными в качестве коэффициентов,
\begin{equation}
 y=U(x);\quad \text{например}\quad
 \begin{cases}
 y_1=x_1,\\
 y_2=u_{21}x_1+x_2,\\
 \cdots\\
 y_n=u_{n1}x_1+\cdots+u_{n,n-1}x_{n-1}+x_n,
 \end{cases}
\tag{12}
\end{equation}
и пусть неопределенные $u_{\mu\nu}$ присоединены к области коэффициентов полиномов, которая тем самым переходит в поле $P(u)$. При этом присоединении $\overline{\mideal}$ переходит в $\overline{\mideal}_{(u)}$; следовательно, $\overline{\mideal}_{(u)}$ содержит, кроме полиномов $\Phi_\lambda(y)$ из $\overline{\mideal}$, также все линейные комбинации $\sum \alpha_\lambda(u)\Phi_\lambda(y)$ конечного числа $\Phi_\lambda$, где под $\alpha_\lambda(u)$ понимаются рациональные функции от $u_{\mu\nu}$ с коэффициентами из $P$. Умножением на подходящий полином от $u$ функции $\alpha_\lambda(u)$ без ограничения общности можно считать степенными произведениями от $u$. Итак, имеем:
\begin{equation}
 \sum \alpha_\lambda(u)\Phi_\lambda(y)\zero{\overline{\mideal}_{(u)}};
 \qquad
 \Phi_\lambda(y)\zero{\overline{\mideal}}\quad\text{и обратно}.
\tag{13}
\end{equation}
Посредством (12) $\overline{\mideal}_{(u)}$ переходит в идеал $\mideal$ из полиномов от $x_1,\ldots,x_n$ с коэффициентами из $P(u)$:
\begin{equation}
 \overline{\mideal}_{(u)}=\mideal\quad \text{посредством}\quad y=U(x).
\tag{14}
\end{equation}
Совмещение соотношений (14) и (13) означает следующее:
\begin{equation}
\begin{aligned}
 &\text{Из } F(x)\zero{\mideal};\\
 &F(x)=\Phi(y)=\sum \alpha_\lambda(u)\Phi_\lambda(y)\\
 &\hspace{2em}=\sum \alpha_\lambda(u)F_\lambda(x)
 \quad\text{следует } F_\lambda(x)\zero{\mideal}.
\end{aligned}
\tag{15}
\end{equation}
а из (14) и (15) в обратном направлении снова получается (13).

\medskip
\noindent\textbf{Определение IV.} \emph{Идеалы $\mideal$ из полиномов от $x_1,\ldots,x_n$ с коэффициентами из $P(u)$, удовлетворяющие соотношениям (15), то есть возникающие из $\overline{\mideal}_{(u)}$ посредством $y=U(x)$, будем называть преобразованными идеалами.}

Преобразованные идеалы выделяются существованием регулярных полиномов; \emph{полином степени $r$ называется регулярным относительно $x_i$, если он содержит член $x_i^r$ с ненулевым коэффициентом}. Видно, что полиномы $F_i(x)$ регулярны относительно $x_1$. Далее существенно будет рассматривать эти преобразованные идеалы как модули из линейных форм по некоторым степенным произведениям от $x$. Для этого надо зафиксировать понятие основного идеала, которое позднее перейдет в основной модуль. Сначала введем обозначение, которого далее будем всюду придерживаться:

\medskip
\noindent\textbf{Обозначение.} \emph{Под $F^{(i)},f^{(i)},a^{(i)},\ldots$ далее всюду понимаются полиномы от $x$, свободные от $x_1,\ldots,x_{i-1}$.}

\medskip
\noindent\textbf{Определение V.} \emph{Для каждого идеала $\mideal$ определены $n$ основных идеалов $\gideal_0,\gideal_1,\ldots,\gideal_{n-1}$ следующим условием: основной идеал $(i-1)$-й ступени $\gideal_{i-1}$ содержит все и только те полиномы $G(x)$, для которых существует полином $b^{(i)}$ -- вообще говоря, меняющийся вместе с $G(x)$, -- такой, что}
\begin{equation}
 b^{(i)}G(x)\zero{\mideal};\qquad b^{(i)}\ne0.
\tag{16}
\end{equation}

То, что $\gideal$ действительно являются идеалами, точно соответствует доказательству модульного свойства для $\Gmod$ после формулы (4), где (16) служит аналогом. Идеал $\gideal_i$ делится на $\gideal_{i-1}$, поскольку каждый полином $b^{(i+1)}$ одновременно является полиномом $b^{(i)}$.

Для основных идеалов верно следующее:

\medskip
\noindent\textbf{Теорема VI.} \emph{Основные идеалы преобразованных идеалов сами являются преобразованными идеалами}\footnote{Теорема VI используется только в §~7.}.

Доказательство опирается на известную теорему Дедекинда--Мертенса\footnote{Dedekind: \emph{Über einen arithmetischen Satz von Gauß}, \emph{Mitt. d. deutsch. math. Ges. zu Prag} 1892. F. Mertens: \emph{Über einen algebraischen Satz}, \emph{Ber. d. Ak. d. Wissensch. Wien} 101 (1892), с. 1560--1566. -- Расширение Кронекером теоремы Гаусса на полиномы с неопределенными коэффициентами является непосредственным следствием этой теоремы.}: пусть $a$ и $b$ являются неопределенными и положено
\[
 \sum a_{i_1,\ldots,i_s}z_1^{i_1}\cdots z_s^{i_s}\cdot
 \sum b_{j_1,\ldots,j_s}z_1^{j_1}\cdots z_s^{j_s}
 =
 \sum c_{k_1,\ldots,k_s}z_1^{k_1}\cdots z_s^{k_s},
\]
\[
 \sum i\le \nu,
 \qquad
 \sum j\le \mu,
 \qquad
 \sum k\le \nu+\mu;
\]
далее пусть $\Amod,\Bmod,\mathfrak C$ обозначают модули, соответственно выведенные из $a,b,c$ при помощи целых рациональных чисел. Тогда существует показатель $q$ такой, что имеет место тождество
\begin{equation}
 \Amod^q\Bmod=\Amod^{q-1}\mathfrak C.
\tag{17}
\end{equation}
Здесь $\Amod^q$ состоит из степенных произведений $q$-й степени от $a$ и их целочисленных линейных комбинаций.

Для доказательства Теоремы VI положим теперь:
\begin{equation}
\begin{cases}
 G(x)=\Gamma(y)=\sum \alpha_\lambda(u)\Gamma_\lambda(y)=\sum \alpha_\lambda(u)G_\lambda(x),\\
 b^{(i)}(x)=\varphi(y)=\sum \beta_\mu(u)\varphi_\mu(y),
\end{cases}
\tag{18}
\end{equation}
где $\alpha_\lambda(u)$ и $\beta_\mu(u)$ обозначают степенные произведения от $u$. Тогда по (16), (14) и (18) имеем:
\begin{equation}
 \sum\beta_\mu(u)\varphi_\mu(y)\cdot
 \sum\alpha_\lambda(u)\Gamma_\lambda(y)\zero{\overline{\mideal}_{(u)}};
 \qquad
 \sum\beta_\mu(u)\varphi_\mu(y)\ne0.
\tag{19}
\end{equation}
Слева стоит произведение двух полиномов от $u$, коэффициентами которых являются полиномы $\varphi_\mu(y)$ и $\Gamma_\lambda(y)$; коэффициенты полинома-произведения от $u$ по правой части (19) являются полиномами из $\overline{\mideal}$. Если, следовательно, в (17) заменить $a,b,c$ этими полиномами от $y$, то получаем
\[
 \left\{\sum\beta_\mu(u)\varphi_\mu(y)\right\}^q\cdot \Gamma_\lambda(y)\zero{\overline{\mideal}_{(u)}},
\]
что по (18) равносильно:
\begin{equation}
 b^{(i)}(x)^q\cdot G_\lambda(x)\zero{\mideal};
 \qquad
 G_\lambda(x)\zero{\gideal_{i-1}}.
\tag{20}
\end{equation}
Соотношения (16), (18) и (20) показывают, что \emph{основные идеалы $\gideal_{i-1}$ действительно являются преобразованными идеалами.}

\section*{22. Обработка K. Hentzelt:\ к теории полиномиальных идеалов и результантов\quad (продолжение)}
\begin{center}
Math. Ann. 88 (1923), с. 53--79
\end{center}

\section*{§ 4. Связь между идеалами из полиномов и модулями из линейных форм.}

Чтобы идеал $\mideal$ из полиномов, который далее всегда предполагается преобразованным, рассматривать как модуль $\Mmod_{i-1}$ $(i=1,\ldots,n)$ из линейных форм, отдельные полиномы $F(x)$ надо рассматривать как линейные формы от степенных произведений $\xi_\lambda$ переменных $x_1\ldots x_{i-1}$:
\[
  F(x)=\sum a^{(i)}_\lambda\xi_\lambda,
\]
где, согласно обозначениям §~3, под $a^{(i)}_\lambda$ понимаются полиномы от $x_i\ldots x_n$. Из определения идеала сразу следует, что $\Mmod_{i-1}$ обладает модульным свойством относительно этих целых величин $a^{(i)},b^{(i)},\ldots$; а так как бесконечно многие степенные произведения $\xi$ от $x_1\ldots x_{i-1}$ входят только \emph{линейно}, то в силу своей линейной независимости они играют для $\Mmod_{i-1}$ роль неопределенных. Каждый модуль $\Mmod_{i-1}$ содержит бесконечно много неопределенных $\xi$, но в каждой отдельной линейной форме встречается лишь конечное число неопределенных. Точно так же каждый основной идеал $\gideal_{i-1}$ следует рассматривать как модуль линейных форм $\Gmod_{i-1}$, где неопределенными снова служат степенные произведения от $x_1\ldots x_{i-1}$. При $i=1$ речь идет только об \emph{одной} неопределенной $\xi_0$, соответствующей единице.

Теперь можно, и на этом основывается все дальнейшее, несмотря на эти бесконечно многие неопределенные $\xi$, ограничиться модулями линейных форм от конечного числа неопределенных, согласно следующей теореме.

\medskip
\noindent\textbf{Теорема VII.} \emph{Система классов вычетов $\Gmod_{i-1}\mid \Mmod_{i-1}$ изоморфна системе классов вычетов $\Gmod^*_{i-1}\mid \Mmod^*_{i-1}$, где $\Mmod^*_{i-1}$ есть модуль от конечного числа неопределенных, а $\Gmod^*_{i-1}$ означает его основной модуль. Модуль $\Mmod_{i-1}$ после присоединения $x_{i+1}\ldots x_n$ к $P(u)$ имеет лишь конечное число элементарных делителей, отличных от $1$, а именно те элементарные делители $\Mmod^*_{i-1}$, которые отличны от $1$.}

При этом элементарные делители $\Mmod_{i-1}$ надо определять так, как в примечании 8). Изоморфизм, входящий в Теорему VII, основывается на следующем определении.

\medskip
\noindent\textbf{Определение V.} \emph{Две системы классов вычетов называются изоморфными, если их можно поставить во взаимно однозначное соответствие так, что разности двух классов соответствует разность соответствующих классов, а произведению класса на целую величину соответствует произведение соответствующего класса на ту же целую величину.}

Доказательство Теоремы VII получается полной индукцией. Для $i=1$ Теорема VII очевидна: $\Mmod_0$ есть модуль линейных форм от одной неопределенной $\xi_0$, соответствующей степенному произведению нулевой степени от $x$, то есть единице; целыми величинами являются все полиномы от $x_1\ldots x_n$. Основной идеал $\gideal_0$ равен единичному идеалу, состоящему из всех полиномов от $x_1\ldots x_n$, следовательно $\Gmod_0$ равен модулю $(\xi_0)$, основному модулю $\Mmod_0$, так что здесь $\Gmod^*_0\mid\Mmod^*_0$ совпадает с $\Gmod_0\mid\Mmod_0$. Наконец, поскольку $\Mmod_0$ имеет ранг $1$, он может иметь не более одного элементарного делителя, отличного от $1$.

Сначала перейдем от $i=1$ к $i=2$, чтобы с помощью полученного здесь понимания строения $\Gmod_1\mid\Mmod_1$ провести общий индукционный шаг. Для этого прежде всего покажем, что для $\Gmod_0\mid\Mmod_0$ можно взять систему представителей, ограниченную по $x_1$; вследствие делимости $\gideal_1$ на $\gideal_0$ это затем приведет к конечному числу неопределенных модуля $\Gmod^*_1$.

По §~3 преобразованный идеал $\mideal$ имеет по крайней мере один полином $C^{(1)}(x)$, регулярный по $x_1$, скажем степени $k$; значит, $\Mmod_0$ содержит линейную форму $C^{(1)}\xi_0$. Благодаря регулярности $C^{(1)}(x)$ каждый полином $H(x)$ допускает представление
\begin{equation}
 H(x)=b^{(2)}_0+b^{(2)}_1x_1+\cdots+b^{(2)}_{k-1}x_1^{k-1}\quad (C^{(1)}(x));
\tag{21}
\end{equation}
поэтому, ввиду принадлежности $C^{(1)}\xi_0$ к $\Mmod_0$, для $\Gmod_0\mid\Mmod_0$ действительно можно выбрать систему представителей, которая не достигает $k$-й степени по $x_1$.

Если теперь, в частности, для полиномов $G(x)$ из $\gideal_1$ и $F(x)$ из $\mideal$ положить:
\begin{equation}
 \left\{\begin{aligned}
  G(x)&=g^{(2)}_0+g^{(2)}_1x_1+\cdots+g^{(2)}_{k-1}x_1^{k-1} &&(C^{(1)}(x)),\\
  F(x)&=a^{(2)}_0+a^{(2)}_1x_1+\cdots+a^{(2)}_{k-1}x_1^{k-1} &&(C^{(1)}(x)),
 \end{aligned}\right.
\tag{22}
\end{equation}
а через $\Gmod^*_1$ обозначить систему всех линейных форм $g(\xi)=g^{(2)}_0\xi_0+\cdots+g^{(2)}_{k-1}\xi_{k-1}$, и соответственно через $\Mmod^*_1$ систему линейных форм $a(\xi)=a^{(2)}_0\xi_0+\cdots+a^{(2)}_{k-1}\xi_{k-1}$, то из свойства идеала для $\gideal_1$ и $\mideal$ следует, что $\Gmod^*_1$ и $\Mmod^*_1$ являются модулями из линейных форм от $\xi_0\ldots\xi_{k-1}$ относительно целых величин $b^{(2)},c^{(2)},\ldots$. Точно так же главный идеал, выведенный из $C^{(1)}(x)$, дает модуль относительно этих целых величин:
\[
 \Gmod_1=(\xi_0,\xi_1,\ldots,\xi_r,\ldots),
\]
где положено $\xi_r=x_1^rC^{(1)}(x)$. Элементы $\xi_r$ линейно независимы относительно этих целых величин $b^{(2)},c^{(2)},\ldots$ как между собой, так и от конечного числа $\xi$. Из делимости $\Gmod_1$ на $\Mmod_1$ и, следовательно, на $\Gmod_1$ следует, что также $\Gmod^*_1$ делится на $\Gmod_1$, соответственно $\Mmod^*_1$ делится на $\Mmod_1$. С учетом (22) получаем:
\begin{equation}
 \Gmod_1=(\Gmod^*_1,\Gmod_1),\qquad
 \Mmod_1=(\Mmod^*_1,\Gmod_1).
\tag{23}
\end{equation}

Из (23) и линейной независимости $\xi$ и $\xi$ Теорема VII для $i=2$ следует так. Благодаря делимости $\Gmod_1$ на $\Mmod_1$ прежде всего имеем $\Gmod_1\mid\Mmod_1=\Gmod^*_1\mid\Mmod_1$. Чтобы доказать изоморфизм с $\Gmod^*_1\mid\Mmod^*_1$, пусть некоторые линейные формы $g^*(\xi)$ из $\Gmod^*_1$ задают систему представителей $\Gmod^*_1\mid\Mmod^*_1$. Из линейной независимости $\xi$ и $\xi$ из $g^*(\xi)\equiv 0(\Mmod_1)$ следует также $g^*(\xi)\equiv 0(\Mmod^*_1)$; значит, формы $g^*(\xi)$ остаются неконгруентными и по модулю $\Mmod_1$, образуют систему представителей $\Gmod_1\mid\Mmod_1$, а соответствие от $\Gmod_1\mid\Mmod_1$ к $\Gmod^*_1\mid\Mmod^*_1$, задаваемое этой системой представителей, изоморфно в смысле Определения V.

Совершенно так же получаем: из $b^{(2)}g(\xi)\equiv0(\Mmod_1)$, $b^{(2)}\ne0$, следует $b^{(2)}g(\xi)\equiv0(\Mmod^*_1)$, $b^{(2)}\ne0$. Так как это выполнено для всех $g(\xi)$ из $\Gmod^*_1$ и только для них -- ведь по (23) $\Gmod^*_1$ состоит из всех полиномов основного идеала $\gideal_1$, которые одновременно являются линейными формами от $\xi$ -- то $\Gmod^*_1$ тем самым характеризуется как основной модуль $\Mmod^*_1$.

Наконец, присоединяя $x_3\ldots x_n$ к $P(u)$, получаем:
\begin{equation}
\left\{\begin{aligned}
 \Gmod^*_1&=(\eta_1,\ldots,\eta_\rho),&
 \Mmod^*_1&=(e_1\eta_1,\ldots,e_\rho\eta_\rho),\\
 \Gmod_1&=(\eta_\rho,\ldots,\eta_1,\xi_0,\ldots,\xi_r,\ldots),&
 \Mmod_1&=(e_\rho\eta_\rho,\ldots,e_1\eta_1,\xi_0,\ldots,\xi_r,\ldots),
\end{aligned}\right.
\tag{24}
\end{equation}
и значит:
\[
 (e_\rho)=\Mmod_1/\Gmod_1=\Mmod_1/(\Mmod_1,\eta_\rho),\qquad
 (e_{\rho-1})=(\Mmod_1,\eta_\rho)/\Gmod_1\quad\hbox{и т. д.}
\]
Тем самым, учитывая примечание 8), $e_\rho,e_{\rho-1},\ldots,e_1,1,\ldots,1,\ldots$ распознаются как элементарные делители $\Mmod_1$. Итак, для $i=2$ Теорема VII доказана во всех частях.

Заметим, что (22) и (23) используют только регулярность $C^{(1)}(x)$ по $x_1$. Поэтому эти формулы и связанные с ними рассуждения, ведущие к Теореме VII, сохраняются, если вместо (12) взять в основу специальное преобразование
\begin{equation}
 y_1=x_1;\quad y_2=u_{21}x_1+z_2;\quad\ldots;\quad y_n=u_{n1}x_1+z_n
\tag{25}
\end{equation}
которое при составлении с
\begin{equation}
\left\{\begin{array}{l}
 z=U_1(x)\quad\hbox{или}\\[2mm]
 z_2=x_2;\quad z_3=u_{32}x_2+x_3;\quad\ldots;\quad z_n=u_{n2}x_2+\cdots+u_{n,n-1}x_{n-1}+x_n
\end{array}\right.
\tag{26}
\end{equation}
снова дает (12). Если $\widetilde{\mideal}_{(u)}$ посредством (25) переходит в $\widetilde{\mideal}$, а $\widetilde{\gideal}_1,\widetilde{\Gmod}_1,\ldots$ имеют для $\widetilde{\mideal}$ такое же значение, как $\gideal_1,\Gmod_1,\ldots$ для $\mideal$, то имеем:
\begin{equation}
 \widetilde{\Gmod}_1=(\widetilde{\Gmod}^{*}_1,\widetilde{\Gmod}_1),\qquad
 \widetilde{\Mmod}_1=(\widetilde{\Mmod}^{*}_1,\widetilde{\Gmod}_1).
\tag{27}
\end{equation}

При этом (27) возникает из (23) просто посредством обращения (26). Это очевидно для $\mideal$ и $C^{(1)}(x)$, значит также для $\Gmod_1$ и $\Mmod_1$. Но это верно и для $\gideal_1$, ибо
\[
 b^{(2)}(x)G(x)\equiv0(\mideal),\quad b^{(2)}\ne0,
 \qquad
 \widetilde b^{(2)}(z)\widetilde G(x,z)\equiv0(\widetilde\mideal),\quad \widetilde b^{(2)}\ne0
\]
взаимно обусловливают друг друга через $z=U_1(x)$. Поэтому $\widetilde\gideal_1=\gideal_1$ посредством (26); следовательно, также $\widetilde\Gmod_1=\Gmod_1$, и по определению $\Gmod_1^*$ и $\Mmod_1^*$ через (22) то же самое верно для $\widetilde\Gmod_1^*$ и $\widetilde\Mmod_1^*$, а значит для всего представления (27).

Для общего индукционного шага пусть выполнены предположения:
\begin{equation}
\left\{\begin{aligned}
 \Gmod_{i-1}&=(\Gmod_{i-1}^*,\Gmod_{i-1}),&
 \Mmod_{i-1}&=(\Mmod_{i-1}^*,\Gmod_{i-1}),\\
 \Gmod_{i-1}&=(\xi_0,\xi_1,\ldots,\xi_r,\ldots),
\end{aligned}\right.
\tag{28}
\end{equation}
где $\Gmod^*_{i-1}$ и $\Mmod^*_{i-1}$ означают модули -- относительно целых величин $b^{(i)},c^{(i)},\ldots$ -- из линейных форм от конечного числа неопределенных $\xi$, некоторых степенных произведений от $x_1,\ldots,x_{i-1}$; и где $\xi$ линейно независимы как между собой, так и от $\xi$ относительно этих целых величин. Представление, соответствующее (28), и линейная независимость $\xi$ и $\xi$ должны выполняться также тогда, когда вместо (12) в основу берется специальное преобразование
\begin{equation}
\begin{aligned}
 y_1&=x_1;\quad \ldots;\quad y_{i-1}=u_{i-1,1}x_1+\cdots+x_{i-1};\\
 y_i&=u_{i,1}x_1+\cdots+u_{i,i-1}x_{i-1}+z_i;\quad\ldots;\quad
 y_n=u_{n,1}x_1+\cdots+u_{n,i-1}x_{i-1}+z_n
\end{aligned}
\tag{29}
\end{equation}
которое составляется в (12) посредством
\[
 z=U_{i-1}(x)\quad\hbox{или}\quad
 z_i=x_i;\quad z_{i+1}=u_{i+1,i}x_i+x_{i+1};\quad\ldots;\quad
 z_n=u_{n,i}x_i+\cdots+u_{n,n-1}x_{n-1}+x_n,
\]
и притом это специальное представление должно возникать из (28) просто посредством обращения $z=U_{i-1}(x)$.

Из предположений (28) и линейной независимости $\xi$ и $\xi$ изоморфизм $\Gmod_{i-1}\mid\Mmod_{i-1}$ с $\Gmod^*_{i-1}\mid\Mmod^*_{i-1}$ следует через соответствие той же системе представителей, точно теми же рассуждениями, которые выше опирались на (23). Точно так же получается, что $\Gmod^*_{i-1}$ становится основным модулем $\Mmod^*_{i-1}$ и что $\Mmod_{i-1}$ имеет лишь конечное число элементарных делителей, отличных от $1$, а именно те элементарные делители $\Mmod^*_{i-1}$, которые отличны от $1$.

Для выполнения индукционного шага прежде всего опять надо показать, что для $\Gmod^*_{i-1}\mid\Mmod^*_{i-1}$, а значит и для $\Gmod_{i-1}\mid\Mmod_{i-1}$, можно взять систему представителей, ограниченную по $x_i$. Это следует из представления частного, доказанного в (7), Теорема III:
\begin{equation}
 \Gmod^*_{i-1}=\Mmod^*_{i-1}/(C^{(i)}),
\tag{30}
\end{equation}
где под $C^{(i)}$ понимается любой не тождественно нулевой определитель $\rho$-го порядка из $\Mmod^*_{i-1}$, а $\rho$ означает ранг $\Mmod^*_{i-1}$. Из той части предположений, которая относится к специальному преобразованию (29), следует, что $C^{(i)}$ всегда можно считать регулярным по $x_i$. Действительно, модуль $\widetilde\Mmod^*_{i-1}$, соответствующий $\Mmod^*_{i-1}$, имеет тот же ранг; если $\widetilde C^{(i)}$ означает любой не тождественно нулевой определитель $\rho$-го порядка из $\widetilde\Mmod^*_{i-1}$, который через $z=U_{i-1}(x)$ переходит в регулярный $C^{(i)}$, то по предположению $C^{(i)}$ становится определителем $\rho$-го порядка из $\Mmod^*_{i-1}$.

Из существования $C^{(i)}$, как в (8), при подходящей нумерации $\xi$ следует принадлежность к $\Mmod^*_{i-1}$ таких $\rho$ линейных форм специального вида:
\[
 L_\lambda(\xi)=C^{(i)}\xi_\lambda+c^{(i)}_{\lambda,1}\xi_{\rho+1}+\cdots+c^{(i)}_{\lambda,s-\rho}\xi_s
 \qquad(\lambda=1\ldots\rho).
\]
Теперь каждую линейную форму от $\xi$ можно привести к виду:
\begin{equation}
\begin{aligned}
 H(\xi)&=a^{(i)}_1\xi_1+\cdots+a^{(i)}_\rho\xi_\rho\\
 &\quad+b^{(i)}_1\xi_{\rho+1}+\cdots+b^{(i)}_{s-\rho}\xi_s
 \quad (L_1(\xi),\ldots,L_\rho(\xi)),
\end{aligned}
\tag{31}
\end{equation}
где $a^{(i)}_1\ldots a^{(i)}_\rho$ ограничены по $x_i$ и не достигают степени $k$ полинома $C^{(i)}$. Это представление единственно: действительно, из
\[
 a^{(i)}_1\xi_1+\cdots+a^{(i)}_\rho\xi_\rho+b^{(i)}_1\xi_{\rho+1}+\cdots+b^{(i)}_{s-\rho}\xi_s\equiv0\quad (L_1(\xi),\ldots,L_\rho(\xi))
\]
сравнение коэффициентов при $\xi_1\ldots\xi_\rho$ с учетом степеней по $x_i$ дает тождественное исчезновение всех $a^{(i)}$ и $b^{(i)}$.

Пусть теперь, в частности,
\[
 H(\xi)\equiv0(\Gmod^*_{i-1});\qquad\hbox{следовательно}\qquad C^{(i)}H(\xi)\equiv0(\Mmod^*_{i-1})\quad\hbox{по (30).}
\]
Отсюда, как в доказательстве Теоремы III, получаем
\[
 C^{(i)}H(\xi)-\sum_{\tau=1}^{s-\rho}\{C^{(i)}b^{(i)}_\tau-\sum_\lambda a^{(i)}_\lambda c^{(i)}_{\lambda,\tau}\}\xi_{\rho+\tau}
 \equiv0(\Mmod^*_{i-1}),
\]
и значит, поскольку, как в Теореме III, коэффициенты при $\xi_{\rho+1}\ldots\xi_s$ должны исчезать,
\[
 C^{(i)}b^{(i)}_\tau=\sum_\lambda a^{(i)}_\lambda c^{(i)}_{\lambda,\tau}\qquad(\tau=1\ldots s-\rho),
\]
откуда следует, что и $b^{(i)}_\tau$ ограничены и не достигают максимальной степени $c^{(i)}_{\lambda,\tau}$. Тем самым доказано существование системы представителей для $\Gmod^*_{i-1}\mid\Mmod^*_{i-1}$, то есть для $\Gmod_{i-1}\mid\Mmod_{i-1}$, ограниченной по $x_i$ и не достигающей некоторой фиксированной степени $r$.

Обозначим теперь через $\vartheta_1\ldots\vartheta_t$ конечное число степенных произведений
\[
 x_i^\alpha\xi_\lambda\quad(\alpha=0,1,\ldots,k-1;\ \lambda=1,2,\ldots,\rho),\qquad
 x_i^\beta\xi_{\rho+\tau}\quad(\beta=0,1,\ldots,r-1;\ \tau=1,2,\ldots,s-\rho)
\]
и расположим счетно бесконечное множество выражений
\[
\begin{aligned}
 &x_i^\gamma L_\lambda\quad(\gamma=0,1,\ldots\text{ до бесконечности};\ \lambda=1,2,\ldots,\rho),\\
 &x_i^\gamma\xi_\nu\quad(\gamma,\nu=0,1,\ldots\text{ до бесконечности}).
\end{aligned}
\]
в простую бесконечную последовательность $\omega_0,\omega_1,\ldots,\omega_\mu,\ldots$. Тогда $\vartheta$ и $\omega$ линейно независимы как каждые между собой, так и друг от друга, относительно целых величин $b^{(i+1)},c^{(i+1)},\ldots$. Действительно, из
\[
 \sum c^{(i+1)}_\alpha x_i^\alpha\xi_\lambda+
 \sum c^{(i+1)}_\beta x_i^\beta\xi_{\rho+\tau}+
 \sum d^{(i+1)}_\gamma x_i^\gamma L_\lambda(\xi)+
 \sum e^{(i+1)}_{\gamma\nu}x_i^\gamma\xi_\nu=0,
\]
где каждая сумма распространяется лишь на конечное число членов, в силу предполагаемой линейной независимости $\xi$ и $\xi$, а также $\xi$ между собой относительно целых величин $b^{(i)}$, следует исчезновение всех $e^{(i+1)}_{\gamma\nu}$. Из единственности представления (31) далее следует исчезновение $c^{(i+1)}_\alpha$ и $c^{(i+1)}_\beta$, а отсюда, наконец, в силу линейной независимости $L_\lambda(\xi)$ -- исчезновение $d^{(i+1)}_\gamma$.

Если теперь модуль $(\Gmod_{i-1},L_1(\xi),\ldots,L_\rho(\xi))$ рассматривать как модуль $\Gmod_i$ относительно целых величин $b^{(i+1)}$, то $\Gmod_i$ делится на $\Mmod_i$ благодаря делимости $\Gmod_{i-1},L_1(\xi)\ldots L_\rho(\xi)$ на $\Mmod_{i-1}$, и получается
\[
 \Gmod_i=(\omega_0,\omega_1,\ldots,\omega_\mu,\ldots).
\]
Для каждого $H(x)\equiv0(\gideal_{i-1})$ из (28), (31) и только что доказанного имеем:
\[
 H(x)=b^{(i+1)}_1\vartheta_1+\cdots+b^{(i+1)}_t\vartheta_t\quad(\Gmod_i)
\]
как модульное равенство относительно целых величин $b^{(i+1)},c^{(i+1)}$. Но $\gideal_i$ делится на $\gideal_{i-1}$, а $\mideal$ делится на $\gideal_i$; следовательно, если теперь, в частности, для каждого $G(x)$ из $\Gmod_i$ и каждого $F(x)$ из $\Mmod_i$ положить
\[
 G(x)=g^{(i+1)}_1\vartheta_1+\cdots+g^{(i+1)}_t\vartheta_t\quad(\Gmod_i),\qquad
 F(x)=a^{(i+1)}_1\vartheta_1+\cdots+a^{(i+1)}_t\vartheta_t\quad(\Gmod_i),
\]
и обозначить модуль, состоящий из всех $g(\vartheta)$, через $\Gmod_i^*$, а модуль, состоящий из всех $a(\vartheta)$, через $\Mmod_i^*$, то точно те же рассуждения, которые привели к (23), дают:
\begin{equation}
 \Gmod_i=(\Gmod_i^*,\Gmod_i),\qquad
 \Mmod_i=(\Mmod_i^*,\Gmod_i),\qquad
 \Gmod_i=(\omega_0,\omega_1,\ldots,\omega_\mu,\ldots).
\tag{32}
\end{equation}
При выводе (32) снова использовалась только регулярность $C^{(i)}$ по $x_i$, поэтому все рассуждение сохраняется, если взять в основу специальное преобразование (29) для индекса на единицу больше. Тогда те же рассуждения, которые опираются на (27), показывают, что это специальное представление возникает из (32) просто посредством обращения $z=U_i(x)$.

Тем самым все предположения (28) и т. д. общего индукционного шага доказаны для индекса на единицу больше; а поскольку из этих предположений, как показано выше, непосредственно следует Теорема VII, Теорема VII доказана в общем виде. Одновременно показано, что произвол пары $(\Gmod^*_{i-1},\Mmod^*_{i-1})$, обусловленный произволом выбора $C^{(i)}$, снова устраняется изоморфизмом системы классов вычетов $\Gmod^*_{i-1}\mid\Mmod^*_{i-1}$.

В частности, если поле $P$, положенное в §~3 в качестве области коэффициентов, имеет характеристику нуль -- то есть если простое поле, содержащееся в $P$ и выводимое из единицы, имеет тип рациональных чисел --, то неопределенные $u_{\mu\nu}$ можно специализировать до величин $\bar u_{\mu\nu}$ из $P$ так, что для $i=1\ldots n$ соответственно каждое $C^{(i)}$ остается регулярным по $x_i$. Приведенные рассуждения показывают, что Теорема VII сохраняется и при этой специализации.\footnote{Hentzelt вместо произвольного $P$ берет только поле всех комплексных чисел и рассматривает главным образом этот последний случай, тогда как неопределенные он кладет в основу только в собственно теории элиминации. Далее Hentzelt показывает, причем существенно теми же рассуждениями, что приведены здесь, что для образования результантной формы достаточно в каждом $\Mmod_{i-1}$ доходить лишь до некоторой конечной степени по $x$, превышающей фиксированную границу. Но отсутствует данный в Теореме VII изоморфизм $\Gmod_{i-1}\mid\Mmod_{i-1}$ с $\Gmod^*_{i-1}\mid\Mmod^*_{i-1}$, который объясняет независимость от чисел степеней. (E. N.)} Основной модуль и элементарные делители при этом, однако, не обязаны возникать из общего случая посредством специализации $u_{\mu\nu}=\bar u_{\mu\nu}$.\footnote{Для $\mideal=(x^2,ux+y)$ имеем $\gideal_0=(1)$, $\gideal_1=(1)$. Если выбрать $C^{(1)}=x^2$, то $\Gmod_1^*=(1,x)$, $\Mmod_1^*=(ux+y,xy)$, причем $\Mmod_1^*$ имеет элементарные делители $y^2$ и $1$, а $C^{(2)}=y^2$ можно выбрать. При $u=0$ полиномы $C^{(1)}$ и $C^{(2)}$ остаются регулярными соответственно по $x$ и $y$; но $\Mmod_1^*=(y,xy)$ имеет элементарные делители $y,y$. Значит, $\Gmod_1\mid\Mmod_1$ также становится изоморфной системе классов вычетов конечного модуля, но не изоморфной $\Gmod_1\mid\Mmod_1$. Если бы, напротив, было выбрано $C^{(1)}=ux+y$ и специализация была проведена так, чтобы $C^{(1)}$ оставался регулярным, то изоморфизм также сохранился бы. (E. N.)}

\clearpage
\section*{§ 5. Результантная форма и форма элементарных делителей идеала из полиномов.}

Пусть $x_{i+1}\ldots x_n$ присоединены к $P(u)$, так что $\Gmod_{i-1}$ и $\Mmod_{i-1}$, а следовательно также $\Gmod^*_{i-1}$ и $\Mmod^*_{i-1}$, переходят в модули линейных форм, где целые величины можно рассматривать как полиномы от \emph{одной} неопределенной $x_i$. По Теореме VII в этом случае элементарные делители $\Mmod_{i-1}$, отличные от $1$, и самое большее конечное число элементарных делителей $1$ модуля $\Mmod_{i-1}$ совпадают с соответствующими делителями $\Mmod^*_{i-1}$. Произведение этих элементарных делителей, $\rho$-й детерминантный делитель $\Mmod^*_{i-1}$, было равно детерминанту переходного преобразования от $\Gmod^*_{i-1}$ к $\Mmod^*_{i-1}$, то есть по Определению III равно $N(\Gmod^*_{i-1}\mid\Mmod^*_{i-1})$. По (24), соответственно по аналогичной формуле для общего $i$, вытекающей из (28), видно, что произведение этих элементарных делителей также равно детерминанту переходного преобразования от $\Gmod_{i-1}$ к $\Mmod_{i-1}$; следовательно, его надо обозначать $N(\Gmod_{i-1}\mid\Mmod_{i-1})$. Поэтому получается
\[
 N(\Gmod_{i-1}\mid\Mmod_{i-1})=N(\Gmod^*_{i-1}\mid\Mmod^*_{i-1})=R^{(i)}(x),
\]
где полином $R^{(i)}(x)$, то есть наибольший общий делитель в полиномиальном смысле всех $\rho$-рядных детерминантов из $\Mmod^*_{i-1}$, можно считать целым и примитивным по $x_{i+1}\ldots x_n$ и потому, как делитель $C^{(i)}$, регулярным по $x_i$. Точно так же высший элементарный делитель $E^{(i)}(x)$ модуля $\Mmod_{i-1}$, соответственно $\Mmod^*_{i-1}$, можно считать целым и примитивным по $x_{i+1}\ldots x_n$ и потому регулярным по $x_i$. Это приводит к следующему определению.

\medskip
\noindent\textbf{Определение VI.} \emph{Полином $R^{(i)}(x)$ называется результантом $i$-го уровня идеала $\mideal$, а $E^{(i)}(x)$ -- элементарным делителем $i$-го уровня. Следовательно, результант $i$-го уровня есть норма основного идеала $\gideal_{i-1}$ относительно $\mideal$, истолкованная как норма модуля $\Gmod_{i-1}$ относительно $\Mmod_{i-1}$, причем $x_{i+1}\ldots x_n$ присоединены к $P(u)$; точно так же $E^{(i)}(x)$ при том же присоединении становится равным частному $\Mmod_{i-1}/\Gmod_{i-1}$. Результантной формой, соответственно формой элементарных делителей, идеала $\mideal$ называются произведения}
\[
 R_\mideal=R^{(1)}R^{(2)}\cdots R^{(n)},\qquad
 E_\mideal=E^{(1)}E^{(2)}\cdots E^{(n)}.
\]

Для результантной формы и формы элементарных делителей имеет место следующее утверждение.

\medskip
\noindent\textbf{Теорема VIII.} \emph{Результантная форма делится на форму элементарных делителей; обратно, некоторая степень $E_\mideal$ делится на $R_\mideal$. Формы $E_\mideal$ и $R_\mideal$ делятся на $\mideal$; более общо, $E^{(i)}\cdots E^{(n)}\gideal_{i-1}\equiv0(\mideal)$ и, следовательно, $R^{(i)}\cdots R^{(n)}\gideal_{i-1}\equiv0(\mideal)$.}

Действительно, поскольку норма делится на высший элементарный делитель и, обратно, некоторая степень этого высшего элементарного делителя делится на норму, эта делимость для $R^{(i)}(x)$ и $E^{(i)}(x)$ следует после присоединения $x_{i+1}\ldots x_n$. Но $R^{(i)}(x)$ и $E^{(i)}(x)$ предполагаются примитивными полиномами относительно $x_{i+1}\ldots x_n$, поэтому делимость имеет место относительно всех неопределенных $x_i,x_{i+1},\ldots,x_n$. Тем самым, по определению этих выражений относительно всех неопределенных $x$, $R_\mideal$ делится на $E_\mideal$, а некоторая степень $E_\mideal$ делится на $R_\mideal$.

Далее, по Теоремам VII и II, после присоединения $x_{i+1}\ldots x_n$ высший элементарный делитель $E^{(i)}(x)$ определяется как частное $\Mmod_{i-1}/\Gmod_{i-1}$. Следовательно, если снова перейти к полиномам от всех неопределенных $x$, для каждого
\begin{equation}
 G(x)\equiv0(\gideal_{i-1})
 \quad\hbox{существует}\quad b^{(i+1)}\ne0,
 \quad\hbox{такое что}\quad b^{(i+1)}E^{(i)}(x)G(x)\equiv0(\mideal).
\tag{33}
\end{equation}
В частности, $\gideal_0$ равен единичному идеалу, поэтому получается
\[
 b^{(2)}E^{(1)}(x)\equiv0(\mideal),\quad b^{(2)}\ne0,
 \qquad\hbox{и следовательно}\qquad E^{(1)}(x)\equiv0(\gideal_1).
\]
В общем случае по (33), примененной к
\[
 E^{(1)}(x)\cdots E^{(i-1)}(x)\equiv0(\gideal_{i-1}),
\]
существует $b^{(i+1)}\ne0$ такое, что $b^{(i+1)}E^{(1)}\cdots E^{(i)}\equiv0(\mideal)$; следовательно, $E^{(1)}\cdots E^{(i)}\equiv0(\gideal_i)$. Так как $\gideal_n=\mideal$, отсюда получается
\begin{equation}
 E_\mideal=E^{(1)}\cdots E^{(n)}\equiv0(\mideal)
 \quad\hbox{и следовательно}\quad
 R_\mideal=R^{(1)}\cdots R^{(n)}\equiv0(\mideal).
\tag{34}
\end{equation}

Пусть $G(x)$ в (33) пробегает конечное число полиномов некоторого базиса идеала $\gideal_{i-1}$, а $c^{(i+1)}(x)$ обозначает произведение соответствующих им $b^{(i+1)}(x)$. Тогда получается
\[
 c^{(i+1)}E^{(i)}(x)\gideal_{i-1}\equiv0(\mideal);
 \qquad c^{(i+1)}\ne0
 \quad\hbox{и поэтому}\quad E^{(i)}(x)\gideal_{i-1}\equiv0(\gideal_i),
\]
и отсюда, как выше, конечным числом повторений:
\[
 E^{(n)}(x)\cdots E^{(i)}(x)\gideal_{i-1}\equiv0(\mideal),
\]
а следовательно также
\[
 R^{(n)}(x)\cdots R^{(i)}(x)\gideal_{i-1}\equiv0(\mideal).
\]
Тем самым Теорема VIII доказана во всех частях.

Теорема VIII существенно опирается на свойства формы элементарных делителей; но характеристическое значение результантной формы для $\mideal$ показывает следующая теорема.

\medskip
\noindent\textbf{Теорема IX.} \emph{Если $\nideal$ является делителем $\mideal$ -- делителем в идеальном смысле -- и результантные формы $R_\nideal$ и $R_\mideal$ совпадают, то идеалы $\nideal$ и $\mideal$ совпадают.}

Действительно, пусть $\gideal_0\ldots\gideal_{n-1},\gideal_n=\mideal$ и $\hideal_0\ldots\hideal_{n-1},\hideal_n=\nideal$ обозначают основные идеалы $\mideal$ и $\nideal$. Сначала $\gideal_0$ и $\hideal_0$ оба равны единичному идеалу, значит $\gideal_0=\hideal_0$. Предположим теперь, что $\gideal_{i-1}=\hideal_{i-1}$, так что $\Mmod_{i-1}$ и $\Nmod_{i-1}$ имеют один и тот же основной модуль $\Gmod_{i-1}$. Тогда предположение $R_\nideal=R_\mideal$ после присоединения $x_{i+1}\ldots x_n$ можно записать в виде
\[
 N(\Gmod_{i-1}\mid\Nmod_{i-1})=N(\Gmod_{i-1}\mid\Mmod_{i-1})
 \quad\hbox{или также}\quad
 N(\Gmod^*_{i-1}\mid\Nmod^*_{i-1})=N(\Gmod^*_{i-1}\mid\Mmod^*_{i-1}),
\]
где, соответственно (32), положено $\Nmod_{i-1}=(\Nmod^*_{i-1},\Gmod_{i-1})$, так что $\Nmod^*_{i-1}$ также становится модулем линейных форм от $\xi$, а $\Gmod^*_{i-1}$ является его основным модулем. Из делимости $\Mmod_{i-1}$ на $\Nmod_{i-1}$, которая влечет делимость $\Mmod^*_{i-1}$ на $\Nmod^*_{i-1}$, по Теореме IV следует, что при этом присоединении модули $\Mmod^*_{i-1}$ и $\Nmod^*_{i-1}$, а следовательно также $\Nmod_{i-1}$ и $\Mmod_{i-1}$, совпадают. Но присоединение $x_{i+1}\ldots x_n$ означает, что к $\Mmod_{i-1}$, соответственно к $\mideal$, добавляются все такие полиномы $G(x)$, для которых
\[
 b^{(i+1)}(x)G(x)\equiv0(\mideal),
 \qquad b^{(i+1)}\ne0.
\]
То есть вследствие присоединения $\Mmod_{i-1}$, соответственно $\mideal$, переходит в $\gideal_i$, а соответственно $\nideal$ переходит в $\hideal_i$; следовательно, $\gideal_i=\hideal_i$ и потому $\gideal_n=\hideal_n$, то есть $\mideal=\nideal$.

Наконец, тот же принцип переноса, примененный к Теореме V, дает еще следующее.

\medskip
\noindent\textbf{Теорема X.} \emph{Пусть $R^{(i)}=S^{(i)}T^{(i)}$ есть разложение $R^{(i)}$ на взаимно простые -- в полиномиальном смысле -- полиномы $S^{(i)}$ и $T^{(i)}$. Тогда существует один и только один идеал $\jideal$, а также один и только один идеал $\tideal$, причем оба являются делителями $\mideal$ с тем же основным идеалом $(i-1)$-го уровня $\gideal_{i-1}$, так что $S^{(i)}$ равно $N(\Gmod_{i-1}\mid\Smod_{i-1})$, а $T^{(i)}$ равно $N(\Gmod_{i-1}\mid\mathfrak T_{i-1})$. Идеалы $\jideal$ и $\tideal$ определяются как совокупность полиномов $S(x)$, соответственно $T(x)$, таких что}
\[
 s^{(i+1)}(x)T^{(i)}(x)S(x)\equiv0(\mideal),\quad s^{(i+1)}\ne0,
\]
\[
 t^{(i+1)}(x)S^{(i)}(x)T(x)\equiv0(\mideal),\quad t^{(i+1)}\ne0,
\]
\emph{и $\gideal_i$ равен наименьшему общему кратному $\jideal$ и $\tideal$,}
\[
 \gideal_i=[\jideal,\tideal].
\]

Здесь надо лишь заметить, что $s^{(i+1)},t^{(i+1)}$ для $\jideal$ и $\tideal$ можно выбрать фиксированно, как произведение тех $s^{(i+1)},t^{(i+1)}$, которые соответствуют базисным элементам $\jideal$, соответственно $\tideal$; и что, как отмечено выше, при присоединении $x_{i+1}\ldots x_n$ идеал $\mideal$ преобразуется в $\gideal_i$. В частности, разложение $R^{(i)}$ можно продолжать вплоть до степеней неприводимых полиномов -- первичных множителей; это влечет соответствующее представление $\gideal_i$ как наименьшего общего кратного.

\clearpage
\section*{§ 6. Собственно теория исключения.}

По Теореме VIII, соответственно по формуле (34), результантная форма и форма элементарных делителей делятся на $\mideal$, следовательно обращаются в нуль во всех нулях $\mideal$; при этом под ``нулями $\mideal$'' понимаются такие -- принадлежащие алгебраически замкнутому полю, полученному из $P(u;x_{i+1}\ldots x_n)$\footnote{Что каждое поле, и притом по существу единственным образом, можно расширить до алгебраически замкнутого, показал Steinitz в своей \emph{Algebraische Theorie der Körper} (\mbox{J. f. M.} 137 (1910), с. 167--309) и тем самым дал рациональный эквивалент основной теоремы алгебры. Фактически для каждого $i=1,\ldots,n$ речь идет о конечном расширении поля $P(u,x_{i+1},\ldots,x_n)$. (E. N.)} -- системы значений $x_1\ldots x_i$, что для этих систем значений обращается в нуль каждый полином из $\mideal$, причем $x_{i+1}\ldots x_n$ мыслятся присоединенными к области коэффициентов $(i=1,\ldots,n)$. Эти присоединенные $x_{i+1}\ldots x_n$, как будет показано, можно также заменить величинами из $P(u)$. Содержание теории исключения, наоборот, состоит в выведении нулей $\mideal$ из нулей результантной формы. Это обращение опирается на последовательное исключение, которое будет развито в этом параграфе; когда в следующем параграфе будет доказана делимость возникающей при этом формы $D^{(i)}(x)$ на $E^{(i)}(x)$, то из делимости некоторой степени $E^{(i)}(x)$ на $R^{(i)}(x)$ получится требуемое обращение.

Теория последовательного исключения, которую надо здесь дать, основана на следующем.

\medskip
\noindent\textbf{Теорема XI.} \emph{Пусть $\bideal$ означает идеал из полиномов от одной неопределенной $t$ с коэффициентами из $P$, то есть главный идеал; пусть $h(t)$ есть полином из $\bideal$ степени $k$, а $\Bmod$ -- модуль линейных форм от $1,t,\ldots,t^{k-1}$, в который $\bideal$ переходит по модулю $h(t)$. Тогда $\Bmod$ имеет ранг $k-p$, если $p$ означает степень базисного полинома $f(t)$ идеала $\bideal$.}

По предположению $h(t)=h_1(t)f(t)$, где $h_1(t)$ имеет степень $k-p$. Поэтому классы вычетов идеала $\bideal$ по модулю $h(t)$, представленные полиномами $f(t),tf(t),\ldots,t^{k-p-1}f(t)$, линейно независимы; и каждый класс вычетов идеала $\bideal$ по модулю $h(t)$ линейно, с коэффициентами из $P$, выражается через эти $k-p$ специальных классов. Но $\bideal$ по модулю $h(t)$ переходит в модуль линейных форм от $1,t,\ldots,t^{k-1}$, ранг которого тем самым равен ровно $k-p$.

Пусть теперь $\aideal=\aideal_1$ есть преобразованный идеал из полиномов, $A^{(1)}(x)$ -- полином из $\aideal_1$, регулярный по $x_1$, степени $k_1$, а $\mathfrak A_1$ -- модуль линейных форм от $1,x_1,\ldots,x_1^{k_1-1}$, с полиномами $a^{(2)}(x)$ в качестве коэффициентов, в который $\aideal_1$ переходит по модулю $A^{(1)}(x)$. Далее пусть $\aideal_2$ означает идеал от $x_2\ldots x_n$, полученный из всех $k_1$-рядных детерминантов из $\mathfrak A_1$. По Теореме XI $\aideal_2$ тогда и только тогда равен нулевому идеалу, когда полиномы из $\aideal_1$ имеют общий делитель -- делитель в полиномиальном смысле -- степени $p_1>0$ по $x_1$. Если $\aideal_2$ отличен от нулевого идеала, то, поскольку $\aideal$ предполагался преобразованным идеалом, он содержит полином $A^{(2)}(x)$ степени $k_2$, регулярный по $x_2$, как это видно непосредственно из составления преобразования (12) из специальных преобразований (25) и (26). Пусть снова $\mathfrak A_2$ означает модуль линейных форм от $1,x_2,\ldots,x_2^{k_2-1}$, в который $\aideal_2$ переходит по модулю $A^{(2)}(x)$, а $\aideal_3$ -- идеал от $x_3\ldots x_n$, полученный из всех $k_2$-рядных детерминантов из $\mathfrak A_2$, который должен содержать полином $A^{(3)}(x)$, регулярный по $x_3$.

Таким образом продолжают процедуру, пока не придут к нулевому идеалу или пока $\aideal_{n+1}$ не станет единичным идеалом; поскольку $\aideal_{n+1}$ свободен от $x$, он может быть только нулевым или единичным идеалом. Оказывается, что идеалы $\aideal_1,\aideal_2,\ldots$ однозначно определены идеалом $\aideal$ и не зависят от положенных в основу полиномов $A^{(\lambda)}(x)$. Это ясно для $\aideal_1=\aideal$, значит можно предполагать это для $\aideal_1,\ldots,\aideal_i$. Если теперь $\aideal_i$ по модулю $A^{(i)}(x)$ переходит в модуль линейных форм $\mathfrak A_i$, то $\mathfrak A_i$ можно также охарактеризовать как совокупность полиномов из $\aideal_i$, которые по $x_i$ не достигают степени $k_i$; следовательно $\mathfrak A_i$ зависит только от степени $k_i$ полинома $A^{(i)}(x)$. Пусть теперь $\bar A^{(i)}(x)$ есть полином из $\aideal_i$, регулярный по $x_i$, степени $\bar k_i>k_i$, $\overline{\mathfrak A}_i$ -- соответствующий модуль линейных форм, а $\bar\aideal_{i+1}$ -- идеал, полученный из $\bar k_i$-рядных детерминантов из $\overline{\mathfrak A}_i$. Тогда имеет место модульное равенство
\[
 \overline{\mathfrak A}_i=(\mathfrak A_i,A^{(i)},x_iA^{(i)},\ldots,x_i^{\bar k_i-k_i-1}A^{(i)}).
\]
Так как из-за регулярности $A^{(i)}(x)$ по $x_i$ полиномы $A^{(i)},x_iA^{(i)},\ldots$ можно ввести как новые неопределенные линейных форм, отсюда следует совпадение $k_i$-рядных детерминантов из $\mathfrak A_i$ с $\bar k_i$-рядными из $\overline{\mathfrak A}_i$; следовательно $\bar\aideal_{i+1}=\aideal_{i+1}$.\footnote{В принципе это тот же самый вывод, на котором покоился изоморфизм $\Gmod^*_{i-1}\mid\Mmod^*_{i-1}$ с $\Gmod_{i-1}\mid\Mmod_{i-1}$.}

Кроме того, $\aideal_{i+1}$ всегда делится на $\aideal_i$. Это ясно, когда $\aideal_{i+1}$ есть нулевой идеал; в противоположном случае $\mathfrak A_i$ имеет ранг $k_i$, следовательно $\aideal_{i+1}$ есть по определению идеал, состоящий из $k_i$-рядных детерминантов из $\mathfrak A_i$ и делящийся на $\mathfrak A_i$, а потому и на $\aideal_i$. Тем самым каждый $\aideal_{i+1}$ делится на $\aideal$; значит, если $\aideal_{n+1}$ становится единичным идеалом, то и $\aideal$ есть единичный идеал.

Пусть теперь $\aideal$ отличен от единичного идеала; $\aideal_1\ne0,\ldots,\aideal_i\ne0;\ \aideal_{i+1}=0$. Пусть $D^{(i)}(x)$ есть наибольший общий делитель -- делитель понимается в полиномиальном смысле -- всех полиномов из $\aideal_i$, существующий по Теореме XI; как делитель $A^{(i)}(x)$ сам $D^{(i)}(x)$ регулярен по $x_i$ и имеет степень $p_i>0$. Присоединим $x_{i+1}\ldots x_n$ к $P(u)$; тогда $D^{(i)}(x)$ в алгебраическом расширении поля $P(u;x_{i+1}\ldots x_n)$ допускает разложение на $p_i$ линейных множителей. Пусть $x_i-\bar x_i$ есть один такой линейный множитель, так что при $x_i=\bar x_i$ все полиномы из $\aideal_i$ обращаются в нуль, и следовательно по Теореме XI полиномы из $\aideal_{i-1}$ при этой специализации получают наибольший общий делитель $D^{(i-1)}(x)$. Как делитель $A^{(i-1)}(x)$ сам $D^{(i-1)}(x)$ снова регулярен по $x_{i-1}$ и имеет степень $p_{i-1}>0$; в частности он не может обращаться в нуль тождественно и в расширении поля $P(u,\bar x_i,x_{i+1},\ldots,x_n)$ допускает разложение на $p_{i-1}$ линейных множителей. Если $x_{i-1}-\bar x_{i-1}$ есть такой линейный множитель, то ему соответствует наибольший общий делитель из $\aideal_{i-2}$. Продолжая так, приходят к конечному числу общих нулей $\aideal$, лежащих в алгебраическом расширении поля $P(u;x_{i+1}\ldots x_n)$. Обратно, из-за делимости $\aideal_i$ на $\aideal$ каждый нуль $\aideal$ является также нулем $\aideal_i$. Если, в частности, $x_1=\bar x_1,\ldots,x_i=\bar x_i,x_{i+1}=x_{i+1},\ldots,x_n=x_n$ есть нуль $\aideal$, то отсюда следует $\aideal_{i+1}=0$; и полиномы из $\aideal_i$ получают наибольший общий делитель с линейным множителем $x_i-\bar x_i$. В итоге получается:

\medskip
\noindent\textbf{Теорема XII.} \emph{Пусть $\aideal$ отличен от единичного идеала, $\aideal_1\ne0,\ldots,\aideal_i\ne0;\ \aideal_{i+1}=0$; тогда при присоединении $x_{i+1}\ldots x_n$ к $P(u)$ существует конечное число связанных систем значений $\bar x_1\ldots\bar x_i$, лежащих в алгебраическом расширении поля $P(u,x_{i+1},\ldots,x_n)$, таких, что все полиномы из $\aideal$ обращаются в нуль при $x_1=\bar x_1,\ldots,x_i=\bar x_i,x_{i+1}=x_{i+1},\ldots,x_n=x_n$. Обратно, если имеется такой нуль $\aideal$, то отсюда следует $\aideal_{i+1}=0$ и появление общего делителя $D^{(i)}(x)$ всех полиномов из $\aideal_i$, имеющего линейный множитель $x_i-\bar x_i$.}

Теорема XII показывает, в частности, что каждый идеал $\aideal$, не имеющий нуля, становится единичным идеалом.

Чтобы провести разложение, которое также явно показывает связанность систем значений, вводят новые неопределенные $v$, которые можно присоединить к $P(u,x_{i+1},\ldots,x_n)$. Положим:
\begin{equation}
 z_i=-v_1x_1-\cdots-v_{i-1}x_{i-1}+x_i,
\tag{35}
\end{equation}
так что составление (12) с (35) снова дает подстановку типа (12), где теперь только $u_{ik}$ заменены на $w_{ik}=u_{ik}-v_k$, а в общем случае $u_{i+h,k}$ -- на $w_{i+h,k}=u_{i+h,k}-u_{i+h,i}v_k$. Следовательно, если подстановка (35) переводит по предположению преобразованный идеал $\aideal$ в $\bideal$, то $\bideal$ просто получается из $\aideal$ заменой $u_{\mu\nu}$ на $w_{\mu\nu}$ и присоединением $v$; и тем же процессом идеалы $\bideal_1,\bideal_2,\ldots$, определенные посредством $\bideal$, получаются из $\aideal_1,\aideal_2,\ldots$. Обратно, $\aideal_i$ получается из $\bideal_i$ посредством $v=0$; следовательно идеалы $\aideal_i$ и $\bideal_i$ одновременно являются нулевыми идеалами или одновременно отличны от нулевого идеала.

Пусть теперь снова $\aideal_1\ne0;\ldots;\aideal_i\ne0;\aideal_{i+1}=0$; следовательно также $\bideal_1\ne0;\ldots;\bideal_i\ne0;\bideal_{i+1}=0$; пусть $\bar x_1\ldots\bar x_i,x_{i+1}\ldots x_n$ есть связанная система нулей $\aideal$. Если определить $\bar z_i=\bar x_i+v_1\bar x_1+\cdots+v_{i-1}\bar x_{i-1}$, то $\bar x_1\ldots\bar x_{i-1},\bar z_i,x_{i+1}\ldots x_n$ посредством (35) становится нулем $\bideal$. Итак, если $H^{(i)}(z,x)$ означает наибольший общий делитель всех полиномов из $\bideal_i$, который можно предполагать целым и примитивным по $v$, то по последней части Теоремы XII $H^{(i)}(z,x)$ имеет множитель
\[
 z_i-\bar z_i=z_i-(\bar x_i+v_1\bar x_1+\cdots+v_{i-1}\bar x_{i-1}),
\]
и каждому нулю $\aideal$, возникающему при присоединении $x_{i+1}\ldots x_n$, соответствует такой множитель. Надо показать, что этим разложение $H^{(i)}(z,x)$ исчерпано, то есть что из $H^{(i)}$ нельзя отделить никакого множителя $H^{(i)}_1$, который не разлагается на линейные множители по $z$ \emph{и} $v$. Пусть это возможно. Тогда из-за предполагаемой примитивности по $v$ полином $H^{(i)}$ содержит по меньшей мере один линейный множитель $z_i-\bar z_i$, где $\bar z_i$ принадлежит алгебраическому расширению поля $P(u,v,x_{i+1},\ldots,x_n)$. Если $\bar x_1\ldots\bar x_{i-1},\bar z_i,x_{i+1}\ldots x_n$ есть соответствующий нуль $\bideal$, то он должен совпасть с одним из конечного числа нулей $\aideal$, возникающих при присоединении $x_{i+1}\ldots x_n$, а значит быть независимым от $v$. Тем самым $\bar z_i$ необходимо содержит $v$ линейно, а не алгебраически или рационально-нелинейно; вопреки предположению, из $H^{(i)}_1$ можно отделить еще один линейный множитель $z_i-(\bar x_i+v_1\bar x_1+\cdots+v_{i-1}\bar x_{i-1})$ по $z$ и $v$. Следовательно, получается явное разложение:
\[
 H^{(i)}(z,x)=\prod \{z_i-(\bar x_i+v_1\bar x_1+\cdots+v_{i-1}\bar x_{i-1})\}^{\alpha_\lambda}.
\]
Так как степень $H^{(i)}$ и отдельные показатели $\alpha_\lambda$ должны совпадать с соответствующими для $D^{(i)}(x)$ -- ведь $H^{(i)}$ возникает из $D^{(i)}$ при специализации $u_{\mu\nu}=w_{\mu\nu}$, а $D^{(i)}$ из $H^{(i)}$ при специализации $v=0$ --, то это разложение еще показывает, что вследствие присоединения неопределенных $u_{\mu\nu}$ к $P$ и при исключении, исходящем из $\aideal$, каждому нулю $\bar x_i$ полинома $D^{(i)}$ соответствует соответственно линейный наибольший общий делитель из $\aideal_{i-1}\ldots\aideal_1$ или степень такого; значит, каждый раз получается лишь одна связанная система нулей $\bar x_1\ldots\bar x_i,x_{i+1}\ldots x_n$ идеала $\aideal$.\footnote{Следует указать, что данное здесь последовательное исключение -- в отличие от метода Кронекера -- зависит только от данного идеала $\aideal$ и не зависит ни от какого базиса идеала; в частности, это относится также к показателям $\alpha_\lambda$. Полное проведение исключения этим методом, по Гентцельту, потребовало бы продолжения процедуры на частном $\aideal/D^{(i)}$, что можно опустить ввиду следующего параграфа. (E. N.)}

\clearpage
\section*{§ 7. Результантная форма и теория исключения.}

Для установления связи между результантной формой и теорией исключения применим последовательное исключение из §~6 специально к частному $\aideal=\mideal/\gideal_{i-1}$ идеала по основному идеалу $(i-1)$-го уровня. Сначала надо показать, что это частное представляет собой преобразованный идеал. Поскольку $\mideal$ преобразован, то по Теореме VI из §~3 преобразован и $\gideal_{i-1}$. Из
\[
 H(x)\equiv0(\mideal/\gideal_{i-1});\qquad
 H(x)=\overline H(y)=\sum\alpha_i(u)\overline H_i(y)=\sum\alpha_i(u)H_i(x)
\]
следовательно получается:
\[
 \overline H(y)\overline\gideal_{i-1,(u)}\equiv0(\overline\mideal_{(u)});
 \quad\hbox{а значит также}\quad
 \overline H(y)\overline\gideal_{i-1}\equiv0(\overline\mideal_{(u)}),
\]
или
\[
 \overline H_i(y)\overline\gideal_{i-1}\equiv0(\overline\mideal),
 \quad\hbox{а значит}\quad H_i(x)\gideal_{i-1}\equiv0(\mideal),
\]
чем доказана делимость $H_i(x)$ на $\mideal/\gideal_{i-1}$ и тем самым доказано, что этот идеал преобразован, так что метод исключения из §~6 применим.

Но (30), с учетом (28) -- §~4 -- показывает, что $C^{(i)}(x)$ делится на $\mideal/\gideal_{i-1}$; здесь под $C^{(i)}(x)$ понимается любой не тождественно исчезающий $\rho$-рядный детерминант из $\Mmod^*_{i-1}$. Так как для $i>1$ этот $C^{(i)}$ имеет нулевую степень по $x_1$, то модуль $\mathfrak A_1$, соответствующий идеалу $\aideal=\aideal_1$, содержит полиномы $C^{(i)},x_1C^{(i)},\ldots,x_1^{k_1-1}C^{(i)}$, и потому $(C^{(i)})^{k_1}$ делится на $\aideal_2$; точно так же $(C^{(i)})^{k_1k_2}$ делится на $\aideal_3$, и наконец $(C^{(i)})^{k_1k_2\cdots k_{i-1}}$ делится на $\aideal_i$. Следовательно, $\aideal_1,\ldots,\aideal_i$ отличны от нулевого идеала, что верно также при $i=1$. Пусть теперь $D^{(i)}(x)$ есть наибольший общий делитель всех полиномов из $\aideal_i$, то есть имеет степень $p_i>0$ только тогда, когда $\aideal_{i+1}$ равен нулевому идеалу. По определению $D^{(i)}(x)$, с учетом делимости $\aideal_i$ на $\aideal$, имеем:
\[
 d^{(i+1)}D^{(i)}(x)\equiv0(\mideal/\gideal_{i-1});\qquad d^{(i+1)}\ne0;
\]
следовательно $d^{(i+1)}D^{(i)}$ становится полиномом из $\mideal/\gideal_{i-1}$, свободным от $x_1\ldots x_{i-1}$. Но $E^{(i)}(x)$ был определен как наивысший элементарный делитель $\Mmod_{i-1}$, соответственно $\Mmod^*_{i-1}$, (Теорема II и §~4) как наибольший общий делитель -- в полиномиальном смысле -- всех свободных от $x_1\ldots x_{i-1}$ полиномов из $\mideal/\gideal_{i-1}$. Поэтому $d^{(i+1)}D^{(i)}$, а вследствие регулярности $E^{(i)}(x)$ по $x_i$ также и $D^{(i)}(x)$, делится на $E^{(i)}(x)$.

То же рассуждение можно провести, если с самого начала преобразование (12) было составлено с (35), то есть если неопределенные $u_{\mu\nu}$ были заменены на $w_{\mu\nu}$; это дает делимость $H^{(i)}(z,x)$ на соответствующий $\bar E^{(i)}(z,x)$. Так как, подобно тому как $D^{(i)}(x)$ и $H^{(i)}(z,x)$ переходят друг в друга специализацией, так же переходят друг в друга $E^{(i)}(x)$ и $\bar E^{(i)}(z,x)$, а также $R^{(i)}(x)$ и $\bar R^{(i)}(z,x)$, то и числа кратности при разложении на линейные множители должны взаимно совпадать. Учитывая еще, что некоторая степень $E^{(i)}(x)$ делится на $R^{(i)}(x)$, получаем в итоге следующее.

\medskip
\noindent\textbf{Теорема XIII.} \emph{Если $R^{(i)}(x)$ разложить в алгебраическом расширении поля $P(u,x_{i+1},\ldots,x_n)$ на линейные множители $x_i-\bar x_i$, то $\bar x_i$ можно одним и только одним способом дополнить до связанной системы значений $\bar x_1\ldots\bar x_i,x_{i+1}\ldots x_n$, которая является нулем $\mideal/\gideal_{i-1}$ и, следовательно, нулем $\mideal$. Конечное число нулей $\mideal$, возникающих таким образом из линейных множителей $R^{(i)}$ -- конечное число после присоединения $x_{i+1}\ldots x_n$ -- объединяется явным разложением:}
\begin{equation}
 \bar R^{(i)}(z,x)=\prod\{z_i-(\bar x_i+v_1\bar x_1+\cdots+v_{i-1}\bar x_{i-1})\}^{\alpha_\lambda}.
\tag{36}
\end{equation}
\emph{Благодаря регулярности $\bar R^{(i)}(z,x)$ по $z_i$ это разложение сохраняется, если присоединенные к $P(u)$ величины $x_{i+1}\ldots x_n$ заменить произвольными специальными системами значений $\bar x_{i+1}\ldots\bar x_n$ из $P(u)$ или из полученного из него алгебраически замкнутого поля.}

Тем самым проведено также разделение нулей $\mideal$ по отдельным множителям результантной формы; число неопределенных $x$, присоединенных к $P(u)$, называется также размерностью соответствующего алгебраического образования.

Если в разложении $\bar R^{(i)}(z,x)$ или в соответствующем разложении $R^{(i)}(x)$ объединить множители, сопряженные относительно $P(u,x_{i+1}\ldots x_n)$, которые при общем $P$ могут быть полностью или частично тождественными, то получается разложение $R^{(i)}(x)$ на степени полиномов, неприводимых относительно $P(u,x_{i+1}\ldots x_n)$:
\[
 R^{(i)}(x)=S^{(i)}_1(x)^{\beta_1}\cdots S^{(i)}_\nu(x)^{\beta_\nu}.
\]
Этому разложению по Теореме X соответствует представление $\gideal_i=[\jideal_1\ldots\jideal_\nu]$ такое, что $S^{(i)}_\mu(x)^{\beta_\mu}$ равно норме $\gideal_{i-1}$ относительно идеала $\jideal_\mu$, соответственно норме модуля $\Gmod_{i-1}$ относительно модуля, соответствующего $\jideal_\mu$. Тем самым проясняется и значение показателей $\beta_\mu$; в частности, если $\gamma_\mu$ означает степень $S^{(i)}_\mu$, то кратность $\beta_\mu\gamma_\mu$ равна числу линейно независимых относительно $P(u,x_{i+1}\ldots x_n)$ классов вычетов $\gideal_{i-1}$ по $\jideal_\mu$.

Рассмотрим еще кратко специальный случай, когда $P$ имеет характеристику нуль. Если здесь специализировать $u_{\mu\nu}$ в величины $\bar u_{\mu\nu}$ из $P$ так, чтобы все $n$ детерминантов $C^{(i)}(x)$ $(i=1,2,\ldots,n)$ оставались регулярными по $x_i$,\footnote{Гентцельт с самого начала выполняет эту специализацию для $u_{\mu\nu}$ и поэтому для доказательства разложимости $H^{(i)}(z,x)$, соответственно $\bar R^{(i)}(z,x)$, на линейные множители по $z$ и $v$ должен привлекать гораздо более сложные рассуждения. Возникающие при этом показатели не обязательно совпадают с приведенными выше. (E. N.)} то регулярность $R^{(i)}$ также сохраняется. При этой специализации $R^{(i)}$, а также $\bar R^{(i)}(z,x)$, становится общим делителем всех $\rho$-рядных детерминантов из $\Mmod^*_{i-1}$, но не обязательно наибольшим общим делителем. Однако специализацию всегда можно выбрать так, чтобы сохранилось и это последнее свойство. Ведь $R^{(i)}(z,x)$ можно -- как показывает существование базиса модуля -- также охарактеризовать как наибольший общий делитель конечного числа $\rho$-рядных детерминантов из $\Mmod^*_{i-1}$. Разложение так специализированного $\bar R^{(i)}(z,x)$ тогда получается просто заменой $u_{\mu\nu}$ на $\bar u_{\mu\nu}$ в (36), так что вся теория исключения сохраняется.

\begin{center}
(Поступила 17.3.1922.)
\end{center}
% END UNIT BODY
% END PRODUCER UNIT 106
\endgroup


\clearpage
\begingroup
\setcounter{footnote}{0}
% BEGIN PRODUCER UNIT 107
% Source: C:/Users/Floris/Documents/interlanguage/03_projects/noether/02_slavic_working_corpus/translations/paper23/russian/v001/Noether_Paper23_Russian_v001.tex
% Identity: 34002 B / SHA-256 3FCB8BA89D8BD7A257C480ECBF8F1B779E485C651EF2B9BEB9877DF8D6AE75B4
\setlength{\parindent}{1.25em}
\setlength{\parskip}{0pt}
\makeatletter
\renewcommand\footnotesize{\@setfontsize\footnotesize{7.7}{8.7}\selectfont}
\makeatother
\emergencystretch=35em
\allowdisplaybreaks
\sloppy
\hbadness=10000
\vbadness=10000
\hfuzz=2pt
\vfuzz=10pt
\raggedbottom
\providecommand{\mJ}{\mathfrak J}
\providecommand{\mS}{\mathfrak S}
\providecommand{\mo}{\mathfrak o}
\providecommand{\ma}{\mathfrak a}
\providecommand{\mb}{\mathfrak b}
\providecommand{\dd}{\mathrm d}
\providecommand{\D}{\Delta}
\providecommand{\Om}{\Omega}
\providecommand{\dx}{\dd x}
\providecommand{\dy}{\dd y}
% BEGIN UNIT BODY
% Unit: Paper23 v001. Direct Russian translation from the German final-audited source slice.
\setcounter{footnote}{0}
\section*{23. Алгебраические и дифференциальные инварианты}
\begin{center}
J. Ber. d. DMV 32 (1923), с. 177--184
\end{center}

Если мне нужно дать обзор развития алгебраических и дифференциальных инвариантов, то в обеих областях я хотела бы ограничиться критическим периодом, который, по замечанию Hilbert, следует за наивным и формальным периодом. Этот критический период для алгебраических инвариантов характеризуется именем самого Hilbert, для дифференциальных инвариантов --- именем Riemann; или, по существу: для алгебраических инвариантов --- арифметическими методами алгебры, которые существенно развились вокруг этих вопросов и здесь смогли проявить всю свою остроту; для дифференциальных инвариантов --- методами формального вариационного исчисления. Поэтому я хотела бы говорить об арифметических методах с их значением за пределами специального предмета, тогда как во второй части могу ограничиться подчинением дифференциальных инвариантов алгебраическим; именно это и осуществляется в связи с методами вариационного исчисления.

\textbf{1.} В основе \emph{алгебраических инвариантов}, как известно, лежит \emph{общая линейная группа} в $x_1,\ldots,x_n$:
\[
\tag{1}
x_i=\sum_k s_{ik}x'_k
\qquad\text{или кратко:}\qquad
x=S(x'),
\]
где $s_{ik}$ обозначают неопределенные, так что $\D=|s_{ik}|$ отлично от нуля. Пусть теперь $f(a,x)$, $g(b,x),\ldots$ --- формы по $x$ степеней $\alpha,\beta,\ldots$ с неопределенными $a,b,\ldots$ в качестве коэффициентов. Тогда посредством (1) возникает \emph{индуцированное преобразование} для $a,b,\ldots$; действительно, из
\[
\begin{aligned}
f(a,x)&=\sum a_{i_1\ldots i_n}x_1^{i_1}\cdots x_n^{i_n}
      =\sum a'_{j_1\ldots j_n}{x'_1}^{j_1}\cdots {x'_n}^{j_n}
      =f(a',x'),\\
g(b,x)&=\sum b_{r_1\ldots r_n}x_1^{r_1}\cdots x_n^{r_n}
      =\sum b'_{s_1\ldots s_n}{x'_1}^{s_1}\cdots {x'_n}^{s_n}
      =g(b',x')
\end{aligned}
\]
следует, что $a',b',\ldots$ становятся линейными однородными функциями от $a,b,\ldots$, коэффициенты которых имеют степени $\alpha,\beta,\ldots$ по $s_{ik}$:
\[
\tag{2}
a'=A_\alpha(a);\qquad b'=B_\beta(b)\ldots .
\]
Под \emph{инвариантом} здесь понимают инвариант относительно преобразований (1) и (2), то есть целую рациональную функцию от неопределенных $a,b,\ldots,x,y,\ldots$ --- где также $y=S(y')$ ---, которая при применении преобразований воспроизводится с точностью до множителя, степени детерминанта подстановки:
\[
\tag{3}
I(a',b',\ldots,x',y',\ldots)=\D^\rho I(a,b,\ldots,x,y,\ldots).
\]
Коэффициенты $I$ можно считать рациональными или даже целыми рациональными числами, поскольку всякий инвариант с коэффициентами из произвольного числового поля $P$ линейно, с коэффициентами из $P$, выражается через конечное число таких специальных инвариантов. Точно так же не является ограничением считать $I$ отдельно однородным в каждом ряду, поскольку и здесь всякий инвариант составляется как сумма конечного числа таких специальных инвариантов.

Отметим, что требование (3), наоборот, снова определяет (1) и (2). Как недавно показали Ostrowski и Schur,\footnote{A. Ostrowski и I. Schur, \emph{Über eine fundamentale Eigenschaft der Invarianten einer binären Form}, Math. Zeitschr. 15 (1922), и связанные с этим, еще не опубликованные работы Ostrowski.} индуцированные преобразования можно также охарактеризовать как совокупность линейных преобразований, разделенных в каждом отдельном ряду, которые --- за исключением некоторых указуемых исключений --- допускают произвольный инвариант $I$, свободный от $x,y,\ldots$.

Произведение двух инвариантов снова является инвариантом, а сумма является им тогда и только тогда, когда совпадает вес --- показатель $\rho$ в (3). Если же ограничиться преобразованиями с детерминантом единица, то и произвольная сумма снова является инвариантом и при этом исчерпывает все инварианты относительно так ограниченной группы. Поэтому получается область целостности, а именно \emph{однородная область целостности}, то есть область, в которой принадлежность полинома влечет отдельную принадлежность его однородных составляющих; здесь это снова сводится к отдельно однородным инвариантам согласно (3).

Здесь возникает фундаментальная проблема, на которой развивались теория инвариантов и вообще арифметическая алгебра: \emph{является ли эта область целостности конечной}, то есть обладает ли она \emph{конечным интегральным базисом} $I_1,\ldots,I_k$ таким, что всякий инвариант можно представить целой рациональной функцией от $I_1,\ldots,I_k$, $I=G(I_1,\ldots,I_k)$? Решение этой проблемы у Hilbert --- доказательство Gordan для бинарного случая не допускает расширения из-за необозримых символических вычислений, а доказательство Mertens--Hilbert не допускает его из-за использования разложения бинарной формы на линейные формы --- основано на сведении вопроса об интегральном базисе к вопросу об идеальном базисе. Я хотела бы здесь кратко набросать ход мысли, подчеркивая основные арифметические идеи, а затем перейти еще к трем связанным кругам вопросов: фактическому построению базиса конечным числом шагов, вопросам конечности для подгрупп, вопросам конечности с учетом целочисленности.

\textbf{2.} Напомню \emph{определение идеала} для конечных алгебраических числовых полей. Система $\ma$ чисел поля называется идеалом, если
\[
\begin{array}{ll}
1.& \ma\ \text{принадлежит области }\mo\text{ всех целых чисел поля},\\
2.& \text{вместе с }\alpha\text{ и }\beta\text{ к }\ma\text{ принадлежит также разность }\alpha-\beta,\\
3.& \text{вместе с }\alpha\text{ к }\ma\text{ принадлежит также }\lambda\alpha,\text{ где под }\lambda
\text{ понимается произвольный элемент из }\mo.
\end{array}
\]
И, как известно, имеет место теорема, что каждый идеал имеет \emph{идеальный базис}: $\ma=(\alpha_1,\ldots,\alpha_k)$; то есть существует конечное число элементов $\alpha_1,\ldots,\alpha_k$ из $\ma$ такое, что посредством 2. и 3. вся $\ma$ выводится из $\alpha_1,\ldots,\alpha_k$; следовательно, $\alpha=\lambda_1\alpha_1+\cdots+\lambda_k\alpha_k$ исчерпывает все элементы из $\ma$.\footnote{То, что $k$ можно уменьшить до двух, для дальнейшего несущественно.}

Это определение идеала опирается только на то, что $\mo$ образует область целостности; поэтому понятие идеала, равно как и понятие идеального базиса, буквально сохраняется, если $\mo$ заменить произвольной областью целостности.

Теперь доказательство конечности у Hilbert распадается на три утверждения:
\[
\begin{array}{ll}
1.& \text{В области всех полиномов от }n\text{ неопределенных с коэффициентами из}\\
& \text{области рациональности теорема об идеальном базисе имеет место для каждого идеала}\\
& \text{--- а следовательно, и для всякой произвольной системы }\mS.\\[2pt]
2.& \text{Однородная область целостности }\mJ\text{ из полиномов конечна тогда и только тогда,}\\
& \text{когда в }\mJ\text{ имеет место теорема об идеальном базисе.}\\[2pt]
3.& \text{Для области инвариантов }\mJ\text{ теорема об идеальном базисе имеет место в}\\
& \text{усиленной форме: каждый идеальный базис, образованный относительно области}\\
& \text{всех полиномов, одновременно является идеальным базисом в }\mJ,\text{ так что наряду с}
\end{array}
\]
\[
 I=A_1I_1+\cdots+A_kI_k
 \qquad\text{всегда выполняется}\qquad
 I=i_1I_1+\cdots+i_kI_k,
\]
где $A$ обозначают полиномы от $a,b,\ldots,x,y,\ldots$, а $i$ --- инварианты.

При этом доказательство 1. в принципе опирается на те же рассуждения, что и в алгебраическом числовом поле; в обоих случаях существование идеального базиса получается как прямое следствие общей теоремы о модулях из линейных форм.\footnote{Ср. мой сравнительный обзор арифметической теории алгебраических функций. Jahresber. 28 (1920), с. 187 и прим. 2), с. 188.} Из 1. прямо следует существование идеального базиса для каждой конечной области целостности из полиномов; для однородных областей легко видеть и обратное. Только в 3. используются специальные свойства инвариантов; у Hilbert 3. получается как следствие известной теоремы об $\Omega$-процессе, тогда как недавно E. Fischer показал, что 3. --- при опоре на другие общие теоремы --- можно уже получить из свойства общей линейной группы содержать вместе с каждым преобразованием также контрагредиентное к нему, соответственно его сопряженно-комплексное.\footnote{E. Fischer, \emph{Über die Endlichkeit der Invarianten}. Göttinger Nachrichten 1915, и лейпцигский доклад.}

\textbf{3.} \emph{Фактическое построение базиса} опирается на дальнейшее арифметическое переформулирование понятия конечности с привлечением теории \emph{полей функций}. Имеет место следующее:
\[
\begin{array}{ll}
4.& \text{Область целостности }\mJ\text{ из полиномов конечна тогда и только тогда,}\\
& \text{когда в }\mJ\text{ содержится конечная подобласть }\mJ'\text{ такая, что }\mJ\\
& \text{является алгебраически целой над }\mJ';\text{ тогда каждый интегральный базис }\mJ'\\
& \text{дополняется фундаментальной системой этих алгебраически целых величин}\\
& \text{до базиса }\mJ.
\end{array}
\]
\footnote{Hilbert показывает (\emph{Über die vollen Invariantensysteme}, Math. Annalen 42 (1893), §§ 1 и 2), что эти факты следуют из предполагаемой конечности; наоборот, то, что из алгебраически-целой зависимости следует конечность, вытекает из существования рационального базиса (E. Noether, \emph{Körper und Systeme rationaler Funktionen}, Math. Annalen 76 (1915), § 12).}
Если, в частности, как в случае инвариантов, речь идет об однородных областях $\mJ$, для которых теорема об идеальном базисе имеет место в более сильной форме 3., то такая область $\mJ'$ выводится из каких-либо, всегда существующих, конечного числа полиномов $I_1,\ldots,I_s$ как интегрального базиса, из обращения которых в нуль следует обращение в нуль всех полиномов из $\mJ$. Этот факт снова основан на общей теореме идеальной теории:
\[
\begin{array}{ll}
5.& \text{Каждому идеалу }\ma\text{ из полиномов соответствует целое рациональное число }r\\
& \text{такое, что для каждого идеала }\mb\text{, который обращается в нуль во всех нулях }\ma,\\
& \mb^r\text{ становится делимым на }\ma.
\end{array}
\]

В случае инвариантов можно определить \emph{верхнюю границу для степеней} $I_1,\ldots,I_s$, а тем самым и для их числа $s$, зависящую только от степеней положенных в основу основных форм. И в этом месте снова привлекаются специальные свойства инвариантов; речь идет о теореме, что численно специализированная система основных форм тогда и только тогда имеет отличный от нуля инвариант, когда $\D=|s_{ik}|$ является алгебраически целой над преобразованными коэффициентами.

Из этой верхней границы K. Hentzelt\footnote{Я вскоре издам эту работу из наследия.} определил \emph{верхнюю границу для степеней инвариантов интегрального базиса}, снова зависящую только от степеней основных форм. Ему это удается потому, что для показателя $r$ в 5. он задает верхнюю границу, которая зависит только от числа и наибольшей степени полиномов идеального базиса $\ma$, не зависит от коэффициентов этих полиномов и вычисляется из указанных чисел конечным числом шагов. В принципе поставленная проблема тем самым полностью решена; правда, указанная граница весьма высока: например, для тернарной квадратичной формы, где дискриминант, то есть инвариант третьей степени, уже образует базис всех свободных от $x$ инвариантов, по Hilbert и Hentzelt доходят далеко до миллионов.

\textbf{4.} В вопросе конечности инвариантов \emph{подгрупп} общей линейной группы пока получены только частичные результаты. Метод почти исключительно опирается на сведение инвариантов подгрупп к симультанным инвариантам общей линейной группы, по процедуре, примененной Study для ортогональной группы. Это осуществили Weitzenböck для инвариантов движений и аффинных инвариантов, Deruyts --- для семиинвариантов и инвариантов сдвига;\footnote{Литературные указания см. в энциклопедической статье Weitzenböck по теории инвариантов; III E 1.} вопрос о пределах применимости этого метода еще не решен. Упомянутое в конце 2. замечание Fischer также решает проблему конечности для ряда подгрупп; в частности, тем самым дано простейшее доказательство для ортогональной группы. Однако пример семиинвариантов показывает --- на что недавно указал Fischer\footnote{Лейпцигский доклад и Math. Zeitschrift.} ---, что для подгрупп, несмотря на конечность, более сильная форма 3. для идеального базиса может не выполняться.

Если понимать имеющийся здесь вопрос конечности в смысле теории полей, он ведет к гильбертовой постановке вопроса об \emph{относительно целых функциях}, то есть к вопросу, всегда ли конечны такие области целостности, которые состоят из совокупности полиномов некоторого поля рациональных функций. В этом направлении лежит элементарно доказуемый факт, что для инвариантов конечных групп выделенный интегральный базис можно указать непосредственно --- систему коэффициентов резольвенты Galois.\footnote{E. Noether, \emph{Der Endlichkeitssatz der Invarianten endlicher Gruppen}. Math. Ann. 77 (1916).}

\textbf{5.} Вопрос о том, можно ли усилить теорему конечности инвариантов относительно \emph{целочисленности}, также в общем виде еще не решен. Правда, как показал Hilbert, теорема об идеальном базисе имеет место и для области всех целочисленных полиномов; точно так же сохраняется связь 2. между идеальным базисом и интегральным базисом;\footnote{Из этой связи, в частности, следует, что для всех конечных областей целостности имеют место теоремы разложения общей идеальной теории, как я развила их в Math. Ann. 83 (1921).} но доказательство для --- не обязательной --- более сильной формы 3. отказывает. Для случая бинарных инвариантов целочисленную конечность можно доказать;\footnote{E. Noether, \emph{Die Endlichkeit des Systems der ganzzahligen Invarianten binärer Formen}. Göttinger Nachrichten 1919.} это доказательство, представляющее собой целочисленное усиление бинарного доказательства конечности Mertens--Hilbert, использует наряду с теоремой о целочисленном идеальном базисе также разложение бинарных форм на линейные формы и тем самым не допускает переноса на большее число неопределенных. Напротив, на сходных основаниях получается целочисленная конечность для инвариантов конечных групп как усиление указанного выше доказательства, хотя теперь снова возникает только доказательство существования, а не фактическое указание базиса.\footnote{§ 4 работы, цитированной под 11).} И здесь, вероятно, только дальнейшее развитие теории полей функций и общей идеальной теории приведет к цели.

\textbf{6.} Теперь я перехожу к \emph{дифференциальным инвариантам} и к поставленному во введении вопросу об их сведении к теории алгебраических инвариантов. Положенная в основу группа здесь, как известно, посредством $x_i=x_i(y)$ определяется как:
\[
\tag{4}
\dd x_i=\sum_k \frac{\partial x_i}{\partial y_k}\dd y_k;\qquad
\dd\delta x_i=\sum_k \frac{\partial x_i}{\partial y_k}\dd\delta y_k+
\sum_{k,l}\frac{\partial^2 x_i}{\partial y_k\partial y_l}\dd y_k\,\delta y_l;\ \ldots .
\]
Чтобы перейти к преобразованиям, индуцированным (4), можно положить в основу произвольную функцию $f(x,\dd x)=\varphi(y,\dd y)$; требуя инвариантности $\delta f,\delta^2f,\ldots$ относительно (4), получают преобразования для производных $f$ по $x$ и $\dd x$. Если, например, ограничиться формами, однородными по $\dd x$:
\[
 f(x,\dd x)=\sum a_{i_1\ldots i_n}(x)\dd x_1^{i_1}\cdots\dd x_n^{i_n}
 =\sum a'_{i_1\ldots i_n}(y)\dd y_1^{i_1}\cdots\dd y_n^{i_n}
 =\varphi(y,\dd y),
\]
то $a'$ линейны по $a$, $\frac{\partial a'}{\partial y}$ линейны по $a$ и $\frac{\partial a}{\partial x}$ и т. д.:
\[
\tag{5}
a'=A_0(a),\qquad
\frac{\partial a'}{\partial y}=A_1\!\left(a,\frac{\partial a}{\partial x}\right),\qquad
\frac{\partial^2 a'}{\partial y^2}=A_2\!\left(a,\frac{\partial a}{\partial x},\frac{\partial^2 a}{\partial x^2}\right),\ldots
\]
и речь идет об инвариантности относительно преобразований (4) и (5). При этом все встречающиеся величины
\[
a,\ \frac{\partial a}{\partial x},\ldots,\ \dd x,\ \delta x,\ \dd\delta x,\ldots,\ 
\frac{\partial x}{\partial y},\ \frac{\partial^2 x}{\partial y^2},\ldots
\]
следует рассматривать как неопределенные, поэтому детерминант каждого отдельного преобразования заведомо отличен от нуля. Только при применении формально готовой теории к специальным функциям возникает вопрос об ограничении на такие функции, чтобы преобразования в некоторой области оставались однозначно обратимыми и существовало достаточное число производных, совершенно так же, как в алгебраическом случае при численной специализации нужно предполагать, что детерминант $|s_{ik}|$ отличен от нуля.

При таком рассмотрении всех встречающихся аргументов как неопределенных речь, следовательно, идет о специальной линейной группе, непосредственно недоступной обычным алгебраическим методам, в которой преобразования $\dd x$ и $a'=A_0(a)$ совпадают с (1) и индуцированным (2). Возникает вопрос, нельзя ли, быть может в общем виде, посредством \emph{введения новых аргументов} свести преобразования (4) и (5) к общей линейной группе (1) и тем самым \emph{индуцированным} преобразованиям (2).

Это действительно так. Высшие дифференциалы $\dd^2x,\delta\dd x,\ldots$ нужно заменить лагранжевыми производными, возникающими из вариационной задачи, принадлежащей $f(x,\dd x)$, и дальнейшими комбинациями, образованными из них поляризацией (переносами по Weyl и Schouten), все из которых когредиентны с $\dd x$, так что их преобразование задается общей линейной группой (1). А $\frac{\partial a}{\partial x},\frac{\partial^2a}{\partial x^2},\ldots$ (или, более общо, производные однородно предполагаемой $f$ по $x$ и $\dd x$) нужно заменить коэффициентами некоторых форм, возникающих из "нормальных форм высших вариаций" посредством указанного выше исключения $\dd^2x,\delta\dd x,\ldots$, и их "ковариантных производных"
\[
 [\Omega_i],\quad [\Omega_i^{(1)}],\quad [\Omega_i^{(2)}],\ldots\qquad (i=1,2,\ldots);
\]
и поскольку $[\Omega_i^{(k)}]$ являются инвариантами относительно (4) и (5), зависящими только от $\dd x$ и $\delta x$, эти коэффициенты удовлетворяют алгебраически индуцированным преобразованиям (2).\footnote{E. Noether, \emph{Invarianten beliebiger Differentialausdrücke}. Göttinger Nachrichten 1918.} Посредством этой \emph{теоремы редукции} теория дифференциальных инвариантов подчиняется алгебраической теории инвариантов.

Ряд $[\Omega_i]$ обрывается на $[\Omega_\alpha]$, если $f(x,\dd x)$ является однородной формой $\alpha$-й степени по $\dd x$. Для квадратичных форм $[\Omega_2]$ совпадает с римановой формой кривизны, а указанный здесь процесс построения в точности тот, который Riemann дал в латинской конкурсной работе и который независимо от него, в связи с габилитационной лекцией Riemann, развил также Lipschitz.\footnote{Ср. также мое изложение Riemann и Lipschitz в № 28 второй части энциклопедической статьи Weitzenböck, упомянутой под 7).} Доказательство теоремы редукции, которое Christoffel и Ricci провели вычислительно для квадратичных дифференциальных форм, опирается на введение "римановых нормальных координат", преобразующих экстремали, исходящие из одной точки, в прямые; поскольку благодаря этому произвольному преобразованию $x$ соответствует линейное преобразование нормальных координат, они и обусловливают переход к общей линейной группе.

Возникает вопрос, существует ли такая теорема редукции и тогда, когда, следуя Weyl и Schouten, для определения группы в основу кладется не вариационная задача, а только перенос; то есть когда, соответственно лагранжевым производным, задаются выражения $\psi_i(\dd\delta x,\dd x,\delta x)$, посредством которых достигается исключение высших дифференциалов, причем по аналогии с квадратичными дифференциальными формами добавляется само по себе ненужное ограничение, что $\psi_i$ должны быть линейны по $\dd x$ (для произвольной однородной $f(x,\dd x)$ поляризованные лагранжевы производные являются однородными только первого порядка по $\dd x$). Требованием когредиентности $\psi$ с $\dd x$ индуцируется преобразование их коэффициентов и тем самым, по аналогии с (5), фиксируется группа. Для инвариантов низших порядков Weitzenböck вычислительно подтвердил теорему редукции в случае Weyl; общего подхода еще нет.\footnote{Ср. по этому поводу: Weyl, \emph{Raum, Zeit, Materie}; Schouten, Math. Zeitschr. 13 (1922); Weitzenböck, Sitzungsber. d. Wiener Akademie, 1920 и 1921.}

\begin{center}
(Поступило 26. 10. 22.)
\end{center}
% END UNIT BODY
% END PRODUCER UNIT 107
\endgroup


\clearpage
\begingroup
\setcounter{footnote}{0}
% BEGIN PRODUCER UNIT 108
% Source: C:/Users/Floris/Documents/interlanguage/03_projects/noether/02_slavic_working_corpus/translations/paper24/working/through-section07/russian/Noether_Paper24_Through_Section07_Russian_working.tex
% Identity: 157291 B / SHA-256 8CF1E6765A2CEE69D90D64ECCADE772950921D81E5E47D43B74C3D3FA4E9C054
\setlength{\parindent}{1.25em}
\setlength{\parskip}{0pt}
\makeatletter
\renewcommand\footnotesize{\@setfontsize\footnotesize{7.7}{8.7}\selectfont}
\makeatother
\emergencystretch=35em
\allowdisplaybreaks
\sloppy
\hbadness=10000
\vbadness=10000
\hfuzz=2pt
\vfuzz=10pt
\raggedbottom
\providecommand{\fraka}{\mathfrak a}
\providecommand{\frakb}{\mathfrak b}
\providecommand{\frakc}{\mathfrak c}
\providecommand{\frakg}{\mathfrak g}
\providecommand{\frakm}{\mathfrak m}
\providecommand{\frakn}{\mathfrak n}
\providecommand{\frako}{\mathfrak o}
\providecommand{\frakp}{\mathfrak p}
\providecommand{\frakq}{\mathfrak q}
\providecommand{\frakr}{\mathfrak r}
\providecommand{\frakS}{\mathfrak S}
\providecommand{\frakG}{\mathfrak G}
\providecommand{\frakM}{\mathfrak M}
\providecommand{\barfrakm}{\bar{\mathfrak m}}
\providecommand{\Norm}{N}
% BEGIN UNIT BODY
% Unit: Paper24 through Section01 working. Direct Russian translation from the German final-audited source slice.
\setcounter{footnote}{0}
\section*{24. Теория исключения и общая теория идеалов}
\begin{center}
Math. Ann. 90 (1923), с. 229--261
\end{center}

Далее речь идет о включении теории исключения --- в той арифметической форме, которую я придала изложению Hentzelt\footnote{Hentzelt, Kurt: Zur Theorie der Polynomideale und Resultanten. Обработано E. Noether. Math. Ann. 88 (1922), с. 53--79; цитируется как H.-N.} и которую можно назвать теорией идеалов в полиномиальной области, --- в общую теорию идеалов,\footnote{Noether, E.: Idealtheorie in Ringbereichen, Math. Ann. 83 (1921), с. 24--66; цитируется как Idealtheorie.} а одновременно о новом обосновании частей, относящихся к нулям. Основные понятия обеих теорий кратко собраны в § 1, как и понятия теории полей,\footnote{Steinitz, E.: Algebraische Theorie der Körper, J. f. M. 137 (1910), с. 167--309; цитируется как Steinitz.} так что здесь я могу на них ссылаться.

Это включение и новое обоснование группируются вокруг четырех вопросов: арифметическая формулировка понятия размерности; теория нулей; связь теорем о разложении для норм и элементарных делителей с соответствующими теоремами общей теории идеалов; характеристика простых идеалов и первичных идеалов нормой и формой элементарных делителей.

Арифметическая формулировка \emph{понятия размерности} (§ 4) дается цепями простых идеалов; это арифметический эквивалент того факта, что на поверхностях есть кривые, а на кривых есть точки. Эта арифметическая формулировка, которая будет доказана тождественной параметрическому определению теории исключения, с небольшой модификацией переносится на произвольные кольцевые области и там, вместе с переносом некоторых частей остальных вопросов, дает точное понимание строения идеалов, как я покажу в другом месте.

В H.-N. вопрос о \emph{нулях}, как обычно, прямо рассматривается для заданного идеала, а именно путем возвращения к последовательному исключению, причем характер проверки не вполне устраняется. В противоположность этому стоит путь, всегда применяемый для полиномов одной переменной: заданный полином представляют как произведение степеней простых функций (первичных функций) и для каждой отдельной простой функции строят поле нулей как поле, изоморфное полю классов вычетов. Степень простой функции дает степень этого поля; степень же первичной функции, нули которой совпадают с нулями простой функции, дает степень кольца классов вычетов по этой первичной функции. Если для каждой простой функции перейти к полю Galois, получается разложение на линейные множители; и тем самым для исходного полинома вопрос о нулях, разложении на линейные множители и кратности решен. Совершенно соответственно здесь (§ 3) система классов вычетов по простому идеалу через образование частных расширяется до поля классов вычетов, и строится изоморфное ему поле нулей. Это поле нулей содержит подполе, порожденное присоединением неопределенных, скажем $x_{i+1},\ldots,x_n$; степень нормы дает степень относительно этого подполя; вместе с тем при переходе к полю Galois получается разложение формы элементарных делителей, а значит и нормы, на линейные множители. Для соответствующего первичного идеала нули те же; разложение также уже дано, поскольку норма становится степенью формы элементарных делителей простого идеала (§ 2); степень нормы дает степень кольца классов вычетов по первичному идеалу относительно подполя, выведенного из неопределенных.

Тем самым вопрос о нулях произвольного идеала сводится к связи теорем о разложении для норм и элементарных делителей с теоремами общей теории идеалов (§ 5). Тому факту, что нули идеала составляются из нулей его отдельных ассоциированных простых идеалов, соответствует разложение нормы на наибольшие первичные множители, которые становятся равными степеням формы элементарных делителей отдельных ассоциированных простых идеалов;\footnote{Этот параллелизм теории исключения с общей теорией идеалов в кронекеровской теории исключения не выполняется; ср. H.-N., прим. *).} тем самым разложение нормы на линейные множители выполнено в общем случае. Кратность же в H.-N. получает свое истолкование посредством основных идеалов и их компонент: степень наибольшего первичного множителя нормы приравнивается числу линейно независимых классов вычетов --- после присоединения $x_{i+1},\ldots,x_n$ --- основного идеала $(i-1)$-й ступени по модулю компоненты основного идеала $i$-й ступени. Затем основной идеал $(i-1)$-й ступени доказывается тождественным изолированной компоненте разложения, однозначно определенной всеми простыми идеалами размерности выше $(n-i)$-й; а компонента основного идеала $i$-й ступени --- тождественной изолированной компоненте разложения, определенной простым идеалом, соответствующим множителю нормы, и всеми простыми идеалами большей размерности. Степень соответствующего множителя нормы --- после присоединения $x_{i+1},\ldots,x_n$ --- равна степени того подкольца классов вычетов этой компоненты, которое состоит только из классов основного идеала $(i-1)$-й ступени.

Характеристика простых идеалов и собственных первичных идеалов (§ 6) получается как прямое следствие рассмотрения полей классов вычетов. Если исходная область коэффициентов является совершенным полем, то простым идеалам как норма и форма элементарных делителей соответствуют простые функции, собственным первичным идеалам --- собственные первичные функции, и наоборот. Если область коэффициентов является несовершенным полем, можно указать только условия, которые либо необходимы, либо достаточны; как показывают примеры, норма простого идеала может стать собственно первичной, а форма элементарных делителей собственного первичного идеала --- простой функцией. Точнее, при характеристике $p$ норма простого идеала становится $p^f$-й степенью ($f>0$) его формы элементарных делителей; тогда как для собственных первичных идеалов, форма элементарных делителей которых является простой функцией, норма может стать произвольной степенью, как снова показывают примеры. Наконец, рассматриваются еще абсолютные простые идеалы (§ 7), то есть такие, которые остаются простыми идеалами при расширении области коэффициентов до алгебраически замкнутого поля. Для абсолютных простых идеалов с коэффициентами из (конечного) алгебраического числового поля характеристика приводит к переносу теоремы, впервые установленной A. Ostrowski для абсолютных простых функций: свойство быть простым идеалом сохраняется по модулю каждого простого идеала числового поля, за исключением, самое большее, конечного числа исключений.

Укажем еще на аналогию последнего вопроса с теорией идеалов в алгебраических числовых полях. Формой элементарных делителей такого идеала следует считать наименьшее целое рациональное число, делимое этим идеалом (прим. 10); поэтому для простого идеала оно всегда становится простым числом, тогда как, наоборот, идеал становится простым идеалом, как только его норма является простым числом; доказательства также идут совершенно параллельно. Приведенные выше примеры, таким образом, показывают, что над несовершенным полем возникает аналог простых идеалов высшей степени и идеалов ветвления, тогда как над совершенным полем имеются только простые идеалы первой степени и нет идеалов ветвления.

\subsection*{§ 1. Основные понятия теории исключения, общей теории идеалов и теории полей}

\textbf{1. Область.} Пусть $\frakS$ обозначает область целостности всех полиномов от $y_1,\ldots,y_n$ с коэффициентами из абстрактно заданного поля $P$. Пусть $y$ подвергаются преобразованию с неопределенными $u_{\mu\nu}$ в качестве коэффициентов:\footnote{Это преобразование, с учетом § 3 и далее, несколько более общее, чем у H.-N.; тем самым все результаты оттуда тем более сохраняются.}
\begin{align}
 y_1&=u_{11}x_1+\cdots+u_{1n}x_n,\quad \ldots,\quad
 y_n=u_{n1}x_1+\cdots+u_{nn}x_n, \tag{1}
\end{align}
с обратным преобразованием
\begin{align}
 x_1&=t_{11}y_1+\cdots+t_{n1}y_n,\quad \ldots,\quad
 x_n=t_{1n}y_1+\cdots+t_{nn}y_n, \tag{1'}
\end{align}
и пусть $u_{\mu\nu}$, или --- что ввиду взаимной рациональной выразимости равносильно --- $t_{\mu\nu}$ присоединены к $P$, так что возникает область коэффициентов $P(u)=P(t)$; через $\frakS'$ обозначим область целостности всех полиномов от $x_1,\ldots,x_n$ с коэффициентами из $P(u)$. Области $\frakS$ и $\frakS'$ лежат в основе теории исключения; рассматриваются идеалы $\fraka,\frakb,\ldots$ в $\frakS$ и $\fraka',\frakb',\ldots$ в $\frakS'$. Идеал в произвольном кольце здесь, как обычно, определяется тем, что вместе с $\alpha$ и $\beta$ он содержит также $\alpha-\beta$, а вместе с $\alpha$ также $\lambda\alpha$, где $\lambda$ понимается как произвольный элемент кольца; $\frako$ всегда обозначает единичный идеал, состоящий из всех элементов.

\textbf{2. Преобразованные идеалы.} Идеал $\frakm$ в $\frakS'$ называется преобразованным, если он возникает из идеала $\bar{\frakm}$ в $\frakS$ посредством (1) при присоединении $u_{\mu\nu}$ или $t_{\mu\nu}$. Следовательно, если $\bar f(y)$ или $\bar g(y)$ пробегают все полиномы из $\bar{\frakm}$, то $\frakm$ состоит из всех линейных комбинаций
\[
 f(x)=\sum U_i(u)\bar f_i(y)=\sum U_i(u)f_i(x)
\]
или также
\[
 g(x)=\sum T_i(t)\bar g_i(y)=\sum T_i(t)g_i(x);
\]
в частности, в нем содержатся все $f_i(x)=\bar f_i(y)$, или, что то же самое, все $g_i(x)=\bar g_i(y)$. Так как $f(x)$ соответственно $g(x)$ определены только с точностью до величин из $P(u)=P(t)$, то $U_i,T_i$ можно без ограничения общности считать степенными произведениями от $u$ соответственно $t$. Полиномы $f(x),g(x)$ снова переходят в полиномы из $\frakm$, если $U,T$ заменить любыми другими степенными произведениями, в частности при перестановке столбцов в (1) соответственно строк в (1'); следовательно, вместе с каким-либо полиномом $\frakm$ содержит и все полиномы, возникающие из него перестановкой $x$, если только в зависящих от $u$ соответственно $t$ коэффициентах произведены соответствующие перестановки. Все идеалы, встречающиеся далее, --- если прямо не сказано обратное --- следует считать преобразованными. Непреобразованные идеалы при составлении (1) с
\[
 x_i=v_{i1}z_1+\cdots+v_{in}z_n
\]
переходят в преобразованные идеалы в полиномиальной области переменных $z$ с коэффициентами из $P(u,v)$, тогда как для преобразованных идеалов это составление сводится лишь к замене $u_{\mu\nu}$ билинейными комбинациями от $u,v$.

\textbf{3. Обозначение.} Под $f^{(i)},g^{(i)},\ldots,a^{(i)},b^{(i)},\ldots$ далее всегда понимаются полиномы в $\frakS'$, свободные от $x_i,\ldots,x_{i-1}$.

\textbf{4. Основные идеалы.} Каждому идеалу $\frakm$ сопоставляются $n$ основных идеалов $\frakg_0,\frakg_1,\ldots,\frakg_{n-1}$ следующим условием: основной идеал $(i-1)$-й ступени $\frakg_{i-1}$ содержит все и только те полиномы $G(x)$, для которых существует полином $b^{(i)}\ne 0$ --- вообще говоря, меняющийся вместе с $G(x)$, --- такой, что $b^{(i)}G(x)\equiv 0\pmod{\frakm}$. Из существования идеального базиса отсюда также следует существование по крайней мере одного фиксированного для всех $G(x)$ полинома $B^{(i)}$, так что
\begin{equation}
 B^{(i)}\frakg_{i-1}\equiv0\pmod{\frakm},\qquad B^{(i)}\ne0 . \tag{2}
\end{equation}
Основные идеалы преобразованных идеалов сами становятся преобразованными идеалами (H.-N., теорема VI). Основной идеал $i$-й ступени $\frakg_i$ также определяется как идеал, в который $\frakm$ переходит при присоединении $x_{i+1},\ldots,x_n$ к $P(u)$, если ограничиться --- что тогда не означает ограничения общности --- полиномами, целыми по $x_{i+1},\ldots,x_n$; $\frakg_0$ всегда равен единичному идеалу $\frako$.

\textbf{5. Модульное представление идеалов, элементарные делители и отдельные нормы.} Если каждый полином $f(x)$ рассматривать как линейную форму от степенных произведений $\xi$ переменных $x_1,\ldots,x_{i-1}$ с полиномами $a^{(i)}$ в качестве коэффициентов, то идеал $\frakm$ переходит в модуль $\frakM_{i-1}$ линейных форм относительно области $a^{(i)}$; то есть $\frakM_{i-1}$ вместе с любыми двумя линейными формами содержит также их разность, а вместе с линейной формой $l(\xi)$ также $a^{(i)}l(\xi)$ при произвольном $a^{(i)}$. Основным модулем такого модуля $\frakM$ называется совокупность линейных форм $g(\xi)$, для которых $b^{(i)}g(\xi)\equiv0\pmod{\frakM}$, $b^{(i)}\ne0$; следовательно, основной идеал $\frakg_{i-1}$ при модульном представлении переходит в основной модуль $\frakG_{i-1}$ модуля $\frakM_{i-1}$. $\frakM_{i-1}$ и $\frakG_{i-1}$ зависят от бесконечно многих неопределенных $\xi$ так, что в каждой отдельной линейной форме встречается лишь конечное число этих неопределенных. $\frakM_{i-1}$ обладает характеристическим свойством: после присоединения $x_{i+1},\ldots,x_n$ к $P(u)$ он имеет лишь конечное число элементарных делителей, отличных от единицы; то есть после такого присоединения $\frakG_{i-1}$ и $\frakM_{i-1}$ допускают базисное представление --- где $\zeta$ обозначают новые неопределенные, связанные с $\xi$ обратимыми линейными, целыми по $x_i$ преобразованиями так, что в $\zeta_i=r_i(\xi)$ всякий раз встречается лишь конечное число этих неопределенных, ---
\begin{align}
\frakG_{i-1}&=(\zeta_0,\zeta_1,\ldots,\zeta_\sigma,\zeta_{\sigma+1},\ldots,\zeta_{\sigma+\nu},\ldots), \notag\\
\frakM_{i-1}&=(E_0^{(i)}\zeta_0,E_1^{(i)}\zeta_1,\ldots,E_\sigma^{(i)}\zeta_\sigma,
\zeta_{\sigma+1},\ldots,\zeta_{\sigma+\nu},\ldots), \tag{3}
\end{align}
где каждый $E_\mu^{(i)}$ делит предыдущий (H.-N. (24) и (28)).

Наивысший элементарный делитель $E^{(i)}$ также определяется как наибольший общий делитель --- в полиномиальном смысле --- всех $B^{(i)}$, встречающихся в (2); его можно считать целым и примитивным по $x_{i+1},\ldots,x_n$, и тогда он регулярен по $x_i$, то есть $E^{(i)}$ содержит член $x_i^\lambda$ с ненулевым коэффициентом, если имеет степень $\lambda$ по $x$.

Произведение $R^{(i)}$ всех элементарных делителей, встречающихся в (3), называется нормой $\frakG_{i-1}$ относительно $\frakM_{i-1}$, или также отдельной нормой идеала $\frakm$; в символах, с учетом того, что $\frakm$ при присоединении $x_{i+1},\ldots,x_n$ переходит в $\frakg_i$:
\begin{equation}
 R^{(i)}=E_0^{(i)}E_1^{(i)}\cdots E_\sigma^{(i)}
 =N(\frakG_{i-1}\mid\frakM_{i-1})=N(\frakg_{i-1}\mid\frakg_i).
\end{equation}
Как показывает (3), при присоединении $x_{i+1},\ldots,x_n$ величина $R^{(i)}$ становится равной детерминанту переходной подстановки от $\frakG_{i-1}$ к $\frakM_{i-1}$, а степень $R^{(i)}$ дает число линейно независимых относительно $P(u,x_{i+1},\ldots,x_n)$ классов вычетов $\frakG_{i-1}$ по модулю $\frakM_{i-1}$, следовательно $\frakg_{i-1}$ по модулю $\frakg_i$; $R^{(i)}$ можно считать целым и примитивным по $x_{i+1},\ldots,x_n$, и тогда он регулярен по $x_i$. $R^{(i)}$ можно вычислить как наибольший общий делитель --- в полиномиальном смысле --- всех $\varrho$-рядных детерминантов некоторого всегда существующего модуля $\frakM^*_{i-1}$ ранга $\varrho$, обладающего тем свойством, что система классов вычетов $\frakG^*_{i-1}/\frakM^*_{i-1}$ изоморфна $\frakG_{i-1}/\frakM_{i-1}$, где под $\frakG^*_{i-1}$ понимается основной модуль $\frakM^*_{i-1}$ (H.-N., теорема VII и § 5).

\textbf{6. Форма элементарных делителей и норма (результантная форма) идеалов.} Произведение $E_\frakm=E^{(1)}E^{(2)}\cdots E^{(n)}$ всех наивысших элементарных делителей называется формой элементарных делителей, а произведение $R_\frakm=R^{(1)}R^{(2)}\cdots R^{(n)}$ всех отдельных норм --- нормой (результантной формой)\footnote{В H.-N. употребляется только название результатантная форма; название норма соответствует арифметическим свойствам.} идеала $\frakm$. Форма элементарных делителей и норма делятся идеалом; норма делится формой элементарных делителей, а некоторая степень последней делится нормой. Более общо имеем еще:
\begin{align}
 E^{(i)}\frakg_{i-1}&\equiv0\pmod{\frakg_i},
& E^{(n)}\cdots E^{(i)}\frakg_{i-1}&\equiv0\pmod{\frakm},\notag\\
 R^{(i)}\frakg_{i-1}&\equiv0\pmod{\frakg_i},
& R^{(n)}\cdots R^{(i)}\frakg_{i-1}&\equiv0\pmod{\frakm}, \tag{4}
\end{align}
(H.-N., теорема VIII), а далее: если $\frakm$ делится $\frakn$ и нормы $R_\frakm$ и $R_\frakn$ совпадают, то совпадают и идеалы $\frakm$ и $\frakn$ (H.-N., теорема IX). Учитывая, что для делителя $\frakm$ из совпадения $R^{(1)},R^{(2)},\ldots$ следует совпадение основных идеалов $\frakg_1,\frakg_2,\ldots$, и что всегда $\frakg_0=\frako$, получаем соответственно: если $\frakn$ есть собственный делитель $\frakm$, то первая отдельная норма $\frakn$, отличная от соответствующей нормы $\frakm$, становится собственным делителем.\footnote{Зато более поздние множители $R_\frakn$ могут становиться кратными соответствующих множителей $R_\frakm$; например, всегда тогда, когда $\frakm$ является простым идеалом, а $\frakn$ не является единичным идеалом.} Если $\frakm$ содержит полином $G^{(i)}\ne0$, то по определению $\frakg_0,\frakg_1,\ldots,\frakg_{i-1}$ равны единичному идеалу; следовательно, по свойству отдельных норм, указанному в 5., а именно определять число линейно независимых классов вычетов, $R^{(1)},\ldots,R^{(i-1)}$ равны единице; значит, и $E^{(1)},\ldots,E^{(i-1)}$ равны единице.

\textbf{7. Разложение отдельных норм и элементарных делителей; компоненты основных идеалов.} Под простой функцией понимается полином, неприводимый относительно $P(u)$; под первичной функцией --- степень простой функции. Собственная первичная функция есть такая, которая не является одновременно простой функцией, то есть речь идет о степени выше первой. Разложение полинома на наибольшие первичные множители есть такое разложение, где каждый множитель является первичной функцией, но произведение любых двух множителей уже не является первичным.\footnote{Определение простой функции и первичной функции можно было бы также сформулировать в точной аналогии с 11., заменив только слово идеал словом полином. Наибольшие первичные множители являются аналогом наибольших первичных компонент разложения в 12.} Разложению отдельной нормы $R^{(i)}=N(\frakg_{i-1}\mid\frakg_i)$ на наибольшие первичные множители $Q_\nu^{(i)}$ соответствует представление $\frakg_i$ как наименьшего общего кратного
\[
 \frakg_i=[\frakr_1,\frakr_2,\ldots,\frakr_s],
\]
такое, что $Q_\nu^{(i)}\frakg_{i-1}\equiv0\pmod{\frakr_\nu}$; при этом $\frakr_\nu$ имеет $\frakg_{i-1}$ своим основным идеалом $(i-1)$-й ступени и совпадает со своим основным идеалом $i$-й ступени; $\frakr_\nu$ называются компонентами основных идеалов. Наивысший элементарный делитель $\frakr_\nu$ равен наибольшему общему делителю --- в полиномиальном смысле --- $Q_\nu^{(i)}$ и $E^{(i)}$, следовательно равен наибольшему первичному множителю $E^{(i)}$. $\frakr_\nu$ однозначно определены своим указанным выше поведением относительно основных идеалов и требованием, чтобы либо отдельная норма, либо наивысший элементарный делитель стал наибольшим первичным множителем $R^{(i)}$ соответственно $E^{(i)}$ (H.-N., теорема X).

\textbf{8. Понятие размерности.} Идеалу $\frakm$ приписывается наивысшая размерность $n-i$, если $R^{(i)}$ является первой отдельной нормой, отличной от единицы, --- или, что по 6. равносильно, если $\frakg_i$ является первым основным идеалом, отличным от единичного идеала; единичному идеалу $\frako$ приписывается размерность $-1$. Если имеется только одна отдельная норма $R^{(i)}$, отличная от единицы, то есть $\frakg_i=\frakg_{i+1}=\cdots=\frakm$, то наивысшая размерность называется также просто размерностью $\frakm$. Если $\frakm$ содержит полином $G^{(i)}\ne0$, то он имеет наивысшую размерность не больше $n-i$, как показывает замечание в конце 6. Если $\frakm$ имеет наивысшую размерность $n-i$, а $\frakn$ является делителем $\frakm$, то $\frakn$ также имеет наивысшую размерность не больше $n-i$.

\textbf{9. Норма (результантная форма) и нули.} Под нулями идеала $\frakm$ понимаются все системы значений $x_1,\ldots,x_i$ из подходящего алгебраического расширения поля $P(u;x_{i+1},\ldots,x_n)$ ($i=1,2,\ldots,n$), для которых --- при присоединении $x_{i+1},\ldots,x_n$ к области коэффициентов --- все полиномы из $\frakm$ обращаются в нуль. Каждому множителю $R^{(i)}$ нормы (результантной формы) при этом присоединении соответствует столько нулей $\frakm$, сколько указывает степень $R^{(i)}$; и $R^{(i)}$, записанная в билинейных комбинациях $w_{\mu\nu}$ от $u_{\mu\nu}$ и дальнейших неопределенных $v_{\mu\nu}$, допускает явное разложение:
\begin{equation}
 R^{(i)}=\prod_\nu (z_i-z_i^{(\nu)})
 =\prod_\nu (v_{i1}x_1+\cdots+v_{in}x_n - z_i^{(\nu)}), \tag{5}
\end{equation}
где $x_1^{(\nu)},\ldots,x_i^{(\nu)}$ каждый раз обозначает систему связанных нулей (H.-N., теорема XIII; новое обоснование будет дано в § 3). Следовательно, среди нулей идеала наивысшей размерности $n-i$ есть такие, которые зависят от $(n-i)$ параметров, но нет таких, которые зависят от большего числа параметров; тем самым показано совпадение понятия размерности, данного в 8., с обычно употребляемой формулировкой.

\textbf{10. Кольцевая область общей теории идеалов.} Основной областью пусть будет коммутативное кольцо, в котором справедлива теорема о конечной цепи: каждая цепь идеалов $\fraka_1,\fraka_2,\ldots,\fraka_k,\ldots$, где каждый $\fraka_i$ является собственным делителем --- делителем в идеальном смысле --- непосредственно предыдущего, обрывается после конечного числа шагов. Это требование эквивалентно другому, а именно что каждый идеал имеет идеальный базис (Idealtheorie, теорема I), и поэтому в частности выполняется для полиномиальной области. В следующих пунктах 11., 12., 13. предполагается именно такая общая кольцевая область.

\textbf{11. Простые идеалы и первичные идеалы.} Идеал $\frakp$ называется простым идеалом, если из делимости произведения на $\frakp$ следует делимость по крайней мере одного множителя; $\frakq$ называется первичным идеалом --- или первичным, --- если из делимости произведения на $\frakq$ следует делимость одного множителя или некоторой степени каждого множителя. В символах:
\[
 a\not\equiv0\pmod{\frakp},\quad b\not\equiv0\pmod{\frakp}\quad\hbox{влечет}\quad ab\not\equiv0\pmod{\frakp};
\]
\[
 a\not\equiv0\pmod{\frakq},\quad b^x\not\equiv0\pmod{\frakq}\ \hbox{для каждой степени }x\quad\hbox{влечет}\quad ab\not\equiv0\pmod{\frakq}.
\]
К каждому первичному идеалу $\frakq$ существует один и только один ассоциированный простой идеал $\frakp$, который является делителем $\frakq$ и некоторая степень которого делится $\frakq$, $\frakp\ne\frako$, $\frakp^\rho\equiv0\pmod{\frakq}$; при этом $\frakp$ состоит из всех элементов кольца, некоторая степень которых делится $\frakq$. Как показывает определение, каждый простой идеал одновременно является первичным идеалом; под собственными первичными идеалами понимаются такие, которые не являются одновременно простыми идеалами, то есть для которых $\rho>1$.

\textbf{12. Теорема о разложении.} Представление $\fraka=[\frakq_1,\ldots,\frakq_s]$ идеала как наименьшего общего кратного называется кратчайшим, если ни один $\frakq$ не входит в наименьшее общее кратное остальных --- в свое дополнение; оно называется представлением через наибольшие первичные компоненты, если каждый $\frakq$ первичен, но наименьшее общее кратное любых двух $\frakq$ не первично. Каждый идеал допускает кратчайшее представление через конечное число наибольших первичных компонент; для двух разных таких представлений совпадают число компонент и ассоциированные простые идеалы, все взаимно различные. Так однозначно определенные простые идеалы будут называться простыми идеалами, или ассоциированными простыми идеалами, данного идеала (Idealtheorie, теорема IX).

\textbf{13. Изолированные компоненты разложения.} Идеал $\frakr=[\frakq_1,\ldots,\frakq_i]$ называется изолированной компонентой разложения $\fraka$, если $\frakq_i$ встречаются хотя бы в одном кратчайшем представлении $\fraka$ через наибольшие первичные компоненты и удовлетворяют условию, что ни один из их ассоциированных простых идеалов не входит в один из остальных простых идеалов $\fraka$. Изолированные компоненты разложения однозначно определены своими простыми идеалами; в частности, изолированные наибольшие первичные компоненты также однозначно определены (Idealtheorie, теорема XIII).

\textbf{14. Характеристика, совершенные и несовершенные поля, расширение первого рода.} Поле $\Omega$ имеет характеристику нуль, если простое поле, содержащееся в $\Omega$ и выведенное из единицы, имеет тип поля рациональных чисел; оно имеет характеристику $p$, если это простое поле имеет тип системы классов вычетов по простому числу $p$. Поле $\Omega$ называется совершенным, если каждая простая функция с коэффициентами из $\Omega$ в подходящем поле расширения разлагается на различные линейные множители; в противном случае оно называется несовершенным. Каждое поле характеристики нуль совершенно; поле характеристики $p$ совершенно тогда и только тогда, когда вместе с каждым элементом оно содержит также его $p$-й корень. Простая функция в несовершенном поле является целой функцией степени $r$ от $x^{p^f}$ и в подходящем поле расширения разлагается на $p^f$-е степени $r$ различных линейных множителей; $f$ называется экспонентой. Простая функция называется функцией первого рода, если она в подходящем поле расширения разлагается на различные линейные множители, то есть если экспонента $f$ равна нулю; расширение называется расширением первого рода, если каждый его элемент первого рода, то есть является нулем простой функции первого рода. Каждое расширение, возникающее присоединением элемента первого рода, является расширением первого рода; над совершенными полями существуют только расширения первого рода (Steinitz, § 11 и § 13).

\textbf{15. Приведенная степень и экспонента конечного расширения.} Если $\mathcal K$ является конечным расширением несовершенного поля $\Omega$, то в $\mathcal K$ содержится подполе $\mathcal K_0$ степени $r$, которое является первого рода относительно $\Omega$ и состоит из всех и только тех элементов $\mathcal K$, которые являются первого рода. $p^f$-я степень каждого элемента из $\mathcal K$ принадлежит $\mathcal K_0$, и в $\mathcal K$ имеются элементы, которые имеют точно степень $rp^f$. $r$ называется приведенной степенью, $f$ --- экспонентой конечного расширения (Steinitz, § 14, 1 и § 14, 5).


\subsection*{§ 2. Форма элементарных делителей и норма простых идеалов и первичных идеалов. Делители простых идеалов}

С этого места речь всюду идет о полиномиальной области, определенной в § 1, 1. Сначала нужно показать, что и для простых идеалов и первичных идеалов можно ограничиваться преобразованными идеалами, согласно следующему утверждению.

\textbf{Лемма I.} Преобразованный идеал $\frakp$ простого идеала $\bar{\frakp}$ в $\frakS$ становится простым идеалом, а преобразованный идеал $\frakq$ первичного идеала $\bar{\frakq}$ в $\frakS$ становится первичным идеалом. Иными словами, свойство быть простым идеалом или первичным идеалом сохраняется при присоединении $u_{\mu\nu}$ к $P$.

В самом деле, пусть предположено, что $\fraka\not\equiv0\pmod{\frakp}$, $\frakb\not\equiv0\pmod{\frakp}$, где $\fraka$ и $\frakb$ не обязательно должны быть преобразованными идеалами; пусть, например, $a(x)\not\equiv0\pmod{\frakp}$, $b(x)\not\equiv0\pmod{\frakp}$, где $a(x)$ и $b(x)$ означают полиномы из $\fraka$ и $\frakb$. Обращением (1) по (1') получаем --- понимая под $T$ без ограничения общности степенные произведения $t_{\mu\nu}$ ---
\[
 a(x)=\sum T_i a_i(y),\qquad b(x)=\sum T_i' b_i(y),
\]
и по крайней мере один $a_i(y)$ и один $b_i(y)$ не должны делиться на $\bar{\frakp}$. Пусть, скажем, $a_h(y)\not\equiv0\pmod{\bar{\frakp}}$, $b_k(y)\not\equiv0\pmod{\bar{\frakp}}$ являются соответствующими старшими членами при некотором лексикографическом порядке на $T$; тогда также $a_h(y)b_k(y)\not\equiv0\pmod{\bar{\frakp}}$, а следовательно $a(x)b(x)\not\equiv0\pmod{\frakp}$ и поэтому $\fraka\frakb\not\equiv0\pmod{\frakp}$. Тем самым доказано, что $\frakp$ прост. Доказательство, очевидно, опирается на то, что система классов вычетов по простому идеалу образует кольцо без делителей нуля; это свойство сохраняется при присоединении неопределенных.

Соответственно пусть $\fraka\not\equiv0\pmod{\frakq}$, $\frakb^x\not\equiv0\pmod{\frakq}$ для каждого показателя $x$; тогда из существования идеального базиса следует, что должны существовать $a(x)$ из $\fraka$ и $b(x)$ из $\frakb$ такие, что $a(x)\not\equiv0\pmod{\frakq}$, $b(x)^x\not\equiv0\pmod{\frakq}$ для каждого показателя $x$: ибо при противоположном предположении некоторая степень $\frakb$ делилась бы на $\frakq$.

Положив $a(x)b(x)=c(x)$, пусть $\fraka^*,\frakb^*,\frakc^*$ означают идеалы, соответственно полученные из коэффициентов $a_i(y),b_j(y),c_k(y)$ полиномов $a(x),b(x),c(x)$. Тогда $\frakb^{*x}\not\equiv0\pmod{\bar{\frakq}}$ для каждого показателя $x$, так как иначе $b(x)^x$ делился бы на $\frakq$; ввиду $\fraka^*\not\equiv0\pmod{\bar{\frakq}}$ поэтому должно также выполняться $\fraka^*\frakb^{*\lambda}\not\equiv0\pmod{\bar{\frakq}}$ для каждого показателя $\lambda$. Но по теореме Дедекинда--Мертенса\footnote{Ср. H.-N., с. 63 и прим. 10).} существует показатель $\lambda$ такой, что $\fraka^*\frakb^{*\lambda}\equiv \frakc^*\frakb^{*\lambda-1}$; следовательно, должно быть $\frakc^*\not\equiv0\pmod{\bar{\frakq}}$. Отсюда следует $c(x)\not\equiv0\pmod{\frakq}$, а значит $\fraka\frakb\not\equiv0\pmod{\frakq}$. Тем самым доказано, что $\frakq$ первичен.

Также при представлении идеала как наименьшего общего кратного можно ограничиваться преобразованными идеалами, согласно следующему утверждению.

\textbf{Лемма II.} Из представления $\barfrakm=[\frakc_1,\ldots,\frakc_s]$ в $\frakS$ следует представление $\frakm=[\frakc_1',\ldots,\frakc_s']$ в $\frakS'$, где $\frakc_i'$ означают преобразованные идеалы $\frakc_i$. В частности, в каждом кратчайшем представлении через наибольшие первичные компоненты эти компоненты можно считать преобразованными.

В самом деле, $\frakm$ делится на $[\frakc_1',\ldots,\frakc_s']$, поскольку каждый полином $f(x)$ из $\frakm$ --- возможно, после умножения на подходящую степень $U(u)$ --- имеет вид $\sum U_i(u)f_i(y)$, где каждый $f_i(y)$ как полином из $\barfrakm$ по предположению делится на все $\frakc_i$. Обратно, если $f(x)$ делится на каждый $\frakc_i'$, то он имеет указанный выше вид и поэтому делится на $\frakm$. Следовательно, если исходить из разложения на наибольшие первичные компоненты в $\frakS$, получается разложение $\frakm=[\frakq_1',\ldots,\frakq_s']$ в $\frakS'$; здесь $\frakq_i'$ как преобразованные идеалы первичных идеалов первичны по лемме I, причем речь идет о наибольших первичных компонентах, поскольку различным ассоциированным простым идеалам $\bar{\frakp}$ в $\frakS$ соответствуют также различные $\frakp$ в $\frakS'$.

Теперь речь идет о первой характеристике простых идеалов и первичных идеалов посредством формы элементарных делителей и нормы, еще без полного разделения простого и первичного; именно в точной аналогии с теорией идеалов в алгебраических числовых полях\footnote{В качестве наивысшего элементарного делителя идеалов в алгебраических числовых полях следует рассматривать наименьшее целое рациональное число, делящееся на этот идеал, то есть в частности простое число, делящееся на простой идеал, или степень простого числа, делящуюся на первичный идеал (степень простого идеала). В самом деле, каждый идеал $\fraka^*$ одновременно является модулем $A^*$ линейных форм от элементов $\omega_1,\ldots,\omega_n$ базиса поля, а $(\omega_1,\ldots,\omega_n)$ является основным модулем $A^*$; наименьшее целое рациональное число, делящееся на $\fraka^*$, следовательно становится наивысшим элементарным делителем этого модуля $A^*$, тогда как норма $\fraka^*$ становится произведением всех элементарных делителей $A^*$.} имеет место:

\textbf{Теорема I.} Форма элементарных делителей простого идеала равна простой функции, а норма --- степени этой простой функции. Для первичных идеалов форма элементарных делителей и норма равны первичным функциям, а именно степеням одной и той же простой функции, которая сама является формой элементарных делителей ассоциированного простого идеала. Для простых идеалов и первичных идеалов наивысшая размерность совпадает с размерностью в собственном смысле; эта размерность одинакова для первичного идеала и его ассоциированного простого идеала.

Теорема I очевидно выполняется для единичного идеала, форма элементарных делителей и норма которого равны единице. Пусть теперь простой идеал $\frakp$ имеет наивысшую размерность $(n-i)$; по (4), § 1, 6. получаем
\[
 E^{(i)}\frakg_{i-1}\not\equiv0\pmod{\frakp},\quad\hbox{но}\quad
 E^{(i)}\frakg_i\equiv0\pmod{\frakp},
\]
так как иначе $\frakp$ по § 1, 8. имел бы наивысшую размерность не больше $n-i-1$. Поэтому $\frakg_i\equiv0\pmod{\frakp}$, то есть $\frakp$ совпадает со своим основным идеалом $i$-й ступени и, следовательно, также с $\frakg_{i+1},\frakg_{i+2},\ldots$. Значит, $R^{(i+1)},\ldots,R^{(n)}$ равны единице, а следовательно и $E^{(i+1)},E^{(i+2)},\ldots$ как делители $R^{(i+1)},R^{(i+2)},\ldots$. Поэтому форма элементарных делителей равна $E^{(i)}$, которое должно быть простой функцией $P^{(i)}$. Действительно, по § 1, 5. и с учетом того, что $\frakg_{i-1}$ становится единичным идеалом, $E^{(i)}$ определено как наибольший общий делитель --- в полиномиальном смысле --- всех $B^{(i)}$, делящихся на $\frakp$; из $E^{(i)}=E_1^{(i)}E_2^{(i)}\equiv0\pmod{\frakp}$ следовало бы, что уже один из множителей $E_1^{(i)}$ или $E_2^{(i)}$ должен обращаться в нуль по модулю $\frakp$. Поскольку же некоторая степень $E^{(i)}$ делится на $R^{(i)}$, $R^{(i)}$ равна степени --- возможно, первой --- этой простой функции $P^{(i)}$; одновременно $R^{(i)}$ является нормой $\frakp$. Так как встречается только одна отдельная норма, отличная от единицы, наивысшая размерность $\frakp$ в итоге совпадает с размерностью в собственном смысле.

Соответственно пусть первичный идеал $\frakq$ имеет наивысшую размерность $n-i$; тогда, как выше,
\[
 (E^{(i)}\frakg_{i-1})^x\not\equiv0\pmod{\frakq},\quad\hbox{но}\quad
 (E^{(i)}\frakg_i)^x\equiv0\pmod{\frakq}
\]
для каждого показателя $x$; следовательно, как выше, $\frakq=\frakg_i$, его форма элементарных делителей равна $E^{(i)}$, его норма равна $R^{(i)}$, а наивысшая размерность равна размерности в собственном смысле. Эта размерность совпадает с размерностью ассоциированного простого идеала; ибо из $\frakp^\rho\equiv0\pmod{\frakq}$, $\frakq\equiv0\pmod{\frakp}$ следует, что размерность $(n-i)$ идеала $\frakp$ не больше размерности $\frakq$, а размерность $\frakq$ не больше размерности $\frakp^\rho$; из $(P^{(i)})^\rho\equiv0\pmod{\frakp^\rho}$, где $P^{(i)}$ означает форму элементарных делителей $\frakp$, получается, что последняя наивысшая размерность не больше $n-i$, и потому $n-i$ является размерностью $\frakp$ и $\frakq$. Из $(P^{(i)})^\rho\equiv0\pmod{\frakq}$ далее следует, что $E^{(i)}$ --- как наибольший общий делитель всех $B^{(i)}$ из $\frakq$ --- становится степенью $P^{(i)}$, возможно также первой, и что то же верно для $R^{(i)}$. Тем самым теорема I доказана во всех частях.

Характеризация простых идеалов посредством понятия размерности дается следующим утверждением.

\textbf{Теорема II.} Простой идеал не имеет собственного делителя той же наивысшей размерности; и обратно, каждый идеал с этим свойством является простым идеалом.

Доказательство опирается на следующую лемму.

\textbf{Лемма III.} Если простой идеал $\frakp$ из полиномов с коэффициентами из поля $\Omega$ имеет только конечное число классов вычетов, линейно независимых относительно $\Omega$, то $\frakp$ не имеет собственного делителя, отличного от единичного идеала; иными словами, система классов вычетов по $\frakp$ образует поле.

Предположение означает, что существует конечное число $k$ такое, что между любыми $(k+1)$ классами вычетов имеется линейная зависимость с коэффициентами из $\Omega$ --- точнее, с коэффициентами из поля классов вычетов $(\Omega)$, представленного всеми элементами $\Omega$, --- но что существует по крайней мере одна система из $k$ классов вычетов, линейно независимых относительно $\Omega$; отсюда немедленно следует, что все классы вычетов можно линейно выразить через любую систему из $k$ линейно независимых классов. Пусть теперь $\fraka$ является собственным делителем $\frakp$, а $a(z)\not\equiv0\pmod{\frakp}$ --- полиномом из $\fraka$; тогда из $c_1f_1(z)+\cdots+c_kf_k(z)\equiv0\pmod{\frakp}$ следует также
\[
 c_1a(z)f_1(z)+\cdots+c_ka(z)f_k(z)\equiv0\pmod{\frakp};
\]
это означает, что наряду с классами вычетов $f_1,\ldots,f_k$ также классы вычетов $af_1,\ldots,af_k$ образуют систему из $k$ линейно независимых классов, через которые, в частности, можно линейно представить единичный класс. Следовательно, $\fraka$ равен единичному идеалу; иначе говоря, система классов вычетов образует поле, так как для $a(z)\not\equiv0\pmod{\frakp}$ всегда существует $f(z)$ такое, что $1\equiv a(z)f(z)\pmod{\frakp}$.

Для доказательства первой части теоремы II остается лишь заметить, что каждый простой идеал размерности $(n-i)$ --- более общо, каждый идеал наивысшей размерности $(n-i)$ --- имеет только конечное число классов вычетов, линейно независимых относительно $P(u;x_{i+1},\ldots,x_n)$, а именно (§ 1, 5.) столько, сколько указывает степень $R^{(i)}$, поскольку здесь $\frakg_{i-1}$ становится единичным идеалом. Поэтому каждый собственный делитель $\fraka$ идеала $\frakp$ при присоединении $x_{i+1},\ldots,x_n$ к $P(u)$ переходит в единичный идеал; иначе говоря, основной идеал $i$-й ступени идеала $\fraka$ становится единичным идеалом. Это означает, что наивысшая размерность $\fraka$ не больше $n-i-1$. Чтобы утверждение выполнялось также для непреобразованного идеала $\fraka$ как делителя, достаточно по § 1, 2. перейти к новой преобразованной области.

Пусть теперь, обратно, $\frakp$ не имеет собственного делителя той же наивысшей размерности $n-i$; и пусть $\fraka\not\equiv0\pmod{\frakp}$, $\frakb\not\equiv0\pmod{\frakp}$, где $\fraka$ и $\frakb$ снова предполагаются преобразованными идеалами. Тогда по предположению, поскольку $(\fraka,\frakp)$ и $(\frakb,\frakp)$ как наибольшие общие делители --- в идеальном смысле --- становятся собственными делителями $\frakp$, имеем
\[
 E^{(n)}\cdots E^{(i+1)}\frakg_i\equiv0\pmod{(\fraka,\frakp)},
\qquad
 E^{(n)}\cdots E^{(i+1)}\frakg_i\equiv0\pmod{(\frakb,\frakp)}.
\]
Отсюда получается
\[
 E^{(n)}\cdots E^{(i+1)}\frakg_i\equiv0\pmod{(\fraka\frakb,\frakp)},
\]
и поэтому необходимо $\fraka\frakb\not\equiv0\pmod{\frakp}$, так как иначе $\frakp$ имел бы меньшую наивысшую размерность, чем $n-i$. Поскольку $\fraka,\frakb$ являются преобразованными идеалами, отсюда следует, что $\frakp$, а по лемме I также $\bar{\frakp}$, является простым идеалом. Тем самым теорема II доказана.


\subsection*{§ 3. Нули простых идеалов и примарных идеалов}

Переход к непосредственному обоснованию теории нулей (§ 1, 9), сначала для простых идеалов, дает следующая лемма, представляющая расширение леммы III и верная для простых идеалов в произвольных кольцах.

\textbf{Лемма IV.} Систему классов вычетов по каждому простому идеалу, отличному от единичного идеала, можно образованием частных расширить до поля, поля классов вычетов простого идеала.

Система классов вычетов по произвольному идеалу образует кольцо, поскольку сумма, разность и произведение снова принадлежат этой системе, а ассоциативный, коммутативный и дистрибутивный законы сохраняются при переходе к классам вычетов. Если речь специально идет о простых идеалах, то это кольцо не имеет делителей нуля; действительно, определение простых идеалов (§ 1, 11) утверждает, что произведение ненулевых классов вычетов также не обращается в нуль. Если простой идеал отличен от единичного идеала, то это кольцо содержит по меньшей мере один элемент, отличный от нулевого элемента, и потому может быть расширено образованием частных -- присоединением пар элементов -- до поля;\footnote{Steinitz, § 3. Предположения о ненулевом элементе в первоначальной области целостности достаточно, поскольку тогда существование единицы обеспечивается образованием частных, и тем самым первоначальная область становится подобластью поля; Steinitz предполагает существование единицы в области целостности.} таким образом поле классов вычетов определено.

Для дальнейшего сначала зафиксируем обозначение, которое будет всюду сохраняться.

\textbf{Обозначение.} Классы вычетов вообще будем обозначать, заключая в скобки какой-либо элемент класса: $(x),\ldots,(a(x)),\ldots$; так же поля из классов вычетов будем обозначать скобками. В частности, $(P)$, $(P(u))$, $(P(u,x_{i+1},\ldots,x_n))$ обозначают поля классов вычетов, состоящие из всех и только тех классов вычетов, которые могут быть представлены элементами из $P$, $P(u)$, $P(u,x_{i+1},\ldots,x_n)$.

Сначала напомним также следующее. Поле $\mathcal R$ называется полем алгебраического ранга (степени трансцендентности) $k$ относительно основного поля $\Omega$, если $\mathcal R$ содержит по меньшей мере одну систему из $k$ величин, алгебраически независимых относительно $\Omega$, -- трансцендентных величин, -- но любые $k+1$ величин из $\mathcal R$ алгебраически зависимы относительно $\Omega$. Если $\mathcal R$ можно представить как алгебраическое расширение подполя $\Omega(z_1,\ldots,z_s)$, где $z$ алгебраически независимы -- трансцендентны -- относительно $\Omega$, то необходимо $s=k$;\footnote{Доказательство этого известного факта, верное при произвольной характеристике, см. у Steinitz, § 22, 9. -- При присоединении $u$ алгебраический ранг не может возрасти, поскольку речь идет о рациональных комбинациях первоначальных элементов с коэффициентами из $\Omega(u)$; он не может и убыть, поскольку всякое соотношение между элементами из $\mathcal R$ с коэффициентами из $\Omega(u)$ распадается на конечное число соотношений с коэффициентами из $\Omega$.} точно так же алгебраический ранг $\mathcal R(u)$ относительно $\Omega(u)$ равен $k$, если $u$ означает неопределенные, не содержащиеся в $\mathcal R$.

Построение поля нулей теперь идет в полной аналогии с обычным путем для простых функций одной неопределенной\footnote{Steinitz, § 6.}, согласно следующей теореме.

\textbf{Теорема III.} Полю классов вычетов $(\mathcal R)$ простого идеала $\frakp$ размерности $n-i$ можно изоморфно поставить в соответствие поле нулей $R$ идеала $\frakp$, которое имеет алгебраический ранг $n-i$ -- зависит от $n-i$ параметров -- и, в частности, может рассматриваться как конечное расширение поля $P(u,x_{i+1},\ldots,x_n)$. Изоморфизм между $(\mathcal R)$ и $R$ охватывает изоморфизм между $(P)$ и $P$, между $(P(u))$ и $P(u)$, а при принятии $x_{i+1},\ldots,x_n$ за параметры также между $(P(u,x_{i+1},\ldots,x_n))$ и $P(u,x_{i+1},\ldots,x_n)$. Обратно, всякое поле нулей алгебраического ранга $n-i$ -- зависящее от $n-i$ параметров -- изоморфно полю классов вычетов $(\mathcal R)$.

Из теоремы III как следствие получается известная теорема: если идеал $\fraka$ обращается в нуль во всех нулях простого идеала $\frakp$, то $\fraka$ делится на $\frakp$.

Для доказательства прежде всего заметим, что поле классов вычетов $(\mathcal R)$ имеет алгебраический ранг $n-i$ относительно $(P(u))$. Действительно, классы $(x_{i+1}),\ldots,(x_n)$ алгебраически независимы относительно $(P(u))$, поскольку $\frakp$, будучи размерности $n-i$, не содержит полинома, свободного от $x_1,\ldots,x_i$; всякий дальнейший класс, однако, зависит от них. Ибо из принадлежности формы элементарных делителей $P^{(i)}$ к $\frakp$ по § 1, 2 также для $\lambda=1,2,\ldots,i$ следует принадлежность полиномов, соответственно зависящих от $x_\lambda,x_{i+1},\ldots,x_n$. Поэтому $(\mathcal R)$ становится конечным расширением поля $(P(u,x_{i+1},\ldots,x_n))$; степень относительно этого подполя по § 1, 5 -- с учетом того, что $\frakg_{i-1}$ равен единичному идеалу, а $\frakg_i$ по теореме I становится равным $\frakp$ -- равна степени нормы.

Чтобы перейти к изоморфному полю нулей, нужно далее заметить, что $P(u,x_{i+1},\ldots,x_n)$ и $(P(u,x_{i+1},\ldots,x_n))$ изоморфно соотнесены тем, что каждому элементу из $P(u,x_{i+1},\ldots,x_n)$ сопоставляется класс вычетов, представленный этим элементом. Действительно, сначала при этом сопоставлении области целостности, состоящие из полиномов в $u,x_{i+1},\ldots,x_n$ соответственно в $(u),(x_{i+1}),\ldots,(x_n)$, соответствуют друг другу изоморфно, поскольку $\frakp$ не содержит полинома, свободного от $x_1,\ldots,x_i$, а значит различным полиномам из $P(u,x_{i+1},\ldots,x_n)$ соответствуют и различные классы; из изоморфных областей целостности, однако, образованием частных возникают изоморфные поля. Изоморфное отношение между $P(u,x_{i+1},\ldots,x_n)$ и $(P(u,x_{i+1},\ldots,x_n))$ охватывает отношение между $P$ и $(P)$, между $P(u)$ и $(P(u))$. Чтобы расширить этот изоморфизм до изоморфизма между $(\mathcal R)$ и полем нулей $R$, пусть классам $(x_1),\ldots,(x_i)$ изоморфно сопоставлены некоторые новые элементы $\xi_1,\ldots,\xi_i$; то есть каждому полиному $(a(x))$ соответствует полином $a(\xi_1,\ldots,\xi_i,x_{i+1},\ldots,x_n)$, а сумме и произведению соответствуют сумма и произведение. Образованием частных -- причем можно ограничиться знаменателями из подполя $(P(u,x_{i+1},\ldots,x_n))$ соответственно $P(u,x_{i+1},\ldots,x_n)$ -- полю классов вычетов $(\mathcal R)$ тогда соответствует поле $R$, которое становится конечным расширением $P(u,x_{i+1},\ldots,x_n)$ степени нормы и которое можно назвать полем нулей; ибо это сопоставление утверждает, что $f(\xi_1,\ldots,\xi_i,x_{i+1},\ldots,x_n)$ обращается в нуль, как только $f(x)\equiv0\pmod{\frakp}$. Вместо $(P(u,x_{i+1},\ldots,x_n))$ во всем рассуждении могло бы фигурировать и любое другое подполе $(\mathcal R)$, возникающее присоединением $n-i$ алгебраически независимых величин.

Обратно, легко показать, что всякое поле нулей $R$ алгебраического ранга $n-i$ становится изоморфным полю классов вычетов. Поскольку $R$ рационально, с коэффициентами из $P(u)$, выводится из величин $\xi_1,\ldots,\xi_i$, среди этих величин должны содержаться $n-i$ алгебраически независимых, следовательно трансцендентных относительно $P(u)$, элементов. В частности, это должно выполняться для $\xi_{i+1},\ldots,\xi_n$, поскольку из $P^{(i)}\equiv0\pmod{\frakp}$, как выше, следует, что $R$ становится алгебраическим расширением поля $P(u,\xi_{i+1},\ldots,\xi_n)$. Поскольку по предположению $f(\xi)$ обращается в нуль для всякого полинома $f(x)$ из $\frakp$, каждому классу полиномиальной области элементов $(x)$, содержащейся в $(\mathcal R)$, соответствует один и только один элемент из $R$, являющийся полиномом от $\xi$; и сумме и произведению соответствуют сумма и произведение. Но различным классам соответствуют и различные элементы из $R$; иначе некоторому ненулевому классу, а значит после подходящего умножения также классу из $(P(u,x_{i+1},\ldots,x_n))$, представленному полиномом, соответствовал бы нулевой элемент в $R$, и между величинами $\xi_{i+1},\ldots,\xi_n$ вопреки предположению существовала бы алгебраическая зависимость. Это сопоставление исчерпывает все полиномы от $\xi$, содержащиеся в $R$; а образованием частных -- причем снова можно ограничиться знаменателями из $(P(u,x_{i+1},\ldots,x_n))$ соответственно $P(u,x_{i+1},\ldots,x_n)$ -- из этих изоморфных областей целостности возникают изоморфные поля.

Пусть теперь идеал $\fraka$ обращается в нуль во всех нулях $\frakp$, следовательно в частности и во всех нулях, зависящих от $n-i$ параметров. Если тогда $a(x)$ есть произвольный полином из $\fraka$, то $a(\xi)$ обращается в нуль; вследствие изоморфизма между $R$ и $(\mathcal R)$ класс $(a(x))$ тем самым становится нулевым классом, а $a(x)$ и, значит, $\fraka$ делятся на $\frakp$. Этим доказаны теорема III и следствие.

Чтобы перейти от теоремы III к разложению, данному в § 1, 9, сначала для простых идеалов, и чтобы иметь возможность делать более точные утверждения об экспонентах, требуется дальнейшее исследование формы элементарных делителей, согласно следующей лемме.

\textbf{Лемма V.} Если при характеристике $p$ форма элементарных делителей $E^{(i)}$ простого идеала или примарного идеала -- более общо, идеала $\fraka$, для которого наивысшая размерность совпадает с размерностью в собственном смысле, -- имеет экспоненту $f$ относительно $x_i$, то есть является целой функцией от $x_i^{p^f}$, то она имеет также экспоненту $f$ относительно $x_{i+1},\ldots,x_n$ и относительно неопределенных $u_{\mu\nu}$ соответственно $t_{\mu\nu}$, входящих в $E^{(i)}$. В частности, когда областью коэффициентов является совершенное поле, форма элементарных делителей $P^{(i)}$ простого идеала всегда имеет экспоненту нуль относительно $x_i$, то есть является простой функцией первого рода.

Для этого надо заметить, что $P(u,x_{i+1},\ldots,x_n)$ становится несовершенным, как только $P$ является совершенным полем характеристики $p$, так что из § 1, 14 не следует прямо, что $P^{(i)}$ становится функцией первого рода. Пусть теперь $E^{(i)}$ имеет экспоненту $f$ относительно $x_i$, но содержит другую неопределенную $x_{i+1}$ с экспонентой $f'$. Тогда -- перестановкой $u_{\mu\nu}$ соответственно $t_{\mu\nu}$ -- по § 1, 2 существует также ненулевой полином $G^{(i)}$, делящийся на идеал, который является целой функцией от $x_{i+1}^{p^{f'}}$ и который по определению формы элементарных делителей делится на $E^{(i)}$. Но поскольку $G^{(i)}$ имеет ту же степень, что и $E^{(i)}$, он должен совпадать с $E^{(i)}$ с точностью до множителя из $P(u)$, и потому необходимо $f'=f$. Следовательно, $E^{(i)}$ -- которую можно считать целой и примитивной по $t_{\mu\nu}$ -- имеет вид
\[
 E^{(i)}=\sum c_\rho(t)\,x_i^{\alpha_\rho p^f}x_{i+1}^{\beta_\rho p^f}\cdots x_n^{\gamma_\rho p^f}
       =\sum c_\rho(t)\,X_\rho^{p^f},
\]
где $X_\rho$ означают степенные произведения от $x$, и $X_\rho$ делятся на $\fraka$. Поэтому получается полином, также принадлежащий $\fraka$, если в $X_\rho$ каждый показатель отдельных $t$ заменить на наибольшее кратное $p^f$, содержащееся в нем; как видно из приведенного представления, это сводится к тому, чтобы в $c_\rho(t)$ каждый показатель $t$ заменить на наибольшее кратное $p^f$, содержащееся в нем. При этой подстановке $E^{(i)}$ переходит в полином от $x_i^{p^f},\ldots,x_n^{p^f}$, который по определению делится на $E^{(i)}$, причем из-за предположенной примитивности также относительно $t$; значит, он должен быть тождественным с $E^{(i)}$, поскольку степень не возросла ни в одном аргументе. В частности, если $P$ является совершенным полем, то $E^{(i)}$ как целая функция от $x_i^{p^f}$ становится $p^f$-й степенью; если же речь идет о простом идеале, то необходимо $f=0$, и форма элементарных делителей $P^{(i)}$ становится простой функцией первого рода.

Разложение теперь дает следующая теорема.

\textbf{Теорема IV.} Форма элементарных делителей $P^{(i)}$ простого идеала допускает в поле Галуа, выведенном из поля нулей, явное разложение:
\begin{align*}
 P^{(i)}
 &=\prod_\nu \bigl(x_i-t_{i+1}x_{i+1}^{(\nu)}-\cdots-t_nx_n^{(\nu)}\bigr)^e\\
 &=\prod_\nu \bigl(t_i(y_i-y_i^{(\nu)})+\cdots+t_n(y_n-y_n^{(\nu)})\bigr)^e,
\end{align*}
где $y_1^{(\nu)},\ldots,y_n^{(\nu)}$ каждый раз задают соответствующую систему нулей первоначальных неопределенных $y$, независимую от $t_i,\ldots,t_n$, различным $\nu$ соответствуют различные линейные множители, а экспонента $e$ для совершенного поля $P$ имеет значение один, для несовершенного -- значение $p^f$ ($f>0$). Тем самым при присоединении новых неопределенных $v$ также определяется разложение формы элементарных делителей, записанной в билинейных комбинациях $u$ и $v$ вместо $u_{\mu\nu}$:
\begin{align*}
 P^{(i)}(w)
 &=\prod_\nu \bigl(z_i-z_i^{(\nu)}\bigr)^e\\
 &=\prod_\nu \bigl(v_{i1}(y_1-y_1^{(\nu)})+\cdots+v_{in}(y_n-y_n^{(\nu)})\bigr)^e,
\end{align*}
где $e$ имеет то же значение, что и выше. Так же этим дано разложение нормы простого идеала или примарного идеала с $\frakp$ как ассоциированным простым идеалом -- как степени $P^{(i)}$.

Как следствие из теоремы IV получается еще: если два простых идеала совпадают по своей форме элементарных делителей $P^{(i)}$, то они тождественны.

Доказательство теоремы IV сведется к установлению, что форма элементарных делителей $P^{(i)}$ тождественна фундаментальному уравнению поля расширения $(P(u,x_{i+1},\ldots,x_n),(x_i))$ над $(P(u,x_{i+1},\ldots,x_n))$. Сначала -- чтобы получить представление о зависимости от неопределенных $u_{\mu\nu}$ соответственно $t_{\mu\nu}$ -- перейдем к полю классов вычетов $(\mathcal S)$ идеала $\frakp$, которое по сказанному при определении алгебраического ранга должно иметь алгебраический ранг $n-i$ относительно $(P)$, где $(P)$ означает подполе, представленное элементами из $P$. Ведь $(\mathcal S)$ возникает присоединением неопределенных $u$ соответственно $t$ к подполю, изоморфному некоторому $(\mathcal R)$, которое получается образованием частных из всех и только тех классов, которые могут быть представлены полиномами из $\frakS$, причем $(P)$ и $(P)$ соответствуют друг другу. Поэтому алгебраический ранг $(\mathcal S)$ относительно $(P(u))$ равен рангу этого подполя относительно $(P)$, и значит, по изоморфизму, равен рангу $(\mathcal R)$ относительно $(P)$. Поскольку $(\mathcal S)$ рационально, с коэффициентами из $(P)$, выводится из классов $(y_1),\ldots,(y_n)$, среди этих классов должны быть $n-i$ алгебраически независимых, и $(\mathcal S)$ возникает как конечное алгебраическое расширение подполя, порожденного присоединением этих $n-i$ классов к $(P)$.\footnote{Это замечание показывает, что теорему III со следствием можно было бы прямо сформулировать и доказать для поля нулей $R$ идеала $\frakp$. Но только переход к преобразованному идеалу обеспечивает, что в наибольших примарных компонентах одинаковой размерности произвольного идеала при нулях появляются одни и те же параметры; и только тем самым получается связь с нормой и формой элементарных делителей произвольного идеала.}

Если присоединить только неопределенные $t_{\mu\nu}$, входящие в $x_{i+1},\ldots,x_n$, то снова возникает поле классов вычетов, которое содержит классы $(y_1),\ldots,(y_n)$ и $(x_{i+1}),\ldots,(x_n)$ и становится конечным алгебраическим расширением поля $(P(t_{\mu\nu},x_{i+1},\ldots,x_n))$; при этом речь идет о полях, изоморфных подполям $(\mathcal S)$, а не тождественных с ними, поскольку при различных неопределенных классы вычетов хотя и представлены одними и теми же элементами, но не тождественны, так как не содержат тех же элементов. Значит, в зависимость классов $(y)$ от $(P(t_{\mu\nu},x_{i+1},\ldots,x_n))$ входят только неопределенные $t_{\mu\nu}$. Теперь образуем, присоединив $t_{i1},\ldots,t_{in}$, класс $(x_i)=(t_{i1}y_1+\cdots+t_{in}y_n)$; и пусть $(x_{i\nu})=(t_{i1}y_{1\nu}+\cdots+t_{in}y_{n\nu})$ будут $r$ различных сопряженных значений, лежащих в поле Галуа относительно $(P(t_{\mu\nu},x_{i+1},\ldots,x_n))$. Тогда -- смотря по тому, идет ли речь о расширениях первого рода или нет, -- произведение по множителям $x_i-(x_{i\nu})$ или по их $p^f$-м степеням является простой функцией $(G)$ относительно $(P(t_{\mu\nu},x_{i+1},\ldots,x_n))$; ибо, поскольку существуют элементы степени $r\cdot p^f$ (§ 1, 15), необходимо, чтобы $(x_i)$ был таким элементом, так как всякий элемент возникает специализацией $t_{\mu\nu}$; но $(x_i)$ является нулем $(G)$. Переходя к изоморфному полю нулей идеала $\frakp$ и к выведенному из него полю Галуа, получаем, что $G$ допускает разложение на линейные множители $x_i-t_{i1}y_{1\nu}-\cdots-t_{in}y_{n\nu}$ соответственно на их $p^f$-е степени. Но поскольку $(G)$ обращается в нуль при $x_i=(x_i)$, то $G(x_i)$ делится на $\frakp$ и тем самым на $P^{(i)}$; а как простая функция относительно $P(t_{\mu\nu},x_{i+1},\ldots,x_n)$ и как примитивная относительно $x_i$, $G$ должна совпадать с $P^{(i)}$. Тем самым первое разложение теоремы IV доказано.

Чтобы перейти ко второму разложению, надо заметить, что свойство быть или не быть расширением первого рода сохраняется при присоединении неопределенных. Поэтому после присоединения $v_1,\ldots,v_n$ и всех $t_{\mu\nu}$ образуем класс $(z)=(v_1x_1+\cdots+v_ix_i)$ и сопряженные классы $(z_\nu)=(v_1x_{1\nu}+\cdots+v_ix_{i\nu})$. Тогда произведение по всем $z-(z_\nu)$ соответственно по их $p^f$-м степеням снова является простой функцией $(H)$, обращающейся в нуль при $z=(z)$. Следовательно, $H$, как и выше, представима в виде произведения линейных множителей и тождественна форме элементарных делителей простого идеала, выведенного из $\frakp$ композицией преобразования (1) с $z=v_1x_1+\cdots+v_nx_n$; эта форма элементарных делителей, однако, -- из-за взаимной специализации -- получается из $P^{(i)}$ заменой $u_{\mu\nu}$ некоторыми билинейными комбинациями $u$ и $v$.\footnote{H.-N., § 7.} С этими разложениями дано также разложение нормы как степени $P^{(i)}$,\footnote{Для расширений первого рода, следовательно в частности над совершенными полями, речь идет о первой степени; над несовершенными полями -- о $p^f$-й степени, как будет показано в § 6, теоремах XIII и XIV.} а также разложение нормы и формы элементарных делителей примарного идеала с $\frakp$ как ассоциированным простым идеалом.

Наконец, если два простых идеала совпадают по форме элементарных делителей $P^{(i)}$, то вследствие приведенного разложения они совпадают по своим нулям; значит, они взаимно делятся друг на друга и потому тождественны.


\subsection*{§ 4. Арифметическая формулировка понятия размерности.}

Первая формулировка понятия размерности, независимая от введения преобразованных идеалов, дается следующей теоремой.

\textbf{Теорема V.} \emph{Размерность простого идеала задается алгебраическим рангом его поля классов вычетов. Наивысшая размерность идеала равна наибольшему из чисел размерности его ассоциированных простых идеалов.}

Первая часть теоремы V была доказана в первом замечании к теореме IV, где было показано, что алгебраический ранг поля классов вычетов $(\mathcal R)$ идеала $\frakp$ равен размерности $\frakp$, в согласии с § 1, 8. Для доказательства второй части пусть $\frakm=[\frakq_1,\ldots,\frakq_a]$ будет кратчайшим представлением через наибольшие первичные компоненты; тогда прежде всего ни один $\frakq$, как делитель $\frakm$, не может иметь более высокой размерности, чем $\frakm$. Но произведение всех $\frakq$ делится на $\frakm$, а следовательно, и произведение всех форм элементарных делителей отдельных $\frakq$; если поэтому $n-i$ есть наибольшее из чисел размерности $\frakq$, то $\frakm$ содержит полином $G^{(i)}\ne0$, а значит, не может иметь размерность выше $n-i$. Числа размерности $\frakq$ по теореме I совпадают с числами размерности их ассоциированных простых идеалов. Если без введения преобразованных идеалов определить размерность простого идеала через алгебраический ранг поля классов вычетов, то вторая часть теоремы V дает определение наивысшей размерности произвольного идеала.

Чтобы, с другой стороны, охватить понятие размерности цепочками простых идеалов, что снова не зависит от введения преобразованных идеалов, нужно в дополнение к теореме II показать следующее.

\textbf{Теорема VI.} \emph{Каждый простой идеал, отличный от единичного идеала, имеет --- после возможного присоединения неопределенных --- по крайней мере один простой идеал в качестве делителя, размерность которого уменьшилась ровно на единицу.}

Действительно, если $\frakp$ имеет размерность $(n-i)>0$ и $v$ обозначает новую неопределенную, то
\[
 \frakc=(\frakp,x_i-v)
\]
является идеалом требуемого свойства. Это следует из изоморфного сопоставления; имеем $\frakc=(\frakp^*,x_i-v)$, где $\frakp^*$ получается из $\frakp$ тем, что $x_i$ заменяется неопределенной $v$, а $v$ присоединяется к области коэффициентов; точно так же классы вычетов получаются заменой класса $(x_i)$ классом $(v)$. Но $\frakp$ переходит в себя, когда $x_i$ присоединяется к $P(u)$; из
\[
 h(x_i)f(x)\equiv0\pmod{\frakp},\qquad h(x_i)\ne0
\]
необходимо следует $f(x)\equiv0\pmod{\frakp}$; ведь если бы $\frakp$ содержал полином, зависящий только от $x_i$, то он содержал бы и полином, зависящий только от $x_n$, а потому вопреки предположению имел бы размерность не выше нуля. При соответствии $v\sim x_i$ нулевому классу по $\frakc$ поэтому необходимо соответствует нулевой класс по $\frakp$, и, конечно, верно также обратное. Следовательно, получается изоморфное сопоставление между системой классов вычетов по $\frakp$ и по $\frakc$; последняя не имеет делителей нуля, так что $\frakc$ становится простым идеалом. В силу этого изоморфного сопоставления алгебраической зависимости между $(x_i)$ и $(P(u,x_{i+1},\ldots,x_n))$ соответствует соотношение между $(x_{i+1}),\ldots,(x_n)$ относительно $(P(u,v))$, тогда как для $\lambda=1,\ldots,i-1$ зависимости между $(x_\lambda)$ и $(P(u,x_{i+1},\ldots,x_n))$ сохраняются как таковые и потому дальше не понижают алгебраический ранг; алгебраический ранг уменьшился ровно на единицу. Если $\frakp$ имел размерность нуль, то $\frakc$ становится единичным идеалом; таким образом, приведенный результат сохраняется. Поэтому в исходной области $\bar{\frakS}$
\[
 \bar{\fraka}=(\bar{\frakp},v+v_1y_1+\cdots+v_ny_n),
\]
где $v_\lambda$ положено равным $-t_{i\lambda}$, является идеалом требуемого свойства. Если, в частности, $P$ содержит бесконечно много элементов, то $v,v_1,\ldots,v_n$ всегда можно специализировать до элементов из $P$ так, чтобы норма и форма элементарных делителей, а следовательно, и размерность $\bar{\fraka}$ стали равны соответствующим данным специализированного идеала $\frakb$;\footnote{Ср. H.-N., последний абзац § 7.} правда, форма элементарных делителей при этом не обязана оставаться простой функцией. Однако по теореме V идеал $\frakb$ имеет по крайней мере один ассоциированный простой идеал требуемой размерности; возвращаясь снова к $\bar{\frakS}$, получаем, что здесь теорема VI верна без присоединения каких-либо неопределенных.

Теоремы II и VI непосредственно дают желаемую формулировку в виде следующей теоремы.

\textbf{Теорема VII.} \emph{Простой идеал $\frakp$ имеет размерность $n-i$, если --- после возможного присоединения неопределенных --- существует по крайней мере одна цепочка из $1+(n-i)+1$ простых идеалов}
\[
 \frakp_0=\frakp;\quad \frakp_1,\ldots,\frakp_{n-i};\quad \frakp_{n-i+1},
\]
\emph{каждый из которых является собственным делителем непосредственно предшествующего, но не существует такой цепочки с большим числом членов. Это определение верно также в исходной области, причем без введения каких-либо неопределенных, если $P$ содержит бесконечно много элементов.}\footnote{Это цепочное определение размерности в случае алгебраического числового поля дает размерность нуль для обычных простых идеалов, размерность $-1$ для единичного идеала и размерность $1$ для нулевого идеала; действительно, последний также удовлетворяет определению простых идеалов --- в поле произведение ненулевых величин всегда отлично от нуля. Теорема II при этом определении размерности сохраняется в алгебраическом числовом поле.}

Прежде всего ясно, что последний элемент цепочки должен быть равен единичному идеалу --- который удовлетворяет определению простого идеала, --- так как иначе цепочку можно было бы продолжить добавлением $\frako$; для самого $\frako$ по теореме VII получается размерность $-1$, в согласии с § 1, 8. Пусть теперь $\frakp$ отличен от $\frako$ и имеет размерность $n-i$; тогда число членов цепочки может быть не больше $1+(n-i)+1$, поскольку по теореме II не могут встречаться два простых идеала одинаковой размерности. Но по теореме VI --- после возможного присоединения новых неопределенных --- существует по крайней мере одна цепочка с таким числом членов, так как размерность на каждом члене цепочки уменьшилась ровно на единицу. Если положить теорему VII в основу как определение размерности --- определение, которое, очевидно, точно так же можно сформулировать в исходной области и при котором по теореме VI введение неопределенных становится ненужным, когда $P$ содержит бесконечно много элементов, --- то определение наивысшей размерности произвольного идеала снова дается второй частью теоремы V.


\subsection*{§ 5. Включение основных идеалов и их компонент в общую теорию идеалов. Нули идеалов.}

Речь идет о характеристике основных идеалов и их компонент (§ 1, 7.) через изолированные компоненты разложения (§ 1, 13.) в смысле общей теории идеалов. Эта характеристика покажет параллельность теорем о разложении норм теоремам общей теории идеалов и позволит сформулировать теорию нулей, данную в § 3 для простых и первичных идеалов, для произвольных идеалов. Для основных идеалов справедливо следующее.

\textbf{Теорема VIII.} \emph{Основные идеалы $i$-й ступени $\frakg_i$ являются изолированными компонентами разложения, однозначно определенными совокупностью ассоциированных простых идеалов размерностей $n-1,n-2,\ldots,n-i$. Если все ассоциированные простые идеалы имеют одну и ту же размерность, то и идеал имеет эту размерность; то есть в норме появляется только один множитель $R^{(i)}$. В общем случае норма имеет столько отличных от единицы множителей $R^{(i)},R^{(i+\sigma)},R^{(i+\tau)},\ldots$, сколько различных чисел размерности встречается среди ассоциированных простых идеалов.}

Доказательству предпошлем следующую лемму.

\textbf{Лемма VI.} \emph{Если идеал представлен как наименьшее общее кратное, $\frakm=[\fraka_1,\ldots,\fraka_a]$, то его основной идеал $\frakg_i$ является наименьшим общим кратным основных идеалов $i$-й ступени $\mathfrak h_i$ идеалов $\fraka$.}

Действительно, из
\[
 b^{(i+1)}f\equiv0\pmod{\frakm},\qquad b^{(i+1)}\ne0
\]
следует
\[
 b^{(i+1)}f\equiv0\pmod{\fraka_\lambda},\qquad b^{(i+1)}\ne0;
\]
следовательно, $\frakg_i$ делится на $[\mathfrak h_{i1},\ldots,\mathfrak h_{ia}]$. Обратно, из
\[
 c_\lambda^{(i+1)}f\equiv0\pmod{\fraka_\lambda},\qquad c_\lambda^{(i+1)}\ne0
\]
также следует
\[
 c_1^{(i+1)}\cdots c_a^{(i+1)}f\equiv0\pmod{\frakm},\qquad
 c_1^{(i+1)}\cdots c_a^{(i+1)}\ne0,
\]
а отсюда обратная делимость, чем лемма доказана.

Пусть теперь ассоциированные простые идеалы $\frakm$ упорядочены по размерности (при этом, конечно, некоторые размерности могут и отсутствовать, $j_\lambda=0$):
\[
 \frakp_{1,1},\ldots,\frakp_{1,j_1};\quad
 \frakp_{2,1},\ldots,\frakp_{2,j_2};\quad\ldots;\quad
 \frakp_{n,1},\ldots,\frakp_{n,j_n},
\]
где вообще $\frakp_{\lambda,\kappa}$ обозначает простой идеал размерности $n-\lambda$; пусть $\frakq_{\lambda,\kappa}$ будет соответствующий первичный идеал, входящий в некоторое фиксированное данное разложение $\frakm$. Тогда основной идеал $i$-й ступени по определению равен единичному идеалу для всех $\frakq_{\lambda,\kappa}$, для которых $\lambda>i$; для $\lambda\le i$ же по теореме I этот основной идеал равен $\frakq_{\lambda,\kappa}$. Поэтому по лемме VI получаем:
\[
 \tag{5}
 \frakg_i=[\frakq_{1,1},\ldots,\frakq_{1,j_1};\ldots;\frakq_{i,1},\ldots,\frakq_{i,j_i}],
\]
и это представление распознает $\frakg_i$ как изолированную компоненту разложения, поскольку ни один из простых идеалов, ассоциированных с $\frakg_i$, не может входить в один из остальных простых идеалов, все из которых имеют меньшую размерность.

Если специально среди ассоциированных простых идеалов встречаются только идеалы размерности $n-i$, то по (5) $\frakg_0,\ldots,\frakg_{i-1}$ равны единичному идеалу, а $\frakg_i=\frakg_{i+1}=\cdots=\frakm$; следовательно, $R_\frakm=R^{(i)}$. Если же имеются простые идеалы различных размерностей, например $j_i,j_{i+\sigma},j_{i+\tau},\ldots$ отличны от нуля, то по (5):
\[
 \frakg_0\ne\frakg_i\ne\frakg_{i+\sigma}\ne\frakg_{i+\tau}\ne\cdots,
\]
но
\[
 \frakg_0=\frakg_1=\cdots=\frakg_{i-1}=\frako,
 \qquad
 \frakg_i=\frakg_{i+1}=\cdots=\frakg_{i+\sigma-1},\ldots,
\]
и потому множители $R^{(i)},R^{(i+\sigma)},R^{(i+\tau)}$ и только они отличны от единицы.\footnote{Из теоремы VIII снова следует, с учетом леммы II, что основные идеалы --- как наименьшие общие кратные преобразованных идеалов --- становятся преобразованными идеалами. Так как теорема VIII опирается только на теоремы из H.-N., где этот факт не используется, тем самым получается новое доказательство, одновременно раскрывающее внутреннюю причину. В H.-N. эта теорема используется только в теории нулей, которая здесь развита заново.}

Посредством теоремы VIII теорему I можно усилить следующим образом.

\textbf{Теорема IX.} \emph{Идеал является первичным тогда и только тогда, когда его форма элементарных делителей --- и, следовательно, его норма --- является первичной функцией.}

То, что условие выполнено для первичных идеалов, было показано в теореме I. Пусть теперь, наоборот, форма элементарных делителей $\frakm$ задана как $Q^{(i)}=(P^{(i)})^\rho$, где $P^{(i)}$ означает простую функцию. Тогда по теореме VIII все простые идеалы $\frakm$ имеют одну и ту же размерность $n-i$. Из $\frakm=[\frakq_1,\ldots,\frakq_a]$ получается:
\[
 Q^{(i)}\equiv0\pmod{\frakq_\lambda};\quad\hbox{следовательно}\quad
 (P^{(i)})^\rho\equiv0\pmod{\frakp_\lambda}\quad\hbox{и потому}\quad
 P^{(i)}\equiv0\pmod{\frakp_\lambda}.
\]
Так как все $\frakp_\lambda$ имеют размерность $n-i$, а $P^{(i)}$ означает простую функцию, то $P^{(i)}$ становится формой элементарных делителей каждого $\frakp_\lambda$. Следовательно, по следствию к теореме IV все $\frakp_\lambda$ совпадают; поэтому, поскольку речь идет о наибольших первичных компонентах, может появиться только один $\frakp$; $\frakm$ становится первичным, причем при $\rho=1$ специальный случай простого идеала не исключается. Теоремы VIII и IX приводят к включению компонент основных идеалов в следующей форме.

\textbf{Теорема X.} \emph{Компонента $\frakr_\lambda$ основного идеала $\frakg_i$, соответствующая наибольшему первичному множителю $Q_\lambda^{(i)}=(P_\lambda^{(i)})^{\rho_\lambda}$ нормы $R^{(i)}$, задается изолированной компонентой разложения, однозначно определенной ассоциированными простыми идеалами размерностей $n-1,\ldots,n-i+1$ и тем простым идеалом размерности $n-i$, форма элементарных делителей которого равна $P_\lambda^{(i)}$. Ассоциированные простые идеалы и наибольшие первичные множители нормы взаимно однозначно соответствуют друг другу.}

По § 1, 7. идеалы $\frakr_\lambda$ совпадают со своим основным идеалом $i$-й ступени и имеют $\frakg_{i-1}$ как основной идеал $(i-1)$-й ступени; следовательно, по теореме VIII, соответственно по (5), они допускают кратчайшее представление:
\[
 \frakr_\lambda=[\frakg_{i-1},\frakq_{\lambda,1},\ldots,\frakq_{\lambda,t_\lambda}],
\]
где все $\frakq$ имеют размерность $n-i$ и принадлежат различным простым идеалам. По определению $Q_\lambda^{(i)}$ --- с учетом того, что $\frakr_\lambda$ совпадает со своим основным идеалом $i$-й ступени, --- получается:
\[
 Q_\lambda^{(i)}\frakg_{i-1}\equiv0\pmod{\frakq_{\lambda,\kappa}},
 \quad(\kappa=1,\ldots,t_\lambda)
\]
и следовательно
\[
 (Q_\lambda^{(i)})^{\rho_\lambda}\frakr_\lambda\equiv0\pmod{\frakq_{\lambda,\kappa}};
 \quad\hbox{следовательно}\quad
 (P_\lambda^{(i)})^{\rho_\lambda}\frakr_\lambda\equiv0\pmod{\frakp_{\lambda,\kappa}}
\]
и потому $P_\lambda^{(i)}\equiv0\pmod{\frakp_{\lambda,\kappa}}$; отсюда, как в доказательстве теоремы IX, следует, что в представлении $\frakr_\lambda$ может появиться только один $\frakq_\lambda$ и соответствующий $\frakp_\lambda$ размерности $n-i$, а форма элементарных делителей $\frakp_\lambda$ задается через $P_\lambda^{(i)}$. Нужно доказать, что $\frakp_\lambda$ является ассоциированным простым идеалом $\frakm$, и показать, что $\frakq_\lambda$ действительно входит по крайней мере в одно представление $\frakm$ через наибольшие первичные компоненты; тогда $\frakr_\lambda$ тем самым распознана как изолированная компонента разложения, поскольку никакой простой идеал не может входить в идеал той же или меньшей размерности, а именно как компонента, определенная $\frakp_\lambda$ и всеми ассоциированными простыми идеалами большей размерности. Для этого нужно лишь показать, что $\frakq_\lambda$ входит в кратчайшее представление $\frakg_i$, так как такое представление по теореме VIII всегда можно дополнить до кратчайшего представления $\frakm$; следовательно, нужно только показать, что представление
\[
 \frakg_i=[\frakr_1,\ldots,\frakr_a]=[\frakg_{i-1},\frakq_1,\ldots,\frakq_a]
\]
является кратчайшим. Если бы, скажем, $\frakq_\lambda$ входил в свой комплемент, то есть если бы $\frakg_{i-1}\frakq_1\cdots\frakq_{\lambda-1}\frakq_{\lambda+1}\cdots\frakq_a$ делился на $\frakq_\lambda$, то из-за $\frakg_{i-1}\not\equiv0\pmod{\frakq_\lambda}$ отсюда следовала бы делимость некоторой степени $\frakq_\mu$ на $\frakq_\lambda$ для $\mu\ne\lambda$; тогда, поскольку ассоциированные простые идеалы $\frakp_\mu$ и $\frakp_\lambda$ имеют одинаковую размерность, эти простые идеалы по теореме II совпадали бы. Но это невозможно, так как их формы элементарных делителей $P_\mu^{(i)}$ и $P_\lambda^{(i)}$ по определению различны. Представление $\frakg_i$, которому по теореме VIII соответствуют все ассоциированные простые идеалы размерности $n-i$, далее показывает, что каждому такому простому идеалу действительно соответствует и наибольший первичный множитель нормы; тем самым взаимно однозначное соответствие доказано.

Теорема X непосредственно дает перенос теоремы IV в следующем виде.

\textbf{Теорема XI.} \emph{Отдельные наибольшие первичные множители $Q_\lambda^{(i)}=(P_\lambda^{(i)})^{\rho_\lambda}$ нормы допускают представление произведением:}
\begin{align*}
 Q_\lambda^{(i)}
 &=\prod (x_i-t_{1i}\bar y_{1\nu}-\cdots-t_{ni}\bar y_{n\nu})^{\delta_\lambda\rho_\lambda}\\
 &=\prod (t_{1i}(y_1-\bar y_{1\nu})+\cdots+t_{ni}(y_n-\bar y_{n\nu}))^{\delta_\lambda\rho_\lambda},
\end{align*}
\emph{где $\bar y_{1\nu},\ldots,\bar y_{n\nu}$ каждый раз задают соответствующую систему нулей ассоциированного простого идеала $\frakp_\lambda$ в исходных неопределенных $y$, независимую от $t_{1i},\ldots,t_{ni}$, а $\delta_\lambda$ имеет значение единица или $p^{f_\lambda}$; тем самым полностью задано разложение нормы на линейные множители. Степень $Q_\lambda^{(i)}$ задает степень подкольца классов вычетов по изолированной компоненте $\frakr_\lambda$, представленного классами вычетов из $\frakg_{i-1}$, а именно относительно подполя классов вычетов $(P(u,x_{i+1},\ldots,x_n))$, полученного образованием частных; для первичных идеалов, следовательно, речь идет о кольце всех классов вычетов. При присоединении новых неопределенных $v$ получается еще разложение}
\[
 Q_\lambda^{(i)}(u,v)=\prod (z-v_1\bar x_{1\nu}-\cdots-v_i\bar x_{i\nu})^{\delta_\lambda\rho_\lambda}.
\]

Поскольку $P_\lambda^{(i)}$ по теореме X распознана как форма элементарных делителей $\frakp_\lambda$, здесь фактически речь идет лишь о подстановке разложения из теоремы IV; замечание о степени есть только иная формулировка сказанного в § 1, 5. и 7.\footnote{Точно так же и упомянутая в прим. 2) у H.-N. формулировка кратности является непосредственным следствием теоремы X. Ведь система классов вычетов $\frakg_{i-1}\mid\frakr_\lambda$ --- при присоединении $x_{i+1},\ldots,x_n$ --- изоморфна $\mathfrak t_\lambda\mid\frakr_\lambda$, где положено $\mathfrak t_\lambda=[\frakg_{i-1},\frakr_1,\ldots,\frakr_{\lambda-1},\frakr_{\lambda+1},\ldots,\frakr_a]$, как непосредственно показывает возвращение к модулям из линейных форм (H.-N., теорема V) с учетом взаимной простоты $Q_\mu^{(i)}$; $\mathfrak t_\lambda\mid\frakr_\lambda$ становится изоморфной $\mathfrak t_\lambda\mid\frakq_\lambda$. При этом $\mathfrak t_\lambda$ возникает из комплемента $\frakq_\lambda$ при присоединении $x_{i+1},\ldots,x_n$.}


\subsection*{§ 6. Характеризация простых идеалов и собственных первичных идеалов через форму элементарных делителей и норму.}

В теореме IX первичные идеалы были охарактеризованы через форму элементарных делителей и норму, но только так, что простой случай рассматривался как специальный случай первичного. Теперь нужно отделить простые идеалы от собственных первичных идеалов; рассуждение будет идти параллельно § 3. Условия, которые являются либо необходимыми, либо достаточными, легко записываются в следующем виде.

\textbf{Теорема XII.} \emph{Необходимым условием для простого идеала является то, что его форма элементарных делителей является простой функцией; достаточным условием является то, что его норма является простой функцией. Иными словами, форма элементарных делителей простого идеала всегда является простой функцией, а норма собственного первичного идеала всегда является собственной первичной функцией.}

То, что форма элементарных делителей простого идеала становится простой функцией, уже было показано в теореме I. Пусть теперь норма $R_\frakq$ равна простой функции; нужно показать, что при этом всякий собственный делитель $\fraka$ идеала $\frakq$ имеет меньшую наивысшую размерность, так что $\frakq$ по теореме II распознается как простой идеал. Действительно, по § 1, 6. первый множитель нормы $R_\fraka$, отличный от соответствующего множителя $R_\frakq$, становится собственным делителем; ввиду $R_\frakq=P^{(i)}$ множитель $R^{(i)}$ нормы $R_\fraka$ должен быть равен собственному делителю $P^{(i)}$, то есть единице; поэтому $\fraka$ имеет меньшую наивысшую размерность.

Второе простое доказательство опирается на следующую лемму, которая лежит также в основе следующей теоремы.

\textbf{Лемма VII.} \emph{Система классов вычетов, представленная многочленами $H^{(i)}$, по модулю формы элементарных делителей $E^{(i)}$ первичного идеала --- более общо, идеала, имеющего только ассоциированные простые идеалы размерности $n-i$, --- изоморфна подсистеме классов вычетов по этому идеалу. Степень формы элементарных делителей задает степень этой системы классов вычетов относительно поля частных $(P(u,x_{i+1},\ldots,x_n))$.}

Действительно, можно установить соответствие, сопоставляя классы, представленные теми же многочленами $H^{(i)}$; сумма и произведение при этом соответствуют сумме и произведению. Это соответствие также взаимно однозначно: ведь все многочлены $H^{(i)}$, принадлежащие нулевому классу по идеалу, по определению делятся на $E^{(i)}$; следовательно, нулевому классу соответствует нулевой класс по $E^{(i)}$, и, конечно, наоборот.

Пусть теперь норма идеала является простой функцией $P^{(i)}$, то есть совпадает с формой элементарных делителей. Тогда степени систем классов вычетов по форме элементарных делителей $P^{(i)}$ и по $\frakq$ относительно $(P(u,x_{i+1},\ldots,x_n))$ совпадают; значит, подсистема, изоморфная системе классов вычетов по $P^{(i)}$, исчерпывает классы вычетов по $\frakq$; а так как эта система не может иметь делителей нуля, $\frakq$ является простым идеалом.

В общем случае теорему XII нельзя усилить до необходимых и достаточных условий, как будет показано ниже на примерах. Зато это возможно, когда область коэффициентов $P$ является совершенным полем (§ 1, 14.), по следующей теореме.

\textbf{Теорема XIII.} \emph{Если область коэффициентов $P$ является совершенным полем, то норма и форма элементарных делителей простого идеала становятся простыми функциями и, следовательно, совпадают; норма и форма элементарных делителей собственного первичного идеала становятся собственными первичными функциями. Следовательно, над совершенным полем условие, что форма элементарных делителей является простой функцией, является необходимым и достаточным для простого идеала; вместо формы элементарных делителей в этом условии может также стоять норма.}

С учетом теоремы XII теорема XIII будет доказана, если показать, что над совершенным полем идеал становится простым, как только его форма элементарных делителей является простой функцией, и что тогда норма равна этой простой функции. В лемме V, § 3, было показано, что над совершенным полем форма элементарных делителей становится простой функцией первого рода, как только она вообще является простой функцией. Поэтому пусть $(\mathcal R)$ обозначает поле классов вычетов, возникающее при присоединении $(x_i)$ к $(P(u,x_{i+1},\ldots,x_n))$, то есть $(P(u,x_i,x_{i+1},\ldots,x_n))$ по $P^{(i)}$. Тогда $(x_i)$ является нулем простой функции первого рода $(P^{(i)})$, и поэтому, как показывает, например, формула Лагранжа, $(y_1),\ldots,(y_n)$ содержатся в $(\mathcal R)$. Подсистема классов вычетов по рассматриваемому идеалу $\frakq$, изоморфная по лемме VII системе $(\mathcal R)$ --- причем речь идет только о появлении знаменателей из $(P(u,x_{i+1},\ldots,x_n))$, --- содержит классы вычетов $(y_1),\ldots,(y_n)$ и потому исчерпывает систему классов вычетов по $\frakq$; но поскольку как система, изоморфная $(\mathcal R)$, она не может иметь делителей нуля, $\frakq$ становится простым идеалом. Степень этого поля классов вычетов $(\mathcal R)$ идеала $\frakq$ относительно $(P(u,x_{i+1},\ldots,x_n))$ равна степени нормы $\frakq$; с другой стороны, благодаря изоморфизму с $(\mathcal R)$ она равна степени формы элементарных делителей $P^{(i)}$; значит, норма и форма элементарных делителей должны совпадать.

Для несовершенных областей коэффициентов теорема XII все же допускает следующее усиление.

\textbf{Теорема XIV.} \emph{Если область коэффициентов $P$ является несовершенным полем характеристики $p$, то норма простого идеала становится $p^g$-й степенью формы элементарных делителей; при этом $0\le g\le(i-1)f$, где $f$ обозначает экспоненту формы элементарных делителей, а $n-i$ --- размерность простого идеала.}

Теорема XIV сводится к утверждению, что поле классов вычетов $(\mathcal R)$ идеала $\frakp$ имеет степень $p^g$ ($g\ge0$) относительно подполя $(\mathcal R')=(P(u,x_i,x_{i+1},\ldots,x_n))$, изоморфного $(\mathcal R)$. Это непосредственно следует из теорем Штейница, приведенных в § 1, 15., точнее из следующего утверждения.

\textbf{Лемма VIII.} \emph{Пусть $\Lambda$ есть конечное расширение несовершенного поля $\Omega$ редуцированной степени $r$ и экспоненты $f$, $\Lambda_0$ --- поле первого рода, содержащееся в $\Lambda$, а $M$ --- промежуточное поле между $\Lambda_0$ и $\Lambda$. Тогда степень каждого элемента из $\Lambda$ относительно $M$ является степенью $p$, где $p$ обозначает характеристику $\Omega$.}

Ибо $p^{f'}$-я степень ($f'\le f$) каждого элемента $z$ из $\Lambda$ принадлежит $\Lambda_0$; следовательно, $z$ является нулем простой функции $G(t)=(t-z)^{p^{f'}}$ с коэффициентами из $\Lambda_0$. Пусть $H(t)=(t-z)^\delta$ является простой функцией в $M$ с нулем $z$; тогда $H(t)$ имеет также коэффициенты из $\Lambda_0(z)$, а значит коэффициенты $H(t)$ принадлежат полю пересечения $M$ и $\Lambda_0(z)$, то есть промежуточному полю между $\Lambda_0$ и $\Lambda_0(z)$. Поэтому степень $z$ относительно $M$ равна степени относительно этого промежуточного поля и, как делитель $p^{f'}$, является степенью $p$.

Для доказательства теоремы XIV остается только заметить, что поле $(\mathcal R)=(P(u,x_i,x_{i+1},\ldots,x_n))$ должно иметь степень $r\cdot p^f$ относительно $(P(u,x_{i+1},\ldots,x_n))$, если $r$ обозначает редуцированную степень, а $f$ --- экспоненту $(\mathcal R)$; и что поэтому поле первого рода $(\mathcal R_0)$, содержащееся в $(\mathcal R)$, должно быть подполем $(\mathcal R')$. Ведь, поскольку в $(\mathcal R)$ существуют элементы степени $r\cdot p^f$ (§ 1, 15.), необходимо, чтобы $(x_i)$ имел именно эту степень, потому что все элементы из $(\mathcal R)$ возникают специализацией $t_{\mu\nu}$ в $(x_i)$; экспонента $(\mathcal R)$ совпадает с экспонентой формы элементарных делителей, и $(x_i)^{p^f}$ имеет степень $r$ и принадлежит $(\mathcal R_0)$, следовательно является примитивным элементом $(\mathcal R_0)$. Так как $(\mathcal R)$ возникает присоединением конечного числа элементов --- скажем, $(x_1),\ldots,(x_{i-1})$ --- к $(\mathcal R')$, конечное многократное применение леммы VIII дает требуемое доказательство и одновременно оценку для $g$.\footnote{Ср. параллельные рассуждения у A. Ostrowski, Zur arithmetischen Theorie der algebraischen Größen, Gött. Nachr. 1919, с. 279--298; в частности с. 288, где соответствующая теорема сформулирована без доказательства.}

Теперь добавим примеры, которые показывают, что теорему XII в общем случае нельзя усилить сверх теоремы XIV. Речь идет о простом идеале, норма которого равна квадрату формы элементарных делителей ($g>0$), и о собственных первичных идеалах, форма элементарных делителей которых является простой функцией; последний пример одновременно показывает, что здесь нет аналога теоремы XIV, а норма может стать произвольной степенью формы элементарных делителей.

Пусть область коэффициентов $P$ получается присоединением двух неопределенных $\lambda,\mu$ к полю классов вычетов по модулю $2$, то есть является несовершенным полем характеристики два. Рассмотрим идеал
\[
 \bar{\frakp}=(y_1^2+\lambda,\,y_2^2+\mu),
\]
который должен быть простым идеалом, потому что система классов вычетов $(\mathcal R)$ изоморфна полю $P(\sqrt\lambda,\sqrt\mu)$. $(\mathcal R)$ имеет четвертую степень относительно $(P)$, редуцированную степень $r=1$, экспоненту $f=2$; поэтому форма элементарных делителей $P^{(2)}$ имеет вторую степень, а норма --- четвертую степень и, следовательно, является квадратом $P^{(2)}$, что также легко проверить по модульному представлению (3), § 1, 5.; получаем
\[
 P^{(2)}=(u_{11}u_{22}+u_{12}u_{21})^2x_2^2+(u_{11}^2\mu+u_{21}^2\lambda)
\]
или, после соответствующего умножения,
\[
 x_2^2+(t_{12}^2\lambda+t_{22}^2\mu).
\]

Пусть теперь $P$ получается присоединением одной неопределенной $\lambda$ к полю классов вычетов по модулю $2$, и пусть взят идеал
\[
 \bar{\frakq}=(y_1^2+\lambda,\,y_2^2+\lambda)=(y_1^2+\lambda,(y_1+y_2)^2).
\]
В качестве формы элементарных делителей, через специализацию $\mu=\lambda$ в приведенной выше формуле, получается
\[
 P^{(2)}=x_2^2+\lambda(t_{12}^2+t_{22}^2),
\]
следовательно простая функция, потому что $t_{12}=(0)$, $t_{22}=(1)$ приводит к простой функции $x_2^2+\lambda$. Но $\bar{\frakq}$ является собственным первичным идеалом, потому что он содержит $(y_1+y_2)^2$, но не содержит $(y_1+y_2)$. Норма должна быть равна квадрату, то есть $p$-й степени, $P^{(2)}$, потому что число линейно независимых относительно $P$ классов вычетов по $\bar{\frakq}$ равно четырем.

Что этот последний аналог теоремы XIV не всегда выполняется, показывает следующий пример. Пусть $P$ получается присоединением неопределенной $\lambda$ к полю классов вычетов по модулю $3$, и пусть
\[
 \frakq=(y_1^3+\lambda,(y_1-y_2)^2)
\]
так что, как и выше, речь идет о собственном первичном идеале. Форма элементарных делителей становится --- с учетом того, что также $y_2^3+\lambda$ делится на $\frakq$, ---
\[
 P^{(2)}=x_2^3+\lambda(t_{12}+t_{22})^3,
\]
следовательно простой функцией; но норма равна квадрату $P^{(2)}$, потому что имеется шесть линейно независимых классов вычетов по $\frakq$. Следовательно, здесь не возникает $p^g$-я степень формы элементарных делителей.

Первый пример показывает, что при несовершенном поле как области коэффициентов, так же как и в алгебраических числовых полях, могут встречаться простые идеалы степени выше первой, где $p^g$ следует называть степенью простого идеала. Дальнейшие примеры следует понимать как аналоги идеалов разветвления: простое число может делиться на более высокую степень простого идеала, то есть на первичный идеал; ведь, как показано в примечании 10), форма элементарных делителей соответствует наименьшему целому рациональному числу, делящемуся на идеал в алгебраическом числовом поле.


\subsection*{§ 7. Абсолютные простые идеалы.}

Простой идеал с коэффициентами из $P$ называется \emph{абсолютным простым идеалом}, если он остается простым идеалом в алгебраически замкнутом поле $A$, до которого можно расширить $P$.\footnote{Алгебраически замкнутое поле, как известно, есть такое поле, в котором каждый полином одной неопределенной разлагается на линейные множители. О существовании и существенной единственности алгебраически замкнутого поля, принадлежащего произвольному полю, см. у Steinitz.} Соответственно простая функция $P$ с коэффициентами из $P$ называется \emph{абсолютной простой функцией} --- абсолютно неприводимым полиномом, --- если $P$ остается простой функцией в $A$. Характеризация абсолютных простых идеалов дается следующей теоремой.

\textbf{Теорема XV.} \emph{Необходимое и достаточное условие для абсолютного простого идеала состоит в том, что его форма элементарных делителей --- которую можно считать целой и примитивной по $u$ --- становится абсолютной простой функцией как полином от $x$ и $u$. Форма элементарных делителей становится тождественной норме, так что условие можно сформулировать также для нормы.}

Для доказательства сначала заметим, что в общем случае норма и форма элементарных делителей сохраняются при алгебраическом расширении области коэффициентов. Ведь отдельные множители $R^{(i)}$ соответственно $E^{(i)}$ определены как наибольшие общие делители в полиномиальном смысле, а это свойство сохраняется при алгебраическом расширении области коэффициентов. Следовательно, $P^{(i)}$ остается формой элементарных делителей $\frakp$ при переходе от $P$ к $A$, то есть при переходе к совершенному полю как области коэффициентов; ведь каждое алгебраически замкнутое поле, как непосредственно показывает определение, совершенно. По теореме XIII необходимое и достаточное условие для того, чтобы $\frakp$ был абсолютным простым идеалом, состоит в том, что $P^{(i)}$ становится простой функцией относительно $A(u)$; то есть абсолютной простой функцией как полином от $x$ и $u$, если $P^{(i)}$ считается целым и примитивным по $u$. Одновременно по теореме XIII $P^{(i)}$ становится также тождественным норме $\frakp$, как только $A$ берется в качестве области коэффициентов; значит, по сказанному выше то же самое верно и при исходной области $P$. Теорема XV доказана.

Для абсолютных простых функций $P$ с коэффициентами из (конечного) алгебраического числового поля $\mathfrak K$ теперь имеется теорема,\footnote{A. Ostrowski: там же (прим. 21); лемма, с. 296. --- Более простое доказательство см. E. Noether, Ein algebraisches Kriterium für absolute Irreduzibilität, Math. Ann. 85 (1922), с. 26--33, № 7.} что $P$ остается простой функцией, точнее абсолютной простой функцией, по модулю каждого простого идеала из $\mathfrak K$, самое большее за исключением конечного числа простых идеалов из $\mathfrak K$. Перенос этой теоремы на абсолютные простые идеалы опирается на следующее.

\textbf{Теорема XVI.} \emph{Если в качестве области коэффициентов вместо поля взять кольцо $\mathfrak o^*$ всех целых алгебраических чисел (конечного) алгебраического числового поля $\mathfrak K$, то норма и форма элементарных делителей полиномиального идеала сохраняются как норма и форма элементарных делителей по модулю каждого простого идеала из $\mathfrak K$, самое большее за исключением конечного числа простых идеалов из $\mathfrak K$.}

Для доказательства надо модифицировать лежащую в основе полиномиальную область по сравнению с § 1, 1. Пусть $\bar{\frakS}$ состоит из всех полиномов от $y_1,\ldots,y_n$ с коэффициентами из $\mathfrak o^*$, то есть из целых алгебраических чисел из $\mathfrak K$; областью коэффициентов для $\frakS$ пусть будет $\mathfrak o^*[u]$, то есть все полиномы от $u$ с коэффициентами из $\mathfrak o^*$. Если $\bar{\frakm}$ есть идеал в $\bar{\frakS}$, то посредством (1) он переходит в преобразованный идеал $\frakm$ в $\frakS$. Нужно показать, что при образовании нормы в знаменателе появляется только конечное число полиномов из $\mathfrak o^*[u]$; отсюда теорема XVI будет следовать непосредственно.

Для этого сначала заметим, что модульное представление (3) из § 1, 5. для образования отдельной нормы сводится к редукции степеней $x_{i-1}$ по модулю полинома, регулярного по $x_{i-1}$,
\[
 C^{(i-1)}=U_{i-1}(u)x_{i-1}^{e_{i-1}}+\hbox{младшие члены};
\]
\footnote{Ср. H.-N. § 4.} при этом все коэффициенты $C^{(i-1)}$, в частности $U_{i-1}$, можно считать величинами из $\mathfrak o^*[u]$. Так как речь идет о последовательной редукции $x_1,\ldots,x_{i-1}$, модульное представление, целое по $x$, является также целым по $\mathfrak o^*[u]$ с точностью до произведений степеней
\[
 U_1^{\lambda_1}\cdots U_{i-1}^{\lambda_{i-1}}
\]
в знаменателе. Произведение $U_1U_2\cdots U_n$, рассматриваемое как полином от $u$, задает своими коэффициентами идеал $\mathfrak r^*$ из $\mathfrak K$, который поэтому делится только на конечное число простых идеалов из $\mathfrak K$. Если взять $\frakp^*$ отличным от этих конечных простых идеалов и положить в качестве области коэффициентов поле классов вычетов по $\frakp^*$, то исходное модульное представление сохраняется: каждое число из $\mathfrak o^*$ только заменяется своим классом вычета по $\frakp^*$.

Далее, поскольку норма и форма элементарных делителей определены только с точностью до множителей из $\mathfrak o^*[u]$, их также можно, относительно $\mathfrak o^*[u]$, считать делящимися на $\frakm$. Пусть при таком предположении, например,
\[
 R^{(i)}=V_i(u)x_i^{e_i}+\hbox{младшие члены},
\]
так что при делимости всех детерминантов порядка $\rho$ модуля $\mathfrak M_{i-1}^*$ ранга $\rho$, определенного через $C^{(i-1)}$, в знаменателе, кроме $U_1^{\lambda_1}\cdots U_{i-1}^{\lambda_{i-1}}$, появляются только степени $V_i$. Если поэтому выбрать $\frakp^*$ еще и отличным от конечного числа простых идеалов из $\mathfrak K$, делящих идеал, определенный произведением $V_1V_2\cdots V_n$, то $R^{(i)}$ остается общим делителем этих детерминантов, когда в качестве области коэффициентов берется поле классов вычетов по $\frakp^*$; ранг $\mathfrak M_{i-1}^*$ также сохраняется, поскольку $C^{(i)}$ был детерминантом порядка $\rho$ из $\mathfrak M_{i-1}^*$.

Наконец $\frakp^*$ надо выбрать еще так, чтобы $R^{(i)}$ каждый раз оставался наибольшим общим делителем. Пусть $D^{(i)}$ пробегает все детерминанты порядка $\rho$ из $\mathfrak M_{i-1}^*$, так что по предположению получается
\[
 U_1^{\lambda_1}\cdots U_{i-1}^{\lambda_{i-1}}V_i^{\rho}D^{(i)}=R^{(i)}T^{(i)},
\]
где $T^{(i)}$ не имеют общего делителя, содержащего $x$, в полиномиальном смысле. Следовательно, идеал, выведенный из всех $T^{(i)}$, содержит полином
\[
 G^{(i+1)}=W_i(u)x_{i+1}^{m_i}+\hbox{младшие члены},\qquad W_i\ne0;
\]
что при этом $G^{(i+1)}$ можно считать регулярным по $x_{i+1}$, показывает составление преобразования (1) с таким преобразованием $x_{i+1},\ldots,x_n$, где коэффициентами служат новые неопределенные. Если, значит, выбрать $\frakp^*$ еще и отличным от конечного числа простых идеалов, делящих идеал из $\mathfrak K$, определенный произведением $W_1W_2\cdots W_n$, то при взятии поля классов вычетов по $\frakp^*$ в качестве области коэффициентов $R_\frakm$ действительно сохраняется как норма. Совершенно соответственно $\frakp^*$ можно выбрать еще и так, чтобы $E_\frakm$ также сохранялась как форма элементарных делителей. Теорема XVI доказана.

Из теорем XV и XVI, как обобщение теоремы об абсолютных простых функциях и с использованием этой теоремы, следует следующее.

\textbf{Теорема XVII.} \emph{Если $\frakp$ является абсолютным простым идеалом с целыми алгебраическими числами из (конечного) алгебраического числового поля $\mathfrak K$ в качестве коэффициентов, то $\frakp$ остается простым идеалом, точнее абсолютным простым идеалом, по модулю каждого простого идеала из $\mathfrak K$, самое большее за исключением конечного числа простых идеалов из $\mathfrak K$.}

Для доказательства пусть норма $R_\frakp$ идеала $\frakp$, как в теореме XVI, выбрана так, что она также относительно $\mathfrak o^*[u]$ делится на $\frakp$. Тогда
\[
 R_\frakp=T(u)P^{(i)}(x),
\]
где $P^{(i)}(x)$ можно считать целым и примитивным по $u$. Поэтому $P^{(i)}$, по теореме XV, если рассматривать его как полином от $u$ и $x$ с коэффициентами из $\mathfrak K$, становится абсолютной простой функцией. По теореме об абсолютных простых функциях он сохраняет это свойство по модулю каждого простого идеала из $\mathfrak K$, самое большее за исключением конечного числа простых идеалов из $\mathfrak K$. Если поэтому выбрать $\frakp^*$ отличным от этих конечных простых идеалов и, по теореме XVI, еще отличным от конечного числа дальнейших простых идеалов из $\mathfrak K$, то $P^{(i)}$ --- коэффициенты которого из $\mathfrak o^*$ надо заменить соответствующими классами вычетов по $\frakp^*$ --- становится нормой $\frakp$, когда в качестве области коэффициентов $P$ берется поле классов вычетов по $\frakp^*$, и эта норма становится абсолютной простой функцией. Следовательно, при взятии $P$ в качестве области коэффициентов $\frakp$ по теореме XV остается абсолютным простым идеалом. Теорема XVII доказана.

\begin{center}
Гёттинген, 8 мая 1923 г.\\[1em]
(Поступила 9. 3. 1923.)
\end{center}
% END UNIT BODY
% END PRODUCER UNIT 108
\endgroup


\clearpage
\begingroup
\setcounter{footnote}{0}
% BEGIN PRODUCER UNIT 109
% Source: C:/Users/Floris/Documents/interlanguage/03_projects/noether/02_slavic_working_corpus/translations/paper25/russian/v001/Noether_Paper25_Russian_v001.tex
% Identity: 20534 B / SHA-256 8A555515DA987D124F774D537898C0490AA18A223D2F3375DD8ED26101922B12
\setlength{\parindent}{1.25em}
\setlength{\parskip}{0pt}
\makeatletter
\renewcommand\footnotesize{\@setfontsize\footnotesize{7.8}{8.8}\selectfont}
\makeatother
\emergencystretch=35em
\allowdisplaybreaks
\sloppy
\hbadness=10000
\vbadness=10000
\hfuzz=2pt
\vfuzz=10pt
\raggedbottom
\providecommand{\fraka}{\mathfrak a}
\providecommand{\frakb}{\mathfrak b}
\providecommand{\frakp}{\mathfrak p}
\providecommand{\frakq}{\mathfrak q}
\providecommand{\frakR}{\mathfrak R}
\providecommand{\frakS}{\mathfrak S}
\providecommand{\Norm}{N}
% BEGIN UNIT BODY
\setcounter{footnote}{0}

% Unit: Paper25 v001. Direct Russian translation from the German final-audited source slice.
\section*{25. Теория элиминации и теория идеалов}

\begin{center}
J. Ber. d. DMV 33 (1924), с. 116--120
\end{center}

В дальнейшем я хотела бы показать, как теорию элиминации можно построить вполне параллельно теории нулей для полиномов от одной переменной; это построение основано на соединении теории элиминации, данной Hentzelt, с методами теории полей, как их развил Steinitz, и теории идеалов, как я сама ее разработала.\footnote{Подробное изложение и библиографические указания: Math. Ann. 90 и работа, которая вскоре появится.}

Сначала я набросаю постановку вопроса в случае полиномов от одной переменной. В качестве области коэффициентов раз и навсегда положим фиксированное поле $P$, к которому относятся понятия простой функции и т. д. При этом $P$ можно считать абстрактно определенным в смысле Steinitz; но, разумеется, в специальном случае оно может совпадать с полем рациональных чисел или с полем всех комплексных чисел --- согласно предположению, ранее общепринятому в теории элиминации. Теория нулей для полиномов от одной переменной теперь распадается на четыре шага:

\noindent 1. \emph{Теорема о разложении}, которая позволяет свести вопрос к вопросу о простых функциях --- полиномах, неприводимых в $P$. Ибо из
\[
 a(x)=p_1(x)^{\varrho_1}\cdots p_r(x)^{\varrho_r}=q_1(x)\cdots q_r(x)
\]
--- где $p(x)$ обозначают различные простые функции, а $q(x)$, следовательно, "наибольшие первичные функции", --- следует, что $a(x)$ обращается в нуль во всех и только в нулях простых функций $p(x)$.

\noindent 2. \emph{Построение поля нулей для простой функции $p(x)$.} Такое поле нулей $\frakR=P(\alpha)$ --- где $\alpha$ обозначает нуль $p(x)$ --- по Steinitz всегда можно построить как поле расширения поля $P$, изоморфное полю классов вычетов $(\frakR)$ по $p(x)$. Обратно, всякое поле нулей, лежащее в поле расширения поля $P$, изоморфно полю классов вычетов $(\frakR)$; поэтому поле нулей и поле классов вычетов абстрактно тождественны.

\noindent 3. \emph{Разложение $p(x)$ на линейные множители в поле Galois, определенном $p(x)$.} Повторяя построение, данное в пункте 2, приходят к существенно однозначно определенному полю Galois $P(\alpha_1,\ldots,\alpha_m)$, которого как раз достаточно, чтобы $p(x)$ разложилось на линейные множители.

\noindent 4. \emph{Значение кратности для исходного полинома $a(x)$.} По пунктам 1, 2, 3 для произвольного полинома $a(x)$ вопрос о нулях и соответствующее разложение на линейные множители решены. Кратностью при этом можно считать степень отдельной первичной функции $q(x)$ --- то есть число линейно независимых относительно $P$ классов вычетов по $q(x)$; ибо в случае алгебраически замкнутой области коэффициентов $P$, где $p(x)$ становятся линейными, это число как раз и дает кратность нулей.

Теперь я перехожу к общей теории элиминации, то есть беру за основу область $\frakS$ всех полиномов от $x_1,\ldots,x_n$ с коэффициентами из $P$; тогда теорию элиминации можно определить как \emph{теорию нулей идеала} из этой полиномиальной области. Под идеалом $\fraka$ я понимаю, как обычно, такую систему полиномов, что вместе с $f$ и $g$ к $\fraka$ принадлежит также $f-g$, а вместе с $f$ также $Af$, где $A$ понимается как произвольный полином из $\frakS$. Под нулем $\fraka$ я понимаю всякую систему значений $\alpha_1,\ldots,\alpha_n$, принадлежащую полю расширения поля $P$, такую, что $f(\alpha)$ обращается в нуль для каждого полинома $f$, принадлежащего идеалу. Если $f_1,\ldots,f_t$ обозначает всегда существующий базис идеала $\fraka$ --- то есть $f=A_1f_1+\cdots+A_tf_t$ для каждого $f$ из $\fraka$ ---, то нули $\fraka$, очевидно, тождественны нулям $f_1,\ldots,f_t$; а так как, обратно, всякую конечную систему полиномов можно рассматривать как базис порожденного ею идеала, то данное выше определение нулей для теории элиминации тождественно обычно употребляемому, но избегает трудности, заключенной в произвольном выделении базиса. Для полиномов от одной переменной, где речь идет только о главных идеалах, эта трудность не возникает, поскольку базисный полином определен однозначно с точностью до единицы --- множителя из $P$; однако и там идеальная точка зрения даст более точное понимание.

Теория элиминации теперь, соответственно случаю одной переменной, распадается на четыре шага:

\noindent I. \emph{Теорема о разложении} --- представление идеала через наибольшие первичные компоненты, --- которая позволяет свести вопрос к вопросу о простых идеалах. Идеал называется простым, если он делит произведение лишь тогда, когда делит по крайней мере один из множителей; первичным, если он делит по крайней мере один из множителей или некоторую степень каждого множителя. Каждому первичному идеалу $\frakq$ принадлежит один и только один ассоциированный простой идеал $\frakp$, который является делителем $\frakq$ и некоторая степень которого делится на $\frakq$. Имеет место представление как наименьшее общее кратное:
\[
 \fraka=[\frakq_1,\ldots,\frakq_r];\qquad
 \frakp_1^{\sigma_1}\cdots \frakp_r^{\sigma_r}\equiv 0\pmod{\fraka}.
\]
Здесь $\frakq$ являются наибольшими первичными компонентами; то есть каждое $\frakq$ первично, но наименьшее общее кратное любых двух $\frakq$ не первично. Кроме того, без ограничения общности можно предполагать, что ни одно $\frakq$ не делит наименьшее общее кратное остальных, так что ни одно $\frakq$ нельзя просто опустить; $\frakp$ являются ассоциированными простыми идеалами соответствующих $\frakq$. Итак, не сами $\frakq$, но --- и в этом состоит аналог случая одной переменной --- их ассоциированные простые идеалы $\frakp$ однозначно определены идеалом $\fraka$; и так как каждое $\frakp$ является делителем $\fraka$, а обратно произведение степеней $\frakp$ делится на $\fraka$, то нули идеала совпадают с нулями его ассоциированных простых идеалов.

\noindent II. \emph{Построение поля нулей для простого идеала $\frakp$.} Классы вычетов по простому идеалу по определению образуют кольцо без делителей нуля, содержащее по крайней мере один элемент, отличный от нулевого, как только $\frakp$ отличается от единичного идеала; образованием частных это кольцо, следовательно, можно расширить до поля классов вычетов $(\frakR)$ идеала $\frakp$. Всегда можно построить поле расширения $\frakR$ поля $P$, изоморфное $(\frakR)$, которое становится полем нулей $\frakp$; то есть $\frakR=P(\alpha_1,\ldots,\alpha_n)$, где $\alpha_1,\ldots,\alpha_n$ обозначает нуль $\frakp$. При этом $\frakR$ имеет степень трансцендентности $k$ $(0\leq k<n)$ относительно $P$, следовательно содержит $k$ алгебраически независимых относительно $P$ элементов, тогда как любые $(k+1)$ элемента становятся алгебраически зависимыми относительно $P$. Обратно, всякое поле нулей степени трансцендентности $k$, лежащее в поле расширения поля $P$, изоморфно полю классов вычетов $(\frakR)$; поэтому поле нулей степени трансцендентности $k$ и поле классов вычетов абстрактно тождественны.

\noindent III. \emph{Разложение нормы $\frakp$ на линейные множители в поле Galois, определенном $\frakp$.} Если считать переменные $x_1,\ldots,x_n$ полученными из исходных переменных $y_1,\ldots,y_n$ линейным преобразованием с неопределенными в качестве коэффициентов, то поле нулей $\frakR$ степени трансцендентности $k$ можно рассматривать как конечное алгебраическое поле расширения поля $P(x_{i+1},\ldots,x_n)$ $(k=n-i)$. Поэтому можно перейти к полю Galois $\frakR$ относительно $P(x_{i+1},\ldots,x_n)$, в котором простая функция $P^{(i)}(x_i,x_{i+1},\ldots,x_n)$ с нулем $\alpha_i$ разлагается на линейные множители. Эта простая функция $P^{(i)}$, если вернуться к исходным переменным $y$, играет роль фундаментального уравнения; ибо $\alpha_i$ становится равным фундаментальной форме $u_1\beta_1+\cdots+u_n\beta_n$, где $\beta$ обозначает систему нулей в переменных $y$, алгебраически зависящую от параметров $x_{i+1},\ldots,x_n$; следовательно, разложение на линейные множители дает произведение по всем нулям. Но простую функцию $P^{(i)}$ или ее степень можно также рассматривать как норму $N(\frakp)$ идеала $\frakp$, если рассматривать $\frakp$ как модуль из линейных форм в степенных произведениях $x_1,\ldots,x_{i-1}$, взять $P(x_{i+1},\ldots,x_n)$ вместо $P$ в качестве области коэффициентов и определить норму как произведение по всем элементарным делителям --- из которых только конечное число отлично от единицы. При совершенном поле $P$ в качестве области коэффициентов всегда имеем $N(\frakp)=P^{(i)}$; при несовершенном поле нормой может оказаться степень $P^{(i)}$.

\noindent IV. \emph{Норма исходного идеала и значение кратности.} Если соответственно последовательно рассматривать идеал $\fraka$ как модуль из линейных форм в степенных произведениях $x_1,\ldots,x_{i-1}$ $(i=1,2,\ldots,n)$, то можно показать, что при взятии за основу $P(x_{i+1},\ldots,x_n)$ этот модуль каждый раз имеет только конечное число элементарных делителей, отличных от единицы; это условие конечности следует из существования базиса идеала. Произведение по этим каждый раз конечным элементарным делителям дает отдельные нормы; норма $N(\fraka)$ определяется как произведение всех отдельных норм. Наибольшие первичные множители $Q^{(i)}(x_i,x_{i+1},\ldots,x_n)$ нормы $N(\fraka)$ взаимно однозначно соответствуют ассоциированным простым идеалам $\frakp$ идеала $\fraka$ так, что $Q^{(i)}$ становится степенью $P^{(i)}$, где $P^{(i)}$ (или, если нужно, степень $P^{(i)}$) обозначает норму $N(\frakp)$ идеала $\frakp$. Тем самым получается разложение
\[
 N(\fraka)=N(\frakp_1)^{\varrho_1}\cdots N(\frakp_r)^{\varrho_r}=Q_1\cdots Q_r,
\]
--- где, если нужно, $N(\frakp)$ следует заменить наивысшим элементарным делителем $E(\frakp)=P^{(i)}$, --- и тем самым по I, II, III для всех нулей $\fraka$ выполнено разложение нормы на линейные множители. Кратностью нулей, соответствующих отдельному простому идеалу $\frakp$ идеала $\fraka$, можно считать степень $Q^{(i)}$, которую можно истолковать как число классов вычетов, линейно независимых относительно $P(x_{i+1},\ldots,x_n)$. А именно речь идет о классах вычетов некоторых изолированных компонент разложения, однозначно определенных $\frakp$ и простыми идеалами большей размерности: о классах вычетов основного идеала $(i-1)$-й ступени по компоненте основного идеала $i$-й ступени, определенной $\frakp$; последняя компонента имеет $\frakp$ и все простые идеалы большей размерности как ассоциированные простые идеалы, первая же имеет только простые идеалы большей размерности.

Полное \emph{включение случая одной переменной} получается, если и здесь перейти к главному идеалу. Разложение, данное в пункте 1, тогда --- смотря по тому, рассматривается ли $a(x)$ как базисный полином или как норма, --- распадается на два отдельных разложения: на представление главного идеала как наименьшего общего кратного наибольших первичных компонент по I, что здесь сводится к представлению как произведения степеней простых идеалов; или на разложение нормы как произведения степеней норм отдельных простых идеалов по IV. В построении поля нулей по пункту 2 на самом деле речь идет о поле нулей простого идеала, определенного $p(x)$, тогда как в пункте 3 на линейные множители разлагается норма. Определение кратности по пункту 4 подчиняется данному в IV, поскольку основной идеал нулевой ступени равен единичному идеалу, а основной идеал первой ступени при одной переменной уже равен самому идеалу.

\emph{Включение в обычную теорию элиминации} задается тем, что норма идеала имеет значение результанта; при таком понимании старая проблема теории элиминации становится частью теории идеалов полиномиальной области, а именно той частью, которая наряду с теоремами о разложении идеалов также рассматривает разложение их норм и связь обоих разложений.

\begin{center}
(Поступило 19. 10. 1923.)
\end{center}
% END UNIT BODY
% END PRODUCER UNIT 109
\endgroup


\clearpage
\begingroup
\setcounter{footnote}{0}
% BEGIN PRODUCER UNIT 110
% Source: C:/Users/Floris/Documents/interlanguage/03_projects/noether/02_slavic_working_corpus/translations/paper26/russian/v001/Noether_Paper26_Russian_v001.tex
% Identity: 2212 B / SHA-256 2F88F142D0263E0C4F3C96B257CBC05D34A139C79F183C2500F4E724276FC4D0
\setlength{\parindent}{1.25em}
\setlength{\parskip}{0pt}
\emergencystretch=25em
\sloppy
\hbadness=10000
\vbadness=10000
\hfuzz=2pt
\vfuzz=10pt
\raggedbottom
% BEGIN UNIT BODY
% Unit: Paper26 v001. Direct Russian translation from the German final-audited source slice.
\section*{26. Абстрактное построение теории идеалов в алгебраическом числовом поле}

\begin{center}
J. Ber. d. DMV 33 (1924), с. 102
\end{center}

4. E. Noether, Гёттинген: Абстрактное построение теории идеалов в алгебраическом числовом поле.

Были указаны абстрактные свойства, вполне эквивалентные теории идеалов в алгебраическом числовом поле и тем самым характеризующие все кольца, обладающие такой теорией идеалов. Это кольца с единицей и без делителей нуля, в которых выполняется условие двойной цепи: каждая цепь идеалов по делимости обрывается за конечное число шагов, и так же каждая цепь кратных по модулю каждого идеала, отличного от нулевого идеала; кроме того, эти кольца интегрально замкнуты относительно своего поля частных. Если отбросить последнее предположение, то получаются утверждения о теории идеалов в конечном порядке, которые отчасти уже встречаются без доказательства у Dedekind (3-е изд.).

Подробное изложение появится в Mathematische Annalen.
% END UNIT BODY
% END PRODUCER UNIT 110
\endgroup


\clearpage
\begingroup
\setcounter{footnote}{0}
% BEGIN PRODUCER UNIT 111
% Source: C:/Users/Floris/Documents/interlanguage/03_projects/noether/02_slavic_working_corpus/translations/paper27/russian/v001/Noether_Paper27_Russian_v001.tex
% Identity: 3219 B / SHA-256 8D2CB176C779257E1B2CD1BCA2B0549A1B52476CC0BF7A40509D8627DAFCE460
\setlength{\parindent}{1.25em}
\setlength{\parskip}{0pt}
\emergencystretch=25em
\sloppy
\hbadness=10000
\vbadness=10000
\hfuzz=2pt
\vfuzz=10pt
\raggedbottom
% BEGIN UNIT BODY
% Unit: Paper27 v001. Direct Russian translation from the German final-audited source slice.
\section*{27. Гильбертовы числа в теории идеалов}

\begin{center}
J. Ber. d. DMV 34 (1925), с. 101
\end{center}

E. Noether: Гильбертовы числа в теории идеалов. Под гильбертовыми числами вообще понимаются такие числа, которые относятся к подсчёту классов вычетов идеалов. Сам Гильберт в своей «характеристической функции» задал функцию, которая для идеалов полиномиальной области, начиная с некоторой границы, даёт точное число линейно независимых классов вычетов. Для меньших значений степени Маколей дополнил эту функцию производящей функцией, которую Островский смог явно вычислить. Но Маколей указал также числа другого типа, которые оказались наиболее существенными для общей теории идеалов, поскольку их можно сформулировать абстрактно. В соответствии с общими теоремами о разложении достаточно дать определение для первичных идеалов $\mathfrak q$. Речь идёт о числе линейно независимых классов вычетов относительно поля классов вычетов ассоциированного простого идеала $\mathfrak p$. Эти числа распадаются на подсистемы, каждая из которых задаётся числом неприводимых компонент идеалов $\mathfrak q,\mathfrak q/\mathfrak p,\mathfrak q/\mathfrak p^2,\ldots$. Число $i$-й подсистемы также характеризуется числом членов композиционного ряда, который идёт от $\mathfrak q/\mathfrak p^i$ к $\mathfrak q/\mathfrak p^{i-1}$. Полный композиционный ряд и его гильбертовы числа образуют эквивалент представления первичных идеалов как степеней простых идеалов, которое в общем случае -- например, в конечных порядках алгебраического числового поля -- уже не действительно.
% END UNIT BODY
% END PRODUCER UNIT 111
\endgroup


\clearpage
\begingroup
\setcounter{footnote}{0}
% BEGIN PRODUCER UNIT 112
% Source: C:/Users/Floris/Documents/interlanguage/03_projects/noether/02_slavic_working_corpus/translations/paper28/russian/v001/Noether_Paper28_Russian_v001.tex
% Identity: 2856 B / SHA-256 FC232D7C63234FBDE8A6C3282D5A41305722B89B2640884C86D4370F40947A98
\setlength{\parindent}{1.25em}
\setlength{\parskip}{0pt}
\emergencystretch=25em
\sloppy
\hbadness=10000
\vbadness=10000
\hfuzz=2pt
\vfuzz=10pt
\raggedbottom
% BEGIN UNIT BODY
% Unit: Paper28 v001. Direct Russian translation from the German final-audited source slice.
\section*{28. Характеры групп и теория идеалов}

\begin{center}
J. Ber. d. DMV 34 (1925), с. 144
\end{center}

3. E. Noether, Гёттинген: Характеры групп и теория идеалов.

Фробениусова теория характеров групп -- то есть представлений конечных групп -- понимается как теория идеалов вполне приводимого кольца, группового кольца. Поэтому в общем виде развиваются теоремы об изоморфизме при аддитивном разложении вполне приводимого кольца; представление в виде прямой суммы прямо неразложимых простых идеалов единственно с точностью до изоморфизма, а неразложимые односторонние идеалы, принадлежащие одному и тому же двустороннему идеалу, изоморфны между собой. Следовательно, если изоморфные идеалы объединить в один класс, то различных классов неразложимых идеалов столько же, сколько неразложимых двусторонних идеалов.

При специализации на вполне приводимые системы гиперкомплексных величин -- в частности на групповое кольцо -- классы неразложимых идеалов и неприводимые классы представлений взаимно однозначно соответствуют друг другу. Далее, двусторонние неразложимые идеалы соответствуют характерам; и число аддитивных неразложимых компонент двустороннего идеала равно рангу неразложимых идеалов. Тем самым фробениусова теория включена в эту схему.

Подробное изложение должно появиться в Math. Ann.
% END UNIT BODY
% END PRODUCER UNIT 112
\endgroup


\clearpage
\begingroup
\setcounter{footnote}{0}
% BEGIN PRODUCER UNIT 113
% Source: C:/Users/Floris/Documents/interlanguage/03_projects/noether/02_slavic_working_corpus/translations/paper29/russian/v001/Noether_Paper29_Russian_v001.tex
% Identity: 31944 B / SHA-256 240F3F153DBC410A69A3B5AA4009E5E955DCB849078BC76F77D5B3120DD70B1D
\setlength{\parindent}{1.25em}
\setlength{\parskip}{0pt}
\emergencystretch=25em
\sloppy
\hbadness=10000
\vbadness=10000
\hfuzz=2pt
\vfuzz=10pt
\raggedbottom
% BEGIN UNIT BODY
% Unit: Paper29 v001. Direct Russian translation from the German final-audited source slice.
\section*{29. Теорема конечности для инвариантов конечных линейных групп характеристики $p$}

\begin{center}
Nachr. v. d. Ges. d. Wiss. zu Göttingen 1926, с. 28--35
\end{center}

Представлено R. Courant на заседании 30 июля 1926 года.

Ниже я показываю, что известная теорема конечности для инвариантов конечных групп линейных подстановок\footnote{Элементарное доказательство конечности см.: E. Noether, \emph{Der Endlichkeitssatz der Invarianten endlicher Gruppen}. Math. Ann. 77 (1916), с. 89--92.} сохраняется, если в качестве области коэффициентов вместо обычного числового поля положить в основу произвольное поле, следовательно, в частности поле характеристики $p$.\footnote{О понятиях теории полей см.: E. Steinitz, \emph{Algebraische Theorie der Körper}. J. f. M. 137 (1910), с. 167--309. Характеристика в § 4, поле корней в § 12, поле частных в § 3.} Однако известные доказательства конечности перестают работать, как только порядок группы делится на $p$; приходится привлекать более глубокие арифметические факты, а именно следующий критерий, до сих пор известный лишь для характеристики нуль:

\textbf{Критерий конечности}\footnote{Ср.: E. Noether, \emph{Algebraische und Differentialinvarianten}. Jahresber. d. d. Math.-Ver. 32 (1923), с. 177--184, № 3, теорема 4, и указанную там литературу. В последнем предложении критерия там ошибочно стоит «модульный базис относительно $\mathfrak R$» (фундаментальная система целых величин) вместо «относительно некоторого подкольца $\mathfrak R$».}: область целостности $\mathfrak S$ из полиномов -- более обще, из рациональных функций -- с коэффициентами из поля конечна тогда и только тогда, когда в $\mathfrak S$ содержится конечное подкольцо $\mathfrak R$ такое, что все элементы $\mathfrak S$ алгебраически целы над $\mathfrak R$. Тогда каждый интегральный базис $\mathfrak R$ дополняется модульным базисом $\mathfrak S$ относительно некоторого подкольца $\mathfrak R$ до интегрального базиса $\mathfrak S$.

Этот критерий, сам по себе представляющий самостоятельный интерес, опирается на существование рационального базиса и на тот факт, что теорема о цепях делителей сохраняется при переходе к целым элементам конечного расширения. Эта вторая теорема в предыдущей заметке Artin и van der Waerden доказана для произвольной характеристики лишь при некотором предположении конечности для исходной области, которое выполнено в случае конечных групп (кольцо корней представляет конечное расширение). Для существования рационального базиса я даю доказательство, действительное для произвольных полей как областей коэффициентов; тем самым я могу доказать критерий конечности при указанном выше ограничении.

То, что для инвариантов конечных групп критерий конечности выполнен, опирается на принцип симметрических функций. Но если в характеристике нуль коэффициенты резольвенты Galois уже дают интегральный базис, то в общем случае они порождают лишь конечное подкольцо, требуемое в критерии конечности. Из существования этого конечного подкольца теми же рассуждениями следует конечность «относительных» инвариантов и «модулярных» инвариантов.

К рассматриваемым здесь инвариантам как специальные случаи относятся «проективные инварианты mod $p$», изучавшиеся Dickson и его школой.\footnote{Ср.: Dickson, \emph{On Invariants and the theory of Numbers}. The Madison Colloquium 1913. Новейшая литература содержится в томах Transactions of the American Mathematical Society, вышедших с тех пор.} Здесь теорема конечности была известна для модулярных инвариантов, но не для общих «формальных» инвариантов.

\subsection*{§ 1. Критерий конечности}

1. \textbf{Теорема о существовании рационального базиса.} Первая формулировка: всякая система $S$ рациональных функций от $n$ неопределенных с коэффициентами из поля $P$ имеет конечный рациональный базис; то есть существует конечное число функций системы такое, что всякую функцию системы можно рационально, с коэффициентами из $P$, выразить через эти конечные функции.

Вторая формулировка: пусть $\mathfrak K$ -- промежуточное поле между $P$ и $P(x_1,\ldots,x_n)$, где $x_1,\ldots,x_n$ обозначают неопределенные, степени трансцендентности $t\leq n$; пусть $y_1,\ldots,y_t$ -- неприводимая система\footnote{По Steinitz система называется неприводимой, если она состоит из алгебраически независимых функций. Говорят, что система имеет степень трансцендентности $t$, если она содержит по крайней мере одну неприводимую систему из $t$ элементов, но не содержит такой системы из большего числа элементов. Простые теоремы о неприводимых системах, используемые ниже, см. у Steinitz, § 22.} из $\mathfrak K$. Тогда $\mathfrak K$ является конечным алгебраическим расширением $P(y_1,\ldots,y_t)$.

Обе формулировки фактически эквивалентны. Достаточно доказать существование рационального базиса для всех промежуточных полей; ведь поле, полученное из системы $S$, становится таким промежуточным полем $\mathfrak K$, рациональный базис которого непосредственно дает рациональный базис $S$, поскольку каждый элемент $\mathfrak K$ становится рациональной комбинацией конечного числа функций из $S$. Из второй формулировки, в свою очередь, следует, что произвольная неприводимая система $y_1,\ldots,y_t$ из $\mathfrak K$, дополненная конечным числом функций, характеризующих $\mathfrak K$ как конечное расширение $P(y_1,\ldots,y_t)$, дает такой рациональный базис. Обратно, если по первой формулировке задан рациональный базис $\mathfrak K$, то по определению степени трансцендентности каждый элемент базиса должен быть алгебраически зависим над $P(y_1,\ldots,y_t)$; следовательно, этим $\mathfrak K$ характеризуется как конечное алгебраическое расширение $P(y_1,\ldots,y_t)$.

Теорема будет доказана во второй формулировке. Сначала пусть $\mathfrak K$ имеет ту же степень трансцендентности $n$, что и $P(x_1,\ldots,x_n)$; пусть $y_1,\ldots,y_n$ -- неприводимая система из $\mathfrak K$, а следовательно и из $P(x_1,\ldots,x_n)$. В силу степени трансцендентности $n$ системы $y_1,\ldots,y_n$ каждое $x_i$ становится алгебраическим над $P(y_1,\ldots,y_n)$; значит, $P(x_1,\ldots,x_n)$, а вместе с ним и промежуточное поле $\mathfrak K$, становится конечным алгебраическим расширением $P(y_1,\ldots,y_n)$.

Теперь пусть $t<n$; нумерация $x$ выбрана так, что $y_1,\ldots,y_t,x_{t+1},\ldots,x_n$ образуют неприводимую систему из $P(x_1,\ldots,x_n)$. Тогда $x_{t+1},\ldots,x_n$ неприводимы относительно $P(y_1,\ldots,y_t)$ и потому, по транзитивному закону алгебраической зависимости, также неприводимы относительно $\mathfrak K$ и относительно всякого промежуточного поля $\mathfrak M$ между $P(y_1,\ldots,y_t)$ и $\mathfrak K$. Это означает, что $\overline{\mathfrak K}=\mathfrak K(x_{t+1},\ldots,x_n)$, соответственно $\overline{\mathfrak M}=\mathfrak M(x_{t+1},\ldots,x_n)$, является чисто трансцендентным расширением $\mathfrak K$, соответственно $\mathfrak M$; следовательно, всякий элемент $\overline{\mathfrak K}$, не принадлежащий $\mathfrak K$, трансцендентен относительно $\mathfrak K$, и точно так же всякий элемент $\overline{\mathfrak M}$, не принадлежащий $\mathfrak M$, трансцендентен относительно $\mathfrak M$. Аналогично положим $\overline P$ равным чисто трансцендентному расширению $P(x_{t+1},\ldots,x_n)$; тогда $y_1,\ldots,y_t$ неприводимы относительно $\overline P$.

Теперь доказательство сводится к уже разобранному случаю, если взять $\overline P$ в качестве области коэффициентов. Из
\[
P<\mathfrak K<P(x_1,\ldots,x_t,x_{t+1},\ldots,x_n)
\]
\footnote{$P<\mathfrak K$ означает: $P$ содержится в $\mathfrak K$, то есть является подполем $\mathfrak K$.} следует
\[
\overline P<\overline{\mathfrak K}<\overline P(x_1,\ldots,x_t),
\]
причем $\overline{\mathfrak K}$ имеет ту же степень трансцендентности $t$ относительно $\overline P$, что и $\overline P(x_1,\ldots,x_t)$. Следовательно, $\overline{\mathfrak K}$ имеет конечный рациональный базис относительно $\overline P$ как области коэффициентов, скажем $y_1,\ldots,y_t; z_1,\ldots,z_s$. При этом $z$ можно выбрать из $\mathfrak K$, так как $\overline{\mathfrak K}$ есть поле, порожденное $\mathfrak K$ и $\overline P$. Остается показать, что $\mathfrak K$ равно $\mathfrak M=P(y_1,\ldots,y_t;z_1,\ldots,z_s)$. Поле $\mathfrak M$ является промежуточным между $P(y_1,\ldots,y_t)$ и $\mathfrak K$; поэтому каждый элемент $\mathfrak K$ алгебраичен над $\mathfrak M$. С другой стороны, $\overline{\mathfrak M}$ равно $\overline{\mathfrak K}$, а значит $\mathfrak K$ содержится в $\overline{\mathfrak M}$. Поэтому $\mathfrak K$ является промежуточным полем между $\mathfrak M$ и $\overline{\mathfrak M}$; поскольку каждый элемент $\overline{\mathfrak M}$, не принадлежащий $\mathfrak M$, трансцендентен относительно $\mathfrak M$, необходимо $\mathfrak K$ совпадает с $\mathfrak M$.

Следствие: если $P<\mathfrak K<\mathfrak L<P(x_1,\ldots,x_n)$ и $\mathfrak L$ алгебраично над $\mathfrak K$, то $\mathfrak L$ является конечным алгебраическим расширением $\mathfrak K$. Ведь каждый элемент рационального базиса $\mathfrak L$ тогда алгебраичен над $\mathfrak K$; следовательно, $\mathfrak L$ характеризуется как конечное расширение $\mathfrak K$.

2. \textbf{Доказательство критерия конечности.} Условие явно необходимо: если $\mathfrak S$ является конечной областью целостности, то саму $\mathfrak S$ можно рассматривать как конечное подкольцо $\mathfrak R$, фигурирующее в критерии.

Условие достаточно при предположении, что поле корней $P^{1/p}$ является конечным расширением поля коэффициентов\footnote{См. с. 28, прим. 2.}, как только последнее имеет характеристику $p$; для характеристики нуль никакого ограничительного предположения не нужно. Для доказательства надо показать, что $\mathfrak S$ имеет конечный модульный базис относительно некоторого подкольца $\mathfrak R$.

Пусть $\mathfrak K$ обозначает поле частных\footnote{См. с. 28, прим. 2.} кольца $\mathfrak R$, а $\mathfrak L$ -- поле частных кольца $\mathfrak S$. Эти поля частных существуют, поскольку речь идет о кольцах без делителей нуля; имеем $P<\mathfrak K<\mathfrak L<P(x_1,\ldots,x_n)$. Следовательно, по следствию из п. 1 поле $\mathfrak L$ является конечным алгебраическим расширением $\mathfrak K$; далее, по предположению, каждый элемент $\mathfrak S$ цел над $\mathfrak R$, и тем самым $\mathfrak S$ становится подкольцом кольца $\mathfrak S'$, состоящего из всех элементов $\mathfrak L$, целых над $\mathfrak R$. Но поскольку ничего не предполагается об интегральной замкнутости $\mathfrak R$ в $\mathfrak K$, факт, что теорема о цепях делителей сохраняется при конечном расширении, нельзя применить непосредственно. Сначала надо перейти к также конечному подкольцу $\mathfrak T$ кольца $\mathfrak R$, для которого эта интегральная замкнутость выполнена.

Сначала предположим, что область коэффициентов $P$ содержит бесконечно много элементов. Тогда из $\mathfrak R$ можно выбрать неприводимую систему $g_1(x),\ldots,g_t(x)$ так, что $\mathfrak R$, а значит и $\mathfrak S$, цела над $\mathfrak T=P[g_1(x),\ldots,g_t(x)]$. Действительно, пусть $f_1(x),\ldots,f_k(x)$ -- интегральный базис $\mathfrak R$, существующий по предположению; если он не образует неприводимой системы, то после возможной перенумерации индексов пусть
\[
F(f_k;f_1,\ldots,f_{k-1})=0
\]
будет соотношением, которому удовлетворяет $f_k$. Если заменить $f_1,\ldots,f_{k-1}$ на
\[
 f_1+t_1f_k,\ldots,f_{k-1}+t_{k-1}f_k,
\]
то $f_k$, а следовательно и $f_1,\ldots,f_{k-1}$, целы над этими новыми функциями, когда $t_i$ выбраны как неопределенные. Так как $P$ имеет бесконечно много элементов, $t_i$ можно специализировать к элементам $P$ так, чтобы эта целая зависимость сохранилась. После конечного числа шагов, с учетом транзитивного закона целой зависимости, получаем искомое кольцо $\mathfrak T=P[g_1(x),\ldots,g_t(x)]$. Тогда $\mathfrak L$ является конечным расширением поля частных $P(g_1(x),\ldots,g_t(x))$ кольца $\mathfrak T$, а кольцо $\mathfrak S'$, состоящее из всех элементов $\mathfrak L$, целых над $\mathfrak R$, совпадает с кольцом всех элементов $\mathfrak L$, целых над $\mathfrak T$.

В $\mathfrak T$ выполняется теорема о цепях делителей для системы всех идеалов, и $\mathfrak T$ интегрально замкнуто в своем поле частных; оба утверждения следуют из того, что $\mathfrak T$ изоморфно полиномиальной области от $t$ неопределенных. Из предположения, что $P^{1/p}$ конечно над $P$, далее следует, что $\mathfrak T^{1/p}$ конечно над $\mathfrak T$.\footnote{Ср.: Hilbert, \emph{Über die vollen Invariantensysteme}, § 1. Math. Ann. 42 (1893), с. 313--373.}\footnote{Ср. заметку Artin--van der Waerden, цитированную во введении.} Следовательно, в $\mathfrak S'$ теорема о цепях делителей выполняется для системы всех $\mathfrak T$-модулей; в частности, $\mathfrak S$, как подкольцо $\mathfrak S'$, содержащее $\mathfrak T$, имеет конечный $\mathfrak T$-модульный базис. Существование конечного $\mathfrak T$-модульного базиса для $\mathfrak S$ далее остается верным без какого-либо ограничительного предположения об области коэффициентов $P$, если $\mathfrak L$ является расширением первого рода поля $P(g_1(x),\ldots,g_t(x))$, следовательно, в частности, всегда когда $P$ имеет характеристику нуль.\footnote{Ср., например: E. Noether, \emph{Abstrakter Aufbau der Idealtheorie in algebraischen Zahl- und Funktionenkörpern}, § 3. Math. Ann. 96 (1926), с. 26--61.}

Пусть $h_1(x),\ldots,h_s(x)$ -- $\mathfrak T$-модульный базис $\mathfrak S$. Тогда представление
\[
 f(x)=A_1(g)h_1(x)+\cdots+A_s(g)h_s(x)
\]
для произвольного $f(x)$ из $\mathfrak S$, где $A$, как элементы $\mathfrak T$, обозначают полиномы от $g$, показывает, что $g_1(x),\ldots,g_t(x),h_1(x),\ldots,h_s(x)$ образуют искомый интегральный базис.

Если поле $P$ содержит лишь конечное число элементов, то поле $P(u)$, полученное присоединением неопределенной $u$, то есть элемента, трансцендентного относительно $\mathfrak S$, содержит бесконечно много элементов. Поэтому кольцо $\mathfrak S(u)$ имеет конечный интегральный базис относительно $P(u)$ как области коэффициентов; и поскольку $\mathfrak S(u)$ является кольцом, порожденным $\mathfrak S$ и $P(u)$, этот базис можно выбрать из $\mathfrak S$, разложив по степеням $u$, хотя базисные элементы $g_i$ из $\mathfrak S$ теперь уже не будут образовывать неприводимую систему. Тем самым каждый полином $f(x)$ из $\mathfrak S$ допускает представление $C(u)f(x)=G(u,g_i,h_i)$, где $C(u)$ обозначает ненулевой полином из $P[u]$, а $G$ является целым полиномом по $u,g_i,h_i$. Сравнение коэффициентов при подходящей степени $u$ показывает, что $g_i(x)$ и $h_i(x)$ дают интегральный базис $\mathfrak S$. Тем самым \textbf{критерий конечности доказан в полной общности}.

\subsection*{§ 2. Конечность инвариантов}

1. Под \textbf{конечной линейной группой характеристики $p$} понимается группа $\mathfrak H$ из $h$ линейных подстановок $A_1,\ldots,A_h$ (с ненулевым детерминантом) с коэффициентами из поля $P$ характеристики $p$. При этом $A_k$ задается линейной подстановкой
\[
\begin{aligned}
 x_1^{(k)}&=a_{11}^{(k)}x_1+\cdots+a_{1n}^{(k)}x_n,\\[-0.2em]
 &\ldots,\\
 x_n^{(k)}&=a_{n1}^{(k)}x_1+\cdots+a_{nn}^{(k)}x_n,
\end{aligned}
\]
или сокращенно: $(x^{(k)})=A_k(x)$. Следовательно, группа $\mathfrak H$ переводит ряд $(x)$ с элементами $x_1,\ldots,x_n$ в ряды $(x^{(k)})$ с элементами $x_1^{(k)},\ldots,x_n^{(k)}$. Так как среди $A_1,\ldots,A_h$ должна содержаться тождественная подстановка, среди рядов $(x^{(k)})$ содержится также ряд $(x)$.

Под \textbf{целым рациональным (абсолютным) инвариантом группы} понимается такая целая рациональная функция от $x_1,\ldots,x_n$ с коэффициентами из $P$\footnote{Это не ограничивает общности, поскольку каждый инвариант с коэффициентами из поля характеристики $p$, которое должно иметь с $P$ общее надполе, чтобы инварианты были определены, можно линейно разложить через конечное число инвариантов с коэффициентами из $P$. В общем случае надо лишь выразить коэффициенты через такие, которые линейно независимы над $P$; тогда свойство инвариантности выполнено тождественно по этим независимым коэффициентам.}, которая при применении $A_1,\ldots,A_h$ тождественно остается неизменной. Следовательно, к этим инвариантам относятся все симметрические функции, с коэффициентами из $P$, от рядов $(x^{(k)})$. Обратно, для каждого такого инварианта имеем
\[
 f(x)=f(x^{(1)})=\cdots=f(x^{(h)})
 \quad\text{или}\quad
 h\cdot f(x)=f(x^{(1)})+\cdots+f(x^{(h)}).
\]
Отсюда следует, что в случае, когда порядок $h$ группы не делится на характеристику, в частности в случае характеристики нуль, симметрические функции рядов $(x^{(k)})$ с коэффициентами из $P$ исчерпывают всю совокупность инвариантов. Так как эти симметрические функции имеют конечный интегральный базис, доказательство конечности в этом случае завершено; одновременно интегральный базис получается как система $G_{\alpha\alpha_1\ldots\alpha_n}(x)$ коэффициентов резольвенты Galois\footnote{См. с. 28, прим. 1.}:
\[
 \Phi(z,u)=\prod_{k=1}^{h}\bigl(z-u_1x_1^{(k)}-\cdots-u_nx_n^{(k)}\bigr)
\]
\[
 =z^h+\sum G_{\alpha\alpha_1\ldots\alpha_n}(x)
 z^{\alpha}u_1^{\alpha_1}\cdots u_n^{\alpha_n}
 \quad(\alpha\ne h;\ \alpha+\alpha_1+\cdots+\alpha_n=h).
\]

2. В случае, когда $h$ делится на $p$, не обязательно каждый инвариант группы является симметрической функцией рядов $(x^{(k)})$, поскольку теперь из $h\cdot f(x)$ уже нельзя заключать о $f(x)$. Но кольцо $\mathfrak R$, полученное из конечного числа инвариантов $G_{\alpha\alpha_1\ldots\alpha_n}(x)$, то есть из коэффициентов резольвенты Galois, образует конечное подкольцо системы $\mathfrak S$ всех инвариантов, и надо показать, что $\mathfrak S$ относительно $\mathfrak R$ удовлетворяет предпосылкам критерия конечности.

Сначала без ограничения общности\footnote{См. с. 33, примечание.} в качестве области коэффициентов можно взять подполе $\overline P$ поля $P$, полученное из всех коэффициентов подстановок $a_{\mu\nu}^{(k)}$. Это поле, полученное из конечного числа элементов, возникает присоединением конечного числа алгебраических или трансцендентных элементов к простому полю; а поскольку простое поле совершенно, отсюда следует, что $\overline P^{1/p}$ конечно относительно $\overline P$.\footnote{См. с. 32, прим. 2.}

Далее надо показать, что $\mathfrak S$ алгебраически цела над $\mathfrak R$. По п. 1 среди $h$ рядов $x^{(k)}$ встречается также исходный ряд $x_1,\ldots,x_n$; значит, резольвента Galois $\Phi(z,u)$ для $z=u_1x_1+\cdots+u_nx_n$ тождественно обращается в нуль по $u$. Если каждый раз специализировать все $u$ в нуль, кроме $u_i$, которое полагается равным единице, получаем соотношение
\[
 x_i^h+\sum x_i^\alpha G_{\alpha0\ldots\alpha_i\ldots0}(x)=0
 \quad (\alpha\ne h;\ \alpha+\alpha_i=h),
\]
показывающее, что каждое $x_i$ алгебраически цело над $\mathfrak R$. Тогда то же верно для любого полинома от $x$; в частности, $\mathfrak S$ алгебраически цела над $\mathfrak R$. Тем самым предпосылки критерия конечности выполнены, и \textbf{конечность инвариантов доказана}.

3. Отметим, что те же рассуждения дают \textbf{конечность относительных и модулярных инвариантов}. Полином $f(x)$ называется \textbf{относительным инвариантом} группы, если все $f(x^{(k)})$ отличаются от $f(x)$ лишь ненулевым множителем из $P$. Система абсолютных инвариантов, в частности кольцо $\mathfrak R$, полученное из $G_{\alpha\alpha_1\ldots\alpha_n}(x)$, образует конечное подкольцо кольца, полученного из всех относительных инвариантов; и, как выше, предпосылки критерия конечности относительно $\mathfrak R$ выполнены.

Полином $f(x)$ называется \textbf{модулярным инвариантом} группы, если $f(x)$ становится равным $f(x^{(k)})$ для каждой специальной системы значений $x_1=\alpha_1,\ldots,x_n=\alpha_n$ из $P$. Каждый абсолютный инвариант одновременно является модулярным инвариантом; однако если $P$ содержит лишь конечное число элементов, то обратное утверждение не обязано быть верным. Здесь тоже предпосылки критерия конечности выполнены относительно подкольца $\mathfrak R$, полученного из $G_{\alpha\alpha_1\ldots\alpha_n}(x)$.

\clearpage
% END UNIT BODY
% END PRODUCER UNIT 113
\endgroup


\clearpage
\begingroup
\setcounter{footnote}{0}
% BEGIN PRODUCER UNIT 114
% Source: C:/Users/Floris/Documents/interlanguage/03_projects/noether/02_slavic_working_corpus/translations/paper30/intro/russian/v001/Noether_Paper30_Intro_Russian_v001.tex
% Identity: 15695 B / SHA-256 DD150B2211DC551EC345B706AC4FABD5C6ADE8F414359C261CDCD041DC873E2A
\setlength{\parindent}{1.25em}
\setlength{\parskip}{0pt}
\emergencystretch=25em
\sloppy
\hbadness=10000
\vbadness=10000
\hfuzz=2pt
\vfuzz=10pt
\raggedbottom
% BEGIN UNIT BODY
% Unit: Paper30 introduction v001. Direct Russian translation from the German final-audited source slice, source lines 1--135 / segments P30-S0001--P30-S0013.
\section*{30. Абстрактное построение теории идеалов в алгебраических числовых полях и полях функций}

\begin{center}
\emph{Math. Ann. 96 (1927), с. 26--61}
\end{center}

Далее дается абстрактная характеристика всех тех колец, теория идеалов которых совпадает с теорией идеалов всех целых элементов алгебраического числового поля, то есть таких колец, идеалы которых однозначно представляются как произведения степеней простых идеалов. К наиболее известным примерам таких колец относится также кольцо всех целых элементов алгебраического поля функций от одной неопределенной; более обще, от нескольких неопределенных, если, вслед за Kronecker, ограничиваться идеалами наибольшей размерности, то есть переходить к некоторому кольцу частных, функциональной области.

Аксиомы, полностью эквивалентные этой теории идеалов для положенного в основу коммутативного кольца, таковы:\footnote{Для определения основных понятий теории идеалов ср., например: E. Noether, \emph{Idealtheorie in Ringbereichen}, Math. Ann. 83 (1921), с. 24--66 (далее цитируется как \emph{Теория идеалов}). Кстати, наиболее существенные понятия кратко формулируются и в настоящей работе, особенно в §§1, 4 и 5. Ради полноты включено также кое-что такое, что не требуется для непосредственных целей этой работы.}

\begin{description}
\item[I. Условие цепей делителей.] Всякая цепь идеалов, в которой каждый идеал является собственным делителем предыдущего, обрывается за конечное число шагов; иначе говоря, в каждой цепи делителей идеалов существует индекс, начиная с которого все идеалы равны.
\item[II. Условие цепей кратных по модулю каждого ненулевого идеала.] Всякая цепь идеалов, все члены которой являются делителями фиксированного идеала, отличного от нулевого, и в которой каждый идеал является собственным кратным предыдущего, обрывается за конечное число шагов.
\item[III.] Существование единичного элемента для умножения.
\item[IV.] Кольцо без делителей нуля.
\item[V. Интегральная замкнутость в поле частных.] Каждый элемент поля частных, целый над данным кольцом, принадлежит самому кольцу.
\end{description}

Из этих аксиом я шаг за шагом развиваю все более ограниченную теорию идеалов вплоть до искомой; и, обратно, показываю, что в ней все аксиомы действительно выполнены.

Из условия цепей делителей I следует, как я показала (\emph{Теория идеалов}, §4), представление идеала как наименьшего общего кратного конечного числа первичных идеалов, принадлежащих различным простым идеалам. Я кратко повторяю здесь это доказательство (§6), поскольку его можно упростить, используя понятие частного идеалов, и поскольку я хочу исправить недосмотр: первоначальное доказательство в одном месте использует не предположенное существование единичного элемента умножения.\footnote{На этот недосмотр мне указал P. Urysohn; мое собственное изменение доказательства было несколько более громоздким, чем приведенное здесь, принадлежащее E. Artin. Он уже раньше самостоятельно заполнил этот пробел и также независимо от меня продумал упрощение с помощью частного идеалов. -- Еще одно упрощение, которое избегает введения «редуцированного представления», я беру из родственной постановки вопроса у W. Krull: \emph{Algebraische Theorie der Ringe III}, Math. Ann. 92 (1924), с. 181--213, теорема I.}

Предположение условия цепей кратных приводит (§7) к тому, что в кольце классов вычетов по каждому идеалу, отличному от нулевого, простой идеал не может иметь собственного делителя, кроме единичного идеала, содержащего все элементы. Поэтому первичные компоненты этих идеалов определяются однозначно, за исключением тех, которые принадлежат единичному идеалу. В несколько менее точной форме этот результат уже имеется у Masazo Sono;\footnote{Masazo Sono, \emph{On the Reduction of Ideals}, Memoirs of the College of Science, Kyoto University, Series A, 7 (1924), с. 191--204.} его предположение о существовании композиционного ряда в кольце классов вычетов эквивалентно «двойному условию цепей» в этом кольце, как я кратко показываю в конце (§10). Под двойным условием цепей я понимаю действие условий цепей делителей и кратных без ограничения идеалами, отличными от нулевого.

Если выполнена еще и аксиома III, то есть существует единичный элемент, единичный идеал не возникает как принадлежащий простой идеал; значит, первичные компоненты определяются однозначно, они попарно взаимно просты, а их наименьшее общее кратное равно их произведению. Аксиомы I, II, III выполнены для всех конечных порядков алгебраического числового поля; здесь уже Dedekind (Dirichlet--Dedekind, \emph{Zahlentheorie}, 3-е изд., 11-е дополнение, §172), не давая доказательства во всех деталях, высказал однозначное представление идеала как произведения попарно взаимно простых однотипных идеалов (первичных компонент).

Из предположения IV, то есть отсутствия делителей нуля, следует, что различные степени идеала, отличного от нулевого и единичного идеалов, все различны, а также что существует поле частных. Если наконец выполнена еще и аксиома V интегральной замкнутости в поле частных, то первичные идеалы, однозначно определенные благодаря IV, становятся степенями простых идеалов; тем самым получается обычная теорема разложения (§8). Это последнее доказательство опирается на следствие интегральной замкнутости, которое в случае числового поля уже вывел Dedekind (3-е изд., §172): истинно дробный элемент можно представить как частное целых элементов так, что даже квадрат числителя не делится на знаменатель. В более поздних изложениях вместо этого следствия, также как следствие интегральной замкнутости, появляется гораздо менее прозрачная обобщенная теорема Gauss или соответствующие, несколько более сильные модульные теоремы; все эти теоремы при применении предполагают базисные элементы, тогда как здесь работа ведется с самими идеалами.

В обратном направлении \hbox{\(\S 9\)} аксиомы, отличные от V, следуют почти непосредственно; сама V следует из доказательства того, что истинно дробный элемент всегда можно представить так, что никакая степень числителя не делится на знаменатель, а это является усилением следствия, первоначально полученного из V. Это доказательство удается лишь после того, как из представления идеалов выводятся обычные следствия: представление главными идеалами по модулю фиксированного идеала и теория дробных идеалов, то есть частных идеалов в поле. Факт, что интегральная замкнутость следует из представления степенями простых идеалов, также уже был известен Dedekind, как показывает замечание (3-е изд., §172, с. 522, внизу).

Аксиомы I--V дают, что четыре вообще различные разложения -- на наименьшие попарно взаимно простые идеалы, на взаимно простые идеалы, на наибольшие первичные компоненты и на неприводимые идеалы -- совпадают; и, обратно, из этого совпадения и аксиом I--IV следует интегральная замкнутость. Для колец с делителями нуля выполнение остальных аксиом, причем V следует определить как интегральную замкнутость в кольце частных, не обязательно влечет совпадение разложений; это показывает пример кольца классов вычетов по некоторым идеалам конечного порядка \hbox{\(\S 9\)}.

В §1 я набрасываю теорию элементов, целых над кольцом, вплоть до следствия Dedekind из интегральной замкнутости. В §2 я показываю, как условия цепей переносятся с кольца на модули из линейных форм, для которых это кольцо служит областью коэффициентов и множителей.\footnote{Эта форма модульных теорем принадлежит Prüfer (Neue Begründung der algebraischen Zahlentheorie, §2, Math. Ann. 94 (1924), с. 198--243) и прозрачнее обычного переноса базисных теорем.} Отсюда я вывожу (§3), что при переходе к целым элементам алгебраического расширения поля первого рода аксиомы I--IV сохраняются для каждого порядка, тогда как для порядка, состоящего из всех целых элементов, выполнена также V. Этот факт, для алгебраического числового поля, конечно, известный, позволяет подчинить упомянутые в начале поля функций; в частности, простое доказательство того, что функциональная область, выведенная из целочисленных многочленов от нескольких неопределенных, удовлетворяет аксиомам, полностью включает сюда и теорию идеалов Kronecker.\footnote{Различие, следовательно, появляется лишь в теоремах о дискриминантах, поскольку кольцо классов вычетов функциональной области по простому числу становится несовершенным полем, и потому могут возникать расширения второго рода.}

Теория идеалов, развиваемая далее, не предполагает этих §§2 и 3, служащих только для включения примеров. В §§4 и 5 даются основания, независимые от аксиом I--V; благодаря этому яснее, чем в первоначальном обосновании, видно, какие части теории не зависят от предположений конечности.
% END UNIT BODY
% END PRODUCER UNIT 114
\endgroup


\clearpage
\begingroup
\setcounter{footnote}{0}
% BEGIN PRODUCER UNIT 115
% Source: C:/Users/Floris/Documents/interlanguage/03_projects/noether/02_slavic_working_corpus/translations/paper30/section01/russian/v001/Noether_Paper30_Section01_Russian_v001.tex
% Identity: 24019 B / SHA-256 F60B8556267A431E67CF475EF911E7013853ECE7A96DC4FF50117C0862F23659
\setlength{\parindent}{1.25em}
\setlength{\parskip}{0pt}
\makeatletter
\renewcommand\footnotesize{\@setfontsize\footnotesize{7.7}{8.7}\selectfont}
\makeatother
\emergencystretch=35em
\allowdisplaybreaks
\sloppy
\hbadness=10000
\vbadness=10000
\hfuzz=2pt
\vfuzz=10pt
\raggedbottom
% BEGIN UNIT BODY
% Unit: Paper30 Section01 v001. Direct Russian translation from the German final-audited source slice, source lines 137--389 / segments P30-S0014--P30-S0038. Footnote counter continues after Paper30 introduction.
\setcounter{footnote}{5}

\subsection*{\S~1. Теория целых элементов}

Пусть задано коммутативное кольцо\footnote{Кольцо, как известно, определено, если на множестве задано отношение равенства, а также две операции, сложение и умножение, удовлетворяющие обычным законам: сложение ассоциативно, коммутативно и однозначно обратимо; умножение ассоциативно -- и коммутативно, когда речь идет о коммутативном кольце, -- и дистрибутивно относительно сложения. Если лежащее в основе определение равенства не является теоретико-множественным тождеством, то, чтобы в каждой подсистеме сохранялось то же определение равенства, такая подсистема (подкольцо, модуль, идеал) вместе с любым элементом должна содержать и все элементы, равные ему. Точно так же при каждом расширении нужно предполагать, что новое определение равенства включает старое. Все понятия вроде «конечно много элементов», «однозначное соответствие» и т. п. понимаются в смысле этого определения равенства; например, «существует один и только один элемент» означает, что существует только один элемент с точностью до равенства. Впрочем, от равенства всегда можно перейти к теоретико-множественному тождеству, объединяя равные элементы в один класс и вводя эти классы как новые элементы. Приведенные здесь замечания об определении равенства принадлежат R. Hölzer.} $T$ без делителей нуля, с единичным элементом $e$ для умножения (аксиомы $III$ и $IV$); в $T$ выделено фиксированное подкольцо $R$, содержащее единичный элемент. Понятия модуля, порядка и целости всюду относятся к этому фиксированному кольцу $R$; далее, ради краткости, это будет опускаться в обозначениях.\footnote{Теория целых элементов сохраняется, если допустить делители нуля и вместо этого предположить в $R$ условие цепей делителей; по §2 это также дает доказательство леммы. Так как эта форма теории целых элементов далее не используется, на нее будет указано только в примечаниях. Все рассматриваемые определения сохраняются при наличии делителей нуля; большинство из них, как известно и легко видно, требует еще более слабых предположений. Не сохраняется лишь дедекиндовская формулировка: «Число является целым, если оно имеет оболочку». Ибо при появлении делителей нуля частное двух конечных модулей не обязано оставаться конечным; поэтому здесь порядок определяется непосредственно, а не как частное модулей.} Элементы из $R$ будем обозначать латинскими буквами, а элементы из $T$ вообще -- греческими.

\textbf{1.} Система $M$ элементов из $T$ называется, как обычно, $R$-модулем, или кратко модулем, если вместе с двумя элементами $\alpha$ и $\beta$ она содержит и их разность, а вместе с $\alpha$ также $r\alpha$, где $r$ -- произвольный элемент из $R$. Говорят, что $M$ делится на $N$,
\[
        M\equiv 0 \pmod N,
\]
а также что $M$ является кратным, а $N$ делителем, если каждый элемент из $M$ содержится также в $N$. $R$-модули, все элементы которых принадлежат $R$, называются идеалами в $R$.

Очевидно, пересечение -- наименьшее общее кратное $[\ldots,M_i,\ldots]$ -- произвольного числа модулей из $T$ снова является модулем; так же и $T$ является модулем. Поэтому для произвольной системы $\Sigma$ элементов из $T$ однозначно существует порожденный ею модуль: это пересечение всех модулей, содержащих $\Sigma$. Сумма -- наибольший общий делитель $(\ldots,M_i,\ldots)$ -- произвольной системы модулей есть модуль, порожденный их объединением. Произведение $MN$ двух модулей определяется как модуль, порожденный элементами $\mu\nu$, где $\mu$ пробегает все элементы из $M$, а $\nu$ все элементы из $N$. Поэтому произведение удовлетворяет ассоциативному и коммутативному законам, поскольку они выполняются для элементов кольца. Из дистрибутивного закона в $T$ далее следует дистрибутивный закон для модулей:
\[
        (\ldots,M_i,\ldots)N=(\ldots,M_iN,\ldots),
\]
так как по тому же самому закону в $T$ правая и левая части являются модулем, порожденным всеми элементами $\mu_i\nu$.

Так как само $R$ является модулем, условие, что $R$ есть область множителей, можно записать как $RM\equiv0\pmod M$. А так как по предположению $R$ содержит единичный элемент, имеем также $M\equiv0\pmod{RM}$, и потому $M=RM$.

Модуль $M$ называется конечным, то есть конечно порожденным, если существует хотя бы одна система $\Sigma$ из конечного числа элементов, порождающая $M$; элементы $\mu_1,\ldots,\mu_n$ из $\Sigma$ называются модульным базисом для
\[
        M=(\mu_1,\ldots,\mu_n).
\]
Сумма и произведение двух, а значит и конечного числа, конечных модулей снова конечны, поскольку они порождаются соответственно объединением или произведениями базисных элементов. Последнее следует из представления $r_1\mu_1+\cdots+r_n\mu_n$ всех элементов $\mu$ из $M$ через базис и из коммутативного закона умножения в $T$.

\textbf{2.} Расширяющее кольцо $R$ в $T$, то есть кольцо в $T$, содержащее все элементы из $R$, называется $R$-порядком, или кратко порядком. Пересечение произвольного числа порядков из $T$ снова является порядком; так же и $T$ является порядком. Следовательно, для произвольной системы $\Sigma$ элементов из $T$ однозначно существует порожденный ею порядок. Суммы и произведения порядков определяются соответственно так же, как для модулей; в частности $O^2=O$ для каждого порядка.

Каждый порядок одновременно является $R$-модулем; порядок называется конечным, если он является конечным модулем. Так как сумма и произведение возникают путем образования модульной суммы и модульного произведения, по пункту 1 суммы и произведения конечных порядков снова являются конечными порядками. Кроме того, выполняется транзитивный закон: если $A$ -- конечный $R$-порядок, а $B$ -- конечный $A$-порядок, то $B$ является также конечным $R$-порядком. Действительно, $B$ одновременно становится $R$-порядком, значит $R$-модулем. Если же $\alpha_1,\ldots,\alpha_m$ образуют модульный базис $A$ относительно $R$, а $\beta_1,\ldots,\beta_n$ такой же базис $B$ относительно $A$, то произведения $\alpha_i\beta_j$, очевидно, образуют модульный базис $B$ относительно $R$.

\textbf{3.} Элемент $\alpha$ из $T$ называется целым над $R$, или кратко целым, если порядок $R_\alpha$, порожденный $\alpha$, конечен; иначе говоря, если цепь делителей модулей
\[
        U_n=(\alpha^0,\alpha,\ldots,\alpha^{n-1})
\]
обрывается за конечное число шагов, $U_n=U_{n+1}=\cdots$.

Определение допускает и вторую форму, которая будет использована в пункте 4: элемент $\alpha$ из $T$ называется целым, если существует хотя бы один главный модуль $C$, отличный от нулевого модуля, то есть $C$ порожден элементом $\gamma\ne0$, такой, что произведение порядка $R_\alpha$ с $C$ является конечным модулем; иначе говоря, что цепь модулей $CU_n$ обрывается за конечное число шагов.\footnote{Если в $R$ допускаются делители нуля, нужно требовать существования регулярного главного модуля; то есть $\gamma$ следует предполагать неделителем нуля, регулярным элементом, что для колец без делителей нуля сводится к $\gamma\ne0$.}

Наконец, это определение тождественно и обычной форме: элемент $\alpha$ из $T$ называется целым, если он удовлетворяет хотя бы одному уравнению
\[
        \alpha^n+r_1\alpha^{n-1}+\cdots+r_n=0,
\]
где все $r$ являются элементами из $R$.

Эти разные определения действительно тождественны. Поскольку $R_\alpha$ совпадает с модулем, порожденным всеми степенями $\alpha$, для конечного $R_\alpha$ всегда можно выбрать модульный базис из этих степеней. Если такой базис, например, имеет вид $e=\alpha^0,\alpha,\ldots,\alpha^{n-1}$, это означает, что цепь $U_n$ обрывается на $U_n$; и обратно, из $U_n=U_{n+1}=\cdots$ получается этот базис. Но обрывание цепи $U_n$ тождественно также существованию уравнения
\[
        \alpha^n+r_1\alpha^{n-1}+\cdots+r_n=0.
\]
Если $R_\alpha$ конечен, то конечен и $CR_\alpha$, значит цепь $CU_n$ обрывается. Обратно, из конечности $CR_\alpha$, то есть из обрывания $CU_n$, следует также обрывание $U_n$ и тем самым конечность $R_\alpha$. Действительно, поскольку $C$ предполагается главным модулем, отличным от нулевого модуля, из $CU_n=CU_{n+1}$ следует также $U_n=U_{n+1}$.

Кольцевое свойство целых элементов и транзитивный закон целости опираются на следующую лемму.

\medskip
\noindent\textbf{Лемма.} Пусть $O$ -- конечный порядок в $T$, а $\alpha$ -- произвольный элемент из $O$. Тогда порядок $R_\alpha$, порожденный $\alpha$, конечен; другими словами, все элементы конечного порядка являются целыми.

Доказательство получается обычными рассуждениями. Пусть $\omega_1,\ldots,\omega_n$ -- модульный базис $O$. Из кольцевого свойства $\alpha\omega_i$ также принадлежит $O$. Поэтому
\[
        \alpha\omega_i=r_{i1}\omega_1+\cdots+r_{in}\omega_n,
\]
и отсюда
\[
        \det(r_{ij}-e_{ij}\alpha)=0,
\]
где $e_{ij}=0$ при $i\ne j$, а $e_{ii}=e$; этим доказано, что $\alpha$ целый. Для зануления определителя используется предположение, что $T$, а следовательно и $O$, не имеет делителей нуля.\footnote{Если в $R$ предполагается условие цепей делителей, но допускаются делители нуля, то лемма становится непосредственным следствием модульной теоремы из §2, согласно которой условие цепей переносится на модули из $O$. Это доказательство для числового поля также принадлежит Dedekind (4-е изд., §173, II) и является первым применением условий цепей в литературе. В следствиях, которые будут даны в пункте 4, Dedekind вместо этого использует более специальный факт, что число классов вычетов по идеалу в алгебраическом числовом поле конечно.}

Система $G$ всех целых элементов из $T$ образует порядок. Она содержит $R$, ибо элементы из $R$ целые: $R$ является конечным $R$-модулем с единичным элементом как базисом. Но $G$ имеет и кольцевое свойство: порядок $R_{\alpha,\beta}$, порожденный двумя произвольными целыми элементами $\alpha,\beta$, равен произведению $R_\alpha R_\beta$ и поэтому конечен. Следовательно, по лемме, $\alpha+\beta$ и $\alpha\beta$ целые как элементы конечного порядка.

\medskip
\noindent\textbf{Транзитивный закон целости.} Если $O$ -- порядок, состоящий из целых элементов, то каждый элемент из $T$, целый над $O$, является также целым над $R$. По определению $\alpha$ удовлетворяет уравнению
\[
        \alpha^m+\omega_1\alpha^{m-1}+\cdots+\omega_m=0,
\]
а значит также целый над порядком $O'=R_{\omega_1,\ldots,\omega_m}$, порожденным $\omega_1,\ldots,\omega_m$; этот порядок конечен как произведение $R_{\omega_1}\cdots R_{\omega_m}$. По транзитивному закону из пункта 2 порожденный $\alpha$ конечный $O'$-порядок $O'_\alpha$ поэтому является также конечным $R$-порядком, и $\alpha$ целый как элемент конечного порядка, по лемме.

\textbf{4.} Кольцо $R$ называется интегрально замкнутым в $T$, если каждый элемент из $T$, целый над $R$, принадлежит $R$. По второй форме определения целых элементов в пункте 3 элемент $\alpha$ из $T$ принадлежит $R$ тогда и только тогда, когда существует хотя бы один главный модуль $C$, отличный от нулевого модуля, такой, что $CR_\alpha$ является конечным модулем.\footnote{Если в $R$ допускаются делители нуля, то $C$ снова нужно предполагать регулярным главным модулем; точно так же элемент $c$ в следствии $I$ нужно предполагать регулярным.}

Из понятия интегрально замкнутого кольца Dedekind (3-е изд., §172) вывел два важных следствия, позволяющих использовать это понятие в теории идеалов.

\medskip
\noindent\textbf{Следствие Dedekind $I$.} Пусть $R$ интегрально замкнуто в $T$, и пусть, как дополнительное предположение, в $R$ выполнено условие цепей делителей для идеалов (аксиома $I$). Если $\beta$ -- нецелый элемент из $T$, а $c\ne0$ -- элемент из $R$, то существует показатель $\rho>0$ такой, что в ряду элементов
\[
        c,\;c\beta,\ldots,c\beta^\nu,\ldots
\]
элементы $c,\ldots,c\beta^\rho$ целые, а все последующие -- нецелые.

Если бы все элементы $c\beta^\nu$ были целыми, то из-за интегральной замкнутости $R$ они принадлежали бы $R$. Тогда цепь модулей $C,CB_1,\ldots,CB_\nu,\ldots$, где $C$ обозначает модуль, порожденный $c$, а $B_\nu$ модуль, порожденный $1,\beta,\ldots,\beta^\nu$, перешла бы в цепь идеалов $C_\nu=(c,c\beta,\ldots,c\beta^\nu)$, которая по предположению обрывается за конечное число шагов. Но тогда, вопреки предположению, $R_\beta$ был бы конечен, а $\beta$ целым. Следовательно, среди элементов $c\beta^\nu$ есть нецелые. Далее, вместе с любым нецелым элементом нецелыми являются и все последующие: ибо если $c\beta^r$ целый, а значит принадлежит $R$, то при $\rho<r$ элемент $c\beta^\rho$ удовлетворяет уравнению
\[
        (c\beta^\rho)^r-c^{\,r-\rho}(c\beta^r)^\rho=0
\]
и поэтому тоже целый. Следовательно, если $c\beta^{\rho+1}$ -- первый нецелый элемент, то, поскольку $c$ целый, $\rho>0$, и это искомый показатель.

Если специализировать расширяющее кольцо $T$ до поля частных кольца $R$, то есть до поля, возникающего присоединением всех пар элементов,\footnote{В известной формулировке Steinitz, J. f. M. 137. Если в $R$ допускаются делители нуля, то вместо поля частных нужно брать кольцо частных, присоединяя все пары частных, у которых знаменатель является регулярным элементом из $R$. Здесь следствие Dedekind $II$ уже не выполняется, ибо $n$ вообще не будет регулярным элементом.} то отсюда следует:

\medskip
\noindent\textbf{Следствие Dedekind $II$.} Пусть $R$ интегрально замкнуто в своем поле частных $K$, и пусть в $R$ выполнено условие цепей делителей для идеалов. Если $\beta$ -- нецелый элемент из $K$, то $\beta$ допускает представление в виде частного элементов из $R$,
\[
        \beta=m/n,
\]
такое, что $m^2/n$ тоже нецелый.

Положим $\beta=b/c$ при $c\ne0$. Тогда в ряду $c,c\beta=b,c\beta^2,\ldots$ первые два элемента наверняка целые, следовательно $\rho>1$, где $\rho$ -- показатель из следствия $I$. Если положить
\[
        m=c\beta^\rho,\qquad n=c\beta^{\rho-1},
\]
то $m/n=\beta$, а $m^2/n=c\beta^{\rho+1}$ нецелый.

Транзитивный закон интегральной замкнутости имеет следующую форму: если $R$ -- произвольное кольцо в $T$, а $G$ обозначает систему всех элементов из $T$, целых над $R$, то $G$ состоит также из всех элементов из $T$, целых над $G$. Ибо каждый элемент, целый над $G$, по транзитивному закону из пункта 3 является также целым над $R$, а значит принадлежит $G$.
% END UNIT BODY
% END PRODUCER UNIT 115
\endgroup


\clearpage
\begingroup
\setcounter{footnote}{0}
% BEGIN PRODUCER UNIT 116
% Source: C:/Users/Floris/Documents/interlanguage/03_projects/noether/02_slavic_working_corpus/translations/paper30/section02/russian/v001/Noether_Paper30_Section02_Russian_v001.tex
% Identity: 10009 B / SHA-256 A373540A28151A21EACBEBF9037CAF370DC21E8F8C2C9C63A0E9CCF718DE4780
\setlength{\parindent}{1.25em}
\setlength{\parskip}{0pt}
\makeatletter
\renewcommand\footnotesize{\@setfontsize\footnotesize{7.7}{8.7}\selectfont}
\makeatother
\emergencystretch=35em
\allowdisplaybreaks
\sloppy
\hbadness=10000
\vbadness=10000
\hfuzz=2pt
\vfuzz=10pt
\raggedbottom
% BEGIN UNIT BODY
% Unit: Paper30 Section02 v001. Direct Russian translation from the German final-audited source slice, source lines 391--493 / segments P30-S0039--P30-S0050. Footnote counter continues after Paper30 Section01.
\setcounter{footnote}{11}

\subsection*{\S~2. Условия цепей в конечных модульных областях}

Модульная теорема, согласно которой переносятся условия цепей, в применении, рассматриваемом здесь, появляется только для модулей расширяющего кольца. Но так как доказательство в случае общей модульной области остается тем же самым, положим в основу именно такую область.

Система $M$ элементов $\alpha,\beta,\ldots$ называется модульной областью относительно кольца $R$, если в $M$ заданы две операции, однозначно приводящие к элементам из $M$: аддитивное соединение элементов и умножение элементов на элементы из $R$; если относительно сложения $M$ образует абелеву группу, тогда как для умножения выполнен ассоциативный закон; и если дистрибутивный закон выполнен в обеих формах.\footnote{Ср. \emph{Idealtheorie}, §9.}

Для модульной области, очевидно, сохраняются все определения модулей из §1, 1, которые не относятся к умножению. В частности, $M$ называется конечной модульной областью, если существует конечное число элементов $\xi_1,\ldots,\xi_k$ из $M$ таких, что $M$ равен модулю, порожденному $\xi_1,\ldots,\xi_k$, то есть системе всех линейных форм $r_1\xi_1+\cdots+r_k\xi_k$, когда в $R$ существует единичный элемент, тогда как в противном случае добавляются еще дополнительные члены $n_i\xi_i$.

\medskip
\noindent\textbf{Модульная теорема.} Пусть $M$ -- конечная модульная область относительно коммутативного кольца $R$ с единичным элементом, и пусть в $R$ выполнено условие цепей делителей соответственно условие цепей кратных для идеалов. Тогда в $M$ выполнено условие цепей делителей соответственно условие цепей кратных для модулей.

Доказательство опирается на то, что каждому модулю из $M$ однозначно ставится в соответствие система конечного числа идеалов из $R$ так, что каждому собственному делителю соответственно собственному кратному модуля каждый раз соответствует по меньшей мере один собственный делитель соответственно собственное кратное среди идеалов.

Элемент из $M$ называется элементом длины $i$, если он допускает по меньшей мере одно представление как линейная форма от $\xi_1,\ldots,\xi_i$, но не допускает представления как линейная форма от $\xi_1,\ldots,\xi_{i-1}$. Произвольному модулю $W$ из $M$ теперь поставим в соответствие $k$ идеалов $\mathfrak a_1,\ldots,\mathfrak a_k$ из $R$: под $W_i$ будем понимать модуль всех элементов из $W$ длины $\le i$, а $\mathfrak a_i$ будет обозначать систему всех коэффициентов при $\xi_i$ в $W_i$. Сначала получается следующее.

Если $B$ -- собственный делитель $W$, а $\mathfrak b_1,\ldots,\mathfrak b_k$ -- идеалы, поставленные в соответствие $B$, то среди $\mathfrak b_i$ имеется по меньшей мере один собственный делитель соответствующего $\mathfrak a_i$. Действительно, из $B\equiv0\pmod W$ следует $B_i\equiv0\pmod {W_i}$, а значит $\mathfrak b_i\equiv0\pmod{\mathfrak a_i}$ для $i=1,2,\ldots,k$. По предположению в $B$ должны существовать элементы, не содержащиеся в $W$, а значит и такие элементы наименьшей длины $\ell$. Тогда $\mathfrak b_\ell$ является собственным делителем $\mathfrak a_\ell$, ибо $\mathfrak b_\ell\equiv0\pmod{\mathfrak a_\ell}$, если
\[
        \beta=b_1\xi_1+\cdots+b_\ell\xi_\ell
\]
есть такой элемент наименьшей длины из $B$. Из $\mathfrak b_\ell=\mathfrak a_\ell$ действительно следует существование элемента
\[
        \alpha=a_\ell\xi_\ell+\cdots+a_1\xi_1\equiv0\pmod W,
\]
а тем самым существование элемента $\beta-\alpha\in B$, $\beta-\alpha\notin W$, меньшей длины, чем $\ell$.

Пусть теперь дана цепь делителей соответственно кратных
\[
        W^{(1)},W^{(2)},\ldots,W^{(\nu)},\ldots
\]
и пусть в $R$ выполнено условие цепей делителей соответственно кратных. Если $\mathfrak a_{1\nu},\ldots,\mathfrak a_{k\nu}$ обозначают идеалы, поставленные в соответствие $W^{(\nu)}$, то ряды
\[
        \mathfrak a_{i1},\mathfrak a_{i2},\ldots
\]
образуют цепи делителей соответственно кратных для $i=1,\ldots,k$. Следовательно, по предположению существует индекс $\nu_0$ такой, что идеал $\mathfrak a_{i\nu_0}$ равен всем последующим для каждого $i$. Но тогда модуль $W^{(\nu_0)}$ по доказанному выше равен всем последующим модулям; тем самым условия цепей перенесены.

Отсюда непосредственно получается следующее.

\medskip
\noindent\textbf{Следствие модульной теоремы.} Если в $R$ выполнено условие цепей кратных по модулю каждого идеала, отличного от нулевого идеала, то в $M$ выполнено условие цепей кратных по модулю каждого модуля $C$, все поставленные в соответствие которому идеалы $\mathfrak c_1,\ldots,\mathfrak c_k$ отличны от нулевого идеала. Действительно, при этих предположениях цепи кратных $\mathfrak a_{i1},\ldots,\mathfrak a_{i\nu},\ldots$ обрываются за конечное число шагов.

\medskip
\noindent\textbf{Дополнительное замечание.} Если $R$ -- кольцо без единичного элемента, то все еще переносится условие цепей делителей и условие цепей кратных по модулю идеалов, отличных от нулевого идеала. Именно, вместо $M$ рассматривают систему $M'$ всех элементов вида
\[
        a_1\xi_1+\cdots+a_k\xi_k+n_1\xi_1+\cdots+n_k\xi_k,
\]
которая снова переходит в $M$, если дополнительные целые кратные $n_i\xi_i$ заменить элементами из $M$. Такому $M'$ можно, соответственно сказанному выше, поставить в соответствие $2k$ идеалов, причем $\mathfrak a_i$ означают идеалы из $R$, а $\mathfrak n_i$ -- идеалы в целых числах. Так как для $\mathfrak n_i$ выполнено условие цепей делителей и условие цепей кратных по модулю каждого идеала, отличного от нулевого, доказательство условия цепей делителей сохраняется для $M'$ и тем самым для $M$; для условия цепей кратных, однако, нужно требовать, чтобы все $2k$ идеалов, поставленные в соответствие $C$ соответственно $C'$, были отличны от нулевого идеала.
% END UNIT BODY
% END PRODUCER UNIT 116
\endgroup


\clearpage
\begingroup
\setcounter{footnote}{0}
% BEGIN PRODUCER UNIT 117
% Source: C:/Users/Floris/Documents/interlanguage/03_projects/noether/02_slavic_working_corpus/translations/paper30/section03/russian/v001/Noether_Paper30_Section03_Russian_v001.tex
% Identity: 16292 B / SHA-256 A5315508E5B5DCD877C639F5643A0355A35854E5D6AF54EAF7DD5EB0EE2E419F
\setlength{\parindent}{1.25em}
\setlength{\parskip}{0pt}
\makeatletter
\renewcommand\footnotesize{\@setfontsize\footnotesize{7.7}{8.7}\selectfont}
\makeatother
\emergencystretch=35em
\allowdisplaybreaks
\sloppy
\hbadness=10000
\vbadness=10000
\hfuzz=2pt
\vfuzz=10pt
\raggedbottom

\providecommand{\mR}{\mathfrak{R}}
\providecommand{\mK}{\mathfrak{K}}
\providecommand{\mL}{\mathfrak{L}}
\providecommand{\mS}{\mathfrak{S}}
\providecommand{\mC}{\mathfrak{C}}
\providecommand{\mF}{\mathfrak{F}}
\providecommand{\mc}{\mathfrak{c}}
% BEGIN UNIT BODY
% Unit: Paper30 Section03 v001. Direct Russian translation from the German final-audited source slice, source lines 516--540 / segments P30-S0052--P30-S0064. Footnote counter continues after Paper30 Section02.
\setcounter{footnote}{12}

\subsection*{\S~3. Переход к конечным расширениям полей}

Для включения числовых полей и полей функций нужно показать, как аксиомы $I$--$V$, сформулированные во введении, переносятся при переходе к конечным расширениям полей.

1. \emph{Пусть в $\mR$ выполнены все аксиомы $I$--$V$, кроме аксиомы $II$ -- условия цепей кратных. Пусть $\mK$ обозначает поле частных $\mR$, а $\mL$ -- конечное расширение поля $\mK$ первого рода.\footnote{Алгебраическое расширение поля, по Steinitz, называется «первого рода», если каждый элемент является нулем простой функции, которая в подходящем расширении поля распадается на попарно различные линейные множители.} Тогда в системе $\mS$ всех элементов из $\mL$, целых над $\mR$, выполнены те же аксиомы; в каждом порядке, содержащемся в $\mS$, также выполнены все эти аксиомы, кроме интегральной замкнутости.}

По §1, 3 система $\mS$ образует кольцо, для которого выполнены аксиомы $III$ и $IV$ -- существование единичного элемента и отсутствие делителей нуля. По заключению §1, 4 система $\mS$ интегрально замкнута в $\mL$; одновременно $\mL$ становится полем частных $\mS$, поскольку каждый элемент из $\mL$ после умножения на подходящий элемент из $\mR$ переходит в элемент из $\mS$. Значит, аксиома $V$ выполнена.

Доказательство $I$ получается из обычных рассуждений. Как расширение первого рода, $\mL$ получается присоединением к $\mK$ одного элемента $\alpha$, который, по предыдущему, можно считать целым над $\mR$. Наряду с $\alpha$ целыми являются и сопряженные элементы $\alpha',\alpha'',\ldots$, содержащиеся в галуаовом поле $\mL$ над $\mK$; следовательно, целым является и произведение их разностей, а также квадрат $D$ этого произведения разностей. Так как $\alpha,\alpha',\ldots$ все различны, $D$ является ненулевым элементом из $\mK$, а значит, по интегральной замкнутости $\mR$ в $\mK$, элементом из $\mR$. По обычному рассуждению, в силу интегральной замкнутости $\mR$ система $\mS$ содержится в конечной $\mR$-модульной области $M$, порожденной элементами $\xi_i=\alpha^i/D$; для этого нужно лишь в представлении элемента из $\mS$ через $\xi_i$ перейти к сопряженным элементам и решить линейную систему уравнений для коэффициентов представления. По модульной теореме из §2 условие цепей делителей выполнено для всех $\mR$-модулей из $M$, в частности для всех идеалов из $\mS$. Точно так же условие цепей делителей выполнено для идеалов любого порядка из $\mS$, поскольку и здесь речь идет об $\mR$-модулях в $M$; для каждого порядка, кроме того, выполнены аксиомы $III$ и $IV$.

2. \emph{Если в $\mR$ выполнены все аксиомы $I$--$V$, то то же самое выполнено и в $\mS$. В каждом порядке из $\mS$ также выполнены все аксиомы, кроме интегральной замкнутости.}

С учетом 1 остается доказать только выполнение аксиомы $II$. По следствию модульной теоремы из §2 достаточно следующего доказательства: если ненулевой идеал из $\mS$ рассматривать как $\mR$-модуль $\mC$ в $M$, то связанные с $\mC$ идеалы $\mc_i$ из $\mR$ все отличны от нулевого идеала. Действительно, $\mC$ имеет тот же конечный линейный ранг над $\mK$, что и $M$, поскольку вместе с любым элементом $\gamma$ модуль $\mC$ содержит также $\gamma\alpha^i$. Поэтому для каждого $\xi_i$ существует элемент $c_i\ne0$ из $\mR$ такой, что $c_i\xi_i$ принадлежит $\mC$; из-за единственности представления через $\xi_i$ -- индекс должен пробегать только значения, меньшие степени поля, -- $c_i$ принадлежит $\mc_i$, и потому $\mc_i$ отличен от нулевого идеала. Это рассуждение сохраняется для всех порядков того же ранга, что и $M$; для порядка меньшего ранга переходят к его полю частных, которое как промежуточное поле между $\mK$ и $\mL$ также является первым родом, и степень которого над $\mK$ совпадает с линейным рангом порядка; значит, те же рассуждения сохраняются. Наконец отметим, что порядок в $\mS$, отличный от самой $\mS$, никогда не является интегрально замкнутым: ведь каждый порядок содержит также $\mR$, и при интегральной замкнутости должен был бы содержать и $\mS$.

Для подчинения числовых полей и полей функций, следовательно, достаточно проверить выполнение аксиом $I$--$V$ для основных областей: целых чисел, многочленов от одной неопределенной и функциональной области многочленов от нескольких неопределенных.

3. \emph{Для кольца $\mR$ целых рациональных чисел, соответственно многочленов от одной неопределенной с коэффициентами из поля, выполнены аксиомы $I$--$V$.} Здесь $I$ и $II$ следуют либо из того, что по модулю каждого ненулевого числа, соответственно такого многочлена от одной неопределенной, существует лишь конечное число, соответственно конечное число линейно независимых, классов вычетов; либо также из того факта, что каждый идеал является главным. Ведь отсюда следует единственное разложение идеалов в произведение степеней простых идеалов, и поэтому ненулевой идеал делится лишь на конечное число идеалов. Аксиомы $III$ и $IV$ очевидно выполнены, причем последняя является следствием определения равенства. Интегральная замкнутость $V$ следует из того, что благодаря единственному разложению элементов из $\mR$ на неприводимые с точностью до единиц -- делителей единицы -- каждый элемент поля частных, не содержащийся в $\mR$, можно представить как частное двух взаимно простых элементов из $\mR$; значит, никакая степень числителя не делится на знаменатель. Такой элемент, следовательно, не может удовлетворять никакому уравнению, характеризующему его как целый над $\mR$.

4. \emph{Пусть теперь $\mR$ обозначает кольцо всех многочленов от нескольких неопределенных с коэффициентами из поля или с целыми рациональными коэффициентами.} Тогда аксиомы $III$, $IV$, $V$ выполнены так же, как в 3, потому что и здесь имеет место единственное разложение элементов на неприводимые с точностью до единиц. То, что условие цепей делителей $I$ выполнено, является непосредственным следствием теоремы Hilbert о существовании базиса идеала; условие цепей кратных $II$ здесь, однако, не выполнено.

Однако переходом к \emph{функциональной области}, что соответствует ограничению на идеалы наибольшей размерности, можно получить также действие аксиомы $II$; действие $I$ тогда можно доказать как в 3, без привлечения теоремы Hilbert. Именно, присоединим к $\mR$ неопределенную $u$ и рассмотрим кольцо $\mR^*$ всех многочленов от $u$ с коэффициентами из $\mR$, причем равенство определяется как равенство коэффициентов. Как обычно, многочлен из $\mR^*$ называется примитивным относительно $\mR$, если наибольший общий делитель его коэффициентов -- делитель в смысле многочлена из $\mR$, а не делитель-идеал -- является единицей из $\mR$. Тогда каждый многочлен из $\mR^*$ является произведением элемента из $\mR$ и примитивного многочлена из $\mR^*$, а произведение примитивных многочленов из $\mR^*$ снова примитивно.

\emph{Итак, совокупность элементов поля частных $\mR^*$, знаменатели которых являются примитивными многочленами из $\mR^*$, образует кольцо, функциональную область $\mF$ кольца $\mR$.} В этой функциональной области $\mF$ единицами являются все и только те элементы, числитель и знаменатель которых являются примитивными многочленами из $\mR^*$; значит, каждый элемент из $\mF$ единственным образом представляется как произведение элемента из $\mR$ и единицы из $\mF$. Поэтому теорема о разложении элементов из $\mR$ переносится на элементы из $\mF$; для $\mF$, следовательно, наряду с аксиомами $III$ и $IV$ также выполнена $V$, как в 3.

Кроме того, в $\mF$ каждый идеал является главным. Идеал вместе с любым элементом из $\mF$ содержит также полученный из него умножением на единицу элемент из $\mR$; а вместе с любыми двумя элементами из $\mR$ содержит и их наибольший общий делитель. Ведь вместе с $f(x)=t(x)\overline f(x)$ и $g(x)=t(x)\overline g(x)$ идеалу принадлежит также $t(x)(\overline f(x)+u\overline g(x))$, а значит и $t(x)$, если $\overline f$ и $\overline g$ взаимно просты и их линейная комбинация поэтому является единицей. Если, следовательно, начать с произвольного элемента $f(x)$ идеала и если в идеале есть элемент $g(x)$, который на него не делится, то наибольший общий делитель $t(x)$ также принадлежит идеалу и является собственным делителем $f(x)$; поэтому конечное повторение приводит к базисному элементу идеала, и идеал распознается как главный. Следовательно, в $\mF$ имеет место единственное разложение идеалов в произведения степеней простых идеалов, соответствующее разложению элементов из $\mR$; отсюда, как в 3, следует выполнение аксиом $I$ и $II$. Рассматриваемые здесь расширения поля частных $\mF$ являются только такими, которые можно породить присоединением элемента из $\mS$.\footnote{Систему элементов, целых над $\mF$, в этих расширениях полей получают также, переходя от $\mS$ к соответствующей функциональной области так же, как от $\mR$ к $\mF$. Ср., например, J. König, \emph{Algebraische Größen}, Leipzig 1903, с. 468/69.} С учетом 1 и 2 этим доказано действие аксиом $I$--$V$ для числовых полей и полей функций.
% END UNIT BODY
% END PRODUCER UNIT 117
\endgroup


\clearpage
\begingroup
\setcounter{footnote}{0}
% BEGIN PRODUCER UNIT 118
% Source: C:/Users/Floris/Documents/interlanguage/03_projects/noether/02_slavic_working_corpus/translations/paper30/section04/russian/v001/Noether_Paper30_Section04_Russian_v001.tex
% Identity: 17613 B / SHA-256 6D4728DB8C14F1CD987399CA01298F7C3368216394BBC515B439717F75609DAF
\setlength{\parindent}{1.25em}
\setlength{\parskip}{0pt}
\makeatletter
\renewcommand\footnotesize{\@setfontsize\footnotesize{7.7}{8.7}\selectfont}
\makeatother
\emergencystretch=35em
\allowdisplaybreaks
\sloppy
\hbadness=10000
\vbadness=10000
\hfuzz=2pt
\vfuzz=10pt
\raggedbottom

\providecommand{\mR}{\mathfrak{R}}
\providecommand{\mA}{\mathfrak{A}}
\providecommand{\mB}{\mathfrak{B}}
\providecommand{\mC}{\mathfrak{C}}
\providecommand{\ma}{\mathfrak{a}}
\providecommand{\mb}{\mathfrak{b}}
\providecommand{\mc}{\mathfrak{c}}
\providecommand{\mm}{\mathfrak{m}}
\providecommand{\mo}{\mathfrak{o}}
\providecommand{\mv}{\mathfrak{v}}
% BEGIN UNIT BODY
% Unit: Paper30 Section04 v001. Direct Russian translation from the German final-audited source slice, source lines 542--596 / segments P30-S0065--P30-S0089. Footnote counter continues after Paper30 Section03.
\setcounter{footnote}{14}

\subsection*{\S~4. Теоремы об изоморфизме. Прямые суммы}

Далее предполагается только свойство кольца, соответственно свойство модуля, но никакая другая аксиома.

1. Пусть $M$ и $\overline M$ являются модульными областями относительно $\mR$ (§2). Тогда $M$ называется гомоморфным $\overline M$ (точнее: модульно-гомоморфным), $M\sim\overline M$, если каждому элементу из $M$ соответствует один и только один элемент из $\overline M$ так, что этим исчерпывается $\overline M$;\footnote{Все в смысле определения равенства, действующего в $\overline M$; ср. примечание 6.} и если при этом соответствии разность и умножение на один и тот же элемент из $\mR$ соответствуют друг другу. То есть из $\beta\sim\overline\beta$ и $\gamma\sim\overline\gamma$ всегда следует:
\[
(\beta-\gamma)\sim(\overline\beta-\overline\gamma),\qquad r\beta\sim r\overline\beta .
\]

Следовательно, $\mR$-модулю $\mB$ из $M$ гомоморфно соответствует $\mR$-модуль $\overline{\mB}$ из $\overline M$; а совокупность $\mC^*$ элементов из $M$, которым соответствуют элементы $\mR$-модуля $\overline{\mC}$ из $\overline M$, образует в $M$ однозначно определенный по $\overline{\mC}$ модуль $\mC^*$, ассоциированный с $\overline{\mC}$, причем $\mC^*\sim\overline{\mC}$. Если, в частности, $\mA$ есть модуль, ассоциированный с нулевым элементом из $\overline M$, то $\mC^*$ является делителем $\mA$ и распадается на классы элементов из $M$, сравнимых по модулю $\mA$, так что эти классы взаимно однозначно соответствуют элементам из $\overline{\mC}$. Если перейти от $\mB$ в $M$ к $\overline{\mB}$ в $\overline M$, а затем обратно к $\mB^*$, то получается $\mB^*=(\mB,\mA)$, то есть наибольший общий делитель $\mB$ и $\mA$; следовательно, $\mB^*=\mB$, если $\mB$ является делителем $\mA$.

Если соответствие между элементами из $M$ и $\overline M$ обратимо однозначно, то эти области называются изоморфными (точнее: модульно-изоморфными), $M\cong\overline M$.

Пусть $\mA$ является произвольным $\mR$-модулем из $M$. Тогда модульная область $\overline M$, гомоморфная $M$, а именно модуль классов вычетов $M\mid\mA$, возникает, если сравнимость по модулю $\mA$ принять за новое отношение равенства. Каждому элементу из $M$ при этом сопоставляются все и только те элементы, которые равны ему в $\overline M$. Переход в $\overline M$ от определения равенства к тождеству означает, что все равные в $\overline M$ элементы объединяют в один класс -- класс вычетов -- и рассматривают эти классы вычетов как новые элементы $\overline M$.

\emph{Каждый гомоморфизм порождается переходом к модулю классов вычетов;} действительно, если $M\sim\overline M$, а $\mA$ есть модуль, ассоциированный с нулевым элементом из $\overline M$, то, как показано выше, $\overline M$ изоморфен модулю классов вычетов $M\mid\mA$.

2. \textbf{Первая теорема об изоморфизме.} \emph{Пусть $\overline M$ обозначает модуль классов вычетов $M\mid\mA$, и пусть $\mC$ является делителем $\mA$. Тогда имеет место изоморфизм
\[
\overline M\mid\overline{\mC}\cong M\mid\mC .
\]}
Действительно, сравнимость по модулю $\mA$ одновременно является сравнимостью по модулю $\mC$; элементы, равные по модулю $\mA$, остаются равными по модулю $\mC$. Поэтому можно образовать модуль классов вычетов $M\mid\mC$ -- то есть равенство по модулю $\mC$ -- так: сначала отождествить элементы по модулю $\mA$, то есть перейти к $\overline M$, а затем среди них объединить те, которые равны по модулю $\mC$; в $\overline M$ это соответствует отождествлению по модулю $\overline{\mC}$, то есть образованию $\overline M\mid\overline{\mC}$.

\textbf{Вторая теорема об изоморфизме.} \emph{Если $\mB$ и $\mA$ являются модулями из $M$, то имеет место изоморфизм
\[
(\mB,\mA)\mid\mA\cong \mB\mid[\mB,\mA].
\]}
Действительно, по 1. $\mB$ становится гомоморфным $\overline{\mB}$, если положить $(\mB,\mA)\mid\mA=\overline{\mB}$; а так как при этом нулевому элементу в $\overline{\mB}$ соответствуют все элементы из $[\mB,\mA]$ и только они, то указанный изоморфизм снова следует из 1.

3. Пусть $\mR$ и $\overline{\mR}$ являются (коммутативными) кольцами. Тогда $\mR$ называется гомоморфным $\overline{\mR}$ (точнее: кольцево-гомоморфным), $\mR\sim\overline{\mR}$, если каждому элементу из $\mR$ соответствует один и только один элемент из $\overline{\mR}$ так, что этим исчерпывается $\overline{\mR}$, и если при этом соответствии разности и произведения соответствуют друг другу. Если соответствие элементов обратимо однозначно, то кольца называются изоморфными (точнее: кольцево-изоморфными), $\mR\cong\overline{\mR}$.

Все рассуждения из 1. и 2. сохраняются, если модули заменить идеалами в $\mR$,\footnote{Для некоммутативных колец в основу должны быть положены двусторонние идеалы.} а понятие модульного гомоморфизма и модульного изоморфизма заменить гомоморфизмом и изоморфизмом колец. В частности, если $\mc$ является делителем $\ma$, то кольцо классов вычетов $\mc\mid\ma$ возникает введением сравнимости по модулю $\ma$ как нового отношения равенства; при этом нужно учитывать, что теперь определены также произведения классов вычетов.

Тогда $\mc$ гомоморфен $\mc\mid\ma$; каждый гомоморфизм порождается таким способом, и теоремы об изоморфизме имеют силу так же, как в 2:

\textbf{Первая теорема об изоморфизме.} Если $\overline{\mR}$ обозначает кольцо классов вычетов $\mR\mid\ma$, а $\mc$ является делителем $\ma$, то $\overline{\mR}\mid\overline{\mc}\cong\mR\mid\mc$ в смысле изоморфизма колец.

\textbf{Вторая теорема об изоморфизме.} Если $\mb$ и $\ma$ являются идеалами из $\mR$, то имеет место изоморфизм колец:
\[
(\mb,\ma)\mid\ma\cong\mb\mid[\mb,\ma].
\]\footnote{Эти теоремы об изоморфизме, известные из теории групп, для модулей в несколько более специальной форме впервые встречаются у Dedekind (ср. 3-е изд., с. 484). Приведенная выше форма для идеалов есть у Masazo Sono: \emph{On Congruences. I, II, III, IV}, Memoirs of the College of Science, Kyoto Imperial University, 2 (1917), 3 (1918, 1919), ср. 2, с. 215. В каждом случае явно сформулирована только вторая теорема об изоморфизме.}

4. Далее будут собраны правила вычисления для взаимно простых идеалов,\footnote{Именно в форме, восходящей к Dedekind (4-е изд., §178, III--VIII). Вычисления с единичным элементом, повсеместно встречающиеся в более новых изложениях, значительно громоздче.} из которых в 5. будут следовать теоремы о прямых суммах. В коммутативном кольце $\mR$ предполагается существование единичного элемента $e$ умножения. Следовательно, для каждого идеала из $\mR$ имеем $\ma=\mo\ma$, где под $\mo$ понимается единичный идеал. Два идеала $\ma,\mb$ называются взаимно простыми, если их наибольший общий делитель $(\ma,\mb)$ равен $\mo$.

\emph{4$\alpha$. Если $(\ma,\mb)=\mo$, то $(\ma,\mb\mc)=(\ma,\mc)$.} Действительно,
\[
(\ma,\mc)=(\ma,\mb)\cdot(\ma,\mc)=(\ma^2,\ma\mb,\ma\mc,\mb\mc)
\]
делится на $(\ma,\mb\mc)$ и обратно. Отсюда следует:

\emph{4$\beta$. Из $\mc\mb\equiv0\pmod{\ma}$ и $(\ma,\mb)=\mo$ следует $\mc\equiv0\pmod{\ma}$.} Действительно, $\mc\mb\equiv0\pmod{\ma}$ равносильно $(\ma,\mb\mc)=\ma$; следовательно, из-за $(\ma,\mb)=\mo$ получается также $(\ma,\mc)=\ma$, или $\mc\equiv0\pmod{\ma}$. Далее:

\emph{4$\gamma$. Если каждый из идеалов $\ma_1,\ldots,\ma_r$ взаимно прост с каждым из идеалов $\mb_1,\ldots,\mb_s$, то произведение $\ma_i$ взаимно просто с произведением $\mb_i$.} Действительно, из $(\ma,\mb)=\mo$ и $(\ma,\mc)=\mo$ получается $(\ma,\mb\mc)=\mo$, а конечное повторение дает утверждение.

Из 4$\alpha$ и 4$\gamma$ получаем:

\emph{4$\delta$. Если идеалы $\mb_1,\ldots,\mb_r$ попарно взаимно просты, то есть $(\mb_i,\mb_k)=\mo$ при $i\ne k$, то их наименьшее общее кратное равно их произведению.} Пусть $(\ma,\mb)=\mo$, $\mv=[\ma,\mb]$. Тогда
\[
\mv=\mo\mv=(\ma\mv,\mb\mv)=0(\ma,\mb);
\]
так как $\ma\mb\equiv0\pmod{\mv}$, отсюда следует $\mv=\ma\mb$, а конечное повторение с учетом 4$\gamma$ дает утверждение.

\emph{4$\varepsilon$. Пусть идеалы $\mb_1,\ldots,\mb_r$ попарно взаимно просты, и пусть}
\[
\ma_i=[\mb_1,\ldots,\mb_{i-1},\mb_{i+1},\ldots,\mb_r]
=\mb_1\cdots\mb_{i-1}\mb_{i+1}\cdots\mb_r .
\]
\emph{Тогда $(\ma_1,\ldots,\ma_{i-1},\ma_{i+1},\ldots,\ma_r)=\mb_i$, а следовательно $(\ma_1,\ldots,\ma_r)=\mo$.} Действительно, по дистрибутивному закону
\[
(\ma_1,\ma_2)=\mb_3\cdots\mb_r(\mb_2,\mb_1)=\mb_3\cdots\mb_r,
\]
и поэтому
\[
(\ma_1,\ma_2,\ma_3)=\mb_4\cdots\mb_r(\mb_3,\mb_1\mb_2)=\mb_4\cdots\mb_r,
\]
откуда конечным повторением при подходящей нумерации следует утверждение.

Идеал $\ma_i$ называется комплементом $\mb_i$ в представлении $\mm=[\mb_1,\ldots,\mb_r]$.

5. В коммутативном кольце $\mR$ снова предполагается существование единичного элемента умножения. Кольцо $\mR$ называется прямой суммой идеалов $\ma_1,\ldots,\ma_r$ -- в обозначении
\[
\mR=\ma_1+\ma_2+\cdots+\ma_r,
\]
если каждый элемент $c$ из $\mR$ можно одним и только одним способом представить в виде $c=a_1+a_2+\cdots+a_r$, где каждый $a_i$ является элементом из $\ma_i$.

\emph{Если нулевой идеал является наименьшим общим кратным попарно взаимно простых идеалов $\mb_1,\ldots,\mb_r$, а $\ma_i$ является комплементом $\mb_i$, то $\mR$ является прямой суммой идеалов $\ma_1,\ldots,\ma_r$.} Действительно, из-за $\mo=(\ma_1,\ldots,\ma_r)$ каждый элемент из $\mR$ имеет по меньшей мере одно аддитивное представление через $\ma_i$; из-за $[\ma_i,\mb_i]=(0)$ это представление единственно, поскольку из $0=a_1+a_2+\cdots+a_r=a_i+b_i$ следует $a_i=0$. По второй теореме об изоморфизме идеалы $\ma_i$ изоморфны кольцам классов вычетов $\mR\mid\mb_i$. Имеют место "ортогональные соотношения" $\ma_i\ma_k=(0)$ при $i\ne k$ и $\ma_i^2=\ma_i$; последнее -- из-за $\ma_i=\mo\ma_i$.

\emph{Наоборот, если существует представление $\mR$ как прямой суммы,
\[
\mR=\ma_1+\ma_2+\cdots+\ma_r,
\]
и если положить
\[
\mb_i=(\ma_1,\ldots,\ma_{i-1},\ma_{i+1},\ldots,\ma_r)
=\ma_1+\cdots+\ma_{i-1}+\ma_{i+1}+\cdots+\ma_r,
\]
то $\mb_i$ попарно взаимно просты, а их наименьшее общее кратное равно нулевому идеалу.} Действительно, из $\mo=(\ma_i,\mb_i)$ следует $\mo=(\mb_k,\mb_i)$ при $k\ne i$, так как $\mb_k$ является делителем $\ma_i$. Пусть далее $c$ делится на все $\mb_i$. Тогда
\[
 c=a_2^{(1)}+\cdots+a_r^{(1)}
 =a_1^{(2)}+a_3^{(2)}+\cdots+a_r^{(2)}
 =\cdots=a_1^{(r)}+\cdots+a_{r-1}^{(r)},
\]
где каждый раз $a_i^{(\lambda)}\equiv0\pmod{\ma_i}$. Из единственности представления, поскольку каждый раз одна компонента равна нулю, следует $0=a_1^{(\lambda)},\ldots,0=a_r^{(\lambda)}$ для каждого $\lambda$, а значит $c=0$.\footnote{Теоремы, приведенные в 5., являются специальными случаями общей теоремы о связи между прямой суммой и прямым пересечением. Ср. H. Prüfer, Theorie der Abelschen Gruppen I, Math. Zeitschr. 20 (1924), с. 165--187, §6. Приведенные там рассуждения остаются верными и при такой общей постановке понятия группы, что модули и идеалы подчиняются ей как специальные случаи.}
% END UNIT BODY
% END PRODUCER UNIT 118
\endgroup


\clearpage
\begingroup
\setcounter{footnote}{0}
% BEGIN PRODUCER UNIT 119
% Source: C:/Users/Floris/Documents/interlanguage/03_projects/noether/02_slavic_working_corpus/translations/paper30/section05/russian/v001/Noether_Paper30_Section05_Russian_v001.tex
% Identity: 16760 B / SHA-256 C31E74CAD25E433B034695785931DE76826D8616879EB2682753D087B030CE62
\setlength{\parindent}{1.25em}
\setlength{\parskip}{0pt}
\makeatletter
\renewcommand\footnotesize{\@setfontsize\footnotesize{7.7}{8.7}\selectfont}
\makeatother
\emergencystretch=45em
\allowdisplaybreaks
\sloppy
\hbadness=10000
\vbadness=10000
\hfuzz=2pt
\vfuzz=10pt
\raggedbottom

\providecommand{\mR}{\mathfrak{R}}
\providecommand{\mT}{\mathfrak{T}}
\providecommand{\mA}{\mathfrak{A}}
\providecommand{\mB}{\mathfrak{B}}
\providecommand{\mC}{\mathfrak{C}}
\providecommand{\mD}{\mathfrak{D}}
\providecommand{\ma}{\mathfrak{a}}
\providecommand{\mb}{\mathfrak{b}}
\providecommand{\mc}{\mathfrak{c}}
\providecommand{\md}{\mathfrak{d}}
\providecommand{\mf}{\mathfrak{f}}
\providecommand{\mm}{\mathfrak{m}}
\providecommand{\mo}{\mathfrak{o}}
\providecommand{\mq}{\mathfrak{q}}
\providecommand{\mt}{\mathfrak{t}}
\providecommand{\mpideal}{\mathfrak{p}}
% BEGIN UNIT BODY
% Unit: Paper30 Section05 v001. Direct Russian translation from the German final-audited source slice, source lines 598--628 / segments P30-S0090--P30-S0102. Footnote counter continues after Paper30 Section04.
\setcounter{footnote}{19}

\subsection*{\S~5. Простые идеалы и первичные идеалы}

Пусть снова $\mR$ будет коммутативным кольцом, для которого не предполагается никакая дальнейшая аксиома -- даже существование единичного элемента. Соберем основные факты о простых идеалах и первичных идеалах.\footnote{В \emph{Idealtheorie}, \S~4, предполагается аксиома I, условие цепей делителей; поэтому там не выявлено резко, какие свойства простых и первичных идеалов не зависят от этой аксиомы.}

1. Идеал $\mpideal$ из $\mR$ называется \emph{слабым простым идеалом}, если кольцо классов вычетов $\mR\mid\mpideal$ есть кольцо без делителей нуля; то есть если из $\ma\not\equiv0\pmod{\mpideal}$ и $\mb\not\equiv0\pmod{\mpideal}$ всегда следует $\ma\mb\not\equiv0\pmod{\mpideal}$. Идеал $\mpideal^*$ из $\mR$ называется \emph{сильным простым идеалом}, если кольцо классов вычетов $\mR\mid\mpideal^*$ есть кольцо без идеалов-делителей нуля; то есть если из $\ma\not\equiv0\pmod{\mpideal^*}$ и $\mb\not\equiv0\pmod{\mpideal^*}$ всегда следует $\ma\mb\not\equiv0\pmod{\mpideal^*}$.

\emph{Каждый сильный простой идеал одновременно является слабым простым идеалом}, как показывает специализация $\ma,\mb$ до главных идеалов. Обратно, \emph{каждый слабый простой идеал также является сильным простым идеалом}; поэтому можно говорить просто о простых идеалах. В самом деле, пусть $\mpideal$ -- слабый простой идеал; пусть $\ma\not\equiv0\pmod{\mpideal}$ и $\mb\not\equiv0\pmod{\mpideal}$; выберем элементы $a$ и $b$ из $\ma$ и $\mb$ так, что $a\not\equiv0\pmod{\mpideal}$ и $b\not\equiv0\pmod{\mpideal}$. Тогда $ab\not\equiv0\pmod{\mpideal}$, а потому $\ma\mb\not\equiv0\pmod{\mpideal}$; значит, $\mpideal$ также является сильным простым идеалом.

2. Идеал $\mq$ из $\mR$ называется \emph{слабым первичным идеалом}, если в кольце классов вычетов $\mR\mid\mq$ некоторая степень каждого делителя нуля обращается в нуль; то есть если из $a\not\equiv0\pmod{\mq}$ и $b^x\not\equiv0\pmod{\mq}$ для каждого $x$ всегда следует $ab\not\equiv0\pmod{\mq}$. Система $\mpideal$ всех элементов из $\mR$, которые становятся делителями нуля в $\mR\mid\mq$, образует идеал, причем простой идеал; он является делителем $\mq$ и называется ассоциированным простым идеалом. В самом деле, вместе с $a$ делителем нуля становится и $ra$, вместе с $a$ и $b$ также $a-b$; последнее потому, что вместе с $a^x$ и $b^\lambda$ элемент $(a-b)^{x+\lambda}$ всегда становится делимым на $\mq$. Если, кроме того, $a$ не является делителем нуля, то по определению никакая степень $a$ также не является делителем нуля; из $a\not\equiv0\pmod{\mpideal}$ и $b\not\equiv0\pmod{\mpideal}$ поэтому следует $a^x b^\lambda=(ab)^x\not\equiv0\pmod{\mq}$, а значит, $ab\not\equiv0\pmod{\mpideal}$.

Идеал $\mq^*$ из $\mR$ называется \emph{сильным первичным идеалом}, если в кольце классов вычетов $\mR\mid\mq^*$ некоторая степень каждого идеала-делителя нуля обращается в нуль; то есть если из $\ma\not\equiv0\pmod{\mq^*}$ и $\mb^x\not\equiv0\pmod{\mq^*}$ для каждого $x$ всегда следует $\ma\mb\not\equiv0\pmod{\mq^*}$. Как и для слабых первичных идеалов, показывают: наибольший общий делитель всех идеалов из $\mR$, которые становятся идеалами-делителями нуля в $\mR\mid\mq^*$, образует простой идеал, являющийся делителем $\mq^*$, то есть ассоциированный простой идеал. Обратно, обо всех первичных идеалах, имеющих один и тот же ассоциированный простой идеал $\mpideal$, говорят, что они "принадлежат $\mpideal$".

\emph{Каждый сильный первичный идеал одновременно является слабым первичным идеалом}, как показывает специализация $\ma,\mb$ до главных идеалов. В общем случае, однако, обратное неверно.\footnote{Это показывает следующий пример, сообщенный мне R. Hölzer. Пусть $\mR$ -- полиномиальная область от счетно многих неопределенных $x_i$ с коэффициентами из поля, и пусть
\[
\mq=(x_1^2,x_2^3,\ldots,x_\nu^{\nu+1},\ldots,x_i x_k\,[i\ne k]\ldots).
\]
Делителями нуля в кольце классов вычетов являются все и только те полиномы, которые делятся на $\mpideal=(x_1,x_2,\ldots,x_\nu,\ldots)$, и каждый раз некоторая степень такого делителя нуля делится на $\mq$; следовательно, $\mq$ является слабым первичным идеалом с ассоциированным простым идеалом $\mpideal$. Напротив, $\mq$ не является сильным первичным идеалом: если положить $\ma=(x_1,x_3,x_5,\ldots,x_{2\nu+1},\ldots)$ и $\mb=(x_2,x_4,\ldots,x_{2\nu},\ldots)$, то $\ma\mb\equiv0\pmod{\mq}$, но никакая степень $\ma$ или $\mb$ не делится на $\mq$.} \emph{Если же в $\mR$ предполагается аксиома I условия цепей делителей, то каждый слабый первичный идеал одновременно становится сильным первичным идеалом;} поэтому можно говорить просто о первичных идеалах. Пусть $\mq$ -- слабый первичный идеал; пусть $\ma\not\equiv0\pmod{\mq}$ и $\mb^x\not\equiv0\pmod{\mq}$ для каждого $x$. Тогда существуют элементы $a$ из $\ma$ и $b$ из $\mb$ такие, что $a\not\equiv0\pmod{\mq}$ и $b^x\not\equiv0\pmod{\mq}$ для каждого $x$; следовательно, $ab\not\equiv0\pmod{\mq}$, а потому $\ma\mb\not\equiv0\pmod{\mq}$. Ведь если для каждого $b$ из $\mb$ существовал бы показатель $\lambda$, для которого $b^\lambda\equiv0\pmod{\mq}$, то $\mb^{\lambda_1+\cdots+\lambda_r}\equiv0\pmod{\mq}$, где $\lambda_1,\ldots,\lambda_r$ означают показатели конечного числа базисных элементов.\footnote{Чтобы из аксиомы I вывести существование конечного базиса идеала, надо положить в основу фиксированный хороший порядок в $\mR$ (ср. \S~6).} Последнее рассуждение показывает еще и следующее: существует (наименьший) показатель $\varrho$, для которого $\mpideal^\varrho\equiv0\pmod{\mq}$, где $\mpideal$ означает ассоциированный простой идеал; $\varrho$ называется \emph{экспонентой} $\mq$.

3. Далее условие цепей делителей снова \emph{не} предполагается. Наименьшее общее кратное $[\ma_1,\ldots,\ma_r]$ конечного числа идеалов называется \emph{кратчайшим представлением}, если никакой $\ma_i$ не содержится в наименьшем общем кратном остальных, то есть если никакой $\ma_i$ нельзя опустить.

\emph{Наименьшее общее кратное конечного числа слабых, соответственно сильных, первичных идеалов, принадлежащих одному и тому же простому идеалу $\mpideal$, снова является слабым первичным идеалом с $\mpideal$ как ассоциированным простым идеалом. Кратчайшее представление конечным числом слабых, соответственно сильных, первичных идеалов, принадлежащих различным простым идеалам, не является слабым, соответственно сильным, первичным идеалом.} Пусть $\mf=[\mq_1,\ldots,\mq_r]$, а $\mpideal$ -- ассоциированный простой идеал слабых первичных идеалов $\mq_i$; тогда $\mpideal$ состоит также из совокупности элементов, некоторая степень которых делится на $\mf$. Поэтому из $a\not\equiv0\pmod{\mf}$ и $b^x\not\equiv0\pmod{\mf}$ для каждого $x$ следует, что $b\not\equiv0\pmod{\mpideal}$ и $a\not\equiv0\pmod{\mq_i}$ по крайней мере для одного индекса $i$; значит, $ab\not\equiv0\pmod{\mq_i}$, а потому $ab\not\equiv0\pmod{\mf}$. Если всюду заменить элементы идеалами, то уже нельзя вывести $\mb\not\equiv0\pmod{\mpideal}$.

Пусть теперь $\mm=[\mq_1,\ldots,\mq_r]$ с $\mpideal_i\ne\mpideal_k$ при $i\ne k$ есть кратчайшее представление слабыми первичными идеалами, и пусть $r\ge2$. Тогда существует по крайней мере один ассоциированный простой идеал, скажем $\mpideal_1$, который не содержится ни в одном из остальных; действительно, всякая собственная цепь кратных среди $\mpideal$ должна оборваться не позднее чем через $r$ шагов. Поэтому существуют элементы $a_i$ из $\mpideal_i$ такие, что $a_i^{\varrho_i}\equiv0\pmod{\mq_i}$ и $a_2\not\equiv0\pmod{\mpideal_1},\ldots,a_r\not\equiv0\pmod{\mpideal_1}$. Если далее $q_1\not\equiv0\pmod{\mm}$ есть элемент из $\mq_1$, то $q_1 a_2^{\varrho_2}\cdots a_r^{\varrho_r}\equiv0\pmod{\mm}$, но никакая степень $a_2^{\varrho_2}\cdots a_r^{\varrho_r}$ не делится на $\mm$, поскольку иначе следовала бы делимость на $\mpideal_1$; значит, $\mm$ не является слабым первичным идеалом. Если всюду заменить элементы идеалами, то есть $q_1$ заменить на $\mq_1$, а $a_i$ -- на идеалы $\ma_i$ с
\[
\ma_i\not\equiv0\pmod{\mpideal_1},\qquad
\ma_i^{\varrho_i}\equiv0\pmod{\mq_i},
\]
получается доказательство для сильных первичных идеалов.\footnote{Этой модификацией моего первоначального доказательства, благодаря которой условие цепей делителей становится ненужным, я обязана B. L. van der Waerden.}

4. \emph{Частное модулей и частное идеалов.} Пусть $\mT$ -- кольцо-расширение $\mR$, и пусть $\mA,\mB,\mC,\ldots$ -- $\mR$-модули в $\mT$. Частное $\mC=\mA:\mB$ определяется как модуль, удовлетворяющий условиям $\mC\mB\equiv0\pmod{\mA}$, и из $\mD\mB\equiv0\pmod{\mA}$ следует $\mD\equiv0\pmod{\mC}$; таким образом, $\mC$ есть "наименьший" модуль, для которого $\mC\mB\equiv0\pmod{\mA}$. Тем самым частное однозначно определяется как наибольший общий делитель всех модулей $\mD$, для которых $\mD\mB\equiv0\pmod{\mA}$; такие $\mD$ существуют, поскольку нулевой модуль заведомо является одним из них.

Из определения следует: если $\mB$ есть кратное $\overline{\mB}$, то $\mA:\overline{\mB}$ есть кратное $\mA:\mB$. Если $\mA$ есть кратное $\overline{\mA}$, то $\mA:\mB$ есть кратное $\overline{\mA}:\mB$. Имеем
\[
        \mA:\mB\mC=(\mA:\mB):\mC=(\mA:\mC):\mB .
\]
Если кольцо-расширение $\mT$ совпадает с $\mR$, то частное модулей переходит в частное идеалов $\ma:\mb$; значит, и для него верны приведенные выше правила вычисления. Из $\ma\mb\equiv0\pmod{\ma}$ здесь еще следует $\ma\equiv0\pmod{\ma:\mb}$.

5. \emph{Если $\ma:\mb=\ma$, то $\mb$ называется простым относительно $\ma$.} Итак, идеал $\mb$ прост относительно $\ma$ тогда и только тогда, когда из $\md\mb\equiv0\pmod{\ma}$ всегда следует $\md\equiv0\pmod{\ma}$. Из правил вычисления в пункте 4 следует: если $\mb$ и $\mc$ просты относительно $\ma$, то их произведение и их наименьшее общее кратное также просты относительно $\ma$. Если $\mb$ прост относительно $\ma$ и $\ma$ прост относительно $\mb$, то $\ma$ и $\mb$ называются взаимно простыми в этом смысле. Из \S~4, 4$\beta$, следует: если $\ma$ и $\mb$ взаимно просты как делители, то они также взаимно просты в этом смысле.
% END UNIT BODY
% END PRODUCER UNIT 119
\endgroup


\clearpage
\begingroup
\setcounter{footnote}{0}
% BEGIN PRODUCER UNIT 120
% Source: C:/Users/Floris/Documents/interlanguage/03_projects/noether/02_slavic_working_corpus/translations/paper30/section06/russian/v001/Noether_Paper30_Section06_Russian_v001.tex
% Identity: 11658 B / SHA-256 45D613F518C2F9A40D9276EFDFA957E8AFFF9E28A40BF6BA624B701DAEC4E2F5
\setlength{\parindent}{1.25em}
\setlength{\parskip}{0pt}
\makeatletter
\renewcommand\footnotesize{\@setfontsize\footnotesize{7.7}{8.7}\selectfont}
\makeatother
\emergencystretch=55em
\allowdisplaybreaks
\sloppy
\hbadness=10000
\vbadness=10000
\hfuzz=2pt
\vfuzz=10pt
\raggedbottom

\providecommand{\mR}{\mathfrak{R}}
\providecommand{\mA}{\mathfrak{A}}
\providecommand{\mB}{\mathfrak{B}}
\providecommand{\mC}{\mathfrak{C}}
\providecommand{\ma}{\mathfrak{a}}
\providecommand{\mb}{\mathfrak{b}}
\providecommand{\mc}{\mathfrak{c}}
\providecommand{\mm}{\mathfrak{m}}
\providecommand{\mt}{\mathfrak{t}}
\providecommand{\mpideal}{\mathfrak{p}}
% BEGIN UNIT BODY
% Unit: Paper30 Section06 v001. Direct Russian translation from the German final-audited source slice, source lines 631--680 / segments P30-S0103--P30-S0113. Footnote counter continues after Paper30 Section05.
\setcounter{footnote}{25}

\subsection*{\S~6. Теория идеалов при условии цепей делителей}

В основе пусть лежит (коммутативное) кольцо $\mR$, для которого выполнена аксиома I условия цепей делителей; далее всюду будем считать заданным фиксированный хороший порядок всех элементов из $\mR$. Тем самым задан и хороший порядок всех идеалов из $\mR$. Действительно, эти две предпосылки сразу дают, что каждый идеал из $\mR$ имеет базис идеала, состоящий из конечного числа элементов. Поэтому если от хорошего порядка элементов из $\mR$ перейти к квазилексикографическому порядку всех конечных подмножеств $\mR$,\footnote{Это следует понимать так: элементы $a_1,\ldots,a_n$ конечного подмножества $\mathfrak W$ обозначены так, что порядок индексов совпадает с хорошим порядком, заданным в $\mR$. Каждое подмножество с $n$ элементами предшествует такому подмножеству с $m>n$ элементами. Если два подмножества равномощны, то первое считается более ранним, когда первый элемент, в котором они различаются, предшествует соответствующему элементу другого подмножества в хорошем порядке. Тем самым каждой системе конечных подмножеств сопоставлен первый элемент. При заданном хорошем порядке на $\mR$ не требуется никакого дополнительного постулата выбора; в частности, если $\mR$ счетно, как в случае алгебраического числового поля, то вообще никакой постулат выбора не лежит в основе.} и сопоставить каждому идеалу как выделенный базис первый среди различных возможных базисов в хорошем порядке конечных подмножеств, то вместе с конечными подмножествами хорошо упорядочены и идеалы.

\textbf{Теорема I.} \emph{При условии цепей делителей каждый идеал из $\mR$ допускает представление как наименьшее общее кратное конечного числа неразложимых идеалов, т.е. идеалов, которые нельзя представить как наименьшее общее кратное двух собственных делителей.}

Для доказательства надо показать: если теорема I не выполнена для идеала $\mm$, то $\mm$ имеет собственный делитель, для которого теорема I также не выполнена; отсюда, вопреки предполагаемому условию цепей делителей, можно построить цепь делителей, не обрывающуюся за конечное число шагов. В самом деле, $\mm$ должен быть разложимым, иначе $\mm=[\mm]$ уже давало бы нужное представление. Если $\mm=[\ma,\mb]$, то теорема I не может одновременно выполняться для $\ma$ и $\mb$, поскольку тогда соответствующее представление следовало бы для $\mm$. Значит, существуют собственные делители $\mm$, для которых теорема I не выполнена; пусть $\ma_1$ будет первым из них в хорошем порядке идеалов. Строя тем же образом для $\ma_1$ собственный делитель $\ma_2$, и вообще для $\ma_i$ собственный делитель $\ma_{i+1}$, получаем вполне определенную цепь делителей, не обрывающуюся за конечное число шагов, вопреки предположению.\footnote{Теорема I, очевидно, точно так же верна для областей модулей, если условие цепей делителей предполагается для системы всех модулей. Она имеет чисто теоретико-множественный характер; вместе с тем это единственная теорема, в доказательстве которой используется хороший порядок идеалов.}

Из-за предполагаемого условия цепей делителей по \S~5, 2 можно говорить просто о первичных идеалах. Связь между первичностью и неразложимостью дает следующая теорема:

\textbf{Теорема II.} \emph{При условии цепей делителей каждый неразложимый идеал является первичным; иначе говоря, каждый непервичный идеал разложим.}

Так как при переходе к кольцу классов вычетов $\mR/\mm$ условие цепей делителей сохраняется благодаря гомоморфизму; так как представлению $\mm$ как наименьшего общего кратного соответствует представление нулевого идеала; и так как по первой теореме об изоморфизме (\S~4, 3) первичным делителям $\mm$ снова соответствуют такие же в $\mR/\mm$ и обратно, можно ограничиться разложением нулевого идеала в кольце классов вычетов. Поскольку это снова кольцо той же общности, можно с самого начала считать разлагаемый идеал нулевым идеалом в $\mR$.

Пусть, таким образом, нулевой идеал из $\mR$ не первичен, так что существует по меньшей мере одна пара идеалов $\ma,\mb$, причем $\mb$ еще можно считать главным идеалом, такая, что
\[
\ma\ne(0),\qquad \mb^x\ne(0)\ (x=1,2,\ldots),
\qquad \ma\mb=(0).
\]
Образуем цепь делителей частных идеалов, которая по предположению обрывается за конечное число шагов:
\[
\ma,\ \ma:\mb,\ldots,\ma:\mb^\nu,\ldots;
\]
пусть, например, $\mt=\ma:\mb^{m-1}=\ma:\mb^m=\cdots$ равно всем последующим идеалам этой цепи; значит, $\mt=\mt:\mb$, или $\mb$ прост относительно $\mt$. Далее $\mt$ как делитель $\ma$ отличен от нулевого идеала, и точно так же $\mb^x$ отлично от нуля для каждого показателя $x$. Поэтому для доказательства теоремы II достаточно установить представление
\[
(0)=[\mt,\mb^{m+1}].
\]
По определению
\[
\mb^{m-1}\mt\equiv0\pmod{\ma},\qquad \mb^m\mt=(0),
\]
следовательно, остается лишь показать:
\[
\mc=[\mt,\mb^{m+1}]\equiv0\pmod{\mb^m\mt}.
\]
Это следует из предположений, что $\mb$ является главным идеалом и прост относительно $\mt$. Каждый элемент $c$ из $\mc$, вследствие делимости на $\mb^{m+1}$, допускает представление, где под $n$ понимается символ целого числа,
\[
 c=k b^{m+1}+n b^{m+1}=r b^m,
\]
где $r$ теперь снова означает элемент из $\mR$. Из $c\equiv0\pmod{\mt}$ поэтому следует $r b^m\equiv0\pmod{\mt}$, а значит $r\equiv0\pmod{\mt}$ ввиду $\mt:\mb=\mt$. Тем самым $c\equiv0\pmod{\mt\mb^m}$, и теорема II доказана.

\textbf{Теорема III.} \emph{При условии цепей делителей каждый идеал допускает кратчайшее представление как наименьшее общее кратное конечного числа наибольших первичных компонент, принадлежащих различным простым идеалам.}

Чтобы получить такое представление, надо только, где это возможно, заменить существующее по теореме I представление через неразложимые, а значит по теореме II первичные идеалы, кратчайшим представлением. Если объединить первичные идеалы, принадлежащие одинаковым простым идеалам, то по \S~5, 3 получается искомое представление. Первичные компоненты можно называть наибольшими, поскольку наименьшее общее кратное любых из них уже не является первичным.\footnote{При этом соответствующие простые идеалы и изолированные компоненты определены однозначно. Ср. \emph{Idealtheorie} или W. Krull, Math. Ann. 90 (1923), S. 55--64.}
% END UNIT BODY
% END PRODUCER UNIT 120
\endgroup


\clearpage
\begingroup
\setcounter{footnote}{0}
% BEGIN PRODUCER UNIT 121
% Source: C:/Users/Floris/Documents/interlanguage/03_projects/noether/02_slavic_working_corpus/translations/paper30/section07/russian/v001/Noether_Paper30_Section07_Russian_v001.tex
% Identity: 10779 B / SHA-256 B827DE77153C163CD19F3D7CD25814511661721C80A4E9335EF1A6A93434A27F
\setlength{\parindent}{1.25em}
\setlength{\parskip}{0pt}
\makeatletter
\renewcommand\footnotesize{\@setfontsize\footnotesize{7.7}{8.7}\selectfont}
\makeatother
\emergencystretch=55em
\allowdisplaybreaks
\sloppy
\hbadness=10000
\vbadness=10000
\hfuzz=2pt
\vfuzz=10pt
\raggedbottom

\providecommand{\mR}{\mathfrak{R}}
\providecommand{\ma}{\mathfrak{a}}
\providecommand{\mb}{\mathfrak{b}}
\providecommand{\mm}{\mathfrak{m}}
\providecommand{\mo}{\mathfrak{o}}
\providecommand{\mq}{\mathfrak{q}}
\providecommand{\mpideal}{\mathfrak{p}}
\providecommand{\mbar}[1]{\overline{#1}}
% BEGIN UNIT BODY
% Unit: Paper30 Section07 v001. Direct Russian translation from the German final-audited source slice, source lines 682--730 / segments P30-S0114--P30-S0129. Footnote counter continues after Paper30 Section06.
\setcounter{footnote}{28}

\subsection*{\S~7. Теория идеалов при двойном условии цепей}

Упрощения по сравнению с теорией, развитой до сих пор, основаны на вспомогательных утверждениях, которые будут даны в пунктах 1 и 2.

1. \emph{Если в коммутативном кольце без делителей нуля выполнено условие цепей кратных, то это кольцо одновременно является полем.} Надо показать, что при $a\ne0$ уравнение $ax=b$ всегда имеет решение в кольце. То, что оно не может иметь более одного решения, следует из предположения, что кольцо не имеет делителей нуля.

Пусть $\ma$ -- главный идеал, порожденный $a$. По предположению ряд степеней обрывается за конечное число шагов; пусть, например, $\ma^m$ равно всем последующим степеням. Если $\mb$ обозначает главный идеал, порожденный $b$, то получаем
\[
\ma^m\mb=\ma^{m+1}\mb,
\qquad \ma^{m+1}\mb=(a^{m+1}b).
\]
Отсюда, в частности, для элемента $a^m b$ получается представление, где под $n$ понимается символ целого числа,
\[
 a^m b=k a^{m+1}b+n a^{m+1}b=a^{m+1}c,
\]
где $c$ снова является элементом из $\mR$. Поскольку по предположению делителей нуля нет, отсюда следует $b=ac$, чем доказано свойство поля.

2. Из пункта 1 непосредственно следует вспомогательное утверждение: \emph{если в коммутативном кольце выполнено условие цепей кратных, то простой идеал не имеет собственного делителя, отличного от $\mo$.}

Так же получается дополнение: \emph{если выполнена только аксиома II, то есть условие цепей кратных по модулю каждого идеала, отличного от нулевого идеала, то каждый простой идеал, отличный от нулевого идеала, не имеет собственного делителя, отличного от $\mo$.}

Действительно, кольцо классов вычетов по простому идеалу $\mpideal$ по определению становится кольцом без делителей нуля; а поскольку благодаря гомоморфизму здесь также выполнено условие цепей кратных, то по пункту 1 оно является полем. Поэтому из взаимно однозначного соответствия между делителями $\mpideal$ и идеалами в кольце классов вычетов следует, что $\mpideal$ не имеет собственного делителя, отличного от $\mo$.\footnote{При этом в $\mR$, а значит и в $\mo$, единичном идеале, состоящем из всех элементов $\mR$, не обязан существовать единичный элемент, хотя по пункту 1 в кольце классов вычетов есть единичный класс.}

\textbf{Теорема IV.} \emph{Если в коммутативном кольце выполнено двойное условие цепей, то в каждом кратчайшем представлении идеала однозначно определены те первичные компоненты, которые не принадлежат единичному идеалу $\mo$.}

\emph{Дополнение.} При предположении аксиом I и II теорема IV верна для каждого идеала, отличного от нулевого идеала.

Пусть
\[
\mm=[\mq,\mq_1,\ldots,\mq_r]=[\mbar\mq,\mbar\mq_1,\ldots,\mbar\mq_s]
\]
суть кратчайшие представления $\mm$ через наибольшие первичные компоненты, и пусть $\mo,\mpideal_1,\ldots,\mpideal_r$ соответственно $\mo,\mbar\mpideal_1,\ldots,\mbar\mpideal_s$ -- соответствующие простые идеалы. При этом $\mq$ соответственно $\mbar\mq$ следует опустить, если не появляется первичный идеал, принадлежащий $\mo$. Из
\[
\mo\mpideal_1\cdots\mpideal_r\equiv0\pmod{\mm}
\]
следует, что каждый $\mpideal_i$ должен содержаться в некотором $\mbar\mpideal_j$, а значит по вспомогательному утверждению быть с ним тождественным. Соответственно совпадают и простые идеалы, возникающие в обратном направлении. Далее простота остальных множителей дает взаимную делимость соответствующих первичных компонент. Наконец, если $\mo$-компонента действительно появляется, то это должно происходить в обоих кратчайших представлениях. Тем самым доказана однозначность; доказательство одновременно показывает, что всякое представление без первичной компоненты, принадлежащей $\mo$, уже является кратчайшим.

3. Если к предположению двойного условия цепей добавить еще существование единичного элемента, то теорема IV усиливается до теоремы V.

\textbf{Теорема V.} \emph{Если в коммутативном кольце с единичным элементом выполнено двойное условие цепей, то каждый идеал можно однозначно представить как произведение конечного числа попарно взаимно простых первичных идеалов.}

\emph{Дополнение.} При предположении аксиом I, II, III теорема V верна для каждого идеала, отличного от нулевого идеала.

Из-за существования единичного элемента имеем $\mo^2=\mo$, и тем самым каждый первичный идеал, принадлежащий $\mo$, равен $\mo$, а значит не может появляться в кратчайшем представлении идеала, отличного от $\mo$. Следовательно, по теореме IV первичные компоненты однозначно определены. Одновременно всякое представление после опускания $\mo$ становится кратчайшим. Два различных простых идеала по вспомогательному утверждению из пункта 2 всегда взаимно просты; значит, то же верно для первичных компонент, и наименьшее общее кратное становится произведением.

Из этих предположений далее следует: \emph{в коммутативном кольце с единичным элементом и двойным условием цепей, кроме единичного идеала, простыми относительно идеала $\mm$ являются все и только те простые идеалы, которые не содержатся в $\mm$. Понятия простоты относительно и взаимной простоты совпадают.}

Если $\mpideal$ не содержится в $\mm$, то он отличен от всех простых идеалов, принадлежащих $\mm$, и потому прост относительно каждой первичной компоненты и относительно $\mm$; одновременно он взаимно прост с каждой первичной компонентой, а значит и с $\mm$. Если же, напротив, $\mpideal\ne\mo$ содержится в $\mm$, то он должен совпадать с одним из соответствующих простых идеалов. Если он принадлежит компоненте $\mq$, то $\mq:\mpideal$ становится собственным делителем $\mq$; благодаря однозначной разложимости тогда и произведение, делимое на $\mm:\mpideal$, становится собственным делителем $\mm$. Следовательно, $\mpideal$ не прост относительно $\mm$.
% END UNIT BODY
% END PRODUCER UNIT 121
\endgroup


\clearpage
\begingroup
\setcounter{footnote}{0}
% BEGIN PRODUCER UNIT 122
% Source: C:/Users/Floris/Documents/interlanguage/03_projects/noether/02_slavic_working_corpus/translations/paper30/section08/russian/v001/Noether_Paper30_Section08_Russian_v001.tex
% Identity: 11124 B / SHA-256 3C9868BEBE4AD33E538F43C4C0F75443EEDB1D539E21FB026CAD7157CC53E7D1
\setlength{\parindent}{1.25em}
\setlength{\parskip}{0pt}
\makeatletter
\renewcommand\footnotesize{\@setfontsize\footnotesize{7.7}{8.7}\selectfont}
\makeatother
\emergencystretch=65em
\allowdisplaybreaks
\sloppy
\hbadness=10000
\vbadness=10000
\hfuzz=2pt
\vfuzz=10pt
\raggedbottom

\providecommand{\mR}{\mathfrak{R}}
\providecommand{\ma}{\mathfrak{a}}
\providecommand{\mb}{\mathfrak{b}}
\providecommand{\mc}{\mathfrak{c}}
\providecommand{\mm}{\mathfrak{m}}
\providecommand{\mn}{\mathfrak{n}}
\providecommand{\mo}{\mathfrak{o}}
\providecommand{\mq}{\mathfrak{q}}
\providecommand{\mt}{\mathfrak{t}}
\providecommand{\mpideal}{\mathfrak{p}}
\providecommand{\mbar}[1]{\overline{#1}}
% BEGIN UNIT BODY
% Unit: Paper30 Section08 v001. Direct Russian translation from the German final-audited source slice, source lines 733--858 / segments P30-S0130--P30-S0140. Footnote counter continues after Paper30 Section07.
\setcounter{footnote}{29}

\subsection*{\S~8. Теория идеалов при интегральной замкнутости в поле частных}

Развитая до сих пор теория идеалов теперь должна быть усилена до обычной формы добавлением двух последних аксиом.

1. \emph{Вспомогательное утверждение.} Пусть $\mR$ -- коммутативное кольцо без делителей нуля с единичным элементом, и пусть в $\mR$ предполагается условие цепей делителей (аксиомы I, III, IV). Тогда из $\ma=\ma\mb$ и $\ma\ne(0)$ всегда следует $\mb=\mo$. Следовательно,
\[
        \mc\ne\mc\mt
\]
для всех идеалов, отличных от нулевого и единичного идеалов.
\footnote{Напротив, при тех же предположениях из $\ma\mb=\ma\mc$ нельзя заключать $\mb=\mc$. В качестве примера можно взять кольцо многочленов $K[x,y]$: $\ma=(xy)$, $\mb=(x^3,xy,y^2)$, $\mc=(x^2,y^2)$; действительно, $\ma\mb=\ma\mc=\ma^2$, но $\mb\ne\mc$.}
Доказательство получается так же, как во вспомогательном утверждении §1, 3, поскольку $\ma\mb$ можно рассматривать как конечный модуль относительно $\mb$. Если $\alpha_1,\ldots,\alpha_n$ обозначает базис идеала $\ma$, существующий благодаря I, то из $\ma=\ma\mb$ получается система уравнений
\[
        \alpha_i=b_{i1}\alpha_1+\cdots+b_{in}\alpha_n,
        \qquad b_{ij}\equiv0\pmod{\mb},\quad i=1,\ldots,n.
\]
Так как $\mR$ предполагается без делителей нуля, отсюда следует $|b_{ij}-e_{ij}|=0$, где $e_{ij}=0$ при $i\ne j$, а $e_{ii}=e$. Из уравнения
\[
        e-b_{11}-\cdots\equiv0\pmod{\mb}
\]
следует $e\equiv0\pmod{\mb}$, а значит $\mb=\mo$.

2. Теперь остается только показать, что после добавления предположения интегральной замкнутости в поле частных (аксиома V) к аксиомам I--IV каждый первичный идеал, отличный от нулевого, становится степенью своего соответствующего простого идеала. Единственность представления
\[
        \mm=\mpideal_1^{\varrho_1}\cdots\mpideal_r^{\varrho_r}
\]
для всех идеалов, отличных от нулевого, уже доказана теоремой V и вспомогательным утверждением из пункта 1; одновременно показано, что эти простые идеалы не имеют собственного делителя, отличного от единичного идеала. Доказательство будет проведено сначала для показателя два с использованием следствия Dedekind II (§1, 4), а в общем случае -- полной индукцией.

\emph{Вспомогательное утверждение.} \emph{При предположении аксиом I--V нет других первичных идеалов показателя два, кроме $\mq=\mpideal^2$.}

Нужно показать, что из
\[
        \mpideal^2\equiv0\pmod{\mq},
        \qquad \mq\equiv0\pmod{\mpideal},
        \qquad \mq\not\equiv0\pmod{\mpideal^2}
\]
необходимо следует $\mq=\mpideal^2$. Итак, пусть
\[
        \mc\equiv0\pmod{\mq},
        \qquad \mc\not\equiv0\pmod{\mpideal},
        \qquad \mc\equiv0\pmod{\mpideal^2}.
\]
Тогда из вспомогательного утверждения §7, 3 следует, что $\mc:\mpideal$ становится собственным делителем $\mc$. Значит, существует элемент $b$ такой, что
\[
        b\mpideal\equiv0\pmod{\mc}\equiv0\pmod{\mq},
        \qquad b\not\equiv0\pmod{\mc}.
\]
Следовательно, $y=b/c$ не является целым, то есть $y$ принадлежит полю частных, но не $\mR$. Нужно доказать $b\not\equiv0\pmod{\mpideal}$; отсюда, поскольку $\mq$ первичен и принадлежит $\mpideal$, следует $\mpideal\equiv0\pmod{\mq}$.

По следствию Dedekind II существуют элементы $m,n$ из $\mR$ такие, что
\[
        y=b/c=m/n
\]
и также $m^2/n$ не является целым, то есть
\[
        bn=mc,
        \qquad m\not\equiv0\pmod{\mn},
        \qquad m^2\not\equiv0\pmod{\mn}.
\]
Умножение $b\mpideal$ на $m$ дает
\[
        mb\mpideal\equiv0\pmod{\mn c},
        \qquad\text{следовательно}\qquad m\mpideal\equiv0\pmod{\mn},
\]
последнее потому, что речь идет о кольце без делителей нуля. Отсюда следует
\[
        m\not\equiv0\pmod{\mpideal},
\]
так как $m^2\not\equiv0\pmod{\mn}$ и $\mn\not\equiv0\pmod{\mpideal}$; последнее снова следует из вспомогательного утверждения §7, 3, поскольку $\mn:\mpideal$ является собственным делителем $\mn$. Четыре соотношения
\[
        m\not\equiv0\pmod{\mpideal},
        \quad n\not\equiv0\pmod{\mpideal},
        \quad c\equiv0\pmod{\mpideal},
        \quad c\not\equiv0\pmod{\mpideal^2},
        \quad bn=mc
\]
дают $b\not\equiv0\pmod{\mpideal}$. Действительно, поскольку $\mpideal^2$ первичен, сначала получается $mc\not\equiv0\pmod{\mpideal^2}$, и тогда из $bn=mc$ следует $b\not\equiv0\pmod{\mpideal}$. Тем самым вспомогательное утверждение доказано.

3. \emph{Вспомогательное утверждение.} \emph{При предположении аксиом I--V нет других первичных идеалов показателя $\varrho$, кроме $\mq=\mpideal^\varrho$.}
\footnote{Доказательство вспомогательного утверждения 3 при предположении справедливости вспомогательного утверждения 2 содержится у Masazo Sono, \emph{On Congruences II}, §§9 и 10.}
Это вспомогательное утверждение следует из вспомогательного утверждения из пункта 2 без дальнейшего применения аксиом. Для доказательства сначала нужно показать:
\[
        \text{если }\mc\equiv0\pmod{\mpideal},\ \mc\not\equiv0\pmod{\mpideal^2},
        \quad\text{то}\quad
        \mpideal^\varrho=(\mc^\varrho,\mpideal^{\varrho+1})
\]
для каждого $\varrho$. По вспомогательному утверждению из пункта 2 имеем $\mpideal=(\mc,\mpideal^2)$. Поэтому если предположить
\[
        \mpideal^{\varrho-1}=(\mc^{\varrho-1},\mpideal^\varrho),
\]
то после умножения на $\mpideal$ по дистрибутивному закону получается
\[
        \mpideal^\varrho=(\mc^\varrho,\mpideal^{\varrho+1}).
\]
Доказываемое вспомогательное утверждение можно привести к следующей второй форме:
\[
        \mq\equiv0\pmod{\mpideal^\varrho},
        \quad \mq\not\equiv0\pmod{\mpideal^{\varrho+1}},
        \quad \mpideal^{\varrho+1}\equiv0\pmod{\mq},
        \quad \varrho\ge1
\]
влечет
\[
        \mpideal^\varrho\equiv0\pmod{\mq}.
\]
Действительно, из $\mq\equiv0\pmod{\mpideal^\varrho}$ и $\mpideal^\varrho\ne\mpideal^{\varrho+1}$ по вспомогательному утверждению из пункта 1 выполнено второе условие. Конечное повторное применение этой второй формы дает $\mpideal^\varrho\equiv0\pmod{\mq}$, а значит $\mq=\mpideal^\varrho$.

Из предположения второй формы с учетом представления $\mpideal^\varrho$ для каждого элемента $q$ из $\mq$ получаем:
\[
        q=b c^\varrho+p^{\varrho+1},
        \qquad b\equiv0\pmod{\mpideal},
\]
или, что то же самое,
\[
        bc^\varrho\equiv0\pmod{\mq,\mpideal^{\varrho+1}}.
\]
Так как $(\mq,\mpideal^{\varrho+1})$ первичен, отсюда следует
\[
        c^\varrho\equiv0\pmod{\mq,\mpideal^{\varrho+1}}
        \quad\text{или}\quad
        c^{\varrho+1}\equiv0\pmod{\mq,\mpideal^{\varrho+2}},
\]
а потому $\mpideal^\varrho\equiv0\pmod{\mq}$.

4. Подводя итог, получаем:

\textbf{Теорема VI.} \emph{Если в коммутативном кольце выполнены аксиомы I--V, то каждый идеал, отличный от нулевого и единичного идеалов, можно однозначно представить как произведение степеней конечного числа простых идеалов, отличных от нулевого и единичного идеалов; эти простые идеалы не имеют собственного делителя, отличного от единичного идеала.}
% END UNIT BODY
% END PRODUCER UNIT 122
\endgroup


\clearpage
\begingroup
\setcounter{footnote}{0}
% BEGIN PRODUCER UNIT 123
% Source: C:/Users/Floris/Documents/interlanguage/03_projects/noether/02_slavic_working_corpus/translations/paper30/section09/russian/v001/Noether_Paper30_Section09_Russian_v001.tex
% Identity: 17397 B / SHA-256 AD3BEA6F2DD78F9C9D1E4981B5EDB73552648F54BF95CB69A7412A5216948DF3
\setlength{\parindent}{1.25em}
\setlength{\parskip}{0pt}
\makeatletter
\renewcommand\footnotesize{\@setfontsize\footnotesize{7.7}{8.7}\selectfont}
\makeatother
\emergencystretch=65em
\allowdisplaybreaks
\sloppy
\hbadness=10000
\vbadness=10000
\hfuzz=2pt
\vfuzz=10pt
\raggedbottom

\providecommand{\mR}{\mathfrak{R}}
\providecommand{\mK}{\mathfrak{K}}
\providecommand{\mT}{\mathfrak{T}}
\providecommand{\ma}{\mathfrak{a}}
\providecommand{\mb}{\mathfrak{b}}
\providecommand{\mc}{\mathfrak{c}}
\providecommand{\mm}{\mathfrak{m}}
\providecommand{\mn}{\mathfrak{n}}
\providecommand{\mo}{\mathfrak{o}}
\providecommand{\mq}{\mathfrak{q}}
\providecommand{\mt}{\mathfrak{t}}
\providecommand{\mpideal}{\mathfrak{p}}
\providecommand{\mbar}[1]{\overline{#1}}
% BEGIN UNIT BODY
% Unit: Paper30 Section09 v001. Direct Russian translation from the German final-audited source slice, source lines 860--996 / segments P30-S0141--P30-S0162. Footnote counter continues after Paper30 Section08.
\setcounter{footnote}{31}

\subsection*{\S~9. Аксиомы как следствие предполагаемого разложения}

Чтобы из существования обычного разложения идеалов, которое по теореме VI следует из аксиом I--V, также в обратном направлении вывести эти аксиомы, разложение нужно предположить в следующей форме.

\emph{Предположение.} В коммутативном кольце $\mR$ каждый идеал, не порожденный нулевым или единичным элементом, однозначно представим как произведение степеней простых идеалов. Эти простые идеалы являются элементарными идеалами, не порожденными единичным элементом, то есть они не имеют собственного делителя, отличного от $\mo$; обратно, все элементарные идеалы являются простыми идеалами. Единственность должна выполняться в точной форме:
\[
        \ma^\varrho\ne\ma^{\varrho+1}
\]
для каждого идеала, не порожденного нулевым или единичным элементом.

Существование единичного элемента при этом не предполагается. Если единичного элемента нет, то условие «не порожденный единичным элементом» не является ограничением.

1. \emph{Доказательство аксиомы III.} Пусть аксиома III не выполнена. Нужно показать, что $\mo^2$ становится элементарным идеалом, не будучи простым идеалом. По предположению точной единственности имеем $\mo\ne\mo^2$; значит, $\mo\equiv0\pmod{\mo^2}$, но $\mo^2\not\equiv0\pmod{\mo^2}$, а следовательно, $\mo^2$ не является простым идеалом. Но $\mo^2$ элементарен: $\mo^2$ не может делиться ни на какой простой идеал, отличный от $\mo$, и потому из-за предполагаемого произведенного представления не имеет делителей, отличных от $\mo$ и $\mo^2$.

2. \emph{Доказательство аксиомы IV.} Если $a$ есть делитель нуля, то есть $ab=0$, но $a\ne0$ и $b\ne0$, то умножение представлений главных идеалов, порожденных $a$ и $b$, дает
\[
        (0)=\mpideal_1^{\varrho_1}\cdots\mpideal_r^{\varrho_r},
\]
где $\mpideal_i$ отличны от нулевого и единичного идеалов. Представление считается кратчайшим в том смысле, что удаление любого фактора дает собственный делитель нуля. Если в этом представлении входит только один простой идеал, то
\[
        (0)=\mpideal^\varrho=\mpideal^{\varrho+1},
\]
вопреки точному требованию единственности. Если же входит больше одного простого идеала, то получается
\[
\begin{aligned}
\mpideal_1^{\varrho_1}
 &= (\mpideal_1^{\varrho_1},\mpideal_2^{\varrho_2}\cdots\mpideal_r^{\varrho_r})\mpideal_1^{\varrho_1}  \\
 &= \mpideal_1^{\varrho_1}\mpideal_2^{\varrho_2}\cdots\mpideal_r^{\varrho_r},
\end{aligned}
\]
поскольку из существования единичного элемента $\mpideal_1$ взаимно прост со всеми остальными простыми идеалами. Следовательно, и здесь получается противоречие точному требованию единственности.

3. \emph{Доказательство аксиом I и II.} Предположения дают для каждого идеала, отличного от нулевого: из делимости следует произведенное представление. Если
\[
        \mm=\mpideal_1^{\varrho_1}\cdots\mpideal_r^{\varrho_r}
\]
и $\ma$ является делителем $\mm$, то в произведенном представлении $\ma$ встречаются только такие простые идеалы, которые совпадают с $\mpideal_i$, а показатели лежат между $0$ и $\varrho_i$. Поэтому имеется лишь конечное число делителей $\mm$, соответствующих комбинациям $0\le\sigma_i\le\varrho_i$. Значит, в кольце классов вычетов по $\mm$ выполняется двойное условие цепей. Тем самым аксиомы I и II доказаны; условие цепей делителей выполняется и в самой $\mR$, поскольку каждый собственный делитель нулевого идеала отличен от нулевого идеала.

4. \emph{Кольцо классов вычетов по каждому идеалу, отличному от нулевого идеала, является кольцом главных идеалов.} Идеал можно считать отличным от единичного идеала. Если идеал первичен, $\mq=\mpideal^\varrho$, и если $\mc\equiv0\pmod{\mpideal}$, но $\mc\not\equiv0\pmod{\mpideal^2}$, то из пункта 3 следует, что $\mc=\mpideal\ma$ с $(\ma,\mpideal)=\mo$. Следовательно,
\[
        \mpideal^\sigma=(\mc^\sigma,\mpideal^{\sigma+1})
        \qquad(\sigma<\varrho),
\]
и потому в кольце классов вычетов по первичному идеалу каждый идеал является главным. В общем случае, при
\[
        \mm=\mq_1\cdots\mq_r
\]
и соответствующем представлении кольца классов вычетов как прямой суммы
\[
        \mR/\mm=\ma_1+\cdots+\ma_r,
\]
$\ma_i$ изоморфен $\mR/\mq_i$, а значит, каждый идеал из $\ma_i$ является главным. Если $\mc$ -- произвольный идеал кольца классов вычетов, то $\ma_i\mc$ является главным идеалом $\ma_i c_i$; получается
\[
        \mc=\ma_1c_1+\cdots+\ma_r c_r=\mo c,
        \qquad c=c_1+\cdots+c_r.
\]
Действительно, вследствие $\ma_i\ma_k=(0)$ и вследствие $c_i\equiv0\pmod{\ma_i}$ главные идеалы $\ma_i\mo c_i$ и $\ma_i\mc$ из $\ma_i$ совпадают.

Эта теорема допускает еще следующую формулировку: если $\mc$ -- произвольный идеал, то его можно превратить в главный идеал умножением на идеал, взаимно простой с заданным идеалом $\mb$. В самом деле, положив $\mm=\mb\mc$, получаем, что $\mc$ по модулю $\mm$ становится главным идеалом, то есть
\[
        \mc=(\mo c,\mm)=(\mc\ma,\mc\mb)=\mc(\ma,\mb);
\]
значит, $\mo c=\mc\ma$ и $(\ma,\mb)=\mo$.

5. \emph{Теория дробных идеалов.} Дробным идеалом называется каждый конечный $\mR$-модуль в поле частных $\mK$.

\emph{Отличные от нулевого модуля конечные $\mR$-модули из $\mK$ образуют относительно умножения абелеву группу.}

Произведение двух конечных $\mR$-модулей снова является конечным $\mR$-модулем; умножение ассоциативно и коммутативно; поскольку $\mK=\mo\mK$, единичный идеал является единичным элементом системы. Следовательно, остается только показать, что уравнение
\[
        \mathfrak U X=\mo
\]
всегда имеет решение в системе конечных $\mR$-модулей.

Предварительное замечание. Если уравнение $\mathfrak A T=\mathfrak B$ имеет одно и только одно решение, то $T$ равно частному модулей $\mathfrak B:\mathfrak A$ (§5, 4). В самом деле,
\[
        T\equiv0\pmod{\mathfrak B:\mathfrak A},
        \qquad
        \mathfrak B=\mathfrak A T\equiv0\pmod{\mathfrak A(\mathfrak B:\mathfrak A)}\equiv0\pmod{\mathfrak B}.
\]
Пусть сначала $\mathfrak A$ -- главный модуль, $\mathfrak A=(\alpha)$; тогда $X=(e/\alpha)$ является решением $\mathfrak A X=\mo$, и $X$ снова конечный $\mR$-модуль. Отсюда следует: каждый конечный $\mR$-модуль является частным модулей двух идеалов из $\mR$, чем оправдывается название дробного идеала. Действительно, пусть $t_1,\ldots,t_n$ -- базис модуля $T$, пусть $t_i=\tau_i/a$, и пусть $\mt$ -- идеал в $\mR$, порожденный $\tau_i$. Тогда $\mo a\,T=\mt$, откуда по предварительному замечанию следует $T=(\mt:\mo a)$.

Решение $\mathfrak A X=\mo$ теперь можно указать в общем виде. Пусть $\mathfrak A=\ma:\mo c$, и пусть $\mb$ выбрано так, что $\ma\mb$ равно главному идеалу $\mo a$. Тогда
\[
        \mathfrak A\,\mo c=\ma,
        \qquad \mathfrak A\,\mb\mo c=\mo a,
        \qquad \mathfrak A(\mb\mo c:\mo a)=\mo.
\]
При этом $X$, как частное идеала на главный идеал, снова является конечным $\mR$-модулем. Из группового свойства еще следует, что частное модулей любых двух идеалов $\ma:\mb$, как решение уравнения $\mb X=\ma$, является конечным $\mR$-модулем.

6. Из группового свойства следуют дальнейшие выводы.

6$\alpha$. Можно сокращать; из $\mathfrak T=\mathfrak A\mathfrak M:\mathfrak B\mathfrak M$ следует $\mathfrak T=\mathfrak A:\mathfrak B$, и обратно.

6$\beta$. Каждый дробный идеал допускает представление $\mathfrak C=\ma:\mb$, где $\ma$ и $\mb$ взаимно просты; этим условием $\ma$ и $\mb$ определены однозначно.

6$\gamma$. Каждый главный модуль $\mo n$, порожденный нецелым элементом, допускает представление как частное главных идеалов так, что никакая степень числителя не делится на знаменатель. Пусть в сокращенном представлении $\mo n=\mb:\mc$, где $(\mb,\mc)=\mo$; пусть $\mo b=\mb\ma$ и $(\ma,\mc)=\mo$. Тогда
\[
        \mo n=\ma\mb:\ma\mc=\mo b:\ma\mc,
\]
и поскольку $\ma\mc$ также становится главным идеалом, получаем
\[
        \mo n=\mo b:\mo c.
\]
Но никакая степень $b$ не делится на $c$.

7. \emph{Доказательство аксиомы V.} Из 6$\gamma$ следует, что каждый нецелый элемент из $\mK$ допускает представление частным
\[
        \eta=a/c
\]
так, что в $\mR$ никакая степень числителя не делится на знаменатель. Следовательно, такой элемент не может удовлетворять никакому уравнению, которое характеризует его как целый относительно $\mR$; ведь из
\[
        \eta^n+r_1\eta^{n-1}+\cdots+r_n=0
\]
следует $a^n\equiv0\pmod{\mo c}$ в $\mR$.

8. \emph{Следствие.} Если в коммутативном кольце $\mR$ выполнены аксиомы I--IV и каждый первичный идеал неразложим, то выполнена и аксиома V. В силу аксиом I--III каждая степень простого идеала $\mpideal^\varrho$ одновременно является первичным идеалом; в частности, это верно для $\mpideal^2$, и по предположению он неразложим. Отсюда нужно доказать
\[
        \mpideal=(\mo c,\mpideal^2);
\]
ибо тогда по вспомогательному утверждению §8, 3 следует, что каждый первичный идеал становится степенью простого идеала, а это вместе с остальными предположениями по пункту 7 влечет интегральную замкнутость.

В силу условия цепей делителей получаем
\[
        \mpideal=(\mo c_1,\ldots,\mo c_n),
\]
причем в качестве множителей для $c_i$ рассматриваются только классы вычетов по $\mpideal$, и $c_i$ можно считать линейно независимыми над полем классов вычетов по $\mpideal$. Но из этой линейной независимости при $n>1$ следует представление
\[
        \mpideal^2=[\mq_1,\mq_2],
        \qquad \mq_1=(\mo c_1,\mpideal^2),
        \qquad \mq_2=(\mo c_2,\ldots,\mo c_n,\mpideal^2),
\]
так что $\mq_1$ и $\mq_2$ становятся собственными делителями $\mpideal^2$. Это противоречит неразложимости. Следовательно, $\mpideal=(\mo c,\mpideal^2)$ доказано.

9. Если в кольце допускаются делители нуля, то из аксиом I--III и интегральной замкнутости в кольце частных не следует, что каждый первичный идеал неразложим.
\footnote{Ср. примечание, приведенное в оригинале.}
Пусть $\mT$ -- кольцо, в котором выполнены все аксиомы I--IV, но не интегральная замкнутость. По следствию 8 в $\mT$ существуют разложимые первичные идеалы $\mq$; пусть $\mpideal^\varrho$ означает собственное кратное $\mq$, а $\mR$ -- кольцо классов вычетов $\mT/\mpideal^\varrho$. В $\mR$ выполнены аксиомы I--III, и одновременно выполнена интегральная замкнутость в кольце частных. Действительно, каждый элемент, не делящийся на $\mpideal$, становится взаимно простым с $\mpideal^\varrho$, так что в $\mR$ все регулярные элементы являются единицами. Следовательно, $\mR$ совпадает со своим кольцом частных.
% END UNIT BODY
% END PRODUCER UNIT 123
\endgroup


\clearpage
\begingroup
\setcounter{footnote}{0}
% BEGIN PRODUCER UNIT 124
% Source: C:/Users/Floris/Documents/interlanguage/03_projects/noether/02_slavic_working_corpus/translations/paper30/section10/russian/v001/Noether_Paper30_Section10_Russian_v001.tex
% Identity: 7705 B / SHA-256 C2A8AEADC3B76EEF42A6F5297F86FCF5D6A811C4068E2E9F9A8BCF25F2F86C60
\setlength{\parindent}{1.25em}
\setlength{\parskip}{0pt}
\makeatletter
\renewcommand\footnotesize{\@setfontsize\footnotesize{7.7}{8.7}\selectfont}
\makeatother
\emergencystretch=65em
\allowdisplaybreaks
\sloppy
\hbadness=10000
\vbadness=10000
\hfuzz=2pt
\vfuzz=10pt
\raggedbottom

\providecommand{\mR}{\mathfrak{R}}
\providecommand{\mG}{\mathfrak{G}}
\providecommand{\mU}{\mathfrak{U}}
\providecommand{\mE}{\mathfrak{E}}
\providecommand{\mA}{\mathfrak{A}}
\providecommand{\mB}{\mathfrak{B}}
\providecommand{\mM}{\mathfrak{M}}
\providecommand{\mH}{\mathfrak{H}}
% BEGIN UNIT BODY
% Unit: Paper30 Section10 v001. Direct Russian translation from the German final-audited source slice, source lines 998--1043 / segments P30-S0163--P30-S0171. Footnote counter continues after Paper30 Section09.
\setcounter{footnote}{32}

\subsection*{\S~10. Двойное условие цепей и композиционный ряд}

Далее будет показано, что для произвольных модульных областей, предполагаемых хорошо упорядоченными, предположение о выполнении двойного условия цепей, то есть о том, что каждая цепь делителей и каждая цепь кратных модулей обрывается за конечное число шагов, равносильно предположению о существовании композиционного ряда.
\footnote{Модули области являются абелевыми группами относительно сложения, а именно обобщенными абелевыми группами, поскольку область множителей состоит из элементов из $\mR$. Утверждения этого параграфа остаются верными, если под «модулями» понимать совокупность нормальных делителей некоммутативной группы.}
Специализация модульной области до коммутативного кольца дает соответствующие факты для системы всех идеалов кольца; при этом в пункте 2 всюду нужно заменить изоморфизм модулей изоморфизмом колец.

Модульная область $\mG$ называется простой, если в $\mG$ нет других $\mR$-модулей, кроме самой $\mG$ и нулевого модуля $\mE$. Модуль $\mA$ называется простым в $\mG$, если $\mG/\mA$ прост. Цепь кратных
\[
        \mG,\mU_1,\ldots,\mU_r,\mE
\]
называется композиционным рядом $\mG$ длины $r$, если все модули цепи различны и каждый модуль прост в предыдущем.

1. Если в $\mG$ предполагается выполнение двойного условия цепей, то в $\mG$ существует композиционный ряд. Из предполагаемого хорошего порядка и условия цепей делителей, как в §6, получается существование хотя бы одного модуля, простого в $\mG$. Действительно, пусть $\mG_1$ -- первый в хорошем порядке собственный делитель $\mG$, а вообще $\mG_i$ -- первый собственный делитель $\mG_{i-1}$; тогда вполне определенная цепь
\[
        \mG,\mG_1,\ldots,\mG_i,\ldots
\]
обрывается за конечное число шагов и необходимо приводит к простому модулю. Повторяя этот шаг, получаем цепь кратных
\[
        \mG,\mU_1,\ldots,\mU_i,
\]
которая обрывается по условию цепей кратных и дает композиционный ряд.

2. Если в $\mG$ предполагается существование композиционного ряда, то в $\mG$ выполняется двойное условие цепей. Доказательство идет полной индукцией; в ходе индукционного перехода одновременно получается теорема Jordan--Hölder.
\footnote{Ср. Masazo Sono, \emph{On Congruences I}, §§11--13, для случая колец; далее Dedekind, \emph{Über die von drei Moduln erzeugte Dualgruppe}, Math. Ann. 53 (1900), S. 371--403.}

Пусть сначала $\mG$ проста; тогда обычные утверждения выполняются непосредственно: через каждый простой модуль проходит композиционный ряд; все композиционные ряды имеют одинаковую длину и изоморфные фактор-группы с точностью до порядка; через каждый собственный модуль проходит композиционный ряд; условия цепей кратных и делителей выполняются.

Теперь эти утверждения доказываются для композиционного ряда
\[
        \mG,\mA,\mA_1,\ldots,\mA_r,\mE
\]
длины $r+1$, если они выполнены для более коротких рядов. Если $\mB$ -- еще один модуль, простой в $\mG$, и $\mB\ne\mA$, то $(\mA,\mB)=\mG$. Положив $\mM=[\mA,\mB]$, по второй теореме об изоморфизме получаем
\[
        \mG/\mB\cong\mA/\mM,
        \qquad
        \mG/\mA\cong\mB/\mM.
\]
Отсюда следует существование композиционного ряда, проходящего через $\mB$; сравнение полученных таким образом рядов обмена дает теорему Jordan--Hölder.

Для произвольного модуля $\mH$, отличного от $\mG$, из фиксированного композиционного ряда образуем ряд
\[
        \mG,(\mU,\mH),(\mU_1,\mH),\ldots,(\mU_r,\mH),\mH.
\]
По второй теореме об изоморфизме различные члены этого ряда последовательно просты, то есть образуют композиционный ряд от $\mG$ до $\mH$. Отсюда следует возможность проводить композиционные ряды через каждый собственный модуль. Наконец, каждая цепь кратных после перехода к такому модулю обрывается по индукции, и точно так же каждая цепь делителей после перехода к модулю классов вычетов. Эквивалентность предположений -- композиционного ряда или двойного условия цепей -- тем самым доказана.

\emph{Поступило 13. 8. 1925.}
% END UNIT BODY
% END PRODUCER UNIT 124
\endgroup


\clearpage
\begingroup
\setcounter{footnote}{0}
% BEGIN PRODUCER UNIT 125
% Source: C:/Users/Floris/Documents/interlanguage/03_projects/noether/02_slavic_working_corpus/translations/paper31/introduction/russian/v001/Noether_Paper31_Introduction_Russian_v001.tex
% Identity: 14951 B / SHA-256 727CEFF3F78F57C96FD958533394BEC65D9ECEFB7399BB8BAAA588207B7A6690
\setlength{\parindent}{1.25em}
\setlength{\parskip}{0pt}
\makeatletter
\renewcommand\footnotesize{\@setfontsize\footnotesize{7.7}{8.7}\selectfont}
\makeatother
\emergencystretch=65em
\allowdisplaybreaks
\sloppy
\hbadness=10000
\vbadness=10000
\hfuzz=2pt
\vfuzz=10pt
\raggedbottom
% BEGIN UNIT BODY
% Unit: Paper31 Introduction v001. Direct Russian translation from the German final-audited source slice, source lines 1--24 / segments P31-S0001--P31-S0011.

\section*{31. Теорема о дискриминанте для порядков алгебраического числового или функционального поля}
\begin{center}
Journal f. d. reine u. angew. Math. 157 (1927), с. 82--104
\end{center}

Теорема Dedekind о дискриминанте, как известно, утверждает: простое число входит в дискриминант поля тогда и только тогда, когда главный идеал, порожденный этим числом, делится по крайней мере на квадрат простого идеала.

Ниже я показываю, что эта теорема сохраняется для всех конечных порядков конечного алгебраического числового поля, то есть для всех таких колец целых чисел, которые содержат кольцо всех целых рациональных чисел, в следующей форме: простое число $p$ входит в дискриминант порядка тогда и только тогда, когда в действующем для этого порядка разложении главного идеала $(p)$ на наибольшие первичные компоненты\footnote{Ср. E. Noether, \emph{Abstrakter Aufbau der Idealtheorie in algebraischen Zahl- und Funktionenkörpern}, Math. Ann. 96 (1926), с. 26--61. Во всех определениях я ссылаюсь на эту работу, ниже цитируемую как "Idealtheorie". В дальнейшем под "порядком" всегда понимается "конечный порядок", а под "кольцом" -- коммутативное кольцо.} появляется по крайней мере один собственный первичный идеал, то есть такой, который не является простым идеалом.

Поскольку для главного порядка, системы всех целых чисел поля, нет других первичных идеалов, кроме степеней простых идеалов, теорема Dedekind получается специализацией.

Доказательство теоремы посредством перехода к кольцу классов вычетов порядка по модулю $(p)$ сводится к теории разложения тех колец, которые имеют конечный ранг относительно поля, содержащегося в самом кольце; значит, они образуют коммутативную систему гиперкомплексных величин с главной единицей и коэффициентами из того же поля. Здесь специально в качестве поля коэффициентов выступает поле классов вычетов целых рациональных чисел по модулю $p$, а разложению $(p)$ в кольце классов вычетов соответствует разложение нулевого идеала.

Итак, если такое кольцо конечного ранга назвать вполне приводимым, когда его нулевой идеал можно представить как наименьшее общее кратное конечного числа простых идеалов, то речь идет о критериях полной приводимости. Первый такой критерий таков: если область коэффициентов является совершенным полем, то кольцо вполне приводимо тогда и только тогда, когда его кольцо расширения с алгебраически замкнутой областью коэффициентов вполне приводимо. Для этого кольца расширения, однако, ненулевость дискриминанта очень просто распознается как необходимый и достаточный критерий полной приводимости;\footnote{В одном специальном случае Hilbert рассуждает сходно: Zur Theorie der aus $n$ Haupteinheiten gebildeten komplexen Größen, Gött. Nachr. 1896.} а так как дискриминант порядка по модулю каждого простого числа переходит в дискриминант кольца классов вычетов, тем самым теорема о дискриминанте доказана.\footnote{Родственные исследования, также в некоммутативном случае, но ограниченные главным порядком (максимальной областью) и предполагающие целые рациональные числовые коэффициенты, содержатся в только что опубликованной работе A. Speiser: \emph{Allgemeine Zahlentheorie}. Naturf. Gesellschaft Zürich 71 (1926), с. 8--48. Ср. также указанную там американскую литературу.}

Если вместо числовых полей рассматривать функциональные поля, то порядок теперь должен содержать основную область всех целых рациональных функций, соответственно, для алгебраических функций от нескольких неопределенных или для целых алгебраических функций, основную область всех целых функционалов с целорациональными коэффициентами. Тогда в кольце классов вычетов может возникать и случай несовершенной области коэффициентов, например, когда в случае целочисленной функциональной области рассматривается кольцо классов вычетов по простому числу. Для несовершенного поля как области коэффициентов нужно различать простые идеалы первого и второго рода, в зависимости от того, является ли поле классов вычетов по простому идеалу полем расширения первого или второго рода,\footnote{В известной формулировке Steinitz, J. f. M. 137. Ср. примечание к \S 2, 4 этой работы.} и соответственно различать полную приводимость первого и второго рода. Все указанные выше критерии относятся к полной приводимости первого рода, которая одна только встречается в совершенном случае; следовательно, в самом общем случае теорема о дискриминанте утверждает, что простой элемент основной области входит в дискриминант порядка тогда и только тогда, когда в разложении идеала $(p)$ появляется по крайней мере одна собственно первичная компонента или простой идеал второго рода.

Для главного порядка этот факт впервые был замечен на примерах Kronecker после того, как в Festschrift еще ошибочно была дана формулировка, верная для главного порядка числового поля, хотя и без полного доказательства. Для главного порядка Ostrowski\footnote{A. Ostrowski: Zur arithmetischen Theorie der algebraischen Größen. Nachr. Gött. Ges. d. Wissensch. 1919.} дал доказательство приведенной выше теоремы о дискриминанте и связанные с ней дальнейшие теоремы разветвления; там же содержится подробный обзор литературы. Наконец, как прямое следствие получается теорема о дискриминанте для всех порядков относительных полей, а именно путем перехода к кольцу частных относительно каждого простого идеала основной области. Идею этого перехода, известную в специальном случае, принципиально подчеркнул W. Krull;\footnote{Ср. его данцигский доклад: Jhrbr. d. D. Math. Ver. 34, с. 121.} систематическую теорию кольца частных в рамках более общих исследований дает H. Grell.\footnote{\emph{Beziehungen zwischen den Idealen verschiedener Ringe} (должна появиться в Math. Annalen).}

Данное здесь единое рассмотрение всех порядков, а также совершенной и несовершенной области коэффициентов, опирается, помимо знания теории идеалов, на то, что с самого начала за основу взята теория колец конечного ранга, а значит, общее понятие конечного расширения, а не понятие расширения, порожденного примитивным элементом. Поэтому модульная база порядка вместе с отношениями между ее элементами заменяет степени примитивного элемента поля и определяющее уравнение,\footnote{Эта точка зрения близко соприкасается с представлением Dedekind о высших комплексных числах: Zur Theorie der aus $n$ Haupteinheiten gebildeten komplexen Größen. Gött. Nachr. 1885.} соответственно степени фундаментальной формы и фундаментальное уравнение.

Такая модульная база по модулю каждого простого элемента основной области переходит в базу кольца классов вычетов, тогда как даже для главного порядка это не обязательно верно для базы, состоящей из степеней примитивного элемента; точно так же это может не быть верно и для базы из степеней фундаментальной формы, как только речь идет о целых алгебраических функциях нескольких неопределенных.\footnote{Ср. Ostrowski, a.a.O. "Irregularität erster Art".} Иными словами: кольцо классов вычетов полиномиальной области одной неопределенной с коэффициентами из основной области по произвольному определяющему уравнению или даже по фундаментальному уравнению может по модулю $p$ не быть изоморфным кольцу классов вычетов главного порядка по $(p)$, и главные трудности обычных изложений заключены именно в этом прекращении изоморфизма. Поскольку эта вспомогательная полиномиальная область здесь не возникает, эти трудности тем самым автоматически обходятся; одновременно аддитивную теорию разложения кольца классов вычетов можно понимать как эквивалент мультипликативного разложения уравнения, которому и здесь, при разложении на взаимно простые множители, соответствует аддитивное разложение в полиномиальном кольце классов вычетов.

Отметим, что теорема о дискриминанте может быть только первой теоремой общей теории разветвления порядков, подобно тому как в главном порядке рядом с теоремой о дискриминанте стоит более тонкая теория идеала разветвления, позволяющая там точнее определять показатели. Это определение показателей, основанное на представлении первичных множителей уравнений степенями, можно истолковать как определение длины композиционного ряда от $p$ до $q=p^t$; вероятно, именно в этой форме оно появится в общей теории разветвления.
% END UNIT BODY
% END PRODUCER UNIT 125
\endgroup


\clearpage
\begingroup
\setcounter{footnote}{0}
% BEGIN PRODUCER UNIT 126
% Source: C:/Users/Floris/Documents/interlanguage/03_projects/noether/02_slavic_working_corpus/translations/paper31/section01_entries01_02/russian/v001/Noether_Paper31_Section01_Entries01_02_Russian_v001.tex
% Identity: 14094 B / SHA-256 A7F70FF6B9FAD25E6BCBE57AE534D9C9E2AC3BFEF0E5CE4F9A05061CB8DDF701
\setlength{\parindent}{1.25em}
\setlength{\parskip}{0pt}
\makeatletter
\renewcommand\footnotesize{\@setfontsize\footnotesize{7.7}{8.7}\selectfont}
\makeatother
\emergencystretch=65em
\allowdisplaybreaks
\sloppy
\hbadness=10000
\vbadness=10000
\hfuzz=2pt
\vfuzz=10pt
\raggedbottom
% BEGIN UNIT BODY
% Unit: Paper31 Section01 entries 1--2 v001. Direct Russian translation from the German final-audited source slice, source lines 26--79 / segments P31-S0012--P31-S0024. Footnote counter continues after Paper31 introduction.
\setcounter{footnote}{9}

\subsection*{\S~1. Кольца расширения полей}

\textbf{1.} Сначала приведем несколько замечаний о произвольных кольцах. Если $T$ -- кольцо, а $S$ -- система элементов из $T$, то под "кольцом, выведенным из $S$ в $T$" понимается пересечение всех колец из $T$, которые содержат $S$. Соответственно определяется "идеал, выведенный из $S$ в $T$", или "$P$-модуль, выведенный из $S$ в $T$", если $P$ является подкольцом $T$.\footnote{Ср. об этом Idealtheorie \S 1.}

Если $R$ является подкольцом $T$, а $S$ -- системой элементов из $T$, то кольцо, выведенное из $R$ и $S$ в $T$, обозначается как кольцо $R[S]$, возникающее присоединением $S$ к $R$ как кольцу. Это кольцо состоит из всех полиномов от элементов из $S$ с коэффициентами из $R$, причем отношения равенства и операционные отношения задаются отношениями равенства и операционными отношениями в $T$. В силу этих уже заданных отношений равенства и операций в специальном случае достаточно образовать лишь подсистему всех полиномов, например только все линейные формы от $S$ с коэффициентами из $R$. Такой специальный случай имеет место, например (ср. 2.), когда $R$ -- кольцо с единичным элементом, а $S$ -- поле с тем же единичным элементом; все линейные комбинации $\sum c_i\gamma_i$, где $c_i$ принадлежат $S$, а $\gamma_i$ принадлежат $R$, образуют кольцо $R[S]$.

Но можно также, если задано только кольцо $R$ и, кроме того, система $S$ элементов, не содержащихся в $R$, с отношениями равенства или также взаимными операционными отношениями, сконструировать кольцо расширения $R[S]$, возникающее присоединением $S$ к $R$. Для этого показывают, что всегда существует кольцо, содержащее $R$ и $S$, а именно формально образованное из всех полиномов от $S$ с коэффициентами из $R$ с фиксированными обычными правилами операций. Для определения равенства, помимо правил операций, здесь возможна полная свобода, за исключением ограничения, что должны быть включены определения равенства, действующие в $R$ и $S$. Значит, необходимо отождествлять только те полиномы, которые, в силу равенства в $R$ и $S$ и в силу операций, можно одинаковым образом образовать из равных элементов $R$ и $S$. Сконструированное таким образом кольцо $T$, содержащее $R$ и $S$, тогда совпадает с кольцом, выведенным из $R$ и $S$ в $T$.\footnote{Ср. конструкцию Steinitz поля частных.}

\textbf{2. Кольца расширения полей.} Начиная отсюда предполагаются кольца $\mathfrak R$ с единичным элементом, содержащие поле $P$ как подкольцо, то есть надкольца полей. Единичный элемент $e$ кольца $\mathfrak R$ тогда одновременно является единичным элементом поля $P$; и кольцо становится $P$-модулем.

Пусть $\overline P$ -- надполе поля $P$. Кольцо расширения $\overline{\mathfrak R}=\mathfrak R[\overline P]$, возникающее присоединением $\overline P$ к $\mathfrak R$ как кольцу, определяется следующим образом.

Предполагается, что теоретико-множественное пересечение $[\mathfrak R,\overline P]$ колец $\mathfrak R$ и $\overline P$ равно $P$, и далее, что между элементами $\mathfrak R$ и $\overline P$, не принадлежащими $P$, нет операционных отношений. Этого предположения всегда можно добиться, заменив $\overline P$ эквивалентным полем расширения поля $P$, или также заменив $\mathfrak R$ эквивалентным кольцом расширения.\footnote{Два поля расширения поля $P$ по Steinitz называются "эквивалентными", если между ними можно установить изоморфизм так, что каждый элемент $P$ соответствует самому себе; соответственно следует определять эквивалентные кольца расширения поля $P$. Введением новых символов всегда можно породить эквивалентные расширения; эта свобода в использовании эквивалентных расширений существенна для дальнейшего.}

Элементами $\overline{\mathfrak R}$ объявляются все формально образованные линейные соединения
\[
   \bar c_{i_1}\gamma_{i_1}+\cdots+\bar c_{i_s}\gamma_{i_s},
\]
где $\bar c_i$ пробегают все элементы $\overline P$, а $\gamma_i$ -- все элементы $\mathfrak R$.\footnote{В общем случае элементы из $\mathfrak R$ или $\overline{\mathfrak R}$ обозначаются греческими буквами, а элементы из $P$ или $\overline P$ -- латинскими.} Определение должно быть однозначным: то есть если $\gamma$ заменяются равными им в $\mathfrak R$, а $\bar c$ -- равными им в $\overline P$, то и элемент в $\overline{\mathfrak R}$ объявляется равным.

Сумма определяется однозначно, то есть так, что замена слагаемого в подлежащем фиксации определении равенства равным ему дает тот же результат:
\[
   \sum \bar c_i\gamma_i+\sum \bar d_i\gamma_i
   =\sum(\bar c_i+\bar d_i)\gamma_i .
\]
Среди $\bar c$ и $\bar d$ могут встречаться нули, благодаря чему формально появляются те же $\gamma_i$.

Произведение определяется однозначно в том же смысле:
\[
   \sum \bar c_i\gamma_i\cdot \sum \bar d_k\gamma_k
   =\sum \bar c_i\bar d_k\gamma_i\gamma_k .
\]
Для определения равенства устанавливается: если $\gamma_{i_1},\ldots,\gamma_{i_s}$ линейно независимы относительно $P$, то два элемента
\[
   \bar c_{i_1}\gamma_{i_1}+\cdots+\bar c_{i_s}\gamma_{i_s}
   \quad\text{и}\quad
   \bar d_{i_1}\gamma_{i_1}+\cdots+\bar d_{i_s}\gamma_{i_s}
\]
равны тогда и только тогда, когда $\bar c_{i_k}$ равен $\bar d_{i_k}$ в $\overline P$. Иными словами: элементы из $\mathfrak R$, линейно независимые относительно $P$, объявляются линейно независимыми относительно $\overline P$.

В силу определений суммы и произведения тем самым задано определение равенства для всех элементов $\overline{\mathfrak R}$; действительно, отсюда следует их выразимость через линейно независимые элементы из $\mathfrak R$. Ведь из отношения
\[
   \mu=c_1\gamma_1+\cdots+c_s\gamma_s\quad\text{в }\mathfrak R
\]
умножением на $\bar b=\bar b e$ по определению произведения получаем
\[
   \bar b\mu=\bar b e\cdot(c_1\gamma_1+\cdots+c_s\gamma_s)
      =\bar b c_1\gamma_1+\cdots+\bar b c_s\gamma_s,
\]
и потому, по определению суммы, выразимость всех элементов через линейно независимые элементы из $\mathfrak R$.

Следовательно, когда все элементы выражены через линейно независимые элементы из $\mathfrak R$, определение равенства становится равенством коэффициентов в $\overline P$. Поэтому для элементов, одновременно принадлежащих $\mathfrak R$ и $\overline{\mathfrak R}$ или $\overline P$ и $\overline{\mathfrak R}$, оно совпадает с действующим там равенством; тогда как для остальных элементов из $\overline{\mathfrak R}$ равенство, действующее в $\mathfrak R$ и $\overline P$, не утверждает ничего сверх однозначности, требуемой при определении этих элементов, ввиду предположения о пересечении и дальнейшего условия.\footnote{Без этого предположения существовал бы по крайней мере один элемент $\bar c$ из $\overline P$, не принадлежащий $P$, для которого $\bar c=\bar c e=\gamma$. Тогда, по определению равенства в $\mathfrak R$, элементы $e$ и $\gamma$ были бы линейно независимы относительно $P$, но зависимы относительно $\overline P$. Из-за отношений между элементами $\overline P$ и $\mathfrak R$ приведенное выше определение равенства тогда было бы невозможно; можно вспомнить, например, понижение степени конечного поля расширения, когда основное поле заменяется промежуточным полем. Напротив, приведенное выше расширение характеризуется возможностью появления новых делителей нуля; так, например, кольцо классов вычетов полиномиальной области с рациональными коэффициентами $P$ по модулю $x^2+1$ является кольцом без делителей нуля, изоморфным полю $P(\sqrt{-1})$, тогда как при переходе к полиномиальной области с коэффициентами из $P(\sqrt{-1})$ появляются делители нуля, соответственно тому, что $(x+i)(x-i)=x^2+1$, но ни один множитель не делится на $x^2+1$. Соответствующее утверждение верно, например, для каждого конечного числового поля; это именно та точка зрения, которую Dedekind впервые высказал для высших комплексных чисел. Ср. об этом \S 6, 4, конец.} Точно так же и требование однозначности суммы и произведения не влечет никаких дальнейших отношений, поскольку равные элементы можно считать равными покоэффициентно, и потому формально они дают те же выражения сумм и произведений.

Теперь сразу проверяется, что все аксиомы кольцевых операций выполнены; следовательно, кольцо расширения $\mathfrak R[\overline P]$, порожденное присоединением кольца, всегда существует.
% END UNIT BODY
% END PRODUCER UNIT 126
\endgroup


\clearpage
\begingroup
\setcounter{footnote}{0}
% BEGIN PRODUCER UNIT 127
% Source: C:/Users/Floris/Documents/interlanguage/03_projects/noether/02_slavic_working_corpus/translations/paper31/section01_entry03/russian/v001/Noether_Paper31_Section01_Entry03_Russian_v001.tex
% Identity: 5371 B / SHA-256 843AD6EC7CA356970BC29F9CDB8AD12EF61708A2C69CB05191C4EB7BB79028EA
\setlength{\parindent}{1.25em}
\setlength{\parskip}{0pt}
\makeatletter
\renewcommand\footnotesize{\@setfontsize\footnotesize{7.7}{8.7}\selectfont}
\makeatother
\emergencystretch=65em
\allowdisplaybreaks
\sloppy
\hbadness=10000
\vbadness=10000
\hfuzz=2pt
\vfuzz=10pt
\raggedbottom
% BEGIN UNIT BODY
% Unit: Paper31 Section01 entry 3 v001. Direct Russian translation from the German final-audited source slice, source lines 81--90 / segments P31-S0025--P31-S0027. Footnote counter continues after Paper31 Section01 entries 1--2.
\setcounter{footnote}{14}

\textbf{3. Идеалы в $\mathfrak R$ и $\overline{\mathfrak R}$.} Если $\mathfrak a$ -- идеал в $\mathfrak R$, то модульное произведение $\overline{\mathfrak R}\mathfrak a$ становится идеалом $\overline{\mathfrak a}$ в $\overline{\mathfrak R}$, идеалом расширения идеала $\mathfrak a$. Одновременно $\overline{\mathfrak a}$ является идеалом, выведенным из $\mathfrak a$ в $\overline{\mathfrak R}$; он состоит из всех линейных соединений $\sum \bar c_i a_i$, где сумма каждый раз берется по конечному числу элементов, причем $a_i$ пробегают все элементы из $\mathfrak a$, а $\bar c_i$ -- все элементы из $\overline P$.

Если $\overline{\mathfrak b}$ -- идеал в $\overline{\mathfrak R}$, то пересечение $[\overline{\mathfrak b},\mathfrak R]$ становится идеалом $\mathfrak b$ в $\mathfrak R$, идеалом сужения идеала $\overline{\mathfrak b}$.\footnote{В обозначении H. Grell (в работе, цитированной в примечании 6, с. 83).} Каждый идеал $\mathfrak a$ из $\mathfrak R$ является идеалом сужения, а именно своего идеала расширения $\overline{\mathfrak a}$. Действительно, каждый элемент $a$ из $[\mathfrak R,\overline{\mathfrak a}]$ имеет вид $\sum \bar c_i a_i$, где $a_i$ можно считать линейно независимыми относительно $P$ или $\overline P$. Значит, элементы $a,a_i$ становятся линейно зависимыми в $\overline{\mathfrak R}$, а потому, по определению равенства, также в $\mathfrak R$; следовательно, $a$ принадлежит $\mathfrak a$. Так как верно и обратное включение, получаем $\mathfrak a=[\mathfrak R,\overline{\mathfrak a}]$.

Если $\overline{\mathfrak p}$ -- простой идеал в $\overline{\mathfrak R}$, то $\mathfrak p=[\mathfrak R,\overline{\mathfrak p}]$ является простым идеалом в $\mathfrak R$. Ведь по второй теореме об изоморфизме\footnote{Idealtheorie, \S 4, 3. Там теорема сформулирована только для идеалов одного и того же кольца; однако то же рассуждение показывает ее справедливость в приведенной выше форме для произвольных идеалов $\overline{\mathfrak a}$ из $\overline{\mathfrak R}$. Действительно, кольцевой гомоморфизм $\overline{\mathfrak R}\sim\overline{\mathfrak R}/\overline{\mathfrak a}$ индуцирует для подкольца $\mathfrak R$ кольца $\overline{\mathfrak R}$ кольцевой гомоморфизм $\mathfrak R\sim(\mathfrak R,\overline{\mathfrak a})/\overline{\mathfrak a}$. При этом гомоморфизме нулевому элементу соответствуют все и только элементы из $[\mathfrak R,\overline{\mathfrak a}]$, и потому получается изоморфизм колец: $(\mathfrak R,\overline{\mathfrak a})/\overline{\mathfrak a}\simeq \mathfrak R/[\mathfrak R,\overline{\mathfrak a}]$. Классы справа и слева можно представлять теми же элементами из $\mathfrak R$, поскольку на каждом шаге доказательства элементам сопоставляются представляемые ими классы.} кольцо классов вычетов $\mathfrak R/\mathfrak p$ изоморфно подкольцу кольца $\overline{\mathfrak R}/\overline{\mathfrak p}$, а значит, является кольцом без делителей нуля. Именно, имеет место изоморфизм колец
\[
   (\mathfrak R,\overline{\mathfrak p})/\overline{\mathfrak p}
      \simeq \mathfrak R/[\mathfrak R,\overline{\mathfrak p}],
\]
где под $(\mathfrak R,\overline{\mathfrak p})$ понимается сумма, то есть наибольший общий делитель, модулей $\mathfrak R$ и $\overline{\mathfrak p}$.
% END UNIT BODY
% END PRODUCER UNIT 127
\endgroup


\clearpage
\begingroup
\setcounter{footnote}{0}
% BEGIN PRODUCER UNIT 128
% Source: C:/Users/Floris/Documents/interlanguage/03_projects/noether/02_slavic_working_corpus/translations/paper31/section02_entry01/russian/v001/Noether_Paper31_Section02_Entry01_Russian_v001.tex
% Identity: 5271 B / SHA-256 5B2A115485141212F0C4AA39733AF8BD18B520E1D6D82EA8EB0D8B4546C84D1E
\setlength{\parindent}{1.25em}
\setlength{\parskip}{0pt}
\makeatletter
\renewcommand\footnotesize{\@setfontsize\footnotesize{7.7}{8.7}\selectfont}
\makeatother
\emergencystretch=75em
\allowdisplaybreaks
\sloppy
\hbadness=10000
\vbadness=10000
\hfuzz=2pt
\vfuzz=10pt
\raggedbottom
% BEGIN UNIT BODY
% Unit: Paper31 Section02 entry 1 v001. Direct Russian translation from the German final-audited source slice, source lines 92--100 / segments P31-S0028--P31-S0032. Footnote counter continues after Paper31 Section01 entry 3.
\setcounter{footnote}{16}

\subsection*{\S 2. Кольца конечного ранга. Представление в виде прямой суммы}

\textbf{1. Кольца конечного ранга.} Если $\mathfrak R$ является конечным $P$-модулем, то $\mathfrak R$ имеет конечный ранг, скажем $k$, относительно $P$, то есть в $\mathfrak R$ существует $k$ и не больше линейно независимых относительно $P$ элементов, и всякая система из $k$ линейно независимых относительно $P$ элементов из $\mathfrak R$ тем самым образует модульную базу $\mathfrak R$. Если $P$-модуль $\mathfrak B$ конечного ранга имеет своим делителем другой $\mathfrak A$ того же конечного ранга, $\mathfrak B=O(\mathfrak A)$, то они, следовательно, тождественны.

Отныне будем предполагать, что $\mathfrak R$ имеет конечный ранг $n$ относительно $P$. Тогда каждый $P$-модуль из $\mathfrak R$ имеет конечный ранг; значит, в каждой собственной цепи делителей или кратных $P$-модулей из $\mathfrak R$ ранг каждого модуля должен быть больше, соответственно меньше, чем ранг непосредственно предыдущего; цепи обрываются за конечное число шагов. В частности, для системы всех идеалов из $\mathfrak R$ выполняется "теорема о двойной цепи".

Отсюда следует,\footnote{Idealtheorie, \S 7, 2.} что простой идеал $\mathfrak p$ из $\mathfrak R$ не имеет никакого собственного делителя, отличного от единичного идеала; следовательно, кольцо классов вычетов $\mathfrak R/\mathfrak p$ по простому идеалу, отличному от единичного, становится полем. Если $\mathfrak q$ -- первичный идеал, принадлежащий $\mathfrak p$, то есть $\mathfrak q=O(\mathfrak p)$ и $\mathfrak p^{\rho}=O(\mathfrak q)$, то кольцо классов вычетов $\mathfrak R/\mathfrak q$ становится первичным кольцом, то есть кольцом, в котором некоторая степень (здесь во всяком случае уже $\rho$-я) каждого делителя нуля обращается в нуль, и $\mathfrak R/\mathfrak q$ содержит только делители нуля и единицы. Последнее непосредственно следует из того, что каждый идеал, не делящийся на $\mathfrak p$, становится взаимно простым с $\mathfrak p$, а значит и с $\mathfrak q$.

Из теоремы о двойной цепи и существования единичного элемента далее следует,\footnote{Idealtheorie, \S 7, 3. Теорема V.} что каждый идеал $\mathfrak a$ из $\mathfrak R$ единственным образом представляется как произведение конечного числа попарно взаимно простых первичных идеалов. В силу взаимной простоты это произведение равно наименьшему общему кратному и соответствует представлению кольца классов вычетов $\mathfrak R/\mathfrak a$ в виде прямой суммы первичных колец,\footnote{Idealtheorie, \S 4, 5. То, что добавлено в настоящем сведении, во многих местах рассеяно по литературе.} то есть $\mathfrak R/\mathfrak a$ становится наибольшим общим делителем этих колец, и представление каждого элемента из $\mathfrak R/\mathfrak a$ в виде суммы единственно.
% END UNIT BODY
% END PRODUCER UNIT 128
\endgroup


\clearpage
\begingroup
\setcounter{footnote}{0}
% BEGIN PRODUCER UNIT 129
% Source: C:/Users/Floris/Documents/interlanguage/03_projects/noether/02_slavic_working_corpus/translations/paper31/section02_entry02/russian/v001/Noether_Paper31_Section02_Entry02_Russian_v001.tex
% Identity: 5845 B / SHA-256 65DCB7C20284350887B43D49C4FB45BE780F476881BEDCC55396FD53BB6BB341
\setlength{\parindent}{1.25em}
\setlength{\parskip}{0pt}
\makeatletter
\renewcommand\footnotesize{\@setfontsize\footnotesize{7.7}{8.7}\selectfont}
\makeatother
\emergencystretch=85em
\allowdisplaybreaks
\sloppy
\hbadness=10000
\vbadness=10000
\hfuzz=2pt
\vfuzz=10pt
\raggedbottom
% BEGIN UNIT BODY
% Unit: Paper31 Section02 entry 2 v001. Direct Russian translation from the German final-audited source slice, source lines 102--157 / segments P31-S0033--P31-S0035. Footnote counter continues after Paper31 Section02 entry 1.
\setcounter{footnote}{19}

\textbf{2. Представление колец конечного ранга в виде прямой суммы первичных колец.} Пусть представление нулевого идеала кольца $\mathfrak R$ задано формулой
\[
   (0)=\mathfrak q_1\mathfrak q_2\cdots\mathfrak q_r
      =[\mathfrak q_1,\mathfrak q_2,\ldots,\mathfrak q_r];
\]
и пусть соответствующее представление в виде прямой суммы имеет вид
\[
   \mathfrak R=\mathfrak R_1+\mathfrak R_2+\cdots+\mathfrak R_r,
   \qquad
   \mathfrak R_i\mathfrak R_k=(0)\quad\hbox{для } i\ne k,
   \qquad
   \mathfrak R_i^2=\mathfrak R_i,
   \qquad
   \mathfrak R\mathfrak R_i=\mathfrak R_i.
\]
Тогда $\mathfrak R_i$ становится идеалом из $\mathfrak R$, а именно
\[
   \mathfrak R_i
   =\mathfrak q_1\cdots\mathfrak q_{i-1}\mathfrak q_{i+1}\cdots\mathfrak q_r
   =[\mathfrak q_1,\ldots,\mathfrak q_{i-1},\mathfrak q_{i+1},\ldots,\mathfrak q_r],
\]
\[
   \mathfrak q_i=\mathfrak R_1+\cdots+\mathfrak R_{i-1}+\mathfrak R_{i+1}+\cdots+\mathfrak R_r,
\]
и $\mathfrak R_i$ становится изоморфным кольцу классов вычетов $\mathfrak R/\mathfrak q_i$; при этом каждому элементу $\gamma_i$ из $\mathfrak R_i$ соответствует класс в $\mathfrak R/\mathfrak q_i$, представленный элементом $\gamma_i$.\footnote{Idealtheorie, \S 4, 5.} Если $\mathfrak a$ -- произвольный идеал из $\mathfrak R$, то по дистрибутивному закону получаем
\[
   \mathfrak a=\mathfrak R\mathfrak a
   =\mathfrak R_1\mathfrak a+\cdots+\mathfrak R_r\mathfrak a
   =\mathfrak a_1+\cdots+\mathfrak a_r;
\]
здесь, поскольку $\mathfrak a_i=\mathfrak R_i\mathfrak a=\mathfrak R\mathfrak a_i$, компоненты $\mathfrak a_i$ являются идеалами как в $\mathfrak R_i$, так и в $\mathfrak R$. Как идеалы, $\mathfrak a_i$, подобно самим $\mathfrak R_i$, являются конечными $P$-модулями; из-за прямой суммы их ранги складываются в ранг $\mathfrak a$, соответственно $\mathfrak R$.

Если представление единичного элемента в виде суммы задано формулами
\[
   e=\varepsilon_1+\cdots+\varepsilon_r,
   \qquad
   \varepsilon_i\varepsilon_k=0\quad\hbox{для } i\ne k,
   \qquad
   \varepsilon_i^2=\varepsilon_i,
   \qquad
   \varepsilon_i\equiv e\pmod{\mathfrak q_i},
\]
то представление каждого произвольного элемента $\gamma$ из $\mathfrak R$ получается единственным образом:
\[
   \gamma=\gamma e=\gamma\varepsilon_1+\cdots+\gamma\varepsilon_r
      =\gamma_1+\cdots+\gamma_r,
\]
где, следовательно, компонента $\gamma_i=\gamma\varepsilon_i$ является элементом из $\mathfrak R_i$; одновременно следует $\mathfrak R_i=\mathfrak R\varepsilon_i$. Из "соотношений ортогональности" получается:
\[
   \gamma\delta=(\gamma_1\varepsilon_1+\cdots+\gamma_r\varepsilon_r)
      (\delta_1\varepsilon_1+\cdots+\delta_r\varepsilon_r)
   =\gamma_1\delta_1\varepsilon_1+\cdots+\gamma_r\delta_r\varepsilon_r;
\]
значит, $\gamma_i\delta_i$ становится $i$-й компонентой произведения, и точно так же $\gamma_i+\delta_i$ становится $i$-й компонентой суммы; эти факты следуют также из гомоморфизма из $\mathfrak R$ в $\mathfrak R_i$.

Кольцо $\mathfrak R_i$ содержит подполе $P_i=P\varepsilon_i$, изоморфное $P$, и имеет конечный ранг относительно $P_i$; этот ранг равен рангу, которым $\mathfrak R_i$ обладает как $P$-модуль. Действительно, для каждого элемента $c$ из $P$ имеем $c\varepsilon_i\equiv ce\equiv c\pmod{\mathfrak q_i}$; значит, $c\varepsilon_i\ne 0$, как только $c\ne0$; гомоморфизм между $P$ и $P_i$ становится изоморфизмом. Далее, всякая линейная зависимость относительно $P$ между элементами из $\mathfrak R_i$ посредством умножения на $\varepsilon_i$ переходит в такую же зависимость относительно $P_i$; и обратно, поскольку $P_i=P\varepsilon_i$, всякое соотношение относительно $P_i$ одновременно представляет собой соотношение относительно $P$. Поэтому ранги относительно $P$ и $P_i$ совпадают.
% END UNIT BODY
% END PRODUCER UNIT 129
\endgroup


\clearpage
\begingroup
\setcounter{footnote}{0}
% BEGIN PRODUCER UNIT 130
% Source: C:/Users/Floris/Documents/interlanguage/03_projects/noether/02_slavic_working_corpus/translations/paper31/section02_entry03/russian/v001/Noether_Paper31_Section02_Entry03_Russian_v001.tex
% Identity: 2777 B / SHA-256 2419D851B284FEB83C68B6C8277FEFCF4B84B5F0D1A76FFA6202B4E60890293B
\setlength{\parindent}{1.25em}
\setlength{\parskip}{0pt}
\makeatletter
\renewcommand\footnotesize{\@setfontsize\footnotesize{7.7}{8.7}\selectfont}
\makeatother
\emergencystretch=85em
\allowdisplaybreaks
\sloppy
\hbadness=10000
\vbadness=10000
\hfuzz=2pt
\vfuzz=10pt
\raggedbottom
% BEGIN UNIT BODY
% Unit: Paper31 Section02 entry 3 v001. Direct Russian translation from the German final-audited source slice, source lines 159--172 / segment P31-S0036. Footnote counter continues after Paper31 Section02 entry 2.
\setcounter{footnote}{20}

\textbf{3. Представление колец расширения в виде прямых сумм.} Если
\[
   \mathfrak R=\mathfrak R_1+\cdots+\mathfrak R_r,
\]
то для $\overline{\mathfrak R}$ имеет место представление
\[
   \overline{\mathfrak R}=\overline{\mathfrak R}_1+\cdots+\overline{\mathfrak R}_r,
\]
где $\overline{\mathfrak R}_i$ означает идеал расширения $\mathfrak R_i$ в $\overline{\mathfrak R}$ и одновременно кольцо, выведенное из $\mathfrak R_i$ в $\overline{\mathfrak R}$. Действительно, из $(0)=\mathfrak q_1\cdots\mathfrak q_r$ умножением на $\overline{\mathfrak R}$ получается представление в виде произведения: $(0)=\overline{\mathfrak q}_1\cdots\overline{\mathfrak q}_r$, где, следовательно, $\overline{\mathfrak q}_i$ снова становятся попарно взаимно простыми. Поэтому для $\overline{\mathfrak R}$ получается представление в виде прямой суммы; а так как $i$-я компонента этого представления задается произведением
\[
   \overline{\mathfrak q}_1\cdots\overline{\mathfrak q}_{i-1}\,
   \overline{\mathfrak q}_{i+1}\cdots\overline{\mathfrak q}_r,
\]
она равна идеалу расширения $\mathfrak R_i$. Но этот идеал расширения, по \S 1, 3, тождествен кольцу расширения $\mathfrak R_i[\overline P_i]$, образованному из $\mathfrak R_i$ относительно $\overline P_i=\overline P\varepsilon_i$; следовательно, при расширениях достаточно расширять отдельные компоненты. При этом $\overline{\mathfrak R}_i$ вообще уже не будет первичным кольцом.
% END UNIT BODY
% END PRODUCER UNIT 130
\endgroup


\clearpage
\begingroup
\setcounter{footnote}{0}
% BEGIN PRODUCER UNIT 131
% Source: C:/Users/Floris/Documents/interlanguage/03_projects/noether/02_slavic_working_corpus/translations/paper31/section02_entry04/russian/v001/Noether_Paper31_Section02_Entry04_Russian_v001.tex
% Identity: 2722 B / SHA-256 569055781055F644D83A9E54160FB26977F6B8555FCA3393590C2C3BCA2C0A7C
\setlength{\parindent}{1.25em}
\setlength{\parskip}{0pt}
\makeatletter
\renewcommand\footnotesize{\@setfontsize\footnotesize{7.7}{8.7}\selectfont}
\makeatother
\emergencystretch=85em
\allowdisplaybreaks
\sloppy
\hbadness=10000
\vbadness=10000
\hfuzz=2pt
\vfuzz=10pt
\raggedbottom
% BEGIN UNIT BODY
% Unit: Paper31 Section02 entry 4 v001. Direct Russian translation from the German final-audited source slice, source line 174 / segment P31-S0037. Footnote counter continues after Paper31 Section02 entry 3.
\setcounter{footnote}{20}

\textbf{4. Простые идеалы первого и второго рода.} По пункту 1 кольцо классов вычетов $\mathfrak R/\mathfrak p$ по любому простому идеалу, отличному от единичного идеала, является полем; это поле имеет подполе, изоморфное $P$,
\[
   (P)=(P,\mathfrak p)/\mathfrak p=P/[P,\mathfrak p]\simeq P,
\]
так как $\mathfrak p$ не может содержать никакого ненулевого элемента из $P$. Следовательно, $(P)$ состоит из всех и только тех классов, которые можно представить элементами из $P$.\footnote{Ср. примечание к \S 1, 3 о второй теореме об изоморфизме.} Кольцо $\mathfrak R/\mathfrak p$ имеет конечный ранг относительно $(P)$, то есть является конечным алгебраическим расширением поля $(P)$. В зависимости от того, является ли это расширение расширением первого или второго рода,\footnote{В известной терминологии Штейница: алгебраическое расширение называется расширением первого рода, если каждый его элемент является нулем неприводимого многочлена, который в подходящем поле расширения распадается на различные линейные множители; иначе оно является расширением второго рода.} идеал $\mathfrak p$ будем называть простым идеалом первого или второго рода.
% END UNIT BODY
% END PRODUCER UNIT 131
\endgroup


\clearpage
\begingroup
\setcounter{footnote}{0}
% BEGIN PRODUCER UNIT 132
% Source: C:/Users/Floris/Documents/interlanguage/03_projects/noether/02_slavic_working_corpus/translations/paper31/section03_entry01/russian/v001/Noether_Paper31_Section03_Entry01_Russian_v001.tex
% Identity: 2282 B / SHA-256 A4E69416FF76A570F749440E73C58132195CCB1A4EAC170504CD16E945CE2A67
\setlength{\parindent}{1.25em}
\setlength{\parskip}{0pt}
\makeatletter
\renewcommand\footnotesize{\@setfontsize\footnotesize{7.7}{8.7}\selectfont}
\makeatother
\emergencystretch=85em
\allowdisplaybreaks
\sloppy
\hbadness=10000
\vbadness=10000
\hfuzz=2pt
\vfuzz=10pt
\raggedbottom
% BEGIN UNIT BODY
% Unit: Paper31 Section03 entry 1 v001. Direct Russian translation from the German final-audited source slice, source lines 178--182 / segments P31-S0038--P31-S0040. Footnote counter continues after Paper31 Section02 entry 4.
\setcounter{footnote}{22}

\subsection*{\S 3. Вполне приводимые кольца}

1. Кольцо $\mathfrak R$ называется \emph{вполне приводимым}, если в представлении его нулевого идеала в виде произведения (§2, 2) все первичные компоненты становятся простыми идеалами; или, что по §2, 1 и 2 равносильно этому, если $\mathfrak R$ можно представить в виде прямой суммы конечного числа полей. Вполне приводимое кольцо называется вполне приводимым первого рода, если все простые идеалы являются идеалами первого рода (§2, 4); в противоположном случае --- вполне приводимым второго рода.

Если $P$ является совершенным полем, то возможна только вполне приводимость первого рода; в частности, это всегда так для алгебраически замкнутого поля. В дальнейшем $\overline P$ специально будет обозначать алгебраически замкнутое поле, выведенное из $P$, а $\overline{\mathfrak R}$ будет равно $\mathfrak R[\overline P]$.
% END UNIT BODY
% END PRODUCER UNIT 132
\endgroup


\clearpage
\begingroup
\setcounter{footnote}{0}
% BEGIN PRODUCER UNIT 133
% Source: C:/Users/Floris/Documents/interlanguage/03_projects/noether/02_slavic_working_corpus/translations/paper31/section03_entry02/russian/v001/Noether_Paper31_Section03_Entry02_Russian_v001.tex
% Identity: 2273 B / SHA-256 3CBACCA37C263B357C6CD2BDA5B1CAC4848E9169C3DF3841026473E0984C7DA2
\setlength{\parindent}{1.25em}
\setlength{\parskip}{0pt}
\makeatletter
\renewcommand\footnotesize{\@setfontsize\footnotesize{7.7}{8.7}\selectfont}
\makeatother
\emergencystretch=90em
\allowdisplaybreaks
\sloppy
\hbadness=10000
\vbadness=10000
\hfuzz=2pt
\vfuzz=10pt
\raggedbottom
% BEGIN UNIT BODY
% Unit: Paper31 Section03 entry 2 v001. Direct Russian translation from the German final-audited source slice, source lines 184--196 / segments P31-S0041--P31-S0043. Footnote counter continues after Paper31 Section03 entry 1.
\setcounter{footnote}{22}

2. \textbf{Теорема.} \emph{Кольцо $\mathfrak R$ является вполне приводимым первого рода тогда и только тогда, когда кольцо расширения $\overline{\mathfrak R}=\mathfrak R[\overline P]$ с алгебраически замкнутым $\overline P$ вполне приводимо.}

Доказательство состоит из трех шагов.

2a. Из полной приводимости $\overline{\mathfrak R}$ --- а в более общем виде из полной приводимости любого кольца расширения --- следует полная приводимость (первого или второго рода) кольца $\mathfrak R$. Действительно, по предположению для нулевого идеала $\overline{\mathfrak R}$ имеется представление как наименьшее общее кратное простых идеалов,
\[
        (0)=[\overline{\mathfrak p}_1,\ldots,\overline{\mathfrak p}_s].
\]
Взятием пересечений получается в $\mathfrak R$ представление нулевого идеала
\[
        (0)=[[\mathfrak R,\overline{\mathfrak p}_1],\ldots,[\mathfrak R,\overline{\mathfrak p}_s]],
\]
а идеалы $[\mathfrak R,\overline{\mathfrak p}_i]$ по \S 1, 3 являются простыми идеалами.
% END UNIT BODY
% END PRODUCER UNIT 133
\endgroup


\clearpage
\begingroup
\setcounter{footnote}{0}
% BEGIN PRODUCER UNIT 134
% Source: C:/Users/Floris/Documents/interlanguage/03_projects/noether/02_slavic_working_corpus/translations/paper31/section03_entry03/russian/v001/Noether_Paper31_Section03_Entry03_Russian_v001.tex
% Identity: 3732 B / SHA-256 870AE9E518E20A17FB64B26660314CF43F302501A4EE1FDD6EE50BC56535EB2F
\setlength{\parindent}{1.25em}
\setlength{\parskip}{0pt}
\makeatletter
\renewcommand\footnotesize{\@setfontsize\footnotesize{7.7}{8.7}\selectfont}
\makeatother
\emergencystretch=95em
\allowdisplaybreaks
\sloppy
\hbadness=10000
\vbadness=10000
\hfuzz=2pt
\vfuzz=10pt
\raggedbottom
% BEGIN UNIT BODY
% Unit: Paper31 Section03 entry 3 v001. Direct Russian translation from the German final-audited source slice, source lines 198--200 / segments P31-S0044--P31-S0045. Footnote counter continues after Paper31 Section03 entry 2.
\setcounter{footnote}{22}

2b. Из полной приводимости первого рода кольца $\mathfrak R$ следует полная приводимость первого рода кольца $\overline{\mathfrak R}$; в более общем виде --- полная приводимость первого рода каждого промежуточного кольца между $\mathfrak R$ и $\overline{\mathfrak R}$.

Так как по \S 2, 3 представлению $\mathfrak R$ в виде прямой суммы соответствует такое представление $\overline{\mathfrak R}$, что каждая компонента $\overline{\mathfrak R}_i$ становится кольцом расширения $\mathfrak R_i[\overline P_i]$ --- где $\overline P_i$ является изоморфным $\overline P$ полем расширения подполя $P_i$ компоненты $\mathfrak R_i$, изоморфного $P$, --- достаточно рассмотреть отдельную компоненту $\mathfrak R_i$. Поэтому без ограничения общности можно считать, что $\mathfrak R$ является полем, а именно конечным полем расширения первого рода своего подполя $P$. Тогда $\mathfrak R$ порождается присоединением к $P$ примитивного элемента первого рода; иначе говоря, $\mathfrak R$ изоморфно кольцу классов вычетов $P[x]/f(x)$, где $f(x)$ означает неприводимый многочлен первого рода в полиномиальной области $P[x]$ одной неопределенной $x$. Если $Q$ --- поле расширения $P$, то $Q[x]/f(x)$ изоморфно $\mathfrak R[Q]$; действительно, при переходе от $P[x]$ к $Q[x]$ ранг кольца классов вычетов сохраняется и равен степени $f(x)$, так что речь идет ровно о конструкции кольца расширения, указанной в \S 1, 2. В $Q[x]$ многочлен $f(x)$ распадается на попарно взаимно простые неприводимые многочлены первого рода, каждый из которых порождает простой идеал в $Q[x]$; это означает, что $Q[x]/f(x)$, а из-за изоморфизма также $\mathfrak R[Q]$, становится вполне приводимым первого рода. Так как это верно для каждого поля расширения $Q$, в частности и для алгебраически замкнутого поля $\overline P$, тем самым 2b доказано.
% END UNIT BODY
% END PRODUCER UNIT 134
\endgroup


\clearpage
\begingroup
\setcounter{footnote}{0}
% BEGIN PRODUCER UNIT 135
% Source: C:/Users/Floris/Documents/interlanguage/03_projects/noether/02_slavic_working_corpus/translations/paper31/section03_entry04/russian/v001/Noether_Paper31_Section03_Entry04_Russian_v001.tex
% Identity: 3732 B / SHA-256 CFAF090F2E94532875EF518C76C1CECB2341B930716BD4DBC67C7BB519370C45
\setlength{\parindent}{1.25em}
\setlength{\parskip}{0pt}
\makeatletter
\renewcommand\footnotesize{\@setfontsize\footnotesize{7.7}{8.7}\selectfont}
\makeatother
\emergencystretch=95em
\allowdisplaybreaks
\sloppy
\hbadness=10000
\vbadness=10000
\hfuzz=2pt
\vfuzz=10pt
\raggedbottom
% BEGIN UNIT BODY
% Unit: Paper31 Section03 entry 4 v001. Direct Russian translation from the German final-audited source slice, source lines 202--220 / segments P31-S0046--P31-S0049. Footnote counter continues after Paper31 Section03 entry 3.
\setcounter{footnote}{22}

2c. Если $\mathfrak R$ вполне приводимо второго рода, то $\overline{\mathfrak R}$ не является вполне приводимым.

Как и в 2b, достаточно считать, что $\mathfrak R$ является полем, а именно конечным полем расширения второго рода над $P$. Нужно показать, что в представлении $\overline{\mathfrak R}$ в виде прямой суммы первичных колец возникает по крайней мере одна собственно первичная компонента. Для этого достаточно предъявить ненулевой элемент из $\overline{\mathfrak R}$, некоторая степень которого обращается в нуль. Действительно, если $\bar\beta\ne0$, то хотя бы одна компонента $\bar\beta_i\ne0$; если $\bar\beta^\rho=0$, то $\bar\beta_i^\rho=0$ (\S 2, 2, гомоморфизм из $\overline{\mathfrak R}$ в $\overline{\mathfrak R}_i$), а значит $\overline{\mathfrak R}_i$ не является полем.

Пусть $\gamma$ --- элемент второго рода из $\mathfrak R$, экспоненты $f\ge 1$; пусть
\[
        F(\gamma)=\gamma^{kp^f}+c_1\gamma^{(k-1)p^f}+\cdots+c_k=0
\]
есть уравнение наименьшей степени, которому $\gamma$ удовлетворяет относительно $P$, где $p$ обозначает характеристику $P$. Тогда элементы
\[
        e,\gamma,\gamma^2,\ldots,\gamma^{kp^f-1}
\]
из $\mathfrak R$ линейно независимы относительно $P$ и поэтому по определению равенства остаются линейно независимыми относительно $\overline P$ также в $\overline{\mathfrak R}$. Следовательно, если положить
\[
        G(\gamma)=\gamma^{kp^{f-1}}+\sqrt[p]{c_1}\,\gamma^{(k-1)p^{f-1}}
        +\cdots+\sqrt[p]{c_k},
\]
то $G(\gamma)$ является ненулевым элементом из $\overline{\mathfrak R}$; поскольку при $p>1$ и $f\ge1$ число $kp^{f-1}$ всегда меньше $kp^f$, никакой зависимости между степенями $\gamma$, входящими в $G(\gamma)$, быть не может. Но $G(\gamma)^p=F(\gamma)$ и потому обращается в нуль; тем самым 2c доказано.

Из 2a и 2c следует, что полная приводимость первого рода $\overline{\mathfrak R}$ влечет полную приводимость первого рода $\mathfrak R$; из 2b следует обратное. Тем самым теорема доказана.
% END UNIT BODY
% END PRODUCER UNIT 135
\endgroup


\clearpage
\begingroup
\setcounter{footnote}{0}
% BEGIN PRODUCER UNIT 136
% Source: C:/Users/Floris/Documents/interlanguage/03_projects/noether/02_slavic_working_corpus/translations/paper31/section03_entry05/russian/v001/Noether_Paper31_Section03_Entry05_Russian_v001.tex
% Identity: 4238 B / SHA-256 0EA3B9C76E52CB828AB78FAA14D5AD307DE1DC096C7B544BCB28814B3653A667
\setlength{\parindent}{1.25em}
\setlength{\parskip}{0pt}
\makeatletter
\renewcommand\footnotesize{\@setfontsize\footnotesize{7.7}{8.7}\selectfont}
\makeatother
\emergencystretch=100em
\allowdisplaybreaks
\sloppy
\hbadness=10000
\vbadness=10000
\hfuzz=2pt
\vfuzz=10pt
\raggedbottom
% BEGIN UNIT BODY
% Unit: Paper31 Section03 entry 5 v001. Direct Russian translation from the German final-audited source slice, source lines 222--235 / segments P31-S0050--P31-S0052. Footnote counter continues after Paper31 Section03 entry 4.
\setcounter{footnote}{22}

3. \textbf{Следствие.} \emph{Если $\mathfrak R$ вполне приводимо первого рода, то $\overline{\mathfrak R}$ становится прямой суммой колец ранга один, следовательно полей, изоморфных $\overline P$. Такое разложение на кольца ранга один существует уже в кольце $\mathfrak R[Q]$, где $Q$ является конечным полем расширения $P$.}

Действительно, по 2b кольцо $\overline{\mathfrak R}$ становится прямой суммой полей, которые как конечные расширения подполей, изоморфных $\overline P$, вследствие алгебраической замкнутости $\overline P$ совпадают с этими подполями. Поэтому получаем
\[
        \overline{\mathfrak R}=\overline P\bar e_1+\cdots+\overline P\bar e_n,
        \qquad e=\bar e_1+\cdots+\bar e_n;
\]
компоненты $\bar e_i$ элемента $e$ одновременно образуют линейно независимую модульную базу $\overline{\mathfrak R}$.

Если $\bar e_i$ по \S 1, 2 представить как линейные комбинации элементов $\gamma$ из $\mathfrak R$ с коэффициентами из $\overline P$, то система всех коэффициентов определяет конечное поле расширения $Q$ поля $P$; следовательно, $\bar e_i$ принадлежат $\mathfrak R[Q]$, и из свойства модульной базы получается
\[
        \mathfrak R[Q]=Q\bar e_1+\cdots+Q\bar e_n.
\]
Тем самым уже в $\mathfrak R[Q]$ существует разложение на кольца ранга один, а в каждом промежуточном кольце между $\mathfrak R[Q]$ и $\overline{\mathfrak R}$ неразложимые компоненты возникают как кольца расширения колец $Q\bar e_i$.\footnote{Наименьшее такое поле расширения $Q$ можно охарактеризовать как композит всех полей Галуа, определенных компонентами $\mathfrak R_i=\mathfrak R e_i$ кольца $\mathfrak R$. А именно, если $f_i(x)$ выбрано по 2b так, что $P[x]/f_i(x)$ соответственно изоморфно компоненте $\mathfrak R_i$, и если $Q^{(i)}$ является полем расширения Галуа, лежащим в $\overline P$ и ровно достаточным для разложения $f_i(x)$ на линейные множители, то $Q$ есть композит всех $Q^{(i)}$. Разложению $f_i(x)$ на линейные множители соответствует представление $Q^{(i)}[x]/f_i(x)$ в виде прямой суммы полей ранга один; при этом $Q^{(i)}$ является наименьшим полем расширения, в котором это происходит. В силу изоморфии то же верно для $\mathfrak R_i[Q^{(i)}\bar e_i]$, а следовательно для $\mathfrak R[Q]$.}
% END UNIT BODY
% END PRODUCER UNIT 136
\endgroup


\clearpage
\begingroup
\setcounter{footnote}{0}
% BEGIN PRODUCER UNIT 137
% Source: C:/Users/Floris/Documents/interlanguage/03_projects/noether/02_slavic_working_corpus/translations/paper31/section04_entry01/russian/v001/Noether_Paper31_Section04_Entry01_Russian_v001.tex
% Identity: 3062 B / SHA-256 06B3671DBE42BFAB9496D31A730DB75CF2A4D98A7499F2540D53503690336D1B
\setlength{\parindent}{1.25em}
\setlength{\parskip}{0pt}
\makeatletter
\renewcommand\footnotesize{\@setfontsize\footnotesize{7.7}{8.7}\selectfont}
\makeatother
\emergencystretch=100em
\allowdisplaybreaks
\sloppy
\hbadness=10000
\vbadness=10000
\hfuzz=2pt
\vfuzz=10pt
\raggedbottom
% BEGIN UNIT BODY
% Unit: Paper31 Section04 entry 1 v001. Direct Russian translation from the German final-audited source slice, source lines 237--243 / segments P31-S0053--P31-S0056. Footnote counter continues after Paper31 Section03 entry 5.
\setcounter{footnote}{23}

\subsection*{§4. Матричные представления, следы и дискриминанты для колец конечного ранга}

В этом и следующем параграфах предположения об основном кольце будут немного более общими, чем раньше, чтобы охватить также порядки в числовых и функциональных полях. Речь идет об общей и единообразной формулировке по существу известных фактов.

\emph{Предположение.} Пусть $\mathfrak R$ является кольцом расширения области целостности $\mathfrak Z$; единичный элемент кольца $\mathfrak R$ уже лежит в $\mathfrak Z$. Для каждого $\mathfrak Z$-модуля из $\mathfrak R$ -- в частности для самого $\mathfrak R$ -- существует хотя бы одна конечная модульная база, состоящая из элементов, линейно независимых над $\mathfrak Z$. Помимо $\alpha_1,\ldots,\alpha_s$, элементы $\beta_1,\ldots,\beta_s$ тогда и только тогда образуют линейно независимую модульную базу того же $\mathfrak Z$-модуля, когда определитель преобразования является единицей в $\mathfrak Z$.

Это предположение, очевидно, выполнено для рассмотренных до сих пор колец конечного ранга, где $\mathfrak Z$ становится полем $P$. Оно также выполнено для порядков в числовых и функциональных полях, где $\mathfrak Z$ является соответственно кольцом рациональных целых чисел или функциональной областью.\footnote{См. \emph{Идеальная теория}, §3. Элементы из $\mathfrak Z$ обозначаются латинскими буквами.}
% END UNIT BODY
% END PRODUCER UNIT 137
\endgroup


\clearpage
\begingroup
\setcounter{footnote}{0}
% BEGIN PRODUCER UNIT 138
% Source: C:/Users/Floris/Documents/interlanguage/03_projects/noether/02_slavic_working_corpus/translations/paper31/section04_entry02/russian/v001/Noether_Paper31_Section04_Entry02_Russian_v001.tex
% Identity: 2978 B / SHA-256 70AFD403A25805EF379909C3FC05B70A767375F20C5CF334786FD1F804E41810
\setlength{\parindent}{1.25em}
\setlength{\parskip}{0pt}
\makeatletter
\renewcommand\footnotesize{\@setfontsize\footnotesize{7.7}{8.7}\selectfont}
\makeatother
\emergencystretch=100em
\allowdisplaybreaks
\sloppy
\hbadness=10000
\vbadness=10000
\hfuzz=2pt
\vfuzz=10pt
\raggedbottom
% BEGIN UNIT BODY
% Unit: Paper31 Section04 entry 2 v001. Direct Russian translation from the German final-audited source slice, source lines 245--267 / segment P31-S0057. Footnote counter continues after Paper31 Section04 entry 1.
\setcounter{footnote}{24}

1. \emph{Матричные кольца, гомоморфные к $\mathfrak R$.} Пусть $\alpha_1,\ldots,\alpha_s$ -- линейно независимая $\mathfrak Z$-модульная база идеала $\mathfrak a$, и пусть $\gamma$ -- произвольный элемент из $\mathfrak R$. Тогда $\gamma\alpha_i$ также принадлежит $\mathfrak a$, и имеются равенства с коэффициентами из $\mathfrak Z$:
\[
        \gamma\alpha_i=c_{i1}\alpha_1+\cdots+c_{is}\alpha_s,
        \qquad i=1,\ldots,s.
\]
Или в матричной форме, если всякий раз $(\tau_1,\ldots,\tau_s)$ означает однострочную матрицу, составленную из этих элементов:
\[
 (\gamma\alpha_1,\ldots,\gamma\alpha_s)
 = (\alpha_1,\ldots,\alpha_s)
   \begin{pmatrix}
      c_{11}&\cdots&c_{1s}\\
      \vdots&&\vdots\\
      c_{s1}&\cdots&c_{ss}
   \end{pmatrix}
 = (\alpha_1,\ldots,\alpha_s)C.
\]
Когда $\gamma$ пробегает все элементы из $\mathfrak R$, совокупность матриц $C$, сопоставленных посредством базы $\alpha_i$, образует гомоморфное кольцо $R_\mathfrak a$; при этом $R_\mathfrak a$ изоморфно кольцу классов вычетов $\mathfrak R/((0):\mathfrak a)$, где $(0):\mathfrak a$ означает частное нулевого идеала по $\mathfrak a$. Действительно, разности $\beta-\gamma$ сопоставляется разность $B-C$, и то же верно для произведения, поскольку
\[
(\beta\gamma\alpha_1,\ldots,\beta\gamma\alpha_s)
=(\gamma\alpha_1,\ldots,\gamma\alpha_s)B
=(\alpha_1,\ldots,\alpha_s)CB.
\]
Элементам $c$ из $\mathfrak Z$ в частности сопоставляются диагональные матрицы $cE$. Всем и только тем элементам $\gamma$, для которых $\gamma\mathfrak a$ равен нулевому идеалу, сопоставляется нулевая матрица.
% END UNIT BODY
% END PRODUCER UNIT 138
\endgroup


\clearpage
\begingroup
\setcounter{footnote}{0}
% BEGIN PRODUCER UNIT 139
% Source: C:/Users/Floris/Documents/interlanguage/03_projects/noether/02_slavic_working_corpus/translations/paper31/section04_entry03/russian/v001/Noether_Paper31_Section04_Entry03_Russian_v001.tex
% Identity: 1691 B / SHA-256 1862C8440890DA0AEB72F355485D54E7566B71B1944EA7B3921551C0514627B3
\setlength{\parindent}{1.25em}
\setlength{\parskip}{0pt}
\makeatletter
\renewcommand\footnotesize{\@setfontsize\footnotesize{7.7}{8.7}\selectfont}
\makeatother
\emergencystretch=100em
\allowdisplaybreaks
\sloppy
\hbadness=10000
\vbadness=10000
\hfuzz=2pt
\vfuzz=10pt
\raggedbottom
% BEGIN UNIT BODY
% Unit: Paper31 Section04 entry 3 v001. Direct Russian translation from the German final-audited source slice, source lines 269--273 / segment P31-S0058. Footnote counter continues after Paper31 Section04 entry 2.
\setcounter{footnote}{24}

Другая база $\bar\alpha_1,\ldots,\bar\alpha_s$ того же идеала порождает матричное кольцо, изоморфное $R_{\mathfrak a}$, причем эквивалентное ему:
\[
        R_{\bar{\mathfrak a}}=P^{-1}R_{\mathfrak a}P.
\]
Действительно, пусть $\bar\alpha_1,\ldots,\bar\alpha_s$ -- еще одна линейно независимая $\mathfrak Z$-модульная база; следовательно, $(\bar\alpha_1,\ldots,\bar\alpha_s)=(\alpha_1,\ldots,\alpha_s)P$, где определитель матрицы $P$ является единицей из $\mathfrak Z$. Тогда
\[
\begin{aligned}
(\gamma\bar\alpha_1,\ldots,\gamma\bar\alpha_s)
&=(\alpha_1,\ldots,\alpha_s)CP\\
&=(\bar\alpha_1,\ldots,\bar\alpha_s)P^{-1}CP,
\end{aligned}
\]
то есть $R_{\bar{\mathfrak a}}=P^{-1}R_{\mathfrak a}P$.
% END UNIT BODY
% END PRODUCER UNIT 139
\endgroup


\clearpage
\begingroup
\setcounter{footnote}{0}
% BEGIN PRODUCER UNIT 140
% Source: C:/Users/Floris/Documents/interlanguage/03_projects/noether/02_slavic_working_corpus/translations/paper31/section04_entry04/russian/v001/Noether_Paper31_Section04_Entry04_Russian_v001.tex
% Identity: 2627 B / SHA-256 DEAB0A9D6AA0C7138626D29031D16B5C64CAE441BB5FFDF23D0BDDB3566CFF6E
\setlength{\parindent}{1.25em}
\setlength{\parskip}{0pt}
\makeatletter
\renewcommand\footnotesize{\@setfontsize\footnotesize{7.7}{8.7}\selectfont}
\makeatother
\emergencystretch=100em
\allowdisplaybreaks
\sloppy
\hbadness=10000
\vbadness=10000
\hfuzz=2pt
\vfuzz=10pt
\raggedbottom
% BEGIN UNIT BODY
% Unit: Paper31 Section04 entry 4 v001. Direct Russian translation from the German final-audited source slice, source lines 275--279 / segment P31-S0059, expanded against the printed scan/OCR witness. Footnote counter continues after Paper31 Section04 entry 3.
\setcounter{footnote}{24}

\textbf{2. \emph{Классы идеалов и классы представлений.}} Два идеала $\mathfrak a$ и $\mathfrak b$ из $\mathfrak R$ принадлежат одному и тому же классу идеалов, если они изоморфны как $\mathfrak R$-модули; то есть если их элементы можно взаимно однозначно сопоставить так, чтобы разность и умножение на один и тот же элемент из $\mathfrak R$ соответствовали друг другу. Следовательно, из $\alpha\simeq\beta$ всегда следует $\gamma\alpha\simeq\gamma\beta$ для произвольного $\gamma$ из $\mathfrak R$.

Два кольца $R_\mathfrak a$ и $R_\mathfrak b$, порожденные матричным представлением по 1, принадлежат одному и тому же классу представлений, если они эквивалентны, то есть если существует унимодулярная матрица $P$ такая, что
\[
        R_\mathfrak b=P^{-1}R_\mathfrak a P
\]
выполняется.\footnote{Вопрос, решенный в отдельных случаях, можно ли каждое гомоморфное матричное кольцо породить идеалом из $\mathfrak R$ по 1 -- возможно, допуская линейно зависимые базы, -- здесь не затрагивается. Классы представлений по определению состоят только из матричных колец, порожденных идеалами из $\mathfrak R$ по 1.}
% END UNIT BODY
% END PRODUCER UNIT 140
\endgroup


\clearpage
\begingroup
\setcounter{footnote}{0}
% BEGIN PRODUCER UNIT 141
% Source: C:/Users/Floris/Documents/interlanguage/03_projects/noether/02_slavic_working_corpus/translations/paper31/section04_entry05/russian/v001/Noether_Paper31_Section04_Entry05_Russian_v001.tex
% Identity: 4102 B / SHA-256 C28042EBA60CB5AB7A6D0D4B55DEDE7D919943929348311939B5E0B48959BA6D
\setlength{\parindent}{1.25em}
\setlength{\parskip}{0pt}
\makeatletter
\renewcommand\footnotesize{\@setfontsize\footnotesize{7.7}{8.7}\selectfont}
\makeatother
\emergencystretch=100em
\allowdisplaybreaks
\sloppy
\hbadness=10000
\vbadness=10000
\hfuzz=2pt
\vfuzz=10pt
\raggedbottom
% BEGIN UNIT BODY
% Unit: Paper31 Section04 entry 5 v001. Direct Russian translation from the German final-audited source slice, source line 281 / segment P31-S0060, expanded against printed scan/OCR witness pages 93--94. Footnote counter continues after Paper31 Section04 entry 4.
\setcounter{footnote}{25}

Классы идеалов и классы представлений взаимно однозначно соответствуют друг другу. В пункте 1 было показано, что различные модульные базы одного и того же идеала порождают кольца $R_\alpha$, $R_{\bar\alpha}$, принадлежащие одному и тому же классу представлений; далее, различные идеалы $\mathfrak a$, $\mathfrak b$ одного и того же класса идеалов также порождают кольца одного и того же класса представлений. Действительно, пусть $\beta_1,\ldots,\beta_s$ есть база $\mathfrak b$, поставленная посредством изоморфизма в соответствие базе $\alpha_1,\ldots,\alpha_s$ идеала $\mathfrak a$. Тогда друг другу соответствуют $\gamma\alpha_i$ и $\gamma\beta_i$, а также $c_{1i}\alpha_1+\cdots+c_{si}\alpha_s$ и $c_{1i}\beta_1+\cdots+c_{si}\beta_s$; следовательно, кольца $R_\alpha$ и $R_\beta$ совпадают.

Обратно, если $R_\alpha$ и $R_{\bar\beta}$ принадлежат одному и тому же классу представлений, то есть если
\[
        R_{\bar\beta}=P^{-1}R_\alpha P,
\]
где $\alpha_i$ образуют базу $\mathfrak a$, а $\bar\beta_i$ образуют базу $\mathfrak b$, то
\[
        (\beta_1,\ldots,\beta_s)=(\bar\beta_1,\ldots,\bar\beta_s)P^{-1}
\]
также образует базу $\mathfrak b$, и получается $R_\beta=PR_{\bar\beta}P^{-1}=R_\alpha$. Если теперь каждому элементу $a=c_1\alpha_1+\cdots+c_s\alpha_s$ из $\mathfrak a$ сопоставить элемент $b=c_1\beta_1+\cdots+c_s\beta_s$ из $\mathfrak b$, то это соответствие взаимно однозначно, а сумма и разность сохраняются. Из равенства $R_\alpha=R_\beta$ далее следует,\footnote{Запись устроена так, чтобы она сохранялась для некоммутативных колец $\mathfrak R$, для которых $\Sigma$ предполагается коммутирующей со всеми элементами.} что $\gamma\alpha_i$ и $\gamma\beta_i$ соответствуют друг другу, поскольку они выражаются одними и теми же линейными формами через $\alpha_i$, соответственно через $\beta_i$; следовательно, в общем случае соответствуют друг другу $\gamma a$ и $\gamma b$. Поэтому идеалы $\mathfrak a$ и $\mathfrak b$ принадлежат одному и тому же классу идеалов. Главным классом называется класс единичного идеала; следовательно, всякая база $\mathfrak R$ дает в качестве класса представлений главный класс.
% END UNIT BODY
% END PRODUCER UNIT 141
\endgroup


\clearpage
\begingroup
\setcounter{footnote}{0}
% BEGIN PRODUCER UNIT 142
% Source: C:/Users/Floris/Documents/interlanguage/03_projects/noether/02_slavic_working_corpus/translations/paper31/section04_entry06/russian/v001/Noether_Paper31_Section04_Entry06_Russian_v001.tex
% Identity: 2156 B / SHA-256 941AC17796FE9C4BFD7D4931454908692B789EAFC169014C8E3C77A75F16F07A
\setlength{\parindent}{1.25em}
\setlength{\parskip}{0pt}
\makeatletter
\renewcommand\footnotesize{\@setfontsize\footnotesize{7.7}{8.7}\selectfont}
\makeatother
\emergencystretch=100em
\allowdisplaybreaks
\sloppy
\hbadness=10000
\vbadness=10000
\hfuzz=2pt
\vfuzz=10pt
\raggedbottom
% BEGIN UNIT BODY
% Unit: Paper31 Section04 entry 6 v001. Direct Russian translation from the German final-audited source slice, source lines 283--287 / segment P31-S0061, expanded against printed scan/OCR witness page 94. Footnote counter continues after Paper31 Section04 entry 5.
\setcounter{footnote}{26}

\textbf{3. \emph{След элемента относительно класса.}} Пусть $C$ есть матрица, сопоставленная элементу $\gamma$ из $\mathfrak R$ в $R_{\mathfrak a}$. Переход от $R_{\mathfrak a}$ к эквивалентному кольцу показывает, что детерминант $|C|$, а в более общем виде $|tE-C|$, где $t$ означает неопределенную, является инвариантом класса представлений. Следовательно, по пункту 2 это одновременно инвариант класса идеалов, то есть, коротко говоря, классовый инвариант. Поэтому отдельные коэффициенты по $t$ в $|tE-C|$ являются классовыми инвариантами; в частности, коэффициент при $(-t^{s-1})$ будем называть следом $\gamma$ относительно класса; в обозначениях:
\[
        S_{(\mathfrak a)}(\gamma)=c_{11}+c_{22}+\cdots+c_{ss}.
\]
Свободный от $t$ коэффициент $\pm |C|$ будем называть нормой $\gamma$ относительно класса.
% END UNIT BODY
% END PRODUCER UNIT 142
\endgroup


\clearpage
\begingroup
\setcounter{footnote}{0}
% BEGIN PRODUCER UNIT 143
% Source: C:/Users/Floris/Documents/interlanguage/03_projects/noether/02_slavic_working_corpus/translations/paper31/section04_entry07/russian/v001/Noether_Paper31_Section04_Entry07_Russian_v001.tex
% Identity: 2641 B / SHA-256 ED992D64849921D155EED8AF978A48EE3C4950FA581EC964B2B0794712621117
\setlength{\parindent}{1.25em}
\setlength{\parskip}{0pt}
\makeatletter
\renewcommand\footnotesize{\@setfontsize\footnotesize{7.7}{8.7}\selectfont}
\makeatother
\emergencystretch=100em
\allowdisplaybreaks
\sloppy
\hbadness=10000
\vbadness=10000
\hfuzz=2pt
\vfuzz=10pt
\raggedbottom
% BEGIN UNIT BODY
% Unit: Paper31 Section04 entry 7 v001. Direct Russian translation from the German final-audited source slice, source lines 289--295 / segment P31-S0062, expanded against printed scan/OCR witness page 94. Footnote counter continues after Paper31 Section04 entry 6.
\setcounter{footnote}{26}

Следы $S_{(\mathfrak a)}(\gamma)$ относительно фиксированного класса $(\mathfrak a)$ образуют $\mathfrak Z$-модуль, в который $\mathfrak R$, рассматриваемое как $\mathfrak Z$-модуль, гомоморфно отображается. Действительно, по доказанному в пункте 1 гомоморфизму колец из $\mathfrak R$ в $R_{\mathfrak a}$ получается: $(\beta-\gamma)\sim(B-C)$ и $c\beta\sim cEB=cB$. Отсюда для следов следует:
\[
        S_{(\mathfrak a)}(\beta-\gamma)=S_{(\mathfrak a)}(\beta)-S_{(\mathfrak a)}(\gamma),
        \qquad
        S_{(\mathfrak a)}(c\beta)=cS_{(\mathfrak a)}(\beta).
\]
Следовательно, получаем модульное свойство и гомоморфизм.

\emph{Замечание.} Если $\mathfrak R$ является кольцом без делителей нуля, то след элемента определяется независимо от класса; значит, можно говорить просто о следе; то же верно для нормы. Ведь когда речь идет о кольцах без делителей нуля, всякий идеал, отличный от нулевого идеала, имеет ранг кольца; поэтому всякая база идеала одновременно является базой единичного идеала в поле частных. Все следы и нормы тем самым одновременно порождаются в поле частных единичным идеалом и потому не зависят от класса в $\mathfrak R$.
% END UNIT BODY
% END PRODUCER UNIT 143
\endgroup


\clearpage
\begingroup
\setcounter{footnote}{0}
% BEGIN PRODUCER UNIT 144
% Source: C:/Users/Floris/Documents/interlanguage/03_projects/noether/02_slavic_working_corpus/translations/paper31/section04_entry08/russian/v001/Noether_Paper31_Section04_Entry08_Russian_v001.tex
% Identity: 3331 B / SHA-256 DB5B361E0037561F7556816B36F4AB08E24D98B9BC725D7D47467C7C7F5CCDF2
\setlength{\parindent}{1.25em}
\setlength{\parskip}{0pt}
\makeatletter
\renewcommand\footnotesize{\@setfontsize\footnotesize{7.7}{8.7}\selectfont}
\makeatother
\emergencystretch=100em
\allowdisplaybreaks
\sloppy
\hbadness=10000
\vbadness=10000
\hfuzz=2pt
\vfuzz=10pt
\raggedbottom
% BEGIN UNIT BODY
% Unit: Paper31 Section04 entry 8 v001. Direct Russian translation from the German final-audited source slice, source lines 297--305 / segments P31-S0063--P31-S0064, expanded against printed scan/OCR witness page 94. Footnote counter continues after Paper31 Section04 entry 7.
\setcounter{footnote}{26}

\textbf{4. \emph{Дискриминант идеала относительно класса.}} Инвариантом идеала и класса оказывается не сам дискриминант, а лишь выведенный из него в $\mathfrak Z$ главный идеал.

Действительно, пусть $\omega_1,\ldots,\omega_s$ и $\bar\omega_1,\ldots,\bar\omega_s$ -- две линейно независимые $\mathfrak Z$-модульные базы идеала $\mathfrak t$ из $\mathfrak R$. Детерминанты
\[
        |S_{(\mathfrak a)}(\omega_i\omega_j)|,
        \qquad
        |S_{(\mathfrak a)}(\bar\omega_i\bar\omega_j)|
\]
отличаются только квадратом единицы из $\mathfrak Z$ в качестве множителя. Это действительно верно для двух детерминантов $|(\omega_i\omega_j)|$ и $|(\bar\omega_i\bar\omega_j)|$, которые отличаются квадратом определителя преобразования между $\omega$ и $\bar\omega$. Благодаря модульному гомоморфизму, доказанному в пункте 3, между $S_{(\mathfrak a)}(\omega_i\omega_j)$ и $S_{(\mathfrak a)}(\bar\omega_i\bar\omega_j)$ имеются те же линейные соотношения, что и между $\omega_i\omega_j$ и $\bar\omega_i\bar\omega_j$; следовательно, вследствие когредиентного преобразования детерминант принимает тот же множитель.

Дискриминантом $\mathfrak t$ относительно класса $(\mathfrak a)$, определенным с точностью до квадрата единицы, называется каждый детерминант $|S_{(\mathfrak a)}(\omega_i\omega_j)|$; дискриминантным идеалом -- выведенный из него в $\mathfrak Z$ главный идеал.

\emph{Замечание.} Если $\mathfrak R$ является кольцом без делителей нуля, то дискриминант идеала не зависит от положенного в основу класса. Действительно, по пункту 3 эта независимость выполнена для каждого отдельного элемента детерминанта.
% END UNIT BODY
% END PRODUCER UNIT 144
\endgroup


\clearpage
\begingroup
\setcounter{footnote}{0}
% BEGIN PRODUCER UNIT 145
% Source: C:/Users/Floris/Documents/interlanguage/03_projects/noether/02_slavic_working_corpus/translations/paper31/section05_entry01/russian/v001/Noether_Paper31_Section05_Entry01_Russian_v001.tex
% Identity: 5715 B / SHA-256 9BE9781960EC8B05F1425ABE69F9D1ADA6A60C1CAC56E7E14F6A2464F3BDBFD4
\setlength{\parindent}{1.25em}
\setlength{\parskip}{0pt}
\makeatletter
\renewcommand\footnotesize{\@setfontsize\footnotesize{7.7}{8.7}\selectfont}
\makeatother
\emergencystretch=100em
\allowdisplaybreaks
\sloppy
\hbadness=10000
\vbadness=10000
\hfuzz=2pt
\vfuzz=10pt
\raggedbottom
% BEGIN UNIT BODY
% Unit: Paper31 Section05 entry 1 v001. Direct Russian translation from the German final-audited source slice, source lines 307--317 / segments P31-S0065--P31-S0068, expanded against printed scan/OCR witness page 95. Footnote counter continues after Paper31 Section04 entry 8.
\setcounter{footnote}{26}

\subsection*{§5. Прямые суммы классов}

\textbf{1. \emph{Прямая сумма класса идеалов или класса представлений.}} Если идеал $\mathfrak a$ является прямой суммой, $\mathfrak a=\mathfrak b+\mathfrak c$, то из изоморфизма следует, что каждый идеал этого класса также становится прямой суммой,
$\bar{\mathfrak a}=\bar{\mathfrak b}+\bar{\mathfrak c}$. При этом соответственно $\bar{\mathfrak b}$ изоморфен $\mathfrak b$, а $\bar{\mathfrak c}$ -- $\mathfrak c$, если идеалы рассматривать как $\mathfrak R$-модули; следовательно, можно говорить о прямой сумме класса идеалов и о компонентных классах.

Если кольцо $R_{\mathfrak a}$ некоторого класса представлений является прямой суммой подколец, которые одновременно являются идеалами в $R_{\mathfrak a}$, то благодаря изоморфизму то же верно для каждого кольца $R_{\mathfrak b}$ этого класса представлений, причем компоненты становятся эквивалентными кольцами. Следовательно, можно говорить о прямой сумме класса представлений и о компонентных классах.

Прямой сумме класса идеалов соответствует прямая сумма класса представлений; именно в этом смысле будем говорить о прямой сумме класса.\footnote{Вопрос, в какой мере все прямые суммы классов представлений могут быть порождены прямыми суммами классов идеалов, здесь снова не рассматривается.} Пусть $\beta_1,\ldots,\beta_m$ есть база $\mathfrak b$, а $\gamma_1,\ldots,\gamma_n$ -- база $\mathfrak c$; так как сумма прямая, $\beta$ и $\gamma$ вместе образуют линейно независимую базу $\mathfrak a$. Если $R_{\mathfrak a}$ есть матричное кольцо, порожденное посредством этой базы, то матрица из $R_{\mathfrak a}$, сопоставленная произвольному элементу $\delta$ из $\mathfrak R$, имеет вид
\[
        \begin{pmatrix}D^{\mathfrak b}&0\\0&D^{\mathfrak c}\end{pmatrix}.
\]
Поэтому определим $R^{\mathfrak b}$ как систему всех матриц
\[
        \begin{pmatrix}D^{\mathfrak b}&0\\0&0\end{pmatrix}
\]
и соответственно определим $R^{\mathfrak c}$. Нужно показать, что $R^{\mathfrak b}$ и $R^{\mathfrak c}$ образуют подкольца в $R_{\mathfrak a}$ и что $R_{\mathfrak a}$ является их прямой суммой; одновременно $R^{\mathfrak b}$ и $R^{\mathfrak c}$ становятся изоморфными кольцам $R_{\mathfrak b}$ и $R_{\mathfrak c}$, порожденным соответственно из $\mathfrak b$ и $\mathfrak c$.

Действительно, через гомоморфизм $R^{\mathfrak b}$ соответствует всем таким элементам $\delta$ из $\mathfrak R$, для которых $\delta\mathfrak c$ обращается в нуль, то есть которые принадлежат частному нулевого идеала по $\mathfrak c$, $(0):\mathfrak c$; благодаря гомоморфизму $R^{\mathfrak b}$ тем самым является идеалом в $R_{\mathfrak a}$, и то же верно для $R^{\mathfrak c}$. Кольцо $R_{\mathfrak a}$ является суммой этих идеалов, причем прямой суммой, потому что произведение $R^{\mathfrak b}R^{\mathfrak c}$ двух колец обращается в нуль. Наконец, соответствие, которое матрице $D^{\mathfrak b}$ ставит в соответствие матрицу из $R^{\mathfrak b}$, содержащую $D^{\mathfrak b}$ и нули, дает изоморфизм между $R^{\mathfrak b}$ и $R_{\mathfrak b}$; в этом смысле компонентным классам $(\mathfrak b)$ и $(\mathfrak c)$ класса идеалов $\mathfrak a$ соответствуют компонентные классы кольца $R_{\mathfrak a}$.
% END UNIT BODY
% END PRODUCER UNIT 145
\endgroup


\clearpage
\begingroup
\setcounter{footnote}{0}
% BEGIN PRODUCER UNIT 146
% Source: C:/Users/Floris/Documents/interlanguage/03_projects/noether/02_slavic_working_corpus/translations/paper31/section05_entry02/russian/v001/Noether_Paper31_Section05_Entry02_Russian_v001.tex
% Identity: 3460 B / SHA-256 C2BD18FDAE54583BC78A71E7245853D94A7DFC597F16E236A5FBB5E9F030A0AB
\setlength{\parindent}{1.25em}
\setlength{\parskip}{0pt}
\makeatletter
\renewcommand\footnotesize{\@setfontsize\footnotesize{7.7}{8.7}\selectfont}
\makeatother
\emergencystretch=100em
\allowdisplaybreaks
\sloppy
\hbadness=10000
\vbadness=10000
\hfuzz=2pt
\vfuzz=10pt
\raggedbottom
% BEGIN UNIT BODY
% Unit: Paper31 Section05 entry 2 v001. Direct Russian translation from the German final-audited source slice, source lines 319--330 / segments P31-S0069--P31-S0070, expanded against printed scan/OCR witness pages 95--96. Footnote counter continues after Paper31 Section05 entry 1.
\setcounter{footnote}{26}

\textbf{2. \emph{Поведение следа и дискриминанта при прямой сумме классов.}} Пусть $\mathfrak a=\mathfrak b+\mathfrak c$. Если для определения следа $S_{(\mathfrak a)}(\delta)$ положить в основу блочную матрицу
\[
        \begin{pmatrix}D^{\mathfrak b}&0\\0&D^{\mathfrak c}\end{pmatrix},
\]
то с учетом пункта 1 непосредственно читаем: след элемента, взятый относительно класса, равен сумме следов, взятых относительно компонентных классов.

Ввиду обращения произведения $\mathfrak b\mathfrak c$ в нуль далее получаем: если $\beta$ есть элемент из $\mathfrak b$, то
\[
        S_{(\mathfrak a)}(\beta)=S_{(\mathfrak b)}(\beta);
\]
след $\beta$, взятый относительно класса $(\mathfrak a)$, равен следу, взятому относительно класса $(\mathfrak b)$.

Отсюда, с учетом §4, 4, далее следует: дискриминантный идеал компонентного идеала $\mathfrak b$, взятый относительно класса $(\mathfrak a)$, равен дискриминантному идеалу $\mathfrak b$, взятому относительно компонентного класса $(\mathfrak b)$.

И наконец получаем: дискриминантный идеал $\mathfrak a$, взятый относительно класса $(\mathfrak a)$, равен произведению дискриминантных идеалов компонент, взятых относительно компонентных классов. Действительно, если и для построения дискриминанта, и для построения следа положить в основу базу $\beta,\gamma$, соответствующую прямой сумме, то база дискриминантного идеала $\mathfrak a$, взятого относительно $(\mathfrak a)$, задается через
\[
\left|
\begin{array}{cc}
S_{(\mathfrak a)}(\beta_i\beta_k)&0\\
0&S_{(\mathfrak a)}(\gamma_\mu\gamma_\nu)
\end{array}
\right|
=
\left|
\begin{array}{cc}
S_{(\mathfrak b)}(\beta_i\beta_k)&0\\
0&S_{(\mathfrak c)}(\gamma_\mu\gamma_\nu)
\end{array}
\right|,
\qquad
\begin{matrix}
 i,k=1,\ldots,l,\\
 \mu,\nu=1,\ldots,m.
\end{matrix}
\]
% END UNIT BODY
% END PRODUCER UNIT 146
\endgroup


\clearpage
\begingroup
\setcounter{footnote}{0}
% BEGIN PRODUCER UNIT 147
% Source: C:/Users/Floris/Documents/interlanguage/03_projects/noether/02_slavic_working_corpus/translations/paper31/section05_entry03/russian/v001/Noether_Paper31_Section05_Entry03_Russian_v001.tex
% Identity: 3025 B / SHA-256 5E9E494D98CF7CB2EFD2367BA0FF1656E2A71779BC56E87BEBBCC9024530D9A1
\setlength{\parindent}{1.25em}
\setlength{\parskip}{0pt}
\makeatletter
\renewcommand\footnotesize{\@setfontsize\footnotesize{7.7}{8.7}\selectfont}
\makeatother
\emergencystretch=100em
\allowdisplaybreaks
\sloppy
\hbadness=10000
\vbadness=10000
\hfuzz=2pt
\vfuzz=10pt
\raggedbottom
% BEGIN UNIT BODY
% Unit: Paper31 Section05 entry 3 v001. Direct Russian translation from the German final-audited source slice, source line 332 / segment P31-S0071, expanded against printed scan/OCR witness page 96. Footnote counter continues after Paper31 Section05 entry 2.
\setcounter{footnote}{26}

\textbf{3. \emph{Переход к кольцам расширения.}} Пусть $\mathfrak X$ является кольцом расширения кольца $\mathfrak Z$ и не имеет делителей нуля, причем пересечение $[\mathfrak R,\mathfrak X]$ равно $\mathfrak Z$; пусть кольцо расширения $\overline{\mathfrak R}=\mathfrak R[\mathfrak X]$, возникшее присоединением $\mathfrak X$, определено, как в \S~1, 2, как кольцо того же ранга.

Если $\mathfrak a$ и $\mathfrak t$ являются идеалами в $\mathfrak R$, а $\overline{\mathfrak a}$ и $\overline{\mathfrak t}$ -- их идеалами расширения в $\overline{\mathfrak R}$, то дискриминантный идеал $\overline{\mathfrak t}$, взятый относительно $(\overline{\mathfrak a})$, также является идеалом расширения в $\mathfrak X$ того дискриминантного идеала $\mathfrak t$, который взят относительно $(\mathfrak a)$. Ведь всякая база $\mathfrak a$ или $\mathfrak t$ является также $\mathfrak X$-модульной базой $\overline{\mathfrak a}$ или $\overline{\mathfrak t}$; поэтому для построения дискриминанта $\overline{\mathfrak t}$ относительно $(\overline{\mathfrak a})$ можно положить в основу базу $\mathfrak a$ и $\mathfrak t$.

В частности, дискриминантный идеал $\overline{\mathfrak R}$ относительно главного класса $(\overline{\mathfrak R})$ является идеалом расширения того дискриминантного идеала $\mathfrak R$, который взят относительно главного класса $(\mathfrak R)$; этот так определенный идеал в $\mathfrak Z$, соответственно в $\mathfrak X$, будем кратко называть дискриминантным идеалом кольца.
% END UNIT BODY
% END PRODUCER UNIT 147
\endgroup


\clearpage
\begingroup
\setcounter{footnote}{0}
% BEGIN PRODUCER UNIT 148
% Source: C:/Users/Floris/Documents/interlanguage/03_projects/noether/02_slavic_working_corpus/translations/paper31/section06_opening/russian/v001/Noether_Paper31_Section06_Opening_Russian_v001.tex
% Identity: 1978 B / SHA-256 0881F51350A494B5CAFDC7C2F8DCB16EF4FFD45072808C96D204974F0E73CBEE
\setlength{\parindent}{1.25em}
\setlength{\parskip}{0pt}
\makeatletter
\renewcommand\footnotesize{\@setfontsize\footnotesize{7.7}{8.7}\selectfont}
\makeatother
\emergencystretch=100em
\allowdisplaybreaks
\sloppy
\hbadness=10000
\vbadness=10000
\hfuzz=2pt
\vfuzz=10pt
\raggedbottom
% BEGIN UNIT BODY
% Unit: Paper31 Section06 opening v001. Direct Russian translation from the German final-audited source slice, source lines 334--336 / segments P31-S0072--P31-S0073, expanded against printed scan/OCR witness page 96. Footnote counter continues after Paper31 Section05 entry 3.
\setcounter{footnote}{26}

\subsection*{\S~6. Дискриминантный критерий полной приводимости первого рода}

С этого момента снова предполагается, что $\mathfrak R$ является кольцом расширения поля $P$, а $\overline P$ является полем расширения $P$. Необходимым и достаточным условием полной приводимости первого рода окажется ненулевость дискриминанта. По \S~5, 3 дискриминантный идеал $\mathfrak R$ можно определить посредством кольца расширения $\overline{\mathfrak R}=\mathfrak R[\overline P]$; причем, по \S~5, 2, достаточно ограничиться первичными кольцами, из которых $\overline{\mathfrak R}$ аддитивно составлено. Поэтому прежде всего следует рассмотреть дискриминантные идеалы первичных колец.
% END UNIT BODY
% END PRODUCER UNIT 148
\endgroup


\clearpage
\begingroup
\setcounter{footnote}{0}
% BEGIN PRODUCER UNIT 149
% Source: C:/Users/Floris/Documents/interlanguage/03_projects/noether/02_slavic_working_corpus/translations/paper31/section06_entry01/russian/v001/Noether_Paper31_Section06_Entry01_Russian_v001.tex
% Identity: 4458 B / SHA-256 0D5E39E725B823667E374F09FCB51733220081905DF1BE46CA1AEA232455F874
\setlength{\parindent}{1.25em}
\setlength{\parskip}{0pt}
\makeatletter
\renewcommand\footnotesize{\@setfontsize\footnotesize{7.7}{8.7}\selectfont}
\makeatother
\emergencystretch=100em
\allowdisplaybreaks
\sloppy
\hbadness=10000
\vbadness=10000
\hfuzz=2pt
\vfuzz=10pt
\raggedbottom
% BEGIN UNIT BODY
% Unit: Paper31 Section06 no. 1 v001. Direct Russian translation from the German final-audited source slice, source lines 338--344 / segments P31-S0074--P31-S0075, expanded against printed scan/OCR witness pages 96--97. Footnote counter continues after Paper31 Section06 opening.
\setcounter{footnote}{26}

\textbf{1. \emph{Специальные $P$-модульные базисы.}} Пусть $\mathfrak a$ и $\mathfrak t$ являются идеалами из $\mathfrak R$, и пусть $\mathfrak a$ делится на $\mathfrak t$. Тогда всякий базис $a_1,\ldots,a_s$ идеала $\mathfrak a$ можно, добавив элементы $\tau_{s+1},\ldots,\tau_t$, дополнить до линейно независимого $P$-модульного базиса идеала $\mathfrak t$.\footnote{Матричное представление произвольного элемента $\gamma$, задаваемое этим базисом, тогда имеет вид $\begin{pmatrix}C_{\mathfrak a}&0\\ C_1&C_2\end{pmatrix}$; и обратно, из такого матричного представления, заданного через $\mathfrak t$, следует существование идеала $\mathfrak a$, кратного $\mathfrak t$. Так как при полной приводимости неразложимые компоненты не имеют идеала, отличного от нулевого и единичного, в классе представления этих компонентных идеалов такое представление не может появиться. Если добавить данную в \S~5, 1 связь между прямой суммой классов, то получается обычное в литературе определение полной приводимости.} Это является непосредственным следствием предполагаемого конечного ранга относительно поля $P$.

Пусть теперь $\mathfrak R$ предполагается первичным кольцом; пусть $\mathfrak p$ -- простой идеал, ассоциированный с нулевым идеалом, а $\rho$ -- его показатель, то есть $\mathfrak p^\rho=(0)$ и $\mathfrak p^{\rho-1}\ne(0)$. Тогда специальный $P$-модульный базис $\mathfrak R$ получается так: последовательно дополняют базис $\mathfrak p^{\rho-1}$ до базиса $\mathfrak p^{\rho-2}$ и, вообще, базис $\mathfrak p^i$ до базиса $\mathfrak p^{i-1}$, причем $\mathfrak R$ считается как $\mathfrak p^0$. С помощью такого базиса можно показать: если $\mathfrak R$ является первичным кольцом, то след, взятый относительно главного класса (\S~4, 2), всякого элемента, делящегося на соответствующий простой идеал $\mathfrak p$, равен нулю. Действительно, пусть
\[
\pi_{0,1},\ldots,\pi_{0,i_0};\quad
\pi_{1,1},\ldots,\pi_{1,i_1};\quad \ldots;\quad
\pi_{\rho-1,1},\ldots,\pi_{\rho-1,i_{\rho-1}}
\]
есть такой специальный $P$-модульный базис с
\[
\pi_{\lambda,j}\equiv0\pmod{\mathfrak p^\lambda},\qquad
\pi_{\lambda,j}\not\equiv0\pmod{\mathfrak p^{\lambda+1}}.
\]
Из $\pi\equiv0\pmod{\mathfrak p}$ следует $\pi\pi_{\lambda,j}\equiv0\pmod{\mathfrak p^{\lambda+1}}$; следовательно, матрица, соответствующая $\pi$ по этому базису, имеет на диагонали только нули, и след $\pi$ равен нулю.
% END UNIT BODY
% END PRODUCER UNIT 149
\endgroup


\clearpage
\begingroup
\setcounter{footnote}{0}
% BEGIN PRODUCER UNIT 150
% Source: C:/Users/Floris/Documents/interlanguage/03_projects/noether/02_slavic_working_corpus/translations/paper31/section06_entry02/russian/v001/Noether_Paper31_Section06_Entry02_Russian_v001.tex
% Identity: 1587 B / SHA-256 D6F9CC8939E1DD6D134D52F22FC30BE2F7E0B2CEA4EB681993470EF17C63940E
\setlength{\parindent}{1.25em}
\setlength{\parskip}{0pt}
\makeatletter
\renewcommand\footnotesize{\@setfontsize\footnotesize{7.7}{8.7}\selectfont}
\makeatother
\emergencystretch=100em
\allowdisplaybreaks
\sloppy
\hbadness=10000
\vbadness=10000
\hfuzz=2pt
\vfuzz=10pt
\raggedbottom
% BEGIN UNIT BODY
% Unit: Paper31 Section06 no. 2 v001. Direct Russian translation from the German final-audited source slice, source lines 346--350 / segment P31-S0076, checked against printed scan/OCR witness page 97. Footnote counter continues after Paper31 Section06 no. 1.
\setcounter{footnote}{27}

\textbf{2. \emph{Дискриминантные идеалы первичных колец с алгебраически замкнутым подкольцом $P$.}} Если нулевой идеал кольца $\mathfrak R$ является простым идеалом, а $P$ алгебраически замкнуто, то $\mathfrak R$ имеет ранг один, следовательно совпадает с $P$ (\S~3, 3); единичный элемент $e$ можно взять в качестве модульного базиса. Поэтому получаем
\[
        S_{(\mathfrak R)}(e^2)=S_{(\mathfrak R)}(e)=e,
\]
и дискриминантный идеал равен единичному идеалу.
% END UNIT BODY
% END PRODUCER UNIT 150
\endgroup


\clearpage
\begingroup
\setcounter{footnote}{0}
% BEGIN PRODUCER UNIT 151
% Source: C:/Users/Floris/Documents/interlanguage/03_projects/noether/02_slavic_working_corpus/translations/paper31/section06_entry03/russian/v001/Noether_Paper31_Section06_Entry03_Russian_v001.tex
% Identity: 2422 B / SHA-256 C617621B5B00EADB561ADC114A1E180BF9F197CF3CB3142ACE44EB30F7AAE538
\setlength{\parindent}{1.25em}
\setlength{\parskip}{0pt}
\makeatletter
\renewcommand\footnotesize{\@setfontsize\footnotesize{7.7}{8.7}\selectfont}
\makeatother
\emergencystretch=100em
\allowdisplaybreaks
\sloppy
\hbadness=10000
\vbadness=10000
\hfuzz=2pt
\vfuzz=10pt
\raggedbottom
% BEGIN UNIT BODY
% Unit: Paper31 Section06 no. 2 proper-primary paragraph v001. Direct Russian translation from the German final-audited source slice, source line 352 / segment P31-S0077, expanded against printed scan/OCR witness page 97. Footnote counter continues after Paper31 Section06 no. 1.
\setcounter{footnote}{27}

Далее можно показать, что дискриминантный идеал собственно первичного кольца с алгебраически замкнутым подкольцом $P$ равен нулевому идеалу. Если для $\mathfrak R$ положить в основу специальный модульный базис, данный в пункте 1, или лишь дополнить произвольный базис $\mathfrak p$ до базиса $\mathfrak R$, то $i_0$ равно единице, поскольку факторкольцо $\mathfrak R/\mathfrak p$ становится кольцом ранга один; в качестве $\pi_{0,1}$ можно, например, взять единичный элемент $e$. Тогда все базисные произведения, отличные от $e^2$, делятся на $\mathfrak p$, и их след, взятый относительно главного класса, равен нулю. Следовательно, равен нулю и дискриминант как определитель порядка не менее двух, поскольку $\mathfrak R$ предполагалось собственно первичным и потому должно иметь по меньшей мере один базисный элемент, отличный от $e$; в этом определителе все элементы, отличные от $S_{(\mathfrak R)}(e^2)$, равны нулю.
% END UNIT BODY
% END PRODUCER UNIT 151
\endgroup


\clearpage
\begingroup
\setcounter{footnote}{0}
% BEGIN PRODUCER UNIT 152
% Source: C:/Users/Floris/Documents/interlanguage/03_projects/noether/02_slavic_working_corpus/translations/paper31/section06_entry04/russian/v001/Noether_Paper31_Section06_Entry04_Russian_v001.tex
% Identity: 1399 B / SHA-256 5649D03DDF828345B36B0DF92208C109B473D1071B39A15B4AADC64A3C2D603C
\setlength{\parindent}{1.25em}
\setlength{\parskip}{0pt}
\makeatletter
\renewcommand\footnotesize{\@setfontsize\footnotesize{7.7}{8.7}\selectfont}
\makeatother
\emergencystretch=100em
\allowdisplaybreaks
\sloppy
\hbadness=10000
\vbadness=10000
\hfuzz=2pt
\vfuzz=10pt
\raggedbottom
% BEGIN UNIT BODY
% Unit: Paper31 Section06 no. 3 theorem v001. Direct Russian translation from the German final-audited source slice, source line 354 / segment P31-S0078, checked against printed scan/OCR witness page 97. Footnote counter continues after Paper31 Section06 no. 1.
\setcounter{footnote}{27}

\textbf{3. Теорема.} \emph{Кольцо $\mathfrak R$ конечного ранга относительно подполя $P$ является вполне приводимым первого рода тогда и только тогда, когда его дискриминантный идеал является единичным идеалом в $P$; другими словами, когда дискриминант, определенный с точностью до квадрата единицы, отличен от нуля.}
% END UNIT BODY
% END PRODUCER UNIT 152
\endgroup


\clearpage
\begingroup
\setcounter{footnote}{0}
% BEGIN PRODUCER UNIT 153
% Source: C:/Users/Floris/Documents/interlanguage/03_projects/noether/02_slavic_working_corpus/translations/paper31/section06_entry05/russian/v001/Noether_Paper31_Section06_Entry05_Russian_v001.tex
% Identity: 2860 B / SHA-256 B6CBABC2A20395510E4D8B864F7F29DD92C7167C81B8BA1C91F8AA3394E25CEF
\setlength{\parindent}{1.25em}
\setlength{\parskip}{0pt}
\makeatletter
\renewcommand\footnotesize{\@setfontsize\footnotesize{7.7}{8.7}\selectfont}
\makeatother
\emergencystretch=100em
\allowdisplaybreaks
\sloppy
\hbadness=10000
\vbadness=10000
\hfuzz=2pt
\vfuzz=10pt
\raggedbottom
% BEGIN UNIT BODY
% Unit: Paper31 Section06 no. 3 proof v001. Direct Russian translation from the German final-audited source slice, source line 356 / segment P31-S0079, expanded against printed scan/OCR witness pages 97--98. Footnote counter continues after Paper31 Section06 no. 1.
\setcounter{footnote}{27}

Полная приводимость первого рода по \S~3, 2 равносильна тому, что кольцо расширения $\overline{\mathfrak R}=\mathfrak R[\overline P]$ является вполне приводимым, то есть его неразложимые компоненты становятся полями, изоморфными $\overline P$. По \S~5, 3 дискриминантный идеал $\mathfrak R$, а одновременно и $\overline{\mathfrak R}$, является нулевым или единичным идеалом; а по \S~5, 2 дискриминантный идеал $\overline{\mathfrak R}$ является произведением дискриминантных идеалов компонент, взятых относительно этих компонент. Эти отдельные дискриминантные идеалы возникают из рассмотренных в пункте 2, где компоненты выступают как кольца сами по себе, а не как идеалы $\overline{\mathfrak R}$, если каждый раз поле $P_i=\overline P e_i$, изоморфное $\overline P$, заменить самим $\overline P$. Отсюда следует, что при полной приводимости $\overline{\mathfrak R}$ дискриминантный идеал как произведение единичных идеалов сам является единичным идеалом; напротив, если $\overline{\mathfrak R}$ не является вполне приводимым, то по меньшей мере для одной компоненты дискриминантный идеал равен нулевому идеалу, откуда и для дискриминантного идеала $\overline{\mathfrak R}$ получается нулевой идеал. Тем самым теорема доказана.
% END UNIT BODY
% END PRODUCER UNIT 153
\endgroup


\clearpage
\begingroup
\setcounter{footnote}{0}
% BEGIN PRODUCER UNIT 154
% Source: C:/Users/Floris/Documents/interlanguage/03_projects/noether/02_slavic_working_corpus/translations/paper31/section06_entry06/russian/v001/Noether_Paper31_Section06_Entry06_Russian_v001.tex
% Identity: 1955 B / SHA-256 1BB49AC02B7A8B2B0D469BC1A6BF2C05AB663AA4146EF00ACCBC96080C5A745D
\setlength{\parindent}{1.25em}
\setlength{\parskip}{0pt}
\makeatletter
\renewcommand\footnotesize{\@setfontsize\footnotesize{7.7}{8.7}\selectfont}
\makeatother
\emergencystretch=100em
\allowdisplaybreaks
\sloppy
\hbadness=10000
\vbadness=10000
\hfuzz=2pt
\vfuzz=10pt
\raggedbottom
% BEGIN UNIT BODY
% Unit: Paper31 Section06 no. 4 opening v001. Direct Russian translation from the German final-audited source slice, source lines 358--362 / segment P31-S0080, expanded against printed scan/OCR witness page 98. Footnote counter continues after Paper31 Section06 no. 1.
\setcounter{footnote}{27}

4. \emph{Представление следа, нормы и дискриминанта посредством сопряженных элементов в случае полной приводимости первого рода.} Пусть по \S~3, 3 при полной приводимости первого рода
\[
        \mathfrak R[Q]=Qe_1+\cdots+Qe_n
\]
имеется представление в виде прямой суммы колец ранга один, то есть полей, изоморфных $Q$, и пусть $Q$ является наименьшим полем расширения, в котором такое представление возможно. Пусть $n$ компонент кольца $\mathfrak R$ заданы как $K_i e_i$, где каждое $K_i$ является подполем $Q$ и $\mathfrak R$ гомоморфно отображается в $K_i$; нужно выразить след, норму и дискриминант (относительно главного класса) через элементы из $K_i$.
% END UNIT BODY
% END PRODUCER UNIT 154
\endgroup


\clearpage
\begingroup
\setcounter{footnote}{0}
% BEGIN PRODUCER UNIT 155
% Source: C:/Users/Floris/Documents/interlanguage/03_projects/noether/02_slavic_working_corpus/translations/paper31/section06_entry07/russian/v001/Noether_Paper31_Section06_Entry07_Russian_v001.tex
% Identity: 3974 B / SHA-256 25884F9DB25F56A01ECDA707DF548E9545FB610783FA5724F536A2F2C1C20DC3
\setlength{\parindent}{1.25em}
\setlength{\parskip}{0pt}
\makeatletter
\renewcommand\footnotesize{\@setfontsize\footnotesize{7.7}{8.7}\selectfont}
\makeatother
\emergencystretch=100em
\allowdisplaybreaks
\sloppy
\hbadness=10000
\vbadness=10000
\hfuzz=2pt
\vfuzz=10pt
\raggedbottom
% BEGIN UNIT BODY
% Unit: Paper31 Section06 no. 4 formulas and field case v001. Direct Russian translation from the German final-audited source slice, source lines 364--383 / segment P31-S0081, expanded against printed scan/OCR witness page 98. Footnote counter continues after Paper31 Section06 no. 1.
\setcounter{footnote}{27}

Если взять $e_1,\ldots,e_n$ в качестве модульной базы, то каждому элементу
\[
        \gamma=c_1e_1+\cdots+c_ne_n
\]
ставится в соответствие диагональная матрица с диагональю $c_1,\ldots,c_n$; при этом, если $\gamma$ принадлежит $\mathfrak R$, то $c_i$ лежит в $K_i$. Следовательно,
\[
        S(\gamma)=c_1+\cdots+c_n,
        \qquad
        N(\gamma)=c_1\cdots c_n,
\]
где под $S(\gamma)$ понимается след, а под $N(\gamma)$ норма относительно главного класса. Пусть далее $\alpha_1,\ldots,\alpha_n$ -- модульная база $\mathfrak R$, состоящая из элементов $\mathfrak R$; тогда дискриминант $|S(\alpha_i\alpha_j)|$ нужно выразить через элементы из всех $K_k$. Если
\[
        \alpha_j=\alpha_j^{(1)}e_1+\cdots+\alpha_j^{(n)}e_n,
\]
то получается
\[
        |S(\alpha_i\alpha_j)|
        =
        \left|\sum_{k=1}^n \alpha_i^{(k)}\alpha_j^{(k)}\right|
        =
        \left|\alpha_i^{(k)}\right|^2.
\]
То есть дискриминант представляется как квадрат определителя сопряженных базисных элементов. В частности, если $\mathfrak R$ само является полем, то есть полем расширения первого рода подполя $P$, то гомоморфизмы из $\mathfrak R$ в $K_i$ переходят в изоморфизмы; при этом элементы из $P$ соответствуют самим себе, поскольку компонента $P$ задана как $P e_i$. Поэтому $K_i$, так как все они лежат в общем объемлющем поле $Q$, являются сопряженными относительно $P$ полями; причем эти $n$ изоморфизмов попарно различны, как показывает приведенное выше детерминантное представление ненулевого дискриминанта. Если $\mathfrak R$ является полем, то представление в виде прямой суммы в $K_1,\ldots,K_n$ дает $n$ сопряженных полей расширения $n$-й степени над $P$, лежащих в поле расширения Галуа и изоморфных $\mathfrak R$. Тогда представление следа, нормы и дискриминанта переходит в известное представление посредством сопряженных элементов. Следует еще раз подчеркнуть, что $\mathfrak R$ предполагается только как поле, эквивалентное $K_i$, а не как поле, само уже лежащее в поле расширения Галуа $Q$.
% END UNIT BODY
% END PRODUCER UNIT 155
\endgroup


\clearpage
\begingroup
\setcounter{footnote}{0}
% BEGIN PRODUCER UNIT 156
% Source: C:/Users/Floris/Documents/interlanguage/03_projects/noether/02_slavic_working_corpus/translations/paper31/section07_entry01/russian/v001/Noether_Paper31_Section07_Entry01_Russian_v001.tex
% Identity: 4289 B / SHA-256 C009B8FBFE45F0F2F433B55C06B6FC5E961562038BE9433B7A46C8D21B80B3E2
\newcommand{\mo}{\mathfrak{o}}
\newcommand{\mO}{\mathfrak{O}}
\newcommand{\mT}{\mathfrak{T}}
\setlength{\parindent}{1.25em}
\setlength{\parskip}{0pt}
\makeatletter
\renewcommand\footnotesize{\@setfontsize\footnotesize{7.7}{8.7}\selectfont}
\makeatother
\emergencystretch=100em
\allowdisplaybreaks
\sloppy
\hbadness=10000
\vbadness=10000
\hfuzz=2pt
\vfuzz=10pt
\raggedbottom
% BEGIN UNIT BODY
% Unit: Paper31 Section07 no. 1 opening v001. Direct Russian translation from the German final-audited source slice, source lines 385--389 / segments P31-S0082--P31-S0084, expanded against printed scan/OCR witness pages 98--99.
\setcounter{footnote}{28}

\subsection*{§7. Теорема о дискриминанте для порядков алгебраического числового поля или функционального поля}

\noindent\textbf{1.}\quad Рассматриваемые числовые и функциональные поля -- причем для алгебраических функций от более чем одной неопределенной и для целых алгебраических функций речь идет о расширении функциональной области -- все подчиняются следующему типу поля с фиксированным понятием целых элементов. Пусть $\mo$ -- кольцо главных идеалов без делителей нуля и с единицей (кольцо целых рациональных чисел, область многочленов от одной неопределенной с коэффициентами из поля, функциональная область), а $\Omega$ -- его поле частных. Пусть $K$ -- конечное расширение первого рода поля $\Omega$, а $\mO$ -- кольцо всех элементов из $K$, целых относительно $\mo$. Всякое подкольцо $\mO$, содержащее $\mo$, называется порядком; само $\mO$ называется главным порядком.

Всякий $\mo$-модуль в $\mO$, в частности все идеалы порядка и сам порядок, имеет линейно независимую модульную базу относительно $\mo$. Если $K$ имеет степень $n$ над $\Omega$, то достаточно ограничиться порядками ранга $n$ относительно $\mo$, поскольку всякий порядок меньшего ранга посредством образования частных порождает промежуточное поле между $K$ и $\Omega$ именно этой степени, и для него все предположения сохраняются. Следовательно, всякий порядок $\mT$ является кольцом, для которого выполняются предположения §4 и §5, а именно кольцом без делителей нуля; поэтому дискриминант $D_\mT$ порядка однозначно определен с точностью до множителя, являющегося квадратом единицы из $\mo$. Этот дискриминант одновременно, с точностью до множителя-единицы из $\Omega$, совпадает с дискриминантом поля $K$, рассматриваемого как $\Omega$-модуль, и потому отличен от нуля, так как $K$ предполагалось расширением первого рода (§6, 3). Дискриминант одновременно допускает данное в §6, 4 представление через квадрат определителя сопряженных базисных элементов.
% END UNIT BODY
% END PRODUCER UNIT 156
\endgroup


\clearpage
\begingroup
\setcounter{footnote}{0}
% BEGIN PRODUCER UNIT 157
% Source: C:/Users/Floris/Documents/interlanguage/03_projects/noether/02_slavic_working_corpus/translations/paper31/section07_entry02/russian/v001/Noether_Paper31_Section07_Entry02_Russian_v001.tex
% Identity: 1527 B / SHA-256 A3995FFDB8BB3FECBE88CDFC971391BA34D539561F23CC70665ACB2CAA7834AD
\newcommand{\mo}{\mathfrak{o}}
\newcommand{\mO}{\mathfrak{O}}
\newcommand{\mT}{\mathfrak{T}}
\setlength{\parindent}{1.25em}
\setlength{\parskip}{0pt}
\makeatletter
\renewcommand\footnotesize{\@setfontsize\footnotesize{7.7}{8.7}\selectfont}
\makeatother
\emergencystretch=100em
\allowdisplaybreaks
\sloppy
\hbadness=10000
\vbadness=10000
\hfuzz=2pt
\vfuzz=10pt
\raggedbottom
% BEGIN UNIT BODY
% Unit: Paper31 Section07 no. 1 ideal-theory theorem v001. Direct Russian translation from the German final-audited source slice, source line 391 / segment P31-S0085, checked against printed scan/OCR witness page 99 with footnote 5.
\setcounter{footnote}{28}

\noindent Теория идеалов в $\mT$ задается следующей теоремой:\footnote{Idealtheorie, \S 7, 3.} в $\mT$ каждый идеал, отличный от нулевого идеала, единственным образом представляется как произведение конечного числа попарно взаимно простых первичных идеалов, чьи соответствующие простые идеалы не имеют собственного делителя, отличного от единичного идеала.
% END UNIT BODY
% END PRODUCER UNIT 157
\endgroup


\clearpage
\begingroup
\setcounter{footnote}{0}
% BEGIN PRODUCER UNIT 158
% Source: C:/Users/Floris/Documents/interlanguage/03_projects/noether/02_slavic_working_corpus/translations/paper31/section07_entry03/russian/v001/Noether_Paper31_Section07_Entry03_Russian_v001.tex
% Identity: 3232 B / SHA-256 AF7A339CB5E9B9F6268BCAB74DA4571CB58FD128E655EC1122EAE3F415BAB0E7
\newcommand{\mo}{\mathfrak{o}}
\newcommand{\mO}{\mathfrak{O}}
\newcommand{\mT}{\mathfrak{T}}
\newcommand{\mR}{\mathfrak{R}}
\setlength{\parindent}{1.25em}
\setlength{\parskip}{0pt}
\makeatletter
\renewcommand\footnotesize{\@setfontsize\footnotesize{7.7}{8.7}\selectfont}
\makeatother
\emergencystretch=100em
\allowdisplaybreaks
\sloppy
\hbadness=10000
\vbadness=10000
\hfuzz=2pt
\vfuzz=10pt
\raggedbottom
% BEGIN UNIT BODY
% Unit: Paper31 Section07 no. 2 passage to finite-rank residue rings v001. Direct Russian translation from the German final-audited source slice, lines 393--403 / segments P31-S0086--P31-S0087, checked against printed scan page 100 / PDF page 19 / Claude handoff OCR line 369.
\setcounter{footnote}{28}

\noindent 2. \emph{Переход от порядка к кольцам конечного ранга относительно поля.}
Пусть $p$ -- простой элемент из $\mo$, то есть главный идеал $\mo p$, выведенный из $p$ в $\mo$, является простым идеалом из $\mo$, отличным от нулевого и единичного идеалов; соответственно обозначим через $\mT p$ главный идеал, выведенный из $p$ в $\mT$. Кольцо классов вычетов
\[
        \mR=\mT/\mT p
\]
становится кольцом ранга $n$ относительно подполя $P$, изоморфного полю $\mo/\mo p$; порядок $\mT$ гомоморфно отображается на $\mT/\mT p$. Это гомоморфное отображение является обычным переходом к кольцу классов вычетов.

То, что $\mR$ содержит подполе, изоморфное $\mo/\mo p$, следует из второй теоремы об изоморфизме: подкольцо $(\mo,\mT p)/\mT p$ изоморфно $\mo/[\mo,\mT p]=\mo/\mo p$. Пусть $\alpha_1,\ldots,\alpha_n$ является линейно независимым модульным базисом $\mT$ относительно $\mo$, и пусть $(\alpha_i)$ -- классы вычетов по $\mT p$, определенные элементами $\alpha_i$. Всякая линейная зависимость между $(\alpha_i)$ относительно $P$ возникла бы из конгруенции
\[
        c_1\alpha_1+\cdots+c_n\alpha_n\equiv0\pmod{\mT p},
\]
а значит, из равенства $c_1\alpha_1+\cdots+c_n\alpha_n=py$ с $y\in\mT$. Из представимости $y$ через базис $\alpha_i$ и единственности такого представления следует, что каждое $c_i$ делится на $p$; тем самым элемент $(c_i)$ из $P$ обращается в нуль. Следовательно, $(\alpha_i)$ образуют базис $\mR$ над $P$.
% END UNIT BODY
% END PRODUCER UNIT 158
\endgroup


\clearpage
\begingroup
\setcounter{footnote}{0}
% BEGIN PRODUCER UNIT 159
% Source: C:/Users/Floris/Documents/interlanguage/03_projects/noether/02_slavic_working_corpus/translations/paper31/section07_entry04/russian/v001/Noether_Paper31_Section07_Entry04_Russian_v001.tex
% Identity: 2911 B / SHA-256 D0F43DC8A1DC4BC6E59957104F10625416B0EE4255140F2C67A71D7093966678
\newcommand{\mo}{\mathfrak{o}}
\newcommand{\mO}{\mathfrak{O}}
\newcommand{\mT}{\mathfrak{T}}
\newcommand{\mR}{\mathfrak{R}}
\newcommand{\mq}{\mathfrak{q}}
\setlength{\parindent}{1.25em}
\setlength{\parskip}{0pt}
\makeatletter
\renewcommand\footnotesize{\@setfontsize\footnotesize{7.7}{8.7}\selectfont}
\makeatother
\emergencystretch=100em
\allowdisplaybreaks
\sloppy
\hbadness=10000
\vbadness=10000
\hfuzz=2pt
\vfuzz=10pt
\raggedbottom
% BEGIN UNIT BODY
% Unit: Paper31 Section07 no. 2 discriminant and ideal-decomposition transfer v001. Direct Russian translation from the German final-audited source slice, lines 405--413 / segment P31-S0088, scan-expanded against printed scan page 100 / PDF page 19 / Claude handoff OCR lines 371--373.
\setcounter{footnote}{28}

\noindent При гомоморфизме дискриминант $D_\mT$ порядка $\mT$ переходит в дискриминант $\mR$, а разложение идеала $\mT p$ -- в разложение нулевого идеала $\mR$. Действительно, в силу только что доказанной линейной независимости $(\alpha_i)$ дискриминант $\mR$ можно построить по этому базису как $\left|S((\alpha_i)(\alpha_k))\right|$; следовательно, при гомоморфизме он возникает из дискриминанта $\left|S(\alpha_i\alpha_k)\right|$ порядка $\mT$, построенного по базису $\alpha_i$. При определении следа здесь положено в основу матричное представление, а не представление через сопряженные элементы, справедливое для $\mT$, но не обязательно для $\mR$; поэтому след, а вместе с ним и дискриминант, возникает через кольцевые операции из элементов кольца.

Если
\[
        \mT p=\mq_1\cdots\mq_s=[\mq_1,\ldots,\mq_s]
\]
есть представление $\mT p$ в $\mT$, то при гомоморфизме оно переходит в
\[
        (0)=[(\mq_1),\ldots,(\mq_s)]
\]
в $\mR$. По первой теореме об изоморфизме $\mR/(\mq_i)$ изоморфно $\mT/\mq_i$; следовательно, $\mq_i$ и $(\mq_i)$ одновременно являются собственно первичными идеалами или простыми идеалами, и в силу гомоморфизма $(\mq_i)$ попарно взаимно просты.
% END UNIT BODY
% END PRODUCER UNIT 159
\endgroup


\clearpage
\begingroup
\setcounter{footnote}{0}
% BEGIN PRODUCER UNIT 160
% Source: C:/Users/Floris/Documents/interlanguage/03_projects/noether/02_slavic_working_corpus/translations/paper31/section07_entry05/russian/v001/Noether_Paper31_Section07_Entry05_Russian_v001.tex
% Identity: 4493 B / SHA-256 FDCD69D04ED8D4630224CAC5F654E29CD09E03756D34E6FF79D1C26126D18DAE
\newcommand{\mo}{\mathfrak{o}}
\newcommand{\mO}{\mathfrak{O}}
\newcommand{\mT}{\mathfrak{T}}
\newcommand{\mR}{\mathfrak{R}}
\newcommand{\mq}{\mathfrak{q}}
\setlength{\parindent}{1.25em}
\setlength{\parskip}{0pt}
\makeatletter
\renewcommand\footnotesize{\@setfontsize\footnotesize{7.7}{8.7}\selectfont}
\makeatother
\emergencystretch=100em
\allowdisplaybreaks
\sloppy
\hbadness=10000
\vbadness=10000
\hfuzz=2pt
\vfuzz=10pt
\raggedbottom
% BEGIN UNIT BODY
% Unit: Paper31 Section07 no. 3 discriminant theorem statement plus scan-expanded example footnote v001. Direct Russian translation from the German final-audited source slice, line 415 / segment P31-S0089, with the printed-scan footnote marker and expanded footnote witness on printed scan pages 100--101 / PDF pages 19--20 / Claude handoff OCR lines 375--403.
\setcounter{footnote}{28}

\noindent 3. \textbf{Теорема о дискриминанте.} \emph{Простой элемент $p$ из $\mo$ входит в дискриминант $D_\mT$ порядка $\mT$ тогда и только тогда, когда в разложении главного идеала $\mT p$ на попарно взаимно простые первичные компоненты из $\mT$ появляется по крайней мере одна собственно первичная компонента или по крайней мере один простой идеал второго рода.}\footnote{Пример появления простых идеалов второго рода таков: пусть $\mo$ будет функциональной областью, выведенной из всех многочленов по $x$ с целыми рациональными коэффициентами, $\Omega$ -- ее полем частных, а $K$ -- полем расширения, возникающим присоединением $\sqrt{x}$. Рассматривается порядок $\mT$ с модульным базисом $1,\sqrt{x}$ относительно $\mo$; как показывает подстановка $x-y^2$, речь идет о главном порядке. Дискриминант равен $\left|\begin{smallmatrix}1&\sqrt{x}\\1&-\sqrt{x}\end{smallmatrix}\right|^2=2^2x$. Главный идеал $\mT2$ остается простым идеалом в $\mT$, но второго рода. Действительно, порядок $\mT$ в силу модульного базиса $1,\sqrt{x}$ изоморфен кольцу классов вычетов $F[t]/(t^2-x)$, где $t$ означает новую неопределенную; тем самым $\mT/\mT2=\mR$ изоморфно $P[t]/(t^2-x)$, где $P$ означает поле всех рациональных функций от $x$ и от неопределенной $u$, встречающейся в функциональной области, с коэффициентами по модулю $2$. Так как $t^2-x$ является неприводимым многочленом в $P[t]$, причем неприводимым многочленом второго рода, $\mR$ становится полем, а именно полем расширения второго рода над $P$; значит, $\mT2$ становится простым идеалом второго рода. Напротив, $\mT x=(\mT\sqrt{x})^2$ является собственно первичным. То, что $\mR$ порождается над $P$ присоединением одного элемента, вообще говоря, неверно: достаточно заменить одну неопределенную $x$ на $x$ и $y$, чтобы в порядке, выведенном из $1,\sqrt{x},\sqrt{y},\sqrt{xy}$, который снова является главным порядком, распознать главный идеал $\mT2$ как простой идеал второго рода, для которого $\mR$ порождается над $P$ только после присоединения двух элементов; ср. о примерах Ostrowski, a. a. O.}
% END UNIT BODY
% END PRODUCER UNIT 160
\endgroup


\clearpage
\begingroup
\setcounter{footnote}{0}
% BEGIN PRODUCER UNIT 161
% Source: C:/Users/Floris/Documents/interlanguage/03_projects/noether/02_slavic_working_corpus/translations/paper31/section07_entry06/russian/v001/Noether_Paper31_Section07_Entry06_Russian_v001.tex
% Identity: 2561 B / SHA-256 23D680D6B7B210ED5E0826620BA02489E91954EAEE1048DD04BBCBC5D0546068
\newcommand{\mo}{\mathfrak{o}}
\newcommand{\mO}{\mathfrak{O}}
\newcommand{\mT}{\mathfrak{T}}
\newcommand{\mR}{\mathfrak{R}}
\newcommand{\mq}{\mathfrak{q}}
\setlength{\parindent}{1.25em}
\setlength{\parskip}{0pt}
\makeatletter
\renewcommand\footnotesize{\@setfontsize\footnotesize{7.7}{8.7}\selectfont}
\makeatother
\emergencystretch=100em
\allowdisplaybreaks
\sloppy
\hbadness=10000
\vbadness=10000
\hfuzz=2pt
\vfuzz=10pt
\raggedbottom
% BEGIN UNIT BODY
% Unit: Paper31 Section07 no. 3 perfect-residue-field consequence and proof v001. Direct Russian translation from the German final-audited source slice, lines 417 and 419 / segments P31-S0090--P31-S0091, with the RA34 line-419 compact footnote deliberately omitted because the printed-scan-expanded example footnote was already translated and attached to the theorem statement in Entry05.
\setcounter{footnote}{29}

\noindent Если поле классов вычетов $\mo/\mo p$ является совершенным полем, то каждый простой элемент $p$, входящий в $D_\mT$, и только такой элемент имеет по крайней мере одну собственно первичную компоненту. Если $\mo/\mo p$ совершенно и для $\mT$ выбран главный порядок $\mO$, то $p$ входит в дискриминант тогда и только тогда, когда $p$ делится по крайней мере на квадрат простого идеала из $\mO$.

Действительно, из-за соответствия между разложением идеала $\mT p$ в $\mT$ и разложением нулевого идеала в $\mR=\mT/\mT p$ появление собственно первичной компоненты или простого идеала второго рода в $\mT p$ означает, что $\mR$ не является вполне приводимым первого рода. А по §6, 3 это имеет место тогда и только тогда, когда дискриминант $\mR$ обращается в нуль; ввиду гомоморфизма из пункта 2 это как раз выражает делимость $D_\mT$ на $p$.
% END UNIT BODY
% END PRODUCER UNIT 161
\endgroup


\clearpage
\begingroup
\setcounter{footnote}{0}
% BEGIN PRODUCER UNIT 162
% Source: C:/Users/Floris/Documents/interlanguage/03_projects/noether/02_slavic_working_corpus/translations/paper31/section08_entry01/russian/v001/Noether_Paper31_Section08_Entry01_Russian_v001.tex
% Identity: 3649 B / SHA-256 F1C90ABCCFC438B525606DB15A2DFB4FA23AEA21556FF8053A0CA65D67CC02D4
\newcommand{\mo}{\mathfrak{o}}
\newcommand{\mO}{\mathfrak{O}}
\newcommand{\mT}{\mathfrak{T}}
\newcommand{\mR}{\mathfrak{R}}
\newcommand{\mq}{\mathfrak{q}}
\setlength{\parindent}{1.25em}
\setlength{\parskip}{0pt}
\makeatletter
\renewcommand\footnotesize{\@setfontsize\footnotesize{7.7}{8.7}\selectfont}
\makeatother
\emergencystretch=100em
\allowdisplaybreaks
\sloppy
\hbadness=10000
\vbadness=10000
\hfuzz=2pt
\vfuzz=10pt
\raggedbottom
% BEGIN UNIT BODY
% Unit: Paper31 Section08 opening and multiplication-ring setup v001. Direct Russian translation from the German final-audited source slice, lines 421--425 / segments P31-S0092--P31-S0094, with scan-visible footnotes 1--4 restored from printed page 101 / PDF page 20.
\setcounter{footnote}{29}

\subsection*{§8. Теорема о дискриминанте для порядков относительных полей}

\noindent Если вместо кольца главных идеалов $\mo$ уже взять в качестве исходного кольца главный порядок числового или функционального поля -- более общо, мультипликационное кольцо\footnote{Краткое название "мультипликационное кольцо" я беру из заметки W. Krull: \emph{Über Multiplikationsringe}. Heidelberger Berichte, 1925, 5. Abhandlung, Nr. 3. Впрочем, Krull называет встречающиеся здесь кольца без делителей нуля "регулярными мультипликационными кольцами".}, то есть кольцо, в котором действует теория идеалов главного порядка, -- то на место дискриминанта $D_\mT$ становится дискриминантный идеал порядка $\mT$ относительно $\mo$;\footnote{В литературе в случае "относительных числовых полей" выступает только дискриминантный идеал главного порядка, "относительный дискриминант". Ср. Hilbert, \emph{Zahlbericht}, §14, и Hecke, \emph{Theorie der algebraischen Zahlen}, §38.} теорема о дискриминанте полностью переносится посредством следующих соображений.

\noindent 1. Пусть $\mo$ -- мультипликационное кольцо, то есть кольцо без делителей нуля и с единицей, в котором каждый идеал, отличный от нулевого и единичного, однозначно представляется как произведение степеней простых идеалов, и где эти простые идеалы не имеют собственного делителя, отличного от единичного идеала.\footnote{Об аксиоматике мультипликационных колец ср. \emph{Idealtheorie}, например введение.} Пусть далее $\Omega,K,\mO,\mT$ определены как в §7, 1; тогда в $\mT$ действует указанная там теория идеалов.\footnote{\emph{Idealtheorie}, §3, 2 и §7, 3.}
% END UNIT BODY
% END PRODUCER UNIT 162
\endgroup


\clearpage
\begingroup
\setcounter{footnote}{0}
% BEGIN PRODUCER UNIT 163
% Source: C:/Users/Floris/Documents/interlanguage/03_projects/noether/02_slavic_working_corpus/translations/paper31/section08_entry02/russian/v001/Noether_Paper31_Section08_Entry02_Russian_v001.tex
% Identity: 1775 B / SHA-256 8E5B7755784F5C2380BEA47A72AF7C2A2AC148061DA9AB721966D6D9F2C3B884
\newcommand{\mo}{\mathfrak{o}}
\newcommand{\mO}{\mathfrak{O}}
\newcommand{\mT}{\mathfrak{T}}
\newcommand{\mR}{\mathfrak{R}}
\newcommand{\mq}{\mathfrak{q}}
\setlength{\parindent}{1.25em}
\setlength{\parskip}{0pt}
\makeatletter
\renewcommand\footnotesize{\@setfontsize\footnotesize{7.7}{8.7}\selectfont}
\makeatother
\emergencystretch=100em
\allowdisplaybreaks
\sloppy
\hbadness=10000
\vbadness=10000
\hfuzz=2pt
\vfuzz=10pt
\raggedbottom
% BEGIN UNIT BODY
% Unit: Paper31 Section08 Entry02 principal-ideal-ring consequence v001. Direct Russian translation from the German final-audited source slice, line 427 / segment P31-S0095, with printed scan page 102 / PDF page 21 logged as a variant witness.
\setcounter{footnote}{33}

\noindent В частности, если мультипликационное кольцо $\mo$ имеет только один простой идеал, отличный от нулевого и единичного идеалов, то $\mo$ является кольцом главных идеалов. В самом деле, если выбрать элемент $p$, главный идеал которого делится ровно на первую степень этого простого идеала, то $p$ порождает сам этот простой идеал; тогда из представления в виде произведения степеней следует, что и все остальные идеалы являются главными.
% END UNIT BODY
% END PRODUCER UNIT 163
\endgroup


\clearpage
\begingroup
\setcounter{footnote}{0}
% BEGIN PRODUCER UNIT 164
% Source: C:/Users/Floris/Documents/interlanguage/03_projects/noether/02_slavic_working_corpus/translations/paper31/section08_entry03/russian/v001/Noether_Paper31_Section08_Entry03_Russian_v001.tex
% Identity: 2210 B / SHA-256 58914ECFEAF2B93235A7E2E6F7B24BBCA15AEE671AB0D08B34552F07D31E129B
\newcommand{\mo}{\mathfrak{o}}
\newcommand{\mO}{\mathfrak{O}}
\newcommand{\mT}{\mathfrak{T}}
\newcommand{\mR}{\mathfrak{R}}
\newcommand{\mq}{\mathfrak{q}}
\setlength{\parindent}{1.25em}
\setlength{\parskip}{0pt}
\makeatletter
\renewcommand\footnotesize{\@setfontsize\footnotesize{7.7}{8.7}\selectfont}
\makeatother
\emergencystretch=100em
\allowdisplaybreaks
\sloppy
\hbadness=10000
\vbadness=10000
\hfuzz=2pt
\vfuzz=10pt
\raggedbottom
% BEGIN UNIT BODY
% Unit: Paper31 Section08 Entry03 trace and discriminant-ideal opening v001. Direct Russian translation from the German final-audited source slice, lines 429--433 / segment P31-S0096, with scan-visible §6,4 cross-reference and Idealtheorie footnote restored from printed page 102 / PDF page 21.
\setcounter{footnote}{33}

\noindent 2. \emph{След и дискриминантный идеал.} След каждого элемента $\gamma$ из $\mT$ определяется как след $\gamma$ в $K$, где $K$ рассматривается как $\Omega$-модуль. Любую систему из $n$ линейно независимых над $\Omega$ элементов из $K$ можно выбрать в качестве базиса для построения следа. Если $\alpha_1,\ldots,\alpha_n$ является такой системой, то, поскольку $K$ предполагается расширением первого рода,
\[
        |S(\alpha_i\alpha_j)|
\]
является ненулевым элементом из $\Omega$, а именно квадратом определителя сопряженных базисных элементов (§6, 4). Из целой замкнутости $\mo$ в $\Omega$\footnote{\emph{Idealtheorie}, §9, 7.} следует, что этот определитель является элементом из $\mo$, как только все $\alpha_i$ принадлежат $\mO$.
% END UNIT BODY
% END PRODUCER UNIT 164
\endgroup


\clearpage
\begingroup
\setcounter{footnote}{0}
% BEGIN PRODUCER UNIT 165
% Source: C:/Users/Floris/Documents/interlanguage/03_projects/noether/02_slavic_working_corpus/translations/paper31/section08_entry04/russian/v001/Noether_Paper31_Section08_Entry04_Russian_v001.tex
% Identity: 3267 B / SHA-256 3E1E023028C68F74837019721A8B8CF60AE135C81D2C05499D82898E89ACEA53
\newcommand{\mo}{\mathfrak{o}}
\newcommand{\mO}{\mathfrak{O}}
\newcommand{\mT}{\mathfrak{T}}
\newcommand{\mR}{\mathfrak{R}}
\newcommand{\mq}{\mathfrak{q}}
\setlength{\parindent}{1.25em}
\setlength{\parskip}{0pt}
\makeatletter
\renewcommand\footnotesize{\@setfontsize\footnotesize{7.7}{8.7}\selectfont}
\makeatother
\emergencystretch=100em
\allowdisplaybreaks
\sloppy
\hbadness=10000
\vbadness=10000
\hfuzz=2pt
\vfuzz=10pt
\raggedbottom
% BEGIN UNIT BODY
% Unit: Paper31 Section08 Entry04 expanded discriminant-ideal definition v001. Direct Russian translation from clean RA34 line 435 / segment P31-S0097, with scan-visible finite-basis proof detail and Idealtheorie footnote 2 restored from printed page 102 / PDF page 21.
\setcounter{footnote}{34}

\noindent \emph{Дискриминантный идеал.} Пусть $\tau_1,\ldots,\tau_n$ пробегают все системы из $n$ элементов из $\mT$; для каждой такой системы образуется определитель $|S(\tau_i\tau_j)|$. Следовательно, он является элементом из $\mo$, причем ненулевым, как только $\tau_1,\ldots,\tau_n$ линейно независимы. Идеал в $\mo$, выведенный из совокупности этих определителей, называется дискриминантным идеалом порядка $\mT$ относительно $\mo$; следовательно, он отличен от нулевого идеала.

Пусть далее $\beta_1,\ldots,\beta_m$ при $m\ge n$ является (всегда существующим)\footnote{\emph{Idealtheorie}, §9 и §3, 1.} $\mo$-модульным базисом $\mT$. Тогда конечное число определителей, соответствующих выбору любых $n$ элементов среди $\beta_1,\ldots,\beta_m$, образует базис этого дискриминантного идеала. Это следует из модульной гомоморфности от $\mT$ к системе следов из $\mT$ (§4, 3). Действительно, определители, образованные из произведений $\tau_i\tau_j$, линейно выражаются с коэффициентами из $\mo$ через конечное число определителей $|S(\beta_i\beta_j)|$, соответствующих $n$ элементам $\beta$; поэтому то же верно для $|S(\tau_i\tau_j)|$ относительно $|S(\beta_i\beta_j)|$. Если, в частности, порядок имеет линейно независимый модульный базис $\alpha_1,\ldots,\alpha_n$, то определитель $|S(\alpha_i\alpha_j)|$ можно взять как базис дискриминантного идеала; определение совпадает с определением в §7, 1.
% END UNIT BODY
% END PRODUCER UNIT 165
\endgroup


\clearpage
\begingroup
\setcounter{footnote}{0}
% BEGIN PRODUCER UNIT 166
% Source: C:/Users/Floris/Documents/interlanguage/03_projects/noether/02_slavic_working_corpus/translations/paper31/section08_entry05/russian/v001/Noether_Paper31_Section08_Entry05_Russian_v001.tex
% Identity: 2817 B / SHA-256 0733ABBB8FF4F46053016CD6B478869B63A927737B438AB1DD56F53CA3B5527A
\newcommand{\mo}{\mathfrak{o}}
\newcommand{\mO}{\mathfrak{O}}
\newcommand{\mT}{\mathfrak{T}}
\newcommand{\mR}{\mathfrak{R}}
\newcommand{\mq}{\mathfrak{q}}
\newcommand{\ma}{\mathfrak{a}}
\setlength{\parindent}{1.25em}
\setlength{\parskip}{0pt}
\makeatletter
\renewcommand\footnotesize{\@setfontsize\footnotesize{7.7}{8.7}\selectfont}
\makeatother
\emergencystretch=100em
\allowdisplaybreaks
\sloppy
\hbadness=10000
\vbadness=10000
\hfuzz=2pt
\vfuzz=10pt
\raggedbottom
% BEGIN UNIT BODY
% Unit: Paper31 Section08 Entry05 quotient-ring passage v001. Direct Russian translation from clean RA34 line 437 / segment P31-S0098, with scan-visible heading footnote 3 and denominator-prime footnote 4 restored from printed pages 102--103 / PDF pages 21--22.
\setcounter{footnote}{35}

\noindent 3. \emph{Переход к кольцу частных.}\footnote{Доказательства всех утверждений, приведенных в этом пункте, см. у H. Grell (в работе, цитированной в примечании 6, с. 83).} Пусть $\ma$ -- произвольный идеал из $\mo$. Кольцо частных $\mT_\ma$ определяется как совокупность элементов поля частных $\mT$, имеющих в знаменателе элемент из $\mT$, простой к $\ma$.\footnote{Так как простые идеалы из $\mT$ не имеют собственных делителей, отличных от единичного идеала, здесь "простой к" совпадает с "взаимно простой" -- в смысле идеалов.} Помимо нулевого и единичного идеалов, кольцо $\mT_\ma$ не имеет других простых идеалов, кроме идеалов расширения простых идеалов, принадлежащих $\ma$. Существует взаимно однозначное соответствие и изоморфизм колец классов вычетов между всеми идеалами из $\mT_\ma$, отличными от нулевого и единичного, и теми идеалами из $\mT$, чьи соответствующие простые идеалы содержатся среди простых идеалов, принадлежащих $\ma$; кроме того, нулевой и единичный идеалы взаимно однозначно соответствуют друг другу.
% END UNIT BODY
% END PRODUCER UNIT 166
\endgroup


\clearpage
\begingroup
\setcounter{footnote}{0}
% BEGIN PRODUCER UNIT 167
% Source: C:/Users/Floris/Documents/interlanguage/03_projects/noether/02_slavic_working_corpus/translations/paper31/section08_entry06/russian/v001/Noether_Paper31_Section08_Entry06_Russian_v001.tex
% Identity: 2825 B / SHA-256 A888DA801DD86616BFDB6A81412120E57C972164982CC7E50943C4B831C839DF
\newcommand{\mo}{\mathfrak{o}}
\newcommand{\mO}{\mathfrak{O}}
\newcommand{\mT}{\mathfrak{T}}
\newcommand{\mR}{\mathfrak{R}}
\newcommand{\mq}{\mathfrak{q}}
\newcommand{\ma}{\mathfrak{a}}
\newcommand{\mb}{\mathfrak{b}}
\newcommand{\mc}{\mathfrak{c}}
\newcommand{\mpideal}{\mathfrak{p}}
\setlength{\parindent}{1.25em}
\setlength{\parskip}{0pt}
\makeatletter
\renewcommand\footnotesize{\@setfontsize\footnotesize{7.7}{8.7}\selectfont}
\makeatother
\emergencystretch=100em
\allowdisplaybreaks
\sloppy
\hbadness=10000
\vbadness=10000
\hfuzz=2pt
\vfuzz=10pt
\raggedbottom
% BEGIN UNIT BODY
% Unit: Paper31 Section08 Entry06 multiplication-ring quotient localization v001. Direct Russian translation from clean RA34 line 439 / segment P31-S0099, with scan-visible final equality/proof sentence restored from printed page 103 / PDF page 22.
\setcounter{footnote}{37}

\noindent Те же рассуждения применимы и к мультипликационному кольцу $\mo$. В частности, если $\mpideal$ -- простой идеал, то $\mo_\mpideal$ имеет только один простой идеал, отличный от нулевого и единичного идеалов; благодаря изоморфизму колец классов вычетов кольцо $\mo_\mpideal$, как и $\mo$, является мультипликационным кольцом, а по пункту 1 поэтому $\mo_\mpideal$ является кольцом главных идеалов. Элемент, порождающий простой идеал, выведенный из $\mpideal$, становится простым элементом $p$ из $\mo_\mpideal$. Идеал расширения произвольного идеала $\mc$ из $\mo$ порождается первичной компонентой $\mc$, принадлежащей $\mpideal$; значит, он равен $(p)^\alpha$ или единичному идеалу в зависимости от того, входит ли $\mpideal$ в $\mc$ -- именно в $\alpha$-й степени -- или не входит в $\mc$. Отсюда еще следует: если расширения двух идеалов $\mb$ и $\mc$ из $\mo$ совпадают во всех кольцах частных $\mo_\mpideal$, то $\mb$ и $\mc$ тождественны, поскольку они делятся одними и теми же простыми идеалами $\mpideal$ и в одинаковых степенях.
% END UNIT BODY
% END PRODUCER UNIT 167
\endgroup


\clearpage
\begingroup
\setcounter{footnote}{0}
% BEGIN PRODUCER UNIT 168
% Source: C:/Users/Floris/Documents/interlanguage/03_projects/noether/02_slavic_working_corpus/translations/paper31/section08_entry07/russian/v001/Noether_Paper31_Section08_Entry07_Russian_v001.tex
% Identity: 1866 B / SHA-256 A24BE64158EB9A374D968A878ECB395152FD96F1BAA9C21CD9261327F3C1E506
\newcommand{\mo}{\mathfrak{o}}
\newcommand{\mT}{\mathfrak{T}}
\newcommand{\mP}{\mathfrak{P}}
\newcommand{\mpideal}{\mathfrak{p}}
\setlength{\parindent}{1.25em}
\setlength{\parskip}{0pt}
\makeatletter
\renewcommand\footnotesize{\@setfontsize\footnotesize{7.7}{8.7}\selectfont}
\makeatother
\emergencystretch=100em
\allowdisplaybreaks
\sloppy
\hbadness=10000
\vbadness=10000
\hfuzz=2pt
\vfuzz=10pt
\raggedbottom
% BEGIN UNIT BODY
% Unit: Paper31 Section08 Entry07 discriminant ideal under passage to quotient ring v001. Direct Russian translation from clean RA34 line 441 / segment P31-S0100, with scan-visible H. Grell proof footnote restored from printed page 103 / PDF page 22.
\setcounter{footnote}{37}

\noindent 4. \emph{Дискриминантный идеал при переходе к кольцу частных.} Пусть $\mpideal$ -- простой идеал из $\mo$, и обозначим через $\mP=\mT\mpideal$ идеал расширения в $\mT$. Кольцо частных $\mT_\mP$ обладает модульным базисом из линейно независимых элементов относительно кольца главных идеалов $\mo_\mpideal$; одновременно $\mT_\mP$ является конечным $\mpideal$-порядком в поле расширения $K$ поля частных кольца $\mo_\mpideal$.\footnote{Для доказательства ср. H. Grell.} Поэтому в $\mT_\mP$ относительно $\mo_\mpideal$ выполнена дискриминантная теорема (§7, 3).
% END UNIT BODY
% END PRODUCER UNIT 168
\endgroup


\clearpage
\begingroup
\setcounter{footnote}{0}
% BEGIN PRODUCER UNIT 169
% Source: C:/Users/Floris/Documents/interlanguage/03_projects/noether/02_slavic_working_corpus/translations/paper31/section08_entry08/russian/v001/Noether_Paper31_Section08_Entry08_Russian_v001.tex
% Identity: 3521 B / SHA-256 D9826F0899F7EA7215B0A3A4372A0B1E8020094891F9D0770B30CB582A62C0CE
\newcommand{\mo}{\mathfrak{o}}
\newcommand{\mO}{\mathfrak{O}}
\newcommand{\mT}{\mathfrak{T}}
\newcommand{\mP}{\mathfrak{P}}
\newcommand{\mpideal}{\mathfrak{p}}
\setlength{\parindent}{1.25em}
\setlength{\parskip}{0pt}
\makeatletter
\renewcommand\footnotesize{\@setfontsize\footnotesize{7.4}{8.4}\selectfont}
\makeatother
\emergencystretch=100em
\allowdisplaybreaks
\sloppy
\hbadness=10000
\vbadness=10000
\hfuzz=2pt
\vfuzz=10pt
\raggedbottom
% BEGIN UNIT BODY
% Unit: Paper31 Section08 Entry08 discriminant basis of extended ideal v001. Direct Russian translation from clean RA34 line 443 / segment P31-S0101, with scan-visible "what is always possible", "nach 1", and the full relative-discriminant footnote restored from printed page 103 / PDF page 22.
\setcounter{footnote}{38}

\noindent Дискриминант $D_{\mT_\mP}$ при этом является базисом идеала расширения в $\mo_\mpideal$ дискриминантного идеала порядка $\mT$ относительно $\mo$. Действительно, что всегда возможно, выберем элементы $\alpha_1,\ldots,\alpha_n$ из $\mT$ как $\mo_\mpideal$-модульный базис $\mT_\mP$; тогда $|S(\alpha_i\alpha_k)|$ принадлежит дискриминантному идеалу, и по пункту 1 каждый определитель $|S(\tau_i\tau_k)|$ в $\mo_\mpideal$ делится на $|S(\alpha_i\alpha_k)|$.\footnote{Отсюда следует -- когда $\mo$ является главным порядком алгебраического числового поля -- согласие с относительным дискриминантом по Hilbert (который, как известно, совпадает с определением Hecke). Действительно, пусть $\omega_1,\ldots,\omega_m$ является $\mo$-модульным базисом главного порядка относительного поля -- Hilbert, в частности, выбирает модульный базис относительно кольца целых рациональных чисел. Тогда по Hilbert образуют все определители из каждых $n$ базисных элементов и их сопряженных; "относительный дискриминант" есть идеал в $\mo$, выведенный из произведений каждых двух, одинаковых или различных, из этих определителей. Поэтому относительный дискриминант в $\mo_\mpideal$ переходит в квадрат определителя базиса $\alpha_1,\ldots,\alpha_n$ и его сопряженных элементов, то есть в $|S(\alpha_i\alpha_k)|$. Значит, расширения относительного дискриминанта и дискриминантного идеала (относительного главного порядка) совпадают во всех кольцах частных $\mo_\mpideal$; по заключению пункта 3 эти два идеала тождественны.}
% END UNIT BODY
% END PRODUCER UNIT 169
\endgroup


\clearpage
\begingroup
\setcounter{footnote}{0}
% BEGIN PRODUCER UNIT 170
% Source: C:/Users/Floris/Documents/interlanguage/03_projects/noether/02_slavic_working_corpus/translations/paper31/section08_entry09/russian/v001/Noether_Paper31_Section08_Entry09_Russian_v001.tex
% Identity: 1746 B / SHA-256 6327F45FC77181E834373BFD521BE850A278674C648A3F86A0E525D9726B6C4C
\newcommand{\mo}{\mathfrak{o}}
\newcommand{\mO}{\mathfrak{O}}
\newcommand{\mT}{\mathfrak{T}}
\newcommand{\mP}{\mathfrak{P}}
\newcommand{\mpideal}{\mathfrak{p}}
\setlength{\parindent}{1.25em}
\setlength{\parskip}{0pt}
\makeatletter
\renewcommand\footnotesize{\@setfontsize\footnotesize{7.4}{8.4}\selectfont}
\makeatother
\emergencystretch=100em
\allowdisplaybreaks
\sloppy
\hbadness=10000
\vbadness=10000
\hfuzz=2pt
\vfuzz=10pt
\raggedbottom
% BEGIN UNIT BODY
% Unit: Paper31 Section08 Entry09 general discriminant theorem statement v001. Direct Russian translation from clean RA34 line 445 / segment P31-S0102, checked against printed scan page 103 / PDF page 22 and Claude handoff OCR line 433. No new footnote belongs to this theorem-statement unit.
\setcounter{footnote}{39}

\noindent 5. \textbf{Общая теорема о дискриминанте.} \emph{Простой идеал $\mpideal$ из мультипликационного кольца $\mo$ входит в дискриминантный идеал порядка относительно $\mo$ тогда и только тогда, когда в разложении идеала $\mT\mpideal$ на попарно взаимно простые первичные компоненты из $\mT$ появляется по крайней мере одна собственно первичная компонента или по крайней мере один простой идеал второго рода.}
% END UNIT BODY
% END PRODUCER UNIT 170
\endgroup


\clearpage
\begingroup
\setcounter{footnote}{0}
% BEGIN PRODUCER UNIT 171
% Source: C:/Users/Floris/Documents/interlanguage/03_projects/noether/02_slavic_working_corpus/translations/paper31/section08_entry10/russian/v001/Noether_Paper31_Section08_Entry10_Russian_v001.tex
% Identity: 2514 B / SHA-256 CF288D5DB506E3A660D991F57C0F4004B959763401819628EEAAAB312C1C946A
\newcommand{\mo}{\mathfrak{o}}
\newcommand{\mO}{\mathfrak{O}}
\newcommand{\mT}{\mathfrak{T}}
\newcommand{\mP}{\mathfrak{P}}
\newcommand{\mpideal}{\mathfrak{p}}
\setlength{\parindent}{1.25em}
\setlength{\parskip}{0pt}
\makeatletter
\renewcommand\footnotesize{\@setfontsize\footnotesize{7.4}{8.4}\selectfont}
\makeatother
\emergencystretch=100em
\allowdisplaybreaks
\sloppy
\hbadness=10000
\vbadness=10000
\hfuzz=2pt
\vfuzz=10pt
\raggedbottom
% BEGIN UNIT BODY
% Unit: Paper31 Section08 Entry10 proof and number-field specialization of the general discriminant theorem v001. Direct Russian translation from clean RA34 lines 447 and 449 / segments P31-S0103--P31-S0104, reconciled with printed scan page 104 / PDF page 23 and Claude handoff OCR lines 443--445. The clean RA34 footnote after the proof duplicates the printed page-103 Hilbert/Hecke note restored in Entry08, so it is intentionally not repeated here.
\setcounter{footnote}{39}

\noindent Из пунктов 3 и 4 следует, что простой идеал $\mpideal$ из $\mo$ входит в дискриминантный идеал порядка $\mT$ тогда и только тогда, когда простой элемент $p$ из $\mo_\mpideal$ входит в дискриминант $D_{\mT_\mP}$. Но это тождественно тому, что кольцо классов вычетов $\mT_\mP/\mT_\mP p$ не является вполне приводимым первого рода; поскольку по пункту 3 это кольцо классов вычетов изоморфно $\mT/\mT\mpideal$, тем самым общая теорема о дискриминанте доказана.

\noindent В качестве специализации к относительным дискриминантам числовых полей получается: простой идеал $\mpideal$ главного порядка числового поля входит в относительный дискриминант поля расширения тогда и только тогда, когда $\mpideal$ в главном порядке поля расширения делится на квадрат простого идеала.
% END UNIT BODY
% END PRODUCER UNIT 171
\endgroup


\clearpage
\begingroup
\setcounter{footnote}{0}
% BEGIN PRODUCER UNIT 172
% Source: C:/Users/Floris/Documents/interlanguage/03_projects/noether/02_slavic_working_corpus/translations/paper31/section08_entry11/russian/v001/Noether_Paper31_Section08_Entry11_Russian_v001.tex
% Identity: 887 B / SHA-256 D96473D582F85610452553448778F69D5EFC9DB80C15C40AD76EEF9F454B28EA
\setlength{\parindent}{1.25em}
\setlength{\parskip}{0pt}
\makeatletter
\renewcommand\footnotesize{\@setfontsize\footnotesize{7.4}{8.4}\selectfont}
\makeatother
\emergencystretch=100em
\allowdisplaybreaks
\sloppy
\hbadness=10000
\vbadness=10000
\hfuzz=2pt
\vfuzz=10pt
\raggedbottom
% BEGIN UNIT BODY
% Unit: Paper31 Section08 Entry11 closing place/date v001. Direct Russian rendering of clean RA34 line 451 / segment P31-S0105, checked against printed scan page 104 / PDF page 23 and Claude handoff OCR line 447. This closes Paper31 before Paper32.
\setcounter{footnote}{39}

\noindent Гёттинген, 30 марта 1926 г.
% END UNIT BODY
% END PRODUCER UNIT 172
\endgroup


\clearpage
\begingroup
\setcounter{footnote}{0}
% BEGIN PRODUCER UNIT 173
% Source: C:/Users/Floris/Documents/interlanguage/03_projects/noether/02_slavic_working_corpus/translations/paper32/opening/russian/v001/Noether_Paper32_Opening_Russian_v001.tex
% Identity: 1279 B / SHA-256 FB81BAAF91D296FCCD90F3FA77CC1B6AFE4C2147577C66E0294EDA670573A0A9
\newcommand{\mo}{\mathfrak{o}}
\newcommand{\mO}{\mathfrak{O}}
\newcommand{\mT}{\mathfrak{T}}
\newcommand{\mR}{\mathfrak{R}}
\newcommand{\mq}{\mathfrak{q}}
\setlength{\parindent}{1.25em}
\setlength{\parskip}{0pt}
\makeatletter
\renewcommand\footnotesize{\@setfontsize\footnotesize{7.7}{8.7}\selectfont}
\makeatother
\emergencystretch=100em
\allowdisplaybreaks
\sloppy
\hbadness=10000
\vbadness=10000
\hfuzz=2pt
\vfuzz=10pt
\raggedbottom
% BEGIN UNIT BODY
% Unit: Paper32 opening/title v001. Direct Russian translation from clean RA34 Paper32 lines 1--6 / segments P32-S0001--P32-S0003, checked against the source scan and Claude batch OCR witness.

\setcounter{footnote}{0}

\section*{32. Совместно с Р. Брауэром: О минимальных полях расщепления неприводимых представлений}
\begin{center}
\emph{Sitz. Ber. d. Preuß. Akad. d. Wiss. 1927, с. 221--228.}
\end{center}

\begin{center}\emph{(Представлено г-ном Шуром.)}\end{center}
% END UNIT BODY
% END PRODUCER UNIT 173
\endgroup


\clearpage
\begingroup
\setcounter{footnote}{0}
% BEGIN PRODUCER UNIT 174
% Source: C:/Users/Floris/Documents/interlanguage/03_projects/noether/02_slavic_working_corpus/translations/paper32/schur_paragraph/russian/v001/Noether_Paper32_Schur_Paragraph_Russian_v001.tex
% Identity: 3926 B / SHA-256 A3C2E3AE0AE29923AACE6457A84CC50FA4927EFD705178E64FD0F067C19E2C04
\newcommand{\mo}{\mathfrak{o}}
\newcommand{\mO}{\mathfrak{O}}
\newcommand{\mT}{\mathfrak{T}}
\newcommand{\mR}{\mathfrak{R}}
\newcommand{\mq}{\mathfrak{q}}
\setlength{\parindent}{1.25em}
\setlength{\parskip}{0pt}
\makeatletter
\renewcommand\footnotesize{\@setfontsize\footnotesize{7.7}{8.7}\selectfont}
\makeatother
\emergencystretch=100em
\allowdisplaybreaks
\sloppy
\hbadness=10000
\vbadness=10000
\hfuzz=2pt
\vfuzz=10pt
\raggedbottom
% BEGIN UNIT BODY
% Unit: Paper32 Schur paragraph v001. Direct Russian translation from clean RA34 Paper32 line 8 / segment P32-S0004, checked against source scan page 1 and Claude batch OCR witness.
\setcounter{footnote}{0}

\noindent Вопрос о числовых полях наименьшей степени над основным полем $P$, в которых неприводимое над $P$ представление конечной группы распадается на абсолютно неприводимые составляющие, впервые рассмотрел И. Шур.\footnote{\emph{Arithmetische Untersuchungen über endliche Gruppen}, Sitzungsber. d. Berl. Ak. d. Wiss. 1906, с. 164. \emph{Beiträge zur Theorie der Gruppen linearer Substitutionen}, Transact. of the Am. Math. Soc. Ser. 2, Vol. XV, 1909, с. 159. Во второй работе результаты перенесены на вполне приводимые гиперкомплексные системы. В один класс объединяется представление со всеми его преобразованными (подобными).} Пусть, например, $P$ содержит характер такой составляющей. Тогда И. Шур показал, что степень этих полей над $P$ -- индекс -- совпадает с числом тех абсолютно неприводимых составляющих, которые все принадлежат одному и тому же классу; что, следовательно, то же самое поле уже определяется требованием выделить одну абсолютно неприводимую составляющую; далее, что степень всякого поля над $P$, в котором происходит такой распад, является кратной индексу. Все эти поля будем называть полями расщепления, также и в случае, когда $P$ не содержит характера. На примере И. Шур показал, что поля расщепления наименьшей степени не обязаны быть изоморфными. Если $P$ не содержало соответствующего неприводимого характера, то поля расщепления должны содержать поле такого характера и его сопряженные относительно $P$; абсолютно неприводимые представления, принадлежащие различным сопряженным характерам, хотя и сопряжены, но, очевидно, принадлежат различным классам. Для выделения системы абсолютно неприводимых представлений одного класса и здесь достаточно поля соответствующего характера и одного из соответствующих полей расщепления.
% END UNIT BODY
% END PRODUCER UNIT 174
\endgroup


\clearpage
\begingroup
\setcounter{footnote}{0}
% BEGIN PRODUCER UNIT 175
% Source: C:/Users/Floris/Documents/interlanguage/03_projects/noether/02_slavic_working_corpus/translations/paper32/second_intro/russian/v001/Noether_Paper32_Second_Intro_Russian_v001.tex
% Identity: 5072 B / SHA-256 6F2C8C53B72FC2D039C4EFA473FC5985483BBF48E488C6A94462E7DFDF9FD812
\newcommand{\mo}{\mathfrak{o}}
\newcommand{\mO}{\mathfrak{O}}
\newcommand{\mT}{\mathfrak{T}}
\newcommand{\mR}{\mathfrak{R}}
\newcommand{\mq}{\mathfrak{q}}
\setlength{\parindent}{1.25em}
\setlength{\parskip}{0pt}
\makeatletter
\renewcommand\footnotesize{\@setfontsize\footnotesize{7.4}{8.4}\selectfont}
\makeatother
\emergencystretch=100em
\allowdisplaybreaks
\sloppy
\hbadness=10000
\vbadness=10000
\hfuzz=2pt
\vfuzz=10pt
\raggedbottom
% BEGIN UNIT BODY
% Unit: Paper32 second introductory paragraph v001. Direct Russian translation from clean RA34 Paper32 line 10 / segment P32-S0005, checked against source scan pages 1--2 and Claude batch OCR witness.
\setcounter{footnote}{0}

\noindent В дальнейшем вопрос о полях расщепления будет прослежен дальше. Поля расщепления наименьшей степени можно характеризовать посредством соответствующих некоммутативных тел; общие поля расщепления -- посредством двусторонне простых колец, инвариантно связанных с этими некоммутативными телами. Далее на примере тела кватернионов будет показано, что степени минимальных полей расщепления не ограничены; при этом поле расщепления называется минимальным, если ни одно из его собственных подполей не является полем расщепления. Эта неограниченность существенно следует из теоретико-числовой теоремы существования, доказательство которой дано Г. Хассе в непосредственно следующей заметке. Однако пример также разбирается элементарно-вычислительно -- скорее в проверочном плане; тем самым без привлечения общей теории и теоретико-числовой теоремы существования неограниченность степеней показывается путем элементарного доказательства менее сильной, но достаточной теоремы существования.\footnote{Этот пример принадлежит Э. Нётер; он основан на модификации примера Р. Брауэра, в котором вообще было показано существование минимального поля расщепления, степень которого превышает индекс. Элементарная обработка примера принадлежит Р. Брауэру. Характеризация полей расщепления восходит к исследованиям авторов, которые в остальном преследуют раздельные цели. У Э. Нётер речь идет о построении теории представлений на основе теории модулей и идеалов при общих предположениях конечности; у Р. Брауэра -- о построении некоммутативных тел конечной степени над заданным коммутативным совершенным основным полем с помощью фактор-систем. К характеризации полей расщепления наименьшей степени авторы пришли независимо: Р. Брауэр для совершенного основного поля, Э. Нётер без этого ограничения. Характеризацию общих полей расщепления Р. Брауэр нашел в связи с вопросом Э. Нётер о возможности некоммутативного вложения; Э. Нётер и здесь смогла снять ограничение на совершенную основную область.} Следует еще заметить, что и в этом примере речь идет о полях расщепления конечной группы, а именно группы кватернионов, лежащей также в основе примера Шура. Указанные там представления являются одновременно (изоморфными) представлениями тела кватернионов.
% END UNIT BODY
% END PRODUCER UNIT 175
\endgroup


\clearpage
\begingroup
\setcounter{footnote}{0}
% BEGIN PRODUCER UNIT 176
% Source: C:/Users/Floris/Documents/interlanguage/03_projects/noether/02_slavic_working_corpus/translations/paper32/section01_opening/russian/v001/Noether_Paper32_Section01_Opening_Russian_v001.tex
% Identity: 3159 B / SHA-256 88B42C49A85B45870D2F59B84D2A3E7A9F4D15C9E6005BEFED3352A5776DA4A3
\newcommand{\mS}{\mathfrak{S}}
\newcommand{\mA}{\mathfrak{A}}
\newcommand{\mo}{\mathfrak{o}}
\newcommand{\mO}{\mathfrak{O}}
\newcommand{\mT}{\mathfrak{T}}
\newcommand{\mR}{\mathfrak{R}}
\newcommand{\mq}{\mathfrak{q}}
\setlength{\parindent}{1.25em}
\setlength{\parskip}{0pt}
\makeatletter
\renewcommand\footnotesize{\@setfontsize\footnotesize{7.6}{8.6}\selectfont}
\makeatother
\emergencystretch=100em
\allowdisplaybreaks
\sloppy
\hbadness=10000
\vbadness=10000
\hfuzz=2pt
\vfuzz=10pt
\raggedbottom
% BEGIN UNIT BODY
% Unit: Paper32 section 1 opening v001. Direct Russian translation from the R122 German Paper32/Paper33 source slice, lines 375 and 378 / segments P32-S0006--P32-S0007.
\subsection*{1. Характеризация полей расщепления.}

\noindent Прежде всего кратко сведем основные понятия. Пусть $\mS$ обозначает гиперкомплексную систему над полем $P_0$. Под представлением $\Gamma$ системы $\mS$ -- над $P$ -- понимается система матриц с элементами из $P$ такая, что $\Gamma$ становится гомоморфным образом $\mS$ и что элементам $p_0$ из $P_0$ сопоставлены диагональные матрицы $p_0E$; следовательно, $P$ должно содержать $P_0$. Представление $\Gamma$ вместе со всеми своими преобразованными $P^{-1}\Gamma P$ образует класс представлений. В так определенные представления входят также представления конечных групп $G$; соответствующая гиперкомплексная система $\mS$, “групповое кольцо”, состоит из всех “групповых чисел”, то есть из всех линейных комбинаций элементов группы с коэффициентами из $P_0$, причем исходные элементы группы рассматриваются как линейно независимые, а умножение задается групповым умножением и правилами вычисления. Понятия “приводимых”, “неприводимых”, “вполне приводимых”, “абсолютно неприводимых” представлений определяются обычным образом; эту терминологию распространяем и на классы представлений. Система $\mS$ называется “вполне приводимой над $P$”, если все ее классы представлений над $P$ вполне приводимы.
% END UNIT BODY
% END PRODUCER UNIT 176
\endgroup


\clearpage
\begingroup
\setcounter{footnote}{0}
% BEGIN PRODUCER UNIT 177
% Source: C:/Users/Floris/Documents/interlanguage/03_projects/noether/02_slavic_working_corpus/translations/paper32/section01_reduction_theorem/russian/v001/Noether_Paper32_Section01_Reduction_Theorem_Russian_v001.tex
% Identity: 5093 B / SHA-256 388470E408F4DF3DD7205AEDFB8B1F4B065EC7F06775E8519AFA82AFF80C9D42
\newcommand{\mS}{\mathfrak{S}}
\newcommand{\mA}{\mathfrak{A}}
\newcommand{\mo}{\mathfrak{o}}
\newcommand{\mO}{\mathfrak{O}}
\newcommand{\mT}{\mathfrak{T}}
\newcommand{\mR}{\mathfrak{R}}
\newcommand{\mq}{\mathfrak{q}}
\setlength{\parindent}{1.25em}
\setlength{\parskip}{0pt}
\makeatletter
\renewcommand\footnotesize{\@setfontsize\footnotesize{7.6}{8.6}\selectfont}
\makeatother
\emergencystretch=100em
\allowdisplaybreaks
\sloppy
\hbadness=10000
\vbadness=10000
\hfuzz=2pt
\vfuzz=10pt
\raggedbottom
% BEGIN UNIT BODY
% Unit: Paper32 section 1 reduction theorem v001. Direct Russian translation from the R122 German Paper32/Paper33 source slice, lines 380 and 382 / segments P32-S0008--P32-S0009.
\noindent Пусть теперь $\mS$ предполагается вполне приводимой над $P_0$; пусть $P_0$ считается совершенным полем, так что то же верно для каждого алгебраического расширения $P$ поля $P_0$. Классы представлений при каждом расширении $P$ поля $P_0$ остаются вполне приводимыми.\footnote{Если $P_0$ не является совершенным, полная приводимость сохраняется тогда и только тогда, когда компоненты $K$ центра становятся “расширениями первого рода” поля $P_0$. При этом условии остаются верны все дальнейшие следствия текста.} Центр $\mS$ является прямой суммой конечного числа полей; всякое расширение $P$ поля $P_0$, изоморфное такому полю $K$, становится полем одного (простого) характера. Если расширить $P_0$ до изоморфного $K$ поля $P$, от $\mS$ отделяется двусторонне простое кольцо, которому соответствует абсолютно неприводимый класс представлений этого характера. Это кольцо -- по теореме Веддерберна -- является прямым произведением матричного кольца над $P$ и некоммутативного тела $\mA$, однозначно определенного с точностью до изоморфизма; центр $\mA$ равен $P$, а его степень над $P$ равна $m^2$, где $m$ обозначает индекс Шура, принадлежащий этому характеру. Всякое поле расщепления $\mS$ -- относительно представления, принадлежащего этому характеру -- одновременно является таким для $\mA$ относительно представлений над $P$, и наоборот. Сопряженные представления $\mS$ получают, исходя из соответствующего двусторонне простого кольца над $P_0$; его соответствующее тело $\mA_0$ становится изоморфным $\mA$ и имеет центр $K$. Тем самым задача о полях расщепления полностью сведена к задаче о соответствующем некоммутативном теле $\mA$ и его расщеплениях.\footnote{Это соответствует шуровскому сведению к системе матриц над $P$, перестановочных с неприводимым представлением $\mS$ над $P$. Транспонированные этих матриц дают представление $\mA$ над $P$.} В частности, имеет место теорема:

\medskip

\noindent Все максимальные коммутативные подтела $\mA$ имеют одну и ту же степень $m$ над $P$, где $m$ означает индекс Шура. Все эти максимальные коммутативные подтела являются полями расщепления наименьшей степени, и каждое поле расщепления наименьшей степени содержится среди них. Степень каждого поля расщепления кратна $m$. У $\mA$ имеется лишь один абсолютно неприводимый класс представлений, а также лишь один неприводимый класс представлений над $P$.
% END UNIT BODY
% END PRODUCER UNIT 177
\endgroup


\clearpage
\begingroup
\setcounter{footnote}{0}
% BEGIN PRODUCER UNIT 178
% Source: C:/Users/Floris/Documents/interlanguage/03_projects/noether/02_slavic_working_corpus/translations/paper32/section01_general_splitting_fields/russian/v001/Noether_Paper32_Section01_General_Splitting_Fields_Russian_v001.tex
% Identity: 3611 B / SHA-256 630D93C1AE9B71ED508BF72BE7A07352BC6D82ED86FCDD71D40675735A5A8142
\newcommand{\mS}{\mathfrak{S}}
\newcommand{\mA}{\mathfrak{A}}
\newcommand{\mo}{\mathfrak{o}}
\newcommand{\mO}{\mathfrak{O}}
\newcommand{\mT}{\mathfrak{T}}
\newcommand{\mR}{\mathfrak{R}}
\newcommand{\mq}{\mathfrak{q}}
\setlength{\parindent}{1.25em}
\setlength{\parskip}{0pt}
\makeatletter
\renewcommand\footnotesize{\@setfontsize\footnotesize{7.6}{8.6}\selectfont}
\makeatother
\emergencystretch=100em
\allowdisplaybreaks
\sloppy
\hbadness=10000
\vbadness=10000
\hfuzz=2pt
\vfuzz=10pt
\raggedbottom
% BEGIN UNIT BODY
% Unit: Paper32 section 1 general splitting fields v001. Direct Russian translation from the R122 German Paper32/Paper33 source slice, lines 384, 386, 388, and section-2 heading lines 390--393 / segments P32-S0010--P32-S0013.
\noindent Для характеристики общих полей расщепления обозначим через $\mA_r$ прямое произведение $\mA$ с матричным кольцом степени $r^2$ над $P$ (системой всех $r$-рядных матриц с элементами из $P$). Тогда $\mA_r$ является двусторонне простым, с $\mA$ как соответствующим некоммутативным телом; и всякое двусторонне простое кольцо, гиперкомплексное над $P$, которому соответствует $\mA$, получается именно так. $\mA_r$ также характеризуется как система всех $r$-рядных матриц с элементами из $\mA$.

\medskip

\noindent Всякое поле расщепления $\mA$ степени $mr$ -- если такое существует -- изоморфно одному из максимальных коммутативных подтел $\mA_r$. Обратно, все максимальные коммутативные подтела $\mA_r$ являются полями расщепления, каждое степени $ms$, где $s$ является делителем $r$. В регулярном случае все максимальные коммутативные подтела $\mA_r$ имеют степень $mr$. При этом под регулярным случаем понимается следующее: основное поле $P$ обладает тем свойством, что для всякого коммутативного поля $\Sigma$ над $P$, конечного над $P$, существуют еще коммутативные расширения над $\Sigma$ произвольной степени.\footnote{Так, например, регулярного случая нет, когда $P$ является полем всех действительных чисел.}

\medskip

\noindent Встречаются ли среди этих полей расщепления при $r>1$ также минимальные, само вложение в $\mA_r$ еще не решает. То, что это действительно может происходить для бесконечно многих $r$, показывает следующий пример.

\subsection*{2. Неограниченность степеней минимальных полей расщепления}
% END UNIT BODY
% END PRODUCER UNIT 178
\endgroup


\clearpage
\begingroup
\setcounter{footnote}{0}
% BEGIN PRODUCER UNIT 179
% Source: C:/Users/Floris/Documents/interlanguage/03_projects/noether/02_slavic_working_corpus/translations/paper32/section02_quaternion_idempotent/russian/v001/Noether_Paper32_Section02_Quaternion_Idempotent_Russian_v001.tex
% Identity: 3463 B / SHA-256 6BF5E4252E12CB02023CBA6B01E58D4D6FE0D091C48057180D1E5ADA5A8BEE10
\newcommand{\mS}{\mathfrak{S}}
\newcommand{\mA}{\mathfrak{A}}
\setlength{\parindent}{1.25em}
\setlength{\parskip}{0pt}
\makeatletter
\renewcommand\footnotesize{\@setfontsize\footnotesize{7.6}{8.6}\selectfont}
\makeatother
\emergencystretch=100em
\allowdisplaybreaks
\sloppy
\hbadness=10000
\vbadness=10000
\hfuzz=2pt
\vfuzz=10pt
\raggedbottom
% BEGIN UNIT BODY
% Unit: Paper32 section 2 quaternion idempotent v001. Direct Russian translation from the R122 German Paper32/Paper33 source slice, lines 395--420 / segment P32-S0014.
\noindent Для этой неограниченности теперь возьмем в качестве примера тело кватернионов, рассматриваемое как гиперкомплексная система относительно поля $P$ рациональных чисел; с базисными элементами $i,j,k$ помимо единицы $1$ и с известными соотношениями:
\[
 i^2=j^2=k^2=-1;\qquad ij=k;\quad jk=i;\quad ki=j;\quad ji=-k;\quad kj=-i;\quad ik=-j.
\]
Окажется, что полями расщепления являются все и только те алгебраические числовые поля, при присоединении которых существует идемпотентный элемент $r$, то есть такой, что $r^2=r$, причем $r$ отличен и от нулевого, и от единичного элемента. Эти числовые поля можно также охарактеризовать как те, в которых $(-1)$ представима как сумма трех квадратов, а следовательно и как сумма двух квадратов.\footnote{То, что из представимости $(-1)$ как суммы трех квадратов всегда следует представимость как суммы двух квадратов, показывает тождество
\[
(c^2+d^2)(a^2+b^2+c^2+d^2)=(ac+bd)^2+(ad-bc)^2+(c^2+d^2)^2,
\]
которое, например, получается из теоремы о произведении норм для кватернионов. Вместо того чтобы требовать существования идемпотентного элемента, для расщепления было бы достаточно требовать элемент нулевой нормы.} Действительно, если положить
\[
        r=\alpha+\frac12\beta i+\frac12\gamma j+\frac12\delta k,
\]
то
\[
        r^2=\left(\alpha^2-\frac14\beta^2-\frac14\gamma^2-\frac14\delta^2\right)
          +\alpha\beta i+\alpha\gamma j+\alpha\delta k,
\]
и потому условие $r^2=r$ равносильно системе
\[
4\alpha^2-4\alpha-\beta^2-\gamma^2-\delta^2=0;\quad
2\alpha\beta-\beta=0;\quad 2\alpha\gamma-\gamma=0;\quad 2\alpha\delta-\delta=0,
\]
а значит, поскольку $\beta,\gamma,\delta$ не должны обращаться в нуль одновременно, равносильно условию
\[
        (-1)=\beta^2+\gamma^2+\delta^2.
\]
% END UNIT BODY
% END PRODUCER UNIT 179
\endgroup


\clearpage
\begingroup
\setcounter{footnote}{0}
% BEGIN PRODUCER UNIT 180
% Source: C:/Users/Floris/Documents/interlanguage/03_projects/noether/02_slavic_working_corpus/translations/paper32/section02_idempotent_splitting_facts/russian/v001/Noether_Paper32_Section02_Idempotent_Splitting_Facts_Russian_v001.tex
% Identity: 4710 B / SHA-256 2D0DFE9308753060145E7A30F0146C1492D76126F578118A3BFE6A8F158CA867
\newcommand{\mS}{\mathfrak{S}}
\newcommand{\mA}{\mathfrak{A}}
\setlength{\parindent}{1.25em}
\setlength{\parskip}{0pt}
\makeatletter
\renewcommand\footnotesize{\@setfontsize\footnotesize{7.6}{8.6}\selectfont}
\makeatother
\emergencystretch=100em
\allowdisplaybreaks
\sloppy
\hbadness=10000
\vbadness=10000
\hfuzz=2pt
\vfuzz=10pt
\raggedbottom
% BEGIN UNIT BODY
% Unit: Paper32 section 2 idempotent splitting facts v001. Direct Russian translation from the R122 German Paper32/Paper33 source slice, lines 422 and 424 / segments P32-S0015--P32-S0016.
\noindent То, что каждый такой идемпотентный элемент порождает расщепление, и что существование такого идемпотентного элемента необходимо для расщепления, следует из следующих фактов. Всякое неприводимое представление полностью приводимой системы порождается односторонне простым идеалом, точнее соответствующим классом идеалов, и каждый такой класс идеалов порождает представление. Все эти односторонне простые идеалы возникают при одностороннем прямом разложении в сумму, и это разложение определяет «примитивные» идемпотентные элементы, которые становятся компонентами единицы. Обратно, каждый идемпотентный элемент дает разложение в сумму
\[
        \mS=r\mS+(e-r)\mS,
\]
где $e$ обозначает единицу; «примитивные» дают отделение простых идеалов.\footnote{Связь между примитивными идемпотентными элементами и неприводимыми представлениями имеется уже в диссертации М. Герцбергер о системах гиперкомплексных величин, глава 4, Берлин, 1923. То, что простые идеалы порождают неприводимые представления, для коммутативного случая можно найти, например, у Э. Нетер, теорема о дискриминанте для порядков, Crelle 157 (1926); для некоммутативного случая -- у В. Крулля, теория и применение обобщенных абелевых групп, Heidelberger Berichte 1926. Ср. также заключительное слово групповой теории Шпайзера. Если $a_1,\ldots,a_m$ есть базис идеала относительно поля расщепления, а $c$ -- произвольный элемент кольца, так что
\[
        a_i c=\gamma_{i1}a_1+\cdots+\gamma_{im}a_m,
\]
то $(\gamma_{ik})$ есть матрица, соответствующая $c$; утверждение, что все неприводимые представления порождаются идеалами, является более глубоким фактом.} Некоммутативное тело, рассматриваемое как гиперкомплексная система относительно своего центра, после расширения до гиперкомплексной системы относительно поля расщепления становится прямой суммой $m$ абсолютно простых односторонних идеалов, где $m$ есть индекс Шура.

\medskip

\noindent В случае тела кватернионов $m$ равно двум; следовательно, каждый идемпотентный элемент $\ne0$ и $\ne1$ уже порождает разложение на абсолютно простые идеалы, которому соответствует расщепление на абсолютно неприводимые классы представлений. Тем самым указанные поля охарактеризованы как поля расщепления.
% END UNIT BODY
% END PRODUCER UNIT 180
\endgroup


\clearpage
\begingroup
\setcounter{footnote}{0}
% BEGIN PRODUCER UNIT 181
% Source: C:/Users/Floris/Documents/interlanguage/03_projects/noether/02_slavic_working_corpus/translations/paper32/section02_cyclic_minimal_fields/russian/v001/Noether_Paper32_Section02_Cyclic_Minimal_Fields_Russian_v001.tex
% Identity: 4478 B / SHA-256 99FA1BB9BA80ACE4AB2D74FF20A822567271F43B00762E2B5EF3497DE23D3F54
\newcommand{\mS}{\mathfrak{S}}
\newcommand{\mA}{\mathfrak{A}}
\setlength{\parindent}{1.25em}
\setlength{\parskip}{0pt}
\makeatletter
\renewcommand\footnotesize{\@setfontsize\footnotesize{7.6}{8.6}\selectfont}
\makeatother
\emergencystretch=100em
\allowdisplaybreaks
\sloppy
\hbadness=10000
\vbadness=10000
\hfuzz=2pt
\vfuzz=10pt
\raggedbottom
% BEGIN UNIT BODY
% Unit: Paper32 section 2 cyclic minimal fields v001. Direct Russian translation from the R122 German Paper32/Paper33 source slice, lines 426 and 428--443 / segments P32-S0017--P32-S0019.
\noindent Далее следует: \emph{все циклические поля $\Omega_n$ степени $2^n$ над полем рациональных чисел, в которых $(-1)$ представляется как сумма двух квадратов, являются минимальными полями расщепления тела кватернионов.} Действительно, поле расщепления не может быть вещественным, поскольку в нем $(-1)$ представим как сумма квадратов; с другой стороны, подполе степени $2^{n-1}$ поля $\Omega_n$ -- а значит и всякое собственное подполе $\Omega_n$ -- необходимо вещественно, как пересечение $\Omega_n$ с полем всех вещественных алгебраических чисел. Именно, $\Omega_n$ является галуа-полем степени два относительно этого пересечения, которое поэтому необходимо совпадает с подполем степени $2^{n-1}$.

\medskip

\noindent По доказанной Г. Хассе теореме существования\footnote{Существование $\Omega_n$ можно уже усмотреть из параметрического представления циклических полей 4-й степени, как оно, например, дано у Weber, \emph{Kleines Lehrbuch der Algebra}, § 93.} \emph{такие поля $\Omega_n$ существуют для каждого $n$; степени минимальных полей расщепления не ограничены. Здесь даже каждая степень индекса появляется как степень минимального поля расщепления.}\footnote{Впоследствии Р. Брауер смог еще показать, что существуют даже минимальные поля расщепления тела кватернионов любой четной степени, то есть всех степеней, вообще возможных для полей расщепления. Такое поле $P(\Theta)$ можно для каждого $n\ge2$ задать уравнением
\[
 f(\Theta)=\Theta^{2n}+a^2p^2\Theta^2+b^2p^2\beta^2=0;
 \quad \text{то есть}\quad -1=\left(\frac{ap}{\Theta^{n-1}}\right)^2+
       \left(\frac{bp\beta}{\Theta^n}\right)^2,
\]
где подходящим выбором $a,b,p,\beta$ можно добиться, что 1. $f(x)$ при каждом $n$ становится неприводимым; 2. $P(\Theta)$ имеет только единственное отличное от $P$ подполе $P(\Theta^2)$; 3. среди сопряженных к $P(\Theta^2)$ имеются и вещественные, так что оно не может быть полем расщепления. Один возможный выбор $a,b,p,\beta$, например, таков:
\[
        f(x)=x^{2n}+(2^5\cdot9)4x^2+9\cdot64.
\]
Доказательство основано на модификации метода Фуртвенглера для построения примитивных уравнений (Math. Ann. 85, с. 34, § 3), но в деталях оно довольно трудоемко.}

\begin{center}
\textbf{3. Элементарное рассмотрение полей расщепления\\
тела кватернионов.}
\end{center}
% END UNIT BODY
% END PRODUCER UNIT 181
\endgroup


\clearpage
\begingroup
\setcounter{footnote}{0}
% BEGIN PRODUCER UNIT 182
% Source: C:/Users/Floris/Documents/interlanguage/03_projects/noether/02_slavic_working_corpus/translations/paper32/section03_elementary_quaternion_criterion/russian/v001/Noether_Paper32_Section03_Elementary_Quaternion_Criterion_Russian_v001.tex
% Identity: 3732 B / SHA-256 B8502267E17E4870F144094999D41E60324F83EBE4EE9B28CEFE7B2144E8E9C0
\newcommand{\mS}{\mathfrak{S}}
\newcommand{\mA}{\mathfrak{A}}
\newcommand{\mat}[1]{\begin{pmatrix}#1\end{pmatrix}}
\setlength{\parindent}{1.25em}
\setlength{\parskip}{0pt}
\makeatletter
\renewcommand\footnotesize{\@setfontsize\footnotesize{7.6}{8.6}\selectfont}
\makeatother
\emergencystretch=100em
\allowdisplaybreaks
\sloppy
\hbadness=10000
\vbadness=10000
\hfuzz=2pt
\vfuzz=10pt
\raggedbottom
% BEGIN UNIT BODY
% Unit: Paper32 section 3 elementary quaternion criterion v001. Direct Russian translation from the R122 German Paper32/Paper33 source slice, lines 445--480 / segments P32-S0020--P32-S0022.
\noindent Пусть $\mathfrak{D}$ обозначает представление группы кватернионов, образованное следующими 8 матрицами:\footnote{См. Sylvester, \emph{Math. Papers} Vol. III, с. 647; Vol. IV, с. 122.}
\[
\pm E=\pm\mat{1&0\\0&1},\qquad
\pm I=\pm\mat{i&0\\0&-i},\qquad
\pm J=\pm\mat{0&1\\-1&0},\qquad
\pm K=\pm\mat{0&i\\ i&0}.
\]
Групповое кольцо, порожденное $\mathfrak{D}$, есть как раз представление тела кватернионов, рассматриваемого как гиперкомплексная система; мы элементарно покажем, что $\mathfrak{D}$ реализуемо над полем $K$ тогда и только тогда, когда $-1$ в $K$ представим как сумма двух квадратов.

\medskip

\noindent \emph{a)} Пусть $\mathfrak{D}^*=T^{-1}\mathfrak{D}T$ есть представление, рациональное над полем $K$, и пусть $I^*=T^{-1}IT$, $J^*=T^{-1}JT$. Поскольку $J$ и $J^*$ являются двумя подобными матрицами, рациональными над $K$, существует также рациональное над $K$ преобразование $R$, переводящее $J^*$ в $J$,
\[
        R^{-1}J^*R=J.
\]
Если с самого начала заменить $\mathfrak{D}^*$ также рациональным над $K$ представлением $R^{-1}\mathfrak{D}^*R$, то видно, что без ограничения общности можно считать $J^*=J$. Пусть
\[
        I^*=\mat{a&b\\ c&d}\qquad (a,b,c,d\text{ числа из }K).
\]
Из $IJ=-JI$ следует $I^*J^*=-J^*I^*$, а ввиду $J^*=J$ это дает
\[
        a=-d,
        \qquad b=c.
\]
Из $I^2=-E$ следует $I^{*2}=-E$, то есть
\[
        \mat{-1&0\\0&-1}=\mat{a&b\\ b&-a}^2
        =\mat{a^2+b^2&0\\0&a^2+b^2},
\]
следовательно, $a^2+b^2=-1$, и потому $-1$ действительно представим в $K$ как сумма двух квадратов.

\medskip

\noindent \emph{b)} Пусть в поле $K$ число $-1$ представимо как сумма двух квадратов. Тогда, если, например, $-1=a^2+b^2$, положим
\[
        I^*=\mat{a&b\\ b&-a},\qquad
        J^*=\mat{0&1\\-1&0},\qquad
        K^*=I^*J^*.
\]
Как легко проверить, тогда $\pm E$, $\pm I^*$, $\pm J^*$, $\pm K^*$ образуют рациональную над $K$ группу, подобную $\mathfrak{D}$.\footnote{Характеризацию полей расщепления можно легко получить также с помощью фактор-систем.}
% END UNIT BODY
% END PRODUCER UNIT 182
\endgroup


\clearpage
\begingroup
\setcounter{footnote}{0}
% BEGIN PRODUCER UNIT 183
% Source: C:/Users/Floris/Documents/interlanguage/03_projects/noether/02_slavic_working_corpus/translations/paper32/section03_cyclic_quaternion_fields/russian/v001/Noether_Paper32_Section03_Cyclic_Quaternion_Fields_Russian_v001.tex
% Identity: 5829 B / SHA-256 5B8090E721CD0475418533566DCB2C433447001F9B766859C6952E48813E554E
\newcommand{\mS}{\mathfrak{S}}
\newcommand{\mA}{\mathfrak{A}}
\newcommand{\mat}[1]{\begin{pmatrix}#1\end{pmatrix}}
\setlength{\parindent}{1.25em}
\setlength{\parskip}{0pt}
\makeatletter
\renewcommand\footnotesize{\@setfontsize\footnotesize{7.6}{8.6}\selectfont}
\makeatother
\emergencystretch=100em
\allowdisplaybreaks
\sloppy
\hbadness=10000
\vbadness=10000
\hfuzz=2pt
\vfuzz=10pt
\raggedbottom
% BEGIN UNIT BODY
% Unit: Paper32 section 3 cyclic quaternion fields v001. Direct Russian translation from the R122 German Paper32/Paper33 source slice, lines 482--516 / segments P32-S0023--P32-S0029.
\noindent Далее следует показать, что для бесконечно многих $n$ существуют циклические поля степени $2^n$, являющиеся минимальными полями расщепления группы кватернионов.

\medskip

\noindent Пусть $p$ -- рациональное простое число, для которого при некотором $r>0$
\[
        (*)\qquad 2^r+1\equiv0\pmod p
\]
выполняется; пусть $2^n$ есть наивысшая степень числа $2$, делящая $p-1$.

Сначала покажем, что для бесконечно многих значений $n$ существуют простые числа $p$. Для этого пусть $k$ -- произвольное положительное целое рациональное число, а $p$ -- простой делитель числа $2^{2^k}+1$. Тогда требуемая для $p$ конгруэнция выполнена при $r=2^k$. Кроме того, $2\pmod p$ имеет порядок $2^{k+1}$, и поэтому для наивысшей степени $2^n$ числа $2$, делящей $p-1$, число $n$ больше $k$. Во всяком случае, значит, существуют сколь угодно большие $n$, которым соответствуют простые числа $p$.

Пусть $\varepsilon$ -- первообразный $p$-й корень из единицы. Теперь покажем, что в $P(\varepsilon)$ число $-1$ представляется как сумма двух квадратов
\[
        -1=a^2+b^2=(a+ib)(a-ib)\qquad (a,b\text{ числа из }P(\varepsilon)).
\]
Это равносильно тому, что в поле $(4p)$-х корней из единицы $P(\varepsilon,i)=P(\varepsilon\cdot i)$ существуют числа, относительная норма которых относительно $P(\varepsilon)$ равна как раз $-1$.

Число $1+i\varepsilon^\nu$ $(\nu=1,2,\ldots,p-1)$ имеет относительную норму $(1+i\varepsilon^\nu)(1-i\varepsilon^\nu)=1+\varepsilon^{2\nu}$; следовательно, все числа $1+\varepsilon^\nu$ $(\nu=1,2,\ldots,p-1)$ являются относительными нормами чисел из $P(\varepsilon i)$ относительно $P(\varepsilon)$ как основного поля. Поэтому то же самое верно и для
\[
\Pi=\prod_{\nu=0}^{r-1}(1+\varepsilon^{2^\nu})
=(1+\varepsilon)(1+\varepsilon^2)(1+\varepsilon^4)\cdots(1+\varepsilon^{2^{r-1}})
=\sum_{\lambda=0}^{2^r-1}\varepsilon^\lambda.
\]
Если к $\Pi$ добавить еще $\varepsilon^{2^r}=\varepsilon^{-1}=\varepsilon^{p-1}$ (в силу $(*)$), то $\Pi+\varepsilon^{p-1}$ становится суммой целых частей выражения $1+\varepsilon+\varepsilon^2+\cdots+\varepsilon^{p-1}=0$, и потому
\[
        \Pi=-\varepsilon^{-1}.
\]
Теперь и $\varepsilon$ является относительной нормой числа из $P(\varepsilon i)$, а именно числа $\varepsilon^{(p+1)/2}$. Следовательно, аналогичное верно для
\[
        \Pi\varepsilon=-1.
\]
Поэтому $-1$ в $P(\varepsilon)$ представим как сумма двух квадратов;\footnote{Для простых чисел вида $8n\pm3$ это рассуждение уже имеется у E. Landau, Über die Darstellung definiter Funktionen durch Quadrate, Math. Ann. Bd. 62, с. 272--285.} значит, $P(\varepsilon)$ является полем расщепления для группы кватернионов.

Пусть $K$ -- циклическое поле степени $2^n$, содержащееся в $P(\varepsilon)$; мы утверждаем, что и $K$ является полем расщепления.\footnote{Выбор циклического поля расщепления степени $2^n$ как подполя поля деления круга повторяет пример Хассе.} Пусть $m$ -- индекс тела кватернионов относительно $K$ как основного поля; тогда $m=1$ или $m=2$. Но поскольку существует поле расщепления $P(\varepsilon)$ нечетной относительной степени $(p-1)/2^n$ относительно $K$, второй случай невозможен; значит, $m=1$. Это и означает, что $K$ является полем расщепления.

Итак, для бесконечно многих $n$ существуют циклические поля степени $2^n$, являющиеся минимальными полями расщепления.
% END UNIT BODY
% END PRODUCER UNIT 183
\endgroup


\clearpage
\begingroup
\setcounter{footnote}{0}
% BEGIN PRODUCER UNIT 184
% Source: C:/Users/Floris/Documents/interlanguage/03_projects/noether/02_slavic_working_corpus/translations/paper33/complete/russian/v001/Noether_Paper33_Complete_Russian_v001.tex
% Identity: 14009 B / SHA-256 4BC3CE044C3FFA038734EAD16712A5A61D8FCF9F00704A036639C2A741246546
\setlength{\parindent}{1.15em}
\setlength{\parskip}{0pt}
\makeatletter
\renewcommand\footnotesize{\@setfontsize\footnotesize{7.5}{8.5}\selectfont}
\makeatother
\emergencystretch=120em
\allowdisplaybreaks
\sloppy
\hbadness=10000
\vbadness=10000
\hfuzz=2pt
\vfuzz=10pt
\raggedbottom
% BEGIN UNIT BODY
% Unit: Paper33 complete v001. Direct Russian translation from refreshed Zenodo R122 Paper32/Paper33 source slice, Paper33 lines 525--569 / segments P33-S0001--P33-S0010. The older local Paper33 final-audited slice is superseded for this unit because it omits the R122/source-scan opening of the two-problem paragraph.
\section*{33. Гиперкомплексные величины и теория представлений в арифметическом понимании}
\emph{Atti Congresso Bologna 2 (1928), с. 71--73.}

\begin{center}
\textsc{Эмми Нётер} (Гёттинген -- Германия)\\[-0.1em]
\rule{3em}{0.4pt}\\[0.8em]
\textsc{Гиперкомплексные величины и теория представлений}\\
\textsc{в арифметическом понимании}\footnote{Подробное изложение должно появиться в \emph{Math. Zeitschr.}; в продолжении там будет рассмотрена также некоммутативная теория Галуа. Ср. также о «полях расщепления»: R. Brauer и E. Noether, \emph{Ber. Berl. Ak.} 1927, а об общих структурных теоремах -- доклад G. Köthe в этих \emph{Atti}.}
\end{center}

\subsection*{Введение}

Я хочу показать, как теорию представлений -- в частности, представления групп -- можно понимать как теорию классов модулей и идеалов; и как последние, в свою очередь, подчиняются обобщенному понятию группы, а именно группам с операторами.\footnote{Группы с операторами восходят к W. Krull (\emph{Math. Zeitschr.} 23) и O. Schmidt (\emph{Math. Zeitschr.} 29). Приведенная здесь трактовка задачи редукции заданного представления также принадлежит -- в ограничении на коммутативные тела представления -- W. Krull (\emph{Theorie und Anwendung der verallgemeinerten Abelschen Gruppen}, Heidelberger Ber. 1926).}

Под \emph{представлением группы} -- в заданном поле $P$ --, как известно, понимают следующее: элементам $a,b,\ldots$ группы $G$ сопоставляются матрицы $A,B,\ldots$ с коэффициентами из $P$ так, что произведению $ab$ соответствует произведение $AB$; система матриц становится гомоморфным образом группы. Вместе с этим представлением дано также представление «группового кольца», которое состоит из всех «групповых чисел», то есть всех линейных комбинаций $a\alpha+b\beta+\cdots+c\gamma$, где $a,b,\ldots,c$ пробегают все комбинации из конечного числа элементов группы, а $\alpha,\beta,\ldots,\gamma$ соответственно -- из конечного числа элементов из $P$; при этом $a,b,\ldots,c$ рассматриваются как линейно независимые, а умножение задается групповым умножением и правилами вычисления. Представление группового кольца одновременно гомоморфно относительно сложения; тем самым представления групп подчиняются общей задаче \emph{представления колец}, где требуется гомоморфность относительно сложения и умножения.

Здесь существенны две задачи. Первая -- это \emph{редукция заданного} представления и вопрос об однозначности возникающих при этом неприводимых составляющих. Вторая состоит в вопросе о совокупности всех возможных представлений -- быть может, только неприводимых -- заданного кольца в заданном, не обязательно коммутативном теле.

Первая задача намного проще; это видно уже из того, что ни для представляемого кольца, ни для -- вообще говоря, некоммутативного -- тела представления не нужны никакие специальные предположения. Сам факт, что представление задается матрицами фиксированной степени, дает все необходимые предположения конечности. Задача полностью решается посредством теории групп с операторами, на которой я поэтому кратко остановлюсь.

Группа $\mathfrak{M}$ называется \emph{группой с областью операторов}, если одновременно задана система символов $\Theta,H,\ldots$ -- операторов -- такая, что выражения $\Theta(a),H(a),\ldots$ для каждого $a$ из $\mathfrak{M}$ однозначно определены, дают элемент из $\mathfrak{M}$ и удовлетворяют дистрибутивному закону: $\Theta(ab)=\Theta(a)\Theta(b)$. Две группы с одной и той же областью операторов называются операторно-изоморфными, если они изоморфны в обычном смысле и, кроме того, из соответствия $a$ и $\bar a$ следует также соответствие $\Theta(a)$ и $\Theta(\bar a)$, $H(a)$ и $H(\bar a)$ и так далее. Если затем в качестве подгрупп допускаются только такие, которые сами допускают область операторов, так что группа называется простой, если у нее нет допустимого собственного нормального делителя, кроме единицы, то теорема Жордана--Гёльдера о композиционном ряде сохраняется в форме: если группа с областью операторов вообще имеет композиционный ряд, то система фактор-групп в смысле операторного изоморфизма однозначно определена с точностью до порядка. При том же предположении верна и теорема о разложении, однозначном в смысле операторного изоморфизма, на прямо неразложимые факторы.

Связь с задачей представлений состоит в том, что к группам с операторами относятся, в частности, идеалы и модули: они рассматриваются как абелевы группы относительно сложения, тогда как операторы задаются умножением на элементы кольца. Всякое представление порождается модулем представления, то есть модулем относительно кольца и тела представления. Точнее, всякий класс модулей, то есть класс операторно-изоморфных между собой модулей представления, порождает одно и то же представление. Тем самым вопрос об однозначности редукции представления решается теоремой о композиционном ряде и теоремой о прямом произведении. Композиционный ряд, примененный к модулю представления, дает для каждой матрицы редукцию вида
\[
\begin{pmatrix}
B_1&0&0\\
&B_2&0\\
A_{ik}&\cdots&B_r
\end{pmatrix},
\]
где диагональные матрицы, соответствующие фактор-группам, порождают «неприводимое» представление, то есть такое, которое само уже не допускает такой редукции.

Но поскольку классы этих фактор-групп однозначно определены, то же верно для диагональных представлений, однозначно в смысле эквивалентности представлений. Аналогично теорема о прямом произведении дает однозначность «распада» на «неразложимые» составляющие:
\[
\begin{pmatrix}
C_1&&&\\
&C_2&&\\
&&\cdots&\\
&&&C_s
\end{pmatrix},
\]
причем каждый $C_i$ еще может быть однозначно редуцирован по теореме о композиционном ряде.

\emph{Вторую задачу} я хотел бы наметить для случая представлений гиперкомплексных систем, более общо -- колец, удовлетворяющих двойному условию цепей. Результат таков: всякий неприводимый (простой) класс модулей является классом идеалов; следовательно, композиционные ряды из односторонних идеалов уже через свои фактор-группы дают все неприводимые представления. Точнее, достаточно композиционных рядов в кольце классов вычетов по радикалу, которое становится «вполне приводимым» кольцом. Тем самым задача сведена к исследованию структуры вполне приводимых колец; и потому оказывается соответствующим задаче брать в качестве исходного пункта представления в телах автоморфизмов. Это надо понимать следующим образом. Все простые правые идеалы двусторонне простой компоненты такого кольца принадлежат одному и тому же классу и потому имеют одно и то же кольцо автоморфизмов; в силу простоты идеалов оно становится телом и одновременно является кольцом автоморфизмов простых левых идеалов. Само простое кольцо становится изоморфным системе всех матриц над телом автоморфизмов; и из этого представления возникают все представления относительно подтела, в частности относительно области коэффициентов гиперкомплексной системы, если сами элементы тела автоморфизмов заменить соответствующими матрицами. Так известная для гиперкомплексных систем структурная теорема Веддерберна здесь переосмысляется через понятие тела автоморфизмов, происходящее из общей теории групп, и одновременно распознается как источник теории представлений гиперкомплексных систем. Возникают новые вопросы теории Галуа некоммутативных тел, тесно связанные с вопросом о «полях расщепления». Вместе с тем открывается перспектива структурных теорем для более общих, не вполне приводимых колец посредством колец автоморфизмов неразложимых идеалов. Теория групп, в своем расширении на самые общие группы с операторами, и здесь снова оказывается упорядочивающим принципом.
% END UNIT BODY
% END PRODUCER UNIT 184
\endgroup


\clearpage
\begingroup
\setcounter{footnote}{0}
% BEGIN PRODUCER UNIT 185
% Source: C:/Users/Floris/Documents/interlanguage/03_projects/noether/02_slavic_working_corpus/translations/paper34/introduction/russian/v001/Noether_Paper34_Introduction_Russian_v001.tex
% Identity: 19038 B / SHA-256 627827BB8D48D1B459840B5CD1464BA7B4CCD99DE4F5F17D714B42D63AB2896E
\setlength{\parindent}{1.15em}
\setlength{\parskip}{0pt}
\makeatletter
\renewcommand\footnotesize{\@setfontsize\footnotesize{7.5}{8.5}\selectfont}
\makeatother
\emergencystretch=120em
\allowdisplaybreaks
\sloppy
\hbadness=10000
\vbadness=10000
\hfuzz=2pt
\vfuzz=10pt
\raggedbottom
% BEGIN UNIT BODY
% Unit: Paper34 introduction v001. Direct Russian translation from the audited German Paper34 slice, segments P34-S0001--P34-S0014, checked against the source scan pages 1--5 and Zenodo record 20822156 freshness snapshot 2026-06-24T05:21Z.
\setcounter{footnote}{0}

\section*{34. Гиперкомплексные величины и теория представлений}
\emph{Math. Zs. 30 (1929), с. 641--692.}

\subsection*{Введение}

Важнейшие общие теоремы о гиперкомплексных системах восходят к Молиену (\emph{Math. Annalen}, т. 41, и дорпатские отчеты 1897 г.). По существу независимо от этого Фробениус вскоре затем развил теорию гиперкомплексных систем и их представлений, в частности теорию представлений конечных групп, как единую теорию. Основу составляет восходящее к Дедекинду понятие группового детерминанта, а в более общем виде -- детерминанта произвольной гиперкомплексной системы. Фробениус показывает, что различным неприводимым множителям этого детерминанта (областью коэффициентов всегда считается поле всех комплексных чисел) соответствуют различные классы неприводимых представлений, и что таким образом исчерпываются все классы неприводимых представлений. Такой класс представлений полностью характеризуется своей “системой характеров”; в случае конечных групп эта система характеров возникает путем разложения детерминанта той коммутативной гиперкомплексной системы, которая выводится из классов сопряженных элементов группы. Этот детерминант разлагается на линейные множители, с характерами в качестве коэффициентов, то есть является прямым обобщением результата, впервые полученного Дедекиндом: групповой детерминант конечной абелевой группы разлагается на линейные множители, коэффициентами которых служат различные характеры абелевой группы (переписка с Фробениусом).\footnote{Дальнейший текст представляет собой свободную разработку моей лекции зимнего семестра 1927/28 г., подготовленную Б. Л. ван дер Варденом. Текст к печати мы готовили совместно. Я также обязана Б. Л. ван дер Вардену благодарностью за ряд критических замечаний.}

Однако эти понятийно простые и прозрачные результаты у Фробениуса получаются посредством трудоемкого вычисления. Дальнейшее развитие было направлено на упрощенный вывод этих результатов и одновременно на их распространение при выборе произвольного тела в качестве области коэффициентов. Это дальнейшее развитие для гиперкомплексных систем и для теории представлений пошло совершенно раздельными путями. Гиперкомплексные системы получили арифметическую трактовку у Веддерберна: всякая гиперкомплексная система единственным образом является прямой суммой двусторонне прямо неразложимых колец; для систем без радикала эти неразложимые кольца становятся изоморфными системе всех $n$-рядных матриц с элементами из присоединенного, не обязательно коммутативного тела, определенного по существу единственным образом.\footnote{Ср. библиографические указания у Dickson, \emph{Algebras and their Arithmetics} (немецкое издание: \emph{Algebren und ihre Zahlentheorie}, Zürich 1927).} Теория представлений получила элементарное обоснование у Бернсайда и И. Шура, которые, независимо от гиперкомплексной системы, исходили непосредственно из заданного представления и оперировали матричными теоремами. В частности, Шур рассмотрел вопрос о числовых полях наименьшей степени, в которых представление, неприводимое над заданным телом, абсолютно разлагается. Эти абсолютно неприводимые составляющие принадлежат конечному числу сопряженных классов представлений. Степень искомых числовых полей равна произведению числа этих классов на индекс (число составляющих в каждом из сопряженных классов). Главным вспомогательным средством теории Шура служит система всех матриц, перестановочных с неприводимым представлением.\footnote{Ср. библиографические указания в 10-й главе теории групп Шпайзера.}

В следующем чисто арифметическом обосновании гиперкомплексные величины и теория представлений снова выступают как единое целое: как специальный случай общей теории некоммутативных колец, удовлетворяющих лишь некоторым условиям конечности. А именно, речь идет о теории классов модулей и идеалов относительно этих колец с главным результатом, что неприводимые классы модулей уже исчерпываются соответствующими классами идеалов; в частности, для колец без радикала все классы модулей распадаются на неприводимые (становятся вполне приводимыми). Это арифметический эквивалент результата Фробениуса о том, что неприводимые множители регулярного группового (соответственно системного) детерминанта уже исчерпывают совокупность классов неприводимых представлений, и что для систем без радикала имеет место полная приводимость представлений. Рассмотрение системного детерминанта и его разложения можно именно понимать как переход к норме, тогда как здесь непосредственно рассматриваются разложение идеалов в прямую сумму и их композиционные ряды. Теория представлений же, благодаря понятию модуля представления, становится теорией классов модулей.

Введение классов модулей и идеалов имеет чисто группово-теоретическое обоснование. Модули и идеалы рассматриваются как абелевы группы относительно сложения, ограниченные тем, что они допускают некоторые умножения на элементы кольца: они образуют “группы с операторами” (§1). Для таких групп вместо обычного изоморфизма появляется “операторный изоморфизм” (§2); группы, операторно-изоморфные между собой (в пределах фиксированной области), объединяются в один класс -- понятие класса, которое, например, в случае идеалов числового поля совпадает с обычным. Теория групп с операторами восходит к W. Krull и O. Schmidt (ср. прим. 6); в главе I она развивается систематически. Теоремы, остающиеся справедливыми и здесь, о композиционных рядах, прямом произведении (соответственно сумме) и вполне приводимых группах -- теоремы единственности в смысле упомянутого выше деления на классы -- образуют основу всего дальнейшего.

Сначала они дают общие теоремы единственности неприводимых диагональных составляющих при произвольных представлениях (глава III, §16), на основании факта взаимно однозначного соответствия классов представлений классам модулей представления, то есть классам групп с операторами. Дальнейший вопрос о совокупности представлений требует более точного изучения структуры представляемых колец, следовательно, в специальном случае, гиперкомплексных систем. В главе II результаты Веддерберна получаются заново и развиваются дальше на основе группово-теоретического понимания, согласно которому на переднем плане стоит рассмотрение односторонних идеалов. При этом оказывается, что “теорема о множественных цепях” для правых идеалов, или тождественное с ней “минимальное условие” (в каждом множестве правых идеалов имеется по меньшей мере один минимальный элемент внутри этого множества), достаточна как условие конечности.\footnote{Методы доказательства Веддерберна переносятся, если предполагается “двойное условие цепей”. Ср. E. Artin, \emph{Zur Theorie der hyperkomplexen Zahlen}, Hamb. Abh. 1927. Ср. также A. Suschkewitsch, \emph{Über die endlichen Gruppen ohne das Gesetz der eindeutigen Umkehrbarkeit}, Math. Annalen 99 (1928), с. 30--50, где для (конечных) областей только с одной ассоциативной, необратимой операцией развиваются параллельные структурные теоремы. Расщепление “ядра” у Сушкевича на правые и левые группы соответствует представлению двусторонне простого кольца с единицей в виде суммы правых и левых идеалов.} Получается тождество колец без радикала, удовлетворяющих минимальному условию, с (право) вполне приводимыми кольцами с единицей. В таких кольцах все простые правые идеалы двусторонне неразложимого кольца принадлежат одному и тому же классу и потому, как группы с операторами, имеют одно и то же кольцо автоморфизмов, которое из-за простоты идеалов становится телом. Это тело автоморфизмов класса и есть то, что осуществляет матричное представление, найденное Веддерберном.

Это матричное представление в теле автоморфизмов одновременно решает задачу представления во вполне приводимом случае, как только требуется представление относительно тела автоморфизмов или его подполей, в частности относительно области коэффициентов гиперкомплексной системы. Не существует, и это прямое следствие теоремы о сведении классов модулей к классам идеалов, никаких иных неприводимых представлений, кроме веддерберновского и тех, которые возникают из него, когда элементы тела автоморфизмов естественным образом заменяются матрицами над подполем.

Итак, во вполне приводимом случае существует столько различных классов неприводимых представлений, сколько классов идеалов; их число совпадает с числом различных двусторонне неразложимых, то есть простых, колец. Это число, наконец, снова тождественно числу неразложимых компонент центра, которые становятся коммутативными полями. В случае гиперкомплексных систем с алгебраически замкнутой областью коэффициентов эти последние полностью определяются различными гомоморфизмами колец (неприводимыми представлениями) центра; а это, с точностью до числового множителя, и есть характеры.

Полученные результаты одновременно дают обзор неприводимых представлений систем с радикалом, поскольку они сводятся к представлениям кольца классов вычетов по радикалу, которое становится системой без радикала.

Как будет показано позднее, для гиперкомплексных систем без радикала из тела автоморфизмов следует также теория “полей разложения”, то есть коммутативных расширений области коэффициентов, в которых представления распадаются на абсолютно неприводимые; иначе говоря, в которых имеет место разложение в прямую сумму абсолютно простых правых идеалов. Поля разложения тождественны полям разложения тела автоморфизмов, и все поля разложения наименьшей степени изоморфны максимальным коммутативным подполям тела автоморфизмов. Связь с упомянутыми выше исследованиями Шура устанавливается тем, что транспонированные матрицы, перестановочные с представлением, как раз дают представление тела автоморфизмов.\footnote{Ср. R. Brauer--E. Noether, \emph{Über minimale Zerfällungskörper irreduzibler Darstellungen}, Sitz.-Ber. Preuß. Akad. Wiss. 1927, с. 221.} Эти теоремы должны быть систематически развиты в рамках теории Галуа некоммутативных тел.

Положенные в основу понятия -- операторный гомоморфизм и кольцо автоморфизмов -- дают также структуру общих колец с радикалом, удовлетворяющих лишь минимальному условию; это будет изложено в другом месте.
% END UNIT BODY
% END PRODUCER UNIT 185
\endgroup


\clearpage
\begingroup
\setcounter{footnote}{0}
% BEGIN PRODUCER UNIT 186
% Source: C:/Users/Floris/Documents/interlanguage/03_projects/noether/02_slavic_working_corpus/translations/paper34/section01/russian/v001/Noether_Paper34_Section01_Russian_v001.tex
% Identity: 8862 B / SHA-256 474A5774E3F3803DE1FF82B1CB53C66AF9895E866AA6C1C0F0166EAC098DF6BD
\setlength{\parindent}{1.15em}
\setlength{\parskip}{0pt}
\makeatletter
\renewcommand\footnotesize{\@setfontsize\footnotesize{7.5}{8.5}\selectfont}
\makeatother
\providecommand{\mG}{\mathfrak{G}}
\providecommand{\mH}{\mathfrak{H}}
\providecommand{\mA}{\mathfrak{A}}
\providecommand{\mB}{\mathfrak{B}}
\providecommand{\mM}{\mathfrak{M}}
\providecommand{\mo}{\mathfrak{o}}
\providecommand{\OO}{\Omega}
\providecommand{\Th}{\Theta}
\emergencystretch=120em
\allowdisplaybreaks
\sloppy
\hbadness=10000
\vbadness=10000
\hfuzz=2pt
\vfuzz=10pt
\raggedbottom
% BEGIN UNIT BODY
% Unit: Paper34 section01 v001. Direct russian translation from audited German source lines 49--102, segments P34-S0015--P34-S0027. Source scan pages 5--7 are retained as witness, and Zenodo record 20822156 was checked at 2026-06-24T05:56Z.
\setcounter{footnote}{5}

\section*{Глава I. Группово-теоретические основания}

\subsection*{§ 1. Группы с операторами}

Пусть $\mG$ -- группа (конечная или бесконечная), с элементами $a,b,\ldots$. Под областью операторов $\Omega$ для группы $\mG$ понимается множество новых символов $H,\Theta,\ldots$ такое, что каждому $a$ из $\mG$ и каждому $\Theta$ из $\Omega$ однозначно сопоставлен элемент $\Theta a$ из $\mG$, и выполнен дистрибутивный закон:
\[
        \Theta(ab)=\Theta a\cdot\Theta b.
\]
Следовательно, каждый оператор определяет гомоморфизм группы в себя, причем группа отображается либо на себя, либо на подгруппу. Единичный элемент переходит в единичный элемент, обратный -- в обратный.

Область операторов называется \emph{абсолютной}, если различные операторы определяют также различные гомоморфизмы. Из произвольной области операторов получают абсолютную, отождествляя все элементы, порождающие один и тот же гомоморфизм. Абсолютная область операторов является взаимно однозначным образом подмножества множества всех гомоморфизмов группы в себя.

\emph{Допустимая подгруппа} $\mH$ группы $\mG$ -- относительно фиксированной области операторов $\Omega$ -- это такая подгруппа, которая допускает эту область операторов, то есть $\Theta a$ лежит в $\mH$ для каждого $a$ из $\mH$ и каждого $\Theta$ из $\Omega$.

\emph{Примеры.} 1. Пусть операторами являются внутренние автоморфизмы: $\Theta a=c^{-1}ac$. Допустимые подгруппы суть нормальные делители.

2. Пусть операторами являются все автоморфизмы. Допустимые подгруппы суть “характеристические подгруппы”, переходящие в себя при всех автоморфизмах.\footnote{Объясняемые здесь понятия восходят к W. Krull, \emph{Über verallgemeinerte endliche Abelsche Gruppen}, Math. Zeitschr. 23 (1925), с. 161--196, и O. Schmidt, \emph{Über unendliche Gruppen mit endlicher Kette}, Math. Zeitschr. 29 (1928), с. 34--41.}

3. Пусть $\mG$ -- кольцо, то есть абелева группа относительно сложения, в которой также определено умножение со свойствами
\[
        r(a+b)=ra+rb,\qquad (a+b)r=ar+br,\qquad ab\cdot c=a\cdot bc.
\]
Каждый элемент $r$ одновременно определяет два оператора: операторы $rx$ и $xr$. Допущенными подгруппами являются “идеалы” $\mathfrak a$, а именно: левые, допускающие операции $rx$: $r\mathfrak a\subseteq\mathfrak a$; правые, допускающие операции $xr$: $\mathfrak a r\subseteq\mathfrak a$; двусторонние, допускающие обе операции. Все идеалы становятся тривиальными $(=\{0\}\text{ или }=\mG)$, если кольцо $\mG$ является телом, то есть если $\mG-\{0\}$ является группой относительно умножения.\footnote{Следовательно, ни для колец, ни для тел не предполагается коммутативный закон умножения.}

4. \emph{Модули относительно кольца $\mo$.} Пусть $\mo$ -- кольцо, а $\mM$ -- абелева группа, записанная аддитивно. Пусть задано умножение $r\cdot a$ элементов из $\mo$ на элементы из $\mM$; произведение снова должно содержаться в $\mM$. Требуется:
\[
\left.
\begin{aligned}
 r(a+b)&=ra+rb,\\
 rs\cdot a&=r\cdot sa,\\
 (r+s)a&=ra+sa,
\end{aligned}
\right\}\qquad r,s\in\mo,\quad a,b\in\mM.
\]
Тогда $\mM$ называется $\mo$-модулем; точнее, поскольку множители из $\mo$ пишутся слева, левым $\mo$-модулем. Кольцо $\mo$ одновременно является областью операторов. Если она абсолютна, то есть если различные элементы кольца дают также различные операции, говорят об \emph{абсолютной области множителей} для модуля $\mM$.

Как и выше, от произвольной области множителей можно перейти к абсолютной, отождествляя элементы, порождающие одну и ту же операцию. Абсолютная область множителей снова является кольцом. Допустимые подгруппы модуля $\mM$ называются \emph{подмодулями}. Если специально $\mo=\mM$, то мы возвращаемся к левым идеалам.\footnote{Точно так же для правых модулей требуют:
\[
(a+b)r=ar+br,\qquad a\cdot rs=ar\cdot s,\qquad a(r+s)=ar+as.
\]}

5. \emph{Бимодули.} Если $\mM$ одновременно является левым $\mo$-модулем и правым $\mo'$-модулем, и, кроме того, для $a$ в $\mo$, $\alpha$ в $\mM$, $a'$ в $\mo'$ всегда
\[
        a\cdot\alpha a'=a\alpha\cdot a',
\]
то $\mM$ называется бимодулем.

6. \emph{Кольцо автоморфизмов абелевой группы.}\footnote{Ср. A. Châtelet, \emph{Les Groupes Abéliens finis}, Paris, Gauthier-Villars, 1925, p. 99.} Для гомоморфизмов в себя абелевой группы (записанной аддитивно) можно определить сложение и умножение формулами
\[
        (H+\Theta)a=Ha+\Theta a,\qquad (H\Theta)a=H(\Theta a).
\]
Определенные таким образом операции, очевидно, снова представляют собой гомоморфизмы и потому снова принадлежат этой системе. Система образует кольцо, и абелеву группу, согласно 4, можно рассматривать как модуль относительно этого кольца.

Под “подгруппами” ниже всегда понимаются допустимые подгруппы. Пересечение $\mA\cap\mB$ двух допустимых подгрупп снова является допустимой подгруппой. Так же и произведение $\mA\mB$, если по меньшей мере одна из двух подгрупп $\mA,\mB$ является нормальным делителем.
% END UNIT BODY
% END PRODUCER UNIT 186
\endgroup


\clearpage
\begingroup
\setcounter{footnote}{0}
% BEGIN PRODUCER UNIT 187
% Source: C:/Users/Floris/Documents/interlanguage/03_projects/noether/02_slavic_working_corpus/translations/paper34/section02/russian/v001/Noether_Paper34_Section02_Russian_v001.tex
% Identity: 9029 B / SHA-256 CBA1743A722F758F0CB2C0EEC4C3881890C65891EB4B18EBAD235FA2C0351386
\setlength{\parindent}{1.15em}
\setlength{\parskip}{0pt}
\makeatletter
\renewcommand\footnotesize{\@setfontsize\footnotesize{7.5}{8.5}\selectfont}
\makeatother
\providecommand{\mG}{\mathfrak{G}}
\providecommand{\mA}{\mathfrak{A}}
\providecommand{\mB}{\mathfrak{B}}
\providecommand{\mN}{\mathfrak{N}}
\providecommand{\mM}{\mathfrak{M}}
\providecommand{\mo}{\mathfrak{o}}
\providecommand{\Th}{\Theta}
\providecommand{\iso}{\simeq}
\emergencystretch=120em
\allowdisplaybreaks
\sloppy
\hbadness=10000
\vbadness=10000
\hfuzz=2pt
\vfuzz=10pt
\raggedbottom
% BEGIN UNIT BODY
% Unit: Paper34 section02 v001. Direct russian translation from audited German source lines 104--142, segments P34-S0028--P34-S0039. Source scan pages 7--8 are retained as witness, and Zenodo record 20822156 was checked at 2026-06-24T06:38Z.
\setcounter{footnote}{9}

\subsection*{§ 2. Теоремы об изоморфизме}

Отображение группы $\mG$ на группу $\overline{\mG}$ -- причем $\mG$ и $\overline{\mG}$ должны иметь одну и ту же область операторов -- называется операторным гомоморфизмом, если, во-первых, оно является гомоморфизмом в обычном смысле ($ab\mapsto \bar a\bar b$), и если, во-вторых, как только $a$ отображается в $\bar a$, элемент $\Theta a$ отображается в $\Theta\bar a$.\footnote{Понятие операторного гомоморфизма можно расширить так, чтобы группы $\mG,\overline{\mG}$ имели отдельные области операторов; тогда гомоморфизм отображает не только $\mG$ на $\overline{\mG}$, но и области операторов друг на друга так, что если $a$ отображается в $\bar a$, а $\Theta$ -- в $\overline\Theta$, то $\overline\Theta\bar a$ является образом $\Theta a$. В такой общей форме понятие операторного гомоморфизма включает также понятие гомоморфизма колец; ср. §8.} Обозначение: $\mG\sim\overline{\mG}$. Если соответствие взаимно однозначно, оно называется операторным изоморфизмом. Обозначение: $\mG\simeq\overline{\mG}$. Как и для обычных групп, доказывается теорема о гомоморфизме:

\emph{Если $\overline{\mG}$ является операторно-гомоморфным образом $\mG$, то $\overline{\mG}$ операторно изоморфна фактор-группе $\mG/\mN$, где $\mN$ -- допустимый нормальный делитель, состоящий из всех элементов $\mG$, которым соответствует единица в $\overline{\mG}$. Обратно, каждый допустимый нормальный делитель $\mN$ определяет фактор-группу $\mG/\mN$, которая допускает область операторов $\mG$ и является операторно-гомоморфным образом $\mG$.}

Группа называется \emph{простой}, если она не имеет нормальных делителей, кроме самой себя и единичной группы. Каждый гомоморфизм простой группы либо является изоморфизмом, либо сопоставляет каждому элементу единичный элемент.

\emph{Первая теорема об изоморфизме.} \emph{Пусть $\overline{\mG}$ -- гомоморфный образ $\mG$, пусть $\overline{\mA}$ -- нормальный делитель в $\overline{\mG}$, а $\mA$ -- совокупность элементов $\mG$, которым соответствуют элементы из $\overline{\mA}$. Тогда $\mA$ снова является нормальным делителем, и}
\[
        \overline{\mG}/\overline{\mA}\simeq \mG/\mA. \tag{1}
\]
\emph{Доказательство.} $\mG\sim\overline{\mG}$, $\overline{\mG}\sim\overline{\mG}/\overline{\mA}$, значит $\mG\sim\overline{\mG}/\overline{\mA}$. Поэтому $\overline{\mG}/\overline{\mA}$ изоморфна фактор-группе в $\mG$; соответствующий нормальный делитель есть совокупность элементов, которым соответствует единица в $\overline{\mG}/\overline{\mA}$; это как раз $\mA$.

\emph{Дополнение.} Если по теореме о гомоморфизме положить $\overline{\mG}=\mG/\mN$, то $\mA$ непременно содержит $\mN$. Из $\mA$ можно восстановить $\overline{\mA}$: $\overline{\mA}=\mA/\mN$. Соответствие $\mA\leftrightarrow\overline{\mA}$ взаимно однозначно. Формулу (1) можно записать также так:
\[
        (\mG/\mN)/(\mA/\mN)\simeq \mG/\mA.
\]

\emph{Вторая теорема об изоморфизме.} \emph{Пусть $\mA$ -- подгруппа, а $\mB$ -- нормальный делитель в $\mG$. Тогда $\mA\cap\mB$ является нормальным делителем в $\mA$, и}
\[
        \mA\mB/\mB\simeq \mA/(\mA\cap\mB).
\]
\emph{Доказательство.} При гомоморфизме $\mG\sim\mG/\mB$ элементам $a$ из $\mA$ в частности сопоставляются некоторые классы смежности $a\mB$, которые вместе составляют группу $\mA\mB/\mB$. Следовательно, $\mA\sim\mA\mB/\mB$, откуда утверждение следует из теоремы о гомоморфизме.

Операторно-гомоморфное отображение группы $\mG$ на саму себя или на подгруппу называется эндоморфизмом группы $\mG$. Если $\mG$ специально является абелевой группой (с операторами), записанной аддитивно, и если сумма и произведение гомоморфизмов определяются, как выше (§1, пример 6), то операторные гомоморфизмы группы в себя образуют кольцо, кольцо автоморфизмов.

\emph{Кольцо автоморфизмов простой абелевой группы (с операторами) является телом.}

\emph{Доказательство.} Каждый гомоморфизм отображает группу либо на саму себя, либо на нулевую группу (так как других допустимых подгрупп нет). Гомоморфизмы, которые не отображают всё в нуль, согласно сказанному выше о простых группах являются изоморфизмами, а изоморфные отображения группы на себя образуют группу. Значит, после удаления нулевого оператора кольцо становится группой, а потому само кольцо является телом.

Если специально $\mG$ является левым $\mo$-модулем и операторные гомоморфизмы записываются как правые операторы, то $\mG$ становится \emph{бимодулем относительно $\mo$ слева и кольца автоморфизмов справа}. Действительно, тот факт, что новые операторы $\Gamma$ были выбраны как гомоморфизмы, выражается формулами
\[
        (a+b)\Gamma=a\Gamma+b\Gamma,
        \qquad ra\cdot\Gamma=r\cdot a\Gamma,
\]
которые (вместе с остальными, уже истолкованными раньше) характеризуют бимодуль.

Обратно, если задан бимодуль, то те же самые формулы выражают, что каждый элемент правой области множителей индуцирует операторный гомоморфизм относительно левых операторов, и наоборот.
% END UNIT BODY
% END PRODUCER UNIT 187
\endgroup


\clearpage
\begingroup
\setcounter{footnote}{0}
% BEGIN PRODUCER UNIT 188
% Source: C:/Users/Floris/Documents/interlanguage/03_projects/noether/02_slavic_working_corpus/translations/paper34/section03/russian/v001/Noether_Paper34_Section03_Russian_v001.tex
% Identity: 6260 B / SHA-256 899812541936EF6487F1B020C2EC4D99DEBE220F76E01C572917F288C7196A2D
\setlength{\parindent}{1.15em}
\setlength{\parskip}{0pt}
\makeatletter
\renewcommand\footnotesize{\@setfontsize\footnotesize{7.5}{8.5}\selectfont}
\makeatother
\providecommand{\mG}{\mathfrak{G}}
\providecommand{\mA}{\mathfrak{A}}
\providecommand{\mB}{\mathfrak{B}}
\providecommand{\mE}{\mathfrak{E}}
\providecommand{\mH}{\mathfrak{H}}
\providecommand{\mo}{\mathfrak{o}}
\providecommand{\iso}{\simeq}
\emergencystretch=120em
\allowdisplaybreaks
\sloppy
\hbadness=10000
\vbadness=10000
\hfuzz=2pt
\vfuzz=10pt
\raggedbottom
% BEGIN UNIT BODY
% Unit: Paper34 section03 v001. Direct russian translation from audited German source lines 144--182, segments P34-S0040--P34-S0050. Source scan pages 9--10 retained as witness, and Zenodo record 20822156 was checked at 2026-06-24T07:09Z.
\setcounter{footnote}{10}

\subsection*{§ 3. Композиционные ряды}

Если в $\mG$ имеется конечный ряд допустимых подгрупп
\[
        \mG>\mA_1>\mA_2>\cdots>\mA_r=\mE \tag{2}
\]
($\mE$ есть группа, состоящая только из единицы $e$), такой, что каждая $\mA_i$ является нормальным делителем в предыдущем члене, и между двумя последовательными членами ряда уже невозможно вставить еще один член с тем же свойством, то этот ряд называется композиционным рядом.

Фактор-группы $\mG/\mA_1,\ \mA_1/\mA_2,\ldots,\ \mA_{r-1}/\mE$ называются композиционными факторами. Они являются простыми группами, то есть не имеют допустимых нормальных делителей, кроме самих себя и единичной группы.

Если среди операторов взять, в частности, все внутренние автоморфизмы, то ряд (2) состоит только из нормальных делителей и называется главным рядом. Если взять все автоморфизмы, он называется характеристическим рядом. В дальнейших примерах §1 речь шла бы о композиционных рядах идеалов в $\mo$, соответственно о композиционных рядах модулей или бимодулей.

\emph{Теорема Жордана--Гельдера.} \emph{Если для группы $\mG$ существуют два различных композиционных ряда:}
\[
        \mG>\mA_1>\mA_2>\cdots>\mA_r=\mE,
        \qquad
        \mG>\mB_1>\mB_2>\cdots>\mB_s=\mE,
\]
\emph{то они имеют одну и ту же длину: $r=s$, и композиционные факторы}
\[
        \mG/\mA_1,\ \mA_1/\mA_2,\ldots,\mA_{r-1}/\mE
\]
\emph{изоморфны в некотором порядке}
\[
        \mG/\mB_1,\ \mB_1/\mB_2,\ldots,\mB_{s-1}/\mE.
\]
Доказательство см., например, у E. Noether, \emph{Abstrakter Aufbau} и т. д., Math. Annalen 96 (1926), §10, с. 57. Там одновременно доказывается:

\emph{Если в $\mG$ имеется композиционный ряд, то через каждый нормальный делитель $\mH$ группы $\mG$ можно провести композиционный ряд.}

Существование композиционного ряда ясно для конечных групп, а в более общем виде для тех групп, в которых выполняются условие максимальности и условие минимальности:

\emph{Условие максимальности} означает, что в каждом множестве подгрупп существует максимальная подгруппа, то есть такая, которая уже не содержится ни в какой другой подгруппе этого множества. Эквивалентно этому: всякая возрастающая цепь подгрупп $\mA_1\subset\mA_2\subset\mA_3\cdots$ обрывается после конечного числа членов.

\emph{Условие минимальности} означает, что в каждом множестве подгрупп существует минимальная подгруппа, которая уже не содержит никакой другой подгруппы этого множества; или, что то же самое, всякая убывающая цепь $\mA_1\supset\mA_2\supset\mA_3\cdots$ обрывается после конечного числа членов.

\emph{Если в качестве подгрупп допускаются только нормальные делители} (например, для абелевых групп), \emph{то из существования композиционного ряда, обратно, следуют условие максимальности и условие минимальности.}

\emph{Доказательство.} Каждый нормальный делитель обладает композиционным рядом; его длина называется длиной нормального делителя. Если $\mA\subset\mB$, то длина $\mA$ меньше длины $\mB$, поскольку через $\mA$ можно провести композиционный ряд для $\mB$. Следовательно, всякая подгруппа наименьшей длины одновременно минимальна, а всякая подгруппа наибольшей длины одновременно максимальна.
% END UNIT BODY
% END PRODUCER UNIT 188
\endgroup


\clearpage
\begingroup
\setcounter{footnote}{0}
% BEGIN PRODUCER UNIT 189
% Source: C:/Users/Floris/Documents/interlanguage/03_projects/noether/02_slavic_working_corpus/translations/paper34/section04/russian/v001/Noether_Paper34_Section04_Russian_v001.tex
% Identity: 10208 B / SHA-256 FBF27E1B6F0098621E226AD9A862DCC1A6E35EB43D7D142BFAB0EC96D35C9EC7
\setlength{\parindent}{1.15em}
\setlength{\parskip}{0pt}
\makeatletter
\renewcommand\footnotesize{\@setfontsize\footnotesize{7.5}{8.5}\selectfont}
\makeatother
\providecommand{\mG}{\mathfrak{G}}
\providecommand{\mA}{\mathfrak{A}}
\providecommand{\mB}{\mathfrak{B}}
\providecommand{\mE}{\mathfrak{E}}
\providecommand{\mH}{\mathfrak{H}}
\providecommand{\mC}{\mathfrak{C}}
\providecommand{\mD}{\mathfrak{D}}
\providecommand{\mR}{\mathfrak{R}}
\providecommand{\mS}{\mathfrak{S}}
\providecommand{\mT}{\mathfrak{T}}
\providecommand{\mo}{\mathfrak{o}}
\providecommand{\iso}{\simeq}
\emergencystretch=120em
\allowdisplaybreaks
\sloppy
\hbadness=10000
\vbadness=10000
\hfuzz=2pt
\vfuzz=10pt
\raggedbottom
% BEGIN UNIT BODY
% Unit: Paper34 section04 v001. Direct russian translation from audited German source lines 184--253, segments P34-S0051--P34-S0065. Source scan pages 10--12 retained as witness, and Zenodo record 20822156 was checked at 2026-06-24T07:42Z.
\setcounter{footnote}{10}

\subsection*{§ 4. Прямые произведения и пересечения}

Группа $\mG$ называется прямым произведением двух факторов, $\mG=\mA\times\mB$, если
\[
\begin{array}{ll}
1.& \mA,\mB\text{ являются нормальными делителями в }\mG,\\
2.& \mA\mB=\mG,\\
3.& \mA\cap\mB=\mE.
\end{array}
\]
\emph{Это определение, очевидно, эквивалентно следующему:} каждый элемент $g$ группы $\mG$ однозначно представляется в виде $g=ab$, где $a$ принадлежит $\mA$, $b$ принадлежит $\mB$, и элементы $\mA$ коммутируют с элементами $\mB$: $ab=ba$.

Группа $\mG$ называется прямым произведением $n$ факторов, $\mG=\mA_1\times\cdots\times\mA_n$, если, положив $\mB_i=\mA_1\cdots\mA_{i-1}\mA_{i+1}\cdots\mA_n$, имеем, что $\mG=\mA_i\times\mB_i$ является прямым произведением для каждого $i$.

\emph{Определение снова эквивалентно следующему:} каждый $g$ из $\mG$ однозначно представляется как
\[
        g=a_1a_2\cdots a_n,\qquad a_i\in\mA_i,
\]
и элементы $\mA_i$ коммутируют с элементами $\mA_k$.

Часто используются следующие факты:
\begin{enumerate}
\item Если $\mA_1\times\cdots\times\mA_n=\mR$ и $\mR\times\mH=\mG$, то $\mA_1\times\cdots\times\mA_n\times\mH=\mG$.
\item Если $\mG=\mA_1\times\cdots\times\mA_n=\mC_1\times\cdots\times\mC_n$ и $\mC_i\subseteq\mA_i$, то $\mC_i=\mA_i$. (Это следует из представления элементов $\mA_i$ через $\mC_j$.)
\item Если $\mG=\mA\times\mB$ и $\mR\supseteq\mA$, то $\mR=\mA\times(\mR\cap\mB)$. (Это следует из представления элементов $\mR$ в форме $ab$, где второй множитель принадлежит как $\mR$, так и $\mB$.)
\item Если $\mG=\mA\times\mB$, то $\mA\sim\mG/\mB$ и $\mB\sim\mG/\mA$ (вторая теорема об изоморфизме). Отсюда далее следует:
\item Если $\mG=\mA\times\mB$, а $\mA$ и $\mB$ имеют композиционные ряды длин $m$ и $n$, то $\mG$ имеет композиционный ряд длины $m+n$.
\item Если $\mA\sim\overline{\mA}$ и $\mB\sim\overline{\mB}$, то $\mA\times\mB\sim\overline{\mA}\times\overline{\mB}$.
\end{enumerate}
Группа называется прямо неразложимой, если ее нельзя представить как прямое произведение факторов, отличных от $\mE$.

Ясно, что \emph{каждая группа с условием минимальности является прямым произведением конечного числа прямо неразложимых групп.}\footnote{Как показали W. Krull в коммутативном случае и O. Schmidt в общем случае (ср. прим. 6), при условии максимальности и минимальности представление единственно с точностью до изоморфизма. Так как, не изменяя понятия прямой неразложимости, можно включить все внутренние автоморфизмы в область операторов, достаточно требовать условий конечности для нормальных делителей или существования главного ряда.}

Если соответствующие группы записывать аддитивно, то понятия произведения и прямого произведения переходят в сумму и прямую сумму. Обозначение: для суммы групп $(\mA_1,\mA_2,\ldots)$, для прямой суммы $\mA_1+\mA_2+\cdots$.

Понятие "прямого пересечения" далее хотя и не используется, но следующие теоремы о связи между прямым пересечением и прямым произведением показывают, как теорию идеалов следующих глав можно формулировать также через пересечения вместо сумм; тем самым устанавливается связь с известными коммутативными теориями.

Группа $\mD$ называется прямым пересечением двух групп $\mR$ и $\mS$ относительно $\mG$, если
\[
\begin{array}{ll}
1.& \mR,\mS\text{ являются нормальными делителями в }\mG,\\
2.& \mR\cap\mS=\mD,\\
3.& \mR\mS=\mG.
\end{array}
\]
Точно так же $\mD$ называется прямым пересечением $n$ групп $\mS_1,\ldots,\mS_n$, если, положив $\mR_i=\mS_1\cap\cdots\cap\mS_{i-1}\cap\mS_{i+1}\cap\cdots\cap\mS_n$, имеем, что $\mD=\mR_i\cap\mS_i$ является прямым для каждого $i$.

Из определения следует, что $\mD$ должен быть нормальным делителем в $\mG$. Если $\mD=\mS_1\cap\cdots\cap\mS_n$ является прямым и при гомоморфизме $\mG\sim\mG/\mD$ группы $\mG,\mS,\mD$ переходят в $\overline{\mG},\overline{\mS},\mE$, то $\mE=\overline{\mS}_1\cap\cdots\cap\overline{\mS}_n$ становится прямым. Обратно, от каждого такого представления $\mE=\overline{\mS}_1\cap\cdots\cap\overline{\mS}_n$ возвращаются к прямому представлению $\mD=\mS_1\cap\cdots\cap\mS_n$. Благодаря этому соответствию теоремы о прямых пересечениях при произвольном $\mD$ сводятся к специальному случаю $\mD=\mE$.

В случае $\mD=\mE$ и двух факторов условия для прямого произведения и прямого пересечения совпадают. Для $n$ факторов существует следующее взаимно однозначное соответствие между прямым произведением и пересечением:

\emph{I. Если $\mG=\mA_1\times\cdots\times\mA_n=\mA_i\times\mB_i$, то $\mE=\mB_1\cap\cdots\cap\mB_n=\mB_i\cap\mC_i$ является прямым и $\mC_i=\mA_i$.}

\emph{II. Если $\mE=\mS_1\cap\cdots\cap\mS_n=\mR_i\cap\mS_i$ является прямым, то $\mG=\mR_1\times\cdots\times\mR_n=\mR_i\times\mT_i$ и $\mT_i=\mS_i$.}

\emph{Доказательство I.} Пусть $\mB=\mB_1\cap\cdots\cap\mB_n=\mB_i\cap\mC_i$; тогда $\mC_i\supseteq\mA_i$. Покажем, что $\mC_i\subseteq\mA_i$. Пусть, например, $c$ принадлежит $\mC_1$; тогда $c$ принадлежит $\mB_2,\ldots,\mB_n$, поэтому представление по компонентам относительно $\mA$ дает:
\[
        c=a_1a_2\cdots a_n=a_1ea_3\cdots a_n=a_1a_2e\cdots a_n=\cdots=a_1\cdots a_{n-1}e.
\]
Так как произведение всех $\mA_i$ является прямым, эти представления совпадают; на каждом месте, кроме первого, один раз стоит $e$, следовательно $c=a_1$, или $\mC_1\subseteq\mA_1$, то есть $\mC_1=\mA_1$, и соответственно $\mC_i=\mA_i$. Отсюда $\mB=\mB_i\cap\mC_i=\mB_i\cap\mA_i=\mE$; $\mB_i\mC_i=\mB_i\mA_i=\mG$, следовательно, пересечение прямое.

\emph{Доказательство II.} Пусть $\mH=\mR_1\cdots\mR_n=\mR_i\mT_i$; тогда $\mT_i\subseteq\mS_i$. Покажем, что $\mS_i\subseteq\mT_i$. Пусть, например, $s$ принадлежит $\mS_1$; тогда, в силу $\mG=\mS_i\times\mR_i$,
\[
        s=s_1e=s_2r_2=\cdots=s_nr_n.
\]
Образуем $t_1=er_2\cdots r_n=t_ir_i$; тогда, учитывая поэлементную коммутативность $\mR_i$ и $\mS_i$, получаем
\[
        st_1^{-1}=s_it_i^{-1},
\]
то есть элемент всех $\mS_i$, а потому равный $e$. Поэтому $\mS_1\subseteq\mT_1$, или $\mT_1=\mS_1$, и соответственно $\mT_i=\mS_i$. Отсюда $\mH=\mR_i\mT_i=\mR_i\mS_i=\mG$ и $\mR_i\cap\mT_i=\mR_i\cap\mS_i=\mE$; следовательно, $\mG$ является прямым произведением $\mR_i$.
% END UNIT BODY
% END PRODUCER UNIT 189
\endgroup


\clearpage
\begingroup
\setcounter{footnote}{0}
% BEGIN PRODUCER UNIT 190
% Source: C:/Users/Floris/Documents/interlanguage/03_projects/noether/02_slavic_working_corpus/translations/paper34/section05/russian/v001/Noether_Paper34_Section05_Russian_v001.tex
% Identity: 3979 B / SHA-256 0CA1CF52B383F391A0B6B64989CC927F99BF0D83B273D45DBDDDC0DE05785600
\setlength{\parindent}{1.15em}
\setlength{\parskip}{0pt}
\makeatletter
\renewcommand\footnotesize{\@setfontsize\footnotesize{7.5}{8.5}\selectfont}
\makeatother
\providecommand{\mG}{\mathfrak{G}}
\providecommand{\mA}{\mathfrak{A}}
\providecommand{\mE}{\mathfrak{E}}
\providecommand{\mH}{\mathfrak{H}}
\providecommand{\iso}{\simeq}
\emergencystretch=120em
\allowdisplaybreaks
\sloppy
\hbadness=10000
\vbadness=10000
\hfuzz=2pt
\vfuzz=10pt
\raggedbottom
% BEGIN UNIT BODY
% Unit: Paper34 section05 v001. Direct Russian translation from audited German source lines 294--332, segments P34-S0068--P34-S0073. Source scan page 13 retained as witness, and Zenodo record 20822156 was rechecked for this unit.
\setcounter{footnote}{0}
\section*{34. Гиперкомплексные величины и теория представлений}
\emph{Продолжение: глава I, §§5--7, и глава II, §§8--14.}

\subsection*{§ 5. Вполне приводимые группы}

Группа называется \emph{вполне приводимой}, если она является прямым произведением конечного числа простых групп:
\[
        \mG=\mA_1\times\cdots\times\mA_n.
\]
В этом случае группы
\[
        \mA_1\times\cdots\times\mA_n
        \supset \mA_1\times\cdots\times\mA_{n-1}
        \supset\cdots\supset \mA_1\supset \mE
\]
образуют композиционный ряд. Его длина равна $n$, а композиционные факторы суть
\[
        \sim \mA_n,\ \mA_{n-1},\ldots,\mA_1.
\]

\emph{Если $\mG$ вполне приводима, то каждый нормальный делитель является прямым фактором, а другой фактор можно выбрать как произведение тех простых групп, которые входят в заданное разложение $\mG$ в произведение. Нормальный делитель $\mH$ также вполне приводим.}

\emph{Доказательство.} Имеем
\begin{equation}
        \mG=\mH\mA_1\mA_2\cdots\mA_n.             \tag{3}
\end{equation}
Положим $\mH_i=\mH\mA_1\cdots\mA_i$; тогда $\mH_{i+1}=\mH_i\mA_{i+1}$. Либо $\mA_{i+1}\leq \mH_i$, следовательно $\mH_{i+1}=\mH_i$; тогда в (3) опускаем фактор $\mA_{i+1}$. Либо $\mH_i\cap\mA_{i+1}$ является собственным нормальным делителем в $\mA_{i+1}$, следовательно равен $\mE$, и потому $\mH_i\times\mA_{i+1}$ является прямым произведением. После удаления лишних факторов (3) принимает вид
\[
        \mG=\mH\times\mA_{i_1}\times\cdots\times\mA_{i_r},
\]
следовательно $\mH$ является прямым фактором. Далее имеем
\[
        \mH\sim \mG/(\mA_{i_1}\times\cdots\times\mA_{i_r})
        \sim \mA_{j_1}\times\cdots\times\mA_{j_\mu},
\]
где $\mA_{j_1},\ldots,\mA_{j_\mu}$ являются остальными группами $\mA_i$; следовательно $\mH$ вполне приводим.

\emph{Следствие.} \emph{Всякое разложение вполне приводимой группы на прямо неразложимые факторы является разложением на простые факторы и, следовательно, дает композиционный ряд; отсюда следует однозначная определенность факторов с точностью до изоморфизма.}
% END UNIT BODY
% END PRODUCER UNIT 190
\endgroup


\clearpage
\begingroup
\setcounter{footnote}{0}
% BEGIN PRODUCER UNIT 191
% Source: C:/Users/Floris/Documents/interlanguage/03_projects/noether/02_slavic_working_corpus/translations/paper34/section06/russian/v001/Noether_Paper34_Section06_Russian_v001.tex
% Identity: 7342 B / SHA-256 56622A5D7D1082068CB60898877632A164A025D1BAB237AF725F927B7B55ED58
\setlength{\parindent}{1.15em}
\setlength{\parskip}{0pt}
\makeatletter
\renewcommand\footnotesize{\@setfontsize\footnotesize{7.5}{8.5}\selectfont}
\makeatother
\providecommand{\mG}{\mathfrak{G}}
\providecommand{\mA}{\mathfrak{A}}
\providecommand{\mE}{\mathfrak{E}}
\providecommand{\mH}{\mathfrak{H}}
\providecommand{\mR}{\mathfrak{R}}
\providecommand{\mo}{\mathfrak{o}}
\providecommand{\iso}{\simeq}
\emergencystretch=120em
\allowdisplaybreaks
\sloppy
\hbadness=10000
\vbadness=10000
\hfuzz=2pt
\vfuzz=10pt
\raggedbottom
% BEGIN UNIT BODY
% Unit: Paper34 section06 v001. Direct Russian translation from audited German source lines 334--404, segments P34-S0074--P34-S0084. Source scan pages 13--15 retained as witnesses, and Zenodo record 20822156 was rechecked for this unit.
\setcounter{footnote}{11}
\section*{34. Гиперкомплексные величины и теория представлений}
\emph{Продолжение: глава I, §§5--7, и глава II, §§8--14.}

\subsection*{§ 6. Модули относительно тела. Гиперкомплексные системы}

Почти тривиальный пример вполне приводимых групп с операторами дают модули с конечной базой относительно (не обязательно коммутативного) тела.

Пусть $\mG$ является правым $K$-модулем, и пусть единичный элемент $e$ тела $K$ одновременно является единичным оператором: $ae=a$ для $a$ из $\mG$. Каждый подмодуль $aK$, порожденный элементом $a$, прост, а именно равен модулю, порожденному любым его ненулевым элементом. Следовательно, если $\mH$ является подмодулем, порожденным некоторыми элементами, а $a$ -- элемент, не принадлежащий $\mH$, то $\mH\cap aK=\mE$, и потому $\mH+aK$ является прямой суммой. Так, начиная с любого базисного элемента $a_1$, можно всякий раз присоединять новые базисные элементы $a_2,a_3,\ldots$ и получить $\mG$ как прямую сумму
\[
        \mG=a_1K+a_2K+\cdots+a_nK.
\]
\emph{Следовательно, $\mG$ вполне приводим.}

Число $n$, то есть длина композиционного ряда, называется \emph{рангом} $\mG$ относительно $K$. База $a_1,\ldots,a_n$ линейно независима в силу единственности представления элементов $\mG$ (и потому, что $e$ был предположен единичным оператором).

Каждый подмодуль $\mH$ является прямым слагаемым:
\[
        \mG=\mH+\mR=(h_1K+\cdots+h_sK)+(r_1K+\cdots+r_{n-s}K),
\]
или: \emph{каждую линейно независимую базу $\mH$ можно дополнить до линейно независимой базы $\mG$.} Элементы $r_i$ можно даже выбрать из исходных базисных элементов $a_i$.

Пусть $(a_1,\ldots,a_n)$ и $(c_1,\ldots,c_n)$ -- линейно независимые базы. Тогда
\[
        c_k=\sum_i a_i\pi_{ik},
\]
или, в матричной записи,
\[
        (c_1,\ldots,c_n)=(a_1,\ldots,a_n)P,
        \qquad
        P=\begin{pmatrix}
        \pi_{11}&\cdots&\pi_{1n}\\
        \vdots&&\vdots\\
        \pi_{n1}&\cdots&\pi_{nn}
        \end{pmatrix}.
\]
Обратно,
\[
        (a_1,\ldots,a_n)=(c_1,\ldots,c_n)Q.
\]
Отсюда следует
\[
        (a_1,\ldots,a_n)=(a_1,\ldots,a_n)PQ,
        \qquad
        (c_1,\ldots,c_n)=(c_1,\ldots,c_n)QP,
\]
а значит, поскольку $a_i$ и $c_i$ линейно независимы,
\[
        PQ=QP=E.
\]

Специальные $K$-модули, которые одновременно являются кольцами, суть ``гиперкомплексные системы''.

Кольцо $\mo$, которое одновременно является правым модулем относительно коммутативного поля $K$, называется \emph{гиперкомплексной системой относительно $K$} (в литературе также ``алгеброй над $K$''), если:
\begin{enumerate}
\item ранг конечен (пусть линейно независимой базой является $u_1,\ldots,u_n$);
\item $ab\cdot x=a\cdot bx=ax\cdot b$. Это выражают так: $K$ коммутативно связано с $\mo$;
\item единичный элемент $e$ поля $K$ одновременно является единичным оператором: $ae=a$ для $a$ из $\mo$.
\end{enumerate}
Поскольку
\[
        \Bigl(\sum_i u_i\alpha_i\Bigr)\Bigl(\sum_k u_k\beta_k\Bigr)
        =\sum_{i,k}(u_i u_k)(\alpha_i\beta_k),
\]
система однозначно определена, как только, кроме $K$, задана \emph{таблица умножения}, указывающая, как каждое произведение $u_i u_k$ выражается через $u_j$:
\[
        u_i u_k=\sum_j u_j\gamma^j_{ik}.
\]
Элементы $\gamma^j_{ik}$ должны удовлетворять известным соотношениям, следующим из ассоциативного закона.

Если $\mo$, как мы часто будем предполагать, имеет единичный элемент $e$ ($ea=ae=a$ для всех $a$), то элементы $x$ поля $K$ можно отождествить с $ex$, то есть рассматривать $K$ как подполе в $\mo$. Тогда гиперкомплексную систему можно также описать как \emph{кольцо конечного ранга относительно поля, лежащего в центре.}

Примером гиперкомплексной системы является \emph{групповое кольцо} конечной группы: в качестве базисных элементов $u_i$ берут элементы конечной группы, а в качестве умножения -- групповое умножение. Поле $K$ при этом произвольно.

\footnotetext[12]{Условие 2 означает, что умножение на элемент из $K$ задает гомоморфизм для $\mo$, когда $\mo$ рассматривается и как левый $\mo$-модуль, и как правый $\mo$-модуль. В силу 3 эти гомоморфизмы даже являются изоморфизмами.}
% END UNIT BODY
% END PRODUCER UNIT 191
\endgroup


\clearpage
\begingroup
\setcounter{footnote}{0}
% BEGIN PRODUCER UNIT 192
% Source: C:/Users/Floris/Documents/interlanguage/03_projects/noether/02_slavic_working_corpus/translations/paper34/section07/russian/v001/Noether_Paper34_Section07_Russian_v001.tex
% Identity: 4146 B / SHA-256 6DF009D8AC0AFBDC720581D4204D22321F35500478F81ADBB49E5E837DDFEDC0
\setlength{\parindent}{1.15em}
\setlength{\parskip}{0pt}
\makeatletter
\renewcommand\footnotesize{\@setfontsize\footnotesize{7.5}{8.5}\selectfont}
\makeatother
\providecommand{\mG}{\mathfrak{G}}
\providecommand{\mA}{\mathfrak{A}}
\providecommand{\mE}{\mathfrak{E}}
\providecommand{\mH}{\mathfrak{H}}
\providecommand{\mR}{\mathfrak{R}}
\providecommand{\mo}{\mathfrak{o}}
\providecommand{\iso}{\simeq}
\emergencystretch=120em
\allowdisplaybreaks
\sloppy
\hbadness=10000
\vbadness=10000
\hfuzz=2pt
\vfuzz=10pt
\raggedbottom
% BEGIN UNIT BODY
% Unit: Paper34 section07 v001. Direct Russian translation from audited German source lines 406--441, segments P34-S0085--P34-S0093. Source scan pages 15--16 retained as witnesses, and Zenodo record 20822156 was rechecked for this unit. The following chapter-II heading is segment P34-S0094 and is intentionally reserved for Section08.
\setcounter{footnote}{12}
\section*{34. Гиперкомплексные величины и теория представлений}
\emph{Продолжение: глава I, §§5--7, и глава II, §§8--14.}

\subsection*{§ 7. Матрицы}

Квадратные матрицы $(\pi_{ik})$, где $\pi_{ik}$ принадлежат кольцу $\mo$, образуют кольцо относительно обычного матричного умножения $\sum_j \pi_{ij}\rho_{jk}=\sigma_{ik}$ и сложения $\pi_{ik}+\rho_{ik}=\sigma_{ik}$.

Матрица с элементами из тела $K$, для которой существуют правая и левая обратные матрицы, называется \emph{регулярной}.

В §6 мы видели: \emph{матрица перехода между двумя линейно независимыми базами $K$-модуля регулярна.}

\emph{Для регулярности достаточно правой обратной.} Действительно, образуем правый $K$-модуль с линейно независимой базой $(a_1,\ldots,a_n)$ (модуль линейных форм от неизвестных $a_1,\ldots,a_n$) и положим
\[
        (c_1,\ldots,c_n)=(a_1,\ldots,a_n)P.
\]
Тогда
\[
        (c_1,\ldots,c_n)P^{-1}=(a_1,\ldots,a_n)PP^{-1}=(a_1,\ldots,a_n);
\]
следовательно, $(c_1,\ldots,c_n)$ снова является базой, значит снова линейно независима; следовательно, матрица $P$ регулярна. Кроме того, правая обратная равна левой обратной.

Точно так же достаточно левой обратной.

\emph{Если $P$ имеет левую обратную, то $P$ не является левым делителем нуля.}
\emph{Доказательство.} Из $PA=0$ следует $A=P^{-1}PA=0$.

Следовательно, $P$ также имеет только одну правую обратную, которая по §6 совпадает с левой обратной.

Если, как обычно, под $P_x$ понимать матрицу $(\pi_{ik}x)$, а под $C_{ik}$ -- матрицу, у которой на месте $ik$ стоит единичный элемент, а всюду в остальных местах нули, то
\[
        (\pi_{ik})=\sum_{i,k} C_{ik}\pi_{ik};
\]
следовательно, матрицы над $K$ образуют $K$-модуль ранга $n^2$. Матрицы $C_{ik}$ удовлетворяют соотношениям
\[
        C_{ij}C_{jk}=C_{ik},
        \qquad
        C_{ij}C_{\bar j k}=0,\quad j\ne\bar j.
\]
Тем самым матрицы над $K$ образуют конечный $K$-модуль и, если $K$ коммутативно, гиперкомплексную систему.
% END UNIT BODY
% END PRODUCER UNIT 192
\endgroup


\clearpage
\begingroup
\setcounter{footnote}{0}
% BEGIN PRODUCER UNIT 193
% Source: C:/Users/Floris/Documents/interlanguage/03_projects/noether/02_slavic_working_corpus/translations/paper34/section08/russian/v001/Noether_Paper34_Section08_Russian_v001.tex
% Identity: 6400 B / SHA-256 30C425BA32338956C6E81D4376809C02B62D04B128F4508E8E6B9CF75E18E492
\setlength{\parindent}{1.15em}
\setlength{\parskip}{0pt}
\makeatletter
\renewcommand\footnotesize{\@setfontsize\footnotesize{7.5}{8.5}\selectfont}
\makeatother
\providecommand{\mA}{\mathfrak{A}}
\providecommand{\mB}{\mathfrak{B}}
\providecommand{\mC}{\mathfrak{C}}
\providecommand{\mM}{\mathfrak{M}}
\providecommand{\ma}{\mathfrak{a}}
\providecommand{\md}{\mathfrak{d}}
\providecommand{\mo}{\mathfrak{o}}
\emergencystretch=120em
\allowdisplaybreaks
\sloppy
\hbadness=10000
\vbadness=10000
\hfuzz=2pt
\vfuzz=10pt
\raggedbottom
% BEGIN UNIT BODY
% Unit: Paper34 section08 v001. Direct Russian translation from audited German source lines 443--472, segments P34-S0094--P34-S0102. Source scan pages 16--17 retained as witnesses, and Zenodo record 20822156 was rechecked for this unit. The following Section09 heading is segment P34-S0103 and is intentionally reserved for Section09.
\setcounter{footnote}{12}
\section*{34. Гиперкомплексные величины и теория представлений}
\emph{Продолжение: глава II, §§8--14.}

\section*{Глава II. Некоммутативная теория идеалов}

\subsection*{§ 8. Теорема о гомоморфизме для колец}

Если $\ma$ является двусторонним идеалом в $\mo$, то область классов вычетов $\mo/\ma$ является не только $\mo$-модулем, но и кольцом, как легко видеть. Каждый гомоморфный образ $\mo$ снова является кольцом и изоморфен кольцу классов вычетов $\mo/\ma$, где $\ma$ является двусторонним идеалом. Доказательство такое же, как для групп.

В частности, если $\mo$ является телом, то нет других идеалов, кроме $(0)$ и $\mo$; следовательно, здесь гомоморфизм либо является изоморфизмом, либо сопоставляет каждому элементу нуль.

В частности, пусть кольцо $\mo$ является правым $K$-модулем, где $K$ должно быть кольцом, поэлементно перестановочным с $\mo$:
\[
        ab\cdot x=a\cdot bx=ax\cdot b
\]
(например, когда речь идет о гиперкомплексной системе относительно коммутативного тела $K$). Тогда допустимыми (правыми, левыми или двусторонними) идеалами считаются только те, которые одновременно являются $K$-модулями, и, кроме гомоморфизма колец, требуется также операторный гомоморфизм относительно умножения на $K$. Допустимые идеалы становятся бимодулями относительно $\mo$ и $K$. (В гиперкомплексных системах под ``идеалами'' без дальнейших уточнений всегда понимают только те, которые допустимы относительно лежащего в основе тела коэффициентов $K$.) И при этом расширении приведенная выше теорема о гомоморфизме остается верной.

Примером гомоморфизма колец служит переход от области множителей $\mo$ к абсолютной области множителей (§1, пример 4). Следовательно, абсолютная область множителей изоморфна кольцу классов вычетов $\mo/\md$, где $\md$ является двусторонним идеалом элементов из $\mo$, аннулирующих $\mM$.

Из справедливости теоремы о гомоморфизме, как и в §2, следуют первая и вторая теоремы об изоморфизме для изоморфизма колец и операторного изоморфизма.

Все произведения $\mA\mB$ подмножеств $\mo$ следует понимать как модульные произведения: $\mA\mB$ есть совокупность всех сумм $\sum ab$, где $a$ принадлежит $\mA$, $b$ принадлежит $\mB$, а $a\mB$ есть совокупность всех сумм $\sum ab$, где $b$ принадлежит $\mB$. Если $\mB$ является правым модулем относительно некоторого кольца как области операторов, то и произведение $\mA\mB$, соответственно $a\mB$, является правым модулем. (В частности: правым идеалом, допустимым идеалом относительно области коэффициентов $K$ и т. д.) Если $(\mA,\mB)$ обозначает сумму модулей $\mA,\mB$, то имеют место правила вычисления
\[
        \mA\mB\cdot\mC=\mA\cdot\mB\mC,
\]
\[
        \mA(\mB,\mC)=(\mA\mB,\mA\mC),
        \qquad
        (\mA,\mB)\mC=(\mA\mC,\mB\mC).
\]
Прямая сумма обозначается через $\mA+\mB$ (§4).

Далее мы часто будем рассматривать кольца, удовлетворяющие условию максимальности или минимальности (§3) для правых идеалов. К ним относятся, в частности, гиперкомплексные системы, поскольку по принятому выше соглашению в качестве идеалов допускаются только $K$-модули, а значит их ранг остается ограниченным (§6).
% END UNIT BODY
% END PRODUCER UNIT 193
\endgroup


\clearpage
\begingroup
\setcounter{footnote}{0}
% BEGIN PRODUCER UNIT 194
% Source: C:/Users/Floris/Documents/interlanguage/03_projects/noether/02_slavic_working_corpus/translations/paper34/section09/russian/v001/Noether_Paper34_Section09_Russian_v001.tex
% Identity: 4622 B / SHA-256 CCA1DF06618ACC783799B5B4AA1A5D95B49014EBBFCAED6A666BFE17F12AF274
\setlength{\parindent}{1.15em}
\setlength{\parskip}{0pt}
\makeatletter
\renewcommand\footnotesize{\@setfontsize\footnotesize{7.5}{8.5}\selectfont}
\makeatother
\providecommand{\ml}{\mathfrak{l}}
\providecommand{\mo}{\mathfrak{o}}
\providecommand{\mr}{\mathfrak{r}}
\emergencystretch=120em
\allowdisplaybreaks
\sloppy
\hbadness=10000
\vbadness=10000
\hfuzz=2pt
\vfuzz=10pt
\raggedbottom
% BEGIN UNIT BODY
% Unit: Paper34 section09 v001. Direct Russian translation from audited German source lines 474--562, segments P34-S0103--P34-S0113. Source scan pages 17--18 retained as witnesses, and Zenodo record 20822156 was rechecked for this unit. The following Section10 heading is segment P34-S0114 and is intentionally reserved for Section10.
\setcounter{footnote}{12}
\section*{34. Гиперкомплексные величины и теория представлений}
\emph{Продолжение: глава II, §§8--14.}

\section*{Глава II. Некоммутативная теория идеалов}

\subsection*{§ 9. Идемпотентные элементы. Прямое разложение на правые идеалы}

Пусть $\mo$ является кольцом с единичным элементом.

Если $\mr$ является правым идеалом, то $\mr\mo\subseteq\mr$, а из-за наличия единичного элемента даже $\mr\mo=\mr$. Соответственно для левых идеалов: $\mo\ml=\ml$.

Элемент $c$ называется идемпотентным, если $c^2=c$.

\emph{Если $\mo=\mr_1+\cdots+\mr_n$ является разложением на правые идеалы и при этом}
\[
        e=e_1+\cdots+e_n,
        \qquad e_i\in\mr_i,
\]
\emph{то}
\[
        e_i^2=e_i,
        \qquad e_i e_k=0\quad(i\ne k)
        \quad\hbox{(``соотношения ортогональности'')},
\]
\[
        \mr_i=e_i\mo.
\]

\emph{Доказательство.} Пусть $r$ является элементом $\mr_1$. Тогда
\[
        r=er=e_1r+\cdots+e_nr,
        \qquad e_ir\in\mr_i,
\]
но также
\[
        r=r+0+\cdots+0,
\]
следовательно,
\[
        e_1r=r,
        \qquad e_ir=0\quad(i\ne1).
\]
Из первой формулы следует $\mr_1\subseteq e_1\mo$; с другой стороны, $e_1\mo\subseteq\mr_1$, значит $e_1\mo=\mr_1$. При специализации $r=e_1$ получаем
\[
        e_1^2=e_1,
        \qquad e_i e_1=0\quad(i\ne1).
\]
То же самое имеет место, если индекс $1$ заменить на $k$.

\emph{Обратно, если $e=\sum e_i$, $e_i e_k=0$ $(i\ne k)$, $e_i^2=e_i$ и положить $\mr_i=e_i\mo$, то $\mo=\mr_1+\cdots+\mr_n$.}

\emph{Доказательство.} Каждый элемент $r$ из $\mo$ имеет вид
\[
        r=er=e_1r+\cdots+e_nr,
\]
поэтому
\[
        \mo=(e_1\mo,\ldots,e_n\mo).
\]
Но это представление единственно: действительно, из
\[
        0=e_1a_1+\cdots+e_na_n
\]
после умножения на $e_i$ следует
\[
        e_i^2a_i=e_i a_i=0.
\]

Из правого представления
\[
        \mo=\mr_1+\cdots+\mr_n=e_1\mo+\cdots+e_n\mo
\]
тем самым следует левое представление
\[
        \mo=\ml_1+\cdots+\ml_n=\mo e_1+\cdots+\mo e_n.
\]

Если $\mr_i$ прямо неразложимы, то такими же являются и $\ml_i$; ведь разложение, скажем, $\ml_1$, означало бы
\[
        \ml_1=\mo e_1=\mo e_1'+\mo e_1'',
        \qquad
        \mo=\mo e_1'+\mo e_1''+\mo e_2+\cdots+\mo e_n.
\]
Отсюда возникло бы правое разложение
\[
        \mo=e_1'\mo+e_1''\mo+e_2\mo+\cdots+e_n\mo
        =\mr_1'+\mr_1''+\mr_2+\cdots+\mr_n,
\]
\[
        \mr_1\cong \mo/(\mr_2+\cdots+\mr_n)=\mr_1'+\mr_1'',
\]
то есть $\mr_1$ был бы разложимым.\footnote{Имеет место еще одно обращение: если $e_i e_k=0$, $e_i^2=e_i$, $c=\sum e_i$, то $c$ является левой единицей для кольца $e_1\mo+\cdots+e_n\mo$. То, что сумма является прямой, видно совершенно так же, как выше.}
% END UNIT BODY
% END PRODUCER UNIT 194
\endgroup


\clearpage
\begingroup
\setcounter{footnote}{0}
% BEGIN PRODUCER UNIT 195
% Source: C:/Users/Floris/Documents/interlanguage/03_projects/noether/02_slavic_working_corpus/translations/paper34/section10/russian/v001/Noether_Paper34_Section10_Russian_v001.tex
% Identity: 7507 B / SHA-256 93C236A4BE578592EC337A2C5F9902CEAF4E72C6AE76B9B23EDB55DA79B3C214
\setlength{\parindent}{1.15em}
\setlength{\parskip}{0pt}
\makeatletter
\renewcommand\footnotesize{\@setfontsize\footnotesize{7.5}{8.5}\selectfont}
\makeatother
\providecommand{\ma}{\mathfrak{a}}
\providecommand{\mb}{\mathfrak{b}}
\providecommand{\mc}{\mathfrak{c}}
\providecommand{\mo}{\mathfrak{o}}
\providecommand{\mr}{\mathfrak{r}}
\emergencystretch=120em
\allowdisplaybreaks
\sloppy
\hbadness=10000
\vbadness=10000
\hfuzz=2pt
\vfuzz=10pt
\raggedbottom
% BEGIN UNIT BODY
% Unit: Paper34 section10 v001. Direct Russian translation from audited German source lines 564--659, segments P34-S0114--P34-S0128. Source scan pages 18--20 retained as witnesses, and Zenodo record 20822156 was rechecked for this unit. The following Section11 heading is segment P34-S0129 and is intentionally reserved for Section11.
\section*{34. Гиперкомплексные величины и теория представлений}
\emph{Продолжение: глава II, §§8--14.}

\section*{Глава II. Некоммутативная теория идеалов}

\subsection*{§ 10. Разложение на двусторонние идеалы}

Пусть $\mo$ снова является кольцом с единичным элементом, и пусть
\begin{equation}
        \mo=\ma_1+\cdots+\ma_n                                  \tag{1}
\end{equation}
есть разложение $\mo$ на двусторонние идеалы. Тогда соотношения ортогональности имеют место и для идеалов:
\begin{equation}
        \ma_i\ma_k=0\quad(i\ne k),
        \qquad
        \ma_i^2=\ma_i.                                             \tag{2}
\end{equation}

\emph{Доказательство.} Пусть $\mb_1=\ma_2+\cdots+\ma_n$, $\mo=\ma_1+\mb_1$; тогда $\ma_1\mb_1\leq \ma_1\cap\mb_1=0$, значит $\ma_1\mb_1=0$, и тем более
\[
        \ma_1\ma_i=0\quad(i\ne1);
\]
\[
        \ma_1=\ma_1\mo=\ma_1(\ma_1+\mb_1)=(\ma_1^2,\ma_1\mb_1)=\ma_1^2.
\]

Обратно: \emph{из (1) и (2) следует, что модули $\ma_i$ являются двусторонними идеалами.}
\emph{Доказательство.}
\[
        \mo\ma_i=(\ma_1+\cdots+\ma_n)\ma_i=\ma_i^2=\ma_i,
        \qquad
        \ma_i\mo=\ma_i.
\]

\emph{Правые идеалы в кольце $\ma_i$ являются правыми идеалами в $\mo$.}
\emph{Доказательство.}
\[
        \mr\mo=\mr(\ma_i+\mb_i)=(\mr\ma_i,\mr\mb_i)=(\mr,0)=\mr.
\]

\emph{Если $\mr$ является правым идеалом в $\mo$, то $\mr$ является прямой суммой правых идеалов $\mr\ma_i\leq\ma_i$.}
\emph{Доказательство.} $\mr=\mr\mo=(\mr\ma_1,\ldots,\mr\ma_n)$, и $\mr\ma_i\mo=\mr\ma_i$, значит $\mr\ma_i$ являются правыми идеалами, содержащимися в $\mr$. Далее, $\mr\ma_i\leq\ma_i$, поэтому сумма прямая:
\[
        \mr=\mr\ma_1+\cdots+\mr\ma_n.
\]

\emph{В частности, если $\mr$ прямо неразложим, то $\mr$ может иметь только одну компоненту; следовательно, $\mr$ должен лежать в одном из $\ma_i$.}

На основании этих теорем теория идеалов в $\mo$ становится управляемой, как только известна теория идеалов отдельных $\ma_i$. Поэтому мы в основном ограничиваемся двусторонне неразложимыми кольцами.

\emph{Два правых идеала, содержащиеся в различных $\ma_i$, никогда не могут быть операторно изоморфны.}
\emph{Доказательство.} Идеал, лежащий в $\ma_1$, аннулируется всеми остальными $\ma_k$, но не аннулируется $\ma_1$. Напротив, идеал, лежащий в $\ma_2$, аннулируется $\ma_1$.

Итак, если кольцо $\mo$ можно разложить на двусторонне неразложимые идеалы $\ma_1+\cdots+\ma_n$, и каждое $\ma_i$ снова разложить на односторонне неразложимые правые идеалы $\mr_i'+\mr_i''+\cdots$, то операторно изоморфные $\mr$ надо искать среди тех, которые принадлежат одному и тому же $\ma_i$. Но бывает, что в одном и том же $\ma_i$ есть еще различные классы $\mr$ -- то есть классы, не операторно изоморфные, -- как показывает следующий пример.

Пусть $K$ есть поле рациональных чисел,
\[
        \mo=e_1K+e_2K+uK
\]
гиперкомплексная система, причем $e_i$ и $u$ коммутируют с числами из $K$, а таблица умножения имеет вид
\[
\begin{array}{c|ccc}
     & e_1&e_2&u\\ \hline
 e_1 & e_1&0&u\\
 e_2 & 0&e_2&0\\
 u   & 0&u&0
\end{array}
\]
Единичным элементом является $e=e_1+e_2$. Элементы $e_i$ удовлетворяют соотношениям ортогональности. Следовательно, правое разложение есть
\[
        \mo=e_1\mo+e_2\mo=(e_1,u)+(e_2),
\]
а левое разложение есть
\[
        \mo=\mo e_1+\mo e_2=(e_1)+(e_2,u).
\]
Так как $e_2\mo$ и $\mo e_1$ явно неразложимы, то $\mo e_2$ и $e_1\mo$ также должны быть неразложимы.

Двустороннее разложение невозможно; ведь тогда одна компонента должна была бы содержать $e_1\mo$, а другая -- $e_2\mo$; первая тогда содержала бы $u$, но и вторая содержала бы $ue_2=u$.

Однако $e_1\mo$ и $e_2\mo$ не операторно изоморфны, поскольку они имеют разный ранг относительно $K$.

\emph{Пусть $\mo=\ma_1+\cdots+\ma_n$, и пусть $\ma_i$ двусторонне неразложимы. Тогда они определены однозначно.}

\emph{Доказательство.} Пусть, кроме того, $\mo=\mc_1+\cdots+\mc_n$, то есть
\[
        \ma_i=\ma_i\mo=(\ma_i\mc_1,\ldots,\ma_i\mc_n).
\]
Последняя сумма прямая, потому что $\ma_i\mc_k\leq\mc_k$ и $\leq\ma_i$; но так как $\ma_i$ неразложимы, все $\ma_i\mc_k$ должны быть нулевыми, кроме одного: $\ma_i\mc_{j_i}$. Следовательно,
\[
        \ma_i=\ma_i\mc_{j_i}.
\]
Точно так же:
\[
        \mc_{j_i}=\mo\mc_{j_i}=\ma_1\mc_{j_i}+\cdots+\ma_n\mc_{j_i},
\]
и $\ma_i\mc_{j_i}\ne0$, следовательно все остальные слагаемые равны $0$,
\[
        \mc_{j_i}=\ma_i\mc_{j_i}=\ma_i,
\]
что и требовалось доказать.
% END UNIT BODY
% END PRODUCER UNIT 195
\endgroup


\clearpage
\begingroup
\setcounter{footnote}{0}
% BEGIN PRODUCER UNIT 196
% Source: C:/Users/Floris/Documents/interlanguage/03_projects/noether/02_slavic_working_corpus/translations/paper34/section11/russian/v001/Noether_Paper34_Section11_Russian_v001.tex
% Identity: 4946 B / SHA-256 DBB292E4173F950DC2527B1E5C3FCFF4AE8008046E6D77E8128D9F15F1C2E2BD
\setlength{\parindent}{1.15em}
\setlength{\parskip}{0pt}
\makeatletter
\renewcommand\footnotesize{\@setfontsize\footnotesize{7.5}{8.5}\selectfont}
\makeatother
\providecommand{\ma}{\mathfrak{a}}
\providecommand{\mb}{\mathfrak{b}}
\providecommand{\mc}{\mathfrak{c}}
\providecommand{\mA}{\mathfrak{A}}
\providecommand{\mZ}{\mathfrak{Z}}
\providecommand{\mo}{\mathfrak{o}}
\providecommand{\mr}{\mathfrak{r}}
\emergencystretch=120em
\allowdisplaybreaks
\sloppy
\hbadness=10000
\vbadness=10000
\hfuzz=2pt
\vfuzz=10pt
\raggedbottom
% BEGIN UNIT BODY
% Unit: Paper34 section11 v001. Direct Russian translation from audited German source lines 661--730, segments P34-S0129--P34-S0136. Source scan pages 20--22 retained as witnesses; Section12 heading P34-S0137 is intentionally reserved for Section12. Zenodo record 20822156 was rechecked for this unit at 2026-06-24T22:09Z UTC.
\section*{34. Гиперкомплексные величины и теория представлений}
\emph{Продолжение: глава II, §§8--14.}

\section*{Глава II. Некоммутативная теория идеалов}

\subsection*{§ 11. Центр}

Центр $\mZ$ кольца $\mo$ есть совокупность всех $z$, которые коммутируют со всеми элементами кольца ($za=az$ для всех $a$). $\mZ$ является коммутативным кольцом.

Каждый односторонний идеал $\ma$ в $\mo$ имеет в $\mZ$ ``идеал сужения'' $\ma\cap\mZ$. Действительно, так как и $\ma$, и $\mZ$ являются $\mZ$-модулями, то $\ma\cap\mZ$ также является $\mZ$-модулем.

Каждый идеал $\mA$ в $\mZ$ имеет идеал расширения: идеал $\ma=(\mA,\mo\mA)$, порожденный $\mA$ в $\mo$. Он двусторонний, поскольку
\[
        \ma\mo=(\mo\mA,\mA)\mo=(\mo\mA\mo,\mA\mo)
        =(\mo^2\mA,\mo\mA)\subseteq\mo\mA\subseteq\ma.
\]
Если в $\mo$ предполагается единичный элемент, то можно писать $\ma=\mo\mA$. В этом случае справедлива следующая теорема:

\emph{Каждому разложению $\mo$ на двусторонние идеалы взаимно однозначно соответствует разложение центра: из}
\begin{equation}
        \mo=\ma_1+\cdots+\ma_n,
        \qquad \mA_i=\ma_i\cap\mZ                         \tag{1}
\end{equation}
\emph{следует}
\begin{equation}
        \mZ=\mA_1+\cdots+\mA_n,
        \qquad \mo\mA_i=\ma_i,                              \tag{2}
\end{equation}
\emph{и обратно.}

\emph{Доказательство.} Пусть $z$ является центральным элементом, и пусть
\begin{equation}
        z=z_1+\cdots+z_n                                    \tag{3}
\end{equation}
есть его разложение согласно (1). Тогда $z_i$ являются центральными элементами; в самом деле,
\[
        za=z_1a+\cdots+z_na=az=az_1+\cdots+az_n;
\]
из прямоты суммы и из того, что $\ma_i$ двусторонний, получаем
\[
        z_i a=az_i.
\]
Следовательно, $z_i$ лежит в $\ma_i\cap\mZ=\mA_i$, и потому (3) дает утверждаемое разложение в сумму. В частности, $e=\sum e_i$; следовательно, $e_i$ являются центральными элементами. Наконец,
\[
        z=ze_1+\cdots+ze_n,
\]
а значит,
\[
        \mA_i=\mZ e_i,
        \qquad
        \mo\mA_i=\mo\mZ e_i=\mo e_i=\ma_i.
\]

Обратно, если $\mZ=\mA_1+\cdots+\mA_n$ есть разложение центра, то по §9
\[
        \mA_i=\mZ e_i,
        \qquad e_i^2=e_i,
        \qquad e_i e_k=0\quad(i\ne k).
\]
Значит, по соотношениям ортогональности \hbox{\(\S 9\)}
\[
\begin{aligned}
        \mo&=\mo e=\mo e_1+\cdots+\mo e_n
        =\mo\mZ e_1+\cdots+\mo\mZ e_n\\
        &=\mo\mA_1+\cdots+\mo\mA_n
        =\ma_1+\cdots+\ma_n .
\end{aligned}
\]
Если теперь положить $\ma_i\cap\mZ=\mA_i'$, то из первой части утверждения известно, что
\[
        \mZ=\mA_1'+\cdots+\mA_n'\quad\hbox{прямо}.
\]
Но $\mA_i'\supseteq\mA_i$, следовательно $\mA_i'=\mA_i$ по §4, 2.

\emph{Следствие.} $\ma_i$ и $\mA_i$ одновременно двусторонне прямо разложимы или нет.
% END UNIT BODY
% END PRODUCER UNIT 196
\endgroup


\clearpage
\begingroup
\setcounter{footnote}{0}
% BEGIN PRODUCER UNIT 197
% Source: C:/Users/Floris/Documents/interlanguage/03_projects/noether/02_slavic_working_corpus/translations/paper34/section12/russian/v001/Noether_Paper34_Section12_Russian_v001.tex
% Identity: 5853 B / SHA-256 5F4CD2136C6072509D53D6F5F790293E1B5232149F79329341F79301D25748C6
\setlength{\parindent}{1.15em}
\setlength{\parskip}{0pt}
\makeatletter
\renewcommand\footnotesize{\@setfontsize\footnotesize{7.5}{8.5}\selectfont}
\makeatother
\providecommand{\ma}{\mathfrak{a}}
\providecommand{\mb}{\mathfrak{b}}
\providecommand{\mc}{\mathfrak{c}}
\providecommand{\md}{\mathfrak{d}}
\providecommand{\mC}{\mathfrak{C}}
\providecommand{\mZ}{\mathfrak{Z}}
\providecommand{\mo}{\mathfrak{o}}
\providecommand{\mr}{\mathfrak{r}}
\emergencystretch=120em
\allowdisplaybreaks
\sloppy
\hbadness=10000
\vbadness=10000
\hfuzz=2pt
\vfuzz=10pt
\raggedbottom
% BEGIN UNIT BODY
% Unit: Paper34 section12 v001. Direct Russian translation from audited German source lines 732--770, segments P34-S0137--P34-S0149. Source scan pages 22--23 retained as witnesses; Section13 heading P34-S0150 is intentionally reserved for Section13. Zenodo record 20822156 was rechecked for this unit at 2026-06-24T22:37Z UTC with no file-list change.
\section*{34. Гиперкомплексные величины и теория представлений}
\emph{Продолжение: глава II, §§8--14.}

\section*{Глава II. Некоммутативная теория идеалов}

\subsection*{§ 12. Нильпотентные идеалы}

Идеал $\mc$ называется \emph{нильпотентным}, если $\mc^\rho=0$.

\emph{Пример.} Идеал $(u)$ в §10.

\emph{Если в $\mo$ существует нильпотентный правый идеал $\mr$, то существует также нильпотентный левый идеал (даже двусторонний идеал).}

\emph{Доказательство.} Пусть $\mr^\rho=0$. Тогда $\mo\mr$ является левым идеалом; соответственно, если $\mo$ не содержит единичного элемента и $\mo\mr$ мог бы быть нулем, левым идеалом является $(\mo\mr,\mr)$, и
\[
        (\mo\mr)^\rho=\mo\mr\mo\mr\cdots\mo\mr
        \subseteq\mo\mr\mr\cdots\mr=\mo\mr^\rho=0.
\]

\emph{Сумма двух нильпотентных правых идеалов снова является нильпотентным правым идеалом.}

\emph{Доказательство.} Пусть $\mc^\rho=0$, $\md^\sigma=0$. В
\[
        (\mc,\md)^{\rho+\sigma-1}=(\ldots,\md\mc\cdots\md\cdots\mc,\ldots)
\]
каждый член содержит либо по меньшей мере $\rho$ множителей $\mc$, либо по меньшей мере $\sigma$ множителей $\md$. В первом случае имеем
\[
        \md\mc\cdots\md\cdots\mc\cdots\subseteq\md\mc\mc\cdots\mc=\md\mc^\rho=0;
\]
во втором случае соответственно для множителей $\md$. Следовательно,
\[
        (\mc,\md)^{\rho+\sigma-1}=0.
\]

Если теперь в $\mo$ выполнено условие максимальности для правых идеалов, то существует максимальный нильпотентный правый идеал. Он содержит все остальные нильпотентные правые идеалы, поскольку иначе можно было бы образовать сумму с таким идеалом. Обозначим его через $\mc$. Так как $\mo\mc$ также является нильпотентным правым идеалом, имеем $\mo\mc\subseteq\mc$, значит $\mc$ является левым идеалом. Каждый нильпотентный левый идеал $\md$ также содержится в $\mc$. Действительно, $(\md,\md\mo)$ является нильпотентным правым идеалом, значит $\md\subseteq\mc$. \emph{Итак, существует максимальный (двусторонний) нильпотентный идеал $\mc$, или “радикал”, содержащий все остальные (правые и левые).}

В примере §10 идеал $(u)$ является максимальным нильпотентным идеалом.

Если радикал является нулевым идеалом, то говорят о “кольце без радикала” (“полупростом кольце”, “системе Дедекинда”). Кольцо классов вычетов по радикалу всегда является кольцом без радикала.

\emph{Если $\mZ$ есть центр $\mo$, а $\mc$ -- радикал $\mo$, то $\mC=\mc\cap\mZ$ есть радикал $\mZ$.}

\emph{Доказательство.} Ясно, что $\mC$ нильпотентен. Если бы в $\mZ$ существовал еще больший нильпотентный идеал $\overline{\mC}$, то его можно было бы расширить до нильпотентного идеала $\overline{\mc}$, не целиком содержащегося в $\mc$, поскольку
\[
        (\overline{\mC}\mo)^\rho=\overline{\mC}^{\rho}\mo^\rho=0.
\]

\emph{Следствие.} Если $\mo$ является кольцом без радикала, то и $\mZ$ тоже. Обратное, однако, неверно, как показывает наш пример в §10. Центр $\mo/\mc$ содержит кольцо, изоморфное $\mZ/\mC$; но оно может быть собственным подкольцом (пример в §10).
% END UNIT BODY
% END PRODUCER UNIT 197
\endgroup


\clearpage
\begingroup
\setcounter{footnote}{0}
% BEGIN PRODUCER UNIT 198
% Source: C:/Users/Floris/Documents/interlanguage/03_projects/noether/02_slavic_working_corpus/translations/paper34/section13/russian/v001/Noether_Paper34_Section13_Russian_v001.tex
% Identity: 9481 B / SHA-256 7A2560E719F0E297E4156B99433DEE38FA69B2CB07494327E71DAF5070C52AFC
\setlength{\parindent}{1.15em}
\setlength{\parskip}{0pt}
\makeatletter
\renewcommand\footnotesize{\@setfontsize\footnotesize{7.5}{8.5}\selectfont}
\makeatother
\providecommand{\ma}{\mathfrak{a}}
\providecommand{\mc}{\mathfrak{c}}
\providecommand{\md}{\mathfrak{d}}
\providecommand{\ml}{\mathfrak{l}}
\providecommand{\mo}{\mathfrak{o}}
\providecommand{\mr}{\mathfrak{r}}
\providecommand{\mt}{\mathfrak{t}}
\emergencystretch=120em
\allowdisplaybreaks
\sloppy
\hbadness=10000
\vbadness=10000
\hfuzz=2pt
\vfuzz=10pt
\raggedbottom
% BEGIN UNIT BODY
% Unit: Paper34 section13 v001. Direct Russian translation from audited German source lines 772--882, segments P34-S0150--P34-S0166. Source scan pages 23--25 retained as witnesses; Section14 heading P34-S0167 is intentionally reserved for Section14. Zenodo record 20822156 was rechecked for this unit at 2026-06-24T23:34Z UTC with no file key/size/checksum change.
\section*{34. Гиперкомплексные величины и теория представлений}
\emph{Продолжение: глава II, §§8--14.}

\section*{Глава II. Некоммутативная теория идеалов}

\subsection*{§ 13. Вполне приводимые кольца}

Согласно §5, кольцо называется справа вполне приводимым, если оно является прямой суммой конечного числа простых правых идеалов.

\emph{Справа вполне приводимое кольцо с единичным элементом не имеет нильпотентного идеала $\ne(0)$, а значит, не имеет радикала.}

\emph{Доказательство.} Каждый правый идеал $\mr$ является прямым слагаемым (§5):
\[
        \mo=\mt+\mr=e_1\mo+e_2\mo.
\]
Так как $e_2^2=e_2$, то также $e_2^\rho=e_2$. Если бы $\mr$ был нильпотентным, то было бы $e_2^\rho=0$, значит $e_2=0$, значит $\mr=0$, q.e.d.

Теперь докажем оба обратных утверждения. Во-первых: \emph{кольцо без радикала, удовлетворяющее условию минимальности для правых идеалов, вполне приводимо относительно правых идеалов.}

\emph{Доказательство.} Пусть $\mt$ -- минимальный правый идеал $\ne0$. Мы хотим показать, что $\mt$ содержит идемпотентный элемент.

Так как $\mt^2\subseteq\mt$ и $\mt^2\ne0$, следует $\mt^2=\mt$, поскольку $\mt^2$ является правым идеалом. Следовательно, существует $a$ в $\mt$ такое, что $a\mt\ne0$. Но тогда должно быть $a\mt=\mt$, поскольку $a\mt$ является правым идеалом.

Совокупность всех $b$ в $\mt$, аннулируемых элементом $a$ ($ab=0$), является правым идеалом и не равна $\mt$, следовательно равна $0$. Итак, из $ab=0$ следует $b=0$.

Из $\mt=a\mt$ следует, что $a$ можно представить в виде $ac$: $a=ac$, $c\ne0$. Отсюда
\[
        ac=ac^2,
        \qquad a(c-c^2)=0,
        \qquad c-c^2=0,
        \qquad c^2=c.
\]
Далее покажем, что $\mt$ является прямым слагаемым. $c\mo$ есть правый идеал $\subseteq\mt$, $\ne0$, так как $c^2=c$ принадлежит $c\mo$; значит, $c\mo=\mt$. Каждый элемент $r$ из $\mo$ допускает представление
\[
        r=cr+(r-cr) \qquad \hbox{(одностороннее разложение Пирса).}
\]
Элементы $cr$ образуют идеал $\mt$; элементы $r-cr$ образуют другой правый идеал $\mr$, который аннулируется элементом $c$: $c\mr=0$. Поэтому
\[
        \mo=(\mt,\mr).
\]
Так как элементы из $\mt$ не аннулируются элементом $c$, эта сумма прямая:
\begin{equation}
        \mo=\mt+\mr.                                             \tag{1}
\end{equation}

Если теперь на основании условия минимальности представить $\mo$ как прямую сумму прямо неразложимых идеалов:
\[
        \mo=\mr_1+\cdots+\mr_n,
\]
то идеалы $\mr_i$ должны быть простыми (= минимальными). Действительно, если бы, например, $\mr_1$ не был простым, а $\mt$ был минимальным идеалом в $\mr_1$, то из (1) имели бы
\[
        \mr_1=\mt+\mr\cap\mr_1 \qquad \hbox{\(\S 4,3\)};
\]
значит, $\mr_1$ был бы разложим, что невозможно. Следовательно, $\mo$ вполне приводимо.

Во-вторых: \emph{кольцо без радикала с условием минимальности имеет единичный элемент.}

Чтобы сначала показать существование левой единицы, достаточно доказать следующее: \emph{если правый идеал $\ma=c_1\mo\ne\mo$ имеет левую единицу $c_1$, то существует правый идеал $c\mo$ большей длины, который также имеет левую единицу $c$.} Поскольку полная приводимость, то есть ограниченность всех длин, уже доказана, этим за конечное число шагов приходим к левой единице для самого $\mo$.

Пусть, следовательно, $\ma=c_1\mo$ и $c_1^2=c_1$. Формула
\[
        r=c_1r+(r-c_1r)
\]
дает, как выше, разложение Пирса $\mo=c_1\mo+\mr$. Пусть, как выше, $c_2$ -- идемпотентный элемент из $\mr$, то есть $c_2^2=c_2$, $c_1c_2=0$ (так как $c_1$ аннулирует все элементы из $\mr$). Положив $e_1=c_1$, $e_2=c_2-c_2c_1$, получаем
\[
        e_1^2=e_1,
        \quad e_1e_2=0,
        \quad e_2e_1=0,
        \quad e_2c_2=e_2,
        \quad e_2\ne0,
        \quad e_2^2=e_2.
\]
По замечанию, сделанному в примечании, отсюда $c=e_1+e_2$ становится левой единицей для кольца
\[
        c\mo=e_1\mo+e_2\mo,
\]
которое, поскольку $e_2^2=e_2\ne0$, как правый идеал имеет большую длину, чем $\ma$.

Чтобы показать, что построенная левая единица $e$ является также правой единицей, выполним разложение Пирса в левые идеалы:
\[
        \mo=\mo e+\ml.
\]
Имеем $\ml e=0$, $\ml=e\ml$, значит $\ml^2=\ml e\ml=0$, значит, поскольку по предположению нильпотентного левого идеала существовать не может, $\ml=0$. Следовательно, $e$ является также правой единицей, а тем самым единицей вообще.

\emph{Теорема. Из правосторонней полной приводимости и существования единицы следует двусторонняя:}
\[
        \mo=\md_1+\cdots+\md_s,
        \qquad \md_i \hbox{ двусторонне просты.}
\]

Как и в предыдущем доказательстве, достаточно показать, что каждый минимальный (двусторонний) идеал $\ma$ является прямым слагаемым. Как правый идеал $\ma$ является прямым слагаемым:
\[
        \mo=\ma+\mr=e_1\mo+e_2\mo.
\]
Соответствующее левое разложение имеет вид
\[
        \mo=\mc+\ml=\mo e_1+\mo e_2.
\]
Следует
\[
        \mr\mc\subseteq \mr\ma\leq \mr\cap\ma=0,
        \qquad
        \ml\ma=\mo e_2e_1\mo=0,
\]
\[
        (\ma\ml)^2=\ma\ml\ma\ml=0,
        \qquad \hbox{значит } \ma\ml=0;
\]
\[
        \mr=\mr\mo=(\mr\mc,\mr\ml)=\mr\ml,
        \qquad
        \ml=\mo\ml=(\ma\ml,\mr\ml)=\mr\ml.
\]
Итак, $\mr=\ml$, значит $\mr$ двусторонний; следовательно, $\ma$ является двусторонним прямым слагаемым. Далее доказательство идет так же, как доказательство односторонней полной приводимости.

Идеалы $\md_i$ также двусторонне просты как кольца, поскольку все идеалы в $\md_i$ являются идеалами в $\mo$. Теперь исследуем их структуру.
% END UNIT BODY
% END PRODUCER UNIT 198
\endgroup


\clearpage
\begingroup
\setcounter{footnote}{0}
% BEGIN PRODUCER UNIT 199
% Source: C:/Users/Floris/Documents/interlanguage/03_projects/noether/02_slavic_working_corpus/translations/paper34/section14/russian/v001/Noether_Paper34_Section14_Russian_v001.tex
% Identity: 14595 B / SHA-256 88B08F0B58A7F61E50DF039736E4FDC5E2F1874049D89F98FF732E8879DC377C
\setlength{\parindent}{1.15em}
\setlength{\parskip}{0pt}
\makeatletter
\renewcommand\footnotesize{\@setfontsize\footnotesize{7.5}{8.5}\selectfont}
\makeatother
\providecommand{\mA}{\mathfrak{A}}
\providecommand{\ml}{\mathfrak{l}}
\providecommand{\mo}{\mathfrak{o}}
\providecommand{\mr}{\mathfrak{r}}
\providecommand{\mZ}{\mathfrak{Z}}
\emergencystretch=120em
\allowdisplaybreaks
\sloppy
\hbadness=10000
\vbadness=10000
\hfuzz=2pt
\vfuzz=10pt
\raggedbottom
% BEGIN UNIT BODY
% Unit: Paper34 section14 v001. Direct Russian translation from audited German source lines 884--1116, segments P34-S0167--P34-S0198. Source scan pages 25--29 retained as witnesses; next-chapter heading P34-S0199 and Section15 P34-S0200 are intentionally reserved for the next tranche. Zenodo record 20822156 was rechecked for this unit at 2026-06-25T00:02Z UTC with no file key/size/checksum change.
\section*{34. Гиперкомплексные величины и теория представлений}
\emph{Продолжение: глава II, §§8--14.}

\section*{Глава II. Некоммутативная теория идеалов}

\subsection*{§ 14. Двусторонне простые вполне приводимые кольца с единичным элементом}

Пусть $\mo$ -- такое кольцо:
\[
        \mo=\mr_1+\cdots+\mr_n=e_1\mo+\cdots+e_n\mo,
        \qquad \mr_i\hbox{ просты}.
\]
Тогда следует
\[
        \mo=\ml_1+\cdots+\ml_n=\mo e_1+\cdots+\mo e_n,
        \qquad \ml_i\hbox{ прямо неразложимы }\hbox{\(\S 9\)}.
\]
$\ml_i\mr_i$ является двусторонним идеалом $\ne0$ (так как $e_i^2$ лежит в $\ml_i\mr_i$), значит, равен $\mo$.

Положив $\mA_{ik}=\mr_i\ml_k$, получаем прямую сумму
\[
        \mo=(\mr_1,\ldots,\mr_n)\cdot(\ml_1,\ldots,\ml_n)
        = (\ldots,\mA_{ij},\ldots).
\]
Действительно, если $0=\sum_{i,j}a_{ij}$, то, так как $a_{ij}$ лежит в $\mr_i$, а $\sum\mr_i$ является прямой суммой,
\[
        0=\sum_j a_{ij},
\]
и далее, поскольку $a_{ij}$ лежит в $\ml_j$,
\[
        0=a_{ij}.
\]

Далее: $\mA_{ij}\mr_j=\mr_i\ml_j\mr_j=\mr_i\mo=\mr_i$, значит $\mA_{ij}\ne0$. Пусть $a_{ij}\ne0$ лежит в $\mA_{ij}$. Тогда
\[
        a_{ij}\mr_j\subseteq\mr_i,
        \qquad
        a_{ij}\mr_j\ne0 \hbox{, так как } a_{ij}e_j=a_{ij},
\]
следовательно
\[
        a_{ij}\mr_j=\mr_i.
\]
Если для каждого $r_j$ из $\mr_j$ положить
\[
        a_{ij}r_j=r_i,
\]
то отображение $r_j\mapsto r_i$ является операторным гомоморфизмом из $\mr_j$ в $\mr_i$. Элементы, которым сопоставляется нуль, образуют идеал $\subseteq\mr_j$, значит нулевой идеал; следовательно, это \emph{изоморфизм}. Тем самым:

\emph{Любые два простых правых идеала $\mr_i$ и $\mr_j$, встречающиеся в одном представлении $\mo$, операторно изоморфны; изоморфизмы осуществляются элементами из $\mA_{ij}$.} Так как любые два различных простых правых идеала $\mr,\mr'$ всегда встречаются вместе хотя бы в одном представлении $\mo$, то любые два таких идеала также изоморфны.

\emph{Все гомоморфизмы из $\mr_j$ в $\mr_i$ осуществляются элементами из $\mA_{ij}$.}

\emph{Доказательство.} Если $e_j\mapsto a_{ij}$, то из $e_j^2\mapsto a_{ij}e_j$ следует, что $a_{ij}=a_{ij}e_j$ лежит в $\mr_i\ml_j=\mA_{ij}$; далее
\[
        \mr_j=e_j\mr_j\mapsto a_{ij}\mr_j=\mr_i.
\]

\emph{В частности, гомоморфизмы $\mr_i$ в себя осуществляются через $\mA_{ii}$.}

Произведению двух элементов из $\mA_{ii}$ соответствует произведение гомоморфизмов, а сумме -- сумма. (Определение см. §1, 6.) Различные $a_{ii}$ дают различные автоморфизмы, поскольку $e_i$ переходит в $a_{ii}e_i=a_{ii}$. Следовательно, кольцо $\mA_{ii}$ изоморфно кольцу автоморфизмов $\mr_i$.

Но кольцо автоморфизмов простого идеала является телом (§2), значит $\mA_{ii}$ является телом. Так как далее все $\mr_i$ операторно изоморфны, то и их кольца автоморфизмов изоморфны как кольца: $\mA_{ii}$ однозначно определяется кольцом $\mo$ с точностью до изоморфизма колец. Различные возможные $\mA_{ii}$ являются лишь различными конкретными реализациями абстрактно определенного кольца автоморфизмов простых правых идеалов.

\emph{Основная теорема. Каждое вполне приводимое двусторонне простое кольцо $\mo$ с единичным элементом изоморфно кольцу матриц $n$-й степени над телом $K$. (Тело $K$ изоморфно телу автоморфизмов правых идеалов кольца $\mo$.)}

\emph{Доказательство. Построение матричных единиц $c_{ik}$.} Пусть $\Gamma_{11}$ -- тождественный автоморфизм $\mr_1$, а $\Gamma_{i1}$ -- произвольный изоморфизм $\Gamma_{i1}\mr_1=\mr_i$; наконец положим
\[
        \Gamma_{ik}=\Gamma_{i1}\Gamma_{k1}^{-1}.
\]
Тогда вообще
\[
        \Gamma_{ik}\Gamma_{kl}=\Gamma_{il}.
\]
Если $\Gamma_{ik}$ осуществляется элементом $c_{ik}$ из $\mA_{ik}$, то далее следует:
\[
        c_{ik}=e_i c_{ik}=c_{ik}e_k,
        \qquad
        c_{ik}c_{kl}e_l=c_{il}e_l,
\]
значит
\[
        c_{ik}c_{kl}=c_{il},
        \qquad
        c_{ij}c_{kl}=c_{ij}e_j e_k c_{kl}=0\quad(j\ne k).
\]
Следовательно, $c_{ik}$ являются матричными единицами (§7).

\footnotetext{Недавно Б. Л. ван дер Варден нашел более простое доказательство, работающее непосредственно с абстрактным телом автоморфизмов; оно должно быть приведено в его книге по алгебре (Grundlehren d. Math. Wiss., Berlin: Julius Springer). [11. 6. 29.]}

\emph{Построение $K$.} Соответствие
\[
        a_{ij}=c_{i1}a_{11}c_{1i}
\]
ставит каждому $a_{11}$ из $\mA_{11}$ в соответствие ``сопряженный элемент'' $a_{ii}$ из $\mA_{ii}$.

Это соответствие является изоморфизмом колец, поскольку изоморфизмы $\Delta,\Gamma$, соответствующие $a$ и $c$, связаны формулой
\[
        \Delta_{ii}=\Gamma_{i1}\Delta_{11}\Gamma_{i1}^{-1},
\]
а такое соотношение, как известно, задает изоморфизм колец. Теперь образуем элементы
\[
        \alpha=a_{11}+\cdots+a_{nn}=a_{11}+c_{21}a_{11}c_{12}+\cdots+c_{n1}a_{11}c_{1n}.
\]
Каждому $a_{11}$ соответствует единственный $\alpha$ (взаимно однозначно, поскольку сумма прямая). Совокупность всех $\alpha$ является телом $K\sim\mA_{11}$, ибо из
\[
        \alpha=a_{11}+\cdots+a_{nn},
        \qquad
        \beta=b_{11}+\cdots+b_{nn},
\]
следует
\[
        \alpha+\beta=(a_{11}+b_{11})+\cdots+(a_{nn}+b_{nn}),
        \qquad
        \alpha\beta=a_{11}b_{11}+\cdots+a_{nn}b_{nn};
\]
значит, соответствие $a_{11}\mapsto\alpha$ является изоморфизмом.

Далее
\[
        c_{ik}\alpha=c_{ik}a_{kk}=c_{ik}c_{k1}a_{11}c_{1k}=c_{i1}a_{11}c_{1k},
\]
\[
        \alpha c_{ik}=a_{ii}c_{ik}=c_{i1}a_{11}c_{1i}c_{ik}=c_{i1}a_{11}c_{1k},
\]
следовательно, каждое $\alpha$ перестановочно со всеми $c_{ik}$.

Наконец, из тех же формул следует
\[
        c_{ik}K=c_{ik}\mA_{kk}=c_{ik}\mr_k\ml_k=\mr_i\ml_k=\mA_{ik},
\]
а значит
\[
        \mo=\sum c_{ik}K.
\]

\emph{Соответствие.} Если каждый элемент $\mo$ представить в виде $\sum c_{ik}\alpha_{ik}$ и сопоставить этому элементу матрицу $(\alpha_{ik})$, то соответствие становится изоморфизмом, что сразу следует из определения умножения матриц. Тем самым основная теорема доказана.

\emph{Обратное утверждение. Кольцо матриц $n$-й степени над телом $A$ является двусторонне простым и вполне приводимым; $A$ изоморфно кольцу автоморфизмов правых идеалов, а $Ae$ есть то $K$, которое построено в предыдущей теореме.}

\emph{Доказательство.} Пусть матричные единицы (§7) равны $c_{ik}$. Положим
\[
        \mr_i=\sum_k c_{ik}A.
\]
Тогда очевидно $\mo=\mr_1+\cdots+\mr_n$. Идеалы $\mr_i$ являются простыми правыми идеалами, поскольку каждый ненулевой элемент $\sum c_{ik}\alpha_k$ порождает весь $\mr_i$. Действительно, если $\alpha_j\ne0$, то
\[
        \Bigl(\sum c_{ik}\alpha_k\Bigr)\cdot(c_{jl}\alpha_j^{-1})=c_{il},
\]
и $\sum c_{il}\beta_l$ пробегает все элементы $\mr_i$.

Идеалы $\mr_i$ операторно изоморфны: $\mr_i=c_{ik}\mr_k$.

Отсюда следует: $\mo$ двусторонне неразложимо, следовательно, по теореме §13 и на основании уже доказанной полной приводимости, а также поскольку существует единичный элемент, оно двусторонне просто (это видно и из того, что любой $\sum\sum c_{ik}\alpha_{ik}\ne0$ порождает весь двусторонний идеал $\mo$).

Далее
\[
        \mr_i=c_{ii}\mo,
        \qquad
        \ml_i=\mo c_{ii},
\]
\[
        e=\sum c_{ii},
        \qquad c_{ii}=e_i,
\]
\[
        \mA_{ii}=\mr_i\cap\ml_i=c_{ii}A\sim A.
\]
Наконец, положив $a_{11}=c_{11}\lambda$, получаем
\[
        \alpha=a_{11}+c_{21}a_{11}c_{12}+\cdots
        =c_{11}\lambda+c_{22}\lambda+\cdots=e\lambda,
\]
что и доказывает все утверждение.

\emph{Центр $\mo$ является центром $K$, а значит телом.}

\emph{Доказательство.} $\mZ(K)$ состоит из всех $\zeta$, перестановочных со всеми $x$. Но они перестановочны также со всеми $c_{ik}$, а значит и с суммами $\sum c_{ik}a_{ik}$. Следовательно:
\[
        \mZ(K)\subseteq\mZ(\mo).
\]
Обратно, пусть $z$ лежит в $\mZ(\mo)$:
\[
        z=\sum_\mu\sum_\nu c_{\mu\nu}\gamma_{\mu\nu}.
\]
\[
        zc_{ik}=c_{ik}z,
        \qquad
        \sum_\mu c_{\mu k}\gamma_{\mu i}=\sum_\nu c_{i\nu}\gamma_{k\nu}.
\]
Значит,
\[
        \gamma_{\mu i}=0\ (\mu\ne i),
        \qquad
        c_{ik}\gamma_{ii}=c_{ik}\gamma_{kk},
        \qquad
        \gamma_{ii}=\gamma_{kk}=\gamma;
\]
\[
        z=\sum c_{ii}\gamma=e\gamma=\gamma.
\]
Итак, $z$ лежит в $K$ и перестановочен со всеми $x$, значит $z$ лежит в $\mZ(K)$.

Поскольку матричное кольцо, конечно, также вполне приводимо слева, а каждое справа вполне приводимое кольцо с единичным элементом является прямой суммой матричных колец, следует:

\emph{Каждое справа вполне приводимое кольцо с единичным элементом является также слева вполне приводимым, и наоборот.}

И еще: \emph{центр вполне приводимого кольца с единичным элементом является прямой суммой коммутативных тел, соответствующих двусторонне простым матричным кольцам.}

Остается следующая теорема:

\emph{Любые два различных разложения $\mo=\sum c_{ik}K$, $\mo=\sum c'_{ik}K'$ переходят друг в друга посредством внутренних автоморфизмов $\alpha'\mapsto x^{-1}\alpha'x$, $c'_{ik}\mapsto x^{-1}c_{ik}x$.}

\emph{Доказательство.} Пусть
\[
        \mo=\mr_1+\cdots+\mr_n=\mr_1'+\cdots+\mr_n',
\]
и пусть $a,b$ -- элементы, осуществляющие два взаимно обратных изоморфизма идеалов $\mr_1,\mr_1'$:
\[
        \mr_1=a\mr_1',
        \qquad
        \mr_1'=b\mr_1,
        \qquad
        ab=c_{11},
        \qquad
        ba=c'_{11}.
\]
Положим
\[
        x=\sum_i c_{i1}ac'_{1i},
        \qquad
        y=\sum_i c'_{i1}bc_{1i}.
\]
Тогда
\[
        xy=e,
        \qquad
        y=x^{-1},
        \qquad
        x^{-1}c_{ik}x=c'_{ik},
        \qquad
        x^{-1}\alpha x=\alpha'.
\]
% END UNIT BODY
% END PRODUCER UNIT 199
\endgroup


\clearpage
\begingroup
\setcounter{footnote}{0}
% BEGIN PRODUCER UNIT 200
% Source: C:/Users/Floris/Documents/interlanguage/03_projects/noether/02_slavic_working_corpus/translations/paper34/section15/russian/v001/Noether_Paper34_Section15_Russian_v001.tex
% Identity: 9949 B / SHA-256 1929C9DF145A8D213E56460B2079E418EC16ED27CE7C0075DD39BED640D038CF
\setlength{\parindent}{1.15em}
\setlength{\parskip}{0pt}
\makeatletter
\renewcommand\footnotesize{\@setfontsize\footnotesize{7.5}{8.5}\selectfont}
\makeatother
\providecommand{\mo}{\mathfrak{o}}
\providecommand{\mM}{\mathfrak{M}}
\emergencystretch=120em
\allowdisplaybreaks
\sloppy
\hbadness=10000
\vbadness=10000
\hfuzz=2pt
\vfuzz=10pt
\raggedbottom
% BEGIN UNIT BODY
% Unit: Paper34 section15 v001. Direct Russian translation from audited German source lines 1121--1245, segments P34-S0199--P34-S0211; P34-S0220/Section16 reserved for next tranche. Zenodo record 20822156 checked for this unit at 2026-06-25T00:45Z UTC with no file key/size/checksum change.
\section*{34. Гиперкомплексные величины и теория представлений}
\emph{Продолжение и завершение: глава III, §§15--19, и глава IV, §§20--26.}

\section*{Глава III. Теория модулей и представлений}

\subsection*{§ 15. Представления и модули представлений}

Пусть $\mo$ -- кольцо, а $K$ -- кольцо с единичным элементом. (В последующих приложениях $K$ всегда является телом.)

Представлением $n$-й степени кольца $\mo$ в $K$ называется гомоморфия
\[
        \mo\sim\mathfrak D,
\]
где $\mathfrak D$ есть кольцо матриц $n$-й степени над $K$.

Под модулем представления кольца $\mo$ относительно $K$ понимается бимодуль $\mM$ (§1, 5), который является левым $\mo$-модулем и правым $K$-модулем:
\[
        \mo\mM\subseteq\mM,
        \qquad
        \mM K\subseteq\mM,
\]
и который, далее, является прямой суммой конечного числа однопорождённых $K$-модулей:
\[
        \mM=x_1K+\cdots+x_nK;
\]
причём единичный элемент $K$ является единичным оператором.

Каждый модуль представления приводит к представлению. Пусть $c$ принадлежит $\mo$ и
\[
        cx_k=\sum_i x_i\gamma_{ik},
\]
или
\[
        c(x_1,\ldots,x_n)=(x_1,\ldots,x_n)C,
        \qquad C=(\gamma_{ik}).
\]
Тогда матрицы $C$ образуют представление элемента $c$, ибо
\[
        (b+c)x_k=\sum_i x_i(\beta_{ik}+\gamma_{ik}),
\]
\[
\begin{aligned}
        bcx_k
        &=b\sum_jx_j\gamma_{jk}
          =\sum_j b x_j\gamma_{jk}
          =\sum_j\sum_i x_i\beta_{ij}\gamma_{jk}  \\
        &=\sum_i x_i\biggl(\sum_j\beta_{ij}\gamma_{jk}\biggr),
\end{aligned}
\]
или
\[
        bc(x_1,\ldots,x_n)=(x_1,\ldots,x_n)BC.
\]

Обратно: каждое представление $\mathfrak D$ кольца $\mo$ принадлежит модулю представления, а именно некоторому определённому базису этого модуля. А именно, под $\mM$ понимается совокупность формальных линейных форм от $x_1,\ldots,x_n$,
\[
        y=\sum_i x_i\alpha_i.
\]
Тогда $\mM$ является правым $K$-модулем. Далее, если элементу $c$ поставлена в соответствие матрица $C=(\gamma_{ik})$, положим
\[
        cx_k=\sum_i x_i\gamma_{ik},
        \qquad
        c\biggl(\sum_k x_k\alpha_k\biggr)=\sum_i x_i\sum_k\gamma_{ik}\alpha_k.
\]
Из гомоморфных соотношений
\[
        c+d\sim C+D,
        \qquad cd\sim CD
\]
следует
\[
        (c+d)x_i=cx_i+dx_i,
        \qquad cdx_i=c(dx_i),
\]
и то же самое для сумм $\sum_i x_i\alpha_i$. Остальные свойства бимодуля
\[
        c(y+z)=cy+cz,
        \qquad c(y\alpha)=(cy)\alpha
\]
тривиальны. Следовательно, $\mM$ есть модуль представления, который по самой конструкции точно принадлежит представлению $\mo\sim\mathfrak D$.

Пусть теперь $(y_1,\ldots,y_n)$ -- другой базис того же модуля представления:
\[
        (y_1,\ldots,y_n)=(x_1,\ldots,x_n)P,
        \qquad
        (x_1,\ldots,x_n)=(y_1,\ldots,y_n)P^{-1}.
\]
Если $b\sim B$ относительно базиса $x$, то
\[
        b(y_1,\ldots,y_n)=b(x_1,\ldots,x_n)P
        =(x_1,\ldots,x_n)BP
        =(y_1,\ldots,y_n)P^{-1}BP.
\]
Представления $b\sim B$ и $b\sim P^{-1}BP$ называются эквивалентными и относятся к одному и тому же классу представлений. Так как всякая регулярная матрица $P$ снова переводит базис $\mM$ в базис, доказано:

\emph{Каждый модуль представления однозначно приводит к определённому классу представлений.}\textsuperscript{15a)}

{\footnotesize\noindent 15) Модули представлений относительно коммутативных тел впервые встречаются у В. Крулля, \emph{Theorie und Anwendung der verallgemeinerten Abelschen Gruppen}, Sitzungsber. Heidelberger Ak. 1926, 1-я работа.\par
15a) Добавлено при корректуре (14 апреля 1929 г.). Как сообщил мне Б. Л. ван дер Варден, связь, не зависящую от специального выбора базиса, то есть инвариантную связь, можно получить непосредственно путём разделения понятий: линейное преобразование и матрица. Линейное преобразование есть гомоморфизм двух модулей линейных форм; матрица есть выражение, то есть представление, этого гомоморфизма при определённом выборе базиса. I. Каждый модуль представления однозначно сопоставляет области мультипликаторов $\mo$ гомоморфную систему линейных преобразований модуля в себя. Действительно, по концу §2 левые мультипликаторы бимодуля $\mM$ порождают операторные гомоморфизмы $\mM$ в себя относительно правых операторов (здесь умножения на $K$), и это сопоставление гомоморфно. II. Каждая гомоморфная $\mo$ система линейных преобразований модуля линейных форм в себя приводит к модулю представления.\par}

Ясно: двум операторно изоморфным модулям представлений соответствует один и тот же класс представлений. Но и обратно: если два модуля представлений $(x_1,\ldots,x_n)$ и $(\bar x_1,\ldots,\bar x_n)$ порождают одно и то же представление, то соответствие
\[
        x_i\mapsto\bar x_i,
        \qquad
        \sum_i x_i\alpha_i\mapsto\sum_i\bar x_i\alpha_i
\]
задаёт изоморфизм. Ведь правило умножения полностью определяется матрицами $C$.

Специальный случай: если два базиса $\mM$ порождают одно и то же представление, то они переходят друг в друга через операторный изоморфизм $\mM$ на себя.

Или: если $C=PCP^{-1}$ для всех $C$ и фиксированного $P$, то соответствие
\[
        y_i\mapsto x_i,
        \qquad y_i=\sum_kx_k\pi_{ik},
\]
является изоморфизмом $\mM$ на себя. Обратно: каждый изоморфизм $y_i\mapsto x_i$ бимодуля $\mM$ на себя осуществляется матрицей $P$, перестановочной со всеми $C$.

Понятие представления можно распространить и на такие кольца, которые ещё коммутативно связаны с коммутативной областью мультипликаторов $P$ (§8). В этом случае берут только такие кольца-носители представлений $K$, которые содержат $P$ в своём центре, и требуют от представления не только того, чтобы оно было гомоморфизмом колец $\mo\sim\mathfrak D$, но даже того, чтобы оно было операторным гомоморфизмом: из $a\sim A$ должно следовать
\[
        a\rho\sim A\rho\qquad(\rho\in P).
\]
Для модулей представлений это требование означает дополнительное правило вычисления
\[
        am\cdot\rho=a(m\rho)=a\rho\cdot m.
\]
Модуль и кольцо тогда, как говорят, ``коммутативно связаны с $P$''.
% END UNIT BODY
% END PRODUCER UNIT 200
\endgroup


\clearpage
\begingroup
\setcounter{footnote}{0}
% BEGIN PRODUCER UNIT 201
% Source: C:/Users/Floris/Documents/interlanguage/03_projects/noether/02_slavic_working_corpus/translations/paper34/section16/russian/v001/Noether_Paper34_Section16_Russian_v001.tex
% Identity: 7581 B / SHA-256 952E7844A3C9C82F316F2300169C8AFFFD418249D9A04D746521A014E0AB4A1A
\setlength{\parindent}{1.15em}
\setlength{\parskip}{0pt}
\makeatletter
\renewcommand\footnotesize{\@setfontsize\footnotesize{7.5}{8.5}\selectfont}
\makeatother
\providecommand{\mo}{\mathfrak{o}}
\providecommand{\mM}{\mathfrak{M}}
\providecommand{\mU}{\mathfrak{U}}
\providecommand{\mA}{\mathfrak{A}}
\providecommand{\mB}{\mathfrak{B}}
\emergencystretch=120em
\allowdisplaybreaks
\sloppy
\hbadness=10000
\vbadness=10000
\hfuzz=2pt
\vfuzz=10pt
\raggedbottom
% BEGIN UNIT BODY
% Unit: Paper34 section16 v001. Direct Russian translation from audited German source lines 1247--1310, segments P34-S0212--P34-S0219; P34-S0220/Section17 reserved for next tranche. Zenodo record 20822156 checked for this unit at 2026-06-25T01:59Z UTC with no file key/size/checksum change against the Section15 preflight.
\section*{34. Гиперкомплексные величины и теория представлений}
\emph{Продолжение и завершение: глава III, §§15--19, и глава IV, §§20--26.}

\section*{Глава III. Теория модулей и представлений}

\subsection*{§ 16. Приводимые представления}

Если $\mU$ является подмодулем модуля представления $\mM$ и если можно выбрать базис для $\mM$, состоящий из базиса $z_1,\ldots,z_t$ модуля $\mU$, дополненного элементами $y_1,\ldots,y_r$, то есть
\[
        \mM=y_1K+\cdots+y_rK+z_1K+\cdots+z_tK,\textsuperscript{16)}
\]
то представления имеют вид
\[
        C=\begin{pmatrix}R&0\\ S&T\end{pmatrix},
\]
где матрицы $T$ сами по себе образуют представление $t$-й степени, порождённое $\mU$, а матрицы $R$ образуют представление $r$-й степени, порождённое $\mM/\mU$.

Доказательство. Имеем
\[
        (cy_1,\ldots,cy_r,cz_1,\ldots,cz_t)
        =(y_1,\ldots,y_r,z_1,\ldots,z_t)C.
\]
Так как $cz_i$, будучи элементами $\mU$, выражаются только через $z_i$, в матрице $C$ справа вверху стоит нуль. Если остальные части $C$ в указанном порядке назвать $R,S,T$, то, в частности,
\[
        (cz_1,\ldots,cz_t)=(z_1,\ldots,z_t)T;
\]
следовательно, $T$ образуют представление, осуществлённое через $\mU$. Далее
\[
        (cy_1,\ldots,cy_r)\equiv (y_1,\ldots,y_r)R\pmod{\mU},
\]
между тем как $y_i$ образуют базис, линейно независимый по модулю $\mU$. Следовательно, $R$ образуют представление, порождённое $\mM/\mU$.

Обратно, если дано ``приводимое'' представление
\[
        C=\begin{pmatrix}R&0\\ S&T\end{pmatrix},
\]
где $R$ и $T$ квадратные матрицы, то в соответствующем модуле представления произведения каждого $c$ с последними $t$ базисными элементами $z_1,\ldots,z_t$ выражаются только через них, то есть $\mU=(z_1,\ldots,z_t)$ является подмодулем.

{\footnotesize\noindent 16) Иначе говоря: когда $\mU$ как $K$-модуль является прямым слагаемым. Если $K$ является телом, это предположение всегда выполнено.\par}

Следствие. Пусть
\[
        \mM\supset \mU_1\supset \mU_2\supset\cdots\supset \mU_e=0
\]
есть композиционный ряд для $\mM$, и пусть каждое $\mU_i$ в предыдущем является прямым слагаемым как $K$-модуль, так что можно выбрать базис
\[
        z_{11},\ldots,z_{1r_1},\ z_{21},\ldots,z_{2r_2},\ldots,
        z_{e1},\ldots,z_{er_e}
\]
так, что $z_{ik}$ с $i>\nu$ образуют базис для $\mU_\nu$. Тогда представления, осуществлённые через $\mM$, имеют вид
\[
 C=\begin{pmatrix}
 R_{11}&0&\cdots&0\\
 R_{12}&R_{22}&\cdots&0\\
 \vdots&\vdots&\ddots&\vdots\\
 R_{1e}&R_{2e}&\cdots&R_{ee}
 \end{pmatrix},                                      \tag{2}
\]
и $R_{\nu\nu}$ образуют представления, осуществлённые через композиционные факторы $\mU_{\nu-1}/\mU_\nu$.

Отдельные представления $R_{\nu\nu}$ неприводимы, так как композиционные факторы $\mU_{\nu-1}/\mU_\nu$ просты. Если предположить, что $K$ является телом, то и обратно всякое представление, максимально приведённое к виду (2), приводит к композиционному ряду. По теореме Жордана--Гёльдера композиционные факторы однозначны с точностью до операторного изоморфизма; значит, $R_{\nu\nu}$ однозначно определены с точностью до эквивалентных представлений и до порядка.

Если $\mM=\mA+\mB$ есть прямая сумма двух модулей представления $\mA=(y_1,\ldots,y_r)$, $\mB=(z_1,\ldots,z_s)$, то представление, осуществлённое через базис $(y_1,\ldots,y_r,z_1,\ldots,z_s)$, очевидно имеет вид
\[
        C=\begin{pmatrix}R&0\\0&S\end{pmatrix},
\]
где $R$ означает представление элемента $c$, осуществлённое через $\mA$, а $S$ -- представление, осуществлённое через $\mB$, и обратно. Если сначала представить модуль $\mM$ как прямую сумму прямо неразложимых составляющих и для них построить композиционные ряды, как выше, то при подходящем выборе базиса представление имеет соответствующий блочный вид. Если снова $K$ является телом, то и обратно всякое максимально приведённое представление приводит к представлению модуля через прямо неразложимые слагаемые и композиционные факторы внутри них. По теореме Крулля и Отто Шмидта классы прямо неразложимых составляющих представления однозначно определены с точностью до порядка.

Если модуль $\mM$ вполне приводим, то все блоки состоят из одного неприводимого представления, и представление называется вполне приводимым.
% END UNIT BODY
% END PRODUCER UNIT 201
\endgroup


\clearpage
\begingroup
\setcounter{footnote}{0}
% BEGIN PRODUCER UNIT 202
% Source: C:/Users/Floris/Documents/interlanguage/03_projects/noether/02_slavic_working_corpus/translations/paper34/section17/russian/v001/Noether_Paper34_Section17_Russian_v001.tex
% Identity: 4496 B / SHA-256 6C396C8C6734E9D78D4CDEAEC0EA87FEABE3D0333FD032DFC9CBD26D6F72E8C1
\setlength{\parindent}{1.15em}
\setlength{\parskip}{0pt}
\makeatletter
\renewcommand\footnotesize{\@setfontsize\footnotesize{7.5}{8.5}\selectfont}
\makeatother
\providecommand{\mo}{\mathfrak{o}}
\providecommand{\ma}{\mathfrak{a}}
\providecommand{\mb}{\mathfrak{b}}
\providecommand{\mM}{\mathfrak{M}}
\emergencystretch=120em
\allowdisplaybreaks
\sloppy
\hbadness=10000
\vbadness=10000
\hfuzz=2pt
\vfuzz=10pt
\raggedbottom
% BEGIN UNIT BODY
% Unit: Paper34 section17 v001. Direct Russian translation from audited German source lines 1311--1352, segments P34-S0220--P34-S0224; P34-S0225/Section18 reserved for next tranche. Zenodo record 20822156 checked for this unit at 2026-06-25T03:04Z UTC with no file key/size/checksum change against the Section16 preflight.
\section*{34. Гиперкомплексные величины и теория представлений}
\emph{Продолжение и завершение: глава III, §§15--19, и глава IV, §§20--26.}

\section*{Глава III. Теория модулей и представлений}

\subsection*{§ 17. Прямое разложение модулей представлений для колец с единицей}

Пусть $\mM$ является $\mo$-модулем, а $\mo$ -- кольцом с единицей. Тогда
\[
        \mM=\mM_e+\mM_0,
\]
где на $\mM_e$ единица является единичным оператором, а на $\mM_0$ -- нулевым оператором. ($\mM_e$ и $\mM_0$ также являются правыми $K$-модулями, если $\mM$ им является.)

Доказательство. Всякое $m$ из $\mM$ представимо в виде
\[
        m=em+(m-em).
\]
Множество элементов $em$ обозначим через $\mM_e$, а множество элементов $m-em$ -- через $\mM_0$. Тогда $\mM_e$ и $\mM_0$, очевидно, являются аддитивными группами и правыми $K$-модулями, если $\mM$ им является. Далее
\[
        r em=erm,
\]
так что $\mM_e$ также является $\mo$-модулем; аналогично
\[
        r(m-em)=rm-rem=rm-rm,
\]
так что $\mM_0$ является $\mo$-модулем. Для элементов $em$ элемент $e$ является единичным оператором, а для элементов $(m-em)$ -- нулевым оператором. Поэтому только нуль принадлежит одновременно $\mM_e$ и $\mM_0$; сумма $\mM_e+\mM_0$ прямая.

Если речь идет о модулях представлений, то $\mM_0$ порождает тривиальное представление, в котором каждому элементу сопоставляется нуль. Отщепив слагаемое $\mM_0$, получаем представление, осуществленное через $\mM_e$, в котором единичному элементу сопоставлена единичная матрица. Поэтому на таких представлениях мы отныне можем и будем ограничиваться.

Пусть
\[
        \mo=\ma_1+\cdots+\ma_s
\]
есть прямая сумма двусторонних идеалов,
\[
        e=e_1+\cdots+e_s
\]
-- разложение $e$, и пусть $\mM$ является $\mo$-модулем, где $e$ является единичным оператором. Тогда
\[
        \mM=\ma_1\mM+\cdots+\ma_s\mM
\]
прямо. В самом деле, очевидно
\[
        \mM=\mo\mM=(\ma_1\mM,\ldots,\ma_s\mM).
\]
Положив $\mb_i=\ma_1+\cdots+\widehat{\ma_i}+\cdots+\ma_s$, получаем
\[
        \mM=(\ma_i\mM,\mb_i\mM),
\]
и эта сумма прямая, поскольку $e_i$ является единичным оператором на $\ma_i\mM$ и нулевым оператором на $\mb_i\mM$.

На основании этой теоремы мы обычно ограничиваемся представлениями двусторонне неразложимых колец.
% END UNIT BODY
% END PRODUCER UNIT 202
\endgroup


\clearpage
\begingroup
\setcounter{footnote}{0}
% BEGIN PRODUCER UNIT 203
% Source: C:/Users/Floris/Documents/interlanguage/03_projects/noether/02_slavic_working_corpus/translations/paper34/source_fidelity_section18_19/russian/v001/Noether_Paper34_Section18_19_SourceFidelity_Russian_v001.tex
% Identity: 16613 B / SHA-256 054F073565F7BA51EE52C826C1C60DF3CCA0501A94931F3AB9C61D49D719C555
\setlength{\parindent}{1.15em}
\setlength{\parskip}{0pt}
\makeatletter
\renewcommand\footnotesize{\@setfontsize\footnotesize{7.5}{8.5}\selectfont}
\makeatother
\providecommand{\mo}{\mathfrak{o}}
\providecommand{\ml}{\mathfrak{l}}
\providecommand{\mc}{\mathfrak{c}}
\providecommand{\mM}{\mathfrak{M}}
\emergencystretch=120em
\allowdisplaybreaks
\sloppy
\hbadness=10000
\vbadness=10000
\hfuzz=2pt
\vfuzz=10pt
\raggedbottom
% BEGIN UNIT BODY
% Unit: Paper34 Section18/19 source-fidelity remediation v001. Direct Russian translation from existing safe Section18 opening plus original scan-witness late Section18 tail and full Section19; supersedes compressed Section18 ending for this boundary.
\section*{34. Гиперкомплексные величины и теория представлений}
\emph{Продолжение и завершение: глава III, §§15--19, и глава IV, §§20--26.}

\section*{Глава III. Теория модулей и представлений}

\subsection*{§ 18. Теория модулей и представлений вполне приводимых колец}

Пусть $\mo$ является вполне приводимым и двусторонне простым, а $\mM$ -- конечным $\mo$-модулем, для которого единичный элемент $\mo$ является единичным оператором. Тогда $\mM$ вполне приводим, а простые составляющие операторно изоморфны простым левым идеалам $\ml_i$.

Доказательство. Пусть
\[
        \mo=\ml_1+\cdots+\ml_n,
        \qquad
        \mM=(\mo m_1,\ldots,\mo m_k),
\]
следовательно,
\[
        \mM=(\ldots,\ml_i m_k,\ldots).                     \tag{4}
\]
Те $\ml_i m_k$, которые отличны от нуля, изоморфны $\ml_i$, поскольку отображение $a\mapsto am_k$ является операторным изоморфизмом. Следовательно, $\ml_i m_k$ просты. Если из выражения (4) опустить те $\ml_i m_k$, которые уже содержатся в сумме предыдущих, то сумма становится прямой.

Следствие. Если, кроме того, $\mM$ прямо неразложим, то $\mM$ прост и изоморфен некоторому $\ml_i$.

Пусть теперь $\Delta$ -- тело автоморфизмов этого простого $\mM$, а $\Lambda$ -- тело автоморфизмов $\ml_i$; тогда на основании операторного изоморфизма между $\mM$ и $\ml_i$ имеем изоморфизм колец $\Delta\simeq\Lambda$. По §1 можно рассматривать $\mM$ и $\ml_i$ как бимодули с $\Delta$ соответственно $\Lambda$ в качестве правой области скаляров, причем эти бимодули также являются модулями представлений. Представления, осуществляемые через них, переходят друг в друга посредством изоморфизма колец $\Delta\simeq\Lambda$. Поэтому для изучения этого класса представлений можно ограничиться представлением, осуществляемым через $\ml_i$ в его теле автоморфизмов.

Специально для $\ml_1$ положим в основу базис из матричных единиц
\[
        \ml_1=(c_{11},\ldots,c_{n1})
\]
и в качестве реализации тела автоморфизмов $\Lambda$ возьмем ранее (§14) обозначенное через $K$ подтело $\mo$. При этом по сравнению с прежними рассуждениями правые идеалы следует заменить левыми, а автоморфизмы соответственно записывать справа. Если представить
\[
        c=\sum_{i,k} c_{ik}\alpha_{ik},
\]
то получаем
\[
        (cc_{11},\ldots,cc_{n1})=(c_{11},\ldots,c_{n1})(\alpha_{ik}).
\]

\paragraph{Восстановленный первоначальный конец §18.}
Если затем представить каждый элемент $a$ кольца $\mo$ в виде
\[
        a=\sum c_{ik}\alpha_{ik},
\]
то $a\mapsto(\alpha_{ik})$ есть представление в $K$, осуществляемое через $\ml_1$.

Пусть $\Gamma$ -- подтело $K$, такое что $K$ имеет конечный ранг $t$ относительно $\Gamma$:
\[
        K=\chi_1\Gamma+\cdots+\chi_t\Gamma .
\]
Тогда $K$ становится собственным модулем представления относительно $\Gamma$. Элементы $\beta$ из $K$ представляются матрицами $B$ над $\Gamma$, которые получают так:
\[
        \beta\chi_j=\sum_i\chi_i\beta_{ij},\qquad B=(\beta_{ij}).
\]
Если теперь рассматривать $\ml_1$ как модуль представления относительно подтела $\Gamma$, то получаем:

\emph{Представление $\mo$ в $\Gamma$, осуществляемое через $\ml_1$, получают так: в указанном выше представлении $\mo$ в $K$
\[
        a\mapsto(\alpha_{ik})
\]
элементы $\alpha_{ik}$ заменяют матрицами $A_{ik}$, соответствующими им в представлении $K$ в $\Gamma$.}

Доказательство.
\[
        \ml_1=\sum_i c_{i1}K=\sum_i\sum_j c_{i1}\chi_j\Gamma .
\]
Посредством базиса $(\ldots,c_{i1}\chi_j,\ldots)$ элемент $c_{ik}$ представляется блочной матрицей, которая в $(i,k)$-м блоке содержит единичную матрицу $E_t$, а во всех остальных блоках -- нулевые матрицы $0$ степени $t$. Точно так же элемент $\alpha$ из $K$ представляется блочно-диагональной матрицей с соответствующими $t$-рядными матрицами $A$. Следовательно, элементу
\[
        a=\sum c_{ik}\alpha_{ik}
\]
соответствует блочная матрица $(A_{ik})$. Переход к произвольному $\mM$ снова дает изоморфизм колец для представления.

Замечание. Изучаемые здесь представления в теле автоморфизмов левых идеалов и в его конечных подтелах, с точностью до изоморфизма, являются единственными, которые могут осуществляться простыми $\mo$-модулями или левыми идеалами. Действительно, согласно сделанному ранее (§2) замечанию, если рассматривать $\mo$-модуль $\mM$ как бимодуль относительно какого-либо тела $K$, то элементы $K$ порождают гомоморфизмы модулей. Следовательно, тело $K$ необходимо должно гомоморфно отображаться в подтело тела автоморфизмов этого $\mo$-модуля. Этот гомоморфизм, поскольку единичный элемент является единичным оператором, является изоморфизмом, а всё тело автоморфизмов должно иметь конечную степень относительно соответствующего подтела; иначе и модуль представления имел бы бесконечный ранг, что невозможно.

В заключение отметим еще, что рассматриваемые здесь неприводимые представления еще не исчерпывают всех неприводимых представлений, а являются только теми, которые осуществляются простыми $\mo$-модулями. Вообще существуют еще модули представления, например тело $\Omega$, рассматриваемое как бимодуль относительно подтела $\Omega_0$ как левой области и самого себя как правой области; они как $\mo$-модули приводимы, но как бимодули неприводимы, и поэтому всё же приводят к неприводимым представлениям.

\subsection*{§ 19. Простые композиционные факторы в модулях и модулях представления}

Пусть $\mo$ -- кольцо с условием максимальности и условием минимальности для левых идеалов, а $\mc$ -- максимальный нильпотентный идеал. Тогда:

\emph{Каждый простой $\mo$-модуль $\mM$ либо аннулируется $\mc$, либо изоморфен простому левому идеалу $\ml$ из кольца без радикала $\mo/\mc$. В последнем случае модуль аннулируется всеми теми двусторонними идеалами из $\mo/\mc$, которые не содержат $\ml$, а абсолютная область множителей (§1) изоморфна двусторонне простому кольцу, содержащему $\ml$.}

Доказательство. Пусть $\mc^\rho=0$. Должно быть $\mc\mM=0$; ведь иначе было бы $\mc\mM=\mM$, значит $\mc^\rho\mM=\mM$, откуда из $\mc^\rho=0$ следовало бы противоречие $\mM=0$. Поэтому $\mM$ можно также рассматривать как $\mo/\mc$-модуль: все элементы одного класса вычетов по $\mc$ дают одно и то же при умножении на элемент из $\mM$.

Как кольцо без радикала, $\mo/\mc$ обладает единичным элементом. Поэтому $\mM$ является прямой суммой модуля, аннулируемого $\mo/\mc$, и модуля, для которого единичный элемент становится единичным оператором (§17). Так как $\mM$ прост, может появиться только одно из этих двух слагаемых. Если единичный элемент является единичным оператором, то, как в §18, следует изоморфность простому левому идеалу. Действительно, если $m\ne0$ -- элемент из $\mM$, то $(\mo/\mc)m\ne0$; следовательно, $\ml m\ne0$ для по крайней мере одного простого левого идеала $\ml$ из $\mo/\mc$, и поэтому он должен исчерпывать $\mM$. Остальные утверждения для таких левых идеалов ясны и немедленно переносятся на все изоморфные им модули.

Следствие. Если $\mM$ -- $\mo$-модуль, обладающий композиционным рядом, то простые композиционные факторы либо аннулируются $\mc$, либо изоморфны простым левым идеалам из $\mo/\mc$.

Если $P$ -- кольцо, коммутативно связанное с $\mo$, то все результаты сохраняются. Нужно лишь вместо модулей рассматривать бимодули относительно $P$ справа и $\mo$ слева, удовлетворяющие условию
\[
        (am)\rho=a(m\rho)=(a\rho)m \qquad(\S 15).
\]
Кроме того, единичный элемент $P$ также должен быть единичным оператором. В качестве подмодулей всегда рассматриваются только такие, которые допускают умножение на $P$, как в случае идеалов. Построенные в наших доказательствах подмодули $\mc\mM,\mc^2\mM,\ldots$ удовлетворяют этому условию.

В частности, если $P$ является телом, а $\mM$ -- $P$-модулем конечного ранга, то $\mM$ одновременно является модулем представления. Представление, осуществляемое через $\mM$, является не только гомоморфизмом колец, но и операторным гомоморфизмом относительно $P$ (§15). Все представления в $P$ с этим свойством доставляются такими модулями. Неприводимые представления принадлежат простым модулям, и наоборот. Поэтому теоремы этого параграфа сразу можно толковать как теоремы о представлениях $\mo$ в $P$: все неприводимые представления, кроме нулевых, осуществляются простыми левыми идеалами $\mo/\mc$. Поэтому если произвольное представление по §16 редуцировать с помощью композиционного ряда, то на главной диагонали, кроме нулей, появляются только такие представления, которые эквивалентны представлениям, осуществляемым простыми левыми идеалами $\mo/\mc$.

Замечание. При предположениях, сделанных относительно $\mo$ и $P$, для $\mo$ не нужно предполагать никакого условия конечности. Действительно, все представления одновременно являются изоморфными представлениями абсолютной области множителей, а она, как изоморфный прообраз кольца матриц конечной степени, сама является $P$-модулем конечного ранга. Так как по определению допустимыми идеалами считаются только $P$-модули, условие максимальности и условие минимальности для этой абсолютной области множителей выполнены.

Эта абсолютная область множителей становится гиперкомплексной системой относительно $P$, если $P$ предполагается коммутативным телом. Это предположение всегда выполнено, как только на главной диагонали появляются не только нули. Тогда ведь $P$, по замечанию в §18, изоморфно подтелу тела автоморфизмов $\ml$, которое при реализации через подтело $K$ кольца $\mo/\mc$ (§14), вследствие коммутативной связанности, должно лежать в центре $K$.
% END UNIT BODY
% END PRODUCER UNIT 203
\endgroup


\clearpage
\begingroup
\setcounter{footnote}{0}
% BEGIN PRODUCER UNIT 204
% Source: C:/Users/Floris/Documents/interlanguage/03_projects/noether/02_slavic_working_corpus/translations/paper34/source_fidelity_section20/russian/v001/Noether_Paper34_Section20_SourceFidelity_Russian_v001.tex
% Identity: 14260 B / SHA-256 EE1EAA514F022ED1FF10233A1F4FDE2D2FE242CD8C14BE3F8557C5C89384DA4A
\setlength{\parindent}{1.15em}
\setlength{\parskip}{0pt}
\makeatletter
\renewcommand\footnotesize{\@setfontsize\footnotesize{7.4}{8.4}\selectfont}
\makeatother
\providecommand{\mo}{\mathfrak{o}}
\providecommand{\ma}{\mathfrak{a}}
\providecommand{\ml}{\mathfrak{l}}
\providecommand{\mc}{\mathfrak{c}}
\providecommand{\mD}{\mathfrak{D}}
\emergencystretch=140em
\allowdisplaybreaks
\sloppy
\hbadness=10000
\vbadness=10000
\hfuzz=2pt
\vfuzz=10pt
\raggedbottom
% BEGIN UNIT BODY
% Unit: Paper34 Section20 source-fidelity remediation v001. Direct Russian translation from the curated original German scan witness, not from the compressed local TeX.
\section*{34. Гиперкомплексные величины и теория представлений}
\emph{Продолжение и завершение: глава III, §§15--19, и глава IV, §§20--26.}

\section*{Глава IV. Представления групп и гиперкомплексных систем}

\subsection*{§ 20. Встраивание гиперкомплексных систем в теорию}

Рассматриваем гиперкомплексные системы $\mo$ над коммутативным полем $P$. Так как, по соглашению, идеалами в $\mo$ считаются только $P$-модули, для левых и правых идеалов выполняются условие максимальности и условие минимальности (ср. §8). Следовательно, применима вся теория колец, изложенная в главе II.

Под представлениями гиперкомплексной системы $\mo$ далее всегда понимаются представления в поле $P$, причем такие, для которых из $a\sim A$ следует
\[
        a\rho\sim A\rho \qquad(\rho\in P),
\]
то есть модульная гомоморфность относительно $P$. Для модулей представления это означает
\[
        (a\rho)m=a(m\rho)
\]
(ср. конец §15). Поэтому для модулей представления действуют теоремы §19; из них следует, что все неприводимые представления осуществляются простыми левыми идеалами кольца
\[
        \mo_0=\mo/\mc,
\]
и, следовательно, имеется столько неэквивалентных неприводимых представлений, сколько двусторонне простых слагаемых $\ma_\nu$ встречается в разложении
\[
        \mo_0=\ma_1+\cdots+\ma_s.
\]

Чтобы действительно определить представление $\mo$, задаваемое таким левым идеалом $\ml_\nu$, сначала переходят от элементов $\mo$ к соответствующим классам вычетов в $\mo_0$. Так как умножение на элемент $P$ задает изоморфизм для всех идеалов в $\mo_0$, поле $P$ гомоморфно подполю $e^{(\nu)}P$ тела автоморфизмов $K_\nu$ каждого простого левого идеала, а поскольку единичный элемент является единичным оператором, это даже изоморфизм. Представление над $P$, осуществляемое простым левым идеалом $\ml_\nu$, можно получить так: сначала исследовать представление, осуществляемое $\ml_\nu$ над этим подполем $e^{(\nu)}P$, а затем посредством изоморфизма $e^{(\nu)}P\simeq P$ перейти к $P$. Тело автоморфизмов $K_\nu$ имеет конечный ранг относительно $P$; значит, речь идет о представлении в подполе конечной степени тела автоморфизмов $\ml_\nu$.

Итак, чтобы получить искомое представление, сначала представим каждый элемент $c$ из $\mo_0$ через его двусторонние компоненты:
\[
        c=c_1+\cdots+c_s.
\]
Теперь надо найти только представляющую матрицу для $c_\nu$, поскольку остальные $c_i$ аннулируют идеал $\ml_\nu$ и потому представляются нулевой матрицей. По §18 эта матрица находится так: представляют $c_\nu$ через матричные единицы $\alpha_{ik}^{(\nu)}$ из $\ma_\nu$:
\[
        c_\nu=\sum c_{ik}^{(\nu)}\alpha_{ik}^{(\nu)}
\]
и в матрице $(\alpha_{ik}^{(\nu)})$ заменяют каждый элемент $\alpha_{ik}^{(\nu)}$ матрицей $A_{ik}^{(\nu)}$, которая соответствует ему в представлении тела $K_\nu$ над $P$, осуществляемом самим $K_\nu$.

Замечание. Тело $K_\nu$ имеет конечный ранг относительно $P$. Каждый элемент $x$ удовлетворяет уравнению $f(x)=0$ с коэффициентами из $e^{(\nu)}P$, так как между степенями $x$ заведомо есть линейная зависимость. Если, в частности, $P$ алгебраически замкнуто, это уравнение распадается на линейные множители; поскольку $K_\nu$ является телом, $x$ уже является корнем одного из линейных множителей и, следовательно, принадлежит $e^{(\nu)}P$. То есть $K_\nu=e^{(\nu)}P$. В этом случае сами матрицы $(\alpha_{ik}^{(\nu)})$ образуют неприводимое представление.

Из этого замечания вместе с концом §19 получается "теорема Бернсайда":

Пусть $P$ алгебраически замкнуто и коммутативно связано с кольцом $\mo$. Тогда каждое неприводимое представление $n$-й степени над $P$ содержит ровно $n^2$ линейно независимых матриц.\footnote{Тем самым $\mo$ не предполагается гиперкомплексной системой и может, например, быть групповым кольцом, состоящим из всех линейных комбинаций конечного числа элементов бесконечной группы.}

Из этой формулировки сразу следует обычная: "Система матриц $n$-й степени над $P$, которая вместе с каждыми двумя матрицами содержит и их произведение, неприводима тогда и только тогда, когда не существует нетривиального однородного линейного соотношения
\[
        \sum \alpha_{ik}x_{ik}=0
\]
с коэффициентами из $P$, которому удовлетворяют элементы $x_{ik}$ каждой матрицы системы." Действительно, из всякой такой системы матриц, присоединив все линейные комбинации, можно получить кольцо и $P$-модуль и рассматривать это кольцо как его собственное представление.

Доказательство. По §19 абсолютная область множителей $\mD$ изоморфна двусторонне простому и вполне приводимому кольцу $\mo_0$ с единичным элементом; такое кольцо является полным матричным кольцом (например, кольцом всех матриц $m$-й степени) относительно тела автоморфизмов его левых идеалов, которое из-за алгебраической замкнутости совпадает с $P$. Каждое неприводимое представление $\mD$, а значит и данное представление, порождается простым левым идеалом. Левый идеал имеет ранг $m$, представление имеет степень $n$, откуда $m=n$; следовательно, в $\mo_0$, а значит и в $\mD$, имеется ровно $n^2$ линейно независимых элементов.

Так же легко следует "обобщенная теорема Бернсайда":\footnote{G. Frobenius und I. Schur, \emph{Über die Äquivalenz der Gruppen linearer Substitutionen}, Sitzungsber. Berlin 1906.}

Пусть при тех же предположениях $\mD$ является вполне приводимым представлением $\mo$, распадающимся на попарно неэквивалентные составляющие степеней
\[
        n_1,n_2,\ldots,n_s;
\]
тогда $\mD$ имеет ранг
\[
        n_1^2+\cdots+n_s^2
\]
относительно $P$.

Доказательство. Пусть снова $\mo$ является абсолютной областью множителей, которая необходимо есть гиперкомплексная система. Различные неэквивалентные представления порождаются левыми идеалами $\ml_1,\ldots,\ml_s$, принадлежащими различным двусторонне простым компонентам $\ma_i$ кольца $\mo/\mc$, где $\mc$ есть радикал. Поэтому данное представление порождается $\ml_1+\cdots+\ml_s$ и имеет $\ma_1+\cdots+\ma_s$ в качестве абсолютной области множителей. Отсюда с помощью тех же рассуждений о ранге, что и выше, следует утверждение.

Пусть теперь $\mo$ является гиперкомплексной системой без радикала над произвольным полем. По §13, $\mo$ вполне приводима и имеет единичный элемент. Из теорем §17 и §18 следует: каждое представление $\mo$ вполне приводимо.

Обратно, каждое вполне приводимое представление кольца $\mo$, коммутативно связанного с $P$, имеет в качестве абсолютной области множителей гиперкомплексную систему без радикала.

Доказательство. Абсолютная область множителей $\mD$ изоморфна кольцу и $P$-модулю, состоящему из матриц над $P$, то есть является гиперкомплексной системой. Элементы радикала $\mc$ по §19 представляются нулем в каждом неприводимом представлении, а значит и во всем представлении; поэтому $\mc=0$.

(Если в представлении встречаются только эквивалентные составляющие, то есть все они осуществляются одним левым идеалом $\ml$, абсолютная область множителей становится двусторонне простой.)

Регулярным представлением гиперкомплексной системы $\mo$ называют представление, осуществляемое самим единичным идеалом $\mo$. (В случае группового кольца в качестве базиса специально берут групповые элементы $u_1,\ldots,u_h$; ср. §6.) Так как в $\mo$ в качестве композиционных факторов присутствуют все левые идеалы из $\mo/\mc$, все неприводимые представления уже содержатся в регулярном представлении как диагональные матрицы $R_{ii}$, если представить, что регулярное представление по формуле (2) §16 посредством композиционного ряда преобразовано к подходящему базису.

Представление известно, когда известны представляющие матрицы базисных элементов $u_1,\ldots,u_h$. Если речь специально идет о групповом кольце, где $u_1,\ldots,u_h$ являются элементами группы, то всякое гомоморфное матричное представление группы порождает также представление группового кольца; значит, задача представления группы является частным случаем задачи представления гиперкомплексных систем.
% END UNIT BODY
% END PRODUCER UNIT 204
\endgroup


\clearpage
\begingroup
\setcounter{footnote}{0}
% BEGIN PRODUCER UNIT 205
% Source: C:/Users/Floris/Documents/interlanguage/03_projects/noether/02_slavic_working_corpus/translations/paper34/source_fidelity_section21/russian/v001/Noether_Paper34_Section21_SourceFidelity_Russian_v001.tex
% Identity: 10700 B / SHA-256 86F95B80D6270B7E95B26B59F4F0DD891A0A9685597FABA813CAEC077000FED9
\setlength{\parindent}{1.15em}
\setlength{\parskip}{0pt}
\makeatletter
\renewcommand\footnotesize{\@setfontsize\footnotesize{7.4}{8.4}\selectfont}
\makeatother
\providecommand{\mo}{\mathfrak{o}}
\providecommand{\ma}{\mathfrak{a}}
\providecommand{\ml}{\mathfrak{l}}
\providecommand{\mc}{\mathfrak{c}}
\providecommand{\mD}{\mathfrak{D}}
\providecommand{\mZ}{\mathfrak{Z}}
\emergencystretch=140em
\allowdisplaybreaks
\sloppy
\hbadness=10000
\vbadness=10000
\hfuzz=2pt
\vfuzz=10pt
\raggedbottom
% BEGIN UNIT BODY
% Unit: Paper34 Section21 source-fidelity continuation v001. Direct Russian translation from curated German scan witness.
\section*{Э. Нётер. Гиперкомплексные величины и теория представлений}
\emph{Источник-точный перевод §21 с немецкого скан-свидетеля}

\subsection*{§ 21. Расширение основного поля. Представления центра}

Пусть
\[
        \mo=a_1P+\cdots+a_hP
\]
есть гиперкомплексная система, а $\Omega$ — поле расширения основного поля $P$. Тогда можно образовать
\[
        \mo\Omega=a_1\Omega+\cdots+a_h\Omega
\]
со старыми правилами умножения для базисных элементов $a_i$.\textsuperscript{18a)} Система $\mo\Omega$ снова гиперкомплексна относительно $\Omega$.

Каждое представление $\mo$ приводит к представлению $\mo\Omega$, потому что представление всех элементов известно, как только известны представления базисных элементов $a_i$.\textsuperscript{19)}

{\footnotesize\noindent 18a) Точнее: вводятся новые символы $a_i$ со старыми правилами умножения; $a_1\Omega+\cdots+a_h\Omega$ тогда содержит подкольцо, изоморфное $\mo$, так что затем эти $a_i$ можно отождествить со старыми.\par
19) Конечно, здесь опять рассматриваются только такие представления, которые являются $P$-модульными гомоморфизмами: из $a\sim A$ должно следовать $a\rho\sim A\rho$ для $\rho\in P$, и соответственно для $\Omega$.\par}

Идеал или модуль представления, а также соответствующее представление, называются абсолютно неприводимыми, если они остаются неприводимыми после перехода к алгебраически замкнутому полю.

Если $\mo\Omega$ не имеет радикала, то и $\mo$ его не имеет; ведь нильпотентный идеал $\mc$ в $\mo$ дал бы расширенный идеал $\mc\Omega$ в $\mo\Omega$. Обратное неверно, как мы увидим ниже.

Центр $\mo\Omega$ равен расширенному центру $\mZ\Omega$. Если $\mo$ вполне приводима, то $\mZ$ является прямой суммой полей:
\[
        \mZ=\mZ_1+\cdots+\mZ_s.
\]
Если $\mo$ при переходе к $\mo\Omega$ остается вполне приводимой, то есть если не добавляется никакого нильпотентного идеала, то новый центр распадается:
\[
        \mZ\Omega=\mZ_1\Omega+\cdots+\mZ_s\Omega
\]
снова в прямую сумму полей над $\Omega$; когда $\Omega$ алгебраически замкнуто, эти поля, по замечанию в §20, имеют ранг $1$ относительно $\Omega$ и изоморфны $\Omega$.

Для представлений центра $\mZ$ справедливы следующие теоремы.

Всякое неприводимое представление коммутативной гиперкомплексной системы $\mZ$ в алгебраически замкнутом поле $\Omega$ имеет первую степень; иначе говоря, неприводимые представления тождественны гомоморфизмам из $\mZ$ в $\Omega$.

Доказательство. Каждое неприводимое представление $\mZ$ дает такое же представление $\mZ\Omega$, а значит и представление $\mZ\Omega/\mc$, где $\mc$ — радикал $\mZ\Omega$. Кольцо $\mZ\Omega/\mc$ является коммутативной системой без радикала, следовательно, по концу §14, прямой суммой полей. Эти поля имеют первую степень, поскольку $\Omega$ предполагается алгебраически замкнутым, и они порождают все неприводимые представления (§20).

Вместе с тем, так как эквивалентные представления первой степени необходимо становятся равными ($A^{-1}\alpha A=\alpha$), число различных гомоморфизмов из $\mZ$ в $\Omega$ равно рангу $\mZ\Omega/\mc$.

Последнее замечание дает для специального случая, когда $\mZ$ является полем над $P$, такую теорему.

Коммутативное поле $\mZ$ тогда и только тогда вполне приводимо, то есть не имеет радикала, когда $\mZ$ является расширением первого рода поля $P$.

Действительно, для поля гомоморфизмы становятся изоморфизмами. Их число равно рангу $\mZ\Omega/\mc$, а значит равно степени поля, то есть $\mZ$ является расширением первого рода, тогда и только тогда, когда $\mc=0$.

Если $\mZ$ вполне приводим,
\[
        \mZ=\mZ_1+\cdots+\mZ_s,
\]
то в каждом представлении отдельные поля $\mZ_i$ также отображаются гомоморфно; в неприводимом представлении одно из этих полей отображается изоморфно, а остальные — нулем (§19).

Применение этих теорем к гиперкомплексным системам прежде всего дает для систем без радикала следующее.

Если $\mo\Omega$, а следовательно и $\mo$, является системой без радикала, то во всяком абсолютно неприводимом представлении $\mo$ элементы центра $z$ представляются диагональными матрицами $E\xi$, и $z\mapsto\xi$ является представлением первой степени $\mZ$ в $\Omega$. Полученная так связь между классами абсолютно неприводимых представлений $\mo$ и классами представлений $\mZ$ взаимно однозначна. Поэтому число этих классов представлений равно рангу центра.\textsuperscript{20)}

{\footnotesize\noindent 20) Соответствующие теоремы справедливы не только для представлений в алгебраически замкнутом поле $\Omega$, но и для представлений в отдельных телах автоморфизмов; это будет изложено в другом месте.\par}

Доказательство. Разложению на двусторонне неразложимые идеалы
\[
        \mo\Omega=\ma_1+\cdots+\ma_s
\]
взаимно однозначно соответствует разложение центра на поля:
\[
        \mZ\Omega=\mZ_1+\cdots+\mZ_s,
\]
где $\mZ_i$ есть центр $\ma_i$. При неприводимом представлении $\mo$ один $\ma_i$ отображается изоморфно, остальные — нулем, и класс представления однозначно определяется этим $\ma_i$. Идеал $\ma_i$ представляется полным матричным кольцом, а центр такого полного матричного кольца состоит из диагональных матриц $E\xi$. Соответствие $z\mapsto\xi$ является гомоморфизмом $\mZ$: одно $\mZ_i$ отображается изоморфно, остальные — нулем. Тем самым все доказано.

Для вопроса, когда системе $\mo$ без радикала соответствует такая же система $\mo\Omega$, получаем еще следующее.

Если $\mZ$ имеет компоненту $\mZ_i$, являющуюся расширением второго рода поля $P$, то $\mo\Omega$ заведомо имеет радикал.\textsuperscript{21)} Ведь $\mZ_i\Omega$ в этом случае уже имеет радикал.

{\footnotesize\noindent 21) Если $\mZ$ имеет только компоненты первого рода, то $\mo\Omega$ не имеет радикала; это также будет доказано позднее. Для характеристики нуль И. Шур уже доказал эквивалентную теорему: неприводимые представления над $P$ остаются вполне приводимыми при каждом расширении области коэффициентов (\emph{Beiträge zur Theorie der Gruppen linearer Substitutionen}, Transact. Am. Math. Soc. 15 (1909), S. 159).\par}
% END UNIT BODY
% END PRODUCER UNIT 205
\endgroup


\clearpage
\begingroup
\setcounter{footnote}{0}
% BEGIN PRODUCER UNIT 206
% Source: C:/Users/Floris/Documents/interlanguage/03_projects/noether/02_slavic_working_corpus/translations/paper34/source_fidelity_section22/russian/v001/Noether_Paper34_Section22_SourceFidelity_Russian_v001.tex
% Identity: 3279 B / SHA-256 AFA93A3D5F660E823959EC26BA844319E0C9AA639111273AA4EBAF7C432568CE
\setlength{\parindent}{1.15em}
\setlength{\parskip}{0pt}
\providecommand{\mZ}{\mathfrak{Z}}
\emergencystretch=120em
\allowdisplaybreaks
\sloppy
\hbadness=10000
\vbadness=10000
\hfuzz=2pt
\vfuzz=10pt
\raggedbottom
% BEGIN UNIT BODY
% Unit: Paper34 Section22 source-fidelity continuation v001.
% Language lane: Russian.
% Direct source: sources/paper34/source_fidelity/Noether_Paper34_Section22_ORIGINAL_SCAN_WITNESS_v001.tex.
\section*{Эмми Нётер: гиперкомплексные величины и теория представлений}
\emph{Работа 34, §22. Источниково-точный перевод со сканированного немецкого свидетеля.}

\subsection*{§ 22. Применение к абелевым группам}

Пусть
\[
        G=G_1\times G_2\times\cdots\times G_r
\]
есть конечная абелева группа, разложенная в прямое произведение циклических групп $G_i$ порядков $h_i$. Ее элементы имеют вид
\[
        a=a_1^{\alpha_1}\cdots a_r^{\alpha_r}.
\]
Пусть $P$ --- поле, характеристика которого не делит порядок группы
\[
        h=h_1h_2\cdots h_r.
\]

Групповое кольцо $\mZ$ состоит из всех сумм
\[
        \sum_{\alpha_1,\ldots,\alpha_r}
        A_{\alpha_1\ldots\alpha_r}a_1^{\alpha_1}\cdots a_r^{\alpha_r},
        \qquad A_{\alpha_1\ldots\alpha_r}\in P,
\]
и является гомоморфным образом кольца многочленов $P[x_1,\ldots,x_r]$ при отображении $x_i\mapsto a_i$. Поэтому
\[
        \mZ\simeq P[x_1,\ldots,x_r]/\mathfrak m,
        \qquad
        \mathfrak m=(x_1^{h_1}-e,\ldots,x_r^{h_r}-e).
\]
В поле расширения $\Omega$ многочлены $x_i^{h_i}-e$ распадаются на попарно различные множители $x_i-\xi_i$; поэтому $\mathfrak m$ становится произведением (или пересечением) попарно различных простых идеалов
\[
        (x_1-\xi_1^{(\alpha_1)},\ldots,x_r-\xi_r^{(\alpha_r)}),
\]
каждому из которых соответствует единственный нуль
\[
        (\xi_1^{(\alpha_1)},\ldots,\xi_r^{(\alpha_r)}).
\]
Образованию пересечения соответствует прямое суммарное разложение (§ 4):
\[
        \mZ\Omega=\mZ_1+\cdots+\mZ_h,
\]
где каждый раз $\mZ_i\sim\Omega$; следовательно, получаем представление
\[
        a_i\mapsto \xi_i^{(\alpha_i)},
        \qquad\hbox{или более общо}\qquad
        a\mapsto \chi^{(\alpha)}(a).
\]
Эти представления являются характерами. Их число равно
\[
        h_1h_2\cdots h_r=h,
\]
как и должно быть по общей теории.
% END UNIT BODY
% END PRODUCER UNIT 206
\endgroup


\clearpage
\begingroup
\setcounter{footnote}{0}
% BEGIN PRODUCER UNIT 207
% Source: C:/Users/Floris/Documents/interlanguage/03_projects/noether/02_slavic_working_corpus/translations/paper34/source_fidelity_section23/russian/v001/Noether_Paper34_Section23_SourceFidelity_Russian_v001.tex
% Identity: 7388 B / SHA-256 4ECB2A43CB85C0482DC5572BB010589067EDBA6973741ED3EDA6364FAEA0E6EA
\setlength{\parindent}{1.15em}
\setlength{\parskip}{0pt}
\providecommand{\mo}{\mathfrak{o}}
\providecommand{\mc}{\mathfrak{c}}
\emergencystretch=140em
\allowdisplaybreaks
\sloppy
\hbadness=10000
\vbadness=10000
\hfuzz=2pt
\vfuzz=10pt
\raggedbottom
% BEGIN UNIT BODY
% Unit: Paper34 Section23 source-fidelity continuation v001.
% Language lane: Russian.
% Direct source: sources/paper34/source_fidelity/Noether_Paper34_Section23_ORIGINAL_SCAN_WITNESS_v001.tex.
\section*{Эмми Нётер: гиперкомплексные величины и теория представлений}
\emph{Работа 34, §23. Источниково-точный перевод со сканированного немецкого свидетеля.}

\subsection*{§ 23. Определитель гиперкомплексной системы}

Пусть
\[
        \mo=a_1P+\cdots+a_hP
\]
есть гиперкомплексная система. Присоединим к $P$ $h$ неопределенных $x_1,\ldots,x_h$ и образуем
\[
        \mo^*=a_1P(x)+\cdots+a_hP(x).
\]
Неопределенные $x_i$ должны перестановочно умножаться с $a_i$; этим уже задаются правила вычислений в $\mo^*$.

В $\mo^*$ лежит «общий элемент» системы $\mo$:
\[
        w=a_1x_1+\cdots+a_hx_h.
\]
Если в некотором представлении $a_i\mapsto A_i$, то элементу $w$ соответствует
\[
        W=A_1x_1+\cdots+A_hx_h.
\]
$W$ называется системной матрицей, принадлежащей этому представлению; в частности, если $a_i$ образуют группу, а $\mo$ поэтому является групповым кольцом, это групповая матрица. В случае регулярного представления имеем регулярную системную матрицу.

Элементы $w_{ij}$ матрицы $W$ являются линейными формами от $x_i$. Поэтому системный определитель $|W|$ имеет степень $n$, если речь идет о представлениях $n$-й степени. В частности регулярный системный определитель имеет степень $h$.

Системный определитель не меняется при переходе к эквивалентным представлениям, поскольку
\[
        |PWP^{-1}|=|P|\,|W|\,|P^{-1}|=|W|.
\]
При переходе от базиса $(a_1,\ldots,a_h)$ к новому базису $(b_1,\ldots,b_h)$ и от $w=\sum a_ix_i$ к $w=\sum b_jy_j$ новые элементы $W$, а значит и новый определитель, получают из старых подстановкой
\[
        x_i=\sum_j\rho_{ij}y_j
\]
с регулярной матрицей подстановки.

Если задан композиционный ряд модуля представления, то при подходящем выборе базиса матрица $W$ имеет вид
\[
        W=\begin{pmatrix}
        W_1&0&\cdots&0\\
        *&W_2&\cdots&0\\
        \vdots&\vdots&\ddots&\vdots\\
        *&*&\cdots&W_s
        \end{pmatrix}.
\]
Среди определителей $|W_i|$ произвольного представления не встречается ничего иного, чем в регулярном представлении системы $\mo$, или даже системы $\mo/\mc$, где $\mc$ --- радикал.

Если $P$ алгебраически замкнуто, то определитель $|W_i|$, принадлежащий неприводимому представлению, является неприводимой функцией от $x_i$, а неэквивалентным представлениям соответствуют различные неприводимые множители.

Доказательство. Так как все неприводимые представления $\mo$ являются также представлениями $\mo/\mc$, можно ограничиться кольцом без радикала $\mo/\mc$. В нем возьмем за базис матричные единицы $c_{ik}^{(\nu)}$; тогда общий элемент имеет вид
\[
        w=\sum_{\nu,i,k}c_{ik}^{(\nu)}x_{ik}^{(\nu)}.
\]
Матрицы неприводимых представлений имеют вид
\[
        W_\nu=(x_{ik}^{(\nu)}).
\]
Функции $|W_\nu|=|x_{ik}^{(\nu)}|$ заведомо неприводимы и очевидно различны.

Чтобы вычислить $|W_i|$, можно разложить регулярную системную матрицу $\mo/\mc$ на ее неприводимые множители. Каждый неприводимый множитель появляется столько раз, какова степень соответствующего неприводимого представления, потому что соответствующий идеал столько раз встречается в композиционном ряду.

Можно также исходить из регулярной системной матрицы $\mo$, но тогда каждый неприводимый множитель получается чаще, а именно столько раз, сколько соответствующий простой идеал встречается как композиционный множитель среди левых идеалов. В этом регулярном представлении $\mo$ рассматривалась как левый идеал; если рассматривать $\mo$ как правый идеал, появляется вторая регулярная системная матрица, «антистрофная матрица» у Фробениуса. Она содержит те же неприводимые множители, а именно системные определители всех неприводимых представлений, но, возможно, с другими показателями; см. пример в § 10.

Системный определитель коммутативной системы распадается на линейные множители, потому что все неприводимые представления имеют первую степень. Эти линейные множители сами являются неприводимыми представлениями, а значит для абелевых групп дают характеры. Этот факт был исходным пунктом Дедекинда в исследовании группового определителя неабелевых групп.
% END UNIT BODY
% END PRODUCER UNIT 207
\endgroup


\clearpage
\begingroup
\setcounter{footnote}{0}
% BEGIN PRODUCER UNIT 208
% Source: C:/Users/Floris/Documents/interlanguage/03_projects/noether/02_slavic_working_corpus/translations/paper34/source_fidelity_section24/russian/v001/Noether_Paper34_Section24_SourceFidelity_Russian_v001.tex
% Identity: 7363 B / SHA-256 6BDCB77143B7C4C2AB26FB49244D8AB6B2F092BA41E5392467795DB01DEF4270
\setlength{\parindent}{1.15em}
\setlength{\parskip}{0pt}
\providecommand{\mo}{\mathfrak{o}}
\providecommand{\mc}{\mathfrak{c}}
\providecommand{\mM}{\mathfrak{M}}
\emergencystretch=140em
\allowdisplaybreaks
\sloppy
\hbadness=10000
\vbadness=10000
\hfuzz=2pt
\vfuzz=10pt
\raggedbottom
% BEGIN UNIT BODY
% Unit: Paper34 Section24 source-fidelity continuation v001.
% Language lane: Russian.
% Direct source: sources/paper34/source_fidelity/Noether_Paper34_Section24_ORIGINAL_SCAN_WITNESS_v001.tex.
\section*{Эмми Нётер: гиперкомплексные величины и теория представлений}
\emph{Работа 34, §24. Источниково-точный перевод со сканированного немецкого свидетеля.}

\subsection*{§ 24. Следы и характеры}

Если в представлении $\mo\sim\mathfrak D$ гиперкомплексной системы $\mo$ элементу $a$ сопоставлена матрица $A$, то полагают
\[
        \operatorname{Sp}_{D}(a)=\operatorname{Sp} A.
\]
Следы являются линейными функциями:
\[
        \operatorname{Sp}(c+d)=\operatorname{Sp}(c)+\operatorname{Sp}(d),
        \qquad
        \operatorname{Sp}(c\alpha)=\alpha\operatorname{Sp}(c).
\]
Эквивалентные представления имеют одни и те же следы.

След в приводимом представлении равен сумме следов в представлениях, задаваемых композиционными факторами.

Доказательство. При подходящем выборе базиса матрица имеет блочный вид
\[
        \begin{pmatrix}
        A_{11}&0&\cdots&0\\
        *&A_{22}&\cdots&0\\
        \vdots&\vdots&\ddots&\vdots\\
        *&*&\cdots&A_{ss}
        \end{pmatrix},
\]
так что ее след равен
\[
        \operatorname{Sp}(A_{11})+\operatorname{Sp}(A_{22})+\cdots+\operatorname{Sp}(A_{ss}).
\]

Главным следом называется след в регулярном представлении. Приведенным следом называется сумма следов в различных неприводимых представлениях.

Если $c$ является элементом максимального нильпотентного идеала $\mc$, то $\operatorname{Sp}(c)=0$ для любого представления.

Доказательство. Достаточно рассматривать представления через простые левые идеалы $\mo/\mc$, потому что в каждом другом представлении встречаются только эти композиционные факторы, а след составляется аддитивно. Но $c$ аннулирует все элементы $\mo/\mc$; значит, каждому $c$ из $\mc$ сопоставляется нулевая матрица, поэтому $\operatorname{Sp}(c)=0$.

Если $\mo$ двусторонне проста, а $P$ алгебраически замкнуто, то есть $\mo$ является матричным кольцом
\[
        \mo=\sum c_{ik}P,
\]
то в неприводимом представлении $\mo$ элементу
\[
        a=\sum c_{ik}\alpha_{ik}
\]
сопоставляется матрица $(\alpha_{ik})$. Следовательно,
\[
        \operatorname{Sp}_{\rm red}(a)=\sum_i\alpha_{ii},
        \qquad
        \operatorname{Sp}_{\rm pr}(a)=n\,\operatorname{Sp}_{\rm red}(a)=n\sum_i\alpha_{ii}.
\]
В частности,
\[
        \operatorname{Sp}_{\rm red}(c_{ik})=0\quad(i\ne k),\qquad
        \operatorname{Sp}_{\rm red}(c_{ii})=e,
\]
а также
\[
        \operatorname{Sp}_{\rm pr}(c_{ik})=0\quad(i\ne k),\qquad
        \operatorname{Sp}_{\rm pr}(c_{ii})=n e.
\]

При переходе от $P$ к алгебраически замкнутому полю $\Omega$ главный след не меняется, поскольку он вычисляется по тому же базису. Для приведенного следа это уже неверно.

Следы элементов $a$ системы без радикала в абсолютно неприводимых представлениях называются характерами и обозначаются $\chi(a)$ или, если нужно указать конкретное представление, $\chi^{(\nu)}(a)$.

В неприводимом представлении степени $n_\nu$ элементы центра, согласно §21, представляются диагональными матрицами $E\theta^{(\nu)}$, где $z\mapsto\theta^{(\nu)}(z)$ есть представление первой степени центра, или гомоморфизм центра в $\Omega$. Поэтому след имеет значение $n_\nu\theta^{(\nu)}(z)$. Следовательно, гомоморфизмы $\theta$ центра связаны с характерами $\chi$ соотношением
\[
        \chi^{(\nu)}(z)=n_\nu\theta^{(\nu)}(z).                \tag{1}
\]

В коммутативном случае $n_\nu=1$, и сами характеры дают гомоморфизмы (ср. §22).

Если поле $\Omega$ имеет характеристику нуль, что мы всюду будем предполагать в дальнейшем, то (1) можно разделить на $n_\nu$:
\[
        \theta^{(\nu)}(z)=\frac{\chi^{(\nu)}(z)}{n_\nu}.
\]
Свойство гомоморфизма для $\theta^{(\nu)}$ выражается формулами
\[
        \frac{\chi^{(\nu)}(z+z')}{n_\nu}
        =
        \frac{\chi^{(\nu)}(z)}{n_\nu}
        +
        \frac{\chi^{(\nu)}(z')}{n_\nu},
\]
\[
        \frac{\chi^{(\nu)}(zz')}{n_\nu}
        =
        \frac{\chi^{(\nu)}(z)}{n_\nu}
        \frac{\chi^{(\nu)}(z')}{n_\nu}.
\]

Класс представлений уже однозначно определяется одними только следами матриц. Чтобы знать следы всех матриц, конечно, достаточно знать следы базисных элементов в данном представлении.

Доказательство. Модуль представления $R$ известен, если известно, сколько раз каждый неприводимый модуль представления $\mM_\nu$ встречается в нем как прямое слагаемое. Если это число равно $\mu_\nu$, то след $e^{(\nu)}$ в представлении равен $\mu_\nu n_\nu$, где $n_\nu$ есть степень неприводимого представления. Следовательно,
\[
        \mu_\nu=\frac{\operatorname{Sp}(e^{(\nu)})}{n_\nu}.
\]
% END UNIT BODY
% END PRODUCER UNIT 208
\endgroup


\clearpage
\begingroup
\setcounter{footnote}{0}
% BEGIN PRODUCER UNIT 209
% Source: C:/Users/Floris/Documents/interlanguage/03_projects/noether/02_slavic_working_corpus/translations/paper34/source_fidelity_section25/russian/v001/Noether_Paper34_Section25_SourceFidelity_Russian_v001.tex
% Identity: 8620 B / SHA-256 97FD918CDEC41D49BFD350FDBDF35E3C0396B7B213329DB0DA8F80431CE07C21
\setlength{\parindent}{1.15em}
\setlength{\parskip}{0pt}
\providecommand{\mo}{\mathfrak{o}}
\providecommand{\mc}{\mathfrak{c}}
\providecommand{\ma}{\mathfrak{a}}
\emergencystretch=160em
\allowdisplaybreaks
\sloppy
\hbadness=10000
\vbadness=10000
\hfuzz=2pt
\vfuzz=10pt
\raggedbottom
% BEGIN UNIT BODY
% Unit: paper34_section25_source_fidelity_v001.
% Language lane: Russian.
% Direct source: sources/paper34/source_fidelity/Noether_Paper34_Section25_ORIGINAL_SCAN_WITNESS_R124plus_v001.tex.
\section*{Эмми Нётер: гиперкомплексные величины и теория представлений}
\emph{Работа 34, §25. Источниково-точный перевод со сканированного немецкого свидетеля с R124plus-коррекцией P34.}

\subsection*{§ 25. Дискриминанты}

Пусть
\[
        \mo=a_1P+\cdots+a_hP
\]
есть гиперкомплексная система. Дискриминантной матрицей называется матрица, элементы которой являются главными следами $\operatorname{Sp}(a_i a_j)$. Приведенная дискриминантная матрица — это то же самое для приведенных следов, построенное в алгебраически замкнутом поле. Определитель этой матрицы называется дискриминантом, соответственно приведенным дискриминантом.

Дискриминант не меняется при расширении основного поля. При переходе к другому базису он умножается на квадрат определителя преобразования.

Доказательство. Пусть
\[
        (b_1,\ldots,b_h)=(a_1,\ldots,a_h)P.
\]
Тогда
\[
        (\operatorname{Sp}(a_i a_j))=(\operatorname{Sp}(a_i b_j))P.             \tag{1}
\]
Так же, если $P^t$ — транспонированная матрица,
\[
        (\operatorname{Sp}(a_i b_j))=P^t(\operatorname{Sp}(b_i b_j)).           \tag{1a}
\]
Из (1) и (1a) следует
\[
        (\operatorname{Sp}(a_i a_j))=P^t(\operatorname{Sp}(b_i b_j))P.
\]
Переход к определителям дает утверждение.

Следовательно, определитель задан лишь с точностью до квадрата из $P$; но его обращение или необращение в нуль является инвариантным свойством. Из (1) также следует, что дискриминант, с точностью до ненулевого множителя, можно вычислять по двум различным базисам, а именно как $|\operatorname{Sp}(a_i b_j)|$.

Дискриминант обращается в нуль, если $\mo$ имеет нильпотентный идеал $\mc$.

Доказательство. В качестве базисных элементов для $\mo$ выберем
\[
        c_1,\ldots,c_r,d_1,\ldots,d_s,
\]
где $c_1,\ldots,c_r$ образуют базис $\mc$. Тогда дискриминантная матрица имеет вид
\[
        \begin{pmatrix}
        \operatorname{Sp}(c_i c_j)&\operatorname{Sp}(c_i d_j)\\
        \operatorname{Sp}(d_i c_j)&\operatorname{Sp}(d_i d_j)
        \end{pmatrix}
        =
        \begin{pmatrix}
        0&0\\
        0&\operatorname{Sp}(d_i d_j)
        \end{pmatrix}.
\]
(То же верно и для приведенного дискриминанта.) Поэтому определитель равен нулю.

Следствие. Дискриминант обращается в нуль и тогда, когда нильпотентный идеал появляется после перехода к алгебраически замкнутому полю.

Пусть
\[
        \mo=\ma_1+\cdots+\ma_s
\]
является прямой двусторонней суммой, и пусть $M,M_1,\ldots,M_s$ — дискриминантные матрицы колец $\mo,\ma_1,\ldots,\ma_s$. Тогда
\[
        M=\begin{pmatrix}
        M_1&0&\cdots&0\\
        0&M_2&\cdots&0\\
        \vdots&\vdots&\ddots&\vdots\\
        0&0&\cdots&M_s
        \end{pmatrix},
\]
откуда $|M|=|M_1|\cdots|M_s|$.

Доказательство. Выберем базис для $\mo$, составленный из базисов для $\ma_1,\ldots,\ma_s$. Если $a\in\ma_i$, то
\[
        \operatorname{Sp}_{\mo}(a)=\operatorname{Sp}_{\ma_i}(a),
\]
где оба раза имеются в виду главные следы, но соответственно в кольцах $\mo$ и $\ma_i$. Далее, если множители лежат в различных двусторонних компонентах, то $a_i a a_j=0$; значит, исчезают и смешанные следы. Отсюда следует утверждение; тем же образом оно следует и для приведенного дискриминанта.

Чтобы построить дискриминант матричного кольца $\sum c_{ik}P$, выбирают два различных базиса: элементы $c_{ik}$ один раз упорядочивают по первым индексам, другой раз по вторым индексам. Таблица умножения имеет вид
\[
\begin{array}{c|ccc}
& \overbrace{c_{11}\;\cdots\;c_{1n}}^{r_1}
& \overbrace{c_{21}\;\cdots\;c_{2n}}^{r_2}
& \cdots \\ \hline
\left.\begin{array}{c}c_{11}\\[-1mm]\vdots\\[-1mm]c_{n1}\end{array}\right\}l_1
& (c_{ik}) & 0 & 0 \\[1.3em]
\left.\begin{array}{c}c_{12}\\[-1mm]\vdots\\[-1mm]c_{n2}\end{array}\right\}l_2
& 0 & (c_{ik}) & 0 \\[1.3em]
\vdots & 0 & 0 & (c_{ik})
\end{array}
\]
Дискриминант равен
\[
        D=\left|\begin{array}{cccc}
        \operatorname{Sp}(c_{11})&0&\cdots&0\\
        0&\operatorname{Sp}(c_{22})&\cdots&0\\
        \vdots&\vdots&\ddots&\vdots\\
        0&0&\cdots&\operatorname{Sp}(c_{nn})
        \end{array}\right|^{n}
        =\begin{cases}
        e, & \hbox{для приведенных следов},\\
        n^{n^2}\cdot e, & \hbox{для главных следов}.
        \end{cases}
\]

Следствие. Если над алгебраически замкнутым полем расширения $\mo$ распадается на матричные кольца степеней $n_1,\ldots,n_s$, то дискриминант равен
\[
        D=n_1^{n_1^2}n_2^{n_2^2}\cdots n_s^{n_s^2}\cdot e
\]
и приведенный дискриминант равен
\[
        D_{\rm red}=e.
\]
Следовательно, приведенный дискриминант для систем без радикала всегда ненулевой; другой дискриминант ненулевой лишь тогда, когда ни одно $n_i$ не делится на характеристику поля. Поэтому:

\emph{Ненулевость приведенного дискриминанта необходима и достаточна для систем без радикала после алгебраического замыкания поля $P$.}

\emph{Ненулевость дискриминанта в случае характеристики нуль необходима и достаточна для невозникновения радикала после алгебраического замыкания поля $P$.}\textsuperscript{22)}

{\footnotesize\noindent 22) Для коммутативных систем ср. E. Noether, \emph{Diskriminantensatz für Ordnungen ...}, J. f. M. 157 (1927), с. 82--104, §§4--6. Метод доказательства там тот же, но в деталях сложнее; переход от одного базиса к другому (§4, 4 там) следует заменить данным здесь в начале этого параграфа, потому что возникающий там определитель обращается в нуль.\par}
% END UNIT BODY
% END PRODUCER UNIT 209
\endgroup


\clearpage
\begingroup
\setcounter{footnote}{0}
% BEGIN PRODUCER UNIT 210
% Source: C:/Users/Floris/Documents/interlanguage/03_projects/noether/02_slavic_working_corpus/translations/paper34/source_fidelity_section26/russian/v001/Noether_Paper34_Section26_SourceFidelity_Russian_v001.tex
% Identity: 4443 B / SHA-256 9935036233D00F9029530C93406062A43263021EA94E844A84A676BAFFC0BF09
\setlength{\parindent}{1.15em}
\setlength{\parskip}{0pt}
\providecommand{\mo}{\mathfrak{o}}
\providecommand{\mc}{\mathfrak{c}}
\providecommand{\ma}{\mathfrak{a}}
\emergencystretch=140em
\allowdisplaybreaks
\sloppy
\hbadness=10000
\vbadness=10000
\hfuzz=2pt
\vfuzz=10pt
\raggedbottom
% BEGIN UNIT BODY
% Unit: paper34_section26_source_fidelity_v001.
% Language lane: Russian.
% Direct source: sources/paper34/source_fidelity/Noether_Paper34_Section26_ORIGINAL_SCAN_WITNESS_R124plus_v001.tex.
\section*{Эмми Нётер: гиперкомплексные величины и теория представлений}
\emph{Работа 34, §26. Источниково-точный перевод со сканированного немецкого свидетеля с R124plus-контролем.}

\subsection*{§ 26. Включение группового кольца}

Пусть $a_1,\ldots,a_h$ — элементы конечной группы, и образуем групповое кольцо над полем, характеристика которого не делит $h$. Тогда дискриминант $D\ne0$, а значит групповое кольцо — это кольцо без радикала.

Доказательство. Сначала покажем, что далее $\operatorname{Sp}$ означает главный след:
\[
        \operatorname{Sp}(e)=he,
        \qquad
        \operatorname{Sp}(a_i)=0\quad(a_i\ne e).
\]
В регулярном представлении имеем
\[
        e\mapsto
        \begin{pmatrix}
        e&0&\cdots&0\\
        0&e&\cdots&0\\
        \vdots&\vdots&\ddots&\vdots\\
        0&0&\cdots&e
        \end{pmatrix},
        \qquad \operatorname{Sp}(e)=he.
\]
Для $a_i\ne e$ произведения $a_i a_1,\ldots,a_i a_h$ дают перестановку базисных элементов без неподвижной точки; соответствующая матрица имеет на главной диагонали только нули, поэтому $\operatorname{Sp}(a_i)=0$.

Теперь возьмем базисы $a_1,\ldots,a_h$ и $a_1^{-1},\ldots,a_h^{-1}$. Матрица $(\operatorname{Sp}(a_i a_j^{-1}))$ диагональна с $he$ на главной диагонали, следовательно
\[
        D=h^h e\ne0.
\]

Классом элемента группы называется сумма, образованная в групповом кольце,
\[
        K_a=\sum_s s^{-1}as,                         \tag{1}
\]
где суммирование ведется только по различным элементам $s^{-1}as$. Классы $K_a$ являются центральными элементами, так как они коммутируют со всеми элементами группы. Они порождают центр: если элемент $\sum_i\alpha_i a_i$ группового кольца коммутирует со всеми $s$, то
\[
        \sum_i\alpha_i a_i=s^{-1}\biggl(\sum_i\alpha_i a_i\biggr)s
        =\sum_i\alpha_i(s^{-1}a_i s),
\]
и значит коэффициенты постоянны на каждом классе сопряженности.

Следовательно, ранг центра равен числу классов сопряженных элементов группы, а потому и число абсолютно неприводимых представлений равно этому числу классов. Матрицы $A$ и $S^{-1}AS$ имеют один и тот же след; значит элементы группы $a$ и $s^{-1}as$ имеют один и тот же след. Беря след в обеих частях (1), получаем
\[
        \chi(K_a)=h_a\chi(a),
\]
где $h_a$ — число элементов класса $K_a$. Иными словами: характер элемента группы равен характеру класса, деленному на число элементов этого класса.

\begin{center}
(Поступила 12 августа 1928 г.)
\end{center}
% END UNIT BODY
% END PRODUCER UNIT 210
\endgroup


\clearpage
\begingroup
\setcounter{footnote}{0}
% BEGIN PRODUCER UNIT 211
% Source: C:/Users/Floris/Documents/interlanguage/03_projects/noether/02_slavic_working_corpus/translations/paper35/source_fidelity/russian/v001/Noether_Paper35_SourceFidelity_Russian_v001.tex
% Identity: 46020 B / SHA-256 07A769966B288FD755DCF8E320193DA9F74552C75AF1520EBA5888FD5A03BC1E
\setlength{\parindent}{1.15em}
\setlength{\parskip}{0pt}
\emergencystretch=180em
\allowdisplaybreaks
\sloppy
\hbadness=10000
\vbadness=10000
\hfuzz=3pt
\vfuzz=10pt
\raggedbottom
% BEGIN UNIT BODY
% Unit: paper35_source_fidelity_v001.
% Language lane: Russian.
% Direct source: sources/paper35/source_fidelity/Noether_Paper35_ORIGINAL_MathNet600_R124plusP40_repaired_witness_v001.tex.
\section*{35. О максимальных областях целочисленных функций}
\emph{Rec. Soc. Math. Moscou 36 (1929), с. 65--72. Источнико-верный перевод по исправленному немецкому свидетелю R124+P40/MathNet600.}

Следующие исследования близко соприкасаются с гильбертовой проблемой относительно целых функций\footnote{D. Hilbert, Mathematische Probleme (Gött. Nachr. 1900, с. 253--297), проблема 14.}; они восходят к случайному вопросу Э. Гекке, который смог прямо вычислительно разобрать специальный случай, нужный ему как лемма\footnote{E. Hecke, Über die Konstruktion relativ-Abelscher Zahlkörper durch Modulfunktionen von zwei Variabeln (Math. Ann. 74 (1913), с. 465--510), §2. У Гекке конечная область целостности \(\mathfrak{I}\) порождается тремя частными степенных рядов от двух переменных, между которыми имеется алгебраическое соотношение.}.

Постановка вопроса такова. Под \textit{максимальной областью} целочисленных функций от \(x_1,\ldots,x_n\), то есть многочленов, степенных рядов или частных таких, внутри данной области \(\mathfrak{Z}\), я понимаю такую область, которую уже нельзя расширить присоединением целочисленных знаменателей: если \(a\cdot g(x)\) принадлежит ей, где \(g(x)\in\mathfrak{Z}\), то всегда принадлежит и \(g(x)\); здесь \(a\) означает ненулевое целое рациональное число. Каждая область целостности \(\mathfrak{I}\) в \(\mathfrak{Z}\) однозначно порождает максимальную область \(M(\mathfrak{I})\), состоящую из всех функций \(g(x)\) из \(\mathfrak{Z}\), для которых существует хотя бы одно \(a\ne0\) такое, что \(a\cdot g(x)\in\mathfrak{I}\).

Вопрос состоит в том, что можно сказать о знаменателях \(a\) этой максимальной области \(M(\mathfrak{I})\), когда исходная область целостности \(\mathfrak{I}\) была конечной, то есть состояла из всех целочисленных многочленов от конечного числа функций \(f_1(x),\ldots,f_r(x)\) из \(\mathfrak{Z}\). При этом в качестве \(\mathfrak{Z}\) следует брать все целочисленные многочлены или степенные ряды от \(x_1,\ldots,x_n\), соответственно такие частные, которые не содержат никаких знаменателей, кроме уже встречающихся в \(f(x)\), и их степеней.

Я показываю, что при дополнительном предположении для степенных рядов или их частных, а именно что между \(f(x)\) имеются только алгебраические зависимости, числа \(a\) всегда можно выбирать так, что в них входит лишь конечное число различных простых чисел, хотя вообще говоря в произвольных степенях.

Последнее непосредственно ясно: если \(a\cdot g(x)\in\mathfrak{I}\), но \(a^\varrho g(x)\) не принадлежит \(\mathfrak{I}\) ни для какого показателя \(\varrho\), то возникают все степени \(a^\varrho\); например, при \(\mathfrak{I}=[2x]\) имеем \(x=(2x)/2\), \(x^2=(2x)^2/2^2\).

Вопрос об ограниченности можно сформулировать только так: можно ли указать, вообще говоря бесконечную, систему образующих для \(M(\mathfrak{I})\), такую что в соответствующих знаменателях показатели простых чисел остаются ограниченными? Так как речь идет лишь о конечном числе простых, этот вопрос по известным рассуждениям\footnote{Ср. Hilbert, там же, или более подробное изложение у E. Noether, Die Endlichkeit des Systems der ganzzahligen Invarianten binärer Formen (Gött. Nachr. 1919, с. 138--156), §5.} тождествен вопросу о конечности \(M(\mathfrak{I})\), то есть гильбертовой проблеме относительно целых функций в целочисленной форме. Этого вопроса применяемыми здесь методами я решить не могу\footnote{Так, например, в случае целочисленных (проективных) инвариантов основной системы форм показано, что они выражаются через полную систему инвариантов в обычном смысле, причем в знаменателях встречается лишь конечное число различных простых чисел. Этот результат не следует из обычного метода \(\Omega\)-процесса. Но о самой целочисленной конечности здесь ничего не сказано; до сих пор она известна только для бинарных основных форм. Ср. заметку, цитированную в примечании 3.}.

Доказательство того, что в знаменателях \(a\) встречается лишь конечное число различных простых \(p\), основано на том, что каждому такому \(p\) соответствует по крайней мере одно соотношение по модулю \(p\) между функциями \(f(x)\), порождающими \(\mathfrak{I}\), причем это соотношение с целочисленными коэффициентами не выполняется тождественно по \(x\). Иначе говоря: если \(\mathfrak{o}\) обозначает область целочисленных многочленов от неопределенных \(u_1,\ldots,u_r\), то \(\mathfrak{I}\) посредством соответствия \(u_i\mapsto f_i(x)\) является гомоморфным образом \(\mathfrak{o}\); этот гомоморфизм опосредуется простым идеалом \(\mathfrak{a}\) из \(\mathfrak{o}\), то есть \(\mathfrak{I}\) изоморфна кольцу классов \(\mathfrak{o}/\mathfrak{a}\). Соответственно \(\mathfrak{I}_p\) является гомоморфным образом \(\mathfrak{o}_p\), где индекс обозначает переход к классам вычетов по модулю \(p\), возможный за исключением конечного числа простых, входящих в знаменатели \(\mathfrak{I}\). Этот гомоморфизм снова опосредуется простым идеалом \(\mathfrak{c}_p\) из \(\mathfrak{o}_p\), содержащим \(\mathfrak{a}_p\). Тогда и только тогда \(p\) входит в знаменатель \(M(\mathfrak{I})\), когда \(\mathfrak{c}_p\) становится собственным надидеалом \(\mathfrak{a}_p\). Это может произойти только если либо степень трансцендентности \(\mathfrak{I}_p\) по сравнению с \(\mathfrak{I}\) уменьшается, либо \(\mathfrak{a}_p\) теряет свойство быть простым идеалом\footnote{Под "степенью трансцендентности" системы, как известно, понимается инвариантное число алгебраически независимых функций, от которых система алгебраически зависит; система, состоящая из алгебраически независимых функций, называется неприводимой. Под степенью трансцендентности, или размерностью, простого идеала понимается степень трансцендентности его поля вычетов. Простой идеал не имеет собственного простого делителя той же степени трансцендентности; ср. B. L. van der Waerden, Zur Nullstellentheorie der Polynomideale, Math. Ann. 96 (1926), с. 183--208.}. Что это происходит лишь для конечного числа \(p\), следует из того, что и по модулю \(p\) алгебраическая зависимость влечет обращение функционального определителя в нуль (§1), а также из того, что \(\mathfrak{a}\) оказывается абсолютным простым идеалом и поэтому теряет простоту только по модулю конечного числа \(p\)\footnote{E. Noether, Eliminationstheorie und allgemeine Idealtheorie, Math. Ann. 90 (1923), с. 229--261, §7. Утверждение следует с помощью теории исключения из более простой теоремы: абсолютно неприводимый многочлен теряет это свойство только по модулю конечного числа простых; A. Ostrowski, Zur arithmetischen Theorie der algebraischen Größen, Gött. Nachr. 1919, с. 279--298, там с. 296. Более простой доказательство см. у E. Noether, Ein algebraisches Kriterium für absolute Irreduzibilität, Math. Ann. 85 (1922), с. 26--33, §7. Если \(\mathfrak{I}\) имеет интегральный базис, содержащий ровно на одну функцию больше, чем его степень трансцендентности, как в случае Гекке, упомянутом в примечании 2, то \(\mathfrak{a}\) становится главным идеалом, и достаточно этой более простой теоремы без привлечения теории исключения. Вместо простых чисел всюду могут стоять простые идеалы конечного алгебраического числового поля. Утверждение, очевидно, верно и тогда, когда \(\mathfrak{a}\) является нулевым идеалом.}.

Заметим еще, что все рассуждения с незначительными изменениями сохраняются, если вместо целых рациональных чисел взять целые числа конечного алгебраического числового поля и его простые идеалы (§4).

\bigskip
\noindent\textbf{§ 1. Функциональные определители над областями коэффициентов характеристики \(p\).}

\smallskip
\noindent\textbf{1.}
Пусть \(f_1(x_1,\ldots,x_n),\ldots,f_t(x_1,\ldots,x_n)\) являются рациональными функциями или, более общо, формально заданными степенными рядами либо частными таких с коэффициентами из поля \(P\) характеристики \(p\). Производные \(\partial f_i/\partial x_k\) также должны быть формально заданы, причем сохраняются все обычные правила вычисления. Пусть \(\Omega\) обозначает алгебраически замкнутое расширение поля \(P\).

Тогда, как и в характеристике нуль, верно:

\noindent\textbf{Теорема.}
\textit{Если между \(f(x)\) имеется алгебраическая зависимость с коэффициентами из \(\Omega\), то ранг функциональной матрицы \((\partial f_i/\partial x_k)\), \(i=1,\ldots,t\), \(k=1,\ldots,n\), меньше \(t\)}\footnote{В первоначальной формулировке я предполагала, что область коэффициентов \(P\) является совершенным полем. Распространением на произвольные, в том числе несовершенные, поля через переход от \(P\) к \(\Omega\) я обязана B. L. van der Waerden. Что обратное утверждение в характеристике \(p\) даже для совершенного \(P\) неверно, показывает пример \(f=x^p\), \(\partial f/\partial x=0\).}.

\noindent\textit{Доказательство.}
В многочленной области \(\Omega[u_1,\ldots,u_t]\) пусть \(\mathfrak{m}\) обозначает простой идеал всех многочленов \(G(u)\), для которых \(G(f)=0\) тождественно по \(x\). По предположению \(\mathfrak{m}\) не является нулевым идеалом; пусть \(G(u)\ne0\) и, следовательно, \(G(u)\ne\mathrm{const.}\) является многочленом наименьшей степени из \(\mathfrak{m}\). Отсюда, в частности, следует, что \(G(u)\) неприводим. Тогда, как обычно, получается
\[
\frac{\partial G}{\partial u_1}\cdot\frac{\partial f_1}{\partial x_k}
+\cdots+
\frac{\partial G}{\partial u_t}\cdot\frac{\partial f_t}{\partial x_k}=0
\quad(k=1,\ldots,n),
\]
откуда следует, что либо ранг функциональной матрицы меньше \(t\), либо все \(\partial G/\partial u_i\) исчезают. Но так как \(G(u)\) был выбран многочленом наименьшей степени, отсюда следует исчезновение всех \(\partial G/\partial u_i\).

Итак, \(G(u)=H(u^p,\ldots,u^p)\); а так как область коэффициентов \(\Omega\) была алгебраически замкнутым, следовательно совершенным, полем, далее \(G(u)=H(u)^p=(H(u))^p\), что противоречит выбору \(G(u)\). Доказательство, конечно, годится и для характеристики нуль, где последний шаг отпадает.

\medskip\noindent\textbf{2. Следствие.}
Пусть \(F_1(x_1,\ldots,x_n),\ldots,F_t(x_1,\ldots,x_n)\) являются целочисленными функциями, то есть многочленами, степенными рядами или частными таких, с функциональной матрицей ранга ровно \(t\). Пусть простое \(p\) выбрано так, что оно не входит ни в знаменатели \(F(x)\), ни во все \(t\)-рядные определители функциональной матрицы.

Тогда \(F_1(x),\ldots,F_t(x)\) остаются алгебраически независимыми и по модулю \(p\).

\noindent\textit{Доказательство.}
Пусть \(\bar P\) обозначает поле вычетов по \(p\), а \(f_1(x),\ldots,f_t(x)\) --- функции над \(\bar P\), которые получаются из \(F_1(x),\ldots,F_t(x)\) заменой целочисленных коэффициентов их классами вычетов по \(p\). Функции \(f_i(x)\) определены, так как \(p\) не входило в знаменатели \(F_i(x)\). Их функциональная матрица \((\partial f_i/\partial x_k)\) также получается из \((\partial F_i/\partial x_k)\) заменой целочисленных коэффициентов классами вычетов по \(p\); в \((\partial F_i/\partial x_k)\) не входят другие знаменатели, кроме тех, что уже есть в \(F(x)\). Благодаря выбору \(p\) матрица \((\partial f_i/\partial x_k)\) сохраняет ранг \(t\); значит, \(f(x)\) алгебраически независимы в \(\Omega\), а тем более в \(\bar P\).

\bigskip
\noindent\textbf{§ 2. Формулировка теоремы.}

\smallskip
\noindent\textbf{1.}
Как во введении, пусть \(\mathfrak{I}\) обозначает область целостности, то есть кольцо без делителей нуля, из целочисленных функций от \(x_1,\ldots,x_n\): многочленов или степенных рядов с целыми рациональными коэффициентами либо частных таких. В случае степенных рядов они могут быть формально заданными или сходящимися в некоторой области; используется только то, что возникающие функции образуют область целостности.

Пусть \(\mathfrak{Q}\) обозначает поле частных \(\mathfrak{I}\), а \(\mathfrak{Z}\) --- область целостности между \(\mathfrak{I}\) и \(\mathfrak{Q}\). Как выше, \(M\) называется \textit{максимальной} в \(\mathfrak{Z}\), если из принадлежности \(a\cdot g(x)\), где \(a\ne0\) целое и \(g(x)\in\mathfrak{Z}\), всегда следует принадлежность \(g(x)\). Для \(\mathfrak{I}\) существует единственная максимальная область \(M(\mathfrak{I})\), порожденная \(\mathfrak{I}\) в \(\mathfrak{Z}\), как пересечение всех максимальных областей в \(\mathfrak{Z}\), содержащих \(\mathfrak{I}\): ведь сама \(\mathfrak{Z}\) такова, а пересечение любого числа таких областей снова максимально и содержит \(\mathfrak{I}\). \(M(\mathfrak{I})\) состоит из всех функций \(g(x)\) из \(\mathfrak{Z}\), для которых существует хотя бы одно \(a\ne0\) с \(a\cdot g(x)\in\mathfrak{I}\).

Действительно, эти \(g(x)\) образуют содержащую \(\mathfrak{I}\) максимальную область \(N\) в \(\mathfrak{Z}\), потому что из \(b\cdot h(x)\in\mathfrak{I}\), \(b\ne0\), \(h(x)\in\mathfrak{Z}\), следует \(ab\cdot h(x)\in\mathfrak{I}\) и \(ab\ne0\), значит \(h(x)\in N\). Любая максимальная область \(M\) в \(\mathfrak{Z}\), содержащая \(\mathfrak{I}\), содержит и \(N\): если \(a\cdot g(x)\in\mathfrak{I}\), то \(a\cdot g(x)\in M\), так как \(\mathfrak{I}\subset M\), следовательно \(g(x)\in M\).

\medskip\noindent\textbf{2.}
Теперь предположим, что \(\mathfrak{I}\) является конечной областью целостности с единицей, с интегральным базисом \(f_1(x),\ldots,f_r(x)\), то есть \(\mathfrak{I}\) состоит из всех целочисленных многочленов от \(f_i(x)\). Далее, пусть \(\mathfrak{Z}\) состоит из всех целочисленных многочленов или степенных рядов от \(x\), соответственно из таких частных, которые содержат только знаменатели, встречающиеся в конечном числе \(f_i(x)\), конечно, в любых степенях.

Предположим также, что хотя бы в одной \(f_i(x)\), а вообще говоря во всех, переменные \(x\) действительно входят; иначе \(\mathfrak{I}\) и \(\mathfrak{Z}\) просто состоят из всех целочисленных многочленов от дробного числа \(1/e\), и доказывать нечего.

Если речь идет о многочленах или рациональных функциях, то поле частных \(\mathfrak{Q}\) имеет степень трансцендентности \(t\), \(1\le t\le r\), над полем \(K\) рациональных чисел. Если, например, \(f_1(x),\ldots,f_t(x)\) образуют неприводимую систему\footnote{Под "неприводимой системой", как известно, понимается система из алгебраически независимых функций.}, то \(\mathfrak{Q}\) является конечным алгебраическим расширением чисто трансцендентного поля \(K(f_1(x),\ldots,f_t(x))\). При этом функциональная матрица \(f_1(x),\ldots,f_r(x)\), а также функциональная матрица \(f_1(x),\ldots,f_t(x)\), имеет ранг ровно \(t\)\footnote{Это дополнительное условие выполнено, когда, как в рассмотренном Гекке случае, речь идет об аналитических функциях, \(f_1(x),\ldots,f_t(x)\) аналитически независимы, а \(f_{t+1}(x),\ldots,f_r(x)\) алгебраически зависят от них. Но условие намеренно сформулировано формально --- ранг функциональной матрицы, --- чтобы оно имело ясный смысл и для формально заданных степенных рядов. Тем самым все рассуждения остаются чисто формально-алгебраическими, а теоремы об аналитических функциях нужны только для проверки условия в специальном случае.}.

Чтобы и в случае степенных рядов или их частных эта структура \(\mathfrak{Q}\) сохранялась, нужно сделать дополнительное предположение: если \(t\) является рангом функциональной матрицы \(f_i(x)\) и одновременно, скажем, рангом функциональной матрицы \(f_1(x),\ldots,f_t(x)\), так что они по §1, 1 образуют неприводимую систему, то \(\mathfrak{Q}\) снова должно быть конечным алгебраическим расширением чисто трансцендентного \(K(f_1(x),\ldots,f_t(x))\). То есть \(f_{t+1}(x),\ldots,f_r(x)\) должны быть алгебраическими над \(K(f_1(x),\ldots,f_t(x))\)\footnote{Это дополнительное условие выполнено, когда, как в случае Гекке, речь идет об аналитических функциях, \(f_1(x),\ldots,f_t(x)\) аналитически независимы, а \(f_{t+1}(x),\ldots,f_r(x)\) алгебраически зависят от них.}. При этом снова \(t\ge1\), так как в характеристике нуль не все \(\partial f_i/\partial x_k\) исчезают.

Одновременно условие следует для всякой системы из \(t\) функций интегрального базиса, функциональная матрица которых имеет ранг ровно \(t\): они образуют неприводимую систему, от которой остальные функции алгебраически зависят.

\medskip\noindent\textbf{3.}
Пусть теперь \(\mathfrak{o}\) обозначает область целочисленных многочленов от неопределенных \(u_1,\ldots,u_r\). Тогда \(\mathfrak{I}\) является кольцево-гомоморфным образом \(\mathfrak{o}\), что непосредственно видно из соответствия \(u_i\mapsto f_i(x)\). Значит, \(\mathfrak{I}\) кольцево изоморфна \(\mathfrak{o}/\mathfrak{a}\), где \(\mathfrak{a}\) --- простой идеал, так как \(\mathfrak{I}\) кольцо без делителей нуля; \(\mathfrak{a}\) состоит из всех многочленов \(G(u)\), для которых \(G(f)\) тождественно исчезает по \(x\).

Как уже сказано во введении, речь идет о следующей теореме.

\medskip\noindent\textbf{Теорема.}
\textit{При дополнительном условии из пункта 2 пусть \(\mathfrak{o}_p,\mathfrak{a}_p,\mathfrak{I}_p\) обозначают области, возникающие из \(\mathfrak{o},\mathfrak{a},\mathfrak{I}\) заменой целочисленных коэффициентов их классами вычетов по модулю \(p\). Тогда, за исключением не более чем конечного числа простых чисел, включая входящие в знаменатели \(f(x)\), чтобы \(\mathfrak{I}_p\) была определена, гомоморфизм из \(\mathfrak{o}_p\) в \(\mathfrak{I}_p\) опосредуется \(\mathfrak{a}_p\). Следовательно, \(\mathfrak{I}_p\) кольцево изоморфна \(\mathfrak{o}_p/\mathfrak{a}_p\), а тем самым \(\mathfrak{o}/(\mathfrak{a},p\mathfrak{o})\).}

Из теоремы факт о конечном числе простых в знаменателях \(a\) области \(M(\mathfrak{I})\) получается посредством следующего следствия.

\medskip\noindent\textbf{Следствие.}
\textit{Если \(p\) выбрано так, что \(\mathfrak{I}_p\) изоморфна \(\mathfrak{o}_p/\mathfrak{a}_p\), то из \(p\cdot g(x)\in\mathfrak{I}\) и \(g(x)\in\mathfrak{Z}\) всегда следует \(g(x)\in\mathfrak{I}\).}

Иначе говоря, \(\mathfrak{I}\) максимальна в \(\mathfrak{Z}\) относительно всех простых, кроме конечного числа исключительных простых.

Имеем \(\mathfrak{o}_p/\mathfrak{a}_p\simeq\mathfrak{o}/(\mathfrak{a},p\mathfrak{o})\). Действительно, по определению \(\mathfrak{o}_p=\mathfrak{o}/p\mathfrak{o}\), соответственно \(\mathfrak{a}_p=(\mathfrak{a},p\mathfrak{o})/p\mathfrak{o}\). Значит, по первой теореме об изоморфизме \(\mathfrak{o}_p/\mathfrak{a}_p\), а по предположению и \(\mathfrak{I}_p\), изоморфны \(\mathfrak{o}/(\mathfrak{a},p\mathfrak{o})\).

Это означает: из \(H(f_i(x))\equiv0\pmod p\) в \(\mathfrak{I}\) следует \(H(u)\equiv0\pmod{(\mathfrak{a},p\mathfrak{o})}\) в \(\mathfrak{o}\).

Это можно преобразовать так: из \(H(f(x))=p\cdot g(x)\), где \(H(f(x))\in\mathfrak{I}\) и \(g(x)\in\mathfrak{Z}\), следует \(H(u)-pF(u)\equiv0\pmod{\mathfrak{a}}\), а потому также \(H(f_i(x))=pF(f_i(x))\) и, значит, \(g(x)=F(f_i(x))\). Следовательно, \(g(x)\in\mathfrak{I}\).

Наконец, конечное повторение дает: если \(a\) не делится ни на одно из исключительных простых, то из \(a\cdot g(x)\in\mathfrak{I}\) и \(g(x)\in\mathfrak{Z}\) всегда следует \(g(x)\in\mathfrak{I}\). Следствие доказано.

\bigskip
\noindent\textbf{§ 3. Доказательство теоремы.}

\smallskip
\noindent\textbf{1. Лемма.}
\textit{Простой идеал \(\mathfrak{a}\), опосредующий гомоморфизм между \(\mathfrak{I}\) и \(\mathfrak{o}\) (§2, 3), является абсолютным простым идеалом.}

Пусть \(\mathfrak{o}^*\) обозначает многочленную область от \(u_1,\ldots,u_r\) с произвольными, в том числе дробными, алгебраическими числами в качестве коэффициентов, а \(\mathfrak{a}^*\) --- идеал в \(\mathfrak{o}^*\), порожденный \(\mathfrak{a}\). Значит, \(\mathfrak{a}^*\) состоит из всех линейных комбинаций \(a_1G_1(u)+\cdots+a_sG_s(u)\), где \(a_i\) --- произвольные алгебраические числа, а \(G_i(u)\) --- многочлены из \(\mathfrak{a}\). Нужно доказать, что \(\mathfrak{a}^*\) является простым идеалом.

Это будет доказано, если показать, что \(\mathfrak{a}^*\) опосредует гомоморфизм из \(\mathfrak{o}^*\) в \(\mathfrak{I}^*\), где под \(\mathfrak{I}^*\) понимается область целостности всех многочленов от \(f_i(x)\) с произвольными алгебраическими коэффициентами. То есть надо показать, что \(H(u)\in\mathfrak{a}^*\), если \(H(u)\in\mathfrak{o}^*\) и \(H(f_i(x))\) исчезает.

Пусть \(\omega_1,\ldots,\omega_s\) является линейно независимым базисом конечного числового поля, порожденного конечным числом коэффициентов \(H(u)\). Тогда
\[
a\cdot H(u)=\omega_1H_1(u)+\cdots+\omega_sH_s(u),
\]
где \(H_i(u)\in\mathfrak{o}\), а \(a\ne0\) --- целое рациональное число. Из \(H(f)=0\) следует \(\omega_1H_1(f)+\cdots+\omega_sH_s(f)=0\), а значит \(H_1(f)=0,\ldots,H_s(f)=0\). Поэтому \(H_1(u),\ldots,H_s(u)\) принадлежат \(\mathfrak{a}\), и, следовательно, \(H(u)\in\mathfrak{a}^*\).

\medskip\noindent\textbf{2.}
Доказательство теоремы теперь можно провести в следующей форме.

Если \(p\) выбрано так, что
\begin{itemize}
\item[I.] \(\mathfrak{I}_p\) определена,
\item[II.] \(\mathfrak{a}_p\) остается простым, притом абсолютным, идеалом,
\item[III.] степень трансцендентности \(\mathfrak{I}_p\) по сравнению со степенью трансцендентности \(\mathfrak{I}\) не уменьшилась,
\end{itemize}
то гомоморфизм из \(\mathfrak{o}_p\) в \(\mathfrak{I}_p\) опосредуется \(\mathfrak{a}_p\). Эти условия, за исключением не более чем конечного числа исключительных простых, выполнены для всех простых чисел.

То, что относительно условий I, II, III имеется лишь конечное число исключительных простых, непосредственно следует из предыдущего. I ясно: исключительными являются лишь конечное число простых, входящих в знаменатели \(f_i(x)\). II следует из того, что \(\mathfrak{a}\) по лемме является абсолютным простым идеалом и потому сохраняет это свойство по модулю всех простых, за исключением не более чем конечного числа. III показано в следствии §1, 2, поскольку функциональная матрица \(f_1(x),\ldots,f_t(x)\) по §2, 2 имеет ранг ровно \(t\).

Теперь пусть \(p\) выбрано отличным от этих исключительных простых. Гомоморфизм из \(\mathfrak{o}_p\) в \(\mathfrak{I}_p\) опосредуется идеалом \(\mathfrak{c}_p\), который является простым, так как \(\mathfrak{I}_p\) также кольцо без делителей нуля, и степень трансцендентности которого задается степенью трансцендентности \(\mathfrak{I}_p\), следовательно не меньше \(t\). Ведь \(\mathfrak{I}_p\simeq\mathfrak{o}_p/\mathfrak{c}_p\), поэтому поле вычетов \(\mathfrak{c}_p\) изоморфно полю частных \(\mathfrak{I}_p\).

Так как \(\mathfrak{a}_p\) также простой идеал и кратен, то есть является подидеалом, \(\mathfrak{c}_p\), то \(\mathfrak{a}_p\) и \(\mathfrak{c}_p\) совпадают тогда и только тогда, когда совпадают их степени трансцендентности\footnote{Ср., например, B. L. van der Waerden, Zur Nullstellentheorie der Polynomideale, Math. Ann. 96 (1926), с. 183--208, §3.}. Степень трансцендентности \(\mathfrak{a}_p\), как кратного \(\mathfrak{c}_p\), также не меньше \(t\); классы \(u_1,\ldots,u_t\) по модулю \(\mathfrak{a}_p\) образуют неприводимую систему. Нужно показать, что эта степень трансцендентности в точности равна \(t\), чем все будет доказано.

Здесь в случае степенных рядов или их частных используется дополнительное предположение о структуре поля частных \(\mathfrak{Q}\) области \(\mathfrak{I}\) (§2, 2). По нему простой идеал \(\mathfrak{a}\) во всяком случае имел степень трансцендентности \(t\), и нумерация была выбрана так, что классы \(u_1,\ldots,u_t\) по модулю \(\mathfrak{a}\) образовывали неприводимую систему, от которой классы \(u_{t+1},\ldots,u_r\) алгебраически зависели. Следовательно, в \(\mathfrak{a}\) содержались целочисленные и примитивные многочлены
\[
G_1(u_{t+1};u_1,\ldots,u_t),\ldots,G_{r-t}(u_r;u_1,\ldots,u_t),\tag{1}
\]
переходящие в многочлены
\[
G^*_1(u_{t+1};u_1,\ldots,u_t),\ldots,G^*_{r-t}(u_r;u_1,\ldots,u_t),\tag{2}
\]
из \(\mathfrak{a}_p\), не исчезающие благодаря примитивности. Но в каждом \(G^*_j\) соответствующая переменная \(u_{t+j}\) действительно должна входить, иначе, вопреки выбору \(p\), имелась бы алгебраическая зависимость между классами \(u_1,\ldots,u_t\) по модулю \(\mathfrak{a}_p\). Тем самым показано, что \(\mathfrak{a}_p\) имеет степень трансцендентности \(t\), поэтому совпадает с \(\mathfrak{c}_p\) и опосредует гомоморфизм. Теорема доказана.

\bigskip
\noindent\textbf{§ 4. Распространение на целые алгебраические коэффициенты.}

\smallskip
Как уже отмечено в конце введения, все рассуждения с небольшими изменениями сохраняются, если вместо целых рациональных чисел взять целые числа конечного алгебраического числового поля \(K\), а вместо простых чисел \(p\) --- простые идеалы \(\mathfrak{p}\) этого числового поля. При этом §1 полностью сохраняется, как и рассуждения, ведущие к формулировке теоремы в §2; в доказательстве (§3, 2) и в следствии из §2, 3 нужны некоторые добавления.

\medskip\noindent\textbf{1.}
В §3, 2 к исключительным простым идеалам надо добавить еще конечное число тех, которые входят в многочлены \(G_1(u),\ldots,G_{r-t}(u)\) из §3, 2, (1), потому что теперь их уже нельзя считать примитивными. Тем самым обеспечивается ненулевость многочленов \(G^*_1(u),\ldots,G^*_{r-t}(u)\) из §3, 2, (2); это является усилением условия III. Условие II сохраняется, поскольку теорема об абсолютной неприводимости не меняется; равно как, конечно, сохраняется и I.

\medskip\noindent\textbf{2.}
Следствие из §2, 3, представляющее конечный результат, согласно тамошней окончательной формулировке следует сформулировать так.

\noindent\textbf{Следствие.}
\textit{Если целое число \(a\) из \(K\) не делится ни на один из конечного числа исключительных простых идеалов, то из \(a\cdot g(x)\in\mathfrak{I}\) и \(g(x)\in\mathfrak{Z}\) всегда следует \(g(x)\in\mathfrak{I}\).}

Если \(\lambda_1,\lambda_2,\ldots\) --- целые числа из \(K\), выбранные так, что каждое из них делится на один исключительный простой идеал, но не на его квадрат и не на остальные, то в качестве знаменателей \(a\) достаточно брать произведения степеней этих \(\lambda\).

\noindent\textit{Доказательство.}
Пусть
\[
(a)=\mathfrak{p}_1^{\alpha_1}\cdots\mathfrak{p}_g^{\alpha_g}
\]
есть разложение \((a)\) на простые идеалы, причем \(\mathfrak{p}_1,\ldots,\mathfrak{p}_g\) отличны от исключительных простых идеалов; тогда \(\mathfrak{I}_{\mathfrak{p}_i}\) изоморфна \(\mathfrak{o}_{\mathfrak{p}_i}/\mathfrak{a}_{\mathfrak{p}_i}\) для каждого \(\mathfrak{p}_i\). От кольца всех целых чисел из \(K\) переходим к кольцу частных \(R_a\), определенному идеалом \((a)\), присоединяя к знаменателям все и только те целые числа, которые взаимно просты с \((a)\), то есть не делятся ни на один из идеалов \(\mathfrak{p}_1,\ldots,\mathfrak{p}_g\). В \(R_a\) идеалы \(\mathfrak{p}_1,\ldots,\mathfrak{p}_g\) остаются простыми, но становятся главными; все остальные становятся единичным идеалом; разложение \((a)\) остается тем же\footnote{Ср., например, H. Grell, Zur Theorie der Ordnungen in algebraischen Zahl- und Funktionenkörpern, Math. Ann. 97 (1927), с. 524--558, §2, 1.}. Поэтому при переходе к \(R_a\) как области коэффициентов следствие и его доказательство сохраняются в точной форме §2, 3.

Это означает: если \(a\cdot g(x)\in\mathfrak{I}\), то \(g(x)=F(f_i(x))\), где многочлен \(F(u)\) имеет коэффициенты из \(R_a\). Если \(d\) --- общий знаменатель этих коэффициентов, то есть \(d\) взаимно прост с \(a\), то это можно сформулировать и так: из \(a\cdot g(x)\in\mathfrak{I}\) следует \(d\cdot g(x)\in\mathfrak{I}\) с \(d\), взаимно простым с \(a\). Значит, \((d,a)\) является единичным идеалом, и поэтому из \(a\cdot g(x)\in\mathfrak{I}\) следует \(g(x)\in\mathfrak{I}\). Первая часть следствия доказана.

Для доказательства второй части пусть \(R^*\) --- кольцо частных в \(K\), определенное конечным числом исключительных простых идеалов, а \(\lambda_1,\ldots,\lambda_s\) --- базисные элементы этих простых идеалов в \(R^*\), которые можно считать целыми числами. В \(R^*\) каждое \(a\) с точностью до единицы из \(R^*\) является произведением степеней \(\lambda\). Возвращаясь к целым числам, получаем
\[
\gamma a=\beta\lambda_1^{\alpha_1}\cdots\lambda_s^{\alpha_s},
\]
где \(\gamma\) и \(\beta\) не делятся ни на один исключительный простой идеал. Поэтому из \(a\cdot g(x)\in\mathfrak{I}\) следует, что \(\beta\lambda_1^{\alpha_1}\cdots\lambda_s^{\alpha_s}g(x)\), а значит, по первой части следствия, и \(\lambda_1^{\alpha_1}\cdots\lambda_s^{\alpha_s}g(x)\), лежит в \(\mathfrak{I}\).

\begin{flushright}
\textit{(Поступила 11 ноября 1931 г.)}
\end{flushright}

\bigskip
\noindent\rule{\textwidth}{0.4pt}

\medskip
\noindent\textbf{Аннотация} (переведена по русскому резюме оригинала):

\smallskip
\textit{Каждая область целостности \(\mathfrak{I}\) определяет некоторую максимальную область \(M(\mathfrak{I})\) того же типа, состоящую из всех функций \(g(x)\), для которых существует хотя бы одно \(a\ne0\) такое, что \(a\cdot g(x)\) содержится в \(\mathfrak{I}\).}

\textit{В этой работе исследуется вопрос о "знаменателях" \(a\) этой максимальной области \(M(\mathfrak{I})\) при предположении, что область \(\mathfrak{I}\) была конечной. Доказано, что эти \(a\) всегда можно выбирать так, чтобы в них в совокупности входило лишь конечное число простых чисел как множителей.}

\textit{При этом в случае, когда \(g(x)\) являются степенными рядами или частными таких рядов, предполагается, что функции \(f_i(x)\), задающие область \(\mathfrak{I}\), могут быть связаны между собой только алгебраическими соотношениями.}

\textit{Доказательство основано на сведении задачи к некоторым вопросам общей теории идеалов. Это сведение осуществляется через рассмотрение появления простого числа \(p\) в знаменателе \(a\) как соотношения по модулю \(p\) между функциями \(f_i\).}
% END UNIT BODY
% END PRODUCER UNIT 211
\endgroup


\clearpage
\begingroup
\setcounter{footnote}{0}
% BEGIN PRODUCER UNIT 212
% Source: C:/Users/Floris/Documents/interlanguage/03_projects/noether/02_slavic_working_corpus/translations/paper36/source_fidelity/russian/v001/Noether_Paper36_SourceFidelity_Russian_v001.tex
% Identity: 1591 B / SHA-256 6291544941EF02B09696E43145941141C2B2FDA3810658AAAAC3F768403E1A36
\setlength{\parindent}{1.15em}
\setlength{\parskip}{0pt}
\emergencystretch=180em
\allowdisplaybreaks
\sloppy
\hbadness=10000
\vbadness=10000
\hfuzz=3pt
\vfuzz=10pt
\raggedbottom
% BEGIN UNIT BODY
% Unit: paper36_source_fidelity_v001.
% Language lane: Russian.
% Direct source: sources/paper36/source_fidelity/Noether_Paper36_ORIGINAL_JDMV39_GDZ600_R124plusP40_repaired_witness_v001.tex.
\section*{36. Дифференцирование идеалов и дифферента}

\begin{center}
\emph{J. Ber. d. DMV39 (1930), с. 17}
\end{center}

Э. Нётер, Гёттинген: Дифференцирование идеалов и дифферента.

Показано, как дифференту алгебраического числового поля можно понимать как \emph{дифференциальное частное} определяющего \emph{идеала}. То, что это определение совпадает с обычным, следует из структурных теорем, интересных самих по себе, которые в аналогичном смысле дают обобщение интерполяционной формулы Лагранжа.

Подробное изложение должно появиться в \emph{Mathematische Annalen}.
% END UNIT BODY
% END PRODUCER UNIT 212
\endgroup


\clearpage
\begingroup
\setcounter{footnote}{0}
% BEGIN PRODUCER UNIT 213
% Source: C:/Users/Floris/Documents/interlanguage/03_projects/noether/02_slavic_working_corpus/translations/paper37/source_fidelity/russian/v001/Noether_Paper37_SourceFidelity_Russian_v001.tex
% Identity: 34190 B / SHA-256 5A71A3021C7D0B3A5D592C8023E9647024DBEAC0678B1A9BED0DE434D5A5D9BA
\setlength{\parindent}{1.1em}
\setlength{\parskip}{0pt}
\providecommand{\frakO}{\mathfrak{O}}
\providecommand{\frako}{\mathfrak{o}}
\providecommand{\frakp}{\mathfrak{p}}
\providecommand{\Gg}{\mathfrak{G}}
\emergencystretch=220em
\allowdisplaybreaks
\sloppy
\hbadness=10000
\vbadness=10000
\hfuzz=4pt
\vfuzz=10pt
\raggedbottom
% BEGIN UNIT BODY
% Unit: paper37_source_fidelity_v001.
% Language lane: Russian.
% Direct source: sources/paper37/source_fidelity/Noether_Paper37_ORIGINAL_JReineAngewMath167_R124plus_source_fidelity_witness_v001.tex.
\section*{37. Нормальный базис в полях без высшего ветвления}

\begin{center}
\emph{Journal f. d. reine u. angew. Math. 167 (1932), с. 147--152}
\end{center}

\begin{center}
Von \emph{Emmy Noether} in Göttingen.
\end{center}

Под нормальным базисом числового поля Галуа \(K/k\) понимают базис максимального порядка \(\frakO\) поля \(K\) относительно максимального порядка \(\frako\) поля \(k\), состоящий из сопряженных элементов. Необходимым условием существования такого базиса, как известно\footnote{Ср. A. Speiser, Gruppendeterminante und Körperdiskriminante, § 6. Math. Ann. 77 (1916), с. 546--562.}, является то, что в \(K/k\) не возникают высшие поля ветвления. Это остается верным и тогда, когда ограничиваются одним местом, то есть переходят к \(\frakp\)-адическому расширению \(k_\frakp\) поля \(k\) и соответствующему \(K_\frakp\) поля \(K\). Ниже я доказываю обратное утверждение:

\emph{Для всех мест \(\frakp\), где \(\frakp\) не делит степень \(K/k\), существует нормальный базис. В частности: если \(K/k\) не имеет высшего ветвления, то в каждом разветвленном месте существует нормальный базис; дискриминант там становится квадратом группового определителя.}

К этому последнему утверждению надо добавить, что, вопреки предположению Э. Артина, переход к его проводникам\footnote{E. Artin, Über die gruppentheoretische Struktur der Diskriminante. J. f. Math. 164 (1931), с. 1--11.} требует еще дальнейших соображений: разложение группового определителя по неприводимым характерам дает обобщенные корневые числа; подходящее объединение дает разложение в основном поле, расширенном характерами. В простейших случаях я могу отождествить эти новые множители с проводниками Артина посредством идеал-теоретического разложения корневых чисел.\footnote{[6. 10. 31.] Что и вообще еще будут необходимы идеал-теоретические соображения, показывает следующий, сообщенный мне М. Дойрингом и произвольно варьируемый, пример немаксимального порядка, дискриминант которого представляется квадратом группового определителя, тогда как сопряженным характерам соответствуют различные идеальные множители. Пусть \(k\) есть поле пятых корней из единицы, \(K=k(\sqrt[5]{2})\). Рассматривается порядок \([1,2\sqrt[5]{2},\sqrt[5]{2^2},\sqrt[5]{2^3},\sqrt[5]{2^4}]\), именно в месте \((2)\), где он по рассуждениям этой заметки имеет нормальный базис. Разложение по характерам дает для соответствующего группового определителя именно приведенные выше элементы базиса как множители. Поэтому ``проводники'', кроме \(1\), становятся: \(2\sqrt[5]{2}/\sqrt[5]{2^4}\), \(\sqrt[5]{2^2}/\sqrt[5]{2^3}\), \(\sqrt[5]{2^3}/\sqrt[5]{2^2}\), \(\sqrt[5]{2^4}/(2\sqrt[5]{2})\), то есть \(4,2,2,4\), и потому различны для сопряженных характеров.} Далее замечу, что и в общем случае, когда нормального базиса не существует, расщепление на обобщенные корневые числа всегда возможно посредством ``композиционного ряда по модулям Галуа''; и что отсюда, аналогично сказанному выше, получается разложение в расширенном основном поле; здесь снова начинаются идеал-теоретические вопросы.

Доказательство существования нормального базиса при указанных предположениях основано на том, что \(\frakO/\frako\), рассматриваемое как модуль Галуа, то есть как \(\frako\)-модуль с подстановками группы Галуа как областью операторов, становится операторно изоморфным одностороннему идеалу целочисленного группового кольца, расширенного с помощью \(\frako\). Это кольцо во всех обычных местах ветвления является максимальным порядком; но идеалы в максимальных порядках \(\frakp\)-адических полупростых систем являются главными идеалами\footnote{Ср. H. Hasse, Über \(\frakp\)-adische Schiefkörper und ihre Bedeutung für die Arithmetik hyperkomplexer Zahlsysteme, теоремы 46 и 47. Math. Ann. 104 (1931), с. 495--534. Для коммутативных систем свойство главных идеалов выполняется уже при переходе к кольцу частных по \(\frakp\) в \(k\); следовательно, для определения места можно ограничиться этим более слабым расширением, если речь идет об абелевых группах и полях.}, и потому здесь они возникают из одного элемента подстановками группы или группового кольца. Возвращаясь отсюда к модулю Галуа, получаем как раз существование нормального базиса. В местах высшего ветвления целочисленное групповое кольцо переходит в немаксимальный порядок, а идеал, соответствующий \(\frakO/\frako\), переходит в неглавный идеал.

Все эти рассуждения, очевидно, сохраняются и для функциональных полей одной переменной, если характеристика поля констант не делит степень \(K/k\).

\subsection*{§1. \(\frakp\)-адически расширенное целочисленное групповое кольцо}

Пусть \(\frakO\) и соответственно \(\frako_\frakp\) обозначают максимальные порядки алгебраического числового поля \(k\) и его \(\frakp\)-адического расширения \(k_\frakp\), где \(\frakp\) есть простой идеал из \(k\). Далее пусть \((\Gg)\) обозначает рациональное, а \([\Gg]\) целочисленное групповое кольцо группы \(\Gg\) с числом элементов \(n\). Под \((\Gg)_k\) и соответственно \((\Gg)_{k_\frakp}\) понимается групповое кольцо с коэффициентами из \(k\) и соответственно из \(k_\frakp\); аналогично под \([\Gg]_\frako\) и \([\Gg]_{\frako_\frakp}\) понимается расширение целочисленного группового кольца до такого с коэффициентами из \(\frako\) и соответственно из \(\frako_\frakp\).

Следовательно, \((\Gg)_k\) и \((\Gg)_{k_\frakp}\) являются системами без радикала, то есть полупростыми системами, а \([\Gg]_\frako\) и \([\Gg]_{\frako_\frakp}\) являются порядками в \((\Gg)_k\) и соответственно в \((\Gg)_{k_\frakp}\).

\paragraph{Теорема 1.}
\emph{Если простой идеал \(\frakp\) не делит число элементов \(n\) группы \(\Gg\), то \([\Gg]_{\frako_\frakp}\) становится максимальным порядком в \((\Gg)_{k_\frakp}\); если же \(\frakp\) делит \(n\), то \([\Gg]_{\frako_\frakp}\) не является максимальным.}

Дискриминант целочисленного группового кольца, а значит также \([\Gg]_\frako\) и \([\Gg]_{\frako_\frakp}\), равен степени числа элементов \(n\).\footnote{Ср., например, E. Noether, Hyperkomplexe Größen und Darstellungstheorie, § 26. Math. Ztschr. 30 (1929), с. 641--692.} Поэтому если \(\frakp\) не делит \(n\), то дискриминант \([\Gg]_{\frako_\frakp}\) является единицей. Это означает, что \([\Gg]_{\frako_\frakp}\), при фиксированном \(\frako_\frakp\), нельзя расширить до большего порядка: такому расширению соответствовало бы отделение квадратного неединичного множителя от дискриминанта. Но \(\frako_\frakp\) уже предполагался максимальным порядком, то есть максимальным в \(k_\frakp\). Тем самым первая часть теоремы 1, единственная используемая далее, доказана.

Если же \(\frakp\) делит \(n\), то прямая сумма, соответствующая единичному представлению,
\[
        \mathfrak C=E^{(1)}[\Gg]_{\frako_\frakp}+(1-E^{(1)})[\Gg]_{\frako_\frakp},
        \qquad
        E^{(1)}=\frac1n\sum_{S\in\Gg} S,
\]
является собственным расширением \([\Gg]_{\frako_\frakp}\). Действительно, вследствие \(E^{(1)}=\frac1n\sum S\) порядок \(\mathfrak C\) становится собственным расширением; \(\mathfrak C\) есть порядок с зависимыми образующими \(E^{(1)}\), \((1-E^{(1)})S\) при \(S\in\Gg\), причем \(E^{(1)}S=E^{(1)}\).

\paragraph{Теорема 2.}
\emph{Если \(\frakp\) не делит \(n\), то каждый целый или дробный идеал относительно \([\Gg]_{\frako_\frakp}\) является главным идеалом.}

Ибо по теореме 1 \([\Gg]_{\frako_\frakp}\) является максимальным порядком \(\frakp\)-адической полупростой системы; следовательно, его идеалы являются главными идеалами.\footnote{Hasse, цит. соч. Свойство главных идеалов там доказано только для простых систем, но известными рассуждениями отсюда оно следует и для полупростых. Действительно, компоненты максимального порядка снова максимальны; компоненты рассматриваемого идеала поэтому являются главными идеалами, скажем с базисами \(a_i\). Тогда \(a=\sum a_i\) является базисом исходного идеала. Для абелевых групп см. также примечание 3.}

\subsection*{§2. Модули Галуа, операторная изоморфность, нормальный базис}

Пусть \(K/k\) есть расширение Галуа, а \(\Gg\) его группа; пусть \(z^S\) обозначает элемент \(K\), возникающий из \(z\) подстановкой \(S\) из \(\Gg\). Подстановки из \(\Gg\) порождают на \(K/k\), рассматриваемом как \(k\)-модуль, то есть как аддитивная абелева группа с \(k\) как областью операторов, операторные автоморфизмы. Поэтому \(K/k\) можно определить как модуль относительно \((\Gg)_k\) правилом
\[
        z\Bigl(\sum_i S_i c_i\Bigr)=\sum_i z^{S_i}c_i,
        \qquad z\in K,\; S_i\in\Gg,\; c_i\in k.
\]

\paragraph{Определение.}
\emph{Каждый \((\Gg)_k\)-модуль из \(K/k\) называется рациональным модулем Галуа. Аналогично \([\Gg]_\frako\)-модули из \(K/k\) называются целыми модулями Галуа; в частности, \(\frakO/\frako\) является целым модулем Галуа.}

\paragraph{Теорема 3.}
\emph{Как рациональный модуль Галуа, \(K/k\) операторно изоморфен \((\Gg)_k\).}

Это следует из того, что \(K/k\) всегда имеет нормальный базис поля, то есть \(k\)-базис из сопряженных элементов \(z^S\).\footnote{Действительно, если \(a_1,\ldots,a_n\) есть базис \(K/k\), а \(u_i\) обозначают неопределенные, то \(\sum u_i a_i^S\) линейно независимы, когда \(S\) пробегает элементы \(\Gg\); следовательно, \(u_i\) можно специализировать до элементов из \(k\) так, чтобы линейная независимость сохранялась.} Соответствие
\[
        S\longmapsto z^S,
        \qquad
        \sum_i S_i c_i\longmapsto \sum_i z^{S_i}c_i
        \quad(c_i\in k),
        \quad\hbox{то есть } ST\longmapsto z^{ST},
\]
дает операторный гомоморфизм, который из-за одинакового ранга над \(k\) становится изоморфизмом.

\paragraph{Теорема 4.}
\emph{Как целый модуль Галуа, \(\frakO/\frako\) операторно изоморфен некоторому \([\Gg]_\frako\)-модулю из \((\Gg)_k\) максимального ранга \(n\). Аналогично \(\frakO_\frakp/\frako_\frakp\) операторно изоморфен \([\Gg]_{\frako_\frakp}\)-модулю ранга \(n\) из \((\Gg)_{k_\frakp}\).}

Указанный в теореме 3 \((\Gg)_k\)-изоморфизм \(K/k\to(\Gg)_k\) сопоставляет \([\Gg]_\frako\)-модулю \(\frakO/\frako\) \([\Gg]_\frako\)-модуль ранга \(n\) в \((\Gg)_k\). Далее этот \((\Gg)_k\)-изоморфизм можно расширить до \((\Gg)_{k_\frakp}\)-изоморфизма от \(K_\frakp/k_\frakp\) к \((\Gg)_{k_\frakp}\), просто позволяя коэффициентам \(c_i\) в соответствии из теоремы 3 пробегать все элементы \(k_\frakp\). Этот изоморфизм тогда индуцирует \([\Gg]_{\frako_\frakp}\)-изоморфизм \(\frakO_\frakp/\frako_\frakp\). Здесь \(K_\frakp/k_\frakp\) следует рассматривать как гиперкомплексную систему над \(k_\frakp\). Оно возникает из \(K/k\) расширением коэффициентов и вообще становится системой с делителями нуля, а именно суммой изоморфных полей, то есть полупростой системой.

\paragraph{Теорема 5.}
\emph{В каждом месте \(\frakp\), которое не делит степень \(n\) расширения \(K/k\), модуль \(\frakO_\frakp/\frako_\frakp\) имеет нормальный базис.}\footnote{В случае абелевых полей место при этом, по примечаниям 3 и 5, можно определять также кольцом частных.}

Действительно, по теореме 1 \([\Gg]_{\frako_\frakp}\) там является максимальным порядком; \([\Gg]_{\frako_\frakp}\)-модули ранга \(n\) из \((\Gg)_{k_\frakp}\) поэтому становятся целыми или дробными идеалами, причем главными идеалами по теореме 2. Если значит идеал, соответствующий \(\frakO_\frakp/\frako_\frakp\), равен \(W[\Gg]_{\frako_\frakp}\), то он имеет \(\frako_\frakp\)-базис \(WS_i\); операторный изоморфизм говорит, что \(w^{S_i}\) является \(\frako_\frakp\)-базисом \(\frakO_\frakp/\frako_\frakp\), если \(w\) обозначает элемент из \(\frakO_\frakp/\frako_\frakp\), соответствующий \(W\). Следовательно, \(n\) сопряженных элементов \(w^{S_i}\) образуют нормальный базис \(\frakO_\frakp/\frako_\frakp\).

Если \(\frakp\) делит \(n\), то модуль, соответствующий \(\frakO_\frakp/\frako_\frakp\), не может быть главным модулем, поскольку в этом случае нормальный базис не может существовать.

\paragraph{Дополнительное замечание к теореме 3.}
Из доказанной в теореме 3 операторной изоморфности \((\Gg)_k\) с \(K/k\)\footnote{Доказательство, очевидно, предполагает только, что \(K\) является расширением Галуа первого рода над \(k\), и что \(k\) имеет бесконечно много элементов. Последнее ограничение не нужно: М. Дойринг нашел доказательство теоремы 3, действительное также для полей с конечным числом элементов; из него, наоборот, следует существование нормального базиса поля. Одновременно так получается доказательство существования примитивного элемента, одинаково работающее для полей с конечным и бесконечным числом элементов.} следуют результаты Шпайзера о проблеме форм Клейна,\footnote{Цит. соч. Понятие модуля Галуа абстрагировано оттуда. Если характеристика \(k\) делит \(n\), то речь идет о композиционных рядах по модулям Галуа и их неприводимых факторах.} которые можно сформулировать так:

\emph{Каждое неприводимое над \(k\) представление группы Галуа \(\Gg\) расширения \(K/k\) порождается неприводимым рациональным модулем Галуа из \(K/k\); каждое абсолютно неприводимое порождается модулем Галуа из \(K_Z/Z\).}

Здесь \(Z\) обозначает поле разложения, содержащее \(k\), для соответствующей компоненты группового кольца; \(K_Z/Z\) следует рассматривать как гиперкомплексную систему над \(Z\), возникающую из \(K/k\) расширением коэффициентов. В частности, \(K_Z\) остается полем, если \(Z\) можно выбрать взаимно простым с \(K\) относительно \(k\), то есть если \(K_Z/Z\) становится акцессорным расширением, как в случае, рассмотренном Шпайзером.

Само утверждение следует посредством операторной изоморфности из соответствующего утверждения для \((\Gg)_k\) или \((\Gg)_Z\), где неприводимые идеалы порождают представление.

Если \(v_1,\ldots,v_t\) есть базис такого модуля Галуа, а \(S\mapsto \bar S\) соответствующее представление, то
\[
        \begin{pmatrix}\vdots\\ v_i^S\\ \vdots\end{pmatrix}
        =\bar S\begin{pmatrix}\vdots\\ v_i\\ \vdots\end{pmatrix}.
\]
Это и есть формулировка Клейна и Шпайзера.

\subsection*{§3. Дискриминант как групповой определитель в полях без высшего ветвления}

Для разложения дискриминанта, как известно, достаточно рассматривать разложение в отдельных местах ветвления; в полях без высшего ветвления, следовательно, речь идет только о местах с нормальным базисом.

\paragraph{Теорема 6.}
\emph{Для полей Галуа \(K/k\) без высшего ветвления дискриминант в каждом месте ветвления равен квадрату группового определителя группы Галуа, когда неопределенные заменены нормальным базисом. Соответственно неприводимым представлениям групповой определитель разлагается на обобщенные корневые числа; дискриминант разлагается на идеалы основного поля, расширенного характерами, причем комплексно-сопряженным характерам соответствуют одни и те же идеалы.}

Действительно, с \(w^S\) как нормальным базисом дискриминант \(\Delta\) равен квадрату
\[
        D=\bigl|w^{ST^{-1}}\bigr|,\qquad S,T\in\Gg.
\]
А \(D\) есть групповой определитель с \(w^S\) вместо неопределенных \(x_S\).

Разложение на корневые числа получается из рассуждений Шпайзера, цит. соч. Пусть
\[
        D=D_1^{f_1}\cdots D_t^{f_t},
        \qquad
        M=f_1M_1+\cdots+f_tM_t
\]
есть разложение группового определителя, соответственно групповой матрицы, по абсолютно неприводимым представлениям. Если \(S\mapsto\bar S\) обозначает представление, соответствующее \(M_\lambda\), где \(\bar S\) лежат в расширенном основном поле, то получается
\[
        M_\lambda=\sum_S w^S\bar S,
        \qquad
        M_\lambda^{T^{-1}}=\sum_S w^{ST^{-1}}\bar S
          =\sum_R w^R\bar R\bar T=M_\lambda\bar T,
\]
и при переходе к определителю
\[
        D_\lambda^{T^{-1}}=D_\lambda|\bar T|=D_\lambda\varepsilon_T.
\]
Тем самым \(D_\lambda\) распознаются как обобщенные корневые числа; \(\varepsilon_T\), как определитель \(\bar T\), является корнем из единицы, то есть неприводимым представлением фактор-группы \(\Gg\) по коммутаторной группе. В случае циклических групп \(D_\lambda\) являются просто корневыми числами Лагранжа \(\sum w^S\chi_\lambda(S)\).

Вообще \(D_\lambda\) лежат в гиперкомплексной системе, возникающей из \(K_\frakp/k_\frakp\) расширением области коэффициентов \(k_\frakp\) соответствующим характером; при этом речь идет о \emph{целых} величинах этой системы. Действительно, определитель, образованный с неопределенными \(x_S\), имеет коэффициенты, являющиеся целыми числами из поля характера; а \(w^S\) являются целыми элементами из \(K_\frakp\).

Чтобы перейти от разложения группового определителя к разложению дискриминанта, каждое представление объединяется с сопряженным к нему. Если \(\lambda,\bar\lambda\) обозначают сопряженные представления, то одновременно
\[
        D_\lambda^{T^{-1}}=D_\lambda\varepsilon_T,
        \qquad
        D_{\bar\lambda}^{T^{-1}}=D_{\bar\lambda}\varepsilon_T^{-1};
\]
вместе с тем определители, образованные для неопределенных \(x_S\) и принадлежащие \(\lambda,\bar\lambda\), комплексно сопряжены, тогда как вообще сопряженным характерам при неопределенных \(x_S\) соответствуют сопряженные определители.

Итак, если положить \(\Delta_\lambda=D_\lambda D_{\bar\lambda}\), то надо присоединить только вещественное подполе \(\lambda\)-го характера, и
\[
        \Delta_\lambda^T=\Delta_\lambda
        \qquad\hbox{для всех }T\hbox{ из }\Gg;
\]
отсюда\footnote{Поскольку речь идет о гиперкомплексном расширении коэффициентов поля, обычный вывод теории Галуа надо модифицировать. Пусть \(P\) есть рассматриваемое поле расширения поля \(k_\frakp\), а
\[
        (K_\frakp)_P=(K_\frakp)_\mathfrak P e^{S_1}+\cdots+(K_\frakp)_\mathfrak P e^{S_r}
\]
прямое разложение на поля. Тогда \((K_\frakp)_\mathfrak P e^{S_i}/Pe^{S_i}\) является расширением Галуа; группа является подгруппой группы разложения одного простого множителя \(\frakp\), а \(e^{S_i}\) сопряжены относительно смежных классов. Из \(\Delta_\lambda^T=\Delta_\lambda\) сначала следует, что \(i\)-я компонента \(\Delta_\lambda\) лежит в \(Pe^{S_i}\), то есть равна, скажем, \(\gamma_i e^{S_i}\) с \(\gamma_i\in P\). Далее из сопряженности \(e^{S_i}\) следует, что все \(\gamma_i\) равны, значит на самом деле \(\Delta_\lambda=\gamma\sum e^{S_i}=\gamma\) лежит в \(P\).} \(\Delta_\lambda\) лежит в основном поле, расширенном вещественным подполем \(\lambda\)-го характера.

\emph{Следовательно, разложение дискриминанта}
\[
        \Delta=\Delta_1^{f_1}\cdots\Delta_t^{f_t}
\]
\emph{является таким разложением в основном поле, расширенном вещественными подполями характеров, причем комплексно-сопряженным характерам соответствуют одни и те же множители. Для остальных сопряженных характеров без дополнительных доводов ничего не следует [ср. 2a)].}

Тем самым теорема 6 доказана во всех частях. Если удается доказать, что идеалы, порожденные \(\Delta_\lambda\), на самом деле уже лежат в основном поле \(k\) и становятся равными для сопряженных характеров, то они совпадают с проводниками Артина, как только составным характерам сопоставляются соответствующие произведения степеней \(\Delta_\lambda\). Действительно, выполняется как функциональное уравнение, так и факт, непосредственно следующий из нормального базиса (ср. Шпайзер, цит. соч.): единичным представлениям подгрупп соответствуют дискриминанты соответствующих подполей. Из совпадения в каждом месте следует совпадение полных идеалов. Если ограничиться рационально неприводимыми характерами, то приведенные факты дают это совпадение. Для абсолютно неприводимых характеров совпадение будет еще прямо подтверждено в самом простом случае:

\paragraph{Теорема 7.}
\emph{Пусть \(k\) есть поле рациональных чисел, а \(K/k\) циклично простой степени \(l\), взаимно простой с дискриминантом, то есть без высшего ветвления. Тогда идеалы, порожденные \(\Delta_\lambda\) для неединичных представлений, попарно равны и совпадают с проводником \(K/k\).}

Это следует из известного разложения корневых чисел.\footnote{Hilbert, Zahlbericht, § 108. Поскольку по теореме 5 существование нормального базиса установлено, а место по примечанию 7 здесь можно просто определить кольцом частных, нет необходимости привлекать использованный в Zahlbericht факт, что абсолютно циклические поля являются круговыми полями.} В разветвленном месте \(\frakp\) расширения \(K/k\), а именно для \(k_\varepsilon\), поля \(l\)-х корней из единицы, и для \(K_\varepsilon\), имеем (\(K_\varepsilon\) остается полем, поскольку \(k_\varepsilon\) взаимно просто с \(K\), а \(\frakp\)-адическое расширение по примечанию 7 здесь излишне):
\[
        (p)=\frakp_1\cdots\frakp_{l-1},
        \qquad
        \frakp_i=\mathfrak P_i^{\,l},
\]
где \(\frakp_i\) лежит в \(k_\varepsilon\), а \(\mathfrak P_i\) в \(K_\varepsilon\). Далее для корневых чисел
\[
        (\Omega_\lambda)=\mathfrak P_1^{r_1}\cdots\mathfrak P_{l-1}^{r_{l-1}},
        \qquad
        (\Omega_{\bar\lambda})=
        \mathfrak P_1^{l-r_1}\cdots \mathfrak P_{l-1}^{l-r_{l-1}},
\]
где \(r_1,\ldots,r_{l-1}\) является перестановкой чисел \(1,\ldots,l-1\). Но корневые числа \(\Omega_\lambda\) здесь тождественны множителям \(D_\lambda\) группового определителя; поэтому
\[
        (\Delta_\lambda)=(D_\lambda D_{\bar\lambda})
        =(\Omega_\lambda\Omega_{\bar\lambda})
        =\mathfrak P_1^{l}\cdots\mathfrak P_{l-1}^{l}
        =\frakp_1\cdots\frakp_{l-1}=(p);
\]
то есть получается совпадение с проводником \(K/k\).

\begin{center}
Поступило 24 августа 1931 г.
\end{center}
% END UNIT BODY
% END PRODUCER UNIT 213
\endgroup


\clearpage
\begingroup
\setcounter{footnote}{0}
% BEGIN PRODUCER UNIT 214
% Source: C:/Users/Floris/Documents/interlanguage/03_projects/noether/02_slavic_working_corpus/translations/paper38/source_fidelity/russian/v001/Noether_Paper38_SourceFidelity_Russian_v001.tex
% Identity: 38172 B / SHA-256 8558D18DEBDD607E3DB7E6A8C420F320D90839802875C9BB83497D3DBD11C2C3
\setlength{\parindent}{1.1em}
\setlength{\parskip}{0pt}
\providecommand{\Om}{\Omega}
\providecommand{\N}{\mathrm{N}}
\newcommand{\foreign}[1]{#1}
\emergencystretch=220em
\allowdisplaybreaks
\sloppy
\hbadness=10000
\vbadness=10000
\hfuzz=5pt
\vfuzz=12pt
\raggedbottom
% BEGIN UNIT BODY
\section*{38. Совместно с Р. Брауэром и Г. Гассе: доказательство главной теоремы в теории алгебр}

\begin{center}
\emph{Journal f. d. reine u. angew. Math. 167 (1932), с. 399--404}
\end{center}

\begin{center}
\textbf{Доказательство главной теоремы в теории алгебр.}\\[0.5em]
Р. Брауэр, Кёнигсберг; Г. Гассе, Марбург; Э. Нётер, Гёттинген\footnote{Составление этой заметки взял на себя Г. Гассе.}.
\end{center}

Наконец нашими совместными усилиями удалось доказать справедливость следующей теоремы, имеющей основное значение для структурной теории алгебр над алгебраическими числовыми полями, а также и за ее пределами:

\paragraph{Главная теорема.}
\emph{Всякая нормальная алгебра с делением над алгебраическим числовым полем циклична (или, как также говорят, имеет тип Диксона).}

Для нас особая радость представить этот результат, как успех, существенно обязанный \(p\)-адическому методу, господину Курту Хензелю, основателю этого метода, к его семидесятилетию.

Наше доказательство состоит из трех редукций; каждый из нас внес одну из них.\footnote{Они приводятся в порядке их возникновения, противоположном систематическому порядку.}

\paragraph{1.}
Первую редукцию дал Г. Гассе на основании недавно развитой им теории циклических алгебр над алгебраическими числовыми полями.\footnote{H. Hasse, \emph{Theorie der zyklischen Algebren über einem algebraischen Zahlkörper}, Gött. Nachr. 1931. Подробное изложение доказательств, где в частности развивается также теория полей разложения и скрещенных произведений, развитая Э. Нётер в лекции и здесь, как и там, основная, вскоре появится в Trans. Amer. Math. Soc. под названием \emph{Theory of cyclic algebras over an algebraic number field}. Ниже эта работа цитируется как H. H, 1--6 составляют, за исключением языка, первую заметку.}

\paragraph{Редукция 1.}
Главная теорема доказана, если показано:
\[
\text{I. Всякая всюду расщепленная алгебра над }\Omega\text{ есть } \sim \Omega .
\]
Для краткости здесь и ниже выражение ``алгебра \(A\) над \(\Omega\)'' всегда означает ``нормальная простая алгебра \(A\) над алгебраическим числовым полем \(\Omega\)'', то есть простая гиперкомплексная система с центром \(\Omega\). Далее, \(A\sim\Omega\) означает, что \(A\) является полной матричной алгеброй над \(\Omega\), то есть принадлежит специальному классу подобия, определенному самим \(\Omega\) как алгеброй с делением над \(\Omega\). Наконец, под всюду расщепленной алгеброй над \(\Omega\) понимается такая алгебра, у которой для каждого простого места \(\mathfrak p\) поля \(\Omega\) \(\mathfrak p\)-адическое расширение, то есть алгебра, возникающая при расширении поля коэффициентов \(\Omega\) до соответствующего \(\mathfrak p\)-адического поля \(\Omega_{\mathfrak p}\), расщепляется:
\[
A_{\mathfrak p}\sim \Omega_{\mathfrak p}.
\]

\paragraph{Доказательство.}
Пусть \(D\) есть алгебра с делением над \(\Omega\). Если \(Z\) есть циклическое поле над \(\Omega\) такое, что для каждого простого места \(\mathfrak p\) поля \(\Omega\) \(\mathfrak p\)-степень поля \(Z\) кратна \(\mathfrak p\)-индексу \(D\), то для соответствующих простых делителей \(\mathfrak P\) в \(Z\) каждое \(\mathfrak P\)-адическое поле \(Z_{\mathfrak P}\) является полем разложения для \(D_{\mathfrak p}\) [H, 18.1]. Отсюда получается, что алгебра \(D_Z\), возникающая из \(D\) расширением поля коэффициентов до \(Z\), всюду расщепляется над \(Z\). Если I уже установлено для всякого \(\Omega\), то следует \(D_Z\sim Z\). Это означает, что \(D\) имеет циклическое поле разложения \(Z\), следовательно допускает циклическое представление [H, 5]. Тогда \(D\) имеет и такое циклическое поле разложения, степень которого совпадает со степенью самой \(D\); то есть \(D\) циклична [H, 6, теорема 6]. Итак, при предположении I главная теорема доказана.

\paragraph{2.}
Для второй редукции решающий толчок дал Р. Брауэр, когда письменно сообщил Г. Гассе, как родственный вопрос о точном значении экспоненты алгебры с делением (который, впрочем, также решается главной теоремой, см. ниже) с помощью силовской теоремы о группах сводится к случаю разрешимого поля разложения. Затем эту мысль удалось использовать и для соответствующей дальнейшей редукции вопроса цикличности.

\paragraph{Редукция 2.}
Утверждение I доказано, если показано:
\[
\text{II. Всякая всюду расщепленная алгебра с разрешимым полем разложения над }\Omega\text{ есть }\sim\Omega .
\]

\paragraph{Доказательство.}
Пусть \(A\) есть всюду расщепленная алгебра над \(\Omega\). Пусть далее \(K\) есть галуа-поле разложения для \(A\), \(p\) произвольное простое число, а \(\Sigma\) одно из силовских полей поля \(K\) над \(\Omega\), относящихся к \(p\) (то есть поле инвариантов силовской \(p\)-подгруппы группы Галуа). Тогда \(A_{\Sigma}\) есть также всюду расщепленная алгебра над \(\Sigma\), имеющая разрешимое поле разложения \(K\). Если II уже установлено для всякого \(\Omega\), то следует \(A_{\Sigma}\sim \Sigma\). Это означает, что \(A\) имеет поле разложения \(\Sigma\), степень которого взаимно проста с \(p\). Поэтому индекс \(A\), как делитель этой степени, взаимно прост с \(p\). Так как это верно для каждого простого \(p\), индекс равен 1, то есть \(A\sim\Omega\). Итак, при предположении II доказано I.\footnote{Идея редукции к разрешимому полю разложения с помощью силовской теоремы была уже раньше применена Р. Брауэром, а именно для доказательства того, что каждый простой делитель индекса встречается также в экспоненте (\emph{Über den Zusammenhang von arithmetischen und invariantentheoretischen Eigenschaften von Gruppen linearer Substitutionen}, Berl. Akad.-Ber. 1926). Недавно А. А. Альберт для этой идеи, а также вообще для ряда общих теорем теории Р. Брауэра и Э. Нётер разработал простые доказательства, независимые от теории представлений: 1. \emph{On direct products, cyclic division algebras, and pure Riemann matrices}; 2. \emph{On direct products}; обе работы в Trans. Amer. Math. Soc. 33 (1931); для рассматриваемой здесь редукции см. особенно теорему 23 в 2.

\emph{Добавление во время печати.} Кроме того, А. А. Альберт, располагая письмом Г. Гассе о том, что главная теорема доказана им для абелевых алгебр (см. далее в тексте), непосредственно и независимо от нас вывел из этого следующие факты:

a) главную теорему для степеней вида \(2^e\),

b) следующую ниже теорему 1 (экспонента = индекс),

c) помимо основной мысли редукции 2, также и мысль последующей редукции 3, естественно без отношения к редукции 1, и соответственно с результатом: для алгебр с делением \(D\) степени \(p^e\), равной степени простого числа, над \(\Omega\) существует расширение \(\Omega'\) степени, взаимно простой с \(p\), над \(\Omega\), такое что \(D_{\Omega'}\) циклична.

Все три результата, разумеется, перекрыты нашим тем временем полученным доказательством главной теоремы. Однако они показывают, что А. А. Альберту также принадлежит независимая доля в доказательстве главной теоремы. Наконец, А. А. Альберт (после знакомства с нашим доказательством главной теоремы) еще заметил, что наша центральная теорема I выводится в несколько строк из теорем 13, 10, 9 его находящейся в печати работы (Bull. Amer. Math. Soc. 37 (1931)). Доказательство этих теорем по существу основано на тех же рассуждениях, что и наши редукции 2 и 3.}

\paragraph{3.}
Третью редукцию, которая затем, как понял Г. Гассе, уже располагая первыми двумя редукциями, ведет к окончательному доказательству, дала Э. Нётер, побужденная сообщением Г. Гассе; в нем главная теорема была доказана для специального случая абелева поля разложения, именно с помощью общей редукции 1 и формульной редукции соответствующей фактор-системы. Ее ядро Э. Нётер и выделила как третью редукцию. Независимо от Э. Нётер эту третью редукцию, впрочем, уже обдумал и Р. Брауэр.

\paragraph{Редукция 3.}
Утверждение II доказано, если показано:
\[
\begin{gathered}\text{III. Всякая всюду расщепленная алгебра с циклическим полем разложения}\\\text{простой степени над }\Omega\text{ есть }\sim\Omega .\end{gathered}
\]

\paragraph{Доказательство.}
Пусть \(A\) есть всюду расщепленная алгебра с разрешимым полем разложения \(K\) над \(\Omega\). Пусть
\[
K=\Lambda_0>\Lambda_1>\cdots>\Lambda_r=\Omega
\]
есть цепочка полей между \(K\) и \(\Omega\), такая что каждое \(\Lambda_i\) над \(\Lambda_{i+1}\) циклично простой степени. Тогда сначала \(A_{\Lambda_1}\) является также всюду расщепленной алгеброй над \(\Lambda_1\), имеющей циклическое поле разложения простой степени \(K=\Lambda_0\). Если III уже установлено для всякого \(\Omega\), то следует \(A_{\Lambda_1}\sim \Lambda_1\). Это означает, что \(\Lambda_1\) есть поле разложения для \(A\). Теперь, исходя из \(\Lambda_1\) вместо \(\Lambda_0\), точно так же доказывают \(A_{\Lambda_2}\sim\Lambda_2\), и так далее, пока наконец \(A_{\Lambda_r}\sim\Lambda_r\), то есть \(A\sim\Omega\). Итак, при предположении III доказано II.

\paragraph{4.}
Но III уже установлено по результатам Г. Гассе [H, 24]. Поэтому, проходя назад через редукции 3, 2, 1, получаем справедливость главной теоремы.

Э. Нётер при этом еще указывает на следующее.
Справедливость III, более общо для циклических полей разложения произвольной степени, сводится к норменной теореме Гассе\footnote{\label{fn:p38-hasse-normensatz}Указано в H, 3.11; доказано в: H. Hasse, \emph{Beweis eines Satzes und Widerlegung einer Vermutung über das allgemeine Normenrestsymbol}, Gött. Nachr. 1931.}. Для редукции 3 нужен, однако, только специальный случай простой степени, то есть первоначальная норменная теорема Гильберта--Фуртвенглера\footnote{См. H. Hasse, \emph{Bericht über neuere Untersuchungen und Probleme in der Theorie der algebraischen Zahlkörper II}, Jahresber. der D. M.-V., Erg.-Bd. 6 (1930), § 8, а также указанную там литературу.}. Тем самым редукция III дает новое простое доказательство норменной теоремы Гассе. Это доказательство шло бы так:

Пусть \(Z\) есть циклическое поле над \(\Omega\), а \(\alpha\) число из \(\Omega\), для которого символ норменного вычета
\[
\left(\frac{\alpha,Z}{\mathfrak p}\right)=1
\]
для каждого простого места \(\mathfrak p\) поля \(\Omega\). Тогда всякая циклическая алгебра
\[
A=(\alpha,Z)
\]
\footnote{В обозначении H, 1. Указание связанного с \(\alpha\) автоморфизма \(S\) поля \(Z\) здесь опущено как несущественное.} всюду расщепляется [H, 17.7]. По редукции 3, следовательно, используя только норменную теорему Гильберта--Фуртвенглера, получаем \(A\sim\Omega\). Но тогда \(\alpha\) есть норма некоторого элемента из \(Z\) [H, 15.4].

В связи с норменной теоремой заметим далее, что теорему I следует рассматривать как ее правильное обобщение на высшие, в том числе неабелевы, случаи, тогда как буквальное обобщение, как показал Г. Гассе\textsuperscript{\ref{fn:p38-hasse-normensatz}}, вообще неверно.

\paragraph{5.}
Из главной теоремы сразу следует ответ на вопрос, к которому первоначально относилась вторая редукционная мысль:

\paragraph{Теорема 1.}
Экспонента нормальной простой алгебры над алгебраическим числовым полем равна ее индексу.

Далее из теоремы I получается примечательный факт:

\paragraph{Теорема 2.}
Основной идеал собственной нормальной алгебры с делением над \(\Omega\) делится по меньшей мере на одно простое место \(\Omega\), и даже по меньшей мере на два.

При этом основной идеал можно определить как (приведенную) норму (приведенной) дифференты.\footnote{В специальном случае рациональных кватернионных алгебр это сводится к введенному Г. Брандтом понятию ``основного числа'' (\emph{Idealtheorie in Quaternionenalgebren}, Math. Ann. 99 (1928)). Так как здесь речь идет не о числе, а об идеале, надо было сказать ``основной идеал''; это не следует смешивать с дедекундовым обозначением ``основной идеал'' и ``основное число'' для дифференты и дискриминанта, которое именно по указанной причине оказывается непригодным для относительных полей. О дефиниции (приведенной) дифференты см. следующую сноску. (Приведенную) норму идеала Э. Нётер также определяет ``по простым местам'', а именно, в соответствии с тем фактом, что для отдельных простых мест всякий идеал главный, просто образованием (приведенных) числовых норм чисел базиса главного идеала для отдельных простых мест.} Кроме того, бесконечное простое место включается в основной идеал тогда и только тогда, когда алгебра для него сводится к кватернионной алгебре, а не к действительному или комплексному числовому полю.

\paragraph{Доказательство.}
По результатам Г. Гассе\footnote{\label{fn:p38-hasse-padisch}H. Hasse, \emph{Über \(p\)-adische Schiefkörper und ihre Bedeutung für die Arithmetik hyperkomplexer Zahlsysteme}, Math. Ann. 104 (1931); см. там теоремы 42, 59.} простое место \(\mathfrak p\) поля \(\Omega\) входит в основной идеал тогда и только тогда, когда \(\mathfrak p\)-индекс отличен от \(1\). Если бы все \(\mathfrak p\)-индексы были равны \(1\), то алгебра всюду расщеплялась бы, и по теореме I не было бы собственной алгебры с делением. (То, что и один-единственный \(\mathfrak p\)-индекс не может быть отличен от \(1\), следует из того, что для символов норменного вычета, порядки которых суть \(\mathfrak p\)-индексы [H, 17.7], действует теорема о произведении, то есть закон взаимности [H, 3.8].)

\paragraph{6.}
Кроме того, теорема I позволяет определить (арифметически охарактеризовать) все поля разложения \(K\), принадлежащие алгебре \(A\) над \(\Omega\). В точном обобщении обстоятельства, уже установленного Г. Гассе в специальном случае циклических \(A,K\) [H, 6, теорема 2], получается:

\paragraph{Теорема 3.}
Для алгебры \(A\) над \(\Omega\) алгебраическое числовое поле \(K\) над \(\Omega\) является полем разложения тогда и только тогда, когда для всех простых делителей \(\mathfrak P_i\) в \(K\) всех простых мест \(\mathfrak p\) поля \(\Omega\) соответствующая \(\mathfrak P_i\)-степень \(n_{\mathfrak P_i}\) поля \(K\) кратна \(\mathfrak p\)-индексу \(m_{\mathfrak p}\) алгебры \(A\).

При этом \(\mathfrak P_i\)-степень поля \(K\) есть степень соответствующего \(\mathfrak P_i\)-адического поля \(K_{\mathfrak P_i}\) над \(\mathfrak p\)-адическим полем \(\Omega_{\mathfrak p}\), то есть, если
\[
\mathfrak p=\prod_i \mathfrak P_i^{\,e_{\mathfrak P_i}},
        \qquad N_{K/\Omega}(\mathfrak P_i)=\mathfrak p^{\,f_{\mathfrak P_i}},
\]
есть разложение \(\mathfrak p\) в \(K\), это произведение степени \(f_{\mathfrak P_i}\) и порядка ветвления \(e_{\mathfrak P_i}\) делителя \(\mathfrak P_i\) относительно \(\mathfrak p\), \(n_{\mathfrak P_i}=f_{\mathfrak P_i}e_{\mathfrak P_i}\) [H, 4]. А \(\mathfrak p\)-индекс \(A\) есть индекс \(m_{\mathfrak p}\) \(\mathfrak p\)-адического расширения \(A_{\mathfrak p}\) [H, 6].

\paragraph{Доказательство.}
То, что \(K\) есть поле разложения для \(A\), то есть \(A_K\sim K\), по теореме I равносильно тому, что \(A_{K_{\mathfrak P_i}}\sim K_{\mathfrak P_i}\) для каждого \(\mathfrak P_i\).

Для конечного простого места \(\mathfrak p\) пусть
\[
        A_{\mathfrak p}\sim(\pi,W_{\mathfrak p})
\]
есть арифметически выделенное циклическое представление \(A_{\mathfrak p}\) [H, 16; кроме того, l. c.\textsuperscript{\ref{fn:p38-hasse-padisch}}, теорема 38], то есть \(W_{\mathfrak p}\) есть неразветвленное поле степени \(m_{\mathfrak p}\) над \(\Omega_{\mathfrak p}\), а \(\pi\) число из \(\Omega_{\mathfrak p}\), точно делящееся на \(\mathfrak p^1\). Далее пусть \(K_{\mathfrak P_i}W_{\mathfrak p}\) есть композит, а
\[
        \Delta_{\mathfrak P_i}=K_{\mathfrak P_i}\cap W_{\mathfrak p}
\]
пересечение \(W_{\mathfrak p}\) и \(K_{\mathfrak P_i}\). Так как наибольшее неразветвленное подполе, содержащееся в \(K_{\mathfrak P_i}\), имеет степень \(f_{\mathfrak P_i}\) над \(\Omega_{\mathfrak p}\), и так как над \(\Omega_{\mathfrak p}\) для каждой степени существует только одно неразветвленное поле, то \(\Delta_{\mathfrak P_i}\)
\[
        \text{имеет степень }d_{\mathfrak P_i}=(m_{\mathfrak p},f_{\mathfrak P_i})\text{ над }\Omega_{\mathfrak p}.
\]
Поэтому \(K_{\mathfrak P_i}W_{\mathfrak p}\) имеет степень \(m_{\mathfrak p}/d_{\mathfrak P_i}\) над \(K_{\mathfrak P_i}\), причем неразветвлено.

Теперь
\[
        A_{K_{\mathfrak P_i}}=(A_{\mathfrak p})_{K_{\mathfrak P_i}}\sim(\pi,K_{\mathfrak P_i}W_{\mathfrak p})
\]
[H, 15.5]. Поэтому \(A_{K_{\mathfrak P_i}}\sim K_{\mathfrak P_i}\) равносильно тому, что \(\pi\) является нормой из \(K_{\mathfrak P_i}W_{\mathfrak p}\) относительно \(K_{\mathfrak P_i}\). Но ввиду неразветвленности \(K_{\mathfrak P_i}W_{\mathfrak p}\) над \(K_{\mathfrak P_i}\) это имеет место тогда и только тогда, когда показатель \(e_{\mathfrak P_i}\) числа \(\pi\) в \(\mathfrak P_i\) кратен степени \(m_{\mathfrak p}/d_{\mathfrak P_i}\) поля \(K_{\mathfrak P_i}W_{\mathfrak p}\) над \(K_{\mathfrak P_i}\), или, что ввиду
\[
\left(\frac{m_{\mathfrak p}}{d_{\mathfrak P_i}},\frac{f_{\mathfrak P_i}}{d_{\mathfrak P_i}}\right)=1
\]
сводится к тому же, когда
\[
\frac{f_{\mathfrak P_i}}{d_{\mathfrak P_i}}e_{\mathfrak P_i}=\frac{n_{\mathfrak P_i}}{d_{\mathfrak P_i}}
\]
кратно \(m_{\mathfrak p}/d_{\mathfrak P_i}\), то есть когда \(n_{\mathfrak P_i}\) кратно \(m_{\mathfrak p}\), как и утверждалось.

Если \(\mathfrak p\) есть бесконечное простое место и не имеет места тривиальный случай \(m_{\mathfrak p}=1\), то \(m_{\mathfrak p}=2\), а \(A_{\mathfrak p}\) есть кватернионная алгебра над действительным числовым полем \(\Omega_{\mathfrak p}\). То, что \(A_{K_{\mathfrak P_i}}=(A_{\mathfrak p})_{K_{\mathfrak P_i}}\sim K_{\mathfrak P_i}\), тогда равносильно тому, что \(K_{\mathfrak P_i}\) есть комплексное числовое поле, то есть \(n_{\mathfrak P_i}=f_{\mathfrak P_i}=2=m_{\mathfrak p}\) (здесь \(e_{\mathfrak P_i}\) не появляется), что снова дает утверждение (\(n_{\mathfrak P_i}\), как и \(m_{\mathfrak p}\), может принимать только значения 1 или 2).

\paragraph{7.}
К значительным результатам приходят, если вопрос, решенный в теореме 3, обо всех полях разложения \(K\) фиксированной алгебры \(A\) обратить: при фиксированном алгебраическом поле \(K\) над \(\Omega\) рассмотреть совокупность алгебр \(A\) над \(\Omega\), расщепляемых полем \(K\). Эти алгебры \(A\) образуют определенную \(K\) подгруппу \(\mathfrak K\) группы Р. Брауэра \(\mathfrak A\) всех алгебр (точнее, всех классов подобных алгебр) над \(\Omega\).\footnote{R. Brauer, \emph{Über Systeme hyperkomplexer Größen}, Jahresber. d. D. M.-V. 38 (1929), с. 47/48. См. также H, 13.1.} Если предположить, что \(K\) галуа над \(\Omega\), то приходят к теоремам, которые следует рассматривать как обобщение главных теорем теории классов полей (теории относительно-абелевых числовых полей) на общие относительно-галуа числовые поля:

\paragraph{Теорема 4.}
\emph{(Теорема разложения).} Относительная степень \(f\) простых делителей в \(K\) простого идеала \(\mathfrak p\) поля \(\Omega\), не входящего в относительный дискриминант \(K\), равна наименьшей экспоненте, для которой
\[
        A_{\mathfrak p}^{\,f}\sim\Omega_{\mathfrak p}
\]
становится верным для всех алгебр \(A\) из группы \(\mathfrak K\), сопоставленной \(K\).

\paragraph{Доказательство.}
Так как \(f\) есть \(\mathfrak p\)-степень поля \(K\), одинаковая для всех простых делителей \(\mathfrak p\) в \(K\), а с другой стороны \(\mathfrak p\)-индекс \(A\) равен индексу, то есть экспоненте, \(A_{\mathfrak p}\), то по теореме 3 \(f\) во всяком случае кратно указанной наименьшей экспоненте; для доказательства остается показать, что в группе \(\mathfrak K\) существуют алгебры \(A\) с точным \(\mathfrak p\)-индексом \(f\).

По теореме Фробениуса о плотности наверняка существует еще один простой идеал \(\mathfrak p'\) в \(\Omega\), не входящий в относительный дискриминант \(K\), простые делители которого в \(K\) имеют относительную степень \(f\). Пусть \(Z\) есть циклическое поле над \(\Omega\), у которого \(\mathfrak p\)-степень и \(\mathfrak p'\)-степень имеют значение \(f\). Далее пусть \(\alpha\) есть число в \(\Omega\), для которого символы норменного вычета
\[
\left(\frac{\alpha,Z}{\mathfrak p}\right)\quad\text{и}\quad
\left(\frac{\alpha,Z}{\mathfrak p'}\right)
\]
имеют взаимно обратные значения порядка \(f\), тогда как для всех остальных делителей \(\mathfrak q\) проводника \(Z\)
\[
\left(\frac{\alpha,Z}{\mathfrak q}\right)=1.
\]
По обобщенной теореме об арифметической прогрессии \(\alpha\) можно к тому же выбрать так, чтобы, кроме, возможно, \(\mathfrak p,\mathfrak p'\) и \(\mathfrak q\), оно содержало еще только один простой идеал \(\mathfrak r\) поля \(\Omega\), причем точно в первой степени. Для него тогда по теореме о произведении для символа норменного вычета (закону взаимности) также
\[
\left(\frac{\alpha,Z}{\mathfrak r}\right)=1
\]
[H, 3.8--10]. Всякая алгебра
\[
        (\alpha,Z)=A
\]
тогда имеет \(\mathfrak p\)-индекс и \(\mathfrak p'\)-индекс \(f\), а для всех других простых мест \(\Omega\) индекс \(1\) [H, 17.7]. Из теоремы 3, ввиду предположения о \(\mathfrak p\) и выбора \(\mathfrak p'\), следует, что \(K\) есть поле разложения для \(A\). Тем самым в группе \(\mathfrak K\) действительно обнаружены алгебры \(A\) с \(\mathfrak p\)-индексом \(f\).

\paragraph{Теорема 5.}
\emph{(Теорема единственности и упорядочения).} Если \(K\leqq K'\), то \(\mathfrak K\leqq\mathfrak K'\), и обратно.

В частности, сопоставление групп алгебр \(\mathfrak K\) галуа-полям \(K\) взаимно однозначно.

\paragraph{Доказательство.}
a) Из \(K\leqq K'\) тривиально следует \(\mathfrak K\leqq\mathfrak K'\), ибо всякая алгебра \(A\), расщепленная \(K\), a fortiori является алгеброй, расщепленной \(K'\).

b) Пусть обратно \(\mathfrak K\leqq\mathfrak K'\). Тогда по теореме разложения для каждого неделителя \(\mathfrak p\) относительного дискриминанта имеем \(f_{\mathfrak p}\mid f'_{\mathfrak p}\). В частности, \(f_{\mathfrak p}=1\) для почти всех \(\mathfrak p\) с \(f'_{\mathfrak p}=1\). Отсюда по обычному аналитическому способу вывода\footnote{См. об этом, например, H. Hasse [l. c. прим. 6], § 25, III.} следует \(K\leqq K'\).

\paragraph{8.}
Наконец, отметим еще, что главная теорема существенно продвигает вопрос, рассмотренный И. Шуром\footnote{I. Schur, \emph{Arithmetische Untersuchungen über endliche Gruppen linearer Substitutionen}, Berl. Akad.-Ber. 1906.}, о числовых полях, в которых возможны абсолютно неприводимые представления конечной группы:

\paragraph{Теорема 6.}
Все абсолютно неприводимые представления конечной группы \(\mathfrak G\) возможны в круговых полях; например, всегда во всяком случае в поле \(n^h\)-х корней из единицы, если \(n\) есть порядок \(\mathfrak G\), а \(h\) достаточно велико.

\paragraph{Доказательство.}
Переходя к рациональному групповому кольцу \(G\) группы \(\mathfrak G\), абсолютно неприводимые представления \(\Gamma_i\) группы \(\mathfrak G\) становятся абсолютно неприводимыми представлениями простых составляющих \(G_i\) полупростой алгебры \(G\), а их центры суть соответствующие поля \(\Omega_i\) характеров,\footnote{См. об этом: a) R. Brauer und E. Noether, \emph{Über minimale Zerfällungskörper irreduzibler Darstellungen}, Berl. Akad.-Ber. 1927; § 1. b) R. Brauer, \emph{Über Systeme hyperkomplexer Zahlen}, Math. Zeitschr. 30 (1929); теорема 3. c) E. Noether, \emph{Hyperkomplexe Größen und Darstellungstheorie}, Math. Zeitschr. 30 (1929); §§ 21, 24, 26.} следовательно, из-за конечности \(\mathfrak G\), во всяком случае круговые поля, причем заведомо подполя поля \(n\)-х корней из единицы.

По теореме 3 циклическое поле \(Z_i\) над \(\Omega_i\) является полем разложения для \(G_i\), если для каждого простого места \(\mathfrak p\) поля \(\Omega_i\) его \(\mathfrak p\)-степень \(n_{i\mathfrak p}\) кратна \(\mathfrak p\)-индексу \(m_{i\mathfrak p}\) алгебры \(G_i\). \(m_{i\mathfrak p}\) отличен от 1 только для простых делителей основного идеала \(G_i\) относительно \(\Omega_i\), значит заведомо только для простых делителей абсолютного дискриминанта \(G_i\); поскольку тот делит \(n^n\), а \(n^n\) есть дискриминант одного (немаксимального) порядка в \(G\)\footnote{См. E. Noether [l. c. прим. 12c], § 26.}, это возможно самое большее для простых делителей \(\mathfrak p\) числа \(n\). Чтобы сделать \(n_{i\mathfrak p}\) для этих \(\mathfrak p\) кратным \(m_{i\mathfrak p}\), достаточно предписать, чтобы рассматриваемые \(\mathfrak p\) в \(Z_i\) разветвлялись с порядком, каждый раз кратным \(m_{i\mathfrak p}\). Как легко видеть, это обеспечивает поле \(n^h\)-х корней из единицы при достаточно большом \(h\).

В последней теореме указанной работы И. Шур устанавливает, что во всех до сих пор известных случаях уже достаточно поля \(n\)-х корней из единицы. Всегда ли это так и достаточны ли развитые здесь методы, чтобы решить этот вопрос, остается предметом дальнейших исследований.

\begin{center}
Поступило 11 ноября 1931 г.
\end{center}
% END UNIT BODY
% END PRODUCER UNIT 214
\endgroup


\clearpage
\begingroup
\setcounter{footnote}{0}
% BEGIN PRODUCER UNIT 215
% Source: C:/Users/Floris/Documents/interlanguage/03_projects/noether/02_slavic_working_corpus/translations/paper39/source_fidelity/russian/v001/Noether_Paper39_SourceFidelity_Russian_v001.tex
% Identity: 27956 B / SHA-256 78E0D078CA7CB4BC6E504D7D212AA9A076B0493284A0E647BD6E018B52AAC7C6
\setlength{\parindent}{1.1em}
\setlength{\parskip}{0pt}
\newcommand{\foreign}[1]{#1}
\emergencystretch=220em
\allowdisplaybreaks
\sloppy
\hbadness=10000
\vbadness=10000
\hfuzz=5pt
\vfuzz=12pt
\raggedbottom
% BEGIN UNIT BODY
\section*{39. Гиперкомплексные системы в их связях с коммутативной алгеброй и теорией чисел}

\begin{center}
\emph{\foreign{Verhandl. Intern. Math.-Kongreß Zürich 1 (1932), с. 189--194}}
\end{center}

\begin{center}
\textsc{Эмми Нётер, Гёттинген}
\end{center}

\paragraph{1.}
Теория гиперкомплексных систем, то есть алгебр, за последние годы сильно оживилась; но только совсем недавно стало ясно значение этой теории для коммутативных вопросов. Об этом значении некоммутативного для коммутативного я сегодня хочу сообщить, а именно проследить его подробно на двух классических вопросах, восходящих к Гауссу: на теореме о главном роде и тесно связанной с ней норменной теореме. Эти вопросы с течением времени снова и снова меняли свою формулировку: у Гаусса они выступают как завершение его теории квадратичных форм; затем играют существенную роль в характеристике относительно циклических и абелевых числовых полей посредством теории полей классов; наконец, их можно высказать как теоремы об автоморфизмах и о расщеплении алгебр, и эта последняя формулировка вместе с тем дает перенос теорем на произвольные относительно галуа-числовые поля.

Этим очерком, который я позднее разверну, я одновременно хочу пояснить \emph{принцип} применения некоммутативного к коммутативному: \emph{посредством теории алгебр ищут инвариантные и простые формулировки для известных фактов о квадратичных формах или циклических полях, то есть такие формулировки, которые зависят только от структурных свойств алгебр. Если эти инвариантные формулировки уже доказаны -- а именно так обстоит дело в указанных выше примерах --, то тем самым автоматически получен перенос этих фактов на произвольные галуа-поля.}

\paragraph{2.}
Прежде чем перейти к подробному изложению, я хочу дать еще общий обзор различных методов и дальнейших результатов. Прежде всего нужно отметить, что главная трудность состоит в нахождении формулировки для общих галуа-полей; без гиперкомплексного метода для этого вообще не было исходной точки. В приведенных примерах лежащий в основе переход к некоммутативному получен через \emph{одновременное рассмотрение поля и группы} посредством ``скрещенного произведения'' и его мультипликативных постоянных, ``фактор-систем'' (см. 3). Так получают нормальную простую алгебру над основным полем, и всякая такая алгебра, по существу, порождается именно так. Такие скрещенные произведения впервые рассматривал Диксон,\footnote{См. его книгу \emph{Algebren und ihre Zahlentheorie}, Zürich 1927, § 34.} тогда как теория фактор-систем развивалась у Шпайзера, Шура и Р. Брауэра\footnote{См., например, R. Brauer, \emph{Untersuchungen über die arithmetischen Eigenschaften von Gruppen linearer Substitutionen}, и приведенную там в прим. 2 литературу, Math. Z. 28 (1928).} из совсем другой исходной точки, а именно из вопроса об абсолютно неприводимых представлениях. Только слияние обеих теорий дало достаточно простое и широкое построение, чтобы с его помощью можно было подступиться к коммутативным вопросам.\footnote{Это построение я впервые развила в лекции зимой 1929/30 г.; оно воспроизведено в гл. 2 работы H. Hasse, \emph{Theory of cyclic algebras}, Trans. Amer. Math. Soc. 34 (1932). Во всей области, рассматриваемой в докладе, ориентирует обзор М. Дойринга о гиперкомплексных числах и числово-теоретических приложениях, который должен выйти в серии \emph{Ergebnisse der Mathematik}.}

При этом для известных фактов одновременно возникают новые прозрачные доказательства. Я хочу указать здесь на гиперкомплексное доказательство закона взаимности для циклических полей, которое Г. Гассе дал и которое вскоре появится в Math. Ann.,\footnote{H. Hasse, \emph{Die Struktur der R. Brauerschen Algebrenklassengruppe über einem algebraischen Zahlkörper (insbesondere Normenrestsymbol und Reziprozitätsgesetz)}. Math. Ann. 107 (1932/33).} посредством инвариантной формулировки своего символа норменного вычета, основанной на теории скрещенного произведения. Далее -- на гиперкомплексное обоснование локальной теории полей классов, тоже на той же основе, которое недавно дал К. Шевалле; при этом, однако, пришлось разработать еще новые алгебраические теоремы о фактор-системах.\footnote{C. Chevalley, \emph{Sur la théorie du symbole de restes normiques}; должно появиться в J. f. Math. 169.} Вместе с тем я должна сделать ограничение: метод скрещенных произведений сам по себе, по-видимому, \emph{не дает полной теории галуа-числовых полей}. Это следует из новых, еще неопубликованных результатов Артина, которые продолжают приведенное выше доказательство Гассе в духе указанного принципа, но дают только равенства чисел вместо полных теорем об изоморфизме.

Методы, дающие полный изоморфизм -- правда, операторный изоморфизм, -- теперь уже имеются в самой алгебре. Речь идет о продолжении подходов А. Шпайзера,\footnote{A. Speiser, \emph{Gruppendeterminante und Körperdiskriminante}. Math. Ann. 77 (1916).} а именно о понимании галуа-поля как ``галуа-модуля'', то есть как модуля над основным полем, допускающего подстановки галуа-группы в качестве операторов. И существует операторный изоморфизм между полем и групповым кольцом (групповой алгеброй) в том смысле, что между элементами имеется взаимно однозначное соответствие так, что линейные формы над основным полем соответствуют друг другу, а подстановкам галуа-группы в поле соответствует умножение в групповом кольце. Эту сформулированную мной теорему\footnote{E. Noether, \emph{Normalbasis bei Körpern ohne höhere Verzweigung}, Satz 3, J. f. Math. 167 (1932) (доказательство содержит пробел).} доказал М. Дойринг,\footnote{M. Deuring, \emph{Galoissche Theorie und Darstellungstheorie}. Math. Ann. 107 (1932).} и на ней он основал доказательство теории Галуа, где операторный изоморфизм реализует соответствие между группой и полем. Дальнейшие структурные теоремы -- также Дойринга -- параллельны формальным фактам об \(L\)-рядах Артина и дают структурный подход к проводникам Артина. Эти \(L\)-ряды и проводники Артина,\footnote{E. Artin, \emph{Über eine neue Art von \(L\)-Reihen}, Math. Sem. Hamburg 3 (1924). \emph{Zur Theorie der \(L\)-Reihen mit allgemeinen Gruppencharakteren}, там же 8 (1931). \emph{Die gruppentheoretische Struktur der Diskriminanten algebraischer Zahlkörper}, J. f. M. 164 (1931).} образованные с общими характерами групп, наряду со Шпайзером\textsuperscript{6)} составляют первую связь между теорией чисел и теорией представлений, первый выход за пределы абелевых полей. Они дали всему развитию сильный импульс; в частности, теория галуа-модулей ориентировалась именно на них.

\paragraph{3.}
Теперь я хочу подробно проследить поставленные в начале вопросы: норменную теорему и теорему о главном роде. Сначала -- \emph{определение скрещенного произведения}. Пусть \(K/k\) есть галуа-поле \(n\)-й степени, а \(\mathfrak G\) -- его группа. Скрещенное произведение означает одновременное вложение \(K\) и \(\mathfrak G\) в алгебру \(A\) так, что автоморфизмы \(K\) становятся внутренними. Пусть \(u_{S_1},\ldots,u_{S_n}\) обозначают символы, соответствующие \(n\) элементам группы; тогда сначала задают \(A\) как модуль линейных форм ранга \(n\) над \(K\):
\begin{equation}
\tag{1}
\begin{gathered}
A=u_{S_1}K+\cdots+u_{S_n}K\\
\text{\normalfont(то есть \(A\) состоит из всех линейных форм}\\
\text{\normalfont\(u_{S_1}a_1+\cdots+u_{S_n}a_n\), где \(a_i\) произвольны в \(K\)).}
\end{gathered}
\end{equation}
В силу требования внутреннего автоморфизма -- порожденного \(u_S\), а в общем виде \(u_SK^\ast\)\footnote{\(K^\ast\) получается из \(K\) удалением нуля; это обозначение употребляется вообще.} -- \(A\) становится кольцом, то есть алгеброй ранга \(n^2\) над \(k\). Это требование выражается равенствами
\begin{equation}
\tag{2}
u_S^{-1}zu_S=z^S,\qquad \text{или}\qquad zu_S=u_Sz^S
\end{equation}
для всякого \(z\in K\),\footnote{\(z^S\), как обычно, означает элемент, получающийся из \(z\) подстановкой \(S\).} и равенством
\begin{equation}
\tag{3}
u_Su_T=u_{ST}a_{S,T}\quad \text{с }a_{S,T}\text{ в }K^\ast .
\end{equation}
Наконец,
\begin{equation}
\tag{4}
\begin{gathered}
a_{ST,R}\,a_{S,T}^{\,R}=a_{S,TR}\,a_{T,R}\\
\text{\normalfont(ассоциативный закон из \([u_Su_T]u_R=u_S[u_Tu_R]\)).}
\end{gathered}
\end{equation}
\(A\) называется скрещенным произведением \(K\) с \(\mathfrak G\) при фактор-системе \(a_{S,T}\). Доказывается, что \(A\) становится нормальной простой алгеброй над \(k\), то есть матричным кольцом \(r\)-й степени \(D_r\) над соответствующей алгеброй с делением \(D\), и что \(K\) становится максимальным коммутативным подполем, следовательно полем разложения (то есть расширение поля коэффициентов \(k\) полем, изоморфным \(K\), дает расщепленную алгебру, полное матричное кольцо над центром). Обратно, для заданной алгебры с делением \(D\) всегда существуют матричные кольца \(D_r\), которые могут быть порождены указанным способом как скрещенное произведение.

Если перейти от \(u_S\) к \(v_S=u_Sc_S\), где \(c_S\in K^\ast\), что порождает тот же автоморфизм, то возникают ``ассоциированные'' фактор-системы:
\begin{equation}
\tag{5}
a'_{S,T}=a_{S,T}\frac{c_S^{\,T}c_T}{c_{ST}}.
\end{equation}
Ассоциированные фактор-системы объединяются в класс \((a)\), а все алгебры, подобные \(A\), -- в алгебровый класс \(\mathfrak A\). Эти классы соответствуют друг другу взаимно однозначно. Классы \(\mathfrak A\) с фиксированным полем разложения \(K\) образуют относительно прямого произведения абелеву группу, изоморфную покомпонентному произведению классов фактор-систем. Единичный элемент есть класс расщепленных алгебр, соответственно фактор-системы из всех величин преобразования \(c_S^{\,T}c_T/c_{ST}\).

\paragraph{4.}
Теперь через \emph{специализацию} к \emph{циклическим полям разложения} я хочу перейти к связи с \emph{понятием нормы}, а тем самым -- к формулировке обобщенной \emph{норменной теоремы} согласно принципу, изложенному в начале. Если \(Z\) циклично, а \(S\) -- порождающая подстановка его группы, то соответствующая алгебра называется циклической. Тогда степеням \(S\) можно поставить в соответствие степени \(u\), то есть
\[
A=Z+uZ+\cdots+u^{n-1}Z,
\]
\[
zu=uz^S,\qquad u^n=a,\qquad a\in k^\ast,
\]
и при \(v=uc\)
\[
a'=a\,N(c).
\]
Следовательно, всякая фактор-система состоит из одного-единственного элемента \(a\in k^\ast\), и пишут
\[
A=(a,Z).
\]
Единичный класс фактор-систем задается нормами из \(Z^\ast\), а группа классов алгебр изоморфна \(k^\ast/N(Z^\ast)\). Поэтому циклическая алгебра \((a,Z)\) расщепляется тогда и только тогда, когда \(a\) есть норма некоторого элемента из \(Z\).

Эта связь между нормой и расщеплением дает формулировку ``норменной теоремы'', а именно \emph{теоремы о расщепленных алгебрах: если алгебра расщепляется в каждом месте, то она расщепляется безусловно}. Здесь ``место'', как обычно в теории чисел, определяется тем, что основное поле \(k\) заменяется его \(p\)-адическим расширением \(k_{\mathfrak p}\), где \(\mathfrak p\) есть простой идеал из \(k\); соответственно конечному числу бесконечных мест, \(k\) и его сопряженные расширяются полем вещественных чисел.

Действительно, в этом содержится норменная теорема для циклических полей. Ибо для циклических алгебр \((a,Z)\), по сказанному выше, эта теорема равносильна утверждению: если \(a\) в каждом (конечном или бесконечном) месте является \(p\)-адической нормой, то \(a\) есть норма числа из \(Z\); или, без перехода к \(p\)-адическому, \emph{если \(a\) есть норменный вычет по каждому простому идеалу \(\mathfrak p\) из \(k\) (и удовлетворяет некоторым условиям знаков), то \(a\) есть норма числа из \(Z\).} Последняя форма и есть норменная теорема, доказываемая в теории полей классов с использованием известных аналитических средств. А доказательство общей теоремы о расщепленных алгебрах можно получить из этого циклического специального случая чисто алгебраико-арифметическими рассуждениями.\footnote{R. Brauer, H. Hasse, E. Noether, \emph{Beweis eines Hauptsatzes in der Theorie der Algebren}. J. f. M. 167 (1932).} Первое важное следствие отсюда вывел Гассе: \emph{всякая нормальная простая алгебра над алгебраическим числовым полем циклична.} В поиске доказательства этого давно предполагаемого факта и возникла общая формулировка.

\paragraph{5.}
Второе следствие из теоремы о расщепленных алгебрах -- опять чисто алгебраико-арифметическими рассуждениями -- это поставленная в начале теорема о главном роде.\footnote{Доказательство должно появиться в Math. Ann.} Ее инвариантная формулировка основана на том факте, что соотношения (2)--(5), определяющие скрещенное произведение, чисто мультипликативны и поэтому сохраняют смысл, если \(K^\ast\) заменить абелевой группой \(\mathfrak J\), удовлетворяющей только условию: ее группа автоморфизмов содержит подгруппу, изоморфную \(\mathfrak G\). Вместо (1) тогда возникает ``расширение \(\mathfrak G\) посредством \(\mathfrak J\)'' в смысле теории групп. Если взять за \(\mathfrak J\) группу всех идеалов из \(K\), то фактор-системы становятся системами идеалов; разбиение на классы в \(\mathfrak J\) индуцирует разбиение на классы для фактор-систем, и вообще это более тонкое разбиение на идеальные классы, чем исходное. Ведь требование, чтобы умножение \(u_S\) с (абсолютными) идеальными классами было однозначным, говорит только, что величины преобразования \(c_S^{\,T}c_T/c_{ST}\) -- единичный класс элементных фактор-систем -- лежат в единичном классе идеальных фактор-систем. В действительности же, как показывает специализация к известным случаям, достаточно уже несколько менее тонкого разбиения. Я определяю: \emph{к единичному классу фактор-систем относятся те элементы \(a_{S,T}\), которые во всех (конечных и бесконечных) местах ветвления \(K\) порождают расщепленные алгебры.}

Возникающее таким образом расширение \(\mathfrak G\) обозначим \(\mathfrak G^\ast\); пусть \((c)\) означает абсолютный идеальный класс \(c\). Тогда

\emph{инвариантная формулировка теоремы о главном роде} такова: \emph{если при определенном так разбиении на классы подстановка \(v_S=u_S(c_S)\) порождает автоморфизм \(\mathfrak G^\ast\) -- все эти \((c_S)\) образуют главный род --, то этот автоморфизм внутренний и порождается некоторым идеальным классом \((\mathfrak b)\).}

Известные специальные случаи следуют из равносильной, несколько более явной формулировки: \emph{если величины преобразования \((c_S^{\,T})(c_T)/(c_{ST})\), образованные из идеальных классов \((c_S)\), принадлежат единичному идеальному классу фактор-систем, то существует идеальный класс \((\mathfrak b)\) такой, что \((c_S)\) становятся символическими \((1-S)\)-ми степенями: \((c_S)=(\mathfrak b)/(\mathfrak b^{S})\) для всех \(S\) из \(\mathfrak G\).} Ибо предположение как раз выражает свойство автоморфизма; а то, что он внутренний, выражается равенством \(v_S=(\mathfrak b)^{-1}u_S(\mathfrak b)=u_S(\mathfrak b)^{1-S}=u_S(c_S)\).

Специализация к циклическому случаю дает, с учетом нормировки: если \(N([c])\) лежит в единичном идеальном классе фактор-систем, то \((c)\) становится символической \((1-S)\)-й степенью:
\[
(c)=(\mathfrak b)^{1-S}.
\]

\paragraph{6.}
Чтобы отсюда перейти к известной форме для циклических полей и квадратичных форм, прежде всего замечу, что теорема высказана для полных идеальных классов, но ее содержание не меняется, если, как обычно, ограничиться идеалами, взаимно простыми с местами ветвления.

Тогда определенный здесь главный род для квадратичных полей переходит в гауссов. Действительно, идеальным классам соответствуют квадратичные формы, а нормам классов -- числа, представимые этими формами. То, что единичный класс порождает алгебры, расщепленные в местах ветвления \(K\), означает, что эти представимые числа становятся квадратичными вычетами в местах ветвления; следовательно, соответствующие формы имеют полный характер главной формы, то есть образуют гауссов главный род. Известно, что символическая \((1-S)\)-я степень переходит в дупликацию.

Для циклических полей формулировка переходит в следующую: главный род состоит из всех идеальных классов, нормы которых в (конечных и бесконечных) местах ветвления становятся норменными вычетами. Но это, как известно, равносильно условию ``норменный вычет по проводнику'', и тем самым обычная теорема возникает как специализация.

Кстати, и в общем случае произвольных галуа-полей можно ввести проводник, составленный только из мест ветвления, так, что при нормировке фактор-систем для каждого из этих мест луч содержит только элементы единичного класса.

Здесь возникает вопрос о связи с проводниками Артина, упомянутыми в обзоре 2, ведь они составлены из тех же простых идеалов; а вместе с этим -- вопрос о связи с теорией галуа-модулей, вторым гиперкомплексным методом. Насколько далеко будут простираться эти два метода, покажет только будущее.
% END UNIT BODY
% END PRODUCER UNIT 215
\endgroup


\clearpage
\begingroup
\setcounter{footnote}{0}
% BEGIN PRODUCER UNIT 216
% Source: C:/Users/Floris/Documents/interlanguage/03_projects/noether/02_slavic_working_corpus/translations/paper40/source_fidelity/russian/v001/Noether_Paper40_SourceFidelity_Russian_v001.tex
% Identity: 122913 B / SHA-256 68B6768B77A1EDFA28D8810FF498D77ED930D762646EA04ECD21BCEA6E102661
\setlength{\parindent}{1.1em}
\setlength{\parskip}{0pt}
\newcommand{\foreign}[1]{#1}
\providecommand{\frA}{\mathfrak A}
\providecommand{\frB}{\mathfrak B}
\providecommand{\frC}{\mathfrak C}
\providecommand{\frG}{\mathfrak G}
\providecommand{\frL}{\mathfrak L}
\providecommand{\frM}{\mathfrak M}
\providecommand{\frN}{\mathfrak N}
\providecommand{\frR}{\mathfrak R}
\providecommand{\frS}{\mathfrak S}
\providecommand{\frT}{\mathfrak T}
\providecommand{\frU}{\mathfrak U}
\providecommand{\frako}{\mathfrak o}
\providecommand{\mA}{\mathfrak A}
\providecommand{\mbarA}{\overline{\mathfrak A}}
\providecommand{\Gcal}{\mathcal G}
\providecommand{\Sp}{\operatorname{Sp}}
\providecommand{\glqq}{``}
\providecommand{\grqq}{''}
\emergencystretch=220em
\allowdisplaybreaks
\sloppy
\hbadness=10000
\vbadness=10000
\hfuzz=5pt
\vfuzz=12pt
\raggedbottom
% BEGIN UNIT BODY
\section*{40. Некоммутативные алгебры}

\begin{center}
\emph{\foreign{Math. Zs. 37 (1933), 514--541}}
\end{center}

\begin{center}
\textbf{Некоммутативная алгебра.}\\[0.5em]
Автор\\[0.25em]
Эмми Нётер, Гёттинген.
\end{center}

Главные теоремы коммутативной алгебры, как известно, содержатся в теории Галуа; ей предшествует теория полей отщепления и полей разложения, то есть полей, достаточных для того, чтобы заданный многочлен отщеплял линейный множитель или, соответственно, полностью разлагался на линейные множители.

Соответствующие части алгебры я развиваю здесь в некоммутативной, в частности гиперкомплексной, постановке. При этом я принципиально работаю некоммутативными методами, а именно представлениями в некоммутативных телах. В заключение я показываю, как названные выше теоремы коммутативной алгебры можно обосновать совершенно параллельно с помощью представлений в коммутативных полях.

Лежащая в основе теория представлений, которой предшествует краткая теория автоморфизмов (§1), есть дальнейшее развитие теории, основанной на модулях представлений.\footnote{E. Noether, \foreign{Hyperkomplexe Größen und Darstellungstheorie}, Math. Zeitschr. 30 (1929), S. 641--692; далее цит. как \foreign{Darstellungstheorie}. Ср. изложение у van der Waerden, \foreign{Moderne Algebra II}.} В частности, наряду с прямым представлением я рассматриваю реципрокное, причем одно сводится к другому; теперь оно порождается посредством реципрокного модуля представления. Преимущество состоит в том, что реципрокный модуль представления, то есть бимодуль, можно также понимать как односторонний модуль над кольцом расширения (§2). Тем самым представление в некоммутативных телах сводится к теории идеалов в кольце расширения; неприводимые классы представлений соответствуют неприводимым идеальным классам кольца расширения, в точной аналогии с фактами, которые при представлении гиперкомплексных систем в их коммутативной области коэффициентов имеют место для самой системы (§3 и 4).

Отсюда получаются структурные теоремы для матричных колец над некоммутативными телами благодаря замечанию (§5), что \emph{каждое подкольцо означает представление этим телом, соответственно реципрокное представление реципрокно изоморфным телом}. Это реципрокно изоморфное тело является первым аналогом минимального поля отщепления и поля разложения в коммутативном случае: оно осуществляет все реципрокные представления первой степени и, при конечном ранге над центром, который до сих пор не предполагался, также дает полное разложение в прямую сумму слагаемых ранга \(1\). Отсюда получается теория Галуа для тел (§6) и одновременно выражается факт (§7), что реципрокно изоморфные алгебры с делением, то есть тела конечного ранга над центром, порождают обратные классы в группе классов алгебр Р. Брауэра.

Второе обоснование теории Галуа, более общее для простых систем, является почти непосредственным следствием теоремы о коммутирующих подкольцах, которая сама почти прямо примыкает к предварительному рассмотрению автоморфизмов.

До этого места рассуждения идут чисто некоммутативно. Но и вопрос о \emph{коммутативных} полях разложения рассматривается некоммутативно, посредством представления поля разложения алгеброй с делением согласно замечанию, подготовленному в §5 (§7). Завершение составляет упомянутый выше перенос на коммутативный случай (§8) и теория полей разложения для произвольных систем (§9), что одновременно дает связь с обычной теорией представлений в коммутативном случае.

На этой теории представлений в коммутативном случае, а также на ``иррациональных'' фактор-системах Р. Брауэр основал теорию полей разложения.\footnote{Ср. нашу совместную заметку: \foreign{Über minimale Zerfällungskörper irreduzibler Darstellungen}. Sitz. Ber. d. Preuß. Ak. d. Wiss. 1927, S. 221--228. Примечание на с. 222 содержит своего рода отчет о положении дел. Главные теоремы возникли независимо и почти одновременно. Ср. далее R. Brauer, \foreign{Über Systeme hyperkomplexer Zahlen}, §3. Math. Zeitschr. 30 (1929), S. 79--107. A. A. Albert позднее независимо вновь нашел эти теоремы.} Из-за этого коммутативного обоснования он, так же как позднее Альберт, должен предполагать центр совершенным полем, что с точки зрения некоммутативного обоснования является ненужным ограничением. С тем же ограничением, также посредством теории представлений в коммутативном случае, Р. Брауэр и К. Шода, уже зная мою теорию Галуа для некоммутативных тел, развили теорию дальше: Р. Брауэр дал упомянутую выше теорему о коммутирующих подкольцах, а К. Шода независимо дал полную теорию также для полупростых систем.\footnote{R. Brauer, \foreign{Über die algebraische Struktur von Schiefkörpern}, §2. Journ. f. Math. 166 (1932), S. 241--252. K. Shoda, \foreign{Über die galoissche Theorie der halbeinfachen hyperkomplexen Systeme}. Math. Ann. 107 (1932), S. 252--258. Эти результаты также развил J. Levitzki.}

Наконец следует упомянуть, что для случая тел краткое изложение имеется у van der Waerden II (§128), в связи с моей лекцией лета 1928 г., когда я впервые докладывала об этих вопросах. Изложение van der Waerden содержит ряд упрощений по сравнению с этой лекцией; часть их я переняла во второй лекции (зима 1929/30),\footnote{Первую лекцию обработал G. Köthe, вторую M. Deuring.} а часть только здесь приняла и продолжила. В частности, перенос гиперкомплексного метода доказательства теории Галуа на коммутативный случай (§8) происходит лишь из второй лекции, где некоммутативное доказательство (§6, 3.) было существенно упрощено по сравнению с первой лекцией. Теорема о коммутирующих подкольцах (§5) и основанный на ней второй метод доказательства некоммутативной теории Галуа (§6, 1 и 2) добавились только после второй лекции, после ознакомления с Р. Брауэром (прим. 3) и van der Waerden §128; в §5, 3, однако, сделаны существенно более слабые предположения конечности, чем там.

Некоторые ходы первой лекции, а именно метод образования пересечений идеалов разностей, хотя и более сложны, имеют то преимущество, что переносятся на системы бесконечного ранга и на целочисленные системы.\footnote{Метод по существу воспроизведен у G. Köthe, \foreign{Schiefkörper unendlichen Ranges über dem Zentrum} §5, Math. Ann. 105 (1931), S. 15--39. Ср. также прим. 10. Этот метод пересечений теперь заменен почти тривиальным фактом (§4, 1), впервые замеченным van der Waerden, что двусторонние идеалы принадлежат инвариантным модулям. Благодаря этому все можно обосновать более простым разложением в прямую сумму, что, однако, не работает при бесконечном ранге и в целочисленном случае.} Для коммутативных целочисленных систем так приходят к связи между дифференцированием идеалов и дифферентой,\footnote{Ср. доклад в Праге, Jahresber. d. Deutsch. Math. Ver. 39 (1930), S. 17 (косая пагинация).} к чему я при случае вернусь подробнее.

Главное применение теория полей разложения находит в теории скрещенных произведений и их фактор-систем, которые, в свою очередь, дают основу теоретико-числовых приложений; здесь мы больше к этому не обращаемся.\footnote{Теория скрещенных произведений развита во второй лекции; она, с малыми изменениями, чтобы не предполагать данную здесь теорию, воспроизведена у H. Hasse: \foreign{Theory of cyclic algebras over an algebraic numberfield}, Chap. II. Transactions of the Amer. Math. Soc. 34 (1932), S. 171--214. Изложение, теснее связанное с лекцией, должно появиться в обзоре M. Deuring в \foreign{Ergebnisse der Mathematik}.}

\subsection*{§1. Автоморфизмы, модули и бимодули}

Теория представлений, основанная на модулях представлений, опирается на теорию кольца автоморфизмов абелевых групп с операторами или без них. Неявно лежащие при этом в основе рассуждения, то есть отношения между отображением и законами вычисления, будут, во избежание повторений, сформулированы в нескольких простых утверждениях.

\paragraph{1. Мультипликативное отображение. Закон ассоциативности.}
Пусть сначала \(\frG\) есть группа без операторов, а \(\frA\) ее абсолютная область автоморфизмов, то есть система всех гомоморфизмов \(\frG\) в себя. \(\frA\) мультипликативно замкнута, потому что произведение \( \sigma\tau \) определено через
\[
        g(\sigma\tau)=(g\sigma)\tau,
        \qquad g\in\frG,\quad \sigma,\tau\in\frA                         \tag{1}
\]
и, очевидно, удовлетворяет закону ассоциативности.

\(\frG\), как известно, становится группой с операторами, если задано множество \(\frB\) символов \(O,H,\ldots\) так, что сочетания \(gO,gH,\ldots\) дают однозначно определенные элементы из \(\frG\) и порождают гомоморфизмы \(\frG\) в себя:
\[
        (gh)O=gO\cdot hO .
\]
Если область операторов \(\frB\) мультипликативно замкнута, то однозначное отображение \(\frB\) на соответствующую часть области автоморфизмов является мультипликативно-гомоморфным тогда и только тогда, когда выполнено ассоциативное соотношение
\[
        g(OH)=(gO)H, \qquad g\in\frG,\ O,H\in\frB,                       \tag{1a}
\]
Соответственно определяется реципрокный гомоморфизм, когда речь идет о левых операторах.

Действительно, отображение \(\frB\) на \(\overline{\frB}\) задается равенством \(g\Theta=g\sigma\) для всех \(g\) из \(\frG\), а значит также равенством
\[
        (g\Theta)H=(g\sigma)\tau=g(\sigma\tau),
\]
так что (1a) становится необходимым и достаточным условием мультипликативной гомоморфности. То, что левые операторы порождают реципрокный гомоморфизм, происходит оттого, что записанные слева автоморфизмы \(\sigma^\ast g,\tau^\ast\sigma^\ast g\) следует читать справа налево: \(\tau^\ast\sigma^\ast g=g\sigma\tau\), тогда как операторы всегда читаются слева направо.

\paragraph{2. Операторно-гомоморфное отображение.}
Если \(\frG\) есть группа с операторами, то операторная гомоморфность определяется равенством
\[
        (gO)\sigma=(g\sigma)O,                                          \tag{2}
\]
а для левых операторов соответственно
\[
        (Og)\sigma=O(g\sigma).                                          \tag{2*}
\]
то есть коммутативной связанностью или сквозным законом ассоциативности. Из рассуждения пункта 1, то есть из отображения на область автоморфизмов, для двух различных областей операторов получается:

Если заданы две области операторов \(\frB\) и \(\frC\), то элементы \(\frC\) порождают \(\frB\)-автоморфизмы и одновременно элементы \(\frB\) порождают \(\frC\)-автоморфизмы тогда и только тогда, когда эти области коммутативно связаны с \(\frG\), соответственно когда выполняется сквозной закон ассоциативности:
\[
        (gO)H=(gH)O,                                                    \tag{2a}
\]
соответственно
\[
        (OH)g=O(Hg).                                                    \tag{2a*}
\]

\paragraph{3. Модули и бимодули над кольцами.}
Правый модуль \(\frM\) над кольцом \(\frR\) есть аддитивная абелева группа с элементами из \(\frR\) в качестве правых операторов, причем помимо ассоциативного соотношения (1a) выполнены также дистрибутивные соотношения:
\[
        (g+h)\rho=g\rho+h\rho,\qquad
        g(\rho+\sigma)=g\rho+g\sigma.                                  \tag{3a}
\]
Соответствующее верно для левых модулей.

Среди бимодулей нужно различать два вида. Правые модули над двумя кольцами \(\frR\) и \(\frS\) называются бимодулями, если кроме связей, действующих отдельно для \(\frR\) и \(\frS\), еще выполнено
\[
        (m\rho)\sigma=(m\sigma)\rho .
\]
\(\frR\)-левые, \(\frS\)-правые модули называются бимодулями, если добавляется сквозной закон ассоциативности
\[
        \rho(m\sigma)=(\rho m)\sigma .
\]

Значение этих связей вытекает из того, что в случае абелевых групп абсолютная область автоморфизмов образует кольцо, абсолютное кольцо автоморфизмов\footnote{Для неабелевых групп речь идет об \glqq{}обобщенном\grqq{} кольце; ср. работу H. Fitting, которая должна появиться в Math. Annalen. [Math. Annalen 107, S. 514--542.]}. В частности, рассуждение через отображение дает:

Если область операторов \(\frR\) аддитивно записанной абелевой группы \(\frM\) есть кольцо, то однозначное отображение \(\frR\) на подмножество \(\overline{\frR}\) абсолютного кольца автоморфизмов является кольцевым гомоморфизмом тогда и только тогда, когда \(\frM\) есть правый модуль над \(\frR\), то есть когда выполнены (1a) и (3a); для левых модулей получается реципрокная кольцевая гомоморфность.

Если \(\frM\) есть правый модуль над кольцами \(\frR\) и \(\frS\), то \(\frM\) становится бимодулем, то есть выполняется (2a), тогда и только тогда, когда \(\frR\) можно кольцево-гомоморфно отобразить на подкольцо кольца \(\frS\)-автоморфизмов и одновременно \(\frS\) на подкольцо кольца \(\frR\)-автоморфизмов. Если \(\frM\) есть правый модуль над \(\frS\) и левый модуль над \(\frR\), то бимодуль, то есть (2a*), означает гомоморфное отображение \(\frS\) на подкольцо кольца \(\frR\)-автоморфизмов и реципрокно гомоморфное отображение \(\frR\) на подкольцо кольца \(\frS\)-автоморфизмов.

Отсюда, в частности, следует:

1) Если \(\frR\) есть кольцо с единицей, то \(\frR\) прямо изоморфно своему кольцу автоморфизмов как \(\frR\)-левый модуль и реципрокно изоморфно своему кольцу автоморфизмов как \(\frR\)-правый модуль.

2) Если \(\frR\) есть полное матричное кольцо над вообще говоря некоммутативным телом \(A\),
\[
        \frR=\sum c_{ik}A=\sum A c_{ik},
\]
то \(A\) становится реципрокно изоморфным телу автоморфизмов простых правых идеалов и прямо изоморфным телу автоморфизмов простых левых идеалов.

\paragraph{4. Переход от правого бимодуля к одностороннему модулю над кольцом произведения.}
На этом переходе существенно основано значение правого бимодуля.

\paragraph{Теорема.}
Если \(\frM\) есть правый модуль над кольцом \(\frT\), содержащим два поэлементно коммутирующих подкольца \(\frR\) и \(\frS\), то \(\frM\) можно также рассматривать как бимодуль над \(\frR\) и \(\frS\). Обратно, если задан бимодуль над \(\frR\) и \(\frS\) и существует кольцо произведения \(\frT\), в котором \(\frR\) и \(\frS\) поэлементно коммутируют, то \(\frM\) можно также рассматривать как правый модуль над \(\frT\), причем \(\frT\)-связь продолжает данные \(\frR\)- и \(\frS\)-связи.

Первое утверждение следует из того, что \(\frT\) по определению можно гомоморфно отобразить на подкольцо \(\overline{\frT}\) абсолютного кольца автоморфизмов, причем \(\frR,\frS\) переходят в поэлементно коммутирующие подкольца \(\overline{\frR},\overline{\frS}\). Но это как раз означает соотношение (2a), характеризующее бимодуль.

Обратно, если \(\frM\) задан как \(\frR,\frS\)-бимодуль, то есть как правый модуль, то \(\frR,\frS\) снова можно гомоморфно отобразить на поэлементно коммутирующие подкольца \(\overline{\frR},\overline{\frS}\) абсолютного кольца автоморфизмов. В абсолютном кольце автоморфизмов для \(\overline{\frR},\overline{\frS}\) существует кольцо произведения\footnote{Кольцо произведения здесь означает лишь кольцо, порожденное двумя кольцами с общим надкольцом; произведение не обязано быть прямым. То, что при заданном пересечении кольцо произведения не всегда обязано существовать, показывает следующий пример. Пусть \(\frako\) есть кольцо целых чисел и положено \(\frR=\frako[x]\), \(\frS=\frako[y]\) с \(2x-1=0\), \(2y-1=0\), так что \(\frako\) есть пересечение \(\frR\) и \(\frS\), если между \(x\) и \(y\) не задана никакая связь. Но если \(\frR\) и \(\frS\) лежат в общем надкольце, то получается \((2x-1)y-x(2y-1)=0\), значит \(x=y\). Следовательно, кольца произведения с пересечением \(\frako\) не существует.} \(\overline{\frT}\), которое из-за поэлементной коммутативности состоит из всех элементов вида
\[
        \sum_i \bar r_i\bar s_i+\bar r+\bar s
\]
(если \(\overline{\frR},\overline{\frS}\) имеют единицы, то добавочные члены \(\bar r,\bar s\) отпадают). Если теперь существует также кольцо произведения \(\frT\), в котором \(\frR\) и \(\frS\) поэлементно коммутируют, то оно соответственно состоит из всех элементов вида \(\sum_i r_i s_i+r+s\). Гомоморфизмы
\[
        \frR\to\overline{\frR},\qquad \frS\to\overline{\frS}
\]
можно, следовательно, через соответствие
\[
        \sum_i r_i s_i+r+s
        \longmapsto
        \sum_i \bar r_i\bar s_i+\bar r+\bar s
\]
дополнить до гомоморфизма \(\frT\to\overline{\frT}\); поэтому связь
\[
        m\Bigl(\sum_i r_i s_i+r+s\Bigr)
        =\sum_i (m r_i)s_i+m r+m s
\]
однозначна и определяет по §1, 3 \(\frT\)-модуль, охватывающий данный \(\frR,\frS\)-модуль.

\subsection*{§2. Реципрокное и прямое представление}

\paragraph{1. Модули линейных форм.}
Переход от теории, развитой в §1, к теории представлений основан на специализации тамошних модулей к модулям линейных форм и на рассмотрении их колец автоморфизмов.

\(\frS\)-правый модуль \(\frM\), где \(\frS\) есть кольцо с единицей, называется модулем линейных форм в \(\frS\), если \(\frM\) есть прямая сумма \(n\) одночленных \(\frS\)-модулей:
\[
        \frM=m_1\frS+\cdots+m_n\frS,
\]
так что \(m_i\frS\) операторно изоморфен \(\frS\), то есть из \(m_i s=0\) всегда следует \(s=0\). Отсюда получается, что \(\frS\) можно отобразить на подкольцо абсолютного кольца автоморфизмов не только гомоморфно, но и изоморфно.

\paragraph{Теорема 1.}
Если \(\frM\) есть модуль линейных форм в \(\frS\), то кольцо \(\frS\)-автоморфизмов \(\mA\) модуля \(\frM\) реципрокно изоморфно, не только гомоморфно, кольцу \(\mbarA\) всех \(n\)-рядных матриц над \(\frS\).

Произвольный \(\frS\)-автоморфизм \(\alpha\) модуля \(\frM\) полностью описывается соответствием
\[
        m_i\longmapsto \overline m_i=m_i\alpha;
\]
потому что из этого по определению следует
\[
        \sum m_i s_i\longmapsto
        \sum \overline m_i s_i
        =\sum_i (m_i\alpha)s_i
        =\sum (m_i s_i)\alpha.
\]
Если теперь в матрической записи
\[
        (m_1\alpha,\ldots,m_n\alpha)=(m_1,\ldots,m_n)A,
\]
то, когда \(\alpha\) пробегает все автоморфизмы, соответствие \(\alpha\mapsto A\) дает взаимно однозначную связь между \(\mA\) и \(\mbarA\). Ведь, поскольку \(\frM\) предположен модулем линейных форм, \(A\) однозначно определяется через \(\alpha\), и каждая матрица \(A\) \(n\)-й степени в \(\frS\) порождает автоморфизм. Далее следует:
\[
        \alpha+\beta\mapsto A+B,\qquad
        \alpha\beta\mapsto BA.
\]
Последнее следует из
\[
        (m_1\alpha\beta,\ldots,m_n\alpha\beta)
        =(m_1,\ldots,m_n)AB
        =(m_1,\ldots,m_n)BA
\]
в силу \(ms\beta=m\beta s\).

\paragraph{2. Представление и модуль представления.}
Если под реципрокным, соответственно прямым, представлением \(n\)-й степени кольца \(\frR\) в \(\frS\) понимать реципрокный, соответственно прямой, кольцевой гомоморфизм \(\frR\) в подкольцо кольца всех \(n\)-рядных матриц над \(\frS\), то теорему 1 можно сформулировать также так:

\paragraph{Теорема 1'.}
Кольцо автоморфизмов \(\mA\) модуля линейных форм \(n\)-й степени в \(\frS\) допускает изоморфное, то есть точное, реципрокное представление полным матричным кольцом \(\mbarA\) всех \(n\)-рядных матриц над \(\frS\).

Отсюда обычные теоремы для прямых представлений получают также формулировки для реципрокных.

\paragraph{Определение.}
Модуль линейных форм \(\frM\) в \(\frS\) называется реципрокным модулем представления \(\frR\) в \(\frS\), если \(\frM\) есть бимодуль над \(\frR\) и \(\frS\) и одновременно \(\frR\)- и \(\frS\)-правый модуль. Он называется прямым модулем представления, если \(\frM\) есть бимодуль и одновременно \(\frR\)-левый, \(\frS\)-правый модуль.

\paragraph{Теорема 2.}
Каждый реципрокный, соответственно прямой, модуль представления порождает класс эквивалентных реципрокных, соответственно прямых, представлений \(\frR\) в \(\frS\), и все реципрокные, соответственно прямые, представления порождаются так.

Сначала пусть \(\frM\) есть реципрокный модуль представления. Тогда по §1, 3 \(\frR\) прямо гомоморфно отображается на подкольцо \(\overline{\frR}\) кольца \(\frS\)-автоморфизмов модуля \(\frM\); по §2, теорема 1', \(\overline{\frR}\) допускает реципрокное, точное, представление, которое тем самым становится и реципрокным представлением \(\frR\). Обратно, если дано реципрокное представление \(\frR\to\frR^*\), то его композиция с реципрокным изоморфизмом \(\frR^*\) к подкольцу \(\overline{\frR}\) кольца \(\frS\)-автоморфизмов дает прямой гомоморфизм \(\frR\) на \(\overline{\frR}\); следовательно, \(\frM\) становится \(\frR\)-модулем, именно бимодулем, то есть модулем представления по §1, 3. Переход к другим \(\frS\)-базисам дает класс эквивалентных представлений.

Если речь идет о прямом модуле представления, то \(\frR\) реципрокно гомоморфно отображается на \(\overline{\frR}\). Композиция с реципрокным представлением \(\overline{\frR}\) поэтому дает прямое представление \(\frR\). Обратно, из прямого представления следует реципрокно гомоморфное отображение на \(\overline{\frR}\); после этого \(\frM\) становится прямым модулем представления по §1, 3.

\paragraph{Замечание.}
Разумеется, оба вида представлений можно сводить друг к другу: прямое представление \(\frR\) есть реципрокное представление кольца, реципрокного к \(\frR\), и обратно. Другое сведение состоит в переходе от \(\frS\) к реципрокному кольцу и замене матриц транспонированными. Это соответствует переходу от \(\frR\)-левого, \(\frS\)-правого модуля к \(\frR\)- и \(\frS\)-левому модулю.

\paragraph{3. Переход от реципрокного модуля представления к модулю над кольцом расширения.}
Преимущество реципрокного модуля представления основано прежде всего на возможном по §1, 4 переходе к кольцу расширения. Для настоящего специального случая, где \(\frM\) есть модуль линейных форм над \(\frS\), речь, следовательно, идет о теореме перехода для модулей представлений:

\emph{Если \(\frM\) есть правый модуль над кольцом \(\frT\), содержащим два поэлементно коммутирующих подкольца \(\frR\) и \(\frS\) так, что \(\frM\) становится модулем линейных форм над \(\frS\), то \(\frM\) можно также рассматривать как реципрокный модуль представления \(\frR\) в \(\frS\).}

\emph{Обратно, если задан реципрокный модуль представления \(\frR\) в \(\frS\) и существует кольцо произведения \(\frT\), в котором \(\frR\) и \(\frS\) поэлементно коммутируют, то \(\frM\) можно также рассматривать как \(\frT\)-модуль так, что \(\frT\)-связь является продолжением данных \(\frR\)- и \(\frS\)-связей.}

В теории представлений существенно используется следующее

\paragraph{Следствие из теоремы перехода для модулей представлений.}
Если \(\frM\) есть реципрокный модуль представления \(\frR\) в \(\frS\) и одновременно \(\frT\)-модуль с \(\frT\) как кольцом произведения, то \(\frR\)-, \(\frS\)-изоморфизм \(\frM\) с модулем \(\frN\) равносилен \(\frT\)-изоморфизму. \emph{Следовательно, каждому классу изоморфных \(\frT\)-модулей соответствует класс реципрокных представлений \(\frR\) в \(\frS\), и обратно.}

Связь этих теорем с §1, 2 и 3 дает основной для дальнейшего ход связь между матрицами, коммутирующими с данным представлением, и \(\frR\)-, \(\frS\)-автоморфизмами.

\paragraph{Теорема о коммутирующих матрицах.}
\emph{Если \(\frR\to\frR^*\) есть реципрокное представление \(n\)-й степени кольца \(\frR\) в \(\frS\), а \(\frB^*\) обозначает кольцо всех матриц \(n\)-й степени над \(\frS\), поэлементно коммутирующих с \(\frR^*\), то \(\frB^*\) дает реципрокно-изоморфное представление \(\frR\)-, \(\frS\)-автоморфизмов, то есть тех \(\frS\)-автоморфизмов, которые одновременно являются \(\frR\)-автоморфизмами, порождающего модуля представления или также, если существует кольцо произведения \(\frT\), \(\frT\)-автоморфизмов.}

Пусть \(\frB\) есть кольцо всех \(\frR\)-, \(\frS\)-автоморфизмов модуля \(\frM\). Как подкольцо кольца \(\frS\)-автоморфизмов \(\mA\), \(\frB\) допускает реципрокно изоморфное представление в \(\frS\). По §1, 2 и 3 \(\frB\) состоит из всех таких автоморфизмов из \(\mA\), которые коммутативно связаны с \(\frR\), то есть удовлетворяют соотношению (2a). Следовательно, \(\frB\) состоит из совокупности тех автоморфизмов, которые поэлементно коммутируют с элементами \(\overline{\frR}\), где \(\overline{\frR}\) обозначает образ \(\frR\) в \(\mA\). Поскольку в представлении \(\frB\) и \(\overline{\frR}\) речь идет о реципрокном изоморфизме, а не только о гомоморфизме, это и есть утверждение, подлежащее доказательству.

\paragraph{Замечание.}
\(\frR\)-, \(\frS\)-автоморфизмы означают соответствие \(m_i\mapsto m_i'\) такое, что посредством \(m_i'\) возникает то же представление \(\frR^*\); при этом \(m_i'\) не обязательно должны давать \(\frS\)-базис модуля \(\frM\), соответственно тому факту, что матрицы из \(\frB^*\) не обязательно обратимы. Если сумма \(\sum m_i'\frS\) уже не является прямой, ``представление'' следует понимать в переносном смысле:
\[
        (m_1',m_2',\ldots,m_n')a=(m_1',\ldots,m_n')A
\]
остается верным, но \(A\) уже не определяется однозначно через \(a\).

Еще одно замечание к теореме перехода таково. Диагональные матрицы \(E\cdot s\) дают \emph{прямо изоморфное} представление \(\frS\), ``тождественное'' представление \(\frS\). То, что здесь речь идет о \emph{прямом} изоморфизме, следует из того, что кольцо \(\frS\)-автоморфизмов модуля \(\frM\) имеет подкольцо, реципрокно изоморфное \(\frS\), и речь идет именно о его реципрокном представлении. Ведь каждый одночленный слагаемый \(m_i\frS\) модуля \(\frM\) операторно изоморфен \(\frS\) как правый модуль, значит его кольцо \(\frS\)-автоморфизмов реципрокно изоморфно \(\frS\) (§1, 3, следствие 1).

Соответствие от \(\frT\) к матрицам над \(\frS\), однозначно определенное через \(\frR\) и \(\frS\), следовательно, является \emph{не представлением кольца \(\frT\), а продолжением реципрокного представления \(\frR\) до операторно-гомоморфного относительно \(\frS\):}
\[
        \sum r_i s_i\longmapsto \sum R_i s_i.
\]
В частности, реципрокное представление \(\frR\) одновременно \emph{операторно-гомоморфно относительно пересечения \([\frR,\frS]\) колец \(\frR\) и \(\frS\), лежащего в центре \(\frR\).}

\subsection*{§3. Модули относительно тела}

Далее речь будет идти только о представлениях в вообще говоря некоммутативных телах. Поэтому нужно собрать несколько простых теорем о модулях линейных форм над телом.

\paragraph{1. Нормальный базис подмодуля относительно заданного базиса полного модуля.}
Пусть \(A\) есть тело и
\[
        \frN=x_1A+\cdots+x_nA
\]
есть модуль линейных форм ранга \(n\); далее пусть
\[
        \frL=z_1A+\cdots+z_\ell A
\]
есть подмодуль ранга \(\ell\le n\). \(z_i\) называются нормальным базисом относительно \(x\), если они имеют вид
\[
        z_i=x_i-(x_{\ell+1}\alpha_{i,\ell+1}+\cdots+x_n\alpha_{i,n}),
        \qquad i=1,\ldots,\ell.
\]
Каждый модуль \(\frL\) после подходящей нумерации \(x_i\) имеет нормальный базис. Действительно, после подходящей нумерации получается
\[
        \frN=\frL+x_{\ell+1}A+\cdots+x_nA,
\]
то есть
\[
        x_i\equiv x_{\ell+1}\alpha_{i,\ell+1}+\cdots+x_n\alpha_{i,n}
        \pmod{\frL},\qquad i=1,\ldots,\ell,
\]
и потому соответствующие
\[
        z_i=x_i-(x_{\ell+1}\alpha_{i,\ell+1}+\cdots+x_n\alpha_{i,n})
\]
лежат в \(\frL\) и исчерпывают \(\frL\).

\paragraph{2. Модуль расширения.}
Снова пусть \(A\) есть тело, \(P\) подтело. \(A\)-модуль \(N\) ранга \(n\) называется модулем расширения \(P\)-модуля \(M\), \(N=M_A\), если \(M\) есть подмодуль того же ранга в \(N\). Если
\[
        M=y_1P+\cdots+y_nP,
\]
то, следовательно,
\[
        N=M_A=y_1A+\cdots+y_nA.
\]
Обратно, для каждого \(P\)-модуля \(M=y_1P+\cdots+y_nP\) с точностью до изоморфии модулей однозначно существует модуль расширения \(M_A\). Ведь модуль \(y_1A+\cdots+y_nA\) содержит подмодуль, изоморфный \(M\), который можно отождествить с \(M\).

\paragraph{Следствие a).}
Если \(z_1,\ldots,z_s\) являются линейно независимыми относительно \(P\) элементами из \(M\), то они также линейно независимы над \(A\) в \(N=M_A\), потому что их можно дополнить до базиса \(M\), а тем самым и до базиса \(N\). Каждый \(P\)-модуль
\[
        T=z_1P+\cdots+z_sP
\]
из \(M\) порождает, следовательно, модуль расширения
\[
        T_A=z_1A+\cdots+z_sA\subset M_A.
\]

\paragraph{Следствие b).}
Каждый \(P\)-модуль \(T=z_1P+\cdots+z_sP\) из \(M\) является модулем сужения, то есть пересечением своего модуля расширения \(T_A\) с \(M\):
\[
        T=T_A\cap M .
\]
Ведь \(T\subseteq T_A\cap M\); обратно, если \(a\) лежит в \(T_A\cap M\), то
\[
        a=\sum z_i\alpha_i,
\]
значит элементы \(M\) \(a,z_1,\ldots,z_s\) зависимы над \(A\), а потому по следствию a) также зависимы над \(P\); следовательно, \(a\) лежит в \(T\).

\paragraph{3. Теорема об инвариантных модулях.}
Снова пусть \(N=M_A\) есть модуль расширения \(P\)-модуля \(M\) ранга \(n\), а \(P\) есть полное тело инвариантов относительно группы \(\frG\) кольцевых автоморфизмов \(A\). Группа \(\frG\) определяется как область операторов модуля \(N=M_A\) условиями
\[
        Gm=m\quad\text{для }m\text{ из }M\text{ и }G\text{ из }\frG
\]
и
\[
        G\sum m_i\alpha_i=\sum m_i\,G(\alpha_i)
        \quad\text{для }m_i\text{ из }M\text{ и }\alpha_i\text{ из }A.
\]

\paragraph{Лемма.}
При этих условиях \(M\) состоит из всей совокупности элементов из \(N\), остающихся неподвижными при \(\frG\). Ведь если \(x_1,\ldots,x_n\) есть \(P\)-базис \(M\), то есть \(w=x_1\alpha_1+\cdots+x_n\alpha_n\), то из \(G(w)=w\) для каждого \(G\) из \(\frG\) следует, что также \(G(\alpha_i)=\alpha_i\), а значит по предположению \(\alpha_i\) лежит в \(P\).

\paragraph{Теорема.}
Модули расширения \(L_A\) подмодулей \(L\) модуля \(M\), и только они, являются допустимыми подмодулями относительно \(\frG\) как области операторов.

Первая половина ясна. Обратно, пусть \(L\) допустим относительно \(\frG\), и пусть \(z_1,\ldots,z_\ell\) есть нормальный базис \(L\) (ср. 1) относительно \(x_1,\ldots,x_n\), где \(x_1,\ldots,x_n\) выбран как \(P\)-базис \(M\), так что \(\frG\) поэлементно его допускает. По предположению
\[
        G(z_i)=x_i-\bigl(x_{\ell+1}G(\alpha_{i,\ell+1})
              +\cdots+x_nG(\alpha_{i,n})\bigr)
\]
есть элемент \(L\) для каждого \(G\) из \(\frG\), значит при фиксированном \(G\) он линейно выражается через \(z\), и поэтому, как видно из сравнения коэффициентов при \(x_1,\ldots,x_n\), имеем \(G(z_i)=z_i\) для каждого \(G\) из \(\frG\). Следовательно, по лемме, \(z\) лежат в \(M\), и \(L\) является расширением \(P\)-модуля
\[
        T=z_1P+\cdots+z_\ell P
\]
из \(M\), причем \(T=L\cap M\) (по 2).\footnote{Теоремы этого параграфа благодаря простым рассуждениям с хорошим упорядочением остаются верными и тогда, когда ранг \(M\) над \(P\) становится бесконечным; ср. G. Köthe, \foreign{Ein Beitrag zur Theorie der kommutativen Ringe ohne Endlichkeitsvoraussetzung} §1, Gött. Nachr. 1931, S. 195--207.}

\subsection*{§4. Гиперкомплексные системы и их классы представлений}

\paragraph{1. Теорема расширения.}
Теперь мы соединяем теоремы о представлениях из §2 с модульными теоремами из §3. Вместо общих колец \(\frR\) рассматриваем гиперкомплексные системы с единицей, \(S,T,\ldots\), над коммутативным полем \(P\), то есть
\[
        S=x_1P+\cdots+x_nP;
\]
вместо произвольных колец представления \(\frS\) берем только тела \(A,B,\ldots\), и, если не оговорено иное, такие, что имеют центр \(P\). В этом случае всегда существует кольцо произведения, в котором \(S\) и \(A\) поэлементно коммутируют, с пересечением \(P\); речь именно о прямом произведении \(S\times A\), которое одновременно становится модулем расширения \(S_A\) системы \(S\) в смысле §3,
\[
        S_A=x_1A+\cdots+x_nA,
\]
и поэтому также называется кольцом расширения.

При этой специализации теорема об инвариантных модулях (§3, 3) переходит в

\paragraph{Теорема расширения.}
Каждый двусторонний идеал из \(S_A\) есть идеал расширения некоторого идеала из \(S\), а именно расширение своего пересечения с \(S\).

Действительно, если положенная в основу группа \(\frG\) есть группа внутренних автоморфизмов \(A\), то \(P\) как центр \(A\) становится полным телом инвариантов, тогда как поэлементная коммутативность \(A\) с \(S\) означает, что \(\frG\) по §3, 3 продолжена на \(S_A\) так, что \(S\) становится полной областью инвариантов. Каждый двусторонний идеал \(\mathfrak a\) из \(S_A\) является допустимой подгруппой, поскольку \(\alpha^{-1}\mathfrak a\alpha\subseteq\mathfrak a\), а значит по §3, 3 есть расширение своего модуля пересечения \(\mathfrak a\cap S\), который становится идеалом в \(S\).

Добавление: если \(S\) коммутативна, то \(S\) становится центром \(S_A\). Ведь \(S\) лежит в центре и является полной областью инвариантов относительно внутренних автоморфизмов, порожденных \(A\). Более общо, центр \(S\) также становится центром \(S_A\).

Обозначение: под \(A_r,B_n,\ldots\) всегда будем понимать матричные кольца указанной степени над \(A,B,\ldots\), то есть
\[
        A_r=\sum A c_{ik}=\sum c_{ik}A.
\]

В дальнейшем существенно используется одно следствие, теорема расширения для простых систем: если \(S\) простая система\footnote{Простая система означает, как обычно, \glqq{}двусторонне простая\grqq{}.}, то \(S_A\) также проста:
\[
        S_A=\sum B c_{ik}=\sum c_{ik}B=B_t
\]
с соответствующим телом \(B\), центр которого содержит \(P\). При этом \(A\), а значит и \(B\), может иметь бесконечный ранг над \(P\); только \(S\) предполагается гиперкомплексной. Если \(S\) и \(A\) имеют общий центр \(P\), то \(P\) также становится центром \(S_A\).

Более общо еще верно: если \(S\) простая гиперкомплексная система, то прямое произведение \(S\times A_r\) также просто. Ведь \(S\times A_r\) совпадает с \(S_A\times P_r\), а потому равно \(B_{tr}\), если \(S_A=B_t\).

\paragraph{2. Классы представлений.}
Под реципрокным или прямым представлением гиперкомплексной системы будем, как обычно, понимать такое представление, которое операторно-гомоморфно относительно области коэффициентов \(P\) (§2, конец) и отлично от нулевого представления. Соединение теоремы перехода и ее следствия (§2, 3) с теоремой расширения из пункта 1 дает

\paragraph{Теорема о классах представлений.}
Если \(S\) гиперкомплексна с единицей, с \(P\) как областью коэффициентов, а \(A\) тело с центром \(P\), то \(S\) имеет столько различных неприводимых реципрокных классов представлений в \(A\), сколько есть классов простых правых идеалов в фактор-кольце \(S_A\) по его радикалу. Если \(S\) простая система, то она имеет ровно один неприводимый реципрокный и один неприводимый прямой класс представлений в \(A\).\footnote{Ср. \foreign{Darstellungstheorie}, прим. 20, для случая, когда \(A\) является соответствующим телом автоморфизмов \(S\).}

Действительно, по следствию из теоремы перехода (§2, 3) простые реципрокные классы представлений и простые модульные классы над \(S_A\) взаимно однозначно соответствуют друг другу. Поскольку \(S_A\) имеет конечный ранг над телом \(A\), оно удовлетворяет максимальному и минимальному условиям для односторонних идеалов; следовательно, первая часть утверждения непосредственно вытекает из \foreign{Darstellungstheorie}, §19. Если \(S\) проста, то по теореме расширения \(S_A\) также проста. Благодаря максимальному и минимальному условиям она одновременно вполне приводима, значит имеет лишь один простой класс односторонних идеалов; то есть \(S\) имеет лишь один неприводимый реципрокный класс представлений в \(A\). Вместе с \(S\) простой является и реципрокно изоморфная ей система, имеющая лишь один неприводимый реципрокный класс представлений; следовательно, \(S\) имеет также лишь один простой прямой класс представлений.

\paragraph{3. Ранговое соотношение.}
Факт, что, когда \(S\) проста,
\[
        S_A=\sum Bc_{ik}
\]
имеет конечный ранг как над \(A\), так и над \(B\), дает ранговое соотношение:
\[
\begin{gathered}
        n=rt,\qquad n=(S:P),\qquad t^2=(S_A:B),\\
        r=\text{степень неприводимого реципрокного представления}.
\end{gathered}                                                            \tag{1}
\]

Ведь \(S_A\) само является реципрокным модулем представления для \(S\) в \(A\) (представление уже лежит в \(P\)). Как матричное кольцо над \(B\), \(S_A\) распадается на \(t\) простых правых идеалов, то есть \(S_A\)-модулей, ранг которых над \(A\) равен степени \(r\) неприводимого реципрокного представления. Поскольку \(n\) одновременно есть ранг \(S_A\) над \(A\), получается ранговое соотношение.

\paragraph{4. Усиление рангового соотношения, когда \(A\) имеет конечный ранг над \(P\).}
В этом случае выполняются три соотношения:
\[
        (S:P)=n=rt,                                                     \tag{1}
\]
\[
        (B:P)\,t=(A:P)\,r,                                               \tag{2}
\]
\[
        (S:P)(B:P)=(A:P)\,r^2.                                           \tag{3}
\]
Здесь (2) означает ранг над \(P\) простого правого идеала из \(S_A\), выраженный по пункту 3 через ранг над \(B\), соответственно над \(A\), а (3) означает ранг простого правого идеала в матричном кольце \(B_n\) и получается из (2) умножением на \(r\) согласно (1).

\subsection*{§5. Применение теорем о представлениях к матричным кольцам}

Реципрокное представление \(S\) в \(A\), рассмотренное в §4, можно также понимать как прямой гомоморфизм \(S\) в подкольцо \(A_r\), где \(\overline A\) обозначает тело, реципрокно изоморфное \(A\). Принцип следующих параграфов состоит в том, чтобы, наоборот, сводить исследование матричных колец \(A_r\) к теории представлений в реципрокно изоморфном теле \(A\).

\paragraph{1. Приводимое и неприводимое вложение в \(A_r\).}
Пусть \(A\) и \(\overline A\) являются реципрокно-изоморфными телами конечного или бесконечного ранга над своим центром \(P\).\footnote{Вообще реципрокные тела или системы будут обозначаться соответствующими греческими и латинскими буквами.} Простая система \(S\) над \(P\) называется неприводимо, соответственно приводимо, вложимой в \(A_r\), если \(S\) имеет неприводимое, соответственно приводимое, реципрокное представление \(r\)-й степени в \(A\). Если \(S\) неприводимо вложима в \(A_r\), то она приводимо вложима во все матричные кольца \(A_{rs}\) с \(s>1\), потому что она имеет приводимые представления именно этих и только этих степеней (§4, 2). Следовательно, \(S\) прямо изоморфна некоторым подкольцам \(A_r\) и \(A_{rs}\). Изоморфизм является продолжением тождественного на \(P\);\footnote{Как уже для представлений, во всех далее встречающихся изоморфизмах речь всегда идет о таких, которые являются продолжением тождественного на \(P\), то есть изоморфны над \(P\), даже если это не оговаривается каждый раз отдельно.} по §4, 3 при этом \(r\) является делителем ранга \(n\) системы \(S\) над \(P\).

Таким образом получаются все простые подкольца конечного ранга над \(P\), содержащие \(P\), в различных \(A_f\) (\(f=1,2,\ldots\)); ведь каждое такое подкольцо реципрокно-изоморфно подкольцу \(\overline A_f\), а значит допускает приводимое или неприводимое реципрокное представление в \(A\).

\paragraph{2. Теорема о внутренних автоморфизмах.}
Если \(S^{(1)}\) и \(S^{(2)}\) два простых подкольца \(A_f\), содержащих \(P\) и имеющих конечный ранг над \(P\), и если \(S^{(1)}\) и \(S^{(2)}\) изоморфны над \(P\), то этот изоморфизм порождается внутренним автоморфизмом \(A_f\).

Ведь \(S^{(1)}\) и \(S^{(2)}\) означают, вообще говоря приводимые, вложения простой системы \(S\) в \(A_f\), то есть прямые представления \(f\)-й степени в \(A\); поэтому они принадлежат одному и тому же приводимому прямому классу представлений в \(A\), и переводящая матрица порождает автоморфизм.

Если \(A\) само имеет конечный ранг над \(P\), то есть гиперкомплексно, то теорема переходит в известную: две простые подсистемы \(A_f\), содержащие центр \(P\) и изоморфные над \(P\), переводятся друг в друга внутренним автоморфизмом \(A_f\). В частности, каждый автоморфизм \(A_f\) внутренний.\footnote{Это уже неверно, когда \(A\) имеет бесконечный ранг над \(P\); ср. Köthe в работе, цитированной в прим. 5.}

\paragraph{3. Теорема о поэлементно коммутирующих подкольцах.}
Эта теорема снова следует из теоремы о коммутирующих матрицах из §2, 3 по принципу, упомянутому в начале параграфа:

Если \(S\) есть простая система из \(A_f\), содержащая \(P\), и \(R\) обозначает совокупность всех элементов \(A_f\), поэлементно коммутирующих с \(S\), то \(R\) также является матричным кольцом: \(R=\overline B_s\). Соответствующие тела \(B\) для \(S_A\) и \(\overline B\) для \(R\) реципрокно изоморфны. Пересечение \(R\) и \(S\) становится центром \(S\). \(R\) является телом тогда и только тогда, когда \(S\) вложено в \(A_f\) неприводимо.

Ведь по §2, 3 \(R\) прямо изоморфно кольцу автоморфизмов реципрокного модуля представления \(\frM\) системы \(S\) в \(A\); но это кольцо автоморфизмов \(\frM\), рассматриваемого как \(S_A\)-модуль. Если \(\frM\) распадается на \(s\) простых \(S_A\)-модулей и, как в §4, 1, положено \(S_A=\sum Bc_{ik}\), то \(R\) изоморфно \(\sum \overline Bc_{ik}\), где \(B\) и \(\overline B\) реципрокно изоморфны (поскольку \(B\) реципрокно изоморфно телу автоморфизмов простых правых идеалов из \(S_A\), то есть простых \(S_A\)-модулей, по концу §1, 3). \(R\) становится телом, то есть \(R=\overline B\), тогда и только тогда, когда \(s=1\), то есть когда \(S\) вложено неприводимо. То, что пересечение \(R\) и \(S\) есть центр \(S\), непосредственно следует из определения \(R\).

\paragraph{4. Усиленная теорема о коммутантах в гиперкомплексном матричном кольце.}
В этом случае ранговые соотношения из §4, 4 дают следующее усиление:

Простые подсистемы \(A_f\), содержащие \(P\), где \(A\) имеет конечный ранг над своим центром \(P\), распадаются на пары \(S,\overline S\) так, что \(\overline S\) состоит из всей совокупности элементов, поэлементно коммутирующих с \(S\), и обратно. Соответствующие тела для \(S_A\) и \(\overline S\) реципрокно изоморфны, так же как соответствующие тела для \(S\) и \(\overline S_A\). Пересечение \(S\) и \(\overline S\) равно общему центру. Одно подкольцо является телом тогда и только тогда, когда другое вложено неприводимо. Произведение ранговых чисел над \(P\) систем \(S\) и \(\overline S\) равно рангу \(A_f\). Если \(f=rs=\overline r\,\overline s\), где \(S\) неприводимо вкладывается в \(A_r\), а \(\overline S\) неприводимо вкладывается в \(A_{\overline r}\), то соответствующие тела для \(S\) и \(\overline S\) изоморфны паре коммутирующих подколец из \(A_g\) с \(f=g s\overline s\).

По существу нужно доказать только утверждение о произведении ранговых чисел. Сначала пусть \(S\) вложено неприводимо, так что \(\overline S\) есть тело и реципрокно изоморфно \(B\), где \(S_A=B_t\). Тогда ранговое соотношение (3) из §4, 4 дает утверждение о произведении, потому что теперь \(f=r\), а ранг \(\overline S\) равен рангу \(B\). Если \(S\) вложено приводимо, то есть \(\overline S=B_s\), то \(f=rs\), и умножение (3) на \(s^2\) дает утверждение. Отсюда непосредственно следует, что \(\overline S\) состоит из всей совокупности элементов, поэлементно коммутирующих с \(S\), и затем все дальнейшее из пункта 3, кроме последнего утверждения. Последнее получается двукратным переходом к неприводимому вложению. В \(A_r\), где \(f=rs\), система \(S\) вложена неприводимо, и \(S\) и \(R\) становятся попарно коммутирующими системами, причем \(R\) изоморфно телу \(B\). Далее, вследствие реципрокности \(S\) и \(\overline S\), \(S\) равно \(\overline B_s\), где \(\overline B\) имеет для \(\overline S\) то же значение, что \(B\) для \(S\); значит \(R\) приводимо вложено в \(A_r\), но неприводимо в \(A_g\) с \(r=g\overline s\). Здесь попарно коммутирующие системы становятся изоморфными \(B\) и \(\overline B\).

\subsection*{§6. Теория Галуа простых систем}

Усиленная теорема о коммутантах из §5, 4 выражает теорию Галуа простых систем относительно центра как основной области. Сначала рассмотрим тела, затем общие простые системы.

\paragraph{1. Теория Галуа тел конечного ранга над центром.}
Группа Галуа \(\Gcal\) тела \(A\) определяется как группа всех автоморфизмов, являющихся продолжением тождественного автоморфизма центра \(P\), то есть по §5, 2 как группа всех внутренних автоморфизмов. Следовательно,
\[
        \Gcal \simeq A^*/P^*,
\]
где \(A^*\), соответственно \(P^*\), обозначают мультипликативную группу ненулевых элементов из \(A\), соответственно из \(P\). Подгруппам \(\mathfrak H\) группы \(\Gcal\) поэтому взаимно однозначно соответствуют подгруппы \(H^*\) группы \(A^*\), содержащие \(P^*\), посредством
\[
        \mathfrak H \sim H^*/P^* .
\]
Подгруппа \(\mathfrak H\) группы \(\Gcal\) называется замкнутой, если \(H^*\) после добавления нуля переходит в тело \(H\). Факт, что \(C\) допускает группу \(\mathfrak H\), равносилен тому, что \(C\) поэлементно коммутирует с \(H\). Тем самым теорема о коммутантах §5, 4 переходит в

\paragraph{Главная теорема теории Галуа некоммутативных тел.}
Тела \(C\) между \(P\) и \(A\) и замкнутые подгруппы \(\mathfrak H\) группы \(\Gcal\) взаимно однозначно соответствуют друг другу так, что \(C\) становится полным телом инвариантов для \(\mathfrak H\), а \(\mathfrak H\) полной группой инвариантов для \(C\).

\paragraph{2. Теория Галуа простых систем.}
Если \(A_f\) есть простая система с центром \(P\), то группа Галуа \(\Gcal\) снова является группой внутренних автоморфизмов, то есть изоморфна
\[
        A_f^*/P^*,
\]
где теперь \(A_f^*\) обозначает мультипликативную группу регулярных элементов из \(A_f\). Подгруппа \(\mathfrak H\sim H^*/P^*\) группы \(\Gcal\) называется просто-замкнутой, если подкольцо \(H\), порожденное \(H^*\), есть простая система и если \(H^*\) состоит из всей совокупности регулярных элементов \(H\). Переход от теоремы о коммутантах к теории Галуа здесь дается следующей

\paragraph{Лемма.}
Каждая простая система из \(A_f\), содержащая \(P\), порождается группой \(H^*\) своих регулярных элементов. Это ясно, если \(P\) имеет бесконечно много элементов; ведь ``общий элемент'', образованный с неопределенными, регулярен, и подходящей специализацией можно образовать регулярные базисные элементы над \(P\). Но по Шоде лемма верна вообще.\footnote{Работа Shoda, цитированная в прим. 3. Там введение неопределенных избегается.}

Благодаря лемме коммутативность \(S\) с простой системой \(\mathfrak S\) снова равносильна тому, что \(S\) допускает автоморфизмы, индуцированные регулярными элементами \(\mathfrak S\). Следовательно, и здесь теорема о коммутантах §5, 4 переходит в

\paragraph{Главная теорема теории Галуа простых систем.}
Простые системы \(S\) из \(A_f\), содержащие \(P\), и просто-замкнутые подгруппы \(\mathfrak H\) группы \(\Gcal\) взаимно однозначно соответствуют друг другу так, что \(S\) становится полной областью инвариантов для \(\mathfrak H\), а \(\mathfrak H\) полной группой инвариантов для \(S\).

\paragraph{3. Доказательство посредством принципа продолжения.}
Для случая тела я даю еще второе доказательство главной теоремы, основанное на принципе продолжения изоморфизмов, то есть представлений первой степени, и, поскольку оно обходится без подсчета рангов, оно требует меньше предположений конечности. В частности, так видно, в каком направлении будет возможен перенос на тела \(S\) бесконечного ранга. Доказательство также не использует факта, что существует лишь один неприводимый реципрокный класс представлений, и потому так же верно в коммутативном случае; ср. §8, 2. Сначала я привожу несколько лемм.

\paragraph{Предположения.}
Пусть простая система \(S\) над \(P\) допускает реципрокное представление первой степени в \(A\), то есть является телом; при этом \(A\) может иметь конечный или бесконечный ранг над своим центром \(P\). Пусть разложение \(S_A\) на \(n\) простых, операторно изоморфных правых идеалов по §4, 3 задано через
\[
        S_A=r_1+\cdots+r_n
            =e_1S_A+\cdots+e_nS_A
            =e_1A+\cdots+e_nA ;
\]
реципрокное представление, порожденное \(e_i\), следовательно, определяется через
\[
        e_i\sigma=e_i\sigma_i .
\]

\paragraph{Лемма 1.}
\(n\) реципрокных представлений, порожденных \(e_1,\ldots,e_n\), различны, следовательно это разные представления одного и того же класса представлений. Поэтому \(S\) имеет по меньшей мере столько разных представлений, каков ее ранг.

Для доказательства продолжим представления из \(S\) по концу §2 до соответствий из \(S_A\), операторно-гомоморфных над \(A\). Эти соответствия здесь определены через
\[
        e_i w=e_i\omega
\]
с \(\omega\in A\) для каждого \(w\in S_A\). Если бы \(e_i\) и \(e_j\), \(i\ne j\), порождали одно и то же представление, то должно было бы быть \(e_iw=e_jw\) для каждого \(w\in S_A\). Это невозможно из-за \(e_ie_i=e_i\) и \(e_je_i=0\).

\paragraph{Лемма 2.}
Если \(T\) есть тело между \(P\) и \(S\), а \(s\) ранг \(S\) над \(T\), то каждый реципрокный изоморфизм \(T\) в \(A\) допускает по меньшей мере \(s\) различных продолжений.

Пусть реципрокный изоморфизм, то есть реципрокное представление \(T\) в \(A\), осуществляется разложением
\[
        T_A=E_1T_A+\cdots+E_hT_A
            =E_1A+\cdots+E_hA,
        \qquad E_it=E_i\tau ,
\]
где \(E_i\) идемпотенты. Это не ограничивает общности, потому что представление, порожденное базисным элементом \(E_i\alpha\), порождается также \(\alpha^{-1}E_i\alpha\), то есть идемпотентом. Правое умножение на \(S\) дает
\[
        S_A=E_1S_A+\cdots+E_hS_A .
\]
При этом все \(E_iS_A\) имеют ранг \(s\) над \(A\). Ведь ранг \(E_iS_A\) не превосходит \(s\), как показывает подстановка \(T\)-базиса в \(S\) благодаря коммутативности \(S\) и \(A\):
\[
        S=Tw_1+\cdots+Tw_s,
\]
\[
        E_iS_A=E_iT_Aw_1+\cdots+E_iT_Aw_s
              =E_iw_1A+\cdots+E_iw_sA .
\]
Поскольку сумма \(S_A\) имеет ранг \(n=hs\), каждое \(E_iS_A\) имеет точно ранг \(s\). Следовательно,
\[
        E_iS_A=r_{i1}+\cdots+r_{is}
              =e_{i1}A+\cdots+e_{is}A,
\]
где \(s\) представлений, порожденных \(e_{i1},\ldots,e_{is}\), по лемме 1 все различны. Все они являются продолжениями представления \(T\), порожденного \(E_i\); ведь из \(E_it=E_i\tau\) следует
\[
        (e_{i1}+\cdots+e_{is})t=(e_{i1}+\cdots+e_{is})\tau
\]
и тем самым \(e_{ij}t=e_{ij}\tau\) из-за разложения в прямую сумму.

\paragraph{Определение.}
Разбиение на классы \(\mathfrak I_T\) реципрокных изоморфизмов \(\mathfrak I\) системы \(S\) на \(A\), индуцированное \(T\), определяется так: все и только те изоморфизмы из \(\mathfrak I\) считаются эквивалентными, которые на \(T\) индуцируют один и тот же изоморфизм.

\paragraph{Главная теорема.}
Если \(\mathfrak I_T\) есть разбиение на классы, индуцированное \(T\), то \(T\) является максимальным подтелом относительно этого разбиения.

Ведь если \(Z\) есть промежуточное тело между \(T\) и \(S\), степени \(l\) над \(T\), то по лемме 2 каждый класс из \(\mathfrak I_T\) распадается по меньшей мере на \(l\) классов из \(\mathfrak I_Z\). Переход от этой формулировки главной теоремы к обычной возникает так: множество изоморфизмов \(\mathfrak I\), после умножения его элементов на обратный к какому-нибудь фиксированному элементу, переходит в группу автоморфизмов, причем класс из \(\mathfrak I_T\) переходит в группу инвариантов \(T\). Утверждение, что замкнутые подгруппы являются полными группами инвариантов, содержится здесь; нужно только для \(\mathfrak H\sim H^*/P^*\) взять тело \(H\) как промежуточное тело \(T\).

\subsection*{§7. Класс подобных алгебр. Поля разложения}

С этого момента речь идет только о телах \(A,\overline A\) и матричных кольцах \(A_f\), имеющих конечный ранг над своим центром \(P\); то есть о простых нормальных алгебрах над \(P\), которые кратко будем называть алгебрами над \(P\) или просто алгебрами. Тогда \(A,\overline A\) являются алгебрами с делением. В один класс подобных алгебр \(A,A_f\) объединяются все \(A_f\) с одним и тем же соответствующим \(A\). Изоморфные \(A\) при этом отождествляются.

\paragraph{1. Группа классов алгебр.}
Если \(A,B\) алгебры над \(P\), то то же верно для их произведения, прямого произведения над \(P\); ведь по концу §3 получается
\[
        A_r\times_P B_s=C_t,
\]
где \(C\) есть с точностью до изоморфии однозначно определенная алгебра с делением с центром \(P\). Умножение \(A_r,B_s\) на матричные кольца над \(P\) показывает, что это умножение однозначно определено для классов, которые тем самым образуют мультипликативно замкнутую коммутативную систему. Кроме того верна

\paragraph{Теорема Р. Брауэра.}
Классы подобных алгебр относительно прямого произведения образуют абелеву группу.

Действительно, система имеет единичный класс, состоящий из матричных колец над \(P\), то есть алгебр, подобных \(P\). Она имеет также для каждого класса \((A)\) обратный \((A)^{-1}\), состоящий из алгебр, подобных \(\overline A\), где \(\overline A\) реципрокно изоморфна \(A\). Это следует из теоремы о коммутантах §5, 3 через специализацию \(S=A\). Ведь \(A\) можно неприводимо вложить в саму \(A\); тогда получается \(B=P\) и \(A_A=P_t\).\footnote{Более тонкие теоремы о группе классов алгебр используют теорию фактор-систем; здесь мы на этом не останавливаемся. Отметим, что до этого места не возникало никаких соображений о коммутативных полях; например, тот факт, что ранг \((A:P)\) есть квадрат, нигде не использовался и не доказывался.}

\paragraph{2. Поле разложения класса алгебр.}
Теория полей разложения снова полностью основана на принципе, высказанном в начале §5: коммутативные поля расширения представляются в \(A\), соответственно в \(\overline A\). Предварительно приведем

\paragraph{Лемма.}
Если \(A\) тело с центром \(P\), а \(Z\) конечное коммутативное поле расширения \(P\), то \(A_Z\) просто с центром \(Z\). Ведь это по §4, 1 верно для системы \(Z_{\overline A}\), реципрокно изоморфной \(A_Z\). Для каждого класса \((A)\) алгебр над \(P\), подобных \(A\), поле \(Z\) поэтому порождает класс расширения \((A)_Z\) простых алгебр над \(Z\), подобных \(A_Z\).

Соответствующая алгебра с делением \(D\) класса расширения немедленно получается из теоремы о коммутантах §5, 3:

\paragraph{Теорема о соответствующей алгебре с делением класса расширения.}
Пусть \((A)_Z\) есть класс расширения, а \(Z\) обозначает неприводимое вложение, то есть представление, \(Z\) в \(A_f\). Тогда
\[
        (A)_Z=(D),
\]
где \(D\) состоит из всей совокупности элементов \(A_f\), поэлементно коммутирующих с \(Z\). Нужно лишь в §5, 3 заменить тамошнюю простую систему \(S\) на \(Z\), учитывая, что \(Z_A\) и \(A_Z\) реципрокно изоморфны. Если речь идет о приводимом вложении \(Z\), то совокупность элементов, поэлементно коммутирующих с \(Z\), дает матричное кольцо над \(D\).

Коммутативное поле расширения \(Z\) поля \(P\), как известно, называется полем разложения класса \((A)\), если класс расширения \((A)_Z\) полностью расщепляется, то есть становится равным единичному классу над \(Z\), классу всех матричных колец над \(Z\).

\paragraph{Теорема о характеристике полей разложения.}
Конечное коммутативное поле расширения \(Z\) поля \(P\) является полем разложения класса \((A)\) тогда и только тогда, когда его неприводимое вложение в \(A_f\) дает максимальное коммутативное подтело \(A_f\).

Ведь свойство \(Z\) быть полем разложения для \((A)\) равносильно \((D)=(Z)\). Следовательно, по теореме о соответствующей алгебре с делением неприводимое вложение \(Z\) должно быть максимальным коммутативным подтелом. Если это условие выполнено, то необходимо \(D=Z\), потому что каждый \(a\) из \(D\), присоединенный к \(Z\), порождает коммутативное поле расширения.

\paragraph{Замечания.}
1. Отсюда одновременно следует, что максимальное коммутативное подтело \(Z\) при неприводимом вложении порождает максимальное коммутативное подкольцо. Ведь совокупность элементов, поэлементно коммутирующих с \(Z\), образует тело.

2. Теорема одновременно дает существование полей разложения; ведь соответствующая алгебра с делением \(A\) наверняка имеет максимальные коммутативные подтела.

\paragraph{3. Ранговые соотношения, индекс, бесконечные коммутативные расширения.}
Ранговое утверждение из §5, 4 сначала дает известный факт, что ранг простой алгебры над центром есть квадрат. Ведь если \(Z\) максимально коммутативно в \(A\), то, поскольку в \(A\) каждое вложение неприводимо, получается
\[
        (Z:P)(Z:P)=(A:P),
\]
следовательно
\[
        (A:P)=m^2,\qquad (Z:P)=m,
\]
где \(m\) есть индекс Шура, который одновременно равен степени всех полей разложения, вложимых в саму алгебру с делением \(A\). Далее для степени \(n\) общего поля разложения получаем
\[
        n=mr,
\]
например из-за \((A_r:P)=m^2r^2\), или также по ранговому соотношению §4, 3; последнее потому, что индекс \(m\) также равен числу простых компонент \(A_Z\), которое становится матричным кольцом над \(Z\) ранга \(m^2\). Более общо, если \(A_L\sim D\), \(L\) есть неприводимое вложение \(A\) в \(A_r\), а \(d^2\) есть ранг \((D:L)\) алгебры с делением \(D\) над ее центром \(L\), то
\[
        l\,d=mr .
\]
Ведь по §5, 4 и теореме о соответствующей алгебре с делением из пункта 2 имеем
\[
        (A_r:P)=(L:P)(D:P),
\]
следовательно
\[
        m^2r^2=l^2d^2 .
\]

Из существования полей разложения для бесконечных коммутативных, алгебраических или трансцендентных, полей расширения \(\Omega\) поля \(P\) следуют факты: класс расширения \((A)_\Omega\) есть класс простых алгебр с центром \(\Omega\). В частности, если \(\Omega\) алгебраически замкнуто, то \(A_\Omega\) становится матричным кольцом степени \(m\). Следовательно, индекс равен ``абсолютному числу компонент''. Ведь если \(Z\) поле разложения \(A\), а \(\Omega'\) композитум \(\Omega\) и \(Z\), то \(A_{\Omega'}\) есть матричное кольцо над \(\Omega'\), то есть без радикала и с центром \(\Omega'\). Поэтому также \(A_\Omega\) не имеет радикала и имеет центр \(\Omega\), потому что элемент \(A_\Omega\), отличный от \(\Omega\), лежал бы в центре \(A_{\Omega'}\), а значит был бы телом. Следовательно, \(A_\Omega\) есть простая алгебра над \(\Omega\), наверняка матричное кольцо над \(\Omega\), когда \(\Omega\) алгебраически замкнуто.\footnote{Вообще существуют также \glqq{}минимальные\grqq{} поля разложения, то есть такие, для которых никакое собственное подполе не является полем разложения, всех возможных степеней. Ср. Brauer--Noether, работу, цитированную в прим. 2.}

\paragraph{4. Существование сепарабельных полей разложения.}
Каждый класс \((A)\) имеет сепарабельные поля разложения, даже такие, которые вкладываются в саму \(A\).\footnote{Ср. G. Köthe, \foreign{Über Schiefkörper mit Unterkörpern zweiter Art über dem Zentrum}, Journ. f. Math. 166 (1932), S. 182--184. Приведенное здесь значительно более простое доказательство восходит к замечанию M. Zorn.}

Пусть основное поле \(P\) имеет характеристику \(p\), а индекс
\[
        m=s p^f
\]
с \(s\), взаимно простым с \(p\). Если тогда \(Z\) есть поле разложения \(A\) степени \(m\), а \(Z_0\) содержащееся в нем расширение первого рода, то есть
\[
        (Z:Z_0)=p^f,
\]
то индекс \(A_{Z_0}\sim D\) равен \(p^f\). Нужно доказать существование сепарабельного поля расширения \(Z_1\) поля \(Z_0\), такого, что индекс \(D_{Z_1}\) становится собственным делителем \(p^f\). Но \(D_\Omega\) с алгебраически замкнутым \(\Omega\) по пункту 3 не имеет радикала; следовательно, приведенная дискриминанта \(D\) не исчезает. Поэтому существует по крайней мере один элемент \(d\) из \(D\) с ненулевым приведенным следом; при этом \(d\) не может уже лежать в основном поле \(Z_0\), потому что след каждого элемента \(\alpha\) из \(Z_0\) равен \(p^f\alpha\), а значит исчезает. Поэтому \(Z_1=Z_0(d)\) есть собственное, и притом сепарабельное, поле расширения; индекс \(D_{Z_1}\) становится собственным делителем \(p^f\). Конечное повторение приводит к сепарабельному полю разложения класса \((A)\), степень которого \(m\) совпадает с индексом \(A\).

\subsection*{§8. Поля разложения, поля отщепления и теория Галуа в коммутативном случае}

Чтобы перейти от полей разложения систем, простых над своим центром, к полям разложения произвольных простых систем, нужно еще развить теорию разложения центра и добавить параллельную §6, 3 теорию Галуа. Все поля и системы, встречающиеся в этом параграфе, коммутативны.

\paragraph{1. Поля разложения и поля отщепления коммутативных простых систем.}
Пусть коммутативная система \(Z\) проста над \(P\), то есть является полем, а \(\Omega\) алгебраически замкнутое поле расширения \(P\). Тогда, как известно, \(Z_\Omega\) остается системой без радикала тогда и только тогда, когда \(Z\) сепарабельно, то есть является расширением первого рода, над \(P\). Ведь тогда и только тогда оно имеет столько изоморфных отображений первой степени в \(\Omega\), каков его ранг над \(P\), а значит столько же различных представлений, которые здесь совпадают с классами представлений.

Факт, что в \(Z_\Omega\) может появиться радикал, то есть не обязательно происходит разложение в прямую сумму абсолютно простых компонент, приводит к следующим определениям. Поле расширения \(\Lambda\) поля \(P\) называется полем разложения, если \(Z_\Lambda\) распадается на композиционные факторы первой степени; другими словами, если все абсолютно неприводимые представления уже лежат в \(\Lambda\). Поле расширения \(T\) поля \(P\) называется полем отщепления, если \(Z_T\) отщепляет по крайней мере один композиционный фактор первой степени; другими словами, если существует по крайней мере одно абсолютно неприводимое представление \(Z\) в \(T\).

Поля разложения и поля отщепления характеризуются следующей

\paragraph{Теорема.}
Поле \(T\) является полем отщепления тогда и только тогда, когда оно содержит подполе \(Z'\), изоморфное \(Z\); поле \(\Lambda\) является полем разложения тогда и только тогда, когда оно содержит подполе \(\Gamma\), изоморфное полю Галуа, принадлежащему \(Z\).

Это прямо следует из определения полей разложения и отщепления через абсолютно неприводимые представления. В частности, по аналогии с некоммутативным случаем получается следствие: поле \(Z'\) является минимальным полем отщепления, то есть никакое собственное подполе не является полем отщепления, тогда и только тогда, когда оно изоморфно \(Z\). Минимальное поле отщепления является одновременно полем разложения тогда и только тогда, когда \(Z\) нормально, то есть является полем Галуа над \(P\).

В этом причина называть простую алгебру нормальной над ее центром. Факт, что внутри фиксированного алгебраически замкнутого \(\Omega\) над \(P\) существует лишь одно минимальное поле разложения, именно поле Галуа, принадлежащее \(Z\), в отличие от вообще говоря бесконечно многих неизоморфных в некоммутативном случае,\footnote{R. Brauer--E. Noether, указанная работа.} имеет свою причину в единственном разложении в прямую сумму в коммутативном случае, в отличие от разложения, единственного лишь с точностью до операторной изоморфии в некоммутативном случае; или, что то же самое, в лишь конечном числе различных абсолютно неприводимых представлений вместо бесконечно многих различных в некоммутативном случае, которые, однако, принадлежат одному и тому же классу.

\paragraph{2. Гиперкомплексное обоснование теории Галуа в случае сепарабельных полей \(Z/P\).}
В точной аналогии с §6, 3 здесь теория изоморфизмов действует для произвольных \(Z\), сепарабельных над \(P\); в случае поля Галуа \(Z\) можно затем перейти к обычной группе автоморфизмов.

\paragraph{Предположения.}
Пусть простая система \(Z\) над \(P\) есть конечное сепарабельное поле расширения, \(\Omega\) обозначает конечное или бесконечное поле разложения над \(P\), \(Z'=Z^{(1)}\) подполе, изоморфное \(Z\), а \(\Gamma\) соответствующее поле Галуа. По пункту 1 имеем прямое разложение
\[
        Z_\Omega=r^{(1)}+\cdots+r^{(n)}
        =e^{(1)}Z_\Omega+\cdots+e^{(n)}Z_\Omega
        =e^{(1)}\Omega+\cdots+e^{(n)}\Omega .
\]
Представление \(Z\to Z^{(i)}\) с \(Z^{(i)}\subseteq\Gamma\), порожденное \(e^{(i)}\), следовательно, определяется через
\[
        e^{(i)}z=e^{(i)}z^{(i)}.
\]

\paragraph{Лемма 1.}
\(n\) представлений, порожденных \(e^{(1)},\ldots,e^{(n)}\), все различны. Доказательство такое же, как в §6, 3, или также по пункту 1; ведь они принадлежат различным классам представлений.

\paragraph{Лемма 2.}
Если \(T\) есть промежуточное поле между \(P\) и \(Z\), а \(s\) ранг \(Z\) над \(T\), то каждый изоморфизм \(T\) в \(\Omega\) допускает по меньшей мере, а потому ровно, \(s\) продолжений.

Доказательство такое же, как в §6, 3. То, что здесь ``по меньшей мере'' усиливается до ``ровно'', следует из однозначного разложения \(Z_\Omega\), по которому \(Z\) имеет в \(\Omega\) не по меньшей мере, а ровно \(n\) изоморфизмов.

Как в §6, 3, теперь определяют разбиение на классы \(\mathfrak I_T\), индуцированное \(T\), конечного множества \(\mathfrak I\) изоморфизмов \(Z\): все и только те изоморфизмы из \(\mathfrak I\) считаются эквивалентными в \(\mathfrak I_T\), которые на \(T\) индуцируют один и тот же изоморфизм. Доказывается

\paragraph{Главная теорема.}
Если \(\mathfrak I_T\) есть разбиение на классы, индуцированное \(T\), то \(T\) является максимальным подполем относительно этого разбиения.

Отсюда для случая поля Галуа \(Z\) можно, как в §6, 3, перейти композицией с реципрокным к одному изоморфизму к группе автоморфизмов и обычной формулировке. Соответствующее утверждение для подгрупп вытекает из рассмотрения идемпотентов, которое интересно само по себе.

\paragraph{3. Идемпотенты \(Z_\Omega\).}
Пусть \(S\) пробегает \(n\) подстановок, переводящих \(Z\) в отдельные \(Z^{(i)}\), то есть композицию изоморфизмов \(Z\to Z_1\) и \(Z\to Z^{(i)}\), где первый означает изоморфизм, реципрокный к \(Z\to Z_1\). \(Z\) не обязано быть полем Галуа. Для элементов \(a\) из \(Z_\Omega\) подстановки \(S\) определяются как операторы условием, что они индуцируют тождественность на \(Z\); для каждого \(a\) тем самым определен сопряженный элемент \(a^S\), лежащий в \(Z_{(i)}\). При этих условиях верна

\paragraph{Теорема.}
\(n\) идемпотентов \(e^{(i)}\) из \(Z_\Omega\) сопряжены:
\[
        e^{(i)}=e^S,\qquad e=e^{(1)};
\]
они лежат в \(n\) сопряженных системах расширения \(Z\Omega\).

Действительно, применение \(S\) переводит определяющие соотношения
\[
        ez=ez^{(1)},\qquad e^2=e
\]
в
\[
        e^S z=e^S z^{(i)},\qquad (e^S)^2=e^S .
\]
Из-за однозначности разложения идемпотенты также однозначно определены; следовательно, \(e^S=e^{(i)}\). Кроме того, поскольку \(Z\) есть поле отщепления, то есть уже осуществляет представление \(ez=ez^{(1)}\), \(e\) уже лежит в \(Z\Omega\), а тем самым \(e^S\) лежит в \(Z^S\Omega\).

Если специально \(Z\) есть поле Галуа, то отсюда непосредственно следует часть главной теоремы теории Галуа, касающаяся подгрупп. Ведь если \(\mathfrak H\) есть подгруппа группы Галуа \(\Gcal\), то, из-за однозначной определенности компонент прямой суммы, элемент
\[
        \sum_{H\in\mathfrak H} e^H
\]
допускает все подстановки из \(\mathfrak H\), но никакие другие. Поскольку группа \(\Gcal\) была определена для \(Z\), речь идет о теории Галуа \(Z\); благодаря изоморфии одновременно закончен и случай для \(Z'\).

\paragraph{4. Связь идемпотентов с комплементарными базисами.}
\(Z\) снова не обязательно предполагается полем Галуа.

\paragraph{Теорема.}
Пусть \(a_1,\ldots,a_n\) есть базис \(Z\) над \(P\), и пусть идемпотент задан как
\[
        e=a_1\beta_1+\cdots+a_n\beta_n
\]
то есть
\[
        e^S=a_1\beta_1^S+\cdots+a_n\beta_n^S,
\]
Тогда
\[
        \beta_1^S,\ldots,\beta_n^S
\]
является образом комплементарного базиса \(b_1,\ldots,b_n\) к \(a_1,\ldots,a_n\) в отображении, осуществленном \(e^S\).

Действительно, для отображения, осуществленного \(e^S\), \(e^Sz=e^S\zeta^S\), в частности \(e^Sa_i=e^S\alpha_i^S\), благодаря
\[
        e^Se^S=e^S,\qquad e^Se^R=0\quad(S\ne R)
\]
действуют соответствия \(e^S\mapsto1\), \(e^R\mapsto0\). Следовательно, получается
\[
        1=\alpha_1^S\beta_1^S+\cdots+\alpha_n^S\beta_n^S,\qquad
        0=\alpha_1^S\beta_1^R+\cdots+\alpha_n^S\beta_n^R\quad(S\ne R).
\]
В матрической записи это есть определяющее уравнение комплементарных базисов:
\[
\begin{pmatrix}
\alpha_1 & \cdots & \alpha_n\\
\cdots & \cdots & \cdots\\
\alpha_1^S & \cdots & \alpha_n^S\\
\cdots & \cdots & \cdots
\end{pmatrix}
\cdot
\begin{pmatrix}
\beta_1 & \cdots & \beta_1^R & \cdots\\
\cdots & \cdots & \cdots & \cdots\\
\beta_n & \cdots & \beta_n^R & \cdots
\end{pmatrix}
=
E.
\]
Переход к определению следа следует из того, что для
\[
        a_i=\sum_S e^S\alpha_i^S,\qquad b_i=\sum_S e^S\beta_i^S
\]
матричное соотношение после перестановки строк переходит в
\[
        \Sp(a_i b_i)=1,\qquad \Sp(a_i b_k)=0\quad\text{для }i\ne k,
\]
причем \(a_i\) и \(b_i\) лежат в \(Z\), потому что они допускают все подстановки \(S\).\footnote{Ср. \foreign{Darstellungstheorie} §25.}\footnote{Связи между комплементарными базисами и идемпотентами впервые встречаются у Dedekind, \foreign{Zur Theorie der aus \(n\) Haupteinheiten gebildeten komplexen Größen}; ср. Bd. II der Gesammelten Werke, S. 4--5.}

Контрагредиентность комплементарных базисов здесь следует из того факта, что идемпотенты \(e^S\) являются инвариантами \(Z\), не зависящими от положенного в основу базиса, то есть что все
\[
        a_1\beta_1^S+\cdots+a_n\beta_n^S
\]
переходят в себя при переходе к другому базису
\[
        a_1',\ldots,a_n'
\]
поля \(Z/P\).

\subsection*{§9. Поля разложения и поля отщепления произвольных систем}

\paragraph{1. Определения и сведение к простому случаю.}
Данным в §8 определениям в общем случае соответствуют следующие определения. Пусть \(S\) гиперкомплексна над \(P\). Коммутативное поле расширения \(\Lambda\) поля \(P\) называется полем разложения \(S\), если \(S_\Lambda\) по композиционному ряду из односторонних идеалов распадается на абсолютно простые композиционные факторы; другими словами, если все неприводимые представления \(S\) в \(\Lambda\) уже абсолютно неприводимы. Поле расширения \(T\) поля \(P\) называется полем отщепления, если \(S_T\) отщепляет по крайней мере один абсолютно простой композиционный фактор, то есть если по крайней мере одно представление, неприводимое в \(T\), уже абсолютно неприводимо.

Эти определения, очевидно, содержат в себе определения, данные в §8 для коммутативного случая и в §7 для простых алгебр; последнее потому, что здесь \(A_\Lambda\) при каждом коммутативном расширении центра вполне приводимо, и потому композиционные факторы совпадают с компонентами прямого разложения. Полные матричные кольца над центром, и только они, здесь дают абсолютно простые компоненты, а значит и абсолютно неприводимые представления. Поскольку далее \(A_\Lambda\) двусторонне просто, то есть распадается на операторно изоморфные компоненты, отщепление абсолютно простой компоненты уже дает полное расщепление: поля отщепления и поля разложения совпадают.

Из определений непосредственно следуют факты: поля разложения \(S\) задаются как поля объединения по одному полю разложения для каждого представления, неприводимого в \(P\); полем отщепления \(S\) является каждое поле отщепления одного представления, неприводимого в \(P\). Благодаря взаимно однозначному соответствию между простыми системами, компонентами фактор-кольца по радикалу, и представлениями, неприводимыми в \(P\), достаточно рассматривать простые системы. Здесь результаты будут получены комбинированием §7 и §8.

\paragraph{2. Поля разложения и поля отщепления простых систем.}
Сначала рассмотрим случай простой системы \(A\) с сепарабельным центром \(Z\) над \(P\). Из §8, 1 и §8, 3 получается: двустороннее разложение \(A_\Omega\), соответствующее прямому разложению \(Z_\Omega\), где \(\Omega\) обозначает поле разложения \(Z\), дает \(n\) сопряженных, изоморфных \(A\), компонент
\[
        e^S A_\Omega
\]
системы \(A_\Omega\), которые становятся простыми алгебрами над своим центром
\[
        e^S Z_\Omega=e^S Z .
\]
При этом подстановки \(S\) определяются как операторы для \(A_\Omega\) условием, что они индуцируют тождественность на \(A\).

Действительно, по §8, 3 идемпотенты сопряжены; поэтому то же по определению подстановок для \(A\) верно для компонент \(e^S A_\Omega\). Поскольку \(A\) двусторонне проста, \(e^S A_\Omega\) кольцево изоморфна \(A\). При этом \(e^SA_\Omega=e^S A_{Z^{S}}\), поскольку \(e^S Z_\Omega=e^S Z\); значит центр совпадает с областью коэффициентов, и компоненты являются нормальными алгебрами.

Результаты §7, 2 далее дают: если \(A\) есть простая система над \(P\) с сепарабельным центром \(Z\) над \(P\), то \(A_\Omega\) также не имеет радикала, то есть вполне приводимо; при этом \(\Omega\) может обозначать любое конечное или бесконечное поле расширения. Ведь по §7, 2 \(e^S A_\Omega\) остается простым при каждом расширении коэффициентов, то есть без радикала. Следовательно, то же верно для прямой суммы; она и есть \(A_\Omega\), или, если \(\Omega\) не является полем разложения \(Z\), возникает из \(A_\Omega\) расширением коэффициентов.

Далее из §7, 2 следует

\paragraph{Характеристика полей отщепления и полей разложения.}
Если \(A\) есть алгебра с делением с сепарабельным центром \(Z\), то неприводимые вложения полей отщепления задаются всеми и только теми максимальными коммутативными подтелами \(A_f\), которые содержат \(Z\); поля разложения суть все и только поля объединения поля отщепления с полем разложения центра, в частности с полем Галуа, принадлежащим центру. Минимальное поле отщепления может быть полем разложения даже тогда, когда центр не является полем Галуа.

Ведь поле отщепления должно содержать поле, изоморфное \(Z\); утверждение о вложении для полей отщепления теперь следует из предыдущих фактов и из §7, 2. Далее поле разложения должно содержать поле разложения \(\Omega\) для \(Z\). Но каждое поле отщепления над \(\Omega\) одновременно является полем разложения, поскольку компоненты \(e^S A_\Omega\) просты и имеют центр, изоморфный \(\Omega\). Если \(Z\) есть поле Галуа, то минимальные поля отщепления и разложения совпадают. Но и в негалуа-случае, в отличие от коммутативного, могут существовать минимальные поля отщепления, которые одновременно являются полями разложения, а именно как только такое минимальное поле отщепления содержит поле Галуа, принадлежащее \(Z\). Пример: тело кватернионов с центром, изоморфным \(P(\sqrt[3]{2})\), где \(P\) есть поле рациональных чисел; поле Галуа, принадлежащее \(P(\sqrt[3]{2})\), является минимальным полем отщепления и разложения.

Если, наконец, речь идет о несепарабельном центре, то \(Z_\Omega\), а с ним и \(A_\Omega\), становится системой с радикалом. Пусть, например, \(\mathfrak C\) радикал \(Z_\Omega\), а \(\mathfrak c\) идеал расширения в \(A_\Omega\). Тогда для \(Z_\Omega/\mathfrak C\) существует прямое разложение на сопряженные компоненты, в точной аналогии с §8, 2; это, как выше, дает прямое разложение \(A_\Omega/\mathfrak c\) так, что компоненты \(\overline e^{(i)}A\) изоморфны \(A\) с центром \(\overline e^{(i)}Z^{(i)}\). Следовательно, теоремы о полях отщепления и разложения остаются верными. Далее подсчет рангов показывает, что композиционные факторы, соответствующие композиционному ряду \(\mathfrak c\), становятся операторно изоморфными \(A_\Omega/\mathfrak c\), а не только гомоморфными.

В терминах неприводимых представлений итог таков: если задано представление гиперкомплексной системы, неприводимое в \(P\), то оно становится абсолютно вполне приводимым тогда и только тогда, когда соответствующий центр сепарабелен над \(P\). В каждом случае поля разложения и отщепления представления можно охарактеризовать через вложение в соответствующую простую систему, как выше. В частности, наименьшая степень поля отщепления равна произведению степени центра на индекс соответствующего класса алгебр над центром; степень каждого поля отщепления кратна этому числу. Представление, в смысле композиционных рядов, распадается на столько различных, но сопряженных, абсолютно неприводимых классов представлений, какова степень наибольшего сепарабельного поля, содержащегося в центре. Каждый класс встречается столько раз, каково произведение индекса и степени центра над этим сепарабельным расширением.

Для совершенного основного поля, где существуют только сепарабельные поля расширения, мы так возвращаемся к известным результатам И. Шура, добавляя характеристику через вложение, которая уже встречается у Р. Брауэра.

\begin{center}
(Поступило 8 июня 1932 г.)
\end{center}
% END UNIT BODY
% END PRODUCER UNIT 216
\endgroup


\clearpage
\begingroup
\setcounter{footnote}{0}
% BEGIN PRODUCER UNIT 217
% Source: C:/Users/Floris/Documents/interlanguage/03_projects/noether/02_slavic_working_corpus/translations/paper41/source_fidelity/russian/v001/Noether_Paper41_SourceFidelity_Russian_v001.tex
% Identity: 32094 B / SHA-256 487D43247BBC0FDA1A4843039BB7ABD765999EB205B13DA3B53DC1B70F959174
\setlength{\parindent}{1.1em}
\setlength{\parskip}{0pt}
\newcommand{\foreign}[1]{#1}
\providecommand{\Norm}{\operatorname{N}}
\providecommand{\Gg}{\mathfrak G}
\providecommand{\Ii}{\mathfrak I}
\providecommand{\PP}{\mathfrak P}
\providecommand{\glqq}{``}
\providecommand{\grqq}{''}
\emergencystretch=220em
\allowdisplaybreaks
\sloppy
\hbadness=10000
\vbadness=10000
\hfuzz=5pt
\vfuzz=12pt
\raggedbottom
% BEGIN UNIT BODY
\section*{41. Теорема о главном роде для относительно галуаовых числовых полей}

\begin{center}
\emph{\foreign{Math. Ann. 108 (1933), S. 411--419}}
\end{center}

В последнее время выяснилось, что некоммутативные методы, в особенности теория алгебр, дают возможность формулировать и доказывать известные теоремы об относительно циклических и относительно абелевых числовых полях для произвольных относительно галуаовых полей.\footnote{Ср. мой цюрихский доклад, \foreign{\emph{Hyperkomplexe Systeme in ihren Beziehungen zur kommutativen Algebra und Zahlentheorie}}, Verhandlungen des Internationalen Mathematiker-Kongresses Zürich 1932, Bd. I. Для теоремы о расщепляющихся алгебрах ср. также H. Hasse, \foreign{\emph{Die Struktur der R. Brauerschen Algebrenklassengruppe}}, Math. Ann. 107 (1932).} Я прежде всего напомню теорему норм, которая в общем виде выражается как теорема о расщепляющихся алгебрах. Ниже я показываю, что соответственно и теорему о главном роде можно в общем гиперкомплексно сформулировать как теорему о внутренних автоморфизмах или также о скрещенных представлениях. Доказательство получается из названной теоремы о расщепляющихся алгебрах чисто алгебро-арифметическими соображениями, тогда как в первой теореме трансцендентным ядром остается теорема норм в циклическом случае, уже для простой степени.

Для лучшего понимания формулировки я сначала привожу теорему о главном роде в минимальном смысле, которая ниже не используется,\footnote{Я говорю ``в минимальном смысле'', поскольку речь идет об алгебраическом аналоге теоремы о главном роде в произвольных полях, тогда как ``в малом'' уже имеет значение переход к \(p\)-адическому основному полю.} элементарную алгебраическую теорему,\footnote{Теорема имеется у A. Speiser, \foreign{Zahlentheoretische Sätze aus der Gruppentheorie}, Math. Z. 5 (1919), S. 1--6; там она уже высказана как теорема о скрещенных представлениях.} которая в циклическом специальном случае переходит в известную теорему 90 из гильбертова \foreign{Zahlbericht}: \(\Norm(a)\) тогда и только тогда равно единице, когда \(a=b^{1-S}\). И уже эту теорему я высказываю как теорему о внутренних автоморфизмах или о скрещенных представлениях и доказываю ее посредством общих теорем о внутренних автоморфизмах и скрещенном произведенном представлении алгебр.

В собственно теореме о главном роде вместо последней конструкции выступает скрещенное произведение группы идеальных классов и группы Галуа, но с более тонким индуцированным разбиением фактор-систем на идеальные классы; тем самым охватывается аналог разбиения на лучевые классы из теории полей классов. Так как общих теорем об автоморфизмах и представлениях здесь еще нет, теорему о главном роде можно понимать как первый шаг в этом направлении. Поэтому, как сказано выше, я даю сведение посредством теоремы о расщепляющихся алгебрах. Тем самым теорема сводится к соответствующей теореме для идеалов вместо идеальных классов; здесь теоремы Брандта о связи максимальных порядков алгебры\footnote{Теория Брандта о максимальных порядках изложена и заново обоснована у H. Hasse, \foreign{Über \(p\)-adische Schiefkörper und ihre Bedeutung für die Arithmetik hyperkomplexer Zahlsysteme}, Math. Ann. 104 (1931).} можно рассматривать как эквивалент теорем об автоморфизмах. Вместе с рассмотрениями максимальных порядков скрещенных произведений\footnote{Было объявлено по одной заметке C. Chevalley, H. Hasse и E. Noether; две последние -- в томе памяти Herbrand.} они дают доказательство, существенно параллельное минимальному доказательству, хотя и более сложное из-за возвращения к отдельным местам. Поэтому я предпочитаю почти тривиальное доказательство, указанное Э. Артином; здесь оно является еще более простым аналогом доказательства Шпайзера.

\subsection*{§1. Теорема о главном роде в минимальном случае}

В этом параграфе \(k\) обозначает произвольное коммутативное поле, не обязательно числовое; \(K/k\) есть сепарабельное галуаово расширение \(n\)-й степени, а \(\Gg\) -- соответствующая группа Галуа.

Сначала соберем известные факты о скрещенных произведениях.\footnote{Эта теория, в связи с моей лекцией (1929/30), изложена у H. Hasse, \foreign{\emph{Theory of cyclic algebras}}, Trans. Amer. Math. Soc. 34 (1932), Chap. 2; ср. также M. Deuring о гиперкомплексных числах.} Скрещенное произведение означает одновременное вложение \(K\) и \(\Gg\) в алгебру \(A\) так, что автоморфизмы \(K\) становятся внутренними. Если \(u_{S_1},\ldots,u_{S_n}\) -- символы, соответствующие \(n\) элементам группы, то сначала полагают \(A\), как модуль линейных форм ранга \(n\) над \(K\), равным
\[
        A=u_{S_1}K+\cdots+u_{S_n}K .                         \tag{1}
\]
В силу требования внутреннего автоморфизма, порожденного \(u_S\), более общо \(u_S K^*\), \(A\) становится кольцом, следовательно алгеброй ранга \(n^2\) над \(k\). Требование выражается соотношениями
\[
        u_S^{-1}zu_S=z^S,\qquad\text{или}\qquad
        z u_S=u_S z^S,                                      \tag{2}
\]
для каждого \(z\in K\), а также
\[
        u_Su_T=u_{ST}a_{S,T},\qquad a_{S,T}\in K^*,           \tag{3}
\]
и законом ассоциативности
\[
        a_{S,T}^{\,R}a_{ST,R}=a_{S,TR}a_{T,R}.                \tag{4}
\]
Если произведение двух произвольных элементов определить формулой
\[
 \Bigl(\sum_S u_S b_S\Bigr)\Bigl(\sum_T u_T c_T\Bigr)
   =\sum_{S,T}u_{ST}\,a_{S,T}\,b_S^{\,T}c_T ,
\]
то \(A\) становится скрещенным произведением \(K\) с \(\Gg\) при фактор-системе \(a_{S,T}\); обозначение
\[
        A=(a_{S,T},K,\Gg).
\]
Доказывают, что \(A\) является простой нормальной алгеброй над \(k\), то есть матричным кольцом \(D_r\) \(r\)-й степени над присоединенной алгеброй с делением \(D\), и что \(K\) является максимальным коммутативным подполем, следовательно полем расщепления. Обратно, для каждой такой нормальной алгебры с делением \(D\) существуют матричные кольца \(D_r\), которые таким образом порождаются как скрещенные произведения.

Если перейти от \(u_S\) к \(v_S=u_Sc_S\), где \(c_S\in K^*\), возникают ассоциированные фактор-системы
\[
        \bar a_{S,T}
        =a_{S,T}\frac{c_S^{\,T}c_T}{c_{ST}} .                 \tag{5}
\]
В частности, по (5) фактор-системы
\[
        \frac{c_S^{\,T}c_T}{c_{ST}}
\]
ассоциированы с единицей; эти величины называются трансформационными величинами. Ассоциированные фактор-системы объединяются в класс \((a)\). Классы \((a)\) взаимно однозначно соответствуют прямым произведенным построениям одной абелевой группы; относительно прямого произведения они образуют абелеву группу, изоморфную покомпонентному произведению классов фактор-систем. Единичный элемент -- это класс расщепляющихся алгебр, то есть система всех трансформационных величин. Речь идет о группе Р. Брауэра классов алгебр.

Для циклических алгебр получается специальный случай: если \(Z/k\) циклично и \(S\) -- порождающая подстановка, то степеням \(S\) соответствуют степени \(u\), и выходит
\[
        A=Z+uZ+\cdots+u^{n-1}Z,                              \tag{1'}
\]
\[
        zu=uz^S,\qquad u^n=\alpha,                            \tag{2'--3'}
\]
где
\[
        \alpha\in k^*,\qquad
        \bar\alpha=\alpha\,\Norm(c)\quad\text{для }v=uc .      \tag{4'--5'}
\]
Каждая фактор-система здесь, таким образом, состоит из одного элемента основного поля; обозначение \(A=(\alpha,Z,S)\). Единичный класс задается нормами из \(Z^*\), а группа классов алгебр изоморфна \(k^*/\Norm(Z^*)\).

Формулировка теоремы о главном роде в минимальном случае использует тот факт, что соотношения (2)--(5) чисто мультипликативны и тем самым одновременно определяют группу расширения \(\Gg^*\) группы \(\Gg\), состоящую из элементов \(n\) комплексов \(u_S K^*\):
\[
        \Gg^*=\{\,u_{S_1}K^*,\ldots,u_{S_n}K^*\,\}.            \tag{6}
\]
\(\Gg^*\) имеет \(K^*\) в качестве абелевой нормальной подгруппы, тогда как \(\Gg^*/K^*\simeq\Gg\). Группа \(\Gg^*\) состоит из всех регулярных элементов \(g^*\), которые переводят \(K\) в себя, то есть \(g^{*-1}Kg^*=K\). Поэтому кольцевые автоморфизмы \(K\) порождаются всеми и только элементами \(\Gg^*\). Действительно, если \(v=uw\) порождает внутренний автоморфизм \(K\), то \(w\) перестановочно со всеми элементами \(K\), а значит лежит в \(K^*\), поскольку \(K\) является максимальным коммутативным подкольцом.

\paragraph{Теорема о главном роде в минимальном случае.}
Первая формулировка: всякий групповой автоморфизм \(\Gg^*\), продолжающий тождественный автоморфизм \(K^*\), является внутренним и порождается элементом из \(K^*\). Другими словами: если трансформационные величины
\[
        \frac{c_S^{\,T}c_T}{c_{ST}}
\]
все имеют значение единица, то \(c_S\) являются символическими \((1-S)\)-ми степенями,
\[
        c_S=b^{1-S}=\frac{b}{b^S}\qquad(S\in\Gg).
\]
Третья формулировка: группа \(\Gg\) имеет только один класс скрещенных представлений первой степени в \(K^*\), принадлежащий фактор-системе единица.

Представление \(u_S\mapsto C_S\) называется скрещенным с фактор-системой \(a_{S,T}\), если
\[
        C_S^{\,T}C_T=C_{ST}a_{S,T}
\]
выполнено. Два представления принадлежат одному классу, если
\[
        C_S=B^{-S}D_SB .
\]

Я доказываю первую формулировку, а затем показываю ее совпадение со второй и третьей. Доказательство основано на том, что всякий групповой автоморфизм указанного вида продолжается до кольцевого автоморфизма \(A\), который, как известно, является внутренним автоморфизмом. Пусть при этом автоморфизме \(v_S\) -- образы \(u_S\). Тогда \(\sum v_SK\) снова образует скрещенное произведение \(\Gg\) с \(K\) для той же фактор-системы, потому что по предположению все соотношения (2)--(4) сохраняются. Соответствие
\[
        \sum u_Sb_S\longmapsto \sum v_Sb_S
\]
задает кольцевой автоморфизм, при котором элементы \(K\) остаются неподвижными. Поэтому он порождается элементом из \(K^*\).

Переход ко второй формулировке основан на том, что по (2) элементы \(v_S\) должны иметь вид \(v_S=u_Sc_S\); поскольку фактор-системы также сохраняются, по (5) соответствующие трансформационные величины должны быть равны единице. Обратно, (2)--(5) показывают, что подстановка \(v_S=u_Sc_S\) порождает автоморфизм указанного вида. Следовательно, из первой формулировки получаем
\[
        v_S=b^{-1}u_Sb=u_S\,\frac{b}{b^S},
        \qquad c_S=b^{1-S}.
\]
В циклическом случае предположение, в частности, означает \(\Norm(c)=1\), и \(c=b^{1-S}\) есть известная теорема. Третья формулировка непосредственно равносильна второй, ибо предположение говорит, что \(u_S\mapsto c_S\) образует скрещенное представление первой степени с фактор-системой единица; \(c_S=b^{1-S}\) означает эквивалентность единичному представлению.

\subsection*{§2. Теорема о главном роде}

\paragraph{Предварительные замечания о разбиении на классы.}
В теореме о главном роде из теории полей классов для циклических относительных полей, как известно, встречаются два различных разбиения на классы: абсолютные идеальные классы в верхнем поле и лучевые классы в нижнем поле. В общем случае этому соответствует то, что разбиение идеалов верхнего поля на классы индуцирует более тонкое разбиение фактор-систем на идеальные классы.

Пусть сначала \(\Ii\) обозначает группу всех идеалов из \(K\). Так как \(\Gg\) порождает автоморфизмы в \(\Ii\), расширение с \(\Gg\) определяется соотношениями (2)--(5) и состоит из \(n\) комплексов \(\{\,\cdots u_S\Ii\cdots\,\}\). Если в \(\Ii\) перейти посредством установления эквивалентности к группе абсолютных идеальных классов, то и теперь определены \(n\) комплексов \(u_S\bar\Ii\). Ведь \(\Gg\) порождает также автоморфизмы группы классов, поскольку главный класс переходит в себя. Тем самым отношение эквивалентности, ведущее от \(\Ii\) к \(\bar\Ii\), переносится с \(\{\,\cdots u_S\Ii\cdots\,\}\) на \(\{\,\cdots u_S\bar\Ii\cdots\,\}\). В частности, \(u_S\) эквивалентно всем произведениям \(u_S(c_S)\), где \((c_S)\) пробегают главные идеалы.

Если теперь потребовать, чтобы произведенные соотношения (3) были однозначными в смысле этой эквивалентности, то это значит, что все трансформационные величины из главных идеалов становятся однозначно эквивалентными. Иными словами:

В системе \(n\) комплексов \(\{\,\cdots u_S\bar\Ii\cdots\,\}\), где \(\bar\Ii\) -- группа абсолютных идеальных классов из \(K\), а \(H\) в частности обозначает главный класс, произведенные соотношения (3) для \(u_SH\) тогда и только тогда однозначно определены, когда в индуцированном разбиении фактор-систем на идеальные классы главный класс содержит все трансформационные величины из главных идеалов.

\paragraph{Определение.}
Индуцированное разбиение фактор-систем на идеальные классы задается так: главный класс состоит из всех главных идеалов \((a_{S,T})\), для которых при по меньшей мере одной системе базисных элементов \(a_{S,T}\) алгебра
\[
        (a_{S,T},K,\Gg)
\]
расщепляется во всех разветвленных местах \(K\).

Эти главные идеалы образуют группу относительно композиции фактор-систем; она содержит трансформационные величины
\[
        \frac{(c_S)^{\,T}(c_T)}{(c_{ST})},
\]
ибо фактор-системы \(c_S^{\,T}c_T/c_{ST}\) порождают всюду расщепляющиеся алгебры.

\paragraph{Формулировка теоремы о главном роде.}
Теорема, соответственно минимальной теореме, высказывается в трех равносильных формах.

Первая формулировка: если при положенном в основу индуцированном разбиении фактор-систем на идеальные классы подстановка
\[
        v_S=u_S c_S
\]
порождает автоморфизм \(n\) комплексов \(\{\,\cdots u_S\bar\Ii\cdots\,\}\), то есть для \(v_S\) в смысле этого разбиения на классы выполняются те же соотношения (3), то этот автоморфизм является внутренним и порождается идеальным классом \(b\).

Вторая формулировка: если трансформационные величины, образованные из идеальных классов \(c_S\),
\[
        \frac{c_S^{\,T}c_T}{c_{ST}}
\]
все принадлежат главному классу индуцированного разбиения фактор-систем на классы, то векторы \(\{\,\cdots c_S\cdots\,\}\) из идеальных классов образуют главный род, а классы \(c_S\) являются символическими \((1-S)\)-ми степенями; то есть существует идеальный класс \(b\) такой, что
\[
        c_S=b^{1-S}=\frac{b}{b^S}\qquad(S\in\Gg).
\]

Третья формулировка: при положенном в основу индуцированном разбиении фактор-систем на идеальные классы группа \(\Gg\) имеет только один класс скрещенных представлений первой степени в группе идеальных классов, принадлежащий единичному классу фактор-систем.

Равносильность этих формулировок следует так же, как в §1. Переход от первой ко второй формулировке здесь даже проще, поскольку автоморфизм уже предполагается порожденным \(v_S=u_Sc_S\). Для третьей формулировки нужно лишь заметить, что представления в группе идеальных классов определены только мультипликативно; утверждение о представлениях первой степени от этого не получает никакого ограничения.

\paragraph{Доказательство теоремы о главном роде.}
Доказательство второй формулировки основано на двух вспомогательных теоремах.

\paragraph{Вспомогательная теорема 1.}
Если трансформационные величины, образованные из идеалов \(c_S\),
\[
        \frac{c_S^{\,T}c_T}{c_{ST}}
\]
дают единичный идеал, то \(c_S\) являются символическими \((1-S)\)-ми степенями:
\[
        c_S=b^{1-S}\qquad(S\in\Gg).
\]
Это равносильно первой формулировке, если там предположить, что автоморфизм порождается \(v_S=u_Sc_S\). Доказательство: \(b=\sum c_S\) имеет требуемое свойство, причем сумма есть модульная сумма, то есть наибольший общий делитель. Действительно,
\[
        b=\sum_S c_S,\qquad
        b^T=\sum_S c_S^{\,T},\qquad
        b^T c_T=\sum_S c_S^{\,T}c_T=\sum_S c_{ST}=b,
\]
следовательно \(c_T=b^{1-T}\).

\paragraph{Вспомогательная теорема 2.}
Если алгебра
\[
        A=(a_{S,T},K,\Gg)
\]
расщепляется во всех конечных и бесконечных разветвленных местах \(K\), а главные идеалы \((a_{S,T})\) являются трансформационными величинами идеалов,
\[
        (a_{S,T})=\frac{c_S^{\,T}c_T}{c_{ST}},
\]
то \(A\) расщепляется всюду, то есть абсолютно. Тогда \((a_{S,T})\) становятся трансформационными величинами главных идеалов,
\[
        (a_{S,T})=\frac{(d_S)^{\,T}(d_T)}{(d_{ST})}.
\]
Обратно, всякая всюду расщепляющаяся алгебра удовлетворяет этому условию.

Доказательство основано на известном поведении \(A\) при \(p\)-адическом расширении. Пусть \(k_p\) обозначает \(p\)-адическое расширение \(k\), соответственно \(K_\PP\) -- \(\PP\)-адическое расширение \(K\), где \(\PP\) есть простой идеал над \(p\); пусть \(A_p\) -- расширение \(A\), полученное переходом от \(k\) к \(k_p\). Далее пусть \(\mathfrak Z\) -- группа разложения при \(\PP\), а \(a_{\mathfrak Z}\) -- соответствующий срез фактор-системы. Тогда
\[
        A_p\sim (a_{\mathfrak Z},K_\PP/k_p,\mathfrak Z).
\]
Если \(p\), в частности, неразветвлено, то есть \(\mathfrak Z\) циклична, то речь идет о выделенном представлении \(A_p\) посредством неразветвленного поля инерции \(K_\PP/k_p\).

Нужно показать, что при предположении вспомогательной теоремы 2 алгебра \(A_p\) расщепляется и в этом неразветвленном случае. Для бесконечных мест здесь \(K_\PP=k_p\), следовательно \(A_p\) расщепляется. Для конечного \(p\) сначала следует, что фактор-система \(a_{\mathfrak Z}\) ассоциирована с фактор-системой, состоящей из единиц. При переходе к нормированному циклическому представлению в \(A_p=(\alpha,K_\PP/k_p,R)\), где \(R\) -- порождающий элемент \(\mathfrak Z\), фактор-система задается соотношениями
\[
        u_{\PP_i}u_{\PP_j}=u_{\PP_{i+j}}\varepsilon_{\PP_i,\PP_j}
\]
Полагая \(u_R=u\), получаем
\[
        u^i=u_R^i\,\varepsilon_R\varepsilon_{R^2}\cdots \varepsilon_{R^{i-1}},
\]
и в частности \(u_E=\varepsilon_E\), если \(E\) есть порядок \(\mathfrak Z\). Так как в неразветвленном расширении \(K_\PP/k_p\) все единицы являются нормами, следует, что \(A_p\) расщепляется. Значит, \(A\) расщепляется всюду, а потому по теореме о расщепляющихся алгебрах -- абсолютно.

Иными словами: из предположения следует, что \(a_{S,T}\) являются трансформационными величинами из элементов,
\[
        a_{S,T}=\frac{d_S^{\,T}d_T}{d_{ST}},
        \qquad d_S\in K^*.
\]
В ослабленной форме это означает: \((a_{S,T})\) становятся трансформационными величинами главных идеалов. Если, наоборот, \(A\) всюду расщепляется, то и главные идеалы \((a_{S,T})\) становятся трансформационными величинами идеалов в каждом месте; и \(A\) действительно расщепляется по другой формулировке предположения теоремы о расщепляющихся алгебрах.

Из двух вспомогательных теорем непосредственно следует доказательство теоремы о главном роде. Пусть \(c_S\) означает некоторый идеал из класса \(c_S\). По определению индуцированного разбиения на классы предположение говорит, что трансформационные величины \(c_S\) становятся главными идеалами \((a_{S,T})\), удовлетворяющими предположениям вспомогательной теоремы 2. Следовательно, существуют главные идеалы \((d_S)\), такие что
\[
        \frac{c_S^{\,T}c_T}{c_{ST}}
\]
посредством \((d_S)\) становится единицей. Идеалы \(c_S/(d_S)\), лежащие в тех же классах \(c_S\), тем самым удовлетворяют предположению вспомогательной теоремы 1. Значит, существует идеальный класс \(b\) с
\[
        c_S=b^{1-S}.
\]
Это означает, что классы \(c_S\) являются символическими \((1-S)\)-ми степенями класса \(b\).

В циклическом случае теорема о главном роде переходит в известную теорему, как только главный класс индуцированного разбиения на классы при положенной в основу нормировке переходит в группу тех главных идеалов \((\alpha)\), для которых \(\alpha\) во всех разветвленных местах становится норменным вычетом.

\begin{center}
(Поступила 27. 10. 1932.)
\end{center}
% END UNIT BODY
% END PRODUCER UNIT 217
\endgroup


\clearpage
\begingroup
\setcounter{footnote}{0}
% BEGIN PRODUCER UNIT 218
% Source: C:/Users/Floris/Documents/interlanguage/03_projects/noether/02_slavic_working_corpus/translations/paper42/source_fidelity/russian/v001/Noether_Paper42_SourceFidelity_Russian_v001.tex
% Identity: 29718 B / SHA-256 C175D59D151847B4A53857068328151ACB8FD4616B65C56218212395D7A7CDA1
\setlength{\parindent}{1.1em}
\setlength{\parskip}{0pt}
\newcommand{\foreign}[1]{#1}
\providecommand{\Norm}{\operatorname{N}}
\providecommand{\Sp}{\operatorname{Sp}}
\providecommand{\Gg}{\mathfrak G}
\providecommand{\Hh}{\mathfrak H}
\providecommand{\Ii}{\mathfrak I}
\providecommand{\PP}{\mathfrak P}
\providecommand{\oo}{\mathfrak o}
\providecommand{\glqq}{``}
\providecommand{\grqq}{''}
\emergencystretch=220em
\allowdisplaybreaks
\sloppy
\hbadness=10000
\vbadness=10000
\hfuzz=5pt
\vfuzz=12pt
\raggedbottom
% BEGIN UNIT BODY
\section*{42. Расщепляющиеся скрещенные произведения и их максимальные порядки}

\begin{center}
\emph{\foreign{Actualités scientifiques et industrielles 148 (1934), S. 5--15}}
\end{center}

В простейшем случае расщепляющихся скрещенных произведений,\footnote{Теория скрещенных произведений, в связи с одной моей лекцией, изложена в разд. II работы H. Hasse, \foreign{\emph{Theory of cyclic algebras over an algebraic number field}}, Trans. Amer. Math. Soc. 34 (1932). Здесь я использую только первые основные понятия. Основания Брандта у Гассе и в последующих рамках обоснованы заново.} то есть тех, которые дают расщепляющиеся алгебры и потому имеют фактор-системы, ассоциированные с единицей, соотношения можно проследить в значительной мере явно. Это остается верным и для негалуаовых полей расщепления, то есть максимальных коммутативных подполей. Относящиеся сюда простые алгебраические факты развиваются в §1. В §2, при алгебраическом числовом поле в качестве основного поля, я даю явные представления всех максимальных порядков и их идеалов посредством модулей и комплементарных модулей лежащего в основе поля расщепления \(k\).\footnote{Эти представления были мне давно известны; толчок к последующим рассуждениям дало сообщение одной теоремы H. Hasse.}

Далее я даю разбиение на области: в одну область помещаются все такие максимальные порядки, которые имеют одно и то же пересечение с \(k\). Тем самым области взаимно однозначно соответствуют порядкам из \(k\); главному порядку соответствует главная область. Максимальные порядки одной области преобразуются друг в друга идеалами, образованными из модулей соответствующего порядка; в частности, в главной области речь идет о расширениях идеалов из \(k\). В §3 я перехожу к произвольным простым алгебрам. Уточнив разбиение на области так, что должны совпадать не только пересечение, но и \(p\)-адические компоненты максимальных порядков в местах ветвления, можно перенести теоремы, полученные в расщепляющемся случае. В частности, и тогда максимальные порядки главной области переходят друг в друга преобразованием с помощью расширений идеалов из \(k\).

Теоремы о главной области имеются у Chevalley и Hasse;\footnote{C. Chevalley, \foreign{Sur certains idéaux d'une algèbre simple}, Abh. Math. Sem. Hamburg, 1934; H. Hasse, \foreign{Über gewisse Ideale in einer einfachen Algebra}, работа, предшествующая настоящей.} Chevalley также показывает, что без уточненного разбиения, относящегося к местам ветвления, теоремы вообще уже не выполняются.

\subsection*{§I. Матричные единицы расщепляющихся скрещенных произведений и комплементарные базы полей расщепления}

В этом параграфе \(\Omega\) обозначает произвольное основное поле. Пусть \(k/\Omega\) -- галуаово сепарабельное расширение \(n\)-й степени, \(\Gg\) -- группа Галуа с элементами \(S,T,\ldots\), а \(u_S,u_T,\ldots\) -- соответствующие операторы. Рассматриваются скрещенные произведения с фактор-системой единица, то есть расщепляющиеся алгебры
\[
        K=\Gg\times k=u_{S_1}k+\cdots+u_{S_n}k,\qquad
        u_Su_T=u_{ST},
\]
с
\[
        z u_S=u_S z^S \qquad(z\in k).
\]
Полагая
\[
        E_\Gg=\sum_{S\in\Gg}u_S,
\]
из фактор-системы единица получаем соотношения
\[
        E_\Gg u_S=u_S E_\Gg=E_\Gg,                            \tag{1}
\]
\[
        zE_\Gg=\sum_Su_Sz^S,                                  \tag{2a}
\]
и
\[
        E_\Gg z E_\Gg
        =E_\Gg\sum_Su_Sz^S
        =E_\Gg\,\Sp(z),                                      \tag{2b}
\]
где \(\Sp(z)\) означает след \(z\) как элемента \(k/\Omega\).

\paragraph{Теорема.}
Если \(a_1,\ldots,a_n\) -- база \(k/\Omega\), а \(\bar a_1,\ldots,\bar a_n\) -- комплементарная к ней база, то
\[
        c_{ik}=\bar a_iE_\Gg a_k
\]
задает систему матричных единиц в \(K\). Алгебра \(K\) представляется как произведение модулей в форме
\[
        K=kE_\Gg k
          =k\Bigl(\sum_Su_S\Bigr)k
          =\sum_{i,k}(\bar a_iE_\Gg a_k)\,\Omega .
\]
Действительно, по (2b) получаем
\[
        c_{ij}c_{hk}
        =\bar a_iE_\Gg a_j\bar a_hE_\Gg a_k
        =\bar a_iE_\Gg a_k\,\Sp(a_j\bar a_h)
        =c_{ik}\,\Sp(a_j\bar a_h).
\]
Это как раз свойство матричных единиц, ввиду определяющих равенств
\[
        \Sp(a_j\bar a_h)=\delta_{jh}
\]
для комплементарных баз.

\paragraph{Замечание 1.}
Если \(Z\) -- какое-либо коммутативное поле расширения поля \(\Omega\), то все сохраняется, когда \(a_i\) означает базу \(k_Z/Z\), а \(K_Z\) становится прямой суммой полей. Подстановки группы тогда определяются тем, что на \(Z\) они индуцируют тождество.

\paragraph{Замечание 2.}
То, что \(K=kE_\Gg k\), следует также непосредственно из того, что правая часть является ненулевым идеалом простой системы \(K\).

Пусть теперь \(\Omega\) имеет характеристику нуль. Представление \(K=kE_\Gg k\) расщепляющейся алгебры остается верным и тогда, когда \(k\) является негалуаовым полем расщепления для \(K/\Omega\). Переходим к галуаову полю расщепления \(\bar k\). Пусть \(\bar k\) -- галуаово поле, соответствующее \(k\), \(\bar\Gg\) -- его группа, \(\bar u_S\) -- соответствующие операторы, а \(\bar K=\bar kE_{\bar\Gg}\bar k\) -- расщепляющееся скрещенное произведение. Подполе \(k\) соответствует подгруппе \(\Hh\) порядка \(h\). Положим
\[
        \bar\Gg=\Hh S_1+\cdots+\Hh S_n
              =S_1^{-1}\Hh+\cdots+S_n^{-1}\Hh,
        \qquad
        E_\Hh=\frac1h\sum_{H\in\Hh}u_H,
\]
и
\[
        E_\Gg=\frac1h E_{\bar\Gg}
              =\sum_i u_{S_i},\qquad
        u_{S_i}=E_\Hh\bar u_{S_i}.
\]
Тогда, в частности, идемпотент \(E_\Hh\) порождает единичное представление \(\Hh\), а \(E_\Gg\), как и \(E_{\bar\Gg}\), порождает единичное представление \(\bar\Gg\). Для так определенных \(E_\Gg\) и \(u_{S_i}\) соотношение (1) выполняется с одной стороны, а (2) -- полностью. Именно, прежде всего
\[
        E_\Gg=E_\Gg E_\Hh=E_\Hh E_\Gg=E_\Hh E_\Gg E_\Hh,       \tag{3}
\]
\[
        zE_\Hh=E_\Hh z,                                      \tag{4}
\]
и потому
\[
        z u_{S_i}=u_{S_i}z^{S_i}.                            \tag{5}
\]
Из (5) суммированием следует (2a), затем (2b), а (1) следует с одной стороны:
\[
        E_\Gg u_{S_i}=E_\Gg .
\]
Доказательство \(K=kE_\Gg k\) теперь проходит точно так же, как выше. Как полное матричное кольцо над \(\Omega\), \(K\) содержит каждое поле \(n\)-й степени как подполе, следовательно, в частности, и \(k\).

В связи с §3 заметим еще, что переход от \(K\) к \(\bar K\) возможен и тогда, когда \(K\) не расщепляется. Помимо \(k\), поле \(\bar k\) тогда также является полем расщепления, следовательно \(K\) подобна \(\bar K\), скрещенному произведению \(\bar\Gg\) с \(\bar k\). При этом, поскольку \(k\) является полем расщепления, подгруппе \(\Hh\), соответствующей \(k\), соответствует фактор-система единица; значит, идемпотент \(E_\Hh\) определен как выше. Отсюда в обратную сторону получается, что \(K\) изоморфна \(E_\Hh\bar K E_\Hh\). Эта алгебра изоморфна кольцу автоморфизмов левого идеала \(\bar K E_\Hh\), а значит подобна \(\bar K\); ранг совпадает с рангом \(K\). В расщепляющемся случае получаем
\[
        E_\Hh\bar K E_\Hh
        =E_\Hh\bar kE_{\bar\Gg}\bar kE_\Hh
        =(E_\Hh\bar kE_\Hh)E_{\bar\Gg}(E_\Hh\bar kE_\Hh)
        =kE_\Gg k .
\]

\subsection*{§II. Максимальные порядки расщепляющихся скрещенных произведений и их разбиение на области}

С этого момента пусть \(\Omega\) -- конечное алгебраическое числовое поле, \(K\) -- расщепляющаяся алгебра \(n\)-й степени над \(\Omega\), а \(k\) -- галуаово или негалуаово максимальное коммутативное подполе, так что \(K=kE_\Gg k\). Пусть \(a,\bar a\) обозначают пару комплементарных конечных \(\oo\)-модулей из \(k\), где \(\oo\) -- главный порядок поля \(\Omega\), а \(o\) -- главный порядок поля \(k\). Структура максимальных порядков в \(K\) относительно \(k\) описывается следующими теоремами.

\paragraph{Теорема 1.}
Произведение модулей
\[
        \mathcal O_a=\bar a E_\Gg a
\]
является максимальным порядком в \(K\); в частности,
\[
        \mathcal O=\mathcal O_o=oE_\Gg o
\]
является максимальным порядком.

\paragraph{Теорема 2.}
Преобразование \(\mathcal O\) в \(\mathcal O_a\) осуществляется левым идеалом
\[
        \mathcal L=oE_\Gg a
\]
и взаимно обратным ему правым идеалом
\[
        \mathcal R=\bar a E_\Gg o
\]
относительно \(\mathcal O\).

\paragraph{Замечание.}
Совпадение двух комплементарных модулей \(a,\bar a\) в произведении \(\mathcal L\mathcal R\) как раз благодаря образованию следа, вызванному \(E_\Gg\), приводит к взаимно обратному соотношению.

\paragraph{Теорема 3.}
Все левые и правые идеалы относительно \(\mathcal O\) исчерпываются указанными в теореме 2. Следовательно, максимальные порядки \(\mathcal O_a\) исчерпывают все максимальные порядки в \(K\). Любые два максимальных порядка \(\mathcal O_a\) и \(\mathcal O_b\) переходят друг в друга преобразованием с помощью
\[
        \mathcal L_{ab}=\bar aE_\Gg b,
        \qquad
        \mathcal R_{ba}=\bar bE_\Gg a
\]
и эти \(\mathcal L_{ab}\) и \(\mathcal R_{ba}\) исчерпывают левые и правые идеалы относительно \(\mathcal O_a\).

\paragraph{Теорема 4.}
Когда \(\mathcal O_a\) пробегает все максимальные порядки в \(K\), пересечение
\[
        [\mathcal O_a,k]
\]
пробегает все порядки в \(k\); это пересечение каждый раз равно порядку модуля \(a\). Если объединить все максимальные порядки, принадлежащие одному и тому же порядку в \(k\), в одну область, то преобразование внутри области осуществляется левыми идеалами \(\bar aE_\Gg b\) относительно \(\mathcal O_a\) и взаимно обратными правыми идеалами, где \(a\) и \(b\) являются модулями соответствующего порядка в \(k\).

\paragraph{Теорема 4a.}
В частности, преобразование максимальных порядков главной области, то есть области всех максимальных порядков, содержащих главный порядок \(o\) поля \(k\), осуществляется идеалами
\[
        aE_\Gg b,
\]
где \(a\) и \(b\) -- идеалы из \(k\). Иными словами: если \(\mathcal L\) -- левый идеал максимального порядка главной области и \(o\) содержится в правом порядке \(\mathcal L\), то \(\mathcal L\) является расширением \(o\)-идеала.

Доказательства основаны на переходе к отдельным местам. Достаточно каждый раз переходить к кольцу частных при \(p\), где \(p\) пробегает все простые идеалы из \(\Omega\). Если \(\mathcal C\) -- произвольный \(\oo\)-модуль в \(K\), то его компонентой \(\mathcal C_p\) будем называть расширенный модуль \(\mathcal C\oo_p\). Тогда:

\paragraph{Лемма.}
\(\mathcal C\) является пересечением всех расширенных модулей \(\mathcal C_p\). Действительно, если \(w\in\mathcal C_p\), то существует не делящийся на \(p\) элемент \(\sigma\in\oo\), такой что \(\sigma w\in\mathcal C\). Если же \(w\) лежит в пересечении всех \(\mathcal C_p\), и \(g\) обозначает идеал всех \(\sigma\in\oo\) с \(\sigma w\in\mathcal C\), то \(g\) не может делиться ни на какой \(p\); следовательно, \(g=\oo\). Лемма не предполагает максимального ранга модуля.

\paragraph{Следствия.}
Каждый модуль определяется своими компонентами. Если \(\mathcal D\) -- собственный надмодуль \(\mathcal C\), то по меньшей мере один \(\mathcal D_p\) является собственным надмодулем \(\mathcal C_p\). В частности, если \(\mathcal M\) -- порядок в \(K\), а все его компоненты \(\mathcal M_p\) являются максимальными порядками над \(\oo_p\), то \(\mathcal M\) тоже является максимальным порядком.

\paragraph{Доказательство теоремы 1.}
\(\mathcal O_a\) является порядком, поскольку \(\mathcal O_a\) -- конечный \(\oo\)-модуль максимального ранга. Далее
\[
        \mathcal O_a^2
        =\bar aE_\Gg a\,\bar aE_\Gg a
        =\bar aE_\Gg a\,\Sp(a\bar a)
        \subseteq \bar aE_\Gg a=\mathcal O_a.
\]
Чтобы показать максимальность, по следствию достаточно показать максимальность каждой компоненты. Так как \(\oo_p\) -- кольцо главных идеалов, \(a_p\) и \(\bar a_p\) имеют независимые комплементарные базы \(a_1,\ldots,a_n\) и \(\bar a_1,\ldots,\bar a_n\). Произведения \(\bar a_iE_\Gg a_k\) по §1 образуют матричные единицы; значит,
\[
        \mathcal O_{a,p}=\sum_{i,k}(\bar a_iE_\Gg a_k)\oo_p
\]
максимален. Так как \(\mathcal O_a\) максимален, отсюда также следует \(\Sp(a\bar a)=\oo\).

\paragraph{Доказательство теоремы 2.}
Из правила следа следует:
\[
\begin{aligned}
 \mathcal O\mathcal L&=oE_\Gg o\,oE_\Gg a=oE_\Gg a=\mathcal L,\\
 \mathcal R\mathcal O&=\bar aE_\Gg o\,oE_\Gg o=\bar aE_\Gg o=\mathcal R,\\
 \mathcal L\mathcal R&=oE_\Gg a\,\bar aE_\Gg o=oE_\Gg o=\mathcal O,\\
 \mathcal R\mathcal L&=\bar aE_\Gg o\,oE_\Gg a=\bar aE_\Gg a=\mathcal O_a.
\end{aligned}
\]

\paragraph{Доказательство теоремы 3.}
Пусть \(\mathcal L\) -- левый идеал относительно \(\mathcal O\). Тогда, в частности,
\[
        \mathcal L=oE_\Gg\mathcal L
        =(oE_\Gg l_1,\ldots,oE_\Gg l_n)
\]
с \(l_1,\ldots,l_n\), образующими \(\oo\)-базу \(\mathcal L\). Полагая \(l_\mu=\sum_i u_{S_i}z_{i\mu}\), получаем \(E_\Gg l_\mu=E_\Gg a_\mu\), значит \(\mathcal L=oE_\Gg a\). Для регулярного \(\mathcal L\) модуль \(a\) должен иметь максимальный ранг \(n\); тогда \(\bar aE_\Gg o\) является взаимно обратным правым идеалом. Общее утверждение следует композицией.

\paragraph{Доказательство теоремы 4.}
Пересечение \(t=[k,\mathcal O_a]\) является порядком. Как подмножество \(\mathcal O_a\), оно является областью правых мультипликаторов и является наибольшим порядком из \(o\), удовлетворяющим этому условию; поэтому оно наверняка содержит порядок \(o_a\) модуля \(a\). Чтобы показать \(t=o_a\), достаточно равенства компонент. Если \(\tau\in t_p\), то \(\tau c_{ik}\) для каждого \(c_{ik}=\bar a_iE_\Gg a_k\) является линейной комбинацией \(c_{ij}\) с коэффициентами из \(\oo_p\). С другой стороны, \(a_k\tau\) является линейной комбинацией \(a_j\) с коэффициентами из \(\oo_p\). В силу единственности представления через \(c_{ij}\) эти коэффициенты должны лежать в \(\oo_p\); следовательно, \(\tau\) лежит в порядке модуля \(a_p\). Тем самым \(t\) равен порядку \(a\).

\paragraph{Доказательство теоремы 4a.}
Это специализация теоремы 4 на главную область. Если сначала \(\mathcal O=oE_\Gg o\), а \(o\) также содержится в правом порядке \(\mathcal L=oE_\Gg a\), то \(\bar aE_\Gg a\) принадлежит главной области; значит, \(a\) является идеалом из \(k\), и \(\mathcal L=\mathcal O a\). Если исходить из \(\mathcal O_a\), то соответственно \(\mathcal L=\bar aE_\Gg b\), а \(a^{-1}b\) является идеалом из \(k\); получается \(\mathcal L=\mathcal O_a a^{-1}b\).

\subsection*{§III. Максимальные порядки произвольных скрещенных произведений}

Теоремы, выведенные для расщепляющихся скрещенных произведений, остаются в силе и в общем случае, как только разбиение на области уточняется. Факты о главной области тогда выполняются точно; остальные должны формулироваться как утверждения об отдельных местах. При этом место теперь следует понимать как \(p\)-адическое расширение.

\paragraph{Определение.}
Пусть \(K\) -- простая нормальная алгебра над \(\Omega\), а \(k\) -- максимальное коммутативное подполе. В область относительно \(k\) зачисляются все максимальные порядки из \(K\), которые имеют одно и то же пересечение с \(k\) и, кроме того, имеют одинаковые \(p\)-адические компоненты в местах ветвления \(K\). Если пересечение с \(k\) есть главный порядок, то речь идет о главной области.

Для главных областей тогда:

\paragraph{Теорема 1.}
Максимальные порядки одной главной области переходят друг в друга преобразованием с помощью расширенных идеалов идеалов из \(k\). Иными словами: левые идеалы максимальных порядков главной области, взаимно простые с дифферентой и содержащие \(o\) в своем правом порядке, являются расширениями идеалов из \(k\).

\paragraph{Теорема 2.}
Если \(K\) имеет простую степень, то существует только одна главная область; она переходит в себя под действием расширенных идеалов \(k\)-идеалов. Соответствующие идеалы фиксированного максимального порядка \(\mathcal O\) -- это идеалы из \(k\), расширенные двусторонними идеалами из \(\mathcal O\).

\paragraph{Доказательство.}
Два максимальных порядка одной главной области отличаются лишь в конечном числе мест; по определению это места, в которых \(K\) расщепляется. Здесь же \(K\) имеет форму, рассмотренную в §§1--2. Это ясно, если \(k\) галуаово; если фактор-система ассоциирована с единицей, ее можно заменить системой единица. По §1 это верно и в негалуаовом случае. Следовательно, в каждом таком месте и тем самым вообще речь идет о расширенных идеалах.

Если \(K\) имеет простую степень, то она либо расщепляется, либо является алгеброй с делением. Во втором случае она остается алгеброй с делением во всех местах ветвления и поэтому имеет только одну главную область, так что снова действует преобразование расширенными идеалами. Самое общее преобразование получается композицией расширенных идеалов с идеалами, которые преобразуют максимальный порядок в себя, то есть с двусторонними идеалами; как известно, от центральных идеалов они отличаются только в местах ветвления, следовательно, только расширенными идеалами.

Далее:

\paragraph{Теорема 3.}
Максимальные порядки произвольной области переходят друг в друга идеалами, взаимно простыми с дифферентой, которые в каждом месте составлены из модулей соответствующего порядка так, как указано в §2, теорема 3. Компоненты этих модулей при фиксированном идеале отличаются от порядка только в конечном числе мест.

Это непосредственно следует из §2; надо лишь учесть, что операторы \(u_S\) в отдельных местах должны быть заменены соответствующими эквивалентными операторами.

Вопрос о том, существуют ли и здесь для каждого порядка из \(k\) максимальные порядки, имеющие этот порядок пересечением, требует более подробного исследования в местах ветвления. Главная область всегда существует, как видно из скрещенного, возможно обобщенного, произведенного представления: оно приводит к порядку, содержащему \(o\), как только фактор-системы взяты целыми.

Следующий вопрос состоит в том, в какой мере можно арифметическим разбиением на области характеризовать алгебраические поля расщепления. Иными словами: всегда ли различные поля расщепления порождают различные разбиения на области, или уже различные главные области? И исчерпывают ли, быть может, главные области всех полей расщепления всю совокупность максимальных порядков?

Что такая характеристика посредством одного-единственного максимального порядка заведомо невозможна, показывают простейшие примеры. Так, максимальный порядок
\[
        \left[\,1,i,j,\frac{1+i+j+k}{2}\,\right]
\]
кватернионного тела содержит, кроме главного порядка \([1,i]\), также главный порядок
\[
        \left[\,1,\frac{1+i+j+k}{2}\,\right]
\]
поля третьих корней из единицы.

\begin{center}
Göttingen, август 1932.
\end{center}
\clearpage
% END UNIT BODY
% END PRODUCER UNIT 218
\endgroup


\clearpage
\begingroup
\setcounter{footnote}{0}
% BEGIN PRODUCER UNIT 219
% Source: C:/Users/Floris/Documents/interlanguage/03_projects/noether/02_slavic_working_corpus/translations/paper43/source_fidelity/russian/v001/Noether_Paper43_SourceFidelity_Russian_v001.tex
% Identity: 47630 B / SHA-256 382A94CE12A005BE789B549E55CC42C16D72C0DA12BB942150B164A08CD000EA
\setlength{\parindent}{1.1em}
\setlength{\parskip}{0pt}
\newcommand{\foreign}[1]{#1}
\providecommand{\Norm}{\operatorname{N}}
\providecommand{\Sp}{\operatorname{Sp}}
\providecommand{\Gg}{\mathfrak G}
\providecommand{\Hh}{\mathfrak H}
\providecommand{\Ii}{\mathfrak I}
\providecommand{\PP}{\mathfrak P}
\providecommand{\oo}{\mathfrak o}
\providecommand{\glqq}{``}
\providecommand{\grqq}{''}
\emergencystretch=220em
\allowdisplaybreaks
\sloppy
\hbadness=10000
\vbadness=10000
\hfuzz=5pt
\vfuzz=12pt
\raggedbottom
% BEGIN UNIT BODY
\section*{43. Дифференцирование идеалов и дифферента}

\begin{center}
\emph{\foreign{Journal f. d. reine u. angew. Math. 188 (1950), S. 1--21}}
\end{center}

\subsection*{Введение}

Главная теорема теории ветвления алгебраического числового поля, как известно, утверждает, что дифферента, то есть идеал ветвления, делится по крайней мере на \((e-1)\)-ю степень простого идеала, входящего в простое число \(p\) в \(e\)-й степени; ровно на \((e-1)\)-ю степень, если \(e\) не делится на \(p\), и на более высокую степень, если \(e\) делится на \(p\).

Эта теорема является аналогом того факта, что дифференциальное частное \(f'(x)\) многочлена \(f(x)\) делится по крайней мере на \((e-1)\)-ю степень линейного множителя, входящего в \(f(x)\) в \(e\)-й степени; ровно на \((e-1)\)-ю степень, если \(e\) не делится на характеристику области коэффициентов, и на более высокую степень, если \(e\) делится на эту характеристику.

Ниже я показываю, что здесь речь идет о большем, чем формальная аналогия. Дифференту можно понимать как дифференциальное частное определяющего идеала числового поля \(K\) в соответствующей целочисленной области многочленов от \(x_1,\ldots,x_n\), причем дифференциальное частное берется в точке \(x=\omega\), то есть при \(x_1=\omega_1,\ldots,x_n=\omega_n\), где \(\omega_1,\ldots,\omega_n\) означает модульный базис системы целых чисел из \(K\), главного порядка \(\mathfrak o\). При этом определяющий идеал \(\mathfrak M\) соответствует совокупности соотношений между \(\omega_i\), то есть состоит из всех целочисленных многочленов \(f(x)\), для которых \(f(\omega)\) обращается в нуль. Дифференциальное частное \(\mathfrak M'[x\to\omega]\) определяется как разностное частное, взятое в точке \(x=\omega\).

Именно, вследствие \(f(\omega)=0\), будем рассматривать идеал \(\mathfrak M\) как состоящий из всех разностей \(f(x)-f(\omega)\). Далее образуем разностный идеал
\[
        \mathfrak B=(x_1-\omega_1,\ldots,x_n-\omega_n)
\]
и разностное частное
\[
        \mathfrak A=\mathfrak M:\mathfrak B,
\]
где взятие частного означает обычное частное идеалов, взятое в области многочленов от \(x\) с целыми числами из \(K\), то есть числами из \(\mathfrak o\), в качестве коэффициентов. Тогда дифферента \(\mathfrak D\) определяется формулой
\[
        \mathfrak D=\mathfrak A[x\to\omega].
\]
Расширение области коэффициентов необходимо уже для того, чтобы разностный идеал и разностное частное вообще были определены. Тем самым это прямое обобщение дифференциального частного многочлена от одной переменной,
\[
        f'(\xi)=\frac{f(x)-f(\xi)}{x-\xi}\bigg|_{x=\xi},
\]
где и здесь область многочленов должна быть расширена присоединением \(\xi\), чтобы разности были определены. В самом деле, идеальное дифференциальное частное переходит в многочленное дифференциальное частное, как только базис \(\omega_i\) состоит из степеней одного элемента.

Данное здесь определение дифференты непосредственно примыкает к известным фактам. Однако оно вводит неинвариантные промежуточные объекты, поскольку идеал \(\mathfrak M\) зависит от выбранного базиса. Равносильное, полностью инвариантное определение получается, если вместо целочисленной области многочленов ввести кольцо \(\mathfrak O\), изоморфное главному порядку \(\mathfrak o\), которое становится изоморфным кольцу классов вычетов по \(\mathfrak M\), и расширить его область коэффициентов посредством \(\mathfrak o\) до \(\mathfrak O_{\mathfrak o}\). Здесь разностный идеал выводится из всех разностей \((x-\xi)\), где \(x\) и \(\xi\) соответствуют друг другу по изоморфизму от \(\mathfrak O\) к \(\mathfrak o\); разностное частное становится частным нулевого идеала по \(\mathfrak B\); дифферента получается из \(\mathfrak A\), когда каждый раз \(x\) заменяется изоморфным ему \(\xi\).

Это инвариантное определение ниже ставится в начало. Так непосредственно определяется дифферента произвольного коммутативного кольца относительно подкольца, причем должны выполняться лишь некоторые предпосылки, касающиеся возможности расширения коэффициентов. Например, это всегда имеет место при существовании модульного базиса, возможно после перехода к подходящим кольцам частных. Тем самым, в частности, определяется и дифферента произвольного порядка числового поля, а также дифферента кольца классов вычетов порядка по степени простого числа, относительная дифферента и дифферента для алгебраических функций от нескольких переменных.

Совпадение дифференциального определения дифференты с обычным следует из точного структурного исследования галуаова кольца расширения, сопоставленного числовому полю посредством прямого произведения. Это исследование основано на разложении в прямую сумму. В специальном случае кольца классов вычетов по многочлену от одной переменной оно сводится к анализу интерполяционной формулы Лагранжа.

Точнее, вводится кольцо \(\mathfrak K\), изоморфное числовому полю \(K\) и содержащее \(\mathfrak o\); его область коэффициентов, поле рациональных чисел, расширяется галуаовым полем \(F\) поля \(K\). Эта расширенная система \(\mathfrak K_F\) становится прямой суммой \(n\) простых компонент, соответствующих \(n\) изоморфизмам \(K\). Эти компоненты являются расширениями разностного частного и его сопряженных; одновременно нулевой идеал \(\mathfrak K_F\) является пересечением \(n\) идеалов, выведенных из разностного идеала и его сопряженных. Простые компоненты имеют вид \(Fe^{(i)}\), где \(e^{(i)}\) означают компоненты единицы. Поскольку при \(x=\xi\) элемент \(e^{(i)}\) переходит в единицу, разностное частное становится равным \(\mathfrak d e^{(i)}\), где под \(\mathfrak d\) понимается дифферента \(\mathfrak o\). Дифферента \(\mathfrak d\) состоит из всех элементов поля, которые после умножения на \(e\) лежат в прямом произведении порядка с самим собой. То же справедливо для сопряженных.

Чтобы отсюда получить определение через дополнительный модуль, надо заметить, что \(e^{(i)}\) становится линейной формой в базисе \(x_1,\ldots,x_n\), соответствующем \(\omega\); коэффициенты при \(x\) в \(e^{(i)}\) как раз дают базис дополнительного модуля \(\mathfrak o\) и его сопряженных. Из того, что разностное частное имеет коэффициенты из \(\mathfrak o\), сразу следует, что дифферента равна частному \(\mathfrak o\) по его дополнительному модулю; так получается дедекиндово определение. Рассуждения не ограничены главным порядком; они характеризуют дифференту каждого порядка как частное порядка по его дополнительному модулю.

В случае главного порядка совпадение дифференциального определения с определением через фундаментальное уравнение следует из указанного выше факта: идеальное дифференциальное частное переходит в многочленное дифференциальное частное, как только базис состоит из степеней одного элемента, фундаментальной формы. Тем самым задана и понятийная связь этого определения с дедекиндовым определением.

Если исходить из фундаментального уравнения, как известно, непосредственно следует главная теорема, упомянутая в начале. К точному определению показателей приводит рассмотрение дифференты \(\mathfrak d[p^t]\) кольца классов вычетов главного порядка по \(p^t\) при достаточно большом \(t\); достаточно \(t=n\), что для полей второй степени также необходимо. Существенно, что дифферента \(\mathfrak d\), рассматриваемая modulo \(p^t\), точно переходит в \(\mathfrak d[p^t]\). Тогда можно ограничиться кольцом классов вычетов по \(\mathfrak p^t\). Оно изоморфно кольцу классов вычетов по многочлену от одной переменной с коэффициентами modulo \(p^t\); дифференцирование этого многочлена дает, как другим путем показал O. Ore, точный обзор возможных показателей при помощи дополнительных чисел.

Наконец, структурное исследование дает также известную теорему: дифферента надполя над \(K\) равна произведению относительной дифференты на дифференту \(K\). Это следует из того, что то же мультипликативное свойство выполняется для единиц разложения в прямую сумму, а значит и для дополнительных модулей. Соответствующие факты справедливы и для относительных дифферент при переходе к кольцам частных, а также для алгебраических функциональных полей с полем коэффициентов характеристики нуль. При характеристике \(p\) и при целых алгебраических функциях структурные свойства, правда, сохраняются; но теория ветвления меняется из-за появления несепарабельных полей расширения.

\subsection*{§1. Расширение области коэффициентов кольца, или образование прямого произведения. Определяющий идеал}

В этом параграфе рассматриваются общие кольца; коммутативность умножения не предполагается. Однако, если явно не указано обратное, предполагается существование единичного элемента.

\paragraph{1. Определение прямого произведения.}
Кольцо \(\mathfrak O\times_{\mathfrak h}\mathfrak o\), или короче \(\mathfrak O_{\mathfrak o}\), называется прямым произведением колец \(\mathfrak O\) и \(\mathfrak o\) относительно \(\mathfrak h\), или кольцом, полученным из \(\mathfrak O\) расширением области коэффициентов \(\mathfrak h\) до \(\mathfrak o\), если выполнены следующие условия:

\begin{enumerate}
\item \(\mathfrak O_{\mathfrak o}\) содержит \(\mathfrak O\) и \(\mathfrak o\) как подкольца и является произведением этих колец.
\item Пересечение \(\mathfrak O\) и \(\mathfrak o\) равно \(\mathfrak h\); при этом \(\mathfrak h\) содержит единичный элемент.
\item \(\mathfrak O\) и \(\mathfrak o\) поэлементно перестановочны. Поэтому \(\mathfrak O_{\mathfrak o}\) состоит из билинейных сочетаний
\[
        x_1y_1+\cdots+x_my_m,
\]
где \(x_i\in\mathfrak O\), \(y_i\in\mathfrak o\).
\item Два элемента равны тогда и только тогда, когда они становятся формально равными вследствие соотношений в \(\mathfrak O\), с одной стороны, и в \(\mathfrak o\), с другой. Поэтому соотношение
\[
        \sum_i x_i y_i=0
\]
должно сводиться к соотношениям в обоих множителях путем добавления формально исчезающих сочетаний вида \((yh)\delta-y(h\delta)\).
\end{enumerate}

Этими требованиями равенство, сложение и умножение однозначно определяются через соответствующие операции в множителях. Если \(\mathfrak O\), \(\mathfrak o\) и их пересечение изоморфны соответствующим кольцам, то и прямое произведение изоморфно. Каждое кольцо, являющееся произведением \(\mathfrak O\) и \(\mathfrak o\) и в котором они поэлементно перестановочны, является гомоморфным образом прямого произведения.

\paragraph{2. Условие существования.}
Если \(\mathfrak O\) и \(\mathfrak o\) являются кольцами с общим центральным подкольцом \(\mathfrak h\), прямое произведение может не существовать. Например, пусть \(\mathfrak h=\mathbb Z\), \(\mathfrak O=\mathfrak h[x]\) с соотношением \(1-2x=0\), а \(\mathfrak o=\mathfrak h[y]\) с соотношением \(1-2y=0\). Общего объемлющего кольца с пересечением \(\mathfrak h\) существовать не может, ибо в нем было бы
\[
        0=(1-2x)y-x(1-2y)=y-x.
\]

Необходимое и достаточное условие существования прямого произведения таково: должно существовать хотя бы одно кольцо \(R\), объемлющее \(\mathfrak O\) и \(\mathfrak o\), в котором их пересечение есть \(\mathfrak h\) и в котором они поэлементно перестановочны. Если прямое произведение существует, оно само является таким кольцом. Обратно, прямое произведение строится как совокупность формальных билинейных сочетаний с равенством, сложением и умножением, заданными в пункте 1. Кольцо вложения \(R\) тогда обеспечивает условие пересечения.

\paragraph{3. Определяющий идеал.}
Система \(S\) элементов \(z\) из \(\mathfrak O\) называется системой образующих \(\mathfrak O\) относительно \(\mathfrak h\), если \(\mathfrak O=\mathfrak h[S]\). Если \(\mathfrak h\) лежит в центре \(\mathfrak O\), то \(\mathfrak O\) становится гомоморфным образом некоммутативной области многочленов \(\mathfrak h\{Z\}\), где переменные \(Z\) соответствуют системе \(S\). Имеем
\[
        \mathfrak O\simeq \mathfrak h\{Z\}/\mathfrak M,
\]
где \(\mathfrak M\) состоит из всех многочленов \(F(Z)\), для которых \(F(z)=0\) в \(\mathfrak O\). Этот двусторонний идеал \(\mathfrak M\) называется определяющим идеалом. Если \(S\) состоит из одного элемента \(z\) и \(\mathfrak M\) является главным идеалом, то базисный многочлен \(G(Z)\) с \(G(z)=0\) называется определяющим уравнением.

\paragraph{4. Определяющий идеал прямого произведения.}
Пусть \(\mathfrak O\simeq\mathfrak h\{Z\}/\mathfrak M\). Наряду с \(\mathfrak h\{Z\}\) рассмотрим \(\mathfrak o\{Z\}\). Идеал расширения \(\mathfrak M_{\mathfrak o}\) состоит из линейных сочетаний элементов \(\mathfrak M\) с коэффициентами из \(\mathfrak o\). Если прямое произведение \(\mathfrak O_{\mathfrak o}\) существует, то
\[
        \mathfrak O_{\mathfrak o}\simeq \mathfrak o\{Z\}/\mathfrak M_{\mathfrak o}.
\]
Следовательно, соотношения в прямом произведении суть ровно соотношения, возникающие из \(\mathfrak M\) расширением коэффициентов.

\subsection*{§2. Кольца с независимым модульным базисом. Построение прямого произведения}

\paragraph{1. Построение.}
Система \(T\) элементов \(t\) из \(\mathfrak O\) называется \(\mathfrak h\)-модульным базисом \(\mathfrak O\), если \(\mathfrak O\) состоит из всех линейных сочетаний
\[
        h_1t_1+\cdots+h_mt_m,
        \qquad h_i\in\mathfrak h.
\]
Модульный базис называется независимым, если из \(\sum h_it_i=0\) всегда следует, что все \(h_i=0\). Если \(\mathfrak O\) обладает независимым модульным базисом, содержащим единичный элемент, то прямое произведение \(\mathfrak O_{\mathfrak o}\) существует для каждого кольца расширения \(\mathfrak o\) кольца \(\mathfrak h\), чье теоретико-множественное пересечение с \(\mathfrak O\) равно \(\mathfrak h\) и элементы которого перестановочны с \(\mathfrak O\). За элементы берутся все линейные сочетания
\[
        y_1t_1+\cdots+y_mt_m,
        \qquad y_i\in\mathfrak o,
\]
с равенством коэффициентов в базисе \(T\). Тем самым построено кольцо вложения; оно одновременно является прямым произведением. Базис \(T\) при этом становится \(\mathfrak o\)-базисом \(\mathfrak O_{\mathfrak o}\).

\paragraph{2. Модули расширения и сужения.}
Если \(\mathfrak B\) является \(\mathfrak h\)-модулем в \(\mathfrak O\), то \(\mathfrak B_{\mathfrak o}\) обозначает порожденный им \(\mathfrak o\)-модуль в \(\mathfrak O_{\mathfrak o}\). Если \(\mathfrak C\) является \(\mathfrak o\)-модулем в \(\mathfrak O_{\mathfrak o}\), то его модуль сужения \([\mathfrak C,\mathfrak O]\) есть пересечение с \(\mathfrak O\). Если \(\mathfrak B\) имеет независимый \(\mathfrak h\)-модульный базис, дополняемый до независимого \(\mathfrak h\)-модульного базиса \(\mathfrak O\), то \(\mathfrak B\) является модулем сужения своего расширения:
\[
        \mathfrak B=[\mathfrak B_{\mathfrak o},\mathfrak O].
\]
Доказательство есть непосредственное сравнение коэффициентов в продолженном базисе.

\paragraph{3. Достаточный критерий.}
Если существует кольцо расширения \(\mathfrak f\) кольца \(\mathfrak h\), такое что прямые произведения \(\mathfrak O_{\mathfrak f}\) и \(\mathfrak o_{\mathfrak f}\) относительно \(\mathfrak h\), а затем и \(\mathfrak O_{\mathfrak f}\times_{\mathfrak f}\mathfrak o_{\mathfrak f}\) существуют, то существует \(\mathfrak O\times_{\mathfrak h}\mathfrak o\). Ведь последнее прямое произведение дает кольцо вложения, и
\[
        [\mathfrak O,\mathfrak o]\subseteq [\mathfrak O_{\mathfrak f},\mathfrak o_{\mathfrak f}]=\mathfrak f,
        \qquad [\mathfrak O,\mathfrak f]=\mathfrak h.
\]

\subsection*{§3. Дифферента кольца относительно подкольца. Дифференциальное частное определяющего идеала}

С этого места все рассматриваемые кольца коммутативны. Кроме того, кольца \(\mathfrak O\) и \(\mathfrak o\) являются изоморфными, а именно эквивалентными, расширениями подкольца \(\mathfrak h\), и предполагается прямое произведение \(\mathfrak O_{\mathfrak o}\) относительно \(\mathfrak h\).

\paragraph{1. Определение дифференты.}
Пусть \(x\) пробегает все элементы \(\mathfrak O\), а \(\xi\) каждый раз обозначает изоморфно соответствующий ему элемент из \(\mathfrak o\). Идеал, порожденный всеми разностями \(x-\xi\),
\[
        \mathfrak B=\{\ldots,x-\xi,\ldots\}
\]
называется разностным идеалом. Разностное частное есть частное нулевого идеала по \(\mathfrak B\):
\[
        \mathfrak A=(0):\mathfrak B.
\]
Дифферента \(\mathfrak d\) кольца \(\mathfrak o\) относительно \(\mathfrak h\) есть
\[
        \mathfrak d=\mathfrak A[x\to\xi],
\]
а дифферента \(\mathfrak D\) кольца \(\mathfrak O\) соответственно есть \(\mathfrak A[\xi\to x]\). Обе являются изоморфно соответствующими друг другу идеалами.

\paragraph{2. Дифференциальное частное определяющего идеала.}
Пусть \(\mathfrak M\) является определяющим идеалом \(\mathfrak O\) относительно \(\mathfrak h\), соответствующим системе образующих \(z\). Тогда в \(\mathfrak o[Z]\)
\[
        \mathfrak B_Z=\{\ldots,Z-\xi,\ldots\},
        \qquad
        \mathfrak C=\mathfrak M_{\mathfrak o}:\mathfrak B_Z.
\]
Следовательно,
\[
        \mathfrak C=\{\ldots,F(Z)-F(\xi),\ldots\}:
        \{\ldots,Z-\xi,\ldots\}.
\]
Дифференциальное частное \(\mathfrak M'\) есть \(\mathfrak C[\xi\to Z]\). Оно содержит \(\mathfrak M\) и при \(Z\to z\), соответственно \(Z\to\xi\), переходит в дифференту \(\mathfrak O\), соответственно \(\mathfrak o\):
\[
        \mathfrak M'[Z\to z]=\mathfrak D,
        \qquad
        \mathfrak M'[Z\to \xi]=\mathfrak d.
\]

\paragraph{3. Случай определяющего уравнения.}
Пусть \(\mathfrak O\) имеет определяющее уравнение
\[
        G(z)=0,\qquad
        G(Z)=Z^n+h_1Z^{n-1}+\cdots+h_n,
        \quad h_i\in\mathfrak h.
\]
Тогда дифферента определена и является главным идеалом с базисным элементом
\[
        G'(\xi) \quad\text{соответственно}\quad G'(z),
\]
где \(G'(Z)\) есть обычное многочленное дифференциальное частное. Идеальное дифференциальное частное есть
\[
        \mathfrak M'=(G'(Z),G(Z)).
\]
Действительно, \(1,z,\ldots,z^{n-1}\) образуют независимый модульный базис. Далее
\[
        C(Z,\xi)=\frac{G(Z)-G(\xi)}{Z-\xi}
\]
является базисом разностного частного; подстановкой \(\xi\to Z\) получается \(G'(Z)\).

\subsection*{§4. Дифферента прямой суммы}

Пусть
\[
        \mathfrak O=R_1+\cdots+R_r,
        \qquad
        e=e_1+\cdots+e_r,
        \qquad R_iR_j=0\;(i\ne j)
\]
представлено как прямая сумма идеалов, где \(e_i\) являются компонентами единицы. Далее пусть каждое \(R_i\) имеет независимый \(\mathfrak h_i\)-модульный базис, содержащий \(e_i\), где \(\mathfrak h_i=\mathfrak h e_i\), и пусть
\[
        [\mathfrak h,\mathfrak O_i]=0
\]
для соответствующих дополнительных слагаемых. Тогда и \(\mathfrak O\) имеет независимый \(\mathfrak h\)-модульный базис, образованный из базисов \(R_i\), а дифферента \(\mathfrak O\) относительно \(\mathfrak h\) есть прямая сумма дифферент \(R_i\) относительно \(\mathfrak h_i\).

Точнее: пусть то же разложение выполнено для \(\mathfrak o\) с компонентами \(r_i\) и идемпотентами \(\varepsilon_i\). Тогда в прямом произведении
\[
        \mathfrak O_{\mathfrak o}=\sum_{i,j} R_i\varepsilon_j,
\]
и разностный идеал разлагается так, что компонента \(R_i\varepsilon_i\) содержит разностный идеал прямого произведения \(R_i\times r_i\). Разностное частное является прямой суммой соответствующих частных. Поэтому, обозначая через \(\mathfrak A\) разностное частное, а через \(\mathfrak D,\mathfrak d\) дифференты, получаем
\[
        \mathfrak D=\mathfrak D_1+\cdots+\mathfrak D_r,
        \qquad
        \mathfrak d=\mathfrak d_1+
        \cdots+\mathfrak d_r.
\]
Отдельные части доказательства состоят в сравнении коэффициентов в независимых базисах, проверке более сильных соотношений пересечения
\[
        [\mathfrak o,\mathfrak O_i]=0,
\]
изоморфизмов
\[
        \mathfrak o\sim\mathfrak o e_i,
        \qquad
        R_i\sim R_i\varepsilon_i,
        \qquad
        \mathfrak h_i\sim \mathfrak h_i e_i,
\]
и отождествления
\[
        \mathfrak A e_i\varepsilon_i=(0):\mathfrak B e_i\varepsilon_i,
\]
то есть разностного частного компоненты. Отсюда следует утвержденная форма дифференты как прямой суммы.

\subsection*{§5. Галуаово кольцо расширения поля. Структурные теоремы}

\paragraph{1. Основные кольца.}
Пусть \(K\) является сепарабельным полем \(n\)-й степени над основным полем \(P\). Далее пусть \(F\) является соответствующим галуаовым полем, которое содержит \(n\) сопряженных полей
\[
        K^{(1)},\ldots,K^{(n)},\qquad K=K^{(1)},
\]
и является их произведением. Пусть \(\mathfrak K\) обозначает изоморфное \(K\) эквивалентное расширение \(P\), такое что \([\mathfrak K,F]=P\). Так как \(\mathfrak K\) имеет конечный независимый \(P\)-модульный базис, существуют прямые произведения
\[
        \mathfrak K_K,
        \qquad \mathfrak K_F,
        \qquad K_F.
\]
Кольцо \(\mathfrak K_F\) называется галуаовым кольцом расширения \(\mathfrak K\). Оно является произведением своих \(n\) подколец \(\mathfrak K_{K^{(i)}}\).

Каждый идеал из \(\mathfrak K\) является идеалом сужения своего расширения в \(\mathfrak K_K\) или \(\mathfrak K_F\); каждый идеал из \(\mathfrak K_K\) является идеалом сужения своего расширения в \(\mathfrak K_F\). Это следует из §2, поскольку над полями каждый модуль имеет дополняемый независимый базис.

Обозначим через
\[
        \mathfrak B^{(i)}=\{\ldots,x-\xi^{(i)},\ldots\}
\]
сопряженные разностные идеалы и через
\[
        \mathfrak A^{(i)}=(0):\mathfrak B^{(i)}
\]
сопряженные разностные частные. Их расширения в \(\mathfrak K_F\) обозначаются \(\mathfrak B_F^{(i)}\) и \(\mathfrak A_F^{(i)}\).

\paragraph{2. Разложение в прямую сумму.}
Структурная теорема гласит: нулевой идеал галуаова кольца расширения \(\mathfrak K_F\) является пересечением расширений \(n\) сопряженных разностных идеалов; само \(\mathfrak K_F\) является прямой суммой расширений \(n\) сопряженных разностных частных:
\[
        (0)=\bigcap_{i=1}^n \mathfrak B_F^{(i)},
        \qquad
        \mathfrak K_F=\mathfrak A_F^{(1)}+\cdots+\mathfrak A_F^{(n)}.
\]
В частности,
\[
        \mathfrak K_K=\mathfrak A+\mathfrak B
\]
как прямая сумма, а его нулевой идеал есть пересечение соответствующего разностного идеала с разностным частным.

Доказательство основано на том, что кольцо классов вычетов по каждому \(\mathfrak B_F^{(i)}\) имеет ранг один над \(F\), и что благодаря сепарабельности идеалы \(\mathfrak B_F^{(i)}\) попарно различны и взаимно просты. Отсюда следует единственное представление единицы в виде прямой суммы
\[
        e=e^{(1)}+\cdots+e^{(n)},
\]
где
\[
        e^{(i)}\in \bigcap_{j\ne i}\mathfrak B_F^{(j)},
        \qquad
        e^{(i)}\equiv 1\pmod{\mathfrak B_F^{(i)}}.
\]
Соответствующие простые компоненты суть \(Fe^{(i)}\).

\paragraph{3. Представление компонент через сопряженные элементы.}
Если
\[
        x=\xi^{(1)}e^{(1)}+\cdots+\xi^{(n)}e^{(n)}
\]
является компонентным представлением элемента \(x\) из \(\mathfrak K\), то компоненты сопряжены. В частности, \(\xi^{(1)},\ldots,\xi^{(n)}\) суть \(n\) сопряженных значений, которые принимает \(x\) при изоморфизмах \(\mathfrak K\) в \(K^{(i)}\).

Отсюда следует: если \(B\) является абсолютным модулем в \(\mathfrak K\), а \(B^{(i)}\) его компоненты в \(\mathfrak K_F\), то
\[
        B=[\mathfrak K,B^{(1)}+\cdots+B^{(n)}].
\]

\paragraph{4. Дополнительные базисы.}
Каждому \(P\)-базису \(t_1,\ldots,t_n\) поля \(\mathfrak K\) соответствует дополнительный базис \(T_1,\ldots,T_n\). Изоморфные ему базисы \(A_1^{(i)},\ldots,A_n^{(i)}\) полей \(K^{(i)}\) определяются коэффициентами компонент единицы:
\[
        e^{(i)}=A_1^{(i)}t_1+
        \cdots+A_n^{(i)}t_n.
\]
Если \(\alpha_1^{(i)},\ldots,\alpha_n^{(i)}\) являются сопряженными базисами, соответствующими \(t_j\), то матрицы
\[
        (\alpha_j^{(i)})_{i,j}
        \quad\text{и}\quad
        (A_j^{(i)})_{i,j}
\]
взаимно обратны:
\[
        (\alpha_j^{(i)})(A_j^{(i)})^{t}=E_n.
\]
Дополнительный базис дополнительного базиса снова является исходным. Если два базиса связаны соотношением \((s)=(t)P\), то дополнительные связаны контраградиентным соотношением
\[
        (S)=P^{-1}(T).
\]

Если
\[
        c=C_1t_1+
        \cdots+C_nt_n
        =c_1T_1+
        \cdots+c_nT_n,
\]
то при \(\operatorname{Sp}\), обозначающем след, то есть сумму сопряженных, имеем
\[
        C_i=\operatorname{Sp}(cT_i),
        \qquad
        c_i=\operatorname{Sp}(ct_i).
\]

\paragraph{5. Определяющее уравнение.}
Предусмотренное в рукописи изложение специального случая определяющего уравнения и интерполяционной формулы Лагранжа осталось незаполненным. Предыдущие теоремы дают инвариантную форму соответствующих утверждений.

\subsection*{§6. Дифферента порядка}

Теперь структурные теоремы для полей надо уточнить в смысле целости. Пусть \(\mathfrak h\) является целозамкнутым подкольцом основного поля \(P\), поле частных которого есть \(P\). Далее пусть \(\mathfrak O\) является \(\mathfrak h\)-порядком в \(\mathfrak K\), то есть подкольцом, содержащим \(\mathfrak h\), имеющим поле частных \(\mathfrak K\) и обладающим конечным, не обязательно независимым, \(\mathfrak h\)-модульным базисом. Изоморфный порядок в \(K\) пусть будет \(\mathfrak o\), его сопряженные -- \(\mathfrak o^{(i)}\); кольцо в \(F\), порожденное ими, пусть будет \(\mathfrak g\).

Прямые произведения \(\mathfrak O_{\mathfrak o}\) и \(\mathfrak O_{\mathfrak g}\) существуют. Следовательно, дифферента \(\mathfrak O\), соответственно \(\mathfrak o\), относительно \(\mathfrak h\) определена. Обозначим через \(\mathfrak B\) разностный идеал, через \(\mathfrak A=(0):\mathfrak B\) разностное частное, а через
\[
        \mathfrak D=\mathfrak A[\xi\to x],
        \qquad
        \mathfrak d=\mathfrak A[x\to\xi]
\]
дифференты.

Структурная теорема для галуаова кольца расширения порядков гласит: разностный идеал и разностное частное порядка переходят при расширении в соответствующие идеалы полевого случая; разностное частное является сужением расширенного частного. Далее
\[
        \mathfrak A=\mathfrak d e^{(1)},
        \qquad
        \mathfrak A^{(i)}=\mathfrak d^{(i)}e^{(i)},
\]
и поэтому
\[
        \mathfrak d=[\mathfrak o,\mathfrak A^{(1)}+
        \cdots+\mathfrak A^{(n)}].
\]
Итак, дифферента \(\mathfrak d\) состоит из всех элементов \(x\) из \(K\), для которых \(xe^{(1)}\) лежит в \(\mathfrak O_{\mathfrak o}\). Она отлична от нулевого идеала.

\paragraph{Дополнительный модуль.}
Если \(\mathfrak O\) имеет независимый \(\mathfrak h\)-базис \(t_1,\ldots,t_n\), то дополнительные элементы \(T_1,\ldots,T_n\) образуют базис \(\mathfrak h\)-модуля \(\mathfrak C\), дополнительного модуля \(\mathfrak O\). Тогда
\[
        \mathfrak d=\mathfrak o:\mathfrak C.
\]
Ибо, записав
\[
        e^{(1)}=A_1t_1+
        \cdots+A_nt_n,
\]
получаем, что \(\mathfrak d\) состоит из элементов \(d\), для которых все \(dA_i\) лежат в \(\mathfrak o\). В специальном случае регулярного порядка с базисом
\[
        1,z,\ldots,z^{n-1}
\]
имеем, следовательно,
\[
        \mathfrak d=(f'(z)).
\]

Эквивалентно, дополнительный модуль \(\mathfrak C\) состоит из всех элементов \(c\) из \(\mathfrak K\), для которых следы
\[
        \operatorname{Sp}(co),\qquad o\in\mathfrak O,
\]
целы, то есть лежат в \(\mathfrak h\).

\paragraph{Дополнение.}
Если в порядке существует независимый \(\mathfrak h\)-модульный базис, то разностное частное \(\mathfrak A\) является пересечением всех сопряженных разностных идеалов, отличных от соответствующего разностного идеала:
\[
        \mathfrak A=[\mathfrak O_{\mathfrak o},
        \mathfrak B^{(2)}\cap\cdots\cap\mathfrak B^{(n)}].
\]
Это следует из того, что тогда и разностный идеал является идеалом сужения своего расширения.

\paragraph{Фундаментальное уравнение.}
Если существует фундаментальное уравнение, так что после присоединения переменных и перехода к кольцу частных порядок имеет базис \(1,u,\ldots,u^{n-1}\), то дифферента задается многочленным дифференциальным частным определяющего уравнения. Если
\[
        G(u)=0,
\]
то дифферента считывается из коэффициентов \(G'(u)\). При целозамкнутости отсюда следует необходимое умножение содержаний коэффициентных идеалов; тем самым становится возможной теория ветвления посредством рассмотрения фундаментального уравнения modulo \(p^t\).

\subsection*{§7. Теория ветвления главного порядка modulo \(p^t\)}

Пусть задано числовое поле, а \(\mathfrak h\) -- его кольцо целых чисел; или относительное поле относительно кольца частных главного порядка по простому идеалу, где \(p\) означает базисный элемент этого простого идеала. После присоединения переменных \(u_1,\ldots,u_n\) как главный порядок, так и кольцо классов вычетов по \(p^t\) имеют \(\mathfrak h\)-модульный базис
\[
        1,u,\ldots,u^{n-1}.
\]
В обоих случаях дифферента задается \(G'(u)\). Дифферента главного порядка modulo \(p^t\) переходит в дифференту кольца классов вычетов по \(p^t\).

Если
\[
        p=\mathfrak p_1^{e_1}
        \cdots\mathfrak p_r^{e_r},
\]
то кольцо классов вычетов по \(p^t\) является прямой суммой компонент, соответствующих кольцам классов вычетов по \(\mathfrak p_i^{te_i}\):
\[
        R=Re_1+
        \cdots+Re_r.
\]
По §4 дифферента этой прямой суммы есть прямая сумма покомпонентных дифферент. Поэтому достаточно рассматривать кольцо классов вычетов по степени \(\mathfrak p^N\).

Выберем \(\xi\) так, что
\[
        \varphi(\xi)\equiv0\pmod p,
        \qquad
        \varphi(\xi)\not\equiv0\pmod {p^2},
\]
где \(\varphi(x)\) является простой функцией modulo \(p\) степени \(f\). Для \(t\ge2\) заменим \(\xi\) на \(\xi^*=\xi e_1\), где \(e_1\) является первым идемпотентом разложения кольца классов вычетов modulo \(p^t\). Тогда даже
\[
        \mathfrak p=(p,\varphi(\xi^*)).
\]
Базис modulo \(p^t\) задается элементами
\[
        \xi^{\mu}\varphi(\xi)^{\nu},
        \qquad
        \mu=0,1,\ldots,f-1,
        \quad
        \nu=0,1,\ldots,e-1.
\]
Ибо из \((\varphi(\xi),p)=\mathfrak p\) и \((p)=\mathfrak p^e\) в кольце классов вычетов следует
\[
        \varphi(\xi)=pM(\xi),
        \qquad M(\xi)\not\equiv0\pmod p,
\]
так что более высокие степени можно последовательно выражать через более низкие. Определяющее уравнение относительно этого базиса имеет вид
\[
        F(x)=\varphi(x)^e+pM(x),
\]
где \(M(x)\) modulo \(p\) не делится на \(\varphi(x)\). Тем самым дифферента кольца классов вычетов задается многочленным дифференциальным частным \(F'(\xi)\). Так определение показателей сводится к результатам O. Ore, в частности к дополнительным числам и утверждениям дедекиндово-гензелевского типа. Одновременно эти результаты переходом к кольцу частных распространяются на относительные дифференты; добавляется только выразимость дифференты надполя как произведения дифференты подполя и относительной дифференты.

\begin{center}
Поступило 25 октября 1949.
\end{center}
% END UNIT BODY
% END PRODUCER UNIT 219
\endgroup

\end{document}
